\IfFileExists{stacks-project-book.cls}{%
\documentclass{stacks-project-book}
}{%
\documentclass{amsbook}
}

\usepackage{verbatim}
\newenvironment{reference}{\comment}{\endcomment}
\newenvironment{slogan}{\comment}{\endcomment}
\newenvironment{history}{\comment}{\endcomment}

\usepackage[all]{xy}

\xyoption{2cell}
\UseAllTwocells

\usepackage{multicol}

\usepackage{lmodern}
\usepackage[T1]{fontenc}
\usepackage[french]{babel}
\AtBeginDocument{\shorthandoff{?}}


\usepackage{hyperref}


\theoremstyle{plain}
\newtheorem{theorem}[subsection]{Théorème}
\newtheorem{proposition}[subsection]{Proposition}
\newtheorem{lemma}[subsection]{Lemme}

\theoremstyle{definition}
\newtheorem{definition}[subsection]{Définition}
\newtheorem{example}[subsection]{Exemple}
\newtheorem{exercise}[subsection]{Exercice}
\newtheorem{situation}[subsection]{Situation}

\theoremstyle{remark}
\newtheorem{remark}[subsection]{Remarque}
\newtheorem{remarks}[subsection]{Remarques}

\renewcommand{\proofname}{Démonstration}

\numberwithin{equation}{subsection}

\def\lim{\mathop{\mathrm{lim}}\nolimits}
\def\colim{\mathop{\mathrm{colim}}\nolimits}
\def\Spec{\mathop{\mathrm{Spec}}}
\def\Hom{\mathop{\mathrm{Hom}}\nolimits}
\def\Ext{\mathop{\mathrm{Ext}}\nolimits}
\def\SheafHom{\mathop{\mathcal{H}\!\mathit{om}}\nolimits}
\def\SheafExt{\mathop{\mathcal{E}\!\mathit{xt}}\nolimits}
\def\Sch{\mathit{Sch}}
\def\Mor{\mathop{\mathrm{Mor}}\nolimits}
\def\Ob{\mathop{\mathrm{Ob}}\nolimits}
\def\Sh{\mathop{\mathit{Sh}}\nolimits}
\def\NL{\mathop{N\!L}\nolimits}
\def\CH{\mathop{\mathrm{CH}}\nolimits}
\def\proetale{{pro\text{-}\acute{e}tale}}
\def\etale{{\acute{e}tale}}
\def\QCoh{\mathit{QCoh}}
\def\Ker{\mathop{\mathrm{Ker}}}
\def\Im{\mathop{\mathrm{Im}}}
\def\Coker{\mathop{\mathrm{Coker}}}
\def\Coim{\mathop{\mathrm{Coim}}}

\DeclareMathSymbol{\boxtimes}{\mathbin}{AMSa}{"02}

\def\QCohstack{\mathcal{QC}\!\mathit{oh}}
\def\Cohstack{\mathcal{C}\!\mathit{oh}}
\def\Spacesstack{\mathcal{S}\!\mathit{paces}}
\def\Quotfunctor{\mathrm{Quot}}
\def\Hilbfunctor{\mathrm{Hilb}}
\def\Curvesstack{\mathcal{C}\!\mathit{urves}}
\def\Polarizedstack{\mathcal{P}\!\mathit{olarized}}
\def\Complexesstack{\mathcal{C}\!\mathit{omplexes}}
\def\Pic{\mathop{\mathrm{Pic}}\nolimits}
\def\Picardstack{\mathcal{P}\!\mathit{ic}}
\def\Picardfunctor{\mathrm{Pic}}
\def\Deformationcategory{\mathcal{D}\!\mathit{ef}}
\begin{document}
\title{Le projet Stacks --- édition française, partie 8}
\maketitle
\tableofcontents

% OK, start here.
%

\chapter{Cohomologie cristalline}



\phantomsection
\label{crystalline-section-phantom}




\section{Introduction}
\label{crystalline-section-introduction}

\noindent
Ce chapitre s'appuie sur une série de cours donnés par Johan de Jong
en 2012 à l'université Columbia.
L'objectif de ce chapitre est de donner une introduction rapide à la
cohomologie cristalline. Le livre \cite{Berthelot} fournit une référence.

\medskip\noindent
Nous avons placé l'étude purement algébrique et plus élémentaire des anneaux
à puissances divisées dans un chapitre préliminaire, car elle sert aussi
à étudier les résolutions de Tate en algèbre commutative.
Voir Algèbre à puissances divisées, section \ref{dpa-section-introduction}.













\section{Enveloppe à puissances divisées}
\label{crystalline-section-divided-power-envelope}

\noindent
La construction du lemme suivant sera appelée l'enveloppe à
puissances divisées. Elle jouera plus loin un rôle important.

\begin{lemma}
\label{crystalline-lemma-divided-power-envelope}
Soit $(A, I, \gamma)$ un anneau à puissances divisées.
Soit $A \to B$ un homomorphisme d'anneaux. Soit $J \subset B$ un idéal
tel que $IB \subset J$. Il existe un homomorphisme d'anneaux
à puissances divisées
$$
(A, I, \gamma) \longrightarrow (D, \bar J, \bar \gamma)
$$
tel que
$$
\Hom_{(A, I, \gamma)}((D, \bar J, \bar \gamma), (C, K, \delta)) =
\Hom_{(A, I)}((B, J), (C, K))
$$
fonctoriellement en l'algèbre à puissances divisées $(C, K, \delta)$ sur
$(A, I, \gamma)$. Ici, le membre de gauche désigne les morphismes d'anneaux
à puissances divisées sur $(A, I, \gamma)$, et celui de droite les morphismes
de couples (anneau, idéal) sur $(A, I)$.
\end{lemma}

\begin{proof}
Notons $\mathcal{C}$ la catégorie des anneaux à puissances divisées
$(C, K, \delta)$. Considérons le foncteur
$F : \mathcal{C} \longrightarrow \textit{Ens}$ défini par
$$
F(C, K, \delta) =
\left\{
(\varphi, \psi)
\middle|
\begin{matrix}
\varphi : (A, I, \gamma) \to (C, K, \delta)
\text{ homomorphisme d'anneaux à puissances divisées} \\
\psi : (B, J) \to (C, K)\text{ un homomorphisme de }
A\text{-algèbres tel que }\psi(J) \subset K
\end{matrix}
\right\}
$$
Nous allons montrer que
le lemme \ref{dpa-lemma-a-version-of-brown} d'Algèbre à puissances divisées
s'applique à ce foncteur, ce qui
démontrera le lemme. Supposons $(\varphi, \psi) \in F(C, K, \delta)$.
Soit $C' \subset C$ le sous-anneau engendré par $\varphi(A)$,
$\psi(B)$ et les $\delta_n(\psi(f))$ pour tout $f \in J$.
Soit $K' \subset K \cap C'$ l'idéal de $C'$ engendré par
$\varphi(I)$ et les $\delta_n(\psi(f))$, pour $f \in J$.
Alors $(C', K', \delta|_{K'})$ est un anneau à puissances divisées et
le cardinal de $C'$ est majoré par le cardinal
$\kappa = |A| \otimes |B|^{\aleph_0}$.
De plus, $\varphi$ se factorise sous la forme $A \to C' \to C$ et $\psi$
sous la forme $B \to C' \to C$. Cela démontre que l'hypothèse (1) du
lemme \ref{dpa-lemma-a-version-of-brown} d'Algèbre à puissances divisées
est satisfaite. L'hypothèse (2) est claire,
car le foncteur d'oubli de la catégorie des anneaux à puissances divisées,
$(C, K, \delta) \mapsto (C, K)$, commute aux limites ; voir le
lemme \ref{dpa-lemma-limits} d'Algèbre à puissances divisées et sa démonstration.
\end{proof}

\begin{definition}
\label{crystalline-definition-divided-power-envelope}
Soit $(A, I, \gamma)$ un anneau à puissances divisées.
Soit $A \to B$ un homomorphisme d'anneaux. Soit $J \subset B$ un idéal
tel que $IB \subset J$. L'algèbre à puissances divisées $(D, \bar J, \bar\gamma)$
construite dans le lemme \ref{crystalline-lemma-divided-power-envelope}
est appelée l'{\it enveloppe à puissances divisées de $J$ dans $B$
relative à $(A, I, \gamma)$} et se note $D_B(J)$ ou $D_{B, \gamma}(J)$.
\end{definition}

\noindent
Soit $(A, I, \gamma) \to (C, K, \delta)$ un homomorphisme d'anneaux à
puissances divisées. La propriété universelle de
$D_{B, \gamma}(J) = (D, \bar J, \bar \gamma)$ est
$$
\begin{matrix}
\text{homomorphismes d'anneaux }B \to C \\
\text{ qui envoient }J\text{ dans }K
\end{matrix}
\longleftrightarrow
\begin{matrix}
\text{homomorphismes à puissances divisées} \\
(D, \bar J, \bar \gamma) \to (C, K, \delta)
\end{matrix}
$$
et la correspondance s'obtient par précomposition avec l'application $B \to D$
qui correspond à $\text{id}_D$. Voici quelques propriétés de
$(D, \bar J, \bar \gamma)$ qui découlent directement de la propriété
universelle. Il existe des homomorphismes de $A$-algèbres
\begin{equation}
\label{crystalline-equation-divided-power-envelope}
B \longrightarrow D \longrightarrow B/J
\end{equation}
La première flèche envoie $J$ dans $\bar J$, et $\bar J$ est le noyau
de la seconde. Les éléments $\bar\gamma_n(x)$, où $n > 0$
et où $x$ appartient à l'image de $J \to D$, engendrent $\bar J$
comme idéal de $D$ et engendrent $D$ comme $B$-algèbre.

\begin{lemma}
\label{crystalline-lemma-divided-power-envelop-quotient}
Soit $(A, I, \gamma)$ un anneau à puissances divisées.
Soit $\varphi : B' \to B$ une surjection de $A$-algèbres de noyau $K$.
Soit $IB \subset J \subset B$ un idéal. Soit $J' \subset B'$
l'image réciproque de $J$. Écrivons
$D_{B', \gamma}(J') = (D', \bar J', \bar\gamma)$.
Alors $D_{B, \gamma}(J) = (D'/K', \bar J'/K', \bar\gamma)$
où $K'$ est l'idéal engendré par les éléments $\bar\gamma_n(k)$
pour $n \geq 1$ et $k \in K$.
\end{lemma}

\begin{proof}
Écrivons $D_{B, \gamma}(J) = (D, \bar J, \bar \gamma)$.
La propriété universelle de $D'$ fournit un homomorphisme $D' \to D$
d'algèbres à puissances divisées. Puisque $B' \to B$ et $J' \to J$ sont surjectifs,
l'homomorphisme $D' \to D$ est surjectif (voir les remarques ci-dessus). Il est clair que
$\bar\gamma_n(k)$ appartient au noyau pour $n \geq 1$ et $k \in K$ ; autrement dit,
nous obtenons un homomorphisme $D'/K' \to D$. Réciproquement, il existe une structure de
puissances divisées sur $\bar J'/K' \subset D'/K'$ ; voir le
lemme \ref{dpa-lemma-kernel} d'Algèbre à puissances divisées.
La propriété universelle de $D$ fournit donc un inverse $D \to D'/K'$, ce qui achève la démonstration.
\end{proof}

\noindent
Dans la situation de la définition \ref{crystalline-definition-divided-power-envelope},
nous pouvons choisir une surjection $P \to B$, où $P$ est une algèbre
polynomiale sur $A$, et prendre pour $J' \subset P$ l'image réciproque de $J$.
Le lemme précédent décrit $D_{B, \gamma}(J)$ en fonction de
$D_{P, \gamma}(J')$. Notons que $\gamma$ se prolonge en une structure de puissances
divisées $\gamma'$ sur $IP$ d'après le
lemme \ref{dpa-lemma-gamma-extends} d'Algèbre à puissances divisées. Ainsi,
$D_{P, \gamma}(J') = D_{P, \gamma'}(J')$ est un exemple du cas particulier
d'enveloppes à puissances divisées que décrit le lemme suivant.

\begin{lemma}
\label{crystalline-lemma-describe-divided-power-envelope}
Soit $(B, I, \gamma)$ une algèbre à puissances divisées. Soit $I \subset J \subset B$
un idéal. Soit $(D, \bar J, \bar \gamma)$ l'enveloppe à puissances divisées
de $J$ relativement à $\gamma$. Choisissons des éléments $f_t \in J$, $t \in T$, tels
que $J = I + (f_t)$. Il existe alors une surjection
$$
\Psi :
(B\langle x_t \rangle, IB\langle x_t \rangle + B\langle x_t \rangle_+, \delta)
\longrightarrow
(D, \bar J, \bar \gamma)
$$
d'anneaux à puissances divisées qui envoie $x_t$ sur l'image de $f_t$ dans $D$.
Le noyau de $\Psi$ est engendré par les éléments $x_t - f_t$ et par tous les
éléments
$$
\delta_n\left(\sum r_t x_t - r_0\right)
$$
dès que $\sum r_t f_t = r_0$ dans $B$ pour certains $r_t \in B$, $r_0 \in I$.
\end{lemma}

\begin{proof}
Dans l'énoncé du lemme, nous considérons $B\langle x_t \rangle$
comme un anneau à puissances divisées d'idéal
$J' = IB\langle x_t \rangle + B\langle x_t \rangle_{+}$ ; voir la
remarque \ref{dpa-remark-divided-power-polynomial-algebra} d'Algèbre à puissances divisées.
L'existence de $\Psi$ résulte de la propriété universelle des
algèbres polynomiales à puissances divisées. La surjectivité de $\Psi$ résulte
du fait que son image est un sous-anneau à puissances divisées de $D$, donc égale à $D$
par la propriété universelle de $D$. Il est clair que
$x_t - f_t$ appartient au noyau. Posons
$$
\mathcal{R} = \{(r_0, r_t) \in I \oplus \bigoplus\nolimits_{t \in T} B
\mid \sum r_t f_t = r_0 \text{ dans }B\}
$$
Si $(r_0, r_t) \in \mathcal{R}$, il est clair que
$\sum r_t x_t - r_0$ appartient au noyau.
Puisque $\Psi$ est un homomorphisme d'anneaux à puissances divisées
et que $\sum r_tx_t - r_0 \in J'$,
il s'ensuit que $\delta_n(\sum r_t x_t - r_0)$ appartient aussi au noyau.
Soit $K \subset B\langle x_t \rangle$ l'idéal engendré par
$x_t - f_t$ et les éléments $\delta_n(\sum r_t x_t - r_0)$ pour
$(r_0, r_t) \in \mathcal{R}$.
Pour montrer que $K = \Ker(\Psi)$, il suffit de montrer que
$\delta$ se prolonge à $B\langle x_t \rangle/K$. En effet, la propriété
universelle de $D$ fournit alors une application $D \to B\langle x_t \rangle/K$
inverse de $\Psi$. Nous devons donc montrer que $K \cap J'$ est
stable par $\delta_n$ ; voir le
lemme \ref{dpa-lemma-kernel} d'Algèbre à puissances divisées.
Soit $K' \subset B\langle x_t \rangle$ l'idéal
engendré par les éléments suivants :
\begin{enumerate}
\item $\delta_m(\sum r_t x_t - r_0)$, où $m > 0$ et
$(r_0, r_t) \in \mathcal{R}$,
\item $x_{t'}^{[m]}(x_t - f_t)$, où $m > 0$ et $t', t \in T$.
\end{enumerate}
Nous affirmons que $K' = K \cap J'$. Cette affirmation prouve que $K \cap J'$
est stable par $\delta_n$, $n > 0$, d'après le critère du
lemme \ref{dpa-lemma-kernel} (2)(c) d'Algèbre à puissances divisées
et le calcul de $\delta_n$ sur
les éléments énumérés, que nous laissons au lecteur.
Pour démontrer l'affirmation, notons que $K' \subset K \cap J'$.
Réciproquement, si $h \in K \cap J'$, nous pouvons écrire modulo $K'$
$$
h = \sum r_t (x_t - f_t)
$$
pour certains $r_t \in B$. Puisque $h \in K \cap J' \subset J'$,
nous avons $r_0 = \sum r_t f_t \in I$. Ainsi $(r_0, r_t) \in \mathcal{R}$,
et nous constatons que
$$
h = \sum r_t x_t - r_0
$$
appartient à $K'$, comme voulu.
\end{proof}

\begin{lemma}
\label{crystalline-lemma-divided-power-envelope-add-variables}
Soit $(A, I, \gamma)$ un anneau à puissances divisées.
Soit $B$ une $A$-algèbre et $IB \subset J \subset B$ un idéal.
Soit $x_i$ une famille d'indéterminées. Alors
$$
D_{B[x_i], \gamma}(JB[x_i] + (x_i)) = D_{B, \gamma}(J) \langle x_i \rangle
$$
\end{lemma}

\begin{proof}
Une démonstration possible consiste à déduire ceci du
lemme \ref{crystalline-lemma-describe-divided-power-envelope},
car toute relation entre les $x_i$ dans $B[x_i]$ est triviale.
D'autre part, le lemme résulte de la propriété universelle
de l'algèbre polynomiale à puissances divisées et de celle des
enveloppes à puissances divisées.
\end{proof}

\noindent
Les conditions (1) et (2) du lemme suivant sont satisfaites si $B \to B'$ est plat
en tout point de $V(IB') \subset \Spec(B')$ et sont très étroitement liées
à cette condition ; voir
Algèbre, lemme \ref{algebra-lemma-what-does-it-mean}.
Il affirme notamment que la formation de l'enveloppe à
puissances divisées commute à la localisation.

\begin{lemma}
\label{crystalline-lemma-flat-base-change-divided-power-envelope}
Soit $(A, I, \gamma)$ un anneau à puissances divisées.
Soit $B \to B'$ un homomorphisme de $A$-algèbres.
Supposons que
\begin{enumerate}
\item $B/IB \to B'/IB'$ soit plat, et
\item $\text{Tor}_1^B(B', B/IB) = 0$.
\end{enumerate}
Alors, pour tout idéal $IB \subset J \subset B$, l'application canonique
$$
D_B(J) \otimes_B B' \longrightarrow D_{B'}(JB')
$$
est un isomorphisme.
\end{lemma}

\begin{proof}
Posons $D = D_B(J)$ et notons $\bar J \subset D$ son idéal à puissances divisées,
muni de la structure de puissances divisées $\bar\gamma$. La propriété universelle de
$D$ fournit un homomorphisme de $B$-algèbres $D \to D_{B'}(JB')$, donc l'application du
lemme. Il suffit de montrer que
les puissances divisées $\bar\gamma$ se prolongent à $D \otimes_B B'$, car alors
la propriété universelle de $D_{B'}(JB')$ fournira une application
$D_{B'}(JB') \to D \otimes_B B'$ inverse de celle du lemme.

\medskip\noindent
Choisissons une surjection $P \to B'$, où $P$ est une algèbre polynomiale sur $B$.
En particulier, $B \to P$ est plat, donc $D \to D \otimes_B P$ est plat d'après le
lemme \ref{algebra-lemma-flat-base-change} d'Algèbre.
Alors $\bar\gamma$ se prolonge à $D \otimes_B P$ d'après le
lemme \ref{dpa-lemma-gamma-extends} d'Algèbre à puissances divisées ; nous
noterons encore ce prolongement
$\bar\gamma$. Posons $\mathfrak a = \Ker(P \to B')$, de sorte que
nous avons la suite exacte courte
$$
0 \to \mathfrak a \to P \to B' \to 0
$$
Ainsi, $\text{Tor}_1^B(B', B/IB) = 0$ implique
$\mathfrak a \cap IP = I\mathfrak a$.
Nous avons maintenant le diagramme commutatif suivant :
$$
\xymatrix{
B/J \otimes_B \mathfrak a \ar[r]_\beta &
B/J \otimes_B P \ar[r] &
B/J \otimes_B B' \\
D \otimes_B \mathfrak a \ar[r]^\alpha \ar[u] &
D \otimes_B P \ar[r] \ar[u] &
D \otimes_B B' \ar[u] \\
\bar J \otimes_B \mathfrak a \ar[r] \ar[u] &
\bar J \otimes_B P \ar[r] \ar[u] &
\bar J \otimes_B B' \ar[u]
}
$$
Ce diagramme est exact, même après ajout de $0$ en haut et à droite.
Nous devons montrer que les puissances divisées sur l'idéal
$\bar J \otimes_B P$ préservent l'idéal
$\Im(\alpha) \cap \bar J \otimes_B P$ ; voir le
lemme \ref{dpa-lemma-kernel} d'Algèbre à puissances divisées.
Considérons la suite exacte
$$
0 \to \mathfrak a/I\mathfrak a \to P/IP \to B'/IB' \to 0
$$
(qui utilise $\mathfrak a \cap IP = I\mathfrak a$, comme on l'a vu ci-dessus).
Puisque $B'/IB'$ est plat sur $B/IB$, cette suite reste exacte après
application de $B/J \otimes_{B/IB} -$ ; voir le
lemme \ref{algebra-lemma-flat-tor-zero} d'Algèbre. Ainsi,
$$
\Ker(B/J \otimes_{B/IB} \mathfrak a/I\mathfrak a \to
B/J \otimes_{B/IB} P/IP) =
\Ker(\mathfrak a/J\mathfrak a \to P/JP)
$$
est nul. Ainsi $\beta$ est injectif. Il s'ensuit que
$\Im(\alpha) \cap \bar J \otimes_B P$ coïncide avec
l'image de $\bar J \otimes \mathfrak a$. Si maintenant
$f \in \bar J$ et $a \in \mathfrak a$, alors
$\bar\gamma_n(f \otimes a) = \bar\gamma_n(f) \otimes a^n$
et le résultat est donc clair.
\end{proof}

\noindent
Le lemme suivant est un cas particulier de
\cite[Proposition 2.1.7]{dJ-crystalline}, qui est elle-même une généralisation de
\cite[Proposition 2.8.2]{Berthelot}.

\begin{lemma}
\label{crystalline-lemma-flat-extension-divided-power-envelope}
Soit $(B, I, \gamma) \to (B', I', \gamma')$ un homomorphisme d'anneaux à
puissances divisées. Soient $I \subset J \subset B$ et
$I' \subset J' \subset B'$ des idéaux. Supposons
\begin{enumerate}
\item $B/I \to B'/I'$ soit plat, et
\item $J' = JB' + I'$.
\end{enumerate}
Alors l'application canonique
$$
D_{B, \gamma}(J) \otimes_B B' \longrightarrow D_{B', \gamma'}(J')
$$
est un isomorphisme.
\end{lemma}

\begin{proof}
Posons $D = D_{B, \gamma}(J)$. Choisissons des éléments $f_t \in J$ engendrant $J/I$.
Posons $\mathcal{R} = \{(r_0, r_t) \in I \oplus \bigoplus\nolimits_{t \in T} B
\mid \sum r_t f_t = r_0 \text{ dans }B\}$ comme dans la démonstration du
lemme \ref{crystalline-lemma-describe-divided-power-envelope}. Ce lemme montre que
$$
D = B\langle x_t \rangle/ K
$$
où $K$ est engendré par les éléments $x_t - f_t$ et
$\delta_n(\sum r_t x_t - r_0)$ pour $(r_0, r_t) \in \mathcal{R}$.
Nous voyons ainsi que
\begin{equation}
\label{crystalline-equation-base-change}
D \otimes_B B' = B'\langle x_t \rangle/K'
\end{equation}
où $K'$ est engendré par les images dans $B'\langle x_t \rangle$
des générateurs de $K$ énumérés ci-dessus. Soit $f'_t \in B'$ l'image
de $f_t$. Par l'hypothèse (1), les éléments $f'_t \in J'$
engendrent $J'/I'$, et $x_t - f'_t \in K'$. Posons
$$
\mathcal{R}' =
\{(r'_0, r'_t) \in I' \oplus \bigoplus\nolimits_{t \in T} B'
\mid \sum r'_t f'_t = r'_0 \text{ dans }B'\}
$$
Pour achever la démonstration, nous devons montrer que
$\delta'_n(\sum r'_t x_t - r'_0) \in K'$
pour $(r'_0, r'_t) \in \mathcal{R}'$, car alors la présentation
(\ref{crystalline-equation-base-change}) de $D \otimes_B B'$ est identique
à la présentation de $D_{B', \gamma'}(J')$ obtenue dans le
lemme \ref{crystalline-lemma-describe-divided-power-envelope} à partir des générateurs $f'_t$.
Supposons que $(r'_0, r'_t) \in \mathcal{R}'$. Alors
$\sum r'_t f'_t = 0$ dans $B'/I'$. Puisque $B/I \to B'/I'$ est plat par
l'hypothèse (1), nous pouvons appliquer le critère équationnel de platitude
(Algèbre, lemme \ref{algebra-lemma-flat-eq}) et constater
qu'il existe un entier $m > 0$ ainsi que des éléments
$r_{jt} \in B$ et $c_j \in B'$, $j = 1, \ldots, m$, tels
que
$$
r_{j0} = \sum\nolimits_t r_{jt} f_t  \in I \text{ pour } j = 1, \ldots, m
$$
et
$$
i'_t  = r'_t - \sum\nolimits_j c_j r_{jt} \in I' \text{ pour tout }t
$$
Notons que cela implique aussi
$r'_0 = \sum_t i'_t f_t + \sum_j c_j r_{j0}$.
Nous avons alors
\begin{align*}
\delta'_n(\sum\nolimits_t r'_t x_t  - r'_0)
& =
\delta'_n(
\sum\nolimits_t i'_t x_t +
\sum\nolimits_{t, j} c_j r_{jt} x_t -
\sum\nolimits_t i'_t f_t -
\sum\nolimits_j c_j r_{j0}) \\
& =
\delta'_n(
\sum\nolimits_t i'_t(x_t - f_t) +
\sum\nolimits_j c_j (\sum\nolimits_t r_{jt} x_t  - r_{j0}))
\end{align*}
Puisque $\delta_n(a + b) = \sum_{m = 0, \ldots, n} \delta_m(a) \delta_{n - m}(b)$
et que $\delta_m(\sum i'_t(x_t - f_t))$ appartient à l'idéal
engendré par les $x_t - f_t \in K'$ pour $m > 0$, il suffit de prouver que
$\delta_n(\sum c_j (\sum r_{jt} x_t  - r_{j0}))$ appartient à $K'$.
Pour cela, nous utilisons
$$
\delta_n(\sum\nolimits_j c_j (\sum\nolimits_t r_{jt} x_t  - r_{j0}))
=
\sum c_1^{n_1} \ldots c_m^{n_m}
\delta_{n_1}(\sum r_{1t} x_t  - r_{10}) \ldots
\delta_{n_m}(\sum r_{mt} x_t  - r_{m0})
$$
où la somme porte sur $n_1 + \ldots + n_m = n$. Cela démontre l'assertion voulue.
\end{proof}





\section{Quelques épaississements explicites à puissances divisées}
\label{crystalline-section-explicit-thickenings}

\noindent
Les constructions de cette section nous aideront à définir la connexion
sur un cristal en modules sur le site cristallin.

\begin{lemma}
\label{crystalline-lemma-divided-power-first-order-thickening}
Soit $(A, I, \gamma)$ un anneau à puissances divisées. Soit $M$ un $A$-module.
Soit $B = A \oplus M$ comme $A$-algèbre, où $M$ est un idéal de carré nul,
et posons $J = I \oplus M$. Posons
$$
\delta_n(x + z) = \gamma_n(x) + \gamma_{n - 1}(x)z
$$
pour $x \in I$ et $z \in M$.
Alors $\delta$ est une structure de puissances divisées et
$A \to B$ est un homomorphisme d'anneaux à puissances divisées de
$(A, I, \gamma)$ vers $(B, J, \delta)$.
\end{lemma}

\begin{proof}
Nous devons vérifier les conditions (1) -- (5) de la
définition \ref{dpa-definition-divided-powers} d'Algèbre à puissances divisées.
Nous le démontrerons directement dans ce cas ; la démonstration du
lemme suivant donne toutefois une méthode qui évite les calculs.
Les conditions (1) et (3) sont claires. La condition (2) résulte de
\begin{align*}
\delta_n(x + z)\delta_m(x + z)
& =
(\gamma_n(x) + \gamma_{n - 1}(x)z)(\gamma_m(x) + \gamma_{m - 1}(x)z) \\
& = \gamma_n(x)\gamma_m(x) + \gamma_n(x)\gamma_{m - 1}(x)z +
\gamma_{n - 1}(x)\gamma_m(x)z \\
& =
\frac{(n + m)!}{n!m!} \gamma_{n + m}(x) +
\left(\frac{(n + m - 1)!}{n!(m - 1)!} +
\frac{(n + m - 1)!}{(n - 1)!m!}\right)
\gamma_{n + m - 1}(x) z \\
& =
\frac{(n + m)!}{n!m!}\delta_{n + m}(x + z)
\end{align*}
La condition (5) résulte de
\begin{align*}
\delta_n(\delta_m(x + z))
& =
\delta_n(\gamma_m(x) + \gamma_{m - 1}(x)z) \\
& =
\gamma_n(\gamma_m(x)) + \gamma_{n - 1}(\gamma_m(x))\gamma_{m - 1}(x)z \\
& =
\frac{(nm)!}{n! (m!)^n} \gamma_{nm}(x) +
\frac{((n - 1)m)!}{(n - 1)! (m!)^{n - 1}}
\gamma_{(n - 1)m}(x) \gamma_{m - 1}(x) z \\
& = \frac{(nm)!}{n! (m!)^n}(\gamma_{nm}(x) + \gamma_{nm - 1}(x) z)
\end{align*}
par un calcul élémentaire d'arithmétique. Pour prouver (4), nous devons voir que
$$
\delta_n(x + x' + z + z')
=
\gamma_n(x + x') + \gamma_{n - 1}(x + x')(z + z')
$$
est égal à
$$
\sum\nolimits_{i = 0}^n
(\gamma_i(x) + \gamma_{i - 1}(x)z)
(\gamma_{n - i}(x') + \gamma_{n - i - 1}(x')z')
$$
Cela résulte aisément du regroupement des coefficients de
$1$, $z$ et $z'$, et de l'emploi de la condition (4) pour $\gamma$.
\end{proof}

\begin{lemma}
\label{crystalline-lemma-divided-power-second-order-thickening}
Soit $(A, I, \gamma)$ un anneau à puissances divisées. Soient $M$, $N$ des $A$-modules.
Soit $q : M \times M \to N$ une application $A$-bilinéaire.
Munissons $B = A \oplus M \oplus N$ de la structure de $A$-algèbre dont le produit est
$$
(x, z, w)\cdot (x', z', w') = (xx', xz' + x'z, xw' + x'w + q(z, z') + q(z', z))
$$
et posons $J = I \oplus M \oplus N$. Posons
$$
\delta_n(x, z, w) = (\gamma_n(x), \gamma_{n - 1}(x)z,
\gamma_{n - 1}(x)w + \gamma_{n - 2}(x)q(z, z))
$$
pour $(x, z, w) \in J$.
Alors $\delta$ est une structure de puissances divisées et
$A \to B$ est un homomorphisme d'anneaux à puissances divisées de
$(A, I, \gamma)$ vers $(B, J, \delta)$.
\end{lemma}

\begin{proof}
Supposons que nous voulions démontrer la propriété (4) de la
définition \ref{dpa-definition-divided-powers} d'Algèbre à puissances divisées.
Choisissons $(x, z, w)$ et $(x', z', w')$ dans $J$.
Choisissons une application
$$
A_0 = \mathbf{Z}\langle s, s'\rangle \longrightarrow A,\quad
s \longmapsto x,
s' \longmapsto x'
$$
ce qui est possible par la propriété universelle des algèbres
polynomiales à puissances divisées. Posons $M_0 = A_0 \oplus A_0$ et
$N_0 = A_0 \oplus A_0 \oplus M_0 \otimes_{A_0} M_0$.
Soit $q_0 : M_0 \times M_0 \to N_0$ l'application évidente.
Définissons $M_0 \to M$ comme l'application $A_0$-linéaire qui envoie
les vecteurs de base de $M_0$ sur $z$ et $z'$. Définissons $N_0 \to N$
comme l'application $A_0$-linéaire qui envoie les deux premiers vecteurs de base
de $N_0$ sur $w$ et $w'$ et qui utilise
$M_0 \otimes_{A_0} M_0 \to M \otimes_A M \xrightarrow{q} N$
sur le dernier facteur direct. Il suffit alors de démontrer
l'identité (4) dans la situation $(A_0, M_0, N_0, q_0)$.
On procède de même pour les autres identités. Nous sommes ainsi ramenés au cas
d'un anneau $A$ sans $\mathbf{Z}$-torsion et de modules sans torsion sur $A$.
Dans ce cas, il suffit de montrer que
$$
n! \delta_n(x, z, w) = (x, z, w)^n
$$
dans l'anneau $A$ ; voir Algèbre à puissances divisées,
lemme \ref{dpa-lemma-silly}. Pour le voir, notons que
$$
(x, z, w)^2 = (x^2, 2xz, 2xw + 2q(z, z))
$$
et, par récurrence,
$$
(x, z, w)^n = (x^n, nx^{n - 1}z, nx^{n - 1}w + n(n - 1)x^{n - 2}q(z, z))
$$
D'autre part,
$$
n! \delta_n(x, z, w) = (n!\gamma_n(x), n!\gamma_{n - 1}(x)z,
n!\gamma_{n - 1}(x)w + n!\gamma_{n - 2}(x) q(z, z))
$$
ce qui coïncide avec l'expression précédente. La démonstration est achevée.
\end{proof}




\section{Compatibilité}
\label{crystalline-section-compatibility}

\noindent
La lecture de cette section n'est pas indispensable ; elle explique comment notre étude
s'articule avec celle de \cite{Berthelot}.
Considérons la notion technique suivante.

\begin{definition}
\label{crystalline-definition-compatible}
Soient $(A, I, \gamma)$ et $(B, J, \delta)$ des anneaux à puissances divisées.
Soit $A \to B$ un homomorphisme d'anneaux. Nous disons que
{\it $\delta$ est compatible avec $\gamma$}
s'il existe une structure de puissances divisées $\bar\gamma$ sur
$J + IB$ telle que
$$
(A, I, \gamma) \to (B, J + IB, \bar \gamma)\quad\text{et}\quad
(B, J, \delta) \to (B, J + IB, \bar \gamma)
$$
soient des homomorphismes d'anneaux à puissances divisées.
\end{definition}

\noindent
Soit $p$ un nombre premier. Soit $(A, I, \gamma)$ un anneau à puissances divisées.
Soit $A \to C$ un homomorphisme d'anneaux tel que $p$ soit nilpotent dans $C$.
Supposons que $\gamma$ se prolonge à $IC$ (voir le
lemme \ref{dpa-lemma-gamma-extends} d'Algèbre à puissances divisées).
Dans cette situation, le (gros) site cristallin affine de
$\Spec(C)$ sur $\Spec(A)$,
tel qu'il est défini dans \cite{Berthelot},
est l'opposé de la catégorie des systèmes
$$
(B, J, \delta, A \to B, C \to B/J)
$$
où
\begin{enumerate}
\item $(B, J, \delta)$ est un anneau à puissances divisées tel que $p$ soit nilpotent dans $B$,
\item $\delta$ est compatible avec $\gamma$, et
\item le diagramme
$$
\xymatrix{
B \ar[r] & B/J \\
A \ar[u] \ar[r] & C \ar[u]
}
$$
est commutatif.
\end{enumerate}
Les conditions
« $\gamma$ se prolonge à $C$ et $\delta$ est compatible avec $\gamma$ »
sont utilisées dans \cite{Berthelot} afin d'assurer que
la cohomologie cristalline de $\Spec(C)$ est la même que celle de
$\Spec(C/IC)$. Nous éviterons cette difficulté
en travaillant exclusivement avec des $C$ tels que $IC = 0$\footnote{Il y aura bien sûr
un prix à payer.}. Dans ce cas,
pour un système $(B, J, \delta, A \to B, C \to B/J)$ comme ci-dessus,
la commutativité du diagramme affiché implique $IB \subset J$, et
la compatibilité équivaut à la condition que
$(A, I, \gamma) \to (B, J, \delta)$ soit un homomorphisme d'anneaux à
puissances divisées.




\section{Site cristallin affine}
\label{crystalline-section-affine-site}

\noindent
Dans cette section, nous étudions la variante algébrique du site cristallin.
La situation de base dans laquelle nous mènerons cette étude sera la
suivante.

\begin{situation}
\label{crystalline-situation-affine}
Ici, $p$ est un nombre premier, $(A, I, \gamma)$ est un anneau à puissances
divisées tel que $A$ soit une $\mathbf{Z}_{(p)}$-algèbre, et $A \to C$ est un
homomorphisme d'anneaux tel que $IC = 0$ et que $p$ soit nilpotent dans $C$.
\end{situation}

\noindent
En général, le nombre premier $p$ appartiendra à l'idéal
à puissances divisées $I$.

\begin{definition}
\label{crystalline-definition-affine-thickening}
Dans la situation \ref{crystalline-situation-affine} :
\begin{enumerate}
\item Un {\it épaississement à puissances divisées} de $C$ sur $(A, I, \gamma)$
est un homomorphisme d'algèbres à puissances divisées $(A, I, \gamma) \to (B, J, \delta)$
tel que $p$ soit nilpotent dans $B$, muni d'un homomorphisme d'anneaux $C \to B/J$ tel que
$$
\xymatrix{
B \ar[r] & B/J \\
& C \ar[u] \\
A \ar[uu] \ar[r] & A/I \ar[u]
}
$$
soit commutatif.
\item Un {\it homomorphisme d'épaississements à puissances divisées}
$$
(B, J, \delta, C \to B/J) \longrightarrow (B', J', \delta', C \to B'/J')
$$
est un homomorphisme $\varphi : B \to B'$ de $A$-algèbres à puissances divisées tel
que $C \to B/J \to B'/J'$ soit l'application donnée $C \to B'/J'$.
\item Nous notons $\text{CRIS}(C/A, I, \gamma)$, ou simplement $\text{CRIS}(C/A)$,
la catégorie des épaississements à puissances divisées de $C$ sur $(A, I, \gamma)$.
\item Nous notons $\text{Cris}(C/A, I, \gamma)$, ou simplement $\text{Cris}(C/A)$,
la sous-catégorie pleine formée des $(B, J, \delta, C \to B/J)$ tels que
$C \to B/J$ soit un isomorphisme. Nous notons souvent un tel objet
$(B \to C, \delta)$, en sous-entendant $J = \Ker(B \to C)$.
\end{enumerate}
\end{definition}

\noindent
Notons que, pour un épaississement à puissances divisées $(B, J, \delta)$ comme ci-dessus,
l'idéal $J$ est localement nilpotent ; voir le
lemme \ref{dpa-lemma-nil} d'Algèbre à puissances divisées.
Il existe un foncteur canonique
\begin{equation}
\label{crystalline-equation-forget-affine}
\text{CRIS}(C/A) \longrightarrow C\text{-algèbres},\quad
(B, J, \delta) \longmapsto B/J
\end{equation}
Cette catégorie n'admet en général ni égalisateurs ni produits fibrés.
Elle n'admet pas non plus, en général, d'objet initial ($=$ colimite vide).

\begin{lemma}
\label{crystalline-lemma-affine-thickenings-colimits}
Dans la situation \ref{crystalline-situation-affine}.
\begin{enumerate}
\item $\text{CRIS}(C/A)$ admet les produits finis (mais non les produits infinis),
\item $\text{CRIS}(C/A)$ admet toutes les colimites finies non vides et
(\ref{crystalline-equation-forget-affine}) commute à celles-ci, et
\item $\text{Cris}(C/A)$ admet toutes les colimites finies non vides et
$\text{Cris}(C/A) \to \text{CRIS}(C/A)$ commute à celles-ci.
\end{enumerate}
\end{lemma}

\begin{proof}
Le produit vide, c'est-à-dire l'objet final de la catégorie des épaississements à
puissances divisées de $C$ sur $(A, I, \gamma)$, est l'anneau nul considéré
comme une $A$-algèbre, muni de l'idéal nul et de l'unique structure de puissances divisées
sur cet idéal, puis de l'unique homomorphisme de $C$ vers
l'anneau nul. Si $(B_t, J_t, \delta_t)_{t \in T}$ est une famille d'objets de
$\text{CRIS}(C/A)$, nous pouvons former le produit
$(\prod_t B_t, \prod_t J_t, \prod_t \delta_t)$ comme dans le
lemme \ref{dpa-lemma-limits} d'Algèbre à puissances divisées.
L'application $C \to \prod B_t/\prod J_t = \prod B_t/J_t$ est évidente.
Toutefois, $p$ n'est nécessairement nilpotent dans $\prod_t B_t$
que si $T$ est fini.

\medskip\noindent
Étant donnés deux objets $(B, J, \gamma)$ et $(B', J', \gamma')$ de
$\text{CRIS}(C/A)$, nous pouvons former un diagramme cocartésien
$$
\xymatrix{
(B, J, \delta) \ar[r] & (B'', J'', \delta'') \\
(A, I, \gamma) \ar[r] \ar[u] & (B', J', \delta') \ar[u]
}
$$
dans la catégorie des anneaux à puissances divisées. Nous avons alors
$$
B''/J'' = B/J \otimes_{A/I} B'/J' \longleftarrow C \otimes_{A/I} C
$$
voir Algèbre à puissances divisées, remarque \ref{dpa-remark-pushout}.
Notons $J'' \subset K \subset B''$
l'idéal tel que
$$
\xymatrix{
B''/J'' \ar[r] & B''/K \\
C \otimes_{A/I} C \ar[r] \ar[u] & C \ar[u]
}
$$
soit cocartésien, c'est-à-dire $B''/K \cong B/J \otimes_C B'/J'$.
Soit $D_{B''}(K) = (D, \bar K, \bar \delta)$
l'enveloppe à puissances divisées de $K$ dans $B''$ relativement à
$(B'', J'', \delta'')$. On vérifie alors aisément que
$(D, \bar K, \bar \delta)$ est un coproduit de $(B, J, \delta)$ et
$(B', J', \delta')$ dans $\text{CRIS}(C/A)$.

\medskip\noindent
Passons maintenant aux coégalisateurs. Soient
$\alpha, \beta : (B, J, \delta) \to (B', J', \delta')$ des morphismes de
$\text{CRIS}(C/A)$. Considérons $B'' = B'/ (\alpha(b) - \beta(b))$. Soit
$J'' \subset B''$ l'image de $J'$. Soit
$D_{B''}(J'') = (D, \bar J, \bar\delta)$ l'enveloppe à puissances divisées de
$J''$ dans $B''$ relativement à $(B', J', \delta')$. On vérifie alors aisément
que $(D, \bar J, \bar \delta)$ est le coégalisateur de $(B, J, \delta)$ et
$(B', J', \delta')$ dans $\text{CRIS}(C/A)$.

\medskip\noindent
D'après Catégories, lemme \ref{categories-lemma-almost-finite-colimits-exist},
$\text{CRIS}(C/A)$ admet toutes les colimites finies non vides. Les constructions
ci-dessus montrent que (\ref{crystalline-equation-forget-affine}) commute à celles-ci.
Cela implique formellement (3), puisque $\text{Cris}(C/A)$ est la catégorie-fibre
de (\ref{crystalline-equation-forget-affine}) au-dessus de $C$.
\end{proof}

\begin{remark}
\label{crystalline-remark-completed-affine-site}
Dans la situation \ref{crystalline-situation-affine}, nous notons
$\text{Cris}^\wedge(C/A)$ la catégorie dont les objets sont
les couples $(B \to C, \delta)$ tels que
\begin{enumerate}
\item $B$ est une $A$-algèbre $p$-adiquement complète,
\item $B \to C$ est une surjection de $A$-algèbres,
\item $\delta$ est une structure de puissances divisées sur $\Ker(B \to C)$,
\item $A \to B$ est un homomorphisme d'anneaux à puissances divisées.
\end{enumerate}
Les morphismes sont définis comme dans la définition \ref{crystalline-definition-affine-thickening}.
Alors $\text{Cris}(C/A) \subset \text{Cris}^\wedge(C/A)$ est la sous-catégorie
pleine formée des $B$ tels que $p$ soit nilpotent dans $B$.
Réciproquement, tout objet $(B \to C, \delta)$ de $\text{Cris}^\wedge(C/A)$
est égal à la limite
$$
(B \to C, \delta) = \lim_e (B/p^eB \to C, \delta)
$$
où, pour $e \gg 0$, l'objet $(B/p^eB \to C, \delta)$ appartient
à $\text{Cris}(C/A)$ ; voir le
lemme \ref{dpa-lemma-extend-to-completion} d'Algèbre à puissances divisées.
En particulier, $\text{Cris}^\wedge(C/A)$ est une sous-catégorie pleine
de la catégorie des pro-objets de $\text{Cris}(C/A)$ ; voir
Catégories, remarque \ref{categories-remark-pro-category}.
\end{remark}

\begin{lemma}
\label{crystalline-lemma-list-properties}
Dans la situation \ref{crystalline-situation-affine}.
Soit $P \to C$ une surjection de $A$-algèbres de noyau $J$.
Écrivons $D_{P, \gamma}(J) = (D, \bar J, \bar\gamma)$.
Soit $(D^\wedge, J^\wedge, \bar\gamma^\wedge)$ le
complété $p$-adique de $D$ ; voir le
lemme \ref{dpa-lemma-extend-to-completion} d'Algèbre à puissances divisées.
Pour tout $e \geq 1$, posons $P_e = P/p^eP$ et désignons par $J_e \subset P_e$
l'image de $J$, puis écrivons
$D_{P_e, \gamma}(J_e) = (D_e, \bar J_e, \bar\gamma)$.
Alors, pour tout $e$ assez grand, nous avons :
\begin{enumerate}
\item $p^eD \subset \bar J$ et $p^eD^\wedge \subset \bar J^\wedge$
sont stables par les puissances divisées,
\item $D^\wedge/p^eD^\wedge = D/p^eD = D_e$ comme anneaux à puissances divisées,
\item $(D_e, \bar J_e, \bar\gamma)$ est un objet de $\text{Cris}(C/A)$,
\item $(D^\wedge, \bar J^\wedge, \bar\gamma^\wedge)$ est égal à
$\lim_e (D_e, \bar J_e, \bar\gamma)$, et
\item $(D^\wedge, \bar J^\wedge, \bar\gamma^\wedge)$ est un objet de
$\text{Cris}^\wedge(C/A)$.
\end{enumerate}
\end{lemma}

\begin{proof}
L'assertion (1) résulte du
lemme \ref{dpa-lemma-extend-to-completion} d'Algèbre à puissances divisées.
Une propriété générale de la complétion $p$-adique donne
$D/p^eD = D^\wedge/p^eD^\wedge$. Puisque $D/p^eD$ est un anneau à puissances divisées
et que $P \to D/p^eD$ se factorise par $P_e$, la propriété universelle de
$D_e$ fournit une application $D_e \to D/p^eD$. Réciproquement, la propriété universelle
de $D$ fournit une application $D \to D_e$ qui se factorise par $D/p^eD$. Nous omettons
la vérification que ces applications sont inverses l'une de l'autre. Cela prouve (2).
Si $e$ est assez grand, alors $p^eC = 0$, d'où (3).
L'assertion (4) résulte du
lemme \ref{dpa-lemma-extend-to-completion} d'Algèbre à puissances divisées.
L'assertion (5) est claire d'après les définitions.
\end{proof}

\begin{lemma}
\label{crystalline-lemma-set-generators}
Dans la situation \ref{crystalline-situation-affine}.
Soit $P$ une algèbre polynomiale sur $A$ et soit
$P \to C$ une surjection de $A$-algèbres de noyau $J$.
Avec $(D_e, \bar J_e, \bar\gamma)$ comme dans le lemme \ref{crystalline-lemma-list-properties},
pour tout objet $(B, J_B, \delta)$ de $\text{CRIS}(C/A)$, il
existe un $e$ et un morphisme $D_e \to B$ dans $\text{CRIS}(C/A)$.
\end{lemma}

\begin{proof}
Nous pouvons trouver un homomorphisme de $A$-algèbres $P \to B$
relevant l'application $C \to B/J_B$. D'après notre définition de
$\text{CRIS}(C/A)$, nous avons $p^eB = 0$ pour
un certain $e$ ; ainsi $P \to B$ se factorise sous la forme $P \to P_e \to B$.
La propriété universelle de l'enveloppe à puissances divisées montre alors
que $P_e \to B$ se factorise par $D_e$.
\end{proof}

\begin{lemma}
\label{crystalline-lemma-generator-completion}
Dans la situation \ref{crystalline-situation-affine}.
Soit $P$ une algèbre polynomiale sur $A$ et soit
$P \to C$ une surjection de $A$-algèbres de noyau $J$.
Soit $(D, \bar J, \bar\gamma)$ le complété $p$-adique de
$D_{P, \gamma}(J)$. Pour tout objet $(B \to C, \delta)$ de
$\text{Cris}^\wedge(C/A)$, il
existe un morphisme $D \to B$ dans $\text{Cris}^\wedge(C/A)$.
\end{lemma}

\begin{proof}
Nous pouvons trouver un homomorphisme de $A$-algèbres $P \to B$ compatible
avec les applications vers $C$. D'après notre définition de
$\text{Cris}^\wedge(C/A)$, l'application $P \to B$ se factorise sous la forme
$P \to D_{P, \gamma}(J) \to B$. Puisque $B$ est $p$-adiquement complet,
nous pouvons factoriser cette application par $D$.
\end{proof}



\section{Module des différentielles}
\label{crystalline-section-differentials}

\noindent
Dans cette section, nous développons une théorie des modules de différentielles
pour les anneaux à puissances divisées.

\begin{definition}
\label{crystalline-definition-derivation}
Soit $A$ un anneau. Soit $(B, J, \delta)$ un anneau à puissances divisées.
Soit $A \to B$ un homomorphisme d'anneaux. Soit $M$ un $B$-module.
Une {\it $A$-dérivation à puissances divisées} à valeurs dans $M$ est une application
$\theta : B \to M$ additive, qui annule les éléments
de $A$, vérifie la règle de Leibniz
$\theta(bb') = b\theta(b') + b'\theta(b)$ et satisfait
$$
\theta(\delta_n(x)) = \delta_{n - 1}(x)\theta(x)
$$
pour tout $n \geq 1$ et tout $x \in J$.
\end{definition}

\noindent
Dans la situation de la définition, comme pour les dérivations
usuelles, il existe une {\it $A$-dérivation universelle à puissances divisées}
$$
\text{d}_{B/A, \delta} : B \to \Omega_{B/A, \delta}
$$
telle que toute $A$-dérivation à puissances divisées $\theta : B \to M$ s'écrive
$\theta = \xi \circ d_{B/A, \delta}$ pour une unique application $B$-linéaire
$\xi : \Omega_{B/A, \delta} \to M$. Si $(A, I, \gamma) \to (B, J, \delta)$
est un homomorphisme d'anneaux à puissances divisées, nous pouvons oublier les
puissances divisées sur $A$ et considérer les dérivations à puissances divisées de
$B$ sur $A$. Voici quelques propriétés fondamentales du module universel
des différentielles (à puissances divisées).

\begin{lemma}
\label{crystalline-lemma-omega}
Soit $A$ un anneau. Soit $(B, J, \delta)$ un anneau à puissances divisées et
$A \to B$ un homomorphisme d'anneaux. 
\begin{enumerate}
\item Considérons $B[x]$ muni de l'idéal à puissances divisées $(JB[x], \delta')$,
où $\delta'$ est le prolongement de $\delta$ à $B[x]$. Alors
$$
\Omega_{B[x]/A, \delta'} =
\Omega_{B/A, \delta} \otimes_B B[x] \oplus B[x]\text{d}x.
$$
\item Considérons $B\langle x \rangle$ muni de l'idéal à puissances divisées
$(JB\langle x \rangle + B\langle x \rangle_{+}, \delta')$. Alors
$$
\Omega_{B\langle x\rangle/A, \delta'} =
\Omega_{B/A, \delta} \otimes_B B\langle x \rangle \oplus
B\langle x\rangle \text{d}x.
$$
\item Soit $K \subset J$ un idéal stable par $\delta_n$ pour
tout $n > 0$. Posons $B' = B/K$ et notons $\delta'$ la structure de
puissances divisées induite sur $J/K$. Alors $\Omega_{B'/A, \delta'}$ est le quotient
de $\Omega_{B/A, \delta} \otimes_B B'$ par le sous-$B'$-module engendré
par les $\text{d}k$, pour $k \in K$.
\end{enumerate}
\end{lemma}

\begin{proof}
Ces assertions se déduisent directement de la construction de $\Omega_{B/A, \delta}$
comme le $B$-module libre sur les éléments $\text{d}b$, modulo les relations
\begin{enumerate}
\item $\text{d}(b + b') = \text{d}b + \text{d}b'$, $b, b' \in B$,
\item $\text{d}a = 0$, $a \in A$,
\item $\text{d}(bb') = b \text{d}b' + b' \text{d}b$, $b, b' \in B$,
\item $\text{d}\delta_n(f) = \delta_{n - 1}(f)\text{d}f$, $f \in J$, $n > 1$.
\end{enumerate}
Notons que la dernière relation explique pourquoi nous obtenons « la même » réponse pour
l'algèbre polynomiale à puissances divisées et l'algèbre polynomiale usuelle :
dans le premier cas, $x$ appartient à l'idéal à puissances divisées, donc
$\text{d}x^{[n]} = x^{[n - 1]}\text{d}x$.
\end{proof}

\noindent
Soit $(A, I, \gamma)$ un anneau à puissances divisées. Dans ce cadre, la
bonne notion de puissances de $I$ est donnée par les puissances divisées
$$
I^{[n]} = \text{idéal engendré par }
\gamma_{e_1}(x_1) \ldots \gamma_{e_t}(x_t)
\text{ avec }\sum e_j \geq n\text{ et }x_j \in I.
$$
Bien entendu, $I^n \subset I^{[n]}$. Notons que $I^{[1]} = I$.
Nous posons parfois aussi $I^{[0]} = A$.

\begin{lemma}
\label{crystalline-lemma-diagonal-and-differentials}
Soit $(A, I, \gamma) \to (B, J, \delta)$ un homomorphisme
d'anneaux à puissances divisées. Soit $(B(1), J(1), \delta(1))$ le coproduit
de $(B, J, \delta)$ avec lui-même sur $(A, I, \gamma)$, c'est-à-dire
tel que
$$
\xymatrix{
(B, J, \delta) \ar[r] & (B(1), J(1), \delta(1)) \\
(A, I, \gamma) \ar[r] \ar[u] & (B, J, \delta) \ar[u]
}
$$
est cocartésien. Notons $K = \Ker(B(1) \to B)$.
Alors $K \cap J(1) \subset J(1)$ est stable par la structure de
puissances divisées et
$$
\Omega_{B/A, \delta} = K/ \left(K^2 + (K \cap J(1))^{[2]}\right)
$$
canoniquement.
\end{lemma}

\begin{proof}
Le fait que $K \cap J(1) \subset J(1)$ soit stable par la structure de
puissances divisées résulte de ce que $B(1) \to B$ est un homomorphisme
d'anneaux à puissances divisées.

\medskip\noindent
Rappelons que $K/K^2$ possède une structure canonique de $B$-module.
Notons $s_0, s_1 : B \to B(1)$ les deux coprojections et considérons
l'application $\text{d} : B \to K/(K^2 +(K \cap J(1))^{[2]})$ donnée par
$b \mapsto s_1(b) - s_0(b)$. Il est clair que $\text{d}$ est additive,
s'annule sur $A$ et satisfait à la règle de Leibniz.
Nous affirmons que $\text{d}$ est une $A$-dérivation à puissances divisées.
Soit $x \in J$. Posons $y = s_1(x)$ et $z = s_0(x)$.
Nous écrirons $\delta$ au lieu de $\delta(1)$ pour
la structure de puissances divisées sur $J(1)$.
Il faut montrer que $\delta_n(y) - \delta_n(z) = \delta_{n - 1}(y)(y - z)$
modulo $K^2 +(K \cap J(1))^{[2]}$ pour $n \geq 1$.
L'égalité vaut pour $n = 1$. Supposons $n > 1$.
Notons que $\delta_i(y - z)$ appartient à $(K \cap J(1))^{[2]}$ si $i > 1$.
En calculant modulo $K^2 + (K \cap J(1))^{[2]}$, nous obtenons
$$
\delta_n(z) = \delta_n(z - y + y) =
\sum\nolimits_{i = 0}^n \delta_i(z - y)\delta_{n - i}(y) =
\delta_{n - 1}(y) \delta_1(z - y) + \delta_n(y)
$$
Cela démontre l'égalité voulue.

\medskip\noindent
Soit $M$ un $B$-module. Soit $\theta : B \to M$ une
$A$-dérivation à puissances divisées.
Posons $D = B \oplus M$, où $M$ est un idéal de carré nul. Définissons une
structure de puissances divisées sur $J \oplus M \subset D$ en posant
$\delta_n(x + m) = \delta_n(x) + \delta_{n - 1}(x)m$ pour $n > 1$ ; voir le
lemme \ref{crystalline-lemma-divided-power-first-order-thickening}.
Il existe deux homomorphismes d'algèbres à puissances divisées $B \to D$ : le premier
est l'inclusion et le second est l'application $b \mapsto b + \theta(b)$.
On obtient donc un homomorphisme canonique $B(1) \to D$ d'algèbres à
puissances divisées sur $(A, I, \gamma)$. Il induit une application $K \to M$
qui s'annule sur $K^2$ (car $M$ est un idéal de carré nul) ainsi que sur
$(K \cap J(1))^{[2]}$, puisque $M^{[2]} = 0$. Par construction, le composé
$B \to K/K^2 + (K \cap J(1))^{[2]} \to M$ est égal à $\theta$.
Il s'ensuit que $\text{d}$
est une $A$-dérivation universelle à puissances divisées, ce qui achève la démonstration.
\end{proof}

\begin{remark}
\label{crystalline-remark-filtration-differentials}
Soit $A \to B$ un homomorphisme d'anneaux et soit $(J, \delta)$ une structure de
puissances divisées sur $B$. Le module universel $\Omega_{B/A, \delta}$
est muni d'une structure supplémentaire : le $B$-sous-module
$N$ de $\Omega_{B/A, \delta}$ engendré par $\text{d}_{B/A, \delta}(J)$.
Au moyen de l'isomorphisme donné dans le
lemme \ref{crystalline-lemma-diagonal-and-differentials},
celui-ci correspond à l'image de
$K \cap J(1)$ dans $\Omega_{B/A, \delta}$. Considérons l'$A$-algèbre
$D = B \oplus \Omega^1_{B/A, \delta}$, munie de l'idéal $\bar J = J \oplus N$
et des puissances divisées $\bar \delta$ définies comme dans la démonstration du lemme.
Alors $(D, \bar J, \bar \delta)$ est un anneau à puissances divisées
et les deux applications $B \to D$ données par $b \mapsto b$ et
$b \mapsto b + \text{d}_{B/A, \delta}(b)$
sont des homomorphismes d'anneaux à puissances divisées sur $A$. De plus, $N$
est le plus petit sous-module de $\Omega_{B/A, \delta}$ pour lequel il en soit ainsi.
\end{remark}

\begin{lemma}
\label{crystalline-lemma-diagonal-and-differentials-affine-site}
Dans la situation \ref{crystalline-situation-affine}.
Soit $(B, J, \delta)$ un objet de $\text{CRIS}(C/A)$.
Soit $(B(1), J(1), \delta(1))$ le coproduit de $(B, J, \delta)$
avec lui-même dans $\text{CRIS}(C/A)$. Notons
$K = \Ker(B(1) \to B)$. Alors $K \cap J(1) \subset J(1)$
est stable par la structure de puissances divisées et
$$
\Omega_{B/A, \delta} = K/ \left(K^2 + (K \cap J(1))^{[2]}\right)
$$
canoniquement.
\end{lemma}

\begin{proof}
La démonstration est mot pour mot celle du
lemme \ref{crystalline-lemma-diagonal-and-differentials}.
Le seul point à vérifier est que l'anneau à
puissances divisées $D = B \oplus M$ est un objet de $\text{CRIS}(C/A)$
et que les deux applications $B \to D$ sont des morphismes de $\text{CRIS}(C/A)$.
Comme $D/(J \oplus M) = B/J$, l'application $C \to B/J$ permet de considérer
$D$ comme un objet de $\text{CRIS}(C/A)$,
et l'assertion concernant les morphismes résulte immédiatement de la construction.
\end{proof}

\begin{lemma}
\label{crystalline-lemma-module-differentials-divided-power-envelope}
Soit $(A, I, \gamma)$ un anneau à puissances divisées. Soit $A \to B$ un homomorphisme
d'anneaux et soit $IB \subset J \subset B$ un idéal. Soit
$D_{B, \gamma}(J) = (D, \bar J, \bar \gamma)$ l'enveloppe à puissances divisées.
Alors on a
$$
\Omega_{D/A, \bar\gamma} = \Omega_{B/A} \otimes_B D
$$
\end{lemma}

\begin{proof}[Première démonstration]
Soit $M$ un $D$-module. Nous affirmons que la donnée d'une $A$-dérivation
$\vartheta : B \to M$ équivaut à celle d'une
$A$-dérivation à puissances divisées $\theta : D \to M$. Le lemme de Yoneda
montre que cette assertion implique l'énoncé.

\medskip\noindent
Considérons l'épaississement de carré nul $D \oplus M$ de $D$.
Il existe une structure de puissances divisées $\delta$ sur $\bar J \oplus M$
si l'on annule sur $M$ les opérations supérieures de puissances divisées.
Autrement dit, on pose
$\delta_n(x + m) = \bar\gamma_n(x) + \bar\gamma_{n - 1}(x)m$ pour
tout $x \in \bar J$ et tout $m \in M$ ; voir le
lemme \ref{crystalline-lemma-divided-power-first-order-thickening}.
Considérons l'homomorphisme de $A$-algèbres $B \to D \oplus M$ dont la première
composante est l'application $B \to D$ et dont la seconde composante
est $\vartheta$. La propriété universelle fournit un
homomorphisme correspondant $D \to D \oplus M$ d'algèbres à puissances divisées,
dont la seconde composante est la
$A$-dérivation à puissances divisées $\theta$ correspondant à $\vartheta$.
\end{proof}

\begin{proof}[Deuxième démonstration]
Démontrons d'abord l'assertion lorsque $B$ est plat sur $A$. Dans ce cas, $\gamma$
se prolonge en une structure de puissances divisées $\gamma'$ sur $IB$ ; voir
Algèbre à puissances divisées, lemme \ref{dpa-lemma-gamma-extends}.
Ainsi $D = D_{B, \gamma'}(J)$ est un quotient de
l'anneau à puissances divisées $(D', J', \delta)$, où $D' =  B\langle x_t \rangle$
et $J' = IB\langle x_t \rangle + B\langle x_t \rangle_{+}$,
par les éléments $x_t - f_t$ et $\delta_n(\sum r_t x_t - r_0)$ ; voir le
lemme \ref{crystalline-lemma-describe-divided-power-envelope} pour les notations
et les explications. Notons $\text{d} : D' \to \Omega_{D'/A, \delta}$
la dérivation universelle. On a
$$
\Omega_{D'/A, \delta} =
\Omega_{B/A} \otimes_B D' \oplus \bigoplus D' \text{d}x_t,
$$
voir le lemme \ref{crystalline-lemma-omega}. On en déduit que $\Omega_{D/A, \bar\gamma}$
est le quotient de $\Omega_{D'/A, \delta} \otimes_{D'} D$ par le sous-module
engendré par les images sous $\text{d}$ des générateurs du
noyau de $D' \to D$ énumérés ci-dessus ; voir le lemme \ref{crystalline-lemma-omega}.
Comme $\text{d}(x_t - f_t) = - \text{d}f_t + \text{d}x_t$,
on a $\text{d}x_t = \text{d}f_t$ dans le quotient.
En particulier, $\Omega_{B/A} \otimes_B D \to \Omega_{D/A, \gamma}$
est surjectif, de noyau engendré par les images sous $\text{d}$
des éléments $\delta_n(\sum r_t x_t - r_0)$.
Or, étant donnée une relation $\sum r_tf_t - r_0 = 0$ dans $B$, avec
$r_t \in B$ et $r_0 \in IB$, on a
\begin{align*}
\text{d}\delta_n(\sum r_t x_t - r_0)
& =
\delta_{n - 1}(\sum r_t x_t - r_0)\text{d}(\sum r_t x_t - r_0)
\\
& =
\delta_{n - 1}(\sum r_t x_t - r_0)
\left(
\sum r_t\text{d}(x_t - f_t) + \sum (x_t - f_t)\text{d}r_t
\right)
\end{align*}
car $\sum r_tf_t - r_0 = 0$ dans $B$. Cette expression est donc déjà nulle dans
$\Omega_{B/A} \otimes_A D$, ce qui achève la démonstration lorsque $B$ est plat sur $A$.

\medskip\noindent
Dans le cas général, écrivons $B$ comme quotient d'un anneau de polynômes
$P \to B$ et soit $J' \subset P$ l'image réciproque de $J$. Alors
$D = D'/K'$, avec les notations du
lemme \ref{crystalline-lemma-divided-power-envelop-quotient}.
D'après le cas traité au premier paragraphe de la démonstration, on a
$\Omega_{D'/A, \bar\gamma'} = \Omega_{P/A} \otimes_P D'$. Alors
$\Omega_{D/A, \bar \gamma}$ est le quotient de $\Omega_{P/A} \otimes_P D$
par le sous-module engendré par les $\text{d}\bar\gamma_n'(k)$, où $k$
parcourt le noyau de $P \to B$ ; voir le
lemme \ref{crystalline-lemma-omega} et la description de $K'$ donnée au
lemme \ref{crystalline-lemma-divided-power-envelop-quotient}. Comme
$\text{d}\bar\gamma_n'(k) = \bar\gamma'_{n - 1}(k)\text{d}k$, on voit
qu'il suffit encore de quotienter par le sous-module engendré par les
$\text{d}k$ avec $k \in \Ker(P \to B)$ ; puisque $\Omega_{B/A}$
est le quotient de $\Omega_{P/A} \otimes_A B$ par ces éléments
(Algèbre, lemme \ref{algebra-lemma-differential-seq}), ce qui achève la démonstration.
\end{proof}

\begin{remark}
\label{crystalline-remark-divided-powers-de-rham-complex}
Soit $A \to B$ un homomorphisme d'anneaux et soit $(J, \delta)$ une structure de
puissances divisées sur $B$. Posons
$\Omega_{B/A, \delta}^i = \wedge^i_B \Omega_{B/A, \delta}$
où $\Omega_{B/A, \delta}$ est le but de la dérivation universelle à puissances divisées
sur $A$, $\text{d} = \text{d}_{B/A} : B \to \Omega_{B/A, \delta}$.
Notons que $\Omega_{B/A, \delta}$ est le quotient de $\Omega_{B/A}$ par le
$B$-sous-module engendré par les éléments
$\text{d}\delta_n(x) - \delta_{n - 1}(x)\text{d}x$ pour $x \in J$.
Nous affirmons que le lemme \ref{algebra-lemma-de-rham-complex} d'Algèbre s'applique.
Pour le voir, il suffit de vérifier que les éléments
$\text{d}\delta_n(x) - \delta_{n - 1}(x)\text{d}x$
de $\Omega_B$ ont une image nulle dans $\Omega^2_{B/A, \delta}$.
On observe que
$$
\text{d}(\delta_{n - 1}(x)) \wedge \text{d}x
= \delta_{n - 2}(x) \text{d}x \wedge \text{d}x = 0
$$
dans $\Omega^2_{B/A, \delta}$, ce qui donne l'annulation voulue. On obtient donc un
{\it complexe de de Rham à puissances divisées}
$$
\Omega^0_{B/A, \delta} \to \Omega^1_{B/A, \delta} \to
\Omega^2_{B/A, \delta} \to \ldots
$$
qui jouera un rôle important dans la suite.
\end{remark}

\begin{remark}
\label{crystalline-remark-connection}
Soit $A \to B$ un homomorphisme d'anneaux. Soit $\Omega_{B/A} \to \Omega$
un quotient satisfaisant aux hypothèses du
lemme \ref{algebra-lemma-de-rham-complex} d'Algèbre.
Soit $M$ un $B$-module. Une {\it connexion} est une application additive
$$
\nabla : M \longrightarrow M \otimes_B \Omega
$$
telle que $\nabla(bm) = b \nabla(m) + m \otimes \text{d}b$
pour $b \in B$ et $m \in M$. Dans cette situation, on peut définir des applications
$$
\nabla : M \otimes_B \Omega^i \longrightarrow M \otimes_B \Omega^{i + 1}
$$
par la règle $\nabla(m \otimes \omega) = \nabla(m) \wedge \omega +
m \otimes \text{d}\omega$. Celle-ci est bien définie car, si $b \in B$, alors
\begin{align*}
\nabla(bm \otimes \omega) - \nabla(m \otimes b\omega)
& =
\nabla(bm) \wedge \omega + bm \otimes \text{d}\omega
- \nabla(m) \wedge b\omega - m \otimes \text{d}(b\omega) \\
& =
b\nabla(m) \wedge \omega + m \otimes \text{d}b \wedge \omega
+ bm \otimes \text{d}\omega \\
& \ \ \ \ \ \ - b\nabla(m) \wedge \omega - bm \otimes \text{d}(\omega)
- m \otimes \text{d}b \wedge \omega = 0
\end{align*}
Comme il est d'usage, on dit que la connexion est {\it intégrable} si et
seulement si le composé
$$
M \xrightarrow{\nabla} M \otimes_B \Omega^1
\xrightarrow{\nabla} M \otimes_B \Omega^2
$$
est nul. Dans ce cas, on obtient un complexe
$$
M \xrightarrow{\nabla} M \otimes_B \Omega^1
\xrightarrow{\nabla} M \otimes_B \Omega^2
\xrightarrow{\nabla} M \otimes_B \Omega^3
\xrightarrow{\nabla} M \otimes_B \Omega^4 \to \ldots
$$
appelé complexe de de Rham de la connexion.
\end{remark}

\begin{remark}
\label{crystalline-remark-base-change-connection}
Considérons un diagramme commutatif d'anneaux
$$
\xymatrix{
B \ar[r]_\varphi & B' \\
A \ar[u] \ar[r] & A' \ar[u]
}
$$
Soient $\Omega_{B/A} \to \Omega$ et $\Omega_{B'/A'} \to \Omega'$
des quotients satisfaisant aux hypothèses du
lemme \ref{algebra-lemma-de-rham-complex} d'Algèbre.
Supposons donnée une application $\varphi : \Omega \to \Omega'$ qui
s'insère dans un diagramme commutatif
$$
\xymatrix{
\Omega_{B/A} \ar[r] \ar[d] &
\Omega_{B'/A'} \ar[d] \\
\Omega \ar[r]^{\varphi} &
\Omega'
}
$$
où la flèche horizontale supérieure est l'application canonique
$\Omega_{B/A} \to \Omega_{B'/A'}$ induite par $\varphi : B \to B'$.
Dans cette situation, pour toute paire $(M, \nabla)$ où $M$ est un $B$-module
et $\nabla : M \to M \otimes_B \Omega$ une connexion,
on obtient par {\it changement de base} la paire $(M \otimes_B B', \nabla')$, où
$$
\nabla' :
M \otimes_B B'
\longrightarrow
(M \otimes_B B') \otimes_{B'} \Omega' = M \otimes_B \Omega'
$$
est définie par la règle
$$
\nabla'(m \otimes b') =
\sum m_i \otimes b'\text{d}\varphi(b_i) + m \otimes \text{d}b' 
$$
si $\nabla(m) = \sum m_i \otimes \text{d}b_i$. Si $\nabla$ est intégrable,
alors $\nabla'$ l'est aussi et il existe dans ce cas une application canonique de
complexes de de Rham (remarque \ref{crystalline-remark-connection})
\begin{equation}
\label{crystalline-equation-base-change-map-complexes}
M \otimes_B \Omega^\bullet
\longrightarrow
(M \otimes_B B') \otimes_{B'} (\Omega')^\bullet =
M \otimes_B (\Omega')^\bullet
\end{equation}
qui envoie $m \otimes \eta$ sur $m \otimes \varphi(\eta)$.
\end{remark}

\begin{lemma}
\label{crystalline-lemma-differentials-completion}
Soit $A \to B$ un homomorphisme d'anneaux et soit $(J, \delta)$ une structure de
puissances divisées sur $B$. Soit $p$ un nombre premier. Supposons que $A$ soit une
$\mathbf{Z}_{(p)}$-algèbre et que $p$ soit nilpotent dans $B/J$. Alors
on a
$$
\lim_e \Omega_{B_e/A, \bar\delta} =
\lim_e \Omega_{B/A, \delta}/p^e\Omega_{B/A, \delta} =
\lim_e \Omega_{B^\wedge/A, \delta^\wedge}/p^e \Omega_{B^\wedge/A, \delta^\wedge}
$$
Pour les notations et les explications, voir la démonstration.
\end{lemma}

\begin{proof}
D'après Algèbre à puissances divisées, lemme \ref{dpa-lemma-extend-to-completion},
la structure $\delta$ se prolonge
à $B_e = B/p^eB$ pour tout $e$ assez grand. La première limite est donc
bien définie. Le lemme fournit aussi une structure de puissances divisées $\delta^\wedge$
sur le complété $B^\wedge = \lim_e B_e$ ; la dernière limite est donc
elle aussi bien définie. D'après le lemme \ref{crystalline-lemma-omega}
et le fait que $\text{d}p^e = 0$ (toujours),
on voit que la surjection
$\Omega_{B/A, \delta} \to \Omega_{B_e/A, \bar\delta}$ a pour noyau
$p^e\Omega_{B/A, \delta}$. Il en va de même du noyau de
$\Omega_{B^\wedge/A, \delta^\wedge} \to \Omega_{B_e/A, \bar\delta}$.
Le lemme en résulte.
\end{proof}



\section{Schémas à puissances divisées}
\label{crystalline-section-divided-power-schemes}

\noindent
Quelques remarques sur la globalisation des notions précédentes.

\begin{definition}
\label{crystalline-definition-divided-power-structure}
Soit $\mathcal{C}$ un site. Soit $\mathcal{O}$ un faisceau d'anneaux
sur $\mathcal{C}$. Soit $\mathcal{I} \subset \mathcal{O}$ un
faisceau d'idéaux. Une {\it structure de puissances divisées $\gamma$} sur $\mathcal{I}$
est une suite d'applications $\gamma_n : \mathcal{I} \to \mathcal{I}$, $n \geq 1$,
telle que, pour tout objet $U$ de $\mathcal{C}$, le triplet
$$
(\mathcal{O}(U), \mathcal{I}(U), \gamma)
$$
soit un anneau à puissances divisées.
\end{definition}

\noindent
Bien entendu, cela s'applique en particulier aux faisceaux d'anneaux sur
les espaces topologiques. Il est toutefois utile de se placer dans ce cadre plus général,
car le faisceau structural du site cristallin vit sur un... site !
Dans ce chapitre, un triplet $(\mathcal{C}, \mathcal{I}, \gamma)$ comme dans la
définition précédente est parfois appelé {\it topos à puissances divisées}.
Étant donnés un second triplet $(\mathcal{C}', \mathcal{I}', \gamma')$ et
un morphisme de topos annelés
$(f, f^\sharp) : (\Sh(\mathcal{C}), \mathcal{O}) \to
(\Sh(\mathcal{C}'), \mathcal{O}')$
on dit que $(f, f^\sharp)$ induit un {\it morphisme de topos à
puissances divisées} si $f^\sharp(f^{-1}\mathcal{I}') \subset \mathcal{I}$
et si les diagrammes
$$
\xymatrix{
f^{-1}\mathcal{I}' \ar[d]_{f^{-1}\gamma'_n} \ar[r]_{f^\sharp} &
\mathcal{I} \ar[d]^{\gamma_n} \\
f^{-1}\mathcal{I}' \ar[r]^{f^\sharp} & \mathcal{I}
}
$$
sont commutatifs pour tout $n \geq 1$. Si $f$ provient d'un morphisme de
sites induit par un foncteur $u : \mathcal{C}' \to \mathcal{C}$, cela signifie
simplement que
$$
(\mathcal{O}'(U'), \mathcal{I}'(U'), \gamma')
\longrightarrow
(\mathcal{O}(u(U')), \mathcal{I}(u(U')), \gamma)
$$
est un homomorphisme d'anneaux à puissances divisées pour tout $U' \in \Ob(\mathcal{C}')$.

\medskip\noindent
Dans le cas des schémas, on exige que l'idéal à puissances divisées soit
{\bf quasi-cohérent}. À cela près, la définition est exactement
la même que dans le cas des topos. La voici.

\begin{definition}
\label{crystalline-definition-divided-power-scheme}
Un {\it schéma à puissances divisées} est un triplet $(S, \mathcal{I}, \gamma)$
où $S$ est un schéma, $\mathcal{I}$ un faisceau quasi-cohérent
d'idéaux et $\gamma$ une structure de puissances divisées sur $\mathcal{I}$.
Un {\it morphisme de schémas à puissances divisées}
$(S, \mathcal{I}, \gamma) \to (S', \mathcal{I}', \gamma')$ est
un morphisme de schémas $f : S \to S'$ tel que
$f^{-1}\mathcal{I}'\mathcal{O}_S \subset \mathcal{I}$ et que
$$
(\mathcal{O}_{S'}(U'), \mathcal{I}'(U'), \gamma')
\longrightarrow
(\mathcal{O}_S(f^{-1}U'), \mathcal{I}(f^{-1}U'), \gamma)
$$
soit un homomorphisme d'anneaux à puissances divisées pour tout ouvert $U' \subset S'$.
\end{definition}

\noindent
Rappelons qu'il existe une correspondance bijective entre les faisceaux
quasi-cohérents d'idéaux et les immersions fermées ; voir
Morphismes, section \ref{morphisms-section-closed-immersions}.
Ainsi, tout schéma à puissances divisées $(T, \mathcal{J}, \gamma)$
détermine une immersion fermée canonique $U \to T$ définie par $\mathcal{J}$.
Réciproquement, une immersion fermée $U \to T$ et une structure de puissances
divisées $\gamma$ sur le faisceau d'idéaux $\mathcal{J}$ associé
à $U \to T$ déterminent un schéma à puissances divisées $(T, \mathcal{J}, \gamma)$.
Dans nombre de situations, on ne considère de tels triplets
$(U, T, \gamma)$ que lorsque le morphisme $U \to T$ est un épaississement ; voir
Compléments sur les morphismes, définition \ref{more-morphisms-definition-thickening}.

\begin{definition}
\label{crystalline-definition-divided-power-thickening}
Un triplet $(U, T, \gamma)$ comme ci-dessus est appelé {\it épaississement à puissances divisées}
si $U \to T$ est un épaississement.
\end{definition}

\noindent
Les produits fibrés de schémas à puissances divisées existent lorsque l'un des
trois est un épaississement à puissances divisées. Voici un énoncé précis.

\begin{lemma}
\label{crystalline-lemma-fibre-product}
Soient $f : (T, \mathcal{J}, \delta) \to (S, \mathcal{I}, \gamma)$ et
$f' : (T', \mathcal{J}', \delta') \to (S, \mathcal{I}, \gamma)$
des morphismes de schémas à puissances divisées. Il existe un schéma à puissances divisées
$(T'', \mathcal{J}'', \delta'')$ et un diagramme cartésien
$$
\xymatrix{
T \ar[d]_f & T'' \ar[d] \ar[l] \\
S & T' \ar[l]_{f'}
}
$$
dans la catégorie des schémas à puissances divisées. Le morphisme
$T'' \to T \times_S T'$ est une immersion fermée et le morphisme
$T''_0 \to T_0 \times_{S_0} T'_0$ est un isomorphisme.
\end{lemma}

\begin{proof}
Esquisse. Notons que, par l'intermédiaire de $\Spec(-)$, les deux dernières assertions
sont compatibles avec ce que nous avons vu des coproduits d'algèbres à puissances divisées dans
Algèbre à puissances divisées, remarque \ref{dpa-remark-pushout}.
On construit donc $T''$ comme sous-schéma fermé de $T \times_S T'$
de la façon suivante : pour tous ouverts affines $U \subset S$, $V \subset T$ et
$V' \subset T'$ tels que $f(V), f'(V') \subset U$, on considère le sous-schéma
fermé de $V \times_U V'$ déterminé par la construction de la
remarque \ref{dpa-remark-pushout} d'Algèbre à puissances divisées. Comme les
schémas $V \times_U V'$ sont les membres d'un recouvrement ouvert
de $T \times_S T'$, on procède ainsi : (1) on montre que ces
sous-schémas fermés se recollent, (2) que les
structures de puissances divisées ainsi obtenues se recollent, et (3) que le résultat du recollement
est le produit fibré dans la catégorie des schémas à puissances divisées.
Pour établir (1), il suffit de montrer
que la construction du coproduit dans la catégorie des
anneaux à puissances divisées commute à la localisation (avec une formulation
appropriée) ; cela résulte du
lemme \ref{dpa-lemma-gamma-extends} d'Algèbre à puissances divisées, qui implique
que les puissances divisées se prolongent aux localisations (par platitude).
\end{proof}

\noindent
Faisons l'observation suivante. Supposons que $(U, T, \gamma)$
soit un schéma à puissances divisées. Supposons que $T$ soit un schéma sur $\mathbf{Z}_{(p)}$
et que $p$ soit localement nilpotent sur $U$. Alors
\begin{enumerate}
\item $p$ est localement nilpotent sur $T \Leftrightarrow U \to T$
est un épaississement (voir Algèbre à puissances divisées, lemme \ref{dpa-lemma-nil}), et
\item $p^e\mathcal{O}_T$ est, localement sur $T$, stable par $\gamma$
pour $e \gg 0$ (voir
Algèbre à puissances divisées, lemme \ref{dpa-lemma-extend-to-completion}).
\end{enumerate}
Cela suggère que l'on disposera de bons résultats sur les épaississements à
puissances divisées sous les hypothèses suivantes.

\begin{situation}
\label{crystalline-situation-global}
Ici, $p$ est un nombre premier et $(S, \mathcal{I}, \gamma)$ est un schéma à
puissances divisées sur $\mathbf{Z}_{(p)}$. On pose $S_0 = V(\mathcal{I}) \subset S$.
Enfin, $X \to S_0$ est un morphisme de schémas tel que $p$ soit
localement nilpotent sur $X$.
\end{situation}

\noindent
C'est dans cette situation que nous définirons le gros et le petit
site cristallin.









\section{Le gros site cristallin}
\label{crystalline-section-big-site}

\noindent
Définissons d'abord le gros site. Étant donné un schéma à puissances divisées
$(S, \mathcal{I}, \gamma)$, on dit que $(T, \mathcal{J}, \delta)$ est
un schéma à puissances divisées sur $(S, \mathcal{I}, \gamma)$ si
$T$ est muni d'un morphisme $T \to S$ de schémas à puissances
divisées. De même, on dit qu'un épaississement à puissances divisées $(U, T, \delta)$
est un épaississement à puissances divisées sur $(S, \mathcal{I}, \gamma)$
si $T$ est muni d'un morphisme $T \to S$ de schémas à puissances
divisées.

\begin{definition}
\label{crystalline-definition-divided-power-thickening-X}
Plaçons-nous dans la situation \ref{crystalline-situation-global}.
\begin{enumerate}
\item Un {\it épaississement à puissances divisées de $X$ relativement à
$(S, \mathcal{I}, \gamma)$} est la donnée d'un épaississement à puissances divisées
$(U, T, \delta)$ sur $(S, \mathcal{I}, \gamma)$
et d'un $S$-morphisme $U \to X$.
\item Un {\it morphisme d'épaississements à puissances divisées de $X$
relativement à $(S, \mathcal{I}, \gamma)$} se définit de la manière
évidente.
\end{enumerate}
La catégorie des épaississements à puissances divisées de $X$ relativement à
$(S, \mathcal{I}, \gamma)$ est notée $\text{CRIS}(X/S, \mathcal{I}, \gamma)$
ou simplement $\text{CRIS}(X/S)$.
\end{definition}

\noindent
Pour tout $(U, T, \delta)$ dans $\text{CRIS}(X/S)$,
$p$ est localement nilpotent sur $T$ ; voir la discussion précédant la
situation \ref{crystalline-situation-global}.
On peut représenter commodément toutes les données associées à $(U, T, \delta)$
par le diagramme commutatif
$$
\xymatrix{
T \ar[dd] & U \ar[l] \ar[d] \\
& X \ar[d] \\
S & S_0 \ar[l]
}
$$
où $S_0 = V(\mathcal{I}) \subset S$. Les morphismes de $\text{CRIS}(X/S)$
se représentent de même par de grands diagrammes commutatifs. En particulier,
il existe un foncteur d'oubli canonique
\begin{equation}
\label{crystalline-equation-forget}
\text{CRIS}(X/S) \longrightarrow \Sch/X,\quad
(U, T, \delta) \longmapsto U
\end{equation}
ainsi que son inverse à droite (et adjoint à gauche)
\begin{equation}
\label{crystalline-equation-endow-trivial}
\Sch/X \longrightarrow \text{CRIS}(X/S),\quad
U \longmapsto (U, U, \emptyset)
\end{equation}
qui est parfois utile.

\begin{lemma}
\label{crystalline-lemma-divided-power-thickening-fibre-products}
Dans la situation \ref{crystalline-situation-global}.
La catégorie $\text{CRIS}(X/S)$ admet toutes les limites finies non vides,
en particulier les produits de deux objets et les produits fibrés.
Le foncteur (\ref{crystalline-equation-forget}) commute aux limites.
\end{lemma}

\begin{proof}
Omis. Indication : voir le lemme \ref{crystalline-lemma-affine-thickenings-colimits}
pour le cas affine. Voir également
Algèbre à puissances divisées, remarque \ref{dpa-remark-pushout}.
\end{proof}

\begin{lemma}
\label{crystalline-lemma-divided-power-thickening-base-change-flat}
Dans la situation \ref{crystalline-situation-global}. Soit
$$
\xymatrix{
(U_3, T_3, \delta_3) \ar[d] \ar[r] & (U_2, T_2, \delta_2) \ar[d] \\
(U_1, T_1, \delta_1) \ar[r] & (U, T, \delta)
}
$$
un carré cartésien dans la catégorie des épaississements à puissances divisées de
$X$ relativement à $(S, \mathcal{I}, \gamma)$. Si $T_2 \to T$ est
plat et $U_2 = T_2 \times_T U$, alors $T_3 = T_1 \times_T T_2$ (comme schémas).
\end{lemma}

\begin{proof}
Cela résulte de ce qu'une structure de puissances divisées se prolonge de manière unique
le long d'un homomorphisme d'anneaux plat. Voir
Algèbre à puissances divisées, lemme \ref{dpa-lemma-gamma-extends}.
\end{proof}

\noindent
Le lemme précédent signifie que le changement de base d'un morphisme plat
d'épaississements à puissances divisées est encore un morphisme plat et qu'il
s'agit en fait du changement de base « usuel » du morphisme. Il en résulte
que la définition suivante est bien définie.

\begin{definition}
\label{crystalline-definition-big-crystalline-site}
Dans la situation \ref{crystalline-situation-global}.
\begin{enumerate}
\item Une famille de morphismes $\{(U_i, T_i, \delta_i) \to (U, T, \delta)\}$
d'épaississements à puissances divisées de $X/S$ est un
{\it recouvrement de Zariski, étale, lisse, syntomique ou fppf}
si et seulement si
\begin{enumerate}
\item $U_i = U \times_T T_i$ pour tout $i$, et
\item $\{T_i \to T\}$ est un recouvrement de Zariski, étale, lisse, syntomique ou fppf.
\end{enumerate}
\item Le {\it gros site cristallin} de $X$ sur $(S, \mathcal{I}, \gamma)$
est la catégorie $\text{CRIS}(X/S)$ munie de la topologie de Zariski.
\item Le topos des faisceaux sur $\text{CRIS}(X/S)$ est noté
$(X/S)_{\text{CRIS}}$ ou parfois
$(X/S, \mathcal{I}, \gamma)_{\text{CRIS}}$\footnote{Cela entre en conflit avec
notre convention consistant à noter le topos associé à un site $\mathcal{C}$
par $\Sh(\mathcal{C})$.}.
\end{enumerate}
\end{definition}

\noindent
On dispose de quelques fonctorialités évidentes pour ces topos.

\begin{remark}[Fonctorialité]
\label{crystalline-remark-functoriality-big-cris}
Soit $p$ un nombre premier.
Soit $(S, \mathcal{I}, \gamma) \to (S', \mathcal{I}', \gamma')$ un
morphisme de schémas à puissances divisées sur $\mathbf{Z}_{(p)}$.
Posons $S_0 = V(\mathcal{I})$ et $S'_0 = V(\mathcal{I}')$.
Soit
$$
\xymatrix{
X \ar[r]_f \ar[d] & Y \ar[d] \\
S_0 \ar[r] & S'_0
}
$$
un diagramme commutatif de morphismes de schémas et supposons que $p$ soit
localement nilpotent sur $X$ et $Y$. On obtient alors un foncteur continu et
cocontinu
$$
\text{CRIS}(X/S) \longrightarrow \text{CRIS}(Y/S')
$$
en faisant correspondre à $(U, T, \delta)$ le triplet $(U, T, \delta)$
où $U \to X \to Y$ est le $S'$-morphisme de $U$ vers $Y$.
On obtient donc un morphisme de topos
$$
f_{\text{CRIS}} : (X/S)_{\text{CRIS}} \longrightarrow (Y/S')_{\text{CRIS}}
$$
voir Sites, section \ref{sites-section-cocontinuous-morphism-topoi}.
\end{remark}

\begin{remark}[Comparaison avec le site de Zariski]
\label{crystalline-remark-compare-big-zariski}
Dans la situation \ref{crystalline-situation-global}.
Le foncteur (\ref{crystalline-equation-forget}) est cocontinu (détails omis) et
commute aux produits et aux produits fibrés
(lemme \ref{crystalline-lemma-divided-power-thickening-fibre-products}).
On obtient donc un morphisme de topos
$$
U_{X/S} : (X/S)_{\text{CRIS}} \longrightarrow \Sh((\Sch/X)_{Zar})
$$
du gros topos cristallin de $X/S$ vers le gros topos de Zariski de $X$.
Voir Sites, section \ref{sites-section-cocontinuous-morphism-topoi}.
\end{remark}

\begin{remark}[Morphisme structural]
\label{crystalline-remark-big-structure-morphism}
Dans la situation \ref{crystalline-situation-global}.
Considérons le sous-schéma fermé $S_0 = V(\mathcal{I}) \subset S$.
Si l'on suppose que $p$ est localement nilpotent sur $S_0$ (ce qui est toujours
le cas en pratique), on obtient une situation comme dans la
définition \ref{crystalline-definition-divided-power-thickening-X}, avec $S_0$ au lieu
de $X$. On obtient donc un site $\text{CRIS}(S_0/S)$. Si $f : X \to S_0$ est
le morphisme structural de $X$ sur $S$, on obtient un diagramme commutatif
de morphismes de topos annelés
$$
\xymatrix{
(X/S)_{\text{CRIS}}
\ar[r]_{f_{\text{CRIS}}} \ar[d]_{U_{X/S}} &
(S_0/S)_{\text{CRIS}} \ar[d]^{U_{S_0/S}} \\
\Sh((\Sch/X)_{Zar}) \ar[r]^{f_{big}} & \Sh((\Sch/S_0)_{Zar}) \ar[rd] \\
& & \Sh((\Sch/S)_{Zar})
}
$$
d'après la remarque \ref{crystalline-remark-functoriality-big-cris}. On considère le composé
$(X/S)_{\text{CRIS}} \to \Sh((\Sch/S)_{Zar})$ comme le morphisme structural du
gros site cristallin. Même si $p$ n'est pas localement nilpotent sur $S_0$,
le morphisme structural
$$
(X/S)_{\text{CRIS}} \longrightarrow \Sh((\Sch/S)_{Zar})
$$
est défini, puisque l'on peut emprunter la voie inférieure du diagramme ci-dessus. C'est donc
le morphisme de topos correspondant au foncteur cocontinu
$\text{CRIS}(X/S) \to (\Sch/S)_{Zar}$ donné par la règle
$(U, T, \delta)/S \mapsto U/S$ ; voir
Sites, section \ref{sites-section-cocontinuous-morphism-topoi}.
\end{remark}

\begin{remark}[Compatibilités]
\label{crystalline-remark-compatibilities-big-cris}
Les morphismes définis ci-dessus satisfont à de nombreuses compatibilités. Par exemple,
dans la situation de la remarque \ref{crystalline-remark-functoriality-big-cris},
on obtient un diagramme commutatif de topos annelés
$$
\xymatrix{
(X/S)_{\text{CRIS}} \ar[d] \ar[r] & (Y/S')_{\text{CRIS}} \ar[d] \\
\Sh((\Sch/S)_{Zar}) \ar[r] & \Sh((\Sch/S')_{Zar})
}
$$
où les flèches verticales sont les morphismes structuraux.
\end{remark}




\section{Le site cristallin}
\label{crystalline-section-site}

\noindent
Comme (\ref{crystalline-equation-forget}) commute aux produits et aux produits
fibrés, les $(U, T, \delta)$ tels que
$U \to X$ soit une immersion ouverte forment une
sous-catégorie pleine stable par produits fibrés (et, plus généralement,
par limites finies non vides). La
définition suivante est donc bien définie.

\begin{definition}
\label{crystalline-definition-crystalline-site}
Dans la situation \ref{crystalline-situation-global}.
\begin{enumerate}
\item Le (petit) {\it site cristallin} de $X$ sur
$(S, \mathcal{I}, \gamma)$, noté $\text{Cris}(X/S, \mathcal{I}, \gamma)$
ou simplement $\text{Cris}(X/S)$, est la sous-catégorie pleine de $\text{CRIS}(X/S)$
formée des $(U, T, \delta)$ de $\text{CRIS}(X/S)$ tels que
$U \to X$ soit une immersion ouverte. Elle est munie de la topologie de Zariski.
\item Le topos des faisceaux sur $\text{Cris}(X/S)$ est noté
$(X/S)_{\text{cris}}$ ou parfois
$(X/S, \mathcal{I}, \gamma)_{\text{cris}}$\footnote{Cela entre en conflit avec
notre convention consistant à noter le topos associé à un site $\mathcal{C}$
par $\Sh(\mathcal{C})$.}.
\end{enumerate}
\end{definition}

\noindent
Pour tout $(U, T, \delta)$ dans $\text{Cris}(X/S)$, le morphisme $U \to X$
définit un objet du petit site de Zariski $X_{Zar}$ de $X$. D'où
On obtient donc un foncteur d'oubli canonique
\begin{equation}
\label{crystalline-equation-forget-small}
\text{Cris}(X/S) \longrightarrow X_{Zar},\quad
(U, T, \delta) \longmapsto U
\end{equation}
et un adjoint à gauche
\begin{equation}
\label{crystalline-equation-endow-trivial-small}
X_{Zar} \longrightarrow \text{Cris}(X/S),\quad
U \longmapsto (U, U, \emptyset)
\end{equation}
qui est parfois utile.

\medskip\noindent
On peut comparer le petit et le gros site cristallin, de même
que l'on compare le petit et le gros site de Zariski d'un schéma ; voir
Topologies, lemme \ref{topologies-lemma-at-the-bottom}.

\begin{lemma}
\label{crystalline-lemma-compare-big-small}
Sous les hypothèses de la définition \ref{crystalline-definition-divided-power-thickening-X}.
Le foncteur d'inclusion
$$
\text{Cris}(X/S) \to \text{CRIS}(X/S)
$$
commute aux limites finies non vides et est pleinement fidèle, continu
et cocontinu. Il existe des morphismes de topos
$$
(X/S)_{\text{cris}} \xrightarrow{i} (X/S)_{\text{CRIS}}
\xrightarrow{\pi} (X/S)_{\text{cris}}
$$
dont le composé est l'identité et dont le premier est induit
par le foncteur d'inclusion. De plus, $\pi_* = i^{-1}$.
\end{lemma}

\begin{proof}
Pour la première assertion, voir le
lemme \ref{crystalline-lemma-divided-power-thickening-fibre-products}.
Cela fournit un morphisme de topos
$i : (X/S)_{\text{cris}} \to (X/S)_{\text{CRIS}}$ et un adjoint à gauche
$i_!$ tel que $i^{-1}i_! = i^{-1}i_* = \text{id}$ ; voir
Sites, lemmes \ref{sites-lemma-when-shriek},
\ref{sites-lemma-preserve-equalizers} et
\ref{sites-lemma-back-and-forth}.
Nous affirmons que $i_!$ est exact. Si tel est le cas, on peut définir
$\pi$ par les formules $\pi^{-1} = i_!$ et $\pi_* = i^{-1}$,
et tout est clair. Pour démontrer l'assertion, rappelons que l'on sait déjà
que $i_!$ est exact à droite et préserve les produits fibrés (voir les références
indiquées). Il suffit donc de montrer que $i_! * = *$, où $*$ désigne
l'objet final de la catégorie des faisceaux d'ensembles.
Pour cela, il suffit de produire un ensemble d'objets
$(U_i, T_i, \delta_i)$, $i \in I$, de $\text{Cris}(X/S)$ tel que
$$
\coprod\nolimits_{i \in I} h_{(U_i, T_i, \delta_i)} \to *
$$
soit surjectif dans $(X/S)_{\text{CRIS}}$ (détails omis ; indication : utiliser le fait que
$\text{Cris}(X/S)$ admet des produits et que le foncteur
$\text{Cris}(X/S) \to \text{CRIS}(X/S)$ commute aux produits).
Dans le cas affine, cela
résulte du lemme \ref{crystalline-lemma-set-generators}. Nous omettons la démonstration
dans le cas général.
\end{proof}

\begin{remark}[Fonctorialité]
\label{crystalline-remark-functoriality-cris}
Soit $p$ un nombre premier.
Soit $(S, \mathcal{I}, \gamma) \to (S', \mathcal{I}', \gamma')$
un morphisme de schémas à puissances divisées sur $\mathbf{Z}_{(p)}$.
Soit
$$
\xymatrix{
X \ar[r]_f \ar[d] & Y \ar[d] \\
S_0 \ar[r] & S'_0
}
$$
un diagramme commutatif de morphismes de schémas et supposons que $p$ soit
localement nilpotent sur $X$ et $Y$. Par analogie avec
Topologies, lemme \ref{topologies-lemma-morphism-big-small}, on définit
$$
f_{\text{cris}} : (X/S)_{\text{cris}} \longrightarrow (Y/S')_{\text{cris}}
$$
par la formule $f_{\text{cris}} = \pi_Y \circ f_{\text{CRIS}} \circ i_X$,
où $i_X$ et $\pi_Y$ sont ceux du lemme \ref{crystalline-lemma-compare-big-small}
pour $X$ et $Y$, et où $f_{\text{CRIS}}$ est celui de la
remarque \ref{crystalline-remark-functoriality-big-cris}.
\end{remark}

\begin{remark}[Comparaison avec le site de Zariski]
\label{crystalline-remark-compare-zariski}
Dans la situation \ref{crystalline-situation-global}.
Le foncteur (\ref{crystalline-equation-forget-small}) est continu, cocontinu et
commute aux produits et aux produits fibrés.
On obtient donc un morphisme de topos
$$
u_{X/S} : (X/S)_{\text{cris}} \longrightarrow \Sh(X_{Zar})
$$
reliant le petit topos cristallin de $X/S$ au
petit topos de Zariski de $X$.
Voir Sites, section \ref{sites-section-cocontinuous-morphism-topoi}.
\end{remark}

\begin{lemma}
\label{crystalline-lemma-localize}
Dans la situation \ref{crystalline-situation-global}.
Soient $X' \subset X$ et $S' \subset S$ des sous-schémas ouverts tels que
l'image de $X'$ soit contenue dans $S'$. Il existe alors un foncteur pleinement fidèle
$\text{Cris}(X'/S') \to \text{Cris}(X/S)$
qui donne naissance à un morphisme de topos s'insérant dans le diagramme
commutatif
$$
\xymatrix{
(X'/S')_{\text{cris}} \ar[r] \ar[d]_{u_{X'/S'}} &
(X/S)_{\text{cris}} \ar[d]^{u_{X/S}} \\
\Sh(X'_{Zar}) \ar[r] & \Sh(X_{Zar})
}
$$
De plus, ce diagramme est un exemple de localisation de morphismes de
topos au sens de Sites, lemme \ref{sites-lemma-localize-morphism-topoi}.
\end{lemma}

\begin{proof}
Le foncteur pleinement fidèle consiste à considérer les
objets de $\text{Cris}(X'/S')$ comme des épaississements à puissances
divisées $(U, T, \delta)$ de $X$ pour lesquels $U \to X$
se factorise par $X' \subset X$ (alors $T \to S$ se
factorise automatiquement par $S'$). Ce foncteur est manifestement cocontinu,
d'où le morphisme de topos indiqué.
Soit $h_{X'} \in \Sh(X_{Zar})$ le faisceau représentable associé
à $X'$, considéré comme objet de $X_{Zar}$. Il est clair que
$\Sh(X'_{Zar})$ est le topos localisé $\Sh(X_{Zar})/h_{X'}$.
D'autre part, la catégorie $\text{Cris}(X/S)/u_{X/S}^{-1}h_{X'}$
(voir Sites, lemme \ref{sites-lemma-localize-topos-site})
s'identifie canoniquement à $\text{Cris}(X'/S')$ par le foncteur précédent.
Cela achève la démonstration.
\end{proof}

\begin{remark}[Morphisme structural]
\label{crystalline-remark-structure-morphism}
Dans la situation \ref{crystalline-situation-global}.
Considérons le sous-schéma fermé $S_0 = V(\mathcal{I}) \subset S$.
Si l'on suppose que $p$ est localement nilpotent sur $S_0$ (ce qui est toujours
le cas en pratique), on obtient une situation comme dans la
définition \ref{crystalline-definition-divided-power-thickening-X}, avec $S_0$ au lieu
de $X$. On obtient donc un site $\text{Cris}(S_0/S)$. Si $f : X \to S_0$
est le morphisme structural de $X$ sur $S$, on obtient un
diagramme commutatif de topos
$$
\xymatrix{
(X/S)_{\text{cris}}
\ar[r]_{f_{\text{cris}}} \ar[d]_{u_{X/S}} &
(S_0/S)_{\text{cris}} \ar[d]^{u_{S_0/S}} \\
\Sh(X_{Zar}) \ar[r]^{f_{small}} & \Sh(S_{0, Zar}) \ar[rd] \\
& & \Sh(S_{Zar})
}
$$
voir la remarque \ref{crystalline-remark-functoriality-cris}. On considère le composé
$(X/S)_{\text{cris}} \to \Sh(S_{Zar})$ comme le morphisme structural du
site cristallin. Même si $p$ n'est pas localement nilpotent sur $S_0$,
le morphisme structural
$$
\tau_{X/S} : (X/S)_{\text{cris}} \longrightarrow \Sh(S_{Zar})
$$
est défini, car on peut suivre le chemin inférieur du diagramme ci-dessus.
\end{remark}

\begin{remark}[Compatibilités]
\label{crystalline-remark-compatibilities}
Les morphismes définis ci-dessus satisfont à de nombreuses compatibilités. Par exemple,
dans la situation de la remarque \ref{crystalline-remark-functoriality-cris},
on obtient un diagramme commutatif de topos annelés
$$
\xymatrix{
(X/S)_{\text{cris}} \ar[d] \ar[r] & (Y/S')_{\text{cris}} \ar[d] \\
\Sh((\Sch/S)_{Zar}) \ar[r] & \Sh((\Sch/S')_{Zar})
}
$$
où les flèches verticales sont les morphismes structuraux.
\end{remark}



\section{Faisceaux sur le site cristallin}
\label{crystalline-section-sheaves}

\noindent
On reprend les notations et les hypothèses de la situation \ref{crystalline-situation-global}.
Afin de traiter simultanément les petits et gros sites cristallins de $X/S$
dans cette section, posons
$$
\mathcal{C} = \text{CRIS}(X/S)
\quad\text{ou}\quad
\mathcal{C} = \text{Cris}(X/S).
$$
Un faisceau $\mathcal{F}$ sur $\mathcal{C}$ détermine une
{\it restriction} $\mathcal{F}_T$ pour tout objet $(U, T, \delta)$
de $\mathcal{C}$. Plus précisément, $\mathcal{F}_T$ est le faisceau de Zariski sur
le schéma $T$ défini par la règle
$$
\mathcal{F}_T(W) = \mathcal{F}(U \cap W, W, \delta|_W)
$$
pour tout ouvert $W \subset T$. De plus, si $f : T \to T'$ est un morphisme
entre les objets
$(U, T, \delta)$ et $(U', T', \delta')$ de $\mathcal{C}$, il existe
un morphisme canonique de {\it comparaison}
\begin{equation}
\label{crystalline-equation-comparison}
c_f : f^{-1}\mathcal{F}_{T'} \longrightarrow \mathcal{F}_T.
\end{equation}
En effet, si $W' \subset T'$ est ouvert, alors $f$ induit un morphisme
$$
f|_{f^{-1}W'} :
(U \cap f^{-1}(W'), f^{-1}W', \delta|_{f^{-1}W'})
\longrightarrow
(U' \cap W', W', \delta|_{W'})
$$
de $\mathcal{C}$ ; on peut donc utiliser l'application de restriction
$(f|_{f^{-1}W'})^*$ de $\mathcal{F}$ pour définir une application
$\mathcal{F}_{T'}(W') \to \mathcal{F}_T(f^{-1}W')$.
Ces applications sont manifestement compatibles avec les restrictions ultérieures et
définissent donc un $f$-morphisme de $\mathcal{F}_{T'}$ vers $\mathcal{F}_T$ (voir
Faisceaux, section \ref{sheaves-section-presheaves-functorial},
et en particulier
Faisceaux, définition \ref{sheaves-definition-f-map}).
On obtient ainsi un morphisme $c_f$ comme dans (\ref{crystalline-equation-comparison}).
Notons que, si $f$ est une immersion ouverte, alors $c_f$ est un
isomorphisme, car $\mathcal{F}_T$ n'est dans ce cas que
la restriction de $\mathcal{F}_{T'}$ à $T$.

\medskip\noindent
Réciproquement, si l'on se donne des faisceaux de Zariski $\mathcal{F}_T$ pour tout objet
$(U, T, \delta)$ de $\mathcal{C}$ et des morphismes de comparaison
$c_f$ comme ci-dessus qui (a) sont des isomorphismes pour les immersions ouvertes et (b)
satisfont à une condition de cocycle convenable, on obtient un faisceau sur
$\mathcal{C}$. La démonstration est exactement celle de
Topologies, lemme \ref{topologies-lemma-characterize-sheaf-big}.

\medskip\noindent
Le {\it faisceau structural} sur $\mathcal{C}$ est le faisceau
$\mathcal{O}_{X/S}$ défini par la règle
$$
\mathcal{O}_{X/S} :
(U, T, \delta)
\longmapsto
\Gamma(T, \mathcal{O}_T)
$$
C'est un faisceau par définition des recouvrements dans $\mathcal{C}$.
Supposons que $\mathcal{F}$ soit un faisceau de $\mathcal{O}_{X/S}$-modules.
Dans ce cas, les applications de comparaison (\ref{crystalline-equation-comparison})
définissent un morphisme de comparaison
\begin{equation}
\label{crystalline-equation-comparison-modules}
c_f : f^*\mathcal{F}_{T'} \longrightarrow \mathcal{F}_T
\end{equation}
de $\mathcal{O}_T$-modules.

\medskip\noindent
Un autre type d'exemple s'obtient en partant d'un faisceau
$\mathcal{G}$ sur $(\Sch/X)_{Zar}$ ou sur $X_{Zar}$ (suivant que
$\mathcal{C} = \text{CRIS}(X/S)$ ou $\mathcal{C} = \text{Cris}(X/S)$).
Alors $\underline{\mathcal{G}}$, défini par la règle
$$
\underline{\mathcal{G}} :
(U, T, \delta)
\longmapsto
\mathcal{G}(U)
$$
est un faisceau sur $\mathcal{C}$. En particulier, si l'on prend
$\mathcal{G} = \mathbf{G}_a = \mathcal{O}_X$, on obtient
$$
\underline{\mathbf{G}_a} :
(U, T, \delta)
\longmapsto
\Gamma(U, \mathcal{O}_U)
$$
Il existe un morphisme surjectif de faisceaux
$\mathcal{O}_{X/S} \to \underline{\mathbf{G}_a}$ défini par les
morphismes canoniques $\Gamma(T, \mathcal{O}_T) \to \Gamma(U, \mathcal{O}_U)$
pour les objets $(U, T, \delta)$. Le noyau de ce morphisme est noté
$\mathcal{J}_{X/S}$, d'où une suite exacte courte
$$
0 \to
\mathcal{J}_{X/S} \to
\mathcal{O}_{X/S} \to
\underline{\mathbf{G}_a} \to 0
$$
Notons que $\mathcal{J}_{X/S}$ est muni d'une structure canonique de
puissances divisées. En effet, pour chaque objet $(U, T, \delta)$,
la troisième composante $\delta$ {\it est} une structure de puissances divisées sur le
noyau de $\mathcal{O}_T \to \mathcal{O}_U$. Ainsi, le (gros)
topos cristallin est un topos à puissances divisées.





\section{Cristaux en modules}
\label{crystalline-section-crystals}

\noindent
Un cristal est en fait un objet très général. Toutefois, sa
définition peut être un peu délicate à analyser ; nous la donnons donc d'abord
dans le cas des modules sur les sites cristallins.

\begin{definition}
\label{crystalline-definition-modules}
Dans la situation \ref{crystalline-situation-global}.
Soit $\mathcal{C} = \text{CRIS}(X/S)$ ou $\mathcal{C} = \text{Cris}(X/S)$.
Soit $\mathcal{F}$ un faisceau de $\mathcal{O}_{X/S}$-modules sur $\mathcal{C}$.
\begin{enumerate}
\item On dit que $\mathcal{F}$ est {\it localement quasi-cohérent} si, pour tout
objet $(U, T, \delta)$ de $\mathcal{C}$, la restriction $\mathcal{F}_T$
est un $\mathcal{O}_T$-module quasi-cohérent.
\item On dit que $\mathcal{F}$ est {\it quasi-cohérent} s'il est quasi-cohérent
au sens de
Modules sur les sites, définition \ref{sites-modules-definition-site-local}.
\item On dit que $\mathcal{F}$ est un {\it cristal en $\mathcal{O}_{X/S}$-modules}
si tous les morphismes de comparaison (\ref{crystalline-equation-comparison-modules}) sont des
isomorphismes.
\end{enumerate}
\end{definition}

\noindent
Ces notions sont reliées de la manière suivante.

\begin{lemma}
\label{crystalline-lemma-crystal-quasi-coherent-modules}
Avec les notations $X/S, \mathcal{I}, \gamma, \mathcal{C}, \mathcal{F}$
de la définition \ref{crystalline-definition-modules}, les conditions suivantes sont équivalentes :
\begin{enumerate}
\item $\mathcal{F}$ est quasi-cohérent, et
\item $\mathcal{F}$ est localement quasi-cohérent et est un cristal en
$\mathcal{O}_{X/S}$-modules.
\end{enumerate}
\end{lemma}

\begin{proof}
Supposons (1). Soit $f : (U', T', \delta') \to (U, T, \delta)$ un morphisme de
$\mathcal{C}$. Il faut montrer (a) que $\mathcal{F}_T$ est un
$\mathcal{O}_T$-module quasi-cohérent et (b) que $c_f : f^*\mathcal{F}_T \to \mathcal{F}_{T'}$
est un isomorphisme. L'hypothèse signifie qu'il existe un recouvrement
$\{(U_i, T_i, \delta_i) \to (U, T, \delta)\}$ tel que, pour chaque $i$,
la restriction de $\mathcal{F}$ à $\mathcal{C}/(U_i, T_i, \delta_i)$
admette une présentation globale. Comme il suffit de prouver (a) et (b)
localement pour la topologie de Zariski, on peut remplacer $f : (U', T', \delta') \to (U, T, \delta)$
par son changement de base à $(U_i, T_i, \delta_i)$ et supposer que la restriction de $\mathcal{F}$
à $\mathcal{C}/(U, T, \delta)$ admet une
présentation globale
$$
\bigoplus\nolimits_{j \in J}
\mathcal{O}_{X/S}|_{\mathcal{C}/(U, T, \delta)} \longrightarrow
\bigoplus\nolimits_{i \in I}
\mathcal{O}_{X/S}|_{\mathcal{C}/(U, T, \delta)} \longrightarrow
\mathcal{F}|_{\mathcal{C}/(U, T, \delta)}
\longrightarrow 0
$$
Il est clair que cela donne une présentation
$$
\bigoplus\nolimits_{j \in J} \mathcal{O}_T \longrightarrow
\bigoplus\nolimits_{i \in I} \mathcal{O}_T \longrightarrow
\mathcal{F}_T
\longrightarrow 0
$$
et donc (a) est vérifiée. De plus, cette présentation se restreint à $T'$
en une présentation analogue de $\mathcal{F}_{T'}$, ce qui établit (b).

\medskip\noindent
Supposons (2). Soit $(U, T, \delta)$ un objet de $\mathcal{C}$.
Il faut trouver un recouvrement de $(U, T, \delta)$ tel que $\mathcal{F}$ admette une
présentation globale après restriction aux localisations de $\mathcal{C}$
associées aux membres du recouvrement. On peut donc supposer $T$ affine.
Dans ce cas, on peut choisir une présentation
$$
\bigoplus\nolimits_{j \in J} \mathcal{O}_T \longrightarrow
\bigoplus\nolimits_{i \in I} \mathcal{O}_T \longrightarrow
\mathcal{F}_T
\longrightarrow 0
$$
puisque $\mathcal{F}_T$ est supposé être un $\mathcal{O}_T$-module quasi-cohérent.
La propriété de cristal de $\mathcal{F}$ montre alors que l'image réciproque de cette présentation donne
une présentation de $\mathcal{F}_{T'}$ pour tout morphisme
$f : (U', T', \delta') \to (U, T, \delta)$ de $\mathcal{C}$.
On obtient ainsi la présentation voulue de $\mathcal{F}|_{\mathcal{C}/(U, T, \delta)}$.
\end{proof}

\begin{definition}
\label{crystalline-definition-crystal-quasi-coherent-modules}
Si $\mathcal{F}$ satisfait aux conditions équivalentes du
lemme \ref{crystalline-lemma-crystal-quasi-coherent-modules}, on dit
que $\mathcal{F}$ est un
{\it cristal en modules quasi-cohérents}.
On dit que $\mathcal{F}$ est un {\it cristal en modules localement libres de type fini}
si, de plus, $\mathcal{F}$ est localement libre de type fini.
\end{definition}

\noindent
Bien entendu, comme le montre le lemme \ref{crystalline-lemma-crystal-quasi-coherent-modules}, cette
terminologie est quelque peu lourde, puisqu'un module quasi-cohérent est toujours un cristal.
Mais c'est la terminologie usuelle dans la littérature.

\begin{remark}
\label{crystalline-remark-crystal}
Pour formuler la notion générale de cristal, on utilise le langage
des champs et des morphismes fortement cartésiens ; voir
Champs, définition \ref{stacks-definition-stack} et
Catégories, définition \ref{categories-definition-cartesian-over-C}.
Dans la situation \ref{crystalline-situation-global}, soit
$p : \mathcal{C} \to \text{Cris}(X/S)$ un champ.
Un {\it cristal en objets de $\mathcal{C}$ sur $X$ relativement à $S$}
est une {\it section cartésienne} $\sigma : \text{Cris}(X/S) \to \mathcal{C}$,
c'est-à-dire un foncteur $\sigma$ tel que $p \circ \sigma = \text{id}$
et tel que $\sigma(f)$ soit fortement cartésien pour tout
morphisme $f$ de $\text{Cris}(X/S)$. De même pour le gros site cristallin.
\end{remark}





\section{Faisceau des différentielles}
\label{crystalline-section-differentials-sheaf}

\noindent
Dans cette section, nous nous en tiendrons au (petit) site cristallin,
qui paraît plus naturel. On globalise
la définition \ref{crystalline-definition-derivation} comme suit.

\begin{definition}
\label{crystalline-definition-global-derivation}
Dans la situation \ref{crystalline-situation-global}, soit
$\mathcal{F}$ un faisceau de $\mathcal{O}_{X/S}$-modules sur
$\text{Cris}(X/S)$. Une
{\it $S$-dérivation $D : \mathcal{O}_{X/S} \to \mathcal{F}$}
est un morphisme de faisceaux tel que, pour tout objet $(U, T, \delta)$ de
$\text{Cris}(X/S)$, l'application
$$
D : \Gamma(T, \mathcal{O}_T) \longrightarrow \Gamma(T, \mathcal{F})
$$
soit une $\Gamma(V, \mathcal{O}_V)$-dérivation à puissances divisées, où $V \subset S$
est un ouvert quelconque tel que $T \to S$ se factorise par $V$.
\end{definition}

\noindent
Cela signifie que $D$ est additive, satisfait à la règle de Leibniz, s'annule sur les
fonctions provenant de $S$ et vérifie $D(f^{[n]}) = f^{[n - 1]}D(f)$
pour toute section locale $f$ de l'idéal à puissances divisées $\mathcal{J}_{X/S}$.
C'est un cas particulier d'une notion très générale que nous allons maintenant décrire.

\medskip\noindent
Comparons la discussion qui suit avec
Modules sur les sites, section \ref{sites-modules-section-differentials}. Soit
$\mathcal{C}$ un site, soit $\mathcal{A} \to \mathcal{B}$ un
morphisme de faisceaux d'anneaux sur $\mathcal{C}$, soit $\mathcal{J} \subset \mathcal{B}$
un faisceau d'idéaux, soit $\delta$ une structure de puissances divisées sur
$\mathcal{J}$ et soit $\mathcal{F}$ un faisceau de $\mathcal{B}$-modules.
Il existe alors une notion de {\it $\mathcal{A}$-dérivation à puissances divisées}
$D : \mathcal{B} \to \mathcal{F}$. Cela signifie que $D$ est $\mathcal{A}$-linéaire,
satisfait à la règle de Leibniz et vérifie
$D(\delta_n(x)) = \delta_{n - 1}(x)D(x)$ pour les sections locales $x$ de
$\mathcal{J}$. Dans cette situation, il existe une
{\it $\mathcal{A}$-dérivation universelle à puissances divisées}
$$
\text{d}_{\mathcal{B}/\mathcal{A}, \delta} :
\mathcal{B}
\longrightarrow
\Omega_{\mathcal{B}/\mathcal{A}, \delta}
$$
De plus, $\text{d}_{\mathcal{B}/\mathcal{A}, \delta}$ est le composé
$$
\mathcal{B}
\longrightarrow
\Omega_{\mathcal{B}/\mathcal{A}}
\longrightarrow
\Omega_{\mathcal{B}/\mathcal{A}, \delta}
$$
où le premier morphisme est la dérivation universelle construite dans la démonstration
de Modules sur les sites, lemme \ref{sites-modules-lemma-universal-module},
et le second est le morphisme quotient par le sous-module engendré par
les sections locales
$\text{d}_{\mathcal{B}/\mathcal{A}}(\delta_n(x)) -
\delta_{n - 1}(x)\text{d}_{\mathcal{B}/\mathcal{A}}(x)$.

\medskip\noindent
Passons à la version relative. Supposons que
$(f, f^\sharp) : (\Sh(\mathcal{C}), \mathcal{O}) \to
(\Sh(\mathcal{C}'), \mathcal{O}')$ soit un morphisme de topos annelés,
que $\mathcal{J} \subset \mathcal{O}$ soit un faisceau d'idéaux, que $\delta$ soit une
structure de puissances divisées sur $\mathcal{J}$ et que $\mathcal{F}$ soit un faisceau
de $\mathcal{O}$-modules. Dans cette situation, on dit que
$D : \mathcal{O} \to \mathcal{F}$ est une $\mathcal{O}'$-dérivation à puissances divisées
si $D$ est une $f^{-1}\mathcal{O}'$-dérivation à puissances divisées au sens défini ci-dessus.
En outre, on écrit
$$
\Omega_{\mathcal{O}/\mathcal{O}', \delta} =
\Omega_{\mathcal{O}/f^{-1}\mathcal{O}', \delta}
$$
qui reçoit la
$\mathcal{O}'$-dérivation universelle à puissances divisées.

\medskip\noindent
En appliquant cette construction au morphisme structural
$$
(X/S)_{\text{Cris}} \longrightarrow \Sh(S_{Zar})
$$
(voir la remarque \ref{crystalline-remark-structure-morphism}), on retrouve la notion
définie dans la définition \ref{crystalline-definition-global-derivation}.
En particulier, il existe une dérivation universelle à puissances divisées
$$
d_{X/S} : \mathcal{O}_{X/S} \to \Omega_{X/S}
$$
Notons que nous omettons dans la notation l'indication de la compatibilité du
module des différentielles avec les puissances divisées (il paraît en effet
peu probable que l'on considère jamais le module usuel des différentielles du
faisceau structural sur le site cristallin).

\begin{lemma}
\label{crystalline-lemma-module-differentials-divided-power-scheme}
Soit $(T, \mathcal{J}, \delta)$ un schéma à puissances divisées.
Soit $T \to S$ un morphisme de schémas.
Le quotient $\Omega_{T/S} \to \Omega_{T/S, \delta}$
décrit ci-dessus est un $\mathcal{O}_T$-module quasi-cohérent.
Pour tout ouvert affine $W \subset T$ dont l'image est contenue dans un
ouvert affine $V \subset S$, on a
$$
\Gamma(W, \Omega_{T/S, \delta}) =
\Omega_{\Gamma(W, \mathcal{O}_W)/\Gamma(V, \mathcal{O}_V), \delta}
$$
où le second membre est
celui construit dans la section \ref{crystalline-section-differentials}.
\end{lemma}

\begin{proof}
Omis.
\end{proof}

\begin{lemma}
\label{crystalline-lemma-module-of-differentials}
Dans la situation \ref{crystalline-situation-global}.
Pour $(U, T, \delta)$ dans $\text{Cris}(X/S)$, la restriction
$(\Omega_{X/S})_T$ à $T$ est $\Omega_{T/S, \delta}$ et la restriction
$\text{d}_{X/S}|_T$ est égale à $\text{d}_{T/S, \delta}$.
\end{lemma}

\begin{proof}
Omis.
\end{proof}

\begin{lemma}
\label{crystalline-lemma-module-of-differentials-on-affine}
Dans la situation \ref{crystalline-situation-global}.
Pour tout objet affine $(U, T, \delta)$ de $\text{Cris}(X/S)$
s'envoyant dans un ouvert affine $V \subset S$, on a
$$
\Gamma((U, T, \delta), \Omega_{X/S}) =
\Omega_{\Gamma(T, \mathcal{O}_T)/\Gamma(V, \mathcal{O}_V), \delta}
$$
où le second membre est
celui construit dans la section \ref{crystalline-section-differentials}.
\end{lemma}

\begin{proof}
Cela résulte des lemmes \ref{crystalline-lemma-module-differentials-divided-power-scheme} et
\ref{crystalline-lemma-module-of-differentials}.
\end{proof}

\begin{lemma}
\label{crystalline-lemma-describe-omega-small}
Dans la situation \ref{crystalline-situation-global}.
Soit $(U, T, \delta)$ un objet de $\text{Cris}(X/S)$.
Soit
$$
(U(1), T(1), \delta(1)) = (U, T, \delta) \times (U, T, \delta)
$$
dans $\text{Cris}(X/S)$. Soit $\mathcal{K} \subset \mathcal{O}_{T(1)}$
le faisceau quasi-cohérent d'idéaux correspondant à l'immersion
fermée $\Delta : T \to T(1)$. Alors
$\mathcal{K} \subset \mathcal{J}_{T(1)}$ est stable par la
structure de puissances divisées sur $\mathcal{J}_{T(1)}$ et l'on a
$$
(\Omega_{X/S})_T = \mathcal{K}/\mathcal{K}^{[2]}
$$
\end{lemma}

\begin{proof}
Notons que $U = U(1)$, puisque $U \to X$ est une immersion ouverte et que
le foncteur (\ref{crystalline-equation-forget-small}) commute aux produits. On en déduit que
$\mathcal{K} \subset \mathcal{J}_{T(1)}$. Cela étant, le lemme s'obtient
en travaillant localement sur des ouverts affines de $T$ et en utilisant les
lemmes \ref{crystalline-lemma-module-of-differentials-on-affine} et
\ref{crystalline-lemma-diagonal-and-differentials-affine-site}.
\end{proof}

\noindent
Le faisceau $\Omega_{X/S}$ n'est pas un cristal en
$\mathcal{O}_{X/S}$-modules quasi-cohérents. Il satisfait néanmoins à deux propriétés
étroitement apparentées (comparer avec le
lemme \ref{crystalline-lemma-crystal-quasi-coherent-modules}).

\begin{lemma}
\label{crystalline-lemma-omega-locally-quasi-coherent}
Dans la situation \ref{crystalline-situation-global}.
Le faisceau des différentielles $\Omega_{X/S}$ possède les deux
propriétés suivantes :
\begin{enumerate}
\item $\Omega_{X/S}$ est localement quasi-cohérent, et
\item pour tout morphisme $(U, T, \delta) \to (U', T', \delta')$
de $\text{Cris}(X/S)$ tel que $f : T \to T'$ soit une immersion fermée,
le morphisme $c_f : f^*(\Omega_{X/S})_{T'} \to (\Omega_{X/S})_T$ est surjectif.
\end{enumerate}
\end{lemma}

\begin{proof}
L'assertion (1) résulte de la combinaison des
lemmes \ref{crystalline-lemma-module-differentials-divided-power-scheme} et
\ref{crystalline-lemma-module-of-differentials}.
L'assertion (2) résulte du fait que
$(\Omega_{X/S})_T = \Omega_{T/S, \delta}$
est un quotient de $\Omega_{T/S}$ et que $f^*\Omega_{T'/S} \to \Omega_{T/S}$
est surjectif.
\end{proof}







\section{Deux épaississements universels}
\label{crystalline-section-universal-thickenings}

\noindent
Les constructions de cette section nous aideront à définir une connexion sur
un cristal en modules sur le site cristallin. En quelque sorte, les constructions
présentées ici sont les versions « faisceautisées et universelles » de celles de la
section \ref{crystalline-section-explicit-thickenings}.

\begin{remark}
\label{crystalline-remark-first-order-thickening}
Dans la situation \ref{crystalline-situation-global}.
Soit $(U, T, \delta)$ un objet de $\text{Cris}(X/S)$.
Notons $\Omega_{T/S, \delta} = (\Omega_{X/S})_T$ ; voir le
lemme \ref{crystalline-lemma-module-of-differentials}.
Décrivons explicitement un épaississement du premier ordre $T'$ de
$T$. Plus précisément, posons
$$
\mathcal{O}_{T'} = \mathcal{O}_T \oplus \Omega_{T/S, \delta}
$$
muni de la structure d'algèbre pour laquelle $\Omega_{T/S, \delta}$ est un
idéal de carré nul. Soit $\mathcal{J} \subset \mathcal{O}_T$
le faisceau d'idéaux de l'immersion fermée $U \to T$. Posons
$\mathcal{J}' = \mathcal{J} \oplus \Omega_{T/S, \delta}$.
Définissons une structure de puissances divisées sur $\mathcal{J}'$ en posant
$$
\delta_n'(f, \omega) = (\delta_n(f), \delta_{n - 1}(f)\omega),
$$
voir le lemme \ref{crystalline-lemma-divided-power-first-order-thickening}.
Il existe deux homomorphismes d'anneaux
$$
p_0, p_1 : \mathcal{O}_T \to \mathcal{O}_{T'}
$$
Le premier est donné par $f \mapsto (f, 0)$ et le second par
$f \mapsto (f, \text{d}_{T/S, \delta}f)$. Notons que tous deux sont compatibles
avec les structures de puissances divisées sur $\mathcal{J}$ et $\mathcal{J}'$,
de même que le morphisme quotient $\mathcal{O}_{T'} \to \mathcal{O}_T$.
On obtient ainsi un objet $(U, T', \delta')$ de $\text{Cris}(X/S)$
et un diagramme commutatif
$$
\xymatrix{
& T \ar[ld]_{\text{id}} \ar[d]^i \ar[rd]^{\text{id}} \\
T & T' \ar[l]_{p_0} \ar[r]^{p_1} & T
}
$$
dans $\text{Cris}(X/S)$ tel que $i$ soit un épaississement du premier ordre dont le faisceau
d'idéaux s'identifie à $\Omega_{T/S, \delta}$ et tel que
$p_1 - p_0 : \mathcal{O}_T \to \mathcal{O}_{T'}$
s'identifie à la dérivation universelle $\text{d}_{T/S, \delta}$
composée avec l'inclusion $\Omega_{T/S, \delta} \to \mathcal{O}_{T'}$.
\end{remark}

\begin{remark}
\label{crystalline-remark-second-order-thickening}
Dans la situation \ref{crystalline-situation-global}.
Soit $(U, T, \delta)$ un objet de $\text{Cris}(X/S)$.
Notons $\Omega_{T/S, \delta} = (\Omega_{X/S})_T$ ; voir le
lemme \ref{crystalline-lemma-module-of-differentials}.
Notons également $\Omega^2_{T/S, \delta}$ sa seconde puissance
extérieure. Décrivons explicitement un épaississement du second ordre $T''$ de $T$.
Plus précisément, posons
$$
\mathcal{O}_{T''} =
\mathcal{O}_T \oplus \Omega_{T/S, \delta} \oplus \Omega_{T/S, \delta}
\oplus \Omega^2_{T/S, \delta}
$$
muni de la structure d'algèbre définie de la manière suivante :
$$
(f, \omega_1, \omega_2, \eta) \cdot
(f', \omega_1', \omega_2', \eta') =
(ff', f\omega_1' + f'\omega_1, f\omega_2' + f'\omega_2,
f\eta' + f'\eta + \omega_1 \wedge \omega_2' + \omega_1' \wedge \omega_2).
$$
Soit $\mathcal{J} \subset \mathcal{O}_T$
le faisceau d'idéaux de l'immersion fermée $U \to T$. Soit
$\mathcal{J}''$ l'image réciproque de $\mathcal{J}$ par la
projection $\mathcal{O}_{T''} \to \mathcal{O}_T$.
Définissons une structure de puissances divisées sur $\mathcal{J}''$ en posant
$$
\delta_n''(f, \omega_1, \omega_2, \eta) =
(\delta_n(f), \delta_{n - 1}(f)\omega_1, \delta_{n - 1}(f)\omega_2,
\delta_{n - 1}(f)\eta + \delta_{n - 2}(f)\omega_1 \wedge \omega_2)
$$
voir le lemme \ref{crystalline-lemma-divided-power-second-order-thickening}.
Il existe trois homomorphismes d'anneaux
$q_0, q_1, q_2 : \mathcal{O}_T \to \mathcal{O}_{T''}$
donnés par
\begin{align*}
q_0(f) & = (f, 0, 0, 0), \\
q_1(f) & = (f, \text{d}f, 0, 0), \\
q_2(f) & = (f, \text{d}f, \text{d}f, 0)
\end{align*}
où $\text{d} = \text{d}_{T/S, \delta}$.
Notons que tous trois sont compatibles avec les structures de puissances divisées
sur $\mathcal{J}$ et $\mathcal{J}''$. Il existe trois homomorphismes d'anneaux
$q_{01}, q_{12}, q_{02} : \mathcal{O}_{T'} \to \mathcal{O}_{T''}$
où $\mathcal{O}_{T'}$ est défini comme dans la remarque \ref{crystalline-remark-first-order-thickening}.
Plus précisément, posons
\begin{align*}
q_{01}(f, \omega) & = (f, \omega, 0, 0), \\
q_{12}(f, \omega) & =
(f, \text{d}f, \omega, \text{d}\omega), \\
q_{02}(f, \omega) & = (f, \omega, \omega, 0)
\end{align*}
Ces homomorphismes sont eux aussi compatibles avec les structures de puissances
divisées données. Vérifions ces assertions pour $q_{12}$. Notons
que $q_{12}$ est un homomorphisme d'anneaux, puisque
\begin{align*}
q_{12}(f, \omega)q_{12}(g, \eta) & =
(f, \text{d}f, \omega, \text{d}\omega)(g, \text{d}g, \eta, \text{d}\eta) \\
& =
(fg, f\text{d}g + g \text{d}f, f\eta + g\omega,
f\text{d}\eta + g\text{d}\omega + \text{d}f \wedge \eta +
\text{d}g \wedge \omega) \\
& = q_{12}(fg, f\eta + g\omega) = q_{12}((f, \omega)(g, \eta))
\end{align*}
Notons que $q_{12}$ est compatible avec les puissances divisées, car
\begin{align*}
\delta_n''(q_{12}(f, \omega)) & =
\delta_n''((f, \text{d}f, \omega, \text{d}\omega)) \\
& =
(\delta_n(f), \delta_{n - 1}(f)\text{d}f, \delta_{n - 1}(f)\omega,
\delta_{n - 1}(f)\text{d}\omega + \delta_{n - 2}(f)\text{d}(f) \wedge \omega)
\\
& = q_{12}((\delta_n(f), \delta_{n - 1}(f)\omega)) =
q_{12}(\delta'_n(f, \omega))
\end{align*}
Les vérifications pour $q_{01}$ et $q_{02}$ sont plus faciles.
Notons que $q_0 = q_{01} \circ p_0$, $q_1 = q_{01} \circ p_1$,
$q_1 = q_{12} \circ p_0$, $q_2 = q_{12} \circ p_1$,
$q_0 = q_{02} \circ p_0$ et $q_2 = q_{02} \circ p_1$.
Ainsi $(U, T'', \delta'')$ est un objet de $\text{Cris}(X/S)$
et l'on obtient des morphismes
$$
\xymatrix{
T''
\ar@<2ex>[r]
\ar@<0ex>[r]
\ar@<-2ex>[r]
&
T'
\ar@<1ex>[r]
\ar@<-1ex>[r]
&
T
}
$$
de $\text{Cris}(X/S)$ satisfaisant aux relations décrites ci-dessus.
Dans les applications, les notations $q_i : T'' \to T$ et
$q_{ij} : T'' \to T'$ désigneront les morphismes associés aux
homomorphismes d'anneaux décrits ci-dessus.
\end{remark}






\section{Le complexe de de Rham}
\label{crystalline-section-de-Rham}

\noindent
Dans la situation \ref{crystalline-situation-global}.
Sur le (petit) site cristallin, on définit
$\Omega^i_{X/S} = \wedge^i_{\mathcal{O}_{X/S}} \Omega_{X/S}$
pour $i \geq 0$. La $S$-dérivation universelle $\text{d}_{X/S}$ définit
le {\it complexe de de Rham}
$$
\mathcal{O}_{X/S} \to \Omega^1_{X/S} \to \Omega^2_{X/S} \to \ldots
$$
sur $\text{Cris}(X/S)$ ; voir le
lemme \ref{crystalline-lemma-module-of-differentials-on-affine} et la
remarque \ref{crystalline-remark-divided-powers-de-rham-complex}.


\section{Connexions}
\label{crystalline-section-connections}

\noindent
Dans la situation \ref{crystalline-situation-global}.
Étant donné un $\mathcal{O}_{X/S}$-module $\mathcal{F}$ sur $\text{Cris}(X/S)$,
une {\it connexion} est un morphisme de faisceaux abéliens
$$
\nabla :
\mathcal{F}
\longrightarrow
\mathcal{F} \otimes_{\mathcal{O}_{X/S}} \Omega_{X/S}
$$
tel que $\nabla(f s) = f\nabla(s) + s \otimes \text{d}f$
pour les sections locales $s, f$ de $\mathcal{F}$ et de $\mathcal{O}_{X/S}$, respectivement.
Une connexion détermine des morphismes canoniques
$
\nabla :
\mathcal{F} \otimes_{\mathcal{O}_{X/S}} \Omega^i_{X/S}
\longrightarrow
\mathcal{F} \otimes_{\mathcal{O}_{X/S}} \Omega^{i + 1}_{X/S}
$
définis par la règle $\nabla(s \otimes \omega) =
\nabla(s) \wedge \omega + s \otimes \text{d}\omega$
comme dans la remarque \ref{crystalline-remark-connection}. On dit que la connexion est
{\it intégrable} si $\nabla \circ \nabla = 0$. Si $\nabla$ est intégrable,
on obtient le {\it complexe de de Rham}
$$
\mathcal{F} \to
\mathcal{F} \otimes_{\mathcal{O}_{X/S}} \Omega^1_{X/S} \to
\mathcal{F} \otimes_{\mathcal{O}_{X/S}} \Omega^2_{X/S} \to \ldots
$$
sur $\text{Cris}(X/S)$. Tout cristal en
$\mathcal{O}_{X/S}$-modules est en fait muni d'une connexion
intégrable canonique.

\begin{lemma}
\label{crystalline-lemma-automatic-connection}
Dans la situation \ref{crystalline-situation-global}.
Soit $\mathcal{F}$ un cristal en $\mathcal{O}_{X/S}$-modules
sur $\text{Cris}(X/S)$. Alors $\mathcal{F}$ est muni d'une
connexion intégrable canonique.
\end{lemma}

\begin{proof}
Soit $(U, T, \delta)$ un objet de $\text{Cris}(X/S)$.
Soit $(U, T', \delta')$ l'épaississement infinitésimal de $T$
par $(\Omega_{X/S})_T = \Omega_{T/S, \delta}$
construit dans la remarque \ref{crystalline-remark-first-order-thickening}.
Il est muni de projections $p_0, p_1 : T' \to T$
et d'une diagonale $i : T \to T'$. Par hypothèse, on obtient des
isomorphismes
$$
p_0^*\mathcal{F}_T \xrightarrow{c_0}
\mathcal{F}_{T'} \xleftarrow{c_1}
p_1^*\mathcal{F}_T
$$
de $\mathcal{O}_{T'}$-modules. L'image réciproque de $c = c_1^{-1} \circ c_0$
sur $T$ par $i$ est l'identité
de $\mathcal{F}_T$. Par conséquent, si $s \in \Gamma(T, \mathcal{F}_T)$,
alors $\nabla(s) = p_1^*s - c(p_0^*s)$ est une section de
$p_1^*\mathcal{F}_T$ dont l'image réciproque par $i$ est nulle. Ainsi,
$\nabla(s)$ est une section de
$$
\mathcal{F}_T
\otimes_{\mathcal{O}_T}
\Omega_{T/S, \delta}
$$
car il s'agit du noyau de $p_1^*\mathcal{F}_T \to \mathcal{F}_T$,
puisque $\mathcal{O}_{T'} = \mathcal{O}_T \oplus \Omega_{T/S, \delta}$
par construction. On vérifie aisément que $\nabla(fs) =
f\nabla(s) + s \otimes \text{d}(f)$ en utilisant la description de
$\text{d}$ dans la remarque \ref{crystalline-remark-first-order-thickening}.

\medskip\noindent
La famille d'applications
$$
\nabla : \Gamma(T, \mathcal{F}_T) \to
\Gamma(T, \mathcal{F}_T \otimes_{\mathcal{O}_T} \Omega_{T/S, \delta})
$$
ainsi obtenue est fonctorielle en $T$, car la construction de $T'$
est fonctorielle en $T$. On obtient donc une connexion.

\medskip\noindent
Pour montrer que la connexion est intégrable, considérons
l'objet $(U, T'', \delta'')$ construit dans la
remarque \ref{crystalline-remark-second-order-thickening}.
Comme $\mathcal{F}$ est un faisceau, on voit que
$$
\xymatrix{
q_0^*\mathcal{F}_T \ar[rr]_{q_{01}^*c} \ar[rd]_{q_{02}^*c} & &
q_1^*\mathcal{F}_T \ar[ld]^{q_{12}^*c} \\
& q_2^*\mathcal{F}_T
}
$$
est un diagramme commutatif de $\mathcal{O}_{T''}$-modules. Pour
$s \in \Gamma(T, \mathcal{F}_T)$, on a
$c(p_0^*s) = p_1^*s - \nabla(s)$. Écrivons
$\nabla(s) = \sum p_1^*s_i \cdot \omega_i$, où $s_i$ est une section locale
de $\mathcal{F}_T$ et $\omega_i$ une section locale de $\Omega_{T/S, \delta}$.
On considère $\omega_i$ comme une section locale du faisceau structural
$\mathcal{O}_{T'}$ et l'on écrit donc un produit plutôt qu'un produit
tensoriel. D'une part,
\begin{align*}
q_{12}^*c \circ q_{01}^*c(q_0^*s) & = 
q_{12}^*c(q_1^*s - \sum q_1^*s_i \cdot q_{01}^*\omega_i) \\
& =
q_2^*s - \sum q_2^*s_i \cdot q_{12}^*\omega_i -
\sum q_2^*s_i \cdot q_{01}^*\omega_i +
\sum q_{12}^*\nabla(s_i) \cdot q_{01}^*\omega_i
\end{align*}
et, d'autre part,
$$
q_{02}^*c(q_0^*s) = q_2^*s - \sum q_2^*s_i \cdot q_{02}^*\omega_i.
$$
Les formules de la remarque \ref{crystalline-remark-second-order-thickening} montrent
que
$q_{01}^*\omega_i + q_{12}^*\omega_i - q_{02}^*\omega_i = \text{d}\omega_i$.
La différence des deux expressions ci-dessus est donc
$$
\sum q_2^*s_i \cdot \text{d}\omega_i -
\sum q_{12}^*\nabla(s_i) \cdot q_{01}^*\omega_i
$$
Notons que
$q_{12}^*\omega \cdot q_{01}^*\omega' = \omega' \wedge \omega =
- \omega \wedge \omega'$ par définition de la multiplication sur
$\mathcal{O}_{T''}$. L'expression ci-dessus est donc $\nabla^2(s)$, considérée
comme une section du sous-faisceau $\mathcal{F}_T \otimes \Omega^2_{T/S, \delta}$ de
$q_2^*\mathcal{F}$. On obtient ainsi la condition d'intégrabilité.
\end{proof}




\section{Algèbre cosimpliciale}
\label{crystalline-section-cosimplicial}

\noindent
Cette section devrait être déplacée ailleurs. Un
{\it anneau cosimplicial} est un objet cosimplicial
de la catégorie des anneaux. Étant donné un anneau $R$, une
{\it $R$-algèbre cosimpliciale} est un objet cosimplicial de la
catégorie des $R$-algèbres. Un {\it idéal cosimplicial} d'un anneau
cosimplicial $A_*$ est la donnée d'idéaux $I_n \subset A_n$, pour tout $n$, tels
que $A(f)(I_n) \subset I_m$ pour tout $f : [n] \to [m]$ dans $\Delta$.

\medskip\noindent
Soit $A_*$ un anneau cosimplicial. Soit $\mathcal{C}$ la catégorie
des couples $(A, M)$ où $A$ est un anneau et $M$ un module sur $A$.
Un morphisme $(A, M) \to (A', M')$ consiste en un homomorphisme d'anneaux $A \to A'$ et
un homomorphisme $A$-linéaire $M \to M'$, où $M'$ est considéré comme un $A$-module
par l'intermédiaire de $A \to A'$ et de la structure de $A'$-module sur $M'$. Cela étant,
on définit un {\it module cosimplicial $M_*$ sur $A_*$} comme un objet
cosimplicial $(A_*, M_*)$ de $\mathcal{C}$ dont la première composante est $A_*$.
Un {\it homomorphisme $\varphi_* : M_* \to N_*$ de modules cosimpliciaux sur
$A_*$} est un morphisme $(A_*, M_*) \to (A_*, N_*)$ d'objets cosimpliciaux
de $\mathcal{C}$ dont la première composante est $1_{A_*}$.

\medskip\noindent
Une {\it homotopie} entre des homomorphismes $\varphi_*, \psi_* : M_* \to N_*$
de modules cosimpliciaux sur $A_*$ est une homotopie entre les
morphismes associés $(A_*, M_*) \to (A_*, N_*)$ dont la première composante est
l'homotopie triviale (duale de
Objets simpliciaux, exemple \ref{simplicial-example-trivial-homotopy}).
Explicitons cette définition. Une telle homotopie est une homotopie
$$
h : M_* \longrightarrow \Hom(\Delta[1], N_*)
$$
entre $\varphi_*$ et $\psi_*$ comme homomorphismes de groupes abéliens
cosimpliciaux, telle que, pour chaque $n$, l'application
$h_n : M_n \to \prod_{\alpha \in \Delta[1]_n} N_n$ soit $A_n$-linéaire.
Le lemme suivant est, pour les modules cosimpliciaux,
une version du lemme \ref{simplicial-lemma-functorial-homotopy}
d'Objets simpliciaux.

\begin{lemma}
\label{crystalline-lemma-homotopy-tensor}
Soit $A_*$ un anneau cosimplicial. Soient $\varphi_*, \psi_* : K_* \to M_*$
des homomorphismes de modules cosimpliciaux sur $A_*$.
\begin{enumerate}
\item
\label{crystalline-item-tensor}
Si $\varphi_*$ et $\psi_*$ sont homotopes, alors
$$
\varphi_* \otimes 1, \psi_* \otimes 1 :
K_* \otimes_{A_*} L_* \longrightarrow M_* \otimes_{A_*} L_*
$$
sont homotopes pour tout $A_*$-module cosimplicial $L_*$.
\item
\label{crystalline-item-wedge}
Si $\varphi_*$ et $\psi_*$ sont homotopes, alors
$$
\wedge^i(\varphi_*), \wedge^i(\psi_*) :
\wedge^i(K_*) \longrightarrow \wedge^i(M_*)
$$
sont homotopes.
\item
\label{crystalline-item-base-change}
Si $\varphi_*$ et $\psi_*$ sont homotopes et si $A_* \to B_*$
est un homomorphisme d'anneaux cosimpliciaux, alors
$$
\varphi_* \otimes 1, \psi_* \otimes 1 :
K_* \otimes_{A_*} B_* \longrightarrow M_* \otimes_{A_*} B_*
$$
sont homotopes comme homomorphismes de modules cosimpliciaux sur $B_*$.
\item
\label{crystalline-item-completion}
Si $I_* \subset A_*$ est un idéal cosimplicial, alors les applications
induites
$$
\varphi^\wedge_*, \psi^\wedge_* :
K_*^\wedge \longrightarrow M_*^\wedge
$$
entre les complétés sont homotopes.
\item Ajouter ici d'autres assertions au besoin, par exemple sur les puissances symétriques.
\end{enumerate}
\end{lemma}

\begin{proof}
Soit $h : K_* \longrightarrow \Hom(\Delta[1], M_*)$ l'homotopie
donnée. En degré $n$, on a
$$
h_n = (h_{n, \alpha}) :
K_n \longrightarrow
\prod\nolimits_{\alpha \in \Delta[1]_n} M_n
$$
voir Objets simpliciaux, section \ref{simplicial-section-homotopy-cosimplicial}.
Pour qu'une famille de $h_{n, \alpha}$ constitue une homotopie,
il faut et il suffit que, pour tout $f : [n] \to [m]$, on
ait
$$
h_{m, \alpha} \circ M_*(f) = N_*(f) \circ h_{n, \alpha \circ f}
$$
voir
Objets simpliciaux, équation (\ref{simplicial-equation-property-homotopy-cosimplicial}).
On doit aussi avoir $\psi_n = h_{n, 0 : [n] \to [1]}$ et
$\varphi_n = h_{n, 1 : [n] \to [1]}$.

\medskip\noindent
Dans chacun des cas du lemme, on peut construire les morphismes correspondants.
Cas (\ref{crystalline-item-tensor}). On peut utiliser l'homotopie $h \otimes 1$ définie
en degré $n$ en posant
$$
(h \otimes 1)_{n, \alpha} = h_{n, \alpha} \otimes 1_{L_n} :
K_n \otimes_{A_n} L_n
\longrightarrow
M_n \otimes_{A_n} L_n.
$$
Cas (\ref{crystalline-item-wedge}). On peut utiliser l'homotopie $\wedge^ih$ définie
en degré $n$ en posant
$$
\wedge^i(h)_{n, \alpha} = \wedge^i(h_{n, \alpha}) :
\wedge^i_{A_n}(K_n)
\longrightarrow
\wedge^i_{A_n}(M_n).
$$
Cas (\ref{crystalline-item-base-change}). On peut utiliser l'homotopie $h \otimes 1$ définie
en degré $n$ en posant
$$
(h \otimes 1)_{n, \alpha} = h_{n, \alpha} \otimes 1 :
K_n \otimes_{A_n} B_n
\longrightarrow
M_n \otimes_{A_n} B_n.
$$
Cas (\ref{crystalline-item-completion}). On peut utiliser l'homotopie $h^\wedge$ définie
en degré $n$ en posant
$$
(h^\wedge)_{n, \alpha} = h_{n, \alpha}^\wedge :
K_n^\wedge
\longrightarrow
M_n^\wedge.
$$
Cette construction convient, car chaque $h_{n, \alpha}$ est $A_n$-linéaire.
\end{proof}






\section{Cristaux en modules quasi-cohérents}
\label{crystalline-section-quasi-coherent-crystals}

\noindent
Dans la situation \ref{crystalline-situation-affine}.
Posons $X = \Spec(C)$ et $S = \Spec(A)$. Nous allons
classifier les cristaux en modules quasi-cohérents sur $\text{Cris}(X/S)$.
Auparavant, fixons quelques notations.

\medskip\noindent
Choisissons un anneau de polynômes $P = A[x_i]$ sur $A$ et une surjection $P \to C$
de $A$-algèbres, de noyau $J = \Ker(P \to C)$. Posons
\begin{equation}
\label{crystalline-equation-D}
D = \lim_e D_{P, \gamma}(J) / p^eD_{P, \gamma}(J)
\end{equation}
pour désigner l'enveloppe à puissances divisées complétée $p$-adiquement.
Cet anneau est muni d'un idéal à puissances divisées $\bar J$ et d'une structure de
puissances divisées $\bar \gamma$ ; voir le lemme \ref{crystalline-lemma-list-properties}.
Posons $D_e = D/p^eD$ et notons $\bar J_e$ l'image de $\bar J$ dans $D_e$.
Nous emploierons l'abréviation
\begin{equation}
\label{crystalline-equation-omega-D}
\Omega_D = \lim_e \Omega_{D_e/A, \bar\gamma} =
\lim_e \Omega_{D/A, \bar\gamma}/p^e\Omega_{D/A, \bar\gamma}
\end{equation}
pour le complété $p$-adique du module des différentielles à puissances divisées ;
voir le lemme \ref{crystalline-lemma-differentials-completion}.
C'est aussi le complété $p$-adique de
$\Omega_{D_{P, \gamma}(J)/A, \bar\gamma}$
qui est libre, ayant pour base les $\text{d}x_i$ ; voir le
lemme \ref{crystalline-lemma-module-differentials-divided-power-envelope}.
Ainsi, tout élément de $\Omega_D$ s'écrit de manière unique comme une somme
$\sum f_i\text{d}x_i$ telle que, pour tout $e$, seuls un nombre fini de $f_i$
n'appartiennent pas à $p^eD$. De plus, les applications
$\text{d}_{D_e/A, \bar\gamma} : D_e \to \Omega_{D_e/A, \bar\gamma}$
sont compatibles et définissent une $A$-dérivation à puissances divisées
\begin{equation}
\label{crystalline-equation-derivation-D}
\text{d} : D \longrightarrow \Omega_D
\end{equation}
après complétion $p$-adique.

\medskip\noindent
Nous aurons également besoin des « produits $\Spec(D(n))$ de $\Spec(D)$ » ; voir la
proposition \ref{crystalline-proposition-compute-cohomology} et sa démonstration pour une
explication. Formellement, ils sont définis comme suit. Pour $n \geq 0$, soit
$J(n) = \Ker(P \otimes_A \ldots \otimes_A P \to C)$, où
le produit tensoriel comporte $n + 1$ facteurs. On pose
\begin{equation}
\label{crystalline-equation-Dn}
D(n) = \lim_e
D_{P \otimes_A \ldots \otimes_A P, \gamma}(J(n))/
p^eD_{P \otimes_A \ldots \otimes_A P, \gamma}(J(n))
\end{equation}
qui est le complété $p$-adique de l'enveloppe à puissances divisées.
On note $\bar J(n)$ son idéal à puissances divisées et $\bar \gamma(n)$
sa structure de puissances divisées. On introduit aussi $D(n)_e = D(n)/p^eD(n)$ ainsi
que le module des différentielles complété $p$-adiquement
\begin{equation}
\label{crystalline-equation-omega-Dn}
\Omega_{D(n)} = \lim_e \Omega_{D(n)_e/A, \bar\gamma} =
\lim_e \Omega_{D(n)/A, \bar\gamma}/p^e\Omega_{D(n)/A, \bar\gamma}
\end{equation}
et la dérivation
\begin{equation}
\label{crystalline-equation-derivation-Dn}
\text{d} : D(n) \longrightarrow \Omega_{D(n)}
\end{equation}
Bien entendu, $D = D(0)$. Notons que les anneaux $D(0), D(1), D(2), \ldots$
forment un objet cosimplicial de la catégorie des anneaux à puissances divisées.

\begin{lemma}
\label{crystalline-lemma-structure-Dn}
Soient $D$ et $D(n)$ comme dans (\ref{crystalline-equation-D}) et (\ref{crystalline-equation-Dn}).
La coprojection $P \to P \otimes_A \ldots \otimes_A P$,
$f \mapsto f \otimes 1 \otimes \ldots \otimes 1$
induit un isomorphisme
\begin{equation}
\label{crystalline-equation-structure-Dn}
D(n) = \lim_e D\langle \xi_i(j) \rangle/p^eD\langle \xi_i(j) \rangle
\end{equation}
de $D$-algèbres, où
$$
\xi_i(j) = x_i \otimes 1 \otimes \ldots \otimes 1 -
1 \otimes \ldots \otimes 1 \otimes x_i \otimes 1 \otimes \ldots \otimes 1
$$
pour $j = 1, \ldots, n$, le second $x_i$ étant placé dans le facteur numéro $j + 1$
facteur ; rappelons que la construction de $D(n)$ part du
produit tensoriel de $n + 1$ copies de $P$ sur $A$.
\end{lemma}

\begin{proof}
On a
$$
P \otimes_A \ldots \otimes_A P = P[\xi_i(j)]
$$
et $J(n)$ est engendré par $J$ et les éléments $\xi_i(j)$.
Le lemme résulte donc du
lemme \ref{crystalline-lemma-divided-power-envelope-add-variables}.
\end{proof}

\begin{lemma}
\label{crystalline-lemma-property-Dn}
Soient $D$ et $D(n)$ comme dans (\ref{crystalline-equation-D}) et (\ref{crystalline-equation-Dn}).
Alors $(D, \bar J, \bar\gamma)$ et $(D(n), \bar J(n), \bar\gamma(n))$
sont des objets de $\text{Cris}^\wedge(C/A)$ (voir la
remarque \ref{crystalline-remark-completed-affine-site}), et
$$
D(n) = \coprod\nolimits_{j = 0, \ldots, n} D
$$
dans $\text{Cris}^\wedge(C/A)$.
\end{lemma}

\begin{proof}
La première assertion est claire. Pour la seconde, si $(B \to C, \delta)$ est un
objet de $\text{Cris}^\wedge(C/A)$, alors on a
$$
\Mor_{\text{Cris}^\wedge(C/A)}(D, B) = 
\Hom_A((P, J), (B, \Ker(B \to C)))
$$
et de même pour $D(n)$ en remplaçant $(P, J)$ par
$(P \otimes_A \ldots \otimes_A P, J(n))$. La propriété relative aux coproduits en résulte,
puisque $P \otimes_A \ldots \otimes_A P$ est un coproduit.
\end{proof}

\noindent
Dans le lemme ci-dessous, nous considérerons des couples $(M, \nabla)$ satisfaisant aux
conditions suivantes :
\begin{enumerate}
\item
\label{crystalline-item-complete}
$M$ est $p$-adiquement complet en tant que $D$-module,
\item
\label{crystalline-item-connection}
$\nabla : M \to M \otimes^\wedge_D \Omega_D$ est une connexion, c'est-à-dire
$\nabla(fm) = m \otimes \text{d}f + f\nabla(m)$,
\item
\label{crystalline-item-integrable}
$\nabla$ est intégrable
(voir la remarque \ref{crystalline-remark-connection}), et
\item
\label{crystalline-item-topologically-quasi-nilpotent}
$\nabla$ est {\it topologiquement quasi-nilpotente} : si l'on écrit
$\nabla(m) = \sum \theta_i(m)\text{d}x_i$ pour des opérateurs
$\theta_i : M \to M$, alors, pour tout $m \in M$, il n'existe qu'un nombre fini
de couples $(i, k)$ tels que $\theta_i^k(m) \not \in pM$.
\end{enumerate}
Dans la littérature, les opérateurs $\theta_i$ sont parfois notés
$\nabla_{\partial/\partial x_i}$.
Dans le lemme suivant, nous construisons un foncteur de la catégorie des cristaux en modules
quasi-cohérents sur $\text{Cris}(X/S)$ vers la catégorie de ces couples. Nous montrerons que
ce foncteur est une équivalence dans la
proposition \ref{crystalline-proposition-crystals-on-affine}.

\begin{lemma}
\label{crystalline-lemma-crystals-on-affine}
Dans la situation ci-dessus, il existe un foncteur
$$
\begin{matrix}
\text{cristaux en} \\
\mathcal{O}_{X/S}\text{-modules quasi-cohérents sur }\text{Cris}(X/S)
\end{matrix}
\longrightarrow
\begin{matrix}
\text{couples }(M, \nabla)\text{ vérifiant} \\
\text{(\ref{crystalline-item-complete}), (\ref{crystalline-item-connection}),
(\ref{crystalline-item-integrable}) et (\ref{crystalline-item-topologically-quasi-nilpotent})}
\end{matrix}
$$
\end{lemma}

\begin{proof}
Soit $\mathcal{F}$ un cristal en modules quasi-cohérents sur $X/S$.
Posons $T_e = \Spec(D_e)$, de sorte que $(X, T_e, \bar\gamma)$ est un objet
de $\text{Cris}(X/S)$ pour $e \gg 0$. On dispose de morphismes
$$
(X, T_e, \bar\gamma) \to (X, T_{e + 1}, \bar\gamma) \to \ldots
$$
qui sont des immersions fermées. Posons
$$
M =
\lim_e \Gamma((X, T_e, \bar\gamma), \mathcal{F}) =
\lim_e \Gamma(T_e, \mathcal{F}_{T_e}) = \lim_e M_e
$$
Notons que, puisque $\mathcal{F}$ est localement quasi-cohérent, on a
$\mathcal{F}_{T_e} = \widetilde{M_e}$. Puisque $\mathcal{F}$ est un
cristal, on a $M_e = M_{e + 1}/p^eM_{e + 1}$. On voit donc que
$M_e = M/p^eM$ et que $M$ est complet $p$-adiquement ; voir
Algèbre, lemme \ref{algebra-lemma-limit-complete}.

\medskip\noindent
Par le lemme \ref{crystalline-lemma-automatic-connection}, on sait que $\mathcal{F}$
est muni d'une connexion intégrable canonique
$\nabla : \mathcal{F} \to \mathcal{F} \otimes \Omega_{X/S}$.
En évaluant cette connexion sur les objets $T_e$ construits ci-dessus,
on obtient une connexion intégrable canonique
$$
\nabla : M \longrightarrow M \otimes^\wedge_D \Omega_D
$$
Pour voir qu'elle est topologiquement quasi-nilpotente, explicitons ce que cela signifie.

\medskip\noindent
On peut maintenant appliquer le même procédé aux anneaux $D(n)$.
On obtient ainsi un module $p$-adiquement complet sur $D(n)$, noté $M(n)$. En utilisant encore
la propriété de cristal de $\mathcal{F}$, on obtient des isomorphismes
$$
M \otimes^\wedge_{D, p_0} D(1) \rightarrow M(1)
\leftarrow M \otimes^\wedge_{D, p_1} D(1)
$$
voir la démonstration du lemme \ref{crystalline-lemma-automatic-connection}.
Notons $c$ le composé de gauche à droite. Prenons $m \in M$.
Posons $\xi_i = x_i \otimes 1 - 1 \otimes x_i$.
D'après (\ref{crystalline-equation-structure-Dn}), on peut écrire de manière unique
$$
c(m \otimes 1) = \sum\nolimits_K \theta_K(m) \otimes \prod \xi_i^{[k_i]}
$$
avec $\theta_K(m) \in M$, où la somme porte sur les multi-indices
$K = (k_i)$ tels que $k_i \geq 0$ et $\sum k_i < \infty$. Posons
$\theta_i = \theta_K$, où $K$ a un $1$ à la $i$-ième place et
des zéros ailleurs. On a
$$
\nabla(m) = \sum \theta_i(m) \text{d}x_i.
$$
Cela se voit en comparant avec la définition de
$\nabla$. En effet, l'équation qui le définit est
$p_1^*m = \nabla(m) - c(p_0^*m)$ dans le lemme \ref{crystalline-lemma-automatic-connection},
mais le signe convient parce que, dans le projet Champs, on utilise systématiquement
$\text{d}f = p_1(f) - p_0(f)$ modulo le carré de l'idéal de la diagonale,
et que $\xi_i = x_i \otimes 1 - 1 \otimes x_i$ s'envoie donc sur $-\text{d}x_i$
modulo le carré de l'idéal de la diagonale.

\medskip\noindent
Notons $q_i : D \to D(2)$ et $q_{ij} : D(1) \to D(2)$ les coprojections
correspondant aux indices $i, j$. Comme dans le dernier paragraphe de la démonstration du
lemme \ref{crystalline-lemma-automatic-connection},
on voit que
$$
q_{02}^*c = q_{12}^*c \circ q_{01}^*c.
$$
Cela signifie que
$$
\sum\nolimits_{K''} \theta_{K''}(m) \otimes \prod {\zeta''_i}^{[k''_i]}
=
\sum\nolimits_{K', K} \theta_{K'}(\theta_K(m))
\otimes \prod {\zeta'_i}^{[k'_i]} \prod \zeta_i^{[k_i]}
$$
dans $M \otimes^\wedge_{D, q_2} D(2)$, où
\begin{align*}
\zeta_i & = x_i \otimes 1 \otimes 1 - 1 \otimes x_i \otimes 1,\\
\zeta'_i & = 1 \otimes x_i \otimes 1 - 1 \otimes 1 \otimes x_i,\\
\zeta''_i & = x_i \otimes 1 \otimes 1 - 1 \otimes 1 \otimes x_i.
\end{align*}
En particulier, $\zeta''_i = \zeta_i + \zeta'_i$, et
$D(2)$ est le complété $p$-adique de l'anneau de polynômes
à puissances divisées en $\zeta_i, \zeta'_i$ sur $q_2(D)$ ; voir le lemme \ref{crystalline-lemma-structure-Dn}.
En comparant les coefficients dans l'expression ci-dessus, on obtient immédiatement
$\theta_i \circ \theta_j = \theta_j \circ \theta_i$
(ce qui fournit une autre démonstration de l'intégrabilité de $\nabla$) et
$$
\theta_K(m) = (\prod \theta_i^{k_i})(m).
$$
En particulier, puisque la somme exprimant $c(m \otimes 1)$ ci-dessus doit converger
$p$-adiquement, on conclut que, pour tout $i$ et tout $m \in M$, les termes
$\theta_i^k(m)$ non nuls modulo $p$ sont en nombre fini.
\end{proof}

\begin{proposition}
\label{crystalline-proposition-crystals-on-affine}
Le foncteur
$$
\begin{matrix}
\text{cristaux en} \\
\mathcal{O}_{X/S}\text{-modules quasi-cohérents sur }\text{Cris}(X/S)
\end{matrix}
\longrightarrow
\begin{matrix}
\text{couples }(M, \nabla)\text{ vérifiant} \\
\text{(\ref{crystalline-item-complete}), (\ref{crystalline-item-connection}),
(\ref{crystalline-item-integrable}) et (\ref{crystalline-item-topologically-quasi-nilpotent})}
\end{matrix}
$$
du lemme \ref{crystalline-lemma-crystals-on-affine}
est une équivalence de catégories.
\end{proposition}

\begin{proof}
Soit $(M, \nabla)$ donné. Construisons
un cristal en modules quasi-cohérents $\mathcal{F}$.
Écrivons $\nabla(m) = \sum \theta_i(m)\text{d}x_i$.
Alors $\theta_i \circ \theta_j = \theta_j \circ \theta_i$, et l'on
peut poser $\theta_K(m) = (\prod \theta_i^{k_i})(m)$ pour tout multi-indice
$K = (k_i)$ tel que $k_i \geq 0$ et $\sum k_i < \infty$.

\medskip\noindent
Soit $(U, T, \delta)$ un objet quelconque de $\text{Cris}(X/S)$ avec $T$ affine.
Écrivons $T = \Spec(B)$ ; l'idéal de $U \to T$ sera noté $J_B \subset B$.
Par le lemme \ref{crystalline-lemma-set-generators}, il existe un entier $e$ et un morphisme
$$
f : (U, T, \delta) \longrightarrow (X, T_e, \bar\gamma)
$$
où $T_e = \Spec(D_e)$, comme dans la démonstration du
lemme \ref{crystalline-lemma-crystals-on-affine}.
Choisissons de tels $e$ et $f$ ; notons encore $f : D \to B$ l'homomorphisme
correspondant de $A$-algèbres à puissances divisées. On définit $\mathcal{F}_T$ comme le
faisceau quasi-cohérent de $\mathcal{O}_T$-modules associé au $B$-module
$$
M \otimes_{D, f} B.
$$
Il faut toutefois montrer que cette définition ne dépend pas du choix de $f$.
Supposons que $g : D \to B$ soit un second morphisme de ce type. Puisque $f$ et $g$
sont des morphismes de $\text{Cris}(X/S)$, on voit que l'image de
$f - g : D \to B$ est contenue dans l'idéal à puissances divisées $J_B$.
Posons $\xi_i = f(x_i) - g(x_i) \in J_B$. Par analogie avec la démonstration
du lemme \ref{crystalline-lemma-crystals-on-affine}, on définit un isomorphisme
$$
c_{f, g} : M \otimes_{D, f} B \longrightarrow M \otimes_{D, g} B
$$
par la formule
$$
m \otimes 1 \longmapsto 
\sum\nolimits_K \theta_K(m) \otimes \prod \xi_i^{[k_i]}
$$
qui a un sens d'après les remarques ci-dessus et le fait que $\nabla$
est topologiquement quasi-nilpotente (la somme est donc finie !).
Un calcul montre que
$$
c_{g, h} \circ c_{f, g} = c_{f, h}
$$
si l'on se donne un troisième morphisme
$h : (U, T, \delta) \longrightarrow (X, T_e, \bar\gamma)$.
On a aussi $c_{f, f} = 1$.
Ces applications sont donc toutes des isomorphismes, et l'on voit que
le module $\mathcal{F}_T$ ne dépend pas du choix de $f$.

\medskip\noindent
Si $a : (U', T', \delta') \to (U, T, \delta)$ est un morphisme entre objets affines
de $\text{Cris}(X/S)$, alors, en choisissant $f' = f \circ a$, il est clair
qu'il existe un isomorphisme canonique
$a^*\mathcal{F}_T \to \mathcal{F}_{T'}$. On omet de vérifier que cette
application ne dépend pas du choix de $f$. En prenant ces applications comme applications de restriction,
il est clair que l'on obtient un cristal en modules quasi-cohérents
sur la sous-catégorie pleine de $\text{Cris}(X/S)$ formée des objets affines.
On omet de montrer que celui-ci s'étend en un cristal sur tout
$\text{Cris}(X/S)$. On omet aussi de montrer que ce procédé est un foncteur
et qu'il est quasi-inverse du foncteur construit dans le
lemme \ref{crystalline-lemma-crystals-on-affine}.
\end{proof}

\begin{lemma}
\label{crystalline-lemma-crystals-on-affine-smooth}
Dans la situation \ref{crystalline-situation-affine}.
Soient $A \to P' \to C$ des homomorphismes d'anneaux tels que $A \to P'$ soit lisse et $P' \to C$
surjectif de noyau $J'$. Soit $D'$ le complété $p$-adique de
$D_{P', \gamma}(J')$. Alors il existe un choix de données $A \to P \to C$
comme ci-dessus et des homomorphismes de $A$-algèbres à puissances divisées
$$
a : D \longrightarrow D',\quad b : D' \longrightarrow D
$$
compatibles avec les applications $D \to C$ et $D' \to C$ tels que
$a \circ b = \text{id}_{D'}$. Ces applications induisent
une équivalence entre la catégorie des couples $(M, \nabla)$ vérifiant
(\ref{crystalline-item-complete}), (\ref{crystalline-item-connection}),
(\ref{crystalline-item-integrable}) et (\ref{crystalline-item-topologically-quasi-nilpotent})
sur $D$ et celle des couples $(M', \nabla')$ vérifiant
(\ref{crystalline-item-complete}), (\ref{crystalline-item-connection}),
(\ref{crystalline-item-integrable}) et
(\ref{crystalline-item-topologically-quasi-nilpotent})\footnote{Cette condition
est délicate à formuler pour $(M', \nabla')$ sur $D'$. Voir la démonstration.} sur $D'$.
En particulier, l'équivalence de catégories de la
proposition \ref{crystalline-proposition-crystals-on-affine}
vaut aussi pour le foncteur correspondant à valeurs dans la catégorie des couples sur $D'$.
\end{lemma}

\begin{proof}
Choisissons une algèbre polynomiale $P = A[y_1, \ldots, y_m]$ sur $A$ et une
surjection $P \to P'$. Définissons $P \to C$ comme le composé $P \to P' \to C$.
On obtient un homomorphisme surjectif $a : D \to D'$ par fonctorialité des
enveloppes à puissances divisées et de la complétion. Prenons $e$ assez grand pour que
$D_e$ soit un épaississement à puissances
divisées de $C$ sur $A$. Alors $D_e \to C$ est une surjection dont le noyau
est localement nilpotent ; voir Algèbre à puissances divisées, lemme \ref{dpa-lemma-nil}.
En posant $D'_e = D'/p^eD'$,
on voit que le noyau de $D_e \to D'_e$ est localement nilpotent.
Par conséquent, d'après Algèbre, lemme \ref{algebra-lemma-smooth-strong-lift},
on peut trouver un relèvement $\beta_e : P' \to D_e$ de l'application $P' \to D'_e$.
Notons que $D_{e + i + 1} \to D_{e + i} \times_{D'_{e + i}} D'_{e + i + 1}$
est surjectif de noyau de carré nul pour tout $i \geq 0$, puisque
$p^{e + i}D \to p^{e + i}D'$ est surjectif. En appliquant successivement la propriété
usuelle de relèvement (Algèbre, proposition \ref{algebra-proposition-smooth-formally-smooth})
aux diagrammes
$$
\xymatrix{
P' \ar[r] & D_{e + i} \times_{D'_{e + i}} D'_{e + i + 1} \\
A \ar[u] \ar[r] & D_{e + i + 1} \ar[u]
}
$$
on voit qu'il existe un homomorphisme de $A$-algèbres $\beta : P' \to D$ dont
le composé avec $a$ est l'application donnée $P' \to D'$.
Par la propriété universelle de l'enveloppe à puissances divisées, on obtient un
homomorphisme $D_{P', \gamma}(J') \to D$. Comme $D$ est $p$-adiquement complet, on
obtient $b : D' \to D$ tel que $a \circ b = \text{id}_{D'}$.

\medskip\noindent
Considérons les foncteurs de changement de base
$$
F : (M, \nabla) \longmapsto
(M \otimes^\wedge_{D, a} D', \nabla')
\quad\text{et}\quad
G : (M', \nabla') \longmapsto
(M' \otimes^\wedge_{D', b} D, \nabla)
$$
sur les modules à connexion vérifiant (\ref{crystalline-item-complete}),
(\ref{crystalline-item-connection}) et (\ref{crystalline-item-integrable}).
Voir la remarque \ref{crystalline-remark-base-change-connection}.
Puisque $a \circ b = \text{id}_{D'}$, on voit que
$F \circ G$ est le foncteur identité. Disons que $(M', \nabla')$
possède la propriété (\ref{crystalline-item-topologically-quasi-nilpotent}) si elle
est satisfaite par $G(M', \nabla')$. Un argument formel montre alors que, pour achever
la démonstration, il suffit de montrer que $G(F(M, \nabla))$ est isomorphe
à $(M, \nabla)$ lorsque $(M, \nabla)$ satisfait les quatre
conditions (\ref{crystalline-item-complete}), (\ref{crystalline-item-connection}),
(\ref{crystalline-item-integrable}) et (\ref{crystalline-item-topologically-quasi-nilpotent}).
Pour cela, on utilise l'isomorphisme fonctoriel
$$
c_{\text{id}_D, b \circ a} :
M \otimes_{D, \text{id}_D} D
\longrightarrow
M \otimes_{D, b \circ a} D
$$
de la démonstration de la proposition \ref{crystalline-proposition-crystals-on-affine}
(qui requiert la quasi-nilpotence topologique de $\nabla$,
supposée ici).
Il reste à montrer que ce morphisme est horizontal, c'est-à-dire
compatible avec les connexions, ce que l'on omet.

\medskip\noindent
La dernière assertion de l'énoncé en résulte.
\end{proof}

\begin{remark}
\label{crystalline-remark-equivalence-more-general}
L'équivalence de la proposition \ref{crystalline-proposition-crystals-on-affine}
vaut si l'on part d'une surjection $P \to C$ telle que $P/A$ satisfasse à la
propriété forte de relèvement du
lemme \ref{algebra-lemma-smooth-strong-lift} d'Algèbre.
Pour le montrer, on peut raisonner comme dans la démonstration du
lemme \ref{crystalline-lemma-crystals-on-affine-smooth}.
(On ajoutera ici les détails si le besoin s'en fait sentir.)
Il existe vraisemblablement aussi une démonstration directe de ce résultat, mais l'avantage
d'utiliser des anneaux de polynômes est que les anneaux $D(n)$ sont les complétés $p$-adiques
d'anneaux de polynômes à puissances divisées, ce qui simplifie l'algèbre.
\end{remark}



\section{Remarques générales sur la cohomologie}
\label{crystalline-section-cohomology-lqc}

\noindent
Dans cette section, on effectue quelques préparatifs pour ramener la cohomologie
des modules sur le site cristallin d'un schéma affine à
une question algébrique.

\begin{lemma}
\label{crystalline-lemma-vanishing-lqc}
Dans la situation \ref{crystalline-situation-global}.
Soit $\mathcal{F}$ un $\mathcal{O}_{X/S}$-module localement quasi-cohérent
sur $\text{Cris}(X/S)$. Alors on a
$$
H^p((U, T, \delta), \mathcal{F}) = 0
$$
pour tout $p > 0$ et tout $(U, T, \delta)$ tel que $T$ ou $U$ soit affine.
\end{lemma}

\begin{proof}
Comme $U \to T$ est un épaississement, on voit que $U$ est affine si et seulement si $T$
est affine ; voir Limites, lemme \ref{limits-lemma-affine}.
Cela étant, appliquons le
lemme \ref{sites-cohomology-lemma-cech-vanish-collection} de Cohomologie sur les sites
à la collection $\mathcal{B}$ des objets affines $(U, T, \delta)$ et à la
collection $\text{Cov}$ des recouvrements ouverts affines
$\mathcal{U} = \{(U_i, T_i, \delta_i) \to (U, T, \delta)\}$. Le
complexe de {\v C}ech
${\check C}^*(\mathcal{U}, \mathcal{F})$ associé à un tel recouvrement est simplement
le complexe de {\v C}ech du $\mathcal{O}_T$-module quasi-cohérent
$\mathcal{F}_T$
(on utilise ici l'hypothèse selon laquelle $\mathcal{F}$ est localement quasi-cohérent)
relativement au recouvrement ouvert affine $\{T_i \to T\}$ du
schéma affine $T$. La cohomologie de {\v C}ech est donc nulle par
Cohomologie des schémas, lemmes
\ref{coherent-lemma-cech-cohomology-quasi-coherent} et
\ref{coherent-lemma-quasi-coherent-affine-cohomology-zero}.
Ainsi, les hypothèses du
lemme \ref{sites-cohomology-lemma-cech-vanish-collection} de Cohomologie sur les sites
sont satisfaites, ce qui conclut.
\end{proof}

\begin{lemma}
\label{crystalline-lemma-compare}
Dans la situation \ref{crystalline-situation-global}.
Supposons en outre que $X$ et $S$ soient des schémas affines.
Considérons la sous-catégorie pleine $\mathcal{C} \subset \text{Cris}(X/S)$
formée des épaississements à puissances divisées $(X, T, \delta)$
et munie de la topologie chaotique (voir
Sites, exemple \ref{sites-example-indiscrete}).
Pour tout $\mathcal{O}_{X/S}$-module localement quasi-cohérent $\mathcal{F}$,
on a
$$
R\Gamma(\mathcal{C}, \mathcal{F}|_\mathcal{C}) =
R\Gamma(\text{Cris}(X/S), \mathcal{F})
$$
\end{lemma}

\begin{proof}
Notons $\text{AffineCris}(X/S)$ la sous-catégorie pleine de $\text{Cris}(X/S)$
formée des objets $(U, T, \delta)$ tels que $U$ et $T$ soient affines.
On en fait un site en déclarant qu'une famille de morphismes
$\{(U_i, T_i, \delta_i) \to (U, T, \delta)\}_{i \in I}$ de
$\text{AffineCris}(X/S)$ est un recouvrement si et seulement si elle en est un
dans $\text{Cris}(X/S)$. Avec cette définition, le foncteur d'inclusion
$$
\text{AffineCris}(X/S) \longrightarrow \text{Cris}(X/S)
$$
est un foncteur cocontinu spécial au sens de
Sites, définition \ref{sites-definition-special-cocontinuous-functor}.
La démonstration est exactement la même que celle du
lemme \ref{topologies-lemma-affine-big-site-Zariski} de Topologies.
On voit ainsi que le topos des faisceaux sur $\text{Cris}(X/S)$
est le même que celui des faisceaux sur $\text{AffineCris}(X/S)$
par restriction le long du foncteur d'inclusion affiché.
Il faut donc montrer l'énoncé correspondant pour
l'inclusion $\mathcal{C} \subset \text{AffineCris}(X/S)$.

\medskip\noindent
On utilisera sans autre mention le fait que $\mathcal{C}$ et
$\text{AffineCris}(X/S)$ admettent des produits et des produits fibrés
(on omet les détails ; voir le
lemme \ref{crystalline-lemma-divided-power-thickening-fibre-products}).
Le foncteur d'inclusion $u : \mathcal{C} \to \text{AffineCris}(X/S)$
est pleinement fidèle, continu et commute aux produits et aux produits fibrés.
On affirme qu'il définit un morphisme de sites annelés
$$
f :
(\text{AffineCris}(X/S), \mathcal{O}_{X/S})
\longrightarrow
(\Sh(\mathcal{C}), \mathcal{O}_{X/S}|_\mathcal{C})
$$
Pour le voir, on utilisera Sites, lemme \ref{sites-lemma-directed-morphism}.
Notons que $\mathcal{C}$ admet des produits fibrés et que $u$ y commute ;
les catégories $\mathcal{I}^u_{(U, T, \delta)}$ sont donc des réunions disjointes
de catégories filtrantes (par Sites, lemme \ref{sites-lemma-almost-directed}, et
Catégories, lemme \ref{categories-lemma-split-into-directed}). Il
suffit donc de montrer que $\mathcal{I}^u_{(U, T, \delta)}$ est connexe.
Le fait qu'elle soit non vide résulte du lemme \ref{crystalline-lemma-set-generators} : puisque $U$ et $T$
sont affines, ce lemme affirme qu'il existe au moins un objet
$(X, T', \delta')$ de $\mathcal{C}$ et un morphisme
$(U, T, \delta) \to (X, T', \delta')$ d'épaississements à puissances divisées.
La connexité résulte du fait que $\mathcal{C}$ admet des produits
et que $u$ y commute (comparer avec la démonstration du
lemme \ref{sites-lemma-directed} de Sites).

\medskip\noindent
Notons que $f_*\mathcal{F} = \mathcal{F}|_\mathcal{C}$. Pour démontrer le lemme, il suffit donc
d'établir que $R^pf_*\mathcal{F} = 0$ pour $p > 0$ ; voir
Cohomologie sur les sites, lemme \ref{sites-cohomology-lemma-apply-Leray}. Par
Cohomologie sur les sites, lemme \ref{sites-cohomology-lemma-higher-direct-images},
il suffit de montrer que
$H^p(\text{AffineCris}(X/S)/(X, T, \delta), \mathcal{F}) = 0$
pour tout $(X, T, \delta)$.
Cela résulte du lemme \ref{crystalline-lemma-vanishing-lqc}, car le
topos du site $\text{AffineCris}(X/S)/(X, T, \delta)$
est équivalent au topos du site
$\text{Cris}(X/S)/(X, T, \delta)$ utilisé dans ce lemme.
\end{proof}

\begin{lemma}
\label{crystalline-lemma-complete}
Dans la situation \ref{crystalline-situation-affine}.
Posons $\mathcal{C} = (\text{Cris}(C/A))^{opp}$ et
$\mathcal{C}^\wedge = (\text{Cris}^\wedge(C/A))^{opp}$
munis de la topologie chaotique ; pour les notations, voir la
remarque \ref{crystalline-remark-completed-affine-site}.
Il existe un morphisme de topos
$$
g : \Sh(\mathcal{C}) \longrightarrow \Sh(\mathcal{C}^\wedge)
$$
tel que, si $\mathcal{F}$ est un faisceau de groupes abéliens sur
$\mathcal{C}$, alors
$$
R^pg_*\mathcal{F}(B \to C, \delta) =
\left\{
\begin{matrix}
\lim_e \mathcal{F}(B_e \to C, \delta) & \text{si }p = 0 \\
R^1\lim_e \mathcal{F}(B_e \to C, \delta) & \text{si }p = 1 \\
0 & \text{sinon}
\end{matrix}
\right.
$$
où $B_e = B/p^eB$ pour $e \gg 0$.
\end{lemma}

\begin{proof}
Tout foncteur entre catégories définit, dans le même sens, un morphisme de
topos chaotiques ; en effet, un tel foncteur
peut être considéré comme un foncteur cocontinu entre sites ; voir
Sites, section \ref{sites-section-cocontinuous-morphism-topoi}.
On omet la démonstration de la description de $g_*\mathcal{F}$.
Notons que, dans l'énoncé, $(B_e \to C, \delta)$
n'est un objet de $\text{Cris}(C/A)$ que pour $e$ assez grand.
Soit $\mathcal{I}$ un faisceau abélien injectif sur $\mathcal{C}$.
Alors les applications de transition
$$
\mathcal{I}(B_e \to C, \delta) \leftarrow
\mathcal{I}(B_{e + 1} \to C, \delta)
$$
sont surjectives, car les morphismes
$$
(B_e \to C, \delta)
\longrightarrow
(B_{e + 1} \to C, \delta)
$$
sont des monomorphismes dans la catégorie $\mathcal{C}$. Ainsi, pour un faisceau
abélien injectif, les deux membres de la formule affichée dans le lemme coïncident.
En prenant une résolution injective de $\mathcal{F}$, on obtient aisément
le résultat (les faisceaux sont des préfaisceaux, de sorte que l'exactitude se mesure au
niveau des groupes de sections sur les objets).
\end{proof}

\begin{lemma}
\label{crystalline-lemma-category-with-covering}
Soit $\mathcal{C}$ une catégorie munie de la topologie chaotique.
Soit $X$ un objet de $\mathcal{C}$ tel que tout objet de
$\mathcal{C}$ admette un morphisme vers $X$. Supposons que $\mathcal{C}$
admette les produits de deux objets.
Alors, pour tout faisceau abélien $\mathcal{F}$ sur $\mathcal{C}$,
la cohomologie totale $R\Gamma(\mathcal{C}, \mathcal{F})$ est représentée
par le complexe
$$
\mathcal{F}(X) \to \mathcal{F}(X \times X) \to
\mathcal{F}(X \times X \times X) \to \ldots
$$
associé au groupe abélien cosimplicial $[n] \mapsto \mathcal{F}(X^{n + 1})$.
\end{lemma}

\begin{proof}
Notons que $H^q(X^p, \mathcal{F}) = 0$ pour tout $q > 0$, puisque tout préfaisceau est un
faisceau sur $\mathcal{C}$. L'hypothèse sur $X$ signifie que $h_X \to *$
est surjectif. En utilisant $H^q(X, \mathcal{F}) = H^q(h_X, \mathcal{F})$ et
$H^q(\mathcal{C}, \mathcal{F}) = H^q(*, \mathcal{F})$, on voit que notre
énoncé résulte de
Cohomologie sur les sites,
lemme \ref{sites-cohomology-lemma-cech-to-cohomology-sheaf-sets}.
\end{proof}









\section{Préparatifs cosimpliciaux}
\label{crystalline-section-cohomology}

\noindent
Dans cette section, on compare la cohomologie cristalline à la cohomologie
de de Rham. On suit \cite{Bhatt}.

\begin{example}
\label{crystalline-example-cosimplicial-module}
Supposons que $A_*$ soit un anneau cosimplicial quelconque.
Considérons le module cosimplicial $M_*$ défini par la règle
$$
M_n = \bigoplus\nolimits_{i = 0, ..., n} A_n e_i
$$
Pour une application $f : [n] \to [m]$, définissons $M_*(f) : M_n \to M_m$
comme l'unique application $A_*(f)$-linéaire envoyant $e_i$ sur $e_{f(i)}$.
On affirme que l'identité de $M_*$ est homotope à $0$.
En effet, une homotopie est donnée par un morphisme de modules cosimpliciaux
$$
h : M_* \longrightarrow \Hom(\Delta[1], M_*)
$$
voir la section \ref{crystalline-section-cosimplicial}.
Pour $j \in \{0, \ldots, n + 1\}$, notons $\alpha^n_j : [n] \to [1]$ l'application
définie par $\alpha^n_j(i) = 0 \Leftrightarrow i < j$. Alors
$\Delta[1]_n = \{\alpha^n_0, \ldots, \alpha^n_{n + 1}\}$ et, par conséquent,
$\Hom(\Delta[1], M_*)_n = \prod_{j = 0, \ldots, n + 1} M_n$ ; voir
Objets simpliciaux, sections \ref{simplicial-section-homotopy} et
\ref{simplicial-section-homotopy-cosimplicial}. Plutôt que d'utiliser
cette écriture comme produit, on considère un élément
de $\Hom(\Delta[1], M_*)_n$ comme une fonction $\Delta[1]_n \to M_n$.
Avec cette notation, on définit $h$ en degré $n$ par la règle
$$
h_n(e_i)(\alpha^n_j) =
\left\{
\begin{matrix}
e_{i} & \text{si} & i < j \\
0 & \text{sinon} 
\end{matrix}
\right.
$$
Vérifions d'abord que $h$ est un morphisme de modules cosimpliciaux. Pour
$f : [n] \to [m]$, on va montrer que
\begin{equation}
\label{crystalline-equation-cosimplicial-morphism}
h_m \circ M_*(f) = \Hom(\Delta[1], M_*)(f) \circ h_n
\end{equation}
Le membre de gauche de (\ref{crystalline-equation-cosimplicial-morphism}), évalué en
$e_i$, puis en $\alpha^m_j$, vaut
$$
h_m(e_{f(i)})(\alpha^m_j) =
\left\{
\begin{matrix}
e_{f(i)} & \text{si} & f(i) < j \\
0 & \text{sinon}
\end{matrix}
\right.
$$
Notons que $\alpha^m_j \circ f = \alpha^n_{j'}$, où
$0 \leq j' \leq n + 1$ est l'unique indice tel que $f(i) < j$
si et seulement si $i < j'$. Ainsi, le membre de droite de
(\ref{crystalline-equation-cosimplicial-morphism}), évalué en $e_i$,
puis en $\alpha^m_j$, vaut
$$
M_*(f)(h_n(e_i)(\alpha^m_j \circ f)) =
M_*(f)(h_n(e_i)(\alpha^n_{j'})) =
\left\{
\begin{matrix}
e_{f(i)} & \text{si} & i < j' \\
0 & \text{sinon} 
\end{matrix}
\right.
$$
Notre description de $j'$ montre que les deux résultats sont égaux.
Ainsi, $h$ est un morphisme de modules cosimpliciaux.
Soient $0 : \Delta[0] \to \Delta[1]$ et
$1 : \Delta[0] \to \Delta[1]$ les applications évidentes, et notons
$ev_0, ev_1 : \Hom(\Delta[1], M_*) \to M_*$ les applications
d'évaluation correspondantes. Le lecteur vérifie aisément que
les composés
$$
ev_0 \circ h, ev_1 \circ h : M_* \longrightarrow M_*
$$
sont respectivement $1$ et $0$ ; ainsi, $h$ est l'homotopie cherchée entre
$1$ et $0$.
\end{example}

\begin{lemma}
\label{crystalline-lemma-vanishing-omega-1}
Avec les notations de (\ref{crystalline-equation-omega-Dn}), le complexe
$$
\Omega_{D(0)} \to \Omega_{D(1)} \to \Omega_{D(2)} \to \ldots
$$
est homotope à zéro comme $D(*)$-module cosimplicial.
\end{lemma}

\begin{proof}
On va utiliser le principe de
Objets simpliciaux, lemme \ref{simplicial-lemma-functorial-homotopy},
et, plus précisément, le
lemme \ref{crystalline-lemma-homotopy-tensor},
qui affirme que tout foncteur transforme des applications homotopes entre objets
(co)simpliciaux en applications homotopes.
Le complexe du lemme est égal au complété $p$-adique du
changement de base du module cosimplicial
$$
M_* = \left(
\Omega_{P/A} \to
\Omega_{P \otimes_A P/A} \to
\Omega_{P \otimes_A P \otimes_A P/A} \to \ldots
\right)
$$
par l'homomorphisme d'anneaux cosimpliciaux $P\otimes_A \ldots \otimes_A P \to D(n)$. Cela
résulte du lemme \ref{crystalline-lemma-module-differentials-divided-power-envelope} ;
voir les commentaires après (\ref{crystalline-equation-omega-D}). Il
suffit donc de montrer que le module cosimplicial $M_*$ est homotope à zéro
(on utilise le changement de base et la complétion $p$-adique).
On peut même supposer que $A = \mathbf{Z}$ et $P = \mathbf{Z}[\{x_i\}_{i \in I}]$,
car on peut effectuer le changement de base par $\mathbf{Z} \to A$.
Dans ce cas, $P^{\otimes n + 1}$ est l'algèbre polynomiale engendrée par les éléments
$$
x_i(e) = 1 \otimes \ldots \otimes x_i \otimes \ldots \otimes 1
$$
où $x_i$ occupe la $e$-ième place. Les modules du complexe sont libres, ayant
pour base les $\text{d}x_i(e)$. Notons que, si $f : [n] \to [m]$ est une
application, alors
$$
M_*(f)(\text{d}x_i(e)) = \text{d}x_i(f(e))
$$
On voit donc que $M_*$ est une somme directe, indexée par $I$, de copies du module
étudié dans l'exemple \ref{crystalline-example-cosimplicial-module}, ce qui conclut.
\end{proof}

\begin{lemma}
\label{crystalline-lemma-vanishing}
Avec les notations de (\ref{crystalline-equation-Dn}) et (\ref{crystalline-equation-omega-Dn}),
pour tout module cosimplicial $M_*$ sur $D(*)$ et tout
$i > 0$, le module cosimplicial
$$
M_0 \otimes^\wedge_{D(0)} \Omega^i_{D(0)} \to
M_1 \otimes^\wedge_{D(1)} \Omega^i_{D(1)} \to
M_2 \otimes^\wedge_{D(2)} \Omega^i_{D(2)} \to \ldots
$$
est homotope à zéro, où $\Omega^i_{D(n)}$ est le complété $p$-adique
de la $i$-ième puissance extérieure de $\Omega_{D(n)}$.
\end{lemma}

\begin{proof}
Par le lemme \ref{crystalline-lemma-vanishing-omega-1}, les endomorphismes $0$ et $1$
de $\Omega_{D(*)}$ sont homotopes.
En appliquant le foncteur $\wedge^i$, on voit qu'il
en va de même pour le module cosimplicial $\wedge^i\Omega_{D(*)}$ ; voir le
lemme \ref{crystalline-lemma-homotopy-tensor}.
Une autre application du même lemme montre que le complété $p$-adique
$\Omega^i_{D(*)}$ est homotopiquement équivalent à zéro.
En tensorisant par $M_*$, on voit que $M_* \otimes_{D(*)} \Omega^i_{D(*)}$
est homotope à zéro ; voir encore le lemme \ref{crystalline-lemma-homotopy-tensor}.
Une dernière application du foncteur de complétion $p$-adique achève la démonstration.
\end{proof}



\section{Lemme de Poincar\'e à puissances divisées}
\label{crystalline-section-poincare}

\noindent
Nous ne considérons que la version la plus simple possible.

\begin{lemma}
\label{crystalline-lemma-poincare}
Soit $A$ un anneau. Soit $P = A\langle x_i \rangle$ un anneau de polynômes
à puissances divisées sur $A$. Pour tout $A$-module $M$, le complexe
$$
0 \to M \to
M \otimes_A P \to
M \otimes_A \Omega^1_{P/A, \delta} \to
M \otimes_A \Omega^2_{P/A, \delta} \to \ldots
$$
est exact. Soit $D$ le complété $p$-adique de $P$.
Soit $\Omega^i_D$ le complété $p$-adique de la $i$-ième puissance
extérieure de $\Omega_{D/A, \delta}$. Pour tout module complet pour la topologie $p$-adique sur
$A$, noté $M$, le complexe
$$
0 \to M \to
M \otimes^\wedge_A D \to
M \otimes^\wedge_A \Omega^1_D \to
M \otimes^\wedge_A \Omega^2_D \to \ldots
$$
est exact.
\end{lemma}

\begin{proof}
Il suffit de montrer que le complexe
$$
E :
(0 \to A \to P \to \Omega^1_{P/A, \delta} \to
\Omega^2_{P/A, \delta} \to \ldots)
$$
est homotopiquement équivalent à zéro comme complexe de $A$-modules.
Pour tout multi-indice $K = (k_i)$, on peut considérer le sous-complexe $E(K)$
qui, en degré $j$, est formé de
$$
\bigoplus\nolimits_{I = \{i_1, \ldots, i_j\} \subset \text{Supp}(K)}
A
\prod\nolimits_{i \not \in I} x_i^{[k_i]}
\prod\nolimits_{i \in I} x_i^{[k_i - 1]}
\text{d}x_{i_1} \wedge \ldots \wedge \text{d}x_{i_j}
$$
Puisque $E = \bigoplus E(K)$, il suffit de prouver que chacun des
complexes $E(K)$ est homotope à zéro. Si $K = 0$, alors
$E(K) : (A \to A)$ est homotope à zéro. Si $K$ a un support (fini) non vide
$S$, alors le complexe $E(K)$ est isomorphe au complexe
$$
0 \to A \to
\bigoplus\nolimits_{s \in S} A \to
\wedge^2(\bigoplus\nolimits_{s \in S} A) \to
\ldots \to \wedge^{\# S}(\bigoplus\nolimits_{s \in S} A) \to 0
$$
qui est homotope à zéro, par exemple par
Compléments d'algèbre, lemme \ref{more-algebra-lemma-homotopy-koszul-abstract}.
\end{proof}

\noindent
Une autre approche, plus directe, du lemme suivant est
expliquée dans l'exemple \ref{crystalline-example-integrate}.

\begin{lemma}
\label{crystalline-lemma-relative-poincare}
Soit $A$ un anneau. Soit $(B, I, \delta)$ un anneau à puissances divisées
tel que $B$ soit une $A$-algèbre.
Soit $P = B\langle x_i \rangle$ un anneau de polynômes à puissances divisées
sur $B$, muni comme d'habitude de l'idéal à puissances divisées $J = IP + B\langle x_i \rangle_{+}$.
Soit $M$ un $B$-module muni d'une connexion intégrable
$\nabla : M \to M \otimes_B \Omega^1_{B/A, \delta}$. Alors le morphisme de
complexes de de Rham
$$
M \otimes_B \Omega^*_{B/A, \delta}
\longrightarrow
M \otimes_P \Omega^*_{P/A, \delta}
$$
est un quasi-isomorphisme. Soient $D$, resp.\ $D'$, les complétés $p$-adiques de
$B$, resp.\ $P$, et soient $\Omega^i_D$, resp.\ $\Omega^i_{D'}$, les complétés $p$-adiques
de $\Omega^i_{B/A, \delta}$,
resp.\ $\Omega^i_{P/A, \delta}$. Soit $M$ un module $p$-adiquement complet
sur $D$, muni d'une connexion intégrable
$\nabla : M \to M \otimes^\wedge_D \Omega^1_D$.
Alors le morphisme de complexes de de Rham
$$
M \otimes^\wedge_D \Omega^*_D
\longrightarrow
M \otimes^\wedge_D \Omega^*_{D'}
$$
est un quasi-isomorphisme.
\end{lemma}

\begin{proof}
Considérons la filtration décroissante $F^*$ de $\Omega^*_{B/A, \delta}$
donnée par les sous-complexes
$F^i(\Omega^*_{B/A, \delta}) = \sigma_{\geq i}\Omega^*_{B/A, \delta}$.
Voir Homologie, section \ref{homology-section-truncations}.
Elle induit une filtration décroissante $F^*$ de $\Omega^*_{P/A, \delta}$
en posant
$$
F^i(\Omega^*_{P/A, \delta}) =
F^i(\Omega^*_{B/A, \delta}) \wedge \Omega^*_{P/A, \delta}.
$$
On a une suite exacte courte scindée
$$
0 \to \Omega^1_{B/A, \delta} \otimes_B P \to
\Omega^1_{P/A, \delta} \to
\Omega^1_{P/B, \delta} \to 0
$$
et le dernier module est libre, ayant pour base les $\text{d}x_i$. Il en résulte que
$F^i(\Omega^*_{P/A, \delta}) \to \Omega^*_{P/A, \delta}$ est, en chaque degré, une
injection scindée et que
$$
\text{gr}^i_F(\Omega^*_{P/A, \delta}) =
\Omega^i_{B/A, \delta} \otimes_B \Omega^*_{P/B, \delta}
$$
comme complexes. On peut donc définir une filtration $F^*$ de
$M \otimes_B \Omega^*_{P/A, \delta}$ en posant
$$
F^i(M \otimes_B \Omega^*_{P/A, \delta}) =
M \otimes_B F^i(\Omega^*_{P/A, \delta})
$$
et l'on a
$$
\text{gr}^i_F(M \otimes_B \Omega^*_{P/A, \delta}) =
M \otimes_B \Omega^i_{B/A, \delta} \otimes_B \Omega^*_{P/B, \delta}
$$
comme complexes.
Par le lemme \ref{crystalline-lemma-poincare}, chacun de ces complexes est
quasi-isomorphe à $M \otimes_B \Omega^i_{B/A, \delta}$ placé en degré $0$.
On voit donc que la première flèche affichée dans l'énoncé est un morphisme de
complexes filtrés qui induit un quasi-isomorphisme sur les gradués associés. Cela
implique qu'il s'agit d'un quasi-isomorphisme, par exemple par la suite spectrale
associée à un complexe filtré ; voir
Homologie, section \ref{homology-section-filtered-complex}.

\medskip\noindent
La démonstration du second quasi-isomorphisme est exactement la même.
\end{proof}



\section{Cohomologie dans le cas affine}
\label{crystalline-section-cohomology-affine}

\noindent
Revenons à la situation étudiée dans la
section \ref{crystalline-section-quasi-coherent-crystals}. On
part de $(A, I, \gamma)$ et de $A/I \to C$, et l'on pose
$X = \Spec(C)$ et $S = \Spec(A)$. On choisit ensuite
un anneau de polynômes $P$ sur $A$ et une surjection $P \to C$ de
noyau $J$. On obtient $D$ et $D(n)$ ; voir
(\ref{crystalline-equation-D}) et (\ref{crystalline-equation-Dn}).
Posons $T(n)_e = \Spec(D(n)/p^eD(n))$, de sorte que
$(X, T(n)_e, \delta(n))$ soit un objet de $\text{Cris}(X/S)$.
Soit $\mathcal{F}$ un faisceau de $\mathcal{O}_{X/S}$-modules, et posons
$$
M(n) = \lim_e \Gamma((X, T(n)_e, \delta(n)), \mathcal{F})
$$
pour $n = 0, 1, 2, 3, \ldots$. La famille ainsi obtenue forme un module cosimplicial
sur l'anneau cosimplicial $D(0), D(1), D(2), \ldots$.

\begin{proposition}
\label{crystalline-proposition-compute-cohomology}
Avec les notations ci-dessus, supposons que
\begin{enumerate}
\item $\mathcal{F}$ soit localement quasi-cohérent, et
\item pour tout morphisme $(U, T, \delta) \to (U', T', \delta')$
de $\text{Cris}(X/S)$ tel que $f : T \to T'$ soit une immersion fermée,
le morphisme $c_f : f^*\mathcal{F}_{T'} \to \mathcal{F}_T$ soit surjectif.
\end{enumerate}
Alors le complexe
$$
M(0) \to M(1) \to M(2) \to \ldots
$$
calcule $R\Gamma(\text{Cris}(X/S), \mathcal{F})$.
\end{proposition}

\begin{proof}
L'hypothèse (1) et le lemme \ref{crystalline-lemma-compare} montrent que
$R\Gamma(\text{Cris}(X/S), \mathcal{F})$ est isomorphe à
$R\Gamma(\mathcal{C}, \mathcal{F})$. Notons que les catégories
$\mathcal{C}$ utilisées dans les lemmes \ref{crystalline-lemma-compare} et \ref{crystalline-lemma-complete}
coïncident. Soit $f : T \to T'$ une immersion fermée comme en (2). La surjectivité
de $c_f : f^*\mathcal{F}_{T'} \to \mathcal{F}_T$ équivaut à celle
de $\mathcal{F}_{T'} \to f_*\mathcal{F}_T$. Ainsi, si
$\mathcal{F}$ satisfait (1) et (2), on obtient une suite exacte courte
$$
0 \to \mathcal{K} \to \mathcal{F}_{T'} \to f_*\mathcal{F}_T \to 0
$$
de $\mathcal{O}_{T'}$-modules quasi-cohérents sur $T'$ ; voir
Schémas, section \ref{schemes-section-quasi-coherent}, et en particulier le
lemme \ref{schemes-lemma-push-forward-quasi-coherent}.
Ainsi, si $T'$ est affine, on conclut que l'application de restriction
$\mathcal{F}(U', T', \delta') \to \mathcal{F}(U, T, \delta)$
est surjective par l'annulation de $H^1(T', \mathcal{K})$ ; voir
Cohomologie des schémas, lemme
\ref{coherent-lemma-quasi-coherent-affine-cohomology-zero}.
Les applications de transition des systèmes inverses du lemme \ref{crystalline-lemma-complete}
sont donc surjectives. On en déduit que
$R^pg_*(\mathcal{F}|_\mathcal{C}) = 0$ pour tout $p \geq 1$,
où $g$ est comme dans le lemme \ref{crystalline-lemma-complete}.
L'objet $D$ de la catégorie $\mathcal{C}^\wedge$
satisfait l'hypothèse du lemme \ref{crystalline-lemma-category-with-covering} par le
lemme \ref{crystalline-lemma-generator-completion},
avec
$$
D \times \ldots \times D = D(n)
$$
dans $\mathcal{C}$, car $D(n)$ est le coproduit de $n + 1$ copies de
$D$ dans $\text{Cris}^\wedge(C/A)$ ; voir le lemme \ref{crystalline-lemma-property-Dn}.
Cela conclut.
\end{proof}

\begin{lemma}
\label{crystalline-lemma-cohomology-is-zero}
On se place sous les hypothèses et avec les notations de la
proposition \ref{crystalline-proposition-compute-cohomology}.
Alors
$$
H^j(\text{Cris}(X/S), \mathcal{F} \otimes_{\mathcal{O}_{X/S}} \Omega^i_{X/S})
= 0
$$
pour tout $i > 0$ et tout $j \geq 0$.
\end{lemma}

\begin{proof}
Il résulte du lemme \ref{crystalline-lemma-omega-locally-quasi-coherent} que
$\mathcal{H} = \mathcal{F} \otimes_{\mathcal{O}_{X/S}} \Omega^i_{X/S}$
satisfait aussi les hypothèses (1) et (2) de la
proposition \ref{crystalline-proposition-compute-cohomology}.
Posons $M(n)_e = \Gamma((X, T(n)_e, \delta(n)), \mathcal{F})$,
de sorte que $M(n) = \lim_e M(n)_e$. Alors
\begin{align*}
\lim_e \Gamma((X, T(n)_e, \delta(n)), \mathcal{H}) & =
\lim_e M(n)_e \otimes_{D(n)_e} \Omega^i_{D(n)}/p^e\Omega^i_{D(n)} \\
& = \lim_e M(n)_e \otimes_{D(n)} \Omega^i_{D(n)}
\end{align*}
Par le
lemme \ref{crystalline-lemma-vanishing},
les modules cosimpliciaux
$$
M(0)_e \otimes_{D(0)} \Omega^i_{D(0)} \to
M(1)_e \otimes_{D(1)} \Omega^i_{D(1)} \to
M(2)_e \otimes_{D(2)} \Omega^i_{D(2)} \to \ldots
$$
sont homotopes à zéro. Comme les applications de transition
$M(n)_{e + 1} \to M(n)_e$ sont surjectives, on voit que
la limite inverse des complexes associés est acyclique\footnote{En fait,
ils sont même homotopes à zéro, car les homotopies sont compatibles, mais cela n'est pas
nécessaire. Cet argument détourné tient au fait que
la limite $\lim_e M(n)_e \otimes_{D(n)} \Omega^i_{D(n)}$ n'est pas le
complété $p$-adique de $M(n) \otimes_{D(n)} \Omega^i_{D(n)}$, car, sous
les hypothèses du lemme, on ne sait pas si
$M(n)_e = M(n)_{e + 1}/p^eM(n)_{e + 1}$. Si $\mathcal{F}$ est un cristal,
cette égalité est vraie.}.
La cohomologie de $\mathcal{H}$ s'annule donc, d'après la
proposition \ref{crystalline-proposition-compute-cohomology}.
\end{proof}

\begin{proposition}
\label{crystalline-proposition-compute-cohomology-crystal}
Sous les hypothèses de la proposition \ref{crystalline-proposition-compute-cohomology},
supposons maintenant que $\mathcal{F}$ soit un cristal en modules quasi-cohérents.
Soit $(M, \nabla)$ le $D$-module à connexion correspondant ; voir la
proposition \ref{crystalline-proposition-crystals-on-affine}. Alors le complexe
$$
M \otimes^\wedge_D \Omega^*_D
$$
calcule $R\Gamma(\text{Cris}(X/S), \mathcal{F})$.
\end{proposition}

\begin{proof}
On va le démontrer à l'aide des deux suites spectrales associées au
complexe double $K^{*, *}$ de termes
$$
K^{a, b} = M \otimes_D^\wedge \Omega^a_{D(b)}
$$
Que sait-on jusqu'ici ? Le lemme \ref{crystalline-lemma-vanishing}
affirme que chaque colonne $K^{a, *}$, pour $a > 0$, est acyclique.
La proposition \ref{crystalline-proposition-compute-cohomology} affirme que
la première colonne $K^{0, *}$ est quasi-isomorphe à
$R\Gamma(\text{Cris}(X/S), \mathcal{F})$.
La première suite spectrale associée au complexe double
montre donc qu'il existe un quasi-isomorphisme canonique de
$R\Gamma(\text{Cris}(X/S), \mathcal{F})$ vers
$\text{Tot}(K^{*, *})$.

\medskip\noindent
Considérons ensuite les lignes $K^{*, b}$. Par le
lemme \ref{crystalline-lemma-structure-Dn},
chacun des $b + 1$ homomorphismes $D \to D(b)$ présente $D(b)$ comme le complété $p$-adique
d'une algèbre polynomiale à puissances divisées sur $D$.
Le lemme \ref{crystalline-lemma-relative-poincare} montre donc que le morphisme
$$
M \otimes^\wedge_D\Omega^*_D
\longrightarrow
M \otimes^\wedge_{D(b)} \Omega^*_{D(b)} = K^{*, b}
$$
est un quasi-isomorphisme. Notons que chacun de ces morphismes induit le {\it même}
morphisme en cohomologie (et même le même morphisme dans la catégorie dérivée), car
l'inverse est donné par le morphisme codiagonal $D(b) \to D$ (correspondant
au morphisme de multiplication $P \otimes_A \ldots \otimes_A P \to P$).
Ainsi, en considérant la page $E_1$ de la seconde suite spectrale,
on obtient
$$
E_1^{a, b} = H^a(M \otimes^\wedge_D\Omega^*_D)
$$
avec les différentielles
$$
E_1^{a, 0} \xrightarrow{0}
E_1^{a, 1} \xrightarrow{1}
E_1^{a, 2} \xrightarrow{0}
E_1^{a, 3} \xrightarrow{1} \ldots
$$
car chacune est la somme alternée des identifications données
$H^a(M \otimes^\wedge_D\Omega^*_D) = E_1^{a, 0} = E_1^{a, 1} = \ldots$.
On voit donc que la page $E_2$ est égale à $H^a(M \otimes^\wedge_D\Omega^*_D)$
sur la première ligne et à zéro ailleurs. Il s'ensuit que l'identification
de $M \otimes^\wedge_D\Omega^*_D$ avec la première ligne induit un
quasi-isomorphisme de $M \otimes^\wedge_D\Omega^*_D$ vers
$\text{Tot}(K^{*, *})$.
\end{proof}

\begin{lemma}
\label{crystalline-lemma-compute-cohomology-crystal-smooth}
On se place sous les hypothèses de la proposition \ref{crystalline-proposition-compute-cohomology-crystal}.
Soient $A \to P' \to C$ des homomorphismes d'anneaux tels que $A \to P'$ soit lisse et $P' \to C$
surjectif de noyau $J'$. Soit $D'$ le complété $p$-adique de
$D_{P', \gamma}(J')$. Soit $(M', \nabla')$ le couple sur $D'$
correspondant à $\mathcal{F}$ ; voir le
lemme \ref{crystalline-lemma-crystals-on-affine-smooth}. Alors le complexe
$$
M' \otimes^\wedge_{D'} \Omega^*_{D'}
$$
calcule $R\Gamma(\text{Cris}(X/S), \mathcal{F})$.
\end{lemma}

\begin{proof}
Choisissons $a : D \to D'$ et $b : D' \to D$ comme dans le
lemme \ref{crystalline-lemma-crystals-on-affine-smooth}.
Notons que le module obtenu par le changement de base $M = M' \otimes^\wedge_{D', b} D$, muni de sa
connexion $\nabla$, correspond à $\mathcal{F}$. Nous savons donc
que $M \otimes^\wedge_D \Omega_D^*$ calcule la cohomologie cristalline
de $\mathcal{F}$, voir la
proposition \ref{crystalline-proposition-compute-cohomology-crystal}.
Il suffit donc de montrer que les morphismes de changement de base (induits
par $a$ et $b$)
$$
M' \otimes^\wedge_{D'} \Omega^*_{D'}
\longrightarrow
M \otimes^\wedge_D \Omega^*_D
\quad\text{et}\quad
M \otimes^\wedge_D \Omega^*_D
\longrightarrow
M' \otimes^\wedge_{D'} \Omega^*_{D'}
$$
sont des quasi-isomorphismes. Puisque $a \circ b = \text{id}_{D'}$, on voit
que la composée dans un sens est l'identité du complexe
$M' \otimes^\wedge_{D'} \Omega^*_{D'}$. Il suffit donc de montrer que
le morphisme
$$
M \otimes^\wedge_D \Omega^*_D
\longrightarrow
M \otimes^\wedge_D \Omega^*_D
$$
induit par $b \circ a : D \to D$ est un quasi-isomorphisme. (Notons que nous
avons le même complexe des deux côtés puisque $M = M' \otimes^\wedge_{D', b} D$,
d'où $M \otimes^\wedge_{D, b \circ a} D =
M' \otimes^\wedge_{D', b \circ a \circ b} D =
M' \otimes^\wedge_{D', b} D = M$.) En fait, nous affirmons que, pour tout
homomorphisme de $A$-algèbres à puissances divisées $\rho : D \to D$ compatible
avec l'augmentation vers $C$, le morphisme induit
$M \otimes^\wedge_D \Omega^*_D \to M \otimes^\wedge_{D, \rho} \Omega^*_D$
est un quasi-isomorphisme.

\medskip\noindent
Écrivons $\rho(x_i) = x_i + z_i$. Les éléments $z_i$ appartiennent à
l'idéal à puissances divisées de $D$ car $\rho$ est compatible avec
l'augmentation vers $C$. On peut donc factoriser le morphisme $\rho$
en une composée
$$
D \xrightarrow{\sigma} D\langle \xi_i \rangle^\wedge \xrightarrow{\tau} D
$$
où le premier morphisme est donné par $x_i \mapsto x_i + \xi_i$ et le
second est le morphisme de $D$-algèbres à puissances divisées qui envoie $\xi_i$ sur $z_i$.
(Cela utilise les propriétés universelles des algèbres polynomiales, des algèbres
polynomiales à puissances divisées, des enveloppes à puissances divisées et de la complétion $p$-adique.)
Notons qu'il existe un {\it automorphisme} $\alpha$ de
$D\langle \xi_i \rangle^\wedge$ avec $\alpha(x_i) = x_i - \xi_i$
et $\alpha(\xi_i) = \xi_i$. En appliquant le lemme \ref{crystalline-lemma-relative-poincare}
à $\alpha \circ \sigma$ (qui envoie $x_i$ sur $x_i$) et en utilisant le fait que
$\alpha$ est un isomorphisme, on conclut que $\sigma$ induit un
quasi-isomorphisme de $M \otimes^\wedge_D \Omega^*_D$ vers
$M \otimes^\wedge_{D, \sigma} \Omega^*_{D\langle \xi_i \rangle^\wedge}$.
D'autre part, le morphisme $\tau$ admet pour inverse à gauche le
morphisme $D \to D\langle \xi_i \rangle^\wedge$, $x_i \mapsto x_i$,
et l'on conclut (en utilisant encore une fois le lemme \ref{crystalline-lemma-relative-poincare})
que $\tau$ induit un quasi-isomorphisme de
$M \otimes^\wedge_{D, \sigma} \Omega^*_{D\langle \xi_i \rangle^\wedge}$
vers $M \otimes^\wedge_{D, \tau \circ \sigma} \Omega^*_D$. En composant ces
deux quasi-isomorphismes, on obtient que $\rho$ induit un quasi-isomorphisme
$M \otimes^\wedge_D \Omega^*_D \to M \otimes^\wedge_{D, \rho} \Omega^*_D$
ce qu'il fallait démontrer.
\end{proof}










\section{Deux contre-exemples}
\label{crystalline-section-examples}

\noindent
Avant de passer à quelques résultats positifs de la cohomologie cristalline,
donnons deux exemples qui expliquent pourquoi la cohomologie cristalline
fonctionne assez mal lorsque les schémas considérés ne sont pas
propres sur la base ou sont singuliers. Le premier exemple se trouve
dans \cite{BO}.

\begin{example}
\label{crystalline-example-torsion}
Soit $A = \mathbf{Z}_p$, avec l'idéal à puissances divisées $(p)$ muni de
son unique structure de puissances divisées $\gamma$. Soit
$C = \mathbf{F}_p[x, y]/(x^2, xy, y^2)$. Choisissons la présentation
$$
C = P/J = \mathbf{Z}_p[x, y]/(x^2, xy, y^2, p)
$$
Soit $D = D_{P, \gamma}(J)^\wedge$, avec pour idéal à puissances divisées
$(\bar J, \bar \gamma)$ comme à la section \ref{crystalline-section-quasi-coherent-crystals}.
Nous noterons encore $x, y$ les images de $x$ et $y$ dans $D$.
Considérons l'élément
$$
\tau = \bar\gamma_p(x^2)\bar\gamma_p(y^2) - \bar\gamma_p(xy)^2 \in D
$$
On remarque que $p\tau = 0$ car
$$
p! \bar\gamma_p(x^2) \bar\gamma_p(y^2) =
x^{2p} \bar\gamma_p(y^2) = \bar\gamma_p(x^2y^2) =
x^py^p \bar\gamma_p(xy) = p! \bar\gamma_p(xy)^2
$$
dans $D$. On remarque aussi que $\text{d}\tau = 0$ dans $\Omega_D$ car
\begin{align*}
\text{d}(\bar\gamma_p(x^2) \bar\gamma_p(y^2))
& =
\bar\gamma_{p - 1}(x^2)\bar\gamma_p(y^2)\text{d}x^2 +
\bar\gamma_p(x^2)\bar\gamma_{p - 1}(y^2)\text{d}y^2 \\
& =
2 x \bar\gamma_{p - 1}(x^2)\bar\gamma_p(y^2)\text{d}x +
2 y \bar\gamma_p(x^2)\bar\gamma_{p - 1}(y^2)\text{d}y \\
& =
2/(p - 1)!( x^{2p - 1} \bar\gamma_p(y^2)\text{d}x +
y^{2p - 1} \bar\gamma_p(x^2)\text{d}y ) \\
& =
2/(p - 1)!
(x^{p - 1} \bar\gamma_p(xy^2)\text{d}x +
y^{p - 1} \bar\gamma_p(x^2y)\text{d}y) \\
& =
2/(p - 1)!
(x^{p - 1}y^p \bar\gamma_p(xy)\text{d}x +
x^py^{p - 1} \bar\gamma_p(xy)\text{d}y) \\
& =
2 \bar\gamma_{p - 1}(xy) \bar\gamma_p(xy)(y\text{d}x + x \text{d}y) \\
& = 
\text{d}(\bar\gamma_p(xy)^2)
\end{align*}
Enfin, nous affirmons que $\tau \not = 0$ dans $D$. Pour le voir, il suffit de
construire un objet $(B \to \mathbf{F}_p[x, y]/(x^2, xy, y^2), \delta)$
de $\text{Cris}(C/S)$ tel que l'image de $\tau$ dans $B$ ne soit pas nulle.
Pour cela, prenons
$$
B = \mathbf{F}_p[x, y, u, v]/(x^3, x^2y, xy^2, y^3, xu, yu, xv, yv, u^2, v^2)
$$
muni de la surjection évidente vers $C$. Soit $K = \Ker(B \to C)$ et
considérons l'application
$$
\delta_p : K \longrightarrow K,\quad
ax^2 + bxy + cy^2 + du + ev + fuv \longmapsto a^pu + c^pv
$$
On vérifie qu'elle satisfait aux hypothèses (1), (2), (3) du
lemme \ref{dpa-lemma-need-only-gamma-p} d'Algèbre à puissances divisées
et définit donc une structure de puissances divisées. De plus,
l'image de $\tau$ est $uv$, qui est non nul dans $B$.
Posons $X = \Spec(C)$ et $S = \Spec(A)$.
Nous en tirons les conclusions suivantes :
\begin{enumerate}
\item $H^0(\text{Cris}(X/S), \mathcal{O}_{X/S})$ a de la $p$-torsion, et
\item l'endomorphisme d'image réciproque par Frobenius $F^* : H^0(\text{Cris}(X/S), \mathcal{O}_{X/S})
\to H^0(\text{Cris}(X/S), \mathcal{O}_{X/S})$ n'est pas injective.
\end{enumerate}
En effet, $\tau$ définit un élément de torsion non nul de
$H^0(\text{Cris}(X/S), \mathcal{O}_{X/S})$ d'après la
proposition \ref{crystalline-proposition-compute-cohomology-crystal}.
De même, $F^*(\tau) = \sigma(\tau)$, où $\sigma : D \to D$ est le
morphisme induit par tout relèvement de Frobenius sur $P$. Si l'on choisit $\sigma(x) = x^p$
et $\sigma(y) = y^p$, un calcul simple montre que $F^*(\tau) = 0$.
\end{example}

\noindent
L'exemple suivant montre que, même pour l'espace affine de dimension $r$, la cohomologie
cristalline ne donne pas le résultat attendu.

\begin{example}
\label{crystalline-example-affine-n-space}
Soit $A = \mathbf{Z}_p$, avec l'idéal à puissances divisées $(p)$ muni de
son unique structure de puissances divisées $\gamma$. Soit
$C = \mathbf{F}_p[x_1, \ldots, x_r]$. Choisissons la présentation
$$
C = P/J = P/pP\quad\text{avec}\quad P = \mathbf{Z}_p[x_1, \ldots, x_r]
$$
Notons que $pP$ est muni de puissances divisées d'après
Algèbre à puissances divisées, lemme \ref{dpa-lemma-gamma-extends}.
Ainsi, en posant $D = P^\wedge$ avec l'idéal à puissances divisées $(p)$, on obtient une
situation comme à la section \ref{crystalline-section-quasi-coherent-crystals}.
On conclut que $R\Gamma(\text{Cris}(X/S), \mathcal{O}_{X/S})$
est représenté par le complexe
$$
D \to \Omega^1_D \to \Omega^2_D \to \ldots \to \Omega^r_D
$$
voir la proposition \ref{crystalline-proposition-compute-cohomology-crystal}.
En supposant $r > 0$, on obtient les conclusions suivantes :
\begin{enumerate}
\item La cohomologie cristalline du faisceau structural cristallin
de $X = \mathbf{A}^r_{\mathbf{F}_p}$ sur $S = \Spec(\mathbf{Z}_p)$
est nulle sauf en degrés $0, \ldots, r$.
\item On a $H^0(\text{Cris}(X/S), \mathcal{O}_{X/S}) = \mathbf{Z}_p$.
\item Le groupe de cohomologie $H^r(\text{Cris}(X/S), \mathcal{O}_{X/S})$
est infini et n'est pas un groupe abélien de torsion.
\item Le groupe de cohomologie $H^r(\text{Cris}(X/S), \mathcal{O}_{X/S})$
n'est pas séparé pour la topologie $p$-adique.
\end{enumerate}
Si les deux premiers énoncés sont raisonnables, les assertions (3) et (4) sont
déconcertants ! Ces énoncés se vérifient immédiatement en
explicitant le complexe affiché ci-dessus. Faisons-le simplement
dans le cas $r = 1$. Nous considérons alors le complexe à deux termes
de modules $p$-adiquement complets
$$
\text{d} :
D = \left(
\bigoplus\nolimits_{n \geq 0} \mathbf{Z}_p x^n
\right)^\wedge
\longrightarrow
\Omega^1_D = \left(
\bigoplus\nolimits_{n \geq 1} \mathbf{Z}_p x^{n - 1}\text{d}x
\right)^\wedge
$$
Le morphisme est donné par $\text{diag}(0, 1, 2, 3, 4, \ldots)$, à ceci près que
le premier facteur de la somme directe manque au membre de droite. Il est alors clair
que $\bigoplus_{n > 0} \mathbf{Z}_p/n\mathbf{Z}_p$ est un sous-groupe
du conoyau, lequel est donc infini. En fait, l'élément
$$
\omega = \sum\nolimits_{e > 0} p^e x^{p^{2e} - 1}\text{d}x
$$
n'est manifestement pas un élément de torsion du conoyau. Mais la situation est pire.
En effet, considérons l'élément
$$
\eta = \sum\nolimits_{e > 0} p^e x^{p^e - 1}\text{d}x
$$
Pour tout $t > 0$, l'élément $\eta$ est congru à
$\sum_{e > t} p^e x^{p^e - 1}\text{d}x$ modulo l'image de
$\text{d}$ ; cette somme est divisible par $p^t$. Ainsi, la classe de cohomologie
de $\eta$ dans $H^1(\text{Cris}(X/S), \mathcal{O}_{X/S})$
appartient à $\bigcap p^tH^1(\text{Cris}(X/S), \mathcal{O}_{X/S})$.
Mais $\eta$ n'appartient pas à l'image de $\text{d}$ car il devrait
être l'image de $a + \sum_{e > 0} x^{p^e}$ pour un certain $a \in \mathbf{Z}_p$,
mais cette somme n'appartient pas au module figurant à gauche.
Ainsi, la classe de cohomologie de $\eta$ est non nulle et l'on voit que
l'assertion (4) est vérifiée. (En fait, la classe de cohomologie de $\eta$ n'est pas de torsion.)
En revanche, les groupes de cohomologie
$H^i(\text{Cris}(X/S), \mathcal{O}_{X/S})$
sont des $\mathbf{Z}_p$-modules complets au sens dérivé : ce sont les groupes de cohomologie
d'un complexe de modules $p$-adiquement complets, voir
Compléments d'algèbre, section \ref{more-algebra-section-derived-completion}.
\end{example}








\section{Applications}
\label{crystalline-section-applications}

\noindent
Dans cette section, nous rassemblons quelques applications des résultats des
sections précédentes.

\begin{proposition}
\label{crystalline-proposition-compare-with-de-Rham}
Dans la situation \ref{crystalline-situation-global}.
Soit $\mathcal{F}$ un cristal en modules quasi-cohérents sur
$\text{Cris}(X/S)$. La flèche de troncation de complexes
$$
(\mathcal{F} \to
\mathcal{F} \otimes_{\mathcal{O}_{X/S}} \Omega^1_{X/S} \to
\mathcal{F} \otimes_{\mathcal{O}_{X/S}} \Omega^2_{X/S} \to \ldots)
\longrightarrow \mathcal{F}[0],
$$
bien qu'il ne soit pas un quasi-isomorphisme, le devient après application du foncteur
$Ru_{X/S, *}$. En fait, pour tout $i > 0$, on a 
$$
Ru_{X/S, *}(\mathcal{F} \otimes_{\mathcal{O}_{X/S}} \Omega^i_{X/S}) = 0.
$$
\end{proposition}

\begin{proof}
Le lemme \ref{crystalline-lemma-automatic-connection} fournit un complexe de de Rham
comme indiqué dans l'énoncé. Abrégeons
$\mathcal{H} = \mathcal{F} \otimes \Omega^i_{X/S}$.
Soit $X' \subset X$ un sous-schéma ouvert affine
dont l'image est contenue dans un sous-schéma ouvert affine $S' \subset S$.
Alors
$$
(Ru_{X/S, *}\mathcal{H})|_{X'_{Zar}} =
Ru_{X'/S', *}(\mathcal{H}|_{\text{Cris}(X'/S')}),
$$
voir le lemme \ref{crystalline-lemma-localize}. Ainsi, le
lemme \ref{crystalline-lemma-cohomology-is-zero}
montre que $Ru_{X/S, *}\mathcal{H}$ est un complexe de faisceaux sur
$X_{Zar}$ dont la cohomologie est nulle sur tout ouvert affine.
Comme la topologie de $X$ possède une base formée d'ouverts affines,
cela implique que $Ru_{X/S, *}\mathcal{H}$ est quasi-isomorphe à zéro.
\end{proof}

\begin{remark}
\label{crystalline-remark-vanishing}
La démonstration de la proposition \ref{crystalline-proposition-compare-with-de-Rham}
montre que la conclusion
$$
Ru_{X/S, *}(\mathcal{F} \otimes_{\mathcal{O}_{X/S}} \Omega^i_{X/S}) = 0
$$
pour $i > 0$ vaut pour tout $\mathcal{O}_{X/S}$-module
$\mathcal{F}$ qui vérifie les conditions (1) et (2) de la
proposition \ref{crystalline-proposition-compute-cohomology}.
Cela s'applique aux faisceaux suivants, qui ne sont pas des cristaux :
$\Omega^i_{X/S}$ pour tout $i$, ainsi qu'à tout faisceau de la forme
$\underline{\mathcal{F}}$, où $\mathcal{F}$ est un $\mathcal{O}_X$-module
quasi-cohérent. En particulier, cela s'applique au
faisceau $\underline{\mathcal{O}_X} = \underline{\mathbf{G}_a}$.
Notons toutefois qu'un résultat tel que le lemme \ref{crystalline-lemma-automatic-connection}
est nécessaire pour produire un complexe de de Rham, ce qui exige que $\mathcal{F}$ soit un cristal.
Par conséquent, à l'heure actuelle, la classe des faisceaux de modules pour lesquels l'énoncé
complet de la proposition \ref{crystalline-proposition-compare-with-de-Rham} est vérifié
est exactement la catégorie des cristaux en modules quasi-cohérents.
\end{remark}

\noindent
Plaçons-nous dans la situation \ref{crystalline-situation-global}.
Soit $\mathcal{F}$ un cristal en modules quasi-cohérents sur
$\text{Cris}(X/S)$. Soit $(U, T, \delta)$ un objet de
$\text{Cris}(X/S)$. La proposition \ref{crystalline-proposition-compare-with-de-Rham}
permet de construire un morphisme canonique
\begin{equation}
\label{crystalline-equation-restriction}
R\Gamma(\text{Cris}(X/S), \mathcal{F})
\longrightarrow
R\Gamma(T, \mathcal{F}_T \otimes_{\mathcal{O}_T} \Omega^*_{T/S, \delta})
\end{equation}
En effet, on a $R\Gamma(\text{Cris}(X/S), \mathcal{F}) =
R\Gamma(\text{Cris}(X/S), \mathcal{F} \otimes \Omega^*_{X/S})$,
on peut restreindre les classes de cohomologie globale à $T$, et $\Omega_{X/S}$
se restreint à $\Omega_{T/S, \delta}$ par le
lemme \ref{crystalline-lemma-module-of-differentials}.

















\section{Quelques résultats supplémentaires}
\label{crystalline-section-missing}

\noindent
Dans cette section, nous mentionnons certains résultats dont nous ne donnons pas encore la démonstration.
Nous les formulons sous forme d'une suite de remarques et ne les transformerons
en véritables lemmes et propositions que lorsque nous ajouterons des
démonstrations détaillées.

\begin{remark}[Images directes supérieures]
\label{crystalline-remark-compute-direct-image}
Soit $p$ un nombre premier. Soit
$(S, \mathcal{I}, \gamma) \to (S', \mathcal{I}', \gamma')$
un morphisme de schémas à puissances divisées sur $\mathbf{Z}_{(p)}$. Soit
$$
\xymatrix{
X \ar[r]_f \ar[d] & X' \ar[d] \\
S_0 \ar[r] & S'_0
}
$$
un diagramme commutatif de morphismes de schémas, et supposons que $p$ soit
localement nilpotent sur $X$ et $X'$. Soit $\mathcal{F}$ un
$\mathcal{O}_{X/S}$-module sur $\text{Cris}(X/S)$. Alors
$Rf_{\text{cris}, *}\mathcal{F}$ se calcule comme suit.

\medskip\noindent
Étant donné un objet $(U', T', \delta')$ de $\text{Cris}(X'/S')$, posons
$U = X \times_{X'} U' = f^{-1}(U')$ (un sous-schéma ouvert de $X$). Notons
$(T_0, T, \delta)$ le schéma à puissances divisées sur $S$ tel que
$$
\xymatrix{
T \ar[r] \ar[d] & T' \ar[d] \\
S \ar[r] & S'
}
$$
soit cartésien dans la catégorie des schémas à puissances divisées ; voir le
lemme \ref{crystalline-lemma-fibre-product}. Il existe un
morphisme induit $U \to T_0$, et l'on obtient un morphisme
$(U/T)_{\text{cris}} \to (X/S)_{\text{cris}}$ ; voir la
remarque \ref{crystalline-remark-functoriality-cris}.
Soit $\mathcal{F}_U$ l'image inverse de $\mathcal{F}$.
Soit $\tau_{U/T} : (U/T)_{\text{cris}} \to T_{Zar}$ le morphisme structural.
Alors on a
\begin{equation}
\label{crystalline-equation-identify-pushforward}
\left(Rf_{\text{cris}, *}\mathcal{F}\right)_{T'} =
R(T \to T')_*\left(R\tau_{U/T, *} \mathcal{F}_U \right)
\end{equation}
où le membre de gauche est la restriction (voir la
section \ref{crystalline-section-sheaves}).

\medskip\noindent
Indications : montrer d'abord que $\text{Cris}(U/T)$ est le topos obtenu en localisant $\text{Cris}(X/S)$ (au sens
du lemme \ref{sites-lemma-localize-topos-site} de Sites)
par rapport au faisceau d'ensembles $f_{\text{cris}}^{-1}h_{(U', T', \delta')}$. Réduire ensuite
l'énoncé au cas où $\mathcal{F}$ est un module injectif
et où l'on considère l'image directe de modules, en utilisant le fait que l'image inverse d'un
$\mathcal{O}_{X/S}$-module injectif est un $\mathcal{O}_{U/T}$-module injectif sur
$\text{Cris}(U/T)$. Enfin, vérifier le résultat pour l'image directe ordinaire.
\end{remark}

\begin{remark}[Mayer-Vietoris]
\label{crystalline-remark-mayer-vietoris}
Dans la situation de la remarque \ref{crystalline-remark-compute-direct-image},
supposons donné un recouvrement ouvert $X = X' \cup X''$. Notons
$X''' = X' \cap X''$. Soient $f'$, $f''$ et $f'''$ les restrictions de $f$
à $X'$, $X''$ et $X'''$. De plus, soient $\mathcal{F}'$, $\mathcal{F}''$
et $\mathcal{F}'''$ les restrictions de $\mathcal{F}$ aux sites cristallins
de $X'$, $X''$ et $X'''$. Il existe alors un triangle distingué
$$
Rf_{\text{cris}, *}\mathcal{F}
\longrightarrow
Rf'_{\text{cris}, *}\mathcal{F}' \oplus Rf''_{\text{cris}, *}\mathcal{F}''
\longrightarrow
Rf'''_{\text{cris}, *}\mathcal{F}'''
\longrightarrow
Rf_{\text{cris}, *}\mathcal{F}[1]
$$
dans $D(\mathcal{O}_{X'/S'})$.

\medskip\noindent
Indications : c'est une conséquence formelle du fait que les sous-catégories
$\text{Cris}(X'/S)$, $\text{Cris}(X''/S)$, $\text{Cris}(X'''/S)$ correspondent
à des sous-objets ouverts du faisceau final sur $\text{Cris}(X/S)$ et que la
dernière est l'intersection des deux premières.
\end{remark}

\begin{remark}[Complexe de {\v C}ech]
\label{crystalline-remark-cech-complex}
Soit $p$ un nombre premier. Soit $(A, I, \gamma)$ un anneau à puissances
divisées, $A$ étant une $\mathbf{Z}_{(p)}$-algèbre. Posons $S = \Spec(A)$ et
$S_0 = \Spec(A/I)$. Soit $X$ un schéma séparé\footnote{Cette hypothèse n'est
pas strictement nécessaire : à l'aide d'hyperrecouvrements, la construction de la
remarque s'étend au cas général.} sur
$S_0$, tel que $p$ soit localement nilpotent sur $X$. Soit $\mathcal{F}$ un
cristal en $\mathcal{O}_{X/S}$-modules quasi-cohérents.

\medskip\noindent
Choisissons un recouvrement ouvert affine
$X = \bigcup_{\lambda \in \Lambda} U_\lambda$ de $X$.
Écrivons $U_\lambda = \Spec(C_\lambda)$. Choisissons une algèbre de polynômes
$P_\lambda$ sur $A$ et une surjection $P_\lambda \to C_\lambda$.
Ces choix étant fixés, on peut construire un complexe de {\v C}ech qui
calcule $R\Gamma(\text{Cris}(X/S), \mathcal{F})$.

\medskip\noindent
Pour $n \geq 0$ et $\lambda_0, \ldots, \lambda_n \in \Lambda$,
posons $U_{\lambda_0 \ldots \lambda_n} = U_{\lambda_0} \cap \ldots
\cap U_{\lambda_n}$. Ce schéma est affine par hypothèse. Écrivons
$U_{\lambda_0 \ldots \lambda_n} = \Spec(C_{\lambda_0 \ldots \lambda_n})$.
Posons
$$
P_{\lambda_0 \ldots \lambda_n} =
P_{\lambda_0} \otimes_A \ldots \otimes_A P_{\lambda_n}
$$
qui est muni d'une surjection canonique vers $C_{\lambda_0 \ldots \lambda_n}$.
Notons son noyau $J_{\lambda_0 \ldots \lambda_n}$ et notons
$D_{\lambda_0 \ldots \lambda_n}$
l'enveloppe à puissances divisées, complétée $p$-adiquement, de
$J_{\lambda_0 \ldots \lambda_n}$ dans $P_{\lambda_0 \ldots \lambda_n}$
relativement à $\gamma$. Soit $M_{\lambda_0 \ldots \lambda_n}$ le
$P_{\lambda_0 \ldots \lambda_n}$-module correspondant
à la restriction de $\mathcal{F}$ à
$\text{Cris}(U_{\lambda_0 \ldots \lambda_n}/S)$ d'après
la proposition \ref{crystalline-proposition-crystals-on-affine}.
Par construction, on obtient un anneau cosimplicial à puissances divisées $D(*)$
dont le terme de degré $n$ est l'anneau
$$
D(n) =
\prod\nolimits_{\lambda_0 \ldots \lambda_n}
D_{\lambda_0 \ldots \lambda_n}
$$
(on utilise la fonctorialité des enveloppes à puissances divisées et la structure
cosimpliciale triviale sur l'anneau $P(*)$ défini de même).
Comme $M_{\lambda_0 \ldots \lambda_n}$ est la « valeur » de $\mathcal{F}$
sur les objets $\Spec(D_{\lambda_0 \ldots \lambda_n})$, on voit que
$M(*)$ défini par la règle
$$
M(n) = \prod\nolimits_{\lambda_0 \ldots \lambda_n}
M_{\lambda_0 \ldots \lambda_n}
$$
définit un module cosimplicial sur $D(*)$. Nous affirmons maintenant que
$$
R\Gamma(\text{Cris}(X/S), \mathcal{F}) = s(M(*))
$$
Ici, $s(-)$ désigne le complexe de cochaînes associé à un module
cosimplicial (voir
Objets simpliciaux, section \ref{simplicial-section-dold-kan-cosimplicial}).

\medskip\noindent
Indications : la démonstration est semblable à celle de la
proposition \ref{crystalline-proposition-compute-cohomology} (en particulier,
le résultat vaut pour tout module qui vérifie les hypothèses de
cette proposition).
\end{remark}

\begin{remark}[Complexe de {\v C}ech alterné]
\label{crystalline-remark-alternating-cech-complex}
Soit $p$ un nombre premier. Soit $(A, I, \gamma)$ un anneau à puissances
divisées, $A$ étant une $\mathbf{Z}_{(p)}$-algèbre. Posons $S = \Spec(A)$ et
$S_0 = \Spec(A/I)$. Soit $X$ un schéma séparé quasi-compact
sur $S_0$, tel que $p$ soit localement nilpotent sur $X$. Soit
$\mathcal{F}$ un cristal en $\mathcal{O}_{X/S}$-modules quasi-cohérents.

\medskip\noindent
Choisissons un recouvrement ouvert affine fini
$X = \bigcup_{\lambda \in \Lambda} U_\lambda$ de $X$
et un ordre total sur $\Lambda$.
Écrivons $U_\lambda = \Spec(C_\lambda)$. Choisissons une algèbre de polynômes
$P_\lambda$ sur $A$ et une surjection $P_\lambda \to C_\lambda$.
Ces choix étant fixés, on peut construire un complexe
de {\v C}ech alterné qui calcule $R\Gamma(\text{Cris}(X/S), \mathcal{F})$.

\medskip\noindent
Nous allons utiliser les notations introduites dans la
remarque \ref{crystalline-remark-cech-complex}.
Notons $\Omega_{\lambda_0 \ldots \lambda_n}$
le module des différentielles, complété $p$-adiquement, de
$D_{\lambda_0 \ldots \lambda_n}$ sur $A$ compatible avec la structure à puissances
divisées. Soit $\nabla$ la connexion intégrable sur
$M_{\lambda_0 \ldots \lambda_n}$ provenant de la
proposition \ref{crystalline-proposition-crystals-on-affine}.
Considérons le complexe double $M^{\bullet, \bullet}$ dont les
termes sont
$$
M^{n, m} =
\bigoplus\nolimits_{\lambda_0 < \ldots < \lambda_n}
M_{\lambda_0 \ldots \lambda_n}
\otimes^\wedge_{D_{\lambda_0 \ldots \lambda_n}}
\Omega^m_{D_{\lambda_0 \ldots \lambda_n}}.
$$
Pour la différentielle $d_1$ (qui augmente $n$), utilisons la différentielle
de {\v C}ech usuelle et, pour la différentielle $d_2$,
la connexion, c'est-à-dire la différentielle du complexe de de Rham.
Nous affirmons que
$$
R\Gamma(\text{Cris}(X/S), \mathcal{F}) = \text{Tot}(M^{\bullet, \bullet})
$$
Ici, $\text{Tot}(-)$ désigne le complexe total associé à un
complexe double ; voir
Homologie, définition \ref{homology-definition-associated-simple-complex}.

\medskip\noindent
Indications : on a
$$
R\Gamma(\text{Cris}(X/S), \mathcal{F}) = R\Gamma(\text{Cris}(X/S),
\mathcal{F} \otimes_{\mathcal{O}_{X/S}} \Omega_{X/S}^\bullet)
$$
par la proposition \ref{crystalline-proposition-compare-with-de-Rham}.
Le membre de droite de la formule n'est autre que le complexe de {\v C}ech alterné
du recouvrement $X = \bigcup_{\lambda \in \Lambda} U_\lambda$
(qui induit un recouvrement ouvert du faisceau final de $\text{Cris}(X/S)$)
et du complexe $\mathcal{F} \otimes_{\mathcal{O}_{X/S}} \Omega_{X/S}^\bullet$ ;
voir la proposition \ref{crystalline-proposition-compute-cohomology-crystal}.
Le résultat découle alors d'un résultat général sur la cohomologie des sites,
selon lequel le complexe de {\v C}ech alterné calcule la cohomologie
dès lors qu'il la calcule correctement sur tous les morceaux (insérer ici une
référence ultérieure).
\end{remark}

\begin{remark}[Quasi-cohérence]
\label{crystalline-remark-quasi-coherent}
Dans la situation de la remarque \ref{crystalline-remark-compute-direct-image},
supposons que $S \to S'$ soit quasi-compact et quasi-séparé et
que $X \to S_0$ soit quasi-compact et quasi-séparé. Alors, pour un cristal
en $\mathcal{O}_{X/S}$-modules quasi-cohérents $\mathcal{F}$,
les faisceaux $R^if_{\text{cris}, *}\mathcal{F}$ sont localement quasi-cohérents.

\medskip\noindent
Indications : il faut montrer que les restrictions à $T'$ sont des
$\mathcal{O}_{T'}$-modules quasi-cohérents, où $(U', T', \delta')$ est un objet quelconque de
$\text{Cris}(X'/S')$. Il suffit de le faire lorsque $T'$ est affine.
On utilise la formule (\ref{crystalline-equation-identify-pushforward}),
le fait que $T \to T'$ est quasi-compact et quasi-séparé (car $T$
est affine au-dessus du changement de base de $T'$ le long de $S \to S'$), ainsi que le
lemme de Cohomologie des schémas
\ref{coherent-lemma-quasi-coherence-higher-direct-images}
pour voir qu'il suffit de montrer que les faisceaux
$R^i\tau_{U/T, *}\mathcal{F}_U$ sont quasi-cohérents.
Notons que $U \to T_0$ est aussi quasi-compact et quasi-séparé ; voir
Schémas, lemmes \ref{schemes-lemma-quasi-compact-permanence} et
\ref{schemes-lemma-compose-after-separated}.

\medskip\noindent
On est ainsi ramené à prouver que $R^i\tau_{X/S, *}\mathcal{F}$
est quasi-cohérent sur $S$ dans le cas où $p$ est localement nilpotent sur $S$. Ici,
$\tau_{X/S}$ est le morphisme structural ; voir la
remarque \ref{crystalline-remark-structure-morphism}.
On peut travailler localement sur $S$, et donc supposer $S$ affine
(voir le lemme \ref{crystalline-lemma-localize}). Une récurrence sur le nombre
d'ouverts affines recouvrant $X$ et Mayer--Vietoris
(remarque \ref{crystalline-remark-mayer-vietoris}) ramènent la question
au cas où $X$ est lui aussi affine (comme dans la démonstration du
lemme de Cohomologie des schémas
\ref{coherent-lemma-quasi-coherence-higher-direct-images}).
Écrivons $X = \Spec(C)$ et $S = \Spec(A)$, de sorte que $(A, I, \gamma)$ et
$A \to C$ soient comme
dans la situation \ref{crystalline-situation-affine}. Choisissons une algèbre de polynômes
$P$ sur $A$ et une surjection $P \to C$ comme dans la
section \ref{crystalline-section-quasi-coherent-crystals}.
Soit $(M, \nabla)$ le module correspondant à $\mathcal{F}$ ; voir la
proposition \ref{crystalline-proposition-crystals-on-affine}.
En appliquant la 
proposition \ref{crystalline-proposition-compute-cohomology-crystal},
on voit que $R\Gamma(\text{Cris}(X/S), \mathcal{F})$ est représenté par
$M \otimes_D \Omega_D^*$. Notons que la complétion n'est pas nécessaire,
car $p$ est nilpotent dans $A$ ! Il faut montrer que cela est compatible
avec le passage aux ouverts principaux de $S = \Spec(A)$. Supposons que $g \in A$.
On conclut de même que $R\Gamma(\text{Cris}(X_g/S_g), \mathcal{F})$
est calculé par $M_g \otimes_{D_g} \Omega_{D_g}^*$ (on utilise encore le fait que
la complétion $p$-adique n'est pas nécessaire). La conclusion s'ensuit puisque la localisation
est un foncteur exact dans la catégorie des $A$-modules.
\end{remark}

\begin{remark}[Bornitude]
\label{crystalline-remark-bounded-cohomology}
Dans la situation de la remarque \ref{crystalline-remark-compute-direct-image},
supposons que $S \to S'$ soit quasi-compact et quasi-séparé et
que $X \to S_0$ soit de type fini et quasi-séparé. Alors il existe
un entier $i_0$ tel que, pour tout cristal
en $\mathcal{O}_{X/S}$-modules quasi-cohérents $\mathcal{F}$,
on ait $R^if_{\text{cris}, *}\mathcal{F} = 0$ pour tout $i > i_0$.

\medskip\noindent
Indications : en raisonnant comme dans la remarque \ref{crystalline-remark-quasi-coherent} (à l'aide du
lemme de Cohomologie des schémas
\ref{coherent-lemma-quasi-coherence-higher-direct-images})
on se ramène à prouver que $H^i(\text{Cris}(X/S), \mathcal{F}) = 0$ pour $i \gg 0$
dans la situation de la proposition \ref{crystalline-proposition-compute-cohomology-crystal},
lorsque $C$ est une $A$-algèbre de type fini. C'est clair, car on peut
choisir une algèbre de polynômes en un nombre fini de variables et l'on a $\Omega^i_D = 0$
pour $i \gg 0$.
\end{remark}

\begin{remark}[Bornes particulières]
\label{crystalline-remark-bounded-cohomology-over-point}
Dans la situation \ref{crystalline-situation-global}, soit $\mathcal{F}$ un cristal en
$\mathcal{O}_{X/S}$-modules quasi-cohérents. Supposons que $S_0$
ait un unique point et que $X \to S_0$ soit de présentation finie.
\begin{enumerate}
\item Si $\dim X = d$ et si $X/S_0$ est de dimension de plongement $e$, alors
$H^i(\text{Cris}(X/S), \mathcal{F}) = 0$ pour $i > d + e$.
\item Si $X$ est séparé et peut être recouvert par $q$ ouverts affines, et si
$X/S_0$ est de dimension de plongement $e$, alors
$H^i(\text{Cris}(X/S), \mathcal{F}) = 0$ pour $i > q + e$.
\end{enumerate}
Indications : dans le cas (1), on peut utiliser l'égalité
$$
H^i(\text{Cris}(X/S), \mathcal{F}) = H^i(X_{Zar}, Ru_{X/S, *}\mathcal{F})
$$
et le fait que $Ru_{X/S, *}\mathcal{F}$ se calcule localement par un complexe de de Rham
construit à l'aide d'une immersion de $X$ dans un schéma lisse
de dimension $e$ sur $S$
(voir le lemme \ref{crystalline-lemma-compute-cohomology-crystal-smooth}).
Ces complexes de de Rham sont nuls en tout degré $> e$. Ainsi, (1)
résulte de Cohomologie, proposition
\ref{cohomology-proposition-vanishing-Noetherian}.
Dans le cas (2), on utilise le complexe de {\v C}ech alterné (voir la
remarque \ref{crystalline-remark-alternating-cech-complex}) pour se ramener au cas où
$X$ est affine. Dans le cas affine, on démontre le résultat à l'aide du complexe de de Rham
associé à une immersion de $X$ dans un schéma lisse de dimension $e$
sur $S$ (la construction d'un tel objet demande un peu de travail).
\end{remark}

\begin{remark}[Morphisme de changement de base]
\label{crystalline-remark-base-change}
Dans la situation de la remarque \ref{crystalline-remark-compute-direct-image},
supposons que $S = \Spec(A)$ et $S' = \Spec(A')$ soient affines.
Soit $\mathcal{F}'$ un $\mathcal{O}_{X'/S'}$-module.
Soit $\mathcal{F}$ l'image inverse de $\mathcal{F}'$.
Il existe alors un morphisme canonique de changement de base
$$
L(S' \to S)^*R\tau_{X'/S', *}\mathcal{F}'
\longrightarrow
R\tau_{X/S, *}\mathcal{F}
$$
où $\tau_{X/S}$ et $\tau_{X'/S'}$ sont les morphismes structuraux ; voir la
remarque \ref{crystalline-remark-structure-morphism}. Après passage aux sections globales,
on obtient un morphisme de changement de base
\begin{equation}
\label{crystalline-equation-base-change-map}
R\Gamma(\text{Cris}(X'/S'), \mathcal{F}') \otimes^\mathbf{L}_{A'} A
\longrightarrow
R\Gamma(\text{Cris}(X/S), \mathcal{F})
\end{equation}
dans $D(A)$.

\medskip\noindent
Indication : on compose le morphisme très général de changement de base de
Cohomologie sur les sites, remarque \ref{sites-cohomology-remark-base-change},
avec le morphisme canonique
$Lf_{\text{cris}}^*\mathcal{F}' \to
f_{\text{cris}}^*\mathcal{F}' = \mathcal{F}$.
\end{remark}

\begin{remark}[Isomorphisme de changement de base]
\label{crystalline-remark-base-change-isomorphism}
Le morphisme (\ref{crystalline-equation-base-change-map}) est un isomorphisme pourvu que
toutes les conditions suivantes soient satisfaites :
\begin{enumerate}
\item $p$ est nilpotent dans $A'$,
\item $\mathcal{F}'$ est un cristal en
$\mathcal{O}_{X'/S'}$-modules quasi-cohérents,
\item $X' \to S'_0$ est un morphisme quasi-compact et quasi-séparé,
\item $X = X' \times_{S'_0} S_0$,
\item $\mathcal{F}'$ est un $\mathcal{O}_{X'/S'}$-module plat,
\item $X' \to S'_0$ est un morphisme d'intersection complète locale (voir
Compléments sur les morphismes, définition \ref{more-morphisms-definition-lci} ; c'est
par exemple le cas si $X' \to S'_0$ est syntomique ou lisse),
\item $X'$ et $S_0$ sont Tor-indépendants sur $S'_0$ (voir
Compléments d'algèbre, définition \ref{more-algebra-definition-tor-independent} ;
c'est par exemple le cas si $S_0 \to S'_0$ ou $X' \to S'_0$ est plat).
\end{enumerate}
Indications : la condition (1) signifie que, dans les arguments ci-dessous, la complétion $p$-adique
n'a aucun effet et peut être ignorée.
À l'aide de la condition (3) et de Mayer--Vietoris (voir la
remarque \ref{crystalline-remark-mayer-vietoris}), on se ramène au cas
où $X'$ est affine. En fait, par la condition (6), quitte à restreindre encore,
on peut supposer que $X' = \Spec(C')$ et que l'on dispose d'une présentation
$C' = A'/I'[x_1, \ldots, x_n]/(\bar f'_1, \ldots, \bar f'_c)$
où $\bar f'_1, \ldots, \bar f'_c$ est une suite Koszul-régulière dans $A'/I'$.
(Cela signifie que, localement pour la topologie lisse, $\bar f'_1, \ldots, \bar f'_c$ forment
une suite régulière ; voir Compléments d'algèbre,
lemme \ref{more-algebra-lemma-Koszul-regular-flat-locally-regular}.)
Choisissons un relèvement de
$\bar f'_i$ en un élément $f'_i \in A'[x_1, \ldots, x_n]$. D'après (4), on voit que
$X = \Spec(C)$, où $C = A/I[x_1, \ldots, x_n]/(\bar f_1, \ldots, \bar f_c)$,
où $f_i \in A[x_1, \ldots, x_n]$ est l'image de $f'_i$.
La propriété (7) montre que $\bar f_1, \ldots, \bar f_c$ est une suite Koszul-régulière
dans $A/I[x_1, \ldots, x_n]$. L'enveloppe à puissances divisées de
$I'A'[x_1, \ldots, x_n] + (f'_1, \ldots, f'_c)$ dans $A'[x_1, \ldots, x_n]$
relativement à $\gamma'$ est
$$
D' = A'[x_1, \ldots, x_n]\langle \xi_1, \ldots, \xi_c \rangle/(\xi_i - f'_i)
$$
voir le lemme \ref{crystalline-lemma-describe-divided-power-envelope}. On vérifie alors que
$\xi_1 - f'_1, \ldots, \xi_c - f'_c$ est une suite Koszul-régulière dans
l'anneau $A'[x_1, \ldots, x_n]\langle \xi_1, \ldots, \xi_c\rangle$.
De même, l'enveloppe à puissances divisées de
$IA[x_1, \ldots, x_n] + (f_1, \ldots, f_c)$ dans $A[x_1, \ldots, x_n]$
relativement à $\gamma$ est
$$
D = A[x_1, \ldots, x_n]\langle \xi_1, \ldots, \xi_c\rangle/(\xi_i - f_i)
$$
et $\xi_1 - f_1, \ldots, \xi_c - f_c$ est une suite Koszul-régulière dans
l'anneau $A[x_1, \ldots, x_n]\langle \xi_1, \ldots, \xi_c\rangle$.
Il s'ensuit que $D' \otimes_{A'}^\mathbf{L} A = D$. La condition (2)
implique que $\mathcal{F}'$ correspond à un couple $(M', \nabla)$
formé d'un $D'$-module muni d'une connexion ; voir la
proposition \ref{crystalline-proposition-crystals-on-affine}.
Alors $M = M' \otimes_{D'} D$ correspond à l'image inverse $\mathcal{F}$.
La condition (5) montre que $M'$ est un $D'$-module plat, donc
$$
M = M' \otimes_{D'} D = M' \otimes_{D'} D' \otimes_{A'}^\mathbf{L} A
= M' \otimes_{A'}^\mathbf{L} A
$$
Comme les modules des différentielles $\Omega_{D'}$ et $\Omega_D$
(tels que définis dans la section \ref{crystalline-section-quasi-coherent-crystals})
sont respectivement un $D'$-module et un $D$-module libres sur les mêmes générateurs, on voit que
$$
M \otimes_D \Omega^\bullet_D =
M' \otimes_{D'} \Omega^\bullet_{D'} \otimes_{D'} D =
M' \otimes_{D'} \Omega^\bullet_{D'} \otimes_{A'}^\mathbf{L} A
$$
ce qui démontre l'assertion en vertu de la
proposition \ref{crystalline-proposition-compute-cohomology-crystal}.
\end{remark}

\begin{remark}[Rlim]
\label{crystalline-remark-rlim}
Soit $p$ un nombre premier. Soit $(A, I, \gamma)$ un anneau à puissances
divisées tel que $A$ soit une algèbre sur $\mathbf{Z}_{(p)}$ et que $p$ soit nilpotent
dans $A/I$. Posons $S = \Spec(A)$ et $S_0 = \Spec(A/I)$.
Soit $X$ un schéma sur $S_0$, tel que $p$ soit localement
nilpotent sur $X$. Soit $\mathcal{F}$ un
$\mathcal{O}_{X/S}$-module quelconque. Pour $e \gg 0$, l'idéal $(p^e) \subset I$
est stable par $\gamma$ ; voir
Algèbre à puissances divisées, lemme \ref{dpa-lemma-extend-to-completion}.
Posons $S_e = \Spec(A/p^eA)$ pour $e \gg 0$.
Alors $\text{Cris}(X/S_e)$ est une sous-catégorie pleine de $\text{Cris}(X/S)$,
et notons $\mathcal{F}_e$ la restriction de $\mathcal{F}$ à
$\text{Cris}(X/S_e)$. Alors
$$
R\Gamma(\text{Cris}(X/S), \mathcal{F}) =
R\lim_e R\Gamma(\text{Cris}(X/S_e), \mathcal{F}_e)
$$

\medskip\noindent
Indications : il suffit de le prouver lorsque $\mathcal{F}$ est injectif.
Dans ce cas, les faisceaux $\mathcal{F}_e$ sont eux aussi des
modules injectifs, les morphismes de transition
$\Gamma(\mathcal{F}_{e + 1}) \to \Gamma(\mathcal{F}_e)$ sont
surjectives, et l'on a
$\Gamma(\mathcal{F}) = \lim_e \Gamma(\mathcal{F}_e)$, car
tout objet de $\text{Cris}(X/S)$ est localement un objet de l'une
des catégories $\text{Cris}(X/S_e)$, par définition de
$\text{Cris}(X/S)$.
\end{remark}

\begin{remark}[Comparaison]
\label{crystalline-remark-comparison}
Soit $p$ un nombre premier. Soit $(A, I, \gamma)$ un anneau à puissances
divisées, avec $p$ nilpotent dans $A$. Posons $S = \Spec(A)$ et
$S_0 = \Spec(A/I)$. Soit $Y$ un schéma lisse sur $S$, et posons
$X = Y \times_S S_0$. Soit
$\mathcal{F}$ un cristal en $\mathcal{O}_{X/S}$-modules quasi-cohérents.
Alors
\begin{enumerate}
\item $\gamma$ s'étend en une structure à puissances divisées sur l'idéal
de $X$ dans $Y$, de sorte que $(X, Y, \gamma)$ soit un objet de $\text{Cris}(X/S)$,
\item la restriction $\mathcal{F}_Y$ (voir la section \ref{crystalline-section-sheaves})
est munie d'une connexion intégrable canonique
$\nabla : \mathcal{F}_Y \to
\mathcal{F}_Y \otimes_{\mathcal{O}_Y} \Omega_{Y/S}$, et
\item on a
$$
R\Gamma(\text{Cris}(X/S), \mathcal{F}) =
R\Gamma(Y, \mathcal{F}_Y \otimes_{\mathcal{O}_Y} \Omega^\bullet_{Y/S})
$$
dans $D(A)$.
\end{enumerate}
Indications : voir Algèbre à puissances divisées, lemme \ref{dpa-lemma-gamma-extends}, pour (1).
Voir le lemme \ref{crystalline-lemma-automatic-connection} pour (2).
Pour la partie (3), considérons le morphisme de
(\ref{crystalline-equation-restriction}). Ce morphisme est un isomorphisme lorsque
$X$ est affine ; voir le
lemme \ref{crystalline-lemma-compute-cohomology-crystal-smooth}.
Il en résulte que $Ru_{X/S, *}\mathcal{F}$ et
$\mathcal{F}_Y \otimes \Omega^\bullet_{Y/S}$ sont quasi-isomorphes
comme complexes sur $Y_{Zar} = X_{Zar}$.
Puisque $R\Gamma(\text{Cris}(X/S), \mathcal{F}) =
R\Gamma(X_{Zar}, Ru_{X/S, *}\mathcal{F})$, le résultat en découle.
\end{remark}

\begin{remark}[Caractère parfait]
\label{crystalline-remark-perfect}
Soit $p$ un nombre premier. Soit $(A, I, \gamma)$ un anneau à puissances
divisées, avec $p$ nilpotent dans $A$. Posons $S = \Spec(A)$ et
$S_0 = \Spec(A/I)$. Soit $X$ un schéma propre et lisse sur $S_0$.
Soit $\mathcal{F}$ un cristal en
$\mathcal{O}_{X/S}$-modules quasi-cohérents localement libres de type fini.
Alors $R\Gamma(\text{Cris}(X/S), \mathcal{F})$ est un
complexe parfait dans $D(A)$.

\medskip\noindent
Indications : la remarque \ref{crystalline-remark-base-change-isomorphism} donne
$$
R\Gamma(\text{Cris}(X/S), \mathcal{F}) \otimes_A^\mathbf{L} A/I
\cong
R\Gamma(\text{Cris}(X/S_0), \mathcal{F}|_{\text{Cris}(X/S_0)})
$$
La remarque \ref{crystalline-remark-comparison} donne
$$
R\Gamma(\text{Cris}(X/S_0), \mathcal{F}|_{\text{Cris}(X/S_0)}) =
R\Gamma(X, \mathcal{F}_X \otimes \Omega^\bullet_{X/S_0})
$$
En utilisant la filtration bête du complexe de de Rham, on voit que
le dernier complexe affiché est parfait dans $D(A/I)$ dès que les complexes
$$
R\Gamma(X, \mathcal{F}_X \otimes \Omega^q_{X/S_0})
$$
sont des complexes parfaits dans $D(A/I)$ ; voir
Compléments d'algèbre, lemme \ref{more-algebra-lemma-two-out-of-three-perfect}.
Cela résulte des arguments usuels
de cohomologie cohérente, car $\mathcal{F}_X \otimes \Omega^q_{X/S_0}$
est un faisceau localement libre de type fini et $X \to S_0$ est propre et plat
(référence à insérer). En appliquant le
lemme \ref{more-algebra-lemma-perfect-modulo-nilpotent-ideal} de Compléments d'algèbre,
on voit que
$$
R\Gamma(\text{Cris}(X/S), \mathcal{F}) \otimes_A^\mathbf{L} A/I^n
$$
est un complexe parfait dans $D(A/I^n)$ pour tout $n$. Cela ne suffit pas tout à fait,
sauf si $A$ est noethérien. En effet, bien que $I$ soit localement nilpotent
par l'hypothèse de nilpotence de $p$ ; voir
Algèbre à puissances divisées, lemme \ref{dpa-lemma-nil},
on ne peut conclure que $I^n = 0$ pour un certain $n$. Un contre-exemple
est fourni par $\mathbf{F}_p\langle x \rangle$. Pour le démontrer en général lorsque
$\mathcal{F} = \mathcal{O}_{X/S}$, l'argument de
\url{https://math.columbia.edu/~dejong/wordpress/?p=2227}
convient. Lorsque les coefficients $\mathcal{F}$ sont non triviaux,
l'argument de \cite{Faltings-very} semble être le suivant. On se ramène au
cas $pA = 0$ par le lemme de Compléments d'algèbre
\ref{more-algebra-lemma-perfect-modulo-nilpotent-ideal}.
Dans ce cas, le morphisme de Frobenius $A \to A$, $a \mapsto a^p$, se factorise
en $A \to A/I \xrightarrow{\varphi} A$ (car $x^p = 0$ pour $x \in I$). Posons
$X^{(1)} = X \otimes_{A/I, \varphi} A$. Le morphisme de Frobenius absolu
de $X$ se factorise par un morphisme $F_X : X \to X^{(1)}$ (une sorte de
Frobenius relatif). Localement sur les ouverts affines, si $X = \Spec(C)$, alors
$X^{(1)} = \Spec( C \otimes_{A/I, \varphi} A)$
et $F_X$ correspond à $C \otimes_{A/I, \varphi} A \to C$,
$c \otimes a \mapsto c^pa$. Cela définit des morphismes de topos annelés
$$
(X/S)_{\text{cris}}
\xrightarrow{(F_X)_{\text{cris}}}
(X^{(1)}/S)_{\text{cris}}
\xrightarrow{u_{X^{(1)}/S}}
\Sh(X^{(1)}_{Zar})
$$
dont le composé est noté $\text{Frob}_X$. On montre alors que
$R\text{Frob}_{X, *}\mathcal{F}$ est représenté par un
complexe parfait de $\mathcal{O}_{X^{(1)}}$-modules (!)
par un calcul local.
\end{remark}

\begin{remark}[Caractère parfait dans le cas complet]
\label{crystalline-remark-complete-perfect}
Soit $p$ un nombre premier. Soit $(A, I, \gamma)$ un anneau à puissances
divisées tel que $A$ soit complet $p$-adiquement et que $p$ soit nilpotent dans $A/I$. Posons
$S = \Spec(A)$ et $S_0 = \Spec(A/I)$. Soit $X$ un schéma propre
et lisse sur $S_0$. Soit $\mathcal{F}$ un cristal en
$\mathcal{O}_{X/S}$-modules quasi-cohérents localement libres de type fini.
Alors $R\Gamma(\text{Cris}(X/S), \mathcal{F})$ est un
complexe parfait dans $D(A)$.

\medskip\noindent
Indications : on sait que $K = R\Gamma(\text{Cris}(X/S), \mathcal{F})$
est la limite dérivée $K = R\lim K_e$ des cohomologies sur $A/p^eA$ ;
voir la remarque \ref{crystalline-remark-rlim}.
Chaque $K_e$ est un complexe parfait dans $D(A/p^eA)$ d'après la
remarque \ref{crystalline-remark-perfect}.
Comme $A$ est complet $p$-adiquement, le résultat
résulte du
lemme \ref{more-algebra-lemma-Rlim-perfect-gives-complete} de Compléments d'algèbre.
\end{remark}

\begin{remark}[Comparaison dans le cas complet]
\label{crystalline-remark-complete-comparison}
Soit $p$ un nombre premier. Soit $(A, I, \gamma)$ un anneau à puissances
divisées tel que $A$ soit un anneau noethérien complet $p$-adiquement et que $p$ soit nilpotent
dans $A/I$. Posons $S = \Spec(A)$ et
$S_0 = \Spec(A/I)$. Soit $Y$ un schéma propre et lisse sur $S$, et posons
$X = Y \times_S S_0$. Soit $\mathcal{F}$ un cristal de type fini en
$\mathcal{O}_{X/S}$-modules quasi-cohérents. Alors
\begin{enumerate}
\item il existe un $\mathcal{O}_Y$-module cohérent $\mathcal{F}_Y$
muni d'une connexion intégrable
$$
\nabla :
\mathcal{F}_Y
\longrightarrow
\mathcal{F}_Y \otimes_{\mathcal{O}_Y} \Omega_{Y/S}
$$
tel que $\mathcal{F}_Y/p^e\mathcal{F}_Y$ soit le module muni de sa connexion
sur $A/p^eA$ obtenu dans la remarque \ref{crystalline-remark-comparison}, et
\item on a
$$
R\Gamma(\text{Cris}(X/S), \mathcal{F}) =
R\Gamma(Y, \mathcal{F}_Y \otimes_{\mathcal{O}_Y} \Omega^\bullet_{Y/S})
$$
dans $D(A)$.
\end{enumerate}
Indications : l'existence de $\mathcal{F}_Y$ résulte du théorème d'existence de Grothendieck
(référence à insérer). L'isomorphisme des cohomologies résulte
du fait que les deux membres sont calculés par $R\lim$ à partir des versions modulo $p^e$
(voir la remarque \ref{crystalline-remark-rlim} pour le membre de gauche ; utiliser le théorème
des fonctions formelles, voir le
théorème \ref{coherent-theorem-formal-functions} de Cohomologie des schémas
pour le membre de droite).
La remarque \ref{crystalline-remark-comparison} identifie les deux versions modulo $p^e$.
\end{remark}





\section{Image inverse par des morphismes purement inséparables}
\label{crystalline-section-pull-back-along-pth-root}

\noindent
Par revêtement $\alpha_p$, on entend un morphisme de la forme
$$
X' = \Spec(C[z]/(z^p - c)) \longrightarrow \Spec(C) = X
$$
où $C$ est une $\mathbf{F}_p$-algèbre et $c \in C$. De manière équivalente,
$X'$ est un $\alpha_p$-torseur sur $X$. Un {\it revêtement
$\alpha_p$ itéré}\footnote{Cette notation n'est pas standard.}
est un morphisme de schémas en caractéristique
$p$ qui, localement sur la cible, s'écrit comme un composé fini de
revêtements $\alpha_p$. Dans cette section, nous montrons que l'image inverse par
un tel morphisme induit un quasi-isomorphisme en cohomologie cristalline
après inversion du nombre premier $p$. En fait, nous démontrons une version précise
de ce résultat. Commençons par un lemme préliminaire dont la formulation
nécessite quelques notations.

\medskip\noindent
Supposons donné un homomorphisme d'anneaux $B \to B'$ et des quotients $\Omega_B \to \Omega$ et
$\Omega_{B'} \to \Omega'$ vérifiant les hypothèses de la
remarque \ref{crystalline-remark-base-change-connection}.
Ainsi, (\ref{crystalline-equation-base-change-map-complexes}) fournit un
morphisme canonique de complexes
$$
c_M^\bullet :
M \otimes_B \Omega^\bullet
\longrightarrow
M \otimes_B (\Omega')^\bullet
$$
pour tout $B$-module $M$ muni d'une connexion intégrable
$\nabla : M \to M \otimes_B \Omega_B$.

\medskip\noindent
Supposons donnés $a \in B$, $z \in B'$ et une application $\theta : B' \to B'$
vérifiant les hypothèses suivantes :
\begin{enumerate}
\item
\label{crystalline-item-d-a-zero}
$\text{d}(a) = 0$,
\item
\label{crystalline-item-direct-sum}
$\Omega' = B' \otimes_B \Omega \oplus B'\text{d}z$ ; nous écrivons
$\text{d}(f) = \text{d}_1(f) + \partial_z(f) \text{d}z$
où $\text{d}_1(f) \in B' \otimes \Omega$ et $\partial_z(f) \in B'$
pour tout $f \in B'$,
\item
\label{crystalline-item-theta-linear}
$\theta : B' \to B'$ est $B$-linéaire,
\item
\label{crystalline-item-integrate}
$\partial_z \circ \theta = a$,
\item
\label{crystalline-item-injective}
$B \to B'$ est universellement injectif (et donc $\Omega \to \Omega'$
est injectif),
\item
\label{crystalline-item-factor}
$af - \theta(\partial_z(f)) \in B$ pour tout $f \in B'$,
\item
\label{crystalline-item-horizontal}
$(\theta \otimes 1)(\text{d}_1(f)) - \text{d}_1(\theta(f)) \in \Omega$
pour tout $f \in B'$, où
$\theta \otimes 1 : B' \otimes \Omega \to B' \otimes \Omega$
\end{enumerate}
Ces conditions ne sont pas logiquement indépendantes.
Par exemple, l'hypothèse (\ref{crystalline-item-integrate}) implique
que $\partial_z(af - \theta(\partial_z(f))) = 0$.
Ainsi, si l'image de $B \to B'$ est l'ensemble des
éléments annulés par $\partial_z$, alors (\ref{crystalline-item-factor})
en découle. Un raisonnement analogue s'applique à la condition (\ref{crystalline-item-horizontal}).

\begin{lemma}
\label{crystalline-lemma-find-homotopy}
Dans la situation ci-dessus, il existe un morphisme de complexes
$$
e_M^\bullet :
M \otimes_B (\Omega')^\bullet
\longrightarrow
M \otimes_B \Omega^\bullet
$$
tel que $c_M^\bullet \circ e_M^\bullet$
et $e_M^\bullet \circ c_M^\bullet$ soient homotopes à
la multiplication par $a$.
\end{lemma}

\begin{proof}
Dans cette démonstration, tous les produits tensoriels sont pris sur $B$.
L'hypothèse (\ref{crystalline-item-direct-sum}) implique que
$$
M \otimes (\Omega')^i =
(B' \otimes M \otimes \Omega^i)
\oplus
(B' \text{d}z \otimes M \otimes \Omega^{i - 1})
$$
pour tout $i \geq 0$. Le groupe additif sous-jacent à
$M \otimes (\Omega')^i$ est engendré par les éléments de la forme
$f \omega$ et d'éléments de la forme $f \text{d}z \wedge \eta$,
où $f \in B'$, $\omega \in M \otimes \Omega^i$, et
$\eta \in M \otimes \Omega^{i - 1}$.

\medskip\noindent
Pour $f \in B'$, écrivons
$$
\epsilon(f) = af - \theta(\partial_z(f))
\quad\text{et}\quad
\epsilon'(f) = (\theta \otimes 1)(\text{d}_1(f)) - \text{d}_1(\theta(f))
$$
de sorte que $\epsilon(f) \in B$ et $\epsilon'(f) \in \Omega$ en vertu des
hypothèses (\ref{crystalline-item-factor}) et (\ref{crystalline-item-horizontal}).
Définissons $e_M^\bullet$ par les formules
$e^i_M(f\omega) = \epsilon(f) \omega$ et
$e^i_M(f \text{d}z \wedge \eta) = \epsilon'(f) \wedge \eta$.
Nous verrons ci-dessous que les applications $e^i_M$ définissent un morphisme de complexes.

\medskip\noindent
Définissons
$$
h^i : M \otimes_B (\Omega')^i \longrightarrow M \otimes_B (\Omega')^{i - 1}
$$
par les formules $h^i(f \omega) = 0$ et
$h^i(f \text{d}z \wedge \eta) = \theta(f) \eta$
pour les éléments ci-dessus. Nous affirmons que
$$
\text{d} \circ h + h \circ \text{d} = a - c_M^\bullet \circ e_M^\bullet
$$
Notons que la multiplication par $a$ est un morphisme de complexes
d'après (\ref{crystalline-item-d-a-zero}). Ainsi, puisque $c_M^\bullet$ est un morphisme injectif
de complexes en vertu de l'hypothèse (\ref{crystalline-item-injective}), on conclut que
$e_M^\bullet$ est un morphisme de complexes. Pour démontrer l'affirmation, calculons
\begin{align*}
(\text{d} \circ h + h \circ \text{d})(f \omega)
& =
h\left(\text{d}(f) \wedge \omega + f \nabla(\omega)\right)
\\
& =
\theta(\partial_z(f)) \omega
\\
& =
a f\omega - \epsilon(f)\omega 
\\
& =
a f \omega - c^i_M(e^i_M(f\omega))
\end{align*}
La deuxième égalité vient de ce que $\text{d}z$ n'intervient pas dans $\nabla(\omega)$,
et la troisième résulte de l'hypothèse (6). De même, on a
\begin{align*}
(\text{d} \circ h + h \circ \text{d})(f \text{d}z \wedge \eta)
& =
\text{d}(\theta(f) \eta) +
h\left(\text{d}(f) \wedge \text{d}z \wedge \eta -
f \text{d}z \wedge \nabla(\eta)\right)
\\
& =
\text{d}(\theta(f)) \wedge \eta + \theta(f) \nabla(\eta)
- (\theta \otimes 1)(\text{d}_1(f)) \wedge \eta
- \theta(f) \nabla(\eta)
\\
& =
\text{d}_1(\theta(f)) \wedge \eta +
\partial_z(\theta(f)) \text{d}z \wedge \eta -
(\theta \otimes 1)(\text{d}_1(f)) \wedge \eta
\\
& =
a f \text{d}z \wedge \eta - \epsilon'(f) \wedge \eta \\
& = a f \text{d}z \wedge \eta - c^i_M(e^i_M(f \text{d}z \wedge \eta))
\end{align*}
La deuxième égalité vient de ce que
$\text{d}(f) \wedge \text{d}z \wedge \eta =
- \text{d}z \wedge \text{d}_1(f) \wedge \eta$.
La quatrième égalité résulte de l'hypothèse (\ref{crystalline-item-integrate}).
D'autre part, il résulte immédiatement des définitions
que $e^i_M(c^i_M(\omega)) = \epsilon(1) \omega = a \omega$.
Cela démontre le lemme.
\end{proof}

\begin{example}
\label{crystalline-example-integrate}
Un exemple standard de la situation ci-dessus se présente lorsque
$B' = B\langle z \rangle$ est l'algèbre de polynômes à puissances divisées
sur un anneau à puissances divisées $(B, J, \delta)$, munie des puissances divisées
$\delta'$ sur $J' = B'_{+} + JB' \subset B'$. Plus précisément, prenons
$\Omega = \Omega_{B, \delta}$ et $\Omega' = \Omega_{B', \delta'}$.
Dans ce cas, on peut prendre $a = 1$ et
$$
\theta( \sum b_m z^{[m]} ) = \sum b_m z^{[m + 1]}
$$
Notons que
$$
f - \theta(\partial_z(f)) = f(0)
$$
est égal au terme constant. Il s'ensuit que, dans ce cas, le
lemme \ref{crystalline-lemma-find-homotopy}
redonne le lemme de Poincar\'e cristallin
(lemme \ref{crystalline-lemma-relative-poincare}).
\end{example}

\begin{lemma}
\label{crystalline-lemma-computation}
Plaçons-nous dans la situation \ref{crystalline-situation-affine}. Supposons que $D$ et $\Omega_D$ soient comme dans
(\ref{crystalline-equation-D}) et (\ref{crystalline-equation-omega-D}).
Soit $\lambda \in D$. Soit $D'$ le complété $p$-adique de
$$
D[z]\langle \xi \rangle/(\xi - (z^p - \lambda))
$$
et soit $\Omega_{D'}$ le complété $p$-adique du module des
différentielles à puissances divisées de $D'$ sur $A$. Pour tout couple $(M, \nabla)$
sur $D$ vérifiant (\ref{crystalline-item-complete}), (\ref{crystalline-item-connection}),
(\ref{crystalline-item-integrable}) et (\ref{crystalline-item-topologically-quasi-nilpotent}),
le morphisme canonique de complexes (\ref{crystalline-equation-base-change-map-complexes})
$$
c_M^\bullet : M \otimes_D^\wedge \Omega^\bullet_D
\longrightarrow
M \otimes_D^\wedge \Omega^\bullet_{D'}
$$
a la propriété suivante : il existe un morphisme $e_M^\bullet$
dans le sens opposé tel que $c_M^\bullet \circ e_M^\bullet$
et $e_M^\bullet \circ c_M^\bullet$ soient toutes deux homotopes à la multiplication par $p$.
\end{lemma}

\begin{proof}
Nous allons le démontrer à l'aide du lemme \ref{crystalline-lemma-find-homotopy} avec $a = p$.
Il nous faut donc trouver $\theta : D' \to D'$ et vérifier
(\ref{crystalline-item-d-a-zero}), (\ref{crystalline-item-direct-sum}), (\ref{crystalline-item-theta-linear}),
(\ref{crystalline-item-integrate}), (\ref{crystalline-item-injective}), (\ref{crystalline-item-factor}),
(\ref{crystalline-item-horizontal}). Commençons par rassembler quelques informations sur les anneaux
$D$ et $D'$ et les modules $\Omega_D$ et $\Omega_{D'}$.

\medskip\noindent
En écrivant
$$
D[z]\langle \xi \rangle/(\xi - (z^p - \lambda))
=
D\langle \xi \rangle[z]/(z^p - \xi - \lambda)
$$
on voit que $D'$ est le complété $p$-adique du $D$-module libre
$$
\bigoplus\nolimits_{i = 0, \ldots, p - 1}
\bigoplus\nolimits_{n \geq 0}
z^i \xi^{[n]} D
$$
où $\xi^{[0]} = 1$.
Il s'ensuit que $D \to D'$ admet une section $D$-linéaire continue ; en particulier,
$D \to D'$ est universellement injectif, c'est-à-dire que (\ref{crystalline-item-injective}) est vérifiée.
Considérons $D'$ comme une algèbre à puissances divisées
sur $A$, d'idéal à puissances divisées $\overline{J}' = \overline{J}D' + (\xi)$.
Alors $D'$ est aussi le complété $p$-adique de l'enveloppe à puissances divisées
de l'idéal engendré par $z^p - \lambda$ dans $D$, voir le
lemme \ref{crystalline-lemma-describe-divided-power-envelope}. Ainsi,
$$
\Omega_{D'} = \Omega_D \otimes_D^\wedge D' \oplus D'\text{d}z
$$
d'après le lemme \ref{crystalline-lemma-module-differentials-divided-power-envelope}.
Cela démontre (\ref{crystalline-item-direct-sum}). Notons que (\ref{crystalline-item-d-a-zero})
est évidente.

\medskip\noindent
À ce stade, construisons $\theta$. (Nous avons écrit un script PARI/gp theta.gp
qui vérifie certaines des formules de cette démonstration et se trouve dans le
sous-répertoire scripts du projet Champs.) Avant cela, calculons
la différentielle des éléments $z^i \xi^{[n]}$. On a
$\text{d}z^i = i z^{i - 1} \text{d}z$. Pour $n \geq 1$, on a
$$
\text{d}\xi^{[n]} =
\xi^{[n - 1]} \text{d}\xi =
- \xi^{[n - 1]}\text{d}\lambda + p z^{p - 1} \xi^{[n - 1]}\text{d}z
$$
car $\xi = z^p - \lambda$. Pour $0 < i < p$ et $n \geq 1$, on a
\begin{align*}
\text{d}(z^i\xi^{[n]})
& =
iz^{i - 1}\xi^{[n]}\text{d}z + z^i\xi^{[n - 1]}\text{d}\xi \\
& =
iz^{i - 1}\xi^{[n]}\text{d}z + z^i\xi^{[n - 1]}\text{d}(z^p - \lambda) \\
& =
- z^i\xi^{[n - 1]}\text{d}\lambda +
(iz^{i - 1}\xi^{[n]} + pz^{i + p - 1}\xi^{[n - 1]})\text{d}z \\
& =
- z^i\xi^{[n - 1]}\text{d}\lambda +
(iz^{i - 1}\xi^{[n]} + pz^{i - 1}(\xi + \lambda)\xi^{[n - 1]})\text{d}z \\
& =
- z^i\xi^{[n - 1]}\text{d}\lambda +
((i + pn)z^{i - 1}\xi^{[n]} + p\lambda z^{i - 1}\xi^{[n - 1]})\text{d}z
\end{align*}
la dernière égalité résultant de $\xi \xi^{[n - 1]} = n\xi^{[n]}$.
On voit donc que
\begin{align*}
\partial_z(z^i) & = i z^{i - 1} \\
\partial_z(\xi^{[n]}) & = p z^{p - 1} \xi^{[n - 1]} \\
\partial_z(z^i\xi^{[n]}) & =
(i + pn) z^{i - 1} \xi^{[n]} + p \lambda z^{i - 1}\xi^{[n - 1]}
\end{align*}
Motivés par ces formules, définissons $\theta$ par les règles
$$
\begin{matrix}
\theta(z^j)
& = & p\frac{z^{j + 1}}{j + 1}
& j = 0, \ldots p - 1, \\
\theta(z^{p - 1}\xi^{[m]})
& = & \xi^{[m + 1]}
& m \geq 1, \\
\theta(z^j \xi^{[m]})
& = &
\frac{p z^{j + 1} \xi^{[m]} - \theta(p\lambda z^j \xi^{[m - 1]})}{(j + 1 + pm)}
& 0 \leq j < p - 1, m \geq 1
\end{matrix}
$$
où, dans la dernière ligne, nous procédons par récurrence sur $m$ pour définir $\theta$.
En développant, on obtient (pour $0 \leq j < p - 1$ et $1 \leq m$)
$$
\theta(z^j \xi^{[m]}) =
\textstyle{\frac{p z^{j + 1} \xi^{[m]}}{(j + 1 + pm)} -
\frac{p^2 \lambda z^{j + 1} \xi^{[m - 1]}}{(j + 1 + pm)(j + 1 + p(m - 1))} +
\ldots +
\frac{(-1)^m p^{m + 1} \lambda^m z^{j + 1}}
{(j + 1 + pm) \ldots (j + 1)}}
$$
bien que nous n'utilisions pas cette expression ci-dessous. Il est clair que $\theta$
se prolonge de manière unique en une application $p$-adiquement continue et $D$-linéaire sur $D'$.
Par construction, on a (\ref{crystalline-item-theta-linear}) et (\ref{crystalline-item-integrate}).
Il reste à démontrer (\ref{crystalline-item-factor}) et (\ref{crystalline-item-horizontal}).

\medskip\noindent
Démonstration de (\ref{crystalline-item-factor}) et (\ref{crystalline-item-horizontal}).
Puisque $\theta$ est $D$-linéaire et continue, il suffit de montrer que
$p - \theta \circ \partial_z$,
resp.\ $(\theta \otimes 1) \circ \text{d}_1 - \text{d}_1 \circ \theta$
prend ses valeurs dans $D$, resp.\ dans $\Omega_D$, lorsqu'on l'évalue sur les
éléments $z^i\xi^{[n]}$\footnote{Cela peut se vérifier par un calcul direct :
on constate que $p - \theta \circ \partial_z$, évalué en
$z^i\xi^{[n]}$, est nul, sauf en $1$, qui est envoyé sur $p$, et en
$\xi$, qui est envoyé sur $-p\lambda$. On constate que
$(\theta \otimes 1) \circ \text{d}_1 - \text{d}_1 \circ \theta$
évalué en $z^i\xi^{[n]}$ est nul, sauf en $z^{p - 1}\xi$,
qui est envoyé sur $-\lambda$.}.
Posons $D_0 = \mathbf{Z}_{(p)}[\lambda]$ et
$D_0' = \mathbf{Z}_{(p)}[z, \lambda]\langle \xi \rangle/(\xi - z^p + \lambda)$.
Observons que chacune des expressions ci-dessus est un élément de
$D_0'$ ou de $\Omega_{D_0'}$. Il suffit donc de démontrer le résultat
dans le cas de $D_0 \to D_0'$. Notons que $D_0$ et $D_0'$
sont des anneaux sans torsion et que $D_0 \otimes \mathbf{Q} = \mathbf{Q}[\lambda]$
et $D'_0 \otimes \mathbf{Q} = \mathbf{Q}[z, \lambda]$.
Ainsi, $D_0 \subset D'_0$ est le sous-anneau des éléments annulés
par $\partial_z$, et (\ref{crystalline-item-factor})
résulte de (\ref{crystalline-item-integrate}) ; voir la discussion qui précède immédiatement le
lemme \ref{crystalline-lemma-find-homotopy}. De même, on a
$\text{d}_1(f) = \partial_\lambda(f)\text{d}\lambda$, donc
$$
\left((\theta \otimes 1) \circ \text{d}_1 - \text{d}_1 \circ \theta\right)(f)
=
\left(\theta(\partial_\lambda(f)) - \partial_\lambda(\theta(f))\right)
\text{d}\lambda
$$
En appliquant $\partial_z$ au coefficient, on obtient
\begin{align*}
\partial_z\left(
\theta(\partial_\lambda(f)) - \partial_\lambda(\theta(f))
\right)
& =
p \partial_\lambda(f) - \partial_z(\partial_\lambda(\theta(f))) \\
& =
p \partial_\lambda(f) - \partial_\lambda(\partial_z(\theta(f))) \\
& =
p \partial_\lambda(f) - \partial_\lambda(p f) = 0
\end{align*}
de sorte que le coefficient ne dépend pas de $z$, comme souhaité.
Cela achève la démonstration du lemme.
\end{proof}

\noindent
Notons qu'un revêtement $\alpha_p$ itéré $X' \to X$ (tel que défini dans
l'introduction de cette section) est fini localement libre. Ainsi, si $X$ est
connexe, le degré de $X' \to X$ est constant et est une puissance de $p$.

\begin{lemma}
\label{crystalline-lemma-pullback-along-p-power-cover}
Soit $p$ un nombre premier. Soit $(S, \mathcal{I}, \gamma)$ un schéma à puissances
divisées sur $\mathbf{Z}_{(p)}$ tel que $p \in \mathcal{I}$. Posons
$S_0 = V(\mathcal{I}) \subset S$. Soit $f : X' \to X$ un revêtement
$\alpha_p$ itéré de schémas sur $S_0$, de degré constant $q$. Soit
$\mathcal{F}$ un cristal quelconque en faisceaux quasi-cohérents sur $X$, et posons
$\mathcal{F}' = f_{\text{cris}}^*\mathcal{F}$.
Dans le triangle distingué
$$
Ru_{X/S, *}\mathcal{F}
\longrightarrow
f_*Ru_{X'/S, *}\mathcal{F}'
\longrightarrow
E
\longrightarrow
Ru_{X/S, *}\mathcal{F}[1]
$$
les faisceaux de cohomologie de $E$ sont annulés par $q$.
\end{lemma}

\begin{proof}
Notons que $X' \to X$ est un homéomorphisme ; on peut donc identifier les espaces
topologiques sous-jacents à $X$ et $X'$. La question est manifestement locale sur $X$ ;
on peut donc supposer $X$, $X'$ et $S$ affines, et que $X' \to X$ s'écrit comme le
composé
$$
X' = X_n \to X_{n - 1} \to X_{n - 2} \to \ldots \to X_0 = X
$$
où chaque morphisme $X_{i + 1} \to X_i$ est un revêtement $\alpha_p$.
Notons $\mathcal{F}_i$ l'image inverse de $\mathcal{F}$ sur $X_i$.
Il suffit de démontrer que chacune des applications
$$
R\Gamma(\text{Cris}(X_i/S), \mathcal{F}_i)
\longrightarrow
R\Gamma(\text{Cris}(X_{i + 1}/S), \mathcal{F}_{i + 1})
$$
s'insère dans un triangle dont les groupes de cohomologie du troisième terme sont annulés
par $p$. (On utilise ici l'axiome TR4 de la catégorie triangulée $D(X)$. Les détails sont
omis.)

\medskip\noindent
On peut donc supposer que $S = \Spec(A)$, $X = \Spec(C)$, $X' = \Spec(C')$
et $C' = C[z]/(z^p - c)$ pour un certain $c \in C$. Choisissons une algèbre de polynômes
$P$ sur $A$ et une surjection $P \to C$. Soit $D$ le complété $p$-adique de
l'enveloppe à puissances divisées de $\Ker(P \to C)$ dans $P$, comme dans (\ref{crystalline-equation-D}).
Posons $P' = P[z]$, muni de la surjection $P' \to C'$ qui envoie $z$ sur la classe de $z$
dans $C'$. Choisissons un relèvement $\lambda \in D$ de $c \in C$. On voit alors que
le complété $p$-adique, noté $D'$, de l'enveloppe à puissances divisées de
$\Ker(P' \to C')$ dans $P'$ est isomorphe au complété $p$-adique de
$D[z]\langle \xi \rangle/(\xi - (z^p - \lambda))$ ; voir le
lemme \ref{crystalline-lemma-computation} et sa démonstration.
On voit donc que le résultat découle de ce lemme
par le calcul de la cohomologie des cristaux en modules quasi-cohérents de la
proposition \ref{crystalline-proposition-compute-cohomology-crystal}.
\end{proof}

\noindent
La borne du lemme suivant n'est probablement pas optimale.

\begin{lemma}
\label{crystalline-lemma-pullback-along-p-power-cover-cohomology}
Sous les hypothèses et avec les notations du
lemme \ref{crystalline-lemma-pullback-along-p-power-cover},
l'application
$$
f^* :
H^i(\text{Cris}(X/S), \mathcal{F})
\longrightarrow
H^i(\text{Cris}(X'/S), \mathcal{F}')
$$
a un noyau et un conoyau annulés par $q^{i + 1}$.
\end{lemma}

\begin{proof}
Cela résulte du fait que $E$ n'a de faisceaux de cohomologie non nuls qu'en
degrés $-1$ et supérieurs, de sorte que la suite spectrale
$H^a(\mathcal{H}^b(E)) \Rightarrow H^{a + b}(E)$ converge.
En combinant cela avec la suite exacte longue de cohomologie associée
à un triangle distingué, on obtient la borne.
\end{proof}

\noindent
Dans la situation \ref{crystalline-situation-global}, supposons que $p \in \mathcal{I}$.
Posons
$$
X^{(1)} = X \times_{S_0, F_{S_0}} S_0.
$$
Notons $F_{X/S_0} : X \to X^{(1)}$ le morphisme de Frobenius relatif.

\begin{lemma}
\label{crystalline-lemma-pullback-relative-frobenius}
Dans la situation ci-dessus, supposons que $X \to S_0$ soit lisse de dimension
relative $d$. Alors $F_{X/S_0}$ est un revêtement $\alpha_p$ itéré
de degré $p^d$. Les lemmes \ref{crystalline-lemma-pullback-along-p-power-cover} et
\ref{crystalline-lemma-pullback-along-p-power-cover-cohomology} s'appliquent donc à cette
situation. En particulier, pour tout cristal en modules quasi-cohérents
$\mathcal{G}$ sur $\text{Cris}(X^{(1)}/S)$, l'application
$$
F_{X/S_0}^* : H^i(\text{Cris}(X^{(1)}/S), \mathcal{G})
\longrightarrow
H^i(\text{Cris}(X/S), F_{X/S_0, \text{cris}}^*\mathcal{G})
$$
a un noyau et un conoyau annulés par $p^{d(i + 1)}$.
\end{lemma}

\begin{proof}
Il suffit de démontrer le premier énoncé. Pour cela, on peut supposer
que $X$ est \'etale sur $\mathbf{A}^d_{S_0}$ ; voir
Morphismes, lemme \ref{morphisms-lemma-smooth-etale-over-affine-space}.
Notons $\varphi : X \to \mathbf{A}^d_{S_0}$ ce morphisme \'etale.
Dans ce cas, le Frobenius relatif de $X/S_0$ s'insère dans un diagramme
$$
\xymatrix{
X \ar[d] \ar[r] & X^{(1)} \ar[d] \\
\mathbf{A}^d_{S_0} \ar[r] & \mathbf{A}^d_{S_0}
}
$$
où la flèche horizontale inférieure est le morphisme de Frobenius relatif
de $\mathbf{A}^d_{S_0}$ sur $S_0$. C'est le morphisme qui élève
toutes les coordonnées à la puissance $p$-ième ; c'est donc un revêtement
$\alpha_p$ itéré. La démonstration s'achève en observant que le diagramme
est un carré cartésien ; voir
Morphismes \'etales, lemme \ref{etale-lemma-relative-frobenius-etale}.
\end{proof}













\section{Action de Frobenius sur la cohomologie cristalline}
\label{crystalline-section-frobenius}

\noindent
Dans cette section, nous démontrons que l'image inverse par Frobenius induit un quasi-isomorphisme
sur la cohomologie cristalline après inversion de $p$. Mais, ne serait-ce que
pour formuler ce résultat, nous devons nous placer dans une situation particulière.

\begin{situation}
\label{crystalline-situation-F-crystal}
Dans la situation \ref{crystalline-situation-global}, supposons ce qui suit :
\begin{enumerate}
\item $S = \Spec(A)$ pour un certain anneau à puissances divisées $(A, I, \gamma)$
avec $p \in I$,
\item on se donne un homomorphisme d'anneaux à puissances divisées $\sigma : A \to A$
tel que $\sigma(x) = x^p \bmod pA$ pour tout $x \in A$.
\end{enumerate}
\end{situation}

\noindent
Dans la situation \ref{crystalline-situation-F-crystal}, le morphisme
$\Spec(\sigma) : S \to S$ est un relèvement du Frobenius absolu
$F_{S_0} : S_0 \to S_0$, et le diagramme
$$
\xymatrix{
X \ar[d] \ar[r]_{F_X} & X \ar[d] \\
S_0 \ar[r]^{F_{S_0}} & S_0
}
$$
est commutatif, où $F_X : X \to X$ est le morphisme de Frobenius absolu
de $X$. On obtient donc un morphisme de topos cristallins
$$
(F_X)_{\text{cris}} :
(X/S)_{\text{cris}}
\longrightarrow
(X/S)_{\text{cris}}
$$
Ce morphisme est décrit dans la remarque \ref{crystalline-remark-functoriality-cris}. Voici la terminologie relative aux
$F$-cristaux, suivant la notation de Saavedra ; voir
\cite{Saavedra}.

\begin{definition}
\label{crystalline-definition-F-crystal}
Dans la situation \ref{crystalline-situation-F-crystal}, un {\it $F$-cristal sur $X/S$
(relatif à $\sigma$)} est un couple $(\mathcal{E}, F_\mathcal{E})$
formé d'un cristal en modules localement libres de type fini sur $\mathcal{O}_{X/S}$,
noté $\mathcal{E}$, et d'un morphisme
$$
F_\mathcal{E} : (F_X)_{\text{cris}}^*\mathcal{E} \longrightarrow \mathcal{E}
$$
Un $F$-cristal est dit {\it non-dégénéré} s'il existe un entier
$i \geq 0$ et un morphisme $V : \mathcal{E} \to (F_X)_{\text{cris}}^*\mathcal{E}$
tel que $V \circ F_{\mathcal{E}} = p^i \text{id}$.
\end{definition}

\begin{remark}
\label{crystalline-remark-F-crystal-variants}
Soit $(\mathcal{E}, F)$ un $F$-cristal comme dans la
définition \ref{crystalline-definition-F-crystal}.
Dans la littérature, la condition de non-dégénérescence fait souvent partie de la
définition d'un $F$-cristal. En outre, on suppose aussi souvent que
$F \circ V = p^n\text{id}$. Ce qui est nécessaire pour le résultat ci-dessous, c'est
qu'il existe un entier $j \geq 0$ tel que $\Ker(F)$ et
$\Coker(F)$ soient annulés par $p^j$. Si le rang de $\mathcal{E}$
est borné (par exemple si $X$ est quasi-compact), ces deux
conditions résultent de la condition de non-dégénérescence telle qu'elle est formulée dans
la définition. En effet, soient $R$ un anneau, $r \geq 1$ un entier et
$K, L \in \text{Mat}(r \times r, R)$ des matrices telles que
$K L = p^i 1_{r \times r}$. Alors $\det(K)\det(L) = p^{ri}$. 
Soit $L'$ la comatrice de $L$, c'est-à-dire
$L' L = L L' = \det(L)$. Posons $K' = p^{ri} K$ et $j = ri + i$.
On a alors $K' L = p^j 1_{r \times r}$ puisque $K L = p^i$, et
$$
L K' = L K \det(L) \det(K) = L K L L' \det(K) = L p^i L' \det(K) =
p^j 1_{r \times r}
$$
Il s'ensuit que, si $V$ est comme dans la définition \ref{crystalline-definition-F-crystal},
alors, en posant $V' = p^N V$, où $N > i \cdot \text{rang}(\mathcal{E})$,
on obtient $V' \circ F = p^{N + i}$ et $F \circ V' = p^{N + i}$.
\end{remark}

\begin{theorem}
\label{crystalline-theorem-cohomology-F-crystal}
Dans la situation \ref{crystalline-situation-F-crystal}, soit $(\mathcal{E}, F_\mathcal{E})$
un $F$-cristal non dégénéré. Supposons que $A$ soit un anneau complet pour la topologie $p$-adique
et noethérien, et que $X \to S_0$ soit propre et lisse. Alors
le morphisme canonique
$$
F_\mathcal{E} \circ (F_X)_{\text{cris}}^* :
R\Gamma(\text{Cris}(X/S), \mathcal{E}) \otimes^\mathbf{L}_{A, \sigma} A
\longrightarrow
R\Gamma(\text{Cris}(X/S), \mathcal{E})
$$
devient un isomorphisme après inversion de $p$.
\end{theorem}

\begin{proof}
Écrivons d'abord la flèche comme le composé de trois flèches.
Plus précisément, posons
$$
X^{(1)} = X \times_{S_0, F_{S_0}} S_0
$$
et notons $F_{X/S_0} : X \to X^{(1)}$ le morphisme de Frobenius relatif.
Notons $\mathcal{E}^{(1)}$ le cristal déduit de $\mathcal{E}$ par changement de base
par $\Spec(\sigma)$, autrement dit l'image inverse de $\mathcal{E}$
sur $\text{Cris}(X^{(1)}/S)$ par le morphisme de topos cristallins
associé au diagramme commutatif
$$
\xymatrix{
X^{(1)} \ar[r] \ar[d] & X \ar[d] \\
S \ar[r]^{\Spec(\sigma)} & S
}
$$
On a alors le morphisme de changement de base
\begin{equation}
\label{crystalline-equation-base-change-sigma}
R\Gamma(\text{Cris}(X/S), \mathcal{E}) \otimes^\mathbf{L}_{A, \sigma} A
\longrightarrow
R\Gamma(\text{Cris}(X^{(1)}/S), \mathcal{E}^{(1)})
\end{equation}
Ce morphisme est celui de la remarque \ref{crystalline-remark-base-change}. Notons que la composée
de $F_{X/S_0} : X \to X^{(1)}$ avec la projection $X^{(1)} \to X$
est le morphisme de Frobenius absolu $F_X$. On voit donc que
$F_{X/S_0}^*\mathcal{E}^{(1)} = (F_X)_{\text{cris}}^*\mathcal{E}$.
Ainsi, l'image inverse par $F_{X/S_0}$ donne un morphisme
\begin{equation}
\label{crystalline-equation-to-prove}
F_{X/S_0}^* :
R\Gamma(\text{Cris}(X^{(1)}/S), \mathcal{E}^{(1)})
\longrightarrow
R\Gamma(\text{Cris}(X/S), (F_X)^*_{\text{cris}}\mathcal{E})
\end{equation}
Enfin, on peut utiliser $F_\mathcal{E}$ pour obtenir un morphisme
\begin{equation}
\label{crystalline-equation-F-E}
R\Gamma(\text{Cris}(X/S), (F_X)^*_{\text{cris}}\mathcal{E})
\longrightarrow
R\Gamma(\text{Cris}(X/S), \mathcal{E})
\end{equation}
Le morphisme du théorème est la composée des trois morphismes
(\ref{crystalline-equation-base-change-sigma}), (\ref{crystalline-equation-to-prove}) et
(\ref{crystalline-equation-F-E}) ci-dessus. Le premier est un
quasi-isomorphisme modulo chaque puissance de $p$ d'après la
remarque \ref{crystalline-remark-base-change-isomorphism}.
C'est donc un quasi-isomorphisme puisque les complexes en jeu sont parfaits
dans $D(A)$ ; voir la remarque \ref{crystalline-remark-complete-perfect}.
Le troisième morphisme est un quasi-isomorphisme après inversion de $p$, simplement
parce que $F_\mathcal{E}$ admet un inverse à une puissance de $p$ près ; voir la
remarque \ref{crystalline-remark-F-crystal-variants}.
Enfin, le deuxième est un isomorphisme après inversion de $p$
d'après le lemme \ref{crystalline-lemma-pullback-relative-frobenius}.
\end{proof}










% OK, start here.
%

\chapter{Cohomologie pro-\'etale}



\phantomsection
\label{proetale-section-phantom}


\section{Introduction}
\label{proetale-section-introduction}

\noindent
Le contenu de ce chapitre, et davantage encore, se trouve dans la pr\'epublication \cite{BS}.

\medskip\noindent
Le but de ce chapitre est d'introduire la topologie pro-\'etale et de
d\'evelopper les fondements de la cohomologie des faisceaux ab\'eliens dans cette topologie.
Un objectif secondaire est de montrer comment l'emploi de la topologie pro-\'etale simplifie
l'introduction de la cohomologie $\ell$-adique en g\'eom\'etrie alg\'ebrique.

\medskip\noindent
Voici un bref aper\c{c}u de l'histoire de la cohomologie \'etale $\ell$-adique
telle que nous l'avons comprise.
Dans \cite[Expos\'es V et VI]{SGA5}, Grothendieck et ses collaborateurs ont d\'evelopp\'e une th\'eorie
qui traite les faisceaux $\ell$-adiques comme des syst\`emes projectifs de faisceaux de
$\mathbf{Z}/\ell^n\mathbf{Z}$-modules.
Dans son second article sur les conjectures de Weil (\cite{WeilII}), Deligne a introduit
une cat\'egorie d\'eriv\'ee de faisceaux $\ell$-adiques comme une certaine 2-limite de cat\'egories
de complexes de faisceaux de $\mathbf{Z}/\ell^n\mathbf{Z}$-modules sur le
site \'etale d'un sch\'ema $X$. Cette approche sert, dans l'article de
Beilinson, Bernstein et Deligne (\cite{BBD}), de fondement \`a leur
belle th\'eorie des faisceaux pervers. Dans un article intitul\'e ``Cohomologie
\'etale continue'' (\cite{Jannsen}), Uwe Jannsen examine une variante importante
de la cohomologie d'un faisceau $\ell$-adique sur une vari\'et\'e sur un corps.
Son article est prolong\'e par un article de Torsten Ekedahl (\cite{Ekedahl}),
qui \`etudie le formalisme adique n\'ecessaire pour manier commod\'ement les cat\'egories
d\'eriv\'ees d\'efinies comme des limites.

\medskip\noindent
Il se trouve qu'en travaillant avec le site pro-\'etale d'un sch\'ema,
on peut \`eviter certaines des difficult\'es techniques rencontr\'ees par ces auteurs.
Cela oblige en contrepartie \`a travailler avec des sch\'emas non noeth\'eriens,
m\^eme lorsque l'on ne s'int\'eresse qu'aux faisceaux $\ell$-adiques
et \`a leur cohomologie sur des vari\'et\'es d\'efinies sur un corps alg\'ebriquement clos.

\medskip\noindent
Une propri\'et\'e tr\`es importante et remarquable du
(petit) site pro-\'etale d'un sch\'ema est
qu'il poss\`ede assez d'objets quasi-compacts w-contractiles. L'existence
de ces objets entra\^ine
plusieurs cons\'equences utiles et (peut-\^etre) inhabituelles pour la cat\'egorie d\'eriv\'ee
des faisceaux ab\'eliens et pour les syst\`emes projectifs de faisceaux.
C'est pr\'ecis\'ement cette propri\'et\'e qui nous permettra de traiter
les subtilit\'es propres aux faisceaux $\ell$-adiques; mais, comme nous
le verrons, elle pr\'esente aussi plusieurs autres avantages.




\section{Quelques \`el\'ements de topologie}
\label{proetale-section-topology}

\noindent
Quelques pr\'eliminaires. Nous avons d\'efini les {\it espaces spectraux} et les
{\it applications spectrales} entre espaces spectraux dans
Topologie, section \ref{topology-section-spectral}.
Le spectre d'un anneau est un espace spectral; voir
Alg\`ebre, lemme \ref{algebra-lemma-spec-spectral}.

\begin{lemma}
\label{proetale-lemma-spectral-split}
Soit $X$ un espace spectral. Soit $X_0 \subset X$ l'ensemble des points ferm\'es.
Les conditions suivantes sont \`equivalentes:
\begin{enumerate}
\item Tout recouvrement ouvert de $X$ peut \^etre raffin\'e par une d\'ecomposition finie
en union disjointe $X = \coprod U_i$, o\`u les $U_i$ sont
ouverts et ferm\'es dans $X$.
\item L'application compos\'ee $X_0 \to X \to \pi_0(X)$ est bijective.
\end{enumerate}
De plus, si $X_0$ est ferm\'e dans $X$ et si tout point de $X$ se sp\'ecialise
en un unique point de $X_0$, alors ces conditions sont satisfaites.
\end{lemma}

\begin{proof}
Nous utiliserons sans autre mention le fait que
$X_0$ est quasi-compact
(Topologie, lemme \ref{topology-lemma-closed-points-quasi-compact})
et que $\pi_0(X)$ est profini
(Topologie, lemme \ref{topology-lemma-spectral-pi0}).
Consid\'erons le diagramme
$$
\xymatrix{
X_0 \ar[rd]_f \ar[r] & X \ar[d]^\pi \\
& \pi_0(X)
}
$$
Si (2) est satisfaite, l'application continue bijective $f : X_0 \to \pi_0(X)$ est
un hom\'eomorphisme d'apr\`es
Topologie, lemme \ref{topology-lemma-bijective-map}.
\'Etant donn\'e un recouvrement ouvert $X = \bigcup U_i$, on obtient le recouvrement ouvert
$\pi_0(X) = \bigcup f(X_0 \cap U_i)$. D'apr\`es
Topologie, lemme \ref{topology-lemma-profinite-refine-open-covering},
il existe un recouvrement ouvert fini de la forme $\pi_0(X) = \coprod V_j$
qui raffine ce recouvrement.
Puisque $X_0 \to \pi_0(X)$ est bijective, chaque composante connexe de
$X$ poss\`ede un unique point ferm\'e et est donc l'ensemble des points qui se
sp\'ecialisent en ce point ferm\'e. Ainsi $\pi^{-1}(V_j)$ est
l'ensemble des points qui se sp\'ecialisent en des points de $f^{-1}(V_j)$.
Or, si $f^{-1}(V_j) \subset X_0 \cap U_i \subset U_i$, alors
$\pi^{-1}(V_j) \subset U_i$ (car l'ouvert
$U_i$ est stable par g\'en\'erisation). On voit donc
que le recouvrement ouvert $X = \coprod \pi^{-1}(V_j)$ raffine
le recouvrement de d\'epart. Cela montre que
(2) implique (1).

\medskip\noindent
Supposons (1). Soient $x, y \in X$ des points ferm\'es. On a alors le recouvrement ouvert
$X = (X \setminus \{x\}) \cup (X \setminus \{y\})$.
Il r\'esulte de (1) qu'il existe une d\'ecomposition en union disjointe
$X = U \amalg V$, avec $U$ et $V$ ouverts (et ferm\'es), $x \in U$ et
$y \in V$. En particulier, toute composante connexe de $X$
poss\`ede au plus un point ferm\'e. D'apr\`es
Topologie, lemme \ref{topology-lemma-quasi-compact-closed-point},
toute composante connexe (qui est ferm\'ee) poss\`ede aussi un point ferm\'e.
Ainsi $X_0 \to \pi_0(X)$ est bijective. Cela montre que (1) implique (2).

\medskip\noindent
Supposons que $X_0$ soit ferm\'e dans $X$ et que tout point se sp\'ecialise en un unique
point de $X_0$. Alors $X_0$ est un espace spectral
(Topologie, lemme \ref{topology-lemma-spectral-sub})
constitu\'e de points ferm\'es, donc profini
(Topologie, lemme \ref{topology-lemma-characterize-profinite-spectral}).
Soient $x, y \in X_0$ distincts. D'apr\`es
Topologie, lemme \ref{topology-lemma-profinite-refine-open-covering},
il existe une d\'ecomposition en union disjointe
$X_0 = U_0 \amalg V_0$, o\`u $U_0$ et $V_0$ sont ouverts et ferm\'es,
avec $x \in U_0$ et $y \in V_0$.
Soient $U \subset X$, resp.\ $V \subset X$, les ensembles form\'es des points
qui se sp\'ecialisent en un point de $U_0$, resp.\ de $V_0$.
On a $X = U \amalg V$.
D'apr\`es Topologie, lemme \ref{topology-lemma-make-spectral-space},
$U$ est une intersection d'ouverts quasi-compacts.
Ainsi $U$ est ferm\'e pour la topologie constructible.
Puisque $U$ est stable par sp\'ecialisation,
$U$ est ferm\'e d'apr\`es Topologie, lemme
\ref{topology-lemma-constructible-stable-specialization-closed}.
Par sym\'etrie, $V$ est ferm\'e; ainsi $U$ et $V$ sont tous deux
ouverts et ferm\'es.
Cela prouve que $x, y$ n'appartiennent pas \`a la m\^eme composante connexe de $X$.
Autrement dit, $X_0 \to \pi_0(X)$ est injective. Elle est aussi
surjective d'apr\`es
Topologie, lemme \ref{topology-lemma-quasi-compact-closed-point}
et le fait que les composantes connexes sont ferm\'ees.
On voit ainsi que la derni\`ere condition implique (2).
\end{proof}

\begin{example}
\label{proetale-example-not-w-local}
Soit $T$ un espace profini. Soit $t \in T$ un point et supposons
que $T \setminus \{t\}$ ne soit pas quasi-compact.
Posons $X = T \times \{0, 1\}$. Munissons $X$ de la topologie
dont une pr\'ebase est donn\'ee par les ensembles
$U \times \{0, 1\}$ pour $U \subset T$ ouvert, $X \setminus \{(t, 0)\}$,
et $U \times \{1\}$ pour $U \subset T$ ouvert tel que $t \not \in U$.
L'ensemble des points ferm\'es de $X$ est $X_0 = T \times \{0\}$, et
$(t, 1)$ appartient \`a l'adh\'erence de $X_0$.
De plus, $X_0 \to \pi_0(X)$ est bijective.
Cet exemple montre que les conditions (1) et (2) du
lemme \ref{proetale-lemma-spectral-split} n'impliquent pas que l'ensemble des points ferm\'es
soit ferm\'e.
\end{example}

\noindent
Il est en fait plus commode de travailler avec des espaces
spectraux qui poss\`edent la propri\'et\'e un peu plus forte mentionn\'ee dans
le dernier \`enonc\'e du lemme \ref{proetale-lemma-spectral-split}.
Nous donnons un nom \`a cette propri\'et\'e.

\begin{definition}
\label{proetale-definition-w-local}
Un espace spectral $X$ est dit {\it w-local} si l'ensemble $X_0$ de ses points ferm\'es
est ferm\'e et si tout point de $X$ se sp\'ecialise en un unique point ferm\'e.
Une application continue $f : X \to Y$ entre espaces w-locaux est dite {\it w-locale}
si elle est spectrale et envoie tout point ferm\'e de $X$ sur un point ferm\'e de $Y$.
\end{definition}

\noindent
Nous avons vu dans la d\'emonstration du lemme \ref{proetale-lemma-spectral-split}
que, dans ce cas, $X_0 \to \pi_0(X)$ est un hom\'eomorphisme et que
$X_0 \cong \pi_0(X)$ est un espace profini. De plus, une composante
connexe de $X$ est exactement l'ensemble des points qui se sp\'ecialisent en
un $x \in X_0$ donn\'e.

\begin{lemma}
\label{proetale-lemma-closed-subspace-w-local}
Soit $X$ un espace spectral w-local. Si $Y \subset X$ est ferm\'e,
alors $Y$ est w-local.
\end{lemma}

\begin{proof}
Le sous-ensemble $Y_0 \subset Y$ des points ferm\'es est ferm\'e puisque
$Y_0 = X_0 \cap Y$. Comme $X$ est $w$-local, tout $y \in Y$ se sp\'ecialise
en un unique point de $X_0$. Ce point appartient \`a $Y$, et donc
aussi \`a $Y_0$, puisque $\overline{\{y\}}\subset Y$. En conclusion, $Y$
est $w$-local.
\end{proof}

\begin{lemma}
\label{proetale-lemma-silly}
Soit $X$ un espace spectral. Supposons que
$$
\xymatrix{
Y \ar[r] \ar[d] & T \ar[d] \\
X \ar[r] & \pi_0(X)
}
$$
soit un diagramme cart\'esien dans la cat\'egorie des espaces topologiques,
o\`u $T$ est profini. Alors $Y$ est spectral et $T = \pi_0(Y)$.
Si de plus $X$ est w-local, alors $Y$ est w-local, $Y \to X$ est w-locale,
et l'ensemble des points ferm\'es de $Y$ est l'image r\'eciproque de
l'ensemble des points ferm\'es de $X$.
\end{lemma}

\begin{proof}
Notons que $Y$ est un sous-espace ferm\'e de $X \times T$, car $\pi_0(X)$
est un espace profini, donc s\'epar\'e
(utiliser Topologie, lemmes \ref{topology-lemma-spectral-pi0} et
\ref{topology-lemma-fibre-product-closed}).
Comme $X \times T$ est spectral
(Topologie, lemme \ref{topology-lemma-product-spectral-spaces}),
il s'ensuit que $Y$ est spectral
(Topologie, lemme \ref{topology-lemma-spectral-sub}).
Consid\'erons la factorisation canonique $Y \to \pi_0(Y) \to T$
(Topologie, lemme \ref{topology-lemma-space-connected-components}).
Il est clair que $\pi_0(Y) \to T$ est surjective.
Les fibres de $Y \to T$ sont hom\'eomorphes aux fibres de
$X \to \pi_0(X)$. Elles sont donc connexes. Il s'ensuit
que $\pi_0(Y) \to T$ est injective. On en conclut que $\pi_0(Y) \to T$
est un hom\'eomorphisme d'apr\`es
Topologie, lemme \ref{topology-lemma-bijective-map}.

\medskip\noindent
Supposons ensuite $X$ w-local et notons $X_0 \subset X$
l'ensemble des points ferm\'es. L'image r\'eciproque $Y_0 \subset Y$ de $X_0$ dans
$Y$ s'identifie à $T$ par une bijection, car $X_0 \to \pi_0(X)$ est bijective
d'apr\`es le lemme \ref{proetale-lemma-spectral-split}. De plus, $Y_0$ est quasi-compact
comme sous-ensemble ferm\'e de l'espace spectral $Y$. Par cons\'equent,
$Y_0 \to \pi_0(Y) = T$ est un hom\'eomorphisme d'apr\`es
Topologie, lemme \ref{topology-lemma-bijective-map}.
Il s'ensuit que tous les points de $Y_0$ sont ferm\'es dans $Y$.
R\'eciproquement, si $y \in Y$ est un point ferm\'e, alors
il est ferm\'e dans la fibre de $Y \to \pi_0(Y) = T$;
son image $x$ dans $X$ est donc ferm\'e dans la fibre (hom\'eomorphe) de
$X \to \pi_0(X)$. Cela implique $x \in X_0$, puis $y \in Y_0$.
Ainsi $Y_0$ est l'ensemble des points ferm\'es de $Y$,
et, pour tout $y \in Y_0$, l'ensemble des g\'en\'erisations de $y$ est
la fibre de $Y \to \pi_0(Y)$. Le lemme en r\'esulte.
\end{proof}




\section{Isomorphismes locaux}
\label{proetale-section-local-isomorphism}

\noindent
Commen\c{c}ons par une d\'efinition.

\begin{definition}
\label{proetale-definition-local-isomorphism}
Soit $\varphi : A \to B$ un homomorphisme d'anneaux.
\begin{enumerate}
\item On dit que $A \to B$ est un {\it isomorphisme local} si, pour tout id\'eal premier
$\mathfrak q \subset B$, il existe $g \in B$, $g \not \in \mathfrak q$,
tel que $A \to B_g$ induise une immersion ouverte $\Spec(B_g) \to \Spec(A)$.
\item On dit que $A \to B$ {\it identifie les anneaux locaux} si, pour tout id\'eal premier
$\mathfrak q \subset B$, l'application canonique
$A_{\varphi^{-1}(\mathfrak q)} \to B_\mathfrak q$ est un isomorphisme.
\end{enumerate}
\end{definition}

\noindent
\'Enum\'erons quelques propri\'et\'es \`el\'ementaires.

\begin{lemma}
\label{proetale-lemma-base-change-local-isomorphism}
Soient $A \to B$ et $A \to A'$ des homomorphismes d'anneaux. Notons $B' = B \otimes_A A'$,
l'alg\`ebre obtenue \`a partir de $B$ par changement de base.
\begin{enumerate}
\item Si $A \to B$ est un isomorphisme local, alors $A' \to B'$ est un
isomorphisme local.
\item Si $A \to B$ identifie les anneaux locaux, alors $A' \to B'$
identifie les anneaux locaux.
\end{enumerate}
\end{lemma}

\begin{proof}
La d\'emonstration est omise.
\end{proof}

\begin{lemma}
\label{proetale-lemma-composition-local-isomorphism}
Soient $A \to B$ et $B \to C$ des homomorphismes d'anneaux.
\begin{enumerate}
\item Si $A \to B$ et $B \to C$ sont des isomorphismes locaux, alors $A \to C$
est un isomorphisme local.
\item Si $A \to B$ et $B \to C$ identifient les anneaux locaux, alors $A \to C$
identifie les anneaux locaux.
\end{enumerate}
\end{lemma}

\begin{proof}
La d\'emonstration est omise.
\end{proof}

\begin{lemma}
\label{proetale-lemma-local-isomorphism-permanence}
Soit $A$ un anneau. Soit $B \to C$ un homomorphisme de $A$-alg\`ebres.
\begin{enumerate}
\item Si $A \to B$ et $A \to C$ sont des isomorphismes locaux, alors $B \to C$
est un isomorphisme local.
\item Si $A \to B$ et $A \to C$ identifient les anneaux locaux, alors $B \to C$
identifie les anneaux locaux.
\end{enumerate}
\end{lemma}

\begin{proof}
La d\'emonstration est omise.
\end{proof}

\begin{lemma}
\label{proetale-lemma-local-isomorphism-implies}
Soit $A \to B$ un isomorphisme local. Alors
\begin{enumerate}
\item $A \to B$ est \'etale,
\item $A \to B$ identifie les anneaux locaux,
\item $A \to B$ est quasi-fini.
\end{enumerate}
\end{lemma}

\begin{proof}
La d\'emonstration est omise.
\end{proof}

\begin{lemma}
\label{proetale-lemma-structure-local-isomorphism}
Soit $A \to B$ un isomorphisme local. Il existe alors $n \geq 0$,
$g_1, \ldots, g_n \in B$ et $f_1, \ldots, f_n \in A$ tels que
$(g_1, \ldots, g_n) = B$ et $A_{f_i} \cong B_{g_i}$.
\end{lemma}

\begin{proof}
La d\'emonstration est omise.
\end{proof}

\begin{lemma}
\label{proetale-lemma-fully-faithful-spaces-over-X}
Soient $p : (Y, \mathcal{O}_Y) \to (X, \mathcal{O}_X)$ et
$q : (Z, \mathcal{O}_Z) \to (X, \mathcal{O}_X)$
des morphismes d'espaces localement annel\'es.
Si $\mathcal{O}_Y = p^{-1}\mathcal{O}_X$, alors
$$
\Mor_{\text{LRS}/(X, \mathcal{O}_X)}((Z, \mathcal{O}_Z), (Y, \mathcal{O}_Y))
\longrightarrow
\Mor_{\textit{Top}/X}(Z, Y),\quad
(f, f^\sharp) \longmapsto f
$$
est bijective. Ici $\text{LRS}/(X, \mathcal{O}_X)$ est la cat\'egorie des
espaces localement annel\'es au-dessus de $X$, et $\textit{Top}/X$ la cat\'egorie
des espaces topologiques au-dessus de $X$.
\end{lemma}

\begin{proof}
Cela r\'esulte imm\'ediatement des d\'efinitions.
\end{proof}

\begin{lemma}
\label{proetale-lemma-local-isomorphism-fully-faithful}
Soit $A$ un anneau. Posons $X = \Spec(A)$. Le foncteur
$$
B \longmapsto \Spec(B)
$$
de la cat\'egorie des $A$-alg\`ebres $B$ telles que $A \to B$ identifie
les anneaux locaux vers la cat\'egorie des
espaces topologiques au-dessus de $X$ est pleinement fid\`ele.
\end{lemma}

\begin{proof}
Cela r\'esulte du lemme \ref{proetale-lemma-fully-faithful-spaces-over-X}
et du fait que, si $A \to B$ identifie les anneaux locaux, l'image inverse
du faisceau structural de $\Spec(A)$ par $p : \Spec(B) \to \Spec(A)$
est \`egale au faisceau structural de $\Spec(B)$.
\end{proof}




\section{Alg\`ebres ind-Zariski}
\label{proetale-section-ind-zariski}

\noindent
Commen\c{c}ons par une d\'efinition; voir la remarque \ref{proetale-remark-slightly-stronger}
pour une comparaison avec la d\'efinition correspondante donn\'ee dans l'article \cite{BS}.

\begin{definition}
\label{proetale-definition-ind-zariski}
Un homomorphisme d'anneaux $A \to B$ est dit {\it ind-Zariski} si $B$ peut s'\'ecrire
comme une limite inductive filtrante $B = \colim B_i$, chaque $A \to B_i$ \`etant un
isomorphisme local.
\end{definition}

\noindent
Une localisation $A \to S^{-1}A$ est un exemple d'homomorphisme ind-Zariski; voir
Alg\`ebre, lemme \ref{algebra-lemma-localization-colimit}.
La cat\'egorie des alg\`ebres ind-Zariski est stable par plusieurs op\'erations
naturelles.

\begin{lemma}
\label{proetale-lemma-base-change-ind-zariski}
Soient $A \to B$ et $A \to A'$ des homomorphismes d'anneaux. Notons $B' = B \otimes_A A'$,
l'alg\`ebre obtenue \`a partir de $B$ par changement de base.
Si $A \to B$ est ind-Zariski, alors $A' \to B'$ est ind-Zariski.
\end{lemma}

\begin{proof}
La d\'emonstration est omise.
\end{proof}

\begin{lemma}
\label{proetale-lemma-composition-ind-zariski}
Soient $A \to B$ et $B \to C$ des homomorphismes d'anneaux. Si $A \to B$ et $B \to C$
sont ind-Zariski, alors $A \to C$ est ind-Zariski.
\end{lemma}

\begin{proof}
La d\'emonstration est omise.
\end{proof}

\begin{lemma}
\label{proetale-lemma-ind-zariski-permanence}
Soit $A$ un anneau. Soit $B \to C$ un homomorphisme de $A$-alg\`ebres.
Si $A \to B$ et $A \to C$ sont ind-Zariski, alors $B \to C$
est ind-Zariski.
\end{lemma}

\begin{proof}
La d\'emonstration est omise.
\end{proof}

\begin{lemma}
\label{proetale-lemma-ind-ind-zariski}
Une limite inductive filtrante de $A$-alg\`ebres ind-Zariski est ind-Zariski sur $A$.
\end{lemma}

\begin{proof}
La d\'emonstration est omise.
\end{proof}

\begin{lemma}
\label{proetale-lemma-ind-zariski-implies}
Soit $A \to B$ ind-Zariski. Alors $A \to B$ identifie les anneaux locaux.
\end{lemma}

\begin{proof}
La d\'emonstration est omise.
\end{proof}







\section{Construction de sch\'emas affines w-locaux}
\label{proetale-section-construction}

\noindent
Un sch\'ema affine $X$ est dit {\it w-local} si son espace topologique
sous-jacent est w-local (d\'efinition \ref{proetale-definition-w-local}).
Pour tout anneau $A$, il existe en fait un homomorphisme d'anneaux canonique, ind-Zariski
et fid\`element plat, $A \to A_w$, tel que $\Spec(A_w)$ soit
w-local. La construction de $A_w$ repose sur le lemme simple suivant.

\begin{lemma}
\label{proetale-lemma-localization}
Soit $A$ un anneau. Posons $X = \Spec(A)$. Soit $Z \subset X$ un sous-sch\'ema localement
ferm\'e de la forme $D(f) \cap V(I)$ pour un $f \in A$ et un
id\'eal $I \subset A$. Alors
\begin{enumerate}
\item il existe une partie multiplicative $S \subset A$ telle que
$\Spec(S^{-1}A)$ s'identifie par un hom\'eomorphisme \`a l'ensemble des points de $X$
qui se sp\'ecialisent en un point de $Z$,
\item la $A$-alg\`ebre $A_Z^\sim = S^{-1}A$ ne d\'epend que du
sous-ensemble localement ferm\'e sous-jacent $Z \subset X$,
\item $Z$ est un sous-sch\'ema ferm\'e de $\Spec(A_Z^\sim)$,
\end{enumerate}
Si $A \to A'$ est un homomorphisme d'anneaux et si $Z' \subset X' = \Spec(A')$ est un
sous-sch\'ema localement ferm\'e de la m\^eme forme dont l'image est contenue dans $Z$,
alors il existe un unique homomorphisme de $A$-alg\`ebres
$A_Z^\sim \to (A')_{Z'}^\sim$.
\end{lemma}

\begin{proof}
Soit $S \subset A$ la partie multiplicative form\'ee des \`el\'ements qui deviennent
inversibles dans $\Gamma(Z, \mathcal{O}_Z) = (A/I)_f$.
Si $\mathfrak p$ est un id\'eal premier de $A$ qui ne se sp\'ecialise pas en un point de $Z$,
alors $\mathfrak p$ engendre l'id\'eal unit\'e dans $(A/I)_f$. On peut donc
\'ecrire $f^n =  g + h$ pour un $n \geq 0$, avec $g \in \mathfrak p$ et
$h \in I$. Alors $g \in S$, et $\mathfrak p$ n'appartient pas
au spectre de $S^{-1}A$. R\'eciproquement, si $\mathfrak p$ se sp\'ecialise
en un point de $Z$, disons $\mathfrak p \subset \mathfrak q \supset I$ avec
$f \not \in \mathfrak q$, alors il existe un homomorphisme de $S^{-1}A$ vers
$A_\mathfrak q$; par cons\'equent, $\mathfrak p$ appartient au spectre de $S^{-1}A$.
Ceci prouve (1).

\medskip\noindent
La classe d'isomorphisme de la localisation $S^{-1}A$ ne d\'epend que
du sous-ensemble correspondant $\Spec(S^{-1}A) \subset \Spec(A)$, d'o\`u
(2). Par construction, l'application surjective de $S^{-1}A$ vers
$(A/I)_f$, d'o\`u (3). La derni\`ere assertion r\'esulte de ce que la partie multiplicative
$S' \subset A'$ correspondant \`a $Z'$ contient l'image de la
partie multiplicative $S$.
\end{proof}

\noindent
Soit $A$ un anneau. Soit $E \subset A$ une partie finie. On obtient une
stratification de $X = \Spec(A)$ en sous-sch\'emas localement ferm\'es en
distinguant, pour chaque \`el\'ement de $E$, s'il s'annule ou non. Plus pr\'ecis\'ement,
pour toute d\'ecomposition en union disjointe $E = E' \amalg E''$, posons
\begin{equation}
\label{proetale-equation-stratum}
Z(E', E'') =
\bigcap\nolimits_{f \in E'} D(f) \cap \bigcap\nolimits_{f \in E''} V(f) =
D(\prod\nolimits_{f \in E'} f) \cap V( \sum\nolimits_{f \in E''} fA)
\end{equation}
Les points de $Z(E', E'')$ sont exactement les $x \in X$ tels que
tout $f \in E'$ ait une image non nulle dans $\kappa(x)$ et que tout $f \in E''$
ait une image nulle dans $\kappa(x)$. Il est donc clair que
\begin{equation}
\label{proetale-equation-stratify}
X = \coprod\nolimits_{E = E' \amalg E''} Z(E', E'')
\end{equation}
en tant qu'ensembles. Notons que chaque strate est constructible.

\begin{lemma}
\label{proetale-lemma-refine}
Soit $X = \Spec(A)$ comme ci-dessus. Pour toute stratification finie
$X = \coprod T_i$ par des sous-ensembles constructibles, il existe une partie finie
$E \subset A$ telle que la stratification (\ref{proetale-equation-stratify})
raffine $X = \coprod T_i$.
\end{lemma}

\begin{proof}
On peut \`ecrire $T_i = \bigcup_j U_{i, j} \cap V_{i, j}^c$ comme une union finie,
o\`u les $U_{i, j}$ et $V_{i, j}$ sont des ouverts quasi-compacts de $X$.
Puis on peut \`ecrire $U_{i, j} = \bigcup D(f_{i, j, k})$ et
$V_{i, j} = \bigcup D(g_{i, j, l})$. Posons alors
$E = \{f_{i, j, k}\} \cup \{g_{i, j, l}\}$. Cela convient, car
la stratification (\ref{proetale-equation-stratify}) est celle dont les strates sont
index\'ees par les conditions d'annulation ou de non-annulation des \`el\'ements de $E$; elle
raffine manifestement la stratification donn\'ee.
\end{proof}

\noindent
Poursuivons la discussion.
Pour toute partie finie $E \subset A$, posons
\begin{equation}
\label{proetale-equation-ring}
A_E = \prod\nolimits_{E = E' \amalg E''} A_{Z(E', E'')}^\sim
\end{equation}
avec les notations du lemme \ref{proetale-lemma-localization}. Cette expression est bien d\'efinie, car
(\ref{proetale-equation-stratum}) montre que chaque $Z(E', E'')$ est de la forme voulue.
Prenons le spectre de cet anneau et notons-le
\begin{equation}
\label{proetale-equation-spectrum}
X_E = \Spec(A_E) = \coprod\nolimits_{E = E' \amalg E''} X_{E', E''}
\end{equation}
o\`u $X_{E', E''} = \Spec(A_{Z(E', E'')}^\sim)$. Notons que
\begin{equation}
\label{proetale-equation-closed}
Z_E = \coprod\nolimits_{E = E' \amalg E''} Z(E', E'')
\longrightarrow
X_E
\end{equation}
est un sous-sch\'ema ferm\'e. Par construction, le sous-sch\'ema ferm\'e $Z_E$
contient tous les points ferm\'es du sch\'ema affine $X_E$, puisque tout point
de $X_{E', E''}$ se sp\'ecialise en un point de $Z(E', E'')$.

\medskip\noindent
Soit $I(A)$ l'ensemble ordonn\'e de toutes les parties finies de $A$.
C'est un ensemble ordonn\'e filtrant. Pour $E_1 \subset E_2$, il existe
un homomorphisme de transition canonique $A_{E_1} \to A_{E_2}$ de $A$-alg\`ebres.
Plus pr\'ecis\'ement, pour une d\'ecomposition $E_2 = E'_2 \amalg E''_2$, posons
$E'_1 = E_1 \cap E'_2$ et $E''_1 = E_1 \cap E''_2$. On a alors
$Z(E'_2, E''_2) \subset Z(E'_1, E''_1)$, d'o\`u un unique homomorphisme de $A$-alg\`ebres
$A_{Z(E'_1, E''_1)}^\sim \to A_{Z(E'_2, E''_2)}^\sim$ d'apr\`es le
lemme \ref{proetale-lemma-localization}. En rassemblant ces homomorphismes, on obtient
l'homomorphisme d'anneaux voulu $A_{E_1} \to A_{E_2}$. Notons que le morphisme correspondant
de sch\'emas affines
\begin{equation}
\label{proetale-equation-transition}
X_{E_2} \longrightarrow X_{E_1}
\end{equation}
envoie $Z_{E_2}$ dans $Z_{E_1}$. Par unicit\'e, on obtient un syst\`eme de
$A$-alg\`ebres index\'e par $I(A)$, et l'on pose
\begin{equation}
\label{proetale-equation-colimit-ring}
A_w = \colim_{E \in I(A)} A_E
\end{equation}
Cette $A$-alg\`ebre est ind-Zariski et fid\`element plate sur $A$.
Enfin, posons $X_w = \Spec(A_w)$ et consid\'erons-y le sous-sch\'ema ferm\'e
$Z = \lim_{E \in I(A)} Z_E$. Cela s'\'ecrit
\begin{equation}
\label{proetale-equation-final}
X_w = \lim_{E \in I(A)} X_E \supset Z = \lim_{E \in I(A)} Z_E
\end{equation}

\begin{lemma}
\label{proetale-lemma-make-w-local}
Soit $X = \Spec(A)$ un sch\'ema affine. Consid\'erons $A \to A_w$, $X_w = \Spec(A_w)$,
ainsi que $Z \subset X_w$, d\'efinis comme ci-dessus.
\begin{enumerate}
\item $A \to A_w$ est ind-Zariski et fid\`element plat,
\item $X_w \to X$ induit une bijection $Z \to X$,
\item $Z$ est l'ensemble des points ferm\'es de $X_w$,
\item $Z$ est un sch\'ema r\'eduit, et
\item tout point de $X_w$ se sp\'ecialise en un unique point de $Z$.
\end{enumerate}
En particulier, $X_w$ est w-local (d\'efinition \ref{proetale-definition-w-local}).
\end{lemma}

\begin{proof}
L'homomorphisme $A \to A_w$ est ind-Zariski par construction.
Pour tout $E$, le morphisme $Z_E \to X$ est bijectif, d'o\`u (2).
Comme $Z \subset X_w$, on en conclut que $X_w \to X$ est surjectif et que
$A \to A_w$ est fid\`element plat d'apr\`es
Alg\`ebre, lemme \ref{algebra-lemma-ff-rings}. Ceci prouve (1).

\medskip\noindent
Supposons que $y \in X_w$ et que $y \not \in Z$. Il existe alors
$E$ tel que l'image de $y$ dans $X_E$ n'appartienne pas \`a
$Z_E$. D\`es que $E \subset E'$, l'image de $y$ dans $X_{E'}$
n'appartient donc pas \`a $Z_{E'}$. Soit $T_{E'} \subset X_{E'}$ le sous-sch\'ema
ferm\'e r\'eduit qui est l'adh\'erence de l'image de $y$. Il est
clair que $T = \lim_{E \subset E'} T_{E'}$ est l'adh\'erence de $y$ dans $X_w$.
D\`es que $E \subset E'$, le sch\'ema $T_{E'} \cap Z_{E'}$ est non vide
par construction de $X_{E'}$. Par suite, $\lim T_{E'} \cap Z_{E'}$ est non vide,
et l'on en conclut que $T \cap Z$ est non vide. Ainsi $y$ n'est pas un point ferm\'e.
Il s'ensuit que tout point ferm\'e de $X_w$ appartient \`a $Z$.

\medskip\noindent
Supposons que $y \in X_w$ se sp\'ecialise en $z, z' \in Z$. Montrons que
$z = z'$, ce qui ach\`evera la preuve de (3) et impliquera (5).
Soient $x, x' \in X$ les images de $z$ et $z'$. Puisque $Z \to X$ est
bijective, il suffit de montrer que $x = x'$. Si $x \not = x'$, alors
il existe $f \in A$ tel que $x \in D(f)$ et $x' \in V(f)$
(ou inversement). Posons $E = \{f\}$; alors
$$
X_E = \Spec(A_f) \amalg \Spec(A_{V(f)}^\sim)
$$
On voit que $z$ et $z'$ ont pour images $x_E$ et $x'_E$, qui appartiennent \`a deux
parties diff\'erentes de la d\'ecomposition ci-dessus de $X_E$. Il est alors impossible
que $x_E$ et $x'_E$ soient des sp\'ecialisations d'un m\^eme point.
C'est la contradiction recherch\'ee.

\medskip\noindent
Rappelons que, pour toute partie finie $E \subset A$, $Z_E$
est une union disjointe des sous-sch\'emas localement ferm\'es $Z(E', E'')$,
chacun isomorphe au spectre de $(A/I)_f$, o\`u $I$ est l'id\'eal
engendr\'e par $E''$ et $f$ le produit des \`el\'ements de $E'$.
Tout \`el\'ement nilpotent $b$ de $(A/I)_f$ est la classe de $g/f^n$
pour un certain $g \in A$. En posant alors $E' = E \cup \{g\}$, le lecteur
v\'erifie que l'image inverse de $b$ est nulle par le morphisme de transition
$Z_{E'} \to Z_E$ du syst\`eme. Ceci prouve (4).
\end{proof}

\begin{remark}
\label{proetale-remark-size-w}
Soit $A$ un anneau. Soit $\kappa$ un cardinal infini sup\'erieur ou
\'egal au cardinal de $A$. Alors le cardinal de $A_w$
(lemme \ref{proetale-lemma-make-w-local})
est au plus $\kappa$. En effet, chaque $A_E$ est de cardinal au plus
$\kappa$, et l'ensemble des parties finies de $A$ est lui aussi de cardinal au plus $\kappa$.
Le r\'esultat en d\'ecoule puisque $\kappa \otimes \kappa = \kappa$; voir
Ensembles, section \ref{sets-section-cardinals}.
\end{remark}

\begin{lemma}[Propri\'et\'e universelle de la construction]
\label{proetale-lemma-universal}
Soit $A$ un anneau. Soit $A \to A_w$ l'homomorphisme d'anneaux construit au
lemme \ref{proetale-lemma-make-w-local}. Pour tout homomorphisme d'anneaux $A \to B$ tel que
$\Spec(B)$ soit w-local, il existe une unique factorisation $A \to A_w \to B$
telle que $\Spec(B) \to \Spec(A_w)$ soit w-locale.
\end{lemma}

\begin{proof}
Notons $Y = \Spec(B)$ et $Y_0 \subset Y$ l'ensemble des points ferm\'es.
Notons $f : Y \to X$ le morphisme donn\'e.
Rappelons que $Y_0$ est profini; en particulier, tout sous-ensemble constructible
de $Y_0$ est ouvert et ferm\'e. Soit $E \subset A$ une partie finie.
Rappelons que $A_w = \colim A_E$ et que l'ensemble des points ferm\'es de
$\Spec(A_w)$ est la limite des sous-ensembles ferm\'es $Z_E \subset X_E = \Spec(A_E)$.
Il suffit donc de montrer qu'il existe une unique factorisation $A \to A_E \to B$
telle que $Y \to X_E$ envoie $Y_0$ dans $Z_E$.
Puisque $Z_E \to X = \Spec(A)$ est bijective et que les strates
$Z(E', E'')$ sont constructibles, on voit que
$$
Y_0 = \coprod f^{-1}(Z(E', E'')) \cap Y_0
$$
est une d\'ecomposition en sous-ensembles ouverts et ferm\'es disjoints.
Comme $Y_0 = \pi_0(Y)$, on obtient une d\'ecomposition correspondante de
$Y$ en parties ouvertes et ferm\'ees. Il suffit donc de construire
la factorisation lorsque $f(Y_0) \subset Z(E', E'')$ pour
une d\'ecomposition $E = E' \amalg E''$.
Dans ce cas, $f(Y)$ est contenu dans l'ensemble des points de $X$
qui se sp\'ecialisent en un point de $Z(E', E'')$, ensemble hom\'eomorphe \`a $X_{E', E''}$.
On obtient donc une unique application continue $Y \to X_{E', E''}$ au-dessus de $X$. D'apr\`es le
lemme \ref{proetale-lemma-fully-faithful-spaces-over-X},
elle correspond \`a un unique morphisme de sch\'emas
$Y \to X_{E', E''}$ au-dessus de $X$. Ceci ach\`eve la d\'emonstration.
\end{proof}

\noindent
Rappelons que le spectre d'un anneau est profini si et seulement si
tout point est ferm\'e. Il existe en fait toute une s\'erie de conditions
\'equivalentes qui entra\^inent ce fait. Voir
Alg\`ebre, lemme \ref{algebra-lemma-ring-with-only-minimal-primes}, ou
Topologie, lemme \ref{topology-lemma-characterize-profinite-spectral}.

\begin{lemma}
\label{proetale-lemma-profinite-goes-up}
Soit $A$ un anneau tel que $\Spec(A)$ soit profini. Soit $A \to B$ un
homomorphisme d'anneaux. Alors $\Spec(B)$ est profini dans chacun des cas suivants:
\begin{enumerate}
\item si $\mathfrak q,\mathfrak q' \subset B$ sont au-dessus d'un m\^eme
id\'eal premier de $A$, on n'a ni $\mathfrak q \subset \mathfrak q'$, ni
$\mathfrak q' \subset \mathfrak q$,
\item $A \to B$ induit des extensions alg\'ebriques des corps r\'esiduels,
\item $A \to B$ est un isomorphisme local,
\item $A \to B$ identifie les anneaux locaux,
\item $A \to B$ est faiblement \'etale,
\item $A \to B$ est quasi-fini,
\item $A \to B$ est non ramifi\'e,
\item $A \to B$ est \'etale,
\item $B$ est une limite inductive filtrante de $A$-alg\`ebres satisfaisant \`a l'une des conditions (1) -- (8),
\item etc.
\end{enumerate}
\end{lemma}

\begin{proof}
D'apr\`es les r\'ef\'erences mentionn\'ees ci-dessus
(Alg\`ebre, lemme \ref{algebra-lemma-ring-with-only-minimal-primes}, ou
Topologie, lemme \ref{topology-lemma-characterize-profinite-spectral}),
il n'existe pas de sp\'ecialisation entre deux points distincts de $\Spec(A)$, et
$\Spec(B)$ est profini si et seulement s'il n'existe pas de sp\'ecialisation
entre deux points distincts de $\Spec(B)$. De telles sp\'ecialisations ne peuvent
avoir lieu que dans les fibres de $\Spec(B) \to \Spec(A)$. On voit ainsi
que (1) est vraie.

\medskip\noindent
L'hypoth\`ese de (2) implique que tous les id\'eaux premiers de $B$ sont maximaux, d'apr\`es
Alg\`ebre, lemme \ref{algebra-lemma-finite-residue-extension-closed}.
Ainsi (2) est vraie.
Si $A \to B$ est un isomorphisme local ou identifie les anneaux locaux,
alors les extensions des corps r\'esiduels sont triviales; (3) et (4)
r\'esultent donc de (2).
Si $A \to B$ est faiblement \'etale, alors Compl\'ements sur l'alg\`ebre, lemme
\ref{more-algebra-lemma-weakly-etale-residue-field-extensions}
montre qu'il induit des extensions alg\'ebriques s\'eparables des corps r\'esiduels, donc
(5) r\'esulte de (2).
Si $A \to B$ est quasi-fini, alors les fibres sont des espaces topologiques
finis discrets. Ainsi (6) r\'esulte de (1).
Ainsi (3) r\'esulte de (1). Les cas (7) et (8)
en d\'ecoulent, car les homomorphismes d'anneaux non ramifi\'es et \'etales sont quasi-finis
(Alg\`ebre, lemmes
\ref{algebra-lemma-unramified-quasi-finite} et
\ref{algebra-lemma-etale-quasi-finite}).
Si $B = \colim B_i$ est une limite inductive filtrante de $A$-alg\`ebres, alors
$\Spec(B) = \lim \Spec(B_i)$ dans la cat\'egorie des espaces topologiques d'apr\`es
Limites, lemme \ref{limits-lemma-inverse-limit-top}.
Ainsi, si chaque $\Spec(B_i)$ est profini, $\Spec(B)$ l'est aussi d'apr\`es
Topologie, lemme \ref{topology-lemma-directed-inverse-limit-profinite}.
Ceci prouve (9).
\end{proof}

\begin{lemma}
\label{proetale-lemma-localize-along-closed-profinite}
Soit $A$ un anneau. Soit $V(I) \subset \Spec(A)$ un sous-ensemble ferm\'e
qui est un espace topologique profini. Il existe alors un
homomorphisme d'anneaux ind-Zariski $A \to B$ tel que $\Spec(B)$ soit w-local,
que l'ensemble de ses points ferm\'es soit $V(IB)$ et que $A/I \cong B/IB$.
\end{lemma}

\begin{proof}
Soient $A \to A_w$ et $Z \subset Y = \Spec(A_w)$ comme au
lemme \ref{proetale-lemma-make-w-local}.
Soit $T \subset Z$ l'image r\'eciproque de $V(I)$.
Alors $T \to V(I)$ est un hom\'eomorphisme d'apr\`es
Topologie, lemme \ref{topology-lemma-bijective-map}.
Posons $B = (A_w)_T^\sim$; voir le lemme \ref{proetale-lemma-localization}.
Il est clair que $B$ est w-local et a pour ensemble de points ferm\'es $V(IB)$.
L'homomorphisme d'anneaux $A/I \to B/IB$ est ind-Zariski
et induit un hom\'eomorphisme entre les espaces topologiques
sous-jacents. C'est donc un isomorphisme d'apr\`es le
lemme \ref{proetale-lemma-local-isomorphism-fully-faithful}.
\end{proof}

\begin{lemma}
\label{proetale-lemma-w-local-algebraic-residue-field-extensions}
Soit $A$ un anneau tel que $X = \Spec(A)$ soit w-local. Soit $I \subset A$
l'id\'eal radical qui d\'efinit l'ensemble $X_0$ des points ferm\'es de $X$.
Soit $A \to B$ un homomorphisme d'anneaux qui induit des extensions alg\'ebriques
des corps r\'esiduels en chaque id\'eal premier de l'anneau d'arriv\'ee. Alors
\begin{enumerate}
\item tout point de $Z = V(IB)$ est un point ferm\'e de $\Spec(B)$,
\item il existe un homomorphisme d'anneaux ind-Zariski $B \to C$ tel que
\begin{enumerate}
\item $B/IB \to C/IC$ soit un isomorphisme,
\item l'espace $Y = \Spec(C)$ soit w-local,
\item l'application induite $p : Y \to X$ soit w-locale, et
\item $p^{-1}(X_0)$ soit l'ensemble des points ferm\'es de $Y$.
\end{enumerate}
\end{enumerate}
\end{lemma}

\begin{proof}
D'apr\`es le lemme \ref{proetale-lemma-profinite-goes-up} appliqu\'e \`a $A/I \to B/IB$,
tous les points de $Z = V(IB) = \Spec(B/IB)$ sont ferm\'es; en fait, $\Spec(B/IB)$
est un espace profini.
Pour achever la d\'emonstration, appliquons le lemme \ref{proetale-lemma-localize-along-closed-profinite}
\`a $IB \subset B$.
\end{proof}





\section{Identification des anneaux locaux et condition ind-Zariski}
\label{proetale-section-connected-components}

\noindent
Un homomorphisme d'anneaux ind-Zariski $A \to B$ identifie les anneaux locaux
(lemme \ref{proetale-lemma-ind-zariski-implies}). La r\'eciproque est fausse
(Exemples, section \ref{examples-section-not-ind-etale}).
Il existe toutefois une sorte de th\'eor\`eme de structure qui d\'ecrit les
homomorphismes d'anneaux identifiant les anneaux locaux au moyen d'homomorphismes
ind-Zariski; voir la proposition \ref{proetale-proposition-maps-wich-identify-local-rings}.

\medskip\noindent
Soit $A$ un anneau et posons $X = \Spec(A)$. L'espace des composantes
connexes $\pi_0(X)$ est un espace profini d'apr\`es
Topologie, lemme \ref{topology-lemma-spectral-pi0}
(et Alg\`ebre, lemme \ref{algebra-lemma-spec-spectral}).

\begin{lemma}
\label{proetale-lemma-construct}
Soit $A$ un anneau et posons $X = \Spec(A)$. Soit $T \subset \pi_0(X)$ un
sous-ensemble ferm\'e. Il existe un homomorphisme d'anneaux surjectif et ind-Zariski $A \to B$
tel que $\Spec(B) \to \Spec(A)$ induise un hom\'eomorphisme de $\Spec(B)$
sur l'image r\'eciproque de $T$ dans $X$.
\end{lemma}

\begin{proof}
Soit $Z \subset X$ l'image r\'eciproque de $T$. Alors $Z$ est l'intersection
$Z = \bigcap Z_\alpha$ des sous-ensembles ouverts et ferm\'es de $X$ qui contiennent $Z$;
voir Topologie, lemme \ref{topology-lemma-closed-union-connected-components}.
Pour tout $\alpha$, on a $Z_\alpha = \Spec(A_\alpha)$, o\`u
$A \to A_\alpha$ est un isomorphisme local (une localisation en un idempotent).
En posant $B = \colim A_\alpha$, on d\'emontre le lemme.
\end{proof}

\begin{lemma}
\label{proetale-lemma-construct-profinite}
Soit $A$ un anneau et posons $X = \Spec(A)$. Soit $T$ un espace profini et
soit $T \to \pi_0(X)$ une application continue. Il existe un
homomorphisme d'anneaux ind-Zariski $A \to B$ tel que, en posant $Y = \Spec(B)$, le diagramme
$$
\xymatrix{
Y \ar[r] \ar[d] & \pi_0(Y) \ar[d] \\
X \ar[r] & \pi_0(X)
}
$$
soit cart\'esien dans la cat\'egorie des espaces topologiques et que
$\pi_0(Y) = T$ comme espaces au-dessus de $\pi_0(X)$.
\end{lemma}

\begin{proof}
\'Ecrivons $T = \lim T_i$ comme la limite d'un syst\`eme projectif d'espaces
finis discrets index\'e par un ensemble ordonn\'e filtrant (voir
Topologie, lemme \ref{topology-lemma-profinite}). Pour tout $i$, posons
$Z_i = \Im(T \to \pi_0(X) \times T_i)$. C'est un sous-ensemble ferm\'e.
Notons que $X \times T_i$ est le spectre de $A_i = \prod_{t \in T_i} A$
et que $A \to A_i$ est un isomorphisme local. D'apr\`es le lemme \ref{proetale-lemma-construct},
le sous-ensemble $Z_i \subset \pi_0(X \times T_i) = \pi_0(X) \times T_i$
correspond \`a un homomorphisme d'anneaux surjectif et ind-Zariski $A_i \to B_i$
tel que $\Spec(B_i) = X \times_{\pi_0(X)} Z_i$ comme sous-ensemble de
$X \times T_i$. Les applications de transition $T_i \to T_{i'}$ induisent des applications
$Z_i \to Z_{i'}$ et $X \times_{\pi_0(X)} Z_i \to X \times_{\pi_0(X)} Z_{i'}$,
donc des homomorphismes d'anneaux $B_{i'} \to B_i$
(lemmes \ref{proetale-lemma-local-isomorphism-fully-faithful} et
\ref{proetale-lemma-ind-zariski-implies}).
Posons $B = \colim B_i$. Puisque $T = \lim Z_i$, on a
$X \times_{\pi_0(X)} T = \lim  X \times_{\pi_0(X)} Z_i$
et donc $Y = \Spec(B) = \lim \Spec(B_i)$
s'ins\`ere dans le diagramme cart\'esien
$$
\xymatrix{
Y \ar[r] \ar[d] & T \ar[d] \\
X \ar[r] & \pi_0(X)
}
$$
d'espaces topologiques. D'apr\`es le lemme \ref{proetale-lemma-silly},
on en conclut que $T = \pi_0(Y)$.
\end{proof}

\begin{example}
\label{proetale-example-construct-space}
Soit $k$ un corps et soit $T$ un espace topologique profini.
Il existe un homomorphisme d'anneaux ind-Zariski $k \to A$ tel que
$\Spec(A)$ soit hom\'eomorphe \`a $T$. Il suffit en effet d'appliquer le
lemme \ref{proetale-lemma-construct-profinite} \`a $T \to \pi_0(\Spec(k)) = \{*\}$.
Dans ce cas, on a m\^eme
$$
A = \colim \text{Map}(T_i, k)
$$
pour toute \`ecriture $T = \lim T_i$ comme limite projective filtrante o\`u chaque $T_i$ est fini.
\end{example}

\begin{lemma}
\label{proetale-lemma-w-local-morphism-equal-points-stalks-is-iso}
Soit $A \to B$ un homomorphisme d'anneaux tel que
\begin{enumerate}
\item $A \to B$ identifie les anneaux locaux,
\item les espaces topologiques $\Spec(B)$ et $\Spec(A)$ soient w-locaux,
\item $\Spec(B) \to \Spec(A)$ soit w-locale, et
\item $\pi_0(\Spec(B)) \to \pi_0(\Spec(A))$ soit bijective.
\end{enumerate}
Alors $A \to B$ est un isomorphisme.
\end{lemma}

\begin{proof}
Soient $X_0 \subset X = \Spec(A)$ et $Y_0 \subset Y = \Spec(B)$ les
ensembles de points ferm\'es. Par hypoth\`ese, l'image de $Y_0$ est contenue dans $X_0$ et
l'application induite $Y_0 \to X_0$ est bijective.
Topologiquement, $\Spec(A)$ est l'union disjointe des spectres
des anneaux locaux de $A$ en ses points ferm\'es.
Il en va de m\^eme pour $B$. Ainsi $X \to Y$ est bijective.
Puisque $A \to B$ est plat, le th\'eor\`eme de descente des id\'eaux premiers s'applique
(Alg\`ebre, lemme \ref{algebra-lemma-flat-going-down}).
Ainsi, le lemme \ref{algebra-lemma-unique-prime-over-localize-below} d'Alg\`ebre
montre que, pour tout id\'eal premier $\mathfrak q \subset B$ au-dessus de
$\mathfrak p \subset A$, on a $B_\mathfrak q = B_\mathfrak p$.
Comme $B_\mathfrak q = A_\mathfrak p$ par hypoth\`ese, on
obtient $A_\mathfrak p = B_\mathfrak p$ pour tout id\'eal premier $\mathfrak p$
de $A$. Par cons\'equent, $A = B$ d'apr\`es
Alg\`ebre, lemme \ref{algebra-lemma-characterize-zero-local}.
\end{proof}

\begin{lemma}
\label{proetale-lemma-w-local-morphism-equal-stalks-is-ind-zariski}
Soit $A \to B$ un homomorphisme d'anneaux tel que
\begin{enumerate}
\item $A \to B$ identifie les anneaux locaux,
\item les espaces topologiques $\Spec(B)$ et $\Spec(A)$ soient w-locaux, et
\item $\Spec(B) \to \Spec(A)$ soit w-locale.
\end{enumerate}
Alors $A \to B$ est ind-Zariski.
\end{lemma}

\begin{proof}
Posons $X = \Spec(A)$ et $Y = \Spec(B)$. Soient $X_0 \subset X$ et
$Y_0 \subset Y$ les ensembles de points ferm\'es. Soit $A \to A'$ l'homomorphisme d'anneaux ind-Zariski
tel que, en posant $X' = \Spec(A')$, le diagramme
$$
\xymatrix{
X' \ar[r] \ar[d] & \pi_0(X') \ar[d] \\
X \ar[r] & \pi_0(X)
}
$$
soit cart\'esien dans la cat\'egorie des espaces topologiques et que
$\pi_0(X') = \pi_0(Y)$ comme espaces au-dessus de $\pi_0(X)$; voir le
lemme \ref{proetale-lemma-construct-profinite}. D'apr\`es le
lemme \ref{proetale-lemma-silly}, $X'$ est w-local et
l'ensemble $X'_0 \subset X'$ de ses points ferm\'es est l'image r\'eciproque de $X_0$.

\medskip\noindent
On obtient une application continue $Y \to X'$ entre les espaces topologiques sous-jacents
au-dessus de $X$, laquelle induit l'identification de $\pi_0(Y)$ avec $\pi_0(X')$. D'apr\`es le
lemme \ref{proetale-lemma-local-isomorphism-fully-faithful}
(et le lemme \ref{proetale-lemma-ind-zariski-implies}),
elle correspond \`a un morphisme de sch\'emas affines $Y \to X'$
au-dessus de $X$. Puisque $Y \to X$ envoie $Y_0$ dans $X_0$, on voit que
$Y \to X'$ envoie $Y_0$ dans $X'_0$, c'est-à-dire que $Y \to X'$ est w-locale.
D'après le lemme \ref{proetale-lemma-w-local-morphism-equal-points-stalks-is-iso},
nous voyons que $Y \cong X'$ et le résultat s'ensuit.
\end{proof}

\noindent
La proposition suivante est un résultat préparatoire du type de ceux
que nous démontrerons plus loin.

\begin{proposition}
\label{proetale-proposition-maps-wich-identify-local-rings}
Soit $A \to B$ un homomorphisme d'anneaux qui identifie les anneaux locaux.
Alors il existe un homomorphisme d'anneaux fidèlement plat et ind-Zariski
$B \to B'$ tel que $A \to B'$ soit ind-Zariski.
\end{proposition}

\begin{proof}
Soient $A \to A_w$, resp. $B \to B_w$, les homomorphismes d'anneaux fidèlement plats et ind-Zariski
construits dans le lemme \ref{proetale-lemma-make-w-local} pour $A$, resp.\ $B$.
Puisque $\Spec(B_w)$ est w-local, il existe une unique factorisation
$A \to A_w \to B_w$ telle que $\Spec(B_w) \to \Spec(A_w)$ soit w-locale
d'après le lemme \ref{proetale-lemma-universal}. Remarquons que $A_w \to B_w$ identifie
les anneaux locaux; voir
le lemme \ref{proetale-lemma-local-isomorphism-permanence}.
D'après le lemme \ref{proetale-lemma-w-local-morphism-equal-stalks-is-ind-zariski},
cela signifie que $A_w \to B_w$ est ind-Zariski. Puisque $B \to B_w$ est
fidèlement plat et ind-Zariski (lemme \ref{proetale-lemma-make-w-local})
et que la composée $A \to B \to B_w$ est ind-Zariski
(lemme \ref{proetale-lemma-composition-ind-zariski}),
la proposition est démontrée.
\end{proof}

\noindent
La proposition ci-dessus nous permet de caractériser les objets affines
faiblement contractiles du site pro-Zariski d'un schéma affine.

\begin{lemma}
\label{proetale-lemma-w-local-extremally-disconnected}
Soit $A$ un anneau. Les assertions suivantes sont équivalentes :
\begin{enumerate}
\item tout homomorphisme d'anneaux fidèlement plat $A \to B$ qui identifie les anneaux locaux
admet une rétraction,
\item tout homomorphisme d'anneaux fidèlement plat et ind-Zariski $A \to B$ admet une rétraction, et
\item $A$ satisfait aux conditions suivantes :
\begin{enumerate}
\item $\Spec(A)$ est w-local, et
\item $\pi_0(\Spec(A))$ est extrêmement discontinu.
\end{enumerate}
\end{enumerate}
\end{lemma}

\begin{proof}
L'équivalence de (1) et (2) découle immédiatement de la
proposition \ref{proetale-proposition-maps-wich-identify-local-rings}.

\medskip\noindent
Supposons (3)(a) et (3)(b). Soit $A \to B$ un homomorphisme d'anneaux fidèlement plat et ind-Zariski.
Nous utiliserons, sans plus le mentionner, le fait qu'un homomorphisme plat
$A \to B$ est fidèlement plat si et seulement si tout point fermé
de $\Spec(A)$ appartient à l'image de $\Spec(B) \to \Spec(A)$.
Nous allons montrer que $A \to B$ admet une rétraction.

\medskip\noindent
Soit $I \subset A$ un idéal tel que $V(I) \subset \Spec(A)$ soit
l'ensemble des points fermés de $\Spec(A)$.
Nous pouvons remplacer $B$ par l'anneau $C$ construit dans le
lemme \ref{proetale-lemma-w-local-algebraic-residue-field-extensions}
pour $A \to B$ et $I \subset A$.
Ainsi, nous pouvons supposer que $\Spec(B)$ est w-local et que l'ensemble des
points fermés de $\Spec(B)$ est $V(IB)$.

\medskip\noindent
Supposons que $\Spec(B)$ soit w-local et que l'ensemble des points fermés de $\Spec(B)$
soit $V(IB)$. Choisissons une section continue de l'application
continue surjective $V(IB) \to V(I)$. Cela est possible puisque
$V(I) \cong \pi_0(\Spec(A))$ est extrêmement discontinu ; voir la
proposition de Topologie
\ref{topology-proposition-projective-in-category-hausdorff-qc}.
Son image est un sous-espace fermé $T \subset \pi_0(\Spec(B)) \cong V(IB)$
qui est envoyé homéomorphiquement sur $\pi_0(A)$. En remplaçant $B$ par l'anneau quotient ind-Zariski
construit dans le lemme \ref{proetale-lemma-construct},
nous voyons que nous pouvons supposer que $\pi_0(\Spec(B)) \to \pi_0(\Spec(A))$
est bijective. À ce stade, $A \to B$ est un isomorphisme d'après le
lemme \ref{proetale-lemma-w-local-morphism-equal-points-stalks-is-iso}.

\medskip\noindent
Supposons (1), ou de manière équivalente (2). Soit $A \to A_w$ l'homomorphisme d'anneaux construit dans le
lemme \ref{proetale-lemma-make-w-local}. D'après (1), il existe une rétraction $A_w \to A$.
Ainsi, $\Spec(A)$ est homéomorphe à un sous-ensemble fermé de $\Spec(A_w)$. Par le
lemme \ref{proetale-lemma-closed-subspace-w-local}, nous voyons que (3)(a) est vérifiée.
Enfin, soit $T \to \pi_0(A)$ une application surjective, où $T$ est un espace topologique
extrêmement discontinu, quasi-compact et séparé
(Topologie, lemme \ref{topology-lemma-existence-projective-cover}).
Choisissons $A \to B$ comme dans le lemme \ref{proetale-lemma-construct-profinite},
adapté à $T \to \pi_0(\Spec(A))$. D'après (1), il existe une rétraction
$B \to A$. Ainsi, nous voyons que $T = \pi_0(\Spec(B)) \to \pi_0(\Spec(A))$
admet une section. Un argument formel de théorie des catégories, utilisant la
proposition de Topologie
\ref{topology-proposition-projective-in-category-hausdorff-qc},
implique que $\pi_0(\Spec(A))$ est extrêmement discontinu.
\end{proof}

\begin{lemma}
\label{proetale-lemma-find-Zariski-w-contractible}
Soit $A$ un anneau. Il existe un homomorphisme d'anneaux fidèlement plat et ind-Zariski
$A \to B$ tel que $B$ satisfasse aux conditions équivalentes
du lemme \ref{proetale-lemma-w-local-extremally-disconnected}.
\end{lemma}

\begin{proof}
Nous appliquons d'abord le lemme \ref{proetale-lemma-make-w-local} pour voir que nous pouvons supposer que
$\Spec(A)$ est w-local.
Choisissons un espace extrêmement discontinu $T$ et une application continue surjective
$T \to \pi_0(\Spec(A))$, voir le
lemme de Topologie \ref{topology-lemma-existence-projective-cover}.
Remarquons que $T$ est profini. Appliquons le lemme \ref{proetale-lemma-construct-profinite}
pour trouver un homomorphisme d'anneaux ind-Zariski $A \to B$ tel que
$\pi_0(\Spec(B)) \to \pi_0(\Spec(A))$ réalise $T \to \pi_0(\Spec(A))$
et tel que
$$
\xymatrix{
\Spec(B) \ar[r] \ar[d] & \pi_0(\Spec(B)) \ar[d] \\
\Spec(A) \ar[r] & \pi_0(\Spec(A))
}
$$
est cartésien dans la catégorie des espaces topologiques. Remarquons que $\Spec(B)$
est w-local, que $\Spec(B) \to \Spec(A)$ est w-locale et que
l'ensemble des points fermés de $\Spec(B)$ est l'image réciproque de
l'ensemble des points fermés de $\Spec(A)$, voir le lemme \ref{proetale-lemma-silly}.
Ainsi, la condition (3) du
lemme \ref{proetale-lemma-w-local-extremally-disconnected}
est satisfaite pour $B$.
\end{proof}

\begin{remark}
\label{proetale-remark-slightly-stronger}
Dans les lemmes \ref{proetale-lemma-construct}, \ref{proetale-lemma-construct-profinite},
dans la proposition \ref{proetale-proposition-maps-wich-identify-local-rings} et dans le
lemme \ref{proetale-lemma-find-Zariski-w-contractible}, nous construisons un homomorphisme d'anneaux ind-Zariski
possédant certaines propriétés. Dans l'article \cite{BS}, les auteurs utilisent la notion
d'ind-(localisation de Zariski), c'est-à-dire de limite inductive filtrante de produits finis
de localisations principales. Dans chacun des résultats énumérés ci-dessus, on peut remplacer ind-Zariski
par ind-(localisation de Zariski).
Nous n'en avons toutefois pas besoin, et la notion d'homomorphisme ind-Zariski
d'anneaux définie ici possède des propriétés formelles légèrement meilleures. De plus,
la notion d'homomorphisme d'anneaux ind-Zariski est l'analogue naturel de la
notion d'homomorphisme d'anneaux ind-\'etale définie à la section suivante.
\end{remark}




\section{Algèbre ind-\'etale}
\label{proetale-section-ind-etale}

\noindent
Commençons par une définition.

\begin{definition}
\label{proetale-definition-ind-etale}
Un homomorphisme d'anneaux $A \to B$ est dit {\it ind-\'etale} si $B$ peut s'écrire
comme une limite inductive filtrante d'algèbres \'etales sur $A$.
\end{definition}

\noindent
La catégorie des algèbres ind-\'etales est stable par un certain nombre d'opérations
naturelles.

\begin{lemma}
\label{proetale-lemma-base-change-ind-etale}
Soient $A \to B$ et $A \to A'$ des homomorphismes d'anneaux. Soit $B' = B \otimes_A A'$
l'algèbre obtenue à partir de $B$ par changement de base.
Si $A \to B$ est ind-\'etale, alors $A' \to B'$ est ind-\'etale.
\end{lemma}

\begin{proof}
C'est le lemme d'Algèbre \ref{algebra-lemma-base-change-colimit-etale}.
\end{proof}

\begin{lemma}
\label{proetale-lemma-composition-ind-etale}
Soient $A \to B$ et $B \to C$ des homomorphismes d'anneaux. Si $A \to B$ et $B \to C$
sont ind-\'etales, alors $A \to C$ est ind-\'etale.
\end{lemma}

\begin{proof}
C'est le lemme d'Algèbre \ref{algebra-lemma-composition-colimit-etale}.
\end{proof}

\begin{lemma}
\label{proetale-lemma-ind-ind-etale}
Une limite inductive filtrante d'algèbres ind-\'etales sur $A$ est ind-\'etale sur $A$.
\end{lemma}

\begin{proof}
C'est le lemme d'Algèbre \ref{algebra-lemma-colimit-colimit-etale}.
\end{proof}

\begin{lemma}
\label{proetale-lemma-ind-etale-permanence}
Soit $A$ un anneau. Soit $B \to C$ un homomorphisme entre des $A$-algèbres ind-\'etales
sur $A$. Alors $C$ est une $B$-algèbre ind-\'etale.
\end{lemma}

\begin{proof}
C'est le lemme d'Algèbre \ref{algebra-lemma-colimits-of-etale}.
\end{proof}

\begin{lemma}
\label{proetale-lemma-ind-etale-implies}
Soit $A \to B$ ind-\'etale. Alors $A \to B$ est faiblement \'etale
(Compléments d'Algèbre, définition \ref{more-algebra-definition-weakly-etale}).
\end{lemma}

\begin{proof}
Cela résulte du lemme de Compléments d'Algèbre
\ref{more-algebra-lemma-when-weakly-etale}.
\end{proof}

\begin{lemma}
\label{proetale-lemma-lift-ind-etale}
Soient $A$ un anneau et $I \subset A$ un idéal. Le foncteur de changement de base
$$
\text{algèbres ind-\'etales sur }A\text{}
\longrightarrow
\text{algèbres ind-\'etales sur }A/I\text{},\quad
C \longmapsto C/IC
$$
admet un adjoint à droite pleinement fidèle $v$. En particulier, étant donnée
une $A/I$-algèbre ind-\'etale $\overline{C}$, il existe
une $A$-algèbre ind-\'etale $C = v(\overline{C})$ telle que
$\overline{C} = C/IC$.
\end{lemma}

\begin{proof}
Soit $\overline{C}$ une $A/I$-algèbre ind-\'etale.
Considérons la catégorie $\mathcal{C}$ des factorisations
$A \to B \to \overline{C}$ où $A \to B$ est \'etale.
(Nous laissons de côté certaines questions ensemblistes dans cette démonstration.)
Nous allons montrer que cette catégorie est filtrante et que
$C = \colim_\mathcal{C} B$ est une $A$-algèbre ind-\'etale
telle que $\overline{C} = C/IC$.

\medskip\noindent
Montrons d'abord que $\mathcal{C}$ est filtrante
(Catégories, définition \ref{categories-definition-directed}).
La catégorie est non vide, puisque $A \to A \to \overline{C}$ en est un objet.
Supposons que $A \to B \to \overline{C}$ et $A \to B' \to \overline{C}$
soient deux objets de $\mathcal{C}$. Alors $A \to B \otimes_A B' \to \overline{C}$
en est un autre (utiliser le lemme d'Algèbre \ref{algebra-lemma-etale}).
Supposons que $f, g : B \to B'$ soient deux morphismes entre les
objets $A \to B \to \overline{C}$ et $A \to B' \to \overline{C}$
de $\mathcal{C}$. Un coégalisateur est alors
$A \to B' \otimes_{f, B, g} B' \to \overline{C}$.
C'est un objet de $\mathcal{C}$ d'après les
lemmes d'Algèbre \ref{algebra-lemma-etale} et
\ref{algebra-lemma-map-between-etale}.
Ainsi, la catégorie $\mathcal{C}$ est filtrante.

\medskip\noindent
Écrivons $\overline{C} = \colim \overline{B_i}$ comme limite inductive filtrante, avec
$\overline{B_i}$ \'etale sur $A/I$. Pour tout $i$,
il existe un homomorphisme $A \to B_i$ qui est \'etale, tel que $\overline{B_i} = B_i/IB_i$, voir le
lemme d'Algèbre \ref{algebra-lemma-lift-etale}.
Ainsi, $C \to \overline{C}$ est surjectif.
Comme $C/IC \to \overline{C}$ est ind-\'etale
(lemme \ref{proetale-lemma-ind-etale-permanence}),
il est plat. Par conséquent, $\overline{C}$ est une localisation de
$C/IC$ par une partie multiplicative $S \subset C/IC$
(Algèbre, lemme \ref{algebra-lemma-pure}).
Soit $f \in C$ dont l'image appartient à $S \subset C/IC$.
Choisissons $A \to B \to \overline{C}$ dans $\mathcal{C}$ et $g \in B$
dont l'image dans la limite inductive est $f$. Alors $A \to B_g \to \overline{C}$
est encore un objet de $\mathcal{C}$. Ainsi, $f$ est un élément inversible
de $C$. Il s'ensuit que $C/IC = \overline{C}$.

\medskip\noindent
Montrons ensuite que, pour toute algèbre ind-\'etale $D$ sur $A$, on a
$$
\Mor_A(D, C) = \Mor_{A/I}(D/ID, \overline{C})
$$
En effet, soit $D/ID \to \overline{C}$ un homomorphisme de $A/I$-algèbres.
Écrivons $D = \colim_{i \in I} D_i$ comme limite inductive indexée par un ensemble préordonné filtrant $I$,
avec $D_i$ \'etale sur $A$. Par le choix de $\mathcal{C}$,
nous obtenons un foncteur $I \to \mathcal{C}$, et donc un homomorphisme
$D \to C$ compatible avec les homomorphismes vers $\overline{C}$. D'où l'assertion.

\medskip\noindent
Il s'ensuit que le foncteur $v$ défini par la règle
$$
\overline{C} 
\longmapsto
v(\overline{C}) = \colim_{A \to B \to \overline{C}} B
$$
est un adjoint à droite du foncteur de changement de base $u$, comme l'affirme le lemme.
Le foncteur $v$ est pleinement fidèle car
$u \circ v = \text{id}$ par construction, voir le
lemme de Catégories \ref{categories-lemma-adjoint-fully-faithful}.
\end{proof}









\section{Construction d'algèbres ind-\'etales}
\label{proetale-section-construction-ind-etale}

\noindent
Soit $A$ un anneau. Rappelons que tout homomorphisme d'anneaux \'etale $A \to B$ est isomorphe
à un homomorphisme lisse standard de dimension relative $0$. Un tel homomorphisme
est de la forme
$$
A \longrightarrow A[x_1, \ldots, x_n]/(f_1, \ldots, f_n)
$$
où le déterminant de la matrice $n \times n$ dont les coefficients sont les
$\partial f_i/\partial x_j$ est inversible dans l'anneau quotient. Voir le
lemme d'Algèbre \ref{algebra-lemma-etale-standard-smooth}.

\medskip\noindent
Soit $S(A)$ l'ensemble de toutes les algèbres {\it fidèlement plates}\footnote{En présence
de platitude, par exemple pour des homomorphismes d'anneaux lisses ou \'etales,
cela signifie simplement que l'application induite sur les spectres est surjective. Voir le
lemme d'Algèbre \ref{algebra-lemma-ff-rings}.}
lisses standard sur $A$ et de dimension relative $0$.
Soit $I(A)$ l'ensemble, partiellement ordonné par inclusion, des sous-ensembles finis
$E$ de $S(A)$. Remarquons que $I(A)$ est un ensemble partiellement ordonné
filtrant. Pour $E = \{A \to B_1, \ldots, A \to B_n\}$, posons
$$
B_E = B_1 \otimes_A \ldots \otimes_A B_n
$$
Remarquons que $B_E$ est une $A$-algèbre \'etale fidèlement plate.
Pour $E \subset E'$, il existe un homomorphisme de transition canonique $B_E \to B_{E'}$
de $A$-algèbres \'etales. En effet, supposons $E = \{A \to B_1, \ldots, A \to B_n\}$
et $E' = \{A \to B_1, \ldots, A \to B_{n + m}\}$; alors
$B_E \to B_{E'}$ envoie $b_1 \otimes \ldots \otimes b_n$ sur
l'élément $b_1 \otimes \ldots \otimes b_n \otimes 1 \otimes \ldots \otimes 1$
de $B_{E'}$. Cette construction définit un système d'algèbres fidèlement plates
\'etales sur $A$, indexé par $I(A)$, et nous posons
$$
T(A) = \colim_{E \in I(A)} B_E
$$
Remarquons que $T(A)$ est une $A$-algèbre ind-\'etale fidèlement plate
(Algèbre, lemme \ref{algebra-lemma-colimit-faithfully-flat}). Par construction,
pour toute $A$-algèbre \'etale fidèlement plate $B$, il existe un homomorphisme, non unique,
de $A$-algèbres $B \to T(A)$. En effet, choisissons $(A \to B_0) \in S(A)$
et un isomorphisme $B \cong B_0$. Alors la coprojection canonique
$$
B \to B_0 \to 
T(A) = \colim_{E \in I(A)} B_E
$$
est le morphisme recherché.

\begin{lemma}
\label{proetale-lemma-first-construction}
Pour tout anneau $A$, il existe une $A$-algèbre ind-\'etale et fidèlement plate $C$
telle que tout homomorphisme d'anneaux \'etale et fidèlement plat $C \to B$ admette une rétraction.
\end{lemma}

\begin{proof}
Posons $T^1(A) = T(A)$ et $T^{n + 1}(A) = T(T^n(A))$. Soit
$$
C = \colim T^n(A)
$$
Cette algèbre est fidèlement plate sur chaque $T^n(A)$, et en particulier
sur $A$, voir le
lemme d'Algèbre \ref{algebra-lemma-colimit-faithfully-flat}.
De plus, $C$ est ind-\'etale sur $A$ d'après le lemme \ref{proetale-lemma-ind-ind-etale}.
Si $C \to B$ est \'etale, il existe alors un entier $n$ et un homomorphisme d'anneaux \'etale
$T^n(A) \to B'$ tels que $B = C \otimes_{T^n(A)} B'$, voir le
lemme d'Algèbre \ref{algebra-lemma-etale}.
Si $C \to B$ est fidèlement plat, alors $\Spec(B) \to \Spec(C) \to \Spec(T^n(A))$
est surjectif; par conséquent, $\Spec(B') \to \Spec(T^n(A))$ est surjectif.
Autrement dit, $T^n(A) \to B'$ est fidèlement plat.
Par construction, il existe un homomorphisme de $T^n(A)$-algèbres
$B' \to T^{n + 1}(A)$. Il induit un homomorphisme de $C$-algèbres $B \to C$,
ce qui achève la démonstration.
\end{proof}

\begin{remark}
\label{proetale-remark-size-T}
Soit $A$ un anneau. Soit $\kappa$ un cardinal infini supérieur ou
égal au cardinal de $A$. Alors le cardinal de $T(A)$
est au plus $\kappa$. En effet, chaque $B_E$ est de cardinal au plus
$\kappa$, et l'ensemble d'indices $I(A)$ est lui aussi de cardinal au plus $\kappa$.
Le résultat s'ensuit puisque $\kappa \otimes \kappa = \kappa$, voir la
section d'Ensembles \ref{sets-section-cardinals}. Il s'ensuit également que
l'anneau construit dans la démonstration du lemme \ref{proetale-lemma-first-construction}
est de cardinal au plus $\kappa$.
\end{remark}

\begin{remark}
\label{proetale-remark-first-construction-functorial}
La construction $A \mapsto T(A)$ est fonctorielle au sens suivant :
si $A \to A'$ est un homomorphisme d'anneaux, nous pouvons construire un diagramme commutatif
$$
\xymatrix{
A \ar[r] \ar[d] & T(A) \ar[d] \\
A' \ar[r] & T(A')
}
$$
En effet, étant donné $(A \to A[x_1, \ldots, x_n]/(f_1, \ldots, f_n))$ dans
$S(A)$, nous pouvons utiliser l'homomorphisme d'anneaux $\varphi : A \to A'$ pour obtenir l'élément correspondant
$(A' \to A'[x_1, \ldots, x_n]/(f^\varphi_1, \ldots, f^\varphi_n))$
de $S(A')$, où $f^\varphi$ désigne le polynôme obtenu en appliquant
$\varphi$ aux coefficients du polynôme $f$.
De plus, on a un diagramme commutatif
$$
\xymatrix{
A \ar[r] \ar[d] & A[x_1, \ldots, x_n]/(f_1, \ldots, f_n) \ar[d] \\
A' \ar[r] & A'[x_1, \ldots, x_n]/(f^\varphi_1, \ldots, f^\varphi_n)
}
$$
dans la catégorie des anneaux. Pour $E \subset S(A)$ fini, posons
$E' = \varphi(E)$ et définissons $B_E \to B_{E'}$ de la manière évidente.
En passant à la limite inductive, on obtient l'homomorphisme recherché $T(A) \to T(A')$, voir le
lemme de Catégories \ref{categories-lemma-functorial-colimit}.
\end{remark}

\begin{lemma}
\label{proetale-lemma-have-sections-quotient}
Soit $A$ un anneau tel que tout homomorphisme d'anneaux \'etale et fidèlement plat
$A \to B$ admette une rétraction. Il en va alors de même pour tout anneau quotient
$A/I$.
\end{lemma}

\begin{proof}
Soit $A/I \to \overline{B}$ \'etale et fidèlement plat. D'après le lemme d'Algèbre
\ref{algebra-lemma-lift-etale}, on peut écrire $\overline{B} = B/IB$ pour
un homomorphisme d'anneaux \'etale $A \to B$. L'image $U$ de $\Spec(B) \to \Spec(A)$
est ouverte et contient $V(I)$. Le complémentaire $Z = \Spec(A) \setminus U$
est donc quasi-compact et disjoint de $V(I)$. Par conséquent,
$Z \subset D(f_1) \cup \ldots \cup D(f_r)$ pour un certain $r \geq 0$
et des $f_i \in I$. Alors $A \to B' = B \times \prod A_{f_i}$
est \'etale et fidèlement plat, et $\overline{B} = B'/IB'$.
La rétraction $B' \to A$ de $A \to B'$ induit donc
une rétraction de $A/I \to \overline{B}$.
\end{proof}

\begin{lemma}
\label{proetale-lemma-have-sections-strictly-henselian}
Soit $A$ un anneau tel que tout homomorphisme d'anneaux \'etale et fidèlement plat
$A \to B$ admette une rétraction. Alors tout anneau local de $A$ en un idéal
maximal est strictement hensélien.
\end{lemma}

\begin{proof}
Soit $\mathfrak m$ un idéal maximal de $A$. Soit $A \to B$ un homomorphisme d'anneaux
\'etale et soit $\mathfrak q \subset B$ un idéal premier
au-dessus de $\mathfrak m$. D'après la description de l'hensélisé strict
$A_\mathfrak m^{sh}$ donnée dans le
lemme d'Algèbre \ref{algebra-lemma-strict-henselization-different},
il suffit de montrer que $A_\mathfrak m = B_\mathfrak q$.
Remarquons qu'il n'existe qu'un nombre fini d'idéaux premiers
$\mathfrak q = \mathfrak q_1, \mathfrak q_2, \ldots, \mathfrak q_n$
au-dessus de $\mathfrak m$ et qu'il n'y a pas de spécialisations
entre eux, car un homomorphisme d'anneaux \'etale est quasi-fini, voir le
lemme d'Algèbre \ref{algebra-lemma-etale-quasi-finite}.
Ainsi, $\mathfrak q_i$ est un idéal maximal et nous pouvons trouver
$g \in \mathfrak q_2 \cap \ldots \cap \mathfrak q_n$, $g \not \in \mathfrak q$
(Algèbre, lemme \ref{algebra-lemma-silly}).
Après avoir remplacé $B$ par $B_g$, nous voyons que $\mathfrak q$
est le seul idéal premier de $B$ au-dessus de $\mathfrak m$.
L'image $U \subset \Spec(A)$ de $\Spec(B) \to \Spec(A)$ est
ouverte (Algèbre, proposition \ref{algebra-proposition-fppf-open}).
Le complémentaire $\Spec(A) \setminus U$ est donc fermé,
et nous pouvons trouver $f \in A$, $f \not \in \mathfrak p$, tel que
$\Spec(A) = U \cup D(f)$. L'homomorphisme d'anneaux $A \to B \times A_f$
est \'etale et fidèlement plat; il admet donc une rétraction
$\sigma : B \times A_f \to A$ en vertu de l'hypothèse faite sur $A$.
Remarquons que $\sigma$ est \'etale, donc plat comme homomorphisme entre algèbres \'etales
sur $A$ (Algèbre, lemme \ref{algebra-lemma-map-between-etale}).
Puisque $\mathfrak q$ est le seul idéal premier de $B \times A_f$ au-dessus
de $A$, nous trouvons que $A_\mathfrak p \to B_\mathfrak q$ admet
une rétraction qui est elle aussi plate. Ainsi,
$A_\mathfrak p \to B_\mathfrak q \to A_\mathfrak p$
sont des homomorphismes locaux plats d'anneaux locaux dont la composée est l'identité. Comme
un homomorphisme local plat d'anneaux locaux est injectif, nous en concluons que ces
homomorphismes sont les isomorphismes recherchés.
\end{proof}

\begin{lemma}
\label{proetale-lemma-have-sections-localize}
Soit $A$ un anneau tel que tout homomorphisme d'anneaux \'etale et fidèlement plat
$A \to B$ admette une rétraction. Soit $Z \subset \Spec(A)$ un sous-schéma fermé.
Soit $A \to A_Z^\sim$ construit comme dans le lemme \ref{proetale-lemma-localization}.
Alors tout homomorphisme d'anneaux \'etale et fidèlement plat $A_Z^\sim \to C$ admet
une rétraction.
\end{lemma}

\begin{proof}
Il existe un homomorphisme d'anneaux \'etale $A \to B'$ tel que
$C = B' \otimes_A A_Z^\sim$ comme $A_Z^\sim$-algèbres.
L'image $U' \subset \Spec(A)$ de $\Spec(B') \to \Spec(A)$
est ouverte et contient $V(I)$; il existe donc $f \in I$ tel
que $\Spec(A) = U' \cup D(f)$. Alors $A \to B' \times A_f$
est \'etale et fidèlement plat. Par hypothèse, il existe une rétraction
$B' \times A_f \to A$. En localisant, on obtient la rétraction voulue
$C \to A_Z^\sim$.
\end{proof}

\begin{lemma}
\label{proetale-lemma-get-w-local-algebraic-residue-field-extensions}
Soit $A \to B$ un homomorphisme d'anneaux induisant des extensions algébriques des corps résiduels.
Il existe un diagramme commutatif
$$
\xymatrix{
B \ar[r] & D \\
A \ar[r] \ar[u] & C \ar[u]
}
$$
qui possède les propriétés suivantes:
\begin{enumerate}
\item $A \to C$ est fidèlement plat et ind-\'etale,
\item $B \to D$ est fidèlement plat et ind-\'etale,
\item $\Spec(C)$ est w-local,
\item $\Spec(D)$ est w-local,
\item $\Spec(D) \to \Spec(C)$ est w-locale,
\item l'ensemble des points fermés de $\Spec(D)$ est l'image réciproque
de l'ensemble des points fermés de $\Spec(C)$,
\item l'ensemble des points fermés de $\Spec(C)$ a pour image tout $\Spec(A)$,
\item l'ensemble des points fermés de $\Spec(D)$ a pour image tout $\Spec(B)$,
\item pour tout idéal maximal $\mathfrak m \subset C$, l'anneau local
$C_\mathfrak m$ est strictement hensélien.
\end{enumerate}
\end{lemma}

\begin{proof}
Il existe un homomorphisme d'anneaux fidèlement plat et ind-Zariski $A \to A'$ tel que
$\Spec(A')$ soit w-local et que l'ensemble des points fermés de
$\Spec(A')$ a pour image tout $\Spec(A)$; voir le lemme \ref{proetale-lemma-make-w-local}.
Soit $I \subset A'$ l'idéal tel que $V(I)$ soit l'ensemble
des points fermés de $\Spec(A')$.
Choisissons $A' \to C'$ comme dans le lemme \ref{proetale-lemma-first-construction}.
Notons que les anneaux locaux $C'_{\mathfrak m'}$, pour les idéaux maximaux
$\mathfrak m' \subset C'$, sont strictement henséliens d'après le
lemme \ref{proetale-lemma-have-sections-strictly-henselian}.
Appliquons le lemme \ref{proetale-lemma-w-local-algebraic-residue-field-extensions}
à $A' \to C'$ et $I \subset A'$ pour obtenir $C' \to C$ avec $C'/IC' \cong C/IC$.
Puisque $A' \to C'$ est fidèlement plat, $\Spec(C'/IC')$
se projette sur l'ensemble des points fermés de $A'$, et donc en particulier
sur $\Spec(A)$. De plus, comme $V(IC) \subset \Spec(C)$
est l'ensemble des points fermés de $C$ et que $C' \to C$ est ind-Zariski
(et identifie les anneaux locaux), on obtient les propriétés (1), (3), (7) et (9).

\medskip\noindent
Notons $J \subset C$ l'idéal tel que $V(J)$ soit l'ensemble des points fermés
de $\Spec(C)$. Posons $D' = B \otimes_A C$. L'homomorphisme d'anneaux
$C \to D'$ induit des extensions algébriques des corps résiduels. Remarquons que,
puisque $V(J) \to \Spec(A)$ est surjective, l'application $T = V(JD') \to \Spec(B)$
est elle aussi surjective. Appliquons le
lemme \ref{proetale-lemma-w-local-algebraic-residue-field-extensions}
à $C \to D'$ et $J \subset C$ pour obtenir 
$D' \to D$ avec $D'/JD' \cong D/JD$.
Toutes les autres propriétés de l'énoncé résultent
immédiatement du
lemme \ref{proetale-lemma-w-local-algebraic-residue-field-extensions}.
\end{proof}








\section{Faiblement \'etale et pro-\'etale}
\label{proetale-section-weakly-etale}

\noindent
Rappelons qu'un homomorphisme d'anneaux $A \to B$ est {\it faiblement \'etale}
si $A \to B$ est plat et si $B \otimes_A B \to B$ est plat. Nous avons
établi certaines propriétés de ces homomorphismes dans
Compléments d'Algèbre, section \ref{more-algebra-section-weakly-etale}.
En particulier, si $A \to B$ est un homomorphisme local et si $A$ est un
anneau local strictement hensélien, alors $A = B$; voir
Compléments d'Algèbre, théorème \ref{more-algebra-theorem-olivier}.
En utilisant ce théorème et le travail accompli ci-dessus, on obtient
le théorème de structure suivant pour les homomorphismes d'anneaux faiblement \'etales.

\begin{proposition}
\label{proetale-proposition-weakly-etale}
Soit $A \to B$ un homomorphisme d'anneaux faiblement \'etale.
Il existe alors un homomorphisme d'anneaux fidèlement plat et ind-\'etale
$B \to B'$ tel que $A \to B'$ soit ind-\'etale.
\end{proposition}

\begin{proof}
L'homomorphisme d'anneaux $A \to B$ induit des extensions algébriques (séparables) des
corps résiduels; voir Compléments d'Algèbre, lemme
\ref{more-algebra-lemma-weakly-etale-residue-field-extensions}.
On peut donc appliquer le
lemme \ref{proetale-lemma-get-w-local-algebraic-residue-field-extensions}
et choisir un diagramme
$$
\xymatrix{
B \ar[r] & D \\
A \ar[r] \ar[u] & C \ar[u]
}
$$
ayant les propriétés énumérées dans le lemme. Notons que $C \to D$
est faiblement \'etale d'après
Compléments d'Algèbre, lemme \ref{more-algebra-lemma-weakly-etale-permanence}.
Choisissons un idéal maximal $\mathfrak m \subset D$. Par construction,
il est au-dessus d'un idéal maximal $\mathfrak m' \subset C$.
D'après Compléments d'Algèbre, théorème \ref{more-algebra-theorem-olivier},
l'homomorphisme d'anneaux $C_{\mathfrak m'} \to D_\mathfrak m$ est un isomorphisme.
Comme tout point de $\Spec(C)$ se spécialise en un point fermé, on en conclut que
$C \to D$ identifie les anneaux locaux.
La proposition \ref{proetale-proposition-maps-wich-identify-local-rings}
s'applique donc à l'homomorphisme d'anneaux $C \to D$. Choisissons $D \to D'$ fidèlement plat
et ind-Zariski tel que $C \to D'$ soit ind-Zariski. Alors
$B \to D'$ résout le problème posé dans la proposition.
\end{proof}





\section{La topologie V et la topologie pro-h}
\label{proetale-section-V-versus-pro-h}

\noindent
La topologie V a été introduite dans
Topologies, section \ref{topologies-section-V}.
La topologie h a été introduite dans
Compléments sur la platitude, section \ref{flat-section-h}.
Une sorte de topologie intermédiaire, la topologie ph,
a été introduite dans Topologies, section \ref{topologies-section-ph}.

\medskip\noindent
Étant donnée une topologie $\tau$ sur une catégorie convenable $\mathcal{C}$
de schémas, on peut introduire une ``topologie pro-$\tau$''
sur $\mathcal{C}$ comme suit. Rappelons que, pour tout objet $X$ de $\mathcal{C}$,
on note $h_X$ le préfaisceau représentable associé à $X$.
Disons provisoirement qu'un morphisme $X \to Y$ de $\mathcal{C}$
est $\tau$-couvrant\footnote{Il ne faut pas confondre ceci avec
la notion de recouvrement. Par exemple, si $\tau = \etale$,
tout morphisme $X \to Y$ qui admet une section est $\tau$-couvrant. Mais notre
définition des recouvrements \'etales $\{V_i \to Y\}_{i \in I}$
impose que chaque $V_i \to Y$ soit \'etale.}
si la $\tau$-faisceautisation du morphisme $h_X \to h_Y$
est surjective. On peut alors définir la topologie pro-$\tau$
comme la topologie la moins fine telle que
\begin{enumerate}
\item la topologie pro-$\tau$ soit plus fine que la topologie $\tau$, et
\item $X \to Y$ soit pro-$\tau$-couvrant si
$Y$ est affine et si $X = \lim X_\lambda$ est une limite projective
filtrante de schémas affines $X_\lambda$ au-dessus de $Y$ telle que
$h_{X_\lambda} \to h_Y$ soit $\tau$-couvrant pour tout $\lambda$.
\end{enumerate}
Nous employons cette formulation minutieuse parce que nous ne voulons pas
fixer un choix de recouvrements pro-$\tau$: selon $\tau$,
différents choix de familles de recouvrements conviennent.
Par exemple, nous verrons à la section \ref{proetale-section-proetale} que
pour définir la topologie pro-\'etale, il est commode de considérer des familles
de morphismes faiblement \'etales possédant une certaine propriété de finitude.
Plus généralement, la construction proposée dans ce paragraphe
vise surtout à motiver les résultats de cette section, et nous ne définirons jamais
implicitement une topologie pro-$\tau$ par cette méthode.

\medskip\noindent
Le lemme suivant affirme que la
topologie pro-V est égale à la topologie V.

\begin{lemma}
\label{proetale-lemma-pro-V-V}
Soit $Y$ un schéma affine. Soit $X = \lim X_i$ une limite projective filtrante
de schémas affines au-dessus de $Y$. Les conditions suivantes sont équivalentes:
\begin{enumerate}
\item $\{X \to Y\}$ est un recouvrement V standard
(Topologies, définition \ref{topologies-definition-standard-V-covering}), et
\item $\{X_i \to Y\}$ est un recouvrement V standard pour tout $i$.
\end{enumerate}
\end{lemma}

\begin{proof}
Le singleton $\{X \to Y\}$ est un recouvrement V standard si et seulement si,
pour tout morphisme $g : \Spec(V) \to Y$, il existe une extension
d'anneaux de valuation $V \subset W$ et un diagramme commutatif
$$
\xymatrix{
\Spec(W) \ar[r] \ar[d] & X \ar[d] \\
\Spec(V) \ar[r]^g & Y
}
$$
L'implication (1) $\Rightarrow$ (2) résulte donc immédiatement de la définition.
Réciproquement, supposons (2) et considérons $g : \Spec(V) \to Y$ comme ci-dessus.
Écrivons $\Spec(V) \times_Y X_i = \Spec(A_i)$.
Puisque $\{X_i \to Y\}$ est un recouvrement V standard, on peut choisir
un anneau de valuation $W_i$ et un homomorphisme d'anneaux $A_i \to W_i$ tels que
la composée $V \to A_i \to W_i$ soit une extension d'anneaux de
valuation. En particulier, le quotient $A'_i$ de $A_i$ par sa
$V$-torsion est une $V$-algèbre fidèlement plate. La platitude résulte de
Compléments d'Algèbre, lemme
\ref{more-algebra-lemma-valuation-ring-torsion-free-flat} et
la surjectivité sur les spectres du fait que $A_i \to W_i$ se factorise par $A'_i$.
Ainsi,
$$
A = \colim A'_i
$$
est une $V$-algèbre fidèlement plate
(Algèbre, lemme \ref{algebra-lemma-colimit-faithfully-flat}).
Puisque $\{\Spec(A) \to \Spec(V)\}$ est un recouvrement fpqc standard, c'est
un recouvrement V standard
(Topologies, lemme \ref{topologies-lemma-standard-fpqc-standard-V});
on peut donc choisir $\Spec(W) \to \Spec(A)$ tel que
$V \to W$ soit une extension d'anneaux de valuation. En composant
avec le morphisme $\Spec(A) \to X = \Spec(\colim A_i)$,
on achève la démonstration.
\end{proof}

\noindent
Le lemme suivant affirme que la topologie pro-h est égale à la
topologie pro-ph, elle-même égale à la topologie V.

\begin{lemma}
\label{proetale-lemma-pro-h-V}
Soit $X \to Y$ un morphisme de schémas affines. Les conditions suivantes sont équivalentes:
\begin{enumerate}
\item $\{X \to Y\}$ est un recouvrement V standard
(Topologies, définition \ref{topologies-definition-standard-V-covering}),
\item $X = \lim X_i$ est une limite projective filtrante de schémas affines sur $Y$
telle que $\{X_i \to Y\}$ soit un recouvrement ph pour tout $i$, et
\item $X = \lim X_i$ est une limite projective filtrante de schémas affines sur $Y$
telle que $\{X_i \to Y\}$ soit un recouvrement h pour tout $i$.
\end{enumerate}
\end{lemma}

\begin{proof}
Démonstration de (2) $\Rightarrow$ (1). Rappelons qu'un recouvrement V donné par une
seule flèche entre schémas affines est un recouvrement V standard; voir
Topologies, définition \ref{topologies-definition-V-covering} et
lemme \ref{topologies-lemma-refine-standard-V}.
Rappelons que tout recouvrement ph est un recouvrement V; voir
Topologies, lemme
\ref{topologies-lemma-zariski-etale-smooth-syntomic-fppf-fpqc-ph-V}.
Ainsi, si $X = \lim X_i$ comme en (2), alors $\{X_i \to Y\}$
est un recouvrement V standard pour tout $i$. Le
lemme \ref{proetale-lemma-pro-V-V} montre donc que (1) est vraie.

\medskip\noindent
Démonstration de (3) $\Rightarrow$ (2). C'est clair, car un recouvrement h
est toujours un recouvrement ph; voir
Compléments sur la platitude, définition \ref{flat-definition-h-covering}.

\medskip\noindent
Démonstration de (1) $\Rightarrow$ (3). C'est l'implication intéressante, mais
le fond de la démonstration est contenu dans Compléments sur la platitude, lemme
\ref{flat-lemma-equivalence-h-v-locally-finite-presentation}.
Écrivons $X = \Spec(A)$ et $Y = \Spec(R)$. On peut écrire
$A = \colim A_i$, où $A_i$ est de présentation finie sur $R$; voir
Algèbre, lemme \ref{algebra-lemma-ring-colimit-fp}.
Posons $X_i = \Spec(A_i)$. Alors $\{X_i \to Y\}$ est un recouvrement V standard
pour tout $i$, d'après (1) et
Topologies, lemme \ref{topologies-lemma-refine-standard-V}.
Ainsi $\{X_i \to Y\}$ est un recouvrement h d'après
Compléments sur la platitude, définition \ref{flat-definition-h-covering}.
Ceci achève la démonstration.
\end{proof}

\noindent
Le lemme suivant affirme, en gros, qu'un faisceau h
qui commute aux limites satisfait la condition de faisceau pour les recouvrements V.
Comparer aussi avec la remarque \ref{proetale-remark-h-limit-preserving}.

\begin{lemma}
\label{proetale-lemma-h-limit-preserving}
Soit $S$ un schéma. Soit $F$ un foncteur contravariant défini
sur la catégorie de tous les schémas au-dessus de $S$. Si
\begin{enumerate}
\item $F$ satisfait la condition de faisceau pour la topologie h, et
\item $F$ commute aux limites
(Limites, remarque \ref{limits-remark-limit-preserving}),
\end{enumerate}
alors $F$ satisfait la condition de faisceau pour la topologie V.
\end{lemma}

\begin{proof}
Nous allons le prouver en vérifiant (1) et (2') de
Topologies, lemme \ref{topologies-lemma-sheaf-property-V}.
La condition de faisceau pour les recouvrements de Zariski résulte
du fait que $F$ satisfait la condition de faisceau pour tout recouvrement h.
Enfin, supposons que $X \to Y$ soit un morphisme de schémas affines
au-dessus de $S$ tel que $\{X \to Y\}$ soit un recouvrement V.
D'après le lemme \ref{proetale-lemma-pro-h-V}, on peut écrire $X = \lim X_i$
comme une limite projective filtrante de schémas affines
au-dessus de $Y$, telle que $\{X_i \to Y\}$ soit un recouvrement h pour tout $i$.
On obtient
\begin{align*}
&
\text{Égalisateur}(
\xymatrix{
F(X) \ar@<1ex>[r] \ar@<-1ex>[r] &
F(X \times_Y X)
}
)
\\
& =
\text{Égalisateur}(
\xymatrix{
\colim F(X_i) \ar@<1ex>[r] \ar@<-1ex>[r] &
\colim F(X_i \times_Y X_i)
}
)
\\
& =
\colim
\text{Égalisateur}(
\xymatrix{
F(X_i) \ar@<1ex>[r] \ar@<-1ex>[r] &
F(X_i \times_Y X_i)
}
)
\\
& =
\colim F(Y) = F(Y)
\end{align*}
ce qui est l'assertion voulue.
La première égalité vient de ce que $F$ commute aux limites, $X = \lim X_i$ et
$X \times_Y X = \lim X_i \times_Y X_i$.
La deuxième vient de l'exactitude des limites inductives filtrantes.
La troisième vient de ce que $F$ satisfait la condition de faisceau
pour les recouvrements h.
\end{proof}

\begin{remark}
\label{proetale-remark-h-limit-preserving}
Soit $S$ un schéma contenu dans un grand site $\Sch_h$. Soit $F$ un faisceau
d'ensembles sur $(\Sch/S)_h$ tel que $F(T) = \colim F(T_i)$ chaque fois que
$T = \lim T_i$ est une limite projective filtrante de schémas affines dans $(\Sch/S)_h$.
Dans cette situation, $F$ se prolonge de manière unique en un foncteur contravariant $F'$
sur la catégorie de tous les schémas au-dessus de $S$, tel que (a) $F'$ satisfasse la
condition de faisceau pour la topologie h et (b) $F'$ commute aux limites.
Voir Compléments sur la platitude, lemme \ref{flat-lemma-extend-sheaf-h}.
Dans cette situation, le lemme \ref{proetale-lemma-h-limit-preserving}
affirme que $F'$ satisfait
la condition de faisceau pour la topologie V.
\end{remark}







\section{Construction de recouvrements w-contractiles}
\label{proetale-section-w-contractible}

\noindent
Dans cette section, nous construisons des recouvrements w-contractiles de schémas affines.

\begin{definition}
\label{proetale-definition-w-contractible}
Soit $A$ un anneau. On dit que $A$ est {\it w-contractile} si tout
homomorphisme d'anneaux fidèlement plat et faiblement \'etale $A \to B$ admet une rétraction.
\end{definition}

\noindent
Remarquons que, d'après la proposition \ref{proetale-proposition-weakly-etale},
une définition équivalente consisterait à demander que tout homomorphisme d'anneaux
fidèlement plat et ind-\'etale $A \to B$ admet une rétraction.
Voici une observation essentielle qui nous permettra de construire
des anneaux w-contractiles.

\begin{lemma}
\label{proetale-lemma-w-local-strictly-henselian-extremally-disconnected}
Soit $A$ un anneau. Les conditions suivantes sont équivalentes:
\begin{enumerate}
\item $A$ est w-contractile,
\item tout homomorphisme d'anneaux fidèlement plat et ind-\'etale $A \to B$ admet
une rétraction, et
\item $A$ satisfait les conditions suivantes:
\begin{enumerate}
\item $\Spec(A)$ est w-local,
\item $\pi_0(\Spec(A))$ est extrêmement discontinu, et
\item pour tout idéal maximal $\mathfrak m \subset A$, l'anneau
local $A_\mathfrak m$ est strictement hensélien.
\end{enumerate}
\end{enumerate}
\end{lemma}

\begin{proof}
L'équivalence de (1) et (2) résulte immédiatement de la
proposition \ref{proetale-proposition-weakly-etale}.

\medskip\noindent
Supposons (3)(a), (3)(b) et (3)(c). Soit $A \to B$ un homomorphisme d'anneaux fidèlement plat
et ind-\'etale. Nous utiliserons sans autre mention le fait qu'un homomorphisme plat
$A \to B$ est fidèlement plat si et seulement si tout point fermé
de $\Spec(A)$ appartient à l'image de $\Spec(B) \to \Spec(A)$.
Montrons que $A \to B$ admet une rétraction.

\medskip\noindent
Soit $I \subset A$ un idéal tel que $V(I) \subset \Spec(A)$ soit
l'ensemble des points fermés de $\Spec(A)$. 
On peut remplacer $B$ par l'anneau $C$ construit au
lemme \ref{proetale-lemma-w-local-algebraic-residue-field-extensions}
pour $A \to B$ et $I \subset A$.
On peut ainsi supposer que $\Spec(B)$ est w-local et que l'ensemble des
points fermés de $\Spec(B)$ est égal à $V(IB)$. Dans ce cas, $A \to B$
identifie les anneaux locaux par la condition (3)(c), car il suffit de le vérifier
aux idéaux maximaux de $B$ situés au-dessus d'idéaux maximaux de $A$.
Ainsi $A \to B$ admet une rétraction d'après le
lemme \ref{proetale-lemma-w-local-extremally-disconnected}.

\medskip\noindent
Supposons (1) ou, ce qui revient au même, (2). On a (3)(c) d'après le
lemme \ref{proetale-lemma-have-sections-strictly-henselian}.
Les propriétés (3)(a) et (3)(b) résultent du
lemme \ref{proetale-lemma-w-local-extremally-disconnected}.
\end{proof}

\begin{proposition}
\label{proetale-proposition-find-w-contractible}
Pour tout anneau $A$, il existe un homomorphisme d'anneaux fidèlement plat et ind-\'etale
$A \to D$ tel que $D$ soit w-contractile.
\end{proposition}

\begin{proof}
En appliquant le lemme \ref{proetale-lemma-get-w-local-algebraic-residue-field-extensions}
à $\text{id}_A : A \to A$, on obtient un homomorphisme d'anneaux fidèlement plat et ind-\'etale
$A \to C$ tel que $C$ soit w-local et que tout anneau local en un
idéal maximal de $C$ soit strictement hensélien.
Choisissons un espace extrêmement discontinu $T$ et une application continue
surjective $T \to \pi_0(\Spec(C))$; voir
Topologie, lemme \ref{topology-lemma-existence-projective-cover}.
Notons que $T$ est profini. Appliquons le lemme \ref{proetale-lemma-construct-profinite}
pour obtenir un homomorphisme d'anneaux ind-Zariski $C \to D$ tel que
$\pi_0(\Spec(D)) \to \pi_0(\Spec(C))$ réalise $T \to \pi_0(\Spec(C))$
et que
$$
\xymatrix{
\Spec(D) \ar[r] \ar[d] & \pi_0(\Spec(D)) \ar[d] \\
\Spec(C) \ar[r] & \pi_0(\Spec(C))
}
$$
soit cartésien dans la catégorie des espaces topologiques. Notons que $\Spec(D)$
est w-local, que $\Spec(D) \to \Spec(C)$ est w-locale et que
l'ensemble des points fermés de $\Spec(D)$ est l'image réciproque de
l'ensemble des points fermés de $\Spec(C)$; voir le lemme \ref{proetale-lemma-silly}.
Il reste donc vrai que les anneaux locaux de $D$ en ses idéaux
maximaux sont strictement henséliens (puisqu'ils sont isomorphes aux
anneaux locaux de $C$ aux idéaux maximaux correspondants).
Il résulte du
lemme \ref{proetale-lemma-w-local-strictly-henselian-extremally-disconnected}
que $D$ est w-contractile.
\end{proof}

\begin{remark}
\label{proetale-remark-size-w-contractible}
Soit $A$ un anneau. Soit $\kappa$ un cardinal infini supérieur ou
égal au cardinal de $A$. Alors le cardinal de l'anneau
$D$ construit dans la proposition \ref{proetale-proposition-find-w-contractible}
est au plus
$$
\kappa^{2^{2^{2^\kappa}}}.
$$
En effet, l'homomorphisme d'anneaux $A \to D$
s'obtient comme le composé
$$
A \to A_w = A' \to C' \to C \to D.
$$
Les trois premières étapes de la construction sont effectuées
dans le premier paragraphe de la démonstration du
lemme \ref{proetale-lemma-get-w-local-algebraic-residue-field-extensions}.
Pour la première étape, on a $|A_w| \leq \kappa$ d'après la
remarque \ref{proetale-remark-size-w}.
On a $|C'| \leq \kappa$ d'après la
remarque \ref{proetale-remark-size-T}.
Puis $|C| \leq \kappa$, car $C$ est une localisation de $(C')_w$
(il est construit à partir de $C'$ par une application du
lemme \ref{proetale-lemma-localize-along-closed-profinite}
dans la démonstration du lemme \ref{proetale-lemma-w-local-algebraic-residue-field-extensions}).
Ainsi $C$ possède au plus $2^\kappa$ idéaux maximaux.
Enfin, l'homomorphisme d'anneaux $C \to D$ identifie les anneaux locaux, et le
cardinal de l'ensemble des idéaux maximaux de $D$ est au plus
$2^{2^{2^\kappa}}$ d'après
Topologie, remarque \ref{topology-remark-size-projective-cover}.
Puisque $D \subset \prod_{\mathfrak m \subset D} D_\mathfrak m$, on voit
que le cardinal de $D$ est au plus celui affiché ci-dessus.
\end{remark}

\begin{lemma}
\label{proetale-lemma-finite-finitely-presented-over-extremally-disconnected}
Soit $A \to B$ un homomorphisme d'anneaux quasi-fini et de présentation finie.
Si les corps résiduels de $A$ sont séparablement clos
et si $\Spec(A)$ est séparé et extrêmement discontinu, alors $\Spec(B)$ est
extrêmement discontinu.
\end{lemma}

\begin{proof}
Posons $X = \Spec(A)$ et $Y = \Spec(B)$. Choisissons une partition finie
$X = \coprod X_i$ et des morphismes $X'_i \to X_i$ comme dans
Cohomologie \'etale, lemme
\ref{etale-cohomology-lemma-decompose-quasi-finite-morphism}.
L'application d'espaces topologiques $\coprod X_i \to X$ (dont la source
est la somme disjointe dans la catégorie des espaces topologiques)
admet une section d'après la proposition de Topologie
\ref{topology-proposition-projective-in-category-hausdorff-qc}.
Ainsi, $X$ est, comme espace topologique, la somme disjointe des
strates $X_i$. On peut donc remplacer $X$ par les $X_i$ et supposer qu'il
existe un morphisme fini localement libre et surjectif $X' \to X$ tel que
$(X' \times_X Y)_{red}$ soit isomorphe à une somme disjointe finie de
copies de $X'_{red}$. Représentons la situation par le diagramme
$$
\xymatrix{
\coprod_{i = 1, \ldots, r} X' \ar[r] \ar[d] & Y \ar[d] \\
X' \ar[r] & X
}
$$
L'hypothèse sur les corps résiduels de $A$ implique que
ce diagramme est cartésien au niveau des ensembles de points
sous-jacents (détails omis).
Puisque $X$ est extrêmement discontinu et que $X'$ est séparé
(lemme \ref{proetale-lemma-profinite-goes-up}), l'application continue
$X' \to X$ admet une section continue $\sigma$. Alors
$\coprod_{i = 1, \ldots, r} \sigma(X) \to Y$ est une application continue
bijective. D'après
Topologie, lemme \ref{topology-lemma-bijective-map},
c'est un homéomorphisme, ce qui achève la démonstration.
\end{proof}

\begin{lemma}
\label{proetale-lemma-finite-finitely-presented-over-w-contractible}
Soit $A \to B$ un homomorphisme d'anneaux fini et de présentation finie.
Si $A$ est w-contractile, alors $B$ l'est aussi.
\end{lemma}

\begin{proof}
Nous utiliserons le critère du
lemme \ref{proetale-lemma-w-local-strictly-henselian-extremally-disconnected}.
Posons $X = \Spec(A)$ et $Y = \Spec(B)$, et notons $f : Y \to X$
le morphisme induit.
Comme $f : Y \to X$ est un morphisme fini, l'ensemble des points fermés
$Y_0$ de $Y$ est l'image réciproque de l'ensemble des points fermés
$X_0$ de $X$. Soit $y \in Y$ d'image $x \in X$. Alors
$x$ se spécialise en un unique point fermé $x_0 \in X$.
Écrivons $f^{-1}(\{x_0\}) = \{y_1, \ldots, y_n\}$, où les $y_i$
sont fermés dans $Y$. Puisque $R = \mathcal{O}_{X, x_0}$ est strictement hensélien
et que $f$ est fini, $Y \times_{f, X} \Spec(R)$ est
égal à $\coprod_{i = 1, \ldots, n} \Spec(R_i)$, où
chaque $R_i$ est un anneau local fini sur $R$
dont l'idéal maximal correspond à $y_i$; voir
Algèbre, lemme \ref{algebra-lemma-characterize-henselian}, partie (10).
Alors $y$ appartient à exactement l'un de ces $\Spec(R_i)$,
et l'on voit que $y$ se spécialise en exactement l'un des $y_i$.
Autrement dit, tout point de $Y$ se spécialise en un unique
point de $Y_0$. Ainsi, $Y$ est w-local.
Pour tout $y \in Y_0$ d'image $x \in X_0$,
l'anneau $\mathcal{O}_{Y, y}$ est strictement hensélien d'après le
lemme d'Algèbre \ref{algebra-lemma-finite-over-henselian},
appliqué à $\mathcal{O}_{X, x} \to B \otimes_A \mathcal{O}_{X, x}$.
Il reste à montrer que $Y_0$ est extrêmement discontinu.
Pour cela, considérons $X_0 \times_X Y \to X_0$,
où $X_0 \subset X$ est muni de la structure de sous-schéma réduite induite.
Notons que l'espace topologique sous-jacent à
$X_0 \times_X Y$ coïncide avec $Y_0$. Le résultat voulu découle alors du
lemme \ref{proetale-lemma-finite-finitely-presented-over-extremally-disconnected}.
\end{proof}

\begin{lemma}
\label{proetale-lemma-localization-w-contractible}
Soit $A$ un anneau. Soit $Z \subset \Spec(A)$ un sous-ensemble fermé
de la forme $Z = V(f_1, \ldots, f_r)$. Posons $B = A_Z^\sim$; voir le
lemme \ref{proetale-lemma-localization}. Si $A$ est w-contractile, alors $B$ l'est aussi.
\end{lemma}

\begin{proof}
Soit $A_Z^\sim \to B$ un homomorphisme d'anneaux faiblement \'etale et fidèlement plat.
Considérons l'homomorphisme d'anneaux
$$
A \longrightarrow A_{f_1} \times \ldots \times A_{f_r} \times B
$$
qui est fidèlement plat et faiblement \'etale. Si $A$ est w-contractile,
cet homomorphisme admet une rétraction $\sigma$. Considérons le morphisme
$$
\Spec(A_Z^\sim) \to \Spec(A) \xrightarrow{\Spec(\sigma)}
\coprod \Spec(A_{f_i}) \amalg \Spec(B)
$$
Tout point de $Z \subset \Spec(A_Z^\sim)$ a son image dans la composante
$\Spec(B)$. Puisque tout point de $\Spec(A_Z^\sim)$ se spécialise en un
point de $Z$, on obtient un morphisme $\Spec(A_Z^\sim) \to \Spec(B)$
comme on le voulait.
\end{proof}






\section{Le site pro-\'etale}
\label{proetale-section-proetale}

\noindent
Dans cette section, nous nous bornons à définir et à construire
les différents sites pro-\'etales ainsi que les morphismes entre eux. L'existence
d'objets faiblement contractiles sera établie dans la
section \ref{proetale-section-weakly-contractible}.

\medskip\noindent
La topologie pro-\'etale ressemble quelque peu à
la topologie fpqc (voir Topologies, section \ref{topologies-section-fpqc}),
en ce que le topos des faisceaux sur le petit site pro-\'etale d'un schéma
dépend du choix de la catégorie sous-jacente de schémas. On ne peut donc pas
parler du {\it topos} pro-\'etale d'un schéma. Toutefois, les
groupes de cohomologie d'un faisceau ne changent pas si l'on agrandit
la catégorie sous-jacente de schémas; voir la section \ref{proetale-section-change-universe}.

\medskip\noindent
Nous définirons les recouvrements pro-\'etales au moyen de morphismes de schémas faiblement \'etales;
voir Compléments sur les morphismes, section \ref{more-morphisms-section-weakly-etale}.
La raison est que, d'une part, il est quelque peu malaisé de définir
la notion de morphisme pro-\'etale de schémas et que, d'autre part, la
proposition \ref{proetale-proposition-weakly-etale}
garantit que l'on obtient les mêmes faisceaux\footnote{Plus précisément,
la topologie pro-\'etale obtenue avec notre choix de recouvrements
est la même que celle issue de la procédure générale expliquée
à la section \ref{proetale-section-V-versus-pro-h} en partant de $\tau = \etale$.}
avec la définition suivante.

\begin{definition}
\label{proetale-definition-fpqc-covering}
Soit $T$ un schéma. Un {\it recouvrement pro-\'etale de $T$} est une famille
de morphismes de schémas $\{f_i : T_i \to T\}_{i \in I}$
telle que chaque $f_i$ soit faiblement \'etale et que, pour tout ouvert affine
$U \subset T$, il existe $n \geq 0$, une application
$a : \{1, \ldots, n\} \to I$ et des ouverts affines
$V_j \subset T_{a(j)}$, $j = 1, \ldots, n$,
tels que $\bigcup_{j = 1}^n f_{a(j)}(V_j) = U$.
\end{definition}

\noindent
Bien entendu, cette condition implique que $T = \bigcup f_i(T_i)$.
Voici un lemme qui permettra de reconnaître les recouvrements pro-\'etales.
Il permettra aussi de ramener de nombreux lemmes sur les recouvrements pro-\'etales
aux résultats correspondants sur les recouvrements fpqc.

\begin{lemma}
\label{proetale-lemma-recognize-proetale-covering}
Soit $T$ un schéma. Soit $\{f_i : T_i \to T\}_{i \in I}$ une famille de
morphismes de schémas de but $T$. Les conditions suivantes sont équivalentes :
\begin{enumerate}
\item $\{f_i : T_i \to T\}_{i \in I}$ est un recouvrement pro-\'etale,
\item chaque $f_i$ est faiblement \'etale et $\{f_i : T_i \to T\}_{i \in I}$
est un recouvrement fpqc,
\item chaque $f_i$ est faiblement \'etale et, pour tout ouvert affine $U \subset T$,
il existe des ouverts quasi-compacts $U_i \subset T_i$, presque tous vides,
tels que $U = \bigcup f_i(U_i)$,
\item chaque $f_i$ est faiblement \'etale et il existe un recouvrement ouvert affine
$T = \bigcup_{\alpha \in A} U_\alpha$ tel que, pour chaque $\alpha \in A$,
il existe $i_{\alpha, 1}, \ldots, i_{\alpha, n(\alpha)} \in I$
et des ouverts quasi-compacts $U_{\alpha, j} \subset T_{i_{\alpha, j}}$ tels que
$U_\alpha =
\bigcup_{j = 1, \ldots, n(\alpha)} f_{i_{\alpha, j}}(U_{\alpha, j})$.
\end{enumerate}
Si $T$ est quasi-séparé, ces conditions sont encore équivalentes à la suivante :
\begin{enumerate}
\item[(5)] chaque $f_i$ est faiblement \'etale et, pour tout $t \in T$, il existe
$i_1, \ldots, i_n \in I$ et des ouverts quasi-compacts $U_j \subset T_{i_j}$
tels que $\bigcup_{j = 1, \ldots, n} f_{i_j}(U_j)$ soit un
voisinage de $t$ dans $T$ (non nécessairement ouvert).
\end{enumerate}
\end{lemma}

\begin{proof}
L'équivalence de (1) et (2) découle immédiatement des définitions.
Le lemme résulte donc de
Topologies, lemme \ref{topologies-lemma-recognize-fpqc-covering}.
\end{proof}

\begin{lemma}
\label{proetale-lemma-etale-proetale}
Tout recouvrement \'etale et tout recouvrement de Zariski sont des recouvrements pro-\'etales.
\end{lemma}

\begin{proof}
Cela résulte du résultat correspondant pour les recouvrements fpqc
(Topologies, lemme
\ref{topologies-lemma-zariski-etale-smooth-syntomic-fppf-fpqc}),
du lemme \ref{proetale-lemma-recognize-proetale-covering} et
du fait qu'un morphisme \'etale est faiblement \'etale; voir
Compléments sur les morphismes, lemme \ref{more-morphisms-lemma-when-weakly-etale}.
\end{proof}

\begin{lemma}
\label{proetale-lemma-proetale}
Soit $T$ un schéma.
\begin{enumerate}
\item Si $T' \to T$ est un isomorphisme, alors $\{T' \to T\}$
est un recouvrement pro-\'etale de $T$.
\item Si $\{T_i \to T\}_{i\in I}$ est un recouvrement pro-\'etale et si, pour chaque
$i$, on a un recouvrement pro-\'etale $\{T_{ij} \to T_i\}_{j\in J_i}$, alors
$\{T_{ij} \to T\}_{i \in I, j\in J_i}$ est un recouvrement pro-\'etale.
\item Si $\{T_i \to T\}_{i\in I}$ est un recouvrement pro-\'etale
et si $T' \to T$ est un morphisme de schémas, alors
$\{T' \times_T T_i \to T'\}_{i\in I}$ est un recouvrement pro-\'etale.
\end{enumerate}
\end{lemma}

\begin{proof}
Cela résulte du fait que les composés et les changements de base
de morphismes faiblement \'etales sont faiblement \'etales
(Compléments sur les morphismes, lemmes
\ref{more-morphisms-lemma-composition-weakly-etale} et
\ref{more-morphisms-lemma-base-change-weakly-etale}),
du lemme \ref{proetale-lemma-recognize-proetale-covering} et
des résultats correspondants pour les recouvrements fpqc; voir
Topologies, lemme \ref{topologies-lemma-fpqc}.
\end{proof}

\begin{lemma}
\label{proetale-lemma-proetale-affine}
Soit $T$ un schéma affine. Soit $\{T_i \to T\}_{i \in I}$ un recouvrement pro-\'etale
de $T$. Il existe alors un recouvrement pro-\'etale
$\{U_j \to T\}_{j = 1, \ldots, n}$ qui raffine
$\{T_i \to T\}_{i \in I}$ et dans lequel chaque $U_j$ est un schéma
affine. De plus, on peut choisir chaque $U_j$ comme ouvert affine
de l'un des $T_i$.
\end{lemma}

\begin{proof}
Cela résulte directement de la définition.
\end{proof}

\noindent
Nous définissons donc comme suit les recouvrements standards correspondants dans le cas affine.

\begin{definition}
\label{proetale-definition-standard-proetale}
Soit $T$ un schéma affine. Un {\it recouvrement pro-\'etale standard}
de $T$ est une famille $\{f_i : T_i \to T\}_{i = 1, \ldots, n}$
où chaque $T_j$ est affine, chaque $f_i$ est faiblement \'etale et
$T = \bigcup f_i(T_i)$.
\end{definition}

\noindent
Nous suivons le schéma général donné dans
Topologies, section \ref{topologies-section-procedure},
pour construire le grand site pro-\'etale avec lequel nous travaillerons.
Toutefois, comme nous avons besoin d'anneaux un peu plus grands pour tenir compte de la taille
de certaines constructions, nous modifions légèrement celles-ci.

\begin{definition}
\label{proetale-definition-big-proetale-site}
Un {\it grand site pro-\'etale} est un site $\Sch_\proetale$ quelconque
au sens de Sites, définition \ref{sites-definition-site}, construit comme suit :
\begin{enumerate}
\item Choisir un ensemble quelconque de schémas $S_0$ et un ensemble quelconque de recouvrements pro-\'etales
$\text{Cov}_0$ parmi ces schémas.
\item Remplacer la fonction $Bound$ de
Ensembles, équation (\ref{sets-equation-bound}), par
$$
Bound(\kappa) = \max\{\kappa^{2^{2^{2^\kappa}}}, \kappa^{\aleph_0}, \kappa^+\}.
$$
\item Prendre pour catégorie sous-jacente une catégorie $\Sch_\alpha$ quelconque
construite comme dans Ensembles, lemme \ref{sets-lemma-construct-category},
à partir de l'ensemble $S_0$ et de la fonction $Bound$.
\item Choisir un ensemble quelconque de recouvrements comme dans
Ensembles, lemme \ref{sets-lemma-coverings-site}, à partir de la
catégorie $\Sch_\alpha$, de la classe des recouvrements pro-\'etales
et de l'ensemble $\text{Cov}_0$ choisi ci-dessus.
\end{enumerate}
\end{definition}

\noindent
Voir les remarques qui suivent Topologies,
définition \ref{topologies-definition-big-zariski-site},
pour la motivation et les explications relatives à la définition des grands sites.

\medskip\noindent
Il résultera du lemme \ref{proetale-lemma-proetale-induced} que la
topologie d'un grand site pro-\'etale $\Sch_\proetale$ est, en un certain sens,
induite par la topologie pro-\'etale sur la catégorie de tous les schémas.

\begin{definition}
\label{proetale-definition-big-small-proetale}
Soit $S$ un schéma. Soit $\Sch_\proetale$ un grand site pro-\'etale
contenant $S$.
\begin{enumerate}
\item Le {\it grand site pro-\'etale de $S$}, noté
$(\Sch/S)_\proetale$, est le site $\Sch_\proetale/S$
introduit dans Sites, section \ref{sites-section-localize}.
\item Le {\it petit site pro-\'etale de $S$}, que nous notons
$S_\proetale$, est la sous-catégorie pleine de $(\Sch/S)_\proetale$
dont les objets sont les $U/S$ tels que $U \to S$ soit faiblement \'etale.
Un recouvrement de $S_\proetale$ est tout recouvrement $\{U_i \to U\}$ de
$(\Sch/S)_\proetale$ avec $U \in \Ob(S_\proetale)$.
\item Le {\it grand site pro-\'etale affine de $S$}, noté
$(\textit{Aff}/S)_\proetale$, est la sous-catégorie pleine de
$(\Sch/S)_\proetale$ dont les objets sont les $U/S$ affines.
Un recouvrement de $(\textit{Aff}/S)_\proetale$ est tout recouvrement
$\{U_i \to U\}$ de $(\Sch/S)_\proetale$ qui est un
recouvrement pro-\'etale standard.
\end{enumerate}
\end{definition}

\noindent
Il n'est pas tout à fait évident que le petit site pro-\'etale et
le grand site pro-\'etale affine soient des sites. Vérifions-le maintenant.

\begin{lemma}
\label{proetale-lemma-verify-site-proetale}
Soit $S$ un schéma. Soit $\Sch_\proetale$ un grand site pro-\'etale
contenant $S$. Alors $S_\proetale$ et $(\textit{Aff}/S)_\proetale$ sont tous deux des sites.
\end{lemma}

\begin{proof}
Montrons que $S_\proetale$ est un site. C'est une catégorie munie d'un
ensemble donné de familles de morphismes à cible fixée. Il faut donc
vérifier les propriétés (1), (2) et (3) de
Sites, définition \ref{sites-definition-site}.
Puisque $(\Sch/S)_\proetale$ est un site, il suffit de montrer
que, pour tout recouvrement $\{U_i \to U\}$ de $(\Sch/S)_\proetale$
avec $U \in \Ob(S_\proetale)$, on a aussi $U_i \in \Ob(S_\proetale)$.
Cela découle des définitions,
car la composée de morphismes faiblement \'etales est faiblement \'etale.

\medskip\noindent
Pour montrer que $(\textit{Aff}/S)_\proetale$ est un site, en raisonnant comme ci-dessus,
il suffit de montrer que la famille des recouvrements pro-\'etales standards
de schémas affines satisfait les propriétés (1), (2) et (3) de
Sites, définition \ref{sites-definition-site}.
Cela découle du lemme \ref{proetale-lemma-recognize-proetale-covering}
et du résultat correspondant pour les recouvrements fpqc standards
(Topologies, lemme \ref{topologies-lemma-fpqc-affine-axioms}).
\end{proof}

\begin{lemma}
\label{proetale-lemma-fibre-products-proetale}
Soit $S$ un schéma. Soit $\Sch_\proetale$ un grand site pro-\'etale
contenant $S$. Soit $\Sch$ la catégorie de tous les schémas.
\begin{enumerate}
\item Les catégories $\Sch_\proetale$, $(\Sch/S)_\proetale$,
$S_\proetale$ et $(\textit{Aff}/S)_\proetale$ possèdent des produits fibrés
qui coïncident avec les produits fibrés dans $\Sch$.
\item Les catégories $\Sch_\proetale$, $(\Sch/S)_\proetale$ et
$S_\proetale$ possèdent des égalisateurs qui coïncident avec les égalisateurs dans $\Sch$.
\item Les catégories $(\Sch/S)_\proetale$ et $S_\proetale$ possèdent toutes deux
un objet final, à savoir $S/S$.
\item La catégorie $\Sch_\proetale$ possède un objet final qui coïncide
avec l'objet final de $\Sch$, à savoir $\Spec(\mathbf{Z})$.
\end{enumerate}
\end{lemma}

\begin{proof}
La catégorie $\Sch_\proetale$ contient $\Spec(\mathbf{Z})$ et
est stable par produits et produits fibrés par construction; voir
Ensembles, lemme \ref{sets-lemma-what-is-in-it}.
Supposons donnés des morphismes $U \to S$, $V \to U$, $W \to U$
de schémas avec $U, V, W \in \Ob(\Sch_\proetale)$.
Le produit fibré $V \times_U W$ dans $\Sch_\proetale$
est un produit fibré dans $\Sch$ et
est le produit fibré de $V/S$ et $W/S$ au-dessus de $U/S$ dans
la catégorie de tous les schémas au-dessus de $S$; c'est donc aussi un
produit fibré dans $(\Sch/S)_\proetale$.
Cela démontre le résultat pour $(\Sch/S)_\proetale$.
Si $U \to S$, $V \to U$ et $W \to U$ sont faiblement \'etales, alors
$V \times_U W \to S$ l'est aussi (voir
Compléments sur les morphismes, section \ref{more-morphisms-section-weakly-etale});
on obtient donc les produits fibrés dans $S_\proetale$.
Si $U, V, W$ sont affines, alors $V \times_U W$ l'est aussi; on obtient donc
les produits fibrés dans $(\textit{Aff}/S)_\proetale$.

\medskip\noindent
Soient $a, b : U \to V$ deux morphismes de $\Sch_\proetale$.
Dans ce cas, l'égalisateur de $a$ et $b$ (dans la catégorie des schémas) est
$$
V
\times_{\Delta_{V/\Spec(\mathbf{Z})}, V \times_{\Spec(\mathbf{Z})} V, (a, b)}
(U \times_{\Spec(\mathbf{Z})} U)
$$
qui est un objet de $\Sch_\proetale$ d'après ce que nous venons de voir.
Ainsi, $\Sch_\proetale$ possède des égalisateurs. Si $a$ et $b$ sont des morphismes au-dessus de $S$,
alors l'égalisateur (dans la catégorie des schémas) est aussi donné par
$$
V \times_{\Delta_{V/S}, V \times_S V, (a, b)} (U \times_S U)
$$
et l'on voit donc que $(\Sch/S)_\proetale$ possède des égalisateurs. De plus, si
$U$ et $V$ sont faiblement \'etales sur $S$, alors l'égalisateur ci-dessus
l'est aussi, puisqu'il est un produit fibré de schémas faiblement \'etales sur $S$.
Ainsi, $S_\proetale$ possède des égalisateurs. Les assertions sur les objets finaux
sont claires.
\end{proof}

\noindent
Vérifions ensuite que le grand site pro-\'etale affine définit le même
topos que le grand site pro-\'etale.

\begin{lemma}
\label{proetale-lemma-affine-big-site-proetale}
Soit $S$ un schéma. Soit $\Sch_\proetale$ un grand site pro-\'etale
contenant $S$.
Le foncteur $(\textit{Aff}/S)_\proetale \to (\Sch/S)_\proetale$
est un foncteur spécial cocontinu. Il induit donc une équivalence
de topos de $\Sh((\textit{Aff}/S)_\proetale)$ vers
$\Sh((\Sch/S)_\proetale)$.
\end{lemma}

\begin{proof}
La notion de foncteur spécial cocontinu est introduite dans
Sites, définition \ref{sites-definition-special-cocontinuous-functor}.
Il faut donc vérifier les hypothèses (1) -- (5) de
Sites, lemme \ref{sites-lemma-equivalence}.
Notons le foncteur d'inclusion
$u : (\textit{Aff}/S)_\proetale \to (\Sch/S)_\proetale$.
La cocontinuité signifie simplement que tout recouvrement pro-\'etale de
$T/S$, avec $T$ affine, peut être raffiné par un recouvrement pro-\'etale
standard de $T$. C'est le contenu du
lemme \ref{proetale-lemma-proetale-affine}.
Ainsi, (1) est vérifiée. Le foncteur $u$ est continu simplement parce qu'un recouvrement
pro-\'etale standard est un recouvrement pro-\'etale. Ainsi, (2) est vérifiée.
Les parties (3) et (4) découlent immédiatement du fait que $u$ est
pleinement fidèle. Enfin, la condition (5) résulte du
fait que tout schéma admet un recouvrement ouvert affine.
\end{proof}

\begin{lemma}
\label{proetale-lemma-put-in-T}
Soit $\Sch_\proetale$ un grand site pro-\'etale.
Soit $f : T \to S$ un morphisme de $\Sch_\proetale$.
Le foncteur $T_\proetale \to (\Sch/S)_\proetale$
est cocontinu et induit un morphisme de topos
$$
i_f :
\Sh(T_\proetale)
\longrightarrow
\Sh((\Sch/S)_\proetale)
$$
Pour tout faisceau $\mathcal{G}$ sur $(\Sch/S)_\proetale$,
on a la formule $(i_f^{-1}\mathcal{G})(U/T) = \mathcal{G}(U/S)$.
Le foncteur $i_f^{-1}$ admet aussi un adjoint à gauche $i_{f, !}$ qui commute
aux produits fibrés et aux égalisateurs.
\end{lemma}

\begin{proof}
Notons $u : T_\proetale \to (\Sch/S)_\proetale$ ce foncteur.
Autrement dit, étant donné un morphisme faiblement \'etale $j : U \to T$ correspondant
à un objet de $T_\proetale$, on pose $u(U \to T) = (f \circ j : U \to S)$.
Ce foncteur commute aux produits fibrés ; voir le
lemme \ref{proetale-lemma-fibre-products-proetale}.
De plus, $T_\proetale$ possède des égalisateurs et $u$ commute avec eux
d'après le lemme \ref{proetale-lemma-fibre-products-proetale}.
Il est manifestement cocontinu.
Il est aussi continu, car $u$ transforme les recouvrements en recouvrements et
commute aux produits fibrés. Le lemme résulte donc de
Sites, lemmes \ref{sites-lemma-when-shriek}
et \ref{sites-lemma-preserve-equalizers}.
\end{proof}

\begin{lemma}
\label{proetale-lemma-at-the-bottom}
Soit $S$ un schéma. Soit $\Sch_\proetale$ un grand site pro-\'etale
contenant $S$.
Le foncteur d'inclusion $S_\proetale \to (\Sch/S)_\proetale$
satisfait aux hypothèses de Sites, lemme \ref{sites-lemma-bigger-site},
et induit donc un morphisme de sites
$$
\pi_S : (\Sch/S)_\proetale \longrightarrow S_\proetale
$$
et un morphisme de topos
$$
i_S : \Sh(S_\proetale) \longrightarrow \Sh((\Sch/S)_\proetale)
$$
tels que $\pi_S \circ i_S = \text{id}$. De plus, $i_S = i_{\text{id}_S}$,
où $i_{\text{id}_S}$ est défini au lemme \ref{proetale-lemma-put-in-T}. En particulier, le
foncteur $i_S^{-1} = \pi_{S, *}$ est décrit par la règle
$i_S^{-1}(\mathcal{G})(U/S) = \mathcal{G}(U/S)$.
\end{lemma}

\begin{proof}
Dans ce cas, le foncteur $u : S_\proetale \to (\Sch/S)_\proetale$,
outre les propriétés établies dans la démonstration du
lemme \ref{proetale-lemma-put-in-T} ci-dessus, est pleinement fidèle
et envoie l'objet final sur l'objet final.
Le lemme découle de Sites, lemme \ref{sites-lemma-bigger-site}.
\end{proof}

\begin{definition}
\label{proetale-definition-restriction-small-proetale}
Dans la situation du
lemme \ref{proetale-lemma-at-the-bottom},
le foncteur $i_S^{-1} = \pi_{S, *}$ est souvent
appelé la {\it restriction au petit site pro-\'etale} ; pour un faisceau
$\mathcal{F}$ sur le grand site pro-\'etale, on note
$\mathcal{F}|_{S_\proetale}$ cette restriction.
\end{definition}

\noindent
Avec cette notation, on a, pour un faisceau $\mathcal{F}$ sur le
grand site et un faisceau $\mathcal{G}$ sur le petit site,
\begin{align*}
\Mor_{\Sh(S_\proetale)}(\mathcal{F}|_{S_\proetale}, \mathcal{G})
& =
\Mor_{\Sh((\Sch/S)_\proetale)}(\mathcal{F},
i_{S, *}\mathcal{G}) \\
\Mor_{\Sh(S_\proetale)}(\mathcal{G}, \mathcal{F}|_{S_\proetale})
& =
\Mor_{\Sh((\Sch/S)_\proetale)}(\pi_S^{-1}\mathcal{G}, \mathcal{F})
\end{align*}
De plus, $(i_{S, *}\mathcal{G})|_{S_\proetale} = \mathcal{G}$
et $(\pi_S^{-1}\mathcal{G})|_{S_\proetale} = \mathcal{G}$.

\begin{lemma}
\label{proetale-lemma-morphism-big}
Soit $\Sch_\proetale$ un grand site pro-\'etale.
Soit $f : T \to S$ un morphisme de $\Sch_\proetale$.
Le foncteur
$$
u : (\Sch/T)_\proetale \longrightarrow (\Sch/S)_\proetale, \quad
V/T \longmapsto V/S
$$
est cocontinu et admet un adjoint à droite continu
$$
v : (\Sch/S)_\proetale \longrightarrow (\Sch/T)_\proetale, \quad
(U \to S) \longmapsto (U \times_S T \to T).
$$
Ils induisent le même morphisme de topos
$$
f_{big} :
\Sh((\Sch/T)_\proetale)
\longrightarrow
\Sh((\Sch/S)_\proetale)
$$
On a $f_{big}^{-1}(\mathcal{G})(U/T) = \mathcal{G}(U/S)$.
On a $f_{big, *}(\mathcal{F})(U/S) = \mathcal{F}(U \times_S T/T)$.
De plus, $f_{big}^{-1}$ admet un adjoint à gauche $f_{big!}$ qui commute aux
produits fibrés et aux égalisateurs.
\end{lemma}

\begin{proof}
Le foncteur $u$ est cocontinu, continu et commute aux produits fibrés
et aux égalisateurs (détails omis ; comparer avec la démonstration du
lemme \ref{proetale-lemma-put-in-T}). Par conséquent,
Sites, lemmes \ref{sites-lemma-when-shriek} et
\ref{sites-lemma-preserve-equalizers}
s'appliquent, et l'on en déduit la formule
pour $f_{big}^{-1}$ ainsi que l'existence de $f_{big!}$. De plus,
le foncteur $v$ est adjoint à droite, car, pour $U/T$ et $V/S$,
on a $\Mor_S(u(U), V) = \Mor_T(U, V \times_S T)$,
ce qui est l'assertion voulue. On peut donc appliquer
Sites, lemmes \ref{sites-lemma-have-functor-other-way} et
\ref{sites-lemma-have-functor-other-way-morphism}
pour obtenir la formule pour $f_{big, *}$.
\end{proof}

\begin{lemma}
\label{proetale-lemma-morphism-big-small}
Soit $\Sch_\proetale$ un grand site pro-\'etale.
Soit $f : T \to S$ un morphisme de $\Sch_\proetale$.
\begin{enumerate}
\item On a $i_f = f_{big} \circ i_T$, où $i_f$ est défini au
lemme \ref{proetale-lemma-put-in-T} et $i_T$ au
lemme \ref{proetale-lemma-at-the-bottom}.
\item Le foncteur $S_\proetale \to T_\proetale$,
$(U \to S) \mapsto (U \times_S T \to T)$ est continu et induit
un morphisme de topos
$$
f_{small} : \Sh(T_\proetale) \longrightarrow \Sh(S_\proetale).
$$
On a $f_{small, *}(\mathcal{F})(U/S) = \mathcal{F}(U \times_S T/T)$.
\item On a un diagramme commutatif de morphismes de sites
$$
\xymatrix{
T_\proetale \ar[d]_{f_{small}} &
(\Sch/T)_\proetale \ar[d]^{f_{big}} \ar[l]^{\pi_T}\\
S_\proetale &
(\Sch/S)_\proetale \ar[l]_{\pi_S}
}
$$
de sorte que $f_{small} \circ \pi_T = \pi_S \circ f_{big}$ comme morphismes de topos.
\item On a $f_{small} = \pi_S \circ f_{big} \circ i_T = \pi_S \circ i_f$.
\end{enumerate}
\end{lemma}

\begin{proof}
L'égalité $i_f = f_{big} \circ i_T$ résulte de
l'égalité $i_f^{-1} = i_T^{-1} \circ f_{big}^{-1}$, qui ressort
clairement des descriptions de ces foncteurs ci-dessus.
Cela démontre (1).

\medskip\noindent
Le foncteur $u : S_\proetale \to T_\proetale$,
$u(U \to S) = (U \times_S T \to T)$,
transforme les recouvrements en recouvrements et commute aux produits fibrés ;
voir les lemmes \ref{proetale-lemma-proetale} et \ref{proetale-lemma-fibre-products-proetale}.
De plus, $S_\proetale$ et $T_\proetale$ ont tous deux un objet final,
à savoir $S/S$ et $T/T$, et $u(S/S) = T/T$. D'après
Sites, proposition \ref{sites-proposition-get-morphism},
le foncteur $u$ correspond donc à un morphisme de sites
$T_\proetale \to S_\proetale$. Celui-ci donne à son tour le
morphisme de topos ; voir
Sites, lemme \ref{sites-lemma-morphism-sites-topoi}. La description
de l'image directe résulte clairement de ces références.

\medskip\noindent
La partie (3) résulte du fait que $\pi_S$ et $\pi_T$ sont donnés par les
foncteurs d'inclusion, et $f_{small}$ et $f_{big}$ par les
foncteurs de changement de base $U \mapsto U \times_S T$.

\medskip\noindent
L'assertion (4) résulte de (3) en précomposant par $i_T$.
\end{proof}

\noindent
Dans la situation du lemme, avec la terminologie de la
définition \ref{proetale-definition-restriction-small-proetale},
on a, pour $\mathcal{F}$ un faisceau sur le grand site pro-\'etale de $T$,
\begin{equation}
\label{proetale-equation-compare-big-small}
(f_{big, *}\mathcal{F})|_{S_\proetale} =
f_{small, *}(\mathcal{F}|_{T_\proetale}),
\end{equation}
Cette égalité résulte clairement de la commutativité du diagramme de
sites du lemme, puisque la restriction au petit site pro-\'etale de
$T$, resp.\ $S$, est donnée par $\pi_{T, *}$, resp.\ $\pi_{S, *}$. Une formule
analogue faisant intervenir les images inverses et les restrictions est fausse.

\begin{lemma}
\label{proetale-lemma-composition-proetale}
Étant donnés des schémas $X$, $Y$, $Y$ dans $\Sch_\proetale$
et des morphismes $f : X \to Y$, $g : Y \to Z$, on a
$g_{big} \circ f_{big} = (g \circ f)_{big}$ et
$g_{small} \circ f_{small} = (g \circ f)_{small}$.
\end{lemma}

\begin{proof}
Cela découle de la description explicite des foncteurs image directe
et image inverse sur les grands sites donnée dans le
lemme \ref{proetale-lemma-morphism-big}. Pour les foncteurs
sur les petits sites, cela découle de la description des
foncteurs image directe du lemme \ref{proetale-lemma-morphism-big-small}.
\end{proof}

\begin{lemma}
\label{proetale-lemma-morphism-big-small-cartesian-diagram}
Soit $\Sch_\proetale$ un grand site pro-\'etale. Considérons un diagramme cartésien
$$
\xymatrix{
T' \ar[r]_{g'} \ar[d]_{f'} & T \ar[d]^f \\
S' \ar[r]^g & S
}
$$
dans $\Sch_\proetale$. Alors
$i_g^{-1} \circ f_{big, *} = f'_{small, *} \circ (i_{g'})^{-1}$
et $g_{big}^{-1} \circ f_{big, *} = f'_{big, *} \circ (g'_{big})^{-1}$.
\end{lemma}

\begin{proof}
Puisque le diagramme est cartésien, on a, pour $U'/S'$,
$U' \times_{S'} T' = U' \times_S T$. Ainsi,
$i_g^{-1} \circ f_{big, *}$ et $f'_{small, *} \circ (i_{g'})^{-1}$
associent tous deux à un faisceau $\mathcal{F}$ sur $(\Sch/T)_\proetale$ le faisceau
$U' \mapsto \mathcal{F}(U' \times_{S'} T')$ sur $S'_\proetale$
(d'après les lemmes \ref{proetale-lemma-put-in-T} et \ref{proetale-lemma-morphism-big}).
La seconde égalité se démontre de la même manière ou peut se
déduire du résultat très général suivant :
Sites, lemme \ref{sites-lemma-localize-morphism}.
\end{proof}

\noindent
On peut considérer un faisceau sur le grand site pro-\'etale de $S$ comme une famille
de faisceaux sur les petits sites pro-\'etales des schémas au-dessus de $S$.

\begin{lemma}
\label{proetale-lemma-characterize-sheaf-big}
Soit $S$ un schéma contenu dans un grand site pro-\'etale $\Sch_\proetale$.
Un faisceau $\mathcal{F}$ sur le grand site pro-\'etale $(\Sch/S)_\proetale$
est déterminé par les données suivantes :
\begin{enumerate}
\item pour tout $T/S \in \Ob((\Sch/S)_\proetale)$, un faisceau
$\mathcal{F}_T$ sur $T_\proetale$,
\item pour tout $f : T' \to T$ dans
$(\Sch/S)_\proetale$, un morphisme
$c_f : f_{small}^{-1}\mathcal{F}_T \to \mathcal{F}_{T'}$.
\end{enumerate}
Ces données sont soumises aux conditions suivantes :
\begin{enumerate}
\item[(a)] pour tous $f : T' \to T$ et $g : T'' \to T'$ dans
$(\Sch/S)_\proetale$, le composé
$c_g \circ g_{small}^{-1}c_f$ est égal à $c_{f \circ g}$, et
\item[(b)] si $f : T' \to T$ dans $(\Sch/S)_\proetale$
est faiblement \'etale, alors $c_f$ est un isomorphisme.
\end{enumerate}
\end{lemma}

\begin{proof}
La démonstration est identique à celle du
Topologies, lemme \ref{topologies-lemma-characterize-sheaf-big-etale}.
\end{proof}

\begin{lemma}
\label{proetale-lemma-alternative}
Soit $S$ un schéma. Notons $S_{affine, \proetale}$ la sous-catégorie pleine
de $S_\proetale$ formée des objets affines. Un recouvrement de
$S_{affine, \proetale}$ sera un recouvrement pro-\'etale standard ; voir la
définition \ref{proetale-definition-standard-proetale}.
Alors la restriction
$$
\mathcal{F} \longmapsto \mathcal{F}|_{S_{affine, \etale}}
$$
définit une équivalence de topos
$\Sh(S_\proetale) \cong \Sh(S_{affine, \proetale})$.
\end{lemma}

\begin{proof}
On peut le montrer directement à partir des définitions, et c'est un bon exercice.
Mais cela découle aussi immédiatement de
Sites, lemme \ref{sites-lemma-equivalence},
en vérifiant que le foncteur d'inclusion
$S_{affine, \proetale} \to S_\proetale$
est un foncteur spécial cocontinu (voir
Sites, définition \ref{sites-definition-special-cocontinuous-functor}).
\end{proof}

\begin{lemma}
\label{proetale-lemma-affine-alternative}
Soit $S$ un schéma affine. Notons $S_{app}$ la sous-catégorie pleine
de $S_\proetale$ formée des objets affines $U$ tels que
$\mathcal{O}(S) \to \mathcal{O}(U)$ soit ind-\'etale. Un recouvrement de
$S_{app}$ sera un recouvrement pro-\'etale standard ; voir la
définition \ref{proetale-definition-standard-proetale}.
Alors la restriction
$$
\mathcal{F} \longmapsto \mathcal{F}|_{S_{app}}
$$
définit une équivalence de topos $\Sh(S_\proetale) \cong \Sh(S_{app})$.
\end{lemma}

\begin{proof}
D'après le lemme \ref{proetale-lemma-alternative}, on peut remplacer $S_\proetale$ par
$S_{affine, \proetale}$.
Le lemme découle de Sites, lemme \ref{sites-lemma-equivalence},
en vérifiant que le foncteur d'inclusion $S_{app} \to S_{affine, \proetale}$
est un foncteur spécial cocontinu ; voir
Sites, définition \ref{sites-definition-special-cocontinuous-functor}.
Les conditions énoncées dans Sites, lemme \ref{sites-lemma-equivalence},
découlent immédiatement de la définition et des faits suivants :
(a) tout objet $U$ de $S_{affine, \proetale}$ admet un recouvrement
$\{V \to U\}$ avec $V$ ind-\'etale sur $X$
(proposition \ref{proetale-proposition-weakly-etale})
et (b) le foncteur $u$ est pleinement fidèle.
\end{proof}

\begin{lemma}
\label{proetale-lemma-proetale-subcanonical}
Soit $S$ un schéma. La topologie de chacun des sites pro-\'etales
$\Sch_\proetale$, $S_\proetale$, $(\Sch/S)_\proetale$,
$S_{affine, \proetale}$ et $(\textit{Aff}/S)_\proetale$ est sous-canonique.
\end{lemma}

\begin{proof}
Il suffit de combiner le lemme \ref{proetale-lemma-recognize-proetale-covering} avec
Descente, lemme \ref{descent-lemma-fpqc-universal-effective-epimorphisms}.
\end{proof}





\section{Objets faiblement contractiles}
\label{proetale-section-weakly-contractible}

\noindent
Dans cette section, nous démontrons le fait essentiel que nos sites pro-\'etales
contiennent de nombreux objets faiblement contractiles. En fait, la démonstration du
lemme \ref{proetale-lemma-get-many-weakly-contractible} explique la forme de la fonction
$Bound$ dans la
définition \ref{proetale-definition-big-proetale-site} (bien que, pour les lecteurs
qui ne se préoccupent pas des questions ensemblistes, cette information
soit sans contenu).

\noindent
Commençons par exprimer la notion d'anneau w-contractile en termes de
recouvrements pro-\'etales.

\begin{lemma}
\label{proetale-lemma-w-contractible-proetale-cover}
Soit $T = \Spec(A)$ un schéma affine. Les conditions suivantes sont équivalentes :
\begin{enumerate}
\item $A$ est w-contractile ;
\item tout recouvrement pro-\'etale de $T$ peut être raffiné par
un recouvrement de Zariski de la forme $T = \coprod_{i = 1, \ldots, n} U_i$.
\end{enumerate}
\end{lemma}

\begin{proof}
Supposons $A$ w-contractile. D'après le lemme \ref{proetale-lemma-proetale-affine},
il suffit de montrer que tout recouvrement pro-\'etale standard
$\{f_i : T_i \to T\}_{i = 1, \ldots, n}$ peut être raffiné par un recouvrement de Zariski de $T$.
Le morphisme $\coprod T_i \to T$ est un morphisme faiblement \'etale et surjectif
de schémas affines. D'après la définition \ref{proetale-definition-w-contractible},
il existe donc un morphisme $\sigma : T \to \coprod T_i$ au-dessus de $T$.
Alors le recouvrement de Zariski $T = \coprod \sigma^{-1}(T_i)$
raffine $\{f_i : T_i \to T\}$.

\medskip\noindent
Réciproquement, supposons (2). Si $A \to B$ est fidèlement plat et faiblement \'etale,
alors $\{\Spec(B) \to T\}$ est un recouvrement pro-\'etale.
Il existe donc un recouvrement de Zariski $T = \coprod U_i$
et des morphismes $U_i \to \Spec(B)$ au-dessus de $T$. Puisque $T = \coprod U_i$,
on obtient $T \to \Spec(B)$, c'est-à-dire un homomorphisme de $A$-algèbres $B \to A$.
Cela signifie que $A$ est w-contractile.
\end{proof}

\begin{lemma}
\label{proetale-lemma-w-contractible-is-weakly-contractible}
Soit $\Sch_\proetale$ un grand site pro-\'etale comme dans la
définition \ref{proetale-definition-big-proetale-site}.
Soit $T = \Spec(A)$ un objet affine de $\Sch_\proetale$.
Les conditions suivantes sont équivalentes :
\begin{enumerate}
\item $A$ est w-contractile,
\item $T$ est un objet faiblement contractile
(Sites, définition \ref{sites-definition-w-contractible})
de $\Sch_\proetale$, et
\item tout recouvrement pro-\'etale de $T$ peut être raffiné par
un recouvrement de Zariski de la forme $T = \coprod_{i = 1, \ldots, n} U_i$.
\end{enumerate}
\end{lemma}

\begin{proof}
Nous avons établi l'équivalence de (1) et (3) dans le
lemme \ref{proetale-lemma-w-contractible-proetale-cover}.

\medskip\noindent
Supposons (3), et soit $\mathcal{F} \to \mathcal{G}$ une surjection de faisceaux sur
$\Sch_\proetale$. Soit $s \in \mathcal{G}(T)$. Pour démontrer (2), montrons que
$s$ appartient à l'image de $\mathcal{F}(T) \to \mathcal{G}(T)$. On peut trouver un
recouvrement $\{T_i \to T\}$ de $\Sch_\proetale$ tel que $s$ se relève
en une section de $\mathcal{F}$ sur $T_i$
(Sites, définition \ref{sites-definition-sheaves-injective-surjective}).
D'après (3), on peut supposer donné un recouvrement fini
$T = \coprod_{j = 1, \ldots, m} U_j$ par des sous-ensembles ouverts et fermés
et des $t_j \in \mathcal{F}(U_j)$ dont l'image est $s|_{U_j}$.
Puisque les recouvrements de Zariski sont des recouvrements de $\Sch_\proetale$
(lemme \ref{proetale-lemma-etale-proetale}), on en conclut que
$\mathcal{F}(T) = \prod \mathcal{F}(U_j)$.
Ainsi, $t = (t_1, \ldots, t_m) \in \mathcal{F}(T)$
est une section dont l'image est $s$.

\medskip\noindent
Supposons (2). Soit $A \to D$ comme dans la
proposition \ref{proetale-proposition-find-w-contractible}.
Alors $\{V \to T\}$ est un recouvrement de $\Sch_\proetale$.
(Notons que $V = \Spec(D)$ est un objet de $\Sch_\proetale$
d'après la remarque \ref{proetale-remark-size-w-contractible},
combinée avec notre choix de la fonction
$Bound$ dans la définition \ref{proetale-definition-big-proetale-site}
et le calcul de la taille des schémas affines dans
Ensembles, lemme \ref{sets-lemma-bound-size}.)
Puisque la topologie de $\Sch_\proetale$ est sous-canonique
(lemme \ref{proetale-lemma-proetale-subcanonical}),
$h_V \to h_T$ est un morphisme surjectif de faisceaux
(Sites, lemme \ref{sites-lemma-covering-surjective-after-sheafification}).
Puisque $T$ est supposé faiblement contractile, il existe un élément
$f \in h_V(T) = \Mor(T, V)$ dont l'image dans $h_T(T)$ est $\text{id}_T$.
Ainsi, $A \to D$ admet une rétraction $\sigma : D \to A$.
Si maintenant $A \to B$ est fidèlement plat et faiblement \'etale, alors
$D \to D \otimes_A B$ possède les mêmes propriétés ; il existe donc
une rétraction $D \otimes_A B \to D$ qui, composée avec
$\sigma$, fournit une rétraction $B \to D \otimes_A B \to D \to A$
de $A \to B$. Ainsi, $A$ est w-contractile et (1) est vérifiée.
\end{proof}

\begin{lemma}
\label{proetale-lemma-get-many-weakly-contractible}
Soit $\Sch_\proetale$ un grand site pro-\'etale comme dans la
définition \ref{proetale-definition-big-proetale-site}.
Pour tout objet $T$ de $\Sch_\proetale$, il existe
un recouvrement $\{T_i \to T\}$ dans $\Sch_\proetale$
tel que chaque $T_i$ soit affine et soit le spectre d'un anneau
w-contractile. En particulier, $T_i$ est faiblement contractile dans $\Sch_\proetale$.
\end{lemma}

\begin{proof}
Pour les lecteurs qui ne se préoccupent pas des questions ensemblistes,
ce lemme est une conséquence immédiate du
lemme \ref{proetale-lemma-w-contractible-is-weakly-contractible} et de la
proposition \ref{proetale-proposition-find-w-contractible}.
Voici les détails.
Choisissons un recouvrement ouvert affine $T = \bigcup U_i$. Écrivons $U_i = \Spec(A_i)$.
Choisissons des homomorphismes d'anneaux fidèlement plats et ind-\'etales $A_i \to D_i$
tels que $D_i$ soit w-contractile comme dans la
proposition \ref{proetale-proposition-find-w-contractible}.
La famille de morphismes $\{\Spec(D_i) \to T\}$ est un
recouvrement pro-\'etale.
Si l'on peut montrer que $\Spec(D_i)$ est isomorphe à un objet, disons $T_i$,
de $\Sch_\proetale$, alors $\{T_i \to T\}$ sera combinatoirement
équivalente à un recouvrement de $\Sch_\proetale$ par la construction
de $\Sch_\proetale$ dans la définition \ref{proetale-definition-big-proetale-site}
et, plus précisément, par l'application de
Ensembles, lemme \ref{sets-lemma-coverings-site}, à la dernière étape.
Pour démontrer que $\Spec(D_i)$ est isomorphe à un objet de
$\Sch_\proetale$, il suffit de montrer que
$|D_i| \leq Bound(\text{size}(T))$ par la construction
de $\Sch_\proetale$ dans la définition \ref{proetale-definition-big-proetale-site}
et, plus précisément, par l'application de
Ensembles, lemme \ref{sets-lemma-construct-category}, à l'étape (3).
Puisque $|A_i| \leq \text{size}(U_i) \leq \text{size}(T)$
d'après Ensembles, lemmes \ref{sets-lemma-bound-affine} et
\ref{sets-lemma-bound-finite-type}, on obtient
$|D_i| \leq \kappa^{2^{2^{2^\kappa}}}$, où $\kappa = \text{size}(T)$,
d'après la remarque \ref{proetale-remark-size-w-contractible}.
Le résultat découle donc de notre choix de la fonction $Bound$ dans la
définition \ref{proetale-definition-big-proetale-site}.
\end{proof}

\begin{lemma}
\label{proetale-lemma-proetale-enough-w-contractible}
Soit $S$ un schéma. Les sites pro-\'etales
$S_\proetale$, $(\Sch/S)_\proetale$, $S_{affine, \proetale}$,
$(\textit{Aff}/S)_\proetale$, ainsi que, lorsque le schéma $S$ est affine, $S_{app}$,
possèdent suffisamment d'objets quasi-compacts, affines et faiblement contractiles ;
voir Sites, définition \ref{sites-definition-w-contractible}.
\end{lemma}

\begin{proof}
Cela découle immédiatement du lemme \ref{proetale-lemma-get-many-weakly-contractible}.
\end{proof}

\begin{lemma}
\label{proetale-lemma-weakly-contractible-cover}
Soit $S$ un schéma. Les sites pro-\'etales
$\Sch_\proetale$, $S_\proetale$ et $(\Sch/S)_\proetale$
possèdent la propriété suivante : pour tout objet
$U$, il existe un recouvrement $\{V \to U\}$ où $V$ est un
objet faiblement contractile. Si $U$ est quasi-compact, on peut
choisir $V$ affine et faiblement contractile.
\end{lemma}

\begin{proof}
Supposons que $V = \coprod_{j \in J} V_j$ soit un objet de $(\Sch/S)_\proetale$
qui soit la somme disjointe d'objets faiblement contractiles $V_j$.
Puisqu'une décomposition en somme disjointe est un recouvrement pro-\'etale,
on a $\mathcal{F}(V) = \prod_{j \in J} \mathcal{F}(V_j)$ pour tout
faisceau pro-\'etale $\mathcal{F}$. Soit $\mathcal{F} \to \mathcal{G}$
un morphisme surjectif de faisceaux d'ensembles. Puisque $V_j$ est faiblement contractile,
le morphisme $\mathcal{F}(V_j) \to \mathcal{G}(V_j)$ est surjectif ; voir
Sites, définition \ref{sites-definition-w-contractible}.
Ainsi, $\mathcal{F}(V) \to \mathcal{G}(V)$ est surjectif, car c'est un produit
de morphismes surjectifs d'ensembles ; on en conclut que $V$ est faiblement contractile.

\medskip\noindent
Choisissons un recouvrement $\{U_i \to U\}_{i \in I}$ où les $U_i$ sont affines et
faiblement contractiles, comme dans le lemme \ref{proetale-lemma-get-many-weakly-contractible}.
Prenons $V = \coprod_{i \in I} U_i$ (il y a ici une question ensembliste
que nous traiterons ci-dessous). Alors $\{V \to U\}$ est le recouvrement
pro-\'etale voulu par un objet faiblement contractile
(pour vérifier qu'il s'agit d'un recouvrement, on applique le
lemme \ref{proetale-lemma-recognize-proetale-covering}).
Si $U$ est quasi-compact, le
lemme \ref{proetale-lemma-recognize-proetale-covering}
montre immédiatement que l'on peut choisir un sous-ensemble fini $I' \subset I$ tel que
$\{U_i \to U\}_{i \in I'}$ soit encore un recouvrement ;
alors $\{\coprod_{i \in I'} U_i \to U\}$ est le recouvrement voulu
par un objet affine et faiblement contractile.

\medskip\noindent
Dans ce paragraphe, que nous conseillons vivement au lecteur de ne pas lire, nous traitons
les problèmes ensemblistes. Pour savoir que la somme disjointe
appartient à notre univers partiel, il faut borner le cardinal de
l'ensemble d'indices $I$. Il ressort immédiatement de la construction du
recouvrement $\{U_i \to U\}_{i \in I}$ dans la démonstration du
lemme \ref{proetale-lemma-get-many-weakly-contractible}
que $|I| \leq \text{size}(U)$, où la taille d'un schéma est celle
définie dans Ensembles, section \ref{sets-section-categories-schemes}.
De plus, pour tout $i$, on a
$\text{size}(U_i) \leq Bound(\text{size}(U))$;
cela résulte de la borne du cardinal de
$\Gamma(U_i, \mathcal{O}_{U_i})$ donnée dans
la démonstration du lemme \ref{proetale-lemma-get-many-weakly-contractible}
et d'Ensembles, lemme \ref{sets-lemma-bound-affine}.
Ainsi, $\text{size}(\coprod_{i \in I} U_i)) \leq Bound(\text{size}(U))$
d'après Ensembles, lemme \ref{sets-lemma-bound-size}.
D'après la construction du grand site pro-\'etale donnée dans
Ensembles, lemme \ref{sets-lemma-construct-category},
on voit donc que $\coprod_{i \in I} U_i$ est isomorphe à un objet
de notre site, ce qui achève la démonstration.
\end{proof}






\section{Hyperrecouvrements faiblement contractiles}
\label{proetale-section-hypercovering}

\noindent
Les résultats de la section \ref{proetale-section-weakly-contractible}
entraînent l'existence d'hyperrecouvrements constitués d'objets
faiblement contractiles.

\begin{lemma}
\label{proetale-lemma-w-contractible-hypercovering}
Soit $X$ un schéma.
\begin{enumerate}
\item Pour tout objet $U$ de $X_\proetale$, il existe un hyperrecouvrement
$K$ de $U$ dans $X_\proetale$ tel que chaque terme $K_n$ soit constitué d'un
seul objet faiblement contractile de $X_\proetale$ qui recouvre $U$.
\item Pour tout objet quasi-compact et quasi-séparé $U$ de $X_\proetale$,
il existe un hyperrecouvrement $K$ de $U$ dans $X_\proetale$ tel que chaque
terme $K_n$ soit constitué d'un seul objet affine et faiblement contractile de
$X_\proetale$ qui recouvre $U$.
\end{enumerate}
\end{lemma}

\begin{proof}
Soit $\mathcal{B} \subset \Ob(X_\proetale)$ l'ensemble des objets faiblement contractiles
de $X_\proetale$. Tout objet $T$ de $X_\proetale$ admet un
recouvrement $\{T_i \to T\}_{i \in I}$ où $I$ est fini et $T_i \in \mathcal{B}$, d'après le
lemme \ref{proetale-lemma-weakly-contractible-cover}.
D'après Hyperrecouvrements, lemme \ref{hypercovering-lemma-w-contractible},
on obtient un hyperrecouvrement $K$ de $U$ tel que $K_n = \{U_{n, i}\}_{i \in I_n}$,
où $I_n$ est fini et $U_{n, i}$ faiblement contractile.
On peut alors remplacer $K$ par l'hyperrecouvrement de $U$ donné par $\{U_n\}$
en degré $n$, où $U_n = \coprod_{i \in I_n} U_{n, i}$.
Cette opération est permise d'après Hyperrecouvrements,
remarque \ref{hypercovering-remark-take-unions-hypercovering-X}.

\medskip\noindent
Soit $X_{qcqs, \proetale} \subset X_\proetale$ la sous-catégorie pleine
formée des objets quasi-compacts et quasi-séparés.
Un recouvrement de $X_{qcqs, \proetale}$ sera un recouvrement pro-\'etale fini.
Alors $X_{qcqs, \proetale}$ est un site, admet des produits fibrés et
le foncteur d'inclusion $X_{qcqs, \proetale} \to X_\proetale$ est continu
et commute aux produits fibrés.
En particulier, si $K$ est un hyperrecouvrement d'un objet $U$ dans
$X_{qcqs, \proetale}$, alors $K$ est un hyperrecouvrement de $U$ dans $X_\proetale$
d'après Hyperrecouvrements, lemme
\ref{hypercovering-lemma-hypercovering-continuous-functor}.
Soit $\mathcal{B} \subset \Ob(X_{qcqs, \proetale})$ l'ensemble des objets
affines et faiblement contractiles. D'après le
lemme \ref{proetale-lemma-get-many-weakly-contractible}
et le fait qu'une somme disjointe finie de schémas affines est affine,
pour tout objet $U$ de $X_{qcqs, \proetale}$, il existe un recouvrement
$\{V \to U\}$ de $X_{qcqs, \proetale}$ tel que $V \in \mathcal{B}$.
D'après Hyperrecouvrements, lemme \ref{hypercovering-lemma-w-contractible},
on obtient un hyperrecouvrement $K$ de $U$ tel que $K_n = \{U_{n, i}\}_{i \in I_n}$,
où $I_n$ est fini et $U_{n, i}$ affine et faiblement contractile.
On peut alors remplacer $K$ par l'hyperrecouvrement de $U$
donné par $\{U_n\}$ en degré $n$, où
$U_n = \coprod_{i \in I_n} U_{n, i}$. Cette opération est permise d'après
Hyperrecouvrements, remarque \ref{hypercovering-remark-take-unions-hypercovering-X}.
\end{proof}

\noindent
Dans le lemme suivant, nous utilisons le complexe de {\v C}ech $s(\mathcal{F}(K))$
associé à un hyperrecouvrement $K$ dans un site. Voir
Hyperrecouvrements, section \ref{hypercovering-section-hyper-cech}.
Si $K$ est un hyperrecouvrement de $U$ et $K_n = \{U_n \to U\}$, alors
le complexe de {\v C}ech est de la forme suivante :
$$
s(\mathcal{F}(K)) =
\left(
\mathcal{F}(U_0) \to \mathcal{F}(U_1) \to \mathcal{F}(U_2) \to \ldots
\right)
$$
où $s(\mathcal{F}(U_n))$ est placé en degré cohomologique $n$.

\begin{lemma}
\label{proetale-lemma-compute-cohomology}
Soit $X$ un schéma. Soit $E \in D^+(X_\proetale)$ représenté par
un complexe borné inférieurement $\mathcal{E}^\bullet$ de faisceaux abéliens.
Soit $K$ un hyperrecouvrement de $U \in \Ob(X_\proetale)$ avec
$K_n = \{U_n \to U\}$, où $U_n$ est un objet faiblement contractile de
$X_\proetale$. Alors
$$
R\Gamma(U, E) = \text{Tot}(s(\mathcal{E}^\bullet(K)))
$$
dans $D(\textit{Ab})$.
\end{lemma}

\begin{proof}
Si $\mathcal{E}$ est un faisceau
abélien sur $X_\proetale$, alors la suite spectrale donnée par
Hyperrecouvrements, lemme \ref{hypercovering-lemma-cech-spectral-sequence},
implique que
$$
R\Gamma(X_\proetale, \mathcal{E}) = s(\mathcal{E}(K))
$$
car les groupes de cohomologie supérieure de tout faisceau sur $U_n$ s'annulent; voir
Cohomologie sur les sites, lemme \ref{sites-cohomology-lemma-w-contractible}.

\medskip\noindent
Si $\mathcal{E}^\bullet$ est borné inférieurement, on peut choisir une résolution
injective $\mathcal{E}^\bullet \to \mathcal{I}^\bullet$ et considérer
le morphisme de complexes
$$
\text{Tot}(s(\mathcal{E}^\bullet(K)))
\longrightarrow
\text{Tot}(s(\mathcal{I}^\bullet(K)))
$$
Pour tout $n$, le morphisme $\mathcal{E}^\bullet(U_n) \to \mathcal{I}^\bullet(U_n)$
est un quasi-isomorphisme, car le foncteur des sections sur $U_n$ est exact.
Le morphisme affiché est donc un quasi-isomorphisme d'après l'une des suites
spectrales fournies par
Homologie, lemme \ref{homology-lemma-first-quadrant-ss}.
En utilisant le résultat du premier paragraphe, on voit que pour tout $p$,
le complexe $s(\mathcal{I}^p(K))$ est acyclique en degrés $n > 0$ et
calcule $\mathcal{I}^p(U)$ en degré $0$. L'autre suite
spectrale fournie par
Homologie, lemme \ref{homology-lemma-first-quadrant-ss},
montre donc que $\text{Tot}(s(\mathcal{I}^\bullet(K)))$ calcule
$R\Gamma(U, E) = \mathcal{I}^\bullet(U)$.
\end{proof}

\begin{lemma}
\label{proetale-lemma-quasi-compact-quasi-separated-commutes-direct-sums}
Soit $X$ un schéma quasi-compact et quasi-séparé.
Le foncteur $R\Gamma(X, -) : D^+(X_\proetale) \to D(\textit{Ab})$
commute aux sommes directes et aux colimites homotopiques.
\end{lemma}

\begin{proof}
L'assertion signifie ce qui suit. Supposons donnée une famille d'objets
$E_i$ de $D^+(X_\proetale)$ telle que $\bigoplus E_i$ soit un objet
de $D^+(X_\proetale)$. Alors
$R\Gamma(X, \bigoplus E_i) = \bigoplus R\Gamma(X, E_i)$.
Pour le voir, choisissons un hyperrecouvrement $K$ de $X$ tel que $K_n = \{U_n \to X\}$,
où $U_n$ est un schéma affine et faiblement contractile; voir le
lemme \ref{proetale-lemma-w-contractible-hypercovering}.
Soit $N$ un entier tel que $H^p(E_i) = 0$ pour $p < N$.
Choisissons un complexe de faisceaux abéliens $\mathcal{E}_i^\bullet$
représentant $E_i$ tel que $\mathcal{E}_i^p = 0$ pour $p < N$.
La somme directe terme à terme $\bigoplus \mathcal{E}_i^\bullet$ représente
$\bigoplus E_i$ dans $D(X_\proetale)$; voir
Injectifs, lemme \ref{injectives-lemma-derived-products}.
D'après le lemme \ref{proetale-lemma-compute-cohomology}, on a
$$
R\Gamma(X, \bigoplus E_i) =
\text{Tot}(s((\bigoplus \mathcal{E}^\bullet_i)(K)))
$$
et
$$
R\Gamma(X, E_i) = \text{Tot}(s(\mathcal{E}^\bullet_i(K)))
$$
Puisque chaque $U_n$ est quasi-compact, on voit que
$$
\text{Tot}(s((\bigoplus \mathcal{E}^\bullet_i)(K))) =
\bigoplus \text{Tot}(s(\mathcal{E}^\bullet_i(K)))
$$
d'après Modules sur les sites, lemme
\ref{sites-modules-lemma-sections-over-quasi-compact}.
L'assertion relative aux colimites homotopiques est une conséquence formelle du fait
que $R\Gamma$ est un foncteur exact de catégories triangulées et du
fait (que nous venons de démontrer) qu'il commute aux sommes directes.
\end{proof}

\begin{remark}
\label{proetale-remark-extend-to-all}
Soit $X$ un schéma. Puisque $X_\proetale$ a suffisamment d'objets faiblement contractiles,
pour tout $K$ dans $D(X_\proetale)$, on a $K = R\lim \tau_{\geq -n}K$
d'après
Cohomologie sur les sites, proposition
\ref{sites-cohomology-proposition-enough-weakly-contractibles}.
Puisque $R\Gamma$ commute à $R\lim$ d'après
Injectifs, lemme \ref{injectives-lemma-RF-commutes-with-Rlim},
on voit que
$$
R\Gamma(X, K) = R\lim R\Gamma(X, \tau_{\geq -n}K)
$$
dans $D(\textit{Ab})$. Cela nous permettra parfois d'étendre à tous les complexes
des résultats établis pour les complexes bornés inférieurement.
\end{remark}







\section{Engendrement compact}
\label{proetale-section-compact-generation}

\noindent
Dans cette section, nous démontrons que diverses catégories dérivées associées
à nos sites pro-\'etales sont compactement engendrées au sens de
Catégories dérivées, définition \ref{derived-definition-compactly-generated}.

\begin{lemma}
\label{proetale-lemma-enough-compact-proetale}
Soit $S$ un schéma. Soit $\Lambda$ un anneau.
\begin{enumerate}
\item $D(S_\proetale)$ est compactement engendrée,
\item $D(S_\proetale, \Lambda)$ est compactement engendrée,
\item $D(S_\proetale, \mathcal{A})$ est compactement engendrée
pour tout faisceau d'anneaux $\mathcal{A}$ sur $S_\proetale$,
\item $D((\Sch/S)_\proetale)$ est compactement engendrée,
\item $D((\Sch/S)_\proetale, \Lambda)$ est compactement engendrée,
\item $D((\Sch/S)_\proetale, \mathcal{A})$ est compactement engendrée
pour tout faisceau d'anneaux $\mathcal{A}$ sur $(\Sch/S)_\proetale$.
\end{enumerate}
\end{lemma}

\begin{proof}
Démonstration de (3). Soit $U$ un objet de $S_\proetale$, affine et
faiblement contractile. Alors $j_{U!}\mathcal{A}_U$ est un objet compact
de la catégorie dérivée $D(S_\proetale, \mathcal{A})$; voir
Cohomologie sur les sites, lemme
\ref{sites-cohomology-lemma-quasi-compact-weakly-contractible-compact}.
Choisissons un ensemble $I$ et, pour chaque $i \in I$, un objet
$U_i$ de $S_\proetale$, affine et faiblement contractile, de telle sorte que tout objet affine et faiblement contractile
de $S_\proetale$ soit isomorphe à l'un des $U_i$. Cela est possible
car $\Ob(S_\proetale)$ est un ensemble. Pour achever la démonstration de (3), il suffit
de montrer que $\bigoplus j_{U_i, !}\mathcal{A}_{U_i}$ est un générateur
de $D(S_\proetale, \mathcal{A})$; voir
Catégories dérivées, définition \ref{derived-definition-generators}.
Pour le voir, soit $K$ un objet non nul
de $D(S_\proetale, \mathcal{A})$. Il existe alors un objet $T$
de notre site $S_\proetale$ et un élément non nul $\xi$ de $H^n(K)(T)$.
Autrement dit, $\xi$ est une section non nulle du $n$-ième faisceau de cohomologie
de $K$.
On peut supposer que $K$ est représenté par un complexe $\mathcal{K}^\bullet$
de faisceaux de $\mathcal{A}$-modules et que $\xi$ est la classe d'une
section $s \in \mathcal{K}^n(T)$ telle que $\text{d}(s) = 0$.
En effet, $\xi$ est localement représenté par la classe d'une section
(on obtient donc le résultat après avoir remplacé $T$ par un membre d'un recouvrement de $T$).
Choisissons ensuite un recouvrement $\{T_j \to T\}_{j \in J}$
comme dans le lemme \ref{proetale-lemma-get-many-weakly-contractible}.
Puisque $H^n(K)$ est un faisceau, on voit que, pour un certain $j$, la restriction
$\xi|_{T_j}$ reste non nulle. Ainsi, $s|_{T_j}$ définit un morphisme non nul
$j_{T_j, !}\mathcal{A}_{T_j} \to K$ dans $D(S_\proetale, \mathcal{A})$.
Comme $T_j \cong U_i$ pour un certain $i \in I$, on conclut.

\medskip\noindent
Le même argument s'applique mot pour mot au grand site pro-\'etale de $S$.
\end{proof}







\section{Comparaison des topologies}
\label{proetale-section-compare-topologies}

\noindent
Cette section est l'analogue de
Cohomologie \'etale, section \ref{etale-cohomology-section-compare-topologies}.

\begin{lemma}
\label{proetale-lemma-presheaf-value-weakly-contractible}
Soit $X$ un schéma. Soit $\mathcal{F}$ un préfaisceau d'ensembles sur $X_\proetale$
qui transforme les sommes disjointes finies en produits. Alors
$\mathcal{F}^\#(W) = \mathcal{F}(W)$ si $W$ est un objet affine et faiblement contractile
de $X_\proetale$.
\end{lemma}

\begin{proof}
Rappelons que $\mathcal{F}^\#$ est égal à $(\mathcal{F}^+)^+$ ; voir
Sites, théorème \ref{sites-theorem-plus}, où $\mathcal{F}^+$
est le préfaisceau qui associe à tout objet $U$ de $X_\proetale$
$\colim H^0(\mathcal{U}, \mathcal{F})$, où la colimite est prise
sur tous les recouvrements pro-\'etales $\mathcal{U}$ de $U$.
Il suffit donc de démontrer que (a) $\mathcal{F}^+$ transforme les sommes
disjointes finies en produits et que (b) sa valeur en $W$ est $\mathcal{F}(W)$.
Si $U = U_1 \amalg U_2$, alors, étant donné un recouvrement pro-\'etale
$\mathcal{U} = \{f_j : V_j \to U\}$ de $U$, on obtient des recouvrements pro-\'etales
$\mathcal{U}_i = \{f_j^{-1}(U_i) \to U_i\}$ et l'on a clairement
$$
H^0(\mathcal{U}, \mathcal{F}) =
H^0(\mathcal{U}_1, \mathcal{F}) \times
H^0(\mathcal{U}_2, \mathcal{F})
$$
car $\mathcal{F}$ transforme les sommes disjointes finies en produits (cela comprend
la condition selon laquelle $\mathcal{F}$ associe au schéma vide le singleton).
Cela démontre (a).
Enfin, tout recouvrement pro-\'etale de $W$ peut être raffiné par une décomposition
en somme disjointe finie $W = W_1 \amalg \ldots W_n$ d'après le
lemme \ref{proetale-lemma-w-contractible-is-weakly-contractible}.
Ainsi, $\mathcal{F}^+(W) = \mathcal{F}(W)$ précisément parce que la valeur de
$\mathcal{F}$ sur $W$ est le produit des valeurs de $\mathcal{F}$
sur les $W_j$. Cela démontre (b).
\end{proof}

\begin{lemma}
\label{proetale-lemma-small-pullback-weakly-contractible}
Soit $f : X \to Y$ un morphisme de schémas. Soit $\mathcal{F}$ un faisceau
d'ensembles sur $Y_\proetale$. Si $W$ est un objet affine et faiblement contractile
de $X_\proetale$, alors
$$
f_{small}^{-1}\mathcal{F}(W) = \colim_{W \to V} \mathcal{F}(V)
$$
où la colimite est prise sur les morphismes $W \to V$ au-dessus de $Y$
tels que $V \in Y_\proetale$.
\end{lemma}

\begin{proof}
Rappelons que $f_{small}^{-1}\mathcal{F}$ est le faisceau associé au
préfaisceau
$$
u_p\mathcal{F} : U \mapsto \colim_{U \to V} \mathcal{F}(V)
$$
sur $X_\etale$ ; voir Sites, sections \ref{sites-section-morphism-sites} et
\ref{sites-section-continuous-functors} ; nous avons omis de la notation le fait
que la colimite est prise sur la catégorie opposée de la catégorie
$\{U \to V, V \in Y_\proetale\}$. D'après le
lemme \ref{proetale-lemma-presheaf-value-weakly-contractible},
il suffit de démontrer que $u_p\mathcal{F}$ transforme les sommes disjointes finies
en produits.
Supposons que $U = U_1 \amalg U_2$ soit une somme disjointe de sous-schémas
ouverts et fermés. Il existe un foncteur
$$
\{U_1 \to V_1\} \times \{U_2 \to V_2\} \longrightarrow
\{U \to V\},\quad
(U_1 \to V_1, U_2 \to V_2) \longmapsto (U \to V_1 \amalg V_2)
$$
qui est initial (Catégories, définition \ref{categories-definition-initial}).
Le foncteur correspondant sur les catégories opposées est donc cofinal ;
d'après Catégories, lemme \ref{categories-lemma-cofinal}, on voit
que $u_p\mathcal{F}$ évalué en $U$ est la colimite des valeurs
$\mathcal{F}(V_1 \amalg V_2)$ indexées par la catégorie produit.
Puisque $\mathcal{F}$ est un faisceau, il transforme les sommes disjointes en produits et
on conclut qu'il en est de même pour $u_p\mathcal{F}$.
\end{proof}

\begin{lemma}
\label{proetale-lemma-describe-pullback}
Soit $S$ un schéma. Considérons le morphisme
$$
\pi_S : (\Sch/S)_\proetale \longrightarrow S_\proetale
$$
du lemme \ref{proetale-lemma-at-the-bottom}. Soit $\mathcal{F}$ un faisceau sur
$S_\proetale$. Alors $\pi_S^{-1}\mathcal{F}$ est donné par la règle
$$
(\pi_S^{-1}\mathcal{F})(T) = \Gamma(T_\proetale, f_{small}^{-1}\mathcal{F})
$$
où $f : T \to S$. De plus, $\pi_S^{-1}\mathcal{F}$ satisfait la
condition de faisceau relativement aux recouvrements fpqc.
\end{lemma}

\begin{proof}
Observons que l'on a un morphisme
$i_f : \Sh(T_\proetale) \to \Sh(\Sch/S)_\proetale)$
tel que $\pi_S \circ i_f = f_{small}$ en tant que morphismes
$T_\proetale \to S_\proetale$ ; voir le lemme \ref{proetale-lemma-put-in-T}.
Puisque l'image inverse est transitive, on voit que
$i_f^{-1} \pi_S^{-1}\mathcal{F} = f_{small}^{-1}\mathcal{F}$, comme souhaité.

\medskip\noindent
Soit $\{g_i : T_i \to T\}_{i \in I}$ un recouvrement fpqc. L'assertion finale
signifie ceci. Soient $\mathcal{G}$ un faisceau sur $T_\proetale$ et
des sections $s_i \in \Gamma(T_i, g_{i, small}^{-1}\mathcal{G})$ dont les images inverses
sur $T_i \times_T T_j$ coïncident, il existe une unique section $s$ de $\mathcal{G}$
sur $T$ dont l'image inverse sur $T_i$ coïncide avec $s_i$. Nous démontrerons cette
assertion lorsque $T$ est affine et que le recouvrement est donné par un unique
morphisme plat surjectif $T' \to T$ entre schémas affines et nous omettons la réduction
du cas général à ce cas.

\medskip\noindent
Soit $g : T' \to T$ un morphisme plat surjectif entre schémas affines et soit
$s' \in g_{small}^{-1}\mathcal{G}(T')$ une section telle que
$\text{pr}_0^*s' = \text{pr}_1^*s'$ sur $T' \times_T T'$.
Choisissons un morphisme surjectif faiblement \'etale $W \to T'$ où
$W$ est affine et faiblement contractile ; voir le
lemme \ref{proetale-lemma-weakly-contractible-cover}. D'après le
lemme \ref{proetale-lemma-small-pullback-weakly-contractible},
la restriction $s'|_W$ est un élément de $\colim_{W \to U} \mathcal{G}(U)$.
Choisissons $\phi : W \to U_0$ et $s_0 \in \mathcal{G}(U_0)$ correspondant à $s'$.
Choisissons un morphisme surjectif faiblement \'etale $V \to W \times_T W$ où
$V$ est affine et faiblement contractile.
Notons $a, b : V \to W$ les morphismes induits.
Puisque $a^*(s'|_W) = b^*(s'|_W)$ et que la catégorie
$\{V \to U, U \in T_\proetale\}$ est cofiltrante
(ce qui est clair ; voir néanmoins
Sites, lemme \ref{sites-lemma-directed-morphism}, en cas de doute),
on voit que les deux morphismes $\phi \circ a , \phi \circ b : V \to U_0$
doivent être égaux. D'après les résultats de
Descente, section \ref{descent-section-fpqc-universal-effective-epimorphisms}
(notamment
Descente, lemme \ref{descent-lemma-fpqc-universal-effective-epimorphisms}),
il existe donc un unique morphisme $T \to U_0$ tel que $\phi$
soit le composé de ce morphisme avec le morphisme structural $W \to T$
(nous omettons un petit détail). On peut alors prendre pour $s$ l'image inverse
de $s_0$ par ce morphisme. Nous omettons de vérifier que
$s$ a pour image inverse $s'$ sur $T'$.
\end{proof}









\section{Comparaison des grands et des petits topos}
\label{proetale-section-compare}

\noindent
Cette section est l'analogue de Cohomologie \'etale, section
\ref{etale-cohomology-section-compare}. Dans la suite, nous
noterons souvent $\mathcal{F} \mapsto \mathcal{F}|_{S_\proetale}$
le foncteur image inverse $i_S^{-1}$ correspondant au
morphisme de topos $i_S : \Sh(S_\proetale) \to \Sh((\Sch/S)_\proetale)$
du lemme \ref{proetale-lemma-at-the-bottom}.

\begin{lemma}
\label{proetale-lemma-compare-injectives}
Soit $S$ un schéma. Soit $T$ un objet de $(\Sch/S)_\proetale$.
\begin{enumerate}
\item Si $\mathcal{I}$ est injectif dans $\textit{Ab}((\Sch/S)_\proetale)$, alors
\begin{enumerate}
\item $i_f^{-1}\mathcal{I}$ est injectif dans $\textit{Ab}(T_\proetale)$,
\item $\mathcal{I}|_{S_\proetale}$ est injectif dans $\textit{Ab}(S_\proetale)$,
\end{enumerate}
\item Si $\mathcal{I}^\bullet$ est un complexe K-injectif
dans $\textit{Ab}((\Sch/S)_\proetale)$, alors
\begin{enumerate}
\item $i_f^{-1}\mathcal{I}^\bullet$ est un complexe K-injectif dans
$\textit{Ab}(T_\proetale)$,
\item $\mathcal{I}^\bullet|_{S_\proetale}$ est un complexe K-injectif dans
$\textit{Ab}(S_\proetale)$,
\end{enumerate}
\end{enumerate}
\end{lemma}

\begin{proof}
Démonstration de (1)(a) et (2)(a) : $i_f^{-1}$ est adjoint à droite
au foncteur exact $i_{f, !}$. En effet, rappelons que $i_f$ correspond
à un foncteur cocontinu $u : T_\proetale \to (\Sch/S)_\proetale$
qui est continu et commute aux produits fibrés et aux égalisateurs ; voir le
lemme \ref{proetale-lemma-put-in-T} et sa démonstration. On obtient donc $i_{f, !}$ d'après
Modules sur les sites, lemme \ref{sites-modules-lemma-g-shriek-adjoint}.
Il est montré dans Modules sur les sites, lemme
\ref{sites-modules-lemma-exactness-lower-shriek}, que ce foncteur est exact.
On conclut alors que (1)(a) et (2)(a) sont vérifiées d'après
Homologie, lemme \ref{homology-lemma-adjoint-preserve-injectives}, et
Catégories dérivées, lemme \ref{derived-lemma-adjoint-preserve-K-injectives}.

\medskip\noindent
Les parties (1)(b) et (2)(b) sont des cas particuliers de (1)(a) et (2)(a),
puisque $i_S = i_{\text{id}_S}$.
\end{proof}

\begin{lemma}
\label{proetale-lemma-compare-higher-direct-image}
Soit $f : T \to S$ un morphisme de schémas.
Pour $K$ dans $D((\Sch/T)_\proetale)$, on a
$$
(Rf_{big, *}K)|_{S_\proetale} = Rf_{small, *}(K|_{T_\proetale})
$$
dans $D(S_\proetale)$. Plus généralement, soit $S' \in \Ob((\Sch/S)_\proetale)$
de morphisme structural $g : S' \to S$. Considérons le produit fibré
$$
\xymatrix{
T' \ar[r]_{g'} \ar[d]_{f'} & T \ar[d]^f \\
S' \ar[r]^g & S
}
$$
Alors, pour $K$ dans $D((\Sch/T)_\proetale)$, on a
$$
i_g^{-1}(Rf_{big, *}K) = Rf'_{small, *}(i_{g'}^{-1}K)
$$
dans $D(S'_\proetale)$ et
$$
g_{big}^{-1}(Rf_{big, *}K) = Rf'_{big, *}((g'_{big})^{-1}K)
$$
dans $D((\Sch/S')_\proetale)$.
\end{lemma}

\begin{proof}
La première égalité découle du lemme \ref{proetale-lemma-compare-injectives}
et de (\ref{proetale-equation-compare-big-small}),
en choisissant un complexe K-injectif de faisceaux abéliens représentant $K$.
La deuxième égalité découle du lemme \ref{proetale-lemma-compare-injectives}
et du lemme
\ref{proetale-lemma-morphism-big-small-cartesian-diagram},
en choisissant un complexe K-injectif de faisceaux abéliens représentant $K$.
La troisième égalité découle de même de
Cohomologie sur les sites, lemmes \ref{sites-cohomology-lemma-cohomology-of-open} et
\ref{sites-cohomology-lemma-restrict-K-injective-to-open},
ainsi que du lemme \ref{proetale-lemma-morphism-big-small-cartesian-diagram},
en choisissant un complexe K-injectif de faisceaux abéliens représentant $K$.
\end{proof}

\noindent
Soit $S$ un schéma et soit $\mathcal{H}$ un faisceau abélien sur
$(\Sch/S)_\proetale$. Rappelons que $H^n_\proetale(U, \mathcal{H})$ désigne
la cohomologie de $\mathcal{H}$ sur un objet $U$ de $(\Sch/S)_\proetale$.

\begin{lemma}
\label{proetale-lemma-compare-cohomology-big-small}
Soit $f : T \to S$ un morphisme de schémas. Pour $K$ dans $D(S_\proetale)$,
on a
$$
H^n_\proetale(S, \pi_S^{-1}K) = H^n(S_\proetale, K)
$$
et
$$
H^n_\proetale(T, \pi_S^{-1}K) = H^n(T_\proetale, f_{small}^{-1}K).
$$
Pour $M$ dans $D((\Sch/S)_\proetale)$, on a
$$
H^n_\proetale(T, M) = H^n(T_\proetale, i_f^{-1}M).
$$
\end{lemma}

\begin{proof}
Pour démontrer la dernière égalité, représentons $M$
par un complexe K-injectif de faisceaux abéliens,
puis appliquons le lemme \ref{proetale-lemma-compare-injectives}
et explicitons les définitions. La deuxième égalité en découle,
car $i_f^{-1} \circ \pi_S^{-1} = f_{small}^{-1}$. La première
égalité est un cas particulier
de la deuxième.
\end{proof}

\begin{lemma}
\label{proetale-lemma-cohomological-descent-proetale}
Soit $S$ un schéma. Pour $K \in D(S_\proetale)$, le morphisme
$$
K \longrightarrow R\pi_{S, *}\pi_S^{-1}K
$$
est un isomorphisme.
\end{lemma}

\begin{proof}
Cela résulte du fait que $\pi_S^{-1}$ et $\pi_{S, *} = i_S^{-1}$ sont tous deux
des foncteurs exacts et que le foncteur composé $\pi_{S, *} \circ \pi_S^{-1}$
est le foncteur identité.
\end{proof}






















\section{Points du site pro-\'etale}
\label{proetale-section-points}

\noindent
Nous appliquons d'abord le critère de Deligne pour montrer qu'il y a suffisamment de points.

\begin{lemma}
\label{proetale-lemma-points-proetale}
Soit $S$ un schéma. Les sites pro-\'etales $\Sch_\proetale$,
$S_\proetale$, $(\Sch/S)_\proetale$, $S_{affine, \proetale}$ et
$(\textit{Aff}/S)_\proetale$ ont suffisamment de points.
\end{lemma}

\begin{proof}
Le grand topos pro-\'etale de $S$ est équivalent au topos défini par
$(\textit{Aff}/S)_\proetale$ ; voir le
lemme \ref{proetale-lemma-affine-big-site-proetale}.
Le topos des faisceaux sur $S_\proetale$ est équivalent au topos
associé à $S_{affine, \proetale}$ ; voir le
lemme \ref{proetale-lemma-alternative}.
Le résultat pour les sites $(\textit{Aff}/S)_\proetale$ et
$S_{affine, \proetale}$ découle immédiatement du résultat de Deligne
(Sites, lemme \ref{sites-lemma-criterion-points}).
Le cas de $\Sch_\proetale$ résulte du fait que ce site est égal à
$(\Sch/\Spec(\mathbf{Z}))_\proetale$.
\end{proof}

\noindent
Soit $S$ un schéma. Soit $\overline{s} : \Spec(k) \to S$ un point
géométrique. On appelle {\it voisinage pro-\'etale} de $\overline{s}$ un
diagramme commutatif
$$
\xymatrix{
\Spec(k) \ar[r]_-{\overline{u}} \ar[rd]_{\overline{s}} & U \ar[d] \\
& S
}
$$
avec $U \to S$ faiblement \'etale.

\begin{lemma}
\label{proetale-lemma-cofinal-etale}
Soit $S$ un schéma et soit $\overline{s} : \Spec(k) \to S$ un
point géométrique. La catégorie des voisinages pro-\'etales de
$\overline{s}$ est cofiltrante.
\end{lemma}

\begin{proof}
La démonstration est identique à celle de
Cohomologie \'etale, lemme \ref{etale-cohomology-lemma-cofinal-etale},
mais en utilisant les résultats correspondants sur les morphismes faiblement \'etales
démontrés dans
Compléments sur les morphismes, lemmes
\ref{more-morphisms-lemma-composition-weakly-etale},
\ref{more-morphisms-lemma-base-change-weakly-etale} et
\ref{more-morphisms-lemma-weakly-etale-permanence}.
\end{proof}

\begin{lemma}
\label{proetale-lemma-geometric-lift-to-cover}
Soit $S$ un schéma. Soit $\overline{s}$ un point géométrique de $S$.
Soit $\mathcal{U} = \{\varphi_i : S_i \to S\}_{i\in I}$ un
recouvrement pro-\'etale. Il existe alors $i \in I$ et un point
géométrique $\overline{s}_i$ de $S_i$ dont l'image est $\overline{s}$.
\end{lemma}

\begin{proof}
Cela découle immédiatement du fait que $\coprod \varphi_i$ est surjectif
et que les extensions de corps résiduels induites par des morphismes faiblement \'etales
sont algébriques séparables (voir, par exemple,
Compléments sur les morphismes, lemme
\ref{more-morphisms-lemma-weakly-etale-strictly-henselian-local-rings}).
\end{proof}

\noindent
Soit $S$ un schéma et soit $\overline{s}$ un point géométrique de $S$.
Pour $\mathcal{F}$ dans $\Sh(S_\proetale)$, définissons la
{\it fibre de $\mathcal{F}$ en $\overline{s}$} par la formule
$$
\mathcal{F}_{\overline{s}} = \colim_{(U, \overline{u})} \mathcal{F}(U)
$$
où la colimite est prise sur tous les voisinages pro-\'etales $(U, \overline{u})$
de $\overline{s}$ tels que $U \in \Ob(S_\proetale)$. Il résulte des
deux lemmes précédents que le foncteur
$$
S_\proetale \textit{Ens},\quad
U \longmapsto
\{\overline{u}\text{ point géométrique de }U\text{ dont l'image est }\overline{s}\}
$$
définit un point du site $S_\proetale$ ; voir
Sites, définition \ref{sites-definition-point} et
le lemme \ref{sites-lemma-neighbourhoods-cofiltered}.
Le foncteur $\mathcal{F} \mapsto \mathcal{F}_{\overline{s}}$
définit donc un point du topos $\Sh(S_\proetale)$ ; voir
Sites, définition \ref{sites-definition-point-topos} et
le lemme \ref{sites-lemma-point-site-topos}. En particulier, ce
foncteur est exact et commute aux colimites quelconques.
En fait, ce foncteur admet une autre description.

\begin{lemma}
\label{proetale-lemma-classical-point}
Dans la situation ci-dessus, le schéma $\Spec(\mathcal{O}_{S, \overline{s}}^{sh})$
est un objet de $X_\proetale$ et il existe un isomorphisme canonique
$$
\mathcal{F}(\Spec(\mathcal{O}_{S, \overline{s}}^{sh})) =
\mathcal{F}_{\overline{s}}
$$
fonctoriel en $\mathcal{F}$.
\end{lemma}

\begin{proof}
La première assertion résulte clairement de la construction de l'hensélisé strict
comme limite inductive filtrante d'algèbres \'etales sur $S$, ou de la caractérisation
des morphismes faiblement \'etales donnée dans
Compléments sur les morphismes, lemme
\ref{more-morphisms-lemma-weakly-etale-strictly-henselian-local-rings}.
La deuxième assertion découle du théorème d'Olivier
(Compléments d'Algèbre, théorème \ref{more-algebra-theorem-olivier}),
selon lequel le schéma $\Spec(\mathcal{O}_{S, \overline{s}}^{sh})$
est un objet initial de la catégorie des voisinages pro-\'etales
de $\overline{s}$.
\end{proof}

\noindent
Contrairement au cas du topos \'etale de $S$, il n'est pas vrai
que tout point de $\Sh(S_\proetale)$ soit de cette forme, ni que
la famille des points associés aux points géométriques soit
conservative. En effet, supposons que $S = \Spec(k)$, où $k$ est un
corps algébriquement clos. Soit $A$ un groupe abélien non nul.
Considérons le faisceau $\mathcal{F}$ sur $S_\proetale$ défini par
$$
\mathcal{F}(U) =
\frac{\{\text{applications }U \to A\}}{\{\text{applications localement constantes}\}}
$$
pour $U$ affine et, en général, par passage au faisceau associé ; voir
l'exemple \ref{proetale-example-functions-mod-locally-constant-functions}.
Alors $\mathcal{F}(U) = 0$ si $U = S = \Spec(k)$, mais, en général, $\mathcal{F}$
n'est pas nul. En effet, $S_\proetale$ contient des objets affines
ayant une infinité de points. Par exemple, soit $E = \lim E_n$
une limite projective d'ensembles finis dont les morphismes de transition sont surjectifs ;
on peut prendre $E = \mathbf{Z}_p = \lim \mathbf{Z}/p^n\mathbf{Z}$. Le schéma
$U = \Spec(\colim \text{Map}(E_n, k))$ est un objet de $S_\proetale$
car $\colim \text{Map}(E_n, k)$ est faiblement \'etale (et même ind-Zariski)
sur $k$. Ainsi, $\mathcal{F}(U)$ est non nul, car il existe des applications $E \to A$
qui ne sont pas localement constantes.
Ainsi, $\mathcal{F}$ est un faisceau abélien non nul dont la fibre en l'unique
point géométrique de $S$ est nulle. Puisque nous savons que $S_\proetale$
a suffisamment de points, on conclut qu'il doit exister un point du site pro-\'etale
qui ne provient pas de la construction expliquée ci-dessus.

\medskip\noindent
Pour remplacer les arguments qui utilisent les points, on emploie des objets affines
faiblement contractiles.
D'une part, il y a suffisamment d'objets affines faiblement contractiles d'après le
lemme \ref{proetale-lemma-proetale-enough-w-contractible}.
D'autre part, si $W \in \Ob(S_\proetale)$ est affine et faiblement contractile,
alors le foncteur
$$
\Sh(S_\proetale) \longrightarrow \textit{Ens},\quad
\mathcal{F} \longmapsto \mathcal{F}(W)
$$
est un foncteur exact $\Sh(S_\proetale) \to \textit{Ens}$ qui commute
à toutes les limites. Le foncteur
$$
\textit{Ab}(S_\proetale) \longrightarrow \textit{Ab},\quad
\mathcal{F} \longmapsto \mathcal{F}(W)
$$
est exact et commute aux sommes directes (car $W$ est quasi-compact ; voir
Sites, lemme \ref{sites-lemma-directed-colimits-sections}) ; il
commute donc à toutes les limites et à toutes les colimites. De plus, on peut vérifier l'exactitude d'un
complexe de faisceaux abéliens par évaluation sur ces objets affines
faiblement contractiles de $S_\proetale$ ; voir
Cohomologie sur les sites, proposition
\ref{sites-cohomology-proposition-enough-weakly-contractibles}.

\medskip\noindent
Remarquons enfin que le foncteur $\mathcal{F} \mapsto \mathcal{F}(W)$,
pour $W$ affine et faiblement contractile, n'est en général pas le foncteur fibre associé à un
point de $S_\proetale$, car il ne préserve pas les coproduits de faisceaux
d'ensembles si $W$ est non connexe. En fait, $W$ est non connexe dès que $W$
possède plus de $1$ point fermé, c'est-à-dire lorsque $W$ n'est pas le spectre d'un
anneau local strictement hensélien (qui est le cas particulier examiné ci-dessus).





\section{Comparaison avec le site \'etale}
\label{proetale-section-comparison}

\noindent
Soit $X$ un schéma. Avec des choix convenables de sites\footnote{Choisissons
un grand site pro-\'etale $\Sch_\proetale$ contenant $X$ comme dans la
définition \ref{proetale-definition-big-proetale-site}. Puis munissons $\Sch_\etale$
de la même catégorie sous-jacente que $\Sch_\proetale$, en prenant
pour recouvrements exactement les recouvrements pro-\'etales qui sont aussi des
recouvrements \'etales. Avec ces choix, soient $X_\etale$ et $X_\proetale$
les sous-catégories définies dans
la définition \ref{proetale-definition-big-small-proetale} et
Topologies, définition \ref{topologies-definition-big-small-etale}.
Comparer avec
Topologies, remarque \ref{topologies-remark-choice-sites}.},
le foncteur $u : X_\etale \to X_\proetale$ qui associe
$U/X$ à $U/X$ définit un morphisme de sites
$$
\epsilon : X_\proetale \longrightarrow X_\etale
$$
Cela résulte de Sites, proposition \ref{sites-proposition-get-morphism}.

\begin{lemma}
\label{proetale-lemma-describe-pullback-from-etale}
Avec les notations ci-dessus, soit $\mathcal{F}$ un faisceau sur $X_\etale$.
La règle
$$
X_\proetale \longrightarrow \textit{Ens},\quad
(f : Y \to X) \longmapsto \Gamma(Y_\etale, f_\etale^{-1}\mathcal{F})
$$
définit un faisceau égal à $\epsilon^{-1}\mathcal{F}$.
Ici, $f_\etale : Y_\etale \to X_\etale$ est le morphisme
entre les petits sites \'etales construit dans
Cohomologie \'etale, section \ref{etale-cohomology-section-functoriality}.
\end{lemma}

\begin{proof}
D'après le lemme \ref{proetale-lemma-recognize-proetale-covering}, tout recouvrement pro-\'etale
est un recouvrement fpqc. La formule définit donc un faisceau sur $X_\proetale$
d'après Cohomologie \'etale, lemme \ref{etale-cohomology-lemma-describe-pullback}.
Soit $a : \Sh(X_\etale) \to \Sh(X_\proetale)$ le foncteur
qui associe à $\mathcal{F}$ le faisceau donné par la formule du lemme.
Pour montrer que $a = \epsilon^{-1}$, il suffit de montrer que $a$ est un
adjoint à gauche de $\epsilon_*$.

\medskip\noindent
Soit $\mathcal{G}$ un objet de $\Sh(X_\proetale)$.
Rappelons que $\epsilon_*\mathcal{G}$ n'est autre que la restriction de
$\mathcal{G}$ à la sous-catégorie pleine $X_\etale$.
Soit $f : Y \to X$ un objet de $X_\proetale$.
Nous considérons $Y_\etale$ comme une sous-catégorie de $X_\proetale$.
Les morphismes de restriction du faisceau $\mathcal{G}$ définissent un morphisme
$$
\epsilon_*\mathcal{G} = \mathcal{G}|_{X_\etale}
\longrightarrow
f_{\etale, *}(\mathcal{G}|_{Y_\etale})
$$
En effet, pour $U$ dans $X_\etale$, la valeur de
$f_{\etale, *}(\mathcal{G}|_{Y_\etale})$ en $U$
est $\mathcal{G}(Y \times_X U)$ et il existe un morphisme de restriction
$\mathcal{G}(U) \to \mathcal{G}(Y \times_X U)$.
Par adjonction, cela détermine un morphisme
$$
f_\etale^{-1}(\epsilon_*\mathcal{G}) \to \mathcal{G}|_{Y_\etale}
$$
En réunissant ces morphismes pour tous les $f : Y \to X$ dans $X_\proetale$,
on obtient un morphisme canonique $a(\epsilon_*\mathcal{G}) \to \mathcal{G}$.

\medskip\noindent
Soit $\mathcal{F}$ un objet de $\Sh(X_\etale)$. On a manifestement
$\mathcal{F} = \epsilon_*a(\mathcal{F})$.

\medskip\noindent 
Nous affirmons que les morphismes $\mathcal{F} \to \epsilon_*a(\mathcal{F})$ et
$a(\epsilon_*\mathcal{G}) \to \mathcal{G}$
sont l'unité et la co-unité de l'adjonction (voir
Catégories, section \ref{categories-section-adjoint}).
Pour le voir, il suffit de montrer que les morphismes correspondants
$$
\Mor_{\Sh(X_\proetale)}(a(\mathcal{F}), \mathcal{G}) \to
\Mor_{\Sh(X_\etale)}(\mathcal{F}, \epsilon^{-1}\mathcal{G})
$$
et
$$
\Mor_{\Sh(X_\etale)}(\mathcal{F}, \epsilon^{-1}\mathcal{G}) \to
\Mor_{\Sh(X_\proetale)}(a(\mathcal{F}), \mathcal{G})
$$
sont inverses l'un de l'autre. Nous omettons la vérification détaillée.
\end{proof}

\begin{lemma}
\label{proetale-lemma-fully-faithful}
Soit $X$ un schéma. Pour tout faisceau $\mathcal{F}$ sur $X_\etale$,
le morphisme d'adjonction $\mathcal{F} \to \epsilon_*\epsilon^{-1}\mathcal{F}$ est un
isomorphisme, c'est-à-dire que $\epsilon^{-1}\mathcal{F}(U) = \mathcal{F}(U)$
pour $U$ dans $X_\etale$.
\end{lemma}

\begin{proof}
Cela découle immédiatement de la description de
$\epsilon^{-1}$ donnée dans le lemme \ref{proetale-lemma-describe-pullback-from-etale}.
\end{proof}

\begin{lemma}
\label{proetale-lemma-limit-pullback}
Soit $X$ un schéma. Soit $Y = \lim Y_i$ la limite d'un système projectif filtrant
d'objets quasi-compacts et quasi-séparés de $X_\proetale$
dont les morphismes de transition sont affines. Pour tout faisceau $\mathcal{F}$ sur $X_\etale$,
on a
$$
\epsilon^{-1}\mathcal{F}(Y) = \colim \epsilon^{-1}\mathcal{F}(Y_i)
$$
De plus, si chaque $Y_i$ appartient à $X_\etale$, on a
$\epsilon^{-1}\mathcal{F}(Y) = \colim \mathcal{F}(Y_i)$.
\end{lemma}

\begin{proof}
D'après la description de $\epsilon^{-1}\mathcal{F}$ dans le
lemme \ref{proetale-lemma-describe-pullback-from-etale}, la formule affichée
est un cas particulier de Cohomologie \'etale, théorème
\ref{etale-cohomology-theorem-colimit}.
(Lorsque $X$, $Y$ et les $Y_i$ sont tous affines, voir
Cohomologie \'etale, lemme
\ref{etale-cohomology-lemma-directed-colimit-cohomology}, dont l'énoncé est plus facile à lire.)
La dernière assertion découle immédiatement de ce qui précède
et du lemme \ref{proetale-lemma-fully-faithful}.
\end{proof}

\begin{lemma}
\label{proetale-lemma-affine-vanishing}
Soit $X$ un schéma affine. Pour tout faisceau abélien injectif $\mathcal{I}$ sur
$X_\etale$, on a $H^p(X_\proetale, \epsilon^{-1}\mathcal{I}) = 0$
pour $p > 0$.
\end{lemma}

\begin{proof}
Pour le démontrer, nous allons utiliser
Cohomologie sur les sites, lemme \ref{sites-cohomology-lemma-cech-vanish-collection}.
Soit $\mathcal{B} \subset \Ob(X_\proetale)$ l'ensemble des objets affines
$U$ de $X_\proetale$ tels que $\mathcal{O}(X) \to \mathcal{O}(U)$ soit
ind-\'etale. Soit $\text{Cov}$ l'ensemble des recouvrements pro-\'etales
$\{U_i \to U\}_{i = 1, \ldots, n}$ avec $U \in \mathcal{B}$ tels que
$\mathcal{O}(U) \to \mathcal{O}(U_i)$ soit ind-\'etale pour $i = 1, \ldots, n$.
Les propriétés (1) et (2) énoncées dans
Cohomologie sur les sites, lemme \ref{sites-cohomology-lemma-cech-vanish-collection},
sont vérifiées pour $\mathcal{B}$ et $\text{Cov}$ d'après les
lemmes \ref{proetale-lemma-composition-ind-etale},
\ref{proetale-lemma-base-change-ind-etale} et
\ref{proetale-lemma-proetale-affine}, ainsi que la
proposition \ref{proetale-proposition-weakly-etale}.

\medskip\noindent
Pour vérifier la condition (3), supposons que
$\mathcal{U} = \{U_i \to U\}_{i = 1, \ldots, n}$
soit un élément de $\text{Cov}$. Il faut montrer que les groupes de
cohomologie supérieure de Čech de $\epsilon^{-1}\mathcal{I}$
relativement à $\mathcal{U}$ sont nuls.
Écrivons d'abord $U_i = \lim_{a \in A_i} U_{i, a}$ sous la forme d'une limite projective filtrante,
avec $U_{i, a} \to U$ \'etale et $U_{i, a}$ affine.
Considérons $A_1 \times \ldots \times A_n$ comme un ensemble ordonné filtrant
pour l'ordre $(a_1, \ldots, a_n) \geq (a_1', \ldots, a_n')$
si et seulement si $a_i \geq a_i'$ pour $i = 1, \ldots, n$.
Observons que
$\mathcal{U}_{(a_1, \ldots, a_n)} = \{U_{i, a_i} \to U\}_{i = 1, \ldots, n}$
est un recouvrement \'etale pour tous
$a_1, \ldots, a_n \in A_1 \times \ldots \times A_n$.
Observons que
$$
U_{i_0} \times_U U_{i_1} \times_U \ldots \times_U U_{i_p}
=
\lim_{(a_1, \ldots, a_n) \in A_1 \times \ldots \times A_n}
U_{i_0, a_{i_0}} \times_U
U_{i_1, a_{i_1}} \times_U \ldots \times_U U_{i_p, a_{i_p}}
$$
pour tous $i_0, \ldots, i_p \in \{1, \ldots, n\}$, car les limites
commutent aux produits fibrés. Ainsi, d'après le lemme \ref{proetale-lemma-limit-pullback}
et l'exactitude des colimites filtrantes, on a
$$
\check{H}^p(\mathcal{U}, \epsilon^{-1}\mathcal{I}) =
\colim \check{H}^p(\mathcal{U}_{(a_1, \ldots, a_n)}, \epsilon^{-1}\mathcal{I}) 
$$
Il suffit donc de démontrer l'annulation pour les recouvrements \'etales
de $U$!

\medskip\noindent
Soit $\mathcal{U} = \{U_i \to U\}_{i = 1, \ldots, n}$ un recouvrement \'etale
avec $U_i$ affine. Écrivons $U = \lim_{b \in B} U_b$ comme limite projective filtrante,
avec $U_b$ affine et $U_b \to X$ \'etale.
D'après Limites, lemmes \ref{limits-lemma-descend-finite-presentation},
\ref{limits-lemma-limit-affine} et
\ref{limits-lemma-descend-etale},
on peut choisir $b_0 \in B$ tel que, pour $i = 1, \ldots, n$, il existe un
morphisme \'etale $U_{i, b_0} \to U_{b_0}$ entre schémas affines
tel que $U_i = U \times_{U_{b_0}} U_{i, b_0}$.
Posons $U_{i, b} = U_b \times_{U_{b_0}} U_{i, b_0}$ pour $b \geq b_0$.
Pour $b$ assez grand, la famille
$\mathcal{U}_b = \{U_{i, b} \to U_b\}_{i = 1, \ldots, n}$
est un recouvrement \'etale ; voir
Limites, lemme \ref{limits-lemma-descend-surjective}.
Comme précédemment, on obtient
$$
\check{H}^p(\mathcal{U}, \epsilon^{-1}\mathcal{I}) =
\colim \check{H}^p(\mathcal{U}_b, \epsilon^{-1}\mathcal{I}) =
\colim \check{H}^p(\mathcal{U}_b, \mathcal{I})
$$
où la dernière égalité résulte du lemme \ref{proetale-lemma-fully-faithful}.
Puisque chacun des complexes de {\v C}ech du membre de droite est
acyclique en degrés strictement positifs
(Cohomologie sur les sites, lemme
\ref{sites-cohomology-lemma-injective-trivial-cech}),
celui du membre de gauche l'est également. Cela démontre la condition (3) de
Cohomologie sur les sites, lemme \ref{sites-cohomology-lemma-cech-vanish-collection}.
Puisque $X \in \mathcal{B}$, le lemme en découle.
\end{proof}

\begin{lemma}
\label{proetale-lemma-relative-comparison}
Soit $X$ un schéma.
\begin{enumerate}
\item Pour un faisceau abélien $\mathcal{F}$ sur $X_\etale$,
on a $R\epsilon_*(\epsilon^{-1}\mathcal{F}) = \mathcal{F}$.
\item Pour $K \in D^+(X_\etale)$, le morphisme $K \to R\epsilon_*\epsilon^{-1}K$
est un isomorphisme.
\end{enumerate}
\end{lemma}

\begin{proof}
Soit $\mathcal{I}$ un faisceau abélien injectif sur $X_\etale$.
Rappelons que $R^q\epsilon_*(\epsilon^{-1}\mathcal{I})$ est le faisceau associé
à $U \mapsto H^q(U_\proetale, \epsilon^{-1}\mathcal{I})$ ; voir
Cohomologie sur les sites, lemme \ref{sites-cohomology-lemma-higher-direct-images}.
D'après le lemme \ref{proetale-lemma-affine-vanishing}, on voit qu'il est nul pour $q > 0$
et $U$ affine et \'etale sur $X$. Puisque tout objet de
$X_\etale$ admet un recouvrement par des objets affines, il s'ensuit que
$R^q\epsilon_*(\epsilon^{-1}\mathcal{I}) = 0$ pour $q > 0$.

\medskip\noindent
Soit $K \in D^+(X_\etale)$. Choisissons un complexe borné inférieurement $\mathcal{I}^\bullet$
de faisceaux abéliens injectifs sur $X_\etale$ représentant $K$.
Alors $\epsilon^{-1}K$ est représenté par $\epsilon^{-1}\mathcal{I}^\bullet$.
Le lemme d'acyclicité de Leray (Catégories dérivées, lemme
\ref{derived-lemma-leray-acyclicity}) montre que $R\epsilon_*\epsilon^{-1}K$
est représenté par $\epsilon_*\epsilon^{-1}\mathcal{I}^\bullet$.
D'après le lemme \ref{proetale-lemma-fully-faithful}, on conclut que
$R\epsilon_*\epsilon^{-1}\mathcal{I}^\bullet = \mathcal{I}^\bullet$,
ce qui achève la démonstration de (2). La partie (1) est un cas particulier de (2).
\end{proof}

\begin{lemma}
\label{proetale-lemma-compare-cohomology}
Soit $X$ un schéma.
\begin{enumerate}
\item Pour un faisceau abélien $\mathcal{F}$ sur $X_\etale$, on a
$$
H^i(X_\etale, \mathcal{F}) =
H^i(X_\proetale, \epsilon^{-1}\mathcal{F})
$$
pour tout $i$.
\item Pour $K \in D^+(X_\etale)$, on a
$$
R\Gamma(X_\etale, K) = R\Gamma(X_\proetale, \epsilon^{-1}K)
$$
\end{enumerate}
\end{lemma}

\begin{proof}
Conséquence immédiate du lemme \ref{proetale-lemma-relative-comparison} et de
la suite spectrale de Leray (Cohomologie sur les sites, lemme
\ref{sites-cohomology-lemma-apply-Leray}).
\end{proof}

\begin{lemma}
\label{proetale-lemma-compare-cohomology-nonabelian}
Soit $X$ un schéma. Soit $\mathcal{G}$ un faisceau de groupes (non nécessairement
commutatifs) sur $X_\etale$. On a
$$
H^1(X_\etale, \mathcal{G}) =
H^1(X_\proetale, \epsilon^{-1}\mathcal{G})
$$
où $H^1$ désigne l'ensemble des classes d'isomorphisme de
torseurs (voir
Cohomologie sur les sites, section \ref{sites-cohomology-section-h1-torsors}).
\end{lemma}

\begin{proof}
Comme le foncteur $\epsilon^{-1}$ est pleinement fidèle d'après le
lemme \ref{proetale-lemma-fully-faithful},
il est clair que l'application
$H^1(X_\etale, \mathcal{G}) \to H^1(X_\proetale, \epsilon^{-1}\mathcal{G})$
est injective. Pour montrer la surjectivité, il suffit de montrer que
tout $\epsilon^{-1}\mathcal{G}$-torseur $\mathcal{F}$ est localement trivial pour la topologie \'etale.
Pour le voir, nous pouvons supposer que $X$ est affine.
Nous nous ramenons donc à démontrer la surjectivité lorsque $X$ est affine.

\medskip\noindent
Choisissons un recouvrement $\{U \to X\}$ tel que (a) $U$ soit affine, (b)
que $\mathcal{O}(X) \to \mathcal{O}(U)$ soit ind-\'etale, et (c) que $\mathcal{F}(U)$
soit non vide. C'est possible d'après la proposition \ref{proetale-proposition-weakly-etale}
et parce que les
recouvrements pro-\'etales standards de $X$ sont cofinaux parmi tous les recouvrements pro-\'etales
de $X$ (lemme \ref{proetale-lemma-proetale-affine}).
Écrivons $U = \lim U_i$ comme une limite de schémas affines, chacun étant \'etale sur $X$.
Choisissons $s \in \mathcal{F}(U)$. Soit
$g \in \epsilon^{-1}\mathcal{G}(U \times_X U)$
l'unique section telle que $g \cdot \text{pr}_1^*s = \text{pr}_2^*s$ dans
$\mathcal{F}(U \times_X U)$. Alors $g$ satisfait la condition de cocycle
$$
\text{pr}_{12}^*g \cdot \text{pr}_{23}^*g = \text{pr}_{13}^*g
$$
dans $\epsilon^{-1}\mathcal{G}(U \times_X U \times_X U)$. D'après le
lemme \ref{proetale-lemma-limit-pullback},
on a
$$
\epsilon^{-1}\mathcal{G}(U \times_X U) =
\colim \mathcal{G}(U_i \times_X U_i)
$$
et
$$
\epsilon^{-1}\mathcal{G}(U \times_X U \times_X U) =
\colim \mathcal{G}(U_i \times_X U_i \times_X U_i)
$$
il existe donc un indice $i$ et un élément $g_i \in \mathcal{G}(U_i \times_X U_i)$
dont l'image est $g$ et qui satisfait la condition de cocycle.
Le cocycle $g_i$ définit alors un torseur sous $\mathcal{G}$ sur
$X_\etale$ dont l'image inverse est isomorphe à $\mathcal{F}$
par construction. Nous omettons certains détails, plus précisément le lien
entre les torseurs et les 1-cocycles, qu'il conviendrait d'ajouter au chapitre
Cohomologie sur les sites.
\end{proof}

\begin{lemma}
\label{proetale-lemma-compare-derived}
Soit $X$ un schéma. Soit $\Lambda$ un anneau.
\begin{enumerate}
\item L'image essentielle du foncteur pleinement fidèle
$\epsilon^{-1} : \textit{Mod}(X_\etale, \Lambda) \to
\textit{Mod}(X_\proetale, \Lambda)$
est une sous-catégorie de Serre faible $\mathcal{C}$.
\item Le foncteur $\epsilon^{-1}$ définit une équivalence de catégories
entre $D^+(X_\etale, \Lambda)$ et $D^+_\mathcal{C}(X_\proetale, \Lambda)$
dont un quasi-inverse est donné par $R\epsilon_*$.
\end{enumerate}
\end{lemma}

\begin{proof}
Pour démontrer (1), nous allons vérifier les conditions (1) -- (4) de
Homologie, lemme \ref{homology-lemma-characterize-weak-serre-subcategory}.
Comme $\epsilon^{-1}$ est pleinement fidèle (lemme \ref{proetale-lemma-fully-faithful})
et exact, tout est clair sauf la condition (4).
Cependant, si
$$
0 \to \epsilon^{-1}\mathcal{F}_1 \to \mathcal{G} \to
\epsilon^{-1}\mathcal{F}_2 \to 0
$$
est une suite exacte courte de faisceaux de $\Lambda$-modules sur $X_\proetale$,
alors on obtient
$$
0 \to \epsilon_*\epsilon^{-1}\mathcal{F}_1 \to \epsilon_*\mathcal{G} \to
\epsilon_*\epsilon^{-1}\mathcal{F}_2 \to
R^1\epsilon_*\epsilon^{-1}\mathcal{F}_1
$$
D'après le lemme \ref{proetale-lemma-relative-comparison},
cette suite se réduit à la suite exacte courte
$$
0 \to \mathcal{F}_1 \to \epsilon_*\mathcal{G} \to \mathcal{F}_2 \to 0
$$
En prenant l'image inverse, on trouve que $\mathcal{G} = \epsilon^{-1}\epsilon_*\mathcal{G}$.
Cela démontre (1).

\medskip\noindent
L'assertion (2) découle de l'assertion (1) et de
Cohomologie sur les sites, lemme \ref{sites-cohomology-lemma-equivalence-bounded}.
\end{proof}

\noindent
Soit $\Lambda$ un anneau. Dans
Modules sur les sites, section \ref{sites-modules-section-locally-constant},
nous avons défini la notion de faisceau localement constant de $\Lambda$-modules
sur un site. Si $M$ est un $\Lambda$-module, alors le faisceau $\underline{M}$ est
de présentation finie comme faisceau de $\underline{\Lambda}$-modules si et seulement
si $M$ est de présentation finie comme $\Lambda$-module ; voir
Modules sur les sites, lemme \ref{sites-modules-lemma-locally-constant-finite-type}.

\begin{lemma}
\label{proetale-lemma-compare-locally-constant}
Soit $X$ un schéma. Soit $\Lambda$ un anneau.
Le foncteur $\epsilon^{-1}$ définit une équivalence de catégories
$$
\left\{
\begin{matrix}
\text{faisceaux localement constants}\\
\text{de }\Lambda\text{-modules sur }X_\etale\\
\text{de présentation finie}
\end{matrix}
\right\}
\longleftrightarrow
\left\{
\begin{matrix}
\text{faisceaux localement constants}\\
\text{de }\Lambda\text{-modules sur }X_\proetale\\
\text{de présentation finie}
\end{matrix}
\right\}
$$
\end{lemma}

\begin{proof}
Soit $\mathcal{F}$ un faisceau localement constant de $\Lambda$-modules
sur $X_\proetale$ et de présentation finie. Choisissons un recouvrement pro-\'etale
$\{U_i \to X\}$ tel que $\mathcal{F}|_{U_i}$ soit constant, disons
$\mathcal{F}|_{U_i} \cong \underline{M_i}_{U_i}$.
Remarquons que $U_i \times_X U_j$ est vide si $M_i$ n'est pas isomorphe
à $M_j$.
Pour tout $\Lambda$-module $M$, posons $I_M = \{i \in I \mid M_i \cong M\}$.
Comme tout recouvrement pro-\'etale est un recouvrement fpqc et que
Descente, lemme \ref{descent-lemma-open-fpqc-covering},
on voit que $U_M = \bigcup_{i \in I_M} \Im(U_i \to X)$
est un ouvert de $X$. Alors $X = \coprod U_M$ est une décomposition en ouverts
disjoints de $X$. Nous pouvons remplacer $X$ par $U_M$, pour un certain $M$, et
supposer que $M_i = M$ pour tout $i$.

\medskip\noindent
Considérons le faisceau $\mathcal{I} = \mathit{Isom}(\underline{M}, \mathcal{F})$.
Ce faisceau est un torseur sous
$\mathcal{G} = \mathit{Isom}(\underline{M}, \underline{M})$.
D'après Modules sur les sites, lemme \ref{sites-modules-lemma-locally-constant},
on a $\mathcal{G} = \underline{G}$,
où $G = \mathit{Isom}_\Lambda(M, M)$.
Puisque les torseurs pour la topologie \'etale
et pour la topologie pro-\'etale coïncident d'après le
lemme \ref{proetale-lemma-compare-cohomology-nonabelian},
il s'ensuit que $\mathcal{I}$ admet localement sur $X$ des sections pour la topologie \'etale.
Ainsi, $\mathcal{F}$ est localement, pour la topologie \'etale, un faisceau constant, ce qui
est précisément ce qu'il fallait démontrer.
\end{proof}

\begin{lemma}
\label{proetale-lemma-compare-locally-constant-derived}
Soit $X$ un schéma. Soit $\Lambda$ un anneau noethérien.
Soient $D_{flc}(X_\etale, \Lambda)$, resp.\ $D_{flc}(X_\proetale, \Lambda)$,
les sous-catégories pleines respectives de
$D(X_\etale, \Lambda)$, resp.\ $D(X_\proetale, \Lambda)$,
formées des complexes dont les faisceaux de cohomologie sont localement
constants et de type fini comme faisceaux de $\Lambda$-modules. Alors
$$
\epsilon^{-1} :
D_{flc}^+(X_\etale, \Lambda)
\longrightarrow
D_{flc}^+(X_\proetale, \Lambda)
$$
est une équivalence de catégories.
\end{lemma}

\begin{proof}
Les catégories $D_{flc}(X_\etale, \Lambda)$ et $D_{flc}(X_\proetale, \Lambda)$
sont des sous-catégories triangulées strictement pleines et saturées de
$D(X_\etale, \Lambda)$ et $D(X_\proetale, \Lambda)$
d'après Modules sur les sites, lemme
\ref{sites-modules-lemma-kernel-finite-locally-constant}
et
Catégories dérivées, section \ref{derived-section-triangulated-sub}.
L'énoncé du lemme s'obtient en combinant les
lemmes \ref{proetale-lemma-compare-derived} et
\ref{proetale-lemma-compare-locally-constant}.
\end{proof}

\begin{lemma}
\label{proetale-lemma-compare-truncations}
Soit $X$ un schéma. Soit $\Lambda$ un anneau noethérien.
Soit $K$ un objet de $D(X_\proetale, \Lambda)$.
Posons $K_n = K \otimes_\Lambda^\mathbf{L} \underline{\Lambda/I^n}$.
Si $K_1$ est
\begin{enumerate}
\item dans l'image essentielle de
$\epsilon^{-1} :D(X_\etale, \Lambda/I) \to D(X_\proetale, \Lambda/I)$, et
\item d'amplitude de Tor dans $[a,\infty)$ pour un certain $a \in \mathbf{Z}$,
\end{enumerate}
alors $K_n$, considéré comme objet de $D(X_\proetale, \Lambda/I^n)$, vérifie (1) et (2).
\end{lemma}

\begin{proof}
Pour $K_n$, l'assertion (2) résulte du résultat plus général de
Cohomologie sur les sites, lemme \ref{sites-cohomology-lemma-bounded}.
Pour $K_n$, l'assertion (1) résulte, par récurrence sur $n$, des
triangles distingués
$$
K \otimes_\Lambda^\mathbf{L} \underline{I^n/I^{n + 1}} \to
K_{n + 1} \to
K_n \to
K \otimes_\Lambda^\mathbf{L} \underline{I^n/I^{n + 1}}[1]
$$
et de l'isomorphisme
$$
K \otimes_\Lambda^\mathbf{L} \underline{I^n/I^{n + 1}} =
K_1 \otimes_{\Lambda/I}^\mathbf{L} \underline{I^n/I^{n + 1}}
$$
ainsi que du fait démontré dans le lemme \ref{proetale-lemma-compare-derived},
que l'image essentielle de $\epsilon^{-1}$ est une sous-catégorie triangulée
de $D^+(X_\proetale, \Lambda/I^n)$.
\end{proof}

\begin{example}
\label{proetale-example-functions-mod-locally-constant-functions}
Soit $X$ un schéma. Soit $A$ un groupe abélien. Notons
$fun(-, A)$ le faisceau sur $X_\proetale$ qui associe à $U$
l'ensemble de toutes les applications $U \to A$ (au niveau des ensembles de points).
Considérons la suite de faisceaux
$$
0 \to \underline{A} \to fun(-, A) \to \mathcal{F} \to 0
$$
sur $X_\proetale$. Puisque le faisceau constant provient par image inverse du topos
final, on voit que $\underline{A} = \epsilon^{-1}\underline{A}$.
Cependant, si $A$ possède plus d'un élément, ni $fun(-, A)$
ni $\mathcal{F}$ ne proviennent du site \'etale de $X$ par image inverse.
Pour déterminer les valeurs de $\mathcal{F}$ dans certains cas, supposons que
tous les points de $X$ soient fermés et aient un corps résiduel séparablement clos,
et que $U$ soit affine. Alors tous les points de $U$ sont fermés et ont un corps
résiduel séparablement clos, et l'on a
$$
H^1_\proetale(U, \underline{A}) =
H^1_\etale(U, \underline{A}) = 0
$$
d'après le lemme \ref{proetale-lemma-compare-cohomology} et le
Cohomologie \'etale, lemme
\ref{etale-cohomology-lemma-affine-only-closed-points}.
On a donc, dans ce cas,
$$
\mathcal{F}(U) = fun(U, A)/\underline{A}(U)
$$
\end{example}





\section{Complétion dérivée dans le cas noethérien constant}
\label{proetale-section-derived-completion-noetherian}

\noindent
Nous poursuivons la discussion commencée dans Géométrie algébrique et formelle, section
\ref{algebraization-section-derived-completion};
nous supposons que le lecteur a lu au moins une partie de cette section.

\medskip\noindent
Soit $\mathcal{C}$ un site. Soit $\Lambda$ un anneau noethérien
et soit $I \subset \Lambda$ un idéal. Rappelons, d'après
Modules sur les sites, lemme \ref{sites-modules-lemma-completion-flat},
que
$$
\underline{\Lambda}^\wedge = \lim \underline{\Lambda/I^n}
$$
est une $\underline{\Lambda}$-algèbre plate et que le morphisme
$\underline{\Lambda} \to \underline{\Lambda}^\wedge$ identifie
les quotients modulo $I$. Ainsi,
Géométrie algébrique et formelle, lemme
\ref{algebraization-lemma-restriction-derived-complete-equivalence}
donne
$$
D_{comp}(\mathcal{C}, \Lambda) =
D_{comp}(\mathcal{C}, \underline{\Lambda}^\wedge)
$$
En particulier, les faisceaux de cohomologie $H^i(K)$ d'un objet $K$ de
$D_{comp}(\mathcal{C}, \Lambda)$ sont des faisceaux de
$\underline{\Lambda}^\wedge$-modules.
Pour simplifier les notations, nous travaillons souvent avec
$D_{comp}(\mathcal{C}, \Lambda)$.

\begin{lemma}
\label{proetale-lemma-naive-completion}
Soit $\mathcal{C}$ un site. Soit $\Lambda$ un anneau noethérien
et soit $I \subset \Lambda$ un idéal. L'adjoint à gauche
du foncteur d'inclusion
$D_{comp}(\mathcal{C}, \Lambda) \to D(\mathcal{C}, \Lambda)$
de Géométrie algébrique et formelle, proposition
\ref{algebraization-proposition-derived-completion}, associe à $K$ le complexe
$$
K^\wedge = R\lim(K \otimes_\Lambda^\mathbf{L} \underline{\Lambda/I^n})
$$
En particulier, $K$ est complet au sens dérivé si et seulement si
$K = R\lim(K \otimes_\Lambda^\mathbf{L} \underline{\Lambda/I^n})$.
\end{lemma}

\begin{proof}
Choisissons des générateurs $f_1, \ldots, f_r$ de $I$. D'après
Géométrie algébrique et formelle, lemme
\ref{algebraization-lemma-derived-completion-koszul},
on a
$$
K^\wedge = 
R\lim (K \otimes_\Lambda^\mathbf{L} \underline{K_n})
$$
où $K_n = K(\Lambda, f_1^n, \ldots, f_r^n)$.
Dans Compléments d'algèbre, lemme \ref{more-algebra-lemma-sequence-Koszul-complexes},
nous avons vu que les pro-systèmes $\{K_n\}$ et
$\{\Lambda/I^n\}$ de $D(\Lambda)$
sont isomorphes. Le lemme en résulte.
\end{proof}

\begin{lemma}
\label{proetale-lemma-pushforward-Noetherian-case}
Soit $\Lambda$ un anneau noethérien. Soit $I \subset \Lambda$ un idéal.
Soit $f : \Sh(\mathcal{D}) \to \Sh(\mathcal{C})$ un morphisme de topos.
Alors
\begin{enumerate}
\item $Rf_*$ envoie $D_{comp}(\mathcal{D}, \Lambda)$
dans $D_{comp}(\mathcal{C}, \Lambda)$,
\item le foncteur $Rf_* : D_{comp}(\mathcal{D}, \Lambda) \to
D_{comp}(\mathcal{C}, \Lambda)$ admet un adjoint à gauche
$Lf_{comp}^* : D_{comp}(\mathcal{C}, \Lambda) \to
D_{comp}(\mathcal{D}, \Lambda)$ qui s'obtient en faisant suivre $Lf^*$ de la
complétion dérivée,
\item $Rf_*$ commute à la complétion dérivée,
\item pour $K$ dans $D_{comp}(\mathcal{D}, \Lambda)$, on a
$Rf_*K = R\lim Rf_*(K \otimes^\mathbf{L}_\Lambda \underline{\Lambda/I^n})$.
\item pour $M$ dans $D_{comp}(\mathcal{C}, \Lambda)$, on a
$Lf^*_{comp}M =
R\lim Lf^*(M \otimes^\mathbf{L}_\Lambda \underline{\Lambda/I^n})$.
\end{enumerate}
\end{lemma}

\begin{proof}
Nous avons établi (1) et (2) dans
Géométrie algébrique et formelle, lemme
\ref{algebraization-lemma-pushforward-derived-complete-adjoint}.
Le point (3) résulte de
Géométrie algébrique et formelle, lemme
\ref{algebraization-lemma-pushforward-commutes-with-derived-completion}.
Pour (4), soit $K$ complet au sens dérivé. Alors
$$
Rf_*K = Rf_*( R\lim K \otimes^\mathbf{L}_\Lambda \underline{\Lambda/I^n}) =
R\lim Rf_*(K \otimes^\mathbf{L}_\Lambda \underline{\Lambda/I^n})
$$
où la première égalité résulte du lemme \ref{proetale-lemma-naive-completion},
et la seconde du fait que $Rf_*$ commute à $R\lim$
(Cohomologie sur les sites, lemme
\ref{sites-cohomology-lemma-Rf-commutes-with-Rlim}). Cela démontre (4).
Pour démontrer (5), le lemme \ref{proetale-lemma-naive-completion} donne
$$
Lf_{comp}^*M =
R\lim ( Lf^*M \otimes_\Lambda^\mathbf{L} \underline{\Lambda/I^n})
$$
Comme $Lf^*$ commute au produit tensoriel dérivé d'après
Cohomologie sur les sites, lemme \ref{sites-cohomology-lemma-pullback-tensor-product},
et que $Lf^*\underline{\Lambda/I^n} = \underline{\Lambda/I^n}$,
on obtient (5).
\end{proof}







\section{Complétion dérivée et objets faiblement contractiles}
\label{proetale-section-derived-completion-proetale}

\noindent
Nous poursuivons la discussion entreprise à la
section \ref{proetale-section-derived-completion-noetherian}.
Dans cette section, nous verrons comment l'existence
d'objets faiblement contractiles simplifie l'étude
des modules complets au sens dérivé.

\medskip\noindent
Soit $\mathcal{C}$ un site. Soit $\Lambda$ un anneau noethérien.
Soit $I \subset \Lambda$ un idéal. Bien que la théorie générale
de $D_{comp}(\mathcal{C}, \Lambda)$ soit tout à fait satisfaisante,
il est difficile de donner explicitement des exemples de complexes complets au sens dérivé.
Nous savons que
\begin{enumerate}
\item tout objet $M$ de $D(\mathcal{C}, \Lambda/I^n)$ devient par restriction un objet
complet au sens dérivé de $D(\mathcal{C}, \Lambda)$, et
\item pour tout $K \in D(\mathcal{C}, \Lambda)$, la complétion dérivée
$K^\wedge = R\lim (K \otimes_\Lambda^\mathbf{L} \underline{\Lambda/I^n})$
est complet au sens dérivé.
\end{enumerate}
Les objets du premier type sont trivialement complets et ne sont peut-être pas
intéressants. Le problème de (2) est que la complétion dérivée
est en général assez mystérieuse, même dans le cas $K = \underline{\Lambda}$.
En effet, la définition des limites homotopiques fournit
un triangle distingué
$$
R\lim(\underline{\Lambda/I^n}) \to
\prod \underline{\Lambda/I^n} \to
\prod \underline{\Lambda/I^n} \to
R\lim(\underline{\Lambda/I^n})[1]
$$
dans $D(\mathcal{C}, \Lambda)$, où les produits sont pris dans
$D(\mathcal{C}, \Lambda)$. Ces produits se calculent en prenant les produits
de résolutions injectives
(Injectifs, lemme \ref{injectives-lemma-derived-products}),
et l'on voit donc que le faisceau $H^p(\prod \underline{\Lambda/I^n})$
est le faisceau associé au préfaisceau
$$
U \longmapsto \prod H^p(U, \Lambda/I^n).
$$
Comme exemple explicite, si $X = \Spec(\mathbf{C}[t, t^{-1}])$,
$\mathcal{C} = X_\etale$, $\Lambda = \mathbf{Z}$, $I = (2)$ et
$p = 1$, on obtient le faisceau associé au préfaisceau
$$
U \mapsto \prod H^1(U_\etale, \mathbf{Z}/2^n\mathbf{Z})
$$
pour $U$ \'etale sur $X$. Observons que $H^1(X_\etale, \mathbf{Z}/m\mathbf{Z})$
est cyclique d'ordre $m$ et admet pour générateur $\alpha_m$, donné par le revêtement \'etale fini
de groupe $\mathbf{Z}/m\mathbf{Z}$ défini par l'équation $t = s^m$
(voir Cohomologie \'etale, section \ref{etale-cohomology-section-computation}).
Alors la section
$$
\alpha = (\alpha_{2^n}) \in \prod H^1(X_\etale, \mathbf{Z}/2^n\mathbf{Z})
$$
du préfaisceau ci-dessus ne devient nulle par restriction à aucun schéma \'etale
non vide au-dessus de $X$; ainsi, le faisceau associé au préfaisceau n'est pas nul.

\medskip\noindent
Cependant, ce phénomène ne se produit pas sur le site pro-\'etale.
La raison en est que nous avons suffisamment d'objets (quasi-compacts) faiblement contractiles.
Dans la proposition suivante, nous rassemblons quelques résultats sur la
complétion dérivée dans le cas d'un anneau noethérien constant, pour les sites qui possèdent suffisamment
d'objets faiblement contractiles (voir
Sites, définition \ref{sites-definition-w-contractible}).

\begin{proposition}
\label{proetale-proposition-enough-weakly-contractibles}
Soit $\mathcal{C}$ un site. Supposons que $\mathcal{C}$ possède suffisamment
d'objets faiblement contractiles.
Soit $\Lambda$ un anneau noethérien. Soit $I \subset \Lambda$ un idéal.
\begin{enumerate}
\item La catégorie des faisceaux de $\Lambda$-modules complets au sens dérivé est une
sous-catégorie de Serre faible de $\textit{Mod}(\mathcal{C}, \Lambda)$.
\item Un faisceau $\mathcal{F}$ de $\Lambda$-modules vérifie
$\mathcal{F} = \lim \mathcal{F}/I^n\mathcal{F}$ si et seulement si
$\mathcal{F}$ est complet au sens dérivé et $\bigcap I^n\mathcal{F} = 0$.
\item Le faisceau $\underline{\Lambda}^\wedge$ est complet au sens dérivé.
\item Si $\ldots \to \mathcal{F}_3 \to \mathcal{F}_2 \to \mathcal{F}_1$
est un système projectif de faisceaux de $\Lambda$-modules complets au sens dérivé,
alors $\lim \mathcal{F}_n$ est complet au sens dérivé.
\item Un objet $K \in D(\mathcal{C}, \Lambda)$ est complet au sens dérivé si
et seulement si chaque faisceau de cohomologie $H^p(K)$ est complet au sens dérivé.
\item Un objet $K \in D_{comp}(\mathcal{C}, \Lambda)$ est borné supérieurement
si et seulement si $K \otimes_\Lambda^\mathbf{L} \underline{\Lambda/I}$
est borné supérieurement.
\item Un objet $K \in D_{comp}(\mathcal{C}, \Lambda)$ est borné
si $K \otimes_\Lambda^\mathbf{L} \underline{\Lambda/I}$ est de dimension de
Tor finie.
\end{enumerate}
\end{proposition}

\begin{proof}
Soit $\mathcal{B} \subset \Ob(\mathcal{C})$ un sous-ensemble tel que tout
$U \in \mathcal{B}$ soit faiblement contractile et que tout objet de $\mathcal{C}$
admette un recouvrement par des éléments de $\mathcal{B}$.
Nous utiliserons les résultats de Cohomologie sur les sites,
lemme \ref{sites-cohomology-lemma-w-contractible} et
proposition \ref{sites-cohomology-proposition-enough-weakly-contractibles},
sans autre mention.

\medskip\noindent
Rappelons que $R\lim$ commute à $R\Gamma(U, -)$
(Injectifs, lemme \ref{injectives-lemma-RF-commutes-with-Rlim}).
Soit $f \in I$. Rappelons que $T(K, f)$ est la limite homotopique
du système
$$
\ldots \xrightarrow{f} K \xrightarrow{f} K \xrightarrow{f} K
$$
dans $D(\mathcal{C}, \Lambda)$. Ainsi,
$$
R\Gamma(U, T(K, f)) = T(R\Gamma(U, K), f).
$$
Puisque les isomorphismes entre objets de
$D(\mathcal{C}, \Lambda)$ se détectent par évaluation sur les $U \in \mathcal{B}$,
nous concluons qu'un objet $K$ de $D(\mathcal{C}, \Lambda)$
est complet au sens dérivé si et seulement si, pour tout $U \in \mathcal{B}$,
l'objet $R\Gamma(U, K)$ est complet au sens dérivé en tant qu'objet de $D(\Lambda)$.

\medskip\noindent
La remarque ci-dessus montre que les points (1) et (5) découlent des résultats correspondants
pour les modules sur les anneaux; voir
Compléments d'algèbre, lemmes \ref{more-algebra-lemma-hom-from-Af} et
\ref{more-algebra-lemma-serre-subcategory}.
De même, (2) se déduit de
Compléments d'algèbre, proposition
\ref{more-algebra-proposition-derived-complete-modules},
puisque $(I^n\mathcal{F})(U) = I^n \cdot \mathcal{F}(U)$
pour $U \in \mathcal{B}$ (par exactitude de l'évaluation en $U$).

\medskip\noindent
Démonstration de (4). La limite homotopique $R\lim \mathcal{F}_n$ appartient à
$D_{comp}(X, \Lambda)$ (voir la discussion qui suit
Géométrie algébrique et formelle, définition
\ref{algebraization-definition-derived-complete}).
D'après le point (5) que nous venons de démontrer, nous concluons que
$\lim \mathcal{F}_n = H^0(R\lim \mathcal{F}_n)$
est complet au sens dérivé.
Le point (3) est un cas particulier de (4).

\medskip\noindent
Démonstration de (6) et (7). Cela résulte du
lemme \ref{proetale-lemma-naive-completion},
du
lemme de Cohomologie sur les sites \ref{sites-cohomology-lemma-bounded}
et du calcul des limites homotopiques dans Cohomologie sur les sites,
proposition \ref{sites-cohomology-proposition-enough-weakly-contractibles}.
\end{proof}






\section{Cohomologie d'un point}
\label{proetale-section-cohomology-point}

\noindent
Soit $\Lambda$ un anneau noethérien complet pour la topologie adique définie par un idéal
$I \subset \Lambda$. Soit $k$ un corps. Dans cette section, nous ``calculons''
$$
H^i(\Spec(k)_\proetale, \underline{\Lambda}^\wedge)
$$
où $\underline{\Lambda}^\wedge = \lim_m \underline{\Lambda/I^m}$ comme précédemment.
Soit $k^{sep}$ une clôture séparable de $k$.
Alors
$$
\mathcal{U} = \{\Spec(k^{sep}) \to \Spec(k)\}
$$
est un recouvrement pro-\'etale de $\Spec(k)$. Nous utiliserons, pour ce recouvrement, la suite spectrale de {\v C}ech
aboutissant à la cohomologie.
Posons $U_0 = \Spec(k^{sep})$ et
\begin{align*}
U_n & =
\Spec(k^{sep}) \times_{\Spec(k)}
\Spec(k^{sep}) \times_{\Spec(k)} \ldots
\times_{\Spec(k)} \Spec(k^{sep}) \\
& =
\Spec(k^{sep} \otimes_k k^{sep} \otimes_k \ldots \otimes_k k^{sep})
\end{align*}
($n + 1$ facteurs). Observons que l'espace topologique sous-jacent
$|U_0|$ de $U_0$ est réduit à un point et que, pour $n \geq 1$, on a
$$
|U_n| = G \times \ldots \times G\quad (n\text{ facteurs})
$$
comme espaces profinis, où $G = \text{Gal}(k^{sep}/k)$. En effet, tout
point de $U_n$ a pour corps résiduel $k^{sep}$ et nous identifions
$(\sigma_1, \ldots, \sigma_n)$ au point correspondant à la
surjection
$$
k^{sep} \otimes_k k^{sep} \otimes_k \ldots \otimes_k k^{sep}
\longrightarrow k^{sep}, \quad
\lambda_0 \otimes \lambda_1 \otimes \ldots \lambda_n
\longmapsto \lambda_0 \sigma_1(\lambda_1) \ldots \sigma_n(\lambda_n)
$$
On calcule alors
\begin{align*}
R\Gamma((U_n)_\proetale, \underline{\Lambda}^\wedge)
& =
R\lim_m R\Gamma((U_n)_\proetale, \underline{\Lambda/I^m}) \\
& =
R\lim_m R\Gamma((U_n)_\etale, \underline{\Lambda/I^m}) \\
& =
\lim_m H^0(U_n, \underline{\Lambda/I^m}) \\
& = 
\text{Maps}_{cont}(G \times \ldots \times G, \Lambda)
\end{align*}
La première égalité résulte du fait que $R\Gamma$ commute aux limites dérivées
et que $\Lambda^\wedge$ est la limite dérivée des faisceaux
$\underline{\Lambda/I^m}$ d'après
la proposition \ref{proetale-proposition-enough-weakly-contractibles}.
La deuxième égalité résulte du lemme \ref{proetale-lemma-compare-cohomology}.
La troisième égalité résulte de Cohomologie \'etale, lemme
\ref{etale-cohomology-lemma-affine-only-closed-points}.
Pour la quatrième égalité, on utilise
Cohomologie \'etale, remarque
\ref{etale-cohomology-remark-constant-locally-constant-maps}
pour identifier les sections du faisceau constant $\underline{\Lambda/I^m}$.
Elle utilise ensuite le fait que $\Lambda$ est complet pour la topologie $I$-adique
et donc égal à $\lim_m \Lambda/I^m$ comme espace topologique, pour voir que
$\lim_m \text{Map}_{cont}(G, \Lambda/I^m) = \text{Map}_{cont}(G, \Lambda)$
et de même pour les puissances supérieures de $G$.
À ce stade, Cohomologie sur les sites, lemmes
\ref{sites-cohomology-lemma-cech-cohomology}  et
\ref{sites-cohomology-lemma-cech-spectral-sequence-application}
montrent que
$$
\Lambda \to \text{Maps}_{cont}(G, \Lambda) \to
\text{Maps}_{cont}(G \times G, \Lambda) \to \ldots
$$
calcule la cohomologie pro-\'etale. En d'autres termes, nous voyons que
$$
H^i(\Spec(k)_\proetale, \underline{\Lambda}^\wedge)
=
H^i_{cont}(G, \Lambda)
$$
où le membre de droite est la cohomologie continue de Tate; voir
Cohomologie \'etale, section
\ref{etale-cohomology-section-continuous-group-cohomology}.
Bien entendu, il devait en être ainsi.

\begin{lemma}
\label{proetale-lemma-more-general-point}
Soit $k$ un corps. Soit $G = \text{Gal}(k^{sep}/k)$ son groupe de
Galois absolu. Considérons en outre l'un des deux cas suivants :
\begin{enumerate}
\item soit $M$ un groupe abélien profini muni d'une action continue
de $G$, ou
\item soit $\Lambda$ un anneau noethérien et $I \subset \Lambda$ un idéal,
et soit $M$ un module $I$-adiquement complet sur $\Lambda$, muni d'une action continue
de $G$.
\end{enumerate}
Il existe alors un faisceau canonique $\underline{M}^\wedge$ sur
$\Spec(k)_\proetale$ associé à $M$ tel que
$$
H^i(\Spec(k), \underline{M}^\wedge) = H^i_{cont}(G, M)
$$
comme groupes abéliens ou comme $\Lambda$-modules.
\end{lemma}

\begin{proof}
Démonstration dans le cas (2). Posons $M_n = M/I^nM$. Alors $M = \lim M_n$, car $M$ est
supposé complet pour la topologie $I$-adique. Puisque l'action de $G$ est continue,
on obtient des actions continues de $G$ sur $M_n$. D'après Cohomologie \'etale, théorème
\ref{etale-cohomology-theorem-equivalence-sheaves-point},
cette action correspond à un faisceau (localement constant)
$\underline{M_n}$ de $\Lambda/I^n$-modules sur $\Spec(k)_\etale$.
Prenons son image inverse sur $\Spec(k)_\proetale$ par le morphisme de comparaison
$\epsilon$ et prenons la limite
$$
\underline{M}^\wedge = \lim \epsilon^{-1}\underline{M_n}
$$
pour obtenir le faisceau annoncé dans le lemme. Le même argument que
dans l'introduction de cette section fournit la comparaison
avec la cohomologie galoisienne continue de Tate.
\end{proof}






\section{Fonctorialité du site pro-\'etale}
\label{proetale-section-morphism}

\noindent
Soit $f : X \to Y$ un morphisme de schémas. Le foncteur
$Y_\proetale \to X_\proetale$, $V \mapsto X \times_Y V$
induit un morphisme de sites $f_\proetale : X_\proetale \to Y_\proetale$; voir
Sites, proposition \ref{sites-proposition-get-morphism}.
En fait, on obtient un diagramme commutatif de morphismes de sites
$$
\xymatrix{
X_\proetale \ar[r]_\epsilon \ar[d]_{f_\proetale} &
X_\etale \ar[d]^{f_\etale} \\
Y_\proetale \ar[r]^\epsilon & Y_\etale
}
$$
où $\epsilon$ est défini comme dans la section \ref{proetale-section-comparison}.
En particulier, on a
$\epsilon^{-1} f_\etale^{-1} = f_\proetale^{-1} \epsilon^{-1}$.
Voici le résultat correspondant pour le foncteur image directe.

\begin{lemma}
\label{proetale-lemma-morphism-comparison}
Soit $f : X \to Y$ un morphisme de schémas, supposé quasi-compact et
quasi-séparé.
\begin{enumerate}
\item Soit $\mathcal{F}$ un faisceau d'ensembles sur $X_\etale$. Alors on a
$f_{\proetale, *}\epsilon^{-1}\mathcal{F} =
\epsilon^{-1}f_{\etale, *}\mathcal{F}$.
\item Soit $\mathcal{F}$ un faisceau abélien sur $X_\etale$. Alors on a
$Rf_{\proetale, *}\epsilon^{-1}\mathcal{F} =
\epsilon^{-1}Rf_{\etale, *}\mathcal{F}$.
\end{enumerate}
\end{lemma}

\begin{proof}
Démonstration de (1). Soit $\mathcal{F}$ un faisceau d'ensembles sur $X_\etale$. Il existe
un morphisme canonique $\epsilon^{-1}f_{\etale, *}\mathcal{F} \to
f_{\proetale, *}\epsilon^{-1}\mathcal{F}$; voir
Sites, section \ref{sites-section-pullback}.
Pour montrer que c'est un isomorphisme, on peut travailler localement sur $Y$ pour la topologie de Zariski; ainsi,
on peut supposer $Y$ affine. Dans ce cas,
tout objet de $Y_\proetale$ admet un recouvrement par des objets $V = \lim V_i$
qui sont des limites de schémas affines $V_i$ \'etales sur $Y$ (d'après la
proposition \ref{proetale-proposition-weakly-etale},
par exemple). En évaluant le morphisme
$\epsilon^{-1}f_{\etale, *}\mathcal{F} \to
f_{\proetale, *}\epsilon^{-1}\mathcal{F}$
en $V$, on obtient un morphisme
$$
\colim \Gamma(X \times_Y V_i, \mathcal{F})
\longrightarrow
\Gamma(X \times_Y V, \epsilon^*\mathcal{F})
$$
Le membre de gauche est décrit dans le lemme \ref{proetale-lemma-limit-pullback}.
D'après le lemme \ref{proetale-lemma-describe-pullback-from-etale}, on a
$$
\Gamma(X \times_Y V, \epsilon^*\mathcal{F}) =
\Gamma(X \times_Y V, g_\etale^{-1}\mathcal{F})
$$
où $g : X \times_Y V \to X$ est la projection. Le résultat découle donc de
Cohomologie \'etale, lemme
\ref{etale-cohomology-lemma-directed-colimit-cohomology}.

\medskip\noindent
Démonstration de (2). En raisonnant exactement comme ci-dessus,
on voit qu'il suffit de montrer que le morphisme
$$
\colim H^i_\etale(X \times_Y V_i, \mathcal{F})
\longrightarrow
H^i_\etale(X \times_Y V, \mathcal{F})
$$
est un isomorphisme, ce qui résulte encore de Cohomologie \'etale, lemme
\ref{etale-cohomology-lemma-directed-colimit-cohomology}.
\end{proof}







\section{Morphismes finis et sites pro-\'etales}
\label{proetale-section-finite}

\noindent
Il n'est pas clair qu'un morphisme fini de schémas détermine
un foncteur image directe exact sur les faisceaux abéliens pro-\'etales.

\begin{lemma}
\label{proetale-lemma-finite}
Soit $f : Z \to X$ un morphisme fini de schémas qui est
localement de présentation finie. Alors
$f_{\proetale, *} : \textit{Ab}(Z_\proetale) \to \textit{Ab}(X_\proetale)$
est exact.
\end{lemma}

\begin{proof}
Pour le démontrer, on peut travailler localement sur $X$ pour la topologie de Zariski et supposer que $X$
est affine, écrivons $X = \Spec(A)$. Alors $Z = \Spec(B)$ pour une certaine
$A$-algèbre $B$ finie et de présentation finie. La construction donnée dans la démonstration de la
proposition \ref{proetale-proposition-find-w-contractible}
produit un homomorphisme d'anneaux fidèlement plat et ind-\'etale $A \to D$
avec $D$ w-contractile. On peut vérifier l'exactitude d'une suite de
faisceaux par évaluation sur un tel objet $U = \Spec(D)$. Alors
$f_{\proetale, *}\mathcal{F}$
évalué en $U$ est égal à $\mathcal{F}$ évalué en
$V = \Spec(D \otimes_A B)$. Puisque $D \otimes_A B$ est w-contractile
d'après le lemme \ref{proetale-lemma-finite-finitely-presented-over-w-contractible},
l'évaluation en $V$ est exacte.
\end{proof}








\section{Immersions fermées et sites pro-\'etales}
\label{proetale-section-closed-immersion}

\noindent
Il n'est pas clair (et c'est probablement faux) qu'une immersion fermée de schémas
détermine un foncteur image directe exact sur les faisceaux abéliens pro-\'etales.

\begin{lemma}
\label{proetale-lemma-closed-immersion-affines}
Soit $i : Z \to X$ une immersion fermée de schémas affines.
Notons $X_{app}$ et $Z_{app}$ les sites introduits dans le
lemme \ref{proetale-lemma-affine-alternative}.
Le foncteur de changement de base
$$
u : X_{app} \to Z_{app},\quad U \longmapsto u(U) = U \times_X Z
$$
est continu et admet un adjoint à gauche pleinement fidèle $v$.
Pour $V$ dans $Z_{app}$, le morphisme $V \to v(V)$ est une immersion fermée
qui identifie $V$ à $u(v(V)) = v(V) \times_X Z$, et tout point de
$v(V)$ se spécialise en un point de $V$.
Le foncteur $v$ est cocontinu et transforme les recouvrements en recouvrements.
\end{lemma}

\begin{proof}
L'existence de l'adjoint découle immédiatement du
lemme \ref{proetale-lemma-lift-ind-etale} et des définitions.
Il est clair que $u$ est continu d'après la définition des
recouvrements dans $X_{app}$.

\medskip\noindent
Écrivons $X = \Spec(A)$ et $Z = \Spec(A/I)$. Soit $V = \Spec(\overline{C})$
un objet de $Z_{app}$, et soit $v(V) = \Spec(C)$.
Nous avons vu dans l'énoncé du lemme \ref{proetale-lemma-lift-ind-etale}
que $V$ est égal à $v(V) \times_X Z = \Spec(C/IC)$.
Tout $g \in C$ dont l'image dans
$C/IC = \overline{C}$ est inversible est lui-même inversible dans $C$. En effet, on a les
homomorphismes de $A$-algèbres $C \to C_g \to C/IC$ et, par adjonction,
on obtient un homomorphisme de $C$-algèbres $C_g \to C$.
Ainsi, tout point de $v(V)$ se spécialise en un point de $V$.

\medskip\noindent
Supposons que $\{V_i \to V\}$ soit un recouvrement dans $Z_{app}$.
Alors $\{v(V_i) \to v(V)\}$ est une famille finie de morphismes de
$Z_{app}$ telle que tout point de $V \subset v(V)$ appartienne
à l'image de l'un des morphismes $v(V_i) \to v(V)$. Comme les
morphismes $v(V_i) \to v(V)$ sont plats (puisqu'ils sont faiblement \'etales),
on conclut que $\{v(V_i) \to v(V)\}$ est conjointement surjective.
Cela démontre que $v$ transforme les recouvrements en recouvrements.

\medskip\noindent
Soit $V$ un objet de $Z_{app}$, et soit $\{U_i \to v(V)\}$
un recouvrement dans $X_{app}$. On voit alors que
$\{u(U_i) \to u(v(V)) = V\}$ est un recouvrement dans $Z_{app}$.
Par adjonction, on obtient des morphismes $v(u(U_i)) \to U_i$.
Ainsi, la famille $\{v(u(U_i)) \to v(V)\}$ raffine le
recouvrement donné, et l'on conclut que $v$ est cocontinu.
\end{proof}

\begin{lemma}
\label{proetale-lemma-closed-immersion-affines-apply}
Soit $Z \to X$ une immersion fermée de schémas affines.
Le morphisme de topos correspondant $i = i_\proetale$
est égal au morphisme de topos
associé au foncteur pleinement fidèle et cocontinu
$v : Z_{app} \to X_{app}$ du lemme \ref{proetale-lemma-closed-immersion-affines}.
Il en résulte que
\begin{enumerate}
\item $i^{-1}\mathcal{F}$ est le faisceau associé au préfaisceau
$V \mapsto \mathcal{F}(v(V))$,
\item pour un objet faiblement contractile $V$ de $Z_{app}$, on a
$i^{-1}\mathcal{F}(V) = \mathcal{F}(v(V))$,
\item $i^{-1} : \Sh(X_\proetale) \to \Sh(Z_\proetale)$
admet un adjoint à gauche $i^{Sh}_!$,
\item $i^{-1} : \textit{Ab}(X_\proetale) \to \textit{Ab}(Z_\proetale)$
admet un adjoint à gauche $i_!$,
\item $\text{id} \to i^{-1}i^{Sh}_!$, $\text{id} \to i^{-1}i_!$, et
$i^{-1}i_* \to \text{id}$ sont des isomorphismes, et
\item $i_*$, $i^{Sh}_!$ et $i_!$ sont pleinement fidèles.
\end{enumerate}
\end{lemma}

\begin{proof}
D'après le lemme \ref{proetale-lemma-affine-alternative}, on peut décrire $i_\proetale$ au moyen
du morphisme de sites $u : X_{app} \to Z_{app}$, $V \mapsto V \times_X Z$.
La première assertion du lemme découle de
Sites, lemmes \ref{sites-lemma-have-functor-other-way} et
\ref{sites-lemma-have-functor-other-way-morphism}
(mais en inversant les rôles de $u$ et $v$).

\medskip\noindent
Démonstration de (1). D'après la description de $i$ comme morphisme de topos associé
à $v$, cela résulte de la construction; voir
Sites, lemme \ref{sites-lemma-cocontinuous-morphism-topoi}.

\medskip\noindent
Démonstration de (2). Puisque le foncteur $v$ transforme les recouvrements en recouvrements d'après le
lemme \ref{proetale-lemma-closed-immersion-affines}, on voit que le préfaisceau
$\mathcal{G} : V \mapsto \mathcal{F}(v(V))$ est un préfaisceau séparé
(Sites, définition \ref{sites-definition-separated}). Par conséquent,
le faisceau associé à $\mathcal{G}$ est $\mathcal{G}^+$; voir
Sites, théorème \ref{sites-theorem-plus}. Soit ensuite $V$ un objet
contractile de $Z_{app}$. Soit
$\mathcal{V} = \{V_i \to V\}_{i = 1, \ldots, n}$
un recouvrement quelconque dans $Z_{app}$. Posons $\mathcal{V}' = \{\coprod V_i \to V\}$.
Puisque $v$ commute aux sommes disjointes finies (en tant qu'adjoint à gauche ou par
construction) et que $\mathcal{F}$ transforme les sommes
disjointes finies en produits, on voit que
$$
H^0(\mathcal{V}, \mathcal{G}) = H^0(\mathcal{V}', \mathcal{G})
$$
(notations de Sites, section \ref{sites-section-sheafification};
comparer avec
Cohomologie \'etale, lemme \ref{etale-cohomology-lemma-cech-complex}).
On peut donc supposer que le recouvrement est donné par un unique morphisme, disons
$\{V' \to V\}$. Puisque $V$ est faiblement contractile, ce recouvrement
peut être raffiné par le recouvrement trivial $\{V \to V\}$.
Il en résulte que la valeur de $\mathcal{G}^+ = i^{-1}\mathcal{F}$
en $V$ est simplement $\mathcal{F}(v(V))$, ce qui démontre (2).

\medskip\noindent
Démonstration de (3). Tout objet de $Z_{app}$ admet un recouvrement par des objets
contractiles (lemme \ref{proetale-lemma-proetale-enough-w-contractible}).
D'après ce qui précède, on aurait $i^{Sh}_!h_V = h_{v(V)}$ pour $V$
faiblement contractile si $i^{Sh}_!$ existait. L'existence de
$i^{Sh}_!$ découle alors de
Sites, lemme \ref{sites-lemma-existence-lower-shriek}.

\medskip\noindent
Démonstration de (4). L'existence de $i_!$ s'obtient de la même manière en posant
$i_!\mathbf{Z}_V = \mathbf{Z}_{v(V)}$ pour $V$ faiblement contractile dans $Z_{app}$,
en procédant de même pour les sommes directes, et en appliquant le
lemme d'Homologie \ref{homology-lemma-partially-defined-adjoint}.
Nous omettons les détails.

\medskip\noindent
Démonstration de (5). Soit $V$ un objet contractile de $Z_{app}$.
Alors $i^{-1}i^{Sh}_!h_V = i^{-1}h_{v(V)} = h_{u(v(V))} = h_V$.
(C'est un fait général que $i^{-1}h_U = h_{u(U)}$.) Puisque les
faisceaux $h_V$, pour $V$ contractile, engendrent $\Sh(Z_{app})$
(Sites, lemme \ref{sites-lemma-sheaf-coequalizer-representable}),
on conclut que $\text{id} \to i^{-1}i^{Sh}_!$ est un isomorphisme.
Il en va de même du morphisme $\text{id} \to i^{-1}i_!$. Ensuite,
$(i^{-1}i_*\mathcal{H})(V) = i_*\mathcal{H}(v(V)) =
\mathcal{H}(u(v(V))) = \mathcal{H}(V)$, et l'on trouve que
$i^{-1}i_* \to \text{id}$ est un isomorphisme.

\medskip\noindent
Les assertions de pleine fidélité de (6) résultent alors de
Catégories, lemme \ref{categories-lemma-adjoint-fully-faithful}.
\end{proof}

\begin{lemma}
\label{proetale-lemma-closed-immersion}
Soit $i : Z \to X$ une immersion fermée de schémas. Alors
\begin{enumerate}
\item $i_\proetale^{-1}$ commute aux limites,
\item $i_{\proetale, *}$ est pleinement fidèle, et
\item $i_\proetale^{-1}i_{\proetale, *} \cong \text{id}_{\Sh(Z_\proetale)}$.
\end{enumerate}
\end{lemma}

\begin{proof}
Les assertions (2) et (3) sont équivalentes d'après le
lemme de Sites \ref{sites-lemma-exactness-properties}.
Les parties (1) et (3) sont locales sur $X$ pour la topologie de Zariski; on peut donc supposer que
$X$ est affine. Dans ce cas, le résultat découle du
lemme \ref{proetale-lemma-closed-immersion-affines-apply}.
\end{proof}

\begin{lemma}
\label{proetale-lemma-thickening}
Soit $i : Z \to X$ un morphisme de schémas entier, universellement injectif et surjectif.
Alors
$i_{\proetale, *}$ et $i_\proetale^{-1}$ définissent des équivalences de catégories
quasi inverses entre les topos pro-\'etales.
\end{lemma}

\begin{proof}
On se ramène immédiatement au cas où $X$ est affine.
Alors $Z$ est également affine. Posons $A = \mathcal{O}(X)$ et $B = \mathcal{O}(Z)$.
Les catégories des algèbres \'etales sur
$A$ et $B$ sont alors équivalentes; voir 
Cohomologie \'etale, théorème
\ref{etale-cohomology-theorem-topological-invariance} et la
remarque \ref{etale-cohomology-remark-affine-inside-equivalence}.
Les catégories des algèbres ind-\'etales sur $A$ et $B$ sont donc
équivalentes. Autrement dit, les catégories $X_{app}$ et $Z_{app}$
du lemme \ref{proetale-lemma-affine-alternative} sont équivalentes.
Nous omettons de vérifier
que cette équivalence transforme les recouvrements en recouvrements et réciproquement.
Le résultat découle alors du lemme \ref{proetale-lemma-affine-alternative},
qui identifie le topos pro-\'etale au topos des faisceaux sur $X_{app}$.
\end{proof}

\begin{lemma}
\label{proetale-lemma-compute-i-star}
Soit $i : Z \to X$ une immersion fermée de schémas.
Soit $U \to X$ un objet de $X_\proetale$ tel que
\begin{enumerate}
\item $U$ soit affine et faiblement contractile, et
\item tout point de $U$ se spécialise en un point de $U \times_X Z$.
\end{enumerate}
Alors $i_\proetale^{-1}\mathcal{F}(U \times_X Z) = \mathcal{F}(U)$
pour tout faisceau abélien sur $X_\proetale$.
\end{lemma}

\begin{proof}
Puisque l'image inverse commute à la restriction, on peut remplacer $X$ par $U$.
Nous pouvons donc supposer que $X$ est affine et faiblement contractile
et que tout point de $X$ se spécialise en un point de $Z$.
D'après le lemme \ref{proetale-lemma-closed-immersion-affines-apply}, partie (1),
il suffit de montrer que $v(Z) = X$ dans ce cas.
Nous devons donc montrer ceci : si $A$ est un anneau w-contractile, si $I \subset A$
est un idéal contenu dans le radical de Jacobson de $A$ et si $A \to B \to A/I$
est une factorisation telle que $A \to B$ soit ind-\'etale, alors il existe
une unique rétraction $B \to A$ compatible avec les morphismes vers $A/I$.
Observons que $B/IB = A/I \times R$ comme $A/I$-algèbres.
Après avoir remplacé $B$ par une localisation, nous pouvons supposer que $B/IB = A/I$.
Notons que $\Spec(B) \to \Spec(A)$ est surjectif, puisque son image
contient $V(I)$, donc tous les points fermés, et qu'elle est stable par
spécialisation. Comme $A$ est w-contractile, il existe une rétraction $B \to A$.
Puisque $B/IB = A/I$, cette rétraction est compatible avec le morphisme vers $A/I$.
Nous omettons la démonstration de l'unicité (indication : utiliser le fait que $A$ et $B$ ont
des anneaux locaux isomorphes en tout idéal maximal de $A$).
\end{proof}

\begin{lemma}
\label{proetale-lemma-closed-immersion-complement-retrocompact-exact}
Soit $i : Z \to X$ une immersion fermée de schémas.
Si $X \setminus i(Z)$ est un ouvert rétrocompact de $X$, alors
$i_{\proetale, *}$ est exact.
\end{lemma}

\begin{proof}
La question est locale sur $X$; nous pouvons donc supposer que $X$ est affine.
Écrivons $X = \Spec(A)$ et $Z = \Spec(A/I)$. Il existe
$f_1, \ldots, f_r \in I$ tels que $Z = V(f_1, \ldots, f_r)$
en tant qu'ensembles ; voir Algèbre, lemme \ref{algebra-lemma-qc-open}.
D'après le lemme \ref{proetale-lemma-thickening}, nous pouvons supposer que
$Z = \Spec(A/(f_1, \ldots, f_r))$. Dans ce cas,
le foncteur $i_{\proetale, *}$ est exact d'après le
lemme \ref{proetale-lemma-finite}.
\end{proof}





\section{Prolongement par zéro}
\label{proetale-section-extension-by-zero}

\noindent
Les résultats généraux de
Modules sur les sites, section \ref{sites-modules-section-localize},
nous permettent de poser la définition suivante.

\begin{definition}
\label{proetale-definition-extension-zero}
Soit $j : U \to X$ un morphisme de schémas faiblement \'etale.
\begin{enumerate}
\item Le foncteur de restriction
$j^{-1} : \Sh(X_\proetale) \to \Sh(U_\proetale)$
admet pour adjoint à gauche
$j_!^{Sh} : \Sh(X_\proetale) \to \Sh(U_\proetale)$.
\item Le foncteur de restriction
$j^{-1} : \textit{Ab}(X_\proetale) \to \textit{Ab}(U_\proetale)$
admet un adjoint à gauche, noté
$j_! : \textit{Ab}(U_\proetale) \to \textit{Ab}(X_\proetale)$
et appelé {\it prolongement par zéro}.
\item Soit $\Lambda$ un anneau. Le foncteur
$j^{-1} : \textit{Mod}(X_\proetale, \Lambda) \to
\textit{Mod}(U_\proetale, \Lambda)$
admet un adjoint à gauche, noté
$j_! : \textit{Mod}(U_\proetale, \Lambda) \to
\textit{Mod}(X_\proetale, \Lambda)$
et appelé {\it prolongement par zéro}.
\end{enumerate}
\end{definition}

\noindent
Comme d'habitude, comparons cette situation au cas \'etale.

\begin{lemma}
\label{proetale-lemma-jshriek-comparison}
Soit $j : U \to X$ un morphisme de schémas \'etale.
Soit $\mathcal{G}$ un faisceau abélien sur $U_\etale$.
Alors $\epsilon^{-1} j_!\mathcal{G} = j_!\epsilon^{-1}\mathcal{G}$
comme faisceaux sur $X_\proetale$.
\end{lemma}

\begin{proof}
Cela est vrai parce que les deux foncteurs sont adjoints à gauche de
$j_\proetale^{-1} \epsilon_* = \epsilon_* j_\etale^{-1}$.
L'égalité résulte de la discussion à la section \ref{proetale-section-morphism}.
\end{proof}

\begin{lemma}
\label{proetale-lemma-jshriek-zero}
Soit $j : U \to X$ un morphisme de schémas faiblement \'etale.
Soit $i : Z \to X$ une immersion fermée telle que $U \times_X Z = \emptyset$.
Soit $V \to X$ un objet affine de $X_\proetale$ tel que tout point
de $V$ se spécialise en un point de $V_Z = Z \times_X V$.
Alors $j_!\mathcal{F}(V) = 0$ pour tout faisceau abélien sur $U_\proetale$.
\end{lemma}

\begin{proof}
Soit $\{V_i \to V\}$ un recouvrement pro-\'etale. Il suffit de pouvoir
raffiner ce recouvrement en un recouvrement dont les membres n'admettent aucun
morphisme vers $U$ au-dessus de $X$ (voir la construction de $j_!$ dans
Modules sur les sites, section \ref{sites-modules-section-localize}).
Raffinons d'abord le recouvrement pour obtenir un recouvrement fini dont les $V_i$ sont affines.
Pour tout $i$, écrivons $V_i = \Spec(A_i)$ et notons $Z_i \subset V_i$ le
sous-schéma image inverse de $Z$.
Posons $W_i = \Spec(A_{i, Z_i}^\sim)$ avec les notations du
lemme \ref{proetale-lemma-localization}.
Alors $\coprod W_i \to V$ est faiblement \'etale et son image contient tous les
points de $V_Z$. Elle contient donc tous les points de $V$ d'après
notre hypothèse sur les spécialisations. Ainsi, $\{W_i \to V\}$ est un
recouvrement pro-\'etale qui raffine le recouvrement donné. Mais tout point de $W_i$
se spécialise en un point situé au-dessus de $Z$; il n'existe donc aucun morphisme
$W_i \to U$ au-dessus de $X$.
\end{proof}

\begin{lemma}
\label{proetale-lemma-open-immersion}
Soit $j : U \to X$ une immersion ouverte de schémas.
Alors $\text{id} \cong j^{-1}j_!$ et $j^{-1}j_* \cong \text{id}$,
et les foncteurs $j_!$ et $j_*$ sont pleinement fidèles.
\end{lemma}

\begin{proof}
Voir Modules sur les sites, lemme \ref{sites-modules-lemma-restrict-back}
(et Sites, lemme \ref{sites-lemma-restrict-back} pour le cas
des faisceaux d'ensembles), ainsi que
Catégories, lemme \ref{categories-lemma-adjoint-fully-faithful}.
\end{proof}

\noindent
Voici la relation entre le prolongement par zéro et la restriction
au sous-schéma fermé complémentaire.

\begin{lemma}
\label{proetale-lemma-ses-associated-to-open}
Soit $X$ un schéma. Soit $Z \subset X$ un sous-schéma fermé et soit
$U \subset X$ son complémentaire. Notons $i : Z \to X$ et $j : U \to X$
les morphismes d'inclusion. Supposons que $j$ soit quasi-compact.
Pour tout faisceau abélien sur $X_\proetale$, il existe une suite exacte courte
canonique
$$
0 \to j_!j^{-1}\mathcal{F} \to \mathcal{F} \to i_*i^{-1}\mathcal{F} \to 0
$$
sur $X_\proetale$, où tous les foncteurs se rapportent à la topologie pro-\'etale.
\end{lemma}

\begin{proof}
Les propriétés d'adjonction des foncteurs en jeu fournissent les morphismes.
Il suffit de montrer que $X_\proetale$ admet suffisamment d'objets
(Sites, définition \ref{sites-definition-w-contractible}) sur lesquels
la suite s'évalue en une suite exacte courte.
Soit $V = \Spec(A)$ un objet affine de $X_\proetale$
tel que $A$ soit w-contractile (il y a suffisamment d'objets de ce type).
Alors $V \times_X Z$ est défini par un idéal $I \subset A$.
L'hypothèse que $j$ est quasi-compact implique qu'il existe
$f_1, \ldots, f_r \in I$ tels que $V(I) = V(f_1, \ldots, f_r)$.
Nous obtenons un homomorphisme d'anneaux fidèlement plat et ind-Zariski
$$
A \longrightarrow A_{f_1} \times \ldots \times A_{f_r} \times
A_{V(I)}^\sim
$$
où $A_{V(I)}^\sim$ est défini comme dans le lemme \ref{proetale-lemma-localization}.
Puisque $V_i = \Spec(A_{f_i}) \to X$ se factorise par $U$, on a
$$
j_!j^{-1}\mathcal{F}(V_i) = \mathcal{F}(V_i)
\quad\text{et}\quad
i_*i^{-1}\mathcal{F}(V_i) = 0
$$
D'autre part, pour le schéma $V^\sim = \Spec(A_{V(I)}^\sim)$,
on a
$$
j_!j^{-1}\mathcal{F}(V^\sim) = 0
\quad\text{et}\quad
\mathcal{F}(V^\sim) = i_*i^{-1}\mathcal{F}(V^\sim)
$$
La première égalité résulte du lemme \ref{proetale-lemma-jshriek-zero},
et la seconde des
lemmes \ref{proetale-lemma-compute-i-star} et \ref{proetale-lemma-localization-w-contractible}.
La suite s'évalue donc en une suite exacte sur
$\Spec(A_{f_1} \times \ldots \times A_{f_r} \times A_{V(I)}^\sim)$,
ce qui démontre le lemme.
\end{proof}

\begin{lemma}
\label{proetale-lemma-j-shriek-limits}
Soit $j : U \to X$ une immersion ouverte quasi-compacte
de schémas. Le foncteur
$j_! : \textit{Ab}(U_\proetale) \to \textit{Ab}(X_\proetale)$
commute aux limites.
\end{lemma}

\begin{proof}
Puisque $j_!$ est exact, il suffit de montrer que $j_!$ commute aux produits.
La question est locale sur $X$; nous pouvons donc supposer que $X$ est affine.
Soit $\mathcal{G}$ un faisceau abélien sur $U_\proetale$.
On a $j^{-1}j_*\mathcal{G} = \mathcal{G}$. En appliquant alors
la suite exacte du lemme \ref{proetale-lemma-ses-associated-to-open},
on obtient
$$
0 \to j_!\mathcal{G} \to j_*\mathcal{G} \to i_*i^{-1}j_*\mathcal{G} \to 0
$$
où $i : Z \to X$ est l'inclusion du schéma muni de la structure réduite
induite sur le complément $Z = X \setminus U$.
Les foncteurs $j_*$ et $i_*$ commutent aux produits en tant qu'adjoints à droite.
Le foncteur $i^{-1}$ commute aux produits d'après le
lemme \ref{proetale-lemma-closed-immersion}.
Il en va donc de même pour $j_!$, car, sur le site pro-\'etale, les produits
sont exacts
(Cohomologie sur les sites, proposition
\ref{sites-cohomology-proposition-enough-weakly-contractibles}).
\end{proof}




\section{Faisceaux constructibles sur le site pro-\'etale}
\label{proetale-section-constructible}

\noindent
Nous nous limitons aux faisceaux constructibles de $\Lambda$-modules pour un
anneau noethérien. Nous envisageons d'étudier ultérieurement les faisceaux constructibles
d'ensembles, de groupes, etc.

\begin{definition}
\label{proetale-definition-constructible}
Soit $X$ un schéma.
Soit $\Lambda$ un anneau noethérien. Un faisceau de $\Lambda$-modules
sur $X_\proetale$ est {\it constructible} si, pour tout ouvert affine
$U \subset X$, il existe une décomposition finie
de $U$ en sous-schémas localement fermés constructibles
$U = \coprod_i U_i$ telle que
$\mathcal{F}|_{U_i}$ soit de type fini et localement constant pour tout $i$.
\end{definition}

\noindent
Là encore, cela ne donne rien de ``nouveau''.

\begin{lemma}
\label{proetale-lemma-compare-constructible}
Soit $X$ un schéma. Soit $\Lambda$ un anneau noethérien.
Le foncteur $\epsilon^{-1}$ définit une équivalence de catégories
$$
\left\{
\begin{matrix}
\text{faisceaux constructibles de}\\
\Lambda\text{-modules sur }X_\etale\\
\end{matrix}
\right\}
\longleftrightarrow
\left\{
\begin{matrix}
\text{faisceaux constructibles de}\\
\Lambda\text{-modules sur }X_\proetale\\
\end{matrix}
\right\}
$$
entre les faisceaux constructibles de $\Lambda$-modules sur $X_\etale$
et les faisceaux constructibles de $\Lambda$-modules sur $X_\proetale$.
\end{lemma}

\begin{proof}
D'après le lemme \ref{proetale-lemma-fully-faithful}, le foncteur $\epsilon^{-1}$
est pleinement fidèle et commute à l'image inverse, c'est-à-dire à la restriction aux
strates. Ainsi, l'image par $\epsilon^{-1}$ d'un faisceau \'etale
constructible est un faisceau pro-\'etale constructible. Pour achever la
démonstration, soit $\mathcal{F}$ un faisceau constructible de $\Lambda$-modules
sur $X_\proetale$ comme dans la définition \ref{proetale-definition-constructible}.
Il existe un morphisme canonique
$$
\epsilon^{-1}\epsilon_*\mathcal{F} \longrightarrow \mathcal{F}
$$
Nous allons montrer que ce morphisme est un isomorphisme. Cela prouvera que
$\mathcal{F}$ appartient à l'image essentielle de $\epsilon^{-1}$
et achèvera la démonstration (nous omettons les détails).

\medskip\noindent
Puisqu'il suffit de le démontrer localement sur $X$, on peut supposer que $X$ est
quasi-compact et quasi-séparé et que l'on dispose
d'une partition finie $X = \coprod_{i = 1, \ldots, n} X_i$
en strates localement fermées constructibles telles que $\mathcal{F}|_{X_i}$
soit localement constant de type fini. Nous raisonnons par récurrence sur $n$.
Le cas initial $n = 1$ résulte du lemme \ref{proetale-lemma-compare-locally-constant}.
Prenons un point $x \in X$; alors $x \in X_k$ pour un certain $k$.
Il suffit de montrer que le morphisme affiché est un isomorphisme dans un
voisinage ouvert de $x$. On peut donc supposer que $X_k$ est fermé dans $X$.
Posons $Z = X_k$ et notons $U \subset X$ le complément de $Z$.
Observons que l'hypothèse de récurrence s'applique à la restriction de
$\mathcal{F}$ à $Z$ et à $U$. D'après le lemme \ref{proetale-lemma-ses-associated-to-open},
on a une suite exacte courte
$$
0 \to j_!j^{-1}\mathcal{F} \to \mathcal{F} \to i_*i^{-1}\mathcal{F} \to 0
$$
sur $X_\proetale$. La fonctorialité donne un diagramme commutatif
$$
\xymatrix{
0 \ar[r] &
\epsilon^{-1}\epsilon_*j_!j^{-1}\mathcal{F} \ar[r] \ar[d] &
\epsilon^{-1}\epsilon_*\mathcal{F} \ar[r] \ar[d] &
\epsilon^{-1}\epsilon_*i_*i^{-1}\mathcal{F} \ar[r] \ar[d] & 0 \\
0 \ar[r] &
j_!j^{-1}\mathcal{F} \ar[r] &
\mathcal{F} \ar[r] &
i_*i^{-1}\mathcal{F} \ar[r] & 0
}
$$
Par récurrence, on sait, d'une part, que
$\epsilon^{-1}\epsilon_*i^{-1}\mathcal{F} \to i^{-1}\mathcal{F}$
et
$\epsilon^{-1}\epsilon_*j^{-1}\mathcal{F} \to j^{-1}\mathcal{F}$
sont des isomorphismes et, d'autre part, que
$i^{-1}\mathcal{F} = \epsilon^{-1}\mathcal{A}$
et
$j^{-1}\mathcal{F} = \epsilon^{-1}\mathcal{B}$
pour certains faisceaux constructibles de $\Lambda$-modules
$\mathcal{A}$ sur $Z_\etale$ et $\mathcal{B}$ sur $U_\etale$.
Alors
$$
\epsilon^{-1}\epsilon_*j_!j^{-1}\mathcal{F} =
\epsilon^{-1}\epsilon_*j_!\epsilon^{-1}\mathcal{B} =
\epsilon^{-1}\epsilon_*\epsilon^{-1}j_!\mathcal{B} =
\epsilon^{-1}j_!\mathcal{B} =
j_!\epsilon^{-1}\mathcal{B} =
j_!j^{-1}\mathcal{F}
$$
où la deuxième égalité résulte du lemme \ref{proetale-lemma-jshriek-comparison},
la troisième du lemme \ref{proetale-lemma-fully-faithful}, et la
quatrième, à nouveau, du lemme \ref{proetale-lemma-jshriek-comparison}.
De même, on a
$$
\epsilon^{-1}\epsilon_*i_*i^{-1}\mathcal{F} =
\epsilon^{-1}\epsilon_*i_*\epsilon^{-1}\mathcal{A} =
\epsilon^{-1}\epsilon_*\epsilon^{-1}i_*\mathcal{A} =
\epsilon^{-1}i_*\mathcal{A} =
i_*\epsilon^{-1}\mathcal{A} =
i_*i^{-1}\mathcal{F}
$$
Cette fois, on utilise le lemme \ref{proetale-lemma-morphism-comparison}.
Le lemme des cinq permet de conclure que le
morphisme vertical médian du grand diagramme est un isomorphisme.
\end{proof}

\begin{lemma}
\label{proetale-lemma-constructible-serre}
Soit $X$ un schéma. Soit $\Lambda$ un anneau noethérien.
La catégorie des faisceaux constructibles de $\Lambda$-modules sur $X_\proetale$
est une sous-catégorie de Serre faible de $\textit{Mod}(X_\proetale, \Lambda)$.
\end{lemma}

\begin{proof}
C'est une conséquence formelle des
lemmes \ref{proetale-lemma-compare-constructible} et \ref{proetale-lemma-compare-derived}
et du résultat correspondant pour le site \'etale
(Cohomologie \'etale, lemme \ref{etale-cohomology-lemma-constructible-abelian}).
\end{proof}

\begin{lemma}
\label{proetale-lemma-compare-constructible-derived}
Soit $X$ un schéma. Soit $\Lambda$ un anneau noethérien.
Soient $D_c(X_\etale, \Lambda)$, resp.\ $D_c(X_\proetale, \Lambda)$,
les sous-catégories pleines de
$D(X_\etale, \Lambda)$, resp.\ $D(X_\proetale, \Lambda)$,
formées des complexes dont les faisceaux de cohomologie sont des
faisceaux constructibles de $\Lambda$-modules. Alors
$$
\epsilon^{-1} :
D_c^+(X_\etale, \Lambda)
\longrightarrow
D_c^+(X_\proetale, \Lambda)
$$
est une équivalence de catégories.
\end{lemma}

\begin{proof}
Les catégories $D_c(X_\etale, \Lambda)$ et $D_c(X_\proetale, \Lambda)$
sont des sous-catégories strictement pleines, saturées et triangulées de
$D(X_\etale, \Lambda)$ et $D(X_\proetale, \Lambda)$ d'après
Cohomologie \'etale, lemme \ref{etale-cohomology-lemma-constructible-abelian}
et
le lemme \ref{proetale-lemma-constructible-serre}
et
Catégories dérivées, section \ref{derived-section-triangulated-sub}.
L'énoncé du lemme s'obtient en combinant les
lemmes \ref{proetale-lemma-compare-derived} et
\ref{proetale-lemma-compare-constructible}.
\end{proof}

\begin{lemma}
\label{proetale-lemma-tensor-c}
Soit $X$ un schéma. Soit $\Lambda$ un anneau noethérien.
Soient $K, L \in D_c^-(X_\proetale, \Lambda)$. Alors
$K \otimes_\Lambda^\mathbf{L} L$ appartient à $D_c^-(X_\proetale, \Lambda)$.
\end{lemma}

\begin{proof}
Remarquons que $H^i(K \otimes_\Lambda^\mathbf{L} L)$ s'identifie à
$H^i(\tau_{\geq i - 1}K \otimes_\Lambda^\mathbf{L} \tau_{\geq i - 1}L)$.
On peut donc supposer que $K$ et $L$ sont bornés.
Dans ce cas, le lemme \ref{proetale-lemma-compare-constructible-derived} permet de
se ramener au cas du site \'etale, voir
Cohomologie \'etale, lemme \ref{etale-cohomology-lemma-tensor-c}.
\end{proof}

\begin{lemma}
\label{proetale-lemma-compare-truncations-constructible}
Soit $X$ un schéma. Soit $\Lambda$ un anneau noethérien.
Soit $I \subset \Lambda$ un idéal.
Soit $K$ un objet de $D(X_\proetale, \Lambda)$.
Posons $K_n = K \otimes_\Lambda^\mathbf{L} \underline{\Lambda/I^n}$.
Si $K_1$ appartient à $D^-_c(X_\proetale, \Lambda/I)$, alors
$K_n$ appartient à $D^-_c(X_\proetale, \Lambda/I^n)$ pour tout $n$.
\end{lemma}

\begin{proof}
Considérons les triangles distingués
$$
K \otimes_\Lambda^\mathbf{L} \underline{I^n/I^{n + 1}} \to
K_{n + 1} \to
K_n \to
K \otimes_\Lambda^\mathbf{L} \underline{I^n/I^{n + 1}}[1]
$$
et les isomorphismes
$$
K \otimes_\Lambda^\mathbf{L} \underline{I^n/I^{n + 1}} =
K_1 \otimes_{\Lambda/I}^\mathbf{L} \underline{I^n/I^{n + 1}}
$$
D'après le lemme \ref{proetale-lemma-tensor-c}, ce produit tensoriel possède des
faisceaux de cohomologie constructibles (et nuls lorsque la cohomologie de $K_1$
est nulle). Une récurrence sur $n$ utilisant le
lemme \ref{proetale-lemma-constructible-serre}
montre donc que chaque $K_n$ possède des faisceaux de cohomologie constructibles.
\end{proof}








\section{Faisceaux adiques constructibles}
\label{proetale-section-adic}

\noindent
Dans cette section, nous définissons la notion de faisceau
constructible de $\Lambda$-modules, ainsi que quelques variantes.

\begin{definition}
\label{proetale-definition-adic}
Soit $\Lambda$ un anneau noethérien et soit $I \subset \Lambda$ un idéal.
Soit $X$ un schéma. Soit $\mathcal{F}$ un faisceau de $\Lambda$-modules
sur $X_\proetale$.
\begin{enumerate}
\item Nous disons que $\mathcal{F}$ est un {\it faisceau constructible de $\Lambda$-modules}
si $\mathcal{F} = \lim \mathcal{F}/I^n\mathcal{F}$ et si chaque
$\mathcal{F}/I^n\mathcal{F}$ est un faisceau constructible de $\Lambda/I^n$-modules.
\item Si $\mathcal{F}$ est un faisceau constructible de $\Lambda$-modules, nous disons
que $\mathcal{F}$ est {\it lisse} si chaque $\mathcal{F}/I^n\mathcal{F}$ est
localement constant.
\item Nous disons que $\mathcal{F}$ est {\it lisse adique}\footnote{Cette notation n'est
peut-être pas standard.} s'il existe un module
complet pour la topologie $I$-adique, muni d'une structure de $\Lambda$-module, noté $M$, tel que $M/IM$ soit fini
et que $\mathcal{F}$ soit localement isomorphe à
$$
\underline{M}^\wedge = \lim \underline{M/I^nM}.
$$
\item Nous disons que $\mathcal{F}$ est
{\it adique constructible}\footnote{Cette notation n'est peut-être pas standard.}
si, pour tout ouvert affine $U \subset X$,
il existe une décomposition $U = \coprod U_i$ en
sous-schémas localement fermés constructibles tels que $\mathcal{F}|_{U_i}$
soit lisse adique.
\end{enumerate}
\end{definition}

\noindent
La définition d'un faisceau constructible de $\Lambda$-modules équivaut
à celle de \cite[Expos\'e VI, d\'efinition 1.1.1]{SGA5} lorsque
$\Lambda = \mathbf{Z}_\ell$ et $I = (\ell)$. Il est clair que
l'on a les implications
$$
\xymatrix{
\text{lisse adique} \ar@{=>}[r] \ar@{=>}[d] &
\text{adique constructible} \ar@{=>}[d] \\
\text{faisceau constructible lisse de }\Lambda\text{-modules} \ar@{=>}[r] &
\text{faisceau constructible de }\Lambda\text{-modules}
}
$$
Les flèches verticales peuvent être inversées dans certains cas
(voir les lemmes \ref{proetale-lemma-Noetherian-constructible} et
\ref{proetale-lemma-Noetherian-adic-constructible}). En général,
ni la catégorie des faisceaux adiques constructibles, ni
la catégorie des faisceaux constructibles de $\Lambda$-modules ne sont stables
par noyaux et conoyaux.

\medskip\noindent
En effet, soit $X$ un schéma affine dont l'espace topologique sous-jacent $|X|$
est homéomorphe à $\Lambda = \mathbf{Z}_\ell$; voir
l'exemple \ref{proetale-example-construct-space}. Notons
$f : |X| \to \mathbf{Z}_\ell = \Lambda$
un homéomorphisme. Nous pouvons considérer $f$ comme une section de
$\underline{\Lambda}^\wedge$ sur $X$; la multiplication par $f$
définit alors un complexe à deux termes
$$
\underline{\Lambda}^\wedge \xrightarrow{f} \underline{\Lambda}^\wedge
$$
sur $X_\proetale$. Le faisceau $\underline{\Lambda}^\wedge$ est lisse adique.
Cependant, le conoyau du morphisme ci-dessus n'est pas adique constructible, car
le type d'isomorphisme de ses fibres prend une infinité
de valeurs : $\mathbf{Z}/\ell^n\mathbf{Z}$ et $\mathbf{Z}_\ell$.
Le conoyau est un faisceau constructible de $\mathbf{Z}_\ell$-modules.
Cependant, le noyau n'est même pas un faisceau constructible de $\mathbf{Z}_\ell$-modules,
car il est nul sur un ouvert non quasi-compact, mais n'est pas nul.

\begin{lemma}
\label{proetale-lemma-Noetherian-constructible}
Soit $X$ un schéma noethérien. Soit $\Lambda$ un anneau noethérien et
soit $I \subset \Lambda$ un idéal. Soit $\mathcal{F}$ un
faisceau constructible de $\Lambda$-modules sur $X_\proetale$.
Il existe alors une partition finie $X = \coprod X_i$ en
sous-schémas localement fermés tels que la restriction $\mathcal{F}|_{X_i}$
soit lisse.
\end{lemma}

\begin{proof}
Posons $R = \bigoplus I^n/I^{n + 1}$. Observons que $R$ est un anneau noethérien.
Puisque chacun des faisceaux
$\mathcal{F}/I^n\mathcal{F}$ est un faisceau constructible de
$\Lambda/I^n\Lambda$-modules, $I^n\mathcal{F}/I^{n + 1}\mathcal{F}$
est également un faisceau constructible de $\Lambda/I$-modules, donc l'image inverse
d'un faisceau constructible $\mathcal{G}_n$ sur $X_\etale$, d'après le
lemme \ref{proetale-lemma-compare-constructible}.
Posons $\mathcal{G} = \bigoplus \mathcal{G}_n$. C'est un faisceau
de $R$-modules sur $X_\etale$, et le morphisme
$$
\mathcal{G}_0 \otimes_{\Lambda/I} \underline{R}
\longrightarrow
\mathcal{G}
$$
est surjectif parce que les morphismes
$$
\mathcal{F}/I\mathcal{F} \otimes \underline{I^n/I^{n + 1}} \to
I^n\mathcal{F}/I^{n + 1}\mathcal{F}
$$
sont surjectifs. Ainsi, $\mathcal{G}$ est un faisceau constructible de
$R$-modules d'après Cohomologie \'etale, proposition
\ref{etale-cohomology-proposition-constructible-over-noetherian}.
Choisissons une partition $X = \coprod X_i$ telle que
$\mathcal{G}|_{X_i}$ soit un faisceau localement constant de $R$-modules
de type fini (Cohomologie \'etale, lemme
\ref{etale-cohomology-lemma-constructible-quasi-compact-quasi-separated}).
Nous affirmons que cette partition possède la propriété de l'énoncé.
En effet, en remplaçant $X$ par $X_i$, nous pouvons supposer que $\mathcal{G}$ est localement
constant. Il s'ensuit que chacun des faisceaux
 $I^n\mathcal{F}/I^{n + 1}\mathcal{F}$
est localement constant. À l'aide des suites exactes courtes
$$
0 \to I^n\mathcal{F}/I^{n + 1}\mathcal{F} \to
\mathcal{F}/I^{n + 1}\mathcal{F} \to \mathcal{F}/I^n\mathcal{F} \to 0
$$
et, par récurrence, du lemme de Modules sur les sites
\ref{sites-modules-lemma-kernel-finite-locally-constant}
on obtient le résultat.
\end{proof}

\begin{lemma}
\label{proetale-lemma-weird}
Soit $X$ un schéma affine faiblement contractile. Soit $\Lambda$ un anneau noethérien,
soit $I \subset \Lambda$ un idéal et soit $\mathcal{F}$ un faisceau de
$\Lambda$-modules sur $X_\proetale$ tel que
\begin{enumerate}
\item $\mathcal{F} = \lim \mathcal{F}/I^n\mathcal{F}$,
\item $\mathcal{F}/I^n\mathcal{F}$ soit un faisceau constant de
$\Lambda/I^n$-modules,
\item $\mathcal{F}/I\mathcal{F}$ soit de type fini.
\end{enumerate}
Alors $\mathcal{F} \cong \underline{M}^\wedge$, où $M$ est
un module de type fini sur $\Lambda^\wedge$.
\end{lemma}

\begin{proof}
Choisissons un $\Lambda/I^n$-module $M_n$ tel que
$\mathcal{F}/I^n\mathcal{F} \cong \underline{M_n}$.
Puisque nous disposons des surjections
$\mathcal{F}/I^{n + 1}\mathcal{F} \to \mathcal{F}/I^n\mathcal{F}$,
nous concluons qu'il existe des
surjections $M_{n + 1} \to M_n$ induisant des isomorphismes
$M_{n + 1}/I^nM_{n + 1} \to M_n$. Fixons un choix de telles surjections
et posons $M = \lim M_n$. Alors $M$, complet pour la topologie $I$-adique,
est un $\Lambda$-module vérifiant $M/I^nM = M_n$; voir le
lemme Algèbre, \ref{algebra-lemma-limit-complete}.
Puisque $M_1$ est un $\Lambda$-module de type fini
(Modules sur les sites, lemme
\ref{sites-modules-lemma-locally-constant-finite-type}),
nous voyons que $M$ est un module de type fini sur $\Lambda^\wedge$.
Considérons les faisceaux
$$
\mathcal{I}_n = \mathit{Isom}(\underline{M_n}, \mathcal{F}/I^n\mathcal{F})
$$
sur $X_\proetale$. Le passage au quotient par $I^n$ définit un morphisme de transition
$$
\mathcal{I}_{n + 1} \longrightarrow \mathcal{I}_n
$$
D'après notre choix de $M_n$, le faisceau $\mathcal{I}_n$ est un torseur sous
$$
\mathit{Isom}(\underline{M_n}, \underline{M_n}) =
\underline{\text{Isom}_\Lambda(M_n, M_n)}
$$
(Modules sur les sites, lemme \ref{sites-modules-lemma-locally-constant}),
puisque $\mathcal{F}/I^n\mathcal{F}$ est localement isomorphe pour la topologie \'etale
à $\underline{M_n}$. D'après
Compléments d'algèbre, lemme \ref{more-algebra-lemma-hom-systems-ML},
que le système de faisceaux $(\mathcal{I}_n)$ vérifie la condition de Mittag-Leffler.
Pour tout $n$, soit $\mathcal{I}'_n \subset \mathcal{I}_n$ l'image
de $\mathcal{I}_N \to \mathcal{I}_n$ pour tout $N \gg n$.
Alors
$$
\ldots \to \mathcal{I}'_3 \to \mathcal{I}'_2 \to \mathcal{I}'_1 \to *
$$
est une suite de faisceaux d'ensembles sur $X_\proetale$ dont les morphismes de
transition sont surjectifs. Puisque $*(X)$ est réduit à un élément (et donc non vide)
et que l'évaluation en $X$ transforme les morphismes surjectifs de faisceaux d'ensembles
en surjections d'ensembles, nous pouvons choisir
$s \in \lim \mathcal{I}'_n(X)$. Ces sections définissent des isomorphismes
$\underline{M}^\wedge \to \lim \mathcal{F}/I^n\mathcal{F} = \mathcal{F}$,
ce qui achève la démonstration.
\end{proof}

\begin{lemma}
\label{proetale-lemma-connected-lisse}
Soit $X$ un schéma connexe. Soit $\Lambda$ un anneau noethérien et soit
$I \subset \Lambda$ un idéal. Si $\mathcal{F}$ est un faisceau constructible lisse
de $\Lambda$-modules sur $X_\proetale$, alors $\mathcal{F}$
est lisse adique.
\end{lemma}

\begin{proof}
D'après le lemme \ref{proetale-lemma-compare-locally-constant}, on a
$\mathcal{F}/I^n\mathcal{F} = \epsilon^{-1}\mathcal{G}_n$
pour un certain faisceau localement constant $\mathcal{G}_n$ de $\Lambda/I^n$-modules. D'après
le lemme de Cohomologie \'etale
\ref{etale-cohomology-lemma-connected-locally-constant},
il existe un module de type fini sur $\Lambda/I^n$, noté $M_n$, tel que
$\mathcal{G}_n$ soit localement isomorphe à $\underline{M_n}$.
Choisissons un recouvrement $\{W_t \to X\}_{t \in T}$ où chaque $W_t$ est
affine et faiblement contractile.
Alors $\mathcal{F}|_{W_t}$ satisfait aux hypothèses du
lemme \ref{proetale-lemma-weird},
et donc $\mathcal{F}|_{W_t} \cong \underline{N_t}^\wedge$
pour un module de type fini sur $\Lambda^\wedge$, noté $N_t$. Remarquons que
$N_t/I^nN_t \cong M_n$ pour tous $t$ et $n$. Ainsi,
$N_t \cong N_{t'}$ pour tous $t, t' \in T$; voir le
lemme de Compléments d'algèbre \ref{more-algebra-lemma-isomorphic-completions}.
Cela démontre que $\mathcal{F}$ est lisse adique.
\end{proof}

\begin{lemma}
\label{proetale-lemma-Noetherian-adic-constructible}
Soit $X$ un schéma noethérien. Soit $\Lambda$ un anneau noethérien et
soit $I \subset \Lambda$ un idéal. Soit $\mathcal{F}$ un
faisceau constructible de $\Lambda$-modules sur $X_\proetale$. Alors $\mathcal{F}$
est adique constructible.
\end{lemma}

\begin{proof}
C'est une conséquence des lemmes \ref{proetale-lemma-Noetherian-constructible} et
\ref{proetale-lemma-connected-lisse}, du fait qu'un schéma noethérien
est localement connexe
(Topologie, lemme \ref{topology-lemma-locally-Noetherian-locally-connected}),
et des définitions.
\end{proof}

\noindent
Il sera utile d'identifier les faisceaux constructibles de $\Lambda$-modules
dans la catégorie des faisceaux de $\Lambda$-modules complets au sens dérivé.
Il se trouve que l'analogue naïf de
Compléments d'algèbre, lemme \ref{more-algebra-lemma-derived-complete-finite},
est faux dans ce contexte. Cependant, voici
l'analogue de Compléments d'algèbre, lemme
\ref{more-algebra-lemma-derived-complete-zero}.

\begin{lemma}
\label{proetale-lemma-derived-complete-zero}
Soit $X$ un schéma. Soit $\Lambda$ un anneau et soit
$I \subset \Lambda$ un idéal de type fini.
Soit $\mathcal{F}$ un faisceau de $\Lambda$-modules sur $X_\proetale$.
Si $\mathcal{F}$ est complet au sens dérivé et si $\mathcal{F}/I\mathcal{F} = 0$,
alors $\mathcal{F} = 0$.
\end{lemma}

\begin{proof}
Supposons que $\mathcal{F}/I\mathcal{F}$ soit nul.
Écrivons $I = (f_1, \ldots, f_r)$. Soit $i < r$ le plus grand
entier tel que $\mathcal{G} = \mathcal{F}/(f_1, \ldots, f_i)\mathcal{F}$
soit non nul. Si un tel $i$ n'existe pas, alors $\mathcal{F} = 0$, ce que nous
voulions montrer. Alors $\mathcal{G}$ est complet au sens dérivé, comme conoyau
d'un morphisme entre modules complets au sens dérivé; voir la
proposition \ref{proetale-proposition-enough-weakly-contractibles}.
D'après le choix de $i$, le morphisme $f_{i + 1} : \mathcal{G} \to \mathcal{G}$
est surjectif. Ainsi,
$$
\lim (\ldots \to \mathcal{G} \xrightarrow{f_{i + 1}} \mathcal{G}
\xrightarrow{f_{i + 1}} \mathcal{G})
$$
est non nul, ce qui contredit le fait que $\mathcal{G}$ soit complet au sens dérivé.
\end{proof}

\begin{lemma}
\label{proetale-lemma-derived-complete-limit}
Soit $X$ un schéma affine faiblement contractile.
Soit $\Lambda$ un anneau noethérien et soit $I \subset \Lambda$ un idéal.
Soit $\mathcal{F}$ un faisceau de $\Lambda$-modules complet au sens dérivé
sur $X_\proetale$, tel que $\mathcal{F}/I\mathcal{F}$ soit un faisceau localement
constant de $\Lambda/I$-modules de type fini.
Alors il existe un entier $t$ et un morphisme surjectif
$$
(\underline{\Lambda}^\wedge)^{\oplus t} \to \mathcal{F}
$$
\end{lemma}

\begin{proof}
Puisque $X$ est faiblement contractile, il existe un recouvrement fini par des ouverts disjoints
$X = \coprod U_i$ tel que $\mathcal{F}/I\mathcal{F}|_{U_i}$
soit isomorphe au faisceau constant associé à un module de type fini sur $\Lambda/I$
$M_i$. Choisissons un nombre fini de générateurs $m_{ij}$ de $M_i$. Nous
pouvons trouver des sections $s_{ij} \in \mathcal{F}(X)$ dont les restrictions sont les
$m_{ij}$ considérés comme sections de $\mathcal{F}/I\mathcal{F}$ sur $U_i$. 
Soit $t$ le nombre total des $s_{ij}$. Nous obtenons alors un morphisme
$$
\alpha : \underline{\Lambda}^{\oplus t} \longrightarrow \mathcal{F}
$$
qui est surjectif modulo $I$ par construction. D'après le
lemme \ref{proetale-lemma-naive-completion},
la complétion dérivée de $\underline{\Lambda}^{\oplus t}$ est le
faisceau $(\underline{\Lambda}^\wedge)^{\oplus t}$. Puisque $\mathcal{F}$
est complet au sens dérivé, on voit que $\alpha$ se factorise par un morphisme
$$
\alpha^\wedge :
(\underline{\Lambda}^\wedge)^{\oplus t}
\longrightarrow
\mathcal{F}
$$
Alors $\mathcal{Q} = \Coker(\alpha^\wedge)$ est un faisceau de
$\Lambda$-modules complet au sens dérivé d'après la
proposition \ref{proetale-proposition-enough-weakly-contractibles}.
Par construction, $\mathcal{Q}/I\mathcal{Q} = 0$. Il résulte du
lemme \ref{proetale-lemma-derived-complete-zero}
que $\mathcal{Q} = 0$, ce que nous voulions montrer.
\end{proof}








\section{Une catégorie dérivée appropriée}
\label{proetale-section-suitable-derived}

\noindent
Soit $X$ un schéma. Il s'avérera que, pour de nombreux schémas $X$, on peut obtenir une
catégorie dérivée appropriée de faisceaux $\ell$-adiques en considérant les
objets $K$ de $D(X_\proetale, \Lambda)$ complets au sens dérivé tels
que $K \otimes_\Lambda^\mathbf{L} \mathbf{F}_\ell$ soit borné et possède des
faisceaux de cohomologie constructibles. Voici la définition générale.

\begin{definition}
\label{proetale-definition-Dbc}
Soit $\Lambda$ un anneau noethérien et soit $I \subset \Lambda$ un idéal.
Soit $X$ un schéma. Un objet $K$ de $D(X_\proetale, \Lambda)$ est dit
{\it constructible} si
\begin{enumerate}
\item $K$ est complet au sens dérivé relativement à $I$,
\item $K \otimes_\Lambda^\mathbf{L} \underline{\Lambda/I}$
possède des faisceaux de cohomologie constructibles et a localement une dimension de Tor finie.
\end{enumerate}
Nous notons $D_{cons}(X, \Lambda)$ la sous-catégorie pleine formée des objets constructibles
$K$ de $D(X_\proetale, \Lambda)$.
\end{definition}

\noindent
Rappelons qu'avec nos conventions, un complexe de dimension de Tor finie
est borné (Cohomologie sur les sites, définition
\ref{sites-cohomology-definition-tor-amplitude}).
En fait, rassemblons dans un lemme tout ce que nous avons démontré jusqu'ici.

\begin{lemma}
\label{proetale-lemma-describe-constructible-complexes}
Dans la situation ci-dessus, supposons que $K$ appartienne à $D_{cons}(X, \Lambda)$
et que $X$ soit quasi-compact. Posons
$K_n = K \otimes_\Lambda^\mathbf{L} \underline{\Lambda/I^n}$.
Il existe $a, b$ tels que
\begin{enumerate}
\item $K = R\lim K_n$ et $H^i(K) = 0$ pour $i \not \in [a, b]$,
\item chaque $K_n$ est d'amplitude de Tor dans $[a, b]$,
\item chaque $K_n$ possède des faisceaux de cohomologie constructibles,
\item pour tout indice, $K_n = \epsilon^{-1}L_n$ pour un certain
$L_n \in D_{ctf}(X_\etale, \Lambda/I^n)$
(Cohomologie \'etale, définition \ref{etale-cohomology-definition-ctf}).
\end{enumerate}
\end{lemma}

\begin{proof}
Puisque $K$ est complet au sens dérivé, nous avons $K = R\lim K_n$ d'après le
lemme \ref{proetale-lemma-naive-completion}.
Par définition de la finitude locale de la dimension de Tor, nous pouvons trouver
$a, b$ tels que $K_1$ soit d'amplitude de Tor dans $[a, b]$.
La partie (2) résulte de
Cohomologie sur les sites, lemme \ref{sites-cohomology-lemma-bounded}.
Ensuite, (1) résulte du fait que $K$ est complet au sens dérivé, par la description
des limites dans
Cohomologie sur les sites, proposition
\ref{sites-cohomology-proposition-enough-weakly-contractibles},
et du fait que $H^b(K_{n + 1}) \to H^b(K_n)$ est surjective,
puisque $K_n = K_{n + 1} \otimes^\mathbf{L}_\Lambda \underline{\Lambda/I^n}$.
La partie (3) résulte du
lemme \ref{proetale-lemma-compare-truncations-constructible}.
La partie (4) résulte du
lemme \ref{proetale-lemma-compare-constructible-derived}
et du fait que $L_n$ est de dimension de Tor finie parce que $K_n$ l'est
(nous omettons un court argument).
\end{proof}

\begin{lemma}
\label{proetale-lemma-local-structure-constructible-complex}
Soit $X$ un schéma affine faiblement contractile. Soit $\Lambda$ un
anneau noethérien et soit $I \subset \Lambda$ un idéal. Soit $K$ un objet de
$D_{cons}(X, \Lambda)$ tel que les faisceaux de cohomologie de
$K \otimes_\Lambda^\mathbf{L} \underline{\Lambda/I}$ soient localement
constants. Alors il existe un recouvrement fini par des ouverts disjoints
$X = \coprod U_i$ et, pour chaque $i$, une famille finie de
modules projectifs de type fini sur $\Lambda^\wedge$, notés $M^a, \ldots, M^b$,
telle que $K|_{U_i}$ soit représenté par un complexe
$$
(\underline{M^a})^\wedge \to \ldots \to (\underline{M^b})^\wedge
$$
dans $D(U_{i, \proetale}, \Lambda)$, pour certains morphismes de faisceaux de
$\Lambda$-modules $(\underline{M^i})^\wedge \to (\underline{M^{i + 1}})^\wedge$.
\end{lemma}

\begin{proof}
Nous utilisons librement les résultats du
lemme \ref{proetale-lemma-describe-constructible-complexes}.
Choisissons $a, b$ comme dans ce lemme. Nous allons démontrer le lemme par
récurrence sur $b - a$. Posons $\mathcal{F} = H^b(K)$.
Remarquons que $\mathcal{F}$ est un faisceau de
$\Lambda$-modules complet au sens dérivé, d'après la
proposition \ref{proetale-proposition-enough-weakly-contractibles}.
De plus, $\mathcal{F}/I\mathcal{F}$ est un faisceau localement
constant de $\Lambda/I$-modules de type fini.
Appliquons le lemme \ref{proetale-lemma-derived-complete-limit} pour obtenir une surjection
$\rho : (\underline{\Lambda}^\wedge)^{\oplus t} \to \mathcal{F}$.

\medskip\noindent
Si $a = b$, alors $K = \mathcal{F}[-b]$. Dans ce cas, nous voyons que
$$
\mathcal{F} \otimes_\Lambda^\mathbf{L} \underline{\Lambda/I} =
\mathcal{F}/I\mathcal{F}
$$
Comme $X$ est faiblement contractile et que $\mathcal{F}/I\mathcal{F}$ est
localement constant, nous pouvons trouver une décomposition finie
$X = \coprod U_i$ en ouverts affines disjoints $U_i$
et des $\Lambda/I$-modules $\overline{M}_i$
tels que $\mathcal{F}/I\mathcal{F}$ se restreigne à
$\underline{\overline{M}_i}$ sur $U_i$. Après avoir raffiné le recouvrement,
nous pouvons supposer que le morphisme
$$
\rho|_{U_i} \bmod I :
\underline{\Lambda/I}^{\oplus t}
\longrightarrow
\underline{\overline{M}_i}
$$
est égal à $\underline{\alpha_i}$ pour un certain morphisme surjectif de modules
$\alpha_i : \Lambda/I^{\oplus t} \to \overline{M}_i$ ; voir
Modules sur les sites, lemme \ref{sites-modules-lemma-morphism-locally-constant}.
Remarquons que chaque $\overline{M}_i$ est un module de type fini sur $\Lambda/I$.
Puisque $\mathcal{F}/I\mathcal{F}$ est d'amplitude de Tor dans $[0, 0]$,
nous en concluons que $\overline{M}_i$ est un $\Lambda/I$-module plat.
Ainsi, $\overline{M}_i$ est projectif de type fini
(Algèbre, lemme \ref{algebra-lemma-finite-projective}).
Nous pouvons donc trouver un projecteur
$\overline{p}_i : (\Lambda/I)^{\oplus t} \to (\Lambda/I)^{\oplus t}$
dont l'image est envoyée isomorphiquement sur $\overline{M}_i$ par le morphisme $\alpha_i$.
Nous pouvons relever $\overline{p}_i$ en un projecteur
$p_i : (\Lambda^\wedge)^{\oplus t} \to
(\Lambda^\wedge)^{\oplus t}$\footnote{Démonstration : d'après
Algèbre, lemme \ref{algebra-lemma-lift-idempotents-noncommutative},
nous pouvons relever $\overline{p}_i$ en un système compatible de
projecteurs $p_{i, n} : (\Lambda/I^n)^{\oplus t} \to (\Lambda/I^n)^{\oplus t}$,
puis poser $p_i = \lim p_{i, n}$, ce qui convient parce que
$\Lambda^\wedge = \lim \Lambda/I^n$.}.
Alors $M_i = \Im(p_i)$ est complet pour la topologie $I$-adique et de type fini
comme $\Lambda^\wedge$-module ; de plus, $M_i/IM_i = \overline{M}_i$.
Sur $U_i$, considérons les morphismes
$$
\underline{M_i}^\wedge \to
(\underline{\Lambda}^\wedge)^{\oplus t} \to
\mathcal{F}|_{U_i}
$$
Par construction, leur composé induit un isomorphisme modulo $I$.
La source et le but sont complets au sens dérivé ; il en va donc de même du conoyau
$\mathcal{Q}$ et du noyau $\mathcal{K}$. Nous avons
$\mathcal{Q}/I\mathcal{Q} = 0$ par construction ; ainsi, $\mathcal{Q}$
est nul d'après le lemme \ref{proetale-lemma-derived-complete-zero}.
Alors
$$
0 \to \mathcal{K}/I\mathcal{K} \to
\underline{\overline{M}_i}
\to \mathcal{F}/I\mathcal{F} \to 0
$$
est exacte grâce à l'annulation de $\text{Tor}_1$ rappelée au début de ce
paragraphe. On utilise aussi l'expression
$\underline{\Lambda}^\wedge/I\overline{\Lambda}^\wedge$,
ainsi que Modules sur les sites, lemme \ref{sites-modules-lemma-completion-flat},
pour voir que
$\underline{M_i}^\wedge/I\underline{M_i}^\wedge = \underline{\overline{M}_i}$.
Ainsi, $\mathcal{K}/I\mathcal{K} = 0$ par construction, et nous concluons
que $\mathcal{K} = 0$ comme précédemment. Cela démontre le résultat lorsque $a = b$.

\medskip\noindent
Si $b > a$, nous relevons le morphisme $\rho$ en un morphisme
$$
\tilde \rho : (\underline{\Lambda}^\wedge)^{\oplus t}[-b] \longrightarrow K
$$
dans $D(X_\proetale, \Lambda)$. Cela est possible car nous pouvons considérer
$K$ comme un complexe de $\underline{\Lambda}^\wedge$-modules, d'après la
discussion dans l'introduction de la
section \ref{proetale-section-derived-completion-noetherian},
et parce que $X_\proetale$ est faiblement contractile ;
il n'y a donc aucune obstruction à relever les éléments
$\rho(e_s) \in H^0(X, \mathcal{F})$ en des éléments de $H^b(X, K)$.
En insérant $\tilde \rho$ dans un triangle distingué
$$
(\underline{\Lambda}^\wedge)^{\oplus t}[-b] \to K \to L \to
(\underline{\Lambda}^\wedge)^{\oplus t}[-b + 1]
$$
nous voyons que $L$ est un objet de $D_{cons}(X, \Lambda)$ tel
que $L \otimes_\Lambda^\mathbf{L} \underline{\Lambda/I}$
soit d'amplitude de Tor contenue dans $[a, b - 1]$ (nous omettons les détails).
Par récurrence, nous pouvons décrire $L$ localement comme dans l'énoncé du lemme ; disons que
$L$ est isomorphe à
$$
(\underline{M^a})^\wedge \to \ldots \to (\underline{M^{b - 1}})^\wedge
$$
Le morphisme
$L \to (\underline{\Lambda}^\wedge)^{\oplus t}[-b + 1]$
correspond à un morphisme
$(\underline{M^{b - 1}})^\wedge \to (\underline{\Lambda}^\wedge)^{\oplus t}$,
qui nous permet de prolonger le complexe d'un terme. Le complexe correspondant
est isomorphe à $K$ dans la catégorie dérivée, par les propriétés
des catégories triangulées. Cela achève la démonstration.
\end{proof}

\noindent
Motivés par ce qui se passe pour les faisceaux constructibles de $\Lambda$-modules,
nous introduisons la notion suivante.

\begin{definition}
\label{proetale-definition-adic-constructible}
Soit $X$ un schéma. Soit $\Lambda$ un anneau noethérien et soit
$I \subset \Lambda$ un idéal. Soit $K \in D(X_\proetale, \Lambda)$.
\begin{enumerate}
\item Nous disons que $K$ est {\it lisse adique}\footnote{Cette notation n'est peut-être
pas standard.} s'il existe un complexe borné de modules
projectifs de type fini sur $\Lambda^\wedge$, noté $M^\bullet$, tel que
$K$ soit localement isomorphe à
$$
\underline{M^a}^\wedge \to \ldots \to \underline{M^b}^\wedge
$$
\item Nous disons que $K$ est {\it adique constructible}\footnote{Cette notation n'est peut-être
pas standard.} si, pour tout ouvert affine $U \subset X$,
il existe une décomposition $U = \coprod U_i$ en
sous-schémas localement fermés constructibles telle que $K|_{U_i}$
soit lisse adique.
\end{enumerate}
\end{definition}

\noindent
La différence entre la structure locale obtenue dans le
lemme \ref{proetale-lemma-local-structure-constructible-complex}
et la structure d'un complexe lisse adique tient à ce que
les morphismes $\underline{M^i}^\wedge \to \underline{M^{i + 1}}^\wedge$ du
lemme \ref{proetale-lemma-local-structure-constructible-complex}
ne sont pas nécessairement constants, tandis que, dans la définition ci-dessus, ils
doivent être constants.

\begin{lemma}
\label{proetale-lemma-weakly-contractible-locally-constant-ML}
Soit $X$ un schéma affine faiblement contractile. Soit $\Lambda$ un
anneau noethérien et soit $I \subset \Lambda$ un idéal. Soit $K$ un objet de
$D_{cons}(X, \Lambda)$ tel que
$K \otimes_\Lambda^\mathbf{L} \underline{\Lambda/I^n}$
soit isomorphe dans $D(X_\proetale, \Lambda/I^n)$ à un
complexe de faisceaux constants de $\Lambda/I^n$-modules. Alors
$$
H^0(X, K \otimes_\Lambda^\mathbf{L} \Lambda/I^n)
$$
satisfait à la condition de Mittag-Leffler.
\end{lemma}

\begin{proof}
Supposons que $K \otimes_\Lambda^\mathbf{L} \underline{\Lambda/I^n}$ soit isomorphe
à $\underline{E_n}$ pour un objet $E_n$ de $D(\Lambda/I^n)$.
Puisque $K \otimes_\Lambda^\mathbf{L} \underline{\Lambda/I}$ a une
dimension de Tor finie et possède des faisceaux de cohomologie de type fini,
on voit que $E_1$ est parfait (voir
Compléments d'algèbre, lemme \ref{more-algebra-lemma-perfect}). Les morphismes de transition
$$
K \otimes_\Lambda^\mathbf{L} \underline{\Lambda/I^{n + 1}}
\to
K \otimes_\Lambda^\mathbf{L} \underline{\Lambda/I^n}
$$
proviennent localement de morphismes de complexes (éventuellement nombreux et distincts)
$E_{n + 1} \to E_n$ dans $D(\Lambda/I^{n + 1})$ ; voir
Cohomologie sur les sites, lemme \ref{sites-cohomology-lemma-locally-constant-map}.
Pour chaque $n$, choisissons un tel morphisme et observons qu'il induit
un isomorphisme
$E_{n + 1} \otimes_{\Lambda/I^{n + 1}}^\mathbf{L} \Lambda/I^n \to E_n$
dans $D(\Lambda/I^n)$. D'après
Compléments d'algèbre, lemme \ref{more-algebra-lemma-Rlim-perfect-gives-complete}
il existe un complexe borné $M^\bullet$ de modules projectifs de type fini sur
$\Lambda^\wedge$ et des isomorphismes $M^\bullet/I^nM^\bullet \to E_n$
dans $D(\Lambda/I^n)$ compatibles avec les morphismes de transition.

\medskip\noindent
Observons maintenant que, pour toute famille finie d'indices
$n > m > k$, le triplet de morphismes
$$
H^0(X, K \otimes_\Lambda^\mathbf{L} \Lambda/I^n)
\to
H^0(X, K \otimes_\Lambda^\mathbf{L} \Lambda/I^m)
\to
H^0(X, K \otimes_\Lambda^\mathbf{L} \Lambda/I^k)
$$
est isomorphe à
$$
H^0(X, \underline{M^\bullet/I^nM^\bullet})
\to
H^0(X, \underline{M^\bullet/I^mM^\bullet})
\to
H^0(X, \underline{M^\bullet/I^kM^\bullet})
$$
En effet, tout isomorphisme
$$
\underline{M^\bullet/I^nM^\bullet} \to
K \otimes_\Lambda^\mathbf{L} \Lambda/I^n
$$
induit des isomorphismes analogues modulo $I^m$ et $I^k$; l'assertion est donc
vraie. Ainsi, pour démontrer le lemme, il suffit de montrer que
le système
$H^0(X, \underline{M^\bullet/I^nM^\bullet})$ vérifie la condition de Mittag-Leffler.
Comme la prise des sections sur $X$ est exacte, il suffit de montrer que
le système de $\Lambda$-modules
$$
H^0(M^\bullet/I^nM^\bullet)
$$
vérifie la condition de Mittag-Leffler. Posons $A = \Lambda^\wedge$ et considérons la suite
spectrale
$$
\text{Tor}_{-p}^A(H^q(M^\bullet), A/I^nA) \Rightarrow
H^{p + q}(M^\bullet/I^nM^\bullet)
$$
D'après Compléments d'algèbre, lemme \ref{more-algebra-lemma-tor-strictly-pro-zero}
les pro-systèmes $\{\text{Tor}_{-p}^A(H^q(M^\bullet), A/I^nA)\}$ sont nuls
pour $p > 0$. Le pro-système $\{H^0(M^\bullet/I^nM^\bullet)\}$
est donc égal au pro-système $\{H^0(M^\bullet)/I^nH^0(M^\bullet)\}$
et le lemme est démontré.
\end{proof}

\begin{lemma}
\label{proetale-lemma-connected-adic-lisse}
Soit $X$ un schéma connexe. Soit $\Lambda$ un anneau noethérien et soit
$I \subset \Lambda$ un idéal. Si $K$ appartient à $D_{cons}(X, \Lambda)$
et si $K \otimes_\Lambda \underline{\Lambda/I}$
a des faisceaux de cohomologie localement constants, alors $K$ est lisse adique
(définition \ref{proetale-definition-adic-constructible}).
\end{lemma}

\begin{proof}
Posons $K_n = K \otimes_\Lambda^\mathbf{L} \underline{\Lambda/I^n}$.
Nous utiliserons les résultats du lemme \ref{proetale-lemma-describe-constructible-complexes}
sans autre mention. D'après Cohomologie sur les sites, lemme
\ref{sites-cohomology-lemma-locally-constant-bounded}
on voit que $K_n$ a des faisceaux de cohomologie localement constants pour tout $n$.
On a $K_n = \epsilon^{-1}L_n$ pour un certain $L_n$ appartenant à
$D_{ctf}(X_\etale, \Lambda/I^n)$ à faisceaux de cohomologie localement constants.
D'après Cohomologie étale, lemme
\ref{etale-cohomology-lemma-connected-ctf-locally-constant}
il existe des complexes parfaits $M_n \in D(\Lambda/I^n)$
tels que $L_n$ soit localement isomorphe à $\underline{M_n}$ pour la topologie \'etale.
Les morphismes $L_{n + 1} \to L_n$ correspondant à $K_{n + 1} \to K_n$
induisent des isomorphismes
$L_{n + 1} \otimes_{\Lambda/I^{n + 1}}^\mathbf{L} \underline{\Lambda/I^n}
\to L_n$. En raisonnant localement sur $X$, on conclut qu'il
existe des morphismes $M_{n + 1} \to M_n$ dans $D(\Lambda/I^{n + 1})$
induisant des isomorphismes
$M_{n + 1} \otimes_{\Lambda/I^{n + 1}} \Lambda/I^n \to M_n$ ; voir
Cohomologie sur les sites, lemme \ref{sites-cohomology-lemma-locally-constant-map}.
Fixons un choix de tels morphismes. D'après
Compléments d'algèbre, lemme \ref{more-algebra-lemma-Rlim-perfect-gives-complete}
il existe un complexe borné $M^\bullet$ de modules projectifs de type fini sur
$\Lambda^\wedge$ et des isomorphismes $M^\bullet/I^nM^\bullet \to M_n$
dans $D(\Lambda/I^n)$ compatibles avec les morphismes de transition.
Pour achever la démonstration, nous allons montrer que $K$ est localement isomorphe
à
$$
\underline{M^\bullet}^\wedge =
\lim \underline{M^\bullet/I^nM^\bullet} =
R\lim \underline{M^\bullet/I^nM^\bullet}
$$
Soit $E^\bullet$ le complexe dual de $M^\bullet$ ; voir
Compléments d'algèbre, lemme \ref{more-algebra-lemma-dual-perfect-complex}
ainsi que sa démonstration. Considérons les objets
$$
H_n =
R\SheafHom_{\Lambda/I^n}(\underline{M^\bullet/I^nM^\bullet}, K_n) =
\underline{E^\bullet/I^nE^\bullet} \otimes_{\Lambda/I^n}^\mathbf{L} K_n
$$
de $D(X_\proetale, \Lambda/I^n)$. La réduction modulo $I^n$ définit un
morphisme de transition $H_{n + 1} \to H_n$. Posons $H = R\lim H_n$. Alors $H$ est un
objet de $D_{cons}(X, \Lambda)$ (nous omettons les détails) tel que
$H \otimes_\Lambda^\mathbf{L} \underline{\Lambda/I^n} = H_n$.
Choisissons un recouvrement $\{W_t \to X\}_{t \in T}$ où chaque $W_t$
est affine et faiblement contractile. Le choix de $M^\bullet$
montre que
\begin{align*}
H_n|_{W_t} & \cong 
R\SheafHom_{\Lambda/I^n}(\underline{M^\bullet/I^nM^\bullet},
\underline{M^\bullet/I^nM^\bullet}) \\
& =
\underline{
\text{Tot}(E^\bullet/I^nE^\bullet \otimes_{\Lambda/I^n} M^\bullet/I^nM^\bullet)
}
\end{align*}
Nous pouvons donc appliquer le lemme \ref{proetale-lemma-weakly-contractible-locally-constant-ML}
à $H = R\lim H_n$. On conclut que le système $H^0(W_t, H_n)$ vérifie la condition de
Mittag-Leffler. Puisque, pour tout $n \gg 1$, il existe un élément de $H^0(W_t, H_n)$
dont l'image est un isomorphisme dans
$$
H^0(W_t, H_1) = \Hom(\underline{M^\bullet/IM^\bullet}, K_1)
$$
on obtient dans la limite inverse un élément $(\varphi_{t, n})$
qui fournit un isomorphisme modulo $I$. Alors
$$
R\lim \varphi_{t, n} :
\underline{M^\bullet}^\wedge|_{W_t} =
R\lim \underline{M^\bullet/I^nM^\bullet}|_{W_t}
\longrightarrow
R\lim K_n|_{W_t} = K|_{W_t}
$$
est un isomorphisme. Cela achève la démonstration.
\end{proof}

\begin{proposition}
\label{proetale-proposition-Noetherian-adic-constructible}
Soit $X$ un schéma noethérien. Soit $\Lambda$ un anneau noethérien et
soit $I \subset \Lambda$ un idéal. Soit $K$ un objet de
$D_{cons}(X, \Lambda)$. Alors $K$ est adique constructible
(définition \ref{proetale-definition-adic-constructible}).
\end{proposition}

\begin{proof}
C'est une conséquence du lemme \ref{proetale-lemma-connected-adic-lisse}
ainsi que du fait qu'un schéma noethérien est localement connexe
(Topologie, lemme \ref{topology-lemma-locally-Noetherian-locally-connected}),
et des définitions.
\end{proof}







\section{Changement de base propre}
\label{proetale-section-proper-base-change}

\noindent
Dans cette section, nous expliquons comment démontrer le théorème de changement de base propre
pour les objets complets au sens dérivé sur le site pro-\'etale, à l'aide du théorème de
changement de base propre en cohomologie \'etale, conformément au principe général
selon lequel nous n'utilisons la topologie pro-\'etale que pour traiter les « problèmes de limite »
et faisons appel aux résultats établis pour la topologie \'etale afin de traiter tout
le reste.

\begin{theorem}
\label{proetale-theorem-proper-base-change}
Supposons que $f : X \to Y$ soit un morphisme propre de schémas et que $g : Y' \to Y$ soit
un morphisme de schémas donnant lieu au diagramme de changement de base
$$
\xymatrix{
X' \ar[r]_{g'} \ar[d]_{f'} & X \ar[d]^f \\
Y' \ar[r]^g & Y
}
$$
Soit $\Lambda$ un anneau noethérien et soit $I \subset \Lambda$ un idéal
tel que $\Lambda/I$ soit de torsion. Soit $K$ un objet
de $D(X_\proetale)$ tel que
\begin{enumerate}
\item $K$ est complet au sens dérivé ;
\item $K \otimes_\Lambda^\mathbf{L} \underline{\Lambda/I^n}$ est
borné inférieurement, à faisceaux de cohomologie provenant de $X_\etale$ ;
\item $\Lambda/I^n$ est un $\Lambda$-module parfait\footnote{Cette hypothèse
peut être supprimée si $K$ est un complexe constructible, voir \cite{BS}.}.
\end{enumerate}
Alors le morphisme de changement de base
$$
Lg_{comp}^*Rf_*K \longrightarrow Rf'_*L(g')^*_{comp}K
$$
est un isomorphisme.
\end{theorem}

\begin{proof}
Nous omettons la construction du morphisme de changement de base (elle n'utilise que
les propriétés formelles de l'image directe dérivée et de l'image inverse dérivée complétée ;
comparer avec
Cohomologie sur les sites, remarque \ref{sites-cohomology-remark-base-change}).
Posons $K_n = K \otimes^\mathbf{L}_\Lambda \underline{\Lambda/I^n}$.
D'après le lemme \ref{proetale-lemma-naive-completion}, on a $K = R\lim K_n$
parce que $K$ est complet au sens dérivé.
Les lemmes \ref{proetale-lemma-pushforward-Noetherian-case} et
\ref{proetale-lemma-naive-completion} permettent de réécrire le membre de gauche :
$$
Lg_{comp}^* Rf_* K =
R\lim Lg^*(Rf_*K)\otimes^\mathbf{L}_\Lambda \underline{\Lambda/I^n} =
R\lim Lg^* Rf_* K_n
$$
La dernière égalité résulte du fait que $\Lambda/I^n$ est un module parfait et de
la formule de projection (Cohomologie sur les sites, lemme
\ref{sites-cohomology-lemma-projection-formula}).
Le lemme \ref{proetale-lemma-pushforward-Noetherian-case} permet de réécrire le membre de
droite :
$$
Rf'_* L(g')^*_{comp} K =
Rf'_* R\lim L(g')^* K_n  =
R\lim Rf'_* L(g')^* K_n
$$
La dernière égalité résulte du fait que $Rf'_*$ commute à $R\lim$
(Cohomologie sur les sites, lemme
\ref{sites-cohomology-lemma-Rf-commutes-with-Rlim}).
Il suffit donc de montrer que les morphismes
$$
Lg^* Rf_* K_n \longrightarrow Rf'_* L(g')^* K_n
$$
sont des isomorphismes. D'après le lemme \ref{proetale-lemma-compare-derived} et la deuxième
condition, on peut écrire $K_n = \epsilon^{-1}L_n$ pour un certain
$L_n \in D^+(X_\etale, \Lambda/I^n)$. D'après le lemme \ref{proetale-lemma-morphism-comparison}
et le fait que $\epsilon^{-1}$ commute aux images inverses,
on obtient
$$
Lg^* Rf_* K_n =
Lg^* Rf_* \epsilon^*L_n =
Lg^* \epsilon^{-1} Rf_* L_n =
\epsilon^{-1} Lg^* Rf_* L_n
$$
et
$$
Rf'_* L(g')^* K_n =
Rf'_* L(g')^* \epsilon^{-1} L_n =
Rf'_* \epsilon^{-1} L(g')^* L_n =
\epsilon^{-1} Rf'_* L(g')^* L_n
$$
(on utilise également ici que $L_n$ est borné inférieurement).
Enfin, d'après le théorème de changement de base propre en cohomologie \'etale
(Cohomologie étale, théorème
\ref{etale-cohomology-theorem-proper-base-change}), on a
$$
Lg^* Rf_* L_n = Rf'_* L(g')^* L_n
$$
(en utilisant de nouveau que $L_n$ est borné inférieurement),
ce qui démontre le théorème.
\end{proof}






\section{Changement d'univers partiel}
\label{proetale-section-change-universe}

\noindent
Nous conseillons au lecteur de passer cette section : nous montrons ici que la cohomologie
des faisceaux pour la topologie pro-\'etale est indépendante du
choix de l'univers partiel. Plus précisément, le foncteur $g_*$ du
lemme \ref{proetale-lemma-proetale-cohomology-independent-partial-universe} ci-dessous
est un plongement de petits topos pro-\'etales qui ne modifie pas la cohomologie.
Pour les grands sites pro-\'etales, les lemmes \ref{proetale-lemma-change-alpha} et
\ref{proetale-lemma-cohomology-enlarge-partial-universe} disent essentiellement la
même chose.

\medskip\noindent
Mais d'abord, 
comme promis dans la section \ref{proetale-section-proetale}, démontrons que la topologie d'un
grand site pro-\'etale $\Sch_\proetale$ est, en un certain sens, induite par
la topologie pro-\'etale sur la catégorie de tous les schémas.

\begin{lemma}
\label{proetale-lemma-proetale-induced}
Soit $\Sch_\proetale$ un grand site pro-\'etale comme dans la
définition \ref{proetale-definition-big-proetale-site}.
Soit $T \in \Ob(\Sch_\proetale)$.
Soit $\{T_i \to T\}_{i \in I}$ un recouvrement pro-\'etale quelconque de $T$.
Il existe un recouvrement $\{U_j \to T\}_{j \in J}$ de $T$ dans le site
$\Sch_\proetale$ qui raffine $\{T_i \to T\}_{i \in I}$.
\end{lemma}

\begin{proof}
Prenons d'abord un recouvrement $\{V_k \to T\}$ comme dans le
lemme \ref{proetale-lemma-get-many-weakly-contractible}.
Alors les recouvrements pro-\'etales $\{T_i \times_T V_k \to V_k\}$
peuvent chacun être raffinés par un recouvrement fini par des ouverts disjoints
$V_k = V_{k, 1} \amalg \ldots \amalg V_{k, n_k}$, voir le
lemme \ref{proetale-lemma-w-contractible-proetale-cover}.
Alors $\{V_{k, i} \to T\}$ constitue, dans $\Sch_\proetale$, un recouvrement de l'objet considéré
qui raffine $\{T_i \to T\}_{i \in I}$.
\end{proof}

\noindent
Énonçons et démontrons d'abord la comparaison pour les petits sites
pro-\'etales. Notons que nous n'affirmons pas que le petit topos pro-\'etale
d'un schéma soit indépendant du choix de l'univers partiel ; cela est faux,
contrairement au cas du petit topos \'etale
(Cohomologie étale, lemme
\ref{etale-cohomology-lemma-etale-topos-independent-partial-universe}).

\begin{lemma}
\label{proetale-lemma-proetale-cohomology-independent-partial-universe}
Soit $S$ un schéma. Soient $S_\proetale \subset S_\proetale'$ deux
petits sites pro-\'etales de $S$ construits dans la
définition \ref{proetale-definition-big-small-proetale}. Alors le foncteur d'inclusion
vérifie les hypothèses de
Sites, lemme \ref{sites-lemma-bigger-site}.
Il existe donc des morphismes de topos
$$
\xymatrix{
\Sh(S_\proetale) \ar[r]^g &
\Sh(S_\proetale') \ar[r]^f &
\Sh(S_\proetale)
}
$$
dont la composée est isomorphe à l'identité et qui vérifient $f_* = g^{-1}$.
De plus,
\begin{enumerate}
\item pour $\mathcal{F}' \in \textit{Ab}(S_\proetale')$, on a
$H^p(S_\proetale', \mathcal{F}') = H^p(S_\proetale, g^{-1}\mathcal{F}')$,
\item pour $\mathcal{F} \in \textit{Ab}(S_\proetale)$, on a
$$
H^p(S_\proetale, \mathcal{F}) =
H^p(S_\proetale', g_*\mathcal{F}) =
H^p(S_\proetale', f^{-1}\mathcal{F}).
$$
\end{enumerate}
\end{lemma}

\begin{proof}
Le foncteur d'inclusion est pleinement fidèle et continu.
Nous avons vu que $S_\proetale$ et $S_\proetale'$ admettent des produits fibrés
et des objets finaux, et que notre foncteur commute à ces constructions
(lemme \ref{proetale-lemma-fibre-products-proetale}).
Il résulte du lemme \ref{proetale-lemma-proetale-induced}
que le foncteur d'inclusion est cocontinu.
L'existence de $f$ et $g$ résulte donc de
Sites, lemme \ref{sites-lemma-bigger-site}.
L'égalité de (1) résulte de
Cohomologie sur les sites, lemme \ref{sites-cohomology-lemma-cohomology-bigger-site}.
L'assertion (2) découle de (1), car
$\mathcal{F} = g^{-1}g_*\mathcal{F} = g^{-1}f^{-1}\mathcal{F}$.
\end{proof}

\noindent
Nous démontrons ensuite un résultat analogue pour les grands topos pro-\'etales.

\begin{lemma}
\label{proetale-lemma-change-alpha}
Supposons donnés des grands sites $\Sch_\proetale$ et $\Sch'_\proetale$ comme dans la
définition \ref{proetale-definition-big-proetale-site}.
Supposons que $\Sch_\proetale$ soit contenu dans $\Sch'_\proetale$.
Le foncteur d'inclusion $\Sch_\proetale \to \Sch'_\proetale$ vérifie
les hypothèses de Sites, lemme \ref{sites-lemma-bigger-site}.
Il existe des morphismes de topos
\begin{eqnarray*}
g : \Sh(\Sch_\proetale) &
\longrightarrow &
\Sh(\Sch'_\proetale) \\
f : \Sh(\Sch'_\proetale) &
\longrightarrow &
\Sh(\Sch_\proetale)
\end{eqnarray*}
tels que $f \circ g \cong \text{id}$. Pour tout objet $S$
de $\Sch_\proetale$, le foncteur d'inclusion
$(\Sch/S)_\proetale \to (\Sch'/S)_\proetale$ vérifie
les hypothèses de Sites, lemme \ref{sites-lemma-bigger-site}.
On obtient donc de même des morphismes
\begin{eqnarray*}
g : \Sh((\Sch/S)_\proetale) &
\longrightarrow &
\Sh((\Sch'/S)_\proetale) \\
f : \Sh((\Sch'/S)_\proetale) &
\longrightarrow &
\Sh((\Sch/S)_\proetale)
\end{eqnarray*}
avec $f \circ g \cong \text{id}$.
\end{lemma}

\begin{proof}
Les hypothèses (b), (c) et (e) du
lemme de Sites \ref{sites-lemma-bigger-site}
sont immédiates pour les foncteurs
$\Sch_\proetale \to \Sch'_\proetale$ et
$(\Sch/S)_\proetale \to (\Sch'/S)_\proetale$. La propriété (a) résulte du
lemme \ref{proetale-lemma-proetale-induced}.
La propriété (d) est satisfaite parce que
les produits fibrés existent dans les catégories $\Sch_\proetale$, $\Sch'_\proetale$
et sont compatibles avec les produits fibrés dans la catégorie des schémas.
\end{proof}

\begin{lemma}
\label{proetale-lemma-cohomology-enlarge-partial-universe}
Soit $S$ un schéma. Soient $(\Sch/S)_\proetale$ et $(\Sch'/S)_\proetale$ deux
grands sites pro-\'etales de $S$ comme dans la
définition \ref{proetale-definition-big-small-proetale}.
Supposons que le premier soit contenu dans
le second. Alors
\begin{enumerate}
\item pour tout faisceau abélien $\mathcal{F}'$ défini sur $(\Sch'/S)_\proetale$
et tout objet $U$ de $(\Sch/S)_\proetale$, on a
$$
H^p(U, \mathcal{F}'|_{(\Sch/S)_\proetale}) =
H^p(U, \mathcal{F}')
$$
Autrement dit, la cohomologie de $\mathcal{F}'$ sur $U$ calculée dans le site le plus grand
coïncide avec la cohomologie de la restriction de $\mathcal{F}'$ au site le plus petit
sur $U$.
\item pour tout faisceau abélien $\mathcal{F}$ sur $(\Sch/S)_\proetale$, il existe un
faisceau abélien $\mathcal{F}'$ sur $(\Sch/S)_\proetale'$ dont la restriction à
$(\Sch/S)_\proetale$ est isomorphe à $\mathcal{F}$.
\end{enumerate}
\end{lemma}

\begin{proof}
D'après le lemme \ref{proetale-lemma-change-alpha}, le foncteur d'inclusion
$(\Sch/S)_\proetale \to (\Sch'/S)_\proetale$ vérifie les hypothèses de
Sites, lemme \ref{sites-lemma-bigger-site}. Cela implique (2), tandis que (1)
résulte de
Cohomologie sur les sites, lemme \ref{sites-cohomology-lemma-cohomology-bigger-site}.
\end{proof}














% OK, start here.
%

\chapter{Cycles relatifs} 



\phantomsection
\label{relative-cycles-section-phantom} 




\section{Introduction} 
\label{relative-cycles-section-introduction} 

\noindent
 Une référence fondamentale est \cite{SV}. 

\medskip\noindent
 Dans ce chapitre, nous ne définissons que les cycles relatifs universellement
 entiers de \cite{SV}. Ce choix permet de développer une théorie 
un peu plus simple que dans l'article original, au prix toutefois d'une certaine perte. 

\medskip\noindent
 Fixons un morphisme $X \to S$ de type fini 
entre schémas noethériens. Une famille $\alpha$ de $r$-cycles sur les fibres
 de $X/S$ est simplement une famille $\alpha = (\alpha_s)_{s \in S}$
 , où $\alpha_s \in Z_r(X_s)$. Le changement 
de base $g^*\alpha$ de $\alpha$ se définit immédiatement pour tout morphisme $g : S' \to S$.
 On dit alors que $\alpha$ est un {\it $r$-cycle relatif sur $X/S$}
 si $\alpha$ est compatible avec les spécialisations : pour tout 
morphisme $g : S' \to S$, où $S'$ est le spectre d'un anneau 
de valuation discrète, la fibre générique de $g^*\alpha$
 doit se spécialiser en la fibre fermée de $g^*\alpha$.
 Voir la section \ref{relative-cycles-section-families-specialization}. 






\section{Conventions et notation} 
\label{relative-cycles-section-conventions} 

\noindent
 On consultera le chapitre Homologie de Chow et classes de Chern pour
 nos conventions et notations relatives aux cycles sur les schémas 
localement de type fini sur une base noethérienne fixée ; voir 
Homologie de Chow, section \ref{chow-section-setup} et suivantes. 

\medskip\noindent
 En particulier, si  $X$  est localement de type fini sur un corps $k$, 
alors  $Z_r(X)$  désigne le groupe des cycles de dimension $r$, voir
 Homologie de Chow, exemple  \ref{chow-example-field}  et 
section \ref{chow-section-cycles}. Étant donné un sous-schéma fermé
 intègre  $Z \subset X$  avec $\dim(Z) = r$ , nous avons  $[Z] \in Z_r(X)$
  et si  $X$  est quasi-compact, alors $Z_r(X)$ est le groupe abélien libre engendré par ces classes. 




\section{Cycles relatifs aux corps} 
\label{relative-cycles-section-relative-fields} 

\noindent
 Soit $k$ un corps. Soit  $X$  un schéma localement algébrique sur  $k$. 
Soit  $r \geq 0$  un entier. Dans ce cadre, nous avons le groupe 
$Z_r(X)$  des  $r$-cycles sur  $X$, voir la section  \ref{relative-cycles-section-conventions}. 

\medskip\noindent
 {\bf  Changement de base. }  Pour toute extension de corps  $k'/k$  il existe une application de changement 
de base  $Z_r(X) \to Z_r(X_{k'})$, voir
 Homologie de Chow, section  \ref{chow-section-change-base}. À 
savoir, étant donné un sous-schéma fermé intègre  $Z \subset X$
  de dimension  $r$  on envoie  $[Z] \in Z_r(X)$  au $r$-cycle  
$[Z_{k'}]_r \in Z_r(X_{k'})$  associé au sous-schéma fermé  
$Z_{k'} \subset X_{k'}$  (bien sûr, en général  $Z_{k'}$
  n'est ni irréductible ni réduit). L'application de changement de base  
$Z_r(X) \to Z_r(X_{k'})$  est toujours injective. 

\begin{lemma} 
\label{relative-cycles-lemma-multiplicities-field-extension} 
Soit $K/k$ une extension de corps. Soit $Z$ un schéma intègre localement 
algébrique sur $k$. La multiplicité  $m_{Z', Z_K}$  d'une composante 
irréductible  $Z' \subset Z_K$  est  $1$  ou une puissance de la caractéristique de  $k$. 
\end{lemma} 

\begin{proof}
 Si la caractéristique de  $k$  est zéro, alors  $k$  est parfait et 
la multiplicité est toujours  $1$  puisque $X_K$ est réduit d'après
 Variétés, lemme  \ref{varieties-lemma-geometrically-reduced}. 
Supposons que la caractéristique de  $k$ soit  $p > 0$. 
Soit  $L$  le corps des fonctions de  $Z$. Puisque  $Z$  est localement algébrique
 sur  $k$, l'extension de corps $L/k$ est finiment engendrée. 
L'anneau  $K \otimes_k L$  est noethérien
 (Algèbre, lemme  \ref{algebra-lemma-Noetherian-field-extension}).
 En termes algébriques, il faut montrer que la longueur de l'anneau
 local artinien  $(K \otimes_k L)_\mathfrak q$
  est une puissance de  $p$ pour chaque idéal premier minimal  $\mathfrak q$. 

\medskip\noindent
 Soit  $L'/L$  une extension finie purement inséparable, disons de degré 
$p^n$. Alors  $K \otimes_k L \subset K \otimes_k L'$  est un homomorphisme
 d'anneaux fini et libre de rang $p^n$  qui induit un homéomorphisme sur les 
spectres et des extensions de corps résiduels purement inséparables. 
Ainsi, pour chaque idéal premier minimal $\mathfrak q$  comme ci-dessus,
 il existe un unique idéal premier minimal 
$\mathfrak q' \subset K \otimes_k L'$  au-dessus de celui-ci et 
$$
p^n \text{longueur}((K \otimes_k L)_\mathfrak q) =
[\kappa(\mathfrak q') : \kappa(\mathfrak q)]
\text{longueur}((K \otimes_k L')_{\mathfrak q'})
$$
 par Algèbre, lemme  \ref{algebra-lemma-pushdown-module}  appliqué 
à  $M = (K \otimes_k L')_{\mathfrak q'} \cong
(K \otimes_k L)_{\mathfrak q}^{\oplus p^n}$. 
Puisque  $[\kappa(\mathfrak q') : \kappa(\mathfrak q)]$  est une puissance
 de  $p$ , nous concluons qu'il suffit de prouver 
l'énoncé pour  $L'$  et  $\mathfrak q'$. 

\medskip\noindent
 D'après le paragraphe précédent et Algèbre, lemme  \ref{algebra-lemma-make-separable} , nous 
pouvons supposer qu'il existe une tour de corps $L/k'/k$ telle que $L/k'$ soit séparable et
 que $k'/k$ soit finie purement inséparable. Alors  $K \otimes_k k'$ est 
un anneau local artinien. L'argument du paragraphe précédent 
(appliqué à $L = k$  et  $L' = k'$) montre que  $\text{longueur}(K \otimes_k k')$
  est une puissance de $p$. Comme $L/k'$ est la localisation 
d'une $k'$-algèbre 
lisse (Algèbre, lemme \ref{algebra-lemma-localization-smooth-separable}),
 $S = (K \otimes_k L)_\mathfrak q$ est la localisation d'une 
$R = K \otimes_k k'$-algèbre lisse en un idéal premier minimal.
 Donc  $R \to S$  est un homomorphisme local plat d'anneaux 
locaux artiniens et  $\mathfrak m_R S = \mathfrak m_S$. Il
 découle de Algèbre, lemme  \ref{algebra-lemma-pullback-module}  que  
$\text{longueur}(K \otimes_k k') = \text{longueur}(R) =
\text{longueur}(S) = \text{longueur}((K \otimes_k L)_\mathfrak q)$
  et la preuve est terminée. 
\end{proof} 

\begin{lemma} 
\label{relative-cycles-lemma-how-different} 
Soit  $k$  un corps de caractéristique  $p > 0$  avec clôture parfaite $k^{perf}$. 
Soit  $X$  un schéma algébrique sur  $k$. Soit $r \geq 0$  un entier. Le 
conoyau de l'application injective  $Z_r(X) \to Z_r(X_{k^{perf}})$  est un module de torsion de 
$p$-puissance (Compléments sur l'algèbre, Définition 
\ref{more-algebra-definition-f-power-torsion}). 
\end{lemma} 

\begin{proof} 
Puisque  $X$  est quasi-compact, le groupe abélien  $Z_r(X)$  est libre avec une 
base donnée par les sous-schémas fermés intègres de dimension $r$. De même pour  
$Z_r(X_{k^{perf}})$. 
Puisque  $X_{k^{perf}} \to X$  est un homéomorphisme, il en résulte
 que  $Z_r(X) \to Z_r(X_{k^{perf}})$  est injectif avec conoyau de torsion. Chaque
 élément du conoyau est de torsion de $p$-puissance d'après le
 lemme  \ref{relative-cycles-lemma-multiplicities-field-extension}. 
\end{proof} 





\section{Spécialisation des cycles} 
\label{relative-cycles-section-specialization} 

\noindent
 Soit  $R$  un anneau de valuation discrète avec corps des fractions $K$
  et corps résiduel  $\kappa$. Soit  $X$  un schéma localement de type fini 
sur  $R$. Soit  $r \geq 0$. Il existe une application de spécialisation 
$$
sp_{X/R} : Z_r(X_K) \longrightarrow Z_r(X_\kappa)
$$
 définie comme suit. Pour un sous-schéma fermé et intègre  $Z \subset X_K$
  de dimension $r$ , nous notons $\overline{Z}$ l'image schématique de
  $Z \to X$. On définit alors $sp_{X/R}$ comme l'unique application $\mathbf{Z}$-linéaire
 telle que 
$$
sp_{X/R}([Z]) = [\overline{Z}_\kappa]_r
$$
 Expliquons brièvement pourquoi cette définition est correcte. Tout d'abord, observons que le
 morphisme  $X_K \to X$  est quasi-compact et donc le morphisme  $Z \to X$
 est quasi-compact. La formation de l'image schématique de $Z \to X$
 commute donc aux changements de base plats 
d'après Morphismes, lemme  \ref{morphisms-lemma-flat-base-change-scheme-theoretic-image}. En 
particulier, après changement de base à $X_K$, on obtient $Z = \overline{Z}_K$. 
Puisque  $Z$  est intègre, bien sûr  $\overline{Z}$  est également intègre
 et en fait est égal au sous-schéma fermé intègre unique dont le point
 générique est l'image du point générique de  $Z$. Il résulte de 
Variétés, lemme  \ref{varieties-lemma-dominate-valuation-ring-dimension-fibres}  
que  $Z_\kappa$  est équidimensionnel de dimension $r$. 

\begin{lemma} 
\label{relative-cycles-lemma-specialization-module} 
Soit  $R$  un anneau de valuation discrète avec corps des fractions $K$  et corps résiduel  
$\kappa$. Soit  $X$  un schéma localement de type fini sur  $R$. Soit  $r \geq 0$. 
Soit  $\mathcal{F}$  un $\mathcal{O}_X$-module cohérent plat sur  $R$. Supposons  
$\dim(\text{Supp}(\mathcal{F}_K)) \leq r$. Alors 
$\dim(\text{Supp}(\mathcal{F}_\kappa)) \leq r$  et 
$$
sp_{X/R}([\mathcal{F}_K]_r) = [\mathcal{F}_\kappa]_r
$$
\end{lemma} 

\begin{proof} 
L'énoncé sur la dimension découle de Compléments sur les morphismes, lemme 
\ref{more-morphisms-lemma-relative-dimension-support-flat}. 
Soit  $x$  un point générique d'un sous-schéma fermé intègre  
$Z \subset X_\kappa$ de dimension  $r$. Pour terminer la preuve,
 nous montrerons que le coefficient de $[Z]$
 dans le membre de gauche (L) et le membre de droite (R) de l'égalité est le même. 

\medskip\noindent
 Posons $A = \mathcal{O}_{X, x}$ et $M = \mathcal{F}_x$.
 Remarquons que  $M$  est un $A$-module fini plat sur  $R$. 
Soit  $\pi \in R$  une uniformisante de sorte que  
$A/\pi A = \mathcal{O}_{X_\kappa, x}$.
 D'après Homologie de Chow, lemme \ref{chow-lemma-additivity-divisors-restricted} , nous
 avons 
$$
\sum\nolimits_i \text{longueur}_A(A/(\pi, \mathfrak q_i))
\text{longueur}_{A_{\mathfrak q_i}}(M_{\mathfrak q_i}) =
\text{longueur}_A(M/\pi M)
$$
 où la somme porte sur les idéaux premiers minimaux  
$\mathfrak q_i$  dans le support de $M$. 
Puisque  $\pi$  est un non-diviseur de zéro sur  $M$ , nous voyons
 que  $\pi \not \in \mathfrak q_i$  et donc ces
 idéaux premiers correspondent aux points génériques  $y_i \in X_K$  du 
support de  $\mathcal{F}_K$  qui se spécialisent en notre $x \in X_\kappa$ 
choisi. Ainsi, le membre de gauche est le coefficient de  $[Z]$
  dans (L). Bien sûr,  $\text{longueur}_A(M/\pi M)$  est le coefficient 
de  $[Z]$  dans (R). Cela termine la démonstration. 
\end{proof} 

\begin{lemma} 
\label{relative-cycles-lemma-specialization-closed} 
Soit  $R$  un anneau de valuation discrète avec corps des fractions $K$  et corps résiduel  
$\kappa$. Soit  $X$  un schéma localement de type fini sur  $R$. Soit  $r \geq 0$. 
Soit  $W \subset X$  un sous-schéma fermé plat sur  $R$. Supposons  
$\dim(W_K) \leq r$. Alors $\dim(W_\kappa) \leq r$  et 
$$
sp_{X/R}([W_K]_r) = [W_\kappa]_r
$$
\end{lemma} 

\begin{proof} En
 prenant  $\mathcal{F} = \mathcal{O}_W$ , il s'agit d'un cas particulier du
 lemme  \ref{relative-cycles-lemma-specialization-module}. Voir 
Homologie de Chow, lemme \ref{chow-lemma-cycle-closed-coherent}. 
\end{proof} 

\begin{lemma} 
\label{relative-cycles-lemma-specialization-extension} 
Soit  $R'/R$  une extension d'anneaux de valuation discrète induisant une extension de corps des
 fractions  $K'/K$  et une extension de corps résiduels $\kappa'/\kappa$
  (Compléments sur l'algèbre, définition  
\ref{more-algebra-definition-extension-discrete-valuation-rings}). 
Soit  $X$ localement de type fini sur  $R$. On note  $X' = X_{R'}$.
 Alors le diagramme 
$$
\xymatrix{
Z_r(X'_{K'}) \ar[rr]_{sp_{X'/R'}} & & Z_r(X'_{\kappa'}) \\
Z_r(X_K) \ar[rr]^{sp_{X/R}} \ar[u] & & Z_r(X_\kappa) \ar[u]
}
$$
 est commutatif, où  $r \geq 0$  et les flèches verticales sont des applications de changement de base. 
\end{lemma} 

\begin{proof} 
On observe que  $X'_{K'} = X_{K'} = X_K \times_{\Spec(K)} \Spec(K')$
  et de même pour les fibres fermées, de sorte que les flèches verticales ont effectivement
 un sens (voir la section  \ref{relative-cycles-section-relative-fields}).
 Si  $Z \subset X_K$  est un sous-schéma fermé intègre 
avec image schématique  $\overline{Z} \subset X$, on voit que  
$Z_{K'} \subset X_{K'}$  est un sous-schéma fermé avec image schématique
  $\overline{Z}_{R'} \subset X_{R'}$. Le changement de base de  
$[Z]$ est  $[Z_{K'}]_r = [\overline{Z}_{K'}]_r$  par définition. Nous avons 
$$
sp_{X/R}([Z]) = [\overline{Z}_\kappa]_r
\quad\text{et}\quad
sp_{X'/R'}([\overline{Z}_{K'}]_r) = [(\overline{Z}_{R'})_{\kappa'}]_r
$$
 d'après le lemme \ref{relative-cycles-lemma-specialization-module}. Comme  
$(\overline{Z}_{R'})_{\kappa'} = (\overline{Z}_\kappa)_{\kappa'}$
 , nous concluons. 
\end{proof} 

\begin{lemma} 
\label{relative-cycles-lemma-specialization-flat-pullback} 
Soit  $R$  un anneau de valuation discrète avec corps des fractions $K$  et corps résiduel  
$\kappa$. Soit  $X$  un schéma localement de type fini sur  $R$. 
Soit  $f : X' \to X$  un morphisme qui est localement de type fini, plat 
et de dimension relative  $e$. Alors le diagramme 
$$
\xymatrix{
Z_{r + e}(X'_K) \ar[rr]_{sp_{X'/R}} & & Z_{r + e}(X'_\kappa) \\
Z_r(X_K) \ar[rr]^{sp_{X/R}} \ar[u] & & Z_r(X_\kappa) \ar[u]
}
$$
 est commutatif, où  $r \geq 0$  et les flèches verticales sont données par 
l'image inverse plate. 
\end{lemma} 

\begin{proof} 
Soit  $Z \subset X$  un sous-schéma fermé intègre dominant  $R$.
 Par la construction de  $sp_{X/R}$ , nous avons  $sp_{X/R}([Z_K]) = [Z_\kappa]_r$
  ce qui caractérise l'application de spécialisation. 
Posons  $Z' = f^{-1}(Z) = X' \times_X Z$.
 Comme  $R$  est un anneau de valuation,  $Z$  est plat sur  $R$. 
Par conséquent, $Z'$  est plat sur  $R$  et  
$sp_{X'/R}([Z'_K]_{r + e}) = [Z'_\kappa]_{r + e}$
  d'après le lemme \ref{relative-cycles-lemma-specialization-closed}. 
D'après Homologie de Chow, lemme \ref{chow-lemma-pullback-coherent},
 on a $f_K^*[Z_K] = [Z'_K]_{r + e}$  et  
$f_\kappa^*[Z_\kappa]_r = [Z'_\kappa]_{r + e}$. Nous concluons. 
\end{proof} 

\begin{lemma} 
\label{relative-cycles-lemma-specialization-proper-pushforward} 
Soit  $R$  un anneau de valuation discrète de corps des fractions $K$ et de corps résiduel  
$\kappa$. Soit  $f : X \to Y$  un morphisme propre de schémas localement 
de type fini sur  $R$. Alors le diagramme 
$$
\xymatrix{
Z_r(X_K) \ar[rr]_{sp_{X/R}} \ar[d] & & Z_r(X_\kappa) \ar[d] \\
Z_r(Y_K) \ar[rr]^{sp_{Y/R}} & & Z_r(Y_\kappa)
}
$$
 est commutatif, où  $r \geq 0$  et les flèches verticales sont données 
par l'image directe propre. 
\end{lemma} 

\begin{proof} 
Soit  $Z \subset X$  un sous-schéma fermé intègre dominant  $R$.
 Par la construction de  $sp_{X/R}$ , nous avons  $sp_{X/R}([Z_K]) = [Z_\kappa]_r$
  ce qui caractérise l'application de spécialisation. 
Posons  $Z' = f(Z) \subset Y$. Alors  $Z'$  est un sous-schéma fermé intègre 
de  $Y$  dominant  $R$. Ainsi $sp_{Y/R}([Z'_K]) = [Z'_\kappa]_r$. 

\medskip\noindent
 Nous pouvons considérer  $[Z]$  comme un élément de  $Z_{r + 1}(X)$. Par définition, 
nous avons  $f_*[Z] = 0$  si  $\dim(Z') < r + 1$  et  $f_*[Z] = d[Z']$
  si $Z \to Z'$  est génériquement fini de degré  $d$.
 Comme l'image directe propre commute avec l'image inverse plate par  $Y_K \to Y$
  (Homologie de Chow, lemme \ref{chow-lemma-flat-pullback-proper-pushforward}),
 on obtient respectivement  $f_{K, *}[Z_K] = 0$  ou  $f_{K, *}[Z_K] = d[Z'_K]$.
 Appliquons Homologie de Chow, lemme \ref{chow-lemma-closed-in-X-gysin}, 
au diagramme commutatif 
$$
\xymatrix{
X_\kappa \ar[d] \ar[r]_i & X \ar[d] \\
Y_\kappa \ar[r]^j & Y
}
$$
 Nous obtenons que  $f_{\kappa, *}[Z_\kappa]_r = 0$  ou  
$f_{\kappa, *}[Z_\kappa] = d[Z'_\kappa]_r$  car
  $i^*[Z] = [Z_k]_r$ et  $j^*[Z'] = [Z'_\kappa]_r$.
 En regroupant le tout, nous concluons. 
\end{proof} 









\section{Familles de cycles sur les fibres} 
\label{relative-cycles-section-cycles-fibres} 

\noindent
 Soit $f : X \to S$ un morphisme de schémas localement de type fini et 
soit $r \geq 0$ un entier. Une 
{\it famille $\alpha$ de $r$-cycles sur les fibres de $X/S$} est une famille 
$$
\alpha = (\alpha_s)_{s \in S}
$$
 indexée par les points $s$ de $S$, où $\alpha_s \in Z_r(X_s)$
 est un $r$-cycle sur la fibre schématique $X_s$ de $f$ en $s$.
 Il existe diverses constructions que nous pouvons effectuer sur des familles de  
$r$-cycles sur les fibres. 

\medskip\noindent
 {\bf  Changement de base. }  Soit 
$$
\xymatrix{
X' \ar[r] \ar[d] & X \ar[d]^f \\
S' \ar[r]^g & S
}
$$
 un carré cartésien de morphismes de schémas avec  $f$  localement de type fini. 
Soit  $r \geq 0$  un entier. Étant donné une famille  $\alpha$  de $r$-cycles sur les
 fibres de  $X/S$ , nous définissons le {\it  changement de base }   $g^*\alpha$  de $\alpha$
  comme étant la famille 
$$
g^*\alpha = (\alpha'_{s'})_{s' \in S'}
$$
 où  $\alpha'_{s'} \in Z_r(X'_{s'})$  est le changement de 
base du cycle  $\alpha_s$  avec  $s = g(s')$ comme à 
la section  \ref{relative-cycles-section-relative-fields}  via l'identification  
$X'_{s'} = X_s \times_{\Spec(\kappa(s))} \Spec(\kappa(s'))$
  des fibres schématiques. 

\medskip\noindent
 {\bf  Restriction.}  Soit  $f : X \to S$  un morphisme de schémas qui est localement
 de type fini. Soit  $r \geq 0$  un entier. Soient $U \subset X$  et  
$V \subset S$  des sous-schémas ouverts avec  $f(U) \subset V$. Étant donnée une famille  
$\alpha$  de  $r$-cycles sur les fibres de $X/S$ , nous pouvons définir la 
 {\it restriction}   $\alpha|_U$  de  $\alpha$  comme la 
famille de $r$-cycles sur les fibres de  $U/V$
$$
\alpha|_U = (\alpha_s|_{U_s})_{s \in V}
$$
 des restrictions aux fibres schématiques. 

\medskip\noindent
 {\bf Image inverse plate.}  Soit  $X \to S$  un morphisme de schémas qui est localement 
de type fini. Soient  $r, e \geq 0$  des entiers. Soit $f : X' \to X$  un 
morphisme plat, localement de type fini, et de dimension relative $e$. 
Étant donnée une famille  $\alpha$  de  $r$-cycles 
sur les fibres de  $X/S$ , nous définissons l'{\it image inverse plate}   $f^*\alpha$  de  $\alpha$
  comme la famille de $(r + e)$-cycles sur les fibres 
$$
f^*\alpha = (f_s^*\alpha_s)_{s \in S}
$$
 où  $f_s^*\alpha_s \in Z_{r + e}(X'_s)$  est l'image inverse 
plate du cycle  $\alpha_s$  dans  $Z_r(X_s)$  par le morphisme plat  
$f_s : X'_s \to X_s$  de dimension relative $e$
 induit sur les fibres schématiques. 

\medskip\noindent
 {\bf Image directe propre.}  Soit 
$$
\xymatrix{
X \ar[rr]_f \ar[rd] & & Y \ar[ld] \\
& S
}
$$
 un diagramme commutatif de morphismes de schémas avec  $X$  et $Y$
  localement de type fini sur  $S$  et  $f$  propre. Soit $r \geq 0$  un entier. Étant 
donné une famille  $\alpha$  de $r$-cycles sur les fibres de  $X/S$ , nous définissons 
l'{\it image directe propre}   $f_*\alpha$  de $\alpha$  comme étant la famille de  
$r$-cycles sur les fibres de $Y/S$  par 
$$
f_*\alpha = (f_{s, *}\alpha_s)_{s \in S}
$$
 où  $f_{s, *}\alpha_s \in Z_r(Y_s)$  est l'image directe propre
 du cycle  $\alpha_s$  dans  $Z_r(X_s)$ par le morphisme propre  
$f_s : X_s \to Y_s$  des fibres schématiques. 

\begin{lemma} 
\label{relative-cycles-lemma-compatibilities} Les 
opérations précédentes vérifient les compatibilités suivantes : (1) le changement
 de base est fonctoriel ; (2) la 
restriction est la composée d'un changement de base et d'un cas particulier d'image inverse
 plate ; (3) 
l'image inverse plate commute au changement de base ; 
(4) l'image inverse plate est fonctorielle
 ; (5) l'image directe propre commute au changement de 
base ; (6) l'image directe propre est fonctorielle 
; (7) l'image directe propre commute à l'image inverse plate. 
\end{lemma} 

\begin{proof}
 Chacune de ces compatibilités découle directement des résultats 
correspondants établis dans le chapitre Homologie de Chow et appliqués aux fibres
 sur $S$ des schémas considérés. Nous omettons les énoncés précis et les 
preuves détaillées. Voici quelques références. Partie
 (1) : Homologie de Chow, lemme 
\ref{chow-lemma-compose-base-change}. Partie 
(2) : c'est immédiat. 
Partie (3) : Homologie de Chow, lemme 
\ref{chow-lemma-pullback-base-change-pullback}. Partie
 (4) : Homologie de Chow, lemme 
\ref{chow-lemma-compose-flat-pullback}. Partie
 (5) : Homologie de Chow, lemme 
\ref{chow-lemma-pullback-base-change-pushforward}. Partie
 (6) : Homologie de Chow, lemme 
\ref{chow-lemma-compose-pushforward}. Partie
 (7) : Homologie de Chow, lemme 
\ref{chow-lemma-flat-pullback-proper-pushforward}. 
\end{proof} 

\begin{example} 
\label{relative-cycles-example-family-associated-module} 
Soit  $f : X \to S$  un morphisme de schémas qui est localement de type fini. 
Soit $r \geq 0$  un entier. Soit  $\mathcal{F}$  un  
$\mathcal{O}_X$-module quasi-cohérent de type fini. Pour  $s \in S$  notons  $\mathcal{F}_s$
  l'image inverse de $\mathcal{F}$ sur $X_s$. 
Supposons  $\dim(\text{Supp}(\mathcal{F}_s)) \leq r$  pour tout  $s \in S$.
 Alors nous pouvons associer à  $\mathcal{F}$  la famille  $[\mathcal{F}/X/S]_r$  de 
$r$-cycles sur les fibres de  $X/S$  définie par la formule 
$$
[\mathcal{F}/X/S]_r = ([\mathcal{F}_s]_r)_{s \in S}
$$
 où  $[\mathcal{F}_s]_r$  est donné par Homologie de Chow, définition  
\ref{chow-definition-cycle-associated-to-coherent-sheaf}. 
\end{example} 

\begin{lemma} 
\label{relative-cycles-lemma-family-associated-module} La
 construction dans l'exemple  \ref{relative-cycles-example-family-associated-module}  est 
compatible avec le changement de base, la restriction et l'image inverse plate. 
\end{lemma} 

\begin{proof} 
Voir Homologie de Chow, lemmes  
\ref{chow-lemma-pullback-coherent-base-change}  et  
\ref{chow-lemma-pullback-coherent}. 
\end{proof} 

\begin{example} 
\label{relative-cycles-example-family-associated-closed} 
Soit  $f : X \to S$  un morphisme de schémas qui est localement de type fini. 
Soit $r \geq 0$  un entier. Soit  $Z \subset X$  un sous-schéma fermé. 
Pour  $s \in S$ , notons  $Z_s$  l'image réciproque de  $Z$  dans  $X_s$
  ou, de manière équivalente, la fibre schématique de  $Z$  en  $s$  vue 
comme un sous-schéma fermé de  $X_s$. 
Supposons  $\dim(Z_s) \leq r$  pour tout $s \in S$.
 Alors nous pouvons associer à  $Z$  la famille  $[Z/X/S]_r$
 de  $r$-cycles sur les fibres de  $X/S$  définie par la formule 
$$
[Z/X/S]_r = ([Z_s]_r)_{s \in S}
$$
 où  $[Z_s]_r$  est donné par
 Homologie de Chow, définition  
\ref{chow-definition-cycle-associated-to-closed-subscheme}. 
\end{example} 

\begin{lemma} 
\label{relative-cycles-lemma-family-associated-closed} La
 construction dans l'exemple  \ref{relative-cycles-example-family-associated-closed}  est 
compatible avec le changement de base, la restriction et l'image inverse plate. 
\end{lemma} 

\begin{proof} En
 prenant  $\mathcal{F} = (Z \to X)_*\mathcal{O}_Z$ , il s'agit d'un cas particulier du
 lemme  \ref{relative-cycles-lemma-family-associated-module}. Voir 
Homologie de Chow, lemme \ref{chow-lemma-cycle-closed-coherent}. 
\end{proof} 

\begin{remark}[Support] 
\label{relative-cycles-remark-supports-family} 
Soit $f : X \to S$ un morphisme de schémas qui est localement de type fini. 
Soit $r \geq 0$ un entier. Soit $\alpha$ une famille de $r$-cycles 
sur les fibres de $X/S$. Nous définissons le {\it support} de $\alpha$ comme 
$$
\text{Supp}(\alpha) =
\bigcup\nolimits_{s \in S} \text{Supp}(\alpha_s) \subset X
$$
 Ici,  $\text{Supp}(\alpha_s) \subset X_s$  est 
le support du cycle  $\alpha_s$, voir
 Homologie de Chow, définition  \ref{chow-definition-support-cycle}. 
Le support  $\text{Supp}(\alpha)$  est rarement un sous-ensemble fermé de  $X$. 
\end{remark} 

\begin{lemma} 
\label{relative-cycles-lemma-support-family} La 
formation du support comme dans la remarque  \ref{relative-cycles-remark-supports-family}  est 
compatible avec le changement de base, la restriction et l'image inverse plate. 
\end{lemma} 

\begin{proof} 
Démonstration omise. 
\end{proof} 

\begin{lemma} 
\label{relative-cycles-lemma-coequalizer-dim-r} 
Soit  $f : X \to S$  un morphisme de schémas qui est localement de type fini. 
Soit $r \geq 0$  un entier. Soit  $g : S' \to S$  un morphisme de schémas 
surjectif. Posons  $S'' = S' \times_S S'$  et soient  $f' : X' \to S'$
  et  $f'' : X'' \to S''$  les changements de base de  $f$. 
Soit  $x \in X$  avec  $\text{trdeg}_{\kappa(f(x))}(\kappa(x)) = r$. 
\begin{enumerate} 
\item Il existe un point $x' \in X'$ se projetant sur $x$
  avec  $\text{trdeg}_{\kappa(f'(x'))}(\kappa(x')) = r$. 
\item Si  $x'_1, x'_2 \in X'$  sont tous deux comme dans (1), alors 
il existe un  $x'' \in X''$ avec  
$\text{trdeg}_{\kappa(f''(x''))}(\kappa(x'')) = r$  et  
$\text{pr}_i(x'') = x'_i$. 
\end{enumerate} 
\end{lemma} 

\begin{proof} La
 partie (1) résulte
 de Morphismes, lemme  \ref{morphisms-lemma-dimension-fibre-after-base-change}. 
Soit  $x'_1, x'_2$  comme dans (2). Alors, puisque  $X'' = X' \times_X X'$
 , nous voyons qu'il 
existe un point $x'' \in X''$ se projetant à la fois sur  $x'_1$  et  $x'_2$  (voir par
 exemple Descente, lemme \ref{descent-lemma-equiv-fibre-product}). On
 note  $s'' \in S''$,  $s'_i \in S'$ et  $s \in S$  les images 
de $x''$,  $x'_i$ et  $x$.
 Posons  $k = \kappa(s)$. Soit $Z \subset X_k$  le sous-schéma 
fermé intègre dont le point générique est $x$. Alors  $x'_i$
  est un point générique d'une composante irréductible de  $Z_{\kappa(s'_i)}$. 
Soit  $Z'' \subset Z_{\kappa(s'')}$  une composante irréductible
 contenant  $x''$. On note  $\xi'' \in Z''$  le point générique.
 Comme $\xi'' \leadsto x''$ , nous voyons que  $\xi''$  doit également se
 projeter sur $x'_i$ sous les deux projections. Par ailleurs, nous 
voyons que  $\text{trdeg}_{\kappa(s'')}(\kappa(\xi'')) = r$
  parce que c'est un point 
générique d'une composante irréductible du changement de base de $Z$. 
\end{proof} 

\begin{lemma} 
\label{relative-cycles-lemma-descend-family} 
Soit  $f : X \to S$  un morphisme de schémas qui est localement de type fini. 
Soit $r \geq 0$  un entier. Soit  $g : S' \to S$  un morphisme de 
schémas et  $X' = S' \times_S X$. Supposons que pour tout  $s \in S$  il 
existe un point  $s' \in S'$  avec $g(s') = s$  et tel que  
$\kappa(s')/\kappa(s)$  est une extension séparable de corps. Alors 
\begin{enumerate} 
\item Si  $\alpha_1$  et  $\alpha_2$  sont des familles de  $r$-cycles sur les fibres de  $X/S$
  et si  $g^*\alpha_1 = g^*\alpha_2$, alors  $\alpha_1 = \alpha_2$. 
\item Étant donné une famille  $\alpha'$  de  $r$-cycles sur les fibres de $X'/S'$ , si  
$\text{pr}_1^*\alpha' = \text{pr}_2^*\alpha'$  en tant que familles de  
$r$-cycles sur les fibres de  $(S' \times_S S') \times_S X / (S' \times_S S')$,
 alors il existe une famille unique  $\alpha$  de $r$-cycles sur les fibres de  $X/S$
  telle que  $g^*\alpha = \alpha'$. 
\end{enumerate} 
\end{lemma} 

\begin{proof} 
La partie (1) découle de l'injectivité de l'application de changement de base étudiée à
 la section  \ref{relative-cycles-section-relative-fields}. (Cet argument fonctionne 
tant que $S' \to S$  est surjectif.) 

\medskip\noindent
 Soit  $\alpha'$  comme dans (2). On note  
$\alpha'' = \text{pr}_1^*\alpha' = \text{pr}_2^*\alpha'$
  la valeur commune. 

\medskip\noindent
 Soit  $(X/S)^{(r)}$  l'ensemble des  $x \in X$  avec  
$\text{trdeg}_{\kappa(f(x))}(\kappa(x)) = r$
  et définissons de même  $(X'/S')^{(r)}$  et  $(X''/S'')^{(r)}$
 . En considérant les coefficients, nous pouvons considérer  $\alpha'$  et  $\alpha''$  comme des fonctions  
$\alpha' : (X'/S')^{(r)} \to \mathbf{Z}$ et  
$\alpha'' : (X''/S'')^{(r)} \to \mathbf{Z}$. 
Étant donnée une fonction 
$$
\varphi : (X/S)^{(r)} \to \mathbf{Z}
$$
 nous définissons  $g^*\varphi : (X'/S')^{(r)} \to \mathbf{Z}$  par analogie avec
 notre opération de changement de base. Plus précisément, supposons que  $x' \in (X'/S')^{(r)}$
  se projette sur $x \in X$,  $s' \in S'$ et $s \in Z$.
 On note  $Z' \subset X'_{s'}$  et  $Z \subset X_s$  les sous-schémas
 fermés intègres avec les points génériques  $x'$  et  $x$. Observons
 que $\dim(Z') = r$. Si  $\dim(Z) < r$, on pose  $(g^*\varphi)(x') = 0$. 
Si  $\dim(Z) = r$, alors  $Z'$  est une composante irréductible de  $Z_{s'}$  et 
possède donc une multiplicité  $m_{Z', Z_{s'}}$. Notons cette multiplicité $m(x', g)$.
 On définit alors 
$$
(g^*\varphi)(x') = m(x', g) \varphi(x)
$$
 Les coefficients  $m(x', g)$  sont toujours des entiers positifs
 (voir par exemple le lemme  \ref{relative-cycles-lemma-multiplicities-field-extension}). De même, nous
 avons des applications de changement de base 
$$
\text{pr}_1^*, \text{pr}_2^* :
\text{Map}((X'/S')^{(r)}, \mathbf{Z})
\longrightarrow
\text{Map}((X''/S'')^{(r)}, \mathbf{Z})
$$
 Il découle de l'associativité du changement de base que nous avons  
$\text{pr}_1^* \circ g^* = \text{pr}_2^* \circ g^*$ (petit détail 
omis). Plus explicitement, en termes d'applications ensemblistes, cette 
égalité signifie simplement que pour  $x'' \in (X''/S'')^{(r)}$  nous avons 
$$
m(x'', \text{pr}_1) m(\text{pr}_1(x''), g) =
m(x'', \text{pr}_2) m(\text{pr}_2(x''), g)
$$
 pourvu que  $\text{pr}_1(x'')$  et  $\text{pr}_2(x'')$  soient 
dans  $(X''/S'')^{(r)}$.
 Le lemme \ref{relative-cycles-lemma-coequalizer-dim-r}  et un argument 
élémentaire\footnote{Étant donné $x \in (X/S)^{(r)}$ , choisissez $x' \in (X'/S')^{(r)}$
 se projetant sur  $x$  et définissez $\alpha(x) = \alpha'(x')/m(x', g)$. Cela 
est bien défini par la formule et le lemme.}, 
appliqués à l'équation affichée précédente, 
montrent qu'il existe une application unique 
$$
\alpha : (X/S)^{(r)} \to \mathbf{Q}
$$
 telle que $g^*\alpha = \alpha'$. Pour terminer la preuve, il suffit 
de montrer que $\alpha$  prend des valeurs entières (petit détail omis:
 il faut voir que $\alpha$  détermine une somme localement finie
 sur chaque fibre, ce qui découle du fait correspondant pour  $\alpha'$). 
Étant donné un  $x \in (X/S)^{(r)}$  avec image $s \in S$
 , nous pouvons choisir un point  $s' \in S'$  tel que  $\kappa(s')/\kappa(s)$
  soit séparable. On peut ensuite choisir  $x' \in (X'/S')^{(r)}$  qui se projette
 sur $s$  et  $x$  et nous voyons que  $m(x', g) = 1$  parce que 
$Z_{s'}$  est réduit dans ce cas. Ainsi $\alpha(x) = \alpha'(x')$
 est un entier. 
\end{proof} 

\begin{lemma} 
\label{relative-cycles-lemma-pullback-universally-bijective} 
Soit  $g : S' \to S$  un morphisme bijectif de schémas qui
 induit des isomorphismes des corps résiduels. 
Soit  $f : X \to S$  localement de type fini. Posons $X' = S' \times_S X$. 
Soit  $r \geq 0$. Alors le changement de base par  $g$ détermine une bijection
 entre le groupe des familles de  $r$-cycles sur les fibres de  $X/S$  et
 le groupe des familles de  $r$-cycles sur les fibres de  $X'/S'$. 
\end{lemma} 

\begin{proof} La 
démonstration est omise. 
\end{proof} 








\section{Cycles relatifs} 
\label{relative-cycles-section-families-specialization} 

\noindent
 Voici la définition que nous utiliserons ; voir la section  \ref{relative-cycles-section-compare}  
pour une comparaison avec les définitions de  \cite{SV}. 

\begin{definition} 
\label{relative-cycles-definition-relative-cycles} 
Soit  $S$  un schéma localement noethérien. Soit  $f : X \to S$  un morphisme de 
schémas qui est localement de type fini. Soit  $r \geq 0$  un entier. 
Un {\it $r$-cycle relatif sur $X/S$} est une famille  $\alpha$  de  $r$-cycles 
sur les fibres de  $X/S$ telle que pour tout morphisme  $g : S' \to S$
  où  $S'$  est le spectre d'un anneau de valuation discrète, nous avons 
$$
sp_{X'/S'}(\alpha_\eta) = \alpha_0
$$
 où  $sp_{X'/S'}$  est défini à la section  \ref{relative-cycles-section-specialization}  
et  $\alpha_\eta$  (resp. \ $\alpha_0$) est la valeur du changement de base  
$g^*\alpha$  de $\alpha$  au point générique (resp. fermé \ ) de  $S'$.
 Le groupe de tous les  $r$-cycles relatifs sur  $X/S$  est noté $z(X/S, r)$. 
\end{definition} 

\begin{lemma} 
\label{relative-cycles-lemma-relative-cycle-functoriality} 
Soit  $\alpha$  un  $r$-cycle relatif sur  $X/S$  comme 
dans la définition  \ref{relative-cycles-definition-relative-cycles}. Toute 
restriction, tout changement de base, toute image inverse plate ou toute image directe propre
 de $\alpha$ est alors un $r$-cycle relatif. 
\end{lemma} 

\begin{proof} Pour
 l'image inverse plate, on applique le lemme \ref{relative-cycles-lemma-specialization-flat-pullback}. 
La restriction en est un cas particulier. Pour le changement de base, 
on utilise sa transitivité. Enfin, pour l'image 
directe propre, on applique le lemme \ref{relative-cycles-lemma-specialization-proper-pushforward}. 
\end{proof} 

\begin{lemma} 
\label{relative-cycles-lemma-relative-cycles-h-descent} 
Soit  $f : X \to S$  un morphisme de schémas. Supposons que  $S$  soit localement noethérien et
 que  $f$  soit localement de type fini. Soit  $r \geq 0$  un entier. Soit  $\alpha$
  une famille de  $r$-cycles sur les fibres de $X/S$. Soit  $\{g_i : S_i \to S\}$
  un recouvrement h (voir Compléments sur la platitude, définition  
\ref{flat-definition-h-covering}). Alors  $\alpha$  est un  $r$-cycle 
relatif si et seulement si chaque changement de base  $g_i^*\alpha$  est un  $r$-cycle relatif. 
\end{lemma} 

\begin{proof} 
Si  $\alpha$  est un  $r$-cycle relatif, alors chaque changement de base $g_i^*\alpha$  est
 un  $r$-cycle relatif d'après le lemme  \ref{relative-cycles-lemma-relative-cycle-functoriality}. Supposons
 que chaque  $g_i^*\alpha$  soit un  $r$-cycle relatif.
 Soit $g : S' \to S$  un morphisme où  $S'$  est le spectre d'un anneau de 
valuation discrète. Après avoir remplacé  $S$  par  $S'$,  $X$  par $X' = X \times_S S'$, et  
$\alpha$  par  $\alpha' = g^*\alpha$, puis utilisé le fait que le changement de base d'un
 recouvrement h est un recouvrement h (Compléments sur la platitude, lemme \ref{flat-lemma-h}), 
nous revenons au problème étudié dans le paragraphe suivant. 

\medskip\noindent
 Supposons que  $S$  soit le spectre d'un anneau de valuation discrète avec le point 
fermé  $0$  et le point générique  $\eta$. Nous devons montrer que  
$sp_{X/S}(\alpha_\eta) = \alpha_0$. Puisqu'un recouvrement h est un recouvrement
 V (par définition), il existe un  $i$  et une spécialisation  $s' \leadsto s$
  des points de  $S_i$  avec  $g_i(s') = \eta$  et  $g_i(s) = 0$, voir
 Topologies, lemme \ref{topologies-lemma-refine-qcqs-V}. D'après
 Propriétés, lemme  \ref{properties-lemma-locally-Noetherian-specialization-dvr} , nous 
pouvons trouver un morphisme $h : S' \to S_i$ du 
spectre  $S'$  d'un anneau de valuation discrète qui envoie
 le point générique  $\eta'$  sur  $s'$  et envoie
 le point fermé  $0'$  sur  $s$. Notons  $\alpha' = h^*g_i^*\alpha$.
 Par hypothèse, nous avons  $sp_{X'/S'}(\alpha'_{\eta'}) = \alpha'_{0'}$. 
Puisque  $g = g_i \circ h : S' \to S$  est le morphisme de schémas 
induit par une extension d'anneaux de valuation discrète, nous en déduisons que 
$sp_{X/S}$  et  $sp_{X'/S'}$  sont compatibles avec les applications de changement 
de base sur les fibres, voir le lemme  \ref{relative-cycles-lemma-specialization-extension}. 
Nous concluons que $sp_{X/S}(\alpha_\eta) = \alpha_0$  car 
l'application de changement de base  $Z_r(X_0) \to Z_r(X'_{0'})$  est injective 
comme on l'a vu à la section  \ref{relative-cycles-section-relative-fields}. 
\end{proof} 

\begin{lemma} 
\label{relative-cycles-lemma-families-specialization-fppf-descent} 
Soit  $f : X \to S$  un morphisme de schémas. Supposons que  $S$  soit localement noethérien et
 que  $f$  soit localement de type fini. Soient  $r, e \geq 0$  des entiers. 
Soit  $\alpha$  une famille de  $r$-cycles sur les fibres de $X/S$. 
Soit  $\{f_i : X_i \to X\}$  une famille de morphismes
 plats, localement de type fini, et de dimension relative  $e$, qui 
est conjointement surjective. Alors  $\alpha$  est un  $r$-cycle relatif si et seulement si 
chaque cycle $f_i^*\alpha$ obtenu par image inverse plate  est un  $r$-cycle relatif. 
\end{lemma} 

\begin{proof} 
Si  $\alpha$  est un  $r$-cycle relatif, alors chaque  $f_i^*\alpha$ est à 
son tour un $r$-cycle relatif d'après le lemme  \ref{relative-cycles-lemma-relative-cycle-functoriality}. Supposons
 que chaque  $f_i^*\alpha$  soit un  $r$-cycle relatif.
 Soit  $g : S' \to S$  un morphisme où  $S'$  est le spectre d'un anneau de 
valuation discrète. Après avoir remplacé  $S$  par  $S'$,  $X$  par  $X' = X \times_S S'$, et 
$\alpha$  par  $\alpha' = g^*\alpha$
 , nous réduisons au problème étudié dans le paragraphe suivant. 

\medskip\noindent
 Supposons que  $S$  soit le spectre d'un anneau de valuation discrète avec point 
fermé  $0$  et point générique  $\eta$. Nous devons montrer que 
$sp_{X/S}(\alpha_\eta) = \alpha_0$. Notons  $f_{i, 0} : X_{i, 0} \to X_0$
  le changement de base de  $f_i$  au point fermé de  $S$. De même 
pour  $f_{i, \eta}$.
 Observons que 
$$
f_{i, 0}^*sp_{X/S}(\alpha_\eta) =
sp_{X_i/S}(f_{i, \eta}^*\alpha_\eta) = f_{i, 0}^*\alpha_0
$$
 En effet, la première égalité résulte du 
lemme  \ref{relative-cycles-lemma-specialization-flat-pullback}  et
 la seconde par hypothèse. Puisque la famille de morphismes  
$f_{i, 0}^* : Z_r(X_0) \to Z_r(X_{i, 0})$  est 
conjointement injective (puisque  $f_{i, 0}$  est conjointement
 surjective), nous concluons ce que nous voulons. 
\end{proof} 

\begin{lemma} 
\label{relative-cycles-lemma-check-after-closed} 
Soit  $S$  un schéma localement noethérien. Soit  $i : X \to Y$  une immersion fermée 
de schémas localement de type fini sur  $S$. Soit  $r \geq 0$. 
Soit  $\alpha$  une famille de  $r$-cycles sur les fibres de $X/S$. 
Alors  $\alpha$  est un  $r$-cycle relatif sur $X/S$  si et seulement si  
$i_*\alpha$  est un  $r$-cycle relatif sur $Y/S$. 
\end{lemma} 

\begin{proof} 
Puisque le changement de base commute avec  $i_*$  (lemme \ref{relative-cycles-lemma-compatibilities}),
 il suffit de montrer ceci : si  $S$  est le spectre d'un anneau de 
valuation discrète avec point générique  $\eta$  et point fermé  $0$, 
alors  $sp_{X/S}(\alpha_\eta) = \alpha_0$  si et seulement si  
$sp_{Y/S}(i_{\eta, *}\alpha_\eta) = i_{0, *}\alpha_0$. 
Cela résulte de l'injectivité de  $i_{0, *} : Z_r(X_0) \to Z_r(Y_0)$
 et de l'égalité  $i_{0, *}sp_{X/S}(\alpha_\eta) =
sp_{Y/S}(i_{\eta, *}\alpha_\eta)$  d'après
 le lemme  \ref{relative-cycles-lemma-specialization-proper-pushforward}. 
\end{proof} 

\noindent
 Le lemme suivant sera renforcé dans le 
lemme  \ref{relative-cycles-lemma-relative-cycles-equal}. 

\begin{lemma} 
\label{relative-cycles-lemma-uniqueness-extension} 
Soit  $f : X \to S$  un morphisme de schémas. Supposons que  $S$  soit localement noethérien
 et  $f$  localement de type fini. Soit  $r \geq 0$. Soient  $\alpha$  et  $\beta$
  des  $r$-cycles relatifs sur $X/S$. Les énoncés suivants sont équivalents 
\begin{enumerate} 
\item $\alpha = \beta$, et 
\item $\alpha_\eta = \beta_\eta$  pour tout point générique  $\eta \in S$
  d'une composante irréductible de  $S$. 
\end{enumerate} 
\end{lemma} 

\begin{proof}
 L'implication (1)  $\Rightarrow$  (2) est immédiate. 
Supposons (2). Pour chaque $s \in S$  nous pouvons trouver un point $\eta$ comme en (2)
 qui se spécialise en  $s$. 
D'après Propriétés, lemme  \ref{properties-lemma-locally-Noetherian-specialization-dvr} , nous 
pouvons trouver un morphisme  $g : S' \to S$  depuis le
 spectre  $S'$  d'un anneau de valuation discrète qui envoie
 le point générique  $\eta'$  sur  $\eta$  et qui 
envoie le point fermé  $0$  sur  $s$. Alors  $\alpha_s$  et  $\beta_s$
 sont des éléments de  $Z_r(X_s)$  dont les images par changement de base 
sont égales dans  $Z_r(X_{0'})$, à savoir  $sp_{X_{S'}/S'}(\alpha_{\eta'})$
  où  $\alpha_{\eta'}$  est le changement de base de  $\alpha_\eta$. 
Comme l'application de changement de base $Z_r(X_s) \to Z_r(X_{0'})$  est injective 
comme on l'a vu à la section  \ref{relative-cycles-section-relative-fields} , nous
 en concluons  $\alpha_s = \beta_s$. 
\end{proof} 

\begin{lemma} 
\label{relative-cycles-lemma-family-associated-module-specialization} 
Dans la situation de l'exemple  \ref{relative-cycles-example-family-associated-module} , supposons
 que  $S$  soit localement noethérien et que  
$\mathcal{F}$  soit plat sur  $S$  en dimensions $\geq r$
  (Compléments sur la platitude, définition  \ref{flat-definition-flat-dimension-n}). 
Alors  $[\mathcal{F}/X/S]_r$ est un  $r$-cycle relatif sur  $X/S$. 
\end{lemma} 

\begin{proof} 
D'après Compléments sur la platitude, lemme \ref{flat-lemma-pre-flat-dimension-n} 
, l'hypothèse sur  $\mathcal{F}$ est conservée par tout changement de base. 
De plus, la formation de  $[\mathcal{F}/X/S]_r$ est compatible avec tout changement
 de base d'après le lemme  \ref{relative-cycles-lemma-family-associated-module}. 
Puisque la condition d'être compatible avec les spécialisations est
 vérifiée après un changement de base par le spectre d'un anneau de valuation 
discrète, cela nous ramène au cas où  $S$  est le spectre d'un anneau de valuation. 
Dans ce cas, l'ensemble  
$U = \{x \in X \mid \mathcal{F}\text{ est plat en }x\text{ sur }S\}$
  est ouvert dans  $X$  
d'après Compléments sur la platitude, lemme \ref{flat-lemma-finite-type-flat}. 
Puisque le complément de  $U$  dans $X$  a des fibres de dimension  $< r$  sur  
$S$  par hypothèse, nous voyons que la restriction à l'ouvert  
$U \subset X$ induit un isomorphisme sur les groupes de  $r$-cycles 
sur les fibres après tout changement de base, compatible avec les applications de
 spécialisation et avec la formation du cycle relatif associé à  $\mathcal{F}$.
 Ainsi, il suffit de montrer la compatibilité avec les
 spécialisations pour  $[\mathcal{F}|_U / U /S]_r$.
 Puisque $\mathcal{F}|_U$  est plat sur  $S$, cela découle du 
lemme  \ref{relative-cycles-lemma-specialization-module}  et des définitions. 
\end{proof} 

\begin{lemma} 
\label{relative-cycles-lemma-family-associated-closed-specialization} 
Dans la situation de l'exemple  \ref{relative-cycles-example-family-associated-closed} , supposons
 que $S$ soit localement noethérien et que  $Z$  soit plat sur  $S$  en dimensions 
$\geq r$. Alors  $[Z/X/S]_r$  est un  $r$-cycle relatif sur $X/S$. 
\end{lemma} 

\begin{proof}
 L'hypothèse signifie que  $\mathcal{O}_Z$  est plat sur  $S$  
en dimensions $\geq r$. Ainsi, en appliquant 
le lemme  \ref{relative-cycles-lemma-family-associated-module-specialization}  
avec  $\mathcal{F} = (Z \to X)_*\mathcal{O}_Z$ , nous en concluons. 
\end{proof} 

\noindent
 Soit  $S$  un schéma localement noethérien. Soit  $f : X \to S$  un morphisme
 de type fini. Soit  $r \geq 0$. Notons  $Hilb(X/S, r)$
  l'ensemble des sous-schémas fermés  $Z \subset X$  tels que  $Z \to S$  soit 
plat et de dimension relative $\leq r$. D'après
 le lemme  \ref{relative-cycles-lemma-family-associated-closed-specialization} , pour chaque  
$Z \in Hilb(X/S, r)$ , nous avons un élément  $[Z/X/S]_r \in z(X/S, r)$.
 Ainsi, nous obtenons un homomorphisme de groupe 
\begin{equation} 
\label{relative-cycles-equation-cycle-classes} 
\text{groupe abélien libre engendré par }Hilb(X/S, r) \longrightarrow z(X/S, r) 
\end{equation} 
qui envoie  $\sum n_i[Z_i]$  sur  $\sum n_i[Z_i/X/S]_r$.
 Une propriété essentielle des $r$-cycles relatifs est qu'ils sont localement 
(sur  $X$  et $S$  dans des topologies appropriées) dans l'image de cette application. 

\begin{lemma} 
\label{relative-cycles-lemma-get-cycles} 
Soit  $f : X \to S$  un morphisme de type fini entre schémas, où $S$ est noethérien. 
Soit  $r \geq 0$. Soit  $\alpha$  un  $r$-cycle relatif sur $X/S$. Alors il existe un 
morphisme propre, complètement décomposé 
(Compléments sur les morphismes, définition  \ref{more-morphisms-definition-cd-morphism})
  $g : S' \to S$  tel que  $g^*\alpha$ soit dans l'image de 
(\ref{relative-cycles-equation-cycle-classes}). 
\end{lemma} 

\begin{proof} 
Par induction noethérienne, on peut supposer que le résultat est vrai pour l'image inverse de 
$\alpha$  par toute immersion fermée  $g : S' \to S$  qui n'est pas un isomorphisme. 

\medskip\noindent
 Soit  $S_1 \subset S$  une composante irréductible (vue comme un sous-schéma fermé
 intègre). Soit  $S_2 \subset S$  l'adhérence du complément de  $S'$
  (vue comme un sous-schéma fermé réduit). Si  $S_2 \not = \emptyset$, alors
 le résultat est vrai pour l'image inverse de $\alpha$  par  $S_1 \to S$  et  $S_2 \to S$. 
Si  $g_1 : S'_1 \to S_1$ et  $g_2 : S'_2 \to S_2$
  sont les morphismes propres complètement décomposés correspondants, 
alors  $S' = S'_1 \amalg S'_2 \to S$
  est un morphisme propre complètement décomposé et nous 
voyons que le résultat est valable pour $S$\footnote{Plus précisément, tout 
sous-schéma fermé de $S'_1 \times_S X$ plat et de dimension relative  
$\leq r$  sur  $S'_1$  peut être considéré comme un sous-schéma fermé de  $S' \times_S X$
  plat et de dimension relative  $\leq r$  sur  $S'$.} 
. Ainsi, nous pouvons supposer que  $S' \to S$  est bijectif
 et nous nous ramenons au cas décrit dans le paragraphe suivant. 

\medskip\noindent
 Supposons que  $S$  soit intègre. Soit  $\eta \in S$  le point générique 
et soit  $K = \kappa(\eta)$  le corps des fonctions de  $S$. 
Alors  $\alpha_\eta$  est un  $r$-cycle sur  $X_K$. 
Écrivons  $\alpha_\eta = \sum n_i[Y_i]$.
 En prenant l'adhérence de  $Y_i$ , nous obtenons des sous-schémas fermés intègres 
$Z_i \subset X$  dont le changement de base à  $\eta$  est  $Y_i$.
 Par platitude générique (par exemple Morphismes,
 proposition  \ref{morphisms-proposition-generic-flatness}), nous
 voyons que  $Z_i$  est plat sur un ouvert non vide  $U$  de  $S$  pour chaque  $i$. 
En appliquant Compléments sur la platitude, lemme \ref{flat-lemma-flat-after-blowing-up} , nous pouvons 
trouver un éclatement  $U$-admissible  $g : S' \to S$
 tel que la transformée stricte  $Z'_i \subset X_{S'}$
  de  $Z_i$  soit plate sur $S'$. Alors  $\beta = \sum n_i[Z'_i/X_{S'}/S']_r$
  est dans l'image de (\ref{relative-cycles-equation-cycle-classes}) et $\beta = g^*\alpha$
  par le lemme \ref{relative-cycles-lemma-uniqueness-extension}. 

\medskip\noindent
 Cependant, cela ne termine pas la preuve car  $S' \to S$  peut ne pas 
être complètement décomposé. Cela se corrige facilement : en notant  $T \subset S$
  le complément de  $U$  (vu comme un sous-schéma fermé), par induction 
noethérienne nous pouvons trouver un morphisme propre complètement décomposé 
$T' \to T$  tel que  $(T' \to S)^*\alpha$
  soit dans l'image de (\ref{relative-cycles-equation-cycle-classes}). Alors 
$S' \amalg T' \to S$  convient. 
\end{proof} 

\begin{lemma} 
\label{relative-cycles-lemma-get-cycles-dvr} 
Soit  $f : X \to S$  un morphisme de schémas de type fini avec  $S$
  le spectre d'un anneau de valuation discrète. Soit  $r \geq 0$.
 Alors (\ref{relative-cycles-equation-cycle-classes}) est surjectif. 
\end{lemma} 

\begin{proof}
 Cela découle bien sûr du lemme \ref{relative-cycles-lemma-get-cycles}, mais nous pouvons 
également le voir directement comme suit. Supposons que  $\alpha$  soit un  $r$-cycle relatif
 sur  $X/S$. Écrivons  $\alpha_\eta = \sum n_i[Z_i]$  (la somme est finie). Notons 
$\overline{Z}_i \subset X$  l'adhérence de  $Z_i$  comme 
dans la section \ref{relative-cycles-section-specialization}. Alors 
$\alpha = \sum n_i[\overline{Z}_i/X/S]$. 
\end{proof} 

\begin{lemma} 
\label{relative-cycles-lemma-relative-cycle-smooth} 
Soit  $f : X \to S$  un morphisme de schémas. Soit  $r \geq 0$. Supposons que $S$
 soit localement noethérien et que $f$ soit lisse de dimension relative $r$. Soit  
$\alpha \in z(X/S, r)$. Alors le support de  $\alpha$  est ouvert et fermé dans  $X$
  (voir la preuve pour un résultat plus précis). 
\end{lemma} 

\begin{proof} 
Soit  $x \in X$  d'image  $s \in S$. Puisque  $f$  est lisse, il existe une
 unique composante irréductible  $Z(x)$  de  $X_s$  qui contient  $x$. 
Alors  $\dim(Z(x)) = r$. 
Soit  $n_x$  le coefficient de $Z(x)$  dans le cycle  $\alpha_s$.
 Nous montrerons que la fonction $x \mapsto n_x$  est localement constante sur  $X$. 

\medskip\noindent
 Soit  $g : S' \to S$  un morphisme de schémas localement noethériens. 
Soit $X'$  le changement de base de  $X$  et soit  $\alpha' = g^*\alpha$
  le changement de base de  $\alpha$. Soit  $x' \in X'$  se projetant sur $s' \in S'$,  
$x \in X$ et  $s \in S$. Nous affirmons que $n_{x'} = n_x$. En effet, puisque  $Z(x)$
  est lisse sur  $\kappa(s)$ , nous voyons que  
$Z(x) \times_{\Spec(\kappa(s))} \Spec(\kappa(s'))$
  est réduit. Puisque  $Z(x')$  est une composante irréductible
 de ce schéma, nous voyons que le coefficient  $n_{x'}$  de  $Z(x')$
  dans  $\alpha'_{s'}$  est le même que le coefficient  $n_x$  de  $Z(x)$
  dans  $\alpha_s$  par la définition du changement de base dans 
la section \ref{relative-cycles-section-relative-fields}, 
ce qui prouve l'affirmation. 

\medskip\noindent
 Puisque  $X$  est localement noethérien, pour montrer que  $x \mapsto n_x$
  est localement constant, il suffit de montrer: si  $x' \leadsto x$
  est une spécialisation dans  $X$, alors  $n_{x'} = n_x$.
 Choisissons un morphisme  $S' \to X$  où $S'$  est le spectre d'un anneau de 
valuation discrète envoyant le point générique  $\eta$  vers  $x'$  et le point fermé  
$0$  vers $x$. Voir Propriétés, lemme 
\ref{properties-lemma-locally-Noetherian-specialization-dvr}. Ensuite,
 le changement de base $X' \to S'$  de  $f$  par  $S' \to S$
  a une section  $\sigma : S' \to X'$  telle que $\sigma(\eta) \leadsto \sigma(0)$
  est une spécialisation de points de  $X'$  envoyés vers $x' \leadsto x$  dans  $X$.
 Ainsi, nous nous réduisons à l'affirmation du paragraphe suivant. 

\medskip\noindent
 Soit  $S$  le spectre d'un anneau de valuation discrète dont le 
point générique est  $\eta$  et le point fermé est  $0$, et soit donnée une section 
$\sigma : S  \to X$. Nous affirmons que $n_{\sigma(\eta)} = n_{\sigma(0)}$.
 D'après la discussion dans Compléments sur les morphismes, section 
\ref{more-morphisms-section-connected-components}, et en particulier Compléments
 sur les morphismes, lemme \ref{more-morphisms-lemma-connected-along-section-open} , après
 avoir remplacé  $X$  par un sous-schéma ouvert, nous pouvons supposer que 
les fibres de  $X \to S$  sont connexes. Comme ces fibres sont
 lisses, elles sont irréductibles. Alors nous voyons que  $\alpha_\eta = n[X_\eta]$
  avec $n = n_{\sigma(\eta)}$  et la relation  $sp_{X/S}(\alpha_\eta) = \alpha_0$
  implique  $\alpha_0 = n[X_0]$, c'est-à-dire $n_{\sigma(0)} = n$, comme voulu. 
\end{proof} 

\begin{lemma} 
\label{relative-cycles-lemma-relative-cycles-equal} 
Soit  $f : X \to S$  un morphisme de schémas. Supposons que  $S$  soit localement noethérien
 et  $f$  localement de type fini. Soit  $r \geq 0$ et soient  
$\alpha, \beta \in z(X/S, r)$. L'ensemble  $E = \{s \in S : \alpha_s = \beta_s\}$
  est fermé dans  $S$. 
\end{lemma} 

\begin{proof} 
La question est locale sur  $S$, donc nous pouvons supposer que $S$  est affine. 
Soit  $X = \bigcup U_i$  un recouvrement ouvert affine. Soit 
$E_i = \{s \in S : \alpha_s|_{U_{i, s}} = \beta_s|_{U_{i, s}}\}$. 
Alors  $E = \bigcap E_i$. Par conséquent, il suffit de prouver le lemme pour  
$U_i \to S$  et la restriction de  $\alpha$  et  $\beta$  à $U_i$.
 Cela nous ramène au cas discuté dans le paragraphe suivant. 

\medskip\noindent
 Supposons que  $X$  et  $S$  soient quasi-compacts. 
Posons $\gamma = \alpha - \beta$. Alors  $E = \{s \in S : \gamma_s = 0\}$. 
D'après le lemme \ref{relative-cycles-lemma-family-associated-closed-specialization}, il 
existe une famille finie de morphismes propres conjointement surjective  
$\{g_i : S_i \to S\}$  telle que  $g_i^*\gamma$  soit dans l'image
 de (\ref{relative-cycles-equation-cycle-classes}). On observe que  $E_i = g_i^{-1}(E)$ est
 l'ensemble des points  $t \in S_i$  tels que  $(g_i^*\gamma)_t = 0$. 
Si  $E_i$  est fermé pour tous les  $i$, alors  $E = \bigcup g_i(E_i)$
  est également fermé. Cela nous ramène au cas discuté dans le paragraphe
 suivant. 

\medskip\noindent
 Supposons que  $X$  et  $S$  soient quasi-compacts et  $\gamma = \sum n_i[Z_i/X/S]_r$
 pour un nombre fini de sous-schémas fermés  $Z_i \subset X$
  plats et de dimension relative  $\leq r$  sur  $S$. 
Posons  $X' = \bigcup Z_i$  (union schématique). 
Alors  $i : X' \to X$  est une immersion fermée et  $X'$
  a une dimension relative  $\leq r$  sur  $S$. De plus  
$\gamma = i_*\gamma'$  où  $\gamma' = \sum n_i[Z_i/X'/S]_r$.
 Puisque  $E = E' = \{s \in S : \gamma'_s = 0\}$
 , nous nous réduisons au cas discuté dans le paragraphe suivant. 

\medskip\noindent
 Supposons que  $X$  ait une dimension relative  $\leq r$  sur $S$. 
Soit  $s \in S$, avec  $s \not \in E$. Nous allons montrer qu'il existe 
un voisinage ouvert  $V \subset S$  de  $s$  tel que  $E \cap V$  soit vide.
 L'hypothèse  $s \not \in E$  signifie qu'il existe un sous-schéma fermé
 intègre $Z \subset X_s$  de dimension  $r$  tel que le coefficient  
$n$  de $[Z]$  dans  $\gamma_s$  soit non nul. Soit  $x \in Z$  le point
 générique. Puisque  $\dim(Z) = r$ , nous voyons que  $x$  est un point 
générique d'une composante irréductible (à savoir  $Z$) de  $X_s$.
 Ainsi, après avoir remplacé  $X$  par un voisinage ouvert de  $x$,
 nous pouvons supposer que  $Z$  est la seule composante irréductible de  $X_s$.
 En particulier, nous avons  $\gamma_s = n[Z]$. 

\medskip\noindent
 À ce stade, nous appliquons Compléments sur les morphismes, lemme 
\ref{more-morphisms-lemma-local-structure-finite-type}, et 
nous obtenons un diagramme 
$$
\xymatrix{
X \ar[dd] & X' \ar[l]^g \ar[d]^\pi & x \ar@{|->}[dd] &
x' \ar@{|->}[l]  \ar@{|->}[d] \\
& Y \ar[d]^h & & y \ar@{|->}[d] \\
S \ar@{=}[r] & S & s & s \ar@{=}[l]
}
$$
 avec toutes les propriétés qui y sont énumérées. Soit  $\gamma' = g^*\gamma$
  l'image inverse plate. Notons que  $E \subset E' = \{s \in S: \gamma'_s = 0\}$
  et que  $s \not \in E'$  parce que le coefficient de $Z'$  dans  $\gamma'_s$
  est non nul, où  $Z' \subset X'_s$  est l'adhérence de $x'$.
 De même, posons  $\gamma'' = \pi_*\gamma'$. Alors nous avons  
$E' \subset E'' = \{s \in S: \gamma''_s = 0\}$  et $s \not \in E''$
  parce que le coefficient de  $Z''$  dans  $\gamma''_s$  est non nul, où  
$Z'' \subset Y_s$  est l'adhérence de  $y$. D'après
 le lemme \ref{relative-cycles-lemma-relative-cycle-smooth}  et l'ouverture de  $Y \to S$
 , nous voyons qu'un voisinage ouvert de $s$  est disjoint de  $E''$
  et la démonstration est complète. 
\end{proof} 

\begin{lemma} 
\label{relative-cycles-lemma-descend-through-limit} 
Soit  $S = \lim_{i \in I} S_i$  la limite d'un système inverse dirigé de schémas
 noethériens avec des morphismes de transition affines. 
Soit  $0 \in I$  et soit  $X_0 \to S_0$ un morphisme de schémas de type fini. 
Pour  $i \geq 0$ , posons  $X_i = S_i \times_{S_0} X_0$  et posons  
$X = S \times_{S_0} X_0$. Si  $S$  est également noethérien, alors 
$$
z(X/S, r) = \colim_{i \geq 0} z(X_i/S_i, r)
$$
 où les applications de transition sont données par changement de base des 
$r$-cycles relatifs. 
\end{lemma} 

\begin{proof} 
Supposons que  $i \geq 0$  et  $\alpha_i, \beta_i \in z(X_i/S_i, r)$
  aient la même image dans $z(X/S, r)$. Alors l'image de  $S \to S_i$
  est contenue dans le sous-ensemble fermé $E \subset S_i$  du
 lemme \ref{relative-cycles-lemma-relative-cycles-equal}. Par conséquent,
 pour un certain  $j \geq i$ , l'image du morphisme  $S_j \to S_i$
  est contenue dans  $E$ ; voir Limites, lemme 
\ref{limits-lemma-limit-contained-in-constructible}. Il 
s'ensuit que les changements de base de  
$\alpha_i$  et $\beta_i$  à  $S_j$  coïncident. Ainsi, l'application est injective. 

\medskip\noindent
 Soit  $\alpha \in z(X/S, r)$. En appliquant le lemme \ref{relative-cycles-lemma-get-cycles}, on 
obtient un morphisme propre complètement décomposé  $g : S' \to S$  tel que  $g^*\alpha$
  soit dans l'image de (\ref{relative-cycles-equation-cycle-classes}). 
Posons  $X' = S' \times_S X$. Nous écrivons  $g^*\alpha = \sum n_a [Z_a/X'/S']_r$
  pour certains sous-schémas fermés  $Z_a \subset X'$  plats et 
de dimension relative  $\leq r$  sur $S'$. 

\medskip\noindent
 Nous appliquons maintenant les résultats
 de Limites, section \ref{limits-section-descending-relative} et suivantes. Nous 
pouvons trouver un  $i \geq 0$  tel qu'il existe 
\begin{enumerate} 
\item un morphisme propre complètement décomposé  $g_i : S'_i \to S_i$
  dont le changement de base à  $S$  est  $g : S' \to S$, 
\item en posant  $X'_i = S'_i \times_{S_i} X_i$
 , des sous-schémas fermés  $Z_{ai} \subset X'_i$  plats et de
 dimension relative  $\leq r$  sur  $S'_i$  dont le changement de base à $S'$
  soit  $Z_a$. 
\end{enumerate} 
Pour ce faire, on utilise Limites, lemmes 
\ref{limits-lemma-descend-finite-presentation}, 
\ref{limits-lemma-descend-closed-immersion-finite-presentation}, 
\ref{limits-lemma-descend-flat-finite-presentation}, 
\ref{limits-lemma-eventually-proper} et 
\ref{limits-lemma-limit-dimension}, ainsi que
 Compléments sur les morphismes, lemme 
\ref{more-morphisms-lemma-descend-cd}. 
Considérons  
$\alpha'_i = \sum n_a [Z_{ai}/X'_i/S'_i]_r \in z(X'_i/S'_i, r)$.
 L'image de  $\alpha'_i$  dans  $z(X'/S', r)$ coïncide par construction avec le changement de base  
$g^*\alpha$. 

\medskip\noindent
 Posons  $S''_i = S'_i \times_{S_i} S'_i$  et  $X''_i = S''_i \times_{S_i} X_i$
  et posons  $S'' = S' \times_S S'$  et  $X'' = S'' \times_S X$.
 Nous notons  $\text{pr}_1, \text{pr}_2 : S'' \to S'$  et  
$\text{pr}_1, \text{pr}_2 : S''_i \to S'_i$  les projections. Les
 deux changements de base $\text{pr}_1^*\alpha'_i$  et  $\text{pr}_1^*\alpha'_i$
  ont la même image dans  $z(X''/S'', r)$  parce que  
$\text{pr}_1^*g^*\alpha = \text{pr}_1^*g^*\alpha$.
 Ainsi, quitte à remplacer  $i$  par un indice plus grand, nous pouvons supposer que  
$\text{pr}_1^*\alpha'_i = \text{pr}_1^*\alpha'_i$
  d'après le premier paragraphe de la démonstration. 
Par le lemme \ref{relative-cycles-lemma-descend-family} , 
nous obtenons une famille unique  $\alpha_i$
  de  $r$-cycles sur les fibres de  $X_i/S_i$
  avec  $g_i^*\alpha_i = \alpha'_i$  (cela utilise le fait que $S'_i \to S_i$
  est complètement décomposé). Par
 le lemme \ref{relative-cycles-lemma-relative-cycles-h-descent}, nous
 voyons que  $\alpha_i \in z(X_i/S_i, r)$.
 L'unicité dans le lemme \ref{relative-cycles-lemma-descend-family} implique que 
l'image de $\alpha_i$  dans  $z(X/S, r)$  est  $\alpha$  et la démonstration est complète. 
\end{proof} 

\begin{lemma} 
\label{relative-cycles-lemma-thickening} 
Soit  $S$  un schéma localement noethérien. Soit  $i : X \to X'$  un épaississement
 de schémas localement de type fini sur  $S$. Soit $r \geq 0$. 
Alors  $i_* : z(X/S, r) \to z(X'/S, r)$  est une bijection. 
\end{lemma} 

\begin{proof} 
Puisque  $i_s : X_s \to X'_s$  est un épaississement, il est clair que  $i_*$  induit 
une bijection entre les
 familles de  $r$-cycles sur les fibres de $X/S$  et 
les familles de  $r$-cycles sur les fibres de $X'/S$. 
Étant donnée une famille  $\alpha$  de $r$-cycles sur les fibres de  $X/S$,
 on a $\alpha \in z(X/S, r) \Leftrightarrow i_*\alpha \in z(X'/S, r)$
  par le lemme \ref{relative-cycles-lemma-check-after-closed}. Le lemme s'ensuit. 
\end{proof} 

\begin{lemma} 
\label{relative-cycles-lemma-extend-to-larger} 
Soit  $S$  un schéma localement noethérien. Soit  $X$  un schéma localement
 de type fini sur  $S$. Soit  $r \geq 0$. Soit  $U \subset X$  un ouvert
 tel que  $X \setminus U$  ait une dimension relative  $< r$  sur $S$, c'est-à-dire  
$\dim(X_s \setminus U_s) < r$  pour tout  $s \in S$. Alors
 la restriction définit une bijection  $z(X/S, r) \to z(U/S, r)$. 
\end{lemma} 

\begin{proof} Étant
 donné que  $Z_r(X_s) \to Z_r(U_s)$  est une bijection par l'hypothèse de dimension, 
nous voyons que la restriction induit une bijection entre les
 familles de $r$-cycles sur les fibres de  $X/S$  et 
les familles de $r$-cycles sur les fibres de  $U/S$.
 Ces applications de restriction  $Z_r(X_s) \to Z_r(U_s)$  sont
 compatibles avec le changement de base et avec les spécialisations ; voir les
 lemmes \ref{relative-cycles-lemma-compatibilities} et \ref{relative-cycles-lemma-specialization-flat-pullback}. Le 
lemme s'en déduit facilement ; la démonstration est omise. 
\end{proof} 

\begin{lemma} 
\label{relative-cycles-lemma-seminormalize-base} 
Soit  $g : S' \to S$  un homéomorphisme universel de schémas localement noethériens 
qui induit des isomorphismes des corps résiduels. Soit  $f : X \to S$  localement de 
type fini. Posons  $X' = S' \times_S X$. Soit  $r \geq 0$. Alors le changement de base par $g$
  détermine une bijection  $z(X/S, r) \to z(X'/S', r)$. 
\end{lemma} 

\begin{proof} 
D'après le lemme \ref{relative-cycles-lemma-pullback-universally-bijective}, nous avons une bijection
 entre le groupe des familles de  $r$-cycles sur les fibres de  $X/S$  et
 le groupe des familles de  $r$-cycles sur les fibres de  $X'/S'$. 
Soit  $\alpha$ une famille de  $r$-cycles sur les fibres de  $X/S$  et  
$\alpha' = g^*\alpha$ le changement de base. 
Si  $R$  est un anneau de valuation discrète, alors tout morphisme  
$h : \Spec(R) \to S$  se factorise comme  $g \circ h'$  pour un 
certain morphisme unique  $h' : \Spec(R) \to S'$.
 En effet, le morphisme  $S' \times_S \Spec(R) \to \Spec(R)$  est
 un homéomorphisme universel induisant des bijections sur les corps résiduels,
 et possède donc une section (par exemple parce que  $R$  est un anneau 
semi-normal ; voir Morphismes, section \ref{morphisms-section-seminormalization}). 
Ainsi, la condition que $\alpha$  soit compatible avec les
 spécialisations (c'est-à-dire qu'il s'agisse d'un  $r$-cycle 
relatif) équivaut à la condition que  
$\alpha'$  soit compatible avec les spécialisations. 
\end{proof} 






\section{Cycles relatifs équidimensionnels} 
\label{relative-cycles-section-equidimensional} 

\noindent
 Voici la définition. 

\begin{definition} 
\label{relative-cycles-definition-equidimensional} 
Soit  $f : X \to S$  un morphisme de schémas. Supposons que  $S$  soit localement noethérien et
 que  $f$  soit localement de type fini. Soit $r \geq 0$  un entier. Nous disons
 qu'un  $r$-cycle relatif  $\alpha$ sur  $X/S$  est {\it équidimensionnel} si le support 
de  $\alpha$  (remarque \ref{relative-cycles-remark-supports-family}) est
 contenu dans un sous-ensemble fermé  $W \subset X$  dont la dimension
 relative sur  $S$  est  $\leq r$.
 Le groupe de tous les  $r$-cycles relatifs équidimensionnels sur $X/S$  
est noté  $z_{equi}(X/S, r)$. 
\end{definition} 

\begin{example} 
\label{relative-cycles-example-not-equidimensional} 
\begin{reference} 
\cite[exemple 3.1.9]{SV} 
\end{reference}
 Il existe des  $r$-cycles relatifs qui ne sont pas équidimensionnels. Plus 
précisément, soit  $k$  un corps et soit  $X = \Spec(k[x, y, t])$
  sur  $S = \Spec(k[x, y])$. Soit $s$  un point de  $S$  et 
notons  $a, b \in \kappa(s)$  les images de $x$  et  $y$.
 Considérons la famille  $\alpha$  de $0$-cycles sur  $X/S$  définie par 
\begin{enumerate} 
\item $\alpha_s = 0$  si  $b = 0$  et sinon 
\item $\alpha_s = [p] - [q]$, où  $p$ (respectivement\ $q$) est le point $\kappa(s)$-rationnel
 de  $\Spec(\kappa(s)[t])$  défini par  $t = a/b$ 
(respectivement\ $t = (a + b^2)/b$). 
\end{enumerate} 
Nous laissons au lecteur le soin de montrer que cela est compatible avec les spécialisations
 ; l'idée est que  $a/b$  et  $(a + b^2)/b = a/b + b$  se spécialisent en un même point 
dans  $\mathbf{P}^1$  sur le corps résiduel de toute valuation  $v$  sur $\kappa(s)$
  avec  $v(b) > 0$. D'autre part, l'adhérence du support de $\alpha$
  contient toute la fibre au-dessus de  $(0, 0)$. 
\end{example} 

\begin{lemma} 
\label{relative-cycles-lemma-equidimensional-functoriality} 
Soit  $f : X \to S$  un morphisme de schémas. Supposons que  $S$  soit localement noethérien et
 que  $f$  soit localement de type fini. Soit $r \geq 0$  un entier. Soit  
$\alpha$  un  $r$-cycle relatif sur $X/S$. Si  $\alpha$  est équidimensionnel, 
alors toute restriction, changement de base ou image inverse plate de  $\alpha$  est
 équidimensionnel. 
\end{lemma} 

\begin{proof} La 
démonstration est omise. 
\end{proof} 

\begin{lemma} 
\label{relative-cycles-lemma-check-equidimensional} 
Soit  $f : X \to S$  un morphisme de schémas. Supposons  $S$  localement noethérien
 et  $f$  localement de type fini. Soit  $r \geq 0$  un entier. Soit  $\alpha$
  un  $r$-cycle relatif sur  $X/S$. Alors, pour vérifier que  $\alpha$  est équidimensionnel, nous pouvons
 travailler localement pour la topologie de Zariski sur  $X$  et  $S$. 
\end{lemma} 

\begin{proof}
 La condition que  $\alpha$  soit équidimensionnel signifie simplement
 que l'adhérence du support de  $\alpha$  a une dimension relative $\leq r$
  sur  $S$. Comme prendre les adhérences commute avec la restriction aux ouverts,
 le lemme s'ensuit (un détail mineur est omis). 
\end{proof} 

\begin{lemma} 
\label{relative-cycles-lemma-equidimensional-fppf-descent} 
Soit  $f : X \to S$  un morphisme de schémas. Supposons  $S$  localement noethérien
 et  $f$  localement de type fini. Soit  $r \geq 0$  un entier. Soit  $\alpha$
  un  $r$-cycle relatif sur  $X/S$. Soit $\{g_i : S_i \to S\}$
  un recouvrement fppf. Alors  $\alpha$  est équidimensionnel si 
et seulement si chaque changement de base  $g_i^*\alpha$  est équidimensionnel. 
\end{lemma} 

\begin{proof} 
Si  $\alpha$  est équidimensionnel, alors chaque  $g_i^*\alpha$  l'est aussi d'après
 le lemme \ref{relative-cycles-lemma-equidimensional-functoriality}. Supposons que chaque  $g_i^*\alpha$
  soit équidimensionnel. Notons  $W$  l'adhérence 
de  $\text{Supp}(\alpha)$  dans  $X$. Comme $g_i : S_i \to S$  est universellement ouvert
 (étant plat et localement de présentation finie), il en est de même pour le morphisme  
$f_i : X_i = S_i \times_S X \to X$. Notons  $\alpha_i = g_i^*\alpha$. Nous avons  
$\text{Supp}(\alpha_i) = f_i^{-1}(\text{Supp}(\alpha))$  d'après
 le lemme \ref{relative-cycles-lemma-support-family}. 
Puisque  $f_i$  est ouvert, nous voyons que  $W_i = f_i^{-1}(W)$  est l'adhérence 
de  $\text{Supp}(\alpha_i)$. Donc, d'après l'hypothèse, le morphisme  
$W_i \to S_i$  a une dimension relative $\leq r$.
 D'après Morphismes, lemme \ref{morphisms-lemma-dimension-fibre-after-base-change}
 (et le fait que les morphismes $S_i \to S$  sont conjointement surjectifs)
 nous concluons que  $W \to S$  a une dimension relative  $\leq r$. 
\end{proof} 

\begin{lemma} 
\label{relative-cycles-lemma-equidimensional-descent-pullbacks} 
Soit  $f : X \to S$  un morphisme de schémas. Supposons  $S$  localement noethérien
 et  $f$  localement de type fini. Soient  $r, e \geq 0$  des entiers. 
Soit  $\alpha$  un  $r$-cycle relatif sur  $X/S$.
 Soit $\{f_i : X_i \to X\}$  une famille conjointement surjective
 de morphismes plats, localement de type fini, et de dimension relative  $e$. 
Alors  $\alpha$  est équidimensionnel si et seulement si chaque image 
inverse plate  $f_i^*\alpha$  est équidimensionnel. 
\end{lemma} 

\begin{proof} La 
démonstration est omise. Indication : comme dans la démonstration du lemme \ref{relative-cycles-lemma-equidimensional-fppf-descent},
 on montre que l'image inverse par  $f_i$  de l'adhérence  $W$  du support de  
$\alpha$  est l'adhérence  $W_i$  du support de  $f_i^*\alpha$. Alors, 
$W \to S$  a une dimension relative  $\leq r$  si  $W_i \to S$
  a une dimension relative  $\leq r + e$  pour tout  $i$. 
\end{proof} 





\section{Pondérations et zéro-cycles relatifs} 
\label{relative-cycles-section-weightings} 

\noindent
 Dans cette section, nous montrons que pour un morphisme quasi-fini, les 
zéro-cycles relatifs sont étroitement liés aux pondérations, telles qu'elles sont définies dans 
Compléments sur les morphismes, section \ref{more-morphisms-section-weightings}. 

\medskip\noindent
 Soit  $S$  un schéma localement noethérien. Soit  $f : X \to S$  un morphisme 
localement quasi-fini de schémas. Alors nous avons  $z(X/S, 0) = z_{equi}(X/S, 0)$
  et  $z(X/S, r) = 0$ pour  $r > 0$. Étant donné  $\alpha \in z(X/S, 0)$ , définissons
 une application 
$$
w_\alpha : X \longrightarrow \mathbf{Z},\quad
x \mapsto \alpha(x) [\kappa(x) : \kappa(s)]_i \quad\text{où }s = f(x)
$$
 Ici,  $\alpha(x)$  désigne le coefficient de  $x$  dans le $0$-cycle  
$\alpha_s$  sur la fibre  $X_s$  et  $[K : k]_i$  désigne le degré inséparable
 d'une extension finie de corps. 

\begin{lemma} 
\label{relative-cycles-lemma-weightings-pre} 
Dans la situation ci-dessus, si  $g : S' \to S$  est un morphisme de schémas 
localement noethériens, alors  $w_{g^*\alpha} = w_\alpha \circ g'$  où  
$g' : X' \to X$  est la projection $X' = S' \times_S X \to X$. 
\end{lemma} 

\begin{proof} 
Soit  $x' \in X'$  dont les images sont  $s', s, x$  dans  $S', S, X$.
 Alors le coefficient de  $[x']$  dans
 le changement de base de  $[x]$  par $\kappa(s')/\kappa(s)$  est la longueur de
 l'anneau local  $(\kappa(s') \otimes_{\kappa(s)} \kappa(x))_\mathfrak q$.
 Ici  $\mathfrak q$  est l'idéal premier correspondant à  $x'$. 
Ainsi, la compatibilité avec le changement de base résulte de l'égalité 
$$
[\kappa(x) : \kappa(s)]_i =
\text{longueur}((\kappa(s') \otimes_{\kappa(s)} \kappa(x))_\mathfrak q)
[\kappa(x') : \kappa(s')]_i
$$
 Fixons une clôture algébrique  $k/\kappa(s')$. Choisissons un idéal premier  
$\mathfrak p \subset k \otimes_{\kappa(s)} \kappa(x)$ au-dessus
 de  $\mathfrak q$. Supposons que nous puissions montrer que 
$$
[\kappa(x) : \kappa(s)]_i =
\text{longueur}((k \otimes_{\kappa(s)} \kappa(x))_\mathfrak p)
\quad\text{et}\quad
[\kappa(x') : \kappa(s')]_i =
\text{longueur}((k \otimes_{\kappa(s')} \kappa(x'))_\mathfrak p)
$$
 Le résultat s'ensuit car 
$$
\text{longueur}((\kappa(s') \otimes_{\kappa(s)} \kappa(x))_\mathfrak q)
\text{longueur}((k \otimes_{\kappa(s')} \kappa(x'))_\mathfrak p)
=
\text{longueur}((k \otimes_{\kappa(s)} \kappa(x))_\mathfrak p)
$$
 d'après Algèbre, lemme \ref{algebra-lemma-pullback-module}, et la platitude de  
$\kappa(s') \otimes_{\kappa(s)} \kappa(x) \to k \otimes_{\kappa(s)} \kappa(x)$.
 Pour montrer les deux égalités, il suffit de prouver la première. 
Soit  $\kappa(x)/\kappa/\kappa(s)$  le sous-corps construit dans
 Corps, lemme \ref{fields-lemma-separable-first}. On obtient alors 
$$
k \otimes_{\kappa(s)} \kappa(x) =
\prod\nolimits_{\sigma : \kappa \to k}
k \otimes_{\sigma, \kappa} \kappa(x)
$$
 et chacun des facteurs est local de degré  
$[\kappa(x) : \kappa] = [\kappa(x) : \kappa(s)]_i$
 , comme voulu. 
\end{proof} 

\begin{lemma} 
\label{relative-cycles-lemma-weightings} 
Soit  $S$  un schéma localement noethérien. Soit  $f : X \to S$  un morphisme
 de schémas localement quasi-fini. 
\begin{enumerate} 
\item Pour  $\alpha \in z(X/S, 0)$ , l'application  $w_\alpha : X \to \mathbf{Z}$
  construite ci-dessus est une pondération. 
\item Si  $X$  est quasi-compact, alors étant donnée une pondération 
$w : X \to \mathbf{Z}$ , il existe un entier  $n > 0$  tel 
que  $nw = w_\alpha$  pour un certain  $\alpha \in z(X/S, 0)$. 
\item L'entier  $n$  de (2) peut être choisi égal à une puissance du 
nombre premier  $p$  si  $S$  est un schéma sur  $\mathbf{F}_p$. 
\end{enumerate} 
\end{lemma} 

\begin{proof} 
Soit  $\alpha \in z(X/S, 0)$  et choisissons un diagramme 
$$
\xymatrix{
X \ar[d]_f & U \ar[l]^h \ar[d]^\pi \\
Y & V \ar[l]_g
}
$$
 comme dans Compléments sur les morphismes, définition \ref{more-morphisms-definition-weighting}. 
Notons  $\beta \in z(U/V, 0)$  la restriction du changement de base  $g^*\alpha$.
 Par la compatibilité avec le changement de base (lemme \ref{relative-cycles-lemma-weightings-pre}), 
nous avons  $w_\beta = w_\alpha \circ h$  et il suffit de montrer que  
$\int_\pi w_\beta$  est localement constant sur  $V$. Ensuite, notons que 
\begin{align*} 
\left( \int_\pi w_\beta \right)(v) 
& = 
\sum\nolimits_{u \in U, \pi(u) = v} 
\beta(u) [\kappa(u) : \kappa(v)]_i [\kappa(u) : \kappa(v)]_s \\
& = 
\sum\nolimits_{u \in U, \pi(u) = v} \beta(u)[\kappa(u) : \kappa(v)] 
\end{align*} 
Cette dernière expression est le coefficient de  $v$  dans  $\pi_*\beta \in z(V/V, 0)$. 
D'après le lemme \ref{relative-cycles-lemma-relative-cycle-smooth}, cette fonction est localement 
constante sur $V$. 

\medskip\noindent
 Réciproquement, soit  $w : X \to S$  une pondération et  $X$  quasi-compact.
 Choisissons un entier  $n$ suffisamment divisible. Soit  $\alpha$  la famille 
de  $0$-cycles sur les fibres de  $X/S$  telle que pour $s \in S$  nous avons 
$$
\alpha_s =
\sum\nolimits_{f(x) = s} \frac{n w(x)}{[\kappa(x) : \kappa(s)]_i} [x]
$$
 comme un zéro-cycle sur  $X_s$.
 Cela a un sens puisque les fibres de $f$  sont universellement bornées
 (Morphismes, lemme 
\ref{morphisms-lemma-locally-quasi-finite-qc-source-universally-bounded}), donc 
nous pouvons trouver  $n$  de sorte que le membre de droite soit un entier pour 
tout  $s \in S$. L'énoncé final du lemme en découle également, à condition que nous montrions
 que  $\alpha$  est un  $0$-cycle relatif. Pour ce faire, il faut 
montrer que  $\alpha$  est compatible avec les spécialisations le long des spectres
 d'anneaux de valuation discrète. Notons que ni la construction de $w_\alpha$
  ni la démonstration du lemme \ref{relative-cycles-lemma-weightings-pre} ne supposent que  $\alpha$
 soit un cycle relatif. Par conséquent, nous avons  $w_\alpha = nw$  et cela reste vrai 
après un changement de base. De même, le changement de base d'une pondération est une
 pondération ; voir 
Compléments sur les morphismes, lemme \ref{more-morphisms-lemma-weighting-base-change}. 
Ainsi nous revenons au problème étudié dans le paragraphe suivant. 

\medskip\noindent
 Supposons que  $S$  soit le spectre d'un anneau de valuation discrète avec un point 
générique  $\eta$  et un point fermé  $0$. Soit  $w : X \to S$ une pondération, avec  
$X$  quasi-fini sur  $S$. Soit  $\alpha$  la famille de  $0$-cycles sur les
 fibres de  $X/S$  construite dans le paragraphe précédent (pour un  $n$ 
approprié). Nous devons montrer que $sp_{X/S}(\alpha_\eta) = \alpha_0$. 
Soit  $\beta \in z(X/S, 0)$  le  $0$-cycle relatif sur  $X/S$
 avec  $\beta_\eta = \alpha_\eta$  et  $\beta_0 = sp_{X/S}(\alpha_\eta)$. 
Alors  $w' = w_\beta - nw : X \to \mathbf{Z}$
  est une pondération (en utilisant le résultat ci-dessus) et nulle aux points de  $X$
  qui se projettent sur  $\eta$. Il résulte de Compléments sur les morphismes,
 lemme \ref{more-morphisms-lemma-weighting-specialization} 
que  $w' = 0$. Cela signifie  $w_\beta = nw$, donc  $\alpha = \beta$  comme souhaité. 
\end{proof} 

\begin{example} 
\label{relative-cycles-example-n-must-be-positive-sometimes} 
Soit  $p$  un nombre premier. Soit  $k'/k$  une extension purement inséparable
 de corps de degré  $p^f$. Posons  $X = \Spec(k')$  et  $S = \Spec(k)$. 
L'application  $w : X \to \mathbf{Z}$  envoyant le point unique  $x$  de  $X$ sur  $1$
  est une pondération de  $X \to S$. D'autre part, 
le groupe $z(X/S, 0)$  est librement engendré par le cycle  $\alpha = [x]$.
 Ainsi, nous voyons que le plus petit entier  $n$  qui convient dans 
le lemme \ref{relative-cycles-lemma-weightings} est  $n = p^f$  dans ce cas. 
\end{example} 

\begin{example} 
\label{relative-cycles-example-Merkurjev} La
 discussion ci-dessus peut être utilisée pour « expliquer » un exemple dû 
à A.S.~Merkurjev \cite[exemple 3.5.10]{SV}. Soit  $p$  un nombre premier 
et soit  $k = \mathbf{F}_p(a, b)$  l'extension purement transcendante 
de  $\mathbf{F}_p$ engendrée par  $a$  et  $b$. Posons 
$$
A = k[x, y, z]/(a x^p + b y^p - z^p)
$$
 L'anneau ainsi obtenu est un anneau intègre normal de type fini sur  $k$. Le spectre  
$S = \Spec(A)$ est régulier, à l'exception d'un point singulier unique  $s$
  correspondant à l'idéal maximal  $(x, y, z)$  de  $A$, de corps résiduel  $k$. 
Soient  $k' = k(a^{1/p}, b^{1/p})$  et  $B = k'[x, y]$. Considérons l'homomorphisme  
$A \to B$ induit par l'inclusion  $k \to k'$ et envoyant  $x$  sur  $x$, $y$  sur  $y$,
 et  $z$  sur  $a^{1/p}x + b^{1/p}y$. Posons $X = \Spec(B)$  et considérons le morphisme  
$f : X \to S$  correspondant au morphisme d'anneaux  $A \to B$.
 La pondération  $w : X \to \mathbf{Z}$  de  $f$
  donnée par Compléments sur les morphismes, lemme 
\ref{more-morphisms-lemma-weighting-quasi-finite-Noetherian}, est 
constante de valeur  $1$, puisque le corps des fractions de  $B$  est une extension purement inséparable de 
celui de  $A$. 
Posons  $V = f^{-1}(U) \subset X$. Par le critère de platitude par miracle, le morphisme  
$V \to U$ est plat de degré  $p$. Il s'ensuit que le cycle  
$\alpha = [V/V/U]_0 \in z(V/U, 0)$  satisfait $w_\alpha = pw|_V$.
 Cependant, puisque le point unique  $x$  de  $X$ situé au-dessus de  $s$
  a pour corps résiduel  $k'$  de degré $p^2$  sur  $k$, nous voyons 
qu'il ne peut y avoir d'élément $\beta$  de  $z(X/S, 0)$  se restreignant à  $\alpha$.
 En effet, la pondération $pw - w_\beta$  serait nulle d'après Compléments sur les morphismes,
 lemme \ref{more-morphisms-lemma-weighting-specialization}, ce qui 
impliquerait (d'après la formule dans la démonstration du lemme ci-dessus) que 
$$
\beta_s = \frac{pw(x)}{[\kappa(x) : \kappa(s)]_i} [x] =
\frac{1}{p}[x]
$$
 dans  $Z_0(X_s)$  et c'est impossible\footnote{Dans  \cite{SV}  il existe
 un élément du groupe des cycles relatifs  $Cycl(X/S, 0)$
  correspondant à $pw$  et ce cycle se spécialise en un cycle sur
  $X_s$  avec des coefficients non entiers.}. Bien sûr,
 le cycle  $p\alpha$  est la restriction d'un élément unique de  
$z(X/S, 0)$. 
\end{example} 









\section{Cycles relatifs effectifs} 
\label{relative-cycles-section-effective} 


\noindent
 Voici la définition. 

\begin{definition} 
\label{relative-cycles-definition-effective} 
Soit  $f : X \to S$  un morphisme de schémas. Supposons que  $S$  soit localement noethérien et
 que  $f$  soit localement de type fini. Soit $r \geq 0$  un entier. On dit 
qu'un  $r$-cycle relatif  $\alpha$  sur $X/S$  est {\it effectif} si  $\alpha_s$ est un
 cycle effectif 
(Homologie de Chow, définition \ref{chow-definition-effective-cycle}) 
pour tout $s \in S$. Le monoïde de tous les  $r$-cycles relatifs effectifs
 sur $X/S$  est noté  $z^{eff}(X/S, r)$. 
\end{definition} 

\noindent
 Plus loin, nous montrerons qu'un cycle relatif effectif est équidimensionnel ; 
voir le lemme \ref{relative-cycles-lemma-effective-equidimensional}. 

\begin{lemma} 
\label{relative-cycles-lemma-effective-functoriality} 
Soit  $f : X \to S$  un morphisme de schémas. Supposons que  $S$  soit localement noethérien et
 que  $f$  soit localement de type fini. Soit $r \geq 0$  un entier. Soit  
$\alpha$  un  $r$-cycle relatif sur $X/S$. Si  $\alpha$  est effectif, alors 
l'effectivité est préservée par toute restriction, tout changement de base, toute image inverse plate et toute image directe propre
 de  $\alpha$. 
\end{lemma} 

\begin{proof} La 
démonstration est omise. 
\end{proof} 

\begin{lemma} 
\label{relative-cycles-lemma-check-effective} 
Soit  $f : X \to S$  un morphisme de schémas. Supposons que  $S$  soit localement noethérien et
 que  $f$  soit localement de type fini. Soit $r \geq 0$  un entier. Soit  $\alpha$
  un  $r$-cycle relatif sur $X/S$. Alors pour vérifier que  $\alpha$  est effectif, nous pouvons
 travailler localement pour la topologie de Zariski sur  $X$  et $S$. 
\end{lemma} 

\begin{proof} La 
démonstration est omise. 
\end{proof} 

\begin{lemma} 
\label{relative-cycles-lemma-effective-descent} 
Soit $f : X \to S$ un morphisme de schémas. Supposons $S$ localement noethérien
 et $f$ localement de type fini. Soit $r \geq 0$ un entier. Soit $\alpha$
 un $r$-cycle relatif sur $X/S$. Soit $g : S' \to S$ un morphisme surjectif. 
Alors $\alpha$ est effectif si et seulement si le changement de base $g^*\alpha$
 est effectif. 
\end{lemma} 

\begin{proof} La 
démonstration est omise. 
\end{proof} 

\begin{lemma} 
\label{relative-cycles-lemma-effective-descent-pullbacks} 
Soit $f : X \to S$ un morphisme de schémas. Supposons $S$ localement noethérien
 et $f$ localement de type fini. Soient $r, e \geq 0$ des entiers. 
Soit $\alpha$ un $r$-cycle relatif sur $X/S$.
 Soit $\{f_i : X_i \to X\}$ une famille conjointement surjective
 de morphismes plats, localement de type fini et de dimension relative $e$.
 Alors $\alpha$ est effectif si et seulement si chaque
 cycle $f_i^*\alpha$ obtenu par image inverse plate est effectif. 
\end{lemma} 

\begin{proof} La 
démonstration est omise. 
\end{proof} 

\begin{lemma} 
\label{relative-cycles-lemma-effective-support-closed} 
Soit  $f : X \to S$  un morphisme de schémas. Supposons que  $S$  soit localement noethérien
 et  $f$  localement de type fini. Soient  $r, e \geq 0$ des entiers. 
Soit  $\alpha$  un  $r$-cycle relatif sur  $X/S$.
 Si $\alpha$  est effectif, alors  $\text{Supp}(\alpha)$  est
 fermé dans  $X$. 
\end{lemma} 

\begin{proof} 
Soit  $g : S' \to S$  l'inclusion d'une composante irréductible vue 
comme un sous-schéma fermé intègre. D'après les 
lemmes  \ref{relative-cycles-lemma-effective-functoriality}  et  \ref{relative-cycles-lemma-support-family} ,
 il suffit de montrer que le support du changement de base  
$g^*\alpha$  est fermé dans $S' \times_S S$. 
Ainsi nous pouvons supposer que  $S$  est un schéma intègre de point générique  
$\eta$. Nous montrerons que  $\text{Supp}(\alpha)$  est l'adhérence 
de  $\text{Supp}(\alpha_\eta)$. Pour ce faire, choisissons un point  $s \in S$. 
Nous pouvons trouver un morphisme  $g : S' \to S$  où  $S'$  est le spectre 
d'un anneau de valuation discrète envoyant le point générique  $\eta' \in S'$  vers $\eta$
  et le point fermé  $0 \in S'$  vers  $s$, voir
 Propriétés, lemme \ref{properties-lemma-locally-Noetherian-specialization-dvr}. 
Il suffit alors de prouver que le support de  $g^*\alpha$
  est égal à l'adhérence de  $\text{Supp}((g^\alpha)_{\eta'})$.
 Cela nous ramène au cas discuté dans le paragraphe suivant. 

\medskip\noindent
 Ici  $S$  est le spectre d'un anneau de valuation discrète avec point 
générique  $\eta$  et point fermé  $0$. Nous devons montrer que 
$\text{Supp}(\alpha)$  est l'adhérence de  $\text{Supp}(\alpha_\eta)$. 
Puisque  $\alpha$  est effectif, nous pouvons écrire  $\alpha_\eta = \sum n_i[Z_i]$
  où  $n_i > 0$  et où les  $Z_i \subset X_\eta$ sont des sous-schémas fermés intègres de dimension  $r$. 
Puisque  $\alpha_0 = sp_{X/S}(\alpha_\eta)$, nous savons que 
$\alpha_0 = \sum n_i [\overline{Z}_{i, 0}]_r$, où  $\overline{Z}_i$
  est l'adhérence de  $Z_i$. D'après Variétés, lemme 
\ref{varieties-lemma-dominate-valuation-ring-dimension-fibres}, nous
 voyons que  $\overline{Z}_{i, 0}$  est équidimensionnel de dimension $r$. 
Puisque  $n_i > 0$, nous en concluons que  $\text{Supp}(\alpha_0)$
  est égal à l'union des  $\overline{Z}_{i, 0}$. Cette union
 est la fibre au-dessus de  $0$  de $\bigcup \overline{Z}_i$ ; ce dernier ensemble
 est à son tour l'adhérence de  $\bigcup Z_i$, comme voulu. 
\end{proof} 

\begin{lemma} 
\label{relative-cycles-lemma-effective-equidimensional} 
Soit  $f : X \to S$  un morphisme de schémas. Supposons  $S$  localement noethérien
 et  $f$  localement de type fini. Soient  $r, e \geq 0$  des entiers. 
Soit  $\alpha$  un  $r$-cycle relatif sur  $X/S$.
 Si $\alpha$  est effectif, alors  $\alpha$  est équidimensionnel. 
\end{lemma} 

\begin{proof} 
Supposons que  $\alpha$  soit effectif. D'après le lemme \ref{relative-cycles-lemma-effective-support-closed},
 le support $\text{Supp}(\alpha)$  est fermé dans  $X$. Ainsi,  
$\alpha$  est équidimensionnel car les fibres de  $\text{Supp}(\alpha) \to S$
  sont les supports des cycles  $\alpha_s$  et ont donc dimension  $r$. 
\end{proof} 

\begin{remark} 
\label{relative-cycles-remark-representable} 
Soit  $f : X \to S$  un morphisme de schémas avec  $S$  localement noethérien
 et  $f$  localement de type fini. On peut se demander si le foncteur contravariant 
$$
\begin{matrix}
\text{schémas }S'\text{ localement} \\
\text{de type fini sur }S
\end{matrix}
\longrightarrow
z^{eff}(X'/S', r)\text{ où }X' = S' \times_S X
$$
 est représentable. Comme  $z(X'/S', r) = z(X'_{red}/S'_{red}, r)$
 , cela ne peut pas être vrai (nous laissons au lecteur le soin de donner un véritable
 contre-exemple). Une meilleure question serait de savoir si nous pouvons 
trouver une sous-catégorie du membre de gauche sur laquelle le foncteur est
 représentable. Le lemme  \ref{relative-cycles-lemma-seminormalize-base}  suggère que 
nous devrions au moins nous restreindre à la catégorie des schémas semi-normaux sur $S$. 

\medskip\noindent
 Si  $S/\Spec(\mathbf{Q})$  est Nagata et  $f$  est un morphisme projectif, alors
 il s'avère que  $S' \mapsto z^{eff}(X'/S', r)$  est représentable sur la 
catégorie dont les objets $S'$  sont semi-normaux. Essentiellement, c'est le contenu de 
\cite[théorème 3.21]{KRC}. 

\medskip\noindent
 Si  $S$  possède des points de caractéristique positive, cela ne fonctionne plus même
 si nous remplaçons la semi-normalité par la normalité faible ; un schéma localement 
noethérien  $T$  est faiblement normal si tout homéomorphisme universel birationnel  
$T' \to T$  possède une section. Un exemple consiste à considérer les  $0$-cycles de degré  
$2$  sur  $X = \mathbf{A}^2_k$ au-dessus de  $S = \Spec(k)$  où  $k$  est un corps 
de caractéristique $2$. Plus précisément, au-dessus de  $W = X \times_S X$  nous avons
 un  $0$-cycle relatif canonique  $\alpha \in z^{eff}(X_W/W, 0)$ : pour  
$w = (x_1, x_2) \in W = X^2$  nous avons le cycle $\alpha_w = [x_1] + [x_2]$.
 Ce cycle est invariant sous l'involution  $\sigma : W \to W$  échangeant
 les facteurs. Puisque  $W$  est lisse (donc normal, donc faiblement normal), si 
$z(-/-, r)$  était représentable par  $M$  sur la catégorie des schémas faiblement
 normaux de type fini sur  $k$ , nous obtiendrions un morphisme invariant par  $\sigma$, 
de  $W$  vers  $M$. Cela définirait à son tour un morphisme du 
schéma quotient  $\text{Sym}^2_S(X) = W/\langle \sigma \rangle$  vers  $M$.
 Puisque $\text{Sym}^2_S(X)$  est normal, nous obtiendrions, par la propriété universelle de l'espace de modules $M$
 , un  $0$-cycle relatif  $\beta$  sur  
$X \times_S \text{Sym}^2_S(X) / \text{Sym}^2_S(X)$
  dont l'image inverse sur  $W$  est  $\alpha$. Cependant, il n'existe aucun cycle $\beta$  
de cette sorte. En effet, en écrivant  $X = \Spec(k[u, v])$, le schéma  
$\text{Sym}^2_S(X)$ est le spectre de 
$$
k[u_1 + u_2, u_1u_2, v_1 + v_2, v_1v_2, u_1v_1 + u_2v_2]
\subset
k[u_1, u_2, v_1, v_2]
$$
 L'image de la diagonale  $u_1 = u_2, v_1 = v_2$  dans  $\text{Sym}^2_S(X)$
  est le sous-schéma fermé $V = \Spec(k[u_1^2, v_1^2])$; ici nous utilisons le 
fait que la caractéristique de  $k$ est  $2$. En regardant le 
point générique  $\eta$  de  $V$, le cycle  $\beta_\eta$  serait un 
zéro-cycle de degré  $2$  sur  $\mathbf{A}^2_{k(u_1^2, v_1^2)}$
 dont l'image inverse sur  $\mathbf{A}^2_{k(u_1, u_2)}$  serait  
$2[\text{le point de coordonnées} (u_1, v_2)]$.
 Cela est clairement impossible. 

\medskip\noindent
 La discussion ci-dessus ne contredit pas  \cite[théorème 4.13]{KRC}  car la variété
 de Chow de ce théorème ne représente un foncteur qu'au sens grossier (en
 fait deux foncteurs distincts, dont un seul coïncide avec le nôtre pour  $X$
  projectif comme on peut le voir avec un peu de travail). De même, dans  \cite[section 4.4]{SV}
  il est montré que pour  $X/S$  projectif, la faisceautisation pour la topologie  $h$ du préfaisceau 
$S' \mapsto z^{eff}(S' \times_S X/S', r)$  est égale à la faisceautisation pour la topologie  $h$ d'un foncteur
 représentable. 
\end{remark} 

\begin{remark} 
\label{relative-cycles-remark-equidimensional-over-geometrically-unibranch} 
Soit  $f : X \to S$  un morphisme de schémas. Soit  $r \geq 0$. Soit  $Z \subset X$
  un sous-schéma fermé. Supposons : 
\begin{enumerate} 
\item $S$  est noethérien et géométriquement unibranche, 
\item $f$  est de type fini, et 
\item $Z \to S$  est de dimension relative  $\leq r$. 
\end{enumerate} 
Alors, pour tout entier suffisamment divisible $n \geq 1$, il 
existe un unique $r$-cycle relatif effectif $\alpha$ sur $X/S$ tel que 
$\alpha_\eta = n[Z_\eta]_r$ pour tout point générique $\eta$  de  $S$.
 Il s'agit d'une reformulation de  \cite[théorème 3.4.2]{SV}. 
Si nous avons besoin de ce résultat, nous le préciserons et le prouverons ici. 
\end{remark} 











\section{Cycles relatifs propres} 
\label{relative-cycles-section-proper} 

\noindent
 Dans notre contexte, la définition suivante est probablement la bonne. 

\begin{definition} 
\label{relative-cycles-definition-proper} 
Soit  $f : X \to S$  un morphisme de schémas. Supposons que  $S$  soit localement noethérien et
 que  $f$  soit localement de type fini. Soit $r \geq 0$  un entier. On dit 
qu'un  $r$-cycle relatif  $\alpha$  sur $X/S$  est un  {\it  cycle relatif propre} 
 si le support de  $\alpha$  (remarque \ref{relative-cycles-remark-supports-family}) est
 contenu dans un sous-ensemble fermé  $W \subset X$  propre sur  $S$
  (Cohomologie des schémas, définition \ref{coherent-definition-proper-over-base}).
 Le groupe de tous les  $r$-cycles relatifs propres sur  $X/S$  
est noté  $c(X/S, r)$. 
\end{definition} 

\noindent
 D'après Cohomologie des schémas, lemme 
\ref{coherent-lemma-closed-closed-proper-over-base}, cela 
signifie simplement que l'adhérence du support est propre sur la base. Pour voir que
 ces cycles forment un groupe, on utilise
 Cohomologie des schémas, lemme \ref{coherent-lemma-union-closed-proper-over-base}. 

\begin{lemma} 
\label{relative-cycles-lemma-proper-functoriality} 
Soit  $f : X \to S$  un morphisme de schémas. Supposons que  $S$  soit localement noethérien et
 que  $f$  soit localement de type fini. Soit $r \geq 0$  un entier. Soit  
$\alpha$  un  $r$-cycle relatif sur $X/S$. Si  $\alpha$  est propre, 
alors tout changement de base de $\alpha$  est propre. 
\end{lemma} 

\begin{proof} La 
démonstration est omise. 
\end{proof} 

\begin{lemma} 
\label{relative-cycles-lemma-proper-h-descent} 
Soit  $f : X \to S$  un morphisme de schémas. Supposons que  $S$  soit localement noethérien et
 que  $f$  soit localement de type fini. Soit $r \geq 0$  un entier. Soit  $\alpha$
  un  $r$-cycle relatif sur $X/S$. Soit  $\{g_i : S_i \to S\}$
  un recouvrement h. Alors  $\alpha$  est propre 
si et seulement si chaque changement de base  $g_i^*\alpha$  est propre. 
\end{lemma} 

\begin{proof} 
Si  $\alpha$  est propre, alors chaque  $g_i^*\alpha$  l'est aussi d'après
 le lemme \ref{relative-cycles-lemma-proper-functoriality}. Supposons que chaque  $g_i^*\alpha$
  soit propre. Pour prouver que $\alpha$  est propre, il suffit de travailler 
localement sur des ouverts affines sur  $S$. Nous pouvons donc supposer que  $S$
  est affine. Nous pouvons alors affiner notre recouvrement  $\{S_i \to S\}$
  par une famille $\{T_j \to S\}$  où  $g : T \to S$  est un morphisme propre 
et surjectif et  $T = \bigcup T_j$ est un recouvrement ouvert.
 Il s'ensuit que  $\beta = g^*\alpha$  est propre sur $Y = T \times_S X$
  au-dessus de  $T$. D'après le lemme \ref{relative-cycles-lemma-support-family}, le support de
  $\beta$  est l'image inverse du support de $\alpha$  par le 
morphisme  $f : Y \to X$. Ainsi, l'adhérence  $W \subset Y$
  de $f^{-1}\text{Supp}(\alpha)$  est propre sur  $T$. Puisque le 
morphisme  $T \to S$  est propre, il s'ensuit que  $W$  est propre sur  $S$.
 D'après Cohomologie des schémas, lemme 
\ref{coherent-lemma-functoriality-closed-proper-over-base},
 l'image  $f(W) \subset X$  est un sous-ensemble fermé propre sur  $S$. 
Puisque  $f(W)$  contient $\text{Supp}(\alpha)$ , nous concluons que  $\alpha$
  est propre. 
\end{proof} 



\section{Cycles relatifs propres et équidimensionnels} 
\label{relative-cycles-section-proper-equidimensional} 

\noindent
 Soit  $f : X \to S$  un morphisme de schémas. Supposons que  $S$  soit localement noethérien et
 que  $f$  soit localement de type fini. Soit $r \geq 0$  un entier. On dit 
qu'un  $r$-cycle relatif  $\alpha$  sur $X/S$  est un {\it cycle relatif propre et 
équidimensionnel} si  $\alpha$  est à la fois équidimensionnel
 (définition \ref{relative-cycles-definition-equidimensional})
 et propre (définition \ref{relative-cycles-definition-proper}).
 Le groupe de tous les  $r$-cycles relatifs propres et équidimensionnels sur  $X/S$  
est noté  $c_{equi}(X/S, r)$. 

\medskip\noindent
 De même, nous disons qu'un  $r$-cycle relatif  $\alpha$  sur $X/S$  est un
 {\it cycle relatif propre et effectif} si  $\alpha$  est à la fois effectif
 (définition \ref{relative-cycles-definition-effective}) et propre 
(définition \ref{relative-cycles-definition-proper}).
 Le monoïde de tous les $r$-cycles relatifs propres et effectifs sur  $X/S$  
est noté $c^{eff}(X/S, r)$.
 Ces cycles sont équidimensionnels d'après le
 lemme \ref{relative-cycles-lemma-effective-equidimensional}. 

\medskip\noindent
 Ainsi, nous avons le diagramme suivant des applications d'inclusion 
$$
\xymatrix{
c^{eff}(X/S, r) \ar[r] \ar[d] &
c_{equi}(X/S, r) \ar[r] \ar[d] &
c(X/S, r) \ar[d] \\
z^{eff}(X/S, r) \ar[r] &
z_{equi}(X/S, r) \ar[r] &
z(X/S, r)
}
$$



\section{Action sur les cycles} 
\label{relative-cycles-section-action} 

\noindent
 Soit  $S$  un schéma localement noethérien, universellement caténaire, muni 
d'une fonction de dimension  $\delta$ ; voir
 Homologie de Chow, section \ref{chow-section-setup}. 
Soit  $X \to Y$  un morphisme de schémas sur  $S$, tous deux localement de type
 fini sur  $S$. Soit  $r \geq 0$. Enfin, soit $\alpha$  une famille de  
$r$-cycles sur les fibres de  $X/Y$. Pour  $e \in \mathbf{Z}$
 , nous allons construire une opération 
$$
\alpha \cap - : Z_e(Y) \longrightarrow Z_{r + e}(X)
$$
 Concrètement, étant donné  $\beta \in Z_e(Y)$, écrivons  $\beta = \sum n_i[Z_i]$
  où  $Z_i \subset Y$  est un sous-schéma fermé intègre de  
$\delta$-dimension  $e$  et la 
famille $Z_i$  est localement finie dans le schéma  $Y$. 
Soit  $y_i \in Z_i$ le point générique. Écrivons  
$\alpha_{y_i} = \sum m_{ij} [V_{ij}]$. Ainsi  $V_{ij} \subset X_{y_i}$
  est un sous-schéma fermé intègre de dimension  $r$  et la famille 
$V_{ij}$ est localement finie dans le schéma $X_{y_i}$. 
Nous posons alors 
$$
\alpha \cap \beta = \sum n_i m_{ij} [\overline{V}_{ij}]
\quad\in\quad
Z_{r + e}(X)
$$
 Ici, $\overline{V}_{ij} \subset X$ est l'image schématique 
du morphisme $V_{ij} \to X_{y_i} \to X$ ou, de manière équivalente, 
$\overline{V}_{ij} \subset X$ est un sous-schéma fermé 
intègre qui domine $Z_i \subset Y$ et dont la fibre générique est $V_{ij}$.
 On en déduit aisément que $\dim_\delta(\overline{V}_{ij}) = r + e$
 et que la famille de sous-schémas fermés 
$\overline{V}_{ij} \subset X$ est localement finie (nous omettons les vérifications). 
Ainsi, $\alpha \cap \beta$ est bien un élément de $Z_{r + e}(X)$. 

\begin{lemma} 
\label{relative-cycles-lemma-action-bilinear} La 
construction ci-dessus est bilinéaire, c'est-à-dire que 
$(\alpha_1 + \alpha_2) \cap \beta = \alpha_1 \cap \beta +
\alpha_2 \cap \beta$ et $\alpha \cap (\beta_1 + \beta_2) =
\alpha \cap \beta_1 + \alpha \cap \beta_2$. 
\end{lemma} 

\begin{proof} 
Démonstration omise. 
\end{proof} 

\begin{lemma} 
\label{relative-cycles-lemma-action-opens} 
Si $U \subset X$ et $V \subset Y$ sont ouverts et si $f(U) \subset V$, alors 
$(\alpha \cap \beta)|_U$ est égal à $\alpha|_U \cap \beta|_V$. 
\end{lemma} 

\begin{proof} 
Cela résulte immédiatement de la description explicite de $\alpha \cap \beta$
 donnée ci-dessus. 
\end{proof} 

\begin{lemma} 
\label{relative-cycles-lemma-action-base-change} La 
formation de $\alpha \cap \beta$ est compatible avec le changement de base plat et avec 
l'image inverse plate (voir la démonstration pour plus de précisions). 
\end{lemma} 

\begin{proof} 
Soient $(S, \delta)$, $(S', \delta')$, $g : S' \to S$ et $c \in \mathbf{Z}$
 comme dans Homologie de Chow, situation \ref{chow-situation-setup-base-change}. 
Soit $X \to Y$ un morphisme de schémas localement de type fini sur $S$.
 Notons $X' \to Y'$ le changement de base de $X \to Y$ par $g$.
 Soit $\alpha$ une famille de $r$-cycles sur les fibres de $X/Y$.
 Soit $\beta \in Z_e(Y)$. Notons $\alpha'$ le 
changement de base de $\alpha$ par $Y' \to Y$. Notons 
$\beta' = g^*\beta \in Z_{e + c}(Y')$ l'image inverse de $\beta$ par $g$; voir
 Homologie de Chow, section \ref{chow-section-change-base}. La 
compatibilité avec le changement de base signifie que 
$\alpha' \cap \beta'$ est le changement de base de $\alpha \cap \beta$. 

\medskip\noindent
 Démonstration de la compatibilité au changement 
de base. Puisque nous démontrons une égalité de cycles sur $X'$, nous pouvons travailler localement
 sur $Y$; voir le lemme \ref{relative-cycles-lemma-action-opens}. Nous pouvons donc supposer $Y$
 affine. En particulier, $\beta$ est une combinaison linéaire finie de cycles 
premiers. Puisque $- \cap -$ est linéaire en la seconde variable
 (lemme \ref{relative-cycles-lemma-action-bilinear}), il suffit de 
démontrer l'égalité lorsque $\beta = [Z]$ pour un sous-schéma fermé intègre 
$Z \subset Y$ de dimension $\delta$ égale à $e$. 

\medskip\noindent
 Soit $y \in Z$ le point générique. Écrivons $\alpha_y = \sum m_j [V_j]$.
 Soit $\overline{V}_j$ l'adhérence de $V_j$ dans $X$. Alors 
$$
\alpha \cap \beta = \sum m_j[\overline{V}_j]
$$
 Le changement de base de $\beta$ est $\beta' = \sum [Z \times_S S']_{e + c}$, considéré
 comme un cycle sur $Y' = Y \times_S S'$. Soient $Z'_a \subset Z \times_S S'$
 les composantes irréductibles; notons $y'_a \in Z'_a$ leurs points 
génériques et $n_a$ la multiplicité de $Z'_a$ dans $Z \times_S S'$. 
Nous avons 
$$
\beta' = \sum [Z \times_S S']_{e + c} = \sum n_a[Z'_a]
$$
 Nous avons $\alpha'_{y'_a} = \sum m_j [V_{j, \kappa(y'_a)}]_r$
 car $\alpha'$ est le changement de base de $\alpha$ par $Y' \to Y$. 
Soient $V'_{jab} \subset V_{j, \kappa(y'_a)}$ les composantes irréductibles,
 et notons $m_{jab}$ la multiplicité de $V'_{jab}$
 dans $V_{j, \kappa(y'_a)}$. Nous avons 
$$
\alpha'_{y'_a} = \sum m_j [V_{j, \kappa(y'_a)}]_r =
\sum m_j m_{jab} [V'_{jab}]
$$
 Ainsi, 
$$
\alpha' \cap \beta' = \sum n_a m_j m_{jab} [\overline{V}'_{jab}]
$$
 où $\overline{V}'_{jab}$ est l'adhérence de $V'_{jab}$ dans $X'$.
 Pour démontrer l'égalité voulue, il suffit donc d'établir que 
\begin{enumerate} 
\item les composantes irréductibles de $\overline{V}_j \times_S S'$
 sont les schémas $\overline{V}'_{jab}$, et que 
\item la multiplicité de $\overline{V}'_{jab}$ dans 
$\overline{V}_j \times_S S'$ est égale à $n_a m_{jab}$. 
\end{enumerate} 
Remarquons que $V_j \to \overline{V}_j$ est un morphisme birationnel 
de schémas intègres. Les morphismes $V_j \times_S S' \to V_j$
 et $\overline{V}_j \times_S S' \to \overline{V}_j$ sont plats
 et envoient donc les points génériques des composantes irréductibles sur les
 points génériques (uniques) de $V_j$ et de $\overline{V}_j$.
 Il s'ensuit que $V_j \times_S S' \to \overline{V}_j \times_S S'$
 est un morphisme birationnel, donc induit une bijection entre les 
composantes irréductibles et identifie leurs multiplicités.
 Il suffit par conséquent de démontrer que les composantes irréductibles de 
$V_j \times_S S'$ sont les schémas $V'_{jab}$ et 
que la multiplicité de $V'_{jab}$ dans $V_j \times_S S'$
 est égale à $n_a m_{jab}$. Mais cela revient précisément à dire
 que le diagramme 
$$
\xymatrix{
Z_r(V_j) \ar[r] & Z_{r + c}(V_j \times_S S') \\
Z_0(\Spec(\kappa(y))) \ar[r] \ar[u] &
Z_c(\Spec(\kappa(y)) \times_S S') \ar[u]
}
$$
 est commutatif, où les flèches horizontales sont données par changement de base le long de 
$\Spec(\kappa(y)) \times_S S' \to \Spec(\kappa(y))$ et 
les flèches verticales sont les images inverses plates. Cela a été démontré 
dans Homologie de Chow, lemme \ref{chow-lemma-pullback-base-change-pullback}. 

\medskip\noindent
 L'assertion du lemme concernant l'image inverse plate signifie ce qui suit. 
Soient $(S, \delta)$, $X \to Y$, $\alpha$ et $\beta$ comme dans la
 construction de $\alpha \cap \beta$ ci-dessus. Soit $Y' \to Y$ un morphisme
 plat, localement de type fini et de dimension relative $c$. Nous pouvons
 alors prendre pour $\alpha'$ le changement de base de $\alpha$ par $Y' \to Y$
 et pour $\beta'$ l'image inverse plate de $\beta$. La compatibilité avec l'image inverse
 plate signifie que $\alpha' \cap \beta'$ est l'image inverse plate de 
$\alpha \cap \beta$ par $X \times_Y Y' \to Y$. Il s'agit en fait d'un
 cas particulier de la discussion précédente si nous posons $S = Y$ et $S' = Y'$. 
\end{proof} 

\begin{lemma} 
\label{relative-cycles-lemma-action-coherent} 
Soient $(S, \delta)$ et $f : X \to Y$ comme ci-dessus. 
Soit $\mathcal{F}$ un $\mathcal{O}_X$-module cohérent
 tel que $\dim(\text{Supp}(\mathcal{F}_y)) \leq r$ pour tout $y \in Y$.
 Soit $\mathcal{G}$ un $\mathcal{O}_Y$-module cohérent
 tel que $\dim_\delta(\text{Supp}(\mathcal{G})) \leq e$. 
Posons $\alpha = [\mathcal{F}/X/Y]_r$
 (exemple \ref{relative-cycles-example-family-associated-module}) et 
$\beta = [\mathcal{G}]_e$ (Homologie de Chow, définition 
\ref{chow-definition-cycle-associated-to-coherent-sheaf}). 
Si $\mathcal{F}$ est plat sur $Y$, alors $\alpha \cap \beta =
[\mathcal{F} \otimes_{\mathcal{O}_X} f^*\mathcal{G}]_{r + e}$. 
\end{lemma} 

\begin{proof} 
Observons que 
$$
\text{Supp}(\mathcal{F} \otimes_{\mathcal{O}_X} f^*\mathcal{G}) =
\text{Supp}(\mathcal{F}) \cap f^{-1}\text{Supp}(\mathcal{G}) =
\bigcup\nolimits_{y \in \text{Supp}(\mathcal{G})} \text{Supp}(\mathcal{F}_y)
$$
 Il en résulte qu'il s'agit d'un fermé dont la dimension $\delta$ est $\leq r + e$.
 L'expression 
$[\mathcal{F} \otimes_{\mathcal{O}_X} f^*\mathcal{G}]_{r + e}$
 a donc un sens. 

\medskip\noindent
 Nous utiliserons les notations 
$\beta = \sum n_i[Z_i]$, $y_i \in Z_i$, 
$\alpha_{y_i} = \sum m_{ij} [V_{ij}]$ et 
$\overline{V}_{ij}$ introduites dans la construction 
de $\alpha \cap \beta$. Puisque $\beta = [\mathcal{G}]_e$
 , les $Z_i$ sont les composantes irréductibles de 
$\text{Supp}(\mathcal{G})$ qui sont de dimension $\delta$ égale à $e$.
 De même, les $V_{ij}$ sont les composantes irréductibles 
de $\text{Supp}(\mathcal{F}_{y_i})$ de dimension $r$.
 Il résulte de ceci et de l'égalité du premier 
paragraphe que les $\overline{V}_{ij}$ sont les composantes 
irréductibles de 
$\text{Supp}(\mathcal{F} \otimes_{\mathcal{O}_X} f^*\mathcal{G})$
 de dimension $\delta$ égale à $r + e$.
 Pour démontrer le lemme, il suffit donc désormais de montrer que 
$$
\text{longueur}_{\mathcal{O}_{X, \xi_{ij}}}(
(\mathcal{F} \otimes_{\mathcal{O}_X} f^*\mathcal{G})_{\xi_{ij}})
=
\text{longueur}_{\mathcal{O}_{X_{y_i}, \xi_{ij}}}((\mathcal{F}_{y_i})_{\xi_{ij}})
\cdot
\text{longueur}_{\mathcal{O}_{Y, y_i}}(\mathcal{G}_{y_i})
$$
 D'après le premier paragraphe de la démonstration, le membre de gauche est 
égal à la longueur du $B = \mathcal{O}_{X, \xi_{ij}}$-module 
$$
\mathcal{G}_{y_i}
\otimes_{\mathcal{O}_{Y, y_i}}
\mathcal{F}_{\xi_{ij}} =
M \otimes_A N
$$
 Ici, $M = \mathcal{G}_{y_i}$ est un 
$A = \mathcal{O}_{Y, y_i}$-module de longueur finie, et $N = \mathcal{F}_{\xi_{ij}}$
 est un $B$-module fini tel que $N/\mathfrak m_AN$ soit de longueur finie. 
Puisque $\mathcal{F}$ est plat sur $Y$, le module $N$ est plat sur $A$. Le 
membre de droite de la formule est égal à 
$$
\text{longueur}_B(N/\mathfrak m_A N) \cdot \text{longueur}_A(M)
$$
 Les membres de droite et de gauche de la formule sont donc additifs en $M$
 (on utilise la platitude de $N$ sur $A$). Il suffit ainsi de démontrer la 
formule pour $M = \kappa_A$, le corps résiduel ; dans ce 
cas, elle est immédiate. 
\end{proof} 

\begin{lemma} 
\label{relative-cycles-lemma-action-closed} 
Soient $(S, \delta)$ et $f : X \to Y$ comme ci-dessus. Soit $Z \subset X$
 un sous-schéma fermé de dimension relative $\leq r$ sur $Y$. 
Posons $\alpha = [Z/X/Y]_r$ (exemple \ref{relative-cycles-example-family-associated-closed}). 
Soit $W \subset Y$ un sous-schéma fermé dont la dimension $\delta$ est $\leq e$. 
Posons $\beta = [W]_e$ (Homologie de Chow, définition 
\ref{chow-definition-cycle-associated-to-closed-subscheme}). 
Si $Z$ est plat sur $Y$, alors $\alpha \cap \beta = [Z \times_Y W]_{r + e}$. 
\end{lemma} 

\begin{proof} 
C'est un cas particulier du lemme \ref{relative-cycles-lemma-action-coherent} si
 l'on prend $\mathcal{F} = \mathcal{O}_Z$ et $\mathcal{F} = \mathcal{O}_W$. 
\end{proof} 

\begin{lemma} 
\label{relative-cycles-lemma-flat-pullback-as-action} 
Soient $(S, \delta)$ et $f : X \to Y$ comme ci-dessus. Supposons que 
$f$ soit plat de dimension relative $e$. Pour $\beta \in Z_r(Y)$
 , nous avons $f^*\beta = [X/X/Y]_e \cap \beta$ dans $Z_{e + r}(X)$. 
\end{lemma} 

\begin{proof} 
Il suffit de démontrer l'égalité localement sur $X$. Nous pouvons donc supposer 
$Y$ affine. Dans ce cas, $\beta$ est une combinaison linéaire finie à coefficients
 entiers de cycles premiers. Puisque les deux membres de l'égalité sont additifs en $\beta$
 , nous pouvons supposer $\beta = [W]$, où $W$ est un sous-schéma fermé intègre de 
$Y$ de dimension $\delta$ égale à $r$. Soit $f^{-1}(W) = W \times_Y X$ l'image
 inverse schématique de $W$ dans $X$. D'après
 le lemme \ref{relative-cycles-lemma-action-closed}, nous avons 
$[X/X/Y]_e \cap \beta = [W \times_Y X]_{r + e}$ et, d'après
 Homologie de Chow, définition \ref{chow-definition-flat-pullback}, 
nous avons $f^*\beta = [f^{-1}(W)]_{r + e}$. Nous obtenons ainsi l'égalité
 voulue. 
\end{proof} 

\begin{lemma} 
\label{relative-cycles-lemma-action-push-pull} 
Soit $(S, \delta)$ comme ci-dessus. Soit 
$$
\xymatrix{
X' \ar[r]_f \ar[d] & X \ar[d] \\
Y' \ar[r]^g & Y
}
$$
 un diagramme cartésien de schémas localement de type fini sur $S$
 tel que $g$ soit propre. Soient $r, e \geq 0$. Soit $\alpha$ une famille de 
$r$-cycles sur les fibres de $X/Y$. Soit $\beta' \in Z_e(Y')$.
 Alors $f_*(g^*\alpha \cap \beta') = \alpha \cap g_*\beta'$. 
\end{lemma} 

\begin{proof} 
Puisque nous démontrons une égalité de cycles sur $X$, nous pouvons travailler localement
 sur $Y$; voir le lemme \ref{relative-cycles-lemma-action-opens}. Nous pouvons donc supposer $Y$
 affine. Alors $Y'$ est quasi-compact. En particulier, $\beta'$
 est une combinaison linéaire finie de cycles premiers. 
Puisque $- \cap -$ est linéaire en la seconde variable
 (lemme \ref{relative-cycles-lemma-action-bilinear}), il suffit de 
démontrer l'égalité lorsque $\beta' = [Z']$ pour un sous-schéma fermé intègre 
$Z' \subset Y'$ de dimension $\delta$ égale à $e$. Posons $Z = g(Z')$. C'est 
un sous-schéma fermé intègre de $Y$ dont la dimension $\delta$ est $\leq e$.
 Pour simplifier, nous supposerons que $Z$ est de dimension $\delta$ égale
 à $e$ et laisserons au lecteur l'autre cas (qui est plus facile). 
Soient $y \in Z$ et $y' \in Z'$ les points génériques. 
Écrivons $\alpha_y = \sum m_j[V_j]$, où $V_j \subset X_y$
 sont des sous-schémas fermés intègres de dimension $r$. 

\medskip\noindent
 Supposons d'abord que $g$ soit une immersion fermée. Alors $g_*\beta' = [Z]$
 et $(g^*\alpha)_{y'} = \sum n_j[V_j]$; cette expression a un sens car 
$V_j$ est contenu dans le sous-schéma fermé $X'_{y'}$ de $X_y$.
 Dans ce cas, l'égalité est donc évidente : dans les deux membres, 
nous obtenons $\sum m_j[\overline{V}_j]$, où $\overline{V}_j$
 est l'adhérence de $V_j$ dans le sous-schéma fermé $X' \subset X$. 

\medskip\noindent
 Revenons au cas général où $\beta' = [Z']$ comme ci-dessus. 
Posons $W = Z \times_Y X$ et $W' = Z' \times_{Y'} X'$.
 Considérons les carrés cartésiens 
$$
\xymatrix{
W \ar[r] \ar[d] & X \ar[d] \\
Z \ar[r] & Y
}
\quad
\xymatrix{
W' \ar[r] \ar[d] & X' \ar[d] \\
Z' \ar[r] & Y'
}
\quad
\xymatrix{
W' \ar[r] \ar[d] & W \ar[d] \\
Z' \ar[r] & Z
}
$$
 Puisque le résultat est connu pour les deux premiers carrés d'après
 le paragraphe précédent, un argument formel montre qu'il suffit 
de démontrer le résultat pour le dernier carré et l'élément 
$\beta' = [Z'] \in Z_e(Z')$. Cela nous ramène au cas traité au 
paragraphe suivant. 

\medskip\noindent
 Supposons que $Y' \to Y$ soit un morphisme génériquement fini entre schémas intègres de dimension
 $\delta$ égale à $e$ et que $\beta' = [Y']$. Dans ce cas, 
$f_*(g^*\alpha \cap \beta')$ et $\alpha \cap g_*\beta'$ sont tous 
deux des cycles qui s'écrivent comme des sommes de cycles premiers dominant $Y$.
 Nous pouvons donc remplacer $Y$ par un sous-schéma ouvert non vide afin de 
vérifier l'égalité. Après un tel remplacement, nous pouvons supposer $g$ fini
 et plat, disons de degré $d \geq 1$. Bien entendu, cela signifie
 que $g_*\beta' = g_*[Y'] = d[Y]$. De plus, $\beta' = [Y'] = g^*[Y]$. Ainsi 
$$
f_*(g^*\alpha \cap \beta') =
f_*(g^*\alpha \cap g^*[Y]) =
f_*f^*(\alpha \cap [Y]) =
d (\alpha \cap [Y]) =
\alpha \cap g_*\beta'
$$
 comme souhaité. La deuxième égalité est le lemme \ref{relative-cycles-lemma-action-base-change}
 et la troisième est Homologie de Chow, lemme \ref{chow-lemma-finite-flat}. 
\end{proof}      







\section{Action sur les groupes de Chow} 
\label{relative-cycles-section-action-chow} 

\noindent
 Lorsque $\alpha$ est un $r$-cycle relatif, l'opération $\alpha \cap -$
 de la section \ref{relative-cycles-section-action} passe au quotient par l'équivalence rationnelle
 et définit une classe bivariante. 

\begin{lemma} 
\label{relative-cycles-lemma-closed-in-X-gysin} 
Soit $(S, \delta)$ comme dans la section \ref{relative-cycles-section-action}. 
Soit $f : X' \to X$ un morphisme propre de schémas 
localement de type fini sur $S$.
 Soit $(\mathcal{L}, s, i : D \to X)$ comme dans
 Homologie de Chow, définition \ref{chow-definition-gysin-homomorphism}. 
Formons le diagramme 
$$
\xymatrix{
D' \ar[d]_g \ar[r]_{i'} & X' \ar[d]^f \\
D \ar[r]^i & X
}
$$
 comme dans Homologie de Chow, remarque \ref{chow-remark-pullback-pairs}. 
Si $\mathcal{L}|_D \cong \mathcal{O}_D$, alors 
$i^*f_*\alpha' = g_*(i')^*\alpha'$ dans $Z_k(D)$
 pour tout $\alpha' \in Z_{k + 1}(X')$. 
\end{lemma} 

\begin{proof} 
L'énoncé a un sens puisque toutes les opérations sont définies au niveau des cycles ; pour les 
applications de Gysin, voir Homologie de Chow, remarque \ref{chow-remark-gysin-on-cycles}.
 Supposons
 que $\alpha = [W']$ pour un sous-schéma fermé intègre 
$W' \subset X'$. Posons $W = f(W') \subset X$. Si $W' \not \subset D'$,
 alors $W \not \subset D$ et nous voyons que 
$$
[W' \cap D']_k = \text{div}_{\mathcal{L}'|_{W'}}({s'|_{W'}})
\quad\text{et}\quad
[W \cap D]_k = \text{div}_{\mathcal{L}|_W}(s|_W)
$$
 et donc l'image du premier cycle par $f_*$ est égale au second d'après Homologie
 de Chow, lemme \ref{chow-lemma-equal-c1-as-cycles}. L'égalité est
 donc vraie au niveau des cycles. Si $W' \subset D'$, alors 
$W \subset D$ et les deux membres sont nuls par construction. 
\end{proof} 

\begin{lemma} 
\label{relative-cycles-lemma-action-gysin} 
Soit $(S, \delta)$ comme dans la section \ref{relative-cycles-section-action}. 
Soit $X \to Y$ un morphisme de schémas 
localement de type fini sur $S$. Soit $r \geq 0$ et soit 
$\alpha \in z(X/Y, r)$ un $r$-cycle relatif sur $X/Y$.
 Soit $(\mathcal{L}, s, i : D \to Y)$ comme dans
 Homologie de Chow, définition \ref{chow-definition-gysin-homomorphism}. 
Formons le diagramme cartésien 
$$
\xymatrix{
E \ar[d] \ar[r]_j & X \ar[d] \\
D \ar[r]^i & Y
}
$$
 Voir Homologie de Chow, remarque \ref{chow-remark-pullback-pairs}. 
Si $\mathcal{L}|_D \cong \mathcal{O}_D$, alors, pour $e \in \mathbf{Z}$
 , le diagramme 
$$
\xymatrix{
Z_e(D) \ar[rr]_{i^*\alpha \cap -} & &
Z_{e + r}(E) \\
Z_{e + 1}(Y) \ar[u]^{i^*} \ar[rr]^{\alpha \cap -} & &
Z_{r + e + 1}(X) \ar[u]_{j^*}
}
$$
 est commutatif, les flèches verticales $i^*$ et $j^*$ étant les 
applications de Gysin sur les cycles définies
 dans Homologie de Chow, remarque \ref{chow-remark-gysin-on-cycles}. 
\end{lemma} 

\begin{proof} 
Remarque préliminaire.
 Supposons que $g : Y' \to Y$ soit une enveloppe (Homologie de Chow, définition 
\ref{chow-definition-envelope}). 
Notons $D', i', E', j', X', \alpha'$ les changements de base de 
$D, i, E, j, X, \alpha$ par $g$, et notons $f : X' \to X$ la projection. 
Supposons le lemme vrai pour $D', i', E', j', X', Y', \alpha'$.
 Alors, si $\beta' \in Z_{e + 1}(Y')$, nous avons 
\begin{align*} 
i^*\alpha \cap i^*g_*\beta'
 & =
 i^*\alpha \cap f_*(i')^*\beta' \\
 & =
 f_*(f^*i^*\alpha \cap (i')^*\beta') \\ 
& =
 f_*((i')^*\alpha' \cap (i')^*\beta') \\ 
& =
 f_*((j')^*(\alpha' \cap \beta')) \\ 
& =
 j^*(f_*(f^*\alpha \cap \beta')) \\ 
& =
 j^*(\alpha \cap g_*\beta') 
\end{align*} 
Ici, la 
première égalité est le lemme \ref{relative-cycles-lemma-closed-in-X-gysin},
 la deuxième le lemme \ref{relative-cycles-lemma-action-push-pull},
 la troisième la définition de $\alpha'$,
 la quatrième résulte de l'hypothèse selon laquelle notre lemme est vrai pour 
$D', i', E', j', X', \alpha'$,
 la cinquième est le lemme \ref{relative-cycles-lemma-closed-in-X-gysin},
 et la sixième le lemme \ref{relative-cycles-lemma-action-push-pull}. Nous voyons
 ainsi que notre lemme est vrai pour tout cycle appartenant à l'image de 
$g_* : Z_{e + 1}(Y') \to Z_e(Y)$. Cependant, puisque $g$ est complètement 
décomposé, cette application est surjective ; nous en concluons que le lemme est vrai 
pour $D, i, E, j, X, Y, \alpha$. 

\medskip\noindent
 Soit $\beta \in Z_{e + 1}(Y)$. Nous devons montrer que 
$(D \to Y)^*\alpha \cap i^*\beta = j^*(\alpha \cap \beta)$
 au niveau des cycles sur $E$. Cette question est locale sur $E$ ; nous pouvons
 donc remplacer $X$ et $Y$ par des sous-schémas ouverts. (Nous utilisons ici le fait que la 
formation des opérateurs $i^*$, $j^*$, $\alpha \cap - $ et 
$(D \to Y)^*\alpha \cap -$ commute à la localisation. Cela est 
évident pour les applications de Gysin et résulte du
 lemme \ref{relative-cycles-lemma-action-opens} pour les autres.) Nous
 pouvons donc supposer que $X$ et $Y$ sont affines
 et nous ramener au cas étudié dans le paragraphe suivant. 

\medskip\noindent
 Supposons $X$ et $Y$ quasi-compacts. D'après le premier paragraphe de la preuve 
et le lemme \ref{relative-cycles-lemma-get-cycles}, nous pouvons en outre supposer que $\alpha$
 appartient à l'image de (\ref{relative-cycles-equation-cycle-classes}). Par linéarité 
des opérations considérées, nous pouvons supposer que $\alpha = [Z/X/Y]_r$
 pour un sous-schéma fermé $Z \subset X$, plat et de dimension 
relative $\leq r$ sur $Y$. De plus, puisque $Y$ est quasi-compact, le cycle 
$\beta$ est une combinaison linéaire finie de cycles premiers. Les opérations 
considérées étant linéaires, il suffit de
 démontrer l'égalité lorsque $\beta = [W]$ pour un sous-schéma fermé intègre 
$W \subset Y$ de $\delta$-dimension $e + 1$. 

\medskip\noindent
 Si $W \subset D$, alors, d'une part, $i^*[W] = 0$ et, d'autre
 part, $\alpha \cap [W]$ est supporté sur $E$, de sorte que 
$j^*(\alpha \cap [W]) = 0$. L'égalité est donc vraie dans ce cas. 

\medskip\noindent
 Supposons $W \not \subset D$. Alors $i^*[W] = [D \cap W]_e$.
 Notons que l'image inverse $i^*\alpha$ de $\alpha = [Z/X/Y]_r$
 par $i$ est $[(E \cap Z)/E/D]_r$ et que 
$(E \cap Z) = E \times_Y Z = D \times_Y Z$
 est plat sur $D$. Le lemme \ref{relative-cycles-lemma-action-closed}, appliqué
 deux fois, donne donc 
$$
i^*\alpha \cap i^*[W] =
[(E \cap Z) \times_D (D \cap W)]_{r + e} =
[E \cap (Z \times_Y W)]_{r + e} =
j^*(\alpha \cap [W])
$$
 comme voulu. 
\end{proof} 

\begin{proposition} 
\label{relative-cycles-proposition-get-bivariant-class} 
Soit $(S, \delta)$ comme dans la section \ref{relative-cycles-section-action}. Soit $X \to Y$
 un morphisme de schémas localement de type fini sur $S$. Soit $r \geq 0$
 et soit $\alpha \in z(X/Y, r)$ un $r$-cycle relatif sur $X/Y$.
 La règle qui, à tout morphisme $g : Y' \to Y$ localement de type fini
 et à tout $e \in \mathbf{Z}$, associe l'opération 
$$
g^*\alpha \cap - : Z_e(Y') \to Z_{r + e}(X')
$$
 où $X' = Y' \times_Y X$, passe au quotient par l'équivalence rationnelle 
et définit une classe bivariante $c(\alpha) \in A^{-r}(X \to Y)$. 
\end{proposition} 

\begin{proof} 
L'opération passe au quotient par l'équivalence rationnelle d'après 
le lemme \ref{relative-cycles-lemma-action-gysin} et 
Homologie de Chow, lemme \ref{chow-lemma-factors-through-rational-equivalence}. L'opération
 ainsi obtenue sur les groupes de Chow est une classe bivariante d'après 
Homologie de Chow, lemme \ref{chow-lemma-bivariant-weaker} et 
les
 lemmes \ref{relative-cycles-lemma-action-push-pull}, \ref{relative-cycles-lemma-action-base-change}, et 
\ref{relative-cycles-lemma-action-gysin}. 
\end{proof} 

\begin{remark} 
\label{relative-cycles-remark-characterize-relative-cycles} 
Soit $(S, \delta)$ comme dans la section \ref{relative-cycles-section-action}. Soit $X \to Y$
 un morphisme de schémas localement de type fini sur $S$. Soit $r \geq 0$.
 Soit $c$ une règle qui, à tout morphisme $g : Y' \to Y$ localement de type
 fini et à tout $e \in \mathbf{Z}$, associe une opération 
$$
c \cap - : Z_e(Y') \to Z_{r + e}(X')
$$
 compatible avec l'image directe propre et l'image inverse plate, ainsi qu'avec les applications de Gysin comme dans
 le lemme \ref{relative-cycles-lemma-action-gysin}. Nous affirmons qu'il existe un $r$-cycle relatif 
$\alpha$ sur $X/Y$ tel que $c \cap = g^*\alpha \cap -$ pour tout $g$ comme ci-dessus.
 Si nous en avons besoin, nous en donnerons ici un énoncé précis et une preuve détaillée. 
\end{remark} 









\section{Composition des familles de cycles sur les fibres} 
\label{relative-cycles-section-compose-families} 

\noindent
 Soient $X \to Y \to S$ des morphismes de schémas, tous deux localement de type fini. 
Soient $r, e \geq 0$. Soit $\alpha$ une famille de $r$-cycles 
sur les fibres de  $X/Y$  et  $\beta$  une famille de  $e$-cycles sur les
 fibres de  $Y/S$. Alors nous obtenons une famille 
de  $(r + e)$-cycles $\alpha \circ \beta$  sur les fibres de  $X/S$
  en posant 
$$
(\alpha \circ \beta)_s = (Y_s \to Y)^*\alpha \cap \beta_s
$$
 Plus précisément, l'expression $(Y_s \to Y)^*\alpha$ désigne 
le changement de base de $\alpha$  par  $Y_s \to Y$, qui est une famille de  $r$-cycles 
sur les fibres de  $X_s/Y_s$ ; l'opération  $- \cap -$
  a été définie et étudiée dans la section \ref{relative-cycles-section-action}\footnote{Précisément, 
nous prenons $s = \Spec(\kappa(s))$ comme schéma de base, avec $\delta(s) = 0$.} 

\begin{lemma} 
\label{relative-cycles-lemma-construction-bilinear} La 
construction ci-dessus est bilinéaire, c'est-à-dire que nous avons  
$(\alpha_1 + \alpha_2) \circ \beta \alpha_1 \circ \beta +
\alpha_1 \circ \beta$ et  $\alpha \circ (\beta_1 + \beta_2) =
\alpha \circ \beta_1 + \alpha \circ \beta_2$. 
\end{lemma} 

\begin{proof} 
Démonstration omise. Indication : sur les fibres, la construction est bilinéaire 
d'après le lemme \ref{relative-cycles-lemma-action-bilinear}. 
\end{proof} 

\begin{lemma} 
\label{relative-cycles-lemma-construction-opens} 
Si  $U \subset X$  et  $V \subset Y$  sont ouverts et  $f(U) \subset V$, alors  
$(\alpha \circ \beta)|_U$ est égal à  $\alpha|_U \circ \beta|_V$. 
\end{lemma} 

\begin{proof} 
Démonstration omise. Indication : sur les fibres, appliquer
 le lemme \ref{relative-cycles-lemma-action-opens}. 
\end{proof} 

\begin{lemma} 
\label{relative-cycles-lemma-construction-base-change} 
La formation de  $\alpha \circ \beta$  est compatible avec le changement de base. 
\end{lemma} 

\begin{proof} 
Soit  $g : S' \to S$  un morphisme de schémas. 
Notons  $X' \to Y'$  le changement de base de  $X \to Y$  par  $g$. Notons  
$\alpha'$  le changement de base de  $\alpha$  par rapport à  $Y' \to Y$.
 Notons  $\beta'$  le changement de base de  $\beta$  par rapport à  $S' \to S$.
 L'assertion signifie que $\alpha' \circ \beta'$  est le changement de base 
de  $\alpha \circ \beta$  par  $g : S' \to S$. 

\medskip\noindent
 Soit  $s' \in S'$  un point d'image  $s \in S$. Alors 
$$
(\alpha' \circ \beta')_{s'} = (Y'_{s'} \to Y')^*\alpha' \cap \beta'_{s'}
$$
 Nous observons que 
$$
(Y'_{s'} \to Y')^*\alpha' =
(Y'_{s'} \to Y')^*(Y' \to Y)^*\alpha =
(Y'_{s'} \to Y_s)^*(Y_s \to Y)^*\alpha
$$
 et que  $\beta'_{s'}$  est le changement de base de  $\beta_s$  par  
$s' = \Spec(\kappa(s')) \to \Spec(\kappa(s)) = s$.
 Par conséquent, le résultat découle du lemme \ref{relative-cycles-lemma-action-base-change} 
appliqué à $(Y_s \to Y)^*\alpha$, $\beta_s$,  
$X_s \to Y_s \to s$, et au changement de base par  $s' \to s$. 
\end{proof} 

\begin{lemma} 
\label{relative-cycles-lemma-construction-coherent} 
Soient  $f : X \to Y$  et  $Y \to S$  des morphismes de schémas, tous deux localement 
de type fini. Soient $r, e \geq 0$. Soit $\mathcal{F}$ un 
$\mathcal{O}_X$-module quasi-cohérent de type fini, avec  
$\dim(\text{Supp}(\mathcal{F}_y)) \leq r$  pour tout $y \in Y$. 
Soit  $\mathcal{G}$  un  $\mathcal{O}_Y$-module quasi-cohérent de 
type fini, avec  $\dim(\text{Supp}(\mathcal{G}_s)) \leq e$  pour tout  $s \in S$. 
Si  $\alpha = [\mathcal{F}/X/Y]_r$  et  $\beta = [\mathcal{G}/Y/S]_e$
 (Exemple  \ref{relative-cycles-example-family-associated-module}) et  $\mathcal{F}$
  est plat sur  $Y$, alors $\alpha \circ \beta =
[\mathcal{F} \otimes_{\mathcal{O}_X} f^*\mathcal{G}/X/S]_{r + e}$. 
\end{lemma} 

\begin{proof} 
Tout d'abord, nous observons que  $\mathcal{F} \otimes_{\mathcal{O}_X} f^*\mathcal{G}$
  est un  $\mathcal{O}_X$-module quasi-cohérent de type fini. 
Soit  $s \in S$. Observons que 
$$
(\mathcal{F} \otimes_{\mathcal{O}_X} f^*\mathcal{G})_s =
\mathcal{F}_s \otimes_{\mathcal{O}_{X_s}} f_s^*\mathcal{G}_s
$$
 par l'exactitude à droite du produit tensoriel. De plus, $\mathcal{F}_s$
 est plat sur $Y_s$, puisqu'il s'obtient par changement de base d'un module plat. Ainsi, l'égalité 
$(\alpha \circ \beta)_s =
[(\mathcal{F} \otimes_{\mathcal{O}_X} f^*\mathcal{G})_s]_{r + e}$
 découle du lemme \ref{relative-cycles-lemma-action-coherent}. 
\end{proof} 

\begin{lemma} 
\label{relative-cycles-lemma-construction-closed} 
Soient  $f : X \to Y$  et  $Y \to S$  des morphismes de schémas, tous deux localement 
de type fini. Soient $r, e \geq 0$. Soit $Z \subset X$ un sous-schéma
 fermé de dimension relative  $\leq r$  sur  $Y$. 
Soit  $W \subset Y$  un sous-schéma fermé de dimension relative  $\leq e$
  sur  $S$. Si $\alpha = [Z/X/Y]_r$  et  $\beta = [W/Y/S]_e$
  (Exemple  \ref{relative-cycles-example-family-associated-closed}) et  $Z$  est plat sur $Y$, 
alors  $\alpha \circ \beta = [Z \times_Y W/X/S]_{r + e}$. 
\end{lemma} 

\begin{proof} 
Il s'agit d'un cas particulier du lemme \ref{relative-cycles-lemma-construction-coherent} si 
nous prenons $\mathcal{F} = \mathcal{O}_Z$ et  $\mathcal{F} = \mathcal{O}_W$. 
\end{proof} 

\begin{lemma} 
\label{relative-cycles-lemma-flat-pullback-as-composition} Soient
  $f : X \to Y$  et  $Y \to S$  des morphismes de schémas, tous deux localement de 
type fini. Supposons que  $f$  soit plat de dimension relative  $e$. Soit  $\beta$  une
 famille de  $r$-cycles sur les fibres de  $Y/S$. Alors  
$f^*\beta = [X/X/Y]_e \circ \beta$
  comme familles de $(e + r)$-cycles sur  $X/S$. 
\end{lemma} 

\begin{proof} 
En déroulant les définitions et en utilisant la compatibilité de la formation de 
$[X/X/Y]_e$ avec le changement de base 
(lemme \ref{relative-cycles-lemma-family-associated-closed}), on conclut par le 
lemme \ref{relative-cycles-lemma-flat-pullback-as-action}. 
\end{proof} 

\begin{lemma} 
\label{relative-cycles-lemma-construction-push-pull} 
Soit  $S$  un schéma. Soit 
$$
\xymatrix{
X' \ar[r]_f \ar[d] & X \ar[d] \\
Y' \ar[r]^g & Y
}
$$
 un diagramme cartésien de schémas localement de type fini sur  $S$
  avec $g$ propre. Soient $r, e \geq 0$. Soit $\alpha$ une famille de 
$r$-cycles sur les fibres de  $X/Y$. Soit  $\beta'$  une famille de  
$e$-cycles sur les fibres de  $Y'/S$. Alors nous avons 
$f_*(g^*(\alpha) \circ \beta') = \alpha \circ g_*\beta'$. 
\end{lemma} 

\begin{proof}
 En déroulant les définitions, cela découle du 
lemme \ref{relative-cycles-lemma-action-push-pull}. 
\end{proof} 

\begin{lemma} 
\label{relative-cycles-lemma-construction-composition} 
Soit $(S, \delta)$ comme dans Homologie de Chow, situation \ref{chow-situation-setup}. 
Soient $X \to Y \to Z$  des morphismes de schémas localement de type fini sur  $S$. 
Soient $r, s, e \geq 0$. Alors 
$$
(\alpha \circ \beta) \cap \gamma = \alpha \cap (\beta \cap \gamma)
\quad\text{dans}\quad Z_{r + s + e}(X)
$$
 où  $\alpha$  est une famille de  $r$-cycles sur les fibres de $X/Y$,  
$\beta$  est une famille de  $s$-cycles sur les fibres de  $Y/Z$, et  $\gamma \in Z_e(Z)$. 
\end{lemma} 

\begin{proof} 
Puisque nous prouvons une égalité de cycles sur  $X$, nous pouvons travailler localement
 sur $Z$ ; voir le lemme \ref{relative-cycles-lemma-action-opens}. Ainsi nous pouvons supposer que  $Z$
  est affine. En particulier,  $\gamma$  est une combinaison linéaire finie de cycles 
premiers. Puisque $- \cap -$ est linéaire en la seconde variable
 (lemme \ref{relative-cycles-lemma-action-bilinear}), il suffit de
 prouver l'égalité lorsque $\gamma = [W]$ pour un sous-schéma fermé intègre 
$W \subset Z$ de $\delta$-dimension $e$. 

\medskip\noindent
 Soit $z \in W$ le point générique. Écrivons $\beta_z = \sum m_j[V_j]$
 dans $Z_s(Y_z)$. Alors  $\beta \cap \gamma$  est égal à  $\sum m_j[\overline{V}_j]$
  où  $\overline{V}_j \subset Y$  est un sous-schéma fermé intègre dont 
l'image par $Y \to Z$ est contenue dans $W$ et dont la fibre générique est $V_j$.
 Soit $y_j \in V_j$ le point générique. Nous le considérons aussi
 comme le point générique de $\overline{V}_j$ (envoyé sur $z$ dans $W$). 
Écrivons $\alpha_{y_j} = \sum n_{jk} [W_{jk}]$ dans $Z_r(X_{y_j})$.
 Alors $\alpha \cap (\beta \cap \gamma)$  est égal à 
$$
\sum m_j n_{jk} [\overline{W}_{jk}]
$$
 où $\overline{W}_{jk} \subset X$ est un sous-schéma fermé intègre dont 
l'image par $X \to Y$ est contenue dans $\overline{V}_j$ et dont la fibre générique est $W_{jk}$. 

\medskip\noindent
 D'autre part, considérons 
$$
(\alpha \circ \beta)_z = (Y_z \to Y)^*\alpha \cap \beta_z =
(Y_z \to Y)^*\alpha \cap (\sum m_j [V_j])
$$
 Par la construction de $- \cap -$, ceci est égal au cycle 
$$
\sum m_j n_{jk} [(\overline{W}_{jk})_z]
$$
 sur  $X_z$. Ainsi, par définition, nous obtenons 
$$
(\alpha \circ \beta) \cap [W] =
\sum m_j n_{jk} [\widetilde{W}_{jk}]
$$
 où $\widetilde{W}_{jk} \subset X$ est un sous-schéma fermé intègre dont
 l'image par $X \to Z$ est contenue dans $W$ et dont la fibre générique est $(\overline{W}_{jk})_z$.
 Il est clair que $\widetilde{W}_{jk} = \overline{W}_{jk}$
 , ce qui achève la preuve. 
\end{proof} 









\section{Composition des cycles relatifs} 
\label{relative-cycles-section-compose} 

\noindent
 Soit  $S$  un schéma localement noethérien. Soit  $X \to Y$
  un morphisme de schémas localement de type fini sur  $S$.
 Nous allons définir une application 
$$
z(X/Y, r) \otimes_\mathbf{Z} z(Y/S, e) \longrightarrow z(X/S, r + e),\quad
\alpha \otimes \beta \longmapsto \alpha \circ \beta
$$
 en utilisant la construction de la section \ref{relative-cycles-section-compose-families}. 
Nous savons déjà que la construction est bilinéaire
 (lemme \ref{relative-cycles-lemma-construction-bilinear}) ; 
pour obtenir la flèche affichée, il reste donc à établir le lemme suivant. 

\begin{lemma} 
\label{relative-cycles-lemma-well-defined} 
Si  $\alpha$  et  $\beta$  sont des cycles relatifs, alors il en est de même pour  $\alpha \circ \beta$. 
\end{lemma} 

\begin{proof}
 La formation de $\alpha \circ \beta$ est compatible avec le changement de base d'après
 le lemme \ref{relative-cycles-lemma-construction-base-change}. Ainsi, nous pouvons supposer que 
$S$ est le spectre d'un anneau de valuation discrète avec point générique  
$\eta$  et point fermé $0$  et nous devons montrer que  
$sp_{X/S}((\alpha \circ \beta)_\eta) = (\alpha \circ \beta)_0$.
 Puisque nous essayons de prouver une égalité de cycles, nous pouvons travailler
 localement sur  $Y$ et $X$ (on utilise les
 lemmes \ref{relative-cycles-lemma-construction-opens} et 
\ref{relative-cycles-lemma-specialization-flat-pullback} pour 
voir que les constructions commutent avec la restriction). Ainsi, nous
 pouvons supposer que $X$ et $Y$ sont affines. 
D'après le lemme \ref{relative-cycles-lemma-get-cycles}, 
nous pouvons trouver un morphisme propre complètement décomposé  
$g : Y' \to Y$  tel que $g^*\alpha$  soit dans l'image de 
(\ref{relative-cycles-equation-cycle-classes}). 

\medskip\noindent
 Comme le morphisme $g_\eta : Y'_\eta \to Y_\eta$
 est complètement décomposé, nous pouvons trouver $\beta'_\eta \in Z_e(Y'_\eta)$ tel que 
$\beta_\eta = \sum g_{\eta, *}\beta'_\eta$ ; voir
 Homologie de Chow, lemme \ref{chow-lemma-envelope}. 
Posons $\beta'_0 = sp_{Y'/S}(\beta'_\eta)$, de sorte que 
$\beta' = (\beta'_\eta, \beta'_0)$ soit un $e$-cycle relatif sur  $Y'/S$. Alors  
$g_*\beta'$  et  $\beta$ sont des $e$-cycles relatifs
 sur $Y/S$ (lemme \ref{relative-cycles-lemma-relative-cycle-functoriality})
 qui ont la même valeur en $\eta$ et sont donc égaux 
(lemme \ref{relative-cycles-lemma-uniqueness-extension}). Par linéarité 
(lemme \ref{relative-cycles-lemma-construction-bilinear}),
 il suffit de montrer que $\alpha \circ g_*\beta'$
 est un $(r + e)$-cycle relatif. 

\medskip\noindent
 Posons $X' = X \times_Y Y'$ et notons $f : X' \to X$ la projection. D'après
 le lemme \ref{relative-cycles-lemma-construction-push-pull}, on a 
$\alpha \circ g_*\beta' = f_*(g^*\alpha \circ \beta')$. 
D'après le lemme \ref{relative-cycles-lemma-relative-cycle-functoriality}, 
il suffit de montrer que  $g^*\alpha \circ \beta'$
  est un  $(r + e)$-cycle relatif. En utilisant le lemme \ref{relative-cycles-lemma-get-cycles-dvr} 
et la bilinéarité, cela nous ramène au cas discuté dans le paragraphe suivant. 

\medskip\noindent
 Supposons  $\alpha = [Z/X/Y]_r$  et  $\beta = [W/Y/S]$
  où  $Z \subset X$  est un sous-schéma fermé plat et de dimension
 relative  $\leq r$  sur  $Y$  et  $W \subset Y$  est un 
sous-schéma fermé plat et de dimension relative  $\leq e$  sur  $S$. 
D'après le lemme \ref{relative-cycles-lemma-construction-closed}, on a 
$$
\alpha \circ \beta = [Z \times_X W/X/S]_{r + e}
$$
 et  $Z \times_X W \subset X$  est un sous-schéma fermé plat sur  $S$
  de dimension relative $\leq r + e$. C'est un $(r + e)$-cycle relatif d'après
 le lemme \ref{relative-cycles-lemma-family-associated-closed-specialization}. 
\end{proof} 

\begin{lemma} 
\label{relative-cycles-lemma-composition-fundamental-cycles} Soient
  $f : X \to Y$  et  $g : Y \to S$  des morphismes de schémas. Supposons
 que $S$ soit localement noethérien, que $g$ soit localement de type fini 
et plat de dimension relative  $e \ge 0$, et  $f$  localement de type fini
 et plat de dimension relative  $r \geq 0$. Alors  
$[X/X/Y]_r \circ [Y/Y/S]_e = [X/X/S]_{r + e}$  dans  $z(X/S, r + e)$. 
\end{lemma} 

\begin{proof} 
C'est un cas particulier du lemme \ref{relative-cycles-lemma-construction-closed}. 
\end{proof} 






\section{Comparaison avec Suslin et Voevodsky} 
\label{relative-cycles-section-compare} 

\noindent
 Nous avons cherché à employer les mêmes notations que dans \cite{SV}, sauf que notre
 notation pour les cycles est tirée 
d'Homologie de Chow, section \ref{chow-section-cycles} et suivantes.
 Voici la comparaison : 
\begin{enumerate} 
\item Dans  \cite[section 3.1]{SV} , il existe une notion
 de « cycle relatif », de « cycle relatif de dimension  $r$ » et 
de « cycle relatif équidimensionnel de dimension $r$ ». 
Il n'existe pas de notion correspondante dans ce chapitre. Par conséquent, les groupes  
$Cycl(X/S, r)$,  $Cycl_{equi}(X/S, r)$,  
$PropCycl(X/S, r)$ et  $PropCycl_{equi}(X/S, r)$ 
n'ont pas de correspondants dans ce chapitre. 
\item En bas de  \cite[page 36]{SV} , les groupes  
$z(X/S, r)$,  $c(X/S, r)$,  $z_{equi}(X/S, r)$ et $c_{equi}(X/S, r)$
  sont définis. Ceux-ci correspondent à nos notions lorsque $S$ est séparé et
 noethérien, et $X \to S$ est séparé et de type fini. 
\item Dans  \cite{SV} , le symbole  $z(X/S, r)$  est parfois utilisé pour le 
préfaisceau $S' \mapsto z(S' \times_S X/S', r)$  sur la catégorie des schémas
 de type fini sur  $S$. De même pour  
$c(X/S, r)$,  $z_{equi}(X/S, r)$ et  $c_{equi}(X/S, r)$. 
\item Le changement de base, l'image inverse plate et l'image directe 
propre définis dans \cite{SV} coïncident avec les nôtres lorsque les deux définitions s'appliquent. 
\item Pour $\alpha \in z(X/S, r)$, l'opération 
$\alpha \cap - : Z_e(S) \to Z_{e + r}(X)$ définie dans 
la section \ref{relative-cycles-section-action} coïncide avec
 l'opération $Cor(\alpha, -)$ de \cite[section 3.7]{SV} lorsque
 les deux sont définies. 
\item Pour $X \to Y \to S$, la loi de composition 
$z(X/Y, r) \otimes_\mathbf{Z} z(Y/S, e) \longrightarrow z(X/S, r + e)$
 définie dans la section \ref{relative-cycles-section-compose}
 coïncide avec l'opération $Cor_{X/Y}(-, -)$ de 
\cite[corollaire 3.7.5]{SV}. 
\end{enumerate} 











\section{Cycles relatifs dans le cas non noethérien} 
\label{relative-cycles-section-non-noetherian} 

\noindent
 Nous recommandons vivement au lecteur de passer cette section. 

\medskip\noindent
 Soit  $f : X \to S$  un morphisme de schémas de présentation finie. 
Soit $r \geq 0$. On note  $Hilb(X/S, r)$  l'ensemble des sous-schémas fermés  
$Z \subset X$ tels que  $Z \to S$  soit plat, de présentation finie,
 et de dimension relative $\leq r$. Nous considérons l'homomorphisme de groupes 
\begin{equation} 
\label{relative-cycles-equation-cycle-classes-general} 
\begin{matrix} 
\text{groupe abélien libre} \\ 
\text{engendré par }Hilb(X/S, r) 
\end{matrix} 
\longrightarrow
\begin{matrix} 
\text{familles de }r\text{-cycles}\\ 
\text{sur les fibres de }X/S 
\end{matrix} 
\end{equation} 
qui envoie  $\sum n_i[Z_i]$  sur  $\sum n_i[Z_i/X/S]_r$. 

\begin{lemma} 
\label{relative-cycles-lemma-relative-r-cycle-general} 
Soit  $S$  un schéma quasi-compact et quasi-séparé. 
Soit  $f : X \to S$  un morphisme de présentation finie. 
Soit  $r \geq 0$  et soit  $\alpha$  une famille de $r$-cycles sur les fibres de $X/S$. 
Les assertions suivantes sont équivalentes : 
\begin{enumerate} 
\item il existe un diagramme cartésien 
$$
\xymatrix{
X \ar[r] \ar[d] & X_0 \ar[d] \\
S \ar[r] & S_0
}
$$
 où  $X_0 \to S_0$  est un morphisme de type fini de schémas noethériens
 et $\alpha_0 \in z(X_0/S_0, r)$  tel que  $\alpha$  soit le changement de base 
de  $\alpha_0$  par $S \to S_0$
\item il existe un morphisme propre complètement décomposé  $g : S' \to S$
  de présentation finie tel que  $g^*\alpha$  soit dans l'image de 
(\ref{relative-cycles-equation-cycle-classes-general}). 
\end{enumerate} 
\end{lemma} 

\begin{proof} 
Supposons donnés un diagramme et $\alpha_0 \in z(X_0/S_0, r)$ comme dans (1). D'après
 le lemme \ref{relative-cycles-lemma-get-cycles}, il existe un morphisme propre complètement 
décomposé  $g_0 : S'_0 \to S_0$  tel que $g_0^*\alpha_0$  soit 
dans l'image de (\ref{relative-cycles-equation-cycle-classes-general}). En effet, puisque  
$S'_0$ est noethérien, tout sous-schéma fermé de  $S'_0 \times_{S_0} X_0$
  est de présentation finie sur  $S'_0$. En posant  $S' = S \times_{S_0} S'_0$
  et en utilisant le changement de base par $S' \to S'_0$, on voit que (2) est vérifiée. 

\medskip\noindent
 Réciproquement, supposons (2) vérifiée. Choisissons un morphisme propre complètement 
décomposé  $g : S' \to S$  de présentation finie tel que  $g^*\alpha$ soit dans 
l'image de (\ref{relative-cycles-equation-cycle-classes-general}). Posons $X' = S' \times_S X$. 
Écrivons $g^*\alpha = \sum n_a [Z_a/X'/S']_r$  pour certains sous-schémas fermés  $Z_a \subset X'$
 , plats, de présentation finie et de dimension relative  
$\leq r$  sur  $S'$. 

\medskip\noindent
 Écrivons $S = \lim S_i$, où le système d'indices est filtrant, les morphismes de transition
 sont affines et les $S_i$ sont de type fini sur $\mathbf{Z}$ 
; voir Limites, proposition \ref{limits-proposition-approximate}.
 Pour $i$ suffisamment grand, il existe 
\begin{enumerate} 
\item un morphisme propre complètement décomposé  $g_i : S'_i \to S_i$
  dont le changement de base à  $S$  est  $g : S' \to S$, 
\item en posant $X'_i = S'_i \times_{S_i} X_i$
 , des sous-schémas fermés $Z_{ai} \subset X'_i$, plats et de
 dimension relative $\leq r$ sur $S'_i$, dont le changement de base à $S'$
 est $Z_a$. 
\end{enumerate} 
Pour ce faire, on utilise Limites, lemmes 
\ref{limits-lemma-descend-finite-presentation}, 
\ref{limits-lemma-descend-closed-immersion-finite-presentation}, 
\ref{limits-lemma-descend-flat-finite-presentation}, 
\ref{limits-lemma-descend-surjective}, 
\ref{limits-lemma-eventually-proper} et 
\ref{limits-lemma-limit-dimension}, ainsi 
que 
Compléments sur les morphismes, lemme 
\ref{more-morphisms-lemma-descend-cd}. 
Considérons $\alpha'_i = \sum n_a [Z_{ai}/X'_i/S_i]_r \in z(X'_i/S'_i, r)$.
 Le changement de base de  $\alpha'_i$, vu comme une famille de  $r$-cycles sur les fibres
 de $X'/S'$, coïncide avec le changement de base $g^*\alpha$ par construction. 

\medskip\noindent
 Posons $S''_i = S'_i \times_{S_i} S'_i$ et $X''_i = S''_i \times_{S_i} X_i$
 , puis $S'' = S' \times_S S'$ et $X'' = S'' \times_S X$.
 Notons $\text{pr}_1, \text{pr}_2 : S'' \to S'$ et 
$\text{pr}_1, \text{pr}_2 : S''_i \to S'_i$ les projections.
 Les $r$-cycles relatifs  $\text{pr}_1^*\alpha'_i$  et  $\text{pr}_1^*\alpha'_i$
  sur la base  $X''_i/S''_i$ deviennent après changement de base la même famille de $r$-cycles 
sur les fibres de $X''/S''$  parce que  
$\text{pr}_1^*g^*\alpha = \text{pr}_1^*g^*\alpha$. 
Par conséquent, l'image du morphisme $S'' \to S''_i$ est contenue dans $E =
\{s \in S''_i : (\text{pr}_1^*\alpha'_i)_s = (\text{pr}_1^*\alpha'_i)_s\}$. 
D'après le lemme \ref{relative-cycles-lemma-relative-cycles-equal}, il s'agit d'un sous-ensemble fermé.
 Comme $S'' = \lim_{i' \geq i} S''_{i'}$, le 
lemme 
\ref{limits-lemma-limit-contained-in-constructible} de Limites montre 
que, pour un certain $i' \geq i$, l'image du morphisme $S''_{i'} \to S''_i$
 est contenue dans $E$. Par conséquent, après avoir remplacé  $i$  par  $i'$, nous pouvons supposer que 
$\text{pr}_1^*\alpha'_i = \text{pr}_1^*\alpha'_i$. 
D'après le lemme \ref{relative-cycles-lemma-descend-family}, nous
 obtenons une famille unique  $\alpha_i$
  de  $r$-cycles sur les fibres de  $X_i/S_i$
 avec  $g_i^*\alpha_i = \alpha'_i$  (cela utilise le fait que  $S'_i \to S_i$
  est complètement décomposé). D'après
 le lemme \ref{relative-cycles-lemma-relative-cycles-h-descent},
 on a $\alpha_i \in z(X_i/S_i, r)$.
 L'unicité dans le lemme \ref{relative-cycles-lemma-descend-family} implique que le 
changement de base de $\alpha_i$ est $\alpha$ ; ainsi (1) est vérifiée. 
\end{proof} 

\noindent
 {\bf Discussion.} 
Si $f : X \to S$, $r$ et $\alpha$ sont comme dans
 le lemme \ref{relative-cycles-lemma-relative-r-cycle-general}, il est naturel de 
dire que $\alpha$ est un {\it $r$-cycle relatif sur $X/S$} si
 les conditions équivalentes
 (1) et (2) du lemme \ref{relative-cycles-lemma-relative-r-cycle-general} sont satisfaites.
 Cette définition possède de nombreuses propriétés utiles ; par 
exemple, elle est compatible avec la définition précédente lorsque $S$
 est noethérien, et la plupart des résultats de 
la section \ref{relative-cycles-section-families-specialization} se généralisent
 à ce cadre. 

\medskip\noindent
 Nous pouvons encore généraliser comme suit. 
Supposons $S$ arbitraire et $f : X \to S$ localement de présentation finie. 
Soient $r \geq 0$ et $\alpha$ une famille de $r$-cycles, notée $\alpha$, sur les fibres de 
$X/S$. Alors $\alpha$ est un {\it $r$-cycle relatif sur $X/S$} si, pour tous 
ouverts affines $U \subset X$ et $V \subset S$ tels que $f(U) \subset V$
 , la restriction $\alpha|_U$ est un $r$-cycle relatif sur $U/V$ au
 sens du paragraphe précédent. Là encore, beaucoup des résultats antérieurs
 se généralisent à ce cadre. 

\medskip\noindent
 Si nous avons besoin de ces généralisations, nous les énoncerons et les prouverons 
soigneusement ici. 























% OK, start here.
%

\chapter{Compléments de cohomologie étale}



\phantomsection
\label{more-etale-section-phantom}



\section{Introduction}
\label{more-etale-section-introduction}

\noindent
Ce chapitre est le deuxième d'une série de chapitres consacrés à la cohomologie
étale des schémas. Le premier chapitre se trouve dans
Cohomologie étale, section \ref{etale-cohomology-section-introduction}.

\medskip\noindent
La séparation avec le chapitre précédent est, en gros, la suivante : tout ce
qui concerne les « foncteurs image directe à support propre » (cohomologie à support propre
et son adjoint à droite), ainsi que leurs applications, figure dans ce chapitre.






\section{Prolongement des sections}
\label{more-etale-section-growing}

\noindent
Dans cette section, nous étudions des résultats du type suivant.

\begin{lemma}
\label{more-etale-lemma-section-support-in-locally-closed-pre}
Soit $X$ un schéma et soit $\mathcal{F}$ un faisceau abélien sur $X_\etale$.
Soit $\varphi : U' \to U$ un morphisme de $X_\etale$. Soit $Z' \subset U'$ un
sous-schéma fermé tel que $Z' \to U' \to U$ soit une immersion fermée
d'image $Z \subset U$. Il existe alors une bijection canonique
$$
\{s \in \mathcal{F}(U) \mid \text{Supp}(s) \subset Z\} =
\{s' \in \mathcal{F}(U') \mid \text{Supp}(s') \subset Z'\}
$$
qui est donnée par restriction si $\varphi^{-1}(Z) = Z'$.
\end{lemma}

\begin{proof}
Considérons le sous-schéma fermé $Z'' = \varphi^{-1}(Z)$ de $U'$.
Alors $Z' \subset Z''$ est fermé, puisque $Z'$ est fermé dans $U'$.
D'autre part, $Z' \to Z''$ est un morphisme étale
(comme morphisme entre schémas étales sur $Z$), et il est donc
ouvert. Ainsi $Z'' = Z' \amalg T$ pour un certain fermé $T$.
Le recouvrement ouvert $U' = (U' \setminus T) \cup (U' \setminus Z')$
montre que
$$
\{s' \in \mathcal{F}(U') \mid \text{Supp}(s') \subset Z'\} =
\{s' \in \mathcal{F}(U' \setminus T) \mid \text{Supp}(s') \subset Z'\}
$$
et le recouvrement étale $\{U' \setminus T \to U, U \setminus Z \to U\}$
montre que
$$
\{s \in \mathcal{F}(U) \mid \text{Supp}(s) \subset Z\} =
\{s' \in \mathcal{F}(U' \setminus T) \mid \text{Supp}(s') \subset Z'\}
$$
Ceci achève la démonstration.
\end{proof}

\begin{lemma}
\label{more-etale-lemma-section-support-in-locally-closed}
Soit $X$ un schéma et soit $Z \subset X$ un sous-schéma localement fermé.
Soit $\mathcal{F}$ un faisceau abélien sur $X_\etale$. Étant donnés des ouverts
$U, U' \subset X$ contenant $Z$ comme sous-schéma fermé,
il existe une bijection canonique
$$
\{s \in \mathcal{F}(U) \mid \text{Supp}(s) \subset Z\} =
\{s \in \mathcal{F}(U') \mid \text{Supp}(s) \subset Z\}
$$
qui est donnée par restriction si $U' \subset U$.
\end{lemma}

\begin{proof}
Puisque $Z$ est un sous-schéma fermé de $U \cap U'$, il suffit de
démontrer le lemme lorsque $U' \subset U$. Il s'agit alors d'un cas particulier
du lemme \ref{more-etale-lemma-section-support-in-locally-closed-pre}.
\end{proof}

\noindent
Introduisons une notation quelque peu inhabituelle qui nous sera
utile par la suite. Dans la situation du
lemme \ref{more-etale-lemma-section-support-in-locally-closed} ci-dessus, posons
$$
H_Z(\mathcal{F}) = \{s \in \mathcal{F}(U) \mid \text{Supp}(s) \subset Z\}
$$
où $U \subset X$ est un sous-schéma ouvert quelconque contenant $Z$ comme
sous-schéma fermé. Si ce manque de précision gêne le lecteur, il peut
choisir $U = X \setminus \partial Z$, où
$\partial Z = \overline{Z}\setminus Z$ est la « frontière » de $Z$ dans $X$.
Cependant, dans plusieurs arguments ci-dessous, la possibilité de choisir
différents ouverts jouera un rôle. Voici quelques propriétés de cette
construction :
\begin{enumerate}
\item
\label{more-etale-item-inclusion}
Si $Z \subset Z'$ sont des sous-schémas localement fermés de $X$ et si $Z$ est
fermé dans $Z'$, il existe une application injective naturelle
$$
H_Z(\mathcal{F}) \to H_{Z'}(\mathcal{F}).
$$
\item
\label{more-etale-item-pullback}
Si $f : Y \to X$ est un morphisme de schémas et si $Z \subset X$ est un
sous-schéma localement fermé, il existe une application naturelle
d'image inverse $f^* : H_Z(\mathcal{F}) \to H_{f^{-1}Z}(f^{-1}\mathcal{F})$.
\end{enumerate}
Il sera commode d'étendre notre notation à la situation suivante :
supposons donnés $W \in X_\etale$ et un sous-schéma localement fermé
$Z \subset W$. Nous noterons alors
$$
H_Z(\mathcal{F}) =
\{s \in \mathcal{F}(U) \mid \text{Supp}(s) \subset Z\} =
H_Z(\mathcal{F}|_{W_\etale})
$$
où $U \subset W$ est un sous-schéma ouvert quelconque contenant $Z$
comme sous-schéma fermé, exactement comme ci-dessus\footnote{En fait, le
lemme \ref{more-etale-lemma-section-support-in-locally-closed-pre}
montre que, si $Z$ est un schéma sur $X$ isomorphe à un sous-schéma
localement fermé d'un objet $W$ de $X_\etale$, alors
le choix de $W$ est sans importance.}.







\section{sections à support propre}
\label{more-etale-section-compact-support}

\noindent
Une référence pour cette section est \cite[Exposé XVII, section 6]{SGA4}.
Soit $f : X \to Y$ un morphisme de schémas séparé et
localement de type fini. Dans cette section, nous définissons un foncteur
$f_! : \textit{Ab}(X_\etale) \to \textit{Ab}(Y_\etale)$
en définissant $f_!\mathcal{F} \subset f_*\mathcal{F}$ comme
le sous-faisceau des sections dont le support est propre sur $Y$
(au sens approprié).

\medskip\noindent
Avertissement : le foncteur $f_!$ est le faisceau de cohomologie de degré zéro
d'un foncteur $Rf_!$ sur la catégorie dérivée (insérer une référence future),
mais $Rf_!$ n'est pas le foncteur dérivé de $f_!$.

\begin{lemma}
\label{more-etale-lemma-f-shriek-separated}
Soit $f : X \to Y$ un morphisme de schémas localement de type fini.
Soit $\mathcal{F}$ un faisceau abélien sur $X_\etale$. La règle
$$
Y_\etale \longrightarrow \textit{Ab},\quad
V \longmapsto \{s \in f_*\mathcal{F}(V) = \mathcal{F}(X_V) \mid
\text{Supp}(s) \subset X_V \text{ est propre sur }V\}
$$
définit un sous-faisceau abélien de $f_*\mathcal{F}$.
\end{lemma}

\noindent
Avertissement : ce faisceau n'est pas le « bon » si $f$ n'est pas séparé.

\begin{proof}
Rappelons que le support d'une section est fermé
(Cohomologie étale, lemme \ref{etale-cohomology-lemma-support-section-closed}) ;
les résultats de
Cohomologie des schémas, section \ref{coherent-section-proper-over-base},
s'appliquent donc. Le lemme ci-dessus et le lemme
\ref{coherent-lemma-union-closed-proper-over-base} de Cohomologie des schémas
montrent que notre sous-ensemble de $f_*\mathcal{F}(V)$ est un sous-groupe.
D'après le lemme de Cohomologie des schémas
\ref{coherent-lemma-base-change-closed-proper-over-base}
nous voyons que notre règle définit un sous-préfaisceau.
Enfin, supposons donnés $s \in f_*\mathcal{F}(V)$
et un recouvrement étale $\{V_i \to V\}$ tels que
le support de $s|_{V_i}$ soit propre sur $V_i$.
Le support de $s|_{V_i}$ est l'image réciproque du
support de $s|_V$ (utiliser la caractérisation du support
au moyen des fibres et le lemme
\ref{etale-cohomology-lemma-stalk-pullback} de Cohomologie étale).
Le support de $s$ est donc propre sur $V$ d'après le lemme
\ref{descent-lemma-descending-property-proper-over-base} de Descente.
Notre règle satisfait ainsi la condition de faisceau.
\end{proof}

\begin{lemma}
\label{more-etale-lemma-separated-etale-shriek}
Soit $j : U \to X$ un morphisme étale séparé. Soit $\mathcal{F}$
un faisceau abélien sur $U_\etale$. L'image de l'application injective
$j_!\mathcal{F} \to j_*\mathcal{F}$
du lemme de Cohomologie étale
\ref{etale-cohomology-lemma-shriek-into-star-separated-etale}
est le sous-faisceau du lemme \ref{more-etale-lemma-f-shriek-separated}.
\end{lemma}

\noindent
On pourrait aussi repousser ce lemme et le démontrer
à l'aide de la description des fibres des deux faisceaux.

\begin{proof}
La construction de $j_!\mathcal{F} \to j_*\mathcal{F}$ dans la démonstration du
lemme de Cohomologie étale
\ref{etale-cohomology-lemma-shriek-into-star-separated-etale}
s'obtient par la construction d'une application de préfaisceaux
$j_{p!}\mathcal{F} \to j_*\mathcal{F}$, dont l'image est manifestement
contenue dans le sous-faisceau du
lemme \ref{more-etale-lemma-f-shriek-separated}.
Puisque $j_!\mathcal{F}$ est le faisceau associé à
$j_{p!}\mathcal{F}$, nous en déduisons que l'image de
$j_!\mathcal{F} \to j_*\mathcal{F}$ est contenue dans
ce sous-faisceau. Réciproquement, soit $s \in j_*\mathcal{F}(V)$
de support $Z$ propre sur $V$. Alors $Z \to V$ est
fini, d'image fermée $Z' \subset V$ ; voir le lemme
\ref{more-morphisms-lemma-characterize-finite} de Compléments sur les morphismes.
La restriction de $s$ à $V \setminus Z'$ est nulle, et la section nulle
appartient à l'image de $j_!\mathcal{F} \to j_*\mathcal{F}$.
D'autre part, si $v \in Z'$, nous pouvons trouver
un voisinage étale
$(V', v') \to (V, v)$ tel que l'on ait une décomposition
$U_{V'} = W \amalg U'_1 \amalg \ldots \amalg U'_n$
en sous-schémas ouverts et fermés, où $U'_i \to V'$ est un isomorphisme
et $T_{V'} \subset U'_1 \amalg \ldots \amalg U'_n$ ; voir le lemme
\ref{etale-lemma-etale-etale-local-technical} de Morphismes étales.
En inversant les isomorphismes $U'_i \to V'$,
nous obtenons $n$ morphismes $\varphi'_i : V' \to U$
et des sections $s'_i$ sur $V'$, obtenues par image inverse de $s$.
La section $\sum (\varphi'_i, s'_i)$ de
$j_{p!}\mathcal{F}$ sur $V'$ — voir la formule de $j_{p!}\mathcal{F}(V')$
dans la démonstration du lemme de Cohomologie étale
\ref{etale-cohomology-lemma-shriek-into-star-separated-etale},
a pour image la restriction de $s$ à $V'$ par construction.
Nous concluons que, localement pour la topologie étale, $s$ appartient à l'image
de $j_!\mathcal{F} \to j_*\mathcal{F}$, ce qui achève la démonstration.
\end{proof}

\begin{definition}
\label{more-etale-definition-f-shriek-separated}
Soit $f : X \to Y$ un morphisme de schémas séparé (!) et
localement de type fini. Soit $\mathcal{F}$ un faisceau abélien sur
$X_\etale$. Le sous-faisceau $f_!\mathcal{F} \subset f_*\mathcal{F}$
construit au lemme \ref{more-etale-lemma-f-shriek-separated} est appelé
{\it l'image directe à support propre}.
\end{definition}

\noindent
Le lemme \ref{more-etale-lemma-separated-etale-shriek} montre que ceci ne contredit pas la
définition \ref{etale-cohomology-definition-extension-zero} de Cohomologie étale,
car les deux définitions coïncident lorsqu'elles s'appliquent. Vérifions-le.

\begin{lemma}
\label{more-etale-lemma-proper-f-shriek}
Soit $f : X \to Y$ un morphisme propre de schémas.
Alors $f_! = f_*$.
\end{lemma}

\begin{proof}
Cela résulte immédiatement de la construction de $f_!$.
\end{proof}

\noindent
Voici une observation très utile.

\begin{remark}[Covariance par rapport aux immersions ouvertes]
\label{more-etale-remark-covariance-f-shriek-separated}
Soit $f : X \to Y$ un morphisme de schémas séparé et
localement de type fini. Soit $\mathcal{F}$ un faisceau abélien sur $X_\etale$.
Soit $X' \subset X$ un sous-schéma ouvert. Notons $f' : X' \to Y$
la restriction de $f$. Il existe une application injective canonique
$$
f'_!(\mathcal{F}|_{X'}) \longrightarrow f_!\mathcal{F}
$$
En effet, soit $V \in Y_\etale$ et considérons une section
$s' \in f'_*(\mathcal{F}|_{X'})(V) = \mathcal{F}(X' \times_Y V)$
de support $Z'$ propre sur $V$. Alors $Z'$ est aussi fermé dans $X \times_Y V$ ;
voir le lemme de Cohomologie des schémas
\ref{coherent-lemma-functoriality-closed-proper-over-base}.
Il existe donc une unique section
$s \in \mathcal{F}(X \times_Y V) = f_*\mathcal{F}(V)$
dont la restriction à $X' \times_Y V$ est $s'$ et dont la restriction
à $X \times_Y V \setminus Z'$ est nulle ; voir le
lemme \ref{more-etale-lemma-section-support-in-locally-closed}. Cette construction est
compatible avec les applications de restriction et induit donc l'application voulue
de faisceaux $f'_!(\mathcal{F}|_{X'}) \to f_!\mathcal{F}$, qui est manifestement
injective. Par construction, nous obtenons un diagramme commutatif
$$
\xymatrix{
f'_!(\mathcal{F}|_{X'}) \ar[r] \ar[d] &
f_!\mathcal{F} \ar[d] \\
f'_*(\mathcal{F}|_{X'}) &
f_*\mathcal{F} \ar[l]
}
$$
fonctoriel en $\mathcal{F}$. Il est clair que, pour $X'' \subset X'$ ouvert
et $f'' = f|_{X''} : X'' \to Y$, la composée des applications canoniques
$f''_!\mathcal{F}|_{X''} \to f'_!\mathcal{F}|_{X'} \to f_!\mathcal{F}$
que nous venons de construire est l'application canonique
$f''_!\mathcal{F}|_{X''} \to f_!\mathcal{F}$.
\end{remark}

\begin{lemma}
\label{more-etale-lemma-compactify-f-shriek-separated}
Soit $Y$ un schéma. Soit $j : X \to \overline{X}$ une immersion
ouverte de schémas sur $Y$, avec $\overline{X}$ propre sur $Y$.
Notons $f : X \to Y$ et $\overline{f} : \overline{X} \to Y$
les morphismes structuraux. Pour $\mathcal{F} \in \textit{Ab}(X_\etale)$,
il existe un isomorphisme canonique (voir la démonstration)
$$
f_!\mathcal{F} \longrightarrow \overline{f}_!j_!\mathcal{F}
$$
Comme $\overline{f}_! = \overline{f}_*$ d'après le
lemme \ref{more-etale-lemma-proper-f-shriek}, nous obtenons
$\overline{f}_* \circ j_! = f_!$ comme foncteurs de
$\textit{Ab}(X_\etale) \to \textit{Ab}(Y_\etale)$.
\end{lemma}

\begin{proof}
Nous avons $(j_!\mathcal{F})|_X = \mathcal{F}$ ; voir le
lemme \ref{etale-cohomology-lemma-jshriek-open} de Cohomologie étale.
Ainsi, la flèche affichée est l'application injective
$f_!(\mathcal{G}|_X) \to \overline{f}_!\mathcal{G}$
de la remarque \ref{more-etale-remark-covariance-f-shriek-separated},
pour $\mathcal{G} = j_!\mathcal{F}$. La description explicite
de cette application montre qu'il suffit de prouver ceci : si $V \in Y_\etale$ et
$s \in \overline{f}_!\mathcal{G}(V) = \overline{f}_*\mathcal{G}(V) =
\mathcal{G}(\overline{X}_V)$
est une section, alors le support de $s$ est contenu dans l'ouvert
$X_V \subset \overline{X}_V$. Cela résulte immédiatement de ce que
les fibres de $\mathcal{G}$ sont nulles aux points géométriques
de $\overline{X} \setminus X$.
\end{proof}

\noindent
Nous voulons relier les fibres de $f_!\mathcal{F}$ aux sections à
support propre sur les fibres. Nous avons besoin pour cela d'une définition.

\begin{definition}
\label{more-etale-definition-compact-support}
Soit $X$ un schéma séparé localement de type fini sur un corps $k$.
Soit $\mathcal{F}$ un faisceau abélien sur $X_\etale$. Nous définissons
$H^0_c(X, \mathcal{F}) \subset H^0(X, \mathcal{F})$ comme l'ensemble
des sections dont le support est propre sur $k$. Les éléments de
$H^0_c(X, \mathcal{F})$ sont appelés {\it sections à support propre}.
\end{definition}

\noindent
Avertissement : cette définition n'est pas la « bonne » si $X$ n'est pas
séparé sur $k$.

\begin{lemma}
\label{more-etale-lemma-proper-compact-support}
Soit $X$ un schéma propre sur un corps $k$. Alors
$H^0_c(X, \mathcal{F}) = H^0(X, \mathcal{F})$.
\end{lemma}

\begin{proof}
Cela résulte immédiatement de la construction de $H^0_c$.
\end{proof}

\begin{remark}[Immersions ouvertes et sections à support propre]
\label{more-etale-remark-covariance-compact-support}
Soit $X$ un schéma séparé localement de type fini sur un corps $k$.
Soit $\mathcal{F}$ un faisceau abélien sur $X_\etale$.
Exactement comme dans la remarque \ref{more-etale-remark-covariance-f-shriek-separated},
pour tout ouvert $X' \subset X$, il existe une application injective
$$
H^0_c(X', \mathcal{F}|_{X'}) \longrightarrow H^0_c(X, \mathcal{F})
$$
et ces applications font de $H^0_c$ un « cofaisceau » sur le site de Zariski de $X$.
\end{remark}

\begin{lemma}
\label{more-etale-lemma-compactify-compact-support}
Soit $k$ un corps. Soit $j : X \to \overline{X}$ une immersion
ouverte de schémas sur $k$, avec $\overline{X}$ propre sur $k$.
Pour $\mathcal{F} \in \textit{Ab}(X_\etale)$,
il existe un isomorphisme canonique (voir la démonstration)
$$
H^0_c(X, \mathcal{F}) \longrightarrow
H^0_c(\overline{X}, j_!\mathcal{F}) =
H^0(\overline{X}, j_!\mathcal{F})
$$
où l'égalité de droite résulte du
lemme \ref{more-etale-lemma-proper-compact-support}.
\end{lemma}

\begin{proof}
Nous avons $(j_!\mathcal{F})|_X = \mathcal{F}$ ; voir le
lemme \ref{etale-cohomology-lemma-jshriek-open} de Cohomologie étale.
Ainsi, la flèche affichée est l'application injective
$H^0_c(X, \mathcal{G}|_X) \to H^0_c(\overline{X}, \mathcal{G})$
de la remarque \ref{more-etale-remark-covariance-compact-support},
pour $\mathcal{G} = j_!\mathcal{F}$. La description explicite
de cette application montre qu'il suffit de prouver ceci : si
$s \in H^0(\overline{X}, \mathcal{G})$ est une section, alors le support de
$s$ est contenu dans l'ouvert $X$. Cela résulte immédiatement de ce que
les fibres de $\mathcal{G}$ sont nulles aux points géométriques
de $\overline{X} \setminus X$.
\end{proof}

\begin{lemma}
\label{more-etale-lemma-stalk-f-shriek-separated}
Soit $f : X \to Y$ un morphisme de schémas séparé et
localement de type fini. Soit $\mathcal{F}$ un faisceau abélien sur
$X_\etale$. Il existe alors un isomorphisme canonique
$$
(f_!\mathcal{F})_{\overline{y}}
\longrightarrow
H^0_c(X_{\overline{y}}, \mathcal{F}|_{X_{\overline{y}}})
$$
pour tout point géométrique $\overline{y} : \Spec(k) \to Y$.
\end{lemma}

\begin{proof}
Rappelons que $(f_*\mathcal{F})_{\overline{y}} = \colim f_*\mathcal{F}(V)$,
où la limite inductive porte sur les voisinages étales $(V, \overline{v})$
de $\overline{y}$. Si $s \in f_*\mathcal{F}(V) = \mathcal{F}(X_V)$,
l'image inverse de $s$ est une section de $\mathcal{F}$ sur
$(X_V)_{\overline{v}} = X_{\overline{y}}$. Nous obtenons ainsi une application canonique
$$
c_{\overline{y}} :
(f_*\mathcal{F})_{\overline{y}}
\longrightarrow 
H^0(X_{\overline{y}}, \mathcal{F}|_{X_{\overline{y}}})
$$
Nous affirmons que cette application induit une bijection entre les sous-groupes
$(f_!\mathcal{F})_{\overline{y}}$ et
$H^0_c(X_{\overline{y}}, \mathcal{F}|_{X_{\overline{y}}})$.
Cette affirmation implique le lemme, mais elle est légèrement plus précise,
car elle décrit l'identification du lemme comme étant donnée
par l'image inverse des sections de $\mathcal{F}$ sur la fibre géométrique de $f$.

\medskip\noindent
Observons que tout élément
$s \in (f_!\mathcal{F})_{\overline{y}} \subset (f_*\mathcal{F})_{\overline{y}}$
est envoyé par $c_{\overline{y}}$ sur un élément de
$H^0_c(X_{\overline{y}}, \mathcal{F}|_{X_{\overline{y}}}) \subset
H^0(X_{\overline{y}}, \mathcal{F}|_{X_{\overline{y}}})$.
En effet, la formation du support d'une section
commute au changement de base et la propreté est préservée par
changement de base. Ceci construit déjà l'application de l'énoncé du lemme.
Pour montrer que c'est un isomorphisme, nous pouvons travailler localement
pour la topologie de Zariski sur $Y$ ; nous supposons donc $Y$ affine.

\medskip\noindent
Nous utiliserons ci-dessous l'observation suivante :
étant donné un sous-schéma ouvert $X' \subset X$,
si $f' = f|_{X'}$, nous obtenons un diagramme commutatif
$$
\xymatrix{
(f'_!(\mathcal{F}|_{X'}))_{\overline{y}} \ar[r] \ar[d] &
H^0_c(X'_{\overline{y}}, \mathcal{F}|_{X'_{\overline{y}}}) \ar[d] \\
(f_!\mathcal{F})_{\overline{y}} \ar[r] &
H^0_c(X_{\overline{y}}, \mathcal{F}|_{X_{\overline{y}}})
}
$$
où les flèches horizontales sont les applications construites ci-dessus et
les flèches verticales sont celles des
remarques \ref{more-etale-remark-covariance-f-shriek-separated} et
\ref{more-etale-remark-covariance-compact-support}.
En effet, étant donnés un voisinage étale $(V, \overline{v})$
de $\overline{y}$ et une section $s \in f_*\mathcal{F}(V) = \mathcal{F}(X_V)$
dont le support $Z$ est contenu dans $X'_V$ et propre sur $V$,
de sorte que $s$ détermine un élément de chacun des groupes
$(f'_!(\mathcal{F}|_{X'}))_{\overline{y}}$ et
$(f_!\mathcal{F})_{\overline{y}}$ qui se correspondent par
la flèche verticale du diagramme, les flèches horizontales envoient alors
ces éléments sur l'image inverse $s|_{X_{\overline{y}}}$ de $s$ dans
$X_{\overline{y}}$, dont le support $Z_{\overline{y}}$ est contenu dans
$X'_{\overline{y}}$ ; cette restriction détermine donc
une paire compatible d'éléments de
$H^0_c(X'_{\overline{y}}, \mathcal{F}|_{X'_{\overline{y}}})$
et de
$H^0_c(X_{\overline{y}}, \mathcal{F}|_{X_{\overline{y}}})$.

\medskip\noindent
Supposons que $s \in (f_!\mathcal{F})_{\overline{y}}$ ait une image nulle dans
$H^0_c(X_{\overline{y}}, \mathcal{F}|_{X_{\overline{y}}})$.
Disons que $s$ correspond à $s \in f_*\mathcal{F}(V) = \mathcal{F}(X_V)$,
de support $Z$ propre sur $V$. Nous pouvons supposer $V$ affine,
et $Z$ est alors quasi-compact. Nous pouvons choisir un ouvert quasi-compact
$X' \subset X$ contenant l'image de $Z$. Alors $Z$ est contenu dans
$X'_V$ et $s$ est donc l'image d'un élément
$s' \in f'_!(\mathcal{F}|_{X'})(V)$, où $f' = f|_{X'}$, comme au
paragraphe précédent. L'image de $s'$ dans
$H^0_c(X'_{\overline{y}}, \mathcal{F}|_{X'_{\overline{y}}})$
est donc nulle. Pour démontrer l'injectivité, nous pouvons remplacer $X$ par
$X'$, c'est-à-dire supposer $X$ quasi-compact. Nous démontrerons
ce cas ci-dessous.

\medskip\noindent
Supposons que
$t \in H^0_c(X_{\overline{y}}, \mathcal{F}|_{X_{\overline{y}}})$.
Alors le support de $t$ est contenu dans un sous-schéma ouvert
quasi-compact $W \subset X_{\overline{y}}$.
Nous pouvons donc trouver un sous-schéma ouvert quasi-compact
$X' \subset X$ tel que $X'_{\overline{y}}$ contienne $W$.
Il est alors clair que $t$ appartient à l'image
de l'application injective
$H^0_c(X'_{\overline{y}}, \mathcal{F}|_{X'_{\overline{y}}}) \to
H^0_c(X_{\overline{y}}, \mathcal{F}|_{X_{\overline{y}}})$.
Pour démontrer la surjectivité, nous pouvons donc remplacer $X$
par $X'$, c'est-à-dire supposer $X$ quasi-compact.
Nous démontrerons ce cas ci-dessous.

\medskip\noindent
Dans ce dernier paragraphe, nous démontrons le lemme lorsque
$X$ est quasi-compact et $Y$ affine. D'après le théorème
\ref{flat-theorem-nagata} de Compléments sur la platitude, il existe une compactification
$j : X \to \overline{X}$ sur $Y$. Posons $\mathcal{G} = j_!\mathcal{F}$,
de sorte que $\mathcal{F} = \mathcal{G}|_X$ d'après le
lemme \ref{etale-cohomology-lemma-jshriek-open} de Cohomologie étale.
La discussion ci-dessus donne un diagramme commutatif
$$
\xymatrix{
(f_!\mathcal{F})_{\overline{y}} \ar[r] \ar[d] &
H^0_c(X_{\overline{y}}, \mathcal{F}|_{X_{\overline{y}}}) \ar[d] \\
(\overline{f}_!\mathcal{G})_{\overline{y}} \ar[r] &
H^0_c(\overline{X}_{\overline{y}}, \mathcal{G}|_{\overline{X}_{\overline{y}}})
}
$$
D'après les lemmes \ref{more-etale-lemma-compactify-f-shriek-separated} et
\ref{more-etale-lemma-compactify-compact-support}, les applications verticales
sont des isomorphismes. Nous sommes ramenés au cas du morphisme propre
$\overline{X} \to Y$. Pour un morphisme propre, notre application
est un isomorphisme d'après les
lemmes \ref{more-etale-lemma-proper-f-shriek} et \ref{more-etale-lemma-proper-compact-support},
ainsi que le changement de base propre pour les images directes ; voir le
lemme de Cohomologie étale
\ref{etale-cohomology-lemma-proper-pushforward-stalk}.
\end{proof}

\begin{lemma}
\label{more-etale-lemma-base-change-f-shriek-separated}
Considérons un carré cartésien
$$
\xymatrix{
X' \ar[r]_{g'} \ar[d]_{f'} & X \ar[d]^f \\
Y' \ar[r]^g & Y
}
$$
de schémas, où $f$ est séparé et localement de type fini.
Pour tout faisceau abélien $\mathcal{F}$ sur $X_\etale$, nous avons
$f'_!(g')^{-1}\mathcal{F} = g^{-1}f_!\mathcal{F}$.
\end{lemma}

\begin{proof}
Plus généralement, il existe une application de changement de base
$g^{-1}f_*\mathcal{F} \to f'_*(g')^{-1}\mathcal{F}$ ; voir
Sites, section \ref{sites-section-pullback}.
Nous affirmons que cette application envoie $g^{-1}f_!\mathcal{F}$
dans le sous-faisceau $f'_!(g')^{-1}\mathcal{F}$
et induit l'isomorphisme du lemme.

\medskip\noindent
Choisissons un point géométrique $\overline{y}': \Spec(k) \to Y'$ et notons
$\overline{y} = g \circ \overline{y}'$ son image dans $Y$. Il existe un
diagramme commutatif
$$
\xymatrix{
(f_*\mathcal{F})_{\overline{y}} \ar[r] \ar[d] &
H^0(X_{\overline{y}}, \mathcal{F}|_{X_{\overline{y}}}) \ar[d] \\
(f'_*(g')^{-1}\mathcal{F})_{\overline{y}'} \ar[r] &
H^0(X'_{\overline{y}'}, (g')^{-1}\mathcal{F}|_{X'_{\overline{y}'}})
}
$$
où les applications horizontales sont celles utilisées dans la démonstration du
lemme \ref{more-etale-lemma-stalk-f-shriek-separated},
et les applications verticales sont celles de changement de base ci-dessus.
Le diagramme commute, car chacune des quatre applications
considérées s'obtient en prenant l'image inverse de sections locales par
un morphisme de schémas, et le diagramme sous-jacent de morphismes
de schémas commute. Puisque le diagramme de l'énoncé du lemme
est cartésien, nous avons $X'_{\overline{y}'} = X_{\overline{y}}$.
Le lemme \ref{more-etale-lemma-stalk-f-shriek-separated}
et sa démonstration donnent donc un diagramme commutatif
$$
\xymatrix{
(f_*\mathcal{F})_{\overline{y}} \ar[rrr] \ar[ddd] & & &
H^0(X_{\overline{y}}, \mathcal{F}|_{X_{\overline{y}}}) \ar[ddd] \\
& (f_!\mathcal{F})_{\overline{y}} \ar[r] \ar@{..>}[d] \ar[lu] &
H^0_c(X_{\overline{y}}, \mathcal{F}|_{X_{\overline{y}}}) \ar[d] \ar[ru] \\
& (f'_!(g')^{-1}\mathcal{F})_{\overline{y}'} \ar[r] \ar[ld] &
H^0_c(X'_{\overline{y}'}, (g')^{-1}\mathcal{F}|_{X'_{\overline{y}'}}) \ar[rd]\\
(f'_*(g')^{-1}\mathcal{F})_{\overline{y}'} \ar[rrr] & & &
H^0(X'_{\overline{y}'}, (g')^{-1}\mathcal{F}|_{X'_{\overline{y}'}})
}
$$
où les flèches horizontales du carré intérieur sont des isomorphismes
et les deux flèches verticales de droite sont des égalités. En outre, les
flèches sud-est, sud-ouest, nord-est et nord-ouest sont injectives. Il existe donc
une unique flèche pointillée bijective complétant le diagramme. Nous en concluons que
$g^{-1}f_!\mathcal{F} \subset g^{-1}f_*\mathcal{F} \to f'_*(g')^{-1}\mathcal{F}$
est envoyé dans le sous-faisceau
$f'_!(g')^{-1}\mathcal{F} \subset f'_*(g')^{-1}\mathcal{F}$
car il en est ainsi sur les fibres ; voir le théorème
\ref{etale-cohomology-theorem-exactness-stalks} de Cohomologie étale.
Le même théorème implique alors que l'application induite est un isomorphisme,
ce qui achève la démonstration.
\end{proof}

\begin{lemma}
\label{more-etale-lemma-f-shriek-composition}
Soient $f : X \to Y$ et $g : Y \to Z$ des morphismes composables de schémas,
séparés et localement de type fini. Soit $\mathcal{F}$ un faisceau abélien
sur $X_\etale$. Alors $g_!f_!\mathcal{F} = (g \circ f)_!\mathcal{F}$
comme sous-faisceaux de $(g \circ f)_*\mathcal{F}$.
\end{lemma}

\begin{proof}
Nous recommandons vivement au lecteur de le démontrer lui-même.
Soient $W \in Z_\etale$ et
$s \in (g \circ f)_*\mathcal{F}(W) = \mathcal{F}(X_W)$.
Notons $T \subset X_W$ le support de $s$ ; c'est un sous-ensemble
fermé. La section $s$ appartient à $(g \circ f)_!\mathcal{F}$
si et seulement si $T$ est propre sur $W$. Nous avons
$f_!\mathcal{F} \subset f_*\mathcal{F}$, et donc
$g_!f_!\mathcal{F} \subset g_!f_*\mathcal{F} \subset g_*f_*\mathcal{F}$.
D'autre part, $s$ appartient à $g_!f_!\mathcal{F}$ si et seulement
si (a) $T$ est propre sur $Y_W$ et (b) le support $T'$ de $s$,
considérée comme section de $f_!\mathcal{F}$, est propre sur $W$.
Si (a) est vérifiée, l'image de $T$ dans $Y_W$ est fermée et, puisque
$f_!\mathcal{F} \subset f_*\mathcal{F}$, nous voyons que
$T' \subset Y_W$ est l'image de $T$ (détails omis ; considérer les fibres).

\medskip\noindent
Il reste donc à montrer qu'un sous-ensemble fermé $T \subset X_W$
est propre sur $W$ si et seulement si $T$ est propre sur $Y_W$
et si l'image de $T$ dans $Y_W$ est propre sur $W$. Munissons $T$
de la structure réduite induite de sous-schéma fermé.
Si $T$ est propre sur $W$, alors $T \to Y_W$ est propre d'après le
lemme \ref{morphisms-lemma-image-proper-scheme-closed} de Morphismes,
et l'image de $T$ dans $Y_W$ est propre sur $W$ d'après le
lemme de Cohomologie des schémas
\ref{coherent-lemma-functoriality-closed-proper-over-base}.
Réciproquement, si $T$ est propre sur $Y_W$
et si l'image de $T$ dans $Y_W$ est propre sur $W$,
alors le morphisme $T \to W$ est propre comme composé
de morphismes propres (nous munissons ici l'image fermée de $T$
dans $Y_W$ de sa structure réduite induite, afin de ramener la
question à des morphismes de schémas) ; voir le
lemme \ref{morphisms-lemma-composition-proper} de Morphismes.
\end{proof}

\begin{remark}
\label{more-etale-remark-f-shriek-base-change-composition}
Les isomorphismes entre foncteurs
construits ci-dessus possèdent les deux propriétés suivantes :
\begin{enumerate}
\item Soient $f : X \to Y$, $g : Y \to Z$ et $h : Z \to T$ des morphismes
composables de schémas, séparés et localement de type fini.
Alors le diagramme
$$
\xymatrix{
(h \circ g \circ f)_! \ar[r] \ar[d] &
(h \circ g)_! \circ f_! \ar[d] \\
h_! \circ (g \circ f)_! \ar[r] &
h_! \circ g_! \circ f_!
}
$$
est commutatif, les flèches étant celles du lemme \ref{more-etale-lemma-f-shriek-composition}.
\item Supposons donné un diagramme de schémas
$$
\xymatrix{
X' \ar[d]_{f'} \ar[r]_c & X \ar[d]^f \\
Y' \ar[d]_{g'} \ar[r]_b & Y \ar[d]^g \\
Z' \ar[r]^a & Z
}
$$
dont les deux carrés sont cartésiens, et où $f$ et $g$ sont séparés et
localement de type fini. Alors le diagramme
$$
\xymatrix{
a^{-1} \circ (g \circ f)_! \ar[d] \ar[rr] & &
(g' \circ f')_! \circ c^{-1} \ar[d] \\
a^{-1} \circ g_! \circ f_! \ar[r] &
g'_! \circ b^{-1} \circ f_! \ar[r] &
g'_! \circ f'_! \circ c^{-1}
}
$$
est commutatif, les flèches horizontales étant celles du
lemme \ref{more-etale-lemma-base-change-f-shriek-separated}
et les autres celles du lemme \ref{more-etale-lemma-f-shriek-composition}.
\end{enumerate}
La partie (1) est vraie parce qu'il existe un diagramme commutatif
analogue pour les images directes. La partie (2) résulte de la compatibilité
très générale des applications de changement de base pour les images directes
(Sites, remarque \ref{sites-remark-compose-base-change})
et du fait que les isomorphismes des
lemmes \ref{more-etale-lemma-base-change-f-shriek-separated} et
\ref{more-etale-lemma-f-shriek-composition}
sont construits à l'aide des applications de changement de base correspondantes pour les images directes.
\end{remark}

\begin{lemma}
\label{more-etale-lemma-colim-f-shriek-separated}
Soit $f : X \to Y$ un morphisme de schémas séparé et
localement de type fini. Soit $X = \bigcup_{i \in I} X_i$ un
recouvrement ouvert tel que, pour tous $i, j \in I$, il existe $k$
avec $X_i \cup X_j \subset X_k$. Notons $f_i : X_i \to Y$
la restriction de $f$. Alors
$$
f_!\mathcal{F} = \colim_{i \in I} f_{i, !}(\mathcal{F}|_{X_i})
$$
fonctoriellement en $\mathcal{F} \in \textit{Ab}(X_\etale)$,
où les morphismes de transition sont ceux construits dans la
remarque \ref{more-etale-remark-covariance-f-shriek-separated}.
\end{lemma}

\begin{proof}
Il suffit de montrer que l'application canonique de
droite à gauche est bijective lorsqu'on l'évalue sur un objet
quasi-compact $V$ de $Y_\etale$.
La limite inductive du membre de droite est filtrante
et ses morphismes de transition sont injectifs.
Nous pouvons donc utiliser le
lemme \ref{sites-lemma-directed-colimits-sections} de Sites
pour évaluer cette limite. L'énoncé se ramène ainsi à observer
qu'un fermé $Z \subset X_V$ propre sur $V$ est quasi-compact
et est donc contenu dans $X_{i, V}$ pour un certain $i$.
\end{proof}

\begin{lemma}
\label{more-etale-lemma-f-shriek-separated-direct-sums}
Soit $f : X \to Y$ un morphisme de schémas séparé et
localement de type fini. Alors le foncteur $f_!$ commute aux sommes directes.
\end{lemma}

\begin{proof}
Soit $\mathcal{F} = \bigoplus \mathcal{F}_i$. Pour montrer que l'application
$\bigoplus f_!\mathcal{F}_i \to f_!\mathcal{F}$ est un isomorphisme,
il suffit de montrer que ces faisceaux ont les mêmes sections sur
un objet quasi-compact $V$ de $Y_\etale$. En remplaçant $Y$ par $V$,
il suffit de montrer que
$H^0(Y, f_!\mathcal{F}) \subset H^0(X, \mathcal{F})$
est égal à
$\bigoplus H^0(Y, f_!\mathcal{F}_i)
\subset \bigoplus H^0(X, \mathcal{F}_i)
\subset H^0(X, \bigoplus \mathcal{F}_i)$.
Dans ce cas, en écrivant $X$ comme réunion de ses ouverts quasi-compacts
et en utilisant le lemme \ref{more-etale-lemma-colim-f-shriek-separated},
nous nous ramenons au cas où $X$ est également quasi-compact.
Alors $H^0(X, \mathcal{F}) = \bigoplus H^0(X, \mathcal{F}_i)$
d'après le théorème \ref{etale-cohomology-theorem-colimit} de Cohomologie étale.
Le lecteur conclura aisément en considérant les supports des sections.
\end{proof}

\begin{lemma}
\label{more-etale-lemma-lqf-f-shriek-separated-colimits}
Soit $f : X \to Y$ un morphisme de schémas séparé et
localement quasi-fini. Alors
\begin{enumerate}
\item pour $\mathcal{F}$ dans $\textit{Ab}(X_\etale)$ et un point
géométrique $\overline{y} : \Spec(k) \to Y$, nous avons
$$
(f_!\mathcal{F})_{\overline{y}} =
\bigoplus\nolimits_{f(\overline{x}) = \overline{y}} \mathcal{F}_{\overline{x}}
$$
fonctoriellement en $\mathcal{F}$, et
\item le foncteur $f_!$ est exact.
\end{enumerate}
\end{lemma}

\begin{proof}
Le foncteur $f_!$ est exact à gauche par construction. L'exactitude à droite
se vérifie sur les fibres
(Cohomologie étale, théorème \ref{etale-cohomology-theorem-exactness-stalks}).
Il suffit donc de démontrer la partie (1).

\medskip\noindent
Soit $\overline{y} : \Spec(k) \to Y$ un point géométrique.
L'espace topologique sous-jacent au schéma $X_{\overline{y}}$
est discret
(Morphismes, lemme \ref{morphisms-lemma-locally-quasi-finite-fibres}),
et tous les corps résiduels en ses points sont égaux à $k$
(comme extensions finies de $k$). Par conséquent,
$\{\overline{x} : \Spec(k) \to X : f(\overline{x}) = \overline{y}\}$
est égal à l'ensemble des points de $X_{\overline{y}}$.
Le calcul de la fibre résulte donc du
lemme plus général \ref{more-etale-lemma-stalk-f-shriek-separated}.
\end{proof}









\section{sections à support fini}
\label{more-etale-section-finite-support}

\noindent
Dans cette section, nous étendons la construction de la
section \ref{more-etale-section-compact-support} aux morphismes localement
quasi-finis qui ne sont pas nécessairement séparés.

\medskip\noindent
Soit $f : X \to Y$ un morphisme de schémas localement quasi-fini.
Soit $\mathcal{F}$ un faisceau abélien sur $X_\etale$. Pour $V$ dans
$Y_\etale$, notons $X_V = X \times_Y V$ le changement de base. Nous allons
considérer le groupe des sommes formelles finies
\begin{equation}
\label{more-etale-equation-formal-sum}
s = \sum\nolimits_{i = 1, \ldots, n} (Z_i, s_i)
\end{equation}
où $Z_i \subset X_V$ est un sous-schéma localement fermé tel que le
morphisme $Z_i \to V$ soit fini\footnote{Puisque $f$ est localement quasi-fini,
le morphisme $Z_i \to V$ est fini si et seulement s'il est propre.}
et où $s_i \in H_{Z_i}(\mathcal{F})$. Ici, comme dans la
section \ref{more-etale-section-growing}, nous posons
$$
H_{Z_i}(\mathcal{F}) =
\{s_i \in \mathcal{F}(U_i) \mid \text{Supp}(s_i) \subset Z_i\}
$$
où $U_i \subset X_V$ est un sous-schéma ouvert contenant $Z_i$ comme
sous-schéma fermé. Nous considérons ces sommes formelles modulo les
relations suivantes :
\begin{enumerate}
\item
\label{more-etale-item-sum}
$(Z, s) + (Z, s') = (Z, s + s')$,
\item
\label{more-etale-item-sub}
$(Z, s) = (Z', s)$ si $Z \subset Z'$.
\end{enumerate}
La seconde relation a bien un sens : puisque $Z \to V$ est fini
et $Z' \to V$ séparé, l'inclusion $Z \to Z'$ est fermée et nous
pouvons utiliser l'application décrite en (\ref{more-etale-item-inclusion}).

\medskip\noindent
Notons $f_{p!}\mathcal{F}(V)$ le quotient du groupe abélien
des sommes formelles (\ref{more-etale-equation-formal-sum}) par ces relations.
La première relation dit que $f_{p!}\mathcal{F}(V)$ est un quotient
de la somme directe des groupes abéliens $H_Z(\mathcal{F})$,
où $Z \subset X_V$ parcourt les sous-schémas localement fermés finis sur $V$.
La seconde relation dit qu'il s'agit en réalité de la limite inductive
\begin{equation}
\label{more-etale-equation-colimit-definition}
f_{p!}\mathcal{F}(V) = \colim_Z H_Z(\mathcal{F})
\end{equation}
Cette formule donne une manière abstraite commode de considérer
notre construction.

\medskip\noindent
Observons ensuite que cette construction donne naturellement un préfaisceau
$f_{p!}\mathcal{F}$ de groupes abéliens sur $Y_\etale$.
En effet, un morphisme $V' \to V$ dans $Y_\etale$ fournit le morphisme de changement de base
$X_{V'} \to X_V$. Si $Z \subset X_V$ est un sous-schéma localement fermé
fini sur $V$, alors l'image inverse schématique $Z' \subset X_{V'}$
est finie sur $V'$. En outre, si $U \subset X_V$ est un ouvert tel
que $Z$ soit fermé dans $U$, alors l'image inverse $U' \subset X_{V'}$
est un ouvert tel que $Z'$ soit fermé dans $U'$. L'application de restriction
$\mathcal{F}(U) \to \mathcal{F}(U')$ de $\mathcal{F}$
envoie donc $H_Z(\mathcal{F})$ dans $H_{Z'}(\mathcal{F})$ ; c'est un cas
particulier de la fonctorialité décrite en (\ref{more-etale-item-pullback}) ci-dessus.
Ces applications sont manifestement compatibles avec les inclusions
$Z_1 \subset Z_2$ de tels sous-schémas localement fermés de $X_V$, et
nous obtenons une application
$$
f_{p!}\mathcal{F}(V) = \colim_Z H_Z(\mathcal{F})
\longrightarrow
\colim_{Z'} H_{Z'}(\mathcal{F}) =
f_{p!}\mathcal{F}(V')
$$
Ces applications font bien de $f_{p!}\mathcal{F}$ un préfaisceau de groupes
abéliens sur $Y_\etale$. Nous omettons les détails.

\medskip\noindent
Enfin, observons que la construction de $f_{p!}\mathcal{F}$
est fonctorielle en $\mathcal{F}$ dans $\textit{Ab}(X_\etale)$.
Ainsi, à tout morphisme localement quasi-fini $f : X \to Y$,
nous avons associé un foncteur
$$
f_{p!} :
\textit{Ab}(X_\etale)
\longrightarrow
\textit{PAb}(Y_\etale)
$$
de la catégorie des faisceaux abéliens sur $X_\etale$ vers celle
des préfaisceaux abéliens sur $Y_\etale$. Avant de définir $f_!$ par
passage au faisceau associé à ce foncteur, vérifions qu'il coïncide avec
la construction de la section \ref{more-etale-section-compact-support}
et avec celle de la
section \ref{etale-cohomology-section-extension-by-zero} de Cohomologie étale
lorsqu'elles s'appliquent toutes deux.

\begin{lemma}
\label{more-etale-lemma-finite-support-f-shriek-separated}
Soit $f : X \to Y$ un morphisme de schémas séparé et localement
quasi-fini. Fonctoriellement en $\mathcal{F} \in \textit{Ab}(X_\etale)$,
il existe un isomorphisme canonique (!)
$$
f_{p!}\mathcal{F} \longrightarrow f_!\mathcal{F}
$$
de préfaisceaux abéliens qui identifie le faisceau
$f_!\mathcal{F}$ de la définition \ref{more-etale-definition-f-shriek-separated}
au préfaisceau $f_{p!}\mathcal{F}$ construit ci-dessus.
\end{lemma}

\begin{proof}
Soit $V$ un objet de $Y_\etale$. Si $Z \subset X_V$ est localement fermé
et fini sur $V$, alors, puisque $f$ est séparé, le
morphisme $Z \to X_V$ est une immersion fermée. De plus, si
$Z_i$, $i = 1, \ldots, n$, sont des sous-schémas fermés de $X_V$ finis
sur $V$, alors $Z_1 \cup \ldots \cup Z_n$ (réunion schématique)
est un sous-schéma fermé fini sur $V$. Dans ce cas, la limite inductive
(\ref{more-etale-equation-colimit-definition}) qui définit $f_{p!}\mathcal{F}(V)$
est donc filtrante, et $f_{!p}\mathcal{F}(V)$ est simplement égal
à l'ensemble des sections de $\mathcal{F}(X_V)$ dont le support est fini sur $V$.
Tout fermé de $X_V$ propre sur $V$ étant
en fait fini sur $V$ (car $f$ est localement quasi-fini),
nous concluons que cet ensemble est égal à $f_!\mathcal{F}(V)$
par définition.
\end{proof}

\begin{lemma}
\label{more-etale-lemma-finite-support-stalk}
Soit $f : X \to Y$ un morphisme de schémas localement quasi-fini.
Soit $\overline{y} : \Spec(k) \to Y$ un point géométrique.
Fonctoriellement en $\mathcal{F}$ dans $\textit{Ab}(X_\etale)$, nous avons
$$
(f_{p!}\mathcal{F})_{\overline{y}} =
\bigoplus\nolimits_{f(\overline{x}) = \overline{y}} \mathcal{F}_{\overline{x}}
$$
\end{lemma}

\begin{proof}
Rappelons que la fibre en $\overline{y}$ d'un préfaisceau est définie par la
limite inductive usuelle sur les voisinages étales $(V, \overline{v})$
de $\overline{y}$ ; voir la définition
\ref{etale-cohomology-definition-stalk} de Cohomologie étale. Supposons donc que
$s = \sum_{i = 1, \ldots, n} (Z_i, s_i)$, comme dans
(\ref{more-etale-equation-formal-sum}), soit un élément de $f_{p!}\mathcal{F}(V)$,
où $(V, \overline{v})$ est un voisinage étale de $\overline{y}$.
Alors, puisque
$$
X_{\overline{y}} = (X_V)_{\overline{v}} \supset Z_{i, \overline{v}}
$$
et puisque $s_i$ est une section de $\mathcal{F}$ sur un voisinage ouvert
de $Z_i$ dans $X_V$, nous pouvons envoyer $s$ sur
$$
\sum\nolimits_{i = 1, \ldots, n} 
\sum\nolimits_{\overline{x} \in Z_{i, \overline{v}}}
\left(\text{classe de }s_i\text{ dans }\mathcal{F}_{\overline{x}}\right)
\quad\in\quad
\bigoplus\nolimits_{f(\overline{x}) = \overline{y}} \mathcal{F}_{\overline{x}}
$$
Nous omettons de vérifier la compatibilité avec les applications de restriction
et le fait que les relations (\ref{more-etale-item-sum}) $(Z, s) + (Z, s') - (Z, s + s')$
et (\ref{more-etale-item-sub}) $(Z, s) - (Z', s)$ si $Z \subset Z'$ sont envoyées sur zéro.
Nous obtenons ainsi une application
$$
(f_{p!}\mathcal{F})_{\overline{y}}
\longrightarrow
\bigoplus\nolimits_{f(\overline{x}) = \overline{y}} \mathcal{F}_{\overline{x}}
$$

\medskip\noindent
Montrons que cette flèche est surjective. Il suffit de choisir
$\overline{x}$ tel que $f(\overline{x}) = \overline{y}$ et de montrer qu'un
élément $s$ du facteur direct $\mathcal{F}_{\overline{x}}$ appartient à
l'image. Supposons que $s$ soit représenté par $s \in \mathcal{F}(U)$,
où $(U, \overline{u})$ est un voisinage étale de $\overline{x}$.
Puisque $f$ est localement quasi-fini, le morphisme $U \to Y$
l'est également. D'après le lemme de Compléments sur les morphismes
\ref{more-morphisms-lemma-etale-makes-quasi-finite-finite-multiple-points-var}
nous pouvons trouver un voisinage étale $(V, \overline{v})$ de
$\overline{y}$, un sous-schéma ouvert
$$
W \subset U \times_Y V,
$$
et un point géométrique $\overline{w}$ s'envoyant sur $\overline{u}$ et
$\overline{v}$, tels que $W \to V$ soit fini et que $\overline{w}$ soit l'unique
point géométrique de $W$ s'envoyant sur $\overline{v}$. (Nous omettons le passage
du langage des points géométriques employé ici à celui des points et des
extensions de corps résiduels utilisé dans
l'énoncé du lemme.) Le morphisme $W \to X_V = X \times_Y V$
est étale. Choisissons un voisinage ouvert affine $W' \subset X_V$
de l'image $\overline{w}'$ de $\overline{w}$. Puisque $\overline{w}$
est l'unique point de $W$ au-dessus de $\overline{v}$ et que $W \to V$
est fermé, nous pouvons, après avoir remplacé $V$ par un voisinage ouvert de $\overline{v}$,
supposer que $W \to X_V$ se factorise par $W'$. Alors $W \to W'$ est fini et
étale, et il existe un unique point géométrique $\overline{w}$ de $W$
au-dessus de $\overline{w}'$. Il s'ensuit que $W \to W'$ est une immersion ouverte
au-dessus d'un voisinage ouvert de $\overline{w}'$ dans $W'$ ; voir le
lemme \ref{etale-lemma-finite-etale-one-point} de Morphismes étales.
En rétrécissant $V$ et $W'$, nous pouvons supposer que $W \to W'$ est un isomorphisme.
Ainsi, $s$ peut être regardée comme une section $s'$ de $\mathcal{F}$ sur
le sous-schéma ouvert $W' \subset X_V$, qui est fini sur $V$.
Par définition, $(W', s')$ définit donc un élément de $j_{p!}\mathcal{F}(V)$
dont l'image est $s$, comme souhaité.

\medskip\noindent
Montrons que la flèche est injective. Pour cela, soit
$s = \sum_{i = 1, \ldots, n} (Z_i, s_i)$, comme dans (\ref{more-etale-equation-formal-sum}),
un élément de $f_{p!}\mathcal{F}(V)$, où $(V, \overline{v})$ est un
voisinage étale de $\overline{y}$. Supposons que l'image de $s$ soit nulle
par l'application construite ci-dessus. Après avoir tout d'abord remplacé
$(V, \overline{v})$ par un voisinage étale de lui-même,
nous pouvons supposer qu'il existe des décompositions
$Z_i = Z_{i, 1} \amalg \ldots \amalg Z_{i, m_i}$ en sous-schémas
ouverts et fermés, telles que chaque $Z_{i, j}$ possède exactement un point géométrique
au-dessus de $\overline{v}$. Dans la décomposition évidente en somme directe
$$
H_{Z_i}(\mathcal{F}) = \bigoplus H_{Z_{i, j}}(\mathcal{F})
$$
supposons que l'élément $s_i$ corresponde à $\sum s_{i, j}$. Les relations
(\ref{more-etale-item-sum}) et (\ref{more-etale-item-sub}) permettent de remplacer $s$ par
$\sum_{i = 1, \ldots, n} \sum_{j = 1, \ldots, m_i} (Z_{i, j}, s_{i, j})$.
Autrement dit, nous pouvons supposer que $Z_i$ possède un unique point géométrique
au-dessus de $\overline{v}$. Soient $\overline{x}_1, \ldots, \overline{x}_m$
les points géométriques de $X$ au-dessus de $\overline{y}$ correspondant aux
points géométriques de nos $Z_i$ au-dessus de $\overline{v}$ ; notons que, pour
un même $j \in \{1, \ldots, m\}$, plusieurs indices $i$ peuvent être tels que
$\overline{x}_j$ corresponde à un point de $Z_i$.
D'après le lemme de Compléments sur les morphismes
\ref{more-morphisms-lemma-etale-makes-quasi-finite-finite-multiple-points-var}
appliqué à $X_V \to V$, après avoir remplacé
$(V, \overline{v})$ par un voisinage étale de lui-même,
nous pouvons supposer qu'il existe des sous-schémas ouverts
$$
W_j \subset X \times_Y V,\quad j = 1, \ldots, m
$$
et un point géométrique $\overline{w}_j$ de $W_j$ s'envoyant sur $\overline{x}_j$ et
$\overline{v}$, tels que $W_j \to V$ soit fini et que $\overline{w}_j$ soit l'unique
point géométrique de $W_j$ s'envoyant sur $\overline{v}$.
Après avoir rétréci $V$, nous pouvons supposer que $Z_i \subset W_j$ pour un certain $j$,
et nous disposons de l'application $H_{Z_i}(\mathcal{F}) \to H_{W_j}(\mathcal{F})$.
La relation (\ref{more-etale-item-sub}) montre donc
que notre élément est équivalent à un élément de la forme
$$
\sum\nolimits_{j = 1, \ldots, m} (W_j, t_j)
$$
pour certains $t_j \in H_{W_j}(\mathcal{F})$. Il est clair que cet élément a
simplement pour image la classe de $t_j$ dans le facteur $\mathcal{F}_{\overline{x}_j}$.
Puisque l'image de $s$ est nulle, celle de $t_j$ dans
$\mathcal{F}_{\overline{x}_j}$ l'est aussi. Il s'ensuit que la restriction de $t_j$
est nulle sur un voisinage ouvert de $\overline{w}_j$ dans $W_j$ ; voir le
lemme \ref{etale-cohomology-lemma-zero-over-image} de Cohomologie étale.
En rétrécissant encore $V$, nous obtenons $t_j = 0$ pour tout $j$, comme souhaité.
\end{proof}

\begin{lemma}
\label{more-etale-lemma-finite-support-etale-shriek}
Soit $f = j : U \to X$ un morphisme étale de schémas. Notons $j_{p!}$
la construction donnée par la formule de Cohomologie étale
(\ref{etale-cohomology-equation-j-p-shriek})
et $f_{p!}$ la construction précédente. Fonctoriellement en
$\mathcal{F} \in \textit{Ab}(X_\etale)$, il existe une application canonique
$$
j_{p!}\mathcal{F} \longrightarrow f_{p!}\mathcal{F}
$$
de préfaisceaux abéliens qui identifie le faisceau
$j_!\mathcal{F} = (j_{p!}\mathcal{F})^\#$ de la définition
\ref{etale-cohomology-definition-extension-zero} de Cohomologie étale
à $(f_{p!}\mathcal{F})^\#$.
\end{lemma}

\begin{proof}
On pourra lire la démonstration du lemme de Cohomologie étale
\ref{etale-cohomology-lemma-shriek-into-star-separated-etale}
avant celle du présent lemme.
Soit $V$ un objet de $X_\etale$. Rappelons que
$$
j_{p!}\mathcal{F}(V) =
\bigoplus\nolimits_{\varphi : V \to U} \mathcal{F}(V \xrightarrow{\varphi} U)
$$
À $\varphi$ est associé un sous-schéma ouvert
$Z_\varphi \subset U_V = U \times_X V$, à savoir
l'image du graphe de $\varphi$. Grâce à $\varphi$,
nous obtenons un isomorphisme $V \to Z_\varphi$ sur $U$
et pouvons regarder un élément
$$
s_\varphi \in \mathcal{F}(V \xrightarrow{\varphi} U) =
\mathcal{F}(Z_\varphi) = H_{Z_\varphi}(\mathcal{F})
$$
comme une section de $\mathcal{F}$ sur $Z_{\varphi}$. Puisque
$Z_\varphi \subset U_V$ est ouvert, nous avons en fait
$H_{Z_\varphi}(\mathcal{F}) = \mathcal{F}(Z_\varphi)$
et pouvons regarder $s_\varphi$ comme un élément de $H_{Z_\varphi}(\mathcal{F})$.
Cela étant, notre application $j_{p!}\mathcal{F} \to f_{p!}\mathcal{F}$
est définie par la règle
$$
\sum\nolimits_{i = 1, \ldots, n} s_{\varphi_i}
\longmapsto
\sum\nolimits_{i = 1, \ldots, n} (Z_{\varphi_i}, s_{\varphi_i})
$$
où le membre de droite est une somme comme dans (\ref{more-etale-equation-formal-sum}).
Nous omettons de vérifier la compatibilité avec les applications de restriction
et la fonctorialité en $\mathcal{F}$.

\medskip\noindent
Pour achever la démonstration, affirmons que, pour tout point géométrique
$\overline{y} : \Spec(k) \to Y$, le diagramme suivant est commutatif :
$$
\xymatrix{
(j_{p!}\mathcal{F})_{\overline{y}} \ar[r] \ar[d] &
\bigoplus_{j(\overline{x}) = \overline{y}} \mathcal{F}_{\overline{x}}
\ar@{=}[d] \\
(f_{p!}\mathcal{F})_{\overline{y}} \ar[r] &
\bigoplus_{f(\overline{x}) = \overline{y}} \mathcal{F}_{\overline{x}}
}
$$
où la flèche horizontale supérieure est construite dans la démonstration de la
proposition de Cohomologie étale
\ref{etale-cohomology-proposition-describe-jshriek},
la flèche horizontale inférieure dans celle du
lemme \ref{more-etale-lemma-finite-support-stalk},
la flèche verticale droite est l'égalité évidente, et
la flèche verticale gauche est l'application induite sur les fibres par celle définie au
paragraphe précédent. L'affirmation résulte immédiatement
de la description explicite de toutes les flèches qui interviennent
ici et dans les références indiquées.
Puisque les flèches horizontales sont des isomorphismes,
la flèche verticale gauche est elle aussi un isomorphisme. Notre application
induit donc un isomorphisme après passage aux faisceaux associés, d'après le
théorème \ref{etale-cohomology-theorem-exactness-stalks} de Cohomologie étale.
\end{proof}

\begin{definition}
\label{more-etale-definition-f-shriek-lqf}
Soit $f : X \to Y$ un morphisme de schémas localement quasi-fini.
Nous appelons {\it image directe à support propre} le
foncteur
$$
f_! : \textit{Ab}(X_\etale) \longrightarrow \textit{Ab}(Y_\etale)
$$
défini par la formule $f_!\mathcal{F} = (f_{p!}\mathcal{F})^\#$,
c'est-à-dire que $f_!\mathcal{F}$ est le faisceau associé au préfaisceau
$f_{p!}\mathcal{F}$ construit ci-dessus.
\end{definition}

\noindent
D'après le lemme \ref{more-etale-lemma-finite-support-f-shriek-separated},
ceci ne contredit pas la définition \ref{more-etale-definition-f-shriek-separated}
(lorsque les deux définitions s'appliquent) et, d'après le
lemme \ref{more-etale-lemma-finite-support-etale-shriek},
ceci ne contredit pas non plus la
définition \ref{etale-cohomology-definition-extension-zero} de Cohomologie étale
(lorsque les deux définitions s'appliquent).

\begin{lemma}
\label{more-etale-lemma-lqf-f-shriek-stalk}
Soit $f : X \to Y$ un morphisme de schémas localement quasi-fini. Alors :
\begin{enumerate}
\item pour $\mathcal{F}$ dans $\textit{Ab}(X_\etale)$ et tout point
géométrique $\overline{y} : \Spec(k) \to Y$, on a
$$
(f_!\mathcal{F})_{\overline{y}} =
\bigoplus\nolimits_{f(\overline{x}) = \overline{y}} \mathcal{F}_{\overline{x}}
$$
fonctoriellement en $\mathcal{F}$ ;
\item le foncteur $f_! : \textit{Ab}(X_\etale) \to \textit{Ab}(Y_\etale)$
est exact et commute aux sommes directes.
\end{enumerate}
\end{lemma}

\begin{proof}
La formule pour les fibres découle immédiatement du
lemme \ref{more-etale-lemma-finite-support-stalk} (et lui est en fait équivalente).
L'exactitude du foncteur en résulte aussitôt,
puisque l'exactitude se vérifie sur les fibres ; voir le
théorème \ref{etale-cohomology-theorem-exactness-stalks} de Cohomologie étale.
\end{proof}

\begin{remark}[Covariance par rapport aux immersions ouvertes]
\label{more-etale-remark-covariance-lqf-f-shriek}
Soit $f : X \to Y$ un morphisme de schémas localement quasi-fini. Soit
$\mathcal{F}$ un faisceau abélien sur $X_\etale$.
Soit $X' \subset X$ un sous-schéma ouvert, et notons $f' : X' \to Y$
la restriction de $f$.
Il existe une application canonique
$$
f'_!(\mathcal{F}|_{X'}) \longrightarrow f_!\mathcal{F}
$$
Cette application est la faisceautisation d'une application canonique
$$
f'_{p!}(\mathcal{F}|_{X'}) \to f_{p!}\mathcal{F}
$$
construite comme suit. Soit $V \in Y_\etale$, et considérons une section
$s' = \sum_{i = 1, \ldots, n} (Z'_i, s'_i)$ comme dans
(\ref{more-etale-equation-formal-sum}), définissant un élément de
$f'_{p!}(\mathcal{F}|_{X'})(V)$.
Alors $Z'_i \subset X'_V$ peut aussi être regardé comme un sous-schéma localement fermé
de $X_V$, et l'on a $H_{Z'_i}(\mathcal{F}|_{X'}) = H_{Z'_i}(\mathcal{F})$.
Nous envoyons $s'$ sur exactement la même somme
$s = \sum_{i = 1, \ldots, n} (Z'_i, s'_i)$,
mais regardée maintenant comme un élément de $f_{p!}\mathcal{F}(V)$.
Nous omettons de vérifier que cette construction est compatible avec les
applications de restriction et fonctorielle en $\mathcal{F}$.
Cette construction possède les propriétés suivantes :
\begin{enumerate}
\item Les applications $f'_{p!}\mathcal{F}' \to f_{p!}\mathcal{F}$ et
$f'_!\mathcal{F}' \to f_!\mathcal{F}$ sont compatibles avec
la description des fibres donnée dans les lemmes
\ref{more-etale-lemma-finite-support-stalk} et \ref{more-etale-lemma-lqf-f-shriek-stalk}.
\item Si $f$ est séparé, l'application
$f'_{p!}\mathcal{F}' \to f_{p!}\mathcal{F}$ est celle
construite dans la remarque \ref{more-etale-remark-covariance-f-shriek-separated}
via l'isomorphisme du lemme \ref{more-etale-lemma-finite-support-f-shriek-separated}.
\item Si $X'' \subset X'$ est un autre ouvert, la composée de
$f''_{p!}(\mathcal{F}|_{X''}) \to f'_{p!}(\mathcal{F}|_{X'}) \to
f_{p!}\mathcal{F}$ est l'application
$f''_{p!}(\mathcal{F}|_{X''}) \to f_{p!}\mathcal{F}$ associée à
l'inclusion $X'' \subset X$. Après passage aux faisceaux associés, il en va
de même de
$f''_!(\mathcal{F}|_{X''}) \to f'_!(\mathcal{F}|_{X'}) \to f_!\mathcal{F}$.
\item L'application $f'_!\mathcal{F}' \to f_!\mathcal{F}$ est injective,
car cela se vérifie sur les fibres.
\end{enumerate}
Tous ces énoncés se démontrent aisément en représentant les éléments
par des sommes finies comme ci-dessus et en examinant leurs images.
\end{remark}

\begin{lemma}
\label{more-etale-lemma-lqf-colimit-f-shriek}
Soit $f : X \to Y$ un morphisme de schémas localement quasi-fini.
Soit $X = \bigcup_{i \in I} X_i$ un recouvrement ouvert. Il existe alors
un complexe exact
$$
\ldots \to
\bigoplus\nolimits_{i_0, i_1, i_2} f_{i_0i_1i_2, !}
\mathcal{F}|_{X_{i_0i_1i_2}} \to
\bigoplus\nolimits_{i_0, i_1} f_{i_0i_1, !} \mathcal{F}|_{X_{i_0i_1}} \to
\bigoplus\nolimits_{i_0} f_{i_0, !} \mathcal{F}|_{X_{i_0}}
\to f_!\mathcal{F} \to 0
$$
fonctoriel en $\mathcal{F} \in \textit{Ab}(X_\etale)$ ; voir les détails
dans la démonstration.
\end{lemma}

\begin{proof}
Comme d'habitude, posons $X_{i_0 \ldots i_p} = X_{i_0} \cap \ldots \cap X_{i_p}$
et notons $f_{i_0 \ldots i_p}$ la restriction de $f$ à
$X_{i_0 \ldots i_p}$. Les différentielles du complexe sont les applications
construites dans la remarque \ref{more-etale-remark-covariance-lqf-f-shriek},
avec les règles de signes du complexe de {\v C}ech.
L'exactitude découle aisément de la description des fibres donnée au
lemme \ref{more-etale-lemma-lqf-f-shriek-stalk}. Nous omettons les détails.
\end{proof}

\begin{remark}[Autre construction]
\label{more-etale-remark-alternative-lqf-f-shriek}
Le lemme \ref{more-etale-lemma-lqf-colimit-f-shriek}
donne une autre construction du foncteur $f_!$
pour les morphismes localement quasi-finis $f$.
En effet, étant donné un morphisme de schémas localement quasi-fini $f : X \to Y$,
nous pouvons choisir un recouvrement ouvert $X = \bigcup_{i \in I} X_i$
tel que chaque $f_i : X_i \to Y$ soit séparé ; il suffit par exemple de choisir
un recouvrement ouvert affine de $X$. Nous
pouvons alors définir $f_!\mathcal{F}$ comme le conoyau de l'avant-dernière flèche
du complexe du lemme, c'est-à-dire
$$
f_!\mathcal{F} = \Coker\left(
\bigoplus\nolimits_{i_0, i_1} f_{i_0i_1, !} \mathcal{F}|_{X_{i_0i_1}} \to
\bigoplus\nolimits_{i_0} f_{i_0, !} \mathcal{F}|_{X_{i_0}}
\right)
$$
où l'on peut utiliser les constructions de $f_{i_0, !}$ et de
$f_{i_0i_1, !}$ de la section \ref{more-etale-section-compact-support},
puisque les morphismes $f_{i_0}$ et $f_{i_0 i_1}$ sont séparés.
On peut alors calculer les fibres de $f_!$ (en se ramenant au cas séparé,
c'est-à-dire au lemme \ref{more-etale-lemma-lqf-f-shriek-separated-colimits})
et obtenir le résultat du lemme \ref{more-etale-lemma-lqf-f-shriek-stalk}.
Tous les autres résultats de la présente section s'en
déduisent également.
\end{remark}

\begin{remark}
\label{more-etale-remark-construct-map-presheaves-downstairs}
Soit $g : Y' \to Y$ un morphisme de schémas.
Pour tout préfaisceau abélien $\mathcal{G}'$ sur $Y'_\etale$, notons
$g_*\mathcal{G}'$ le préfaisceau $V \mapsto \mathcal{G}'(Y' \times_Y V)$.
Si $\alpha : \mathcal{G} \to g_*\mathcal{G}'$ est une application de préfaisceaux abéliens
sur $Y_\etale$, il existe une unique application
$\alpha^\# : \mathcal{G}^\# \to g_*((\mathcal{G}')^\#)$
de faisceaux abéliens sur $Y_\etale$ telle que le diagramme
$$
\xymatrix{
\mathcal{G} \ar[d] \ar[r]_\alpha & g_*\mathcal{G}' \ar[d] \\
\mathcal{G}^\# \ar[r]^-{\alpha^\#} & g_*((\mathcal{G}')^\#)
}
$$
soit commutatif, les flèches verticales provenant des applications canoniques
$\mathcal{G} \to \mathcal{G}^\#$ et $\mathcal{G}' \to (\mathcal{G}')^\#$. Si
$\alpha' : g^{-1}\mathcal{G}^\# \to (\mathcal{G}')^\#$
est l'application adjointe à $\alpha^\#$, alors, pour tout point géométrique
$\overline{y}' : \Spec(k) \to Y'$ d'image
$\overline{y} = g \circ \overline{y}'$ dans $Y$, l'application
$$
\alpha'_{\overline{y}'} :
\mathcal{G}_{\overline{y}} =
(\mathcal{G}^\#)_{\overline{y}} =
(g^{-1}\mathcal{G}^\#)_{\overline{y}'}
\longrightarrow
(\mathcal{G}')^\#_{\overline{y}'} =
\mathcal{G}'_{\overline{y}'}
$$
est définie comme suit : une section $s$ de $\mathcal{G}$ sur un voisinage étale
$(V, \overline{v})$ définit une classe dans la fibre, envoyée sur celle de la section
$\alpha(s)$ de $g_*\mathcal{G}'(V) = \mathcal{G}'(Y' \times_Y V)$
sur le voisinage étale $(Y' \times_Y V, (\overline{y}', \overline{v}))$
dans la fibre de $\mathcal{G}'$ en $\overline{y}'$.
\end{remark}

\begin{lemma}
\label{more-etale-lemma-lqf-base-change-f-shriek}
Considérons un carré cartésien
$$
\xymatrix{
X' \ar[r]_{g'} \ar[d]_{f'} & X \ar[d]^f \\
Y' \ar[r]^g & Y
}
$$
de schémas, où $f$ est localement quasi-fini. Il existe un isomorphisme
$g^{-1}f_!\mathcal{F} \to f'_!(g')^{-1}\mathcal{F}$, fonctoriel en
$\mathcal{F}$ dans $\textit{Ab}(X_\etale)$, qui est compatible avec
la description des fibres donnée au lemme \ref{more-etale-lemma-lqf-f-shriek-stalk}
(voir la démonstration pour l'énoncé précis).
\end{lemma}

\begin{proof}
Avec les conventions de la remarque \ref{more-etale-remark-construct-map-presheaves-downstairs},
nous allons construire explicitement une application
$$
c : f_{p!}\mathcal{F} \longrightarrow g_*f'_{p!}(g')^{-1}\mathcal{F}
$$
de préfaisceaux abéliens sur $Y_\etale$. D'après la discussion de la
remarque \ref{more-etale-remark-construct-map-presheaves-downstairs},
elle détermine une application canonique
$g^{-1}f_!\mathcal{F} \to f'_!(g')^{-1}\mathcal{F}$. Enfin, nous
montrerons que cette application induit des isomorphismes sur les fibres et conclurons par le
théorème \ref{etale-cohomology-theorem-exactness-stalks} de Cohomologie étale.

\medskip\noindent
Construction de l'application $c$. Soit $V \in Y_\etale$, et considérons une section
$s = \sum_{i = 1, \ldots, n} (Z_i, s_i)$ comme dans
(\ref{more-etale-equation-formal-sum}), définissant un élément de $f_{p!}\mathcal{F}(V)$.
La valeur de $g_*f'_{p!}(g')^{-1}\mathcal{F}$ en $V$ est
$f'_{p!}(g')^{-1}\mathcal{F}(V')$, où $V' = V \times_Y Y'$.
Notons $Z'_i \subset X'_{V'}$ le changement de base de $Z_i$ à $V'$.
D'après (\ref{more-etale-item-pullback}), il existe une application de changement de base
$H_{Z_i}(\mathcal{F}) \to H_{Z'_i}((g')^{-1}\mathcal{F})$.
Si $s'_i \in H_{Z'_i}((g')^{-1}\mathcal{F})$ désigne l'image de $s_i$
par changement de base, posons $c(s) = \sum_{i = 1, \ldots, n} (Z'_i, s'_i)$ comme dans
(\ref{more-etale-equation-formal-sum}), ce qui définit un élément de
$f'_{p!}(g')^{-1}\mathcal{F}(V')$. Nous omettons de vérifier
que cette construction est compatible avec les relations
(\ref{more-etale-item-sum}) et (\ref{more-etale-item-sub}), ainsi qu'avec les
applications de restriction. La construction est manifestement
fonctorielle en $\mathcal{F}$.

\medskip\noindent
Soit $\overline{y}' : \Spec(k) \to Y'$ un point géométrique d'image
$\overline{y} = g \circ \overline{y}'$ dans $Y$. Remarquons que
$X'_{\overline{y}'} = X_{\overline{y}}$ par transitivité des
produits fibrés. Ainsi, $g'$ induit une bijection
$\{f'(\overline{x}') = \overline{y}'\} \to \{f(\overline{x}) = \overline{y}\}$
et, si $\overline{x}'$ s'envoie sur $\overline{x}$, alors
$((g')^{-1}\mathcal{F})_{\overline{x}'} = \mathcal{F}_{\overline{x}}$
d'après le lemme \ref{etale-cohomology-lemma-stalk-pullback} de Cohomologie étale.
Affirmons maintenant que le diagramme
$$
\xymatrix{
(g^{-1}f_!\mathcal{F})_{\overline{y}'} \ar@{=}[r] \ar[d] &
(f_!\mathcal{F})_{\overline{y}} \ar[r] \ar[ld] &
\bigoplus\nolimits_{f(\overline{x}) = \overline{y}}
\mathcal{F}_{\overline{x}} \ar[d]
\\
(f'_!(g')^{-1}\mathcal{F})_{\overline{y}'} \ar[rr] & &
\bigoplus\nolimits_{f'(\overline{x}') = \overline{y}'}
(g')^{-1}\mathcal{F}_{\overline{x}'}
}
$$
est commutatif, les flèches horizontales étant celles de la démonstration du
lemme \ref{more-etale-lemma-finite-support-stalk}, tandis que la flèche verticale droite
est une égalité d'après ce qui précède. La flèche sud-ouest est
décrite dans la remarque \ref{more-etale-remark-construct-map-presheaves-downstairs}
comme l'application de changement de base, c'est-à-dire qu'elle est
simplement donnée par notre construction $c$ ci-dessus. La description
élémentaire de l'image d'une somme $\sum (Z_i, z_i)$ dans la
fibre en $\overline{x}$, donnée dans la démonstration du
lemme \ref{more-etale-lemma-finite-support-stalk}, montre immédiatement que le
diagramme commute. Ceci achève la démonstration du lemme.
\end{proof}

\begin{lemma}
\label{more-etale-lemma-lqf-separated-shriek-composition}
Soient $f' : X \to Y'$ et $g : Y' \to Y$ des morphismes composables de schémas,
où $f'$ et $f = g \circ f'$ sont localement quasi-finis, et $g$ séparé et
localement de type fini. Il existe alors un isomorphisme canonique de foncteurs
$g_! \circ f'_! = f_!$. Cet isomorphisme est compatible avec
\begin{enumerate}
\item[(a)] la covariance par rapport aux immersions ouvertes des
remarques \ref{more-etale-remark-covariance-f-shriek-separated} et
\ref{more-etale-remark-covariance-lqf-f-shriek},
\item[(b)] les isomorphismes de changement de base des
lemmes \ref{more-etale-lemma-lqf-base-change-f-shriek}
et \ref{more-etale-lemma-base-change-f-shriek-separated} ;
\item[(c)] il coïncide avec l'isomorphisme du lemme \ref{more-etale-lemma-f-shriek-composition}
via les identifications du lemme \ref{more-etale-lemma-finite-support-f-shriek-separated}
lorsque $f'$ est séparé.
\end{enumerate}
\end{lemma}

\begin{proof}
Soit $\mathcal{F}$ un faisceau abélien sur $X_\etale$. Avec les conventions de la
remarque \ref{more-etale-remark-construct-map-presheaves-downstairs}, nous allons
construire explicitement une application
$$
c : f_{p!}\mathcal{F} \longrightarrow g_*f'_{p!}\mathcal{F}
$$
de préfaisceaux abéliens sur $Y_\etale$. D'après la discussion de la
remarque \ref{more-etale-remark-construct-map-presheaves-downstairs},
elle détermine une application canonique
$c^\# : f_!\mathcal{F} \to g_*f'_!\mathcal{F}$.
Nous montrerons que l'image de $c^\#$ est contenue dans le sous-faisceau
$g_!f'_!\mathcal{F}$, obtenant ainsi une application
$c' : f_!\mathcal{F} \to g_!f'_!\mathcal{F}$. Nous démontrerons ensuite
(a), (b) et (c). Enfin, le point (b)
permettra de montrer que $c'$ est un isomorphisme.

\medskip\noindent
Construction de l'application $c$. Soit $V \in Y_\etale$, et
soit $s = \sum (Z_i, s_i)$ une somme comme dans (\ref{more-etale-equation-formal-sum}),
définissant un élément de $f_{p!}\mathcal{F}(V)$.
Rappelons que $Z_i \subset X_V = X \times_Y V$
est un sous-schéma localement fermé fini sur $V$.
En posant $V' = Y' \times_Y V$, on obtient $X_{V'} = X \times_{Y'} V' = X_V$.
Ainsi $Z_i \subset X_{V'}$ est localement fermé, et
$Z_i$ est fini sur $V'$ puisque $g$ est séparé
(lemme \ref{morphisms-lemma-finite-permanence} de Morphismes).
Nous pouvons donc poser $c(s) = \sum (Z_i, s_i)$, mais en la regardant maintenant
comme un élément de $f'_{p!}\mathcal{F}(V') = (g_*f'_{p!}\mathcal{F})(V)$.
La construction est manifestement compatible avec les relations
(\ref{more-etale-item-sum}) et (\ref{more-etale-item-sub}),
ainsi qu'avec les applications de restriction ; nous obtenons donc l'application $c$.

\medskip\noindent
Remarquons que, dans la discussion précédente, notre section $c(s) = \sum (Z_i, s_i)$ de
$f'_!\mathcal{F}$ sur $V'$ a une restriction nulle sur
$V' \setminus \Im(\coprod Z_i \to V')$. Puisque $\Im(\coprod Z_i \to V')$
est propre sur $V$ (par exemple d'après le lemme de Morphismes
\ref{morphisms-lemma-scheme-theoretic-image-is-proper}),
nous concluons que $c(s)$ définit une section de
$g_!f'_!\mathcal{F} \subset g_*f'_!\mathcal{F}$ sur $V$.
Puisque toute section locale de $f_!\mathcal{F}$ provient localement d'une
section locale de $f_{p!}\mathcal{F}$, nous concluons que l'image
de $c^\#$ est contenue dans $g_!f'_!\mathcal{F}$.
Nous obtenons donc une application induite $c' : f_!\mathcal{F} \to g_!f'_!\mathcal{F}$
par laquelle $c^\#$ se factorise, comme annoncé au premier paragraphe de la démonstration.

\medskip\noindent
Démonstration de (a). Soit $Y'_1 \subset Y'$ un sous-schéma ouvert,
et posons $X_1 = (f')^{-1}(W')$. Nous obtenons un diagramme
$$
\xymatrix{
X_1 \ar[d]_{f'_1} \ar[r]_a \ar@/_2em/[dd]_{f_1} &
X \ar[d]^{f'} \ar@/^2em/[dd]^f \\
Y'_1 \ar[d]_{g_1} \ar[r]_{b'} &
Y' \ar[d]^g \\
Y \ar@{=}[r] &
Y
}
$$
où les flèches horizontales sont des immersions ouvertes. Nous affirmons alors que
le diagramme
$$
\xymatrix{
f_{1, !}\mathcal{F}|_{X_1} \ar[r]_{c'_1} \ar[dd] &
g_{1, !}f'_{1, !}\mathcal{F}|_{X_1} \ar@{=}[d] \\
& g_{1, !}(f'_!\mathcal{F})|_{Y'_1} \ar[d] \\
f_!\mathcal{F} \ar[r]^{c'} &
g_!f'_!\mathcal{F} \ar[r] & g_*f'_!\mathcal{F}
}
$$
commute, la flèche verticale gauche étant celle de la
remarque \ref{more-etale-remark-covariance-lqf-f-shriek}, et
la flèche verticale droite celle de la remarque \ref{more-etale-remark-covariance-f-shriek-separated}.
Le signe d'égalité du diagramme vient de ce que $f'_1$
est la restriction de $f'$ à $Y'_1$ et que notre construction
de $f'_!$ est locale sur la base.
Enfin, pour démontrer la commutativité, choisissons un objet $V$ de
$Y_\etale$ et une somme formelle $s_1 = \sum (Z_{1, i}, s_{1, i})$ comme dans
(\ref{more-etale-equation-formal-sum}), définissant un élément de
$f_{1, p!}\mathcal{F}|_{X_1}(V)$. Cela signifie que
$Z_{1, i} \subset X_1 \times_Y V$ est localement fermé et fini sur $V$,
et que $s_{1, i} \in H_{Z_{1, i}}(\mathcal{F})$.
Suivons cette section
le long des applications qui interviennent : il suffit de montrer que nous
obtenons le même élément de
$g_*f'_!\mathcal{F}(V) = f'_!\mathcal{F}(Y' \times_Y V)$.
En parcourant chacun des deux côtés du diagramme, on voit immédiatement
que l'on obtient l'élément $\sum (Z_{1, i}, s_{1, i})$,
où $Z_{1, i}$ est maintenant regardé comme un sous-schéma localement fermé
de $X \times_{Y'} (Y' \times_Y V) = X \times_Y V$, fini sur
$Y' \times_Y V$.

\medskip\noindent
Démonstration de (b). Soit $b : Y_1 \to Y$ un morphisme de schémas. Formons le
diagramme commutatif
$$
\xymatrix{
X_1 \ar[d]_{f'_1} \ar[r]_a \ar@/_2em/[dd]_{f_1} &
X \ar[d]^{f'} \ar@/^2em/[dd]^f \\
Y'_1 \ar[d]_{g_1} \ar[r]_{b'} &
Y' \ar[d]^g \\
Y_1 \ar[r]^b &
Y
}
$$
dont les carrés sont cartésiens. Nous affirmons que notre construction est compatible
avec les applications de changement de base des lemmes \ref{more-etale-lemma-lqf-base-change-f-shriek}
et \ref{more-etale-lemma-base-change-f-shriek-separated}, c'est-à-dire
que le rectangle supérieur du diagramme
$$
\xymatrix{
b^{-1}f_!\mathcal{F} \ar[rr] \ar[d]_{b^{-1}c'} & &
f_{1, !}a^{-1}\mathcal{F} \ar[d]^{c_1'} \\
b^{-1}g_!f'_!\mathcal{F} \ar[r] \ar[d] &
g_{1, !}(b')^{-1}f'_!\mathcal{F} \ar[r] \ar[d] &
g_{1, !}f'_{1, !}a^{-1}\mathcal{F} \ar[d] \\
b^{-1}g_*f'_!\mathcal{F} \ar[r] &
g_{1, *}(b')^{-1}f'_!\mathcal{F} \ar[r] &
g_{1, *}f'_{1, !}a^{-1}\mathcal{F}
}
$$
commute. Cette vérification est élémentaire, et nous
invitons le lecteur à l'omettre. Comme les flèches de la ligne médiane
vers la ligne inférieure sont injectives, il suffit de montrer que
le contour extérieur du diagramme commute.
Pour cela, il suffit de prendre une section locale de
$b^{-1}f_!\mathcal{F}$ et de montrer que les deux chemins donnent la même
section locale de $g_{1, *}f'_{1, !}a^{-1}\mathcal{F}$.
Il suffit même de le vérifier
pour les sections locales qui sont les images inverses par $b$
d'une section $s = \sum (Z_i, s_i)$ de $f_{p!}\mathcal{F}(V)$
comme ci-dessus (puisque de telles images inverses engendrent le faisceau abélien
$b^{-1}f_!\mathcal{F}$). Notons $V_1$, $V'_1$ et $Z_{1, i}$
les changements de base de $V$, $V' = Y' \times_Y V$ et $Z_i$ par $Y_1 \to Y$.
Rappelons que $Z_i$ est un sous-schéma localement fermé de $X_V = X_{V'}$ ;
ainsi $Z_{1, i}$ est un sous-schéma localement fermé
de $(X_1)_{V_1} = (X_1)_{V'_1}$. Alors $b^{-1}c'$ envoie l'image inverse
de $s$ sur l'image inverse de la section locale $c(s) \sum (Z_i, s_i)$, regardée
comme un élément de $f'_{p!}\mathcal{F}(V') = (g_*f'_{p!}\mathcal{F})(V)$.
La composée des deux applications de changement de base inférieures
l'envoie simplement sur $\sum (Z_{i, 1}, s_{1, i})$, regardé comme un
élément de $f'_{1, p!}a^{-1}\mathcal{F}(V'_1) =
g_{1, *}f'_{1, p!}a^{-1}\mathcal{F}(V_1)$.
D'autre part, l'application de changement de base en haut du diagramme
envoie l'image inverse de $s$ sur $\sum (Z_{1, i}, s_{1, i})$, regardé
comme un élément de $f_{1, !}a^{-1}\mathcal{F}(V_1)$.
Enfin, par sa construction même, $c'_1$
envoie bien cet élément sur $\sum (Z_{i, 1}, s_{1, i})$, regardé comme un
élément de $f'_{1, p!}a^{-1}\mathcal{F}(V'_1) =
g_{1, *}f'_{1, p!}a^{-1}\mathcal{F}(V_1)$ ; la commutativité
est ainsi vérifiée.

\medskip\noindent
Démonstration de (c). Elle résulte de la comparaison des définitions
des deux applications ; nous omettons les détails.

\medskip\noindent
Pour achever la démonstration, il suffit de montrer que l'image inverse de
$c'$ par tout point géométrique $\overline{y} : \Spec(k) \to Y$
est un isomorphisme. En effet, prendre l'image inverse par $\overline{y}$
revient à prendre la fibre en $\overline{y}$
(remarque \ref{etale-cohomology-remark-stalk-pullback} de Cohomologie étale),
et nous pouvons donc invoquer le
théorème \ref{etale-cohomology-theorem-exactness-stalks} de Cohomologie étale.
Grâce à la compatibilité (b) que nous venons d'établir, nous 
concluons que nous pouvons supposer que $Y$ est le spectre de $k$,
et il reste à montrer que $c'$ est un isomorphisme.
Il suffit pour cela de montrer que l'application induite
$$
\bigoplus\nolimits_{x \in X} \mathcal{F}_x = H^0(Y, f_!\mathcal{F})
\longrightarrow
H^0(Y, g_!f'_!\mathcal{F}) = H^0_c(Y', f'_!\mathcal{F})
$$
est un isomorphisme. Les égalités résultent des
lemmes \ref{more-etale-lemma-lqf-f-shriek-stalk} et
\ref{more-etale-lemma-stalk-f-shriek-separated}.
Rappelons que $X$ est une somme disjointe de
spectres d'anneaux locaux artiniens de corps résiduel $k$ ; voir le
lemme \ref{varieties-lemma-algebraic-scheme-dim-0} de Variétés.
Comme les deux membres commutent aux sommes directes
(nous omettons les détails), nous pouvons supposer que $\mathcal{F}$ est un faisceau gratte-ciel
$x_*A$ concentré en un certain point $x \in X$.
Alors $f'_!\mathcal{F}$ est le faisceau gratte-ciel au point
$y'$, image de $x$ dans $Y$, d'après le lemme \ref{more-etale-lemma-lqf-f-shriek-stalk}.
Dans ce cas, il est manifeste que notre construction
donne l'application identité $A \to H^0_c(Y', y'_*A) = A$,
comme souhaité.
\end{proof}

\begin{lemma}
\label{more-etale-lemma-lqf-shriek-composition}
Soient $f : X \to Y$ et $g : Y \to Z$ des morphismes de schémas
localement quasi-finis et composables. Il existe alors un isomorphisme canonique de foncteurs
$$
(g \circ f)_! \longrightarrow g_! \circ f_!
$$
Ces isomorphismes vérifient les propriétés suivantes :
\begin{enumerate}
\item Si $f$ et $g$ sont séparés, l'isomorphisme coïncide avec celui du
lemme \ref{more-etale-lemma-f-shriek-composition}.
\item Si $g$ est séparé, l'isomorphisme coïncide avec celui du
lemme \ref{more-etale-lemma-lqf-separated-shriek-composition}.
\item Pour tout point géométrique $\overline{z} : \Spec(k) \to Z$, le diagramme
$$
\xymatrix{
((g \circ f)_!\mathcal{F})_{\overline{z}} \ar[d] \ar[rr] & &
\bigoplus\nolimits_{g(f(\overline{x})) = \overline{z}}
\mathcal{F}_{\overline{x}} \ar@{=}[d] \\
(g_!f_!\mathcal{F})_{\overline{z}} \ar[r] &
\bigoplus\nolimits_{g(\overline{y}) = \overline{z}}
(f_!\mathcal{F})_{\overline{y}} \ar[r] &
\bigoplus\nolimits_{g(f(\overline{x})) = \overline{z}}
\mathcal{F}_{\overline{x}}
}
$$
est commutatif, les flèches horizontales étant données par le
lemme \ref{more-etale-lemma-lqf-f-shriek-stalk}.
\item Soit $h : Z \to T$ un troisième morphisme de schémas
localement quasi-fini. Alors le diagramme
$$
\xymatrix{
(h \circ g \circ f)_! \ar[r] \ar[d] &
(h \circ g)_! \circ f_! \ar[d] \\
h_! \circ (g \circ f)_! \ar[r] &
h_! \circ g_! \circ f_!
}
$$
est commutatif.
\item Supposons donné un diagramme de schémas
$$
\xymatrix{
X' \ar[d]_{f'} \ar[r]_c & X \ar[d]^f \\
Y' \ar[d]_{g'} \ar[r]_b & Y \ar[d]^g \\
Z' \ar[r]^a & Z
}
$$
dont les deux carrés sont cartésiens et où $f$ et $g$ sont
localement quasi-finis. Alors le diagramme
$$
\xymatrix{
a^{-1} \circ (g \circ f)_! \ar[d] \ar[rr] & &
(g' \circ f')_! \circ c^{-1} \ar[d] \\
a^{-1} \circ g_! \circ f_! \ar[r] &
g'_! \circ b^{-1} \circ f_! \ar[r] &
g'_! \circ f'_! \circ c^{-1}
}
$$
est commutatif, les flèches horizontales étant celles du
lemme \ref{more-etale-lemma-lqf-base-change-f-shriek}.
\end{enumerate}
\end{lemma}

\begin{proof}
Si $f$ et $g$ sont séparés, il s'agit d'un cas particulier du
lemme \ref{more-etale-lemma-f-shriek-composition}.
Si $g$ est séparé, il s'agit d'un cas particulier du
lemme \ref{more-etale-lemma-lqf-separated-shriek-composition},
qui coïncide en outre avec le cas où $f$ et $g$ sont séparés.

\medskip\noindent
Construction dans le cas général. Choisissons un recouvrement ouvert $Y = \bigcup Y_i$
tel que la restriction $g_i : Y_i \to Z$ de $g$ soit séparée.
Posons $X_i = f^{-1}(Y_i)$ et notons $f_i : X_i \to Y_i$ la restriction
de $f$. Notons aussi $h = g \circ f$ et $h_i : X_i \to Z$ la restriction
de $h$. Considérons le diagramme suivant :
$$
\xymatrix{
\bigoplus\nolimits_{i_0, i_1}
h_{i_0i_1, !}\mathcal{F}|_{X_{i_0i_1}} \ar[r] \ar[d] &
\bigoplus\nolimits_{i_0} h_{i_0, !}\mathcal{F}|_{X_{i_0}} \ar[r] \ar[d] &
h_!\mathcal{F} \ar[r] \ar@{..>}[dd] &
0 \\
\bigoplus\nolimits_{i_0, i_1}
g_{i_0i_1, !} f_{i_0i_1, !}\mathcal{F}|_{X_{i_0i_1}} \ar[r] \ar[d] &
\bigoplus\nolimits_{i_0}
g_{i_0, !} f_{i_0, !}\mathcal{F}|_{X_{i_0}} \ar[d] \\
\bigoplus\nolimits_{i_0, i_1}
g_{i_0i_1, !} (f_!\mathcal{F})|_{Y_{i_0i_1}} \ar[r] &
\bigoplus\nolimits_{i_0}
g_{i_0, !} (f_!\mathcal{F})|_{Y_{i_0}} \ar[r] &
g_!f_!\mathcal{F} \ar[r] &
0
}
$$
D'après le lemme \ref{more-etale-lemma-lqf-colimit-f-shriek}, les lignes supérieure et inférieure
du diagramme sont exactes. D'après le lemme \ref{more-etale-lemma-lqf-separated-shriek-composition},
le carré supérieur gauche commute. Les flèches verticales du
carré inférieur gauche proviennent des égalités
$(f_!\mathcal{F})|_{Y_{i_0i_1}} = f_{i_0i_1, !}\mathcal{F}|_{X_{i_0i_1}}$ et
$(f_!\mathcal{F})|_{Y_{i_0}} = f_{i_0, !}\mathcal{F}|_{X_{i_0}}$
puisque la construction de $f_!$ est locale sur la base. De plus, ces
égalités sont bien entendu compatibles avec les identifications
$((f_!\mathcal{F})|_{Y_{i_0}})|_{Y_{i_0i_1}} =
(f_!\mathcal{F})|_{Y_{i_0i_1}}$ et
$(f_{i_0, !}\mathcal{F}|_{X_{i_0}})|_{Y_{i_0i_1}} =
f_{i_0i_1, !}\mathcal{F}|_{X_{i_0i_1}}$
qui servent (avec la covariance par rapport aux immersions ouvertes
pour $Y_{i_0i_1} \subset Y_{i_0}$)
à définir les flèches horizontales du carré inférieur gauche.
Ce carré commute donc lui aussi.
Nous en concluons qu'il existe une unique
flèche pointillée comme indiqué dans le diagramme et
que cette flèche est en outre un isomorphisme.

\medskip\noindent
Démonstration des propriétés (1) -- (5). Fixons le recouvrement ouvert $Y = \bigcup Y_i$.
Remarquons que, si $Y \to Z$ est séparé, nous obtenons une flèche pointillée
s'insérant dans le grand diagramme ci-dessus grâce à l'application du
lemme \ref{more-etale-lemma-lqf-separated-shriek-composition}
(par les propriétés mêmes de ce lemme).
Ceci démontre (2), puis (1) par compatibilité des
applications du lemme \ref{more-etale-lemma-lqf-separated-shriek-composition}
et du lemme \ref{more-etale-lemma-f-shriek-composition}.
Ensuite, pour tout schéma $Z'$ sur $Z$, nous obtenons la compatibilité de (5)
pour l'application $(g' \circ f')_! \to g'_! \circ f'_!$
construite au moyen du recouvrement ouvert $Y' = \bigcup b^{-1}(Y_i)$.
Cela résulte clairement de la compatibilité correspondante des applications
construites au lemme \ref{more-etale-lemma-lqf-separated-shriek-composition}.
En particulier, nous pouvons considérer un point géométrique
$\overline{z} : \Spec(k) \to Z$. Comme les applications
$X_{\overline{z}} \to Y_{\overline{z}} \to \Spec(k)$
sont séparées, le changement de base de
$(g \circ f)_!\mathcal{F} \to g_! f_! \mathcal{F}$
par $\overline{z}$ coïncide avec l'application du
lemme \ref{more-etale-lemma-f-shriek-composition}.
Le lecteur voit alors immédiatement que l'on obtient la propriété (3).
La propriété (3) garantit bien sûr que notre transformation de foncteurs
$(g \circ f)_! \to g_! \circ f_!$, construite au moyen du recouvrement ouvert
$Y = \bigcup Y_i$, ne dépend pas du choix de ce recouvrement.
Enfin, la propriété (4) se déduit de la propriété (3), déjà démontrée,
en examinant ce qui se passe sur les fibres.
\end{proof}








\section{Pondérations et morphismes trace pour les morphismes localement quasi-finis}
\sectionmark{Pondérations et morphismes trace}
\label{more-etale-section-weightings}

\noindent
Une référence pour cette section est
\cite[Exposé XVII, Proposition 6.2.5]{SGA4}.

\medskip\noindent
Soit $f : X \to Y$ un morphisme de schémas localement quasi-fini.
Soit $w : X \to \mathbf{Z}$ une pondération de $f$ ; voir la
définition \ref{more-morphisms-definition-weighting} de Compléments sur les morphismes.
Soit $\mathcal{F}$ un faisceau abélien sur $Y_\etale$.
Dans cette section, nous montrerons qu'il existe un morphisme
$$
\text{Tr}_{f, w, \mathcal{F}} :
f_!f^{-1}\mathcal{F}
\longrightarrow
\mathcal{F}
$$
de faisceaux abéliens sur $Y_\etale$, caractérisé par la propriété suivante :
sur les fibres en un point géométrique $\overline{y}$ de $Y$, on obtient l'application
$$
\bigoplus\nolimits_{f(\overline{x}) = \overline{y}} w(\overline{x}) :
(f_!f^{-1}\mathcal{F})_{\overline{y}} =
\bigoplus\nolimits_{f(\overline{x}) = \overline{y}}
\mathcal{F}_{\overline{y}}
\longrightarrow
\mathcal{F}_{\overline{y}}
$$
Comme indiqué, la flèche est donnée par multiplication par l'entier
$w(\overline{x})$ sur le facteur correspondant à $\overline{x}$.
L'égalité située à gauche de la flèche résulte du
lemme \ref{more-etale-lemma-lqf-f-shriek-stalk}, combiné au
lemme \ref{etale-cohomology-lemma-stalk-pullback} de Cohomologie étale.

\medskip\noindent
Si le morphisme $f : X \to Y$ est plat, localement quasi-fini et localement de
présentation finie, il existe une pondération canonique,
et l'on obtient un morphisme trace canonique dont la formation est
compatible au changement de base ; voir
l'exemple \ref{more-etale-example-trace-for-flat-quasi-finite}.
Si $Y$ est un schéma localement noethérien unibranche et si $f : X \to Y$
est localement quasi-fini, on peut également définir une pondération
(naturelle) de $f$ et l'on dispose aussi de morphismes trace dans ce cas ; voir
l'exemple \ref{more-etale-example-trace-for-quasi-finite-over-normal}.

\begin{lemma}
\label{more-etale-lemma-f-shriek-projection}
Soit $f : X \to Y$ un morphisme de schémas localement quasi-fini.
Soit $\Lambda$ un anneau.
Soit $\mathcal{F}$ un faisceau de $\Lambda$-modules sur $X_\etale$
et soit $\mathcal{G}$ un faisceau de $\Lambda$-modules sur $Y_\etale$.
Il existe un isomorphisme canonique
$$
can :
f_!\mathcal{F} \otimes_\Lambda \mathcal{G}
\longrightarrow
f_!(\mathcal{F} \otimes_\Lambda f^{-1}\mathcal{G})
$$
de faisceaux de $\Lambda$-modules sur $Y_\etale$.
\end{lemma}

\begin{proof}
Rappelons que $f_!\mathcal{F} = (f_{p!}\mathcal{F})^\#$
d'après la définition \ref{more-etale-definition-f-shriek-lqf}, où
$f_{p!}\mathcal{F}$ est le préfaisceau
construit à la section \ref{more-etale-section-finite-support}.
Pour construire la flèche, il suffit donc de construire une application
$$
f_{p!}\mathcal{F} \otimes_{p, \Lambda} \mathcal{G}
\longrightarrow
f_{p!}(\mathcal{F} \otimes_\Lambda f^{-1}\mathcal{G})
$$
de préfaisceaux sur $Y_\etale$. Ici, le symbole $\otimes_{p, \Lambda}$
désigne le produit tensoriel de préfaisceaux ; voir la
section \ref{sites-modules-section-tensor-product} de Modules sur les sites.
Soit $V$ un objet de $Y_\etale$. Rappelons que
$$
f_{p!}\mathcal{F}(V) = \colim_Z H_Z(\mathcal{F})
\quad\text{et}\quad
f_{p!}(\mathcal{F} \otimes_\Lambda f^{-1}\mathcal{G})(V) =
\colim_Z H_Z(\mathcal{F} \otimes_\Lambda f^{-1}\mathcal{G})
$$
Voir la section \ref{more-etale-section-finite-support}. Notre application est définie
sur les tenseurs purs par la règle
$$
(Z, s) \otimes t \longmapsto (Z, s \otimes f^{-1}t)
$$
(voir ci-dessous pour les notations), puis prolongée par linéarité à tout
$(f_{p!}\mathcal{F} \otimes_{p, \Lambda} \mathcal{G})(V) =
f_{p!}\mathcal{F}(V) \otimes_\Lambda \mathcal{G}(V)$.
Les notations employées ici sont les suivantes :
\begin{enumerate}
\item $Z \subset X_V$ est un sous-schéma localement fermé fini sur $V$ ;
\item $s \in H_Z(\mathcal{F})$, ce qui signifie que $s \in \mathcal{F}(U)$
avec $\text{Supp}(s) \subset Z$ pour un ouvert $U \subset X_V$ tel que
$Z \subset U$ soit fermé ;
\item $t \in \mathcal{G}(V)$, d'image $f^{-1}t \in f^{-1}\mathcal{G}(U)$.
\end{enumerate}
Puisque le support de $s \in \mathcal{F}(U)$ est contenu dans $Z$, il est clair
que le support de $s \otimes f^{-1}t$ est lui aussi contenu dans $Z$.
Le couple $(Z, s \otimes f^{-1}t)$ a donc bien un sens.
Il est immédiat que la construction commute aux applications de transition
de la limite inductive $\colim_Z H_Z(\mathcal{F})$ et qu'elle est
compatible avec les applications de restriction. Enfin, il est tout aussi clair
que la construction est compatible avec les identifications des
fibres de $f_!$ données au lemme \ref{more-etale-lemma-lqf-f-shriek-stalk}.
Autrement dit, l'application $can$ que nous avons construite sur les fibres en un point géométrique
$\overline{y}$ s'insère dans le diagramme commutatif
$$
\xymatrix{
(f_!\mathcal{F} \otimes_\Lambda \mathcal{G})_{\overline{y}}
\ar[r]_-{can_{\overline{y}}} \ar[d] &
f_!(\mathcal{F} \otimes_\Lambda f^{-1}\mathcal{G})_{\overline{y}} \ar[d] \\
(\bigoplus \mathcal{F}_{\overline{x}})
\otimes_\Lambda \mathcal{G}_{\overline{y}} \ar[r] &
\bigoplus
(\mathcal{F}_{\overline{x}} \otimes_\Lambda \mathcal{G}_{\overline{y}})
}
$$
où les sommes directes portent sur les points géométriques $\overline{x}$
au-dessus de $\overline{y}$, les flèches verticales sont les identifications
du lemme \ref{more-etale-lemma-lqf-f-shriek-stalk}, et la flèche horizontale inférieure
est l'isomorphisme évident.
Nous concluons que $can$ est un isomorphisme, comme souhaité.
\end{proof}

\begin{lemma}
\label{more-etale-lemma-trace-map-exists}
Soit $f : X \to Y$ un morphisme de schémas localement quasi-fini.
Soit $w : X \to \mathbf{Z}$ une pondération de $f$. Pour tout faisceau abélien
$\mathcal{F}$ sur $Y$, il existe un unique morphisme trace
$\text{Tr}_{f, w, \mathcal{F}} : f_!f^{-1}\mathcal{F} \to \mathcal{F}$
ayant le comportement prescrit sur les fibres.
\end{lemma}

\begin{proof}
D'après le lemme \ref{more-etale-lemma-f-shriek-projection}, nous avons une identification
$f_!f^{-1}\mathcal{F} = f_!\underline{\mathbf{Z}} \otimes \mathcal{F}$
compatible avec la description des fibres de ces faisceaux aux
points géométriques. Il suffit donc de construire le morphisme
$$
\text{Tr}_{f, w, \underline{\mathbf{Z}}} :
f_!\underline{\mathbf{Z}}
\longrightarrow
\underline{\mathbf{Z}}
$$
ayant le comportement prescrit sur les fibres. D'après la
définition \ref{more-etale-definition-f-shriek-lqf}, on a
$f_!\underline{\mathbf{Z}} = (f_{p!}\underline{\mathbf{Z}})^\#$,
où $f_{p!}\underline{\mathbf{Z}}$ est le préfaisceau
construit à la section \ref{more-etale-section-finite-support}.
Il suffit donc de construire une application
$$
f_{p!}\underline{\mathbf{Z}} \longrightarrow \underline{\mathbf{Z}}
$$
de préfaisceaux sur $Y_\etale$. Soit $V$ un objet de $Y_\etale$.
Rappelons, d'après la section \ref{more-etale-section-finite-support}, que
$$
f_{p!}\underline{\mathbf{Z}}(V) = \colim_Z H_Z(\underline{\mathbf{Z}})
$$
Ici, la limite inductive est prise sur l'ensemble ordonné des sous-schémas
localement fermés $Z \subset X_V$ qui sont finis sur $V$. Pour chaque tel
$Z$, nous allons définir une application
$$
H_Z(\underline{\mathbf{Z}}) \longrightarrow \underline{\mathbf{Z}}(V)
$$
compatible avec les applications qui définissent la limite inductive.

\medskip\noindent
Soit $Z \subset X_V$ localement fermé et fini sur $V$.
Choisissons un ouvert $U \subset X_V$ contenant $Z$ comme partie fermée.
Un élément $s$ de $H_Z(\underline{\mathbf{Z}})$ est une section
$s \in \underline{\mathbf{Z}}(U)$ dont le support est contenu dans $Z$.
Soit $U_n \subset U$ la partie ouverte et fermée où la valeur
de $s$ est $n \in \mathbf{Z}$. La condition sur le support donne
$Z \cap U_n = U_n$ pour $n \not = 0$. Ainsi, pour $n \not = 0$, l'ouvert
$U_n$ est aussi fermé dans $Z$ (comme complémentaire de tous les autres),
et $U_n \to V$ est donc fini puisque $Z$ est fini sur $V$.
Par définition même d'une pondération, la
fonction $\int_{U_n \to V} w|_{U_n}$ est alors localement constante sur $V$,
et nous pouvons la regarder comme un élément de $\underline{\mathbf{Z}}(V)$.
Notre construction envoie $(Z, s)$ sur l'élément
$$
\sum\nolimits_{n \in \mathbf{Z},\ n \not = 0}
n \left(\int_{U_n \to V} w|_{U_n}\right)
\quad \in \quad \underline{\mathbf{Z}}(V)
$$
La somme est localement finie sur $V$, donc bien définie ; nous omettons les détails
(dans toute la discussion, le lecteur peut commencer par choisir des ouverts affines et
s'assurer que tous les schémas qui interviennent dans l'argument sont quasi-compacts, afin que
la somme soit finie). Nous omettons de vérifier que cette construction
est compatible avec les applications de la limite inductive et avec les applications de
restriction qui définissent $f_{p!}\underline{\mathbf{Z}}$.

\medskip\noindent
Soit $\overline{y}$ un point géométrique de $Y$ au-dessus du point $y \in Y$.
En prenant les fibres en $\overline{y}$, la construction précédente détermine une application
$$
(f_!\underline{\mathbf{Z}})_{\overline{y}}
= \bigoplus\nolimits_{f(\overline{x}) = \overline{y}} \mathbf{Z}
\longrightarrow
\mathbf{Z} = \underline{\mathbf{Z}}_{\overline{y}}
$$
Pour achever la démonstration, montrons que cette application est donnée par multiplication par
$w(\overline{x})$ sur le facteur correspondant à $\overline{x}$.
Choisissons en effet $\overline{x}$ au-dessus de $\overline{y}$.
Il existe un voisinage étale $(V, \overline{v}) \to (Y, \overline{y})$
tel que $X_V$ contienne un ouvert $U$ fini sur $V$,
tel que $\overline{x}$ soit le seul à appartenir à $U$ parmi les points géométriques
de $X$ relevant $\overline{y}$.
Ceci résulte du lemme de Compléments sur les morphismes
\ref{more-morphisms-lemma-etale-makes-quasi-finite-finite-multiple-points-var};
nous omettons certains détails. Alors $(U, 1)$ définit une section de
$f_!\underline{\mathbf{Z}}$ sur $V$, dont l'image vaut $1$ dans le facteur
correspondant à $\overline{x}$ et zéro dans les autres facteurs
(voir la démonstration du lemme \ref{more-etale-lemma-finite-support-stalk}), et
notre construction précédente envoie $(U, 1)$ sur $\int_{U \to V} w|_U$,
qui est constante de valeur $w(\overline{x})$ dans un voisinage
de $\overline{v}$, comme souhaité.
\end{proof}

\begin{lemma}
\label{more-etale-lemma-properties-trace-map}
Soit $f : X \to Y$ un morphisme de schémas localement quasi-fini.
Soit $w : X \to \mathbf{Z}$ une pondération de $f$. Les morphismes trace
construits ci-dessus vérifient les propriétés suivantes :
\begin{enumerate}
\item $\text{Tr}_{f, w, \mathcal{F}}$ est fonctoriel en $\mathcal{F}$ ;
\item $\text{Tr}_{f, w, \mathcal{F}}$ est compatible à tout changement de base ;
\item étant donnés un anneau $\Lambda$ et $K$ dans $D(Y_\etale, \Lambda)$,
on obtient $\text{Tr}_{f, w, K} : f_!f^{-1}K \to K$, fonctoriel en $K$
et compatible à tout changement de base.
\end{enumerate}
\end{lemma}

\begin{proof}
Le point (1) résulte soit de la construction du morphisme trace
dans la démonstration du lemme \ref{more-etale-lemma-trace-map-exists},
soit, plus simplement, de ce que la caractérisation de l'application
impose cette fonctorialité sur toutes les fibres.
Soit
$$
\xymatrix{
X' \ar[r]_{g'} \ar[d]_{f'} & X \ar[d]^f \\
Y' \ar[r]^g & Y
}
$$
un diagramme cartésien de schémas. Alors la fonction
$w' = w \circ g' : X' \to \mathbf{Z}$ est une pondération de $f'$ d'après le
lemme \ref{more-morphisms-lemma-weighting-base-change} de Compléments sur les morphismes.
Le point (2) signifie que le diagramme
$$
\xymatrix{
g^{-1}f_!f^{-1}\mathcal{F}
\ar[rr]_-{g^{-1}\text{Tr}_{f, w, \mathcal{F}}} \ar@{=}[d] & &
g^{-1}\mathcal{F} \ar@{=}[d] \\
f'_!(f')^{-1}g^{-1}\mathcal{F}
\ar[rr]^-{\text{Tr}_{f', w', g^{-1}\mathcal{F}}} & &
g^{-1}\mathcal{F}
}
$$
est commutatif, l'égalité verticale gauche étant donnée par
$$
g^{-1}f_!f^{-1}\mathcal{F} =
f'_!(g')^{-1}f^{-1}\mathcal{F} =
f'_!(f')^{-1}g^{-1}\mathcal{F}
$$
où le premier signe d'égalité est donné par le lemme \ref{more-etale-lemma-lqf-base-change-f-shriek}
(changement de base pour l'image directe à support propre). La commutativité de ce diagramme
résulte de la caractérisation de l'action de nos morphismes trace
sur les fibres et du fait que la flèche de changement de base du
lemme \ref{more-etale-lemma-lqf-base-change-f-shriek} respecte la description des fibres.

\medskip\noindent
Les points (1) et (2) étant acquis, le point (3) résulte de ce que les foncteurs
$f^{-1} : D(Y_\etale, \Lambda) \to D(X_\etale, \Lambda)$ et
$f_! : D(X_\etale, \Lambda) \to D(Y_\etale, \Lambda)$
s'obtiennent en appliquant $f^{-1}$ et $f_!$ à des complexes quelconques
de modules représentant les objets considérés.
\end{proof}

\begin{lemma}
\label{more-etale-lemma-trace-composition}
Soient $f : X \to Y$ et $g : Y \to Z$ des morphismes localement quasi-finis.
Soit $w_f : X \to \mathbf{Z}$ une pondération de $f$, et soit
$w_g : Y \to \mathbf{Z}$ une pondération de $g$. Pour
$K \in D(Z_\etale, \Lambda)$, la composée
$$
(g \circ f)_!(g \circ f)^{-1}K =
g_! f_! f^{-1} g^{-1}K
\xrightarrow{g_! \text{Tr}_{f, w_f, g^{-1}K}}
g_!g^{-1}K
\xrightarrow{\text{Tr}_{g, w_g, K}}
K
$$
est égale à $\text{Tr}_{g \circ f, w_{g \circ f}, K}$, où
$w_{g \circ f}(x) = w_f(x) w_g(f(x))$.
\end{lemma}

\begin{proof}
On a $(g \circ f)_! = g_! \circ f_!$ d'après le
lemme \ref{more-etale-lemma-lqf-shriek-composition}.
Le lemme \ref{more-morphisms-lemma-weighting-composition} de Compléments sur les morphismes
montre que $w_{g \circ f}$ est une pondération de $g \circ f$ ;
l'énoncé a donc un sens. L'égalité se vérifie sur les fibres.
Nous omettons les détails.
\end{proof}

\begin{example}[Morphisme trace dans le cas plat quasi-fini]
\label{more-etale-example-trace-for-flat-quasi-finite}
Soit $f : X \to Y$ un morphisme de schémas plat,
localement quasi-fini et localement de présentation finie.
On obtient une pondération positive canonique $w : X \to \mathbf{Z}$
en posant
$$
w(x) = \text{longueur}_{\mathcal{O}_{X, x}}
(\mathcal{O}_{X, x}/\mathfrak m_{f(x)} \mathcal{O}_{X, x})
[\kappa(x) : \kappa(f(x))]_i
$$
Voir le lemme de Compléments sur les morphismes
\ref{more-morphisms-lemma-weighting-flat-quasi-finite}.
Ainsi, d'après les lemmes \ref{more-etale-lemma-trace-map-exists} et
\ref{more-etale-lemma-properties-trace-map}, on obtient pour $f$ des morphismes trace
$$
\text{Tr}_{f, K} : f_!f^{-1}K \longrightarrow K
$$
fonctoriels en $K$ dans $D(Y_\etale, \Lambda)$ et compatibles
à tout changement de base. Notons que tout changement de base
$f' : X' \to Y'$ de $f$ possède les mêmes propriétés et que la restriction de $w$
est la pondération canonique de $f'$.
\end{example}

\begin{remark}
\label{more-etale-remark-trace-etale-counit}
Soit $j : U \to X$ un morphisme étale de schémas. Alors le morphisme trace
$\text{Tr} : j_!j^{-1}K \to K$ de
l'exemple \ref{more-etale-example-trace-for-flat-quasi-finite} est égal à la
flèche d'adjonction associée à l'adjonction entre $j_!$ et $j^{-1}$.
Nous avons déjà employé le terme « trace » pour cette flèche d'adjonction dans la
section \ref{etale-cohomology-section-trace-method} de Cohomologie étale.
\end{remark}

\begin{example}[Morphisme trace dans le cas quasi-fini sur un schéma normal]
\label{more-etale-example-trace-for-quasi-finite-over-normal}
Soit $Y$ un schéma géométriquement unibranche et localement noethérien ;
par exemple, $Y$ peut être une variété normale. Soit $f : X \to Y$ un morphisme de schémas
localement quasi-fini. Il existe alors une
pondération positive $w : X \to \mathbf{Z}$ de $f$,
définie approximativement en associant à $x$ le
« degré séparable générique »
de $\mathcal{O}_{X, x}^{sh}$ sur $\mathcal{O}_{Y, f(x)}^{sh}$.
Voir le lemme de Compléments sur les morphismes
\ref{more-morphisms-lemma-weighting-quasi-finite-Noetherian}.
Ainsi, d'après les lemmes \ref{more-etale-lemma-trace-map-exists} et
\ref{more-etale-lemma-properties-trace-map}, on obtient pour $f$ et $w$ des morphismes trace
$$
\text{Tr}_{f, w, K} : f_!f^{-1}K \longrightarrow K
$$
fonctoriels en $K$ dans $D(Y_\etale, \Lambda)$ et compatibles
à tout changement de base. Toutefois, dans ce cas, pour un changement de base
$f' : X' \to Y'$ de $f$, la restriction de $w$ à $X'$ ne possède pas en général
d'interprétation « naturelle » en termes du
morphisme $f'$.
\end{example}
















\section{Image inverse extraordinaire pour les morphismes localement quasi-finis}
\sectionmark{Image inverse extraordinaire}
\label{more-etale-section-duality-locally-quasi-finite}

\noindent
Pour tout morphisme de schémas localement quasi-fini $f : X \to Y$, le
foncteur $f_! : \textit{Ab}(X_\etale) \to \textit{Ab}(Y_\etale)$ commute
aux sommes directes et est exact ; voir le lemme \ref{more-etale-lemma-lqf-f-shriek-stalk}.
Ceci suggère qu'il admet un adjoint à droite, que nous noterons $f^!$.

\medskip\noindent
Avertissement : ce foncteur est la version non dérivée !

\begin{lemma}
\label{more-etale-lemma-lqf-f-upper-shriek}
Soit $f : X \to Y$ un morphisme de schémas localement quasi-fini.
\begin{enumerate}
\item Le foncteur $f_! : \textit{Ab}(X_\etale) \to \textit{Ab}(Y_\etale)$
admet un adjoint à droite $f^! : \textit{Ab}(Y_\etale) \to \textit{Ab}(X_\etale)$.
\item On a
$f^!(\overline{y}_*A) = \prod_{f(\overline{x}) = \overline{y}} \overline{x}_*A$.
\item Si $\Lambda$ est un anneau, le foncteur
$f_! : \textit{Mod}(X_\etale, \Lambda) \to \textit{Mod}(Y_\etale, \Lambda)$
admet un adjoint à droite
$f^! : \textit{Mod}(Y_\etale, \Lambda) \to \textit{Mod}(X_\etale, \Lambda)$
qui coïncide avec $f^!$ sur les faisceaux abéliens sous-jacents.
\end{enumerate}
\end{lemma}

\begin{proof}
Démonstration de (1). Soit $E \subset \Ob(\textit{Ab}(Y_\etale))$ la classe
formée des produits de faisceaux gratte-ciel. Nous affirmons ce qui suit :
\begin{enumerate}
\item[(a)] tout $\mathcal{G}$ de $\textit{Ab}(Y_\etale)$ est un sous-faisceau
d'un élément de $E$ ;
\item[(b)] pour tout $\mathcal{G} \in E$, il existe un objet
$\mathcal{H}$ de $\textit{Ab}(X_\etale)$ tel que
$\Hom(f_!\mathcal{F}, \mathcal{G}) = \Hom(\mathcal{F}, \mathcal{H})$
fonctoriellement en $\mathcal{F}$.
\end{enumerate}
Une fois cette affirmation vérifiée, le dual du
lemme \ref{homology-lemma-partially-defined-adjoint} d'Homologie
fournit le foncteur adjoint $f^!$.

\medskip\noindent
Le point (a) est vrai, car on peut envoyer $\mathcal{G}$ dans le faisceau
$\prod \overline{y}_*\mathcal{G}_{\overline{y}}$, où le
produit porte sur tous les points géométriques de $Y$. C'est une injection d'après le
théorème \ref{etale-cohomology-theorem-exactness-stalks} de Cohomologie étale.
(Il s'agit de la première étape de la résolution de Godement lorsque
l'on travaille avec des faisceaux abéliens sur des espaces topologiques.)

\medskip\noindent
Le point (b) et le point (2) du lemme s'établissent comme suit.
Supposons que $\mathcal{G} = \prod \overline{y}_*A_{\overline{y}}$
pour certains groupes abéliens $A_{\overline{y}}$. Alors
$$
\Hom(f_!\mathcal{F}, \mathcal{G}) =
\prod \Hom(f_!\mathcal{F}, \overline{y}_*A_{\overline{y}})
$$
Il suffit donc de trouver des faisceaux abéliens $\mathcal{H}_{\overline{y}}$
sur $X_\etale$ qui représentent les foncteurs
$\mathcal{F} \mapsto \Hom(f_!\mathcal{F}, \overline{y}_*A_{\overline{y}})$
et de prendre $\mathcal{H} = \prod \mathcal{H}_{\overline{y}}$.
Nous sommes ainsi ramenés au cas $\mathcal{H} = \overline{y}_*A$
pour un point géométrique fixé $\overline{y} : \Spec(k) \to Y$
et un groupe abélien fixé $A$. Nous affirmons que, dans ce cas,
$\mathcal{H} = \prod_{f(\overline{x}) = \overline{y}} \overline{x}_*A$ convient.
Ceci achèvera la démonstration des points (1) et (2) du lemme.
En effet, on a
$$
\Hom(f_!\mathcal{F}, \overline{y}_*A) =
\Hom_{\textit{Ab}}((f_!\mathcal{F})_{\overline{y}}, A) =
\Hom_{\textit{Ab}}(\bigoplus\nolimits_{f(\overline{x}) = \overline{y}}
\mathcal{F}_{\overline{x}}, A)
$$
d'une part, d'après la description des fibres donnée au
lemme \ref{more-etale-lemma-lqf-f-shriek-stalk},
et, d'autre part, on a
$$
\Hom(\mathcal{F}, \mathcal{H}) =
\prod\nolimits_{f(\overline{x}) = \overline{y}}
\Hom(\mathcal{F}, \overline{x}_*A) =
\prod\nolimits_{f(\overline{x}) = \overline{y}}
\Hom_{\textit{Ab}}(\mathcal{F}_{\overline{x}}, A)
$$
Nous laissons au lecteur le soin d'identifier ces deux expressions comme foncteurs en $\mathcal{F}$.

\medskip\noindent
Démonstration du point (3). Remarquons qu'un objet de $\textit{Mod}(X_\etale, \Lambda)$
n'est autre qu'un objet $\mathcal{F}$ de $\textit{Ab}(X_\etale)$
muni d'une application $\Lambda \to \text{End}(\mathcal{F})$. Les
foncteurs $f_!$ et $f^!$ de (1) définissent les foncteurs $f_!$ et $f^!$ comme en (3).
Un calcul immédiat montre qu'ils sont adjoints.
\end{proof}

\begin{lemma}
\label{more-etale-lemma-etale-upper-shriek}
Soit $j : U \to X$ un morphisme étale. Alors $j^! = j^{-1}$.
\end{lemma}

\begin{proof}
Cela résulte du fait que le foncteur $j_!$ défini à la
section \ref{more-etale-section-finite-support}
coïncide avec $j_!$ tel qu'il est défini dans Cohomologie étale, à la section
\ref{etale-cohomology-section-extension-by-zero} ; voir le
lemme \ref{more-etale-lemma-finite-support-etale-shriek}. Enfin, dans la
section \ref{etale-cohomology-section-extension-by-zero} de Cohomologie étale,
le foncteur $j_!$ est défini
comme l'adjoint à gauche de $j^{-1}$ ; on conclut donc par
unicité des foncteurs adjoints.
\end{proof}

\begin{lemma}
\label{more-etale-lemma-upper-shriek-restriction}
Soient $f : X \to Y$ et $g : Y \to Z$ des morphismes séparés et localement
quasi-finis. Il existe un isomorphisme canonique $(g \circ f)^! \to f^! \circ g^!$.
Étant donné un troisième morphisme localement quasi-fini $h : Z \to T$,
le diagramme
$$
\xymatrix{
(h \circ g \circ f)^! \ar[r] \ar[d] &
f^! \circ (h \circ g)^! \ar[d] \\
(g \circ f)^! \circ h^! \ar[r] & f^! \circ g^! \circ h^!
}
$$
est commutatif.
\end{lemma}

\begin{proof}
Par unicité des foncteurs adjoints, ceci se traduit immédiatement
par l'énoncé (dual) correspondant pour les foncteurs $f_!$.
Voir le lemme \ref{more-etale-lemma-lqf-shriek-composition}.
\end{proof}

\begin{lemma}
\label{more-etale-lemma-upper-shriek-restriction-etale}
Soient $j : U \to X$ et $j' : V \to U$ des morphismes étales.
L'isomorphisme $(j \circ j')^{-1} = (j')^{-1} \circ j^{-1}$
et l'isomorphisme $(j \circ j')^! = (j')^! \circ j^!$ du
lemme \ref{more-etale-lemma-upper-shriek-restriction}
coïncident via l'isomorphisme du lemme \ref{more-etale-lemma-etale-upper-shriek}.
\end{lemma}

\begin{proof}
La démonstration est omise.
\end{proof}

\begin{lemma}
\label{more-etale-lemma-lqf-base-change-upper-shriek}
Considérons un carré cartésien
$$
\xymatrix{
X' \ar[r]_{g'} \ar[d]_{f'} & X \ar[d]^f \\
Y' \ar[r]^g & Y
}
$$
de schémas, où $f$ est localement quasi-fini. Pour tout faisceau abélien $\mathcal{F}$
sur $Y'_\etale$, on a $(g')_*(f')^!\mathcal{F} = f^!g_*\mathcal{F}$.
\end{lemma}

\begin{proof}
Par unicité des foncteurs adjoints, ceci résulte de
l'énoncé (dual) correspondant pour les foncteurs $f_!$.
Voir le lemme \ref{more-etale-lemma-lqf-base-change-f-shriek}.
\end{proof}

\begin{remark}
\label{more-etale-remark-pointed-sets}
Les résultats de cette section se généralisent aux faisceaux d'ensembles pointés.
En effet, pour un site $\mathcal{C}$, notons $\Sh^*(\mathcal{C})$ la catégorie des
faisceaux d'ensembles pointés. Les constructions de la présente section et de la précédente
s'appliquent, mutatis mutandis, aux faisceaux d'ensembles pointés. Ainsi, étant donné un morphisme
de schémas localement quasi-fini $f : X \to Y$, on obtient
un couple de foncteurs adjoints
$$
f_! : \Sh^*(X_\etale) \longrightarrow \Sh^*(Y_\etale)
\quad\text{et}\quad
f^! : \Sh^*(Y_\etale) \longrightarrow \Sh^*(X_\etale)
$$
tel que, pour tout point géométrique $\overline{y}$ de $Y$, il existe des
isomorphismes
$$
(f_!\mathcal{F})_{\overline{y}} =
\coprod\nolimits_{f(\overline{x}) = \overline{y}}
\mathcal{F}_{\overline{x}}
$$
(le coproduit étant pris dans la catégorie des ensembles pointés) fonctoriels en
$\mathcal{F} \in \Sh^*(X_\etale)$ et des isomorphismes
$$
f^!(\overline{y}_*S) =
\prod\nolimits_{f(\overline{x}) = \overline{y}}
\overline{x}_*S
$$
fonctoriels par rapport à l'ensemble pointé $S$. Si
$F : \textit{Ab}(X_\etale) \to \Sh^*(X_\etale)$ et
$F : \textit{Ab}(Y_\etale) \to \Sh^*(Y_\etale)$
désignent les foncteurs d'oubli, la compatibilité des constructions
garantit l'existence d'applications canoniques
$$
f_!F(\mathcal{F}) \longrightarrow F(f_!\mathcal{F})
$$
fonctorielles en $\mathcal{F} \in \textit{Ab}(X_\etale)$ et
$$
F(f^!\mathcal{G}) \longrightarrow f^!F(\mathcal{G})
$$
fonctorielles en $\mathcal{G} \in \textit{Ab}(Y_\etale)$,
qui induisent les applications évidentes sur les fibres, resp.\ les faisceaux gratte-ciel.
En fait, la transformation $F \circ f^! \to f^! \circ F$ est un isomorphisme
(car $f^!$ commute aux produits).
\end{remark}







\section{Image inverse extraordinaire dérivée pour les morphismes localement quasi-finis}
\sectionmark{Image inverse extraordinaire dérivée}
\label{more-etale-section-derived-duality-locally-quasi-finite}

\noindent
On peut passer aux foncteurs dérivés dans la
section \ref{more-etale-section-duality-locally-quasi-finite},
et obtenir l'énoncé suivant.

\begin{lemma}
\label{more-etale-lemma-lqf-shriek-derived}
Soit $f : X \to Y$ un morphisme de schémas localement quasi-fini.
Soit $\Lambda$ un anneau. Les foncteurs $f_!$ et $f^!$ de la
définition \ref{more-etale-definition-f-shriek-lqf} et du
lemme \ref{more-etale-lemma-lqf-f-upper-shriek}
induisent des foncteurs adjoints $f_! : D(X_\etale, \Lambda) \to D(Y_\etale, \Lambda)$
et $Rf^! : D(Y_\etale, \Lambda) \to D(X_\etale, \Lambda)$
sur les catégories dérivées.
\end{lemma}

\noindent
Dans le cas séparé, le foncteur $f_!$ est défini à la
section \ref{more-etale-section-compact-support}.

\begin{proof}
Ceci résulte immédiatement du
lemme \ref{derived-lemma-derived-adjoint-functors} de Catégories dérivées,
du fait que $f_!$ est exact (lemme \ref{more-etale-lemma-lqf-f-shriek-stalk})
et donc que $Lf_! = f_!$,
et du fait que l'on dispose d'assez de complexes K-injectifs de $\Lambda$-modules
sur $Y_\etale$ pour que $Rf^!$ soit défini.
\end{proof}

\begin{remark}
\label{more-etale-remark-lqf-shriek-derived-torsion}
Soit $f : X \to Y$ un morphisme de schémas localement quasi-fini. Soit
$\Lambda$ un anneau. Le foncteur
$f_! : D(X_\etale, \Lambda) \to D(Y_\etale, \Lambda)$ du
lemme \ref{more-etale-lemma-lqf-shriek-derived}
envoie les complexes dont les faisceaux de cohomologie sont de torsion
sur des complexes dont les faisceaux de cohomologie sont de torsion.
Ceci résulte immédiatement de la description des fibres de $f_!$ ;
voir le lemme \ref{more-etale-lemma-lqf-f-shriek-stalk}.
\end{remark}

\begin{lemma}
\label{more-etale-lemma-mayer-vietoris-shriek}
Soit $X$ un schéma. Supposons $X = U \cup V$, avec $U$ et $V$ ouverts.
Soit $\Lambda$ un anneau. Soit $K \in D(X_\etale, \Lambda)$. Il existe
un triangle distingué
$$
j_{U \cap V!}K|_{U \cap V} \to
j_{U!}K|_U \oplus j_{V!}K|_V \to K \to 
j_{U \cap V!}K|_{U \cap V}[1]
$$
dans $D(X_\etale, \Lambda)$, avec les notations évidentes.
\end{lemma}

\begin{proof}
Comme les foncteurs de restriction et d'image directe à support propre employés
sont exacts, il suffit de montrer que, pour tout faisceau abélien $\mathcal{F}$ sur
$X_\etale$, la suite
$$
0 \to j_{U \cap V!}\mathcal{F}|_{U \cap V} \to
j_{U!}\mathcal{F}|_U \oplus j_{V!}\mathcal{F}|_V \to \mathcal{F} \to 0
$$
est exacte. Il suffit de le vérifier sur les fibres.
\end{proof}

\begin{lemma}
\label{more-etale-lemma-triangle-associated-to-open}
Soit $X$ un schéma. Soit $Z \subset X$ un sous-schéma fermé, et soit
$U \subset X$ son complémentaire. Notons $i : Z \to X$ et $j : U \to X$
les morphismes d'inclusion. Soit $\Lambda$ un anneau.
Soit $K \in D(X_\etale, \Lambda)$. Il existe un triangle distingué
$$
j_!j^{-1}K \to K \to i_*i^{-1}K \to j_!j^{-1}K[1]
$$
dans $D(X_\etale, \Lambda)$.
\end{lemma}

\begin{proof}
Conséquence immédiate du lemme de Cohomologie étale
\ref{etale-cohomology-lemma-ses-associated-to-open}
et du fait que les foncteurs $j_!$, $j^{-1}$, $i_*$, $i^{-1}$
sont exacts ; leurs versions dérivées se calculent donc en
appliquant ces foncteurs à tout complexe de faisceaux représentant $K$.
\end{proof}





\section{Préliminaires à l'image directe dérivée à support propre par compactification}
\sectionmark{Préliminaires à l'image directe dérivée}
\label{more-etale-section-prelim}

\noindent
Dans cette section, nous démontrons quelques lemmes sur l'existence
de certains isomorphismes naturels de foncteurs qui résultent immédiatement
du changement de base propre.

\begin{lemma}
\label{more-etale-lemma-shriek-proper-and-open}
Considérons un diagramme commutatif de schémas
$$
\xymatrix{
X' \ar[r]_{g'} \ar[d]_{f'} & X \ar[d]^f \\
Y' \ar[r]^g & Y
}
$$
où $f$ et $f'$ sont propres, et $g$ et $g'$ séparés et localement quasi-finis.
Soit $\Lambda$ un anneau. Fonctoriellement en $K \in D(X'_\etale, \Lambda)$,
il existe un morphisme canonique
$$
g_!Rf'_*K \longrightarrow Rf_*(g'_!K)
$$
dans $D(Y_\etale, \Lambda)$. Cette application est un isomorphisme si
(a) $K$ est borné inférieurement et ses faisceaux de cohomologie sont de torsion, ou si
(b) $\Lambda$ est un anneau de torsion.
\end{lemma}

\begin{proof}
Représentons $K$ par un complexe K-injectif $\mathcal{J}^\bullet$ de
faisceaux de $\Lambda$-modules sur $X'_\etale$. Choisissons un quasi-isomorphisme
$g'_!\mathcal{J}^\bullet \to \mathcal{I}^\bullet$ vers un
complexe K-injectif $\mathcal{I}^\bullet$ de
faisceaux de $\Lambda$-modules sur $X_\etale$.
Nous pouvons alors considérer l'application
$$
g_!f'_*\mathcal{J}^\bullet =
g_!f'_!\mathcal{J}^\bullet =
f_!g'_!\mathcal{J}^\bullet =
f_*g'_!\mathcal{J}^\bullet \to
f_*\mathcal{I}^\bullet
$$
où la première et la troisième égalité résultent du
lemme \ref{more-etale-lemma-proper-f-shriek},
et la deuxième du
lemme \ref{more-etale-lemma-f-shriek-composition}, qui affirme que
$g_! \circ f'_!$ et $f_! \circ g'_!$ sont tous deux égaux à
$(g \circ f')_! = (f \circ g')_!$ comme sous-faisceaux de
$(g \circ f')_* = (f \circ g')_*$.

\medskip\noindent
Supposons que $\Lambda$ soit de torsion, c'est-à-dire que nous soyons dans le cas (b).
Avec les notations précédentes, il suffit de montrer que
$f_*g'_!\mathcal{J}^\bullet \to f_*\mathcal{I}^\bullet$
est un isomorphisme. La question est locale sur $Y$.
Nous pouvons donc supposer que la dimension
des fibres de $f$ est bornée ; voir le lemme de Morphismes
\ref{morphisms-lemma-morphism-finite-type-bounded-dimension}.
Alors $Rf_*$ est de dimension cohomologique finie ; voir le
lemme de Cohomologie étale
\ref{etale-cohomology-lemma-cohomological-dimension-proper}.
Ainsi, d'après le lemme \ref{derived-lemma-unbounded-right-derived} de Catégories dérivées,
il suffit de montrer que $R^qf_*(g'_!\mathcal{J}) = 0$ pour $q > 0$
et pour tout faisceau injectif de $\Lambda$-modules $\mathcal{J}$
sur $X'_\etale$.

\medskip\noindent
La fibre de $R^qf_*(g'_!\mathcal{J})$ en un point géométrique
$\overline{y}$ est égale à
$H^q(X_{\overline{y}}, (g'_!\mathcal{J})|_{X_{\overline{y}}})$ d'après le
lemme de Cohomologie étale
\ref{etale-cohomology-lemma-proper-base-change-stalk}.
Puisque la formation de $g'_!$ commute au changement de base
(lemme \ref{more-etale-lemma-base-change-f-shriek-separated}),
ce groupe est égal à
$$
H^q(X_{\overline{y}}, g'_{\overline{y}, !}(\mathcal{J}|_{X'_{\overline{y}}}))
$$
où $g'_{\overline{y}} : X'_{\overline{y}} \to X_{\overline{y}}$
est le morphisme induit entre les fibres géométriques.
Puisque $Y' \to Y$ est localement quasi-fini,
$X'_{\overline{y}}$ est la somme disjointe des fibres
$X'_{\overline{y}'}$ aux points géométriques $\overline{y}'$ de $Y'$
au-dessus de $\overline{y}$. Notons
$g'_{\overline{y}'} : X'_{\overline{y}'} \to X_{\overline{y}}$
la restriction de $g'_{\overline{y}}$ à $X'_{\overline{y}'}$.
Le groupe de cohomologie précédent est donc égal à
$$
H^q(X_{\overline{y}},
\bigoplus\nolimits_{\overline{y}'/\overline{y}}
g'_{\overline{y}', !}(\mathcal{J}|_{X'_{\overline{y}'}}))
$$
par exemple d'après le lemme \ref{more-etale-lemma-colim-f-shriek-separated} (mais cela
résulte aussi immédiatement de la définition de $g'_{\overline{y}, !}$
à la section \ref{more-etale-section-compact-support}).
Comme la cohomologie étale sur $X_{\overline{y}}$
commute aux sommes directes
(théorème \ref{etale-cohomology-theorem-colimit} de Cohomologie étale),
il suffit de montrer que
$$
H^q(X_{\overline{y}}, g'_{\overline{y}', !}(\mathcal{J}|_{X'_{\overline{y}'}}))
$$
est nul. Remarquons que
$g_{\overline{y}'} : X'_{\overline{y}'} \to X_{\overline{y}}$
est un morphisme entre schémas propres sur $\overline{y}$ et est donc
lui-même propre. Comme il est aussi localement quasi-fini, on en conclut que
$g_{\overline{y}'}$ est fini. Par conséquent,
$g'_{\overline{y}', !} = g'_{\overline{y}', *} = Rg'_{\overline{y}', *}$.
La suite spectrale de Leray nous ramène à montrer que
$$
H^q(X'_{\overline{y}'}, \mathcal{J}|_{X'_{\overline{y}'}})
$$
est nul. Comme $\Lambda$ est de torsion, cela résulte du changement de base propre
(lemme de Cohomologie étale
\ref{etale-cohomology-lemma-proper-base-change-stalk})
puisque les images directes supérieures de $\mathcal{J}$ par $f'$ sont nulles.

\medskip\noindent
Démonstration dans le cas (a). Nous allons le déduire du cas (b) par des arguments standards.
Nous montrerons que l'application induite $g_! R^pf'_* K \to R^pf_*(g'_!K)$
est un isomorphisme pour tout $p \in \mathbf{Z}$. Fixons un entier
$p_0 \in \mathbf{Z}$. Soit $a$ un entier tel que $H^j(K) = 0$
pour $j < a$. Nous allons démontrer que $g_! R^pf'_* K \to R^pf_*(g'_!K)$
est un isomorphisme pour $p \leq p_0$, par récurrence descendante sur $a$. Si
$a > p_0$, les deux membres de
l'application sont nuls pour $p \leq p_0$ par annulation triviale ; voir le
lemme \ref{derived-lemma-negative-vanishing} de Catégories dérivées
(et utiliser que $g_!$ et $g'_!$ sont des foncteurs exacts).
Supposons $a \leq p_0$. Considérons le triangle distingué
$$
H^a(K)[-a] \to K \to \tau_{\geq a + 1}K
$$
Par récurrence, le résultat vaut pour $\tau_{\geq a + 1}K$.
Au paragraphe suivant, nous le démontrerons pour $H^a(K)[-a]$.
Le lemme des cinq, appliqué au morphisme entre les
suites exactes longues de faisceaux de cohomologie
associées au morphisme de triangles distingués
$$
\xymatrix{
g_! Rf'_*(H^a(K)[-a]) \ar[d] \ar[r] &
g_! Rf'_* K \ar[r] \ar[d] &
g_! Rf'_* \tau_{\geq a + 1} K \ar[d] \\
Rf_*(g'_!(H^a(K)[-a])) \ar[r] &
Rf_*(g'_!K) \ar[r] &
Rf_*(g'_!\tau_{|geq a + 1}K)
}
$$
donne alors le résultat pour $K$. Nous omettons certains détails.

\medskip\noindent
Soit $\mathcal{F}$ un faisceau abélien de torsion sur $X'_\etale$.
Pour achever la démonstration, montrons que
$g_! Rf'_*\mathcal{F} \to R^pf_*(g'_!\mathcal{F})$
est un isomorphisme pour tout $p$.
On peut écrire $\mathcal{F} = \bigcup \mathcal{F}[n]$, où
$\mathcal{F}[n] = \Ker(n : \mathcal{F} \to \mathcal{F})$.
Nous avons l'isomorphisme pour $\mathcal{F}[n]$ d'après le cas (b).
Puisque les foncteurs $g_!$, $g'_!$, $R^pf_*$, $R^pf'_*$ commutent
aux limites inductives filtrantes (cela résulte du
lemme \ref{more-etale-lemma-lqf-f-shriek-separated-colimits} et du
lemme de Cohomologie étale
\ref{etale-cohomology-lemma-relative-colimit-general})
la démonstration est achevée.
\end{proof}

\begin{lemma}
\label{more-etale-lemma-shriek-proper-and-open-compose}
Considérons un diagramme commutatif de schémas
$$
\xymatrix{
X' \ar[r]_k \ar[d]_{f'} & X \ar[d]^f \\
Y' \ar[r]_l \ar[d]_{g'} & Y \ar[d]^g \\
Z' \ar[r]^m & Z
}
$$
où $f$, $f'$, $g$ et $g'$ sont propres, et
$k$, $l$ et $m$ séparés et localement quasi-finis.
Alors les isomorphismes du lemme \ref{more-etale-lemma-shriek-proper-and-open}
associés aux deux carrés se composent pour donner l'isomorphisme
associé au rectangle extérieur (voir la démonstration pour un énoncé précis).
\end{lemma}

\begin{proof}
L'énoncé signifie que, si l'on écrit
$R(g \circ f)_* = Rg_* \circ Rf_*$ et
$R(g' \circ f')_* = Rg'_* \circ Rf'_*$,
alors l'isomorphisme
$m_! \circ Rg'_* \circ Rf'_* \to Rg_* \circ Rf_* \circ k_!$
associé au rectangle extérieur est égal à la composée
$$
m_! \circ Rg'_* \circ Rf'_* \to
Rg_* \circ l_! \circ Rf'_* \to
Rg_* \circ Rf_* \circ k_!
$$
des deux morphismes associés aux carrés du diagramme. Pour le montrer, choisissons
un complexe K-injectif $\mathcal{J}^\bullet$ de $\Lambda$-modules
sur $X'_\etale$ et un quasi-isomorphisme
$k_!\mathcal{J}^\bullet \to \mathcal{I}^\bullet$
vers un complexe K-injectif $\mathcal{I}^\bullet$ de $\Lambda$-modules
sur $X_\etale$. La démonstration du lemme \ref{more-etale-lemma-shriek-proper-and-open}
montre que le morphisme canonique
$$
a : l_!f'_*\mathcal{J}^\bullet \to f_*\mathcal{I}^\bullet
$$
est un quasi-isomorphisme, lequel fournit la
deuxième flèche après application de $Rg_*$. D'après le
lemme \ref{sites-cohomology-lemma-K-injective-flat} de Cohomologie sur les sites,
le complexe $f_*\mathcal{I}^\bullet$, resp.\ $f'_*\mathcal{J}^\bullet$,
est un complexe K-injectif de $\Lambda$-modules sur 
$Y_\etale$, resp.\ $Y'_\etale$.
(Ce recours est un peu abusif et pourrait être évité.)
En particulier, le même raisonnement montre que le morphisme canonique
$$
b : m_!g'_*f'_*\mathcal{J}^\bullet \to g_*f_*\mathcal{I}^\bullet
$$
est un quasi-isomorphisme, et ce quasi-isomorphisme représente
la première flèche. Enfin, la démonstration du lemme \ref{more-etale-lemma-shriek-proper-and-open}
montre que $g_*l_!f'_!\mathcal{J}^\bullet$ représente
$Rg_*(l_!f'_*\mathcal{J}^\bullet)$, car $f'_*\mathcal{J}^\bullet$
est K-injectif. Ainsi $Rg_*(a) = g_*(a)$, et la composée
$g_*(a) \circ b$ est la flèche du lemme \ref{more-etale-lemma-shriek-proper-and-open}
associée au rectangle.
\end{proof}

\begin{lemma}
\label{more-etale-lemma-shriek-proper-and-open-compose-horizontal}
Considérons un diagramme commutatif de schémas
$$
\xymatrix{
X'' \ar[r]_{g'} \ar[d]_{f''} &
X' \ar[r]_g \ar[d]_{f'} &
X \ar[d]^f \\
Y'' \ar[r]^{h'} &
Y' \ar[r]^h &
Y
}
$$
où $f$, $f'$ et $f''$ sont propres, et
$g$, $g'$, $h$ et $h'$ séparés et localement quasi-finis.
Alors les isomorphismes du lemme \ref{more-etale-lemma-shriek-proper-and-open}
associés aux deux carrés se composent pour donner l'isomorphisme
associé au rectangle extérieur (voir la démonstration pour un énoncé précis).
\end{lemma}

\begin{proof}
L'énoncé signifie que, si l'on écrit
$(h \circ h')_! = h_! \circ h'_!$ et
$(g \circ g')_! = g_! \circ g'_!$
au moyen des égalités du lemme \ref{more-etale-lemma-f-shriek-composition},
alors l'isomorphisme
$h_! \circ h'_! \circ Rf''_* \to Rf_* \circ g_! \circ g'_!$
associé au rectangle extérieur est égal à la composée
$$
h_! \circ h'_! \circ Rf''_* \to
h_! \circ Rf'_* \circ g'_! \to
Rf_* \circ g_! \circ g'_!
$$
des deux morphismes associés aux carrés du diagramme. Pour le montrer, choisissons
un complexe K-injectif $\mathcal{I}^\bullet$ de $\Lambda$-modules
sur $X''_\etale$ et un quasi-isomorphisme
$g'_!\mathcal{I}^\bullet \to \mathcal{J}^\bullet$
vers un complexe K-injectif $\mathcal{J}^\bullet$ de $\Lambda$-modules
sur $X'_\etale$. Choisissons ensuite un quasi-isomorphisme
$g_!\mathcal{J}^\bullet \to \mathcal{K}^\bullet$
vers un complexe K-injectif $\mathcal{K}^\bullet$ de $\Lambda$-modules
sur $X_\etale$.
La démonstration du lemme \ref{more-etale-lemma-shriek-proper-and-open} montre que les morphismes
canoniques
$$
h'_!f''_*\mathcal{I}^\bullet \to f'_*\mathcal{J}^\bullet
\quad\text{et}\quad
h_!f'_*\mathcal{J}^\bullet \to f_*\mathcal{K}^\bullet
$$
sont des quasi-isomorphismes, qui définissent respectivement la
première et la deuxième flèche ci-dessus. Puisque $g_!$ est un foncteur exact
(lemme \ref{more-etale-lemma-lqf-f-shriek-separated-colimits}),
on voit que $g_!g'_!\mathcal{I}^\bullet \to \mathcal{K}^\bullet$
est un quasi-isomorphisme ; par conséquent, le morphisme canonique
$$
h_!h'_!f''_*\mathcal{I}^\bullet \to f_*\mathcal{K}^\bullet
$$
est un quasi-isomorphisme et représente le morphisme associé au rectangle
extérieur dans la catégorie dérivée. Il est clair que ce morphisme est la
composée des deux autres, ce qui achève la démonstration.
\end{proof}

\begin{remark}
\label{more-etale-remark-going-around}
Considérons un diagramme commutatif
$$
\xymatrix{
X'' \ar[r]_{k'} \ar[d]_{f''} & X' \ar[r]_k \ar[d]_{f'} & X \ar[d]^f \\
Y'' \ar[r]^{l'} \ar[d]_{g''} & Y' \ar[r]^l \ar[d]_{g'} & Y \ar[d]^g \\
Z'' \ar[r]^{m'} & Z' \ar[r]^m & Z
}
$$
de schémas dont les flèches verticales sont propres et les flèches horizontales
séparées et localement quasi-finies.
Étiquetons les carrés du diagramme $A$, $B$, $C$, $D$
comme suit :
$$
\begin{matrix}
A & B \\
C & D
\end{matrix}
$$
Les morphismes du lemme \ref{more-etale-lemma-shriek-proper-and-open}
associés aux carrés sont alors (où l'on utilise $Rf_* = f_*$, etc.)
$$
\begin{matrix}
\gamma_A : l'_! \circ f''_* \to f'_* \circ k'_! &
\gamma_B : l_! \circ f'_* \to f_* \circ k_! \\
\gamma_C : m'_! \circ g''_* \to g'_* \circ l'_! &
\gamma_D : m_! \circ g'_* \to g_* \circ l_!
\end{matrix}
$$
Pour les rectangles $2 \times 1$ et $1 \times 2$, nous disposons de quatre autres
morphismes
$$
\begin{matrix}
\gamma_{A + B} :
(l \circ l')_! \circ f''_* \to f_* \circ (k \circ k')_*  \\
\gamma_{C + D} :
(m \circ m')_! \circ g''_* \to g_* \circ (l \circ l')_! \\
\gamma_{A + C} :
m'_! \circ (g'' \circ f'')_* \to (g' \circ f')_* \circ k'_! \\
\gamma_{B + D} :
m_! \circ (g' \circ f')_* \to (g \circ f)_* \circ k_!
\end{matrix}
$$
D'après le lemme \ref{more-etale-lemma-shriek-proper-and-open-compose-horizontal}, on a
$$
\gamma_{A + B} = \gamma_B \circ \gamma_A, \quad
\gamma_{C + D} = \gamma_D \circ \gamma_C
$$
et, d'après le lemme \ref{more-etale-lemma-shriek-proper-and-open-compose}, on a
$$
\gamma_{A + C} = \gamma_A \circ \gamma_C, \quad
\gamma_{B + D} = \gamma_B \circ \gamma_D
$$
Il serait ici plus exact d'écrire
$\gamma_{A + B} = (\gamma_B \star \text{id}_{k'_!}) \circ
(\text{id}_{l_!} \star \gamma_A)$ avec les notations de la
section \ref{categories-section-formal-cat-cat} de Catégories
et de même pour les autres.
Cela étant, nous obtenons a priori deux transformations
$$
m_! \circ m'_! \circ g''_* \circ f''_*
\longrightarrow
g_* \circ f_* \circ k_! \circ k'_!
$$
à savoir
$$
\gamma_B \circ \gamma_D \circ \gamma_A \circ \gamma_C =
\gamma_{B + D} \circ \gamma_{A + C}
$$
et
$$
\gamma_B \circ \gamma_A \circ \gamma_D \circ \gamma_C =
\gamma_{A + B} \circ \gamma_{C + D}
$$
Le but de cette remarque est de signaler que ces transformations
sont égales. Pour le voir, il suffit en effet de montrer que
$$
\xymatrix{
m_! \circ g'_* \circ l'_! \circ f''_* \ar[r]_{\gamma_D} \ar[d]_{\gamma_A} &
g_* \circ l_! \circ l'_! \circ f''_* \ar[d]^{\gamma_A} \\
m_! \circ g'_* \circ f'_* \circ k'_! \ar[r]^{\gamma_D} &
g_* \circ l_! \circ f'_* \circ k'_!
}
$$
commute. Cela résulte du fait que les carrés $A$ et $D$ ne se rencontrent qu'en
un point ; plus précisément, cela résulte du
lemme \ref{categories-lemma-properties-2-cat-cats} de Catégories,
ou, plus simplement, de la discussion qui précède la
définition \ref{categories-definition-horizontal-composition} de Catégories.
\end{remark}

\begin{lemma}
\label{more-etale-lemma-shriek-proper-and-open-base-change}
Soit $b : Y_1 \to Y$ un morphisme de schémas. Considérons un diagramme commutatif
de schémas
$$
\vcenter{
\xymatrix{
X' \ar[r]_{g'} \ar[d]_{f'} & X \ar[d]^f \\
Y' \ar[r]^g & Y
}
}
\quad\text{et soit}\quad
\vcenter{
\xymatrix{
X'_1 \ar[r]_{g'_1} \ar[d]_{f'_1} & X_1 \ar[d]^{f_1} \\
Y'_1 \ar[r]^{g_1} & Y_1
}
}
$$
son changement de base par $b$. Supposons $f$ et $f'$ propres, et
$g$ et $g'$ séparés et localement quasi-finis.
Pour un anneau $\Lambda$ et $K$ dans $D(X'_\etale, \Lambda)$,
le diagramme suivant est commutatif :
$$
\xymatrix{
b^{-1}g_!Rf'_*K \ar[d] \ar[r] &
g_{1, !}(b')^{-1}Rf'_*K \ar[r] &
g_{1, !}Rf'_{1, *}(a')^{-1}K \ar[d] \\
b^{-1}Rf_*g'_!K \ar[r] &
Rf_{1, *}a^{-1}g'_!K \ar[r] &
Rf_{1, *}g'_{1, !}(a')^{-1}K
}
$$
dans $D(Y_{1, \etale}, \Lambda)$, où $a : X_1 \to X$, $a' : X'_1 \to X'$,
$b' : Y'_1 \to Y'$ sont les projections, les flèches verticales sont celles
du lemme \ref{more-etale-lemma-shriek-proper-and-open},
et les flèches horizontales sont la flèche de changement de base
(de la section
\ref{etale-cohomology-section-base-change-preliminaries} de Cohomologie étale)
et la flèche de changement de base du lemme \ref{more-etale-lemma-base-change-f-shriek-separated}.
\end{lemma}

\begin{proof}
Représentons $K$ par un complexe K-injectif $\mathcal{J}^\bullet$ de
faisceaux de $\Lambda$-modules sur $X'_\etale$. Choisissons un quasi-isomorphisme
$g'_!\mathcal{J}^\bullet \to \mathcal{I}^\bullet$ vers un
complexe K-injectif $\mathcal{I}^\bullet$ de
faisceaux de $\Lambda$-modules sur $X_\etale$.
La démonstration du lemme \ref{more-etale-lemma-shriek-proper-and-open}
construit $g_!Rf'_*K \to Rf_*g'_!K$ sous la forme
$$
g_!f'_*\mathcal{J}^\bullet =
g_!f'_!\mathcal{J}^\bullet =
f_!g'_!\mathcal{J}^\bullet =
f_*g'_!\mathcal{J}^\bullet \to
f_*\mathcal{I}^\bullet
$$
Choisissons un quasi-isomorphisme
$(a')^{-1}\mathcal{J}^\bullet \to \mathcal{J}_1^\bullet$
vers un complexe K-injectif $\mathcal{J}_1^\bullet$ de
faisceaux de $\Lambda$-modules sur $X'_{1, \etale}$.
On peut alors choisir un diagramme de complexes
$$
\xymatrix{
g'_{1, !}\mathcal{J}_1^\bullet \ar[rr] & &
\mathcal{I}_1^\bullet \\
g'_{1, !}(a')^{-1}\mathcal{J}^\bullet \ar[u] \ar@{=}[r] &
a^{-1}g'_!\mathcal{J}^\bullet \ar[r] &
a^{-1}\mathcal{I}^\bullet \ar[u]
}
$$
commutatif à homotopie près, où toutes les flèches sont des quasi-isomorphismes ;
l'égalité résulte du lemme \ref{more-etale-lemma-proper-f-shriek},
et $\mathcal{I}_1^\bullet$ est un complexe K-injectif de faisceaux de
$\Lambda$-modules sur $X_{1, \etale}$. Le morphisme
$g_{1, !}Rf'_{1, *}(a')^{-1}K \to Rf_{1, *}g'_{1, !}(a')^{-1}K$ est donné par
$$
g_{1, !}f'_{1, *}\mathcal{J}_1^\bullet =
g_{1, !}f'_{1, !}\mathcal{J}_1^\bullet =
f_{1, !}g'_{1, !}\mathcal{J}_1^\bullet =
f_{1, *}g'_{1, !}\mathcal{J}_1^\bullet \to
f_{1, *}\mathcal{I}_1^\bullet
$$
Les identifications correspondant aux $3$ signes d'égalité dans les deux flèches
sont compatibles avec les morphismes de changement de base, c'est-à-dire que le diagramme
$$
\xymatrix{
b^{-1}g_!f'_*\mathcal{J}^\bullet \ar@{=}[d] \ar[r] &
g_{1, !}(b')^{-1}f'_*\mathcal{J}^\bullet \ar[r] &
g_{1, !}f'_{1, *}(a')^{-1}\mathcal{J}^\bullet \ar@{=}[d] \\
b^{-1}f_*g'_!\mathcal{J}^\bullet \ar[r] &
f_{1, *}a^{-1}g'_!\mathcal{J}^\bullet \ar[r] &
f_{1, *}g'_{1, !}(a')^{-1}\mathcal{J}^\bullet
}
$$
de complexes de faisceaux abéliens commute. Pour le montrer, il suffit de
vérifier la commutativité du diagramme après avoir remplacé $g_!, g_{1, !}, g'_!, g'_{1, !}$
par $g_*, g_{1, *}, g'_*, g'_{1, *}$ (car les foncteurs d'image directe à support propre
sont définis comme des sous-foncteurs des foncteurs $*$ et les
morphismes de changement de base sont définis de manière compatible avec cette inclusion ; voir la
démonstration du lemme \ref{more-etale-lemma-base-change-f-shriek-separated}).
Pour ce nouveau diagramme, la commutativité résulte de la compatibilité
des morphismes de changement de base avec la composition horizontale et verticale des diagrammes ; voir les
remarques \ref{sites-remark-compose-base-change} et
\ref{sites-remark-compose-base-change-horizontal}
de Sites, de sorte que chacun des deux parcours du diagramme donne le morphisme de
changement de base de $f \circ g' = g \circ f'$ par $b$.
Puisque, bien sûr,
$$
\xymatrix{
g_{1, !}f'_{1, *}(a')^{-1}\mathcal{J}^\bullet \ar@{=}[d] \ar[r] &
g_{1, !}f'_{1, *}\mathcal{J}_1^\bullet \ar@{=}[d] \\
f_{1, *}g'_{1, !}(a')^{-1}\mathcal{J}^\bullet \ar[r] &
f_{1, *}g'_{1, !}\mathcal{J}_1^\bullet 
}
$$
commute, il suffit, pour achever la démonstration, de montrer que
$$
\xymatrix{
b^{-1}f_*g'_!\mathcal{J}^\bullet \ar[r] \ar[d] &
f_{1, *}a^{-1}g'_!\mathcal{J}^\bullet \ar[r] \ar[d] &
f_{1, *}g'_{1, !}(a')^{-1}\mathcal{J}^\bullet \ar[r] &
f_{1, *}g'_{1, !}\mathcal{J}_1^\bullet \ar[d] \\
b^{-1}f_*\mathcal{I}^\bullet \ar[r] &
f_{1, *}a^{-1}\mathcal{I}^\bullet \ar[rr] & &
f_{1, *}\mathcal{I}_1^\bullet
}
$$
commute dans la catégorie dérivée, ce qui résulte du choix des morphismes
effectué plus haut.
\end{proof}

\begin{lemma}
\label{more-etale-lemma-shriek-lqf-and-proper}
Considérons un diagramme commutatif de schémas
$$
\xymatrix{
X \ar[r]_f \ar[rd]_g & Y \ar[d]^h \\
& Z
}
$$
où $f$ et $g$ sont localement quasi-finis et $h$ propre. Soit $\Lambda$ un anneau.
Fonctoriellement en $K \in D(X_\etale, \Lambda)$, il existe un morphisme canonique
$$
g_!K \longrightarrow Rh_*(f_!K)
$$
dans $D(Z_\etale, \Lambda)$. Ce morphisme est un isomorphisme si
(a) $K$ est borné inférieurement et ses faisceaux de cohomologie sont de torsion, ou si
(b) $\Lambda$ est un anneau de torsion.
\end{lemma}

\begin{proof}
C'est un cas particulier du lemme \ref{more-etale-lemma-shriek-proper-and-open}
si $f$ et $g$ sont séparés. Nous recommandons au lecteur de ne pas lire la démonstration
dans le cas général,
car nous utiliserons principalement le cas où $f$ et $g$ sont séparés.

\medskip\noindent
Représentons $K$ par un complexe $\mathcal{K}^\bullet$ de faisceaux de
$\Lambda$-modules sur $X_\etale$. Choisissons un quasi-isomorphisme
$f_!\mathcal{K}^\bullet \to \mathcal{I}^\bullet$ vers un complexe K-injectif
$\mathcal{I}^\bullet$ de faisceaux de $\Lambda$-modules sur $Y_\etale$.
Considérons le morphisme
$$
g_!\mathcal{K}^\bullet =
h_!f_!\mathcal{K}^\bullet =
h_*f_!\mathcal{K}^\bullet
\longrightarrow
h_*\mathcal{I}^\bullet
$$
où les égalités résultent des
lemmes \ref{more-etale-lemma-lqf-separated-shriek-composition}
et \ref{more-etale-lemma-proper-f-shriek}. Ce morphisme de complexes détermine le morphisme
$g_!K \to Rh_*(f_!K)$ de l'énoncé du lemme.

\medskip\noindent
Supposons que $\Lambda$ soit un anneau de torsion, c'est-à-dire que nous sommes dans le cas (b). Pour vérifier que le morphisme
est un isomorphisme, nous pouvons travailler localement sur $Z$.
Nous pouvons donc supposer que la dimension des fibres de $h$ est bornée ;
voir le lemme de Morphismes
\ref{morphisms-lemma-morphism-finite-type-bounded-dimension}.
Alors $Rh_*$ est de dimension cohomologique finie ; voir le
lemme de Cohomologie étale
\ref{etale-cohomology-lemma-cohomological-dimension-proper}.
Ainsi, d'après le lemme \ref{derived-lemma-unbounded-right-derived} de Catégories dérivées,
si nous montrons que $R^qh_*(f_!\mathcal{F}) = 0$ pour $q > 0$
et pour tout faisceau $\mathcal{F}$ de $\Lambda$-modules sur $X_\etale$, alors
$h_*f_!\mathcal{K}^\bullet \to h_*\mathcal{I}^\bullet$ est un quasi-isomorphisme.

\medskip\noindent
Remarquons que $\mathcal{G} = f_!\mathcal{F}$ est un faisceau de $\Lambda$-modules
sur $Y$ dont les fibres ne sont non nulles qu'aux points $y \in Y$ tels que
$\kappa(y)/\kappa(h(y))$ soit une extension finie. Cela résulte de la
description des fibres de $f_!\mathcal{F}$ au
lemme \ref{more-etale-lemma-lqf-f-shriek-stalk}
et du fait que $f$ et $g$ sont tous deux localement quasi-finis.
Ainsi, d'après le théorème de changement de base propre (lemme de Cohomologie étale
\ref{etale-cohomology-lemma-proper-base-change-stalk})
il suffit de montrer que $H^q(Y_{\overline{z}}, \mathcal{H}) = 0$,
où $\mathcal{H}$ est un faisceau sur le schéma propre $Y_{\overline{z}}$
sur $\kappa(\overline{z})$, dont le support est contenu dans l'ensemble
des points fermés. L'annulation voulue résulte donc du lemme de Cohomologie étale
\ref{etale-cohomology-lemma-supported-in-closed-points}.

\medskip\noindent
Le cas (a) se déduit du cas (b) par exactement le même argument que celui employé
dans la démonstration du lemme \ref{more-etale-lemma-shriek-proper-and-open}
(en utilisant le lemme \ref{more-etale-lemma-lqf-f-shriek-stalk} à la place du
lemme \ref{more-etale-lemma-lqf-f-shriek-separated-colimits}).
\end{proof}







\section{Image directe dérivée à support propre au moyen de compactifications}
\sectionmark{Image directe dérivée par compactification}
\label{more-etale-section-derived-lower-shriek-compactification}

\noindent
Soit $f : X \to Y$ un morphisme de schémas séparé et de type fini,
où $Y$ est quasi-compact et quasi-séparé. Choisissons une compactification
$j : X \to \overline{X}$ sur $Y$ ; voir le
théorème \ref{flat-theorem-nagata} de Compléments sur la platitude.
Soit $\Lambda$ un anneau. Notons $D^+_{tors}(X_\etale, \Lambda)$
la sous-catégorie triangulée strictement pleine et saturée de
$D(X_\etale, \Lambda)$ formée des objets $K$ bornés inférieurement
et dont les faisceaux de cohomologie sont de torsion. Nous considérerons le
foncteur
$$
Rf_! = R\overline{f}_* \circ j_! :
D^+_{tors}(X_\etale, \Lambda)
\longrightarrow
D^+_{tors}(Y_\etale, \Lambda)
$$
où $\overline{f} : \overline{X} \to Y$ est le morphisme structural.
Cette définition a un sens : le foncteur $j_!$ envoie $D^+_{tors}(X_\etale, \Lambda)$ dans
$D^+_{tors}(\overline{X}_\etale, \Lambda)$ d'après la
remarque \ref{more-etale-remark-lqf-shriek-derived-torsion}, et $R\overline{f}_*$ envoie
$D^+_{tors}(\overline{X}_\etale, \Lambda)$ dans $D^+_{tors}(Y_\etale, \Lambda)$
d'après le lemme \ref{etale-cohomology-lemma-torsion-direct-image} de Cohomologie étale.
Si $\Lambda$ est un anneau de torsion, nous définissons alors
$$
Rf_! = R\overline{f}_* \circ j_! :
D(X_\etale, \Lambda)
\longrightarrow
D(Y_\etale, \Lambda)
$$
Voici le lemme indispensable.

\begin{lemma}
\label{more-etale-lemma-shriek-well-defined}
Soit $f : X \to Y$ un morphisme séparé de type fini entre schémas quasi-compacts
et quasi-séparés. Les foncteurs $Rf_!$ construits ci-dessus sont,
à isomorphisme canonique près, indépendants du choix de la
compactification.
\end{lemma}

\begin{proof}
Nous le démontrerons pour le foncteur
$Rf_! : D(X_\etale, \Lambda) \to D(Y_\etale, \Lambda)$
quand $\Lambda$ est un anneau de torsion ; le cas du foncteur
$Rf_! : D^+_{tors}(X_\etale, \Lambda) \to D^+_{tors}(Y_\etale, \Lambda)$
se démontre exactement de la même manière.

\medskip\noindent
Considérons la catégorie des compactifications de $X$ sur $Y$, qui est
cofiltrante d'après le théorème \ref{flat-theorem-nagata} et les
lemmes \ref{flat-lemma-compactifications-cofiltered} et
\ref{flat-lemma-compactifyable} de Compléments sur la platitude.
À toute compactification
$$
j : X \to \overline{X},\quad \overline{f} : \overline{X} \to Y
$$
la construction précédente associe le foncteur $R\overline{f}_* \circ j_! :
D(X_\etale, \Lambda) \to D(Y_\etale, \Lambda)$.
Précisons un peu. Étant donné un complexe $\mathcal{K}^\bullet$
de faisceaux de $\Lambda$-modules sur $X_\etale$, choisissons un quasi-isomorphisme
$j_!\mathcal{K}^\bullet \to \mathcal{I}^\bullet$ vers un complexe K-injectif
de faisceaux de $\Lambda$-modules sur $\overline{X}_\etale$.
Notre foncteur envoie alors $\mathcal{K}^\bullet$ sur
$\overline{f}_*\mathcal{I}^\bullet$.

\medskip\noindent
Supposons donné un morphisme $g : \overline{X}_1 \to \overline{X}_2$
entre les compactifications $j_i : X \to \overline{X}_i$ sur $Y$.
On obtient alors un isomorphisme
$$
R\overline{f}_{2, *} \circ j_{2, !} =
R\overline{f}_{2, *} \circ Rg_* \circ j_{1, !} =
R\overline{f}_{1, *} \circ j_{1, !}
$$
en utilisant le lemme \ref{more-etale-lemma-shriek-lqf-and-proper} pour la première égalité.

\medskip\noindent
Pour achever la démonstration, puisque la catégorie des compactifications de $X$ sur $Y$
est cofiltrante, il suffit de montrer que la composition des morphismes de
compactifications de $X$ sur $Y$ correspond à la composition
des isomorphismes de foncteurs\footnote{En effet, si $\alpha, \beta : F \to G$
sont des morphismes de foncteurs et $\gamma : G \to H$ un isomorphisme
de foncteurs tel que $\gamma \circ \alpha = \gamma \circ \beta$, alors
on conclut que $\alpha = \beta$.}. Pour cela, supposons que
$j_3 : X \to \overline{X}_3$
soit une troisième compactification et que $h : \overline{X}_2 \to \overline{X}_3$
soit un morphisme de compactifications. Il faut alors montrer que la
composée
$$
R\overline{f}_{3, *} \circ j_{3, !} =
R\overline{f}_{3, *} \circ Rh_* \circ j_{2, !} =
R\overline{f}_{2, *} \circ j_{2, !} =
R\overline{f}_{2, *} \circ Rg_* \circ j_{1, !} =
R\overline{f}_{1, *} \circ j_{1, !}
$$
est égale à l'isomorphisme de foncteurs construit en utilisant seulement
$j_3$, $g \circ h$ et $j_1$. Un calcul montre qu'il suffit de
prouver que la composée des morphismes
$$
j_{3, !} \to Rh_* \circ j_{2, !} \to Rh_* \circ Rg_* \circ j_{1, !}
$$
du lemme \ref{more-etale-lemma-shriek-lqf-and-proper} coïncide avec le morphisme correspondant
$j_{3, !} \to R(h \circ g)_* \circ j_{1, !}$
via l'identification $R(h \circ g)_* = Rh_* \circ Rg_*$.
Comme le morphisme du lemme \ref{more-etale-lemma-shriek-lqf-and-proper}
est un cas particulier de celui du lemme \ref{more-etale-lemma-shriek-proper-and-open}
(puisque $j_1$ et $j_2$ sont séparés), cela résulte immédiatement du
lemme \ref{more-etale-lemma-shriek-proper-and-open-compose}.
\end{proof}

\begin{lemma}
\label{more-etale-lemma-shriek-composition}
Soient $f : X \to Y$ et $g : Y \to Z$ des morphismes séparés de type fini
entre schémas quasi-compacts et quasi-séparés. Il existe alors un
isomorphisme canonique $Rg_! \circ Rf_! \to R(g \circ f)_!$.
\end{lemma}

\begin{proof}
Choisissons une compactification $i : Y \to \overline{Y}$ de $Y$ sur $Z$.
Choisissons une compactification $X \to \overline{X}$ de $X$ sur
$\overline{Y}$. On applique deux fois le théorème \ref{flat-theorem-nagata}
et le lemme \ref{flat-lemma-compactifyable} de Compléments sur la platitude.
Soit $U$ l'image inverse de $Y$ dans $\overline{X}$ ;
on obtient ainsi le diagramme commutatif
$$
\xymatrix{
X \ar[r]_j \ar[d]_f &
U \ar[dl]^{f'} \ar[r]_{j'} &
\overline{X} \ar[dl]^{\overline{f}} \\
Y \ar[r]_i \ar[d]_g &
\overline{Y} \ar[dl]^{\overline{g}} \\
Z
}
$$
On a alors
\begin{align*}
R(g \circ f)_!
& =
R(\overline{g} \circ \overline{f})_* \circ (j' \circ j)_! \\
& =
R\overline{g}_* \circ R\overline{f}_* \circ j'_! \circ j_! \\
& =
R\overline{g}_* \circ i_! \circ Rf'_* \circ j_! \\
& =
Rg_! \circ Rf_!
\end{align*}
La première égalité est la définition de $R(g \circ f)_!$.
La deuxième utilise les identifications
$R(\overline{g} \circ \overline{f})_* = R\overline{g}_* \circ R\overline{f}_*$
et $(j' \circ j)_! = j'_! \circ j_!$ du lemme \ref{more-etale-lemma-f-shriek-composition}.
L'identification $i_! \circ Rf'_* \to R\overline{f}_* \circ j_!$
employée dans la troisième égalité est celle du lemme \ref{more-etale-lemma-shriek-proper-and-open}.
La quatrième et dernière égalité est la définition de $Rg_!$ et $Rf_!$.
Pour achever la démonstration, montrons que
cet isomorphisme est indépendant des choix effectués.

\medskip\noindent
Supposons donnés deux diagrammes
$$
\vcenter{
\xymatrix{
X \ar[r]_{j_1} \ar[d] &
U_1 \ar[dl]^{f_1} \ar[r]_{j'_1} &
\overline{X}_1 \ar[dl]^{\overline{f}_1} \\
Y \ar[r]_{i_1} \ar[d] &
\overline{Y}_1 \ar[dl]^{\overline{g}_1} \\
Z
}
}
\quad\text{et}\quad
\vcenter{
\xymatrix{
X \ar[r]_{j_2} \ar[d] &
U_2 \ar[dl]^{f_2} \ar[r]_{j'_2} &
\overline{X}_2 \ar[dl]^{\overline{f}_2} \\
Y \ar[r]_{i_2} \ar[d] &
\overline{Y}_2 \ar[dl]^{\overline{g}_2} \\
Z
}
}
$$
Nous pouvons d'abord choisir une compactification $i : Y \to \overline{Y}$
de $Y$ sur $Z$ qui domine à la fois $\overline{Y}_1$ et $\overline{Y}_2$ ;
voir le lemme \ref{flat-lemma-compactifications-cofiltered} de Compléments sur la platitude.
D'après le lemme \ref{flat-lemma-right-multiplicative-system} de Compléments sur la platitude et les
lemmes \ref{categories-lemma-morphisms-right-localization} et
\ref{categories-lemma-equality-morphisms-right-localization}
de Catégories, nous pouvons choisir une compactification $X \to \overline{X}$ de
$X$ sur $\overline{Y}$, avec des morphismes $\overline{X} \to \overline{X}_1$
et $\overline{X} \to \overline{X}_2$, telle que la composée
$\overline{X} \to \overline{Y} \to \overline{Y}_1$ soit égale à
la composée $\overline{X} \to \overline{X}_1 \to \overline{Y}_1$
et que la composée
$\overline{X} \to \overline{Y} \to \overline{Y}_2$ soit égale à
la composée $\overline{X} \to \overline{X}_2 \to \overline{Y}_2$.
Il suffit donc de comparer les morphismes
déterminés par nos diagrammes lorsque l'on dispose d'un diagramme commutatif
de la forme suivante :
$$
\xymatrix{
X \ar[rr]_{j_1} \ar@{=}[d] & &
U_1 \ar[d]^{h'} \ar[ddll] \ar[rr]_{j'_1} & &
\overline{X}_1 \ar[d]^h \ar[ddll] \\
X \ar'[r][rr]^-{j_2} \ar[d] & &
U_2 \ar'[dl][ddll] \ar'[r][rr]^-{j'_2} & &
\overline{X}_2 \ar[ddll] \\
Y \ar[rr]^{i_1} \ar@{=}[d] & & \overline{Y}_1 \ar[d]^k \\
Y \ar[rr]^{i_2} \ar[d] & & \overline{Y}_2 \ar[dll] \\
Z
}
$$
Chacun des carrés
$$
\xymatrix{
X \ar[r]_{j_1} \ar[d]_{\text{id}} \ar@{}[dr]|A &
U_1 \ar[d]^{h'} \\
X \ar[r]^{j_2} &
U_2
}
\quad
\xymatrix{
U_2 \ar[r]_{j_2'} \ar[d]_{f_2} \ar@{}[dr]|B &
\overline{X}_2 \ar[d]^{\overline{f}_2} \\
Y \ar[r]^{i_2} &
\overline{Y}_2
}
\quad
\xymatrix{
U_1 \ar[r]_{j_1'} \ar[d]_{f_1} \ar@{}[dr]|C &
\overline{X}_1 \ar[d]^{\overline{f}_1} \\
Y \ar[r]^{i_1} &
\overline{Y}_1
}
\quad
\xymatrix{
Y \ar[r]_{i_1} \ar[d]_{\text{id}} \ar@{}[dr]|D &
\overline{Y}_1 \ar[d]^k \\
Y \ar[r]^{i_2} &
\overline{Y}_2
}
\quad
\xymatrix{
X \ar[r]_{j_1' \circ j_1} \ar[d]_{\text{id}} \ar@{}[dr]|E &
\overline{X}_1 \ar[d]^h \\
X \ar[r]^{j_2} &
\overline{X}_2
}
$$
donne lieu à un isomorphisme comme suit :
\begin{align*}
\gamma_A & :
j_{2, !} \to Rh'_* \circ j_{1, !}  \\
\gamma_B & :
i_{2, !} \circ Rf_{2, *} \to R\overline{f}_{2, *} \circ j'_{2, !} \\
\gamma_C & :
i_{1, !} \circ Rf_{1, *} \to R\overline{f}_{1, *} \circ j'_{1, !} \\
\gamma_D & :
i_{2, !} \to Rk_* \circ i_{1, !} \\
\gamma_E & :
j_{2, !} \to Rh_* \circ (j'_1 \circ j_1)_!
\end{align*}
en appliquant la flèche du lemme \ref{more-etale-lemma-shriek-proper-and-open}
(qui est la même que celle du lemme \ref{more-etale-lemma-shriek-lqf-and-proper}
lorsque la flèche verticale gauche est l'identité). Posons
\begin{align*}
F_1 & = Rf_{1, *} \circ j_{1, !} \\
F_2 & = Rf_{2, *} \circ j_{2, !} \\
G_1 & = R\overline{g}_{1, *} \circ i_{1, !} \\
G_2 & = R\overline{g}_{2, *} \circ i_{2, !} \\
C_1 & = R(\overline{g}_1 \circ \overline{f}_1)_* \circ (j'_1 \circ j_1)_! \\
C_2 & = R(\overline{g}_2 \circ \overline{f}_2)_* \circ (j'_2 \circ j_2)_!
\end{align*}
La construction donnée dans le premier paragraphe de la démonstration
et dans le lemme \ref{more-etale-lemma-shriek-well-defined} utilise
\begin{enumerate}
\item $\gamma_C$ pour la flèche $G_1 \circ F_1 \to C_1$,
\item $\gamma_B$ pour la flèche $G_2 \circ F_2 \to C_2 $,
\item $\gamma_A$ pour la flèche $F_2 \to F_1$,
\item $\gamma_D$ pour la flèche $G_2 \to G_1$, et
\item $\gamma_E$ pour la flèche $C_2 \to C_1$.
\end{enumerate}
Il faut donc montrer que le diagramme
$$
\xymatrix{
C_2 \ar[rr]_{\gamma_E} & &
C_1 \\
G_2 \circ F_2 \ar[rr]^{\gamma_D \circ \gamma_A} \ar[u]^{\gamma_B} & &
G_1 \circ F_1 \ar[u]_{\gamma_C}
}
$$
est commutatif. Nous utiliserons les
lemmes \ref{more-etale-lemma-shriek-proper-and-open-compose} et
\ref{more-etale-lemma-shriek-proper-and-open-compose-horizontal}
et reprendrons, avec le même léger abus, les notations de la
remarque \ref{more-etale-remark-going-around} (en omettant notamment, dans la notation,
les produits $\star$ par des transformations
identiques).
On peut écrire $\gamma_E = \gamma_F \circ \gamma_A$, où
$$
\xymatrix{
U_1 \ar[r]_{j'_1} \ar[d]_{h'} \ar@{}[rd]|F &
\overline{X}_1 \ar[d]^h \\
U_2 \ar[r]^{j'_2} &
\overline{X}_2
}
$$
On a donc
$$
\gamma_E \circ \gamma_B = \gamma_F \circ \gamma_A  \circ \gamma_B
= \gamma_F \circ \gamma_B \circ \gamma_A
$$
la dernière égalité venant de ce que les deux carrés $A$ et $B$ ne
se rencontrent qu'en un point (comme dans le dernier argument de la
remarque \ref{more-etale-remark-going-around}). Il suffit donc de prouver que
$\gamma_C \circ \gamma_D = \gamma_F \circ \gamma_B$.
Comme chacune de ces deux flèches est égale à celle associée au carré
$$
\xymatrix{
U_1 \ar[r] \ar[d] & \overline{X}_1 \ar[d] \\
Y \ar[r] & \overline{Y}_2
}
$$
on conclut.
\end{proof}

\begin{lemma}
\label{more-etale-lemma-pseudo-functor}
Soient $f : X \to Y$, $g : Y \to Z$, $h : Z \to T$ des morphismes séparés de
type fini entre schémas quasi-compacts et quasi-séparés. Alors
le diagramme
$$
\xymatrix{
Rh_! \circ Rg_! \circ Rf_! \ar[r]_{\gamma_C} \ar[d]^{\gamma_A} &
R(h \circ g)_! \circ Rf_! \ar[d]_{\gamma_{A + B}} \\
Rh_! \circ R(g \circ f)_! \ar[r]^{\gamma_{B + C}} &
R(h \circ g \circ f)_!
}
$$
d'isomorphismes du lemme \ref{more-etale-lemma-shriek-composition} est commutatif
(pour la signification des $\gamma$, voir la démonstration).
\end{lemma}

\begin{proof}
Pour cela, choisissons une compactification $\overline{Z}$
de $Z$ sur $T$, puis une compactification $\overline{Y}$ de $Y$ sur
$\overline{Z}$, et enfin une compactification $\overline{X}$ de
$X$ sur $\overline{Y}$. On utilise le
théorème \ref{flat-theorem-nagata} et le
lemme \ref{flat-lemma-compactifyable} de Compléments sur la platitude.
Soit $W \subset \overline{Y}$ l'image inverse de $Z$ par
$\overline{Y} \to \overline{Z}$, et soient $U \subset V \subset \overline{X}$
les images inverses respectives des deux sous-schémas $Y \subset W$ par $\overline{X} \to \overline{Y}$.
On obtient le diagramme suivant :
$$
\xymatrix{
X \ar[d]_f \ar[r] & U \ar[r] \ar[d] \ar@{}[dr]|A &
V \ar[d] \ar[r] \ar@{}[rd]|B & \overline{X} \ar[d] \\
Y \ar[d]_g \ar[r] & Y \ar[r] \ar[d] & W \ar[r] \ar[d] \ar@{}[rd]|C &
\overline{Y} \ar[d] \\
Z \ar[d]_h \ar[r] & Z \ar[d] \ar[r] & Z \ar[d] \ar[r] & \overline{Z} \ar[d] \\
T \ar[r] & T \ar[r] & T \ar[r] & T
}
$$
Sans multiplier les notations, mais en raisonnant exactement
comme dans la démonstration du lemme \ref{more-etale-lemma-shriek-composition},
on voit que les flèches du premier diagramme affiché utilisent les
flèches du lemme \ref{more-etale-lemma-shriek-proper-and-open} pour les rectangles
$A + B$, $B + C$, $A$ et $C$, comme l'indique le diagramme de
l'énoncé du lemme. D'après les
lemmes \ref{more-etale-lemma-shriek-proper-and-open-compose} et
\ref{more-etale-lemma-shriek-proper-and-open-compose-horizontal}
on a $\gamma_{A + B} = \gamma_B \circ \gamma_A$ et
$\gamma_{B + C} = \gamma_B \circ \gamma_C$ ; on en conclut
que l'égalité voulue est vraie pourvu que
$\gamma_A \circ \gamma_C = \gamma_C \circ \gamma_A$.
C'est bien le cas, car les deux carrés $A$ et $C$ ne
se rencontrent qu'en un point (comme dans le dernier argument de la
remarque \ref{more-etale-remark-going-around}).
\end{proof}

\begin{lemma}
\label{more-etale-lemma-base-change-shriek}
Considérons un carré cartésien
$$
\xymatrix{
X' \ar[r]_{g'} \ar[d]_{f'} & X \ar[d]^f \\
Y' \ar[r]^g & Y
}
$$
de schémas quasi-compacts et quasi-séparés, où $f$ est séparé et
de type fini. Il existe alors un isomorphisme canonique
$$
g^{-1} \circ Rf_! \to Rf'_! \circ (g')^{-1}
$$
De plus, ces isomorphismes sont compatibles avec ceux
du lemme \ref{more-etale-lemma-shriek-composition}.
\end{lemma}

\begin{proof}
Choisissons une compactification $j : X \to \overline{X}$ sur $Y$
et notons $\overline{f} : \overline{X} \to Y$ le morphisme structural.
Notons $j' : X' \to \overline{X}'$ et $\overline{f}' : \overline{X}' \to Y'$
les morphismes déduits de $j$ et de $\overline{f}$ par changement de base.
Puisque $Rf_! = R\overline{f}_* \circ j_!$ et $Rf'_! =
R\overline{f}'_* \circ j'_!$, l'isomorphisme se construit au moyen de
$$
g^{-1} \circ R\overline{f}_* \circ j_! \to
R\overline{f}'_* \circ (\overline{g}')^{-1} \circ j_! \to
R\overline{f}'_* \circ j'_! \circ (g')^{-1} 
$$
où la première flèche est l'isomorphisme fourni par le
théorème de changement de base propre (lemme de Cohomologie étale
\ref{etale-cohomology-lemma-proper-base-change} dans le cas borné inférieurement
à cohomologie de torsion, et le lemme de Cohomologie étale
\ref{etale-cohomology-lemma-proper-base-change-mod-n} dans le
cas où $\Lambda$ est un anneau de torsion), tandis que la deuxième est l'isomorphisme du
lemme \ref{more-etale-lemma-base-change-f-shriek-separated}.

\medskip\noindent
Pour achever la démonstration, il faut établir deux points : d'abord que
l'isomorphisme de foncteurs ainsi obtenu ne dépend pas du
choix de la compactification ; ensuite que, si l'on empile verticalement
deux diagrammes de changement de base comme dans le lemme,
les isomorphismes de changement de base sont compatibles avec ceux
du lemme \ref{more-etale-lemma-shriek-composition}.
Un argument direct, que nous omettons, montre que ces deux points découlent
de la compatibilité au changement de base des isomorphismes
\begin{enumerate}
\item $Rg_* \circ Rf_* = R(g \circ f)_*$ pour $f : X \to Y$ et $g : Y \to Z$
propres,
\item $g_! \circ f_! = (g \circ f)_!$ pour $f : X \to Y$ et $g : Y \to Z$
séparés et quasi-finis, et
\item $g_! \circ Rf'_* = Rf_* \circ g'_!$ pour $f : X \to Y$ et
$f' : X' \to Y'$ propres, et $g : Y' \to Y$ et $g' : X' \to X$
séparés et quasi-finis, avec $f \circ g' = g \circ f'$,
\end{enumerate}
sont compatibles au changement de base. Cela vaut pour (1) d'après la
remarque \ref{sites-cohomology-remark-compose-base-change} de Cohomologie sur les sites,
pour (2) d'après la remarque \ref{more-etale-remark-f-shriek-base-change-composition}, et pour
(3) d'après le lemme \ref{more-etale-lemma-shriek-proper-and-open-base-change}.
\end{proof}

\begin{remark}
\label{more-etale-remark-shriek-to-star}
Soit $f : X \to Y$ un morphisme de schémas séparé et de type fini,
où $Y$ est quasi-compact et quasi-séparé. Nous construirons plus loin
un morphisme
$$
Rf_!K \longrightarrow Rf_*K
$$
fonctoriel en $K$ dans $D^+_{tors}(X_\etale, \Lambda)$, ou dans
$D(X_\etale, \Lambda)$ si $\Lambda$ est un anneau de torsion. Dans les deux cas, cette transformation
de foncteurs est compatible avec
\begin{enumerate}
\item l'isomorphisme $Rg_! \circ Rf_! \to R(g \circ f)_!$ du
lemme \ref{more-etale-lemma-shriek-composition} et l'isomorphisme
$Rg_* \circ Rf_* \to R(g \circ f)_*$ du
lemme de Cohomologie sur les sites
\ref{sites-cohomology-lemma-derived-pushforward-composition} ;
\item l'isomorphisme $g^{-1} \circ Rf_! \to Rf'_! \circ (g')^{-1}$
du lemme \ref{more-etale-lemma-base-change-shriek}
et la flèche de changement de base de la
remarque \ref{sites-cohomology-remark-base-change} de Cohomologie sur les sites.
\end{enumerate}
En effet, choisissons une compactification $j : X \to \overline{X}$ sur $Y$
et notons $\overline{f} : \overline{X} \to Y$ le morphisme structural.
Puisque $Rf_! = R\overline{f}_* \circ j_!$ et
$Rf_* = R\overline{f}_* \circ Rj_*$, il suffit de construire une
transformation de foncteurs $j_! \to Rj_*$. Nous utilisons pour cela la
transformation canonique $j_! \to j_*$ du
lemme de Cohomologie étale
\ref{etale-cohomology-lemma-shriek-into-star-separated-etale}.
Nous omettons la démonstration de l'indépendance de la transformation obtenue
par rapport au choix de la compactification, ainsi que celle des
compatibilités (1) et (2).
\end{remark}













\section{Propriétés de l'image directe dérivée à support propre}
\label{more-etale-section-derived-lower-shriek-properties}

\noindent
Voici quelques propriétés de l'image directe dérivée à support propre.

\begin{lemma}
\label{more-etale-lemma-derived-lower-shriek-commute-direct-sums}
Soit $f : X \to Y$ un morphisme séparé de type fini entre schémas quasi-compacts
et quasi-séparés. Soit $\Lambda$ un anneau.
\begin{enumerate}
\item Considérons une famille d'objets $K_i \in D^+_{tors}(X_\etale, \Lambda)$, $i \in I$.
Supposons donné $a \in \mathbf{Z}$ tel que
$H^n(K_i) = 0$ pour $n < a$ et $i \in I$. Alors $Rf_!(\bigoplus_i K_i) =
\bigoplus_i Rf_!K_i$.
\item Si $\Lambda$ est un anneau de torsion, le foncteur
$Rf_! : D(X_\etale, \Lambda) \to D(Y_\etale, \Lambda)$
commute aux sommes directes.
\end{enumerate}
\end{lemma}

\begin{proof}
Par construction, il suffit de le démontrer lorsque $f$ est une immersion ouverte
et lorsque $f$ est un morphisme propre. Pour toute immersion ouverte $j : U \to X$
de schémas, le foncteur $j_! : D(U_\etale) \to D(X_\etale)$ est
adjoint à gauche de l'image inverse $j^{-1} : D(X_\etale) \to D(U_\etale)$
et commute donc aux sommes directes ; voir le lemme de Cohomologie sur les sites
\ref{sites-cohomology-lemma-adjoint-lower-shriek-restrict}.
Dans le cas propre, on a $Rf_! = Rf_*$ et le résultat découle du
lemme de Cohomologie étale
\ref{etale-cohomology-lemma-direct-sum-bounded-below-pushforward}
dans le cas borné inférieurement, et du
lemme \ref{etale-cohomology-lemma-proper-mod-n-direct-sums} de Cohomologie étale
lorsque $\Lambda$ est un anneau de torsion.
\end{proof}

\begin{lemma}
\label{more-etale-lemma-derived-lower-shriek-bounded}
Soit $f : X \to Y$ un morphisme séparé de type fini entre schémas quasi-compacts
et quasi-séparés. Soit $\Lambda$ un anneau. Les foncteurs $Rf_!$
construits à la section \ref{more-etale-section-derived-lower-shriek-compactification}
sont bornés au sens suivant : il existe un entier $N$ tel que,
pour $E \in D^+_{tors}(X_\etale, \Lambda)$ ou $E \in D(X_\etale, \Lambda)$
si $\Lambda$ est un anneau de torsion, on ait :
\begin{enumerate}
\item $H^i(Rf_!(\tau_{\leq a}E) \to H^i(Rf_!(E))$ est un isomorphisme
pour $i \leq a$ ;
\item $H^i(Rf_!(E)) \to H^i(Rf_!(\tau_{\geq b - N}E))$ est un isomorphisme
pour $i \geq b$ ;
\item si $H^i(E) = 0$ pour $i \not \in [a, b]$ avec
$-\infty \leq a \leq b \leq \infty$, alors $H^i(Rf_!(E)) = 0$
pour $i \not \in [a, b + N]$.
\end{enumerate}
\end{lemma}

\begin{proof}
Supposons que $\Lambda$ soit un anneau de torsion et considérons le foncteur
$Rf_! : D(X_\etale, \Lambda) \to D(Y_\etale, \Lambda)$.
Par construction, il suffit de le démontrer lorsque $f$ est une immersion ouverte
et lorsque $f$ est un morphisme propre. Pour toute immersion ouverte $j : U \to X$
de schémas, le foncteur $j_! : D(U_\etale) \to D(X_\etale)$ est exact ;
l'énoncé vaut donc avec $N = 0$ dans ce cas.
Si $f$ est propre, alors $Rf_! = Rf_*$, c'est-à-dire qu'il s'agit d'un foncteur
dérivé à droite. La borne à gauche résulte donc du
lemme \ref{derived-lemma-negative-vanishing} de Catégories dérivées.
De plus, dans ce cas, $f_* : \textit{Mod}(X_\etale, \Lambda)
\to \textit{Mod}(Y_\etale, \Lambda)$ est de dimension cohomologique finie d'après le
lemme \ref{morphisms-lemma-morphism-finite-type-bounded-dimension} de Morphismes
et le lemme de Cohomologie étale
\ref{etale-cohomology-lemma-cohomological-dimension-proper}.
On conclut donc par le
lemme \ref{derived-lemma-unbounded-right-derived} de Catégories dérivées.

\medskip\noindent
Supposons maintenant $\Lambda$ quelconque et considérons le foncteur
$Rf_! : D^+_{tors}(X_\etale, \Lambda) \to D^+_{tors}(Y_\etale, \Lambda)$.
On se ramène encore immédiatement au cas où $f$ est propre et
$Rf_! = Rf_*$. Le point (1) est encore immédiat. Pour démontrer le point (3),
on peut raisonner par récurrence sur $b - a$, en utilisant les triangles distingués
de troncation et le lemme de Cohomologie étale
\ref{etale-cohomology-lemma-cohomological-dimension-proper}.
Le point (2) résulte du point (3). Nous omettons les détails.
\end{proof}

\begin{lemma}
\label{more-etale-lemma-Rf-shriek-for-quasi-finite}
Soit $f : X \to Y$ un morphisme quasi-fini séparé entre schémas quasi-compacts et
quasi-séparés. Alors les foncteurs $Rf_!$ construits à la
section \ref{more-etale-section-derived-lower-shriek-compactification}
coïncident avec la restriction du foncteur
$f_! : D(X_\etale, \Lambda) \to D(Y_\etale, \Lambda)$
construit à la section \ref{more-etale-section-derived-duality-locally-quasi-finite}
à leurs domaines de définition communs.
\end{lemma}

\begin{proof}
Par le théorème principal de Zariski (lemme de Compléments sur les morphismes
\ref{more-morphisms-lemma-quasi-finite-separated-pass-through-finite})
on peut trouver une immersion ouverte $j : X \to \overline{X}$ et un morphisme fini
$\overline{f} : \overline{X} \to Y$ tels que $f = \overline{f} \circ j$.
Par construction, on a $Rf_! = R\overline{f}_* \circ j_!$.
Comme $\overline{f}$ est fini, on a $R\overline{f}_* = \overline{f}_*$
d'après la proposition de Cohomologie étale
\ref{etale-cohomology-proposition-finite-higher-direct-image-zero}.
Le lemme en résulte, car $\overline{f}_* \circ j_! = f_!$,
par exemple d'après le lemme \ref{more-etale-lemma-compactify-f-shriek-separated}.
\end{proof}

\begin{lemma}
\label{more-etale-lemma-relative-mayer-vietoris}
Soit $f : X \to Y$ un morphisme séparé de type fini entre schémas quasi-compacts
et quasi-séparés. Soient $U$ et $V$ des ouverts quasi-compacts de $X$
tels que $X = U \cup V$. Notons $a : U \to Y$, $b : V \to Y$ et
$c : U \cap V \to Y$ les restrictions de $f$. Soit $\Lambda$ un anneau.
Pour $K$ dans $D^+_{tors}(X_\etale, \Lambda)$, ou $K \in D(X_\etale, \Lambda)$
si $\Lambda$ est un anneau de torsion, on a un triangle distingué
$$
Rc_!(K|_{U \cap V}) \to
Ra_!(K|_U) \oplus Rb_!(K|_V) \to
Rf_!K \to 
Rc_!(K|_{U \cap V})[1]
$$
dans $D(Y_\etale, \Lambda)$.
\end{lemma}

\begin{proof}
Cela résulte du lemme \ref{more-etale-lemma-mayer-vietoris-shriek},
de l'égalité $Rf_! \circ Rj_{U!} = Ra_!$ donnée par le
lemme \ref{more-etale-lemma-shriek-composition}, et de l'égalité
$Rj_{U!} = j_{U!}$ donnée par le lemme \ref{more-etale-lemma-Rf-shriek-for-quasi-finite}.
\end{proof}

\begin{lemma}
\label{more-etale-lemma-relative-triangle-associated-to-open}
Soit $f : X \to Y$ un morphisme séparé de type fini entre schémas quasi-compacts
et quasi-séparés. Soit $U$ un ouvert quasi-compact de $X$ dont
le complémentaire est $Z \subset X$. Notons $g : U \to Y$ et $h : Z \to Y$
les restrictions de $f$. Soit $\Lambda$ un anneau.
Pour $K$ dans $D^+_{tors}(X_\etale, \Lambda)$, ou $K \in D(X_\etale, \Lambda)$
si $\Lambda$ est un anneau de torsion, on a un triangle distingué
$$
Rg_!(K|_U) \to
Rf_!K \to
Rh_!(K|_Z) \to
Rg_!(K|_U)[1]
$$
dans $D(Y_\etale, \Lambda)$.
\end{lemma}

\begin{proof}
Cela résulte du
lemme \ref{more-etale-lemma-triangle-associated-to-open},
des égalités $Rf_! \circ Rj_! = Rg_!$ et $Rf_! \circ Ri_!$ données par le
lemme \ref{more-etale-lemma-shriek-composition}, ainsi que des égalités
$Rj_! = j_!$ et $Ri_! = i_! = i_*$ données par le
lemme \ref{more-etale-lemma-Rf-shriek-for-quasi-finite}.
\end{proof}

\begin{lemma}
\label{more-etale-lemma-shriek-and-thickening}
Soit $f' : X' \to Y$ un morphisme séparé de type fini entre schémas quasi-compacts
et quasi-séparés. Soit $i : X \to X'$ un épaississement et
notons $f = f' \circ i$. Soit $\Lambda$ un anneau.
Pour $K'$ dans $D^+_{tors}(X'_\etale, \Lambda)$, ou $K' \in D(X'_\etale, \Lambda)$
si $\Lambda$ est un anneau de torsion, on a $Rf_!i^{-1}K' = Rf'_!K'$.
\end{lemma}

\begin{proof}
Cela résulte de ce que $i^{-1}$ et $i_* = i_!$ sont des équivalences de catégories
quasi-inverses l'une de l'autre par invariance topologique du petit topos étale
(théorème de Cohomologie étale
\ref{etale-cohomology-theorem-topological-invariance}),
ainsi que du
lemme \ref{more-etale-lemma-shriek-composition}.
\end{proof}

\begin{lemma}
\label{more-etale-lemma-projection-formula-Rf-lower-shriek}
Soit $f : X \to Y$ un morphisme séparé de type fini entre schémas quasi-compacts
et quasi-séparés. Soit $\Lambda$ un anneau de torsion. Soient
$E \in D(X_\etale, \Lambda)$ et $K \in D(Y_\etale, \Lambda)$. Alors
$$
Rf_!E \otimes_\Lambda^\mathbf{L} K =
Rf_!(E \otimes_\Lambda^\mathbf{L} f^{-1}K)
$$
dans $D(Y_\etale, \Lambda)$.
\end{lemma}

\begin{proof}
Choisissons $j : X \to \overline{X}$ et $\overline{f} : \overline{X} \to Y$
comme dans la construction de $Rf_!$. On a
$j_!E \otimes_\Lambda^\mathbf{L} \overline{f}^{-1}K =
j_!(E \otimes_\Lambda^\mathbf{L} f^{-1}K)$ d'après
le lemme \ref{sites-cohomology-lemma-j-shriek-and-tensor} de Cohomologie sur les sites.
On obtient alors le résultat en appliquant le
lemme de Cohomologie étale
\ref{etale-cohomology-lemma-projection-formula-proper-mod-n}
et en utilisant $f^{-1} = j^{-1} \circ \overline{f}^{-1}$ et
$Rf_! = R\overline{f}_*j_!$.
\end{proof}

\begin{remark}
\label{more-etale-remark-Rf-lower-shriek-change-of-rings}
Soit $\Lambda_1 \to \Lambda_2$ un homomorphisme d'anneaux de torsion.
Soit $f : X \to Y$ un morphisme séparé de type fini entre schémas quasi-compacts
et quasi-séparés. Le diagramme
$$
\xymatrix{
D(X_\etale, \Lambda_2) \ar[r]_{res} \ar[d]_{Rf_!} &
D(X_\etale, \Lambda_1) \ar[d]^{Rf_!} \\
D(Y_\etale, \Lambda_2) \ar[r]^{res} &
D(Y_\etale, \Lambda_1)
}
$$
est commutatif, où $res$ est le foncteur de restriction des scalaires qui considère un
$\Lambda_2$-module comme un $\Lambda_1$-module au moyen du morphisme d'anneaux donné.
Écrivons $Rf_! = R\overline{f}_* \circ j_!$ pour une factorisation
$f = \overline{f} \circ j$ comme dans la
section \ref{more-etale-section-derived-lower-shriek-compactification}. On voit que
le résultat se vérifie directement pour $j_!$ et résulte, pour $R\overline{f}_*$,
du lemme de Cohomologie sur les sites
\ref{sites-cohomology-lemma-modules-abelian-unbounded}.
D'autre part, le diagramme
$$
\xymatrix{
D(X_\etale, \Lambda_1)
\ar[r]_{- \otimes_{\Lambda_1}^\mathbf{L} \Lambda_2} \ar[d]_{Rf_!} &
D(X_\etale, \Lambda_2) \ar[d]^{Rf_!} \\
D(Y_\etale, \Lambda_1)
\ar[r]^{- \otimes_{\Lambda_1}^\mathbf{L} \Lambda_2} &
D(Y_\etale, \Lambda_2)
}
$$
est également commutatif d'après le
lemme \ref{more-etale-lemma-projection-formula-Rf-lower-shriek}.
\end{remark}

\begin{remark}
\label{more-etale-remark-projection-formula-for-internal-hom}
Soit $f : X \to Y$ un morphisme séparé de type fini entre schémas quasi-compacts
et quasi-séparés. Soit $\Lambda$ un anneau de torsion.
Soient $K$ et $L$ des objets de $D(X_\etale, \Lambda)$. Nous affirmons
qu'il existe un morphisme canonique
$$
\alpha :
Rf_*R\SheafHom_\Lambda(K, L)
\longrightarrow
R\SheafHom_\Lambda(Rf_!K, Rf_!L)
$$
fonctoriel en $K$ et $L$. Choisissons en effet $j : X \to \overline{X}$ et
$\overline{f} : \overline{X} \to Y$ comme dans la construction de $Rf_!$.
Définissons d'abord un morphisme
$$
\beta :
Rj_*R\SheafHom_\Lambda(K, L)
\longrightarrow
R\SheafHom_\Lambda(j_!K, j_!L)
$$
Par la construction du Hom interne dans la catégorie dérivée, cela revient
à définir un morphisme
$$
\beta' :
Rj_*R\SheafHom_\Lambda(K, L) \otimes_\Lambda^\mathbf{L} j_!K
\longrightarrow
j_!L
$$
Voir Cohomologie sur les sites, section \ref{sites-cohomology-section-internal-hom}.
La source de $\beta'$ est égale à
$$
j_!\left(R\SheafHom_\Lambda(K, L) \otimes_\Lambda^\mathbf{L} K\right)
$$
par le lemme \ref{sites-cohomology-lemma-j-shriek-and-tensor} de Cohomologie sur les sites.
On peut donc poser $\beta' = j_!\beta''$, où
$\beta'' : R\SheafHom_\Lambda(K, L) \otimes_\Lambda^\mathbf{L} K \to L$
correspond à l'identité de $R\SheafHom_\Lambda(K, L)$
par la propriété universelle du Hom interne mentionnée ci-dessus. D'après la
remarque de Cohomologie sur les sites
\ref{sites-cohomology-remark-projection-formula-for-internal-hom}
on a un morphisme canonique
$$
\gamma :
R\overline{f}_*R\SheafHom_\Lambda(j_!K, j_!L)
\longrightarrow
R\SheafHom_\Lambda(R\overline{f}_*j_!K, R\overline{f}_*j_!L)
$$
Comme $Rf_! = R\overline{f}_*j_!$ et $Rf_* = R\overline{f}_* Rj_*$
(par Leray), on obtient le morphisme recherché
$\alpha = \gamma \circ R\overline{f}_*\beta$.
\end{remark}









\section{Image inverse extraordinaire dérivée}
\label{more-etale-section-derived-upper-shriek}

\noindent
Nous obtenons $Rf^!$ au moyen d'un théorème de représentabilité de Brown.

\begin{lemma}
\label{more-etale-lemma-upper-shriek-derived}
Soit $f : X \to Y$ un morphisme séparé de type fini entre schémas quasi-compacts
et quasi-séparés. Soit $\Lambda$ un anneau de torsion.
Le foncteur
$Rf_! : D(X_\etale, \Lambda) \to D(Y_\etale, \Lambda)$
admet un adjoint à droite
$Rf^! : D(Y_\etale, \Lambda) \to D(X_\etale, \Lambda)$.
\end{lemma}

\begin{proof}
Cela résulte de la
proposition \ref{injectives-proposition-brown} d'Injectifs
et du lemme \ref{more-etale-lemma-derived-lower-shriek-commute-direct-sums} ci-dessus.
\end{proof}

\begin{lemma}
\label{more-etale-lemma-Rf-upper-shriek-for-quasi-finite}
Soit $f : X \to Y$ un morphisme quasi-fini séparé entre schémas quasi-compacts
et quasi-séparés. Soit $\Lambda$ un anneau de torsion.
Le foncteur $Rf^! : D(Y_\etale, \Lambda) \to D(X_\etale, \Lambda)$
du lemme \ref{more-etale-lemma-upper-shriek-derived} s'identifie au foncteur
$Rf^!$ du lemme \ref{more-etale-lemma-lqf-shriek-derived}.
\end{lemma}

\begin{proof}
Cela résulte de l'unicité des adjoints, puisque $Rf_! = f_!$ d'après le
lemme \ref{more-etale-lemma-Rf-shriek-for-quasi-finite}.
\end{proof}

\begin{lemma}
\label{more-etale-lemma-Rf-upper-shriek-for-etale}
Soit $j : U \to X$ un morphisme étale séparé entre schémas quasi-compacts
et quasi-séparés. Soit $\Lambda$ un anneau de torsion.
Le foncteur $Rj^! : D(X_\etale, \Lambda) \to D(U_\etale, \Lambda)$
s'identifie à $j^{-1}$.
\end{lemma}

\begin{proof}
En effet, $Rj^!$ et $j^{-1}$ sont tous deux adjoints à droite du foncteur
$Rj_! = j_!$. Voir par exemple les lemmes
\ref{more-etale-lemma-Rf-upper-shriek-for-quasi-finite} et
\ref{more-etale-lemma-etale-upper-shriek}.
\end{proof}

\begin{lemma}
\label{more-etale-lemma-twisted-inverse-image-bounded-below}
Soit $f : X \to Y$ un morphisme séparé de type fini entre
schémas quasi-compacts et quasi-séparés. Soit $\Lambda$ un anneau de
torsion. Le foncteur $Rf^!$ envoie $D^+(Y_\etale, \Lambda)$ dans
$D^+(X_\etale, \Lambda)$. Plus précisément, il existe un entier
$N \geq 0$ tel que, si $K \in D(Y_\etale, \Lambda)$ vérifie $H^i(K) = 0$
pour $i < a$, alors $H^i(Rf^!K) = 0$ pour $i < a - N$.
\end{lemma}

\begin{proof}
Soit $N$ l'entier fourni par le lemme \ref{more-etale-lemma-derived-lower-shriek-bounded}.
Par construction, pour $K \in D(Y_\etale, \Lambda)$ et
$L \in \in D(X_\etale, \Lambda)$, on a
$\Hom_X(L, Rf^!K) = \Hom_Y(Rf_!L, K)$. Supposons $H^i(K) = 0$
pour $i < a$. Prenons alors $L = \tau_{\leq a - N - 1}Rf^!K$. D'après le
lemme \ref{more-etale-lemma-derived-lower-shriek-bounded},
le complexe $Rf_!L$ a ses faisceaux de cohomologie nuls aux
degrés $\leq a - 1$. Ainsi, $\Hom_Y(Rf_!L, K) = 0$ d'après le
lemme \ref{derived-lemma-negative-exts} de Catégories dérivées.
Le morphisme canonique $\tau_{\leq a - N - 1}Rf^!K \to Rf^!K$
est donc nul, ce qui entraîne $H^i(Rf^!K) = 0$ pour $i \leq a - N - 1$.
\end{proof}

\noindent
Soit $f : X \to Y$ un morphisme séparé de type fini entre schémas quasi-séparés
et quasi-compacts. Soit $\Lambda$ un anneau de torsion.
Pour tous $K \in D(Y_\etale, \Lambda)$ et $L \in D(X_\etale, \Lambda)$,
on obtient un morphisme canonique
\begin{equation}
\label{more-etale-equation-sheafy-trace}
Rf_*R\SheafHom_\Lambda(L, Rf^!K) \longrightarrow R\SheafHom_\Lambda(Rf_!L, K)
\end{equation}
Ce morphisme est construit comme la composée
$$
Rf_*R\SheafHom_\Lambda(L, Rf^!K) \to
R\SheafHom_\Lambda(Rf_!L, Rf_!Rf^!K) \to
R\SheafHom_\Lambda(Rf_!L, K)
$$
où la première flèche est celle de la
remarque \ref{more-etale-remark-projection-formula-for-internal-hom}
et la seconde est le morphisme d'adjonction $Rf_!Rf^!K \to K$.

\begin{lemma}
\label{more-etale-lemma-iso-on-RSheafHom}
Soit $f : X \to Y$ un morphisme séparé de type fini entre schémas quasi-compacts et
quasi-séparés. Soit $\Lambda$ un anneau de torsion.
Pour tous $K \in D(Y_\etale, \Lambda)$ et $L \in D(X_\etale, \Lambda)$,
le morphisme (\ref{more-etale-equation-sheafy-trace})
$$
Rf_*R\SheafHom_\Lambda(L, Rf^!K) \longrightarrow R\SheafHom_\Lambda(Rf_!L, K)
$$
est un isomorphisme.
\end{lemma}

\begin{proof}
Pour démontrer le lemme, il suffit de montrer que, pour tout $M \in D(Y_\etale, \Lambda)$,
le morphisme (\ref{more-etale-equation-sheafy-trace}) induit une bijection
$$
\Hom_Y(M, Rf_*R\SheafHom_\Lambda(L, Rf^!K))
\longrightarrow
\Hom_Y(M, R\SheafHom_\Lambda(Rf_!L, K))
$$
Pour le voir, utilisons la suite d'égalités suivante :
\begin{align*}
\Hom_Y(M, Rf_*R\SheafHom_\Lambda(L, Rf^!K))
& =
\Hom_X(f^{-1}M, R\SheafHom_\Lambda(L, Rf^!K)) \\
& =
\Hom_X(f^{-1}M \otimes_\Lambda^\mathbf{L} L, Rf^!K) \\
& =
\Hom_Y(Rf_!(f^{-1}M \otimes_\Lambda^\mathbf{L} L), K) \\
& =
\Hom_Y(M \otimes_\Lambda^\mathbf{L} Rf_!L, K) \\
& =
\Hom_Y(M, R\SheafHom_\Lambda(Rf_!L, K))
\end{align*}
La première égalité résulte du
lemme \ref{sites-cohomology-lemma-adjoint} de Cohomologie sur les sites.
La deuxième résulte du
lemme \ref{sites-cohomology-lemma-internal-hom} de Cohomologie sur les sites.
La troisième résulte de la construction de $Rf^!$.
La quatrième résulte du lemme \ref{more-etale-lemma-projection-formula-Rf-lower-shriek}
(c'est l'étape essentielle).
La cinquième résulte du
lemme \ref{sites-cohomology-lemma-internal-hom} de Cohomologie sur les sites.
\end{proof}

\begin{lemma}
\label{more-etale-lemma-iso-global-hom}
Soit $f : X \to Y$ un morphisme séparé de type fini entre schémas
quasi-séparés et quasi-compacts. Soit $\Lambda$ un anneau de torsion.
Pour tous $K \in D(Y_\etale, \Lambda)$ et $L \in D(X_\etale, \Lambda)$, le morphisme
(\ref{more-etale-equation-sheafy-trace}) induit un isomorphisme
$$
R\Hom_X(L, Rf^!K) \longrightarrow R\Hom_Y(Rf_!L, K)
$$
entre les Hom dérivés globaux.
\end{lemma}

\begin{proof}
Par la construction donnée dans la
section \ref{sites-cohomology-section-global-RHom} de Cohomologie sur les sites,
on a
$$
R\Hom_X(L, Rf^!K) =
R\Gamma(X, R\SheafHom_\Lambda(L, Rf^!K)) =
R\Gamma(Y, Rf_*R\SheafHom_\Lambda(L, Rf^!K))
$$
(la seconde égalité résultant de Leray) et
$$
R\Hom_Y(Rf_!L, K) = R\Gamma(Y, R\SheafHom_\Lambda(Rf_!L, K))
$$
Le lemme résulte donc du lemme \ref{more-etale-lemma-iso-on-RSheafHom}.
\end{proof}

\begin{lemma}
\label{more-etale-lemma-cartesian-square-Rf-upper-shriek}
Considérons un carré cartésien
$$
\xymatrix{
X' \ar[r]_{g'} \ar[d]_{f'} & X \ar[d]^f \\
Y' \ar[r]^g & Y
}
$$
de schémas quasi-compacts et quasi-séparés, où $f$ est séparé et
de type fini. Alors $Rf^! \circ Rg_* = Rg'_* \circ R(f')^!$.
\end{lemma}

\begin{proof}
Par unicité des foncteurs adjoints, cela résulte du changement de base
pour l'image directe dérivée à support propre : on a
$g^{-1} \circ Rf_! = Rf'_! \circ (g')^{-1}$ d'après
le lemme \ref{more-etale-lemma-base-change-shriek}.
\end{proof}

\begin{remark}
\label{more-etale-remark-Rf-upper-shriek-change-of-rings}
Soit $\Lambda_1 \to \Lambda_2$ un homomorphisme d'anneaux de torsion.
Soit $f : X \to Y$ un morphisme séparé de type fini entre schémas quasi-compacts
et quasi-séparés. Le diagramme
$$
\xymatrix{
D(X_\etale, \Lambda_2) \ar[r]_{res} &
D(X_\etale, \Lambda_1) \\
D(Y_\etale, \Lambda_2) \ar[r]^{res} \ar[u]^{Rf^!} &
D(Y_\etale, \Lambda_1) \ar[u]_{Rf^!}
}
$$
est commutatif, où $res$ est le foncteur de restriction des scalaires qui considère un
$\Lambda_2$-module comme un $\Lambda_1$-module au moyen du morphisme d'anneaux donné.
Cela résulte de l'unicité des adjoints, du second diagramme commutatif
de la remarque \ref{more-etale-remark-Rf-lower-shriek-change-of-rings}, et du fait que
$$
\Hom_{\Lambda_2}(K_1 \otimes_{\Lambda_1}^\mathbf{L} \Lambda_2, K_2) =
\Hom_{\Lambda_1}(K_1, res(K_2))
$$
Cette égalité, tant pour les objets sur $X_\etale$ que pour ceux sur
$Y_\etale$, est un cas très particulier du
lemme \ref{sites-cohomology-lemma-adjoint} de Cohomologie sur les sites.
\end{remark}











\section{Cohomologie à support propre}
\label{more-etale-section-compactly-supported-cohomology}

\noindent
Soit $k$ un corps. Soit $\Lambda$ un anneau. Soit $X$ un schéma séparé
de type fini sur $k$, de morphisme structural $f : X \to \Spec(k)$.
À la section \ref{more-etale-section-derived-lower-shriek-compactification},
nous avons défini le foncteur
$Rf_! : D^+_{tors}(X_\etale, \Lambda) \to D^+_{tors}(\Spec(k), \Lambda)$
et le foncteur $Rf_! : D(X_\etale, \Lambda) \to D(\Spec(k), \Lambda)$
si $\Lambda$ est un anneau de torsion. En composant avec le foncteur des sections globales
sur $\Spec(k)$, on obtient ce que nous appellerons la cohomologie à support propre.

\begin{definition}
\label{more-etale-definition-cohomology-compact-support}
Soit $X$ un schéma séparé de type fini sur un corps $k$.
Soit $\Lambda$ un anneau. Soit $K$ un objet de
$D^+_{tors}(X_\etale, \Lambda)$
ou de $D(X_\etale, \Lambda)$ lorsque $\Lambda$ est un anneau de torsion.
La {\it cohomologie de $K$ à support propre}, ou la
{\it cohomologie à support propre de $K$}, est
$$
R\Gamma_c(X, K) = R\Gamma(\Spec(k), Rf_!K)
$$
où $f : X \to \Spec(k)$ est le morphisme structural. Nous
écrirons $H^i_c(X, K) = H^i(R\Gamma_c(X, K))$.
\end{definition}

\noindent
Nous vérifierons que cette définition est compatible avec la
définition \ref{more-etale-definition-compact-support} grâce au
lemme \ref{more-etale-lemma-compact-support-h0}.
L'intérêt de cette définition réside dans le résultat suivant.

\begin{lemma}
\label{more-etale-lemma-stalk-R-f-shriek}
Soit $f : X \to Y$ un morphisme séparé de type fini entre schémas,
où $Y$ est quasi-compact et quasi-séparé. Soit $K$ un objet de
$D^+_{tors}(X_\etale, \Lambda)$ ou de $D(X_\etale, \Lambda)$ lorsque
$\Lambda$ est un anneau de torsion. Il existe alors un isomorphisme canonique
$$
(Rf_!K)_{\overline{y}}
\longrightarrow
R\Gamma_c(X_{\overline{y}}, K|_{X_{\overline{y}}})
$$
dans $D(\Lambda)$ pour tout point géométrique $\overline{y} : \Spec(k) \to Y$.
\end{lemma}

\begin{proof}
C'est une conséquence immédiate du lemme \ref{more-etale-lemma-base-change-shriek} et des
définitions.
\end{proof}

\begin{lemma}
\label{more-etale-lemma-compact-support-h0}
Soit $X$ un schéma séparé de type fini sur un corps $k$.
Si $\mathcal{F}$ est un faisceau abélien de torsion, alors le groupe abélien
$H^0_c(X, \mathcal{F})$ défini dans la définition \ref{more-etale-definition-compact-support}
coïncide avec le groupe abélien $H^0_c(X, \mathcal{F})$ défini dans la
définition \ref{more-etale-definition-cohomology-compact-support}.
\end{lemma}

\begin{proof}
Choisissons une compactification $j : X \to \overline{X}$ sur $k$.
Dans les deux cas, le groupe est défini comme $H^0(\overline{X}, j_!\mathcal{F})$.
Pour la première version, cela résulte du
lemme \ref{more-etale-lemma-compactify-compact-support} ;
pour la seconde, cela résulte de la construction.
\end{proof}

\begin{lemma}
\label{more-etale-lemma-finiteness-curves}
Soit $k$ un corps algébriquement clos. Soit $X$ un schéma séparé
de type fini sur $k$, de dimension $\leq 1$.
Soit $\Lambda$ un anneau noethérien.
Soit $\mathcal{F}$ un faisceau constructible de $\Lambda$-modules
sur $X$, supposé de torsion. Alors $H^q_c(X, \mathcal{F})$ est un
$\Lambda$-module de type fini.
\end{lemma}

\begin{proof}
Cela résulte du théorème de Cohomologie étale
\ref{etale-cohomology-theorem-vanishing-affine-curves-coefficients}.
Choisissons en effet une compactification $j : X \to \overline{X}$.
Après avoir remplacé $\overline{X}$ par l'adhérence schématique
de $X$, on voit que l'on peut supposer $\dim(\overline{X}) \leq 1$.
Alors $H^q_c(X, \mathcal{F}) = H^q(\overline{X}, j_!\mathcal{F})$
et le théorème s'applique.
\end{proof}

\begin{remark}[Covariance de la cohomologie à support propre]
\label{more-etale-remark-covariance-compactly-supported}
Soit $k$ un corps. Soit $f : X \to Y$ un morphisme de schémas séparés
de type fini sur $k$. Si $X$, $Y$ et $f$ vérifient l'une
des conditions suivantes :
\begin{enumerate}
\item $f$ est étale ;
\item $f$ est plat et quasi-fini ;
\item $f$ est quasi-fini et $Y$ est géométriquement unibranche ;
\item $f$ est quasi-fini et il existe une pondération
$w : X \to \mathbf{Z}$ de $f$,
\end{enumerate}
alors la cohomologie à support propre est covariante vis-à-vis de $f$.
Plus précisément, soit $\Lambda$ un anneau. Soit $K$ un objet de
$D^+_{tors}(Y_\etale, \Lambda)$ ou de $D(Y_\etale, \Lambda)$ lorsque
$\Lambda$ est un anneau de torsion. Sous l'une des hypothèses (1) -- (4),
il existe un morphisme canonique
$$
\text{Tr}_{f, w, K} : f_!f^{-1}K \longrightarrow K
$$
Voir la section \ref{more-etale-section-weightings} pour l'existence du
morphisme trace, et les exemples \ref{more-etale-example-trace-for-flat-quasi-finite} et
\ref{more-etale-example-trace-for-quasi-finite-over-normal} pour les cas (2) et (3).
Si $p : X \to \Spec(k)$ et $q : Y \to \Spec(k)$ désignent
les morphismes structuraux, alors $Rq_! \circ f_! = Rp_!$
d'après le lemme \ref{more-etale-lemma-shriek-composition} et l'égalité
$Rf_! = f_!$ pour le morphisme quasi-fini séparé $f$, donnée par le
lemme \ref{more-etale-lemma-Rf-shriek-for-quasi-finite}. On peut donc considérer
le morphisme
\begin{align*}
R\Gamma_c(X, f^{-1}K)
& =
R\Gamma(\Spec(k), Rp_!f^{-1}K) \\
& =
R\Gamma(\Spec(k), Rq_!f_!f^{-1}K) \\
& \xrightarrow{Rq_!\text{Tr}_{f, w, K}}
R\Gamma(\Spec(k), Rq_!K) \\
& =
R\Gamma_c(Y, K)
\end{align*}
En particulier, si $\Lambda$ est un anneau de torsion, on obtient une flèche
$$
\text{Tr}_f : R\Gamma_c(X, \Lambda) \longrightarrow R\Gamma_c(Y, \Lambda)
$$
Ce morphisme vérifie de nombreuses propriétés supplémentaires ; il est notamment
compatible avec tout changement de corps de base.
\end{remark}










\section{Un résultat de constructibilité}
\label{more-etale-section-family-smooth-curves-cohomology}

\noindent
Nous « calculons » la cohomologie d'une famille projective lisse
de courbes à coefficients constants.

\begin{lemma}
\label{more-etale-lemma-frobenius-linear-on-vb}
Soit $p$ un nombre premier. Soit $S$ un schéma sur $\mathbf{F}_p$.
Soit $\mathcal{E}$ un module localement libre de type fini
sur $\mathcal{O}_S$, vu comme $\mathcal{O}_S$-module sur $S_\etale$.
Soit $F : \mathcal{E} \to \mathcal{E}$ un homomorphisme de faisceaux abéliens
sur $S_\etale$ tel que $F(a e) = a^pF(e)$ pour toutes sections locales $a$, $e$
respectivement de $\mathcal{O}_S$ et $\mathcal{E}$ sur $S_\etale$. Alors
$$
\Coker(F - 1 : \mathcal{E} \to \mathcal{E})
$$
est nul, et
$$
\Ker(F - 1 : \mathcal{E} \to \mathcal{E})
$$
est un faisceau abélien constructible sur $S_\etale$.
\end{lemma}

\noindent
Ce lemme généralise le lemme de Cohomologie étale
\ref{etale-cohomology-lemma-F-1}.

\begin{proof}
Nous pouvons supposer $S = \Spec(A)$, où $A$ est une $\mathbf{F}_p$-algèbre,
et que $\mathcal{E}$ est le module quasi-cohérent associé
au $A$-module libre $Ae_1 \oplus \ldots \oplus Ae_n$.
Écrivons $F(e_i) = \sum a_{ij} e_j$.

\medskip\noindent
Surjectivité de $F - 1$. Il suffit de montrer que tout élément
$\sum a_i e_i$, $a_i \in A$, appartient à l'image de $F - 1$
après remplacement de $A$ par une extension étale fidèlement plate.
Remarquons que
$$
F(\sum x_ie_i) - \sum x_i e_i = \sum x_i^p a_{ij} e_j - \sum x_i e_i
$$
Considérons la $A$-algèbre
$$
A' = A[x_1, \ldots, x_n]/(a_i + x_i - \sum\nolimits_j a_{ji} x_j^p)
$$
Un calcul montre que $\text{d}x_i$ est nul dans $\Omega_{A'/A}$,
donc que $\Omega_{A'/A} = 0$. Comme $A'$ est de type fini sur
$A$, il en résulte que $\Spec(A') \to \Spec(A)$ est non ramifié,
donc quasi-fini. Comme $A'$ est engendré par $n$ éléments
et défini par $n$ équations, on conclut que $A'$
est une intersection complète relative globale sur $A$.
Ainsi, $A'$ est plat sur $A$, et $A \to A'$
est étale (comme morphisme d'anneaux plat et non ramifié). Enfin,
le lecteur peut montrer que $A \to A'$ est fidèlement plat
en vérifiant directement que toutes les fibres géométriques de
$\Spec(A') \to \Spec(A)$ sont non vides ; cela résulte aussi du lemme de
Cohomologie étale
\ref{etale-cohomology-lemma-F-1}.
Enfin, l'élément $\sum x_i e_i \in A'e_1 \oplus \ldots \oplus A'e_n$
a pour image $\sum a_i e_i$ par $F - 1$.

\medskip\noindent
Constructibilité du noyau. Les calculs ci-dessus montrent que
$\Ker(F - 1)$ est représenté par le schéma
$$
\Spec(A[x_1, \ldots, x_n]/(x_i - \sum\nolimits_j a_{ji} x_j^p))
$$
sur $S = \Spec(A)$. Comme il s'agit d'un schéma affine et étale
sur $S$, le résultat découle du
lemme \ref{etale-cohomology-lemma-jshriek-constructible} de Cohomologie étale.
\end{proof}

\begin{lemma}
\label{more-etale-lemma-proper-smooth-family-curves}
Soit $f : X \to S$ un morphisme propre et lisse de schémas, à
fibres géométriquement connexes de dimension $1$. Soit $\ell$
un nombre premier. Alors $R^qf_*\underline{\mathbf{Z}/\ell\mathbf{Z}}$
est constructible.
\end{lemma}

\begin{proof}
Nous pouvons supposer $S$ affine, disons $S = \Spec(A)$. Si l'on écrit
$A = \bigcup A_i$ comme réunion de ses sous-$\mathbf{Z}$-algèbres de type fini,
on peut trouver un $i$ et un morphisme $f_i : X_i \to S_i = \Spec(A_i)$ de type
fini dont le changement de base à $S$ est $f : X \to S$ ; voir le
lemme \ref{limits-lemma-descend-finite-presentation} de Limites.
Quitte à augmenter $i$, on peut supposer $f_i : X_i \to S_i$
lisse, propre et de dimension relative $1$ ; voir les
lemmes de Limites
\ref{limits-lemma-eventually-proper},
\ref{limits-lemma-descend-smooth}, et
\ref{limits-lemma-descend-dimension-d}.
D'après le lemme \ref{more-morphisms-lemma-proper-flat-geom-red} de Compléments sur les morphismes,
on obtient un sous-schéma ouvert $U_i \subset S_i$ tel que les fibres
de $f_i : X_i \to S_i$ au-dessus de $U_i$ soient géométriquement connexes.
Alors $S \to S_i$ se factorise par $U_i$. On peut remplacer $X \to S$
par $f_i : f_i^{-1}(U_i) \to U_i$ pour se ramener au cas
traité dans le paragraphe suivant.

\medskip\noindent
Supposons $S$ noethérien. On peut écrire $S = U \cup Z$, où $U$ est le
sous-schéma ouvert défini par la non-annulation de $\ell$ et
$Z = V(\ell) \subset S$. Comme la formation de
$R^qf_*\underline{\mathbf{Z}/\ell\mathbf{Z}}$ commute à tout
changement de base (théorème de Cohomologie étale
\ref{etale-cohomology-theorem-proper-base-change}),
il suffit de démontrer le résultat sur $U$ et sur $Z$.
On se ramène ainsi aux deux cas suivants : (a) $\ell$
est inversible sur $S$ ; (b) $\ell$ est nul sur $S$.

\medskip\noindent
Cas (a). Nous affirmons que, dans ce cas, les faisceaux
$R^qf_*\underline{\mathbf{Z}/\ell\mathbf{Z}}$ sont
localement constants à fibres finies sur $S$. D'abord, par changement de base propre
(sous la forme du lemme de Cohomologie étale
\ref{etale-cohomology-lemma-proper-base-change-stalk})
et par finitude
(théorème de Cohomologie étale
\ref{etale-cohomology-theorem-vanishing-affine-curves})
on voit que les fibres de $R^qf_*\underline{\mathbf{Z}/\ell\mathbf{Z}}$
sont finies. D'après le lemme de Cohomologie étale
\ref{etale-cohomology-lemma-sp-isom-proper-loc-cst-torsion}
tous les morphismes de spécialisation sont des isomorphismes.
On conclut grâce au
lemme de Cohomologie étale
\ref{etale-cohomology-lemma-characterize-locally-constant-module}.

\medskip\noindent
Cas (b). Ici, $\ell = p$ est premier et $S$ est un schéma sur
$\Spec(\mathbf{F}_p)$. Les mêmes références que ci-dessus montrent déjà
que les fibres de $R^qf_*\underline{\mathbf{Z}/p\mathbf{Z}}$
sont finies et nulles pour $q \geq 2$. Il résulte du
lemme \ref{etale-cohomology-lemma-sections-upstairs} de Cohomologie étale
que $f_*\underline{\mathbf{Z}/p\mathbf{Z}} =
\underline{\mathbf{Z}/p\mathbf{Z}}$. Il reste à démontrer
que $R^1f_*\underline{\mathbf{Z}/p\mathbf{Z}}$ est constructible.
Considérons la suite d'Artin--Schreier
$$
0 \to \underline{\mathbf{Z}/p\mathbf{Z}} \to \mathcal{O}_X
\xrightarrow{F - 1} \mathcal{O}_X
\to 0
$$
Voir Cohomologie étale, section \ref{etale-cohomology-section-artin-schreier}.
Rappelons que $f_*\mathcal{O}_X = \mathcal{O}_S$ et que
$R^1f_*\mathcal{O}_X$ est un $\mathcal{O}_S$-module localement libre de type fini,
de rang égal au genre des fibres de $X \to S$ ; voir le
lemme \ref{curves-lemma-genus-in-nodal-family-of-curves} de Courbes algébriques.
On en déduit une suite exacte courte
$$
0 \to
\Coker(F - 1 : \mathcal{O}_S \to \mathcal{O}_S)
\to R^1f_*\underline{\mathbf{Z}/p\mathbf{Z}} \to
\Ker(F - 1 : R^1f_*\mathcal{O}_X \to R^1f_*\mathcal{O}_X)
\to 0
$$
Le lemme \ref{more-etale-lemma-frobenius-linear-on-vb} permet de conclure.
\end{proof}

\begin{lemma}
\label{more-etale-lemma-proper-smooth-family-curves-modules}
Soit $f : X \to S$ un morphisme propre et lisse de schémas, à
fibres géométriquement connexes de dimension $1$. Soit $\Lambda$
un anneau noethérien. Soit $M$ un $\Lambda$-module de type fini
annulé par un entier $n > 0$. Alors $R^qf_*\underline{M}$
est un faisceau constructible de $\Lambda$-modules sur $S$.
\end{lemma}

\begin{proof}
Si $n = \ell n'$ pour un nombre premier $\ell$, on obtient
une suite exacte courte $0 \to M[\ell] \to M \to M' \to 0$
de $\Lambda$-modules de type fini, et $M'$ est annulé par $n'$.
Elle donne une suite exacte courte correspondante de faisceaux
constants, qui donne à son tour une suite exacte
$$
R^{q - 1}f_*\underline{M'} \to
R^qf_*\underline{M[n]} \to
R^qf_*\underline{M} \to
R^qf_*\underline{M'} \to
R^{q + 1}f_*\underline{M[n]}
$$
Ainsi, si l'on démontre le résultat lorsque $M$ est annulé par
un nombre premier, une récurrence sur $n$ permet de conclure grâce au
lemme \ref{etale-cohomology-lemma-constructible-abelian} de Cohomologie étale.

\medskip\noindent
Soit $\ell$ un nombre premier tel que $\ell$ annule $M$.
On peut alors remplacer $\Lambda$ par la $\mathbf{F}_\ell$-algèbre
$\Lambda/\ell \Lambda$. En effet, le faisceau $R^qf_*\underline{M}$,
où $\underline{M}$ est vu comme faisceau de $\Lambda$-modules,
s'identifie au faisceau $R^qf_*\underline{M}$ calculé en
considérant $\underline{M}$ comme faisceau de $\Lambda/\ell \Lambda$-modules ; voir le
lemme de Cohomologie sur les sites
\ref{sites-cohomology-lemma-modules-abelian-unbounded}.

\medskip\noindent
Supposons que $\ell$ soit un nombre premier tel que $\ell$ annule $M$ et $\Lambda$.
Ramenons-nous au cas où $M$ est un $\Lambda$-module libre de type fini.
Choisissons pour cela une résolution
$$
\ldots \to
\Lambda^{\oplus m_2} \to
\Lambda^{\oplus m_1} \to
\Lambda^{\oplus m_0} \to
M \to 0
$$
Rappelons que $f_*$ est de dimension cohomologique finie sur les faisceaux de
$\Lambda$-modules ; voir le
lemme de Cohomologie étale
\ref{etale-cohomology-lemma-cohomological-dimension-proper}
et le
lemme \ref{derived-lemma-unbounded-right-derived} de Catégories dérivées.
On voit ainsi que $R^qf_*\underline{M}$ est le $q$-ième faisceau
de cohomologie de l'objet
$$
Rf_*(\underline{\Lambda^{\oplus m_a}} \to \ldots \to
\underline{\Lambda^{\oplus m_0}})
$$
de $D(S_\etale, \Lambda)$ pour un entier $a$ assez grand.
En utilisant la première suite spectrale du
lemme \ref{derived-lemma-two-ss-complex-functor} de Catégories dérivées
(ou bien en utilisant un argument par troncations),
on conclut qu'il suffit de démontrer que $R^qf_*\underline{\Lambda})$
est constructible.

\medskip\noindent
Nous pouvons enfin utiliser l'égalité
$$
(Rf_*\underline{\mathbf{Z}/\ell\mathbf{Z}})
\otimes_{\mathbf{Z}/\ell\mathbf{Z}}^\mathbf{L} \underline{\Lambda} =
Rf_*\underline{\Lambda}
$$
du lemme de Cohomologie étale
\ref{etale-cohomology-lemma-projection-formula-proper-mod-n}.
Comme tout module sur le corps $\mathbf{Z}/\ell\mathbf{Z}$
est plat, on obtient
$$
(R^qf_*\underline{\mathbf{Z}/\ell\mathbf{Z}})
\otimes_{\mathbf{Z}/\ell\mathbf{Z}} \underline{\Lambda} =
R^qf_*\underline{\Lambda}
$$
Il suffit donc de démontrer le résultat pour
$R^qf_*\underline{\mathbf{Z}/\ell\mathbf{Z}}$
d'après le lemme \ref{etale-cohomology-lemma-tensor-constructible} de Cohomologie étale.
C'est précisément le lemme \ref{more-etale-lemma-proper-smooth-family-curves}.
\end{proof}






\section{Complexes à cohomologie constructible}
\label{more-etale-section-Dc}

\noindent
Nous poursuivons la discussion commencée dans Cohomologie étale, section
\ref{etale-cohomology-section-Dc}. En particulier, pour un schéma
$X$ et un anneau noethérien $\Lambda$, nous notons
$D_c(X_\etale, \Lambda)$ la sous-catégorie triangulée strictement pleine et saturée
de $D(X_\etale, \Lambda)$ formée des objets
dont les faisceaux de cohomologie sont des faisceaux constructibles de $\Lambda$-modules.

\begin{lemma}
\label{more-etale-lemma-qf-f-shriek-constructible}
Soit $f : X \to Y$ un morphisme de schémas
localement quasi-fini et de présentation finie.
Le foncteur $f_! : D(X_\etale, \Lambda) \to D(Y_\etale, \Lambda)$
du lemme \ref{more-etale-lemma-lqf-shriek-derived}
envoie $D_c(X_\etale, \Lambda)$ dans $D_c(Y_\etale, \Lambda)$.
\end{lemma}

\begin{proof}
Comme le foncteur $f_!$ est exact, il suffit de montrer que
$f_!\mathcal{F}$ est constructible pour tout faisceau constructible
$\mathcal{F}$ de $\Lambda$-modules sur $X_\etale$.
La question est locale sur $Y$ ; nous pouvons donc supposer
$Y$ affine, ce que nous faisons.
Alors $X$ est quasi-compact et quasi-séparé ; voir la
définition \ref{morphisms-definition-finite-presentation} de Morphismes.
Soit $X = \bigcup_{i = 1, \ldots, n} X_i$ un recouvrement ouvert affine
fini. D'après le lemme \ref{more-etale-lemma-lqf-colimit-f-shriek},
il suffit de montrer que
$f_{i, !}\mathcal{F}|_{X_i}$ et $f_{ii', !}\mathcal{F}|_{X_i \cap X_{i'}}$
sont constructibles, où $f_i : X_i \to Y$ et
$f_{ii'} : X_i \cap X_{i'} \to Y$ sont les restrictions de $f$.
Comme $X_i$ et $X_i \cap X_{i'}$ sont quasi-compacts et séparés,
on peut donc supposer $f$ séparé.
D'après le théorème principal de Zariski (sous la forme du
lemme de Compléments sur les morphismes
\ref{more-morphisms-lemma-quasi-finite-separated-pass-through-finite-addendum})
on peut choisir une factorisation $f = g \circ j$, où
$j : X \to X'$ est une immersion ouverte et $g : X' \to Y$ est
fini et de présentation finie. Alors $f_! = g_! \circ j_!$ d'après le
lemme \ref{more-etale-lemma-f-shriek-composition}.
D'après le lemme \ref{etale-cohomology-lemma-jshriek-constructible} de Cohomologie étale,
$j_!\mathcal{F}$ est constructible sur $X'$.
Le morphisme $g$ est fini, donc $g_! = g_*$
d'après le lemme \ref{more-etale-lemma-proper-f-shriek}.
Ainsi, $f_!\mathcal{F} = g_!j_!\mathcal{F} = g_*j_!\mathcal{F}$
est constructible d'après le
lemme de Cohomologie étale
\ref{etale-cohomology-lemma-finite-pushforward-constructible}.
\end{proof}

\begin{lemma}
\label{more-etale-lemma-constant-shriek-rel-dim-1}
Soit $S$ un schéma affine noethérien de dimension finie.
Soit $f : X \to S$ un morphisme séparé, affine et lisse
de dimension relative $1$. Soit $\Lambda$ un anneau noethérien
de torsion. Soit $M$ un $\Lambda$-module de type fini.
Alors $Rf_!\underline{M}$ a ses faisceaux de cohomologie constructibles.
\end{lemma}

\begin{proof}
Nous démontrerons le résultat par récurrence sur $d = \dim(S)$.

\medskip\noindent
Cas initial. Si $d = 0$, il suffit de montrer que les fibres
de $R^qf_!\underline{M}$ sont des $\Lambda$-modules de type fini.
Si $\overline{s}$ est un point géométrique de $S$, alors
$(R^qf_!\underline{M})_{\overline{s}} = H^q_c(X_{\overline{s}}, \underline{M})$
d'après le lemme \ref{more-etale-lemma-stalk-R-f-shriek}.
C'est un $\Lambda$-module de type fini d'après le lemme \ref{more-etale-lemma-finiteness-curves}.

\medskip\noindent
Étape de récurrence. Il suffit de trouver un ouvert dense $U \subset S$ tel
que $Rf_!\underline{M}|_U$ ait ses faisceaux de cohomologie constructibles.
En effet, la restriction de $Rf_!\underline{M}$ au complémentaire
$S \setminus U$ aura ses faisceaux de cohomologie constructibles par récurrence
et parce que la formation de $Rf_!\underline{M}$ commute
à tout changement de base (lemme \ref{more-etale-lemma-base-change-shriek}).
Plus précisément, soit $\eta \in S$ un point générique d'une composante
irréductible de $S$. Il suffit alors de trouver un voisinage ouvert $U$
de $\eta$ tel que la restriction de $Rf_!\underline{M}$
à $U$ soit constructible. C'est l'objet du paragraphe suivant.

\medskip\noindent
Étant donné un point générique $\eta \in S$, choisissons un diagramme
$$
\xymatrix{
\overline{Y}_1 \amalg \ldots \amalg \overline{Y}_n \ar[rd] &
Y_1 \amalg \ldots \amalg Y_n \ar[r]_-\nu \ar[d] \ar[l]^j &
X_V \ar[r] \ar[d] &
X_U \ar[r] \ar[d] &
X \ar[d]^f \\
& T_1 \amalg \ldots \amalg T_n \ar[r] &
V \ar[r] &
U \ar[r] &
S
}
$$
comme dans le lemme \ref{more-morphisms-lemma-make-good-curves} de Compléments sur les morphismes.
Nous allons montrer que $Rf_!\underline{M}|_U$ est constructible.
D'abord, comme $V \to U$ est fini et surjectif, il suffit
de montrer que son image inverse sur $V$ est constructible ; voir le
lemme \ref{etale-cohomology-lemma-check-constructible} de Cohomologie étale.
Comme la formation de $Rf_!$ commute au changement de base,
il suffit de montrer que
$R(X_V \to V)_!\underline{M}$ est constructible.
Soit $W \subset X_V$ le sous-schéma ouvert fourni par la partie (4) du
lemme \ref{more-morphisms-lemma-make-good-curves} de Compléments sur les morphismes.
Soit $Z \subset X_V$ le sous-schéma induit réduit sur
le complémentaire de $W$ dans $X_V$. Alors les fibres de $Z \to V$
sont de dimension $0$ (car $W$ est dense dans les fibres), donc
$Z \to V$ est quasi-fini. Du triangle distingué
$$
R(W \to V)_!\underline{M} \to
R(X_V \to V)_!\underline{M} \to
R(Z \to V)_!\underline{M} \to
\ldots
$$
du lemme \ref{more-etale-lemma-relative-triangle-associated-to-open} et
du lemme \ref{more-etale-lemma-qf-f-shriek-constructible},
on déduit qu'il suffit de montrer que
$R(W \to V)_!\underline{M}$ a ses faisceaux de cohomologie constructibles.
Ensuite, on a
$$
R(W \to V)_!\underline{M} = R(\nu^{-1}(W) \to V)_!\underline{M}
$$
parce que le morphisme $\nu : \nu^{-1}(W) \to W$ est un épaississement
et que l'on peut appliquer le lemme \ref{more-etale-lemma-shriek-and-thickening}.
Notons ensuite $Z' \subset \coprod \overline{Y}_i$ le
complémentaire de l'ouvert $j(\nu^{-1}(W))$. Là encore, $Z' \to V$ est quasi-fini.
Utilisons encore le triangle distingué
$$
R(\nu^{-1}(W) \to V)_!\underline{M} \to
R(\coprod \overline{Y}_i \to V)_!\underline{M} \to
R(Z' \to V)_!\underline{M} \to
\ldots
$$
pour conclure qu'il suffit de démontrer que
$$
R(\coprod \overline{Y}_i \to V)_!\underline{M} =
\bigoplus\nolimits_i R(\overline{Y}_i \to V)_!\underline{M} =
\bigoplus\nolimits_i R(T_i \to V)_!R(\overline{Y}_i \to T_i)_!\underline{M}
$$
a ses faisceaux de cohomologie constructibles (la seconde égalité résultant du
lemme \ref{more-etale-lemma-shriek-composition}).
Le résultat pour $R(\overline{Y}_i \to T_i)_!\underline{M}$
résulte du lemme \ref{more-etale-lemma-proper-smooth-family-curves-modules},
et l'on conclut parce que $T_i \to V$ est fini étale
et que l'on peut appliquer le lemme \ref{more-etale-lemma-qf-f-shriek-constructible}.
\end{proof}

\begin{lemma}
\label{more-etale-lemma-loc-constant-shriek-rel-dim-1}
Soit $Y$ un schéma affine noethérien de dimension finie.
Soit $\Lambda$ un anneau noethérien de torsion.
Soit $\mathcal{F}$ un faisceau de
$\Lambda$-modules localement constant et de type fini sur un sous-schéma ouvert $U \subset \mathbf{A}^1_Y$.
Alors $Rf_!\mathcal{F}$ a ses faisceaux de cohomologie constructibles, où
$f : U \to Y$ est le morphisme structural.
\end{lemma}

\begin{proof}
On peut décomposer $\Lambda$ en un produit
$\Lambda = \Lambda_1 \times \ldots \times \Lambda_r$
où $\Lambda_i$ est $\ell_i$-primaire pour un nombre premier $\ell_i$.
On peut donc supposer qu'il existe un nombre premier $\ell$ et un entier $n > 0$
tels que $\ell^n$ annule $\Lambda$ (et donc $\mathcal{F}$).

\medskip\noindent
Comme $U$ est noethérien, $U$ n'a qu'un nombre fini de composantes
connexes. On peut donc supposer $U$ connexe.
Soit $g : U' \to U$ le revêtement étale fini
construit dans le lemme de Cohomologie étale
\ref{etale-cohomology-lemma-pullback-filtered-modules}.
La discussion de la
section \ref{etale-cohomology-section-trace-method} de Cohomologie étale
donne des morphismes
$$
\mathcal{F} \to g_*g^{-1}\mathcal{F} \to \mathcal{F}
$$
dont la composée est un isomorphisme. Il suffit donc de
démontrer le résultat pour $g_*g^{-1}\mathcal{F}$. D'autre part,
on a $Rf_!g_*g^{-1}\mathcal{F} = R(f \circ g)_!g^{-1}\mathcal{F}$
d'après le lemme \ref{more-etale-lemma-shriek-composition}.
Comme $g^{-1}\mathcal{F}$ admet une filtration finie dont les quotients sont des
faisceaux constants de $\Lambda$-modules de la forme
$\underline{M}$ pour un $\Lambda$-module de type fini $M$
(par notre choix de $g$), on se ramène au cas démontré au
lemme \ref{more-etale-lemma-constant-shriek-rel-dim-1}.
\end{proof}

\begin{lemma}
\label{more-etale-lemma-constructible-shriek-rel-dim-1}
Soit $Y$ un schéma affine. Soit $\Lambda$ un anneau noethérien.
Soit $\mathcal{F}$ un faisceau constructible de $\Lambda$-modules de torsion sur
$\mathbf{A}^1_Y$. Alors $Rf_!\mathcal{F}$ a ses faisceaux de cohomologie
constructibles, où $f : \mathbf{A}^1_Y \to Y$ est le morphisme structural.
\end{lemma}

\begin{proof}
Supposons $\mathcal{F}$ annulé par $n > 0$. On peut alors remplacer
$\Lambda$ par $\Lambda/n\Lambda$ sans changer $Rf_!\mathcal{F}$.
On peut donc supposer, et supposons, que $\Lambda$ est un anneau de torsion.

\medskip\noindent
Écrivons $Y = \Spec(R)$. Si l'on écrit
$R = \bigcup R_i$ comme réunion de ses sous-$\mathbf{Z}$-algèbres de type fini,
on peut trouver un $i$ tel que $\mathcal{F}$ provienne par image inverse
d'un faisceau constructible de $\Lambda$-modules sur $\mathbf{A}^1_{R_i}$ ; voir le
lemme \ref{etale-cohomology-lemma-category-is-colimit} de Cohomologie étale.
On peut donc supposer $Y$ noethérien et de dimension finie.

\medskip\noindent
Supposons $Y$ noethérien de dimension finie $d = \dim(Y)$
et $\Lambda$ un anneau de torsion. Nous démontrerons le résultat par récurrence sur $d$.

\medskip\noindent
Cas initial. Si $d = 0$, il suffit de montrer que les fibres
de $R^qf_!\mathcal{F}$ sont des $\Lambda$-modules de type fini.
Si $\overline{y}$ est un point géométrique de $Y$, alors
$(R^qf_!\mathcal{F})_{\overline{y}} = H^q_c(X_{\overline{y}}, \mathcal{F})$
d'après le lemme \ref{more-etale-lemma-stalk-R-f-shriek}.
C'est un $\Lambda$-module de type fini d'après le lemme \ref{more-etale-lemma-finiteness-curves}.

\medskip\noindent
Étape de récurrence. Il suffit de trouver un ouvert dense $V \subset Y$ tel
que $Rf_!\mathcal{F}|_V$ ait ses faisceaux de cohomologie constructibles. En effet, la
restriction de $Rf_!\mathcal{F}$ au complémentaire $Y \setminus V$
aura ses faisceaux de cohomologie constructibles par récurrence
et parce que la formation de $Rf_!\mathcal{F}$ commute
à tout changement de base (lemme \ref{more-etale-lemma-base-change-shriek}).
Par définition des faisceaux constructibles de $\Lambda$-modules, il existe un
sous-schéma ouvert dense $U \subset \mathbf{A}^1_Y$ tel que $\mathcal{F}|_U$
soit un faisceau de $\Lambda$-modules localement constant et de type fini.
Notons $Z \subset \mathbf{A}^1_Y$ le complémentaire (muni de sa structure de
sous-schéma fermé réduit). Remarquons que $U$ contient tous les points génériques des
fibres de $\mathbf{A}^1_Y \to Y$ au-dessus des points génériques
$\xi_1, \ldots, \xi_n$ des composantes irréductibles de $Y$.
Ainsi, $Z \to Y$ a des fibres finies au-dessus de $\xi_1, \ldots, \xi_n$.
Quitte à remplacer $Y$ par un ouvert dense (ce qui est permis), on peut
supposer $Z \to Y$ fini ; voir le
lemme \ref{morphisms-lemma-generically-finite} de Morphismes.
Le triangle distingué du
lemme \ref{more-etale-lemma-relative-triangle-associated-to-open}
et le résultat pour $Z \to Y$ (lemme \ref{more-etale-lemma-qf-f-shriek-constructible})
nous ramènent à montrer que
$R(U \to Y)_!\mathcal{F}$ a ses faisceaux de cohomologie constructibles.
C'est précisément le lemme \ref{more-etale-lemma-loc-constant-shriek-rel-dim-1}.
\end{proof}

\begin{theorem}
\label{more-etale-theorem-constructible-shriek}
Soit $f : X \to Y$ un morphisme séparé de présentation finie entre
schémas quasi-compacts et quasi-séparés. Soit $\Lambda$ un anneau
noethérien. Soit $K$ un objet de $D^+_{tors, c}(X_\etale, \Lambda)$
ou de $D_c(X_\etale, \Lambda)$ lorsque $\Lambda$ est un anneau de torsion.
Alors $Rf_!K$ a ses faisceaux de cohomologie constructibles, c'est-à-dire que
$Rf_!K$ appartient à $D^+_{tors, c}(Y_\etale, \Lambda)$
ou à $D_c(Y_\etale, \Lambda)$ lorsque $\Lambda$ est un anneau de torsion.
\end{theorem}

\begin{proof}
La question est locale sur $Y$ ; nous pouvons donc supposer $Y$ affine.
Par le principe de récurrence et le lemme \ref{more-etale-lemma-relative-mayer-vietoris},
on se ramène au cas où $X$ est aussi affine.

\medskip\noindent
Supposons $X$ et $Y$ affines. Comme $X$ est de présentation finie,
on peut choisir une immersion fermée $i : X \to \mathbf{A}^n_Y$
de présentation finie. Si $p : \mathbf{A}^n_Y \to Y$ désigne
le morphisme structural, on a $Rf_! = Rp_! \circ Ri_!$
d'après le lemme \ref{more-etale-lemma-shriek-composition}. D'après le
lemme \ref{more-etale-lemma-qf-f-shriek-constructible},
le résultat vaut pour $Ri_! = i_!$. On peut donc supposer que $f$
est le morphisme de projection $\mathbf{A}^n_Y \to Y$.
Comme on peut écrire $f$ comme la composée
$$
X = \mathbf{A}^n_Y \to \mathbf{A}^{n - 1}_Y \to
\mathbf{A}^{n - 2}_S \to \ldots \to
\mathbf{A}^1_Y \to Y
$$
on peut supposer $n = 1$.

\medskip\noindent
Supposons $Y$ affine et $X = \mathbf{A}^1_Y$. Comme $Rf_!$ est de
dimension cohomologique finie
(lemme \ref{more-etale-lemma-derived-lower-shriek-bounded}),
on peut supposer $K$ borné inférieurement. En utilisant la première suite spectrale du
lemme \ref{derived-lemma-two-ss-complex-functor} de Catégories dérivées
(ou bien en utilisant un argument par troncations),
on se ramène au résultat du
lemme \ref{more-etale-lemma-constructible-shriek-rel-dim-1}.
\end{proof}








\section{Applications}
\label{more-etale-section-applications}

\noindent
Dans cette section, nous donnons quelques applications du
théorème \ref{more-etale-theorem-constructible-shriek}.

\begin{lemma}
\label{more-etale-lemma-finiteness-compactly-supported}
Soit $k$ un corps algébriquement clos.
Soit $X$ un schéma séparé de type fini sur $k$.
Soit $\Lambda$ un anneau noethérien. Soit $K$ un objet de
$D^+_{tors, c}(X_\etale, \Lambda)$ ou de $D_c(X_\etale, \Lambda)$
lorsque $\Lambda$ est un anneau de torsion. Alors
$H^i_c(X, K)$ est un $\Lambda$-module de type fini pour tout
$i \in \mathbf{Z}$.
\end{lemma}

\begin{proof}
C'est une conséquence immédiate du théorème \ref{more-etale-theorem-constructible-shriek}
et de la définition de la cohomologie à support propre donnée à la
section \ref{more-etale-section-compactly-supported-cohomology}.
\end{proof}

\begin{proposition}
\label{more-etale-proposition-loc-cst-torsion}
Soit $f : X \to S$ un morphisme propre et lisse entre schémas.
Soit $\Lambda$ un anneau noethérien. Soit $\mathcal{F}$ un
faisceau de $\Lambda$-modules localement constant et de type fini
sur $X_\etale$ tel que, pour tout point géométrique $\overline{x}$ de $X$,
la fibre $\mathcal{F}_{\overline{x}}$ soit annulée par un entier
$n > 0$ premier à la caractéristique résiduelle de $\overline{x}$.
Alors $R^if_*\mathcal{F}$ est un faisceau de
$\Lambda$-modules localement constant et de type fini sur $S_\etale$ pour tout $i \in \mathbf{Z}$.
\end{proposition}

\begin{proof}
La question est locale sur $S$ ; on peut donc supposer $S$ affine.
Pour un point $x$ de $X$, notons $n_x \geq 1$
le plus petit entier annulant $\mathcal{F}_{\overline{x}}$
pour un point géométrique $\overline{x}$ de $X$
situé au-dessus de $x$ (et donc pour tout tel point). Comme $X$ est quasi-compact (étant propre sur un affine),
il existe un recouvrement étale fini $\{U_j \to X\}_{j = 1, \ldots, m}$
tel que $\mathcal{F}|_{U_j}$ soit constant. Comme $U_j \to X$
est ouvert, on en déduit que la fonction $x \mapsto n_x$ est localement
constante et ne prend qu'un nombre fini de valeurs. On obtient donc
une décomposition finie $X = X_1 \amalg \ldots \amalg X_N$
en sous-schémas ouverts et fermés telle que $n_x = n$ si et seulement si $x \in X_n$.
Il suffit alors de démontrer la proposition pour les morphismes induits
$X_n \to S$ et la restriction de $\mathcal{F}$ à $X_n$.
On peut donc supposer qu'il existe un entier $n > 0$ tel
que $\mathcal{F}$ soit annulé par $n$ et que $n$
soit premier aux caractéristiques résiduelles de tous les corps résiduels de $X$.

\medskip\noindent
Comme $f$ est lisse et propre, l'image $f(X) \subset S$ est ouverte et fermée.
On peut donc remplacer $S$ par $f(X)$ et supposer $f(X) = S$. En particulier,
on peut supposer $n$ inversible dans l'anneau de définition du
schéma affine $S$.

\medskip\noindent
Dans ce paragraphe, nous nous ramenons au cas où $S$ est noethérien.
Écrivons $S = \Spec(A)$ pour une $\mathbf{Z}[1/n]$-algèbre $A$.
Écrivons $A = \bigcup A_i$ comme réunion de ses sous-algèbres de type fini
sur $\mathbf{Z}[1/n]$. On peut trouver un $i$ et un morphisme
$f_i : X_i \to S_i = \Spec(A_i)$ de type
fini dont le changement de base à $S$ est $f : X \to S$ ; voir le
lemme \ref{limits-lemma-descend-finite-presentation} de Limites.
Quitte à augmenter $i$, on peut supposer $f_i : X_i \to S_i$
lisse et propre ; voir les lemmes de Limites
\ref{limits-lemma-eventually-proper},
\ref{limits-lemma-descend-smooth} et
\ref{limits-lemma-descend-dimension-d}.
D'après le lemme de Cohomologie étale
\ref{etale-cohomology-lemma-category-loc-cst-is-colimit}
il existe un $i$ et un faisceau
de $\Lambda$-modules localement constant et de type fini
$\mathcal{F}_i$ dont l'image inverse sur $X$ est isomorphe à $\mathcal{F}$.
Comme $\mathcal{F}$ est annulé par $n$, on peut remplacer
$\mathcal{F}_i$ par $\Ker(n : \mathcal{F}_i \to \mathcal{F}_i)$
et supposer qu'il en va de même pour $\mathcal{F}_i$.
On se ramène ainsi au cas traité dans le paragraphe suivant.

\medskip\noindent
Supposons donnés un entier $n \geq 1$, un schéma de base noethérien $S$
sur $\mathbf{Z}[1/n]$, et un faisceau $\mathcal{F}$ de $n$-torsion.
D'après le théorème \ref{more-etale-theorem-constructible-shriek},
les faisceaux $R^if_*\mathcal{F}$ sont des faisceaux constructibles de $\Lambda$-modules.
D'après le lemme de Cohomologie étale
\ref{etale-cohomology-lemma-sp-isom-proper-torsion-loc-ac}
les morphismes de spécialisation de $R^if_*\mathcal{F}$ sont toujours des isomorphismes.
On conclut par le lemme de Cohomologie étale
\ref{etale-cohomology-lemma-characterize-locally-constant-module}.
\end{proof}







\section{Compléments sur l'image inverse extraordinaire dérivée}
\label{more-etale-section-more-derived-upper-shriek}

\noindent
Soit $\Lambda$ un anneau de torsion. Considérons un diagramme commutatif
$$
\xymatrix{
U \ar[rr]_j \ar[rd]_g & & U' \ar[ld]^{g'} \\
& Y
}
$$
de schémas quasi-compacts et quasi-séparés, où $g$ et $g'$
sont séparés et de type fini, et où $j$ est étale. Il induit
un morphisme canonique
$$
Rg_!\Lambda \longrightarrow Rg'_!\Lambda
$$
dans $D(Y_\etale, \Lambda)$. En effet, d'après les lemmes \ref{more-etale-lemma-shriek-composition}
et \ref{more-etale-lemma-Rf-shriek-for-quasi-finite}, on a $Rg_! = Rg'_! \circ j_!$.
D'autre part, comme $j_!$ est adjoint à gauche de $j^{-1}$, on dispose du
morphisme d'adjonction $\text{Tr}_j : j_!\Lambda = j_!j^{-1}\Lambda \to \Lambda$ ;
nous l'appelons aussi le morphisme trace de $j$ ; voir la
remarque \ref{more-etale-remark-trace-etale-counit}.
Le morphisme ci-dessus est la composée
$$
Rg_!\Lambda = Rg'_!j_!\Lambda
\xrightarrow{Rg'_! \text{Tr}_j}
Rg'_!\Lambda
$$
Étant donné un second morphisme étale $j' : U' \to U''$
et un morphisme séparé de type fini $g'' : U'' \to Y$,
la composée
$$
Rg_!\Lambda \longrightarrow Rg'_!\Lambda \longrightarrow Rg''_!\Lambda
$$
des morphismes associés à $j$ et $j'$ est égale au morphisme
$Rg_!\Lambda \longrightarrow Rg''_!\Lambda$ construit pour $j' \circ j$.
Cela résulte de l'énoncé correspondant pour les morphismes trace ;
voir le lemme \ref{more-etale-lemma-trace-composition} pour un cas plus général.

\medskip\noindent
Soit $f : X \to Y$ un morphisme séparé de type fini entre schémas quasi-compacts
et quasi-séparés. On obtient alors un foncteur
$$
X_{affine, \etale}
\longrightarrow
\left\{
\begin{matrix}
\text{schémas séparés de type fini sur }Y\\
\text{et morphismes étales entre eux}
\end{matrix}
\right\}
$$
La construction ci-dessus détermine donc un foncteur
$X_{affine, \etale}^{opp} \to D(Y_\etale, \Lambda)$
qui envoie $U$ sur $R(U \to Y)_!\Lambda$.

\begin{lemma}
\label{more-etale-lemma-describe-Rf-upper-shriek}
Soit $f : X \to Y$ un morphisme séparé de type fini entre schémas quasi-compacts
et quasi-séparés. Soit $\Lambda$ un anneau de torsion.
Soit $K \in D(Y_\etale, \Lambda)$. Pour $n \in \mathbf{Z}$, le
faisceau de cohomologie $H^n(Rf^!K)$ restreint à $X_{affine, \etale}$
est le faisceau associé au préfaisceau
$$
U \longmapsto \Hom_Y(R(U \to Y)_!\Lambda, K[n])
$$
Voir la discussion ci-dessus pour la fonctorialité de $R(U \to Y)_!\Lambda$.
\end{lemma}

\begin{proof}
Soit $j : U \to X$ un objet de $X_{affine, \etale}$, et posons $g = f \circ j$.
Rappelons que $\Hom_X(j_!\Lambda, M[n]) = H^n(U, M)$ pour tout objet $M$ de
$D(X_\etale, \Lambda)$. Alors $H^n(Rf^!K)$ est le faisceau associé au
préfaisceau
$$
U \mapsto H^n(U, Rf^!K) = \Hom_X(j_!\Lambda, Rf^!K[n]) =
\Hom_Y(Rf_!j_!\Lambda, K[n] = \Hom_Y(Rg_!\Lambda, K[n])
$$
Nous omettons de vérifier que les morphismes de transition sont induits par les morphismes
de transition entre les objets $Rg_!\Lambda = R(U \to Y)_!\Lambda$ construits
ci-dessus.
\end{proof}




















% OK, start here.
%

\chapter{La formule des traces}



\phantomsection
\label{trace-section-phantom}



\section{Introduction}
\label{trace-section-introduction}

\noindent
Voici les notes de la seconde partie d'un cours de cohomologie \'etale
donné par Johan de Jong à l'Université Columbia à l'automne 2009. Les
auteurs des notes initiales étaient Thibaut Pugin, Zachary Maddock et Min Lee.
Au fil du temps, nous ajouterons des références aux prérequis dans le reste du
projet Champs et donnerons des démonstrations rigoureuses de tous les énoncés.









\section{La formule des traces}
\label{trace-section-trace-formula}

\noindent
Un cours classique de cohomologie \'etale énoncerait et démontrerait normalement
les théorèmes de changement de base propre et lisse, la pureté et la dualité de
Poincar\'e. Tout cela se trouve dans \cite[Arcata]{SGA4.5}. Nous allons plutôt
étudier la formule des traces pour le Frobenius, en suivant l'exposé de Deligne
dans \cite[Rapport]{SGA4.5}. Nous nous limiterons à la dimension 1, mais le
changement de base propre suffit alors à traiter le cas général. Puisque tous
les groupes de cohomologie considérés seront \'etales, nous omettons l'indice
$_\etale$. Décrivons maintenant
la formule que nous cherchons. Soient $X$ un schéma de type fini de dimension 1
sur un corps fini $k$, $\ell$ un nombre premier et $\mathcal{F}$ un faisceau
constructible et plat de $\mathbf{Z}/\ell^n\mathbf{Z}$-modules. Alors
\begin{equation}
\label{trace-equation-trace-formula-initial}
\sum\nolimits_{x \in X(k)}
\text{Tr}(\text{Frob} | \mathcal{F}_{\bar x}) =
\sum\nolimits_{i = 0}^2
(-1)^i \text{Tr}(\pi_X^* | H^i_c(X \otimes_k \bar k, \mathcal{F}))
\end{equation}
comme éléments de $\mathbf{Z}/\ell^n\mathbf{Z}$. Comme nous le verrons, cette
formulation est légèrement incorrecte telle quelle. Décrivons néanmoins les
symboles qui y figurent.




\section{Les Frobenius}
\label{trace-section-frobenii}

\noindent
Dans cette section, nous allons démontrer un théorème ``déroutant''.
Un analogue topologique de ce théorème déroutant est le suivant.

\begin{exercise}
\label{trace-exercise-baffling}
Soient $X$ un espace topologique et $g : X \to X$ une application continue telle
que $g^{-1}(U) = U$ pour tout ouvert $U$ de $X$. Alors $g$ induit l'identité sur
la cohomologie de $X$ (à coefficients quelconques).
\end{exercise}

\noindent
Passons maintenant à l'énoncé pour le site \'etale.

\begin{lemma}
\label{trace-lemma-baffling}
Soient $X$ un schéma et $g : X \to X$ un morphisme. Supposons que, pour tout
morphisme \'etale $\varphi : U \to X$, il existe un isomorphisme
$$
\xymatrix{
U \ar[rd]_\varphi \ar[rr]^-\sim & & {U
\times_{\varphi, X, g} X} \ar[ld]^{\text{pr}_2} \\
& X
}
$$
fonctoriel en $U$. Alors $g$ induit l'identité en cohomologie (pour tout faisceau).
\end{lemma}

\begin{proof}
La démonstration est formelle et ne présente aucune difficulté.
\end{proof}

\noindent
Voir Variétés, section \ref{varieties-section-frobenius},
pour une discussion des différentes variantes du morphisme de Frobenius.

\begin{theorem}[Le théorème déroutant]
\label{trace-theorem-baffling}
Soit $X$ un schéma de caractéristique $p > 0$. Alors le Frobenius absolu
induit (par image réciproque) l'identité en cohomologie, c'est-à-dire que, pour
tout entier $j\geq 0$,
$$
F_X^* : H^j (X, \underline{\mathbf{Z}/n\mathbf{Z}}) \longrightarrow H^j (X,
\underline{\mathbf{Z}/n\mathbf{Z}})
$$
est l'identité.
\end{theorem}

\noindent
Ce théorème est purement formel. Il est toutefois utile de revoir comment
calculer l'image réciproque d'une classe de cohomologie. Disons simplement que,
lorsque la cohomologie coïncide avec celle de {\v C}ech, il suffit de tirer en
arrière les cocycles de {\v C}ech (au moyen des produits fibrés sur un site). Le
cas général est assez technique ; voir
Hyperrecouvrements, théorème \ref{hypercovering-theorem-cohomology-hypercoverings}.
Pour démontrer le théorème, il nous suffit
de vérifier que l'hypothèse du lemme \ref{trace-lemma-baffling}
est satisfaite par le Frobenius.

\begin{proof}[Démonstration du théorème \ref{trace-theorem-baffling}]
Nous devons vérifier l'existence d'un isomorphisme fonctoriel comme ci-dessus.
Pour un morphisme \'etale $\varphi : U \to X$, considérons le diagramme
$$
\xymatrix{
U \ar@{-->}[rd] \ar@/^1pc/[rrd]^{F_U}
\ar@/_1pc/[rdd]_\varphi \\
& {U \times_{\varphi, X, F_X} X} \ar[r]_-{\text{pr}_1}
\ar[d]^{\text{pr}_2} & U \ar[d]^\varphi \\
& X \ar[r]^{F_X} & X.
}
$$
La flèche en pointillés est un morphisme \'etale et un homéomorphisme
universel ; c'est donc un isomorphisme. Voir
Morphismes étales, lemme \ref{etale-lemma-relative-frobenius-etale}.
\end{proof}

%10.22.09

\begin{definition}
\label{trace-definition-geometric-frobenius}
Soit $k$ un corps fini à $q = p^f$ éléments. Soit $X$ un schéma
sur $k$. Le {\it Frobenius géométrique} de $X$ est le morphisme
$\pi_X : X \to X$ au-dessus de $\Spec(k)$ égal à $F_X^f$.
\end{definition}

\noindent
Puisque $\pi_X$ est un morphisme sur $k$, nous pouvons le changer de base par
n'importe quel schéma sur $k$. En particulier, le changement de base à la
clôture algébrique $\bar k$ donne un morphisme $\pi_X : X_{\bar k} \to
X_{\bar k}$. Employer aussi $\pi_X$ pour ce changement de base ne devrait pas
prêter à confusion, car $X_{\bar k}$ n'a pas de Frobenius géométrique propre.

\begin{lemma}
\label{trace-lemma-sheaf-over-finite-field-has-frobenius-descent}
Soit $\mathcal{F}$ un faisceau sur $X_\etale$.
Il existe alors des isomorphismes canoniques
$\pi_X^{-1} \mathcal{F} \cong \mathcal{F}$ et
$\mathcal{F} \cong {\pi_X}_*\mathcal{F}$.
\end{lemma}

\noindent
Cela est faux pour le site fppf.

\begin{proof}
Soit $\varphi : U \to X$ un morphisme \'etale. Rappelons que
${\pi_X}_* \mathcal{F} (U) = \mathcal{F} (U \times_{\varphi, X, \pi_X} X)$.
Puisque $\pi_X = F_X^f$, la démonstration du
théorème \ref{trace-theorem-baffling} fournit un isomorphisme fonctoriel
$$
\xymatrix{
U \ar[rd]_{\varphi} \ar[rr]_-{\gamma_U}
& & U \times_{\varphi, X, \pi_X} X \ar[ld]^{\text{pr}_2} \\
& X
}
$$
où $\gamma_U = (\varphi, F_U^f)$. Nous définissons maintenant un
isomorphisme
$$
\mathcal{F} (U) \longrightarrow {\pi_X}_* \mathcal{F} (U) =
\mathcal{F} (U \times_{\varphi, X, \pi_X} X)
$$
comme l'application de restriction de $\mathcal{F}$ le long de $\gamma_U^{-1}$.
L'autre isomorphisme se construit de manière analogue.
\end{proof}

\begin{remark}
\label{trace-remark-may-be-confusing}
Il se peut, ou non, que $F^f_U$ soit égal à $\pi_U$.
\end{remark}

\noindent
Poursuivons l'étude de la cohomologie des faisceaux sur notre schéma $X$ défini
sur le corps fini $k$ à $q = p^f$ éléments.
Fixons une clôture algébrique $\bar k$ de $k$ et notons $G_k =
\text{Gal}(\bar k/k)$ le groupe de Galois absolu de $k$.
Soit $\mathcal{F}$ un faisceau abélien sur $X_\etale$.
Nous allons munir d'une structure de $G_k$-module à gauche le
groupe de cohomologie $H^j (X_{\bar k}, \mathcal{F}|_{X_{\bar k}})$
de la manière suivante : si $\sigma \in G_k$, le diagramme
$$
\xymatrix{
X_{\bar k} \ar[rd] \ar[rr]^{\Spec(\sigma) \times \text{id}_X} & &
X_{\bar k} \ar[ld] \\
& X
}
$$
est commutatif. Nous pouvons donc poser, pour $\xi \in H^j (X_{\bar k},
\mathcal{F}|_{X_{\bar k}})$,
$$
\sigma \cdot \xi := (\Spec(\sigma) \times \text{id}_X)^*\xi \in
H^j(X_{\bar k}, (\Spec(\sigma) \times \text{id}_X)^{-1}
\mathcal{F}|_{X_{\bar k}})
= H^j (X_{\bar k}, \mathcal{F}|_{X_{\bar k}}),
$$
où la dernière égalité résulte de la commutativité du diagramme précédent.
Cela munit ce dernier groupe d'une structure de $G_k$-module.

\begin{lemma}
\label{trace-lemma-two-actions-agree}
Dans la situation ci-dessus, notons $\alpha : X \to \Spec(k)$ le morphisme
structural. Considérons la fibre $(R^j\alpha_*\mathcal{F})_{\Spec(\bar k)}$
munie de son action galoisienne naturelle comme dans
Cohomologie étale, section \ref{etale-cohomology-section-galois-action-stalks}.
Alors l'identification
$$
(R^j\alpha_*\mathcal{F})_{\Spec(\bar k)} \cong H^j (X_{\bar k},
\mathcal{F}|_{X_{\bar k}})
$$
du
théorème \ref{etale-cohomology-theorem-higher-direct-images} de Cohomologie étale
est un isomorphisme de $G_k$-modules.
\end{lemma}

\noindent
Un résultat analogue vaut lorsque l'on compare
$(R^j\alpha_!\mathcal{F})_{\Spec(\bar k)}$ à
$H^j_c (X_{\bar k}, \mathcal{F}|_{X_{\bar k}})$.

\begin{proof}
Omise.
\end{proof}

\begin{definition}
\label{trace-definition-arithmetic-frobenius}
Le {\it Frobenius arithmétique} est l'application
$\text{frob}_k : \bar k \to \bar k$, $x \mapsto x^q$ appartenant à $G_k$.
\end{definition}

\begin{theorem}
\label{trace-theorem-geometric-arithmetic-inverse}
Soit $\mathcal{F}$ un faisceau abélien sur $X_\etale$. Alors, pour tout
$j\geq 0$, $\text{frob}_k$ agit sur le groupe de cohomologie $H^j(X_{\bar k},
\mathcal{F}|_{X_{\bar k}})$ comme l'inverse de l'application $\pi_X^*$.
\end{theorem}

\noindent
L'application $\pi_X^*$ est définie par la composée
$$
H^j(X_{\bar k}, \mathcal{F}|_{X_{\bar k}}) \xrightarrow{{\pi_X}_{\bar k}^*}
H^j(X_{\bar k}, (\pi_X^{-1} \mathcal{F})|_{X_{\bar k}}) \cong
H^j(X_{\bar k}, \mathcal{F}|_{X_{\bar k}}).
$$
où le dernier isomorphisme provient de l'isomorphisme canonique
$\pi_X^{-1} \mathcal{F} \cong \mathcal{F}$ du
lemme \ref{trace-lemma-sheaf-over-finite-field-has-frobenius-descent}.

\begin{proof}
La composée $X_{\bar k} \xrightarrow{\Spec(\text{frob}_k)} X_{\bar k}
\xrightarrow{\pi_X} X_{\bar k}$ est égale à $F_{X_{\bar k}}^f$ ; le
résultat découle donc du théorème déroutant convenablement généralisé à des
coefficients non triviaux. Notons que la composée précédente commute au sens où
$F_{X_{\bar k}}^f = \pi_X \circ \Spec(\text{frob}_k) =
\Spec(\text{frob}_k) \circ \pi_X$.
\end{proof}

\begin{definition}
\label{trace-definition-geometric-frobenius-on-stalk}
Si $x \in X(k)$ est un point rationnel et $\bar x : \Spec(\bar k) \to X$
est le point géométrique au-dessus de $x$, nous notons $\pi_x : \mathcal{F}_{\bar x} \to
\mathcal{F}_{\bar x}$ l'action de $\text{frob}_k^{-1}$ et l'appelons le
{\it Frobenius géométrique}\footnote{Cette notation n'est pas standard.
Cet opérateur est noté $F_x$ dans \cite{SGA4.5}. Nous changerons probablement
cette notation à l'avenir.}
\end{definition}

\noindent
Nous pouvons maintenant donner un énoncé plus précis (quoique faux) de la formule des
traces (\ref{trace-equation-trace-formula-initial}). Soit $X$ un schéma de
type fini de dimension 1
sur un corps fini $k$, $\ell$ un nombre premier et $\mathcal{F}$ un faisceau
constructible et plat de $\mathbf{Z}/\ell^n\mathbf{Z}$-modules. Alors
\begin{equation}
\label{trace-equation-trace-formula-second}
\sum\nolimits_{x \in X(k)}
\text{Tr}(\pi_x | \mathcal{F}_{\bar x})
=
\sum\nolimits_{i = 0}^2
(-1)^i \text{Tr}(\pi_X^* | H^i_c(X_{\bar k}, \mathcal{F}))
\end{equation}
comme éléments de $\mathbf{Z}/\ell^n\mathbf{Z}$. Cette égalité est incorrecte
parce que la trace du membre de droite n'a pas de sens pour le type de
faisceaux considéré. Avant d'aborder ce problème, cherchons à motiver
l'apparition du Frobenius géométrique (outre le fait qu'il s'agit d'un
morphisme naturel !).

\medskip\noindent
Considérons le cas où $X = \mathbf{P}^1_k$ et $\mathcal{F} =
\underline{\mathbf{Z}/\ell\mathbf{Z}}$. En tout point, le module galoisien
$\mathcal{F}_{\bar x}$ est trivial ; ainsi, pour tout morphisme $\varphi$ agissant sur
$\mathcal{F}_{\bar x}$, le membre de gauche vaut
$$
\sum\nolimits_{x \in X(k)} \text{Tr}(\varphi | \mathcal{F}_{\bar x}) =
\#\mathbf{P}^1_k(k) = q+1.
$$
Or $\mathbf{P}^1_k$ est propre, donc la cohomologie à support propre coïncide avec la
cohomologie usuelle ; ainsi, pour un morphisme $\pi : \mathbf{P}^1_k \to
\mathbf{P}^1_k$, le membre de droite vaut
$$
\text{Tr}(\pi^* | H^0 (\mathbf{P}^1_{\bar k},
\underline{\mathbf{Z}/\ell\mathbf{Z}})) + \text{Tr}(\pi^* | H^2
(\mathbf{P}^1_{\bar k}, \underline{\mathbf{Z}/\ell\mathbf{Z}})).
$$
Le module galoisien $H^0 (\mathbf{P}^1_{\bar k},
\underline{\mathbf{Z}/\ell\mathbf{Z}}) = \mathbf{Z}/\ell\mathbf{Z}$ est trivial,
car l'image réciproque de l'identité est l'identité. La première trace vaut donc 1,
quel que soit $\pi$. Pour la seconde trace, nous devons calculer l'image réciproque
$\pi^* : H^2(\mathbf{P}^1_{\bar k}, \underline{\mathbf{Z}/\ell\mathbf{Z}}))$
pour une application $\pi : \mathbf{P}^1_{\bar k} \to \mathbf{P}^1_{\bar k}$. C'est un
bon exercice, et la réponse est la multiplication par le degré de $\pi$
(pour une démonstration, voir
Cohomologie étale, lemme \ref{etale-cohomology-lemma-pullback-on-h2-curve}).
Autrement dit, le calcul est le même que dans le cadre familier de la cohomologie complexe.
En particulier, si $\pi$ est le Frobenius géométrique, nous obtenons
$$
\text{Tr}(\pi_X^* | H^2 (\mathbf{P}^1_{\bar k},
\underline{\mathbf{Z}/\ell\mathbf{Z}})) = q
$$
et si $\pi$ est le Frobenius arithmétique, nous obtenons
$$
\text{Tr}(\text{frob}_k^* | H^2 (\mathbf{P}^1_{\bar k},
\underline{\mathbf{Z}/\ell\mathbf{Z}})) = q^{-1}.
$$
Cette dernière possibilité est manifestement incorrecte.

\begin{remark}
\label{trace-remark-compute-degree-lifting}
Le calcul des degrés peut s'effectuer en relevant (en un sens évident)
en caractéristique 0 et en considérant la situation à coefficients complexes.
Cette méthode ne fonctionne presque jamais, puisque le relèvement est en général impossible pour les
schémas autres que l'espace projectif.
\end{remark}

\noindent
Reste à comprendre pourquoi nous devons considérer la cohomologie à support
propre. En fait, au vu de la dualité de Poincar\'e, elle n'est pas absolument
nécessaire pour les variétés lisses, mais cela impose d'introduire certaines puissances
de $q$. Par exemple, considérons le cas où
$X = \mathbf{A}^1_k$ et
$\mathcal{F} = \underline{\mathbf{Z}/\ell\mathbf{Z}}$.
L'action sur les fibres est encore triviale ; nous n'avons donc qu'à examiner l'action
sur la cohomologie. Mais alors $\pi_X^*$ agit comme l'identité sur
$H^0(\mathbf{A}^1_{\bar k}, \underline{\mathbf{Z}/\ell\mathbf{Z}})$
et comme la multiplication par $q$ sur
$H^2_c(\mathbf{A}^1_{\bar k}, \underline{\mathbf{Z}/\ell\mathbf{Z}})$.




\section{Traces}
\label{trace-section-traces}

\noindent
Nous expliquons maintenant comment prendre la trace d'un endomorphisme d'un module sur un
anneau non commutatif. Fixons un anneau fini $\Lambda$ dont le cardinal est premier à $p$.
En général, $\Lambda$ est l'anneau de groupe $(\mathbf{Z}/\ell^n\mathbf{Z})[G]$ d'un
groupe fini $G$. Par convention, tous les $\Lambda$-modules considérés
seront des $\Lambda$-modules à gauche.

\medskip\noindent
Introduisons la notation suivante :
nous notons $\Lambda^\natural$ le quotient de $\Lambda$ par le sous-groupe additif
engendré par les commutateurs (c'est-à-dire les éléments de la forme
$ab-ba$, $a, b \in \Lambda$). Notons que $\Lambda^\natural$ n'est pas un anneau.

\medskip\noindent
Par exemple, le module $(\mathbf{Z}/\ell^n\mathbf{Z})[G]^\natural$ est le
dual de l'espace des fonctions centrales, de sorte que
$$
(\mathbf{Z}/\ell^n\mathbf{Z})[G]^\natural
=
\bigoplus\nolimits_{\text{classes de conjugaison dans }G}
\mathbf{Z}/\ell^n\mathbf{Z}.
$$
Pour un $\Lambda$-module libre, on a $\text{End}_\Lambda(\Lambda^{\oplus m}) =
\text{Mat}_m(\Lambda)$. Puisque les modules sont à gauche,
la représentation des endomorphismes par des matrices est une action à droite : si $a \in
\text{End}(\Lambda^{\oplus m})$ a pour matrice $A$ et si $v \in \Lambda$, alors $a(v)
= v A$.

\begin{definition}
\label{trace-definition-trace}
La {\it trace} de l'endomorphisme $a$ est la somme des termes diagonaux d'une
matrice qui le représente. On obtient ainsi une application additive $\text{Tr} :
\text{End}_\Lambda(\Lambda^{\oplus m}) \to \Lambda^\natural$.
\end{definition}

\begin{exercise}
\label{trace-exercise-trace-is-trace}
Étant données des applications
$$
\Lambda^{\oplus m} \xrightarrow{a} \Lambda^{\oplus n}
\quad\text{et}\quad
\Lambda^{\oplus n} \xrightarrow{b} \Lambda^{\oplus m}
$$
montrer que $\text{Tr}(ab) = \text{Tr}(ba)$.
\end{exercise}

\noindent
Étendons la définition de la trace à un $\Lambda$-module projectif de type fini
$P$ et à un endomorphisme $\varphi$ de $P$. Écrivons $P$ comme facteur direct
d'un $\Lambda$-module libre, c'est-à-dire considérons des applications $P \xrightarrow{a}
\Lambda^{\oplus n} \xrightarrow{b} P$ telles que
\begin{enumerate}
\item
$\Lambda^{\oplus n} = \Im(a) \oplus \Ker(b)$ ; et
\item
$b\circ a = \text{id}_P$.
\end{enumerate}
Nous posons alors $\text{Tr}(\varphi) = \text{Tr}(a\varphi b)$. Il est facile de vérifier
que cette définition ne dépend pas des choix, au moyen de l'exercice précédent.








\section{Pourquoi les catégories dérivées ?}
\label{trace-section-derived-categories-why}

\noindent
Avec cette définition de la trace, examinons un autre problème que pose la
formule telle qu'elle est énoncée. Soit $C$ une courbe projective lisse sur $k$. Il existe
une correspondance entre les faisceaux finis localement constants $\mathcal{F}$ sur
$C_\etale$ dont les fibres sont isomorphes à
${(\mathbf{Z}/\ell^n\mathbf{Z})}^{\oplus m}$ d'une part, et les représentations continues
$\rho : \pi_1 (C, \bar c) \to
\text{GL}_m(\mathbf{Z}/\ell^n\mathbf{Z}))$ (pour un choix fixé de $\bar c$)
d'autre part. Nous notons $\mathcal{F}_\rho$ le faisceau correspondant à
$\rho$. Alors $H^2 (C_{\bar k}, \mathcal{F}_\rho)$ est le groupe des coinvariants
pour l'action de $\rho(\pi_1 (C, \bar c))$ sur
${(\mathbf{Z}/\ell^n\mathbf{Z})}^{\oplus m}$, et l'on a la suite exacte courte
suivante
$$
0 \longrightarrow \pi_1 (C_{\bar k}, \bar c) \longrightarrow \pi_1 (C, \bar c)
\longrightarrow G_k \longrightarrow 0.
$$
Par exemple, faisons agir $\mathbf{Z} = \mathbf{Z} \sigma$ sur
$\mathbf{Z}/\ell^2\mathbf{Z}$ par $\sigma(x) = (1+\ell) x$. Les coinvariants
sont $(\mathbf{Z}/\ell^2\mathbf{Z})_{\sigma} = \mathbf{Z}/\ell\mathbf{Z}$, qui
n'est pas un $\mathbf{Z}/\ell^2\mathbf{Z}$-module plat. Nous ne pouvons donc pas prendre la
trace d'une action sur $H^2(C_{\bar k}, \mathcal{F}_\rho)$, du moins pas au
sens de la section précédente.

\medskip\noindent
En fait, notre objectif est d'établir une formule des traces à coefficients $\ell$-adiques.
Mais $\mathbf{Q}_\ell = \mathbf{Z}_\ell[1/\ell]$ et $\mathbf{Z}_\ell =
\lim \mathbf{Z}/\ell^n\mathbf{Z}$, et même pour un faisceau plat de
$\mathbf{Z}/\ell^n\mathbf{Z}$-modules, les groupes de cohomologie individuels peuvent ne pas
être plats, de sorte que nous ne pouvons calculer les traces. Une solution possible consiste à considérer le complexe
dérivé total $R\Gamma(C_{\bar k}, \mathcal{F}_\rho)$ dans la catégorie dérivée
$D(\mathbf{Z}/\ell^n\mathbf{Z})$ et à montrer qu'il s'agit d'un objet parfait, c'est-à-dire
qu'il est quasi-isomorphe à un complexe fini de modules libres de type fini.
Pour de tels complexes, on peut définir la trace, mais cela exigera un exposé
des catégories dérivées.






\section{Catégories dérivées}
\label{trace-section-derived-categories}

\noindent
Pour fixer les notations, soit $\mathcal{A}$ une catégorie abélienne. Soit
$\text{Comp}(\mathcal{A})$ la catégorie abélienne des complexes de
$\mathcal{A}$. Soit $K(\mathcal{A})$ la catégorie des complexes à
homotopie près, dont les objets sont les complexes de $\mathcal{A}$ et les morphismes
des classes
d'homotopie de morphismes de complexes. Ce n'est pas une catégorie abélienne.
Grossièrement, $D(A)$ est définie comme la catégorie obtenue en inversant
tous les quasi-isomorphismes de $\text{Comp}(\mathcal{A})$ ou, de manière équivalente, de
$K(\mathcal{A})$. On peut en outre définir $\text{Comp}^+(\mathcal{A}),
K^+(\mathcal{A}), D^+(\mathcal{A})$ de manière analogue en ne considérant que des complexes bornés inférieurement.
De même, on peut définir $\text{Comp}^-(\mathcal{A}),
K^-(\mathcal{A}), D^-(\mathcal{A})$ à l'aide des complexes bornés supérieurement, et l'on peut
définir $\text{Comp}^b(\mathcal{A}), K^b(\mathcal{A}), D^b(\mathcal{A})$ à l'aide des
complexes bornés.

\begin{remark}
\label{trace-remarks-derived-categories}
Quelques remarques sur les catégories dérivées.
\begin{enumerate}
\item
Des problèmes ensemblistes se posent lorsque $\mathcal{A}$ est assez
arbitraire ; nous les ignorerons allègrement.
\item
Les catégories $K(A)$ et $D(A)$ sont munies d'une structure de
catégorie triangulée.
\item
Les catégories $\text{Comp}(\mathcal{A})$ et $K(\mathcal{A})$ peuvent aussi être
définies lorsque $\mathcal{A}$ est une catégorie additive.
\end{enumerate}
\end{remark}

\noindent
Le foncteur d'homologie $H^i : \text{Comp}(\mathcal{A}) \to \mathcal{A}$ défini par
$K^\bullet \mapsto H^i(K^\bullet)$ se prolonge en des foncteurs $H^i :
K(\mathcal{A}) \to \mathcal{A}$ et $H^i : D(\mathcal{A}) \to \mathcal{A}$.

\begin{lemma}
\label{trace-lemma-when-in-bounded}
Un objet $E$ de $D(\mathcal{A})$ appartient à $D^+(\mathcal{A})$ si et seulement
si $H^i(E) =0 $ pour tout $i \ll 0$. Des énoncés analogues valent pour $D^-$ et
$D^b$.
\end{lemma}

\begin{proof}
Indication : utiliser les foncteurs de troncation. Voir
Catégories dérivées, lemme \ref{derived-lemma-complex-cohomology-bounded}.
\end{proof}

\begin{lemma}
\label{trace-lemma-derived-categories}
Morphismes entre objets de la catégorie dérivée.
\begin{enumerate}
\item
Soit $I^\bullet \in \text{Comp}^+(\mathcal{A})$ tel que $I^n$ soit injectif pour tout
$n \in \mathbf{Z}$. Alors
$$
\Hom_{D(\mathcal{A})}(K^\bullet, I^\bullet)
=
\Hom_{K(\mathcal{A})}(K^\bullet, I^\bullet).
$$
\item
Soit $P^\bullet \in \text{Comp}^-(\mathcal{A})$ tel que $P^n$ soit projectif pour tout
$n \in \mathbf{Z}$. Alors
$$
\Hom_{D(\mathcal{A})}(P^\bullet, K^\bullet)
=
\Hom_{K(\mathcal{A})}(P^\bullet, K^\bullet).
$$
\item
Si $\mathcal{A}$ possède assez d'injectifs et si $\mathcal{I} \subset \mathcal{A}$
est la sous-catégorie additive des objets injectifs, alors
$
D^+(\mathcal{A})\cong K^+(\mathcal{I})
$
(en tant que catégories triangulées).
\item
Si $\mathcal{A}$ possède assez de projectifs et si $\mathcal{P} \subset \mathcal{A}$
est la sous-catégorie additive des objets projectifs, alors
$
D^-(\mathcal{A}) \cong K^-(\mathcal{P}).
$
\end{enumerate}
\end{lemma}

\begin{proof}
Démonstration omise.
\end{proof}

\begin{definition}
\label{trace-definition-derived-functor}
Soit $F: \mathcal{A} \to \mathcal{B}$ un foncteur exact à gauche et supposons que
$\mathcal{A}$ possède assez d'injectifs. Nous définissons le {\it foncteur dérivé total
à droite de $F$} comme le foncteur $RF: D^+(\mathcal{A}) \to D^+(\mathcal{B})$
s'insérant dans le diagramme
$$
\xymatrix{
D^+(\mathcal{A}) \ar[r]^{RF} & D^+(\mathcal{B}) \\
K^+(\mathcal I) \ar[u] \ar[r]^F & K^+(\mathcal{B}). \ar[u]
}
$$
Cela est possible puisque la flèche verticale de gauche est inversible d'après le lemme
précédent. De même, soit $G: \mathcal{A} \to \mathcal{B}$ un foncteur exact à droite
et supposons que $\mathcal{A}$ possède assez de projectifs. Nous définissons le
{\it foncteur dérivé total à gauche de $G$} comme le foncteur $LG: D^-(\mathcal{A})
\to D^-(\mathcal{B})$ s'insérant dans le diagramme
$$
\xymatrix{
D^-(\mathcal{A}) \ar[r]^{LG} & D^-(\mathcal{B}) \\
K^-(\mathcal{P}) \ar[u] \ar[r]^G & K^-(\mathcal{B}). \ar[u]
}
$$
Cela est possible puisque la flèche verticale de gauche est inversible d'après le lemme
précédent.
\end{definition}

\begin{remark}
\label{trace-remark-cohomology-of-derived-functor}
Dans ces cas, on a $R^iF(K^\bullet) = H^i(RF(K^\bullet))$, où
le membre de gauche est défini comme la $i$-ième homologie du complexe
$F(K^\bullet)$.
\end{remark}




\section{Catégorie dérivée filtrée}
\label{trace-section-filtered-derived-category}

\noindent
Il s'avère que nous devons tout recommencer et construire également la catégorie
dérivée filtrée.

\begin{definition}
\label{trace-definition-filtered}
Soit $\mathcal{A}$ une catégorie abélienne.
\begin{enumerate}
\item Soit $\text{Fil}(\mathcal{A})$ la catégorie des objets filtrés
$(A, F)$ de $\mathcal{A}$, où $F$ est une filtration de la forme
$$
A \supset \ldots \supset F^n A \supset F^{n+1}A \supset \ldots
\supset 0.
$$
Il s'agit d'une catégorie additive.
\item Nous notons $\text{Fil}^f(\mathcal{A})$ la sous-catégorie pleine
de $\text{Fil}(\mathcal{A})$ dont les objets $(A, F)$ ont une filtration
finie. C'est encore une catégorie additive.
\item Un objet $I \in \text{Fil}^f(\mathcal{A})$ est dit
{\it injectif filtré} (respectivement {\it projectif filtré}) si
$\text{gr}^p(I) = \text{gr}_F^p(I) = F^pI/F^{p+1}I$ est injectif
(respectivement projectif) dans $\mathcal{A}$ pour tout $p$.
\item La catégorie des complexes
$\text{Comp}(\text{Fil}^f(\mathcal{A})) \supset
\text{Comp}^+(\text{Fil}^f(\mathcal{A}))$
et sa catégorie homotopique
$K(\text{Fil}^f(\mathcal{A})) \supset K^+(\text{Fil}^f(\mathcal A))$
sont définies comme précédemment.
\item Un morphisme $\alpha : K^\bullet \to L^\bullet$ de complexes de
$\text{Comp}(\text{Fil}^f(\mathcal{A}))$ est appelé un
{\it quasi-isomorphisme filtré} si
$$
\text{gr}^p(\alpha): \text{gr}^p(K^\bullet) \to \text{gr}^p(L^\bullet)
$$
est un quasi-isomorphisme pour tout $p \in \mathbf{Z}$.
\item Nous définissons $DF(\mathcal{A})$ (respectivement $DF^+(\mathcal{A})$)
en inversant les quasi-isomorphismes filtrés dans
$K(\text{Fil}^f(\mathcal{A}))$ (respectivement $K^+(\text{Fil}^f(\mathcal{A}))$).
\end{enumerate}
\end{definition}

\begin{lemma}
\label{trace-lemma-filtered-derived-category}
Si $\mathcal{A}$ possède assez d'injectifs, alors $DF^+(\mathcal{A}) \cong
K^+(\mathcal{I})$, où $\mathcal{I}$ est la sous-catégorie additive pleine de
$\text{Fil}^f(\mathcal{A})$ formée des objets injectifs filtrés.
De même, si $\mathcal{A}$ possède assez de projectifs, alors $DF^-(\mathcal{A})
\cong K^-(\mathcal{P})$, où $\mathcal P$ est la sous-catégorie additive pleine de
$\text{Fil}^f(\mathcal{A})$ formée des objets projectifs filtrés.
\end{lemma}

\begin{proof}
Démonstration omise.
\end{proof}





\section{Foncteurs dérivés filtrés}
\label{trace-section-filtered-derived-functors}

\noindent
Il y a ensuite les foncteurs dérivés filtrés.

\begin{definition}
\label{trace-definition-filtered-derived-functors}
Soit $T: \mathcal{A} \to \mathcal{B}$ un foncteur exact à gauche et supposons que
$\mathcal{A}$ possède assez d'injectifs. Définissons $RT: DF^+(\mathcal{A}) \to D
F^+(\mathcal{B})$ de sorte qu'il s'insère dans le diagramme
$$
\xymatrix{
DF^+(\mathcal{A}) \ar[r]^{RT} & DF^+(\mathcal{B}) \\
K^+(\mathcal{I}) \ar[u] \ar[r]^{T \quad} & K^+(\text{Fil}^f(\mathcal{B})).
\ar[u]}
$$
Ce foncteur est bien défini d'après le lemme précédent. Soit $G: \mathcal{A} \to
\mathcal{B}$ un foncteur exact à droite et supposons que $\mathcal{A}$ possède assez de
projectifs. Définissons $LG: DF^-(\mathcal{A}) \to DF^-(\mathcal{B})$ de sorte qu'il s'insère dans
le diagramme
$$
\xymatrix{
DF^-(\mathcal{A}) \ar[r]^{LG} & DF^-(\mathcal{B}) \\
K^-(\mathcal{P}) \ar[u] \ar[r]^{G \quad} & K^-(\text{Fil}^f(\mathcal{B})).
\ar[u]}
$$
Là encore, ce foncteur est bien défini d'après le lemme précédent.
Les foncteurs $RT$, respectivement\ $LG$, sont appelés le {\it foncteur dérivé
filtré} de $T$, respectivement de $G$.
\end{definition}

\begin{proposition}
\label{trace-proposition-compare-filtered-graded}
Dans la situation ci-dessus, on a
$$
\text{gr}^p \circ RT = RT \circ \text{gr}^p
$$
où le $RT$ de gauche est le foncteur dérivé filtré, tandis que celui de
droite est le foncteur dérivé total. Autrement dit, on a un diagramme commutatif
$$
\xymatrix{
DF^+(\mathcal{A}) \ar[r]^{RT} \ar[d]_{\text{gr}^p} & DF^+(\mathcal{B})
\ar[d]^{\text{gr}^p}\\
D^+(\mathcal{A}) \ar[r]^{RT} & D^+(\mathcal{B}).}
$$
\end{proposition}

\begin{proof}
Démonstration omise.
\end{proof}

\noindent
Pour $K^\bullet \in DF^+(\mathcal{B})$, on obtient une suite spectrale
$$
E_1^{p, q} = H^{p+q}(\text{gr}^p K^\bullet) \Rightarrow H^{p+q}(\text{oubli de la
filtration}(K^\bullet)).
$$






\section{Application des complexes filtrés}
\label{trace-section-applications-filtered}

\noindent
Soit $\mathcal{A}$ une catégorie abélienne possédant assez d'injectifs, et soit
$0 \to L \to M \to N \to 0$ une suite exacte courte dans $\mathcal{A}$.
Considérons $\widetilde M \in \text{Fil}^f(\mathcal{A})$ obtenu en munissant $M$ de la
filtration définie par
$$
F^1M = L, \ F^nM = M
\text{ pour }n \leq 0\text{, et }F^nM = 0\text{ pour }n \geq 2.
$$
Par définition, on a
$$
\text{oubli de la filtration}(\widetilde M) = M, \quad
\text{gr}^0(\widetilde M) = N, \quad
\text{gr}^1(\widetilde M) = L
$$
et $\text{gr}^n(\widetilde M) = 0$ pour tout autre $n \neq 0, 1$. Soit $T:
\mathcal{A} \to \mathcal{B}$ un foncteur exact à gauche. Supposons que $\mathcal{A}$
possède assez d'injectifs. Alors $RT(\widetilde M) \in DF^+(\mathcal{B})$ est un
complexe filtré vérifiant
$$
\text{gr}^p(RT(\widetilde M))
\stackrel{\text{qis}}{=}
\left\{
\begin{matrix}
0 & \text{si} & p \neq 0, 1, \\
RT(N) & \text{si} & p = 0, \\
RT(L) & \text{si} & p = 1.
\end{matrix}
\right.
$$
et $\text{oubli de la filtration}(RT(\widetilde M))\stackrel{\text{qis}}{ = } RT(M)$. La
suite spectrale appliquée à $RT(\widetilde M)$ donne
$$
E_1^{p, q} = R^{p+q}T(\text{gr}^p(\widetilde M)) \Rightarrow
R^{p+q}T(\text{oubli de la filtration}(\widetilde M)).
$$
En développant la suite spectrale, on obtient la suite exacte longue
$$
\xymatrix{
0 \ar[r] & T(L) \ar[r] & T(M) \ar[r] & T(N) \ar@(rd, ul)[rdllllr] \\
& R^1T(L) \ar[r] & R^1T(M) \ar[r] & \ldots
}
$$
Nous l'utiliserons de la manière suivante. Soit $X/k$ un schéma de type fini. Soit
$\mathcal{F}$ un $\mathbf{Z}/\ell^n \mathbf{Z}$-module constructible plat.
Nous voulons alors montrer que la trace
$$
\text{Tr}( \pi_X^\ast | R\Gamma_c(X_{\bar k}, \mathcal{F})) \in
\mathbf{Z}/\ell^n \mathbf{Z}
$$
est additive sur les suites exactes courtes. Pour le voir, il ne suffira pas de
travailler avec $R\Gamma_c(X_{\bar k}, -) \in D^+(\mathbf{Z}/\ell^n \mathbf{Z})$, mais
il faudra utiliser la catégorie dérivée filtrée.







%10.29.09
\section{Complexes parfaits}
\label{trace-section-perfect}

\noindent
Soit $\Lambda$ un anneau (éventuellement non commutatif), $\text{Mod}_{\Lambda}$ la
catégorie des $\Lambda$-modules à gauche, $K(\Lambda) = K(\text{Mod}_\Lambda)$ sa
catégorie homotopique, et $D(\Lambda)= D(\text{Mod}_\Lambda)$ la catégorie dérivée
correspondante.

\begin{definition}
\label{trace-definition-perfect}
Nous notons $K_{perf}(\Lambda)$ la catégorie dont les objets sont les complexes bornés
de $\Lambda$-modules projectifs de type fini, et dont les morphismes sont les
morphismes de complexes à homotopie près. Le foncteur $K_{perf}(\Lambda)\to
D(\Lambda)$ est pleinement fidèle (Catégories dérivées, lemme
\ref{derived-lemma-morphisms-from-projective-complex}).
Notons $D_{perf}(\Lambda)$ son image essentielle.
Un objet de $D(\Lambda)$ est dit {\it parfait} s'il appartient à
$D_{perf}(\Lambda)$.
\end{definition}

\begin{proposition}
\label{trace-proposition-trace-well-defined}
Soient $K\in D_{perf}(\Lambda)$ et $f\in \text{End}_{D(\Lambda)}(K)$. Alors la
trace $\text{Tr}(f)\in \Lambda^\natural$ est bien définie.
\end{proposition}

\begin{proof}
Nous utiliserons le lemme de Catégories dérivées
\ref{derived-lemma-morphisms-from-projective-complex}
sans plus le mentionner dans cette démonstration.
Soit $P^\bullet$ un complexe borné de $\Lambda$-modules projectifs de type fini
et soit $\alpha : P^\bullet \to K$ un isomorphisme dans $D(\Lambda)$. Alors
$\alpha^{-1}\circ f\circ \alpha$ correspond à un morphisme de complexes
$f^\bullet : P^\bullet \to P^\bullet$ bien défini à homotopie près.
Posons
$$
\text{Tr}(f) = \sum_i (-1)^i \text{Tr}(f^i : P^i \to P^i) \in \Lambda^\natural.
$$
Une fois $P^\bullet$ et $\alpha$ fixés, cette définition ne dépend pas du choix de
$f^\bullet$. En effet, tout autre choix est de la forme
$\tilde{f}^\bullet = f^\bullet + dh +hd$ pour un certain
$h^i : P^i \to P^{i-1}(i\in \mathbf{Z})$. Mais
\begin{eqnarray*}
\text{Tr}(dh) & = & \sum_i (-1)^i \text{Tr}(P^i\xrightarrow{dh} P^i) \\
& = & \sum_i (-1)^i \text{Tr}(P^{i-1}\xrightarrow{hd} P^{i-1}) \\
& = & -\sum_i (-1)^{i-1}\text{Tr}(P^{i-1}\xrightarrow{hd} P^{i-1}) \\
& = & - \text{Tr}(hd)
\end{eqnarray*}
donc $\sum_i (-1)^i \text{Tr} ((dh+hd)|_{P^i})=0$.
En outre, cette définition ne dépend pas du choix de $(P^\bullet , \alpha)$ :
supposons que $(Q^\bullet, \beta)$ soit un autre choix. Les composées
$$
Q^\bullet \xrightarrow{\beta} K \xrightarrow{\alpha^{-1}} P^\bullet
\quad\text{et}\quad
P^\bullet \xrightarrow{\alpha} K \xrightarrow{\beta^{-1}} Q^\bullet
$$
sont représentables par des morphismes de complexes $\gamma_1^\bullet$ et
$\gamma_2^\bullet$ respectivement, tels que $\gamma_1^\bullet \circ
\gamma_2^\bullet$ soit homotope à l'identité. Ainsi, le morphisme de complexes
$\gamma_2^\bullet\circ f^\bullet\circ \gamma_1^\bullet : Q^\bullet\to Q^\bullet$
représente le morphisme $\beta^{-1}\circ f\circ\beta$ dans $D(\Lambda)$. Or
\begin{eqnarray*}
\text{Tr}(\gamma_2^\bullet\circ f^\bullet\circ\gamma_1^\bullet|_{Q^\bullet}) &
= & \text{Tr}(\gamma_1^\bullet \circ\gamma_2^\bullet \circ
f^\bullet|_{P^\bullet})\\
& = & \text{Tr}(f^\bullet|_{P^\bullet})
\end{eqnarray*}
car $\gamma_1^\bullet \circ \gamma_2^\bullet$ est homotope à
l'identité et d'après l'indépendance du choix de $f^\bullet$ établie ci-dessus.
\end{proof}




\section{Filtrations et complexes parfaits}
\label{trace-section-filtrations-perfect}

\noindent
Présentons maintenant une version filtrée de la catégorie des complexes parfaits. Un
objet $(M, F)$ de $\text{Fil}^f(\text{Mod}_\Lambda)$ est dit {\it projectif filtré
de type fini} si, pour tout $p$, $\text{gr}^p_F (M)$ est de type fini et
projectif. Nous considérons alors la catégorie homotopique
$KF_{\text{perf}}(\Lambda)$ des complexes bornés d'objets projectifs filtrés
de type fini de $\text{Fil}^f(\text{Mod}_\Lambda)$. On a un diagramme de
catégories
$$
\begin{matrix}
KF(\Lambda) & \supset & KF_{\text{perf}}(\Lambda)\\
\downarrow & & \downarrow\\
DF(\Lambda) & \supset & DF_{\text{perf}}(\Lambda)
\end{matrix}
$$
où le foncteur vertical de droite est pleinement fidèle et où la catégorie
$DF_{\text{perf}}(\Lambda)$ est son image essentielle, comme précédemment.

\begin{lemma}[Additivité]
\label{trace-lemma-additivity}
Soient $K\in DF_{\text{perf}}(\Lambda)$ et $f\in
\text{End}_{DF}(K)$. Alors
$$
\text{Tr}(f|_K) =
\sum\nolimits_{p\in \mathbf{Z}} \text{Tr}(f|_{\text{gr}^p K}).
$$
\end{lemma}

\begin{proof}
D'après la proposition \ref{trace-proposition-trace-well-defined}, on peut supposer disposer
d'un complexe borné
$P^\bullet$ d'objets projectifs filtrés de type fini de
$\text{Fil}^f(\text{Mod}_\Lambda)$ et d'un morphisme $f^\bullet : P^\bullet\to
P^\bullet$ dans $\text{Comp}(\text{Fil}^f(\text{Mod}_\Lambda))$. Le lemme
découle donc du résultat suivant, dont la démonstration est laissée au lecteur.
\end{proof}

\begin{lemma}
\label{trace-lemma-additive-filtered-finite-projective}
Soit $P \in \text{Fil}^f(\text{Mod}_\Lambda)$ un objet projectif filtré de type fini, et
soit $f : P \to P$ un endomorphisme dans $\text{Fil}^f(\text{Mod}_\Lambda)$. Alors
$$
\text{Tr}(f|_P) =
\sum\nolimits_p \text{Tr}(f|_{\text{gr}^p(P)}).
$$
\end{lemma}

\begin{proof}
Démonstration omise.
\end{proof}







\section{Caractérisation des objets parfaits}
\label{trace-section-characterizing-perfect}

\noindent
Pour le cas commutatif, voir
Compléments sur l'algèbre, sections
\ref{more-algebra-section-pseudo-coherent},
\ref{more-algebra-section-tor}, et
\ref{more-algebra-section-perfect}.

\begin{definition}
\label{trace-definition-finite-tor-dimension}
Soit $\Lambda$ un anneau (éventuellement non commutatif).
Un objet $K\in D(\Lambda)$ est de {\it dimension de $\text{Tor}$ finie}
s'il existe $a, b \in \mathbf{Z}$ tels que, pour tout
$\Lambda$-module à droite $N$, on ait
$H^i(N \otimes_{\Lambda}^\mathbf{L} K) = 0$ pour tout
$i \not \in [a, b]$.
\end{definition}

\noindent
Cela signifie en particulier que $K \in D^b(\Lambda)$, comme on le voit en prenant
$N = \Lambda$.

\begin{lemma}
\label{trace-lemma-characterize-perfect}
Soient $\Lambda$ un anneau noethérien à gauche et $K\in D(\Lambda)$. Alors $K$ est
parfait si et seulement si les deux conditions suivantes sont satisfaites :
\begin{enumerate}
\item
$K$ est de dimension de $\text{Tor}$ finie, et
\item
pour tout $i \in \mathbf{Z}$, $H^i(K)$ est un $\Lambda$-module de type fini.
\end{enumerate}
\end{lemma}

\begin{proof}
Voir Compléments sur l'algèbre, lemme \ref{more-algebra-lemma-perfect},
pour la démonstration dans le cas commutatif.
\end{proof}

\noindent
Le lecteur est vivement encouragé à essayer de démontrer cet énoncé. La démonstration repose sur le
fait qu'un module à gauche de présentation finie sur un anneau est plat si et seulement
s'il est projectif.

\begin{remark}
\label{trace-remark-variant}
Une variante de ce lemme consiste à considérer un schéma noethérien $X$
et la catégorie $D_{perf}(\mathcal{O}_X)$ des complexes qui sont localement
quasi-isomorphes à un complexe fini de modules localement libres de type fini sur
$\mathcal{O}_X$. Les objets $K$ de $D_{perf}(\mathcal{O}_X)$
se caractérisent par la cohérence de leurs faisceaux de cohomologie et par une
dimension de Tor bornée.
\end{remark}








\section{Cohomologie des complexes raisonnables}
\label{trace-section-cohomology-ctf}

\noindent
Ce qui suit est un cas particulier d'un résultat plus général concernant la
cohomologie à support propre des objets de $D_{ctf}(X, \Lambda)$.

\begin{proposition}
\label{trace-proposition-projective-curve-constructible-cohomology}
Soient $X$ une courbe projective sur un corps $k$, $\Lambda$ un anneau fini et
$K\in D_{ctf}(X, \Lambda)$. Alors $R\Gamma(X_{\bar k}, K)\in
D_{perf}(\Lambda)$.
\end{proposition}

\begin{proof}[Esquisse de démonstration.]
La première étape consiste à montrer :
\begin{enumerate}
\item[(1)]
{\it La cohomologie de $R\Gamma(X_{\bar k}, K)$ est bornée.}
\end{enumerate}
Considérons la suite spectrale
$$
H^i(X_{\bar k}, \underline H^j(K))
\Rightarrow
H^{i+j} (R\Gamma(X_{\bar k}, K)).
$$
Comme $K$ est borné et que $\Lambda$ est fini, les faisceaux $\underline H^j(K)$
sont de torsion. De plus, $X_{\bar k}$ est de dimension cohomologique finie, donc le
membre de gauche n'est non nul que pour un nombre fini de $i$ et de $j$. Il en va donc
de même du membre de droite.
\begin{enumerate}
\item[(2)]
{\it Les groupes de cohomologie $H^{i+j} (R\Gamma(X_{\bar k}, K))$ sont finis.}
\end{enumerate}
Comme les faisceaux $\underline H^j(K)$ sont constructibles, les groupes
$H^i(X_{\bar k}, \underline H^j(K))$ sont finis
(Cohomologie \'etale, section \ref{etale-cohomology-section-vanishing-torsion})
et l'assertion découle de nouveau de la suite spectrale.
\begin{enumerate}
\item[(3)]
{\it $R\Gamma(X_{\bar k}, K)$ est de dimension de $\text{Tor}$ finie.}
\end{enumerate}
Soit $N$ un $\Lambda$-module à droite (en fait, comme $\Lambda$ est fini, il
suffit de supposer que $N$ est fini). Par la formule de projection (changement de
module),
$$
N \otimes^\mathbf{L}_\Lambda R \Gamma(X_{\bar k}, K) = R\Gamma(X_{\bar k},
N \otimes^\mathbf{L}_\Lambda K).
$$
Par conséquent,
$$
H^i (N \otimes^\mathbf{L}_\Lambda R\Gamma(X_{\bar k}, K)) = H^i(R\Gamma(X_{\bar
k}, N \otimes_{\Lambda}^\mathbf{L} K)).
$$
Considérons maintenant la suite spectrale
$$
H^i (X_{\bar k}, \underline H^j (N \otimes_{\Lambda}^\mathbf{L} K))
\Rightarrow
H^{i+j}(R\Gamma(X_{\bar k}, N \otimes_{\Lambda}^\mathbf{L} K)).
$$
Comme $K$ est de dimension de $\text{Tor}$ finie, $\underline H^j
(N \otimes_{\Lambda}^\mathbf{L} K)$ s'annule universellement pour $j$ assez petit,
et le membre de gauche s'annule dès que $i < 0$. Par conséquent, $R\Gamma(X_{\bar
k}, K)$ est de dimension de $\text{Tor}$ finie, comme annoncé. C'est donc un
complexe parfait d'après le lemme \ref{trace-lemma-characterize-perfect}.
\end{proof}





\section{Nombres de Lefschetz}
\label{trace-section-lefschetz-numbers}

\noindent
Le fait que la cohomologie totale d'un complexe constructible de dimension de Tor
finie soit un complexe parfait est la raison technique essentielle pour laquelle la cohomologie
se comporte bien et permet de définir rigoureusement les traces qui apparaissent dans la
formule des traces.

\begin{definition}
\label{trace-definition-global-lefschetz-number}
Soient $\Lambda$ un anneau fini, $X$ une courbe projective sur un corps fini $k$
et $K \in D_{ctf}(X, \Lambda)$ (par exemple $K = \underline\Lambda$).
Il existe un morphisme canonique $c_K : \pi_X^{-1}K \to K$, et son changement de base
$c_K|_{X_{\bar k}}$ induit une action, notée $\pi_X^*$, sur le complexe parfait
$R\Gamma(X_{\bar k}, K|_{X_{\bar k}})$. Le
{\it nombre global de Lefschetz} de $K$ est la trace
$\text{Tr}(\pi_X^* |_{R\Gamma(X_{\bar k}, K)})$ de cette action.
C'est un élément de $\Lambda^\natural$.
\end{definition}

\begin{definition}
\label{trace-definition-local-lefschetz-number}
Soient $\Lambda, X, k, K$ comme dans la
définition \ref{trace-definition-global-lefschetz-number}.
Comme $K\in D_{ctf}(X, \Lambda)$, pour tout point géométrique $\bar x$ de $X$,
le complexe $K_{\bar x}$ est un complexe parfait (dans $D_{perf}(\Lambda)$). Comme nous
l'avons vu à la section \ref{trace-section-frobenii}, le Frobenius $\pi_X$ agit sur
$K_{\bar x}$. Le {\it nombre local de Lefschetz} de $K$ est la somme
$$
\sum\nolimits_{x\in X(k)} \text{Tr}(\pi_x |_{K_{\overline{x}}})
$$
qui est encore un élément de $\Lambda^\natural$.
\end{definition}

\noindent
Nous pouvons enfin formuler précisément la formule des traces.

\begin{theorem}[Formule des traces de Lefschetz]
\label{trace-theorem-trace}
Soient $X$ une courbe projective sur un corps fini $k$, $\Lambda$ un anneau fini
et $K \in D_{ctf}(X, \Lambda)$. Alors les nombres de Lefschetz global et local
de $K$ sont égaux, c'est-à-dire
\begin{equation}
\label{trace-equation-trace-formula}
\text{Tr}(\pi^*_X |_{R\Gamma(X_{\bar k}, K)})
=
\sum\nolimits_{x\in X(k)} \text{Tr}(\pi_X |_{K_{\bar x}})
\end{equation}
dans $\Lambda^\natural$.
\end{theorem}

\begin{proof}
Voir la discussion ci-dessous.
\end{proof}

%11.5.09
\noindent
Nous utiliserons la formule des traces plutôt que de la démontrer. Néanmoins, nous
donnerons de nombreux détails de la démonstration du théorème telle qu'elle est donnée dans
\cite{SGA4.5} (certaines des choses qui n'y sont pas expliquées de manière satisfaisante
sont énumérées à la section \ref{trace-section-list-skipped}).

\medskip\noindent
Nous n'avons énoncé la formule que pour les courbes et, en un sens assez
faible, elle est une conséquence du résultat suivant.

\begin{theorem}[Weil]
\label{trace-theorem-weil-trace-formula}
Soient $C$ une courbe projective non singulière sur un corps algébriquement clos
$k$, et $\varphi : C \to C$ un $k$-endomorphisme de $C$ distinct de
l'identité. Posons $V(\varphi) = \Delta_C \cdot \Gamma_\varphi$, où $\Delta_C$ est
la diagonale, $\Gamma_\varphi$ est le graphe de $\varphi$, et où le nombre d'intersection
est calculé sur $C \times C$. Soit $J = \underline{\Picardfunctor}^0_{C/k}$
la jacobienne de $C$, et notons $\varphi^* : J \to J$ l'action induite par
$\varphi$ par image réciproque. Alors
$$
V(\varphi) = 1 - \text{Tr}_J(\varphi^*) + \deg \varphi.
$$
\end{theorem}

\begin{proof}
Le nombre $V(\varphi)$ est le nombre de points fixes de $\varphi$ comptés avec multiplicité ; il est égal
à
$$
V(\varphi) =
\sum\nolimits_{c \in |C| : \varphi(c) = c} m_{\text{Fix}(\varphi)} (c)
$$
où $m_{\text{Fix}(\varphi)} (c)$ est la multiplicité de $c$ comme point fixe
de $\varphi$, c'est-à-dire l'ordre d'annulation de l'image d'une uniformisante locale
par $\varphi - \text{id}_C$. On trouvera des démonstrations de ce théorème dans
\cite{Lang} et \cite{Weil}.
\end{proof}

\begin{example}
\label{trace-example-elliptic-curve}
Soit $C = E$ une courbe elliptique et soit $\varphi = [n]$ la multiplication par $n$.
Alors $\varphi^* = \varphi^t$ est la multiplication par $n$ sur la jacobienne, de sorte qu'elle
a pour trace $2n$ et pour degré $n^2$. D'autre part, le schéma des points fixes de
$\varphi$ est formé des points $p \in E$ tels que $n p = p$ ; c'est donc le schéma de
$(n-1)$-torsion, qui est de degré $(n-1)^2$. Le théorème s'écrit donc
$$
(n-1)^2 = 1 - 2n + n^2.
$$
\end{example}

\noindent
{\bf Jacobiennes.}
Nous examinons maintenant, sans démonstration, la correspondance entre une courbe et sa
jacobienne qui intervient dans la démonstration de Weil. Soit $C$ une courbe projective non singulière
sur un corps algébriquement clos $k$ et choisissons un point de base $c_0 \in
C(k)$. Notons $A^1(C \times C)$ (ou $\Pic(C \times C)$, ou
$\text{CaCl}(C \times C)$) le groupe abélien des diviseurs de codimension 1 de
$C \times C$. Alors
$$
A^1(C \times C) = \text{pr}_1^* (A^1(C)) \oplus \text{pr}_2^* (A^1(C)) \oplus R
$$
où
$$
R = \{ Z \in A^1(C \times C) \ |
\ Z|_{C \times \{c_0\}} \sim_\text{rat} 0
\text{ et }
Z|_{\{c_0\} \times C} \sim_\text{rat} 0 \}.
$$
Autrement dit,
$R$ est le sous-groupe des fibrés en droites dont les restrictions aux
deux sous-schémas indiqués sont triviales. Il existe alors un isomorphisme canonique de groupes abéliens $R
\cong \text{End}(J)$ qui envoie un diviseur $Z$ de $R$ sur l'endomorphisme
$$
\begin{matrix}
J & \to & J \\
\left[ \mathcal{O}_C(D) \right] & \mapsto & (\text{pr}_1 |_Z)_* (\text{pr}_2
|_Z)^* (D).
\end{matrix}
$$
La correspondance susmentionnée est la suivante. Notons $\sigma$
l'automorphisme de $C \times C$ qui permute les facteurs.
$$
\begin{matrix}
\hline & \\
\text{End}(J) & R \\
& \\
\hline & \\
\text{composition de }\alpha, \beta &
{\text{pr}_{13}}_* ({\text{pr}_{12}}^*(\alpha) \circ {\text{pr}_{23}}^*(\beta))
\\
& \\
\text{id}_J &
\Delta_C - \{c_0\} \times C - C \times \{c_0\} \\
& \\
\varphi^* &
\Gamma_\varphi - C \times \{\varphi(c_0)\}
- \sum_{\varphi(c) = c_0} \{c\} \times C \\
& \\
{
\begin{matrix}
\text{la forme trace} \\
\alpha, \beta \mapsto \text{Tr}(\alpha \beta)
\end{matrix}
}
&
\alpha, \beta \mapsto - \int_{C \times C} \alpha . \sigma^*\beta
\\
& \\
{
\begin{matrix}
\text{l'involution de Rosati} \\
\alpha \mapsto \alpha^\dagger
\end{matrix}
}
&
\alpha \mapsto \sigma^*\alpha
\\
& \\
{
\begin{matrix}
\text{positivité de Rosati} \\
\text{Tr}(\alpha\alpha^\dagger) > 0
\end{matrix}
}
&
{
\begin{matrix}
\text{théorème de l'indice de Hodge sur }C \times C \\
- \int_{C \times C} \alpha \sigma^*\alpha > 0.
\end{matrix}
}
\\
& \\
\hline 
\end{matrix}
$$
En fait, au vu de la formule de Künneth, le sous-groupe $R$ correspond aux classes de Hodge de type
$1, 1$ dans $H^1(C)\otimes H^1(C)$.

\medskip\noindent
{\bf Démonstration de Weil.} Cette correspondance permet de démontrer la formule
des traces. Nous avons
\begin{eqnarray*}
V(\varphi) & = & \int_{C \times C} \Gamma_\varphi.\Delta \\
& = & \int_{C \times C} \Gamma_\varphi. \left(\Delta_C - \{c_0\} \times C - C
\times \{c_0\}\right) + \int_{C \times C} \Gamma_\varphi. \left(\{c_0\} \times C
+ C \times \{c_0\}\right).
\end{eqnarray*}
D'une part,
$$
\int_{C \times C} \Gamma_\varphi. \left(\{c_0\} \times C + C \times
\{c_0\}\right)
=
1 + \deg \varphi
$$
et, d'autre part, puisque $R$ est l'orthogonal du diviseur ample
$\{c_0\} \times C + C \times \{c_0\}$,
\begin{eqnarray*}
& &
\int_{C \times C} \Gamma_\varphi. \left(\Delta_C - \{c_0\} \times C - C \times
\{c_0\}\right) \\
& = &
\int_{C \times C} \left(\Gamma_\varphi - C \times \{\varphi(c_0)\} -
\sum_{\varphi(c) = c_0} \{c\} \times C \right). \left(\Delta_C - \{c_0\} \times
C - C \times \{c_0\}\right) \\
& = & - \text{Tr}_J (\varphi^* \circ \text{id}_J).
\end{eqnarray*}
En récapitulant, nous obtenons
$$
V(\varphi) = 1 - \text{Tr}_J (\varphi^*) + \deg \varphi
$$
ce qui donne la formule des traces.

\begin{lemma}
\label{trace-lemma-weil-mod}
Considérons la situation du
théorème \ref{trace-theorem-weil-trace-formula}
et soit $\ell$ un nombre premier inversible dans $k$. Alors
$$
\sum\nolimits_{i = 0}^2
(-1)^i
\text{Tr}(\varphi^* |_{H^i (C, \underline{\mathbf{Z}/\ell^n \mathbf{Z}})})
=
V(\varphi) \mod \ell^n.
$$
\end{lemma}

\begin{proof}[Esquisse de démonstration]
Remarquons d'abord que l'énoncé a un sens, car les groupes $H^i(C,
\underline{\mathbf{Z}/\ell^n \mathbf{Z}})$ sont des modules libres sur $\mathbf{Z}/\ell^n
\mathbf{Z}$ pour tout $i$. La trace de $\varphi^*$ sur la
cohomologie de degré zéro vaut 1. Le choix d'une racine primitive $\ell^n$-ième de l'unité dans $k$
fournit un isomorphisme
$$
H^i(C, \underline{\mathbf{Z}/\ell^n \mathbf{Z}}) \cong H^i(C, \mu_{\ell^n})
$$
de manière compatible avec l'action induite par l'endomorphisme considéré. D'autre part,
$H^1(C, \mu_{\ell^n}) = J[\ell^n]$. Par conséquent,
\begin{eqnarray*}
\text{Tr}(\varphi^* |_{H^1 (C, \underline{\mathbf{Z}/\ell^n \mathbf{Z}})})) & =
& \text{Tr}_J (\varphi^*) \mod \ell^n \\
& = & \text{Tr}_{\mathbf{Z}/\ell^n \mathbf{Z}} (\varphi^* : J[\ell^n] \to
J[\ell^n]).
\end{eqnarray*}
En outre, $H^2(C, \mu_{\ell^n}) = \Pic(C)/\ell^n\Pic(C) \cong
\mathbf{Z}/\ell^n \mathbf{Z}$, où $\varphi^*$ agit par multiplication par $\deg
\varphi$. Par suite,
$$
\text{Tr} (\varphi^*|_{H^2 (C, \underline{\mathbf{Z}/\ell^n \mathbf{Z}})}) =
\deg \varphi.
$$
Ainsi,
$$
\sum_{i = 0}^2 (-1)^i
\text{Tr}(\varphi^* |_{H^i (C, \underline{\mathbf{Z}/\ell^n \mathbf{Z}})}) =
1 - \text{Tr}_J(\varphi^*) + \deg \varphi \mod \ell^n
$$
et le lemme découle du théorème \ref{trace-theorem-weil-trace-formula}.
\end{proof}

\noindent
Une autre manière de démontrer ce lemme consiste à montrer que
$$
X \mapsto H^* (X, \mathbf{Q}_\ell) =
\mathbf{Q}_\ell \otimes
\lim_n H^*(X, \mathbf{Z}/\ell^n\mathbf{Z})
$$
définit une théorie cohomologique de Weil sur les variétés projectives lisses sur $k$. Alors
la formule des traces
$$
V(\varphi) = \sum_{i = 0}^2 (-1)^i
\text{Tr}(\varphi^* |_{H^i(C, \mathbf{Q}_\ell)})
$$
est une conséquence formelle des axiomes (il s'agit d'un exercice d'algèbre linéaire, la
démonstration étant la même que dans le cas topologique).




%11.10.09
\section{Préliminaires et sorites}
\label{trace-section-preliminaries}

\noindent
Notations :
Fixons les notations de cette section. On note $A$ un anneau commutatif,
et $\Lambda$ un anneau (éventuellement non commutatif) muni d'un morphisme d'anneaux $A\to \Lambda$ dont
l'image est contenue dans le centre de $\Lambda$. Soient $G$ un groupe fini et $\Gamma$ une
{\it extension de monoïdes de noyau $G$ et de quotient $\mathbf{N}$}, c'est-à-dire qu'il existe une
suite exacte
$$
1\to G\to \tilde\Gamma\to \mathbf{Z}\to 1
$$
et $\Gamma$ est constitué des éléments de $\tilde\Gamma$ dont l'image est
positive ou nulle. Enfin, soit $P$ un $A[\Gamma]$-module qui est de type fini et
projectif en tant que $A[G]$-module, et soit $M$ un $\Lambda[\Gamma]$-module qui est
de type fini et projectif en tant que $\Lambda$-module.

\medskip\noindent
Notre but est de calculer la trace de $1 \in \mathbf{N}$ agissant sur le $\Lambda$-module
des coinvariants de $G$ dans $P \otimes_A M$, c'est-à-dire l'élément
$$
\text{Tr}_{\Lambda}\left(1; \left(P \otimes_A M\right)_G\right) \in
\Lambda^\natural.
$$
L'élément $1\in \mathbf{N}$ correspondra au Frobenius.

\begin{lemma}
\label{trace-lemma-epsilon}
Soit $e\in G$ l'élément neutre. L'application
$$
\begin{matrix}
\Lambda[G] & \longrightarrow & \Lambda^{\natural}\\
\sum \lambda_g\cdot g & \longmapsto & \lambda_e
\end{matrix}
$$
se factorise par $\Lambda[G]^\natural$. On note
$\varepsilon : \Lambda[G]^\natural\to \Lambda^\natural$ l'application induite.
\end{lemma}

\begin{proof}
Il faut montrer que l'application annule les commutateurs. On a
$$
\left(\sum\lambda_g g\right)\left(\sum\mu_g g\right)-\left(\sum \mu_g
g\right)\left(\sum\lambda_g g\right)
= \sum_g\left(\sum_{g_1g_2=g}
\lambda_{g_1}\mu_{g_2}-\mu_{g_1}\lambda_{g_2}\right)g
$$
Le coefficient de $e$ est
$$
\sum_g\left(\lambda_g\mu_{g^{-1}}-\mu_g\lambda_{g^{-1}}\right) =
\sum_g\left(\lambda_g\mu_{g^{-1}}-\mu_{g^{-1}}\lambda_g\right)
$$
qui est une somme de commutateurs et est donc nul dans $\Lambda^\natural$.
\end{proof}

\begin{definition}
\label{trace-definition-trace-G}
Soit $f : P\to P$ un
endomorphisme d'un $\Lambda[G]$-module projectif de type fini
$P$. On pose
$$
\text{Tr}_{\Lambda}^G(f; P) := \varepsilon\left(\text{Tr}_{\Lambda[G]}(f;
P)\right)
$$
et l'on appelle cet élément la {\it $G$-trace de $f$ sur $P$}.
\end{definition}

\begin{lemma}
\label{trace-lemma-lambda-trace}
Soit $f : P\to P$ un endomorphisme du
$\Lambda[G]$-module projectif de type fini $P$. Alors
$$
\text{Tr}_{\Lambda}(f; P) = \# G \cdot \text{Tr}_\Lambda^G(f; P).
$$
\end{lemma}

\begin{proof}
Par additivité, il suffit de traiter le cas $P = \Lambda[G]$.
Dans ce cas, $f$ est donné par
la multiplication à droite par un élément $\sum\lambda_g\cdot g$ de $\Lambda[G]$. Dans
la base $(g)_{g \in G}$, la matrice de $f$ a pour coefficient
$\lambda_{g_2^{-1}g_1}$ à la position $(g_1, g_2)$. En particulier, tous les
coefficients diagonaux sont égaux à $\lambda_e$, et il y en a $\# G$.
\end{proof}

\begin{lemma}
\label{trace-lemma-A-module-structure}
Le morphisme $A\to \Lambda$ définit une structure de $A$-module sur $\Lambda^\natural$.
\end{lemma}

\begin{proof}
C'est clair.
\end{proof}

\begin{lemma}
\label{trace-lemma-diagonal-action-projective-module}
Soient $P$ un $A[G]$-module projectif de type fini et $M$ un $\Lambda[G]$-module,
projectif de type fini en tant que $\Lambda$-module. Alors $P \otimes_A M$ est un
$\Lambda[G]$-module projectif de type fini, pour la structure induite par l'action
diagonale de $G$.
\end{lemma}

\noindent
Notons que $P \otimes_A M$ est naturellement un $\Lambda$-module, puisque $M$ en est un.
Explicitement, avec l'action diagonale, cela s'écrit
$$
\left(\sum\lambda_g g\right)\left(p \otimes m\right)
=
\sum g p \otimes \lambda_g g m.
$$

\begin{proof}
Pour tout $\Lambda[G]$-module $N$, on a
$$
\Hom_{\Lambda[G]}\left(P \otimes_A M, N\right)= \Hom_{A[G]}\left(P,
\Hom_{\Lambda}(M, N)\right)
$$
où l'action de $G$ sur $\Hom_{\Lambda}(M, N)$ est donnée par $(g\cdot
\varphi)(m) = g \varphi (g^{-1} m) $. Il suffit alors de remarquer que le
membre de droite est un composé de foncteurs exacts, grâce à la projectivité
de $P$ et de $M$.
\end{proof}

\begin{lemma}
\label{trace-lemma-multiplicative-trace}
Sous les hypothèses du
lemme \ref{trace-lemma-diagonal-action-projective-module},
soient
$u\in \text{End}_{A[G]}(P)$ et $v\in \text{End}_{\Lambda[G]}(M)$. Alors
$$
\text{Tr}_\Lambda^G \left(u \otimes v; P \otimes_A M\right) = \text{Tr}_A^G(u;
P)\cdot \text{Tr}_\Lambda(v;M).
$$
\end{lemma}

\begin{proof}[Esquisse de démonstration]
On se ramène au cas $P=A[G]$. Dans ce cas, $u$ est la multiplication à droite par un
élément $a = \sum a_gg$ de $A[G]$, et l'on écrit $u = R_a$. Il existe un
isomorphisme de $\Lambda[G]$-modules
$$
\begin{matrix}
\varphi : & A[G]\otimes_A M & \cong & \left(A[G]\otimes_A M\right)'\\
& g \otimes m & \longmapsto & g \otimes g^{-1}m
\end{matrix}
$$
où $\left(A[G]\otimes_A M\right)'$ est muni de la structure de module définie par
l'action à gauche de $G$, ainsi que par la $\Lambda$-linéarité sur $M$. Ce transport
de structure transforme $u \otimes v$ en $\sum_ga_gR_g \otimes g^{-1}v$. En d'autres
termes,
$$
\varphi \circ (u \otimes v) \circ \varphi^{-1}
=
\sum_ga_gR_g \otimes g^{-1}v.
$$
En explicitant les deux membres de l'égalité, il suffit de montrer que
$$
\text{Tr}_\Lambda^G\left(\sum_g a_gR_g \otimes g^{-1}v\right) = a_e\cdot
\text{Tr}_\Lambda(v; M).
$$
Pour cela, on montre que
$$
\text{Tr}_\Lambda^G\left(a_gR_g \otimes g^{-1}v\right) =
\left\{
\begin{matrix}
0 & \text{ si } g\neq e\\
a_e\text{Tr}_\Lambda\left(v; M\right) & \text{ si }g = e
\end{matrix}
\right.
$$
en se ramenant au cas $M=\Lambda$.
\end{proof}

\noindent
Notation :
Considérons l'extension de monoïdes $1 \to G\to \Gamma\to \mathbf{N} \to 1$ et soit
$\gamma\in \Gamma$.
On pose alors $Z_\gamma = \{g\in G | g\gamma = \gamma g\}$.

\begin{lemma}
\label{trace-lemma-gamma-z-gamma-trace}
Soit $P$ un $\Lambda[\Gamma]$-module, projectif de type fini en tant que
$\Lambda[G]$-module, et soit $\gamma \in \Gamma$. Alors
$$
\text{Tr}_{\Lambda}(\gamma, P) =
\# Z_\gamma \cdot \text{Tr}_\Lambda^{Z_\gamma}\left(\gamma, P\right).
$$
\end{lemma}

\begin{proof}
Cela résulte immédiatement du lemme \ref{trace-lemma-lambda-trace}.
\end{proof}

\begin{lemma}
\label{trace-lemma-weak-trace}
Soit $P$ un $A[\Gamma]$-module, projectif de type fini en tant que $A[G]$-module. Soit $M$
un $\Lambda[\Gamma]$-module, projectif de type fini en tant que $\Lambda$-module. Alors
$$
\text{Tr}_{\Lambda}^{Z_\gamma}(\gamma, P \otimes_A M) =
\text{Tr}_A^{Z_\gamma}(\gamma, P)\cdot \text{Tr}_\Lambda(\gamma, M).
$$
\end{lemma}

\begin{proof}
Cela résulte directement du lemme \ref{trace-lemma-multiplicative-trace}.
\end{proof}

\begin{lemma}
\label{trace-lemma-trivial-trace}
Soit $P$ un $\Lambda[\Gamma]$-module, projectif de type fini en tant que
$\Lambda[G]$-module. Alors les coinvariants
$P_G = \Lambda \otimes_{\Lambda[G]} P$
forment un $\Lambda$-module projectif de type fini, muni d'une action de
$\Gamma/G = \mathbf{N}$. De plus, on a
$$
\text{Tr}_\Lambda(1; P_G) =
\sum\nolimits'_{\gamma \mapsto 1} \text{Tr}_\Lambda^{Z_\gamma}(\gamma, P)
$$
où $\sum_{\gamma\mapsto 1}'$ désigne la somme étendue aux classes de $G$-conjugaison
dans $\Gamma$.
\end{lemma}

\begin{proof}[Esquisse de démonstration]
Nous le démontrons d'abord après multiplication par $\# G$.
$$
\# G\cdot \text{Tr}_\Lambda(1; P_G)
= \text{Tr}_\Lambda(\sum\nolimits_{\gamma\mapsto 1} \gamma, P_G)
= \text{Tr}_\Lambda(\sum\nolimits_{\gamma\mapsto 1} \gamma, P)
$$
où la seconde égalité résulte du triangle commutatif
$$
\xymatrix{
P^G \ar[rd]_a & & P_G \ar[ll]^c \\
& P \ar[ur]_b
}
$$
où $a$ est l'inclusion canonique, $b$ la surjection canonique et $c =
\sum_{\gamma \mapsto 1} \gamma$. On a alors
$$
(\sum\nolimits_{\gamma \mapsto 1} \gamma) |_P = a \circ c \circ b
\quad\text{et}\quad
(\sum\nolimits_{\gamma \mapsto 1} \gamma) |_{P_G} = b \circ a \circ c
$$
ils ont donc même trace. On a alors
$$
\# G\cdot \text{Tr}_\Lambda(1; P_G)
=
{\sum_{\gamma\mapsto 1}}'
\frac{\# G}{\# Z_\gamma}\text{Tr}_\Lambda(\gamma, P)
= \# G{\sum_{\gamma\mapsto 1}}' \text{Tr}_\Lambda^{Z_\gamma}(\gamma, P).
$$
Pour achever la démonstration, on se ramène au cas où $\Lambda$ est sans torsion par un
argument d'universalité. Voir \cite{SGA4.5} pour les détails.
\end{proof}

\begin{remark}
\label{trace-remark-content-trivial-trace}
Illustrons le contenu de la formule du
lemme \ref{trace-lemma-weak-trace}.
Supposons que $\Lambda$, vu comme $\Gamma$-module trivial, admette une
résolution finie
$
0\to P_r\to \ldots \to P_1 \to P_0\to \Lambda\to 0
$
par des $\Lambda[\Gamma]$-modules $P_i$ qui sont projectifs de type fini en tant que
$\Lambda[G]$-modules. Dans ce cas,
$$
H_*\left(\left(P_\bullet\right)_G\right) =
\text{Tor}_*^{\Lambda[G]}\left(\Lambda, \Lambda\right) = H_*(G, \Lambda)
$$
et
$$
\text{Tr}_\Lambda^{Z_\gamma}\left(\gamma, P_\bullet\right) =\frac{1}{\#
Z_\gamma}\text{Tr}_\Lambda(\gamma, P_\bullet)=\frac{1}{\#
Z_\gamma}\text{Tr}(\gamma, \Lambda) = \frac{1}{\# Z_\gamma}.
$$
Par conséquent, le lemme \ref{trace-lemma-weak-trace} donne
$$
\text{Tr}_\Lambda (1 , P_G)
= \text{Tr}\left(1 |_{H_*(G, \Lambda)}\right)
= {\sum_{\gamma\mapsto 1}}'\frac{1}{\# Z_\gamma}.
$$
On peut interpréter ceci comme un comptage de points sur le champ $BG$. Si
$\Lambda = \mathbf{F}_\ell$ avec $\ell$ premier à $\# G$, alors
$H_*(G, \Lambda)$ est égal à $\mathbf{F}_\ell$ en degré 0 (et à $0$ aux
autres degrés) et la formule s'écrit
$$
1 =
\sum\nolimits_{
\frac{\sigma\text{-conjugaison}}{\text{classes}\langle\gamma\rangle}
}
\frac{1}{\# Z_\gamma} \mod \ell.
$$
C'est, en un certain sens, une formule des traces ``triviale'' pour $G$.
Nous verrons plus loin que (\ref{trace-equation-trace-formula}) peut, dans
certains cas, être considérée comme une formule des traces hautement non triviale pour un
certain type de groupe ; voir la
section \ref{trace-section-abstract-trace-formula}.
\end{remark}





%11.12.09
\section{Démonstration de la formule des traces}
\label{trace-section-proof-trace-formula}

\begin{theorem}
\label{trace-theorem-trace-formula-again}
Soit $k$ un corps fini et soit $X$ un schéma séparé de type fini de dimension
au plus 1 sur $k$. Soit $\Lambda$ un anneau fini dont le cardinal est premier
à celui de $k$, et soit $K\in D_{ctf}(X, \Lambda)$. Alors
\begin{equation}
\label{trace-equation-trace-formula-again}
\text{Tr}(\pi_X^* |_{R\Gamma_c(X_{\bar k}, K)})
=
\sum\nolimits_{x\in X(k)}
\text{Tr}(\pi_x |_{K_{\bar x}})
\end{equation}
dans $\Lambda^{\natural}$.
\end{theorem}

\noindent
On trouvera dans la remarque \ref{trace-remark-on-trace-formula-again} quelques remarques
sur l'énoncé. Notation : pour abréger, nous écrivons
$$
T'(X, K) =
\sum\nolimits_{x\in X(k)}
\text{Tr}(\pi_x |_{K_{\bar x}})
$$
pour le membre de droite de (\ref{trace-equation-trace-formula-again}) et
$$
T''(X, K)
=\text{Tr}(\pi_x^* |_{R\Gamma_c(X_{\bar k}, K)})
$$
pour le membre de gauche.

\begin{proof}[Démonstration du théorème \ref{trace-theorem-trace-formula-again}]
La démonstration se fait en plusieurs étapes.

\medskip\noindent
Étape 1. {\it Soit $j : \mathcal{U}\hookrightarrow X$ une immersion ouverte de
complément $Y = X - \mathcal{U}$ et soit $i : Y \hookrightarrow X$. Alors
$T''(X, K) = T''(\mathcal{U}, j^{-1} K)+ T''(Y, i^{-1}K)$ et
$T'(X, K) = T'(\mathcal{U}, j^{-1} K)+ T'(Y, i^{-1}K)$.}

\medskip\noindent
C'est clair pour $T'$. Pour $T''$, utilisons la suite exacte
$$
0\to j_!j^{-1} K \to K \to i_* i^{-1} K \to 0
$$
pour munir $K$ d'une filtration. On obtient ainsi un objet
$\widetilde K \in DF(X, \Lambda)$
dont les gradués sont $j_!j^{-1}K$ et $i_*i^{-1}K$;
tous deux appartiennent à $D_{ctf}(X, \Lambda)$. Par les généralités abstraites sur les catégories dérivées filtrées
(INSÉRER UNE RÉFÉRENCE),
$R\Gamma_c(X_{\bar k}, K)\in DF_{perf}(\Lambda)$,
et cet objet est muni de $\pi_x^*$ dans
$DF_{perf}(\Lambda)$.
D'après l'analyse des traces sur les complexes filtrés (INSÉRER UNE RÉFÉRENCE), on obtient
\begin{eqnarray*}
\text{Tr}(\pi_X^* |_{R\Gamma_c(X_{\bar k}, K)})
& = & \text{Tr}(\pi_X^* |_{R\Gamma_c(X_{\bar k}, j_!j^{-1}K)}) +
\text{Tr}(\pi_X^* |_{R\Gamma_c(X_{\bar k}, i_*i^{-1}K)}) \\
& = & T''(U, i^{-1}K) + T''(Y, i^{-1}K).
\end{eqnarray*}

\noindent
Étape 2. {\it Le théorème est vrai si $\dim X\leq 0$. }

\medskip\noindent
En effet, dans ce cas
$$
R\Gamma_c(X_{\bar k}, K) = R\Gamma(X_{\bar k}, K) = \Gamma(X_{\bar k}, K) =
\bigoplus\nolimits_{\bar x \in X_{\bar k}} K_{\bar x}
$$
qui est muni de l'endomorphisme $\pi_X^*$.
Comme les points fixes de $\pi_X : X_{\bar k}\to X_{\bar k}$ sont exactement les
points $\bar x \in X_{\bar k}$ situés au-dessus d'un point $k$-rationnel $x \in X(k)$
nous obtenons
$$
\text{Tr}\big(\pi_X^*|_{R\Gamma_c(X_{\bar k}, K)}\big) =
\sum\nolimits_{x \in X(k)} \text{Tr}(\pi_x|_{K_{\bar x}}).
$$
ce qui est le résultat voulu.

\medskip\noindent
Étape 3. {\it Il suffit de démontrer l'égalité
$T'(\mathcal{U}, \mathcal{F}) = T''(\mathcal{U}, \mathcal{F})$
dans le cas où
\begin{itemize}
\item $\mathcal{U}$ est une courbe affine irréductible lisse sur $k$,
\item $\mathcal{U}(k) = \emptyset$,
\item $K=\mathcal{F}$ est un faisceau fini localement constant de $\Lambda$-modules
sur $\mathcal{U}$ dont les fibres sont des $\Lambda$-modules projectifs de type fini, et
\item $\Lambda$ est annulé par une puissance d'un nombre premier $\ell$ et $\ell \in k^*$.
\end{itemize}
}

\medskip\noindent
En effet, d'après l'étape 2, nous pouvons retrancher tout ensemble fini de points. Mais nous
n'avons qu'un nombre fini de points rationnels, de sorte que nous pouvons supposer qu'il n'y en a
aucun\footnote{À ce stade, un rire maléfique devrait retentir à l'arrière-plan.}.
Nous pouvons supposer que $\mathcal{U}$ est lisse, irréductible et affine en passant aux
composantes irréductibles et en retranchant au besoin les mauvais points. Les
hypothèses sur $\mathcal{F}$ s'obtiennent en explicitant la définition de
$D_{ctf}(X, \Lambda)$ et celles sur $\Lambda$ de l'examen de sa décomposition
primaire.

\medskip\noindent
Dans la suite de la démonstration, nous considérons la situation
$$
\xymatrix{
\mathcal{V} \ar[d]_f \ar[r] & Y \ar[d]^{\bar f} \\
\mathcal{U} \ar[r] & X
}
$$
où $\mathcal{U}$ est comme ci-dessus, $f$ est un revêtement \'etale galoisien fini,
$\mathcal{V}$ est connexe et les flèches horizontales sont des
compactifications projectives. En posant $G=\text{Aut}(\mathcal{V}|\mathcal{U})$, nous supposons également
(comme nous le pouvons) que $f^{-1}\mathcal{F} =\underline M$ est constant, où le
module $M = \Gamma(\mathcal{V}, f^{-1}\mathcal{F})$ est un $\Lambda[G]$-module
projectif de type fini sur $\Lambda$. Cela correspond à l'extension de
monoïdes triviale
$$
1\to G\to \Gamma = G \times \mathbf{N}\to \mathbf{N}\to 1.
$$
Dans ce contexte, compte tenu des réductions ci-dessus, nous devons montrer que
$T''(\mathcal{U}, \mathcal{F}) = 0$.

\medskip\noindent
Étape 4. {\it Il existe une action naturelle de $G$ sur $f_*f^{-1}\mathcal{F}$ et
le morphisme trace $f_*f^{-1}\mathcal{F}\to \mathcal{F}$ définit un isomorphisme}
$$
(f_*f^{-1}\mathcal{F})\otimes_{\Lambda[G]} \Lambda =
(f_*f^{-1}\mathcal{F})_G \cong \mathcal{F}.
$$

\medskip\noindent
Pour le démontrer, il suffit de tout expliciter en un point géométrique.

\medskip\noindent
Étape 5. {\it Soit $A = \mathbf{Z}/\ell^n \mathbf{Z}$ avec $n\gg 0$. Alors
$f_*f^{-1}\mathcal{F} \cong (f_*\underline A)
\otimes_{\underline A} \underline M$ muni de l'action diagonale de $G$.}

\medskip\noindent
Étape 6. {\it Il existe un isomorphisme canonique
$(f_*\underline A \otimes_{\underline A} \underline M)
\otimes_{\Lambda[G]} \underline \Lambda \cong \mathcal{F}$.
}

\medskip\noindent
En fait, il s'agit d'un produit tensoriel dérivé, grâce à l'hypothèse de projectivité
faite sur $\mathcal{F}$.

\medskip\noindent
Étape 7. {\it Il existe un isomorphisme canonique
$$
R\Gamma_c(\mathcal{U}_{\bar k}, \mathcal{F})
= (R\Gamma_c(\mathcal{U}_{\bar k}, f_*A)\otimes_A^\mathbf{L}
M)\otimes_{\Lambda[G]}^\mathbf{L} \Lambda,
$$
compatible avec l'action de $\pi^*_\mathcal{U}$.
}

\medskip\noindent
Cela résulte du théorème des coefficients universels, c'est-à-dire du fait que
$R\Gamma_c$ commute à $\otimes^\mathbf{L}$, et de la platitude de
$\mathcal{F}$ comme $\Lambda$-module.

\medskip\noindent
Nous avons
\begin{eqnarray*}
\text{Tr}(
\pi_\mathcal{U}^* |_{R\Gamma_c(\mathcal{U}_{\bar k}, \mathcal{F})})
& = &
{\sum_{g \in G}}'
\text{Tr}_{\Lambda}^{Z_g}
\left(
(g, \pi_\mathcal{U}^*)
|_{R\Gamma_c(\mathcal{U}_{\bar k}, f_*A)\otimes_A^\mathbf{L} M}
\right) \\
& = &
{\sum_{g\in G}}'
\text{Tr}_A^{Z_g}
(
(g, \pi_\mathcal{U}^*) |_{R\Gamma_c(\mathcal{U}_{\bar k}, f_*A)}
)
\cdot
\text{Tr}_\Lambda(g|_M)
\end{eqnarray*}
où l'action de $\Gamma$ sur $R\Gamma_c(\mathcal{U}_{\bar k}, \mathcal{F})$ prolonge celle de $G$
et $(e, 1)$ agit par $\pi_\mathcal{U}^*$. L'extension de monoïdes est donc donnée
par $\Gamma = G \times \mathbf{N} \to \mathbf{N}$, $\gamma \mapsto 1$. La première
égalité résulte du lemme \ref{trace-lemma-trivial-trace} et la seconde du
lemme \ref{trace-lemma-weak-trace}.

\medskip\noindent
Étape 8. {\it Il suffit de montrer que
$\text{Tr}_A^{Z_g}((g, \pi_\mathcal{U}^*)
|_{R\Gamma_c(\mathcal{U}_{\bar k}, f_*A)}) \in A$
s'envoie sur zéro dans $\Lambda$.
}

\medskip\noindent
Rappelons que
\begin{eqnarray*}
\# Z_g \cdot \text{Tr}_A^{Z_g}((g, \pi_\mathcal{U}^*)
|_{R\Gamma_c(\mathcal{U}_{\bar k}, f_*A)})
& = & \text{Tr}_A((g, \pi_\mathcal{U}^*)
|_{R\Gamma_c(\mathcal{U}_{\bar k}, f_*A)})\\
& = &
\text{Tr}_A((g^{-1}\pi_\mathcal{V})^* |_{R\Gamma_c(\mathcal{V}_{\bar k}, A)}).
\end{eqnarray*}
La première égalité découle du
le lemme \ref{trace-lemma-gamma-z-gamma-trace},
la seconde de la suite spectrale de Leray,
en utilisant la finitude de $f$ et le fait que nous ne prenons que
des traces sur $A$. Or, comme $A=\mathbf{Z}/\ell^n\mathbf{Z}$ avec
$n \gg 0$ et $\# Z_g = \ell^a$ pour un certain entier (fixé) $a$,
il suffit de montrer le résultat suivant.

\medskip\noindent
Étape 9. {\it On a
$\text{Tr}_A((g^{-1}\pi_\mathcal{V})^* |_{R\Gamma_c(\mathcal{V}, A)}) = 0$
dans $A$.}

\medskip\noindent
À nouveau, par additivité, nous avons
\begin{eqnarray*}
& \text{Tr}_A((g^{-1}\pi_\mathcal{V})^* |_{R\Gamma_c(\mathcal{V}_{\bar k} A)})
+
\text{Tr}_A((g^{-1}\pi_\mathcal{V})^*
|_{R\Gamma_c(Y-\mathcal {V})_{\bar k}, A)}) \\
& =
\text{Tr}_A((g^{-1}\pi_Y)^* |_{R\Gamma(Y_{\bar k}, A)})
\end{eqnarray*}
Cette dernière trace est le nombre de points fixes de $g^{-1}\pi_Y$ sur $Y$, d'après
la formule des traces de Weil,
théorème \ref{trace-theorem-weil-trace-formula}.
De plus, d'après le cas de dimension 0 déjà démontré à l'étape 2,
$$
\text{Tr}_A((g^{-1}\pi_\mathcal{V})^*|_{R\Gamma_c(Y-\mathcal{V})_{\bar k}, A)})
$$
est le nombre de points fixes de $g^{-1}\pi_Y$ sur $(Y-\mathcal{V})_{\bar k}$.
Par conséquent,
$$
\text{Tr}_A((g^{-1}\pi_\mathcal{V})^* |_{R\Gamma_c(\mathcal{V}_{\bar k}, A)})
$$
est le nombre de points fixes de $g^{-1}\pi_Y$ sur $\mathcal{V}_{\bar k}$. Mais
il n'y en a aucun : si $\bar y\in Y_{\bar k}$ est fixe par
$g^{-1}\pi_Y$, alors $\bar f(\bar y) \in X_{\bar k}$ est fixe par $\pi_X$. Mais
$\mathcal{U}$ n'a aucun point $k$-rationnel, donc on doit avoir $\bar f(\bar y)\in
(X-\mathcal{U})_{\bar k}$ et ainsi $\bar y\notin \mathcal{V}_{\bar k}$, une
contradiction.
Ceci achève la démonstration.
\end{proof}

\begin{remark}
\label{trace-remark-on-trace-formula-again}
Remarques sur le théorème \ref{trace-theorem-trace-formula-again}.
\begin{enumerate}
\item
Cette formule vaut en toute dimension. Par un lemme de d\'evissage (qui utilise le changement de
base propre, etc.), elle se ramène, même dans cette généralité, à l'énoncé actuel.
\item
Le complexe $R\Gamma_c(X_{\bar k}, K)$ est défini en choisissant une immersion ouverte
$j : X \hookrightarrow \bar X$ avec $\bar X$ projectif sur $k$ et de dimension au
plus 1, et en posant
$$
R\Gamma_c(X_{\bar k}, K) := R\Gamma(\bar X_{\bar k}, j_!K).
$$
L'indépendance de ce choix de $\bar X$ résulte de
(insérer ici une référence). Nous définissons $H^i_c(X_{\bar k}, K)$
comme le $i$-ième groupe de cohomologie de $R\Gamma_c(X_{\bar k}, K)$.
\end{enumerate}
\end{remark}

\begin{remark}
\label{trace-remark-stronger}
Bien que nous n'ayons fait que des réductions et surtout de l'algèbre, la formule des traces
du théorème \ref{trace-theorem-trace-formula-again} est bien plus forte que
la formule géométrique des traces de Weil (théorème \ref{trace-theorem-weil-trace-formula})
car elle s'applique aux systèmes de coefficients
(aux faisceaux), et pas seulement aux coefficients constants.
\end{remark}


%11.17.09
\section{Applications}
\label{trace-section-applications}

\noindent
Bon, après avoir indiqué la démonstration de la formule des traces, essayons d'en tirer
quelque chose.





\section{Sur les faisceaux l-adiques}
\label{trace-section-l-adic-sheaves}

\begin{definition}
\label{trace-definition-l-adic-sheaf}
Soit $X$ un schéma noethérien. Un {\it faisceau de $\mathbf{Z}_\ell$-modules} sur $X$, ou
simplement un {\it faisceau $\ell$-adique} $\mathcal{F}$, est un
système projectif $\left\{\mathcal{F}_n\right\}_{n\geq 1}$ dans lequel
\begin{enumerate}
\item
$\mathcal{F}_n$ est un faisceau constructible de $\mathbf{Z}/\ell^n\mathbf{Z}$-modules sur
$X_\etale$, et
\item
les morphismes de transition $\mathcal{F}_{n+1}\to \mathcal{F}_n$ induisent des isomorphismes
$\mathcal{F}_{n+1} \otimes_{\mathbf{Z}/\ell^{n+1}\mathbf{Z}}
\mathbf{Z}/\ell^n\mathbf{Z} \cong \mathcal{F}_n$.
\end{enumerate}
On dit que $\mathcal{F}$ est {\it lisse} si chaque $\mathcal{F}_n$ est localement
constant. Un {\it morphisme} entre de tels faisceaux n'est autre qu'un morphisme de systèmes projectifs.
\end{definition}

\begin{lemma}
\label{trace-lemma-eventually-constant}
Soit $\{\mathcal{G}_n\}_{n\geq 1}$ un système projectif de faisceaux constructibles de
$\mathbf{Z}/\ell^n\mathbf{Z}$-modules.
Supposons que, pour tout $k\geq 1$, les morphismes
$$
\mathcal{G}_{n+1}/\ell^k \mathcal{G}_{n+1}\to \mathcal{G}_n /\ell^k
\mathcal{G}_n
$$
soient des isomorphismes pour tout $n\gg 0$ (la borne pouvant dépendre de $k$).
En d'autres termes, supposons que le système
$\{\mathcal{G}_n/\ell^k\mathcal{G}_n\}_{n\geq 1}$
soit essentiellement constant, et désignons par $\mathcal{F}_k$ le faisceau correspondant.
Alors le système $\left\{\mathcal{F}_k\right\}_{k\geq 1}$ est un
faisceau de $\mathbf{Z}_\ell$-modules sur $X$.
\end{lemma}

\begin{proof}
La démonstration est immédiate.
\end{proof}

\begin{lemma}
\label{trace-lemma-l-adic-abelian}
La catégorie des faisceaux de $\mathbf{Z}_\ell$-modules sur $X$ est abélienne.
\end{lemma}

\begin{proof}
Soit
$\Phi = \left\{\varphi_n\right\}_{n\geq 1} :
\left\{\mathcal{F}_n\right\}
\to
\left\{\mathcal{G}_n\right\}$
un morphisme de faisceaux de $\mathbf{Z}_\ell$-modules. Posons
$$
\Coker(\Phi) =
\left\{
\Coker\left(\mathcal{F}_n \xrightarrow{\varphi_n} \mathcal{G}_n\right)
\right\}_{n\geq 1}
$$
et $\Ker(\Phi)$ s'obtient en appliquant le
lemme \ref{trace-lemma-eventually-constant}
au système projectif
$$
\left\{
\bigcap_{m\geq n}
\Im\left(\Ker(\varphi_m) \to \Ker(\varphi_n)\right)
\right\}_{n \geq 1}.
$$
On laisse au lecteur le soin de vérifier que ceci définit une catégorie abélienne.
\end{proof}

\begin{example}
\label{trace-example-kernel}
Soit $X=\Spec(\mathbf{C})$ et soit $\Phi : \mathbf{Z}_\ell\to \mathbf{Z}_\ell$
la multiplication par $\ell$. Plus précisément,
$$
\Phi = \left\{ \mathbf{Z}/\ell^n\mathbf{Z} \xrightarrow{\ell}
\mathbf{Z}/\ell^n\mathbf{Z}\right\}_{n \geq 1}.
$$
Pour calculer le noyau, considérons le système projectif
$$
\ldots\to \mathbf{Z}/\ell\mathbf{Z}\xrightarrow{0}
\mathbf{Z}/\ell\mathbf{Z}\xrightarrow{0}\mathbf{Z}/\ell\mathbf{Z}.
$$
Puisque les images sont toujours nulles, $\Ker(\Phi)$ est le système nul.
\end{example}

\begin{remark}
\label{trace-remark-stalk-l-adic-sheaf}
Si $\mathcal{F} = \left\{\mathcal{F}_n\right\}_{n\geq 1}$ est un
faisceau de $\mathbf{Z}_\ell$-modules sur $X$ et si $\bar x$ est un point géométrique, alors
$M_n = \left\{\mathcal{F}_{n, \bar x}\right\}$ est un système projectif de
$\mathbf{Z}/\ell^n\mathbf{Z}$-modules de type fini dont $M_{n+1}\to M_n$ est surjectif
et $M_n = M_{n+1}/\ell^n M_{n+1}$. Il en résulte que
$$
M = \lim_n M_n = \lim \mathcal{F}_{n, \bar x}
$$
est un module de type fini sur $\mathbf{Z}_\ell$. Cela résulte des
lemmes \ref{algebra-lemma-limit-complete} et
\ref{algebra-lemma-finite-over-complete-ring} du chapitre Algèbre, ainsi que du fait que
$M/\ell M = M_1$ est de dimension finie sur $\mathbf{F}_\ell$.
On a donc $M\cong \mathbf{Z}_\ell^{\oplus r} \oplus \oplus_{i = 1}^t
\mathbf{Z}_\ell/\ell^{e_i}\mathbf{Z}_\ell$ pour certains $r, t\geq 0$, $e_i\geq 1$.
Le module $M = \mathcal{F}_{\bar x}$ est appelé la {\it fibre} de
$\mathcal{F}$ en $\bar x$.
\end{remark}

\begin{definition}
\label{trace-definition-torsion-l-adic-sheaf}
Un faisceau de $\mathbf{Z}_\ell$-modules $\mathcal{F}$ est {\it de torsion} si
$\ell^n : \mathcal{F} \to \mathcal{F}$ est le morphisme nul pour un certain $n$.
La catégorie abélienne
des faisceaux de $\mathbf{Q}_\ell$-modules sur $X$ est le quotient de la catégorie abélienne des
faisceaux de $\mathbf{Z}_\ell$-modules par la sous-catégorie de Serre des faisceaux de torsion. En
d'autres termes, ses objets sont les faisceaux de $\mathbf{Z}_\ell$-modules sur $X$ et, si
$\mathcal{F}, \mathcal{G}$ sont deux tels faisceaux, alors
$$
\Hom_{\mathbf{Q}_\ell} \left(\mathcal{F}, \mathcal{G} \right) =
\Hom_{\mathbf{Z}_\ell} \left(\mathcal{F}, \mathcal{G}\right)
\otimes_{\mathbf{Z}_\ell} \mathbf{Q}_\ell.
$$
On désigne par $\mathcal{F} \mapsto \mathcal{F} \otimes \mathbf{Q}_\ell$ le
foncteur de passage au quotient (adjoint à droite du foncteur d'inclusion). Si $\mathcal{F} =
\mathcal{F}' \otimes \mathbf{Q}_\ell$ où $\mathcal{F}'$ est un
faisceau de $\mathbf{Z}_\ell$-modules et $\bar x$ un point géométrique, alors la
{\it fibre} de $\mathcal{F}$ en $\bar x$ est $\mathcal{F}_{\bar x} =
\mathcal{F}'_{\bar x} \otimes \mathbf{Q}_\ell$.
\end{definition}

\begin{remark}
\label{trace-remark-torsion-stalks}
Puisqu'un faisceau de $\mathbf{Z}_\ell$-modules n'est défini que sur un schéma noethérien, il est
de torsion si et seulement si ses fibres sont de torsion.
\end{remark}

\begin{definition}
\label{trace-definition-cohomology-l-adic}
Si $X$ est un schéma séparé de type fini sur un corps algébriquement clos
$k$ et si $\mathcal{F} = \left\{\mathcal{F}_n\right\}_{n\geq 1}$ est un
faisceau de $\mathbf{Z}_\ell$-modules sur $X$, on définit
$$
H^i(X, \mathcal{F}) := \lim_n H^i(X, \mathcal{F}_n)
\quad\text{et}\quad
H_c^i(X, \mathcal{F}) := \lim_n H_c^i(X, \mathcal{F}_n).
$$
Si $\mathcal{F} = \mathcal{F}'\otimes \mathbf{Q}_\ell$ pour un
faisceau de $\mathbf{Z}_\ell$-modules $\mathcal{F}'$, on pose
$$
H_c^i(X , \mathcal{F}) := H_c^i(X,
\mathcal{F}')\otimes_{\mathbf{Z}_\ell}\mathbf{Q}_\ell.
$$
On appelle ces groupes la {\it cohomologie $\ell$-adique} de $X$ à coefficients dans
$\mathcal{F}$.
\end{definition}





\section{Fonctions L}
\label{trace-section-L-function}

\begin{definition}
\label{trace-definition-L-function-finite-ring}
Soit $X$ un schéma de type fini sur un corps fini $k$. Soit $\Lambda$ un
anneau fini d'ordre premier à la caractéristique de $k$ et $\mathcal{F}$ un
faisceau constructible et plat de $\Lambda$-modules sur $X_\etale$. On pose alors
$$
L(X, \mathcal{F}) :=
\prod\nolimits_{x\in |X|}
\det(1 - \pi_x^*T^{\deg x} |_{\mathcal{F}_{\bar x}})^{-1} \in \Lambda [[ T ]]
$$
où $|X|$ est l'ensemble des points fermés de $X$, $\deg x = [\kappa(x): k]$ et
$\bar x$ est un point géométrique au-dessus de $x$. Cette définition se
généralise clairement au cas où l'on remplace $\mathcal{F}$ par
$K \in D_{ctf}(X, \Lambda)$. On appelle cette expression la {\it fonction $L$ de
$\mathcal{F}$}.
\end{definition}

\begin{remark}
\label{trace-remark-T}
Intuitivement, il convient de considérer $T$ comme $T = t^f$, où $p^f = \# k$. Les
définitions ne dépendent alors pas du cardinal du corps de base.
\end{remark}

\begin{definition}
\label{trace-definition-L-function-l-adic}
Supposons maintenant que $\mathcal{F}$ soit un faisceau de $\mathbf{Q}_\ell$-modules sur $X$.
Dans ce cas, on définit
$$
L(X, \mathcal{F}) :=
\prod\nolimits_{x \in |X|}
\det(1 - \pi_x^*T^{\deg x} |_{\mathcal{F}_{\bar x}})^{-1}
\in \mathbf{Q}_\ell[[T]].
$$
Notons que ce produit converge puisqu'il n'y a qu'un nombre fini de points de
degré donné. On appelle ce produit la {\it fonction $L$ de
$\mathcal{F}$}.
\end{definition}




\section{Interprétation cohomologique}
\label{trace-section-L-cohomological}

\noindent
C'est ainsi que Grothendieck a interprété la fonction $L$.

\begin{theorem}[Coefficients finis]
\label{trace-theorem-A}
Soit $X$ un schéma de type fini sur un corps fini $k$. Soit $\Lambda$ un
anneau fini d'ordre premier à la caractéristique de $k$ et $\mathcal{F}$ un
faisceau constructible et plat de $\Lambda$-modules sur $X_\etale$. Alors
$$
L(X, \mathcal{F}) =
\det(1 - \pi_X^*\ T |_{R\Gamma_c(X_{\bar k}, \mathcal{F})})^{-1}
\in \Lambda[[T]].
$$
\end{theorem}

\begin{proof}
Démonstration omise.
\end{proof}

\noindent
À ce stade, nous ne savons même pas si chacun des groupes de cohomologie
$H^i_c(X_{\bar k}, \mathcal{F})$ est libre.

\begin{theorem}[Faisceaux adiques]
\label{trace-theorem-B}
Soit $X$ un schéma de type fini sur un corps fini $k$, et soit $\mathcal{F}$ un
faisceau de $\mathbf{Q}_\ell$-modules sur $X$. Alors
$$
L(X, \mathcal{F}) =
\prod\nolimits_i
\det(1 - \pi_X^*T |_{H_c^i(X_{\bar k} , \mathcal{F})})^{(-1)^{i + 1}}
\in \mathbf{Q}_\ell[[T]].
$$
\end{theorem}

\begin{proof}
La démonstration est esquissée ci-dessous.
\end{proof}

\begin{remark}
\label{trace-remark-which-is-harder}
Comme nous n'avons développé qu'une partie de la théorie des traces, et non
celle des déterminants, le théorème \ref{trace-theorem-A}
est plus difficile à démontrer que le
théorème \ref{trace-theorem-B}.
Nous ne démontrerons que ce dernier; pour le premier, voir \cite{SGA4.5}.
Remarquons aussi qu'il n'existe pas de version plus générale de ce théorème à
coefficients dans $\mathbf{Z}_\ell$, puisqu'il n'y a pas de $\ell$-torsion.
\end{remark}

\noindent
La démonstration du théorème \ref{trace-theorem-B} se ramène à une formule des traces. Puisque
$\mathbf{Q}_\ell$ est de caractéristique 0, il suffit de prouver l'égalité après
avoir pris les dérivées logarithmiques. Plus précisément, nous appliquons $T\frac{d}{dT} \log$
aux deux membres. Nous avons d'une part
\begin{align*}
T \frac{d}{dT} \log L(X, \mathcal{F})
& =
T\frac{d}{dT} \log
\prod_{x \in |X|} \det(1 - \pi_x^* T^{\deg x} |_{\mathcal{F}_{\bar x}})^{-1} \\
& =
\sum_{x \in |X|} T \frac{d}{dT} \log
( \det(1 - \pi_x^* T^{\deg x} |_{\mathcal{F}_{\bar x}})^{-1}) \\
& =
\sum_{x \in |X|} \deg x
\sum_{n \geq 1} \text{Tr}((\pi_x^n)^* |_{\mathcal{F}_{\bar x}}) T^{n\deg x}
\end{align*}
où la dernière égalité résulte de la formule
$$
T\frac{d}{dT}\log\left(\det\left(1-fT|_M\right)^{-1}\right) = \sum_{n\geq 1}
\text{Tr}(f^n|_M)T^n
$$
ce qui vaut pour tout anneau commutatif $\Lambda$ et tout endomorphisme $f$ d'un
$\Lambda$-module projectif de type fini $M$. D'autre part, nous avons
\begin{align*}
& T\frac{d}{dT} \log\left(
\prod\nolimits_i
\det(1-\pi_X^*T |_{H_c^i\left(X_{\bar k} , \mathcal{F}\right)})^{(-1)^{i+1}}
\right) \\
& =
\sum\nolimits_i (-1)^i \sum\nolimits_{n \geq 1}
\text{Tr}\left((\pi_X^n)^* |_{H_c^i(X_{\bar k},\mathcal{F})}\right) T^n
\end{align*}
en appliquant encore une fois la même formule. Maintenant, en comparant les puissances de $T$ et en utilisant la
formule d'inversion de Möbius, nous voyons que le théorème \ref{trace-theorem-B} découle de
l'égalité suivante
$$
\sum_{d | n} d \sum_{x \in |X| \atop \deg x = d}
\text{Tr}((\pi_X^{n/d})^* |_{\mathcal{F}_{\bar x}}) =
\sum_i (-1)^i \text{Tr}((\pi^n_X)^* |_{H^i_c(X_{\bar k}, \mathcal{F})}).
$$
Notons $k_n$ l'extension de degré $n$ de $k$,
$X_n = X \times_{\Spec k} \Spec(k_n)$ et
$_n\mathcal{F} = \mathcal{F}|_{X_n}$, cela se ramène
à
$$
\sum_{x \in X_n(k_n)} \text{Tr}(\pi_X^* |_{_n\mathcal{F}_{\bar x}})
=
\sum_i (-1)^i \text{Tr}((\pi^n_X)^* |_{H^i_c({(X_n)}_{\bar k}, _n\mathcal{F})})
$$
ce qui découle du théorème \ref{trace-theorem-C}.

%11.19.09


\begin{theorem}
\label{trace-theorem-D}
Soit $X/k$ comme ci-dessus, soit $\Lambda$ un anneau fini tel que $\#\Lambda \in k^*$
et soit $K\in D_{ctf}(X, \Lambda)$. Alors $R\Gamma_c(X_{\bar k}, K)\in
D_{perf}(\Lambda)$ et
$$
\sum_{x\in X(k)}\text{Tr}\left(\pi_x |_{K_{\bar x}}\right) =
\text{Tr}\left(\pi_X^* |_{R\Gamma_c(X_{\bar k}, K )}\right).
$$
\end{theorem}

\begin{proof}
Remarquons que nous avons déjà démontré ceci (RÉFÉRENCE) lorsque $\dim X \leq 1$. Le
cas général découle aisément de ce cas et du théorème de changement de base
propre.
\end{proof}

\begin{theorem}
\label{trace-theorem-C}
Soit $X$ un schéma séparé de type fini sur un corps fini $k$ et soit
$\mathcal{F}$ un $\mathbf{Q}_\ell$-faisceau sur $X$. Alors
$\dim_{\mathbf{Q}_\ell}H_c^i(X_{\bar k}, \mathcal{F})$ est finie pour tout $i$,
et ne peut être non nulle que pour $0\leq i \leq 2 \dim X$. De plus, nous avons
$$
\sum_{x\in X(k)} \text{Tr}\left(\pi_x |_{\mathcal{F}_{\bar x}}\right) =
\sum_i (-1)^i\text{Tr}\left(\pi_X^* |_{H_c^i(X_{\bar k}, \mathcal{F})}\right).
$$
\end{theorem}

\begin{proof}
Nous expliquons comment déduire ce résultat du théorème \ref{trace-theorem-D}.
Nous utilisons d'abord quelques arguments de cohomologie \'etale pour ramener la démonstration
à un énoncé algébrique que nous démontrons ensuite.

\medskip\noindent
Soit $\mathcal{F}$ comme dans le théorème. Nous pouvons écrire
$\mathcal{F}$ sous la forme
$\mathcal{F}'\otimes \mathbf{Q}_\ell$ où $\mathcal{F}' =
\left\{\mathcal{F}'_n\right\}$ est un $\mathbf{Z}_\ell$-faisceau sans torsion,
c'est-à-dire que $\ell : \mathcal{F}'\to \mathcal{F}'$ a un noyau nul dans la
catégorie des $\mathbf{Z}_\ell$-faisceaux. Alors chaque $\mathcal{F}_n'$ est un
faisceau constructible et plat de $\mathbf{Z}/\ell^n\mathbf{Z}$-modules sur $X_\etale$, donc
$\mathcal{F}'_n \in D_{ctf}(X, \mathbf{Z}/\ell^n\mathbf{Z})$ et
$\mathcal{F}_{n+1}'
\otimes^{\mathbf{L}}_{\mathbf{Z}/\ell^{n+1}\mathbf{Z}}
\mathbf{Z}/\ell^n\mathbf{Z} = \mathcal{F}_n'$.
Remarquons que la dernière égalité est également valable
pour le produit tensoriel usuel (non dérivé), puisque $\mathcal{F}'_n$ est plat
(il s'agit de la même égalité). Par conséquent,
\begin{enumerate}
\item
le complexe $K_n = R\Gamma_c\left(X_{\bar k}, \mathcal{F}_n'\right)$ est parfait,
et il est muni d'un endomorphisme $\pi_n : K_n\to K_n$ dans
$D(\mathbf{Z}/\ell^n\mathbf{Z})$,
\item
on dispose d'identifications
$$
K_{n+1}
\otimes^{\mathbf{L}}_{\mathbf{Z}/\ell^{n+1}\mathbf{Z}}
\mathbf{Z}/\ell^n\mathbf{Z}
=
K_n
$$
dans $D_{perf}(\mathbf{Z}/\ell^n\mathbf{Z})$, compatibles avec les endomorphismes
$\pi_{n+1}$ et $\pi_n$ (voir \cite[Rapport 4.12]{SGA4.5}),
\item
l'égalité $\text{Tr}\left(\pi_X^* |_{K_n}\right) =
\sum_{x\in X(k)} \text{Tr}\left(\pi_x |_{(\mathcal{F}'_n)_{\bar x}}\right)$
est vérifiée, et
\item
pour tout $x\in X(k)$, les éléments
$\text{Tr}(\pi_x |_{\mathcal{F}'_{n, \bar x}}) \in \mathbf{Z}/\ell^n\mathbf{Z}$
définissent un élément de
$\mathbf{Z}_\ell$ qui est égal à
$\text{Tr}(\pi_x |_{\mathcal{F}_{\bar x}}) \in \mathbf{Q}_\ell$.
\end{enumerate}
Il suffit donc de démontrer le lemme d'algèbre suivant.
\end{proof}

\begin{lemma}
\label{trace-lemma-piece-together}
Supposons donnés
$K_n\in D_{perf}(\mathbf{Z}/\ell^n\mathbf{Z})$, $\pi_n : K_n\to K_n$
et des isomorphismes
$\varphi_n :
K_{n+1} \otimes^\mathbf{L}_{\mathbf{Z}/\ell^{n+1}\mathbf{Z}}
\mathbf{Z}/\ell^n\mathbf{Z}
\to K_n$
compatibles avec $\pi_{n+1}$ et $\pi_n$. Alors
\begin{enumerate}
\item
les éléments $t_n = \text{Tr}(\pi_n |_{K_n})\in \mathbf{Z}/\ell^n\mathbf{Z}$
définissent un élément $t_\infty = \{t_n\}$ de $\mathbf{Z}_\ell$,
\item
le $\mathbf{Z}_\ell$-module $H_\infty^i = \lim_n H^i(k_n)$ est de type fini et
est nul sauf pour un nombre fini d'indices $i$, et
\item
les opérateurs $H^i(\pi_n): H^i(K_n)\to H^i(K_n)$ sont compatibles et définissent
$\pi_\infty^i : H_\infty^i\to H_\infty^i$ qui vérifie
$$
\sum (-1)^i \text{Tr}(
\pi_\infty^i |_{H_\infty^i \otimes_{\mathbf{Z}_\ell}\mathbf{Q}_\ell}) =
t_\infty.
$$
\end{enumerate}
\end{lemma}

\begin{proof}
Puisque $\mathbf{Z}/\ell^n\mathbf{Z}$ est un anneau local et que $K_n$ est parfait, chaque
$K_n$ peut être représenté par un complexe fini $K_n^\bullet$ de
$\mathbf{Z}/\ell^n \mathbf{Z}$-modules libres de type fini tel que l'image du morphisme $K_n^p \to K_n^{p+1}$
soit contenue dans $\ell K_n^{p+1}$. Un tel complexe est en fait
unique à isomorphisme près. De plus, $\pi_n$ peut être représenté par un morphisme de
complexes $\pi_n^\bullet : K_n^\bullet\to K_n^\bullet$ (qui est unique à
homotopie près). De même, l'isomorphisme
$\varphi_n : K_{n+1} \otimes_{\mathbf{Z}/\ell^{n+1}\mathbf{Z}}^{\mathbf{L}}
\mathbf{Z}/\ell^n\mathbf{Z}\to K_n$ est représenté par un morphisme de complexes
$$
\varphi_n^\bullet :
K_{n+1}^\bullet
\otimes_{\mathbf{Z}/\ell^{n+1}\mathbf{Z}}
\mathbf{Z}/\ell^n\mathbf{Z} \to K_n^\bullet.
$$
En fait, $\varphi_n^\bullet$ est un isomorphisme de complexes ; nous voyons donc que
\begin{itemize}
\item
il existe $a, b\in \mathbf{Z}$ indépendants de $n$ tels que $K_n^i = 0$
pour tout $i\notin[a, b]$, et
\item
le rang de $K_n^i$ ne dépend pas de $n$.
\end{itemize}
Par conséquent, le module $K_\infty^i = \lim_n \{K_n^i, \varphi_n^i\}$ est un
$\mathbf{Z}_\ell$-module libre de type fini et $K_\infty^\bullet$ est un complexe fini
de $\mathbf{Z}_\ell$-modules libres de type fini. Par récurrence sur le nombre de termes non nuls,
on peut démontrer que $H^i\left(K_\infty^\bullet\right) = \lim_n
H^i\left(K_n^\bullet\right)$ (ceci est faux pour les complexes non bornés). Nous
en concluons que $H_\infty^i = H^i\left(K_\infty^\bullet\right)$ est un
$\mathbf{Z}_\ell$-module de type fini. Cela démontre {\it ii}. Pour démontrer le reste du
lemme, nous devons remédier au défaut éventuel de commutativité des diagrammes
$$
\xymatrix{
{K_{n+1}^\bullet} \ar[d]_{\pi_{n+1}^\bullet} \ar[r]^{\varphi_n^\bullet} &
{K_n^\bullet} \ar[d]^{\pi_n^\bullet} \\
{K_{n+1}^\bullet} \ar[r]_{\varphi_n^\bullet} & {K_n^\bullet.}
}
$$
Cependant, ce diagramme commute bien dans la catégorie dérivée ; il commute donc
à homotopie près. Nous remplaçons par récurrence $\pi_n^\bullet$ pour $n\geq 2$ par
des morphismes de complexes homotopes rendant ces diagrammes commutatifs. Plus précisément, si $h^i :
K_{n+1}^i \to K_n^{i-1}$ est une homotopie, c'est-à-dire
$$
\pi_n^\bullet \circ \varphi_n^\bullet -
\varphi_n^\bullet \circ \pi_{n + 1}^\bullet = dh + hd,
$$
alors nous choisissons $\tilde h^i : K_{n+1}^i\to K_{n+1}^{i-1}$ relevant $h^i$. Cela est
possible parce que $K_{n+1}^i$ est libre et que le morphisme $K_{n+1}^{i-1}\to K_n^{i-1}$ est
surjectif. Nous remplaçons alors $\pi_n^\bullet$ par $\tilde\pi_n^\bullet$ défini par
$$
\tilde\pi_{n+1}^\bullet = \pi_{n+1}^\bullet + d\tilde h+\tilde hd.
$$
Avec ce choix de $\{\pi_n^\bullet\}$, les diagrammes ci-dessus commutent et les
morphismes sont compatibles et définissent un endomorphisme $\pi_\infty^\bullet =
\lim_n\pi_n^\bullet$ de $K_\infty^\bullet$. La partie {\it i} est alors claire :
les éléments $t_n = \sum(-1)^i \text{Tr}\left(\pi_n^i |_{K_n^i}\right)$
définissent un élément $t_\infty$ de $\mathbf{Z}_\ell$. De plus,
\begin{align*}
t_\infty
& =
\sum (-1)^i \text{Tr}_{\mathbf{Z}_\ell}(\pi_\infty^i |_{K_\infty^i}) \\
& =
\sum (-1)^i \text{Tr}_{\mathbf{Q}_\ell}(
\pi_\infty^i |_{K_\infty^i \otimes_{\mathbf{Z}_\ell}\mathbf{Q}_\ell}) \\
& =
\sum (-1)^i \text{Tr}(
\pi_\infty |_{H^i(K_\infty^\bullet \otimes \mathbf{Q}_\ell)})
\end{align*}
où la dernière égalité découle du fait que $\mathbf{Q}_\ell$ est un
corps, de sorte que le complexe $K_\infty^\bullet \otimes \mathbf{Q}_\ell$ est
quasi-isomorphe à sa cohomologie
$H^i(K_\infty^\bullet \otimes \mathbf{Q}_\ell)$. Celle-ci est également égale à
$H^i(K_\infty^\bullet)\otimes_{\mathbf{Z}}\mathbf{Q}_\ell = H_\infty^i \otimes
\mathbf{Q}_\ell$, ce qui achève la démonstration du lemme ainsi que celle du
théorème \ref{trace-theorem-C}.
\end{proof}




\section{Liste des compléments à apporter ci-dessus}
\label{trace-section-list-skipped}

\noindent
Qu'avons-nous omis de démontrer jusqu'ici dans les cours ?
\begin{enumerate}
\item les courbes et leurs jacobiennes,
\item le théorème de changement de base propre,
\item une présentation insuffisante de $R\Gamma_c$,
\item plus généralement, étant donné un morphisme séparé de type fini $f : X \to S$,
avec $S$ quasi-projectif, une présentation de $Rf_!$ sur les faisceaux \'etales.
\item une présentation de $\otimes^{\mathbf{L}}$
\item une explication des raisons pour lesquelles $R\Gamma_c$ commute à $\otimes^{\mathbf{L}}$
\end{enumerate}






%11.24.09
\section{Exemples de fonctions L}
\label{trace-section-examples-L-functions}

\noindent
Nous utilisons le théorème \ref{trace-theorem-B} pour les courbes afin de donner des exemples de fonctions $L$.





\section{Faisceaux constants}
\label{trace-section-L-function-constant-sheaf}

\noindent
Soit $k$ un corps fini, $X$ une courbe lisse géométriquement irréductible sur
$k$ et $\mathcal{F} = \underline{\mathbf{Q}_\ell}$ le faisceau constant. Si
$\bar x$ est un point géométrique de $X$, le module galoisien
$\mathcal{F}_{\bar x} = \mathbf{Q}_\ell$ est trivial, donc
$$
\det(1-\pi_x^*\ T^{\deg x} |_{\mathcal{F}_{\bar x}})^{-1} =
\frac{1}{1-T^{\deg x}}.
$$
En appliquant le théorème \ref{trace-theorem-B}, nous obtenons
\begin{align*}
L(X, \mathcal{F})
& =
\prod_{i = 0}^2
\det(1 - \pi_X^*T |_{H_c^i(X_{\bar k}, \mathbf{Q}_\ell)})^{(-1)^{i+1}} \\
& =
\frac{\det(1 - \pi_X^*T |_{H_c^1(X_{\bar k}, \mathbf{Q}_\ell)})}{
\det(1 - \pi_X^*T |_{H_c^0(X_{\bar k}, \mathbf{Q}_\ell)})
\cdot \det(1 - \pi_X^*T |_{H_c^2(X_{\bar k}, \mathbf{Q}_\ell)})}.
\end{align*}
Pour calculer ce dernier, distinguons deux cas.


\medskip\noindent
{\bf Cas projectif.}
Supposons que $X$ soit projective, de sorte que $H_c^i(X_{\bar k}, \mathbf{Q}_\ell) =
H^i(X_{\bar k}, \mathbf{Q}_\ell)$, et nous avons
$$
H^i(X_{\bar k}, \mathbf{Q}_\ell) =
\left\{
\begin{matrix}
\mathbf{Q}_\ell & \pi_X^* = 1 & \text{si }i = 0, \\
\mathbf{Q}_\ell^{2g} & \pi_X^* = ? & \text{si }i = 1, \\
\mathbf{Q}_\ell & \pi_X^* = q & \text{si }i = 2.
\end{matrix}
\right.
$$
L'identification de l'action de $\pi_X^*$ sur $H^2$ résulte du
Cohomologie \'etale, lemme \ref{etale-cohomology-lemma-pullback-on-h2-curve}
et du fait que le degré
de $\pi_X$ est $q = \#(k)$.
Nous savons peu de choses sur l'action de $\pi_X^*$ sur la cohomologie de degré 1.
Appelons $\alpha_1, \ldots, \alpha_{2g}$ ses valeurs propres dans
$\bar{\mathbf{Q}}_\ell$. En rassemblant ces informations,
le théorème \ref{trace-theorem-B}
donne l'égalité
$$
\prod\nolimits_{x \in |X|} \frac{1}{1 - T^{\deg x}} =
\frac{\det(1- \pi_X^* T|_{H^1(X_{\bar k}, \mathbf{Q}_\ell)})}{(1-T)(1-qT)} =
\frac{(1 - \alpha_1 T) \ldots (1 - \alpha_{2g}T)}{(1-T)(1-qT)}
$$
d'où nous déduisons le résultat suivant.

\begin{lemma}
\label{trace-lemma-count-points-projective}
Soit $X$ une courbe lisse, projective et géométriquement irréductible
sur un corps fini $k$. Alors
\begin{enumerate}
\item la fonction $L$ notée $L(X, \mathbf{Q}_\ell)$ est rationnelle,
\item les valeurs propres $\alpha_1, \ldots, \alpha_{2g}$ de $\pi_X^*$ sur
$H^1(X_{\bar k}, \mathbf{Q}_\ell)$ sont des entiers algébriques
indépendants de $\ell$,
\item le nombre de points de $X$ rationnels sur $k_n$, où $[k_n : k] = n$, est
$$
\# X(k_n) = 1 - \sum\nolimits_{i = 1}^{2g}\alpha_i^n + q^n,
$$
\item pour tout $i$, $|\alpha_i| < q$.
\end{enumerate}
\end{lemma}

\begin{proof}
La partie (3) est le théorème \ref{trace-theorem-C} appliqué à $\mathcal{F} =
\underline{\mathbf{Q}_\ell}$ sur $X \otimes k_n$. Pour la partie (4), utilisons le
résultat suivant.
\end{proof}

\begin{exercise}
\label{trace-exercise-powers}
Soient $\alpha_1, \ldots, \alpha_n \in \mathbf{C}$. Alors, pour tout secteur angulaire
contenant l'axe réel positif, de la forme $C_\varepsilon = \{ z \in
\mathbf{C} \ | \ |\arg z| < \varepsilon \}$ avec $\varepsilon > 0$, il existe
un entier $k \geq 1$ tel que $\alpha_1^k, \ldots, \alpha_n^k \in
C_\varepsilon$.
\end{exercise}

\noindent
Démontrer alors que $|\alpha_i| \leq q$ pour tout $i$. Puis utiliser des
considérations élémentaires sur les nombres complexes pour démontrer (comme dans la preuve du théorème
des nombres premiers) que $|\alpha_i| < q$. En fait, l'hypothèse de Riemann affirme que
l'on a $|\alpha_i| = \sqrt{q}$ pour tout $i$. Nous y reviendrons plus loin.

\medskip\noindent
{\bf Cas affine.}
Supposons maintenant que $X$ soit affine, disons $X= \bar X-\left\{x_1, \ldots, x_n\right\}$
où $j : X \hookrightarrow \bar X$ est une compactification projective non singulière.
Alors $H_c^0(X_{\bar k}, \mathbf{Q}_\ell) = 0$ et $H_c^2(X_{\bar k},
\mathbf{Q}_\ell) = H^2(\bar X_{\bar k}, \mathbf{Q}_\ell)$, de sorte que
le théorème \ref{trace-theorem-B}
s'écrit
$$
L(X, \mathbf{Q}_\ell) = \prod_{x \in |X|}\frac{1}{1 - T^{\deg x}} =
\frac{\det(1-\pi_X^*T |_{H_c^1(X_{\bar k}, \mathbf{Q}_\ell)})}{1 - qT}.
$$
D'autre part, le cas précédent donne
\begin{eqnarray*}
L(X, \mathbf{Q}_\ell) & = & L(\bar X,
\mathbf{Q}_\ell)\prod_{i = 1}^n\left(1-T^{\deg x_i}\right) \\
& = & \frac{\prod_{i = 1}^n(1-T^{\deg
x_i})\prod_{j = 1}^{2g}(1-\alpha_jT)}{(1-T)(1-qT)}.
\end{eqnarray*}
Par conséquent, nous voyons que $\dim H_c^1(X_{\bar k}, \mathbf{Q}_\ell) =
2g+\sum_{i = 1}^n \deg(x_i)-1$, et que les valeurs propres $\alpha_1, \ldots,
\alpha_{2g}$ de $\pi_{\bar X}^*$ agissant sur la cohomologie de degré 1 sont des racines de
l'unité. Plus précisément, chaque $x_i$ fournit un système complet de racines $\deg(x_i)$-ièmes
de l'unité, et l'on omet une occurrence de 1. Pour le voir directement à l'aide des
faisceaux cohérents, considérons la suite exacte courte sur $\bar X$
$$
0\to j_!\mathbf{Q}_\ell\to \mathbf{Q}_\ell\to\bigoplus_{i = 1}^n
\mathbf{Q}_{\ell, x_i}\to 0.
$$
La suite exacte longue de cohomologie s'écrit
$$
0\to \mathbf{Q}_\ell \to \bigoplus_{i = 1}^n \mathbf{Q}_\ell^{\oplus \deg x_i}
\to H_c^1(X_{\bar k}, \mathbf{Q}_\ell) \to H_c^1(\bar X_{\bar k},
\mathbf{Q}_\ell)\to 0
$$
où l'action du Frobenius sur $\bigoplus_{i = 1}^n \mathbf{Q}_\ell^{\oplus
\deg x_i}$ est donnée par la permutation cyclique de chaque terme ; par ailleurs, $H_c^2(X_{\bar k},
\mathbf{Q}_\ell) = H_c^2(\bar X_{\bar k}, \mathbf{Q}_\ell)$.






\section{La famille de Legendre}
\label{trace-section-legendre-family}

\noindent
Soit $k$ un corps fini de caractéristique impaire,
$X = \Spec(k[\lambda, \frac{1}{\lambda(\lambda - 1)}])$, et
considérons la famille de courbes elliptiques
$f : E \to X$ dans $\mathbf{P}^2_X$, dont l'équation affine est $y^2 =
x(x - 1)(x - \lambda)$. Posons $\mathcal{F} = Rf_*^1\mathbf{Q}_\ell =
\left\{R^1f_*\mathbf{Z}/\ell^n\mathbf{Z}\right\}_{n\geq 1} \otimes
\mathbf{Q}_\ell$. Dans cette situation, on a les propriétés suivantes :
\begin{itemize}
\item pour tout $n \geq 1$, le faisceau $R^1f_*(\mathbf{Z}/\ell^n\mathbf{Z})$ est
fini et localement constant -- en fait, il est libre de rang 2 sur
$\mathbf{Z}/\ell^n\mathbf{Z}$,
\item le système $\{R^1f_*\mathbf{Z}/\ell^n\mathbf{Z}\}_{n\geq 1}$ est un faisceau
$\ell$-adique lisse, et
\item pour tout $x\in |X|$,
$\det(1 - \pi_x\ T^{\deg x} |_{\mathcal{F}_{\bar x}}) =
(1 - \alpha_x T^{\deg x})(1 - \beta_x T^{\deg x})$
où $\alpha_x, \beta_x$ sont les valeurs propres du Frobenius géométrique
de $E_x$ agissant sur $H^1(E_{\bar x}, \mathbf{Q}_\ell)$.
\end{itemize}
Notons que $E_x$ n'est défini que sur $\kappa(x)$ et non sur $k$. La démonstration de
ces faits utilise le théorème de changement de base propre et l'acyclicité locale des
morphismes lisses. Pour plus de détails, voir \cite{SGA4.5}. Il en résulte que
$$
L(E/X) := L(X, \mathcal{F}) = \prod_{x\in |X|}
\frac{1}{(1-\alpha_xT^{\deg x})(1-\beta_xT^{\deg x })} .
$$
En appliquant le théorème \ref{trace-theorem-B}, nous obtenons
$$
L(E/X) =
\prod_{i = 0}^2
\det\left(1 - \pi_X^*T |_{H_c^i(X_{\bar k}, \mathcal{F})}\right)^{(-1)^{i+1}},
$$
et nous voyons en particulier qu'il s'agit d'une fonction rationnelle. De plus, il est
assez facile de montrer que $H_c^0(X_{\bar k}, \mathcal{F}) = H_c^2(X_{\bar
k}, \mathcal{F}) = 0$, de sorte que nous avons simplement
$$
L(E/X) = \det(1 - \pi_X^*T |_{H_c^1(X, \mathcal{F})}).
$$
Pour calculer explicitement ce déterminant, considérons la suite spectrale de Leray
du morphisme propre $f : E \to X$ à coefficients dans $\mathbf{Q}_\ell$, à savoir
$$
H_c^i(X_{\bar k}, R^jf_*\mathbf{Q}_\ell) \Rightarrow H_c^{i+j}(E_{\bar
k}, \mathbf{Q}_\ell)
$$
qui dégénère. Nous avons $f_*\mathbf{Q}_\ell = \mathbf{Q}_\ell$ et
$R^1f_*\mathbf{Q}_\ell = \mathcal{F}$. Le faisceau $R^2f_*\mathbf{Q}_\ell =
\mathbf{Q}_\ell(-1)$ est le {\it twist à la Tate} de $\mathbf{Q}_\ell$, c'est-à-dire
qu'il s'agit du faisceau $\mathbf{Q}_\ell$ sur lequel l'action de Galois est donnée par
la multiplication par $\#\kappa(x)$ dans la fibre en $\bar x$. Il s'ensuit que,
pour tout $n\geq 1$,
\begin{align*}
\# E(k_n)
& =
\sum\nolimits_i (-1)^i
\text{Tr}({\pi_E^n}^* |_{H_c^i(E_{\bar k}, \mathbf{Q}_\ell)}) \\
& =
\sum\nolimits_{i, j} (-1)^{i+j}
\text{Tr}({\pi^n_X}^* |_{H_c^i(X_{\bar k}, R^jf_*\mathbf{Q}_\ell)}) \\
& =
(q^n - 2) +
\text{Tr}({\pi_X^n}^* |_{H_c^1(X_{\bar k}, \mathcal{F})}) + q^n(q^n - 2) \\
& =
q^{2n} - q^n - 2 +
\text{Tr}({\pi_X^n}^* |_{H_c^1(X_{\bar k}, \mathcal{F})})
\end{align*}
où la première égalité résulte du
théorème \ref{trace-theorem-C},
la deuxième de
la suite spectrale de Leray et la troisième du calcul explicite des images directes supérieures
de $\mathbf{Q}_\ell$ par $f$. On pourrait aussi écrire
$$
\# E(k_n) = \sum_{x \in X(k_n)} \# E_x(k_n)
$$
et appliquer la formule des traces à chaque courbe. On peut aussi déterminer le nombre de
points rationnels sur $k_n$ par un simple comptage. La section nulle fournit $q^n -2$ points
(on omet les points où $\lambda = 0, 1$), d'où
$$
\# E(k_n) =
q^n - 2 + \# \{y^2 = x(x - 1)(x - \lambda), \lambda\neq 0, 1\}.
$$
Nous avons maintenant
$$
\begin{matrix}
\# \{y^2 = x(x - 1)(x - \lambda),\ \lambda\neq 0, 1\} \\
\\
\quad =
\# \{y^2 = x(x - 1)(x - \lambda)\text{ dans }\mathbf{A}^3\}
- \# \{y^2 = x^2(x - 1)\} - \# \{y^2 = x(x - 1)^2\}\\
\\
\quad = \# \{\lambda = \frac{-y^2}{x(x - 1)} + x,\ x\neq 0, 1\} +
\# \{y^2 = x(x - 1)(x - \lambda), x = 0, 1\} - 2(q^n - \varepsilon_n) \\
\\
\quad = q^n(q^n - 2)+2q^n - 2(q^n - \varepsilon_n)\\
\\
\quad = q^{2n}-2q^n+2\varepsilon_n
\end{matrix}
$$
où $\varepsilon_n = 1$ si $-1$ est un carré dans $k_n$, et 0 sinon,
c'est-à-dire
$$
\varepsilon_n = \frac{1}{2}\left(1+\left(\frac{-1}{k_n}\right)\right) =
\frac{1}{2}\left(1+(-1)^{\frac{q^n - 1}{2}}\right).
$$
Ainsi, $ \# E(k_n) = q^{2n} - q^n - 2+ 2\varepsilon_n$.
En comparant avec la formule précédente, on obtient
$$
\text{Tr}({\pi_X^n}^* |_{H_c^1(X_{\bar k}, \mathcal{F})}) =
2 \varepsilon_n = 1 + (-1)^{\frac{q^n - 1}{2}},
$$
ce qui implique, par un calcul élémentaire dans les nombres complexes, que si $-1$ est un
carré dans $k_n^*$, alors $\dim H_c^1(X_{\bar k}, \mathcal{F}) = 2$ et les
valeurs propres sont $1$ et $1$. Par conséquent, dans ce cas, nous avons
$$
L(E/X) = (1 - T)^2.
$$




\section{Sommes exponentielles}
\label{trace-section-exponential-sums}

\noindent
Un problème classique de théorie des nombres consiste à évaluer des sommes de la forme
$$
S_{a, b}(p) = \sum_{x\in \mathbf{F}_p - \left\{0, 1\right\}} e^{\frac{2\pi
ix^a(x - 1)^b}{p}}.
$$
Dans notre contexte, on peut l'interpréter comme une somme cohomologique de la manière suivante.
Considérons le schéma de base
$S = \Spec(\mathbf{F}_p[x, \frac{1}{x(x - 1)}])$ et la courbe affine
$f : X \to \mathbf{P}^1-\{0, 1, \infty\}$ sur $S$ donnée par l'équation
$y^{p - 1} = x^a(x - 1)^b$. C'est un revêtement galoisien fini \'etale de groupe
$\mathbf{F}_p^*$ et il existe une décomposition
$$
f_*(\bar{\mathbf{Q}}_\ell^*) =
\bigoplus_{\chi : \mathbf{F}_p^*\to \bar{\mathbf{Q}}_\ell^*} \mathcal{F}_\chi
$$
où $\chi$ parcourt les caractères de $\mathbf{F}_p^*$ et
$\mathcal{F}_\chi$ est un faisceau lisse de rang 1 de $\mathbf{Q}_\ell$-modules sur lequel
$\mathbf{F}_p^*$ agit par $\chi$ sur les fibres. On obtient une décomposition correspondante
$$
H_c^1(X_{\bar k}, \mathbf{Q}_\ell) = \bigoplus_\chi H^1(\mathbf{P}_{\bar
k}^1-\{0, 1, \infty\}, \mathcal{F}_\chi)
$$
et l'interprétation cohomologique de la somme exponentielle est donnée par la
formule des traces appliquée à $\mathcal{F}_\chi$ sur $\mathbf{P}^1 - \{0, 1,
\infty\}$ pour un $\chi$ convenable. Elle s'écrit
$$
S_{a, b}(p) =
-\text{Tr}(\pi_X^*
|_{H^1(\mathbf{P}_{\bar k}^1-\{0, 1, \infty\}, \mathcal{F}_\chi)}).
$$
Le yoga général de Weil suggère qu'il devrait y avoir une certaine compensation dans la
somme. En appliquant (approximativement) la formule de Riemann--Hurwitz, nous voyons que
$$
2g_X-2 \approx -2 (p-1) + 3(p-2) \approx p
$$
donc $g_X\approx p/2$, ce qui suggère également que les composantes en $\chi$ sont petites.



%12.01.09
\section{Formule des traces en termes de groupes fondamentaux}
\label{trace-section-trace-formula-fundamental-group}

\noindent
Dans les sections suivantes, nous reformulons entièrement la formule des traces
en termes du groupe fondamental d'une courbe, sauf lorsque la courbe
se trouve être $\mathbf{P}^1$.




\section{Groupes fondamentaux}
\label{trace-section-fundamental-groups}

\noindent
Ce matériau est présenté plus en détail dans le chapitre consacré aux
groupes fondamentaux. Voir
Groupes fondamentaux, section \ref{pione-section-introduction}.
Soit $X$ un schéma connexe et soit $\overline{x}\to X$ un
point géométrique. Considérons le foncteur
$$
\begin{matrix}
F_{\overline{x}}: &
\text{ schémas finis \'etales } \atop \text{ sur } X &
\longrightarrow & \text{ ensembles finis} \\
&
Y/X &
\longmapsto &
F_{\overline{x}}(Y) =
\left\{\text{ points géométriques }\overline y \atop \text{ de } Y
\text{ au-dessus de }\overline{x}\right\} = Y_{\overline{x}}
\end{matrix}
$$
Posons
$$
\pi_1(X, \overline{x})
=
Aut(F_{\overline{x}})
=
\text{ ensemble des automorphismes du foncteur }F_{\overline{x}}
$$
Notons que, pour tout morphisme fini \'etale $Y \to X$, il existe une action
$$
\pi_1(X, \overline{x}) \times F_{\overline{x}}(Y) \to F_{\overline{x}}(Y)
$$

\begin{definition}
\label{trace-definition-open}
Un sous-groupe de la forme
$\text{Stab}(\overline y\in F_{\overline{x}}(Y))\subset \pi_1(X, \overline{x})$
est dit {\it ouvert}.
\end{definition}

\begin{theorem}[Grothendieck]
\label{trace-theorem-fundamental-group}
Soit $X$ un schéma connexe.
\begin{enumerate}
\item Il existe une topologie sur $\pi_1(X, \overline{x})$ telle que les sous-groupes
ouverts forment un système fondamental de voisinages ouverts de $e\in \pi_1(X, \overline
x)$.
\item Muni de la topologie de (1), le groupe
$\pi_1(X, \overline{x})$ est un groupe profini.
\item Le foncteur
$$
\begin{matrix}
\text{ schémas finis } \atop \text{ \'etales sur }X & \to &
\text{ ensembles finis discrets à action continue } \atop \pi_1(X, \overline{x})\text{-ensembles}\\
Y / X& \mapsto & F_{\overline{x}}(Y) \text{ muni de son action naturelle}
\end{matrix}
$$
est une équivalence de catégories.
\end{enumerate}
\end{theorem}

\begin{proof}
Cf. \cite{SGA1}.
\end{proof}

\begin{proposition}
\label{trace-proposition-integral-normal-fundamental-group}
Soit $X$ un schéma noethérien normal intègre. Soit
$\overline y\to X$ un point géométrique algébrique situé
au-dessus du point générique $\eta\in X$. Alors
$$
\pi_x(X, \overline \eta) = Gal(M/\kappa(\eta))
$$
($\kappa(\eta)$ est le corps des fonctions de $X$), où
$$
\kappa(\overline \eta)\supset M\supset \kappa(\eta) = k(X)
$$
est la sous-extension maximale telle que, pour toute sous-extension finie
$M\supset L\supset \kappa(\eta)$, la normalisation de $X$ dans $L$ soit finie
\'etale sur $X$.
\end{proposition}

\begin{proof}
Omis.
\end{proof}

\noindent
{\bf Changement de point base.} Pour tous points géométriques $\overline{x}_1, \overline{x}_2$
de $X$, il existe un isomorphisme des foncteurs fibres
$$
\mathcal{F}_{\overline{x}_1} \cong \mathcal{F}_{\overline{x}_2}
$$
(C'est un chemin de $\overline{x}_1$ à $\overline{x}_2$.) La conjugaison
par ce chemin donne un isomorphisme
$$
\pi_1(X, \overline{x}_1) \cong \pi_1(X, \overline{x}_2)
$$
bien défini à automorphisme intérieur près.

\medskip\noindent
{\bf Fonctorialité.} Pour tout morphisme $X_1\to X_2$ de schémas connexes
et tout $\overline{x}\in X_1$, il existe un morphisme canonique
$$
\pi_1(X_1, \overline{x}) \to \pi_1(X_2, \overline{x})
$$
(Pourquoi ? Parce que le foncteur fibre ...)

\medskip\noindent
{\bf Corps de base.} Soit $X$ une variété sur un corps $k$. On obtient alors
$$
\pi_1(X, \overline{x}) \to
\pi_1(Spec(k), \overline{x}) =^{\text{prop}} Gal(k^{\text{sep}}/k)
$$
Ce morphisme est surjectif si et seulement si $X$ est géométriquement connexe sur $k$.
Dans le cas géométriquement connexe, nous obtenons donc une suite exacte courte de groupes
profinis
$$
1 \to \pi_1(X_{\overline{k}}, \overline{x}) \to
\pi_1(X, \overline{x}) \to
Gal(k^{\text{sep}}/k) \to 1
$$
($\pi_1(X_{\overline{k}}, \overline{x})$ : groupe fondamental géométrique de
$X$, $\pi_1(X, \overline{x})$ : groupe fondamental arithmétique de $X$)

\medskip\noindent
{\bf Comparaison.} Si $X$ est une variété sur $\mathbf{C}$, alors
$$
\pi_1(X, \overline{x}) =
\text{ complété profini de }
\pi_1(X(\mathbf{C})(\text{ topologie usuelle}), x)
$$
(où $x\in X(\mathbf{C})$)

\medskip\noindent
{\bf Frobenius.} Soit $X$ une variété sur $k$, avec $\# k < \infty$. Pour tout $x \in X$
point fermé, soit
$$
F_x\in \pi_1(x, \overline{x}) =
\text{Gal}(\kappa(x)^{\text{sep}}/\kappa(x))
$$
le Frobenius géométrique.
Soit $\overline\eta$ un point géométrique algébrique au-dessus du point générique. Alors
$$
\pi_1(X, \overline\eta) \leftarrow^{\cong}
\pi_1(X, \overline{x})
{\text{fonctorialité} \atop \leftarrow} \pi_1(x, \overline{x})
$$

\noindent
Fait élémentaire :
$$
\begin{matrix}
\pi_1(X, \overline \eta) & \to^{\deg} \pi_1(\Spec(k), \overline \eta) * &
= Gal(k^{sep}/k) \\
& & || \\
& & \widehat{\mathbf{Z}}\cdot F_{\Spec(k)} \\
F_x & \mapsto & \deg(x)\cdot F_{\Spec(k)}
\end{matrix}
$$
Rappelons que $\deg(x) = [\kappa(x):k]$.

\medskip\noindent
{\bf Groupes fondamentaux et faisceaux lisses.}
Soit $X$ un schéma connexe et $\overline{x}$ un point géométrique. Il existe des
équivalences de catégories
\begin{center}
\begin{tabular}{@{}c@{}}
\((\Lambda\text{ anneau fini})\)\\[2pt]
\begin{tabular}{@{}c@{\quad$\leftrightarrow$\quad}c@{}}
\parbox[c]{0.38\textwidth}{\centering
faisceaux de \(\Lambda\)-modules localement constants constructibles sur
\(X_\etale\)}
&
\parbox[c]{0.38\textwidth}{\centering
modules finis (discrets) sur \(\Lambda\) munis d'une action continue de
\(\pi_1(X, \overline{x})\)}
\end{tabular}\\[10pt]
\((\ell\text{ premier})\)\\[2pt]
\begin{tabular}{@{}c@{\quad$\leftrightarrow$\quad}c@{}}
\parbox[c]{0.38\textwidth}{\centering
faisceaux \(\ell\)-adiques lisses}
&
\parbox[c]{0.38\textwidth}{\centering
un \(\mathbf{Z}_\ell\)-module \(M\) de type fini muni d'une action continue
de \(\pi_1(X, \overline{x})\) ; on utilise la topologie \(\ell\)-adique sur
\(M\)}
\end{tabular}
\end{tabular}
\end{center}
En particulier, les faisceaux lisses de $\mathbf{Q}_l$-modules correspondent aux
homomorphismes continus
$$
\pi_1(X, \overline{x}) \to \text{GL}_r(\mathbf{Q}_l), \quad r\geq 0
$$

\noindent
Notation : un module muni d'une action $(M, \rho)$ correspond au faisceau
$\mathcal{F}_\rho$.

\medskip\noindent
{\bf Formules des traces.} Soit $X$ une variété sur $k$, avec $\# k < \infty$.
\begin{enumerate}
\item $\Lambda$ est un anneau fini tel que $(\# \Lambda, \# k)=1$
$$
\rho : \pi_1(X, \overline{x})\to \text{GL}_r(\Lambda)
$$
continue. Pour tout $n\geq 1$, on a
$$
\sum_{d|n}d\left(
\sum_{x\in |X|, \atop \deg(x)=d}
\text{Tr}(\rho(F_x^{n/d}))\right) =
\text{Tr}\left(
(\pi_x^n)^* |_{R\Gamma_c(X_{\overline{k}}, \mathcal{F}_\rho)}\right)
$$
\item Soit $l\neq char(k)$ un nombre premier et soit $\rho : \pi_1(X, \overline{x})\to
\text{GL}_r(\mathbf{Q}_l)$. Pour tout $n\geq 1$,
$$
\sum_{d|n} d\left(
\sum_{x\in |X| \atop \deg(x)=d}
\text{Tr}
\left(
\rho(F_x^{n/d})
\right)
\right) =
\sum_{i = 0}^{2\dim X}
(-1)^i
\text{Tr}\left(
\pi_X^* |_{H_c^i(X_{\overline{k}}, \mathcal{F}_\rho)}\right)
$$
\end{enumerate}

\noindent
{\bf Conjectures de Weil.} (Deligne--Weil I, 1974) Si $X$ est lisse et projective sur $k$ et
$\# k = q$, alors les valeurs propres de $\pi_X^*$ sur $H^i(X_{\overline{k}},
\mathbf{Q}_l)$ sont des entiers algébriques $\alpha$ tels que $|\alpha|=q^{i/2}$.

\medskip\noindent
{\bf Conjectures de Deligne.} (presque entièrement démontrées par
Lafforgue et d'autres, $\ldots$) Soit $X$ une variété normale sur le corps fini $k$ et soit
$$
\rho : \pi_1(X, \overline{x}) \longrightarrow \text{GL}_r(\mathbf{Q}_l)
$$
continue. Supposons $\rho$ irréductible et $\det(\rho)$ d'ordre fini. Alors :
\begin{enumerate}
\item il existe un corps de nombres $E$ tel que, pour tout $x\in
|X|$ (point fermé), le polynôme caractéristique de $\rho(F_x)$ ait ses coefficients dans $E$ ;
\item pour tout $x\in |X|$, les valeurs propres $\alpha_{x, i}$, $i = 1, \ldots,
r$, de $\rho(F_x)$ ont toutes une valeur absolue complexe égale à $1$
(ce sont des nombres algébriques, pas nécessairement des entiers) ;
\item pour toute place finie $\lambda$ de $E$ ne divisant pas $p$
(quitte à agrandir légèrement $E$), il existe
$$
\rho\lambda : \pi_1(X, \overline{x}) \to \text{GL}_r(E_\lambda)
$$
compatible avec $\rho$ (les polynômes caractéristiques des $F_x$ sont les mêmes).
\end{enumerate}

\begin{theorem}[Deligne, Weil II]
\label{trace-theorem-weil-II}
Pour un faisceau
$\mathcal{F}_\rho$ tel que $\rho$ satisfasse les conclusions de la conjecture
ci-dessus, les valeurs propres de $\pi_X^*$ sur $H_c^i(X_{\overline{k}},
\mathcal{F}_{\rho})$ sont des nombres algébriques $\alpha$ dont les valeurs absolues
$$
|\alpha|=q^{w/2}, \text{ pour }w\in \mathbf{Z},\ w\leq i
$$
De plus, si $X$ est lisse et projective, alors $w = i$.
\end{theorem}

\begin{proof}
Voir \cite{WeilII}.
\end{proof}




%12.03.09
\section{Groupes profinis, cohomologie et homologie}
\label{trace-section-profinite-cohomology}

\noindent
Soit $G$ un groupe profini.

\medskip\noindent
{\bf Cohomologie.}
Considérons la catégorie des modules discrets munis d'une action continue de $G$.
Cette catégorie possède assez d'injectifs, et l'on peut définir
$$
H^i(G, M) = R^iH^0(G, M) = R^i(M\mapsto M^G)
$$
Il existe aussi une version dérivée $RH^0(G, -)$.

\medskip\noindent
{\bf Homologie.}
Considérons la catégorie des groupes abéliens compacts munis d'une action continue de $G$.
Cette catégorie possède assez de projectifs, et l'on peut définir
$$
H_i(G, M) = L_iH_0(G, M)=L_i(M\mapsto M_G)
$$
et il existe également une version dérivée.

\medskip\noindent
{\bf Dualité élémentaire.}
Le foncteur $M\mapsto M^\wedge = \Hom_{cont}(M, S^1)$
échange les catégories précédentes, et
$$
H^i(G, M)^\wedge = H_i(G, M^\wedge)
$$
De plus, ce foncteur envoie les $G$-modules discrets de torsion sur les
$G$-modules profinis continus, et réciproquement ; si $M$ est un module continu
discret ou profini muni d'une action de $G$, alors
$M^\wedge = \Hom(M, \mathbf{Q}/\mathbf{Z})$.

\medskip\noindent
{\bf Notes sur l'homologie.}
\begin{enumerate}
\item Si l'on considère les $\Lambda$-modules pour un anneau fini $\Lambda$,
alors on peut identifier
$$
H_i(G, M)=Tor_i^{\Lambda[[G]]}(M, \Lambda)
$$
où $\Lambda[[G]]$ est la limite des algèbres de groupe des quotients finis
de $G$.
\item Si $G$ est un sous-groupe normal de $\Gamma$, et si $\Gamma$ est également
profini, alors
\begin{itemize}
\item $H^0(G, -)$ : $\Gamma$-modules discrets$\to$ modules discrets
munis d'une action de $\Gamma/G$
\item $H_0(G, -)$ : $\Gamma$-modules compacts $\to$ modules compacts
munis d'une action de $\Gamma/G$
\end{itemize}
et le groupe profini $\Gamma/G$ agit donc sur les groupes de cohomologie
de $G$ à valeurs dans un $\Gamma$-module. Autrement dit, il existe des foncteurs
dérivés
$$
\begin{aligned}
RH^0(G, -) :\quad
&D^{+}(\text{modules discrets }\Gamma\text{-équivariants})\\
&\longrightarrow
D^{+}(\text{modules discrets }\Gamma/G\text{-équivariants})
\end{aligned}
$$
et de même pour $LH_0(G, -)$.
\end{enumerate}








\section{Cohomologie des courbes, revisitée}
\label{trace-section-cohomology-curves-revisited}

\noindent
Soit $k$ un corps et soit $X$ une courbe lisse, géométriquement connexe sur $k$.
On dispose de la suite exacte fondamentale
$$
1 \to
\pi_1(X_{\overline{k}}, \overline \eta) \to
\pi_1(X, \overline\eta) \to
\text{Gal}(k^{^{sep}}/k) \to 1
$$
Si $\Lambda$ est un anneau fini tel que $\#\Lambda\in k^*$, si $M$ est un
$\Lambda$-module fini et si l'on se donne
$$
\rho : \pi_1(X, \overline\eta) \to \text{Aut}_{\Lambda}(M)
$$
continue, alors $\mathcal{F}_\rho$ désigne le faisceau associé sur
$X_\etale$.

\begin{lemma}
\label{trace-lemma-identify-h2c}
Il existe un isomorphisme canonique
$$
H_c^2(X_{\overline{k}}, \mathcal{F}_\rho)=(M)_{\pi_1(X_{\overline{k}},
\overline\eta)}(-1)
$$
en tant que $\text{Gal}(k^{^{sep}}/k)$-modules.
\end{lemma}

\noindent
Ici, l'indice ${}_{\pi_1(X_{\overline{k}}, \overline\eta)}$
désigne les coinvariants, et $(-1)$ le twist à la Tate, c'est-à-dire que
$\sigma\in \text{Gal}(k^{^{sep}}/k)$ agit par
$$
\chi_{cycl}(\sigma)^{-1}.\sigma\text{ sur le membre de droite}
$$
où
$$
\chi_{cycl} :
\text{Gal}(k^{^{sep}}/k)
\to
\prod\nolimits_{l\neq char(k)}\mathbf{Z}_l^*
$$
est le caractère cyclotomique.

\medskip\noindent
Reformulation (Deligne, Weil II, page 338). Pour tout faisceau fini localement
constant $\mathcal{F}$ sur $X$, il existe un quotient maximal $\mathcal{F}\to
\mathcal{F}''$ tel que $\mathcal{F}''/X_{\overline{k}}$ soit un faisceau constant ; ainsi
$$
\mathcal{F}'' = (X\to \Spec(k))^{-1}F''
$$
où $F''$ est un faisceau sur $\Spec(k)$, c'est-à-dire un
$\text{Gal}(k^{^{sep}}/k)$-module. Alors
$$
H_c^2(X_{\overline{k}}, \mathcal{F})\to H_c^2(X_{\overline{k}},
\mathcal{F}'')\to F''(-1)
$$
est un isomorphisme.

\begin{proof}[Démonstration du lemme \ref{trace-lemma-identify-h2c}]
Soit $Y\to^{\varphi}X$ le revêtement galoisien étale fini
correspondant à $\Ker(\rho) \subset \pi_1(X, \overline\eta)$. Ainsi
$$
\text{Aut}(Y/X)=Ind(\rho)
$$
est le groupe de Galois. Alors $\varphi^*\mathcal{F}_\rho =\underline M_Y$ et
$$
\varphi_*\varphi^*\mathcal{F}_\rho\to \mathcal{F}_\rho
$$
ce qui donne
\begin{align*}
& H_c^2(X_{\overline{k}}, \varphi_*\varphi^*\mathcal{F}_\rho) \to
H_c^2(X_{\overline{k}}, \mathcal{F}_\rho)\\
& =H_c^2(Y_{\overline{k}}, \varphi^*\mathcal{F}_\rho)\\
& =H_c^2(Y_{\overline{k}}, \underline M) = \oplus_{\text{composantes irréd.
de } \atop Y_{\overline{k}}}M
\end{align*}
$$
\Im(\rho) \to H_c^2(Y_{\overline{k}}, \underline M) =
\oplus_{\text{composantes irréd. de } \atop Y_{\overline{k}}}
M \to_{\Im(\rho) \text{équivariant}} H_c^2(X_{\overline{k}},
\mathcal{F}_{\rho}) \to^{\text{action }
\Im(\rho) \atop \text{triviale}}
$$
si $C/\overline{k}$ est une courbe irréductible, alors $H_c^2(C, \underline M)=M$.

\medskip\noindent
Puisque
$$
{\text{ensemble des composantes } \atop \text{irréductibles de }Y_k} =
\frac{Im(\rho)}{Im(\rho|_{\pi_1(X_{\overline{k}}, \overline \eta)})}
$$
On en conclut que $H_c^2(X_{\overline{k}}, \mathcal{F}_\rho)$ est un
quotient de $M_{\pi_1(X_{\overline{k}}, \overline \eta)}$. D'autre part,
il existe une surjection
$$
\mathcal{F}_\rho\to \mathcal{F}'' = {\text{ faisceau sur }
X\text{ associé à } \atop (M)_{\pi_1(X_{\overline{k}}, \overline
\eta)}\leftarrow\pi_1(X, \overline \eta)}
$$
$$
H_c^2(X_{\overline{k}}, \mathcal{F}_\rho)\to
M_{\pi_1(X_{\overline{k}}, \overline\eta)}
$$
L'action de Galois ainsi tordue provient du fait que
$H_c^2(X_{\overline{k}}, \mu_n)=^{\text{can}} \mathbf{Z}/n\mathbf{Z}$.
\end{proof}

\begin{remark}
\label{trace-remark-projective}
On en conclut donc que, si $X$ est en outre projective,
on dispose, fonctoriellement en la représentation $\rho$,
des identifications
$$
H^0(X_{\overline{k}}, \mathcal{F}_\rho) =
M^{\pi_1(X_{\overline{k}}, \overline\eta)}
$$
et
$$
H_c^2(X_{\overline{k}}, \mathcal{F}_\rho) =
M_{\pi_1(X_{\overline{k}}, \overline \eta)}(-1)
$$
Bien entendu, si $X$ n'est pas projective, alors
$H^0_c(X_{\overline{k}}, \mathcal{F}_\rho) = 0$.
\end{remark}


\begin{proposition}
\label{trace-proposition-curve-kpi1}
Soit $X/k$ comme précédemment, mais avec $X_{\overline{k}}\neq \mathbf{P}^1_{\overline{k}}$.
Les foncteurs
$
(M, \rho)\mapsto H_c^{2-i}(X_{\overline{k}}, \mathcal{F}_\rho)
$
sont les foncteurs dérivés à gauche de
$(M, \rho)\mapsto H_c^2(X_{\overline{k}}, \mathcal{F}_\rho)$
donc
$$
H_c^{2-i}(X_{\overline{k}}, \mathcal{F}_\rho) =
H_i(\pi_1(X_{\overline{k}}, \overline \eta), M)(-1)
$$
De plus, il existe une version dérivée, à savoir
$$
R\Gamma_c(X_{\overline{k}}, \mathcal{F}_\rho)
=
LH_0(\pi_1(X_{\overline{k}}, \overline \eta), M(-1))
=
M(-1)
\otimes_{\Lambda[[\pi_1(X_{\overline{k}}, \overline \eta)]]}^\mathbf{L}
\Lambda
$$
dans $D(\Lambda[[\widehat{\mathbf{Z}}]])$.
De même, les foncteurs
$(M, \rho)\mapsto H^i(X_{\overline{k}}, \mathcal{F}_\rho)$
sont les foncteurs dérivés à droite de
$(M, \rho)\mapsto M^{\pi_1(X_{\overline{k}}, \overline \eta)}$
donc
$$
H^i(X_{\overline{k}}, \mathcal{F}_\rho) =
H^i(\pi_1(X_{\overline{k}}, \overline \eta), M)
$$
De plus, il existe également une version dérivée dans ce cas.
\end{proposition}

\begin{proof}
(Idée) Montrer que les deux membres sont des $\delta$-foncteurs universels.
\end{proof}

\begin{remark}
\label{trace-remark-poincare-groups}
La proposition et la dualité élémentaire donnent alors
$$
H^{2-i}_c(X_{\overline{k}}, \mathcal{F}_\rho)
\times
H^i(X_{\overline{k}}, \mathcal{F}_\rho^\wedge(1))
\to
\mathbf{Q}/\mathbf{Z}
$$
un accouplement parfait. Si $X$ est projective, c'est la dualité de Poincaré.
\end{remark}





\section{Formule des traces abstraite}
\label{trace-section-abstract-trace-formula}

\noindent
Supposons donnée une extension de groupes profinis
$$
1 \to G \to \Gamma \xrightarrow{\deg} \widehat{\mathbf{Z}} \to 1
$$
On dit que $\Gamma$ {\it possède une formule des traces abstraite} si et seulement s'il
existe
\begin{enumerate}
\item un entier $q\geq 1$, et
\item pour tout $d\geq 1$, un ensemble fini $S_d$ et, pour chaque $x\in S_d$, une
classe de conjugaison $F_x \in \Gamma$ telle que $\deg(F_x) = d$
\end{enumerate}
de sorte que les conditions suivantes soient satisfaites :
\begin{enumerate}
\item pour tout $\ell$ ne divisant pas $q$, on a $\text{cd}_\ell(G)<\infty$, et
\item pour tout anneau fini $\Lambda$ tel que $q\in \Lambda^*$,
pour tout $\Lambda$-module projectif fini $M$ muni d'une action continue de
$\Gamma$, et pour tout $n > 0$, on a
$$
\sum\nolimits_{d|n}d \left(
\sum\nolimits_{x \in S_d}
\text{Tr}( F_x^{n/d} |_M)
\right)
=
q^n \text{Tr}(F^n|_{M \otimes_{\Lambda[[G]]}^{\mathbf{L}}\Lambda})
$$
dans $\Lambda^\natural$.
\end{enumerate}
Ici, $M \otimes_{\Lambda[[G]]}^{\mathbf{L}}\Lambda = LH_0(G, M)$ désigne
l'homologie dérivée, et $F=1$ dans $\Gamma/G = \widehat{\mathbf{Z}}$.

\begin{remark}
\label{trace-remark-abstract-trace-formula}
Voici quelques observations au sujet de cette notion.
\begin{enumerate}
\item Pour modéliser les courbes projectives, on peut utiliser la cohomologie et
le facteur $q^n$ n'est pas nécessaire.
\item Les seuls exemples que je connaisse sont $\Gamma = \pi_1(X, \overline \eta)$,
où $X$ est lisse, géométriquement irréductible et un $K(\pi, 1)$ sur un corps
fini. Dans ce cas, $q = (\# k)^{\dim X}$. Compte tenu de la proposition, nous avons démontré
ce résultat pour les courbes dans ce cours.
\item Une fois l'entier $q$ donné, les ensembles $S_d$ sont déterminés de manière
unique. (On peut multiplier $q$ par un entier $m$, puis remplacer $S_d$ par
$m^d$ copies de $S_d$ sans changer la formule.)
\end{enumerate}
\end{remark}

\begin{example}
\label{trace-example-commutative}
Fixons un entier $q\geq 1$.
$$
\begin{matrix}
1 &
\to &
G = \widehat{\mathbf{Z}}^{(q)} &
\to &
\Gamma &
\to &
\widehat{\mathbf{Z}} &
\to &
1 \\
&
&
= \prod_{l\not \mid q} \mathbf{Z}_l &
&
F &
\mapsto &
1
\end{matrix}
$$
avec $FxF^{-1} = ux$, où $u \in (\widehat{\mathbf{Z}}^{(q)})^*$.
En utilisant seulement les modules triviaux
$\mathbf{Z}/m\mathbf{Z}$, on voit que
$$
q^n - (qu)^n \equiv \sum\nolimits_{d|n} d\# S_d
$$
dans $\mathbf{Z}/m\mathbf{Z}$ pour tout $(m, q)=1$ (à
$u \to u^{-1}$ près) ; cela entraîne $qu = a\in \mathbf{Z}$
et $|a| < q$. Le cas particulier $a = 1$ se présente avec
$$
\Gamma = \pi_1^t(\mathbf{G}_{m, \mathbf{F}_p}, \overline \eta),
\quad
\# S_1 = q - 1,
\quad\text{et}\quad
\# S_2 = \frac{(q^2-1)-(q-1)}{2}
$$
\end{example}



%12.08.09
\section{Formes automorphes et faisceaux}
\label{trace-section-automorphic}

\noindent
Références : voir notamment les remarquables articles
\cite{D1}, \cite{D2} et \cite{D0} de Drinfeld.

\medskip\noindent
{\bf Formes cuspidales non ramifiées.}
Soit $k$ un corps fini de caractéristique $p$.
Soit $X$ une courbe projective lisse géométriquement irréductible sur $k$.
Posons $K = k(X)$, le corps des fonctions de $X$.
Soit $v$ une place de $K$, ce qui revient à se donner un
point fermé $x\in X$. Soit $K_v$ le complété de $K$ en $v$ ;
c'est aussi le corps des fractions du complété de
l'anneau local de $X$ en $x$.
Notons $O_v\subset K_v$ l'anneau des entiers. Posons en outre
$$
O = \prod\nolimits_v O_v \subset \mathbf{A} = \prod_v' K_v
$$
et soit $\Lambda$ un anneau, avec $p$ inversible dans $\Lambda$.

\begin{definition}
\label{trace-definition-unramified}
Une {\it forme cuspidale non ramifiée sur $\text{GL}_2(\mathbf{A})$ à valeurs dans
$\Lambda$}\footnote{Cette notation n'est probablement pas standard.}
est une fonction
$$
f : \text{GL}_2(\mathbf{A}) \to \Lambda
$$
telle que
\begin{enumerate}
\item $f(x\gamma) = f(x)$ pour tout $x\in \text{GL}_2(\mathbf{A})$ et tout
$\gamma\in \text{GL}_2(K)$
\item $f(ux) = f(x)$ pour tout $x\in \text{GL}_2(\mathbf{A})$ et tout
$u\in \text{GL}_2(O)$
\item pour tout $x\in \text{GL}_2(\mathbf{A})$,
$$
\int_{\mathbf{A} \mod K} f
\left(x
\left(
\begin{matrix}
1 & z \\
0 & 1
\end{matrix}
\right)
\right) dz = 0
$$
voir \cite[section 4.1]{dJ-conjecture}
pour une explication de la manière de donner un sens à cette expression
pour un anneau quelconque $\Lambda$ dans lequel $p$ est inversible.
\end{enumerate}
\end{definition}

\noindent
{\bf Opérateurs de Hecke.}
Pour une place $v$ de $K$ et une forme cuspidale non ramifiée $f$, posons
$$
T_v(f)(x) =
\int_{g\in M_v}f(g^{-1}x)dg,
$$
et
$$
U_v(f)(x) =
f\left(
\left(
\begin{matrix}
\pi_v^{-1} & 0 \\
0 & \pi_v^{-1}
\end{matrix}
\right)x\right)
$$
Notations : ici, $\pi_v \in O_v$ est une uniformisante,
$$
M_v =
\left\{
h\in Mat(2\times 2, O_v) | \det h = \pi_vO_v^*\right\}
$$
et $dg$ désigne la mesure de Haar sur $\text{GL}_2(K_v)$ telle que
$\int_{\text{GL}_2(O_v)} dg = 1$. Explicitement, on a
$$
T_v(f)(x) =
f\left(
\left(
\begin{matrix}
\pi_v^{-1}& 0 \\
0 & 1
\end{matrix}
\right)
x\right) +
\sum_{i = 1}^{q_v}
f\left(\left(
\begin{matrix}
1 & 0 \\
-\pi_v^{-1}\lambda_i
& \pi_v^{-1}
\end{matrix}
\right) x\right)
$$
où les $\lambda_i\in O_v$ forment un système de représentants de
$O_v/(\pi_v)=\kappa_v$, $q_v = \#\kappa_v$.

\medskip\noindent
{\bf Formes propres.} Une {\it forme propre} $f$ est une forme cuspidale non ramifiée
telle qu'une valeur de $f$ soit une unité, et telle que $T_vf = t_vf$ et
$U_vf = u_vf$ pour certains $t_v, u_v \in \Lambda$ déterminés de manière unique.

\begin{theorem}
\label{trace-theorem-drinfeld-make-rho}
Étant donnée une forme propre $f$ à valeurs dans
$\overline{\mathbf{Q}}_l$, dont les valeurs propres
$u_v\in \overline{\mathbf{Z}}_l^*$, il existe une représentation
$$
\rho : \pi_1(X)\to \text{GL}_2(E)
$$
continue et absolument irréductible, où
$E$ est une extension finie de $\mathbf{Q}_\ell$ contenue dans
$\overline{\mathbf{Q}}_l$, et cette représentation vérifie
$t_v = \text{Tr}(\rho(F_v))$, et
$u_v = q_v^{-1}\det\left(\rho(F_v)\right)$ pour toute place $v$.
\end{theorem}

\begin{proof}
Voir \cite{D0}.
\end{proof}

\begin{theorem}
\label{trace-theorem-drinfeld-make-f}
Supposons finie l'extension $\mathbf{Q}_l \subset E$, et soit
$$
\rho : \pi_1(X)\to \text{GL}_2(E)
$$
une représentation absolument irréductible et continue. Il existe alors une forme propre $f$ à
valeurs dans $\overline{\mathbf{Q}}_l$, dont les valeurs propres $t_v$, $u_v$
satisfont les égalités
$t_v = \text{Tr}(\rho(F_v))$ et $u_v = q_v^{-1}\det(\rho(F_v))$.
\end{theorem}

\begin{proof}
Voir \cite{D1}.
\end{proof}

\begin{remark}
\label{trace-remark-lafforgue}
Grâce à Lafforgue et à de nombreux autres mathématiciens, nous disposons désormais
d'énoncés complets analogues aux deux théorèmes précédents pour $\text{GL}_n$,
y compris avec ramification !
Autrement dit, la correspondance de Langlands globale complète pour $\text{GL}_n$
est connue pour les corps de fonctions de courbes sur des corps finis. Cela
ne signifie toutefois pas qu'il ne reste pas de nombreuses questions intéressantes
sur les groupes fondamentaux des courbes sur les corps finis, comme
nous le verrons ci-dessous.
\end{remark}

\noindent
{\bf Caractère central.} Si $f$ est une forme propre, alors
$$
\begin{matrix}
\chi_f : &
O^*\backslash \mathbf{A}^*/K^* &
\to &
\Lambda^* \\
&
(1, \ldots, \pi_v, 1, \ldots, 1) &
\mapsto &
u_v^{-1}
\end{matrix}
$$
est appelé le {\it caractère central}. Il correspond, compte tenu des
normalisations ci-dessus, au déterminant de $\rho$. Posons
$$
C(\Lambda) =
\left\{
{\text{formes cuspidales non ramifiées } f \text{ à coefficients dans }\Lambda}
\atop {\text{ telles que } U_v f = \varphi_v^{-1}f\forall v}
\right\}
$$

\begin{proposition}
\label{trace-proposition-cusp-forms-finite}
Si $\Lambda$ est noethérien, alors $C(\Lambda)$ est un
$\Lambda$-module de type fini. De plus, si $\Lambda$ est un corps et si
$\mathbf{F} \subset \Lambda$ est son sous-corps premier, alors
$$
C(\Lambda)=(C(\mathbf{F}))\otimes_{\mathbf{F}}\Lambda
$$
de manière compatible avec l'action de $T_v$.
\end{proposition}

\begin{proof}
Voir \cite[Proposition 4.7]{dJ-conjecture}.
\end{proof}

\noindent
Cette proposition entraîne immédiatement le lemme suivant.

\begin{lemma}
\label{trace-lemma-eigenvalues-algebraic}
Algébricité des valeurs propres.
Si $\Lambda$ est un corps, alors les valeurs propres $t_v$ de $f\in
C(\Lambda)$ sont algébriques sur le sous-corps premier
$\mathbf{F} \subset \Lambda$.
\end{lemma}

\begin{proof}
Cela résulte de la proposition \ref{trace-proposition-cusp-forms-finite}.
\end{proof}

\noindent
En combinant ce qui précède, on obtient l'astuce très utile suivante.

\begin{lemma}
\label{trace-lemma-switch-l}
Changement de $l$. Soit $E$ un corps de nombres.
Partons d'une représentation
$$
\rho : \pi_1(X)\to SL_2(E_\lambda)
$$
absolument irréductible et continue, où $\lambda$ est une place de $E$
ne divisant pas $p$. Pour toute autre place $\lambda'$ de $E$
ne divisant pas $p$, il existe une extension finie $E'_{\lambda'}$
et une représentation continue absolument irréductible
$$
\rho': \pi_1(X)\to SL_2(E'_{\lambda'})
$$
compatible avec $\rho$, au sens où les polynômes caractéristiques
de tous les Frobenius sont les mêmes.
\end{lemma}

\noindent
Notons qu'il s'agit d'un cas particulier de la conjecture de Deligne !

\begin{proof}
Pour démontrer le lemme de changement, on applique
le théorème \ref{trace-theorem-drinfeld-make-f}
pour obtenir une forme propre $f\in C(\overline{\mathbf{Q}}_l)$ associée à $\rho$.
On applique ensuite
la proposition \ref{trace-proposition-cusp-forms-finite}
pour voir que l'on peut choisir $f\in C(E')$, où $E \subset E'$ est une extension finie.
En complétant ensuite $E'$, on obtient
une forme propre $f\in C(E'_{\lambda'})$, où
$E'_{\lambda'}$ est une extension finie de $E_{\lambda'}$.
Enfin, on applique
le théorème \ref{trace-theorem-drinfeld-make-rho}
pour obtenir
$\rho': \pi_1(X) \to SL_2(E_{\lambda'}')$ absolument irréductible et continue,
quitte à agrandir encore légèrement $E'_{\lambda'}$.
\end{proof}

\noindent
Spéculation : si, pour un anneau (topologique) $\Lambda$, on a
$$
\left(
{\rho : \pi_1(X)\to SL_2(\Lambda) \atop \text{ absolument irréductible}}
\right)
\leftrightarrow
\text{ formes propres dans } C(\Lambda)
$$
alors toutes les valeurs propres de $\rho(F_v)$ sont algébriques (cet argument ne fonctionne
pas directement lorsque $\Lambda$ est un anneau fini). En partant de l'idée que la
correspondance de Langlands vaut plus généralement que pour les seuls corps,
on arrive à la conjecture suivante.

\medskip\noindent
{\bf Conjecture.}
(Voir \cite{dJ-conjecture}.)
Pour toute représentation continue
$$
\rho : \pi_1(X)\to \text{GL}_n(\mathbf{F}_l[[t]])
$$
on a $\# \rho(\pi_1(X_{\overline{k}}))<\infty$.

\medskip\noindent
Reformulation dans le langage des faisceaux :
« Pour tout faisceau lisse de $\overline{\mathbf{F}_l((t))}$-modules, la monodromie
géométrique est finie. »

\begin{theorem}
\label{trace-theorem-conjecture-n-2}
La conjecture est vraie si $n\leq 2$.
\end{theorem}

\begin{proof}
Voir \cite{dJ-conjecture}.
\end{proof}

\begin{theorem}
\label{trace-theorem-conjecture-l-bigger-2n}
La conjecture est vraie si $l > 2n$, sous réserve de certains résultats non démontrés.
\end{theorem}

\begin{proof}
Voir \cite{Gaitsgory}.
\end{proof}

\noindent
Il se trouve que la conjecture a une application.
Voir les travaux de Drinfeld sur les conjectures de Kashiwara. Il existe aussi
l'application beaucoup plus élémentaire suivante.

\begin{theorem}
\label{trace-theorem-deformation-rings}
(Voir \cite[Théorème 3.5]{dJ-conjecture}.)
Supposons que
$$
\rho_0: \pi_1(X)\to \text{GL}_n(\mathbf{F}_l)
$$
soit continue, avec $l\neq p$. Supposons :
\begin{enumerate}
\item que la conjecture soit vraie pour $X$ ;
\item que $\rho_0 |_{\pi_1(X_{\overline{k}})}$ soit absolument irréductible ; et
\item que $l$ ne divise pas $n$.
\end{enumerate}
Alors l'anneau de déformation universel $R_{\text{univ}}$ de $\rho_0$ est
fini et plat sur $\mathbf{Z}_l$.
\end{theorem}

\noindent
Explication : il existe une représentation $\rho_{\text{univ}}:
\pi_1(X)\to \text{GL}_n(R_{\text{univ}})$ (anneau de déformation universel),
où $R_{\text{univ}}$ est local et
complet, de corps résiduel $\mathbf{F}_l$, et $(R_{\text{univ}}\to
\mathbf{F}_l)\circ\rho_{\text{univ}}\cong\rho_0$.
Étant donnés $R\to \mathbf{F}_l$, avec $R$ local et complet, et
$\rho : \pi_1(X)\to \text{GL}_n(R)$, il existe alors
$\psi : R_{\text{univ}}\to R$ tel que
$\psi\circ\rho_{\text{univ}}\cong \rho$. Le théorème affirme que le morphisme
$$
\Spec(R_{\text{univ}})
\longrightarrow
\Spec(\mathbf{Z}_l)
$$
est fini et plat. En particulier, une telle représentation $\rho_0$
se relève en une représentation $\rho : \pi_1(X) \to \text{GL}_n(\overline{\mathbf{Q}}_l)$.

\medskip\noindent
Notes :
\begin{enumerate}
\item Le théorème sur les déformations est facile.
\item Tout progrès vers la conjecture semble difficile.
\item Il serait intéressant de disposer de davantage de conjectures sur $\pi_1(X)$ !
\end{enumerate}




%12.10.09
\section{Comptage des points}
\label{trace-section-counting}

\noindent
Soit $X$ une courbe lisse, géométriquement irréductible et
projective sur $k$, et posons $q = \# k$. La formule des traces donne ceci :
il existe des entiers algébriques $w_1, \ldots, w_{2g}$ tels que
$$
\# X(k_n) = q^n - \sum\nolimits_{i = 1}^{2g_X} w_i^n + 1.
$$
Si $\sigma\in \text{Aut}(X)$, alors, pour tout $i$, il existe $j$ tel que
$\sigma(w_i)=w_j$.

\medskip\noindent
{\bf Hypothèse de Riemann.} Pour tout $i$, on a $|\omega_i| = \sqrt{q}$.

\medskip\noindent
Cette hypothèse a été formulée par Emil Artin en 1924 pour les
courbes hyperelliptiques, puis démontrée par Weil en 1940. Weil en a donné deux démonstrations :
\begin{itemize}
\item l'une utilisant la théorie de l'intersection sur $X \times X$ et le
théorème de l'indice de Hodge ;
\item l'autre utilisant la jacobienne de $X$.
\end{itemize}
Il existe une autre démonstration, dont l'idée initiale est due à Stepanov et
que Bombieri a donnée : elle utilise le corps de fonctions $k(X)$ et
son opérateur de Frobenius (1969). Le point de départ est que, pour
$f\in k(X)$, la fonction rationnelle $f^q - f$
s'annule en tous les points $\mathbf{F}_q$-rationnels de $X$, et que l'on
peut essayer d'utiliser cette idée pour majorer le nombre de points.


\section{Forme précise du théorème de Chebotarev}
\label{trace-section-chebotarev}

\noindent
Comme première application, démontrons une forme précise du théorème de Chebotarev
pour un revêtement galoisien étale fini de courbes.
Soit $\varphi : Y \to X$ un revêtement galoisien étale fini de
groupe $G$. Il correspond à un homomorphisme
$$
\pi_1(X) \longrightarrow G = \text{Aut}(Y/X)
$$
Supposons $Y_{\overline{k}}$ irréductible. Si $C\subset G$ est une classe de conjugaison,
alors, pour tout $n > 0$, on a
$$
| \# \{x \in X(k_n) \mid F_x \in C\} - \frac{\# C}{\# G} \cdot \# X(k_n) |
\leq
(\# C)(2g - 2) \sqrt{q^n}
$$
(Avertissement : vérifier soigneusement le coefficient $\# C$ au membre de droite
avant d'utiliser cette estimation.)

\begin{proof}[Esquisse]
Écrivons
$$
\varphi_*(\overline{\mathbf{Q}_l}) =
\oplus_{\pi \in \widehat{G}} \mathcal{F}_{\pi}
$$
où $\widehat{G}$ est l'ensemble des classes d'isomorphisme
des représentations irréductibles de
$G$ sur $\overline{\mathbf{Q}}_l$. Pour $\pi \in \widehat{G}$,
soit $\chi_{\pi}: G \to \overline{\mathbf{Q}}_l$
le caractère de $\pi$. Alors
$$
H^*(Y_{\overline{k}}, \overline{\mathbf{Q}}_l) =
\oplus_{\pi\in \widehat{G}}
H^*(Y_{\overline{k}}, \overline{\mathbf{Q}}_l)_\pi
=_{(\varphi\text{ finie })}
\oplus_{\pi\in \widehat{G}}
H^*(X_{\overline{k}}, \mathcal{F}_\pi)
$$
Si $\pi\neq 1$, alors
$$
H^0(X_{\overline{k}}, \mathcal{F}_\pi) =
H^2(X_{\overline{k}}, \mathcal{F}_\pi) = 0,\quad
\dim H^1(X_{\overline{k}}, \mathcal{F}_\pi) = (2g_X - 2)d_\pi^2
$$
(cela peut se déduire de la formule des traces en faisant agir sur ...) et l'on voit que
$$
|\sum_{x \in X(k_n)} \chi_\pi(\mathcal{F}_x)| \leq
(2g_X - 2) d_\pi^2\sqrt{q^n}
$$
Écrivons $1_C = \sum_\pi a_\pi \chi_\pi$ ; alors
$a_\pi = \langle 1_C, \chi_\pi\rangle$, et
$a_1 = \langle 1_C, \chi_1\rangle = \frac{\# C}{\# G}$, où
$$
\langle f, h\rangle = \frac{1}{\# G}\sum_{g \in G} f(g)\overline{h(g)}
$$
On a donc la relation
$$
\frac{\# C}{\# G} = ||1_C||^2 = \sum|a_\pi|^2
$$
Dernière étape :
\begin{align*}
\#\left\{x \in X(k_n) \mid F_x \in C\right\}
& =
\sum_{x \in X(k_n)} 1_C(x) \\
& =
\sum_{x \in X(k_n)} \sum_\pi a_\pi \chi_\pi(F_x) \\
& =
\underbrace{\frac{\# C}{\# G} \# X(k_n)}_{
\text{terme pour }\pi = 1}
+
\underbrace{\sum_{\pi\neq 1}a_\pi\sum_{x\in X(k_n)}\chi_\pi(F_x)}_{
\text{ terme d'erreur (à majorer par }E)}
\end{align*}
On peut majorer le terme d'erreur par
\begin{align*}
|E|
& \leq
\sum_{\pi \in \widehat{G}, \atop \pi \neq 1}
|a_\pi| (2g - 2) d_\pi^2 \sqrt{q^n} \\
& \leq
\sum_{\pi \neq 1} \frac{\# C}{\# G} (2g_X - 2) d_\pi^3 \sqrt{q^n}
\end{align*}
D'après la conjecture de Weil, $\# X(k_n)\sim q^n$.
\end{proof}



\section{Combien de nombres premiers se décomposent complètement ?}
\sectionmark{Nombres premiers complètement décomposés}
\label{trace-section-how-many}

\noindent
Cette section donne une seconde application de l'hypothèse de Riemann pour les
courbes sur un corps fini. Les théoriciens des nombres pourront consulter
l'article d'Ihara intitulé
« Combien de nombres premiers se décomposent complètement dans une extension de Galois
non ramifiée infinie d'un corps global ? », voir \cite{Ihara}.
Considérons la suite exacte fondamentale
$$
1 \to
\pi_1(X_{\overline{k}}) \to
\pi_1(X) \xrightarrow{\deg}
\widehat{\mathbf{Z}} \to 1
$$

\begin{proposition}
\label{trace-proposition-finite-set-frobenii-generate-topologically}
Il existe des points fermés $x_1, \ldots, x_n$ de $X$
tels que l'ensemble de {\bf tous} les éléments de Frobenius correspondant à ces
points engendre topologiquement $\pi_1(X)$.
\end{proposition}

\noindent
On peut aussi formuler l'énoncé ainsi :
il existe $x_1, \ldots, x_n\in |X|$ tels que
le plus petit sous-groupe normal fermé $\Gamma$ de $\pi_1(X)$
contenant exactement $1$ élément de Frobenius pour chaque $x_i$ soit égal à $\pi_1(X)$, c'est-à-dire
$\Gamma = \pi_1(X)$.

\begin{proof}
Choisissons $N\gg 0$ et posons
$$
\{x_1, \ldots, x_n\} =
{\text{ ensemble de tous les points fermés de}
\atop X \text{ de degré} \leq N\text{ sur } k}
$$
Soit $\Gamma\subset \pi_1(X)$ comme dans la reformulation précédente pour ces
points. Supposons $\Gamma \neq \pi_1(X)$. On peut alors choisir un sous-groupe ouvert normal
$U$ de $\pi_1(X)$ contenant $\Gamma$ et tel que
$U \neq \pi_1(X)$. Par l'hypothèse de Riemann pour $X$, notre ensemble de points contient un
$x_{i_1}$ de degré $N$ et un $x_{i_2}$ de degré $N - 1$. Cela montre que
$\deg : \Gamma \to \widehat{\mathbf{Z}}$ est surjective,
et il en va donc de même pour $U$. Cela signifie exactement que, si
$Y \to X$ est le revêtement galoisien étale fini
correspondant à $U$, alors $Y_{\overline{k}}$ est irréductible.
Posons $G = \text{Aut}(Y/X)$. On a
$$
Y \to^G X,\quad G = \pi_1(X)/U
$$
Par construction, tous les points de $X$ de degré $\leq N$ se décomposent
complètement dans $Y$. En particulier,
$$
\# Y(k_N)\geq (\# G)\# X(k_N)
$$
Appliquons l'hypothèse de Riemann aux deux membres. On obtient
$$
q^N+1+2g_Yq^{N/2}\geq \# G\# X(k_N)\geq \#
G(q^N+1-2g_Xq^{N/2})
$$
Comme $2g_Y-2 = (\# G)(2g_X-2)$, cela donne
$$
q^N + 1 + (\# G)(2g_X - 1) + 1)q^{N/2}\geq
\# G (q^N + 1 - 2g_Xq^{N/2})
$$
On voit donc que $G$ est nécessairement trivial si $N$ est assez grand.
\end{proof}

\noindent
{\bf Question curieuse.}
Posons $W_X = \deg^{-1}(\mathbf{Z})\subset \pi_1(X)$.
Existe-t-il un ensemble fini de points fermés $x_1, \ldots, x_n$ de $X$ tel que
l'ensemble de tous les éléments de Frobenius correspondant à ces points
engendre {\it algébriquement} $W_X$ ?

\medskip\noindent
Un argument de catégorie de Baire ramène cette question à la même question
pour tous les éléments de Frobenius.





\section{Combien y a-t-il réellement de points ?}
\label{trace-section-really}

\noindent
Si le genre de la courbe est grand par rapport à $q$, alors le terme
principal de la formule $\# X(k) = q - \sum \omega_i + 1$ n'est pas le terme $q$,
mais le deuxième terme $\sum \omega_i$, qui peut avoir a priori une
taille de l'ordre de $2g_X\sqrt{q}$. Dans l'article \cite{Drinfeld-number},
Drinfeld et Vlăduţ montrent que ce maximum est en fait au plus de l'ordre de
$g\sqrt{q}$, comme Ihara l'avait prédit.

\medskip\noindent
Fixons $q$ et soit $k$ un corps à $q$ éléments. Posons
$$
A(q) = \limsup_{g_X \to \infty} \frac{\# X(k)}{g_X}
$$
où $X$ parcourt les courbes projectives lisses géométriquement irréductibles
sur $k$. Avec cette définition, on a les résultats suivants :
\begin{itemize}
\item HR $\Rightarrow A(q)\leq 2\sqrt{q}$
\item Ihara $\Rightarrow A(q)\leq \sqrt{2q}$
\item DV $\Rightarrow A(q)\leq \sqrt{q}-1$ (en fait, cette borne est atteinte si $q$
est un carré).
\end{itemize}

\begin{proof}
Étant donné $X$, soient $w_1, \ldots, w_{2g}$ comme précédemment et posons $g = g_X$.
Posons $\alpha_i = \frac{w_i}{\sqrt{q}}$ ; alors $|\alpha_i| = 1$. Si $\alpha_i$
apparaît, alors $\overline{\alpha}_i = \alpha_i^{-1}$ apparaît également. On a
$$
N = \# X(k) \leq X(k_r) = q^r + 1 - (\sum_i \alpha_i^r) q^{r/2}
$$
En réécrivant, on voit que, pour tout $r \geq 1$,
$$
-\sum_i \alpha_i^r \geq Nq^{-r/2} - q^{r/2} - q^{-r/2}
$$
Observons que
$$
0 \leq |\alpha_i^n +\alpha_i^{n-1} +\ldots +\alpha_i +1|^2
= (n + 1) + \sum_{j = 1}^n (n + 1 - j) (\alpha_i^j + \alpha_i^{-j})
$$
Ainsi,
\begin{align*}
2g(n+1) & \geq - \sum_i \left(\sum_{j = 1}^n (n+1-j)(\alpha_i^j
+\alpha_i^{-j})\right)\\
& =-\sum_{j = 1}^n (n+1-j)\left(\sum_i\alpha_i^j
+\sum_i\alpha_i^{-j}\right)
\end{align*}
En prenant la moitié, on obtient
\begin{align*}
g(n+1)& \geq - \sum_{j = 1}^n (n+1-j)(\sum_i\alpha_i^j)\\
& \geq N\sum_{j = 1}^n (n+1-j)q^{-j/2}-\sum_{j = 1}^n
(n+1-j)(q^{j/2}+q^{-j/2})
\end{align*}
On en déduit
$$
\frac{N}{g}\leq \left(\sum_{j = 1}^n \frac{n+1-j}{n+1}q^{-j/2} \right)^{-1}
\cdot
\left(
1 + \frac{1}{g} \sum_{j = 1}^n \frac{n + 1 - j}{n + 1}(q^{j/2} + q^{-j/2})
\right)
$$
Fixons $n$ et faisons tendre $g\to \infty$.
$$
A(q)\leq \left(\sum_{j = 1}^n \frac{n+1-j}{n+1}q^{-j/2}\right)^{-1}
$$
Donc
$$
A(q)\leq \lim_{n\to\infty}(\ldots) = \left(\sum_{j = 1}^\infty
q^{-j/2}\right)^{-1}=\sqrt{q}-1
$$
\end{proof}









% OK, start here.
%

\chapter{Espaces algébriques}



\phantomsection
\label{spaces-section-phantom}


\section{Introduction}
\label{spaces-section-introduction}

\noindent
Les espaces algébriques ont été introduits par Michael Artin,
voir \cite{ArtinI}, \cite{ArtinII},
\cite{Artin-Theorem-Representability},
\cite{Artin-Construction-Techniques},
\cite{Artin-Algebraic-Spaces},
\cite{Artin-Algebraic-Approximation},
\cite{Artin-Implicit-Function},
et \cite{ArtinVersal}.
Une partie des fondements a été élaborée conjointement avec
Knutson, qui en a tiré l'ouvrage \cite{Kn}.
Artin a défini (voir \cite[Définition 1.3]{Artin-Implicit-Function})
un espace algébrique comme un faisceau pour la topologie étale
qui est localement représentable pour la topologie étale.
Dans la plupart des travaux d'Artin, les catégories de schémas
considérées sont formées de schémas localement de type fini sur une base
noethérienne excellente fixée.

\medskip\noindent
Notre définition diffère légèrement de la définition originelle d'Artin.
En effet, nos espaces algébriques sont des faisceaux pour la topologie fppf
dont la diagonale est représentable et qui possèdent un «~revêtement~» étale
par un schéma. Travailler avec la topologie fppf plutôt qu'avec la topologie
étale n'est qu'un point technique et ne change presque
rien; nous montrerons dans
Amorçage, section \ref{bootstrap-section-spaces-etale}
que nous aurions obtenu la même catégorie d'espaces algébriques
en travaillant avec la topologie étale. Dans ce même chapitre,
nous démontrerons que la condition portant sur la diagonale
peut, en un certain sens, être supprimée; voir
Amorçage, section \ref{bootstrap-section-bootstrap}.

\medskip\noindent
Après avoir défini les espaces algébriques, nous formulons quelques
observations fondamentales. Le principal résultat de ce chapitre est qu'avec
nos définitions, la donnée d'un espace algébrique équivaut à celle d'une
relation d'équivalence étale; voir l'exposé de la section
\ref{spaces-section-presentations} et le Théorème \ref{spaces-theorem-presentation}.
L'analogue de ce théorème dans le cadre d'Artin est \cite[Théorème 1.5]{Artin-Implicit-Function}, ou
\cite[Proposition II.1.7]{Kn}. Autrement dit, le faisceau
défini par une relation d'équivalence étale a une diagonale représentable.
Il s'ensuit que notre définition coïncide, au sens large, avec la définition
originelle d'Artin. Cela signifie aussi que l'on peut donner des exemples
d'espaces algébriques en écrivant simplement une relation d'équivalence étale.

\medskip\noindent
Dans la section \ref{spaces-section-separation}, nous introduisons divers axiomes
de séparation pour les espaces algébriques rencontrés dans la littérature.
Enfin, dans la section \ref{spaces-section-examples},
nous donnons des exemples insolites d'espaces algébriques et d'autres qui le sont moins.






\section{Remarques générales}
\label{spaces-section-general}

\noindent
Nous travaillons dans un grand site fppf convenable $\Sch_{fppf}$
comme dans Topologies, Définition \ref{topologies-definition-big-fppf-site}.
Ainsi, sauf mention expresse du contraire, tous les schémas seront des objets
de $\Sch_{fppf}$.
Dans la section \ref{spaces-section-change-big-site}, nous étudions les modifications
entraînées par un changement de grand site fppf.

\medskip\noindent
Nous travaillerons toujours relativement à un schéma de base $S$ appartenant à $\Sch_{fppf}$.
Nous travaillerons alors avec le grand site fppf $(\Sch/S)_{fppf}$,
voir Topologies, Définition \ref{topologies-definition-big-small-fppf}.
On retrouve le cas absolu en prenant
$S = \Spec(\mathbf{Z})$.

\medskip\noindent
Si $U, T$ sont des schémas sur $S$, nous désignons
par $U(T)$ l'ensemble des points à valeurs dans $T$ {\it au-dessus} de $S$.
Autrement dit : $U(T) = \Mor_S(T, U)$.

\medskip\noindent
Notons que tout recouvrement fpqc est un épimorphisme effectif universel, voir
Descente, Lemme \ref{descent-lemma-fpqc-universal-effective-epimorphisms}.
Par conséquent, la topologie sur $\Sch_{fppf}$
est moins fine que la topologie canonique, et tous les préfaisceaux représentables
sont des faisceaux.







\section{Morphismes représentables de préfaisceaux}
\label{spaces-section-representable}

\noindent
Soit $S$ un schéma appartenant à $\Sch_{fppf}$.
Soient $F, G : (\Sch/S)_{fppf}^{opp} \to \textit{Ens}$.
Soit $a : F \to G$ un morphisme représentable de foncteurs, voir
Catégories,
Définition \ref{categories-definition-representable-map-presheaves}.
Cela signifie que pour tout
$U \in \Ob((\Sch/S)_{fppf})$
et tout $\xi \in G(U)$, le produit fibré $h_U \times_{\xi, G} F$ est représentable.
Choisissons un objet $V_\xi$ qui le représente et un isomorphisme
$h_{V_\xi} \to h_U \times_G F$.
D'après le lemme de Yoneda, voir Catégories, Lemme \ref{categories-lemma-yoneda},
la projection $h_{V_\xi} \to h_U \times_G F \to h_U$ provient d'un unique
morphisme de schémas $a_\xi : V_\xi \to U$.
Le diagramme suivant permet de se représenter la situation :
$$
\xymatrix{
V_\xi \ar@{~>}[r] \ar[d]_{a_\xi} & h_{V_\xi} \ar[d] \ar[r] & F \ar[d]^a \\
U \ar@{~>}[r] & h_U \ar[r]^\xi & G
}
$$
où les flèches ondulées représentent le plongement de Yoneda.
Voici quelques lemmes sur cette notion, valables dans un cadre très général.

\begin{lemma}
\label{spaces-lemma-morphism-schemes-gives-representable-transformation}
Soit $S$ un schéma appartenant à $\Sch_{fppf}$ et soient
$X$, $Y$ des objets de $(\Sch/S)_{fppf}$.
Soit $f : X \to Y$ un morphisme de schémas.
Alors
$$
h_f : h_X \longrightarrow h_Y
$$
est un morphisme représentable de foncteurs.
\end{lemma}

\begin{proof}
Cela est formel et repose seulement sur le fait que
la catégorie $(\Sch/S)_{fppf}$ admet des produits fibrés.
\end{proof}

\begin{lemma}
\label{spaces-lemma-composition-representable-transformations}
Soit $S$ un schéma appartenant à $\Sch_{fppf}$.
Soient $F, G, H : (\Sch/S)_{fppf}^{opp} \to \textit{Ens}$.
Soient $a : F \to G$, $b : G \to H$ des morphismes représentables de foncteurs.
Alors
$$
b \circ a : F \longrightarrow H
$$
est un morphisme représentable de foncteurs.
\end{lemma}

\begin{proof}
Cela est entièrement formel et vaut dans toute catégorie.
\end{proof}

\begin{lemma}
\label{spaces-lemma-base-change-representable-transformations}
Soit $S$ un schéma appartenant à $\Sch_{fppf}$.
Soient $F, G, H : (\Sch/S)_{fppf}^{opp} \to \textit{Ens}$.
Soit $a : F \to G$ un morphisme représentable de foncteurs.
Soit $b : H \to G$ un morphisme de foncteurs quelconque.
Considérons le diagramme de produit fibré
$$
\xymatrix{
H \times_{b, G, a} F \ar[r]_-{b'} \ar[d]_{a'} & F \ar[d]^a \\
H \ar[r]^b & G
}
$$
Alors le changement de base $a'$ est un morphisme représentable de foncteurs.
\end{lemma}

\begin{proof}
Cela est entièrement formel et vaut dans toute catégorie.
\end{proof}

\begin{lemma}
\label{spaces-lemma-product-representable-transformations}
Soit $S$ un schéma appartenant à $\Sch_{fppf}$.
Soient $F_i, G_i : (\Sch/S)_{fppf}^{opp} \to \textit{Ens}$, $i = 1, 2$.
Soient $a_i : F_i \to G_i$, $i = 1, 2$,
des morphismes représentables de foncteurs.
Alors
$$
a_1 \times a_2 : F_1 \times F_2 \longrightarrow G_1 \times G_2
$$
est un morphisme représentable de foncteurs.
\end{lemma}

\begin{proof}
Écrivons $a_1 \times a_2$ comme la composée
$F_1 \times F_2 \to G_1 \times F_2 \to G_1 \times G_2$.
La première flèche est le changement de base de $a_1$ par le morphisme
$G_1 \times F_2 \to G_1$, et la seconde flèche
est le changement de base de $a_2$ par le morphisme
$G_1 \times G_2 \to G_2$. Ce lemme est donc une conséquence
formelle des Lemmes \ref{spaces-lemma-composition-representable-transformations}
et \ref{spaces-lemma-base-change-representable-transformations}.
\end{proof}

\begin{lemma}
\label{spaces-lemma-representable-transformation-to-sheaf}
Soit $S$ un schéma appartenant à $\Sch_{fppf}$.
Soient $F, G : (\Sch/S)_{fppf}^{opp} \to \textit{Ens}$.
Soit $a : F \to G$ un morphisme représentable de foncteurs.
Si $G$ est un faisceau, alors $F$ l'est aussi.
\end{lemma}

\begin{proof}
Soit $\{\varphi_i : T_i \to T\}$ une famille couvrante du site
$(\Sch/S)_{fppf}$.
Soient $s_i \in F(T_i)$ satisfaisant à la condition de recollement.
Alors $\sigma_i = a(s_i) \in G(T_i)$ satisfont elles aussi à la condition de recollement.
Il existe donc un unique $\sigma \in G(T)$ tel
que $\sigma_i = \sigma|_{T_i}$. Par hypothèse,
$F' = h_T \times_{\sigma, G, a} F$ est un préfaisceau représentable
et donc (voir les remarques de la section \ref{spaces-section-general}) un faisceau.
Notons que $(\varphi_i, s_i) \in F'(T_i)$ satisfont
elles aussi à la condition de recollement et proviennent donc d'un unique
$(\text{id}_T, s) \in F'(T)$. Manifestement, $s$ est la section de
$F$ que nous cherchions.
\end{proof}

\begin{lemma}
\label{spaces-lemma-representable-transformation-diagonal}
Soit $S$ un schéma appartenant à $\Sch_{fppf}$.
Soient $F, G : (\Sch/S)_{fppf}^{opp} \to \textit{Ens}$.
Soit $a : F \to G$ un morphisme représentable de foncteurs.
Alors $\Delta_{F/G} : F \to F \times_G F$ est représentable.
\end{lemma}

\begin{proof}
Soit $U \in \Ob((\Sch/S)_{fppf})$. Soit
$\xi = (\xi_1, \xi_2) \in (F \times_G F)(U)$.
Posons $\xi' = a(\xi_1) = a(\xi_2) \in G(U)$.
Par hypothèse, il existe un schéma $V$ muni d'un morphisme $V \to U$
qui représente le produit fibré $h_U \times_{\xi', G} F$.
En particulier, les éléments $\xi_1, \xi_2$ définissent des morphismes
$f_1, f_2 : U \to V$ au-dessus de $U$. Puisque $V$ représente le
produit fibré $h_U \times_{\xi', G} F$ et que
$\xi' = a \circ \xi_1 = a \circ \xi_2$
nous voyons que, si $g : U' \to U$ est un morphisme, alors
$$
g^*\xi_1 = g^*\xi_2
\Leftrightarrow
f_1 \circ g = f_2 \circ g.
$$
Autrement dit, nous voyons que $h_U \times_{\xi, F \times_G F} F$
est représenté par $V \times_{\Delta, V \times V, (f_1, f_2)} U$
qui est un schéma.
\end{proof}










\section{Listes de propriétés utiles des morphismes de schémas}
\label{spaces-section-lists}

\noindent
Pour faciliter les renvois, les remarques suivantes énumèrent les
propriétés de morphismes qui vérifient certaines des conditions
requises dans des résultats ultérieurs.

\begin{remark}
\label{spaces-remark-list-properties-stable-base-change}
Voici une liste de propriétés/types de morphismes
qui sont {\it stables par changement de base quelconque} :
\begin{enumerate}
\item immersions fermées, ouvertes et localement fermées, voir
Schémas, Lemme \ref{schemes-lemma-base-change-immersion},
\item quasi-compact, voir
Schémas, Lemme \ref{schemes-lemma-quasi-compact-preserved-base-change},
\item universellement fermé, voir
Schémas, Définition \ref{schemes-definition-universally-closed},
\item (quasi-)séparé, voir
Schémas, Lemme \ref{schemes-lemma-separated-permanence},
\item monomorphisme, voir
Schémas, Lemme \ref{schemes-lemma-base-change-monomorphism}
\item surjectif, voir
Morphismes, Lemme \ref{morphisms-lemma-base-change-surjective},
\item radiciel, voir
Morphismes, Lemme \ref{morphisms-lemma-universally-injective},
\item affine, voir
Morphismes, Lemme \ref{morphisms-lemma-base-change-affine},
\item quasi-affine, voir
Morphismes, Lemme \ref{morphisms-lemma-base-change-quasi-affine},
\item (localement) de type fini, voir
Morphismes, Lemme \ref{morphisms-lemma-base-change-finite-type},
\item (localement) quasi-fini, voir
Morphismes, Lemme \ref{morphisms-lemma-base-change-quasi-finite},
\item (localement) de présentation finie, voir
Morphismes, Lemme \ref{morphisms-lemma-base-change-finite-presentation},
\item localement de type fini de dimension relative $d$, voir
Morphismes, Lemme \ref{morphisms-lemma-base-change-relative-dimension-d},
\item universellement ouvert, voir
Morphismes, Définition \ref{morphisms-definition-open},
\item plat, voir
Morphismes, Lemme \ref{morphisms-lemma-base-change-flat},
\item syntomique, voir
Morphismes, Lemme \ref{morphisms-lemma-base-change-syntomic},
\item lisse, voir
Morphismes, Lemme \ref{morphisms-lemma-base-change-smooth},
\item non ramifié (resp.\ G-non ramifié), voir
Morphismes, Lemme \ref{morphisms-lemma-base-change-unramified},
\item \'etale, voir
Morphismes, Lemme \ref{morphisms-lemma-base-change-etale},
\item propre, voir
Morphismes, Lemme \ref{morphisms-lemma-base-change-proper},
\item H-projectif, voir
Morphismes, Lemme \ref{morphisms-lemma-H-projective-base-change},
\item (localement) projectif, voir
Morphismes, Lemme \ref{morphisms-lemma-base-change-projective},
\item fini ou entier, voir
Morphismes, Lemme \ref{morphisms-lemma-base-change-finite},
\item fini localement libre, voir
Morphismes, Lemme \ref{morphisms-lemma-base-change-finite-locally-free},
\item universellement submersif, voir
Morphismes, Lemme \ref{morphisms-lemma-base-change-universally-submersive},
\item homéomorphisme universel, voir
Morphismes, Lemme \ref{morphisms-lemma-base-change-universal-homeomorphism}.
\end{enumerate}
Compléter au besoin.
\end{remark}

\begin{remark}
\label{spaces-remark-list-properties-stable-composition}
Parmi les propriétés de morphismes qui sont stables par changement de base
(telles qu'elles sont énumérées dans la
Remarque \ref{spaces-remark-list-properties-stable-base-change}),
les suivantes sont aussi {\it stables par composition} :
\begin{enumerate}
\item immersions fermées, ouvertes et localement fermées, voir
Schémas, Lemme \ref{schemes-lemma-composition-immersion},
\item quasi-compact, voir
Schémas, Lemme \ref{schemes-lemma-composition-quasi-compact},
\item universellement fermé, voir
Morphismes, Lemme \ref{morphisms-lemma-composition-proper},
\item (quasi-)séparé, voir
Schémas, Lemme \ref{schemes-lemma-separated-permanence},
\item monomorphisme, voir
Schémas, Lemme \ref{schemes-lemma-composition-monomorphism},
\item surjectif, voir
Morphismes, Lemme \ref{morphisms-lemma-composition-surjective},
\item radiciel, voir
Morphismes, Lemme \ref{morphisms-lemma-composition-universally-injective},
\item affine, voir
Morphismes, Lemme \ref{morphisms-lemma-composition-affine},
\item quasi-affine, voir
Morphismes, Lemme \ref{morphisms-lemma-composition-quasi-affine},
\item (localement) de type fini, voir
Morphismes, Lemme \ref{morphisms-lemma-composition-finite-type},
\item (localement) quasi-fini, voir
Morphismes, Lemme \ref{morphisms-lemma-composition-quasi-finite},
\item (localement) de présentation finie, voir
Morphismes, Lemme \ref{morphisms-lemma-composition-finite-presentation},
\item universellement ouvert, voir
Morphismes, Lemme \ref{morphisms-lemma-composition-open},
\item plat, voir
Morphismes, Lemme \ref{morphisms-lemma-composition-flat},
\item syntomique, voir
Morphismes, Lemme \ref{morphisms-lemma-composition-syntomic},
\item lisse, voir
Morphismes, Lemme \ref{morphisms-lemma-composition-smooth},
\item non ramifié (resp.\ G-non ramifié), voir
Morphismes, Lemme \ref{morphisms-lemma-composition-unramified},
\item \'etale, voir
Morphismes, Lemme \ref{morphisms-lemma-composition-etale},
\item propre, voir
Morphismes, Lemme \ref{morphisms-lemma-composition-proper},
\item H-projectif, voir
Morphismes, Lemme \ref{morphisms-lemma-H-projective-composition},
\item fini ou entier, voir
Morphismes, Lemme \ref{morphisms-lemma-composition-finite},
\item fini localement libre, voir
Morphismes, Lemme \ref{morphisms-lemma-composition-finite-locally-free},
\item universellement submersif, voir
Morphismes, Lemme \ref{morphisms-lemma-composition-universally-submersive},
\item homéomorphisme universel, voir
Morphismes, Lemme \ref{morphisms-lemma-composition-universal-homeomorphism}.
\end{enumerate}
Compléter au besoin.
\end{remark}

\begin{remark}
\label{spaces-remark-list-properties-fpqc-local-base}
Parmi les propriétés mentionnées qui sont stables par changement de base
(énumérées dans la Remarque \ref{spaces-remark-list-properties-stable-base-change}),
les suivantes sont aussi {\it locales sur la base pour la topologie fpqc}
(et a fortiori locales sur la base pour la topologie fppf) :
\begin{enumerate}
\item pour les immersions, c'est le cas pour
\begin{enumerate}
\item les immersions fermées, voir
Descente, Lemme \ref{descent-lemma-descending-property-closed-immersion},
\item les immersions ouvertes, voir
Descente, Lemme \ref{descent-lemma-descending-property-open-immersion}, et
\item les immersions quasi-compactes, voir
Descente,
Lemme \ref{descent-lemma-descending-property-quasi-compact-immersion},
\end{enumerate}
\item quasi-compact, voir
Descente, Lemme \ref{descent-lemma-descending-property-quasi-compact},
\item universellement fermé, voir
Descente, Lemme
\ref{descent-lemma-descending-property-universally-closed},
\item (quasi-)séparé, voir
Descente, Lemmes
\ref{descent-lemma-descending-property-quasi-separated}, et
\ref{descent-lemma-descending-property-separated},
\item monomorphisme, voir
Descente, Lemme \ref{descent-lemma-descending-property-monomorphism},
\item surjectif, voir
Descente, Lemme \ref{descent-lemma-descending-property-surjective},
\item radiciel, voir
Descente, Lemme \ref{descent-lemma-descending-property-universally-injective},
\item affine, voir
Descente, Lemme \ref{descent-lemma-descending-property-affine},
\item quasi-affine, voir
Descente, Lemme \ref{descent-lemma-descending-property-quasi-affine},
\item (localement) de type fini, voir
Descente,
Lemmes \ref{descent-lemma-descending-property-locally-finite-type}, et
\ref{descent-lemma-descending-property-finite-type},
\item (localement) quasi-fini, voir
Descente, Lemme \ref{descent-lemma-descending-property-quasi-finite},
\item (localement) de présentation finie, voir
Descente, Lemmes
\ref{descent-lemma-descending-property-locally-finite-presentation}, et
\ref{descent-lemma-descending-property-finite-presentation},
\item localement de type fini de dimension relative $d$, voir
Descente,
Lemme \ref{descent-lemma-descending-property-relative-dimension-d},
\item universellement ouvert, voir
Descente, Lemme \ref{descent-lemma-descending-property-universally-open},
\item plat, voir
Descente, Lemme \ref{descent-lemma-descending-property-flat},
\item syntomique, voir
Descente, Lemme \ref{descent-lemma-descending-property-syntomic},
\item lisse, voir
Descente, Lemme \ref{descent-lemma-descending-property-smooth},
\item non ramifié (resp.\ G-non ramifié), voir
Descente, Lemme \ref{descent-lemma-descending-property-unramified},
\item \'etale, voir
Descente, Lemme \ref{descent-lemma-descending-property-etale},
\item propre, voir
Descente, Lemme \ref{descent-lemma-descending-property-proper},
\item fini ou entier, voir
Descente, Lemme \ref{descent-lemma-descending-property-finite},
\item fini localement libre, voir
Descente, Lemme \ref{descent-lemma-descending-property-finite-locally-free},
\item universellement submersif, voir
Descente, Lemme \ref{descent-lemma-descending-property-universally-submersive},
\item homéomorphisme universel, voir
Descente, Lemme \ref{descent-lemma-descending-property-universal-homeomorphism}.
\end{enumerate}
Notons que la propriété d'être une ``immersion'' n'est pas nécessairement locale
sur la base pour la topologie fpqc, mais que dans
Descente, Lemme \ref{descent-lemma-descending-fppf-property-immersion},
nous avons démontré qu'elle est locale sur la base pour la topologie fppf.
\end{remark}








\section{Propriétés des morphismes représentables de préfaisceaux}
\label{spaces-section-representable-properties}

\noindent
Voici la définition qui permet de procéder.

\begin{definition}
\label{spaces-definition-relative-representable-property}
On se donne $S$ et $a : F \to G$ représentable comme ci-dessus.
Soit $\mathcal{P}$ une propriété de morphismes de schémas qui
\begin{enumerate}
\item est stable par tout changement de base,
voir Schémas, Définition \ref{schemes-definition-preserved-by-base-change},
et
\item est locale sur la base pour la topologie fppf, voir
Descente, Définition \ref{descent-definition-property-morphisms-local}.
\end{enumerate}
Dans ce cas, nous disons que $a$ possède {\it la propriété $\mathcal{P}$} si, pour tout
$U \in \Ob((\Sch/S)_{fppf})$ et
tout $\xi \in G(U)$, le morphisme de schémas qui en résulte
$V_\xi \to U$ possède la propriété $\mathcal{P}$.
\end{definition}

\noindent
Il importe de noter que nous n'utiliserons cette définition que pour
des propriétés de morphismes stables par changement de base et
locales sur la base pour la topologie fppf. Ce choix ne vient
pas de ce que la définition n'aurait autrement aucun sens ; il vient plutôt
de ce que nous pouvons vouloir donner une autre définition
mieux adaptée à la propriété envisagée.

\begin{remark}
\label{spaces-remark-warning}
Considérons la propriété $\mathcal{P}=$``surjectif''.
Dans ce cas, dire
``soit $F \to G$ un morphisme surjectif'' pourrait être ambigu.
En effet, nous pourrions entendre la notion définie
dans la Définition \ref{spaces-definition-relative-representable-property}
ci-dessus, ou bien un morphisme surjectif de préfaisceaux, voir
Sites, Définition \ref{sites-definition-presheaves-injective-surjective},
ou encore, si $F$ et $G$ sont tous deux des faisceaux,
un morphisme surjectif de faisceaux, voir
Sites, Définition \ref{sites-definition-sheaves-injective-surjective}.
Sauf mention expresse du contraire, lorsque nous parlerons de morphismes d'espaces algébriques,
nous entendrons toujours la première notion. Voir le
Lemme \ref{spaces-lemma-surjective-flat-locally-finite-presentation}
pour un cas où cette propriété entraîne la surjectivité du morphisme de faisceaux sous-jacent.
\end{remark}

\noindent
Voici une vérification de cohérence.

\begin{lemma}
\label{spaces-lemma-morphism-schemes-gives-representable-transformation-property}
Soient $S$, $X$, $Y$ des objets de $\Sch_{fppf}$.
Soit $f : X \to Y$ un morphisme de schémas.
Soit $\mathcal{P}$ comme dans la
Définition \ref{spaces-definition-relative-representable-property}.
Alors $h_X \longrightarrow h_Y$ possède la propriété $\mathcal{P}$ si
et seulement si $f$ possède la propriété $\mathcal{P}$.
\end{lemma}

\begin{proof}
Notons que le lemme a un sens en vertu du
Lemme \ref{spaces-lemma-morphism-schemes-gives-representable-transformation}.
Démonstration omise.
\end{proof}

\begin{lemma}
\label{spaces-lemma-composition-representable-transformations-property}
Soit $S$ un schéma appartenant à $\Sch_{fppf}$.
Soient $F, G, H : (\Sch/S)_{fppf}^{opp} \to \textit{Ens}$.
Soit $\mathcal{P}$ une propriété comme dans la
Définition \ref{spaces-definition-relative-representable-property}
qui est stable par composition.
Soient $a : F \to G$, $b : G \to H$ des morphismes représentables de foncteurs.
Si $a$ et $b$ possèdent la propriété $\mathcal{P}$, il en va de même de
$b \circ a : F \longrightarrow H$.
\end{lemma}

\begin{proof}
Notons que le lemme a un sens en vertu du
Lemme \ref{spaces-lemma-composition-representable-transformations}.
Démonstration omise.
\end{proof}

\begin{lemma}
\label{spaces-lemma-base-change-representable-transformations-property}
Soit $S$ un schéma appartenant à $\Sch_{fppf}$.
Soient $F, G, H : (\Sch/S)_{fppf}^{opp} \to \textit{Ens}$.
Soit $\mathcal{P}$ une propriété comme dans la
Définition \ref{spaces-definition-relative-representable-property}.
Soit $a : F \to G$ un morphisme représentable de foncteurs.
Soit $b : H \to G$ un morphisme quelconque de foncteurs.
Considérons le diagramme de produit fibré
$$
\xymatrix{
H \times_{b, G, a} F \ar[r]_-{b'} \ar[d]_{a'} & F \ar[d]^a \\
H \ar[r]^b & G
}
$$
Si $a$ possède la propriété $\mathcal{P}$, alors son changement de base $a'$
possède également la propriété $\mathcal{P}$.
\end{lemma}

\begin{proof}
Notons que le lemme a un sens en vertu du
Lemme \ref{spaces-lemma-base-change-representable-transformations}.
Démonstration omise.
\end{proof}

\begin{lemma}
\label{spaces-lemma-descent-representable-transformations-property}
Soit $S$ un schéma appartenant à $\Sch_{fppf}$.
Soient $F, G, H : (\Sch/S)_{fppf}^{opp} \to \textit{Ens}$.
Soit $\mathcal{P}$ une propriété comme dans la
Définition \ref{spaces-definition-relative-representable-property}.
Soit $a : F \to G$ un morphisme représentable de foncteurs.
Soit $b : H \to G$ un morphisme quelconque de foncteurs.
Considérons le diagramme de produit fibré
$$
\xymatrix{
H \times_{b, G, a} F \ar[r]_-{b'} \ar[d]_{a'} & F \ar[d]^a \\
H \ar[r]^b & G
}
$$
Supposons que $b$ induise un morphisme surjectif de faisceaux pour la topologie fppf $H^\# \to G^\#$.
Dans ce cas, si $a'$ possède la propriété $\mathcal{P}$, alors $a$
possède également la propriété $\mathcal{P}$.
\end{lemma}

\begin{proof}
Remarquons d'abord que, d'après le
Lemme \ref{spaces-lemma-base-change-representable-transformations},
le morphisme $a'$ est représentable.
Soient $U \in \Ob((\Sch/S)_{fppf})$ et
$\xi \in G(U)$. Par hypothèse, il existe un recouvrement fppf
$\{U_i \to U\}_{i \in I}$ et des éléments $\xi_i \in H(U_i)$
qui sont envoyés sur $\xi|_{U_i}$ par $b$. Il résulte de la théorie générale des catégories que, pour
chaque $i$, on a un diagramme de produit fibré
$$
\xymatrix{
U_i \times_{\xi_i, H, a'} (H \times_{b, G, a} F) \ar[r] \ar[d] &
U \times_{\xi, G, a} F \ar[d] \\
U_i \ar[r] & U
}
$$
Par hypothèse, la flèche verticale de gauche est un morphisme de schémas qui
possède la propriété $\mathcal{P}$. Comme $\mathcal{P}$ est locale sur la base pour la
topologie fppf, il en résulte que la flèche verticale de droite possède aussi la propriété
$\mathcal{P}$, comme souhaité.
\end{proof}

\begin{lemma}
\label{spaces-lemma-product-representable-transformations-property}
Soit $S$ un schéma appartenant à $\Sch_{fppf}$.
Soient $F_i, G_i : (\Sch/S)_{fppf}^{opp} \to \textit{Ens}$,
$i = 1, 2$.
Soient $a_i : F_i \to G_i$, $i = 1, 2$, des morphismes représentables
de foncteurs.
Soit $\mathcal{P}$ une propriété comme dans la
Définition \ref{spaces-definition-relative-representable-property}
qui est stable par composition.
Si $a_1$ et $a_2$ possèdent la propriété $\mathcal{P}$, il en va de même de
$a_1 \times a_2 : F_1 \times F_2 \longrightarrow G_1 \times G_2$.
\end{lemma}

\begin{proof}
Notons que le lemme a un sens en vertu du
Lemme \ref{spaces-lemma-product-representable-transformations}.
Démonstration omise.
\end{proof}

\begin{lemma}
\label{spaces-lemma-representable-transformations-property-implication}
Soit $S$ un schéma appartenant à $\Sch_{fppf}$.
Soient $F, G : (\Sch/S)_{fppf}^{opp} \to \textit{Ens}$.
Soit $a : F \to G$ un morphisme représentable de foncteurs.
Soient $\mathcal{P}$, $\mathcal{P}'$ des propriétés comme dans la
Définition \ref{spaces-definition-relative-representable-property}.
Supposons que, pour tout morphisme de schémas $f : X \to Y$,
on ait $\mathcal{P}(f) \Rightarrow \mathcal{P}'(f)$.
Si $a$ possède la propriété $\mathcal{P}$, alors
$a$ possède la propriété $\mathcal{P}'$.
\end{lemma}

\begin{proof}
Formel.
\end{proof}

\begin{lemma}
\label{spaces-lemma-surjective-flat-locally-finite-presentation}
Soit $S$ un schéma.
Soient $F, G : (\Sch/S)_{fppf}^{opp} \to \textit{Ens}$ des faisceaux.
Soit $a : F \to G$ représentable, plat,
localement de présentation finie et surjectif.
Alors $a : F \to G$ est surjectif en tant que morphisme de faisceaux.
\end{lemma}

\begin{proof}
Soient $T$ un schéma au-dessus de $S$ et $g : T \to G$ un point à valeurs dans $T$ du foncteur
$G$. Par hypothèse, $T' = F \times_G T$ est (représenté par) un schéma et
le morphisme $T' \to T$ est plat, localement de présentation finie et
surjectif. Ainsi, $\{T' \to T\}$ est un recouvrement fppf tel
que $g|_{T'} \in G(T')$ provient d'un élément de $F(T')$, à savoir
le morphisme $T' \to F$. Cela prouve que le morphisme est surjectif en tant que
morphisme de faisceaux, voir
Sites, Définition \ref{sites-definition-sheaves-injective-surjective}.
\end{proof}

\noindent
Voici une caractérisation des foncteurs dont la
diagonale est représentable.

\begin{lemma}
\label{spaces-lemma-representable-diagonal}
Soit $S$ un schéma appartenant à $\Sch_{fppf}$.
Soit $F$ un préfaisceau d'ensembles sur $(\Sch/S)_{fppf}$.
Les assertions suivantes sont équivalentes :
\begin{enumerate}
\item la diagonale $F \to F \times F$ est représentable,
\item pour $U \in \Ob((\Sch/S)_{fppf})$ et tout $a \in F(U)$,
le morphisme $a : h_U \to F$ est représentable,
\item pour toute paire $U, V \in \Ob((\Sch/S)_{fppf})$
et tous $a \in F(U)$ et $b \in F(V)$, le produit fibré
$h_U \times_{a, F, b} h_V$ est représentable.
\end{enumerate}
\end{lemma}

\begin{proof}
C'est purement formel, voir
Catégories, Lemme \ref{categories-lemma-representable-diagonal}.
Cela ne dépend que du fait que la catégorie $(\Sch/S)_{fppf}$
possède les produits de deux objets et les produits fibrés, voir
Topologies, Lemme \ref{topologies-lemma-fibre-products-fppf}.
\end{proof}

\noindent
Dans la situation du lemme, pour tout morphisme
$\xi : h_U \to F$ comme dans son énoncé, il est légitime
de dire que $\xi$ possède la propriété $\mathcal{P}$, pour toute propriété
comme dans la Définition \ref{spaces-definition-relative-representable-property}.
C'est notamment le cas pour $\mathcal{P} = $ ``surjectif''
et $\mathcal{P} = $ ``\'etale'', voir la
Remarque \ref{spaces-remark-list-properties-fpqc-local-base}
ci-dessus. Nous utiliserons cette remarque dans la définition
des espaces algébriques ci-dessous.

\begin{lemma}
\label{spaces-lemma-transformation-diagonal-properties}
Soit $S$ un schéma appartenant à $\Sch_{fppf}$.
Soit $F$ un préfaisceau d'ensembles sur $(\Sch/S)_{fppf}$.
Soit $\mathcal{P}$ une propriété comme dans la
Définition \ref{spaces-definition-relative-representable-property}.
Supposons que, pour tous $U, V \in \Ob((\Sch/S)_{fppf})$, $a \in F(U)$ et
$b \in F(V)$, on ait
\begin{enumerate}
\item $h_U \times_{a, F, b} h_V$ est représentable, disons par le schéma $W$, et
\item le morphisme $W \to U \times_S V$ correspondant au morphisme
$h_U \times_{a, F, b} h_V \to h_U \times h_V$ possède la propriété $\mathcal{P}$,
\end{enumerate}
alors $\Delta : F \to F \times F$ est représentable et possède
la propriété $\mathcal{P}$.
\end{lemma}

\begin{proof}
Observons que $\Delta$ est représentable d'après le
Lemme \ref{spaces-lemma-representable-diagonal}.
Nous pouvons reformuler la condition (2) en disant que
le morphisme $h_U \times_{a, F, b} h_V \to h_{U \times_S V}$
possède la propriété $\mathcal{P}$, voir le Lemme
\ref{spaces-lemma-morphism-schemes-gives-representable-transformation-property}.
Considérons $T \in \Ob((\Sch/S)_{fppf})$ et $(a, b) \in (F \times F)(T)$.
Observons que nous avons le diagramme commutatif
$$
\xymatrix{
F \times_{\Delta, F \times F, (a, b)} h_T \ar[d] \ar[r] &
h_T \ar[d]^{\Delta_{T/S}} \\
h_T \times_{a, F, b} h_T \ar[r] \ar[d] &
h_{T \times_S T} \ar[d]^{(a, b)} \\
F \ar[r]^\Delta & F \times F
}
$$
dont les deux carrés sont cartésiens. Nous voyons ainsi que
le morphisme $F \times_{F \times F} h_T \to h_T$ est obtenu par changement
de base d'un morphisme qui possède la propriété $\mathcal{P}$ au moyen de
$\Delta_{T/S}$. Comme $\mathcal{P}$ est stable par changement de
base, cela achève la démonstration.
\end{proof}


















\section{Espaces algébriques}
\label{spaces-section-algebraic-spaces}

\noindent
Voici la définition.

\begin{definition}
\label{spaces-definition-algebraic-space}
Soit $S$ un schéma appartenant à $\Sch_{fppf}$.
Un {\it espace algébrique au-dessus de $S$} est un préfaisceau
$$
F : (\Sch/S)^{opp}_{fppf} \longrightarrow \textit{Ens}
$$
qui possède les propriétés suivantes :
\begin{enumerate}
\item Le préfaisceau $F$ est un faisceau.
\item Le morphisme diagonal $F  \to F \times F$ est représentable.
\item Il existe un schéma $U \in \Ob((\Sch/S)_{fppf})$
et un morphisme $h_U \to F$ qui est surjectif et \'etale\footnote{Voir le
Lemme \ref{spaces-lemma-representable-diagonal}, la
Définition \ref{spaces-definition-relative-representable-property} et la
Remarque \ref{spaces-remark-warning}.}.
\end{enumerate}
\end{definition}

\noindent
Notre définition diffère en deux points de la définition ``usuelle'', par exemple de la
définition donnée dans l'ouvrage de Knutson \cite{Kn}.

\medskip\noindent
La première est que nous exigeons que $F$ soit un faisceau pour la topologie fppf.
L'une des raisons de ce choix est que de nombreux exemples naturels
d'espaces algébriques vérifient la condition de faisceau pour les recouvrements fppf
(et même pour les recouvrements fpqc). En outre, l'utilité des espaces
algébriques tient notamment aux résultats de Michael Artin à leur sujet.
Sa méthode comporte une condition qui garantit que le résultat est
localement de présentation finie au-dessus de $S$.
Tout bien considéré, il nous semble que la topologie fppf
est la topologie naturelle à employer. Finalement, on obtient
la même catégorie d'espaces algébriques. Voir
Amorçage, section \ref{bootstrap-section-spaces-etale}.

\medskip\noindent
La seconde est que nous exigeons seulement que le morphisme diagonal de $F$ soit
représentable, tandis que \cite{Kn} exige qu'il soit également
quasi-compact. Si $F = h_U$ pour un schéma $U$ au-dessus de $S$,
cela revient à demander que $U$ soit quasi-séparé.
Nous cherchons à démontrer un certain
nombre des résultats qui suivent en supposant seulement que la diagonale
de $F$ soit représentable, et à ajouter simplement une hypothèse supplémentaire chaque fois que
cela est nécessaire. Cette convention a en tout cas pour conséquence agréable que
le lemme suivant est vrai.

\begin{lemma}
\label{spaces-lemma-scheme-is-space}
Tout schéma est un espace algébrique. Plus précisément,
pour tout schéma $T \in \Ob((\Sch/S)_{fppf})$,
le foncteur représentable $h_T$ est un espace algébrique.
\end{lemma}

\begin{proof}
Le foncteur $h_T$ est un faisceau d'après nos remarques de la section \ref{spaces-section-general}.
La diagonale $h_T \to h_T \times h_T = h_{T \times T}$ est
représentable parce que $(\Sch/S)_{fppf}$ possède les produits fibrés.
Le morphisme identité $h_T \to h_T$ est surjectif et \'etale.
\end{proof}

\begin{definition}
\label{spaces-definition-morphism-algebraic-spaces}
Soient $F$, $F'$ des espaces algébriques au-dessus de $S$.
Un {\it morphisme $f : F \to F'$ d'espaces algébriques au-dessus de $S$}
est un morphisme de foncteurs de $F$ vers $F'$.
\end{definition}

\noindent
La catégorie des espaces algébriques au-dessus de $S$ contient la catégorie
$(\Sch/S)_{fppf}$ comme sous-catégorie pleine par l'intermédiaire du
plongement de Yoneda $T/S \mapsto h_T$. Désormais, nous ne distinguerons plus
un schéma $T/S$ de l'espace algébrique qu'il représente.
Ainsi, lorsque nous disons ``Soit $f : T \to F$ un morphisme du schéma
$T$ vers l'espace algébrique $F$'', nous voulons dire que
$T \in \Ob((\Sch/S)_{fppf})$, que $F$ est un
espace algébrique au-dessus de $S$ et que $f : h_T \to F$ est un morphisme
d'espaces algébriques au-dessus de $S$.






\section{Produits fibrés d'espaces algébriques}
\label{spaces-section-fibre-products}

\noindent
La catégorie des espaces algébriques au-dessus de $S$ possède à la fois des produits et des
produits fibrés.

\begin{lemma}
\label{spaces-lemma-product-spaces}
Soit $S$ un schéma appartenant à $\Sch_{fppf}$.
Soient $F, G$ des espaces algébriques au-dessus de $S$.
Alors $F \times G$ est un espace algébrique et c'est un produit
dans la catégorie des espaces algébriques au-dessus de $S$.
\end{lemma}

\begin{proof}
Il est immédiat que $H = F \times G$ est un faisceau.
La diagonale de $H$ est simplement le produit des
diagonales de $F$ et de $G$. Elle est donc représentable d'après le
Lemme \ref{spaces-lemma-product-representable-transformations}.
Enfin, si $U \to F$ et $V \to G$ sont des morphismes étales
surjectifs, où $U, V \in \Ob((\Sch/S)_{fppf})$,
alors $U \times V \to F \times G$ est étale surjectif
d'après le Lemme \ref{spaces-lemma-product-representable-transformations-property}.
\end{proof}

\begin{lemma}
\label{spaces-lemma-fibre-product-spaces-over-sheaf-with-representable-diagonal}
Soit $S$ un schéma appartenant à $\Sch_{fppf}$.
Soit $H$ un faisceau sur $(\Sch/S)_{fppf}$ dont la diagonale
est représentable. Soient $F, G$ des espaces algébriques au-dessus de $S$.
Soient $F \to H$, $G \to H$ des morphismes de faisceaux.
Alors $F \times_H G$ est un espace algébrique.
\end{lemma}

\begin{proof}
Vérifions les trois conditions de la
Définition \ref{spaces-definition-algebraic-space}.
Un produit fibré de faisceaux est un faisceau; ainsi, $F \times_H G$ est un faisceau.
La diagonale de $F \times_H G$ est la flèche verticale de gauche dans le diagramme
$$
\xymatrix{
F \times_H G \ar[r] \ar[d]_\Delta &
F \times G \ar[d]^{\Delta_F \times \Delta_G} \\
(F \times F) \times_{(H \times H)} (G \times G) \ar[r] &
(F \times F) \times (G \times G)
}
$$
qui est cartésien. Ainsi, $\Delta$ est représentable, puisqu'elle s'obtient par changement de base
à partir du morphisme de droite, qui est représentable; voir les
Lemmes \ref{spaces-lemma-product-representable-transformations} et
\ref{spaces-lemma-base-change-representable-transformations}.
Enfin, soient $U, V \in \Ob((\Sch/S)_{fppf})$
et soient $a : U \to F$, $b : V \to G$ des morphismes étales surjectifs.
Comme $\Delta_H$ est représentable, on voit que $U \times_H V$ est un schéma.
Le morphisme
$$
U \times_H V \longrightarrow F \times_H G
$$
est étale surjectif, car il est le composé des morphismes
$U \times_H V \to U \times_H G$ et $U \times_H G \to F \times_H G$,
obtenus par changement de base à partir des morphismes étales surjectifs $U \to F$ et $V \to G$; voir les
Lemmes \ref{spaces-lemma-composition-representable-transformations} et
\ref{spaces-lemma-base-change-representable-transformations}.
Cela montre que la dernière condition de la
Définition \ref{spaces-definition-algebraic-space}
est satisfaite, et l'on conclut que $F \times_H G$ est un espace algébrique.
\end{proof}

\begin{lemma}
\label{spaces-lemma-fibre-product-spaces}
Soit $S$ un schéma appartenant à $\Sch_{fppf}$.
Soient $F \to H$, $G \to H$ des morphismes d'espaces algébriques au-dessus de $S$.
Alors $F \times_H G$ est un espace algébrique et c'est un produit fibré
dans la catégorie des espaces algébriques au-dessus de $S$.
\end{lemma}

\begin{proof}
Il résulte du Lemme plus général
\ref{spaces-lemma-fibre-product-spaces-over-sheaf-with-representable-diagonal}
que $F \times_H G$ est un espace algébrique.
Il est clair que $F \times_H G$
est un produit fibré dans la catégorie des espaces algébriques au-dessus de $S$,
puisque celle-ci est une sous-catégorie pleine de la catégorie
des (pré)faisceaux d'ensembles sur $(\Sch/S)_{fppf}$.
\end{proof}






\section{Recollement d'espaces algébriques}
\label{spaces-section-glueing-algebraic-spaces}

\noindent
Dans cette section, nous commençons véritablement à abuser des notations et à ne plus
distinguer les schémas des espaces qu'ils représentent.

\begin{lemma}
\label{spaces-lemma-coproduct-sheaves-open-and-closed}
Soit $S \in \Ob(\Sch_{fppf})$. Soient $F$ et $G$ des faisceaux sur
$(\Sch/S)_{fppf}^{opp}$ et notons $F \amalg G$ leur somme
dans la catégorie des faisceaux. Le morphisme $F \to F \amalg G$ est représentable par
des immersions à la fois ouvertes et fermées.
\end{lemma}

\begin{proof}
Soient $U$ un schéma et $\xi \in (F \amalg G)(U)$. Rappelons que la
somme dans la catégorie des faisceaux est le faisceau associé au
préfaisceau somme (Sites, Lemme \ref{sites-lemma-colimit-sheaves}).
Il existe donc un recouvrement fppf $\{g_i : U_i \to U\}_{i \in I}$
et une décomposition en somme disjointe $I = I' \amalg I''$ telle que
$U_i \to U \to F \amalg G$ se factorise par $F$, resp.\ par $G$,
si et seulement si $i \in I'$, resp.\ $i \in I''$. Puisque l'intersection de $F$ et de
$G$ dans $F \amalg G$ est vide, on en déduit que
$U_i \times_U U_j$ est vide si $i \in I'$ et $j \in I''$.
Par conséquent, $U' = \bigcup_{i \in I'} g_i(U_i)$ et
$U'' = \bigcup_{i \in I''} g_i(U_i)$ sont des sous-schémas ouverts
disjoints de $U$ (Morphismes, Lemme \ref{morphisms-lemma-fppf-open})
tels que $U = U' \amalg U''$.
Nous omettons la vérification de l'égalité $U' = U \times_{F \amalg G} F$.
\end{proof}

\begin{lemma}
\label{spaces-lemma-representable-sheaf-coproduct-sheaves}
Soit $S \in \Ob(\Sch_{fppf})$.
Soit $U \in \Ob((\Sch/S)_{fppf})$.
Étant donnés un ensemble $I$ et des faisceaux $F_i$ sur $\Ob((\Sch/S)_{fppf})$,
si $U \cong \coprod_{i\in I} F_i$
comme faisceaux, chaque $F_i$ est représentable par un sous-schéma
$U_i$ à la fois ouvert et fermé, et $U \cong \coprod U_i$ comme schémas.
\end{lemma}

\begin{proof}
D'après le Lemme \ref{spaces-lemma-coproduct-sheaves-open-and-closed},
le morphisme $F_i \to U$ est représentable par des immersions à la fois ouvertes et fermées.
Ainsi, $F_i$ est représentable par un sous-schéma $U_i$ à la fois ouvert et fermé de $U$.
On a $U = \coprod U_i$, puisque $U \cong \coprod F_i$
comme faisceaux et que l'égalité peut être vérifiée sur les points.
\end{proof}

\begin{lemma}
\label{spaces-lemma-algebraic-space-coproduct-sheaves}
Soit $S \in \Ob(\Sch_{fppf})$.
Soit $F$ un espace algébrique au-dessus de $S$.
Étant donnés un ensemble $I$ et des faisceaux $F_i$ sur
$\Ob((\Sch/S)_{fppf})$,
si $F \cong \coprod_{i\in I} F_i$ comme faisceaux,
chaque $F_i$ est un espace algébrique au-dessus de $S$.
\end{lemma}

\begin{proof}
La représentabilité de $F \to F \times F$ implique que chacun des morphismes diagonaux
$F_i \to F_i \times F_i$ est représentable (cela résulte immédiatement des définitions
et de l'égalité $F \times_{(F \times F)} (F_i \times F_i) = F_i$).
Choisissons un schéma $U$ dans $(\Sch/S)_{fppf}$ et un morphisme
\'etale surjectif $U \to F$ (dont l'existence résulte de l'hypothèse).
Le morphisme $U \times_F F_i \to F_i$ obtenu par changement de base est \'etale surjectif
d'après le Lemme \ref{spaces-lemma-base-change-representable-transformations-property}.
D'autre part, $U \times_F F_i$ est un schéma d'après le
Lemme \ref{spaces-lemma-coproduct-sheaves-open-and-closed}.
Nous avons donc vérifié toutes les conditions de la
Définition \ref{spaces-definition-algebraic-space},
et $F_i$ est un espace algébrique.
\end{proof}

\noindent
La condition portant sur la taille de $I$ et des $F_i$ dans le
lemme suivant peut être ignorée par ceux que ne préoccupent pas les
questions de théorie des ensembles.

\begin{lemma}
\label{spaces-lemma-coproduct-algebraic-spaces}
Soit $S \in \Ob(\Sch_{fppf})$.
Supposons donnés un ensemble $I$ et des espaces algébriques $F_i$, $i \in I$.
Alors $F = \coprod_{i \in I} F_i$ est un espace algébrique,
pourvu que $I$ et les $F_i$ ne soient pas trop «~grands~» : par exemple, si nous
pouvons choisir des morphismes \'etales surjectifs $U_i \to F_i$ tels que
$\coprod_{i \in I} U_i$ soit isomorphe à un objet de
$(\Sch/S)_{fppf}$, alors $F$ est un espace algébrique.
\end{lemma}

\begin{proof}
Par construction, $F$ est un faisceau. Nous omettons de vérifier que le
morphisme diagonal de $F$ est représentable. Enfin, si $U$ est un
objet de $(\Sch/S)_{fppf}$ isomorphe à $\coprod_{i \in I} U_i$,
alors il est immédiat de vérifier que le morphisme qui en résulte
$U \to \coprod F_i$ est \'etale surjectif.
\end{proof}

\noindent
Voici l'analogue de Schémas, Lemme \ref{schemes-lemma-glue-functors}.

\begin{lemma}
\label{spaces-lemma-glueing-algebraic-spaces}
Soit $S \in \Ob(\Sch_{fppf})$.
Soit $F$ un préfaisceau d'ensembles sur $(\Sch/S)_{fppf}$.
Supposons que
\begin{enumerate}
\item $F$ est un faisceau,
\item il existe un ensemble d'indices $I$
et des sous-foncteurs $F_i \subset F$ tels que
\begin{enumerate}
\item chaque $F_i$ est un espace algébrique,
\item chaque $F_i \to F$ est représentable,
\item chaque $F_i \to F$ est une immersion ouverte (voir la
Définition \ref{spaces-definition-relative-representable-property}),
\item le morphisme $\coprod F_i \to F$ est surjectif comme morphisme de faisceaux, et
\item $\coprod F_i$ est un espace algébrique (condition ensembliste, voir le
Lemme \ref{spaces-lemma-coproduct-algebraic-spaces}).
\end{enumerate}
\end{enumerate}
Alors $F$ est un espace algébrique.
\end{lemma}

\begin{proof}
Soit $T$ un objet de $(\Sch/S)_{fppf}$. Soit $T \to F$ un morphisme.
D'après les hypothèses (2)(b) et (2)(c), le produit fibré $F_i \times_F T$
est représentable par un sous-schéma ouvert $V_i \subset T$. Il s'ensuit que
$(\coprod F_i) \times_F T$ est représenté par le schéma $\coprod V_i$ au-dessus de $T$.
D'après l'hypothèse (2)(d), il existe un recouvrement fppf $\{T_j \to T\}_{j \in J}$
tel que $T_j \to T \to F$ se factorise par $F_i$, où $i = i(j)$.
Ainsi, $T_j \to T$ se factorise par le sous-schéma ouvert $V_{i(j)} \subset T$.
Puisque la famille $\{T_j \to T\}$ est surjective, il s'ensuit que
$T = \bigcup V_i$ est un recouvrement ouvert. En particulier, le morphisme
de foncteurs $\coprod F_i \to F$ est représentable
et surjectif au sens de la
Définition \ref{spaces-definition-relative-representable-property}
(voir la Remarque \ref{spaces-remark-warning} pour une discussion).

\medskip\noindent
Ensuite, soit $T' \to F$ un second morphisme ayant pour source un objet de $(\Sch/S)_{fppf}$.
Écrivons comme ci-dessus $T' = \bigcup V'_i$, où $V'_i = T' \times_F F_i$.
Pour montrer que la diagonale $F \to F \times F$ est représentable,
il faut montrer que $G = T \times_F T'$ est représentable; voir le
Lemme \ref{spaces-lemma-representable-diagonal}.
Considérons les sous-foncteurs $G_i = G \times_F F_i$.
Notons que $G_i = V_i \times_{F_i} V'_i$; il est donc représentable,
car $F_i$ est un espace algébrique.
D'après ce qui précède, les $G_i$ forment un recouvrement de Zariski de $G$.
Ainsi, d'après Schémas, Lemme \ref{schemes-lemma-glue-functors},
on voit que $G$ est représentable.

\medskip\noindent
Choisissons un schéma $U \in \Ob((\Sch/S)_{fppf})$ et un morphisme
\'etale surjectif $U \to \coprod F_i$ (dont l'existence résulte de l'hypothèse).
Nous pouvons écrire $U = \coprod U_i$, où $U_i$ est l'image réciproque de $F_i$;
voir le Lemme \ref{spaces-lemma-representable-sheaf-coproduct-sheaves}.
Nous affirmons que $U \to F$ est \'etale surjectif. La surjectivité découle
de celle de $\coprod F_i \to F$ (voir le premier paragraphe de la démonstration),
en appliquant le
Lemme \ref{spaces-lemma-composition-representable-transformations-property}.
Considérons le produit fibré $U \times_F T$, où $T \to F$ est comme
ci-dessus. Il faut montrer que $U \times_F T \to T$ est \'etale.
Puisque $U \times_F T = \coprod U_i \times_F T$, il suffit de montrer que
chacun des morphismes $U_i \times_F T \to T$ est \'etale. Comme
$U_i \times_F T = U_i \times_{F_i} V_i$, cela résulte du
fait que $U_i \to F_i$ est \'etale et que $V_i \to T$ est une immersion ouverte
(voir aussi Morphismes, Lemmes \ref{morphisms-lemma-open-immersion-etale}
et \ref{morphisms-lemma-composition-etale}).
\end{proof}














\section{Présentations d'espaces algébriques}
\label{spaces-section-presentations}

\noindent
Étant donné un espace algébrique, on peut en trouver une «~présentation~».


\begin{lemma}
\label{spaces-lemma-space-presentation}
Soit $F$ un espace algébrique au-dessus de $S$. Soit $f : U \to F$ un
morphisme \'etale surjectif d'un schéma vers $F$. Posons $R = U \times_F U$.
Alors
\begin{enumerate}
\item $j : R \to U \times_S U$ définit sur
$U$ une relation d'équivalence au-dessus de $S$ (voir
Groupoïdes, Définition \ref{groupoids-definition-equivalence-relation}).
\item les morphismes $s, t : R \to U$ sont \'etales, et
\item le diagramme
$$
\xymatrix{
R \ar@<1ex>[r] \ar@<-1ex>[r] &
U \ar[r] &
F
}
$$
est un diagramme de conoyau dans $\Sh((\Sch/S)_{fppf})$.
\end{enumerate}
\end{lemma}

\begin{proof}
Soit $T/S$ un objet de $(\Sch/S)_{fppf}$.
Alors $R(T) = \{(a, b) \in U(T) \times U(T) \mid f \circ a = f \circ b\}$,
ce qui définit une relation d'équivalence sur $U(T)$.
Les morphismes $s, t : R \to U$ sont \'etales parce que le morphisme
$U \to F$ est \'etale.

\medskip\noindent
Pour démontrer (3), montrons d'abord que
$U \to F$ est un morphisme surjectif de faisceaux; voir
Sites, Définition \ref{sites-definition-sheaves-injective-surjective}.
Soit $\xi \in F(T)$, où $T$ est comme ci-dessus. Posons $V = T \times_{\xi, F, f}U$.
Par hypothèse, $V$ est un schéma et $V \to T$ est un morphisme \'etale surjectif.
Ainsi, $\{V \to T\}$ est un recouvrement pour la topologie fppf.
Puisque, par construction, $\xi|_V$ se factorise par $U$, on
en déduit que $U \to F$ est surjectif. La surjectivité implique que
$F$ est le conoyau du diagramme d'après
Sites, Lemme \ref{sites-lemma-coequalizer-surjection}.
\end{proof}

\noindent
Ce lemme conduit aux définitions suivantes.

\begin{definition}
\label{spaces-definition-etale-equivalence-relation}
Soit $S$ un schéma. Soit $U$ un schéma au-dessus de $S$.
Une {\it relation d'équivalence \'etale} sur $U$ au-dessus de $S$
est une relation d'équivalence $j : R \to U \times_S U$
telle que $s, t : R \to U$ soient des morphismes \'etales de schémas.
\end{definition}

\begin{definition}
\label{spaces-definition-presentation}
Soit $F$ un espace algébrique au-dessus de $S$.
Une {\it présentation} de $F$ est la donnée d'un schéma
$U$ au-dessus de $S$, d'une relation d'équivalence \'etale $R$ sur $U$ au-dessus de $S$, et
d'un morphisme \'etale surjectif $U \to F$ tel que $R = U \times_F U$.
\end{definition}

\noindent
De manière équivalente, on pourrait demander l'existence d'un isomorphisme
$$
U/R \cong F
$$
où le quotient $U/R$ est défini comme dans
Groupoïdes, section \ref{groupoids-section-quotient-sheaves}.
Pour construire des espaces algébriques, nous étudierons la question réciproque, à savoir
pour quelles relations d'équivalence le faisceau quotient $U/R$ est un espace algébrique.
On verra finalement qu'il en est toujours ainsi si $R$ est une relation
d'équivalence \'etale sur $U$ au-dessus de $S$; voir le Théorème \ref{spaces-theorem-presentation}.





























\section{Espaces algébriques et relations d'équivalence}
\label{spaces-section-spaces-from-equivalence-relations}

\noindent
Supposons donnés un schéma $U$ au-dessus de $S$
et une relation d'équivalence \'etale $R$ sur $U$ au-dessus de $S$.
Nous voudrions montrer que ces données définissent un espace algébrique.
Nous allons établir une suite de lemmes montrant que le faisceau quotient $U/R$
(voir Groupoïdes, Définition \ref{groupoids-definition-quotient-sheaf})
possède toutes les propriétés que lui impose la
Définition \ref{spaces-definition-algebraic-space}.

\begin{lemma}
\label{spaces-lemma-pullback-etale-equivalence-relation}
Soit $S$ un schéma. Soit $U$ un schéma au-dessus de $S$.
Soit $j = (s, t) : R \to U \times_S U$
une relation d'équivalence \'etale sur $U$ au-dessus de $S$.
Soit $U' \to U$ un morphisme \'etale.
Soit $R'$ la restriction de $R$ à $U'$; voir
Groupoïdes, Définition \ref{groupoids-definition-restrict-relation}.
Alors $j' : R' \to U' \times_S U'$ est également une relation
d'équivalence \'etale.
\end{lemma}

\begin{proof}
Il résulte de la description de $s', t'$ dans
Groupoïdes, Lemme \ref{groupoids-lemma-restrict-groupoid}
que $s' , t' : R' \to U'$ sont \'etales,
car ce sont des composés de changements de base de morphismes \'etales
(voir Morphismes, Lemme \ref{morphisms-lemma-base-change-etale}
et \ref{morphisms-lemma-composition-etale}).
\end{proof}

\noindent
Nous utiliserons souvent le lemme suivant pour trouver des sous-espaces ouverts des espaces
algébriques. Une légère amélioration de ce lemme (sous des hypothèses plus générales) est
Amorçage, Lemme \ref{bootstrap-lemma-better-finding-opens}.

\begin{lemma}
\label{spaces-lemma-finding-opens}
Soit $S$ un schéma.
Soit $U$ un schéma au-dessus de $S$.
Soit $j = (s, t) : R \to U \times_S U$ une pré-relation.
Soit $g : U' \to U$ un morphisme.
Supposons que
\begin{enumerate}
\item $j$ est une relation d'équivalence,
\item $s, t : R \to U$ sont surjectifs, plats et
localement de présentation finie,
\item $g$ est plat et localement de présentation finie.
\end{enumerate}
Soit $R' = R|_{U'}$ la restriction de $R$ à $U'$. Alors
$U'/R' \to U/R$ est représentable et est une immersion ouverte.
\end{lemma}

\begin{proof}
D'après Groupoïdes, Lemme \ref{groupoids-lemma-restrict-relation},
le morphisme $j' = (s', t') : R' \to U' \times_S U'$
définit une relation d'équivalence. Puisque $g$ est plat et localement de
présentation finie, $g$ est aussi universellement ouvert
(Morphismes, Lemme \ref{morphisms-lemma-fppf-open}).
Pour la même raison, $s, t$ sont également universellement ouverts.
Posons $W^1 = g(U') \subset U$ et $W = t(s^{-1}(W^1))$.
Alors $W^1$ et $W$ sont ouverts dans $U$. De plus, puisque $j$ est une
relation d'équivalence, on a $t(s^{-1}(W)) = W$ (voir
Groupoïdes, Lemme \ref{groupoids-lemma-constructing-invariant-opens},
par exemple).

\medskip\noindent
D'après
Groupoïdes,
Lemme \ref{groupoids-lemma-quotient-pre-equivalence-relation-restrict},
le morphisme de faisceaux $F' = U'/R' \to F = U/R$ est injectif.
Soit $a : T \to F$ un morphisme d'un schéma vers $U/R$.
Il faut montrer que $T \times_F F'$ est représentable
par un sous-schéma ouvert de $T$.

\medskip\noindent
Le morphisme $a$ est donné par les données suivantes :
un recouvrement fppf $\{\varphi_j : T_j \to T\}_{j \in J}$ de $T$ et des
morphismes $a_j : T_j \to U$ tels que les morphismes
$$
a_j \times a_{j'} :
T_j \times_T T_{j'}
\longrightarrow
U \times_S U
$$
se factorisent par $j : R \to U \times_S U$ au moyen de morphismes (uniques)
$r_{jj'} : T_j \times_T T_{j'} \to R$. Le système
$(a_j)$ correspond à $a$ en ce sens que les diagrammes
$$
\xymatrix{
T_j \ar[r]_{a_j} \ar[d] & U \ar[d] \\
T \ar[r]^a & F
}
$$
commutent.

\medskip\noindent
Considérons les ouverts $W_j = a_j^{-1}(W) \subset T_j$.
Puisque $t(s^{-1}(W)) = W$, on voit que
$$
W_j \times_T T_{j'} =
r_{jj'}^{-1}(t^{-1}(W)) = r_{jj'}^{-1}(s^{-1}(W)) =
T_j \times_T W_{j'}.
$$
D'après
Descente, Lemme \ref{descent-lemma-open-fpqc-covering},
cela signifie qu'il existe un ouvert
$W_T \subset T$ tel que $\varphi_j^{-1}(W_T) = W_j$ pour tout $j \in J$.
Nous affirmons que $W_T \to T$ représente $T \times_F F' \to T$.

\medskip\noindent
Montrons d'abord que $W_T \to T \to F$ définit un élément de
$F'(W_T)$. Puisque $\{W_j \to W_T\}_{j \in J}$ est un
recouvrement fppf de $W_T$, il suffit de montrer que
chaque $W_j \to U \to F$ définit un élément de $F'(W_j)$ (car $F'$ est un faisceau
pour la topologie fppf). Considérons le diagramme commutatif
$$
\xymatrix{
W'_j \ar[rr] \ar[dd] \ar[rd] & & U' \ar[d]^g \\
& s^{-1}(W^1) \ar[r]_s \ar[d]^t & W^1 \ar[d] \\
W_j \ar[r]^{a_j|_{W_j}} & W \ar[r] & F
}
$$
où $W'_j = W_j \times_W s^{-1}(W^1) \times_{W^1} U'$.
Puisque $t$ et $g$ sont surjectifs, plats et localement de présentation
finie, il en est de même de $W'_j \to W_j$. Par conséquent, la restriction de
l'élément $W_j \to U \to F$ à $W'_j$ est un élément de $F'$,
comme voulu.

\medskip\noindent
Supposons que $f : T' \to T$ soit un morphisme de schémas
tel que $a|_{T'} \in F'(T')$. Il faut montrer que
$f$ se factorise par l'ouvert $W_T$. Puisque
$\{T' \times_T T_j \to T'\}$ est un recouvrement fppf de $T'$,
il suffit de montrer que chaque $T' \times_T T_j \to T$
se factorise par $W_T$. Nous pouvons donc supposer que $f$ se factorise
sous la forme $\varphi_j \circ f_j : T' \to T_j \to T$ pour un certain $j$.
Dans ce cas, la condition $a|_{T'} \in F'(T')$ signifie qu'il existe
un recouvrement fppf $\{\psi_i : T'_i \to T'\}_{i \in I}$ et des
morphismes $b_i : T'_i \to U'$ tels que
$$
\xymatrix{
T'_i \ar[r]_{b_i} \ar[d]_{f_j \circ \psi_i} & U' \ar[r]_g & U \ar[d] \\
T_j \ar[r]^{a_j} & U \ar[r] & F
}
$$
soit commutatif. Cette commutativité signifie qu'il existe un
morphisme $r'_i : T'_i \to R$ tel que
$t \circ r'_i = a_j \circ f_j \circ \psi_i$ et
$s \circ r'_i = g \circ b_i$. Il en résulte que
$\Im(f_j \circ \psi_i) \subset W_j$, ce qui conclut.
\end{proof}

\noindent
Le lemme suivant n'est pas tout à fait trivial, bien qu'il semble
devoir l'être.

\begin{lemma}
\label{spaces-lemma-when-it-works-it-works}
Soit $S$ un schéma. Soit $U$ un schéma au-dessus de $S$.
Soit $j = (s, t) : R \to U \times_S U$
une relation d'équivalence \'etale sur $U$ au-dessus de $S$.
Si le quotient $U/R$ est un espace algébrique, alors
$U \to U/R$ est \'etale et surjectif. Par conséquent,
$(U, R, U \to U/R)$ est une présentation de l'espace
algébrique $U/R$.
\end{lemma}

\begin{proof}
Notons $c : U \to U/R$ le morphisme en question.
Soit $T$ un schéma et soit $a : T \to U/R$ un morphisme.
Il faut montrer que le morphisme (de schémas)
$\pi : T \times_{a, U/R, c} U \to T$ est \'etale et surjectif.
Le morphisme $a$ correspond à un recouvrement fppf
$\{\varphi_i : T_i \to T\}$ et à des morphismes $a_i : T_i \to U$ tels
que
$a_i \times a_{i'} : T_i \times_T T_{i'} \to U \times_S U$
se factorise par $R$, et tels que $c \circ a_i = a \circ \varphi_i$.
Ainsi,
$$
T_i \times_{\varphi_i, T} T \times_{a, U/R, c} U =
T_i \times_{c \circ a_i, U/R, c} U =
T_i \times_{a_i, U} U \times_{c, U/R, c} U = T_i \times_{a_i, U, t} R.
$$
Puisque $t$ est \'etale et surjectif, on en déduit que
le changement de base de $\pi$ à $T_i$ est surjectif et \'etale.
Comme la propriété d'être surjectif et \'etale est locale
sur la base pour la topologie fpqc (voir la
Remarque \ref{spaces-remark-list-properties-fpqc-local-base}),
cela conclut.
\end{proof}

\begin{lemma}
\label{spaces-lemma-presentation-quasi-compact}
Soit $S$ un schéma.
Soit $U$ un schéma au-dessus de $S$.
Soit $j = (s, t) : R \to U \times_S U$
une relation d'équivalence \'etale sur $U$ au-dessus de $S$.
Supposons $U$ affine. Alors le quotient $F = U/R$
est un espace algébrique, et $U \to F$ est \'etale et surjectif.
\end{lemma}

\begin{proof}
Puisque $j : R \to U \times_S U$ est un monomorphisme, $j$ est séparé
(voir Schémas, Lemme \ref{schemes-lemma-monomorphism-separated}).
Comme $U$ est affine, on voit que $U \times_S U$
(qui est muni d'un monomorphisme vers le schéma affine
$U \times U$) est séparé. Il s'ensuit que $R$ est séparé.
En particulier, les morphismes $s, t$ sont séparés et \'etales.

\medskip\noindent
Puisque le composé $R \to U \times_S U \to U$ est
localement de type fini, on en déduit que
$j$ est localement de type fini (voir
Morphismes, Lemme \ref{morphisms-lemma-permanence-finite-type}).
Comme $j$ est aussi un monomorphisme, ses fibres sont finies, et
l'on voit que $j$ est localement quasi-fini d'après
Morphismes, Lemme \ref{morphisms-lemma-finite-fibre}.
En définitive, $j$ est séparé et localement quasi-fini.

\medskip\noindent
La première étape consiste à montrer que le morphisme de passage au quotient
$c : U \to F$ est représentable.
Considérons un schéma $T$ et un morphisme $a : T \to F$.
Il faut montrer que le faisceau $G = T \times_{a, F, c} U$
est représentable.
Comme on l'a vu dans les démonstrations des Lemmes \ref{spaces-lemma-finding-opens} et
\ref{spaces-lemma-when-it-works-it-works}, il existe un recouvrement fppf
$\{\varphi_i : T_i \to T\}_{i \in I}$ et des morphismes $a_i : T_i \to U$
tels que $a_i \times a_{i'} : T_i \times_T T_{i'} \to U \times_S U$
se factorise par $R$, et tels que $c \circ a_i = a \circ \varphi_i$.
Comme dans la démonstration du Lemme \ref{spaces-lemma-when-it-works-it-works}, on voit que
\begin{eqnarray*}
T_i \times_{\varphi_i, T} G & = &
T_i \times_{\varphi_i, T} T \times_{a, U/R, c} U \\
& = & T_i \times_{c \circ a_i, U/R, c} U \\
& = & T_i \times_{a_i, U} U \times_{c, U/R, c} U \\
& = & T_i \times_{a_i, U, t} R
\end{eqnarray*}
Puisque $t$ est séparé et \'etale, donc en particulier
séparé et localement quasi-fini (d'après Morphismes, Lemmes
\ref{morphisms-lemma-unramified-quasi-finite} et
\ref{morphisms-lemma-flat-unramified-etale}),
on voit que la restriction
de $G$ à chaque $T_i$ est représentable par un morphisme de schémas
$X_i \to T_i$ séparé et localement quasi-fini. D'après
Descente, Lemme \ref{descent-lemma-descent-data-sheaves},
on obtient une donnée de descente $(X_i, \varphi_{ii'})$ relativement
au recouvrement fppf $\{T_i \to T\}$. Comme chaque
$X_i \to T_i$ est séparé et localement quasi-fini, il résulte de
Compléments sur les morphismes, Lemme
\ref{more-morphisms-lemma-separated-locally-quasi-finite-morphisms-fppf-descend}
que cette donnée de descente est effective.
Ainsi, d'après
Descente, Lemme \ref{descent-lemma-descent-data-sheaves} (2),
on en déduit que $G$ est représentable, comme voulu.

\medskip\noindent
La deuxième étape de la démonstration consiste à montrer que $U \to F$ est surjectif et
\'etale. Cela résulte de ce qui précède : dans la première étape, nous avons
vu que $G = T \times_{a, F, c} U$ est un schéma au-dessus de $T$ dont les changements de base
sont les schémas $X_i \to T_i$, qui sont surjectifs et \'etales. Ainsi, $G \to T$
est surjectif et \'etale (voir la
Remarque \ref{spaces-remark-list-properties-fpqc-local-base}).
On peut également reprendre la démonstration du
Lemme \ref{spaces-lemma-when-it-works-it-works} dans la présente
situation.

\medskip\noindent
La troisième et dernière étape consiste à montrer que le morphisme diagonal $F \to F \times F$
est représentable. Observons d'abord que le diagramme
$$
\xymatrix{
R \ar[r] \ar[d]_j & F \ar[d]^\Delta \\
U \times_S U \ar[r] & F \times F
}
$$
est un carré cartésien. D'après le
Lemme \ref{spaces-lemma-product-representable-transformations}, le morphisme
$U \times_S U \to F \times F$ est représentable (notons que
$h_U \times h_U = h_{U \times_S U}$). De plus, d'après le
Lemme \ref{spaces-lemma-product-representable-transformations-property},
le morphisme $U \times_S U \to F \times F$ est surjectif
et \'etale (notons aussi que \'etale et surjectif figurent dans les listes des
Remarques \ref{spaces-remark-list-properties-fpqc-local-base}
et \ref{spaces-remark-list-properties-stable-composition}).
Il résulte soit du
Lemme \ref{spaces-lemma-base-change-representable-transformations}
et du diagramme ci-dessus, soit de l'écriture de $R \to F$ comme $R \to U \to F$ et des
Lemmes
\ref{spaces-lemma-morphism-schemes-gives-representable-transformation} et
\ref{spaces-lemma-composition-representable-transformations}, que
$R \to F$ est également représentable. Soient $T$ un schéma et
$a : T \to F \times F$ un morphisme. Il faut montrer que
$G = T \times_{a, F \times F, \Delta} F$ est représentable.
D'après ce qui précède, le morphisme (de schémas)
$$
T' = (U \times_S U) \times_{F \times F, a} T \longrightarrow T
$$
est surjectif et \'etale. Ainsi, $\{T' \to T\}$ est un recouvrement
\'etale de $T$. Notons également que
$$
T' \times_T G = T' \times_{U \times_S U, j} R
$$
comme on le voit en examinant le cube suivant
$$
\xymatrix{
& R \ar[rr] \ar[dd] & & F \ar[dd] \\
T' \times_T G \ar[rr] \ar[dd] \ar[ru] & & G \ar[dd] \ar[ru] & \\
& U \times_S U \ar'[r][rr] & & F \times F \\
T' \ar[rr] \ar[ru] & & T \ar[ru]
}
$$
On voit ainsi que la restriction de $G$ à $T'$ est représentable
par un schéma $X$, et de plus que le morphisme $X \to T'$ est
un changement de base du morphisme $j$. Par conséquent, $X \to T'$ est
séparé et localement quasi-fini (voir le deuxième paragraphe de la démonstration).
D'après Descente, Lemme \ref{descent-lemma-descent-data-sheaves},
on obtient une donnée de descente $(X, \varphi)$ relativement
au recouvrement fppf $\{T' \to T\}$. Puisque
$X \to T'$ est séparé et localement quasi-fini, il résulte de
Compléments sur les morphismes, Lemme
\ref{more-morphisms-lemma-separated-locally-quasi-finite-morphisms-fppf-descend}
que cette donnée de descente est effective.
Ainsi, d'après
Descente, Lemme \ref{descent-lemma-descent-data-sheaves} (2),
on en déduit que $G$ est représentable, comme voulu.
\end{proof}

\begin{theorem}
\label{spaces-theorem-presentation}
Soit $S$ un schéma. Soit $U$ un schéma au-dessus de $S$.
Soit $j = (s, t) : R \to U \times_S U$
une relation d'équivalence \'etale sur $U$ au-dessus de $S$.
Alors le quotient $U/R$ est un espace algébrique,
et $U \to U/R$ est \'etale et surjectif; autrement dit,
$(U, R, U \to U/R)$ est une présentation de $U/R$.
\end{theorem}

\begin{proof}
D'après le Lemme \ref{spaces-lemma-when-it-works-it-works},
il suffit de prouver que $U/R$ est un espace algébrique.
Soit $U' \to U$ un morphisme surjectif et \'etale.
Alors $\{U' \to U\}$ est en particulier un recouvrement fppf.
Soit $R'$ la restriction de $R$ à $U'$; voir
Groupoïdes, Définition \ref{groupoids-definition-restrict-relation}.
D'après
Groupoïdes, Lemme \ref{groupoids-lemma-quotient-groupoid-restrict},
on a $U/R \cong U'/R'$.
D'après le Lemme \ref{spaces-lemma-pullback-etale-equivalence-relation}, $R'$ est une
relation d'équivalence \'etale sur $U'$. Nous pouvons donc remplacer $U$ par $U'$.

\medskip\noindent
Appliquons la remarque précédente à $U' = \coprod U_i$, où
$U = \bigcup U_i$ est un recouvrement ouvert affine de $U$. Nous
pouvons donc supposer, et nous le ferons, que $U = \coprod U_i$, où
chaque $U_i$ est un schéma affine.

\medskip\noindent
Considérons la restriction $R_i$ de $R$ à $U_i$.
D'après le Lemme \ref{spaces-lemma-pullback-etale-equivalence-relation},
c'est une relation d'équivalence \'etale.
Posons $F_i = U_i/R_i$ et $F = U/R$.
Il est clair que $\coprod F_i \to F$ est surjectif.
D'après le Lemme \ref{spaces-lemma-finding-opens}, chaque $F_i \to F$
est représentable et est une immersion ouverte.
En appliquant le Lemme \ref{spaces-lemma-presentation-quasi-compact}
à $(U_i, R_i)$, on voit que $F_i$ est un espace algébrique.
Puis, d'après le Lemme \ref{spaces-lemma-when-it-works-it-works},
$U_i \to F_i$ est \'etale et surjectif.
Il résulte du Lemme \ref{spaces-lemma-coproduct-algebraic-spaces}
que $\coprod F_i$ est un espace algébrique.
Enfin, nous avons vérifié toutes les
hypothèses du Lemme \ref{spaces-lemma-glueing-algebraic-spaces},
et il s'ensuit que $F = U/R$ est un espace algébrique.
\end{proof}










\section{Espaces algébriques, reformulation}
\label{spaces-section-algebraic-spaces-retrofitted}

\noindent
Nous commençons à constituer notre arsenal de lemmes sur les espaces algébriques.
Le premier résultat affirme que, dans la Définition \ref{spaces-definition-algebraic-space},
on peut affaiblir comme suit la condition portant sur la diagonale.

\begin{lemma}
\label{spaces-lemma-etale-locally-representable-gives-space}
Soit $S$ un schéma appartenant à $\Sch_{fppf}$.
Soit $F$ un faisceau sur $(\Sch/S)_{fppf}$
tel qu'il existe $U \in \Ob((\Sch/S)_{fppf})$ et un morphisme
$U \to F$ représentable, surjectif et \'etale.
Alors $F$ est un espace algébrique.
\end{lemma}

\begin{proof}
Posons $R = U \times_F U$. C'est un schéma, car $U \to F$ est supposé représentable.
Les projections $s, t : R \to U$ sont \'etales, car $U \to F$ est supposé \'etale.
Le morphisme $j = (t, s) : R \to U \times_S U$ est un monomorphisme et définit une relation
d'équivalence, puisque $R = U \times_F U$. D'après le Théorème \ref{spaces-theorem-presentation},
le faisceau quotient $F' = U/R$ est un espace algébrique et $U \to F'$
est surjectif et \'etale. À nouveau, comme $R = U \times_F U$, on obtient
une factorisation canonique $U \to F' \to F$, et $F' \to F$ est un morphisme injectif
de faisceaux. D'autre part, $U \to F$ est surjectif comme morphisme
de faisceaux d'après le Lemme \ref{spaces-lemma-surjective-flat-locally-finite-presentation}.
Ainsi, $F' \to F$ est également surjectif, et l'on conclut que $F' = F$ est un
espace algébrique.
\end{proof}

\begin{lemma}
\label{spaces-lemma-etale-locally-representable-by-space-gives-space}
Soit $S$ un schéma appartenant à $\Sch_{fppf}$. Soit $G$ un espace
algébrique au-dessus de $S$, soit $F$ un faisceau sur $(\Sch/S)_{fppf}$, et soit
$G \to F$ un morphisme représentable de foncteurs qui est
surjectif et \'etale. Alors $F$ est un espace algébrique.
\end{lemma}

\begin{proof}
Choisissons un schéma $U$ et un morphisme surjectif \'etale $U \to G$.
Puisque $G$ est un espace algébrique, $U \to G$ est représentable.
Le composé $U \to G \to F$ est donc représentable,
surjectif et \'etale. Voir les Lemmes
\ref{spaces-lemma-composition-representable-transformations} et
\ref{spaces-lemma-composition-representable-transformations-property}.
Ainsi, $F$ est un espace algébrique d'après le
Lemme \ref{spaces-lemma-etale-locally-representable-gives-space}.
\end{proof}

\begin{lemma}
\label{spaces-lemma-representable-over-space}
\begin{slogan}
Un foncteur qui admet un morphisme représentable vers un espace algébrique est
un espace algébrique.
\end{slogan}
Soit $S$ un schéma appartenant à $\Sch_{fppf}$.
Soit $F$ un espace algébrique au-dessus de $S$.
Soit $G \to F$ un morphisme représentable de foncteurs.
Alors $G$ est un espace algébrique.
\end{lemma}

\begin{proof}
D'après le Lemme \ref{spaces-lemma-representable-transformation-to-sheaf},
on voit que $G$ est un faisceau. Le diagramme
$$
\xymatrix{
G \times_F G \ar[r] \ar[d] & F \ar[d]^{\Delta_F} \\
G \times G \ar[r] & F \times F
}
$$
est cartésien. On voit donc que $G \times_F G \to G \times G$
est représentable d'après le
Lemme \ref{spaces-lemma-base-change-representable-transformations}.
D'après le
Lemme \ref{spaces-lemma-representable-transformation-diagonal},
on voit que $G \to G \times_F G$ est représentable.
Ainsi, $\Delta_G : G \to G \times G$ est représentable comme composé
de morphismes représentables de foncteurs; voir le
Lemme \ref{spaces-lemma-composition-representable-transformations}.
Enfin, soit $U$ un objet de $(\Sch/S)_{fppf}$
et soit $U \to F$ surjectif et \'etale. Par hypothèse,
$U \times_F G$ est représentable par un schéma $U'$. D'après le
Lemme \ref{spaces-lemma-base-change-representable-transformations-property},
le morphisme $U' \to G$ est surjectif et \'etale. Cela vérifie
la dernière condition de la Définition \ref{spaces-definition-algebraic-space}, ce qui conclut.
\end{proof}

\begin{lemma}
\label{spaces-lemma-representable-morphisms-spaces-property}
Soit $S$ un schéma appartenant à $\Sch_{fppf}$.
Soient $F$, $G$ des espaces algébriques au-dessus de $S$.
Soit $G \to F$ un morphisme représentable.
Soient $U \in \Ob((\Sch/S)_{fppf})$ et $q : U \to F$
surjectif et \'etale. Posons $V = G \times_F U$.
Enfin, soit $\mathcal{P}$ une propriété des morphismes
de schémas comme dans la Définition \ref{spaces-definition-relative-representable-property}.
Alors $G \to F$ possède la propriété $\mathcal{P}$ si et seulement si
$V \to U$ possède la propriété $\mathcal{P}$.
\end{lemma}

\begin{proof}
(Ce lemme résulte des
Lemmes \ref{spaces-lemma-base-change-representable-transformations-property} et
\ref{spaces-lemma-descent-representable-transformations-property},
mais nous en donnons aussi ici une démonstration directe.)
Il résulte immédiatement des définitions que, si $G \to F$ possède la propriété
$\mathcal{P}$, alors $V \to U$ possède la propriété $\mathcal{P}$.
Réciproquement, supposons que $V \to U$ possède la propriété $\mathcal{P}$.
Soit $T \to F$ un morphisme d'un schéma vers $F$.
Soit $T' = T \times_F G$, qui est un schéma puisque $G \to F$ est
représentable. Il faut montrer que $T' \to T$ possède la propriété $\mathcal{P}$.
Considérons le diagramme commutatif de schémas
$$
\xymatrix{
V \ar[d] & T \times_F V \ar[d] \ar[l] \ar[r] &
T \times_F G \ar[d] \ar@{=}[r] & T' \\
U & T \times_F U \ar[l] \ar[r] & T
}
$$
où les deux carrés sont cartésiens. On en déduit que
la flèche médiane possède la propriété $\mathcal{P}$ en tant que changement de base
de $V \to U$. Enfin, $\{T \times_F U \to T\}$ est un recouvrement fppf,
car il est surjectif \'etale; on en déduit donc que
$T' \to T$ possède la propriété $\mathcal{P}$, puisque celle-ci est locale sur la
base pour la topologie fppf.
\end{proof}

\begin{lemma}
\label{spaces-lemma-morphism-sheaves-with-P-effective-descent-etale}
Soit $S$ un schéma appartenant à $\Sch_{fppf}$.
Soit $G \to F$ un morphisme de préfaisceaux sur $(\Sch/S)_{fppf}$.
Soit $\mathcal{P}$ une propriété des morphismes de schémas.
Supposons que
\begin{enumerate}
\item $\mathcal{P}$ soit stable par tout changement de base, locale sur la
base pour la topologie fppf, et que les morphismes de type $\mathcal{P}$ satisfassent à la descente pour les recouvrements fppf;
voir Descente, Définition \ref{descent-definition-descending-types-morphisms},
\item $G$ soit un faisceau,
\item $F$ soit un espace algébrique,
\item il existe $U \in \Ob((\Sch/S)_{fppf})$
et un morphisme surjectif \'etale $U \to F$ tels que
$V = G \times_F U$ soit représentable, et
\item $V \to U$ possède la propriété $\mathcal{P}$.
\end{enumerate}
Alors $G$ est un espace algébrique, $G \to F$ est représentable et possède la propriété
$\mathcal{P}$.
\end{lemma}

\begin{proof}
Soit $T$ un schéma et soit $T \to F$ un morphisme. Alors
$U \times_F T \to T$ est surjectif \'etale; ainsi, $\{U \times_F T \to T\}$ est
un recouvrement pour la topologie \'etale. Considérons
$$
W = G \times_F (U \times_F T) = V \times_F T = V \times_U (U \times_F T).
$$
C'est un schéma puisque $F$ est un espace algébrique. Le morphisme
$W \to U \times_F T$ possède la propriété $\mathcal{P}$, car c'est un
changement de base de $V \to U$. Il existe un isomorphisme
\begin{align*}
W \times_T (U \times_F T) & =
(G \times_F (U \times_F T)) \times_T (U \times_F T) \\
& = (U \times_F T) \times_T (G \times_F (U \times_F T)) \\
& = (U \times_F T) \times_T W
\end{align*}
au-dessus de $(U \times_F T) \times_T (U \times_F T)$. L'égalité médiane envoie
$((g, (u_1, t)), (u_2, t))$ sur $((u_1, t), (g, (u_2, t)))$.
Cela définit une donnée de descente pour $W/U \times_F T/T$; voir
Descente, Définition \ref{descent-definition-descent-datum}.
Cela résulte du
Lemme \ref{descent-lemma-descent-data-sheaves} de Descente.
Plus précisément, on dispose du faisceau $G \times_F T$, dont le
changement de base à $U \times_F T$ est représenté par $W$, et l'isomorphisme
ci-dessus est celui de la démonstration du
Lemme \ref{descent-lemma-descent-data-sheaves} de Descente.
Par hypothèse sur $\mathcal{P}$, la donnée de descente ci-dessus est représentable.
La dernière assertion du
Lemme \ref{descent-lemma-descent-data-sheaves} de Descente
montre donc que $G \times_F T$ est représentable. Cela prouve que
$G \to F$ est un morphisme représentable de foncteurs.

\medskip\noindent
Comme $G \to F$ est représentable, on voit que $G$ est un espace algébrique d'après le
Lemme \ref{spaces-lemma-representable-over-space}. Le fait que $G \to F$ possède
la propriété $\mathcal{P}$ résulte maintenant du
Lemme \ref{spaces-lemma-representable-morphisms-spaces-property}.
\end{proof}

\begin{lemma}
\label{spaces-lemma-lift-morphism-presentations}
Soit $S$ un schéma appartenant à $\Sch_{fppf}$.
Soient $F, G$ des espaces algébriques au-dessus de $S$.
Soit $a : F \to G$ un morphisme.
Étant donnés $V \in \Ob((\Sch/S)_{fppf})$
et un morphisme surjectif \'etale $q : V \to G$, il existe
$U \in \Ob((\Sch/S)_{fppf})$
et un diagramme commutatif
$$
\xymatrix{
U \ar[d]_p \ar[r]_\alpha &
V \ar[d]^q \\
F \ar[r]^a & G
}
$$
où $p$ est surjectif et \'etale.
\end{lemma}

\begin{proof}
Choisissons d'abord $W \in \Ob((\Sch/S)_{fppf})$
et un morphisme surjectif \'etale $W \to F$.
Posons ensuite $U = W \times_G V$. Puisque $G$ est un espace algébrique,
on voit que $U$ est isomorphe à un objet de $(\Sch/S)_{fppf}$.
Comme $q$ est surjectif \'etale, on voit que $U \to W$ est surjectif
\'etale (voir le
Lemme \ref{spaces-lemma-base-change-representable-transformations-property}).
Ainsi, $U \to F$ est surjectif \'etale comme composé de morphismes surjectifs
\'etales (voir le
Lemme \ref{spaces-lemma-composition-representable-transformations-property}).
\end{proof}



























\section{Immersions et recouvrements de Zariski des espaces algébriques}
\sectionmark{Immersions et recouvrements de Zariski}
\label{spaces-section-Zariski}

\noindent
À ce stade, un phénomène intéressant se produit. Nous avons déjà défini
la notion d'immersion ouverte d'espaces algébriques (au moyen de la
Définition \ref{spaces-definition-relative-representable-property}),
mais il nous reste à définir la notion de {\it point}\footnote{Nous
associerons un espace topologique à un espace algébrique dans
Propriétés des espaces algébriques, section
\ref{spaces-properties-section-points},
et ses ouverts correspondront exactement aux sous-espaces ouverts définis ci-dessous.}.
Ainsi, la {\it topologie de Zariski} d'un espace algébrique
a déjà été définie, alors qu'il n'y a encore aucun espace !

\medskip\noindent
Peut-être de façon superflue, nous introduisons formellement les immersions comme suit.

\begin{definition}
\label{spaces-definition-immersion}
Soit $S \in \Ob(\Sch_{fppf})$ un schéma.
Soit $F$ un espace algébrique au-dessus de $S$.
\begin{enumerate}
\item Un morphisme d'espaces algébriques au-dessus de $S$
est appelé {\it immersion ouverte} s'il est représentable et si c'est une immersion ouverte
au sens de la Définition \ref{spaces-definition-relative-representable-property}.
\item Un {\it sous-espace ouvert} de $F$ est un sous-foncteur $F' \subset F$
tel que $F'$ soit un espace algébrique et que $F' \to F$ soit une
immersion ouverte.
\item Un morphisme d'espaces algébriques au-dessus de $S$
est appelé {\it immersion fermée} s'il est représentable et si c'est une immersion
fermée au sens de la
Définition \ref{spaces-definition-relative-representable-property}.
\item Un {\it sous-espace fermé} de $F$ est un sous-foncteur $F' \subset F$
tel que $F'$ soit un espace algébrique et que $F' \to F$ soit une
immersion fermée.
\item Un morphisme d'espaces algébriques au-dessus de $S$
est appelé {\it immersion} s'il est représentable et si c'est une immersion
au sens de la Définition \ref{spaces-definition-relative-representable-property}.
\item Un {\it sous-espace localement fermé} de $F$ est un sous-foncteur $F' \subset F$
tel que $F'$ soit un espace algébrique et que $F' \to F$ soit une
immersion.
\end{enumerate}
\end{definition}

\noindent
Notons que ces définitions ont un sens, puisqu'une immersion
est en particulier un monomorphisme (voir
Schémas, Lemme \ref{schemes-lemma-immersions-monomorphisms}
et Lemme \ref{spaces-lemma-representable-transformations-property-implication}),
et que, par conséquent, l'image d'une
immersion $G \to F$ d'espaces algébriques est un sous-foncteur $F' \subset F$
qui est (canoniquement) isomorphe à $G$. Ainsi, une partie de la discussion
de Schémas, section \ref{schemes-section-immersions}, se transpose au
cadre des espaces algébriques.

\begin{lemma}
\label{spaces-lemma-composition-immersions}
Soit $S \in \Ob(\Sch_{fppf})$ un schéma.
Un composé d'immersions (fermées, resp.\ ouvertes)
d'espaces algébriques au-dessus de $S$ est une immersion (fermée, resp.\ ouverte)
d'espaces algébriques au-dessus de $S$.
\end{lemma}

\begin{proof}
Voir le Lemme \ref{spaces-lemma-composition-representable-transformations-property} et les
Remarques \ref{spaces-remark-list-properties-fpqc-local-base} (voir la toute dernière ligne de
cette remarque) et \ref{spaces-remark-list-properties-stable-composition}.
\end{proof}

\begin{lemma}
\label{spaces-lemma-base-change-immersions}
Soit $S \in \Ob(\Sch_{fppf})$ un schéma.
Un changement de base d'une immersion (fermée, resp.\ ouverte)
d'espaces algébriques au-dessus de $S$ est une immersion (fermée, resp.\ ouverte)
d'espaces algébriques au-dessus de $S$.
\end{lemma}

\begin{proof}
Voir le Lemme \ref{spaces-lemma-base-change-representable-transformations-property} et la
Remarque \ref{spaces-remark-list-properties-fpqc-local-base} (voir la toute dernière ligne de
cette remarque).
\end{proof}

\begin{lemma}
\label{spaces-lemma-sub-subspaces}
Soit $S \in \Ob(\Sch_{fppf})$ un schéma.
Soit $F$ un espace algébrique au-dessus de $S$. Soient $F_1$, $F_2$ des
sous-espaces localement fermés de $F$. Si $F_1 \subset F_2$ comme sous-foncteurs
de $F$, alors $F_1$ est un sous-espace localement fermé de $F_2$.
Il en va de même pour les sous-espaces fermés et ouverts.
\end{lemma}

\begin{proof}
Soit $T \to F_2$ un morphisme, où $T$ est un schéma.
Puisque $F_2 \to F$ est un monomorphisme, on voit que
$T \times_{F_2} F_1 = T \times_F F_1$. Le lemme en résulte
formellement.
\end{proof}

\noindent
Définissons formellement la notion de recouvrement ouvert de Zariski
d'espaces algébriques. Notons que, dans le Lemme \ref{spaces-lemma-glueing-algebraic-spaces},
nous avons déjà rencontré de tels recouvrements ouverts comme méthode pour
construire des espaces algébriques.

\begin{definition}
\label{spaces-definition-Zariski-open-covering}
Soit $S \in \Ob(\Sch_{fppf})$ un schéma.
Soit $F$ un espace algébrique au-dessus de $S$.
Un {\it recouvrement de Zariski} $\{F_i \subset F\}_{i \in I}$ de $F$
est constitué d'un ensemble $I$ et d'une famille de sous-espaces ouverts
$F_i \subset F$ telle que $\coprod F_i \to F$ soit un morphisme surjectif
de faisceaux.
\end{definition}

\noindent
Notons que, si $T$ est un schéma
et si $a : T \to F$ est un morphisme, chacun des produits fibrés
$T \times_F F_i$ s'identifie à un sous-schéma ouvert
$T_i \subset T$. La dernière condition de la définition signifie
exactement que $T = \bigcup_{i \in I} T_i$.

\medskip\noindent
Il est clair que la collection $F_{Zar}$ des sous-espaces ouverts de
$F$ est un ensemble (puisque $(\Sch/S)_{fppf}$ est un site, donc un ensemble).
De plus, on peut faire de $F_{Zar}$ une catégorie en prenant pour
morphismes les inclusions de sous-foncteurs (qui sont automatiquement des immersions
ouvertes d'après le Lemme \ref{spaces-lemma-sub-subspaces}). Enfin, la
Définition \ref{spaces-definition-Zariski-open-covering}
fournit la notion de recouvrement de Zariski $\{F_i \to F'\}_{i \in I}$
dans la catégorie $F_{Zar}$. Ainsi, comme dans le cas d'un espace topologique
(voir Sites, Exemple \ref{sites-example-site-topological}),
en choisissant convenablement un ensemble de recouvrements,
on peut obtenir un site de Zariski de l'espace algébrique $F$.

\begin{definition}
\label{spaces-definition-small-Zariski-site}
Soit $S \in \Ob(\Sch_{fppf})$ un schéma. Soit $F$ un espace algébrique au-dessus de
$S$. Un {\it petit site de Zariski $F_{Zar}$} d'un espace algébrique $F$ est l'un
des sites décrits ci-dessus.
\end{definition}

\noindent
Cela donne donc un sens à l'expression : une propriété est vraie
localement pour la topologie de Zariski sur un espace algébrique; c'est ainsi que nous utiliserons cette
notion. En général, la topologie de Zariski n'est pas assez fine pour nos
besoins. Par exemple, on peut considérer la catégorie des faisceaux de Zariski
sur un espace algébrique. Il s'avérera que ce n'est pas la
bonne notion à considérer, même pour les faisceaux quasi-cohérents.
On n'obtient le résultat recherché qu'en utilisant le site \'etale ou fppf de
$F$ pour définir les faisceaux quasi-cohérents.











\section{Conditions de séparation des espaces algébriques}
\label{spaces-section-separation}

\noindent
Une condition de séparation pour un espace algébrique $F$ est une condition
portant sur le morphisme diagonal $F \to F \times F$. Commençons par
énumérer les propriétés que la diagonale possède automatiquement.
Comme la diagonale est représentable par définition, l'énoncé du lemme suivant
a un sens (grâce à la
Définition \ref{spaces-definition-relative-representable-property}).

\begin{lemma}
\label{spaces-lemma-properties-diagonal}
Soit $S$ un schéma appartenant à $\Sch_{fppf}$.
Soit $F$ un espace algébrique sur $S$.
Soit $\Delta : F \to F \times F$ le morphisme diagonal.
Alors
\begin{enumerate}
\item $\Delta$ est localement de type fini,
\item $\Delta$ est un monomorphisme,
\item $\Delta$ est séparé, et
\item $\Delta$ est localement quasi-fini.
\end{enumerate}
\end{lemma}

\begin{proof}
Soit $F = U/R$ une présentation de $F$.
Comme dans la démonstration du Lemme \ref{spaces-lemma-presentation-quasi-compact}, le diagramme
$$
\xymatrix{
R \ar[r] \ar[d]_j & F \ar[d]^\Delta \\
U \times_S U \ar[r] & F \times F
}
$$
est cartésien. Par conséquent, d'après le
Lemme \ref{spaces-lemma-representable-morphisms-spaces-property},
il suffit de montrer que $j$ possède les propriétés énumérées dans le lemme.
(Notons que chacune des propriétés (1) -- (4) figure dans les listes
des Remarques \ref{spaces-remark-list-properties-stable-base-change}
et \ref{spaces-remark-list-properties-fpqc-local-base}.)
Comme $j$ définit une relation d'équivalence, c'est un monomorphisme.
Il est donc séparé d'après
Schémas, Lemme \ref{schemes-lemma-monomorphism-separated}.
Comme $R$ est une relation d'équivalence \'etale, on voit que
$s, t : R \to U$ sont \'etales. Ainsi, $s, t$ sont localement de
type fini. Il résulte alors de
Morphismes, Lemme \ref{morphisms-lemma-permanence-finite-type}, que
$j$ est localement de type fini. Enfin, comme c'est un monomorphisme,
ses fibres sont finies. On en conclut qu'il est localement quasi-fini d'après
Morphismes, Lemme \ref{morphisms-lemma-finite-fibre}.
\end{proof}

\noindent
Voici quelques types usuels de conditions de séparation, relativement au schéma
de base $S$. Il existe aussi une notion absolue de ces conditions, dont nous
discuterons dans
Propriétés des espaces algébriques, section \ref{spaces-properties-section-separation}.
En outre, nous étudierons les conditions de séparation d'un morphisme
d'espaces algébriques dans
Morphismes d'espaces algébriques, section \ref{spaces-morphisms-section-separation-axioms}.

\begin{definition}
\label{spaces-definition-separated}
Soit $S$ un schéma appartenant à $\Sch_{fppf}$.
Soit $F$ un espace algébrique sur $S$.
Soit $\Delta : F \to F \times F$ le morphisme diagonal.
\begin{enumerate}
\item On dit que $F$ est {\it séparé sur $S$} si $\Delta$ est une immersion fermée.
\item On dit que $F$ est {\it localement séparé sur $S$}\footnote{Dans la
littérature, cela désigne souvent des espaces algébriques quasi-séparés et
localement séparés.} si $\Delta$ est une
immersion.
\item On dit que $F$ est {\it quasi-séparé sur $S$} si $\Delta$ est quasi-compact.
\item On dit que $F$ est {\it localement quasi-séparé sur $S$ pour la topologie de Zariski}\footnote{Cette
définition a été proposée par B.\ Conrad.} s'il
existe un recouvrement de Zariski $F = \bigcup_{i \in I} F_i$ tel que
chacun des $F_i$ soit quasi-séparé.
\end{enumerate}
\end{definition}

\noindent
Notons que, si la diagonale est quasi-compacte (lorsque $F$ est séparé ou
quasi-séparé), elle est en fait
quasi-finie et séparée, donc quasi-affine (d'après Compléments sur les morphismes,
Lemme \ref{more-morphisms-lemma-quasi-finite-separated-quasi-affine}).








\section{Exemples d'espaces algébriques}
\label{spaces-section-examples}

\noindent
Dans cette section, nous construisons quelques exemples d'espaces algébriques.
Certains de ces exemples ont été suggérés par B.\ Conrad.
Comme nous ne disposons pas encore de beaucoup de théorie,
l'exposé est quelque peu malaisé par endroits.

\begin{example}
\label{spaces-example-affine-line-involution}
Soit $k$ un corps de caractéristique $\not = 2$. Soit $U = \mathbf{A}^1_k$. Posons
$$
j : R = \Delta \amalg \Gamma \longrightarrow U \times_k U
$$
où $\Delta = \{(x, x) \mid x \in \mathbf{A}^1_k\}$ et
$\Gamma = \{(x, -x) \mid x \in \mathbf{A}^1_k, x \not = 0\}$.
Il est clair que $s, t : R \to U$ sont \'etales, et donc que
$j$ est une relation d'équivalence \'etale.
Le quotient $X = U/R$ est un espace algébrique d'après le
Théorème \ref{spaces-theorem-presentation}.
Comme $R$ est quasi-compact, on voit que $X$ est quasi-séparé.
En revanche, $X$ n'est pas localement séparé, car
le morphisme $j$ n'est pas une immersion.
\end{example}

\begin{example}
\label{spaces-example-non-representable-descent}
Soit $k$ un corps. Soit $k'/k$ une extension galoisienne de degré $2$
avec $\text{Gal}(k'/k) = \{1, \sigma\}$. Soient $S = \Spec(k[x])$
et $U = \Spec(k'[x])$. Notons que
$$
U \times_S U =
\Spec((k' \otimes_k k')[x]) =
\Delta(U) \amalg \Delta'(U)
$$
où $\Delta' = (1, \sigma) : U \to U \times_S U$. Prenons
$$
R = \Delta(U) \amalg \Delta'(U \setminus \{0_U\})
$$
où $0_U \in U$ désigne le point $k'$-rationnel dont la coordonnée $x$ est nulle.
Il est facile de voir que $R$ est une relation d'équivalence \'etale sur $U$ au-dessus de $S$
et donc que $X = U/R$ est un espace algébrique d'après le
Théorème \ref{spaces-theorem-presentation}. Voici quelques propriétés de $X$ (dont certaines
n'auront de sens que plus tard) :
\begin{enumerate}
\item $X \to S$ est un isomorphisme au-dessus de $S \setminus \{0_S\}$,
\item le morphisme $X \to S$ est \'etale (voir
Propriétés des espaces algébriques,
Définition \ref{spaces-properties-definition-etale}),
\item la fibre $0_X$ de $X \to S$ au-dessus de $0_S$ est isomorphe à
$\Spec(k') = 0_U$,
\item $X$ n'est pas un schéma, car, s'il l'était, $\mathcal{O}_{X, 0_X}$
serait un anneau local intègre $(\mathcal{O}, \mathfrak m, \kappa)$ de
corps des fractions $k(x)$, avec $x \in \mathfrak m$, et de corps résiduel
$\kappa = k'$, ce qui est impossible,
\item $X$ n'est pas séparé, mais il est
localement séparé et quasi-séparé,
\item il existe un morphisme surjectif, fini et \'etale $S' \to S$
tel que le changement de base $X' = S' \times_S X$ soit un schéma (en effet, si
l'on effectue le changement de base vers $S' = \Spec(k'[x])$, alors $U$ se scinde en
deux copies de $S'$ et $X'$ devient isomorphe à la droite affine dont le point
$0$ est dédoublé, voir
Schémas, Exemple \ref{schemes-example-affine-space-zero-doubled}), et
\item si l'on considère $X$ comme un espace algébrique de type fini sur
$\Spec(k)$, alors, de même, le changement de base $X_{k'}$ est un schéma,
mais $X$ n'est pas un schéma.
\end{enumerate}
En particulier, on obtient ainsi un exemple de donnée de descente pour des schémas,
relativement au recouvrement $\{\Spec(k') \to \Spec(k)\}$,
qui n'est pas effective.
\end{example}

\noindent
Voir aussi Exemples,
Lemme \ref{examples-lemma-non-effective-descent-projective},
qui montre que les données de descente ne sont pas nécessairement effectives,
même pour un morphisme projectif de schémas. Cet exemple
fournit un espace algébrique lisse et séparé de dimension 3
sur ${\mathbf C}$ qui n'est pas un schéma.

\medskip\noindent
Nous utiliserons le lemme suivant comme moyen commode de construire
des espaces algébriques comme quotients de schémas par des actions libres de groupes.

\begin{lemma}
\label{spaces-lemma-quotient}
Soit $U \to S$ un morphisme de $\Sch_{fppf}$.
Soit $G$ un groupe abstrait. Soit $G \to \text{Aut}_S(U)$
un homomorphisme de groupes. Supposons que
\begin{itemize}
\item[(*)] si $u \in U$ est un point et si $g(u) = u$
pour un élément non neutre $g \in G$, alors $g$
induit un automorphisme non trivial de $\kappa(u)$.
\end{itemize}
Alors
$$
j :
R = \coprod\nolimits_{g \in G} U
\longrightarrow
U \times_S U,
\quad
(g, x) \longmapsto (g(x), x)
$$
est une relation d'équivalence \'etale et, par conséquent,
$$
F = U/R
$$
est un espace algébrique d'après le Théorème \ref{spaces-theorem-presentation}.
\end{lemma}

\begin{proof}
Dans l'énoncé du lemme, le symbole $\text{Aut}_S(U)$ désigne
le groupe des automorphismes de $U$ au-dessus de $S$.
Supposons $(*)$ vérifiée. Montrons que
$$
j :
R = \coprod\nolimits_{g \in G} U
\longrightarrow
U \times_S U,
\quad
(g, x) \longmapsto (g(x), x)
$$
est un monomorphisme. Cela signifie que, si $T$ est un schéma
non vide et si $h : T \to U$ est un point à valeurs dans $T$ tel que
$g \circ h = g' \circ h$, alors $g = g'$. Supposons
$T \not = \emptyset$, $h : T \to U$ et $g \circ h = g' \circ h$.
Soit $t \in T$. Considérons le composé
$\Spec(\kappa(t)) \to \Spec(\kappa(h(t))) \to U$.
On en conclut que $g^{-1} \circ g'$ fixe $u = h(t)$ et
agit comme l'identité sur son corps résiduel. On a donc $g = g'$ par $(*)$.

\medskip\noindent
Ainsi, si $(*)$ est vérifiée, on voit que $j$ est une relation (voir
Schémas en groupoïdes, Définition \ref{groupoids-definition-equivalence-relation}).
De plus, c'est une relation d'équivalence puisque, sur les points à valeurs dans $T$
pour un schéma connexe $T$, on a
$R(T) = G \times U(T) \to U(T) \times U(T)$ (rappelons que nous
travaillons toujours au-dessus de $S$). En outre, les morphismes $s, t : R \to U$ sont \'etales,
puisque $R$ est une somme disjointe de copies de $U$.
Cela prouve que $j : R \to U \times_S U$ est une relation d'équivalence \'etale.
\end{proof}

\noindent
Étant donnés un schéma $U$ et une action d'un groupe $G$ sur $U$, on dit que l'action
de $G$ sur $U$ est {\it libre} si la condition $(*)$ du Lemme \ref{spaces-lemma-quotient}
est vérifiée. Cela équivaut à la notion d'action libre du schéma en groupes
constant $G_S$ sur $U$, telle qu'elle est définie dans
Schémas en groupoïdes, Définition \ref{groupoids-definition-free-action}.
On peut interpréter le lemme comme affirmant que les quotients de schémas par des
actions libres de groupes existent dans la catégorie des espaces algébriques.

\begin{definition}
\label{spaces-definition-quotient}
Reprenons les notations $U \to S$, $G$, $R$ du Lemme \ref{spaces-lemma-quotient}.
Si l'action de $G$ sur $U$ satisfait $(*)$, on dit que $G$ {\it agit librement}
sur le schéma $U$. Dans ce cas, l'espace algébrique $U/R$ est noté
$U/G$ et appelé le {\it quotient de $U$ par $G$}.
\end{definition}

\noindent
Cette notation est compatible avec la notation $U/G$ introduite dans
Schémas en groupoïdes, Définition \ref{groupoids-definition-quotient-sheaf}.
Nous donnerons plus tard un sens au quotient comme champ algébrique sans
aucune hypothèse sur l'action; nous utiliserons alors la
notation $[U/G]$. Avant d'aborder les exemples, démontrons
encore quelques lemmes afin de faciliter la discussion. Le lemme suivant traite des
diverses conditions de séparation de ce quotient lorsque $G$ est fini.

\begin{lemma}
\label{spaces-lemma-quotient-finite-separated}
Reprenons les notations et les hypothèses du Lemme \ref{spaces-lemma-quotient}.
Supposons $G$ fini. Alors
\begin{enumerate}
\item si $U \to S$ est quasi-séparé, alors $U/G$ est quasi-séparé
sur $S$, et
\item si $U \to S$ est séparé, alors $U/G$ est séparé sur $S$.
\end{enumerate}
\end{lemma}

\begin{proof}
Dans la démonstration du Lemme \ref{spaces-lemma-properties-diagonal},
nous avons vu qu'il suffit d'établir les
propriétés correspondantes pour le morphisme $j : R \to U \times_S U$.
Si $U \to S$ est quasi-séparé, alors, pour tout ouvert affine $V \subset U$
dont l'image est contenue dans un ouvert affine de $S$,
les ouverts $g(V) \cap V$ sont quasi-compacts. Il s'ensuit que $j$ est
quasi-compact.
Si $U \to S$ est séparé, la diagonale $\Delta_{U/S}$ est une immersion
fermée. Ainsi, $j : R \to U \times_S U$ est une somme finie
d'immersions fermées aux images disjointes. Ainsi, $j$ est une immersion fermée.
\end{proof}

\begin{lemma}
\label{spaces-lemma-quotient-field-map}
Reprenons les notations et les hypothèses du Lemme \ref{spaces-lemma-quotient}.
Si $\Spec(k) \to U/G$ est un morphisme, alors il existe
\begin{enumerate}
\item une extension galoisienne finie $k'/k$,
\item un sous-groupe fini $H \subset G$,
\item un isomorphisme $H \to \text{Gal}(k'/k)$, et
\item un morphisme $H$-équivariant $\Spec(k') \to U$.
\end{enumerate}
Réciproquement, de telles données déterminent un morphisme $\Spec(k) \to U/G$.
\end{lemma}

\begin{proof}
Considérons le produit fibré $V = \Spec(k) \times_{U/G} U$.
Voici un diagramme
$$
\xymatrix{
V \ar[r] \ar[d] & U \ar[d] \\
\Spec(k) \ar[r] & U/G
}
$$
Alors $V$ est un schéma non vide \'etale sur $\Spec(k)$ et est donc une
réunion disjointe $V = \coprod_{i \in I} \Spec(k_i)$
de spectres de corps $k_i$, extensions finies séparables de $k$
(Morphismes de schémas, Lemme \ref{morphisms-lemma-etale-over-field}).
On a
\begin{align*}
V \times_{\Spec(k)} V
& =
(\Spec(k) \times_{U/G} U) \times_{\Spec(k)}(\Spec(k) \times_{U/G} U) \\
& = 
\Spec(k) \times_{U/G} U \times_{U/G} U \\
& =
\Spec(k) \times_{U/G} U \times G \\
& =
V \times G
\end{align*}
L'action de $G$ sur $U$ induit une action $a : G \times V \to V$.
L'égalité affichée signifie que
$G \times V \to V \times_{\Spec(k)} V$, $(g, v) \mapsto (a(g, v), v)$
est un isomorphisme. En particulier, on voit que, pour tout $i$, on a
un isomorphisme $H_i \times \Spec(k_i) \to \Spec(k_i \otimes_k k_i)$,
où $H_i \subset G$ est le sous-groupe des éléments qui fixent $i \in I$.
Ainsi, $H_i$ est fini et est le groupe de Galois de $k_i/k$.
Nous omettons la construction réciproque.
\end{proof}

\noindent
Il résulte par exemple de ce lemme que,
si $k'/k$ est une extension galoisienne finie, alors
$\Spec(k')/\text{Gal}(k'/k) \cong \Spec(k)$.
Que se passe-t-il si l'extension est infinie ? Voici un exemple.

\begin{example}
\label{spaces-example-Qbar}
Soit $S = \Spec(\mathbf{Q})$.
Soit $U = \Spec(\overline{\mathbf{Q}})$.
Soit $G = \text{Gal}(\overline{\mathbf{Q}}/\mathbf{Q})$, muni de son action
évidente sur $U$. Alors, par construction, la propriété $(*)$ du
Lemme \ref{spaces-lemma-quotient} est vérifiée, et l'on obtient un espace algébrique
$$
X = \Spec(\overline{\mathbf{Q}})/G
\longrightarrow
S = \Spec(\mathbf{Q}).
$$
Bien entendu, c'est une approximation tout à fait ridicule de $S$ !
En effet, d'après le théorème d'Artin-Schreier,
voir \cite[Théorème 17, page 316]{JacobsonIII},
les seuls sous-groupes finis de $\text{Gal}(\overline{\mathbf{Q}}/\mathbf{Q})$
sont $\{1\}$ et les conjugués du groupe d'ordre deux
$\text{Gal}(\overline{\mathbf{Q}}/\overline{\mathbf{Q}} \cap \mathbf{R})$.
Par conséquent, si
$\Spec(k) \to X$ est un morphisme avec $k$ algébrique sur $\mathbf{Q}$,
il résulte du Lemme \ref{spaces-lemma-quotient-field-map} et du théorème
que nous venons de mentionner que $k$ est soit $\overline{\mathbf{Q}}$, soit isomorphe à
$\overline{\mathbf{Q}} \cap \mathbf{R}$.
\end{example}

\noindent
Ce qui ne va pas dans l'exemple ci-dessus, c'est que
le groupe de Galois est muni d'une topologie,
et que celle-ci devrait intervenir d'une manière ou d'une autre dans toute construction
d'un quotient de $\Spec(\overline{\mathbf{Q}})$.
L'exemple suivant est, à mon avis, beaucoup plus raisonnable
et peut effectivement apparaître dans « la nature ».

\begin{example}
\label{spaces-example-affine-line-translation}
Soit $k$ un corps de caractéristique zéro.
Soit $U = \mathbf{A}^1_k$ et soit $G = \mathbf{Z}$.
Nous prenons pour action $n(x) = x + n$, c'est-à-dire l'action de
$\mathbf{Z}$ sur la droite affine par translation.
Le seul point fixe est le point générique, et il
est clair que $\mathbf{Z}$ s'injecte dans
le groupe des automorphismes du corps $k(x)$. (C'est
ici que nous utilisons l'hypothèse de caractéristique zéro.)
Considérons le morphisme
$$
\gamma : \Spec(k(x)) \longrightarrow X = \mathbf{A}^1_k/\mathbf{Z}
$$
du point générique de la droite affine vers le quotient.
Nous affirmons que ce morphisme ne se factorise par aucun
monomorphisme $\Spec(L) \to X$ du spectre d'un
corps vers $X$. (À la différence de ce qui se passe pour les schémas, voir
Schémas, section \ref{schemes-section-points}.) En effet, comme
$\mathbf{Z}$ ne possède aucun sous-groupe fini non trivial, il résulte du
Lemme \ref{spaces-lemma-quotient-field-map} que, pour toute telle
factorisation, $k(x) = L$. Enfin, $\gamma$ n'est pas un monomorphisme,
puisque
$$
\Spec(k(x)) \times_{\gamma, X, \gamma} \Spec(k(x))
\cong
\Spec(k(x)) \times \mathbf{Z}.
$$
\end{example}

\noindent
Cet exemple suggère que, pour définir les points d'un espace algébrique
$X$, il faut considérer les classes d'équivalence de morphismes de spectres
de corps vers $X$, et non l'ensemble des monomorphismes de spectres de corps.

\medskip\noindent
Nous terminons par un exemple véritablement affreux.

\begin{example}
\label{spaces-example-infinite-product}
Soit $k$ un corps.
Soit $A = \prod_{n \in \mathbf{N}} k$ le produit infini.
Posons $U = \Spec(A)$, considéré comme un schéma sur $S = \Spec(k)$.
Notons que les projections $\text{pr}_n : A \to k$ définissent des immersions
ouvertes et fermées $f_n : S \to U$. Posons
$$
R = U \amalg \coprod\nolimits_{(n, m) \in \mathbf{N}^2, \ n \not = m} S
$$
où le morphisme $j$ est égal à $\Delta_{U/S}$ sur la composante $U$
et à $j = (f_n, f_m)$ sur la composante $S$ correspondant à $(n, m)$.
La remarque ci-dessus montre clairement que $s, t$ sont \'etales.
Il est également clair que $j$ est une relation d'équivalence. Nous
obtenons donc un espace algébrique
$$
X = U/R.
$$
Pour comprendre ce que cela signifie, plaçons-nous dans le cas où
le corps $k$ est fini et possède $q$ éléments. Commençons par
examiner quelque peu l'espace topologique $|U|$ associé au schéma $U$.
Tout élément de $A$ vérifie $x^q = x$.
Par conséquent, tout corps résiduel de $A$ est isomorphe à $k$, et
tous les points de $U$ sont fermés. Mais la topologie de $U$ n'est pas
la topologie discrète. Soit $u_n \in |U|$ le point correspondant
à $f_n$. Comme indiqué ci-dessus, les points $u_n$ sont
les points ouverts (et sont donc isolés). Il doit donc exister
d'autres points, puisque nous savons que $U$ est quasi-compact, voir
Algèbre commutative, Lemme \ref{algebra-lemma-quasi-compact}
(et n'est donc pas égal à un ensemble discret infini).
Une autre façon de le voir consiste à remarquer que l'idéal (propre)
$$
I =
\{x = (x_n) \in A \mid \text{tous les }x_n\text{ sauf un nombre fini d'entre eux sont nuls}\}
$$
est contenu dans un idéal maximal. Notons aussi que tout élément
$x$ de $A$ s'écrit $x = ue$, où $u$ est une unité et $e$ un
idempotent. Une base de la topologie de $A$ est donc formée de parties ouvertes et
fermées (voir Algèbre commutative, Lemme \ref{algebra-lemma-idempotent-spec}.)
Ainsi, $|U|$ est totalement discontinu, mais non trivial.
Enfin, notons que $\{u_n\}$ est dense dans $|U|$.

\medskip\noindent
Nous définirons plus tard un espace topologique $|X|$ associé à $X$, voir
Propriétés des espaces algébriques, section \ref{spaces-properties-section-points}.
Que peut-on dire de $|X|$ ?
Il s'avère que l'application $|U| \to |X|$ est surjective et continue.
Tous les points $u_n$ ont pour image le même point $x_0$ de $|X|$, et aucun
des autres points n'est identifié à un autre. Comme $\{u_n\}$ est dense dans $|U|$, on
conclut que l'adhérence de $x_0$ dans $|X|$ est $|X|$. Autrement dit,
$|X|$ est irréductible et $x_0$ est un point générique de $|X|$. Cela paraît
bizarre, puisque $x_0$ est aussi l'image d'une section
$S \to X$ du morphisme structural $X \to S$ (et, dans le cas des
schémas, cela impliquerait qu'il s'agit d'un point fermé, voir
Morphismes de schémas, Lemme
\ref{morphisms-lemma-algebraic-residue-field-extension-closed-point-fibre}).
\end{example}

\noindent
Quoi que l'on pense de ce qui se passe réellement dans cet exemple, il montre assurément
qu'il faut faire preuve d'une certaine prudence lorsque l'on définit les composantes
irréductibles, la connexité, etc., des espaces algébriques.






\section{Changement de grands sites}
\label{spaces-section-change-big-site}

\noindent
Dans cette section, nous examinons brièvement ce qui se passe lorsque l'on
change de grand site. En substance, on peut toujours agrandir à volonté le grand
site; on peut donc supposer que tout ensemble de schémas que l'on souhaite
considérer est contenu dans le grand site fppf sur lequel on considère notre
espace algébrique. Voici un énoncé précis du résultat.

\begin{lemma}
\label{spaces-lemma-change-big-site}
Supposons donnés de grands sites $\Sch_{fppf}$ et $\Sch'_{fppf}$.
Supposons que $\Sch_{fppf}$ soit contenu dans $\Sch'_{fppf}$,
voir Topologies sur les schémas, section \ref{topologies-section-change-alpha}.
Soit $S$ un objet de $\Sch_{fppf}$. Soient
\begin{align*}
g : \Sh((\Sch/S)_{fppf})
\longrightarrow
\Sh((\Sch'/S)_{fppf}), \\
f : \Sh((\Sch'/S)_{fppf})
\longrightarrow
\Sh((\Sch/S)_{fppf})
\end{align*}
les morphismes de topos de
Topologies sur les schémas, Lemme \ref{topologies-lemma-change-alpha}.
Soit $F$ un faisceau d'ensembles sur $(\Sch/S)_{fppf}$. Alors
\begin{enumerate}
\item si $F$ est représentable par un schéma
$X \in \Ob((\Sch/S)_{fppf})$ sur $S$,
alors $f^{-1}F$ est lui aussi représentable; en fait, il est représentable par le
même schéma $X$, considéré maintenant comme objet de $(\Sch'/S)_{fppf}$, et
\item si $F$ est un espace algébrique sur $S$, alors $f^{-1}F$ est lui aussi un
espace algébrique sur $S$.
\end{enumerate}
\end{lemma}

\begin{proof}
Soit $X \in \Ob((\Sch/S)_{fppf})$. Notons $h_X$ le
faisceau représentable sur $(\Sch/S)_{fppf}$ associé à $X$, et
$h'_X$ le faisceau représentable sur $(\Sch'/S)_{fppf}$ associé à
$X$. D'après la description de $f^{-1}$ dans
Topologies sur les schémas, section \ref{topologies-section-change-alpha},
on a $f^{-1}h_X = h'_X$. Cela prouve (1).

\medskip\noindent
Supposons ensuite que $F$
soit un espace algébrique sur $S$. D'après le Lemme \ref{spaces-lemma-space-presentation},
cela signifie que $F = h_U/h_R$ pour une relation d'équivalence \'etale
$R \to U \times_S U$ dans $(\Sch/S)_{fppf}$. Comme $f^{-1}$ est un
foncteur exact, on en conclut que $f^{-1}F = h'_U/h'_R$. Par conséquent,
$f^{-1}F$ est un espace algébrique sur $S$ d'après le Théorème \ref{spaces-theorem-presentation}.
\end{proof}

\noindent
Notons que ce lemme est purement ensembliste et n'a pratiquement aucun contenu.
De plus, il n'est pas vrai en général que la restriction d'un espace
algébrique sur le site plus grand soit un espace algébrique sur le site plus petit
(pour de simples raisons de cardinalité). On ne peut donc utiliser un lemme
de ce genre que pour agrandir la catégorie de base, et jamais pour la rétrécir.

\begin{lemma}
\label{spaces-lemma-fully-faithful}
Supposons que $\Sch_{fppf}$ soit contenu dans $\Sch'_{fppf}$.
Soit $S$ un objet de $\Sch_{fppf}$. Notons
$\textit{Espaces}/S$ la catégorie des espaces algébriques sur $S$
définie à l'aide de $\Sch_{fppf}$. De même, notons
$\textit{Espaces}'/S$ la catégorie des espaces algébriques sur $S$
définie à l'aide de $\Sch'_{fppf}$. La construction du
Lemme \ref{spaces-lemma-change-big-site}
définit un foncteur pleinement fidèle
$$
\textit{Espaces}/S \longrightarrow \textit{Espaces}'/S
$$
dont l'image essentielle est formée des $X' \in \Ob(\textit{Espaces}'/S)$
pour lesquels il existe $U, R \in \Ob((\Sch/S)_{fppf})$\footnote{Exiger
l'existence de $R$ est nécessaire à cause de notre choix de la fonction $Bound$ dans
Théorie des ensembles, Équation (\ref{sets-equation-bound}). La taille du produit fibré
$U \times_{X'} U$ peut croître plus vite que $Bound$ en fonction de la taille de $U$. On
peut l'illustrer en posant $S = \Spec(A)$, $U = \Spec(A[x_i, i \in I])$ et
$R = \coprod_{(\lambda_i) \in A^I} \Spec(A[x_i, y_i]/(x_i - \lambda_i y_i))$.
Dans ce cas, la taille de $R$ croît comme $\kappa^\kappa$, où $\kappa$ est la
taille de $U$.} et des morphismes
$$
U \longrightarrow X'
\quad\text{et}\quad
R \longrightarrow U \times_{X'} U
$$
dans $\Sh((\Sch'/S)_{fppf})$ qui sont surjectifs comme morphismes de faisceaux
(par exemple, si les morphismes affichés sont surjectifs et \'etales).
\end{lemma}

\begin{proof}
Dans Sites et faisceaux, Lemme \ref{sites-lemma-bigger-site}, nous avons vu que le foncteur
$f^{-1} : \Sh((\Sch/S)_{fppf}) \to \Sh((\Sch'/S)_{fppf})$
est pleinement fidèle (voir la discussion dans
Topologies sur les schémas, section \ref{topologies-section-change-alpha}).
Le foncteur affiché dans le lemme est donc pleinement fidèle.

\medskip\noindent
Supposons que $X' \in \Ob(\textit{Espaces}'/S)$ soit tel qu'il existe
$U \in \Ob((\Sch/S)_{fppf})$ et un morphisme $U \to X'$ dans
$\Sh((\Sch'/S)_{fppf})$ qui soit surjectif comme morphisme de faisceaux.
Soit $U' \to X'$ un morphisme \'etale surjectif avec
$U' \in \Ob((\Sch'/S)_{fppf})$. Posons $\kappa = \text{size}(U)$, voir
Théorie des ensembles, section \ref{sets-section-categories-schemes}.
Alors $U$ possède un recouvrement ouvert affine $U = \bigcup_{i \in I} U_i$
avec $|I| \leq \kappa$. Remarquons que $U' \times_{X'} U \to U$ est \'etale
et surjectif. Pour chaque $i$, on peut choisir un ouvert quasi-compact
$U'_i \subset U'$ tel que $U'_i \times_{X'} U_i \to U_i$
soit surjectif (car le schéma $U' \times_{X'} U_i$ est la
réunion des ouverts de Zariski $W \times_{X'} U_i$ pour $W \subset U'$
affine et parce que $U' \times_{X'} U_i \to U_i$ est \'etale, donc ouvert).
Alors $\coprod_{i \in I} U'_i \to X'$ est \'etale et surjectif
en vertu de notre hypothèse selon laquelle $U \to X'$ et donc $\coprod U_i \to X'$
est un morphisme surjectif de faisceaux (détails omis).
Comme $U'_i \times_{X'} U \to U'_i$ est un morphisme surjectif de faisceaux
et que $U'_i$ est quasi-compact,
on peut trouver un ouvert quasi-compact $W_i \subset U'_i \times_{X'} U$
tel que $W_i \to U'_i$ soit surjectif comme morphisme de faisceaux
(détails omis). Alors $W_i \to U$ est \'etale et l'on en conclut
que $\text{size}(W_i) \leq \text{size}(U)$, voir
Théorie des ensembles, Lemme \ref{sets-lemma-bound-finite-type}. D'après
Théorie des ensembles, Lemme \ref{sets-lemma-bound-by-covering}, on en conclut
que $\text{size}(U'_i) \leq \text{size}(U)$.
Ainsi, $\coprod_{i \in I} U'_i$ est isomorphe à un objet de
$(\Sch/S)_{fppf}$ d'après Théorie des ensembles, Lemme \ref{sets-lemma-bound-size}.

\medskip\noindent
Soient maintenant $X'$, $U \to X'$ et $R \to U \times_{X'} U$ comme dans
l'énoncé du lemme. Au paragraphe précédent, nous avons vu que
l'on peut trouver $U' \in \Ob((\Sch/S)_{fppf})$ et un morphisme \'etale surjectif
$U' \to X'$ dans $\Sh((\Sch'/S)_{fppf})$. Alors
$U' \times_{X'} U \to U'$ est un morphisme surjectif de faisceaux, c'est-à-dire
qu'il existe un recouvrement fppf $\{U'_i \to U'\}$ tel que $U'_i \to U'$ se factorise par
$U' \times_{X'} U \to U'$.
D'après Théorie des ensembles, Lemme \ref{sets-lemma-bound-fppf-covering},
on peut trouver
$\tilde U \to U'$
qui soit surjectif, plat et localement de présentation finie,
avec $\text{size}(\tilde U) \leq \text{size}(U')$,
tel que $\tilde U \to U'$ se factorise par
$U' \times_{X'} U \to U'$. Considérons alors
$$
\xymatrix{
U' \times_{X'} U' \ar[d] &
\tilde U \times_{X'} \tilde U  \ar[l] \ar[d] \ar[r] &
U \times_{X'} U \ar[d] \\
U' \times_S U' & \tilde U \times_S \tilde U \ar[l] \ar[r] & U \times_S U
}
$$
Les carrés sont cartésiens. On sait que les objets de la ligne inférieure
sont représentés par des objets de $(\Sch/S)_{fppf}$. D'après le résultat de
l'argument du paragraphe précédent, il en va de même pour $U \times_{X'} U$
(puisque, par hypothèse, on dispose du morphisme surjectif de faisceaux
$R \to U \times_{X'} U$). Comme $(\Sch/S)_{fppf}$ est stable par produits
fibrés (par construction), on voit que $\tilde U \times_{X'} \tilde U$
est représenté par un objet de $(\Sch/S)_{fppf}$. Enfin, le
morphisme $\tilde U \times_{X'} \tilde U \to U' \times_{X'} U'$ est
un morphisme surjectif de faisceaux fppf, puisque $\tilde U \to U'$ l'est.
On peut donc appliquer une fois encore le résultat du paragraphe précédent
pour conclure que $R' = U' \times_{X'} U'$ est représenté par un objet
de $(\Sch/S)_{fppf}$. À ce stade, le
Lemme \ref{spaces-lemma-space-presentation} et le
Théorème \ref{spaces-theorem-presentation} impliquent que $X = h_{U'}/h_{R'}$
est un objet de $\textit{Espaces}/S$ tel que $f^{-1}X \cong X'$,
comme souhaité.
\end{proof}



\section{Changement de schéma de base}
\label{spaces-section-change-base-scheme}

\noindent
Dans cette section, nous examinons brièvement ce qui se passe lorsque l'on
change de schéma de base. Le point essentiel est que, étant donné un morphisme
$S \to S'$ de schémas de base, tout espace algébrique au-dessus de $S$ peut être
considéré comme un espace algébrique au-dessus de $S'$. De plus, étant donné un
espace algébrique $F'$ au-dessus de $S'$, il existe un changement de base $F'_S$, qui est
un espace algébrique au-dessus de $S$. Nous expliquons seulement le cas où $S \to S'$
est un morphisme du grand site fppf considéré. Si seul l'un des schémas $S$ et $S'$
appartient au grand site, on commence par agrandir celui-ci comme dans la
section \ref{spaces-section-change-big-site}.

\begin{lemma}
\label{spaces-lemma-change-base-scheme}
Soit $\Sch_{fppf}$ un grand site.
Soit $g : S \to S'$ un morphisme de $\Sch_{fppf}$.
Soit $j : (\Sch/S)_{fppf} \to (\Sch/S')_{fppf}$ le
foncteur de localisation correspondant.
Soit $F$ un faisceau d'ensembles sur $(\Sch/S)_{fppf}$.
Alors
\begin{enumerate}
\item pour un schéma $T'$ au-dessus de $S'$, on a
$j_!F(T'/S') =
\coprod\nolimits_{\varphi : T' \to S} F(T' \xrightarrow{\varphi} S),$
\item si $F$ est représentable par un schéma
$X \in \Ob((\Sch/S)_{fppf})$,
alors $j_!F$ est représentable par $j(X)$, c'est-à-dire par
$X$ considéré comme un schéma au-dessus de $S'$, et
\item si $F$ est un espace algébrique au-dessus de $S$, alors $j_!F$ est un
espace algébrique au-dessus de $S'$ et, si $F = U/R$ est une présentation,
alors $j_!F = j(U)/j(R)$ est une présentation.
\end{enumerate}
Soit $F'$ un faisceau d'ensembles sur $(\Sch/S')_{fppf}$. Alors
\begin{enumerate}
\item[(4)] pour un schéma $T$ au-dessus de $S$, on a $j^{-1}F'(T/S) = F'(T/S')$,
\item[(5)] si $F'$ est représentable par un schéma
$X' \in \Ob((\Sch/S')_{fppf})$, alors
$j^{-1}F'$ est représentable, à savoir par $X'_S = S \times_{S'} X'$, et
\item[(6)] si $F'$ est un espace algébrique, alors
$j^{-1}F'$ est un espace algébrique et, si $F' = U'/R'$ est une présentation,
alors $j^{-1}F' = U'_S/R'_S$ est une présentation.
\end{enumerate}
\end{lemma}

\begin{proof}
Les foncteurs $j_!$, $j_*$ et $j^{-1}$ sont définis dans
Sites et faisceaux, Lemme \ref{sites-lemma-relocalize},
où l'on montre également que $j = j_{S/S'}$ est la localisation
de $(\Sch/S')_{fppf}$ en l'objet $S/S'$. Tout ce qui concerne
les foncteurs de localisation s'applique donc à $j$.
La formule de (1) est celle de
Sites et faisceaux, Lemme \ref{sites-lemma-describe-j-shriek-good-site}.
Par définition, $j_!$ est adjoint à gauche du foncteur de restriction $j^{-1}$ ;
$j_!$ est donc exact à droite. D'après
Sites et faisceaux, Lemme \ref{sites-lemma-j-shriek-commutes-equalizers-fibre-products},
il commute en outre aux produits fibrés et aux égalisateurs.
D'après
Sites et faisceaux, Lemme \ref{sites-lemma-describe-j-shriek-representable},
on a $j_!h_X = h_{j(X)}$, d'où (2).
Si $F$ est un espace algébrique au-dessus de $S$, écrivons $F = U/R$
(Lemme \ref{spaces-lemma-space-presentation}) ;
on obtient
$$
j_!F = j(U)/j(R)
$$
puisque $j_!$, étant exact à droite, commute aux conoyaux ; de plus,
$j(R) = j(U) \times_{j_!F} j(U)$ puisque $j_!$ commute aux produits fibrés.
Comme les morphismes $j(s), j(t) : j(R) \to j(U)$ ne sont autres que les morphismes
$s, t : R \to U$ (considérés comme des morphismes de schémas au-dessus de $S'$), ils
sont encore \'etales. Ainsi, $(j(U), j(R), s, t)$ est une relation d'équivalence \'etale.
Le
Théorème \ref{spaces-theorem-presentation}
permet donc de conclure que $j_!F$ est un espace algébrique.

\medskip\noindent
Démonstration de (4), (5) et (6). La description de $j^{-1}$ se trouve dans
Sites et faisceaux, section \ref{sites-section-localize}.
La restriction du faisceau représentable associé à $X'/S'$
est le faisceau représentable associé à
$X'_S = S \times_{S'} X'$ d'après
Sites et faisceaux, Lemme \ref{sites-lemma-localize-given-products}.
Le foncteur de restriction $j^{-1}$ est exact, donc $j^{-1}F' = U'_S/R'_S$.
Par exactitude encore, le faisceau $R'_S$ demeure une relation d'équivalence sur
$U'_S$. Enfin, les deux morphismes $R'_S \to U'_S$ sont \'etales, car ils se
déduisent par changement de base des morphismes \'etales $R' \to U'$. Ainsi,
$j^{-1}F' = U'_S/R'_S$ est un espace algébrique d'après le
Théorème \ref{spaces-theorem-presentation},
ce qui achève la démonstration.
\end{proof}

\noindent
Notons que la présentation $j_!F = j(U)/j(R)$ n'est autre que la présentation
de $F$, considérée comme une présentation par des schémas au-dessus de $S'$.
La définition suivante a donc un sens.

\begin{definition}
\label{spaces-definition-base-change}
Soit $\Sch_{fppf}$ un grand site fppf.
Soit $S \to S'$ un morphisme de ce site.
\begin{enumerate}
\item Si $F'$ est un espace algébrique au-dessus de $S'$, alors le
{\it changement de base de $F'$ à $S$} est
l'espace algébrique $j^{-1}F'$ décrit dans le
Lemme \ref{spaces-lemma-change-base-scheme}. Nous le notons $F'_S$.
\item Si $F$ est un espace algébrique au-dessus de $S$, alors $F$,
{\it considéré comme un espace algébrique au-dessus de $S'$},
est l'espace algébrique $j_!F$ au-dessus de $S'$ décrit dans le
Lemme \ref{spaces-lemma-change-base-scheme}. Nous le notons souvent simplement
$F$ ; sinon, nous écrivons $j_!F$.
\end{enumerate}
\end{definition}

\noindent
L'espace algébrique $j_!F$ est muni d'un morphisme canonique
$j_!F \to S$ d'espaces algébriques au-dessus de $S'$. Cela vient simplement
de ce que le faisceau $j_!F$ admet un morphisme vers $h_S$ (voir, par exemple, la
description explicite du
Lemme \ref{spaces-lemma-change-base-scheme}).
En effet, dans
Sites et faisceaux, Lemme \ref{sites-lemma-essential-image-j-shriek},
nous avons vu que la catégorie des faisceaux sur $(\Sch/S)_{fppf}$
est équivalente à la catégorie des couples $(\mathcal{F}', \mathcal{F}' \to h_S)$
formés d'un faisceau sur $(\Sch/S')_{fppf}$ et d'un morphisme de faisceaux
$\mathcal{F}' \to h_S$. L'équivalence associe au faisceau $\mathcal{F}$
le couple $(j_!\mathcal{F}, j_!\mathcal{F} \to h_S)$.
Ceci, joint à ce qui précède, conduit au résultat suivant
pour les catégories d'espaces algébriques.

\begin{lemma}
\label{spaces-lemma-category-of-spaces-over-smaller-base-scheme}
Soit $\Sch_{fppf}$ un grand site fppf.
Soit $S \to S'$ un morphisme de ce site.
La construction ci-dessus donne une équivalence de
catégories
$$
\begin{gathered}
\left\{
\begin{matrix}
\text{catégorie des espaces}\\
\text{algébriques sur }S
\end{matrix}
\right\}\\[4pt]
\leftrightarrow\\[4pt]
\left\{
\begin{matrix}
\text{catégorie des couples }(F', F' \to S)\text{ formés}\\
\text{d'un espace algébrique }F'\text{ sur }S'\\
\text{et d'un morphisme }F' \to S\text{ d'espaces algébriques}\\
\text{sur }S'
\end{matrix}
\right\}
\end{gathered}
$$
\end{lemma}

\begin{proof}
Soit $F$ un espace algébrique au-dessus de $S$. Le foncteur de gauche à droite
associe le couple $(j_!F, j_!F \to S)$ à $F$,
et ce couple appartient au membre de droite d'après le
Lemme \ref{spaces-lemma-change-base-scheme}.
Puisque ceci définit une équivalence de catégories de faisceaux d'après
Sites et faisceaux, Lemme \ref{sites-lemma-essential-image-j-shriek},
il suffit, pour achever la démonstration, de montrer ceci :
si $F$ est un faisceau et $j_!F$ un espace algébrique, alors $F$
est un espace algébrique. Pour cela, écrivons
$j_!F = U'/R'$ comme dans le
Lemme \ref{spaces-lemma-space-presentation},
avec $U', R' \in \Ob((\Sch/S')_{fppf})$.
Les composés $U' \to j_!F \to S$ et $R' \to j_!F \to S$
sont alors des morphismes de schémas au-dessus de $S'$. Notons $U, R$ les
objets correspondants de $(\Sch/S)_{fppf}$. Les deux morphismes
$R' \to U'$ sont des morphismes au-dessus de $S$ et correspondent donc à des
morphismes $R \to U$. Comme ce sont les mêmes
morphismes (considérés au-dessus de $S$), nous obtenons une relation
d'équivalence \'etale au-dessus de $S$. Puisque $j_!$ définit une équivalence de
catégories de faisceaux (voir la référence ci-dessus), on a
$F = U/R$ ; le
Théorème \ref{spaces-theorem-presentation}
montre alors que $F$ est un espace algébrique.
\end{proof}

\noindent
Le lemme suivant est une légère reformulation de ce qui précède.

\begin{lemma}
\label{spaces-lemma-rephrase}
Soit $\Sch_{fppf}$ un grand site fppf.
Soit $S \to S'$ un morphisme de ce site.
Soit $F'$ un faisceau sur $(\Sch/S')_{fppf}$.
Les assertions suivantes sont équivalentes :
\begin{enumerate}
\item la restriction $F'|_{(\Sch/S)_{fppf}}$
est un espace algébrique au-dessus de $S$, et
\item le faisceau $h_S \times F'$ est un espace algébrique au-dessus de $S'$.
\end{enumerate}
\end{lemma}

\begin{proof}
La restriction et le produit se correspondent par
l'équivalence de catégories de
Sites et faisceaux, Lemme \ref{sites-lemma-essential-image-j-shriek},
de sorte que le
Lemme \ref{spaces-lemma-category-of-spaces-over-smaller-base-scheme}
ci-dessus donne le résultat.
\end{proof}

\noindent
Nous terminons cette section par un lemme de compatibilité.

\begin{lemma}
\label{spaces-lemma-viewed-as-properties}
Soit $\Sch_{fppf}$ un grand site fppf.
Soit $S \to S'$ un morphisme de ce site.
Soit $F$ un espace algébrique au-dessus de $S$.
Soit $T$ un schéma au-dessus de $S$ et soit $f : T \to F$
un morphisme au-dessus de $S$.
Soit $f' : T' \to F'$ le morphisme au-dessus de $S'$ obtenu à partir de
$f$ en appliquant l'équivalence de catégories décrite dans le
Lemme \ref{spaces-lemma-category-of-spaces-over-smaller-base-scheme}.
Pour toute propriété $\mathcal{P}$ comme dans la
Définition \ref{spaces-definition-relative-representable-property},
on a $\mathcal{P}(f') \Leftrightarrow \mathcal{P}(f)$.
\end{lemma}

\begin{proof}
Supposons que $U$ soit un schéma au-dessus de $S$ et que $U \to F$ soit un
morphisme \'etale surjectif. Notons $U'$ le schéma $U$ considéré comme un schéma au-dessus de $S'$. Dans le
Lemme \ref{spaces-lemma-change-base-scheme},
nous avons vu que $U' \to F'$ est \'etale surjectif. Puisque
$$
j(T \times_{f, F} U) = T' \times_{f', F'} U'
$$
le morphisme de schémas $T \times_{f, F} U \to U$
s'identifie au morphisme de schémas
$T' \times_{f', F'} U' \to U'$.
Il s'agit du même morphisme, considéré au-dessus de schémas de base différents.
Le lemme résulte donc du
Lemme \ref{spaces-lemma-representable-morphisms-spaces-property}.
\end{proof}














% OK, start here.
%

\chapter{Propriétés des espaces algébriques}



\phantomsection
\label{spaces-properties-section-phantom}


\section{Introduction}
\label{spaces-properties-section-introduction}

\noindent
Veuillez consulter
Espaces, section \ref{spaces-section-introduction},
pour une brève introduction aux espaces algébriques, et lire aussi
une partie de ce chapitre pour nos définitions et conventions de base
concernant les espaces algébriques. Dans le présent chapitre, nous introduisons
quelques notions et propriétés élémentaires des espaces algébriques. Une référence
fondamentale pour le cas des espaces algébriques quasi-séparés est
\cite{Kn}.

\medskip\noindent
L'exposé est parfois quelque peu malaisé, car nous avons choisi
de traiter d'abord des propriétés des espaces algébriques
en eux-mêmes, et seulement plus tard de celles des morphismes d'espaces
algébriques. Nous dérogeons à cette règle pour les
{\it morphismes étales} d'espaces algébriques, que nous introduisons à la
section \ref{spaces-properties-section-etale-morphisms}.
Mais, jusqu'à cette section, lorsque nous disons qu'un morphisme possède une propriété donnée,
cela signifie automatiquement que sa source est un schéma (ou peut-être que
le morphisme est représentable).

\medskip\noindent
Certains résultats de ce chapitre (notamment ceux qui concernent les points)
seront améliorés dans le chapitre consacré aux espaces algébriques décents.


\section{Conventions}
\label{spaces-properties-section-conventions}

\noindent
On suppose une fois pour toutes que tous les schémas appartiennent à
un grand site fppf $\Sch_{fppf}$. De plus, tous les anneaux $A$ considérés
sont tels que $\Spec(A)$ soit (isomorphe à) un
objet de ce grand site.

\medskip\noindent
Soient $S$ un schéma et $X$ un espace algébrique sur $S$.
Dans ce chapitre et le suivant, nous noterons $X \times_S X$
le produit de $X$ par lui-même (dans la catégorie des espaces
algébriques sur $S$), au lieu de $X \times X$. Nous voulons ainsi
éviter toute confusion lorsque nous changeons de schéma de base, comme dans
Espaces, section \ref{spaces-section-change-base-scheme}.


\section{Axiomes de séparation}
\label{spaces-properties-section-separation}

\noindent
Dans cette section, nous réunissons toutes les conditions de séparation « absolues »
des espaces algébriques. Puisque, selon nos conventions, tout espace algébrique est un
espace algébrique sur un schéma de base bien déterminé, toute propriété absolue
de $X$ sur $S$ correspond à une condition imposée à $X$ considéré
comme espace algébrique sur $\Spec(\mathbf{Z})$. En voici la
formulation précise.

\begin{definition}
\label{spaces-properties-definition-separated}
(Comparer avec Espaces, définition \ref{spaces-definition-separated}.)
Considérons un grand site fppf
$\Sch_{fppf} = (\Sch/\Spec(\mathbf{Z}))_{fppf}$.
Soit $X$ un espace algébrique sur
$\Spec(\mathbf{Z})$. Soit $\Delta : X \to X \times X$
le morphisme diagonal.
\begin{enumerate}
\item On dit que $X$ est {\it séparé} si $\Delta$ est une immersion fermée.
\item On dit que $X$ est {\it localement séparé}\footnote{Dans la
littérature, cela désigne souvent des espaces algébriques quasi-séparés et
localement séparés.} si $\Delta$ est une
immersion.
\item On dit que $X$ est {\it quasi-séparé} si $\Delta$ est quasi-compact.
\item On dit que $X$ est {\it localement quasi-séparé pour la topologie de Zariski}\footnote{
Cette notion a été proposée par B.\ Conrad.} s'il
existe un recouvrement de Zariski $X = \bigcup_{i \in I} X_i$ (voir Espaces,
définition \ref{spaces-definition-Zariski-open-covering}) tel que
chaque $X_i$ soit quasi-séparé.
\end{enumerate}
Soit $S$ un schéma appartenant à $\Sch_{fppf}$, et soit
$X$ un espace algébrique sur $S$. On dit alors que $X$ est {\it séparé},
{\it localement séparé}, {\it quasi-séparé} ou
{\it localement quasi-séparé pour la topologie de Zariski}
si $X$, considéré comme espace algébrique sur $\Spec(\mathbf{Z})$ (voir
Espaces, définition \ref{spaces-definition-base-change}),
possède la propriété correspondante.
\end{definition}

\noindent
Un espace algébrique $X$ sur $S$ qui est séparé
(au sens absolu ci-dessus) est bien séparé sur $S$ (et il en va de même
pour les autres propriétés absolues de séparation ci-dessus). Ce point sera étudié
en détail dans
Morphismes d'espaces, section \ref{spaces-morphisms-section-separation-axioms}.
Nous verrons au
lemme \ref{spaces-properties-lemma-quasi-separated-quasi-compact-pieces}
que la propriété d'être localement quasi-séparé pour la topologie de Zariski est indépendante du schéma de base
(et donc équivalente à la notion absolue).

\begin{lemma}
\label{spaces-properties-lemma-trivial-implications}
Soit $S$ un schéma.
Soit $X$ un espace algébrique sur $S$.
On a les implications suivantes entre les axiomes de séparation
de la définition \ref{spaces-properties-definition-separated} :
\begin{enumerate}
\item la propriété d'être séparé entraîne toutes les autres,
\item la propriété d'être quasi-séparé entraîne celle d'être localement quasi-séparé pour la topologie de Zariski.
\end{enumerate}
\end{lemma}

\begin{proof}
La démonstration est omise.
\end{proof}

\begin{lemma}
\label{spaces-properties-lemma-characterize-quasi-separated}
Soit $S$ un schéma. Soit $X$ un espace algébrique sur $S$.
Les conditions suivantes sont équivalentes :
\begin{enumerate}
\item $X$ est un espace algébrique quasi-séparé,
\item pour $U \to X$, $V \to X$, où $U$ et $V$ sont des schémas quasi-compacts,
le produit fibré $U \times_X V$ est quasi-compact,
\item pour $U \to X$, $V \to X$, où $U$ et $V$ sont affines,
le produit fibré $U \times_X V$ est quasi-compact.
\end{enumerate}
\end{lemma}

\begin{proof}
Le lemme
\ref{spaces-lemma-category-of-spaces-over-smaller-base-scheme}
d'Espaces permet de supposer que $S = \Spec(\mathbf{Z})$.
Puisque $U \times_X V = X \times_{X \times X} (U \times V)$
et que $U \times V$ est quasi-compact lorsque $U$ et $V$ le sont,
on voit que (1) entraîne (2). Il est clair que (2) entraîne (3).
Supposons (3). Choisissons un schéma $W$ et un morphisme étale surjectif
$W \to X$. Alors $W \times W \to X \times X$ est étale surjectif.
Il suffit donc de montrer que
$$
j : W \times_X W = X \times_{(X \times X)} (W \times W) \to W \times W
$$
est quasi-compact, voir Espaces, lemme
\ref{spaces-lemma-descent-representable-transformations-property}.
Si $U \subset W$ et $V \subset W$ sont des ouverts affines, alors
$j^{-1}(U \times V) = U \times_X V$ est quasi-compact par hypothèse.
Puisque les ouverts affines $U \times V$ forment un recouvrement ouvert affine de
$W \times W$ (Schémas, lemme \ref{schemes-lemma-affine-covering-fibre-product}),
on conclut par
Schémas, lemme \ref{schemes-lemma-quasi-compact-affine}.
\end{proof}

\begin{lemma}
\label{spaces-properties-lemma-characterize-separated}
Soit $S$ un schéma. Soit $X$ un espace algébrique sur $S$.
Les conditions suivantes sont équivalentes :
\begin{enumerate}
\item $X$ est un espace algébrique séparé,
\item pour $U \to X$, $V \to X$, où $U$ et $V$ sont affines,
le produit fibré $U \times_X V$ est affine et
$$
\mathcal{O}(U) \otimes_\mathbf{Z} \mathcal{O}(V)
\longrightarrow
\mathcal{O}(U \times_X V)
$$
est surjectif.
\end{enumerate}
\end{lemma}

\begin{proof}
Le lemme
\ref{spaces-lemma-category-of-spaces-over-smaller-base-scheme}
d'Espaces permet de supposer que $S = \Spec(\mathbf{Z})$.
Puisque $U \times_X V = X \times_{X \times X} (U \times V)$
et que $U \times V$ est affine lorsque $U$ et $V$ le sont,
on voit que (1) entraîne (2).
Supposons (2). Choisissons un schéma $W$ et un morphisme étale surjectif
$W \to X$. Alors $W \times W \to X \times X$ est étale surjectif.
Il suffit donc de montrer que
$$
j : W \times_X W = X \times_{(X \times X)} (W \times W) \to W \times W
$$
est une immersion fermée, voir Espaces, lemme
\ref{spaces-lemma-descent-representable-transformations-property}.
Si $U \subset W$ et $V \subset W$ sont des ouverts affines, alors
$j^{-1}(U \times V) = U \times_X V$ est affine par hypothèse
et le morphisme $U \times_X V \to U \times V$ est une immersion fermée,
car l'homomorphisme d'anneaux correspondant est surjectif.
Puisque les ouverts affines $U \times V$ forment un recouvrement ouvert affine
de $W \times W$
(Schémas, lemme \ref{schemes-lemma-affine-covering-fibre-product}),
on conclut par
Morphismes, lemme \ref{morphisms-lemma-closed-immersion}.
\end{proof}







\section{Points des espaces algébriques}
\label{spaces-properties-section-points}

\noindent
Comme il ressort clairement de
Espaces, exemple \ref{spaces-example-affine-line-translation},
il ne convient pas de définir un point d'un espace algébrique comme un monomorphisme
ayant pour source le spectre d'un corps.
Nous définissons plutôt les points comme des classes d'équivalence de morphismes de spectres
de corps, exactement comme expliqué dans
Schémas, section \ref{schemes-section-points}.

\medskip\noindent
Soit $S$ un schéma.
Soit $F$ un préfaisceau sur $(\Sch/S)_{fppf}$.
Soit $K$ un corps. Considérons un morphisme
$$
\Spec(K) \longrightarrow F.
$$
D'après le lemme de Yoneda, celui-ci est donné par un
élément $p \in F(\Spec(K))$. On dit que deux tels
couples $(\Spec(K), p)$ et $(\Spec(L), q)$
sont {\it équivalents} s'il existe
un troisième corps $\Omega$ et un diagramme commutatif
$$
\xymatrix{
\Spec(\Omega) \ar[r] \ar[d] &
\Spec(L) \ar[d]^q \\
\Spec(K) \ar[r]^p &
F.
}
$$
Autrement dit, il existe des extensions de corps
$K \to \Omega$ et $L \to \Omega$ telles que
les images de $p$ et de $q$ coïncident dans
$F(\Spec(\Omega))$. Nous omettons la vérification du fait que ceci
définit une relation d'équivalence.

\begin{definition}
\label{spaces-properties-definition-points}
Soit $S$ un schéma. Soit $X$ un espace algébrique sur $S$.
Un {\it point} de $X$ est une classe d'équivalence de morphismes
de spectres de corps vers $X$.
L'ensemble des points de $X$ est noté $|X|$.
\end{definition}

\noindent
Notons que, si $f : X \to Y$ est un morphisme d'espaces algébriques
sur $S$, alors on dispose d'une application induite $|f| : |X| \to |Y|$ qui
envoie un représentant $x : \Spec(K) \to X$ sur le représentant
$f \circ x : \Spec(K) \to Y$.

\begin{lemma}
\label{spaces-properties-lemma-scheme-points}
Soit $S$ un schéma. Soit $X$ un schéma sur $S$.
Les points de $X$ en tant que schéma sont canoniquement en bijection
avec les points de $X$ en tant qu'espace algébrique.
\end{lemma}

\begin{proof}
C'est Schémas, lemme \ref{schemes-lemma-characterize-points}.
\end{proof}

\begin{lemma}
\label{spaces-properties-lemma-points-cartesian}
Soit $S$ un schéma. Soit
$$
\xymatrix{
Z \times_Y X \ar[r] \ar[d] & X \ar[d] \\
Z \ar[r] & Y
}
$$
un diagramme cartésien d'espaces algébriques sur $S$. Alors l'application entre les ensembles
de points
$$
|Z \times_Y X|
\longrightarrow
|Z| \times_{|Y|} |X|
$$
est surjective.
\end{lemma}

\begin{proof}
En effet, soient des corps $K$, $L$ et des morphismes
$\Spec(K) \to X$, $\Spec(L) \to Z$. L'hypothèse
que leurs images dans $|Y|$ coïncident signifie qu'il
existe une extension commune $M/K$ et $M/L$
telle que les morphismes
$\Spec(M) \to \Spec(K) \to X \to Y$ et
$\Spec(M) \to \Spec(L) \to Z \to Y$ coïncident.
C'est précisément la condition qui assure l'existence d'un
morphisme $\Spec(M) \to Z \times_Y X$.
\end{proof}

\begin{lemma}
\label{spaces-properties-lemma-characterize-surjective}
Soit $S$ un schéma.
Soit $X$ un espace algébrique sur $S$.
Soit $f : T \to X$ un morphisme d'un schéma vers $X$.
Les assertions suivantes sont équivalentes :
\begin{enumerate}
\item $f : T \to X$ est surjectif (au sens de
Espaces, définition \ref{spaces-definition-relative-representable-property}),
et
\item $|f| : |T| \to |X|$ est surjectif.
\end{enumerate}
\end{lemma}

\begin{proof}
Supposons (1). Soit $x : \Spec(K) \to X$ un morphisme
du spectre d'un corps vers $X$. Par hypothèse, le morphisme de
schémas $\Spec(K) \times_X T \to \Spec(K)$ est surjectif.
Il existe donc une extension de corps $K'/K$ et un morphisme
$\Spec(K') \to \Spec(K) \times_X T$ tels que le carré
de gauche du diagramme
$$
\xymatrix{
\Spec(K') \ar[r] \ar[d] &
\Spec(K) \times_X T \ar[d] \ar[r] &
T \ar[d] \\
\Spec(K) \ar@{=}[r] &
\Spec(K) \ar[r]^-x & X
}
$$
soit commutatif. Cela montre que $|f| : |T| \to |X|$ est surjectif.

\medskip\noindent
Supposons (2). Soit $Z \to X$ un morphisme, où $Z$ est
un schéma. Il faut montrer que le morphisme de schémas $Z \times_X T \to Z$
est surjectif, c'est-à-dire que $|Z \times_X T| \to |Z|$ est surjectif.
Cela résulte de (2) et du
lemme \ref{spaces-properties-lemma-points-cartesian}.
\end{proof}

\begin{lemma}
\label{spaces-properties-lemma-points-presentation}
Soit $S$ un schéma.
Soit $X$ un espace algébrique sur $S$.
Soit $X = U/R$ une présentation de $X$, voir
Espaces, définition \ref{spaces-definition-presentation}.
Alors l'image de $|R| \to |U| \times |U|$ est une relation d'équivalence
et $|X|$ est le quotient de $|U|$ par cette relation d'équivalence.
\end{lemma}

\begin{proof}
L'hypothèse signifie que $U$ est un schéma, que $p : U \to X$ est un morphisme
étale surjectif, que $R = U \times_X U$ est un schéma et définit une relation
d'équivalence étale sur $U$ telle que $X = U/R$ comme faisceaux. D'après le
lemme \ref{spaces-properties-lemma-characterize-surjective},
on voit que $|U| \to |X|$ est surjective. D'après le
lemme \ref{spaces-properties-lemma-points-cartesian},
l'application
$$
|R| \longrightarrow |U| \times_{|X|} |U|
$$
est surjective. L'image de $|R| \to |U| \times |U|$ est donc
exactement l'ensemble des couples $(u_1, u_2) \in |U| \times |U|$
tels que $u_1$ et $u_2$ aient la même image dans $|X|$.
En réunissant ces deux assertions, on obtient le résultat du lemme.
\end{proof}

\begin{lemma}
\label{spaces-properties-lemma-topology-points}
Soit $S$ un schéma. Il existe une unique topologie sur les ensembles de points
des espaces algébriques sur $S$ qui possède les propriétés suivantes :
\begin{enumerate}
\item si $X$ est un schéma sur $S$, la topologie sur $|X|$ est la topologie usuelle
(via l'identification du lemme \ref{spaces-properties-lemma-scheme-points}),
\item pour tout morphisme d'espaces algébriques $X \to Y$ sur $S$,
l'application $|X| \to |Y|$ est continue, et
\item pour tout morphisme étale $U \to X$, où $U$ est un schéma,
l'application d'espaces topologiques $|U| \to |X|$ est continue et ouverte.
\end{enumerate}
\end{lemma}

\begin{proof}
Soit $X$ un espace algébrique sur $S$. Soit $p : U \to X$ un
morphisme étale surjectif, où $U$ est un schéma sur $S$.
On déclare que $W \subset |X|$ est ouvert si et seulement si $|p|^{-1}(W)$
est une partie ouverte de $|U|$. Cela définit une topologie sur $|X|$
(c'est la topologie quotient sur $|X|$, voir
Topologie, lemme \ref{topology-lemma-quotient}).

\medskip\noindent
Montrons que la topologie est indépendante du choix de
la présentation. Pour cela, il suffit de montrer que, si $U'$ est un schéma
et $U' \to X$ un morphisme étale, alors l'application $|U'| \to |X|$
(la topologie sur $|X|$ étant définie à l'aide de $U \to X$ comme ci-dessus)
est ouverte et continue, ce qui démontrera en outre (3).
Posons $U'' = U \times_X U'$ ; on a alors le diagramme commutatif
$$
\xymatrix{
U'' \ar[r] \ar[d] & U' \ar[d] \\
U \ar[r] & X
}
$$
Puisque $U \to X$ et $U' \to X$ sont étales, on voit que
$U'' \to U$ et $U'' \to U'$ sont tous deux des morphismes étales de schémas.
De plus, $U'' \to U'$ est surjectif. On obtient donc
un diagramme commutatif d'applications d'ensembles
$$
\xymatrix{
|U''| \ar[r] \ar[d] & |U'| \ar[d] \\
|U| \ar[r] & |X|
}
$$
La flèche horizontale inférieure est surjective (voir le
lemme \ref{spaces-properties-lemma-characterize-surjective}
ou le
lemme \ref{spaces-properties-lemma-points-presentation})
et continue par définition de la topologie sur $|X|$.
La flèche horizontale supérieure est surjective, continue et ouverte d'après
Morphismes, lemme \ref{morphisms-lemma-etale-open}.
La flèche verticale gauche est également continue et ouverte d'après
Morphismes, lemme \ref{morphisms-lemma-etale-open}.
Il en résulte formellement que la flèche verticale droite
est continue et ouverte.

\medskip\noindent
Pour achever la démonstration, prouvons (2).
Soit $a : X \to Y$ un morphisme d'espaces algébriques. D'après
Espaces, lemme \ref{spaces-lemma-lift-morphism-presentations},
on peut trouver un diagramme
$$
\xymatrix{
U \ar[d]_p \ar[r]_\alpha & V \ar[d]^q \\
X \ar[r]^a & Y
}
$$
où $U$ et $V$ sont des schémas, et où $p$ et $q$ sont étales et surjectifs.
On en déduit le diagramme
$$
\xymatrix{
|U| \ar[d]_p \ar[r]_\alpha & |V| \ar[d]^q \\
|X| \ar[r]^a & |Y|
}
$$
où toutes les flèches sauf l'horizontale inférieure sont continues et
où les deux flèches verticales sont surjectives et ouvertes. Il s'ensuit que la
flèche horizontale inférieure est continue, comme voulu.
\end{proof}

\begin{definition}
\label{spaces-properties-definition-topological-space}
Soit $S$ un schéma. Soit $X$ un espace algébrique sur $S$.
L'{\it espace topologique sous-jacent} à $X$ est l'ensemble des points
$|X|$ muni de la topologie construite au
lemme \ref{spaces-properties-lemma-topology-points}.
\end{definition}

\noindent
Il se trouve que cet espace topologique porte la même information
que le petit site de Zariski $X_{Zar}$ de
Espaces, définition \ref{spaces-definition-small-Zariski-site}.

\begin{lemma}
\label{spaces-properties-lemma-open-subspaces}
Soit $S$ un schéma.
Soit $X$ un espace algébrique sur $S$.
\begin{enumerate}
\item La règle $X' \mapsto |X'|$ définit une bijection respectant les inclusions
entre les sous-espaces ouverts $X'$ (voir
Espaces, définition \ref{spaces-definition-immersion})
de $X$ et les ouverts de l'espace topologique $|X|$.
\item Une famille $\{X_i \subset X\}_{i \in I}$ de sous-espaces ouverts de $X$
est un recouvrement de Zariski (voir
Espaces, définition \ref{spaces-definition-Zariski-open-covering})
si et seulement si $|X| = \bigcup |X_i|$.
\end{enumerate}
Autrement dit, le petit site de Zariski $X_{Zar}$ de $X$ s'identifie canoniquement
à un site associé à l'espace topologique $|X|$ (voir
Sites, exemple \ref{sites-example-site-topological}).
\end{lemma}

\begin{proof}
Pour démontrer (1), construisons l'inverse de cette règle.
Supposons donc que $W \subset |X|$ soit ouvert. Choisissons une présentation
$X = U/R$ correspondant au morphisme étale surjectif
$p : U \to X$ et aux morphismes étales $s, t : R \to U$.
Par construction, on voit que $|p|^{-1}(W)$ est un
ouvert de $U$. Notons $W' \subset U$ le sous-schéma ouvert correspondant.
Il est clair que $R' = s^{-1}(W') = t^{-1}(W')$ est un ouvert de Zariski
de $R$ qui définit une relation d'équivalence étale sur $W'$.
D'après Espaces, lemme \ref{spaces-lemma-finding-opens}, le morphisme
$X' = W'/R' \to X$ est une immersion ouverte. Ainsi, $X'$ est un espace algébrique
d'après Espaces, lemme \ref{spaces-lemma-representable-over-space}.
Par construction, $|X'| = W$, c'est-à-dire que $X'$ est un sous-espace de $X$
correspondant à $W$. Cela démontre (1).

\medskip\noindent
Pour démontrer (2), notons que, si $\{X_i \subset X\}_{i \in I}$ est une famille
de sous-espaces ouverts, elle est un recouvrement de Zariski si et seulement si
$U = \bigcup U \times_X X_i$ est un recouvrement ouvert. Cela résulte de
la définition d'un recouvrement de Zariski et du fait que le morphisme
$U \to X$ est surjectif comme morphisme de préfaisceaux sur $(\Sch/S)_{fppf}$.
D'autre part, on voit que $|X| = \bigcup |X_i|$ si et seulement si
$U = \bigcup U \times_X X_i$ d'après le lemme \ref{spaces-properties-lemma-points-presentation}
(et le fait que les projections $U \times_X X_i \to X_i$ sont surjectives
et étales). L'équivalence de (2) en résulte.
\end{proof}

\begin{lemma}
\label{spaces-properties-lemma-factor-through-open-subspace}
Soit $S$ un schéma.
Soient $X$, $Y$ des espaces algébriques sur $S$.
Soit $X' \subset X$ un sous-espace ouvert.
Soit $f : Y \to X$ un morphisme d'espaces algébriques sur $S$.
Alors $f$ se factorise par $X'$ si et seulement si $|f| : |Y| \to |X|$
se factorise par $|X'| \subset |X|$.
\end{lemma}

\begin{proof}
D'après Espaces, lemme \ref{spaces-lemma-base-change-immersions},
on voit que $Y' = Y \times_X X' \to Y$ est une immersion ouverte.
Si $|f|(|Y|) \subset |X'|$, alors $|Y'| = |Y|$ manifestement. Ainsi, $Y' = Y$ d'après le
lemme \ref{spaces-properties-lemma-open-subspaces}.
\end{proof}

\begin{lemma}
\label{spaces-properties-lemma-etale-image-open}
Soit $S$ un schéma. Soit $X$ un espace algébrique sur $S$.
Soit $U$ un schéma et soit $f : U \to X$ un morphisme étale.
Soit $X' \subset X$ le sous-espace ouvert correspondant à
l'ouvert $|f|(|U|) \subset |X|$ par le
lemme \ref{spaces-properties-lemma-open-subspaces}.
Alors $f$ se factorise par un morphisme étale surjectif $f' : U \to X'$.
En outre, si $R = U \times_X U$, alors $R = U \times_{X'} U$ et $X'$ admet
la présentation $X' = U/R$.
\end{lemma}

\begin{proof}
L'existence de la factorisation résulte du
lemme \ref{spaces-properties-lemma-factor-through-open-subspace}.
Le morphisme $f'$ est surjectif d'après le
lemme \ref{spaces-properties-lemma-characterize-surjective}.
Pour voir que $f'$ est étale, soit $T \to X'$ un morphisme,
où $T$ est un schéma. Alors $T \times_X U = T \times_{X'} U$,
car $X' \to X$ est un monomorphisme de faisceaux. Ainsi, la projection
$T \times_{X'} U \to T$ est étale puisque $f$ est supposé étale.
On a $U \times_X U = U \times_{X'} U$, car $X' \to X$ est un monomorphisme.
Alors $X' = U/R$ résulte de
Espaces, lemme \ref{spaces-lemma-space-presentation}.
\end{proof}

\begin{lemma}
\label{spaces-properties-lemma-equivalence-class-point-monomorphism}
Soit $S$ un schéma. Soit $X$ un espace algébrique sur $S$.
Soient $p : \Spec(K) \to X$ et $q : \Spec(L) \to X$
des morphismes, où $K$ et $L$ sont des corps. Supposons que $p$ et $q$
déterminent le même point de $|X|$ et que $p$ soit un monomorphisme.
Alors $q$ se factorise de manière unique par $p$.
\end{lemma}

\begin{proof}
Puisque $p$ et $q$ définissent le même point de $|X|$, on voit que le schéma
$$
Y = \Spec(K) \times_{p, X, q} \Spec(L)
$$
est non vide. Comme tout changement de base d'un monomorphisme est un monomorphisme,
cela signifie que le morphisme de projection $Y \to \Spec(L)$
est un monomorphisme. Ainsi, $Y = \Spec(L)$, voir
Schémas, lemme \ref{schemes-lemma-mono-towards-spec-field}.
On en conclut que $q$ se factorise par $p$. L'unicité
résulte du fait que $p$ est un monomorphisme.
\end{proof}

\begin{lemma}
\label{spaces-properties-lemma-points-monomorphism}
Soit $S$ un schéma. Soit $X$ un espace algébrique sur $S$.
Considérons l'application
$$
\{\Spec(k) \to X \text{ monomorphisme où }k\text{ est un corps}\}
\longrightarrow
|X|
$$
Cette application est injective.
\end{lemma}

\begin{proof}
Cela résulte du lemme \ref{spaces-properties-lemma-equivalence-class-point-monomorphism}.
\end{proof}

\noindent
Nous verrons dans Espaces algébriques décents,
lemme \ref{decent-spaces-lemma-decent-points-monomorphism},
que l'application du lemme \ref{spaces-properties-lemma-points-monomorphism}
est bijective lorsque $X$ est décent.

















\section{Espaces quasi-compacts}
\label{spaces-properties-section-quasi-compact}

\begin{definition}
\label{spaces-properties-definition-quasi-compact}
Soit $S$ un schéma.
Soit $X$ un espace algébrique sur $S$.
On dit que $X$ est {\it quasi-compact} s'il existe un morphisme
étale surjectif $U \to X$, où $U$ est un schéma quasi-compact.
\end{definition}

\begin{lemma}
\label{spaces-properties-lemma-quasi-compact-space}
Soit $S$ un schéma.
Soit $X$ un espace algébrique sur $S$.
Alors $X$ est quasi-compact si et seulement si $|X|$ est quasi-compact.
\end{lemma}

\begin{proof}
Choisissons un schéma $U$ et un morphisme étale surjectif $U \to X$.
Nous utiliserons le lemme \ref{spaces-properties-lemma-characterize-surjective}.
Si $U$ est quasi-compact, alors, puisque $|U| \to |X|$ est surjective,
on en conclut que $|X|$ est quasi-compact.
Si $|X|$ est quasi-compact, alors, puisque $|U| \to |X|$ est ouverte,
on voit qu'il existe un ouvert quasi-compact $U' \subset U$
tel que $|U'| \to |X|$ soit surjective (et encore étale).
D'où le résultat.
\end{proof}

\begin{lemma}
\label{spaces-properties-lemma-finite-disjoint-quasi-compact}
Une somme disjointe finie d'espaces algébriques quasi-compacts est
un espace algébrique quasi-compact.
\end{lemma}

\begin{proof}
Cela résulte immédiatement du
lemme \ref{spaces-properties-lemma-quasi-compact-space}
et du fait topologique correspondant.
\end{proof}

\begin{example}
\label{spaces-properties-example-quasi-compact-not-very-reasonable}
L'espace $\mathbf{A}^1_{\mathbf{Q}}/\mathbf{Z}$ est un espace algébrique
quasi-compact.
\end{example}

\begin{lemma}
\label{spaces-properties-lemma-space-locally-quasi-compact}
Soit $S$ un schéma.
Soit $X$ un espace algébrique sur $S$.
Tout point de $|X|$ admet un système fondamental de voisinages
ouverts quasi-compacts.
En particulier, $|X|$ est localement quasi-compact au sens de
Topologie, définition \ref{topology-definition-locally-quasi-compact}.
\end{lemma}

\begin{proof}
Cela découle formellement du fait qu'il existe un schéma
$U$ et une application surjective, ouverte et continue
$U \to |X|$ entre espaces topologiques. Plus précisément, si
$u \in U$ a pour image $x \in |X|$, alors les images des voisinages
affines de $u$ forment un système fondamental de voisinages ouverts
quasi-compacts de $x$.
\end{proof}







\section{Recouvrements spéciaux}
\label{spaces-properties-section-special-coverings}

\noindent
Dans cette section, nous réunissons quelques lemmes immédiats sur l'existence
de recouvrements étales surjectifs d'espaces algébriques.

\begin{lemma}
\label{spaces-properties-lemma-cover-by-union-affines}
Soit $S$ un schéma.
Soit $X$ un espace algébrique sur $S$.
Il existe un morphisme étale surjectif $U \to X$, où
$U$ est une somme disjointe de schémas affines.
On peut en outre supposer que chacun de ces schémas affines
est envoyé dans un ouvert affine de $S$.
\end{lemma}

\begin{proof}
Soit $V \to X$ un morphisme étale surjectif.
Soit $V = \bigcup_{i \in I} V_i$ un recouvrement ouvert de Zariski
tel que chaque $V_i$ soit envoyé dans un ouvert affine de $S$.
Posons alors $U = \coprod_{i \in I} V_i$, muni du morphisme induit
$U \to V \to X$. Celui-ci est étale et surjectif, car il est composé
de transformations représentables de foncteurs qui sont étales et
surjectives (d'après le principe général de
Espaces, lemme
\ref{spaces-lemma-composition-representable-transformations-property}
et de
Morphismes, lemmes \ref{morphisms-lemma-composition-surjective} et
\ref{morphisms-lemma-composition-etale}).
\end{proof}

\begin{lemma}
\label{spaces-properties-lemma-union-of-quasi-compact}
Soit $S$ un schéma.
Soit $X$ un espace algébrique sur $S$.
Il existe un recouvrement de Zariski $X = \bigcup X_i$
tel que chaque espace algébrique $X_i$ admette un recouvrement
étale surjectif par un schéma affine. On peut en outre supposer
que chaque $X_i$ soit envoyé dans un ouvert affine de $S$.
\end{lemma}

\begin{proof}
D'après le lemme \ref{spaces-properties-lemma-cover-by-union-affines}, on peut trouver un morphisme
étale surjectif $U = \coprod U_i \to X$, où chaque $U_i$ est affine et est envoyé
dans un ouvert affine de $S$. Soit $X_i \subset X$ le sous-espace ouvert
de $X$ tel que $U_i \to X$ se factorise par un morphisme étale surjectif
$U_i \to X_i$, voir le
lemme \ref{spaces-properties-lemma-etale-image-open}.
Puisque $U = \bigcup U_i$, on voit que $X = \bigcup X_i$.
Comme $U_i \to X_i$ est surjectif, il s'ensuit que $X_i \to S$ se factorise par
un ouvert affine de $S$.
\end{proof}

\begin{lemma}
\label{spaces-properties-lemma-quasi-compact-affine-cover}
Soit $S$ un schéma.
Soit $X$ un espace algébrique sur $S$.
Alors $X$ est quasi-compact si et seulement s'il
existe un morphisme étale surjectif $U \to X$
où $U$ est un schéma affine.
\end{lemma}

\begin{proof}
S'il existe un morphisme étale surjectif $U \to X$ où $U$
est affine, alors $X$ est quasi-compact par la définition \ref{spaces-properties-definition-quasi-compact}.
Réciproquement, si $X$ est quasi-compact, alors $|X|$ est quasi-compact.
Soit $U = \coprod_{i \in I} U_i$ une somme disjointe de schémas affines
munie d'un morphisme étale surjectif $\varphi : U \to X$
(lemme \ref{spaces-properties-lemma-cover-by-union-affines}).
Alors $|X| = \bigcup \varphi(|U_i|)$ et,
par quasi-compacité, il existe des indices $i_1, \ldots, i_n$
tels que $|X| = \bigcup \varphi(|U_{i_j}|)$. Par conséquent,
$U_{i_1} \cup \ldots \cup U_{i_n}$ est un schéma affine muni d'un
morphisme étale surjectif vers $X$.
\end{proof}

\noindent
Le lemme suivant deviendra caduc lorsque nous traiterons des
morphismes séparés dans le chapitre consacré aux morphismes d'espaces algébriques.

\begin{lemma}
\label{spaces-properties-lemma-separated-cover}
Soit $S$ un schéma.
Soit $X$ un espace algébrique sur $S$.
Soit $U$ un schéma séparé et soit $U \to X$ étale.
Alors $U \to X$ est séparé, et $R = U \times_X U$ est un schéma séparé.
\end{lemma}

\begin{proof}
Soit $X' \subset X$ le sous-espace ouvert tel que $U \to X$ se factorise
par un morphisme étale surjectif $U \to X'$, voir le
lemme \ref{spaces-properties-lemma-etale-image-open}.
Si $U \to X'$ est séparé, alors $U \to X$ l'est aussi, voir
Espaces, lemme
\ref{spaces-lemma-composition-representable-transformations-property}
(car l'immersion ouverte $X' \to X$ est séparée d'après
Espaces, lemme
\ref{spaces-lemma-representable-transformations-property-implication}
et
Schémas, lemme \ref{schemes-lemma-immersions-monomorphisms}).
De plus, puisque $U \times_{X'} U = U \times_X U$, il suffit
de démontrer le résultat après avoir remplacé $X$ par $X'$, c'est-à-dire qu'on peut
supposer $U \to X$ surjectif.
Considérons le diagramme commutatif
$$
\xymatrix{
R = U \times_X U \ar[r] \ar[d] & U \ar[d] \\
U \ar[r] & X
}
$$
Dans la démonstration de
Espaces, lemme \ref{spaces-lemma-properties-diagonal},
nous avons vu que $j : R \to U \times_S U$ est séparé.
Le morphisme de schémas $U \to S$ est séparé, puisque $U$ est un schéma
séparé, voir
Schémas, lemme \ref{schemes-lemma-compose-after-separated}.
Par conséquent, $U \times_S U \to U$ est séparé en tant que changement de base, voir
Schémas, lemme \ref{schemes-lemma-separated-permanence}.
Le schéma $U \times_S U$ est donc séparé (d'après le même lemme).
Puisque $j$ est séparé, on voit de même que $R$ est séparé.
Ainsi, $R \to U$ est un morphisme séparé (d'après
Schémas, lemme \ref{schemes-lemma-compose-after-separated},
là encore). Par conséquent, d'après
Espaces, lemme \ref{spaces-lemma-representable-morphisms-spaces-property}
et le diagramme ci-dessus, on conclut que $U \to X$ est séparé.
\end{proof}

\begin{lemma}
\label{spaces-properties-lemma-quasi-separated}
Soit $S$ un schéma. Soit $X$ un espace algébrique sur $S$.
S'il existe un schéma quasi-séparé $U$ et un morphisme
étale surjectif $U \to X$ tel que l'une ou l'autre des projections
$U \times_X U \to U$ soit quasi-compacte, alors $X$ est quasi-séparé.
\end{lemma}

\begin{proof}
On peut considérer $X$ comme un espace algébrique sur $\mathbf{Z}$.
Considérons le diagramme cartésien
$$
\xymatrix{
U \times_X U \ar[r] \ar[d]_j & X \ar[d]^\Delta \\
U \times U \ar[r] & X \times X
}
$$
Puisque $U$ est quasi-séparé, la projection $U \times U \to U$ est
quasi-séparée (elle s'obtient par changement de base à partir d'un morphisme
quasi-séparé de schémas, voir Schémas, lemme \ref{schemes-lemma-separated-permanence}).
L'hypothèse du lemme implique donc que $j$ est quasi-compact d'après
Schémas, lemme \ref{schemes-lemma-quasi-compact-permanence}.
D'après Espaces, lemme \ref{spaces-lemma-representable-morphisms-spaces-property},
on voit que le morphisme $\Delta$ est quasi-compact, comme voulu.
\end{proof}

\begin{lemma}
\label{spaces-properties-lemma-quasi-separated-quasi-compact-pieces}
Soit $S$ un schéma.
Soit $X$ un espace algébrique sur $S$.
Les assertions suivantes sont équivalentes :
\begin{enumerate}
\item $X$ est localement quasi-séparé pour la topologie de Zariski sur $S$,
\item $X$ est localement quasi-séparé pour la topologie de Zariski,
\item il existe un recouvrement ouvert de Zariski $X = \bigcup X_i$
tel que, pour tout $i$, il existe un schéma affine
$U_i$ et un morphisme étale, surjectif et quasi-compact
$U_i \to X_i$, et
\item il existe un recouvrement ouvert de Zariski $X = \bigcup X_i$
tel que, pour tout $i$, il existe un schéma affine
$U_i$ qui est envoyé dans un ouvert affine de $S$ et un morphisme
étale, surjectif et quasi-compact $U_i \to X_i$.
\end{enumerate}
\end{lemma}

\begin{proof}
Supposons que les morphismes $U_i \to X_i \subset X$ soient comme en (3). Pour démontrer (4),
choisissons, pour tout $i$, un recouvrement ouvert affine fini $U_i =
U_{i1} \cup \ldots \cup U_{in_i}$ tel que chaque $U_{ij}$ soit envoyé
dans un ouvert affine de $S$. Les composés
$U_{ij} \to U_i \to X_i$ sont étales et quasi-compacts (voir
Espaces, lemme
\ref{spaces-lemma-composition-representable-transformations-property}).
Soit $X_{ij} \subset X_i$ le sous-espace ouvert correspondant à
l'image de $|U_{ij}| \to |X_i|$, voir le
lemme \ref{spaces-properties-lemma-etale-image-open}.
Notons que $U_{ij} \to X_{ij}$ est quasi-compact, car $X_{ij} \subset X_i$
est un monomorphisme et que $U_{ij} \to X$ est quasi-compact.
Alors $X = \bigcup X_{ij}$ est un recouvrement comme en (4).
L'implication (4) $\Rightarrow$ (3) est immédiate.

\medskip\noindent
Supposons (4). Pour montrer que $X$ est localement quasi-séparé pour la topologie de Zariski sur $S$,
il suffit de montrer que $X_i$ est quasi-séparé sur $S$.
On peut donc supposer qu'il existe un schéma affine $U$ envoyé dans
un ouvert affine de $S$ et un morphisme étale, surjectif et quasi-compact
$U \to X$. Considérons le carré cartésien
$$
\xymatrix{
U \times_X U \ar[r] \ar[d] & U \times_S U \ar[d] \\
X \ar[r]^-{\Delta_{X/S}} & X \times_S X
}
$$
La flèche verticale droite est étale et surjective (voir
Espaces, lemme
\ref{spaces-lemma-product-representable-transformations-property})
et $U \times_S U$ est affine (car $U$ est envoyé dans un ouvert affine de $S$, voir
Schémas, section \ref{schemes-section-fibre-products}),
et $U \times_X U$ est quasi-compact,
car la projection $U \times_X U \to U$ est quasi-compacte en tant que
changement de base de $U \to X$. Il résulte de
Espaces, lemme \ref{spaces-lemma-representable-morphisms-spaces-property}
que $\Delta_{X/S}$ est quasi-compact, comme voulu.

\medskip\noindent
Supposons (1). Pour démontrer (3), on se ramène immédiatement au cas
où $X$ est quasi-séparé sur $S$. D'après le
lemme \ref{spaces-properties-lemma-union-of-quasi-compact},
on peut trouver un recouvrement ouvert de Zariski $X = \bigcup X_i$ tel que chaque
$X_i$ soit envoyé dans un ouvert affine de $S$ et qu'il existe des schémas affines
$U_i$ et des morphismes étales surjectifs $U_i \to X_i$.
Puisque $U_i \to S$ se factorise par un ouvert affine de $S$, on voit que
$U_i \times_S U_i$ est affine, voir
Schémas, section \ref{schemes-section-fibre-products}.
Comme $X$ est quasi-séparé sur $S$, les morphismes
$$
R_i = U_i \times_{X_i} U_i = U_i \times_X U_i
\longrightarrow
U_i \times_S U_i
$$
sont quasi-compacts en tant que changements de base de $\Delta_{X/S}$. On en conclut
que $R_i$ est un schéma quasi-compact. Cela implique à son tour que chaque
projection $R_i \to U_i$ est quasi-compacte. Par conséquent, en appliquant
Espaces, lemme \ref{spaces-lemma-representable-morphisms-spaces-property}
au recouvrement $U_i \to X_i$ et au morphisme $U_i \to X_i$,
on conclut que les morphismes $U_i \to X_i$ sont quasi-compacts, comme voulu.

\medskip\noindent
On voit maintenant que (1), (3) et (4) sont équivalentes. Puisque (3) ne
fait pas intervenir le schéma de base, on en conclut qu'elles sont également équivalentes
à (2).
\end{proof}

\noindent
Le lemme suivant se révélera très utile.

\begin{lemma}
\label{spaces-properties-lemma-finite-fibres-presentation}
Soit $S$ un schéma.
Soit $X$ un espace algébrique sur $S$.
Soit $U$ un schéma. Soit $\varphi : U \to X$ un morphisme étale tel que
les projections $R = U \times_X U \to U$ soient quasi-compactes ; c'est par exemple le cas si
$\varphi$ est quasi-compact. Alors les fibres de
$$
|U| \to |X|
\quad\text{et}\quad
|R| \to |X|
$$
sont finies.
\end{lemma}

\begin{proof}
Notons $R = U \times_X U$, et $s, t : R \to U$ les projections.
Soit $u \in U$ un point, et soit $x \in |X|$ son image.
La fibre de $|U| \to |X|$ au-dessus de $x$ est égale à
$s(t^{-1}(\{u\}))$ d'après le
lemme \ref{spaces-properties-lemma-points-cartesian},
et la fibre de $|R| \to |X|$ au-dessus de $x$ est $t^{-1}(s(t^{-1}(\{u\})))$.
Puisque $t : R \to U$ est étale et quasi-compact, ses fibres sont finies
(ses fibres sont des sommes disjointes de spectres de corps d'après
Morphismes, lemme \ref{morphisms-lemma-etale-over-field},
et elles sont quasi-compactes). D'où le résultat.
\end{proof}








\section{Propriétés des espaces définies par des propriétés des schémas}
\sectionmark{Propriétés des espaces héritées des schémas}
\label{spaces-properties-section-types-properties}

\noindent
Toute propriété de schémas qui est locale pour la topologie \'etale donne lieu à une
propriété correspondante des espaces algébriques grâce au lemme suivant.

\begin{lemma}
\label{spaces-properties-lemma-type-property}
Soit $S$ un schéma.
Soit $X$ un espace algébrique sur $S$.
Soit $\mathcal{P}$ une propriété de schémas locale pour la topologie \'etale,
voir
Descente, définition \ref{descent-definition-property-local}.
Les assertions suivantes sont équivalentes :
\begin{enumerate}
\item il existe un schéma $U$ et un morphisme \'etale surjectif $U \to X$
tels que le schéma $U$ possède la propriété $\mathcal{P}$, et
\item pour tout schéma $U$ et tout morphisme \'etale $U \to X$,
le schéma $U$ possède la propriété $\mathcal{P}$.
\end{enumerate}
Si $X$ est représentable, ces conditions équivalent à $\mathcal{P}(X)$.
\end{lemma}

\begin{proof}
L'implication (2) $\Rightarrow$ (1) est immédiate.
Réciproquement, choisissons un morphisme \'etale surjectif $U \to X$
où $U$ est un schéma qui possède $\mathcal{P}$, et soit $V$ un schéma \'etale sur
$X$. Alors $U \times_X V \rightarrow V$ est un morphisme \'etale surjectif
de schémas ; ainsi, $V$ hérite de $\mathcal{P}$ de $U \times_X V$, lequel
hérite à son tour de $\mathcal{P}$ de $U$ (voir la discussion qui suit
Descente, définition \ref{descent-definition-property-local}).
La dernière assertion résulte clairement de (1) et de
Descente, définition \ref{descent-definition-property-local}.
\end{proof}

\begin{definition}
\label{spaces-properties-definition-type-property}
Soit $\mathcal{P}$ une propriété de schémas qui est
locale pour la topologie \'etale.
Soit $S$ un schéma.
Soit $X$ un espace algébrique sur $S$.
On dit que $X$ {\it possède la propriété $\mathcal{P}$}
si l'une des conditions équivalentes du
lemme \ref{spaces-properties-lemma-type-property}
est satisfaite.
\end{definition}

\begin{remark}
\label{spaces-properties-remark-list-properties-local-etale-topology}
Voici une liste de propriétés qui sont locales pour la topologie \'etale
(rappelons que les topologies fpqc, fppf, syntomique et lisse sont
plus fines que la topologie \'etale) :
\begin{enumerate}
\item localement noethérien, voir
Descente, lemme \ref{descent-lemma-Noetherian-local-fppf},
\item de Jacobson, voir
Descente, lemme \ref{descent-lemma-Jacobson-local-fppf},
\item localement noethérien et vérifiant $(S_k)$, voir
Descente, lemme \ref{descent-lemma-Sk-local-syntomic},
\item Cohen-Macaulay, voir
Descente, lemme \ref{descent-lemma-CM-local-syntomic},
\item de Gorenstein, voir
Dualité pour les schémas, lemme \ref{duality-lemma-gorenstein-local-syntomic},
\item réduit, voir
Descente, lemme \ref{descent-lemma-reduced-local-smooth},
\item normal, voir
Descente, lemme \ref{descent-lemma-normal-local-smooth},
\item localement noethérien et vérifiant $(R_k)$, voir
Descente, lemme \ref{descent-lemma-Rk-local-smooth},
\item régulier, voir
Descente, lemme \ref{descent-lemma-regular-local-smooth},
\item de Nagata, voir
Descente, lemme \ref{descent-lemma-Nagata-local-smooth}.
\end{enumerate}
\end{remark}

\noindent
Toute propriété de germes de schémas locale pour la topologie \'etale donne lieu à une
propriété correspondante des espaces algébriques. Voici le lemme de rigueur.

\begin{lemma}
\label{spaces-properties-lemma-local-source-target-at-point}
Soit $\mathcal{P}$ une propriété de germes de schémas locale pour la topologie \'etale, voir
Descente, définition \ref{descent-definition-local-at-point}.
Soit $S$ un schéma.
Soit $X$ un espace algébrique sur $S$.
Soit $x \in |X|$ un point de $X$.
Considérons les morphismes \'etales $a : U \to X$, où $U$ est un schéma.
Les assertions suivantes sont équivalentes :
\begin{enumerate}
\item pour tout $U \to X$ comme ci-dessus et tout $u \in U$ tel que $a(u) = x$, on a
$\mathcal{P}(U, u)$, et
\item il existe $U \to X$ comme ci-dessus et $u \in U$ tel que $a(u) = x$ et que l'on ait
$\mathcal{P}(U, u)$.
\end{enumerate}
Si $X$ est représentable, ces conditions équivalent à $\mathcal{P}(X, x)$.
\end{lemma}

\begin{proof}
Démonstration omise.
\end{proof}

\begin{definition}
\label{spaces-properties-definition-property-at-point}
Soit $S$ un schéma. Soit $X$ un espace algébrique sur $S$.
Soit $x \in |X|$. Soit $\mathcal{P}$ une propriété de germes de schémas qui est
locale pour la topologie \'etale.
On dit que $X$ {\it possède la propriété $\mathcal{P}$ en $x$} si l'une des
conditions équivalentes du
lemme \ref{spaces-properties-lemma-local-source-target-at-point}
est satisfaite.
\end{definition}

\begin{remark}
\label{spaces-properties-remark-list-properties-local-ring-local-etale-topology}
Soit $P$ une propriété d'anneaux locaux. Supposons que, pour tout
homomorphisme \'etale d'anneaux $A \to B$ et tout idéal premier $\mathfrak q$ de $B$ au-dessus de
l'idéal premier $\mathfrak p$ de $A$, on ait
$P(A_\mathfrak p) \Leftrightarrow P(B_\mathfrak q)$.
On obtient alors une propriété locale pour la topologie \'etale des germes $(U, u)$ de schémas
en posant $\mathcal{P}(U, u) = P(\mathcal{O}_{U, u})$.
Dans cette situation, nous emploierons l'expression
``l'anneau local de $X$ en $x$ possède $P$'' pour signifier que
$X$ possède la propriété $\mathcal{P}$ en $x$.
Voici une liste de telles propriétés $P$ :
\begin{enumerate}
\item noethérien, voir
Compléments d'algèbre, lemme \ref{more-algebra-lemma-Noetherian-etale-extension},
\item de dimension $d$, voir
Compléments d'algèbre, lemme \ref{more-algebra-lemma-dimension-etale-extension},
\item régulier, voir
Compléments d'algèbre, lemme \ref{more-algebra-lemma-regular-etale-extension},
\item anneau de valuation discrète, cela résulte de (2), (3) et du
lemme \ref{algebra-lemma-characterize-dvr} d'Algèbre,
\item réduit, voir
Compléments d'algèbre, lemme \ref{more-algebra-lemma-henselization-reduced},
\item normal, voir
Compléments d'algèbre, lemme \ref{more-algebra-lemma-henselization-normal},
\item noethérien et de profondeur $k$, voir
Compléments d'algèbre, lemme \ref{more-algebra-lemma-henselization-depth},
\item noethérien et Cohen-Macaulay, voir
Compléments d'algèbre, lemme \ref{more-algebra-lemma-henselization-CM},
\item noethérien et de Gorenstein, voir
Complexes dualisants, lemme \ref{dualizing-lemma-flat-under-gorenstein}.
\end{enumerate}
Il existe d'autres propriétés pour lesquelles ceci vaut, par exemple celles d'être un anneau G ou un
anneau de Nagata. Si nous en avons un jour besoin, nous les ajouterons ici,
ainsi que des références à des démonstrations détaillées des faits
d'algèbre correspondants.
\end{remark}







\section{Ensembles constructibles}
\label{spaces-properties-section-constructible}

\begin{lemma}
\label{spaces-properties-lemma-locally-constructible}
Soit $S$ un schéma. Soit $X$ un espace algébrique sur $S$.
Soit $E \subset |X|$ une partie. Les assertions suivantes sont équivalentes :
\begin{enumerate}
\item pour tout morphisme \'etale $U \to X$, où $U$ est un
schéma, l'image réciproque de $E$ dans $U$ est une partie localement constructible
de $U$,
\item pour tout morphisme \'etale $U \to X$, où $U$ est un
schéma affine, l'image réciproque de $E$ dans $U$ est une partie constructible
de $U$,
\item il existe un morphisme \'etale surjectif $U \to X$, où $U$ est un
schéma, tel que l'image réciproque de $E$ dans $U$ soit une partie localement constructible
de $U$.
\end{enumerate}
\end{lemma}

\begin{proof}
D'après Propriétés, lemme \ref{properties-lemma-locally-constructible},
on voit que (1) et (2) sont équivalentes. Il est immédiat que (1)
implique (3). Supposons donc donné un morphisme \'etale surjectif
$\varphi : U \to X$, où $U$ est un schéma tel que $\varphi^{-1}(E)$
soit localement constructible. Soit $\varphi' : U' \to X$ un autre
morphisme \'etale, où $U'$ est un schéma. On a alors
$$
E'' = \text{pr}_1^{-1}(\varphi^{-1}(E)) = \text{pr}_2^{-1}((\varphi')^{-1}(E))
$$
où $\text{pr}_1 : U \times_X U' \to U$ et
$\text{pr}_2 : U \times_X U' \to U'$ sont les projections.
D'après Morphismes, lemme \ref{morphisms-lemma-inverse-image-constructible},
on voit que $E''$ est localement constructible dans $U \times_X U'$.
Soit $W' \subset U'$ un ouvert affine. Puisque $\text{pr}_2$ est
\'etale et donc ouverte, on peut choisir un ouvert quasi-compact
$W'' \subset U \times_X U'$ tel que $\text{pr}_2(W'') = W'$.
Alors $\text{pr}_2|_{W''} : W'' \to W'$ est quasi-compact.
On a $W' \cap (\varphi')^{-1}(E) = \text{pr}_2(E'' \cap W'')$
puisque $\varphi$ est surjectif, voir le lemme \ref{spaces-properties-lemma-points-cartesian}.
Ainsi, $W' \cap (\varphi')^{-1}(E) = \text{pr}_2(E'' \cap W'')$
est localement constructible d'après
Morphismes, théorème \ref{morphisms-theorem-chevalley}, comme voulu.
\end{proof}

\begin{definition}
\label{spaces-properties-definition-locally-constructible}
Soit $S$ un schéma. Soit $X$ un espace algébrique sur $S$.
Soit $E \subset |X|$ une partie. On dit que $E$ est
{\it localement constructible pour la topologie \'etale} si les conditions
équivalentes du lemme \ref{spaces-properties-lemma-locally-constructible} sont satisfaites.
\end{definition}

\noindent
Bien entendu, si $X$ est représentable, c'est-à-dire si $X$ est un schéma,
cela signifie simplement que $E$ est une partie localement constructible
de l'espace topologique sous-jacent.







\section{Dimension en un point}
\label{spaces-properties-section-dimension}

\noindent
On peut utiliser
Descente, lemme \ref{descent-lemma-dimension-at-point-local},
pour définir la dimension d'un espace
algébrique $X$ en un point $x$. On obtient ainsi une notion différente de la
notion topologique (c'est-à-dire la dimension de $|X|$ en $x$).

\begin{definition}
\label{spaces-properties-definition-dimension-at-point}
Soit $S$ un schéma.
Soit $X$ un espace algébrique sur $S$.
Soit $x \in |X|$ un point de $X$.
On définit la {\it dimension de $X$ en $x$} comme
l'élément $\dim_x(X) \in \{0, 1, 2, \ldots, \infty\}$
tel que $\dim_x(X) = \dim_u(U)$ pour tout (ou, de manière équivalente, pour un)
couple $(a : U \to X, u)$ constitué d'un morphisme \'etale $a : U \to X$
d'un schéma vers $X$ et d'un point $u \in U$ tel que $a(u) = x$.
Voir la
définition \ref{spaces-properties-definition-property-at-point}, le
lemme \ref{spaces-properties-lemma-local-source-target-at-point} et
Descente, lemme \ref{descent-lemma-dimension-at-point-local}.
\end{definition}

\noindent
Attention : en général, on n'a {\bf pas} $\dim_x(X) = \dim_x(|X|)$.
Un contre-exemple est l'espace algébrique $X$ de
Espaces, exemple \ref{spaces-example-infinite-product}.
Plus précisément, soit $x \in |X|$ un point distinct du point générique $x_0$
de $|X|$. Alors $\dim_x(X) = 0$, mais $\dim_x(|X|) = 1$.
En particulier, la dimension de $X$ (telle qu'elle est définie
ci-dessous) est différente de la dimension de $|X|$.

\begin{definition}
\label{spaces-properties-definition-dimension}
Soit $S$ un schéma. Soit $X$ un espace algébrique sur $S$.
La {\it dimension} $\dim(X)$ de $X$ est définie par la formule
$$
\dim(X) = \sup\nolimits_{x \in |X|} \dim_x(X)
$$
\end{definition}

\noindent
D'après
Propriétés, lemme \ref{properties-lemma-dimension},
on voit que c'est la notion usuelle lorsque $X$ est un schéma.
Il existe un autre invariant qui mesure la dimension d'un schéma
en un point, à savoir la dimension de l'anneau local. Cet invariant
est lui aussi compatible avec les morphismes \'etales, voir la
section \ref{spaces-properties-section-dimension-local-ring}.






\section{Dimension des anneaux locaux}
\label{spaces-properties-section-dimension-local-ring}

\noindent
La dimension de l'anneau local d'un espace algébrique est une notion bien
définie.

\begin{lemma}
\label{spaces-properties-lemma-pre-dimension-local-ring}
Soit $S$ un schéma.
Soit $X$ un espace algébrique sur $S$.
Soit $x \in |X|$ un point.
Soit $d \in \{0, 1, 2, \ldots, \infty\}$.
Les assertions suivantes sont équivalentes :
\begin{enumerate}
\item il existe un schéma $U$, un morphisme \'etale $a : U \to X$ et un point
$u \in U$ tel que $a(u) = x$ et $\dim(\mathcal{O}_{U, u}) = d$,
\item pour tout schéma $U$, tout morphisme \'etale $a : U \to X$ et tout point
$u \in U$ tel que $a(u) = x$, on a $\dim(\mathcal{O}_{U, u}) = d$.
\end{enumerate}
Si $X$ est un schéma, ces conditions équivalent à $\dim(\mathcal{O}_{X, x}) = d$.
\end{lemma}

\begin{proof}
Combiner le
lemme \ref{spaces-properties-lemma-local-source-target-at-point} et
Descente, lemme \ref{descent-lemma-dimension-local-ring-local}.
\end{proof}

\begin{definition}
\label{spaces-properties-definition-dimension-local-ring}
Soit $S$ un schéma. Soit $X$ un espace algébrique sur $S$. Soit $x \in |X|$
un point. La {\it dimension de l'anneau local de $X$ en $x$} est
l'élément $d \in \{0, 1, 2, \ldots, \infty\}$ qui satisfait les conditions
équivalentes du lemme \ref{spaces-properties-lemma-pre-dimension-local-ring}. Dans ce cas, on
dira aussi que {\it $x$ est un point de codimension $d$ dans $X$}.
\end{definition}

\noindent
Outre le lemme ci-dessous, signalons aussi au lecteur les
lemmes \ref{spaces-properties-lemma-dimension-local-ring} et
\ref{spaces-properties-lemma-dimension-decent-invariant-under-etale}.

\begin{lemma}
\label{spaces-properties-lemma-dimension}
Soit $S$ un schéma. Soit $X$ un espace algébrique sur $S$.
Les quantités suivantes sont égales :
\begin{enumerate}
\item La dimension de $X$.
\item La borne supérieure des dimensions des anneaux locaux de $X$.
\item La borne supérieure des $\dim_x(X)$ pour $x \in |X|$.
\end{enumerate}
\end{lemma}

\begin{proof}
Les nombres de (1) et (3) sont égaux par la définition \ref{spaces-properties-definition-dimension}.
Soit $U \to X$ un morphisme \'etale surjectif issu d'un schéma $U$.
La borne supérieure des $\dim_x(X)$ pour $x \in |X|$ est égale à la
borne supérieure des $\dim_u(U)$ pour les points $u$ de $U$ par définition.
Celle-ci est égale à la borne supérieure des $\dim(\mathcal{O}_{U, u})$ d'après
Propriétés, lemme \ref{properties-lemma-dimension}. Ce dernier
est à son tour égal à (2) par définition.
\end{proof}




\section{Points génériques}
\label{spaces-properties-section-generic-points}

\noindent
Soit $T$ un espace topologique. D'après la deuxième édition de
l'EGA I, un {\it point maximal de $T$} est un point générique d'une composante
irréductible de $T$. Si $T = |X|$ est l'espace topologique associé à
un espace algébrique $X$, il existe au moins deux notions de points maximaux :
on peut considérer les points maximaux de $T$ vu comme espace topologique, ou
les images des points maximaux de $U$, où $U \to X$ est un
morphisme \'etale et $U$ un schéma. La seconde notion correspond à
l'ensemble des points de codimension $0$ (lemme \ref{spaces-properties-lemma-codimension-0-points}).
Les points de codimension $0$ sont plus faciles
à manier pour les espaces algébriques généraux ; les deux notions
coïncident pour les espaces algébriques quasi-séparés et, plus généralement, décents
(Espaces algébriques décents, lemme \ref{decent-spaces-lemma-decent-generic-points}).

\begin{lemma}
\label{spaces-properties-lemma-codimension-0-points}
Soit $S$ un schéma et soit $X$ un espace algébrique sur $S$.
Soit $x \in |X|$. Considérons les morphismes \'etales $a : U \to X$, où
$U$ est un schéma. Les assertions suivantes sont équivalentes :
\begin{enumerate}
\item $x$ est un point de codimension $0$ dans $X$,
\item il existe $U \to X$ comme ci-dessus et $u \in U$ tel que $a(u) = x$,
et le point $u$ est le point générique d'une composante irréductible de $U$,
\item pour tout $U \to X$ comme ci-dessus et tout $u \in U$ ayant pour image $x$,
le point $u$ est le point générique d'une composante irréductible de $U$.
\end{enumerate}
Si $X$ est représentable, ces conditions équivalent à ce que $x$ soit un point générique
d'une composante irréductible de $|X|$.
\end{lemma}

\begin{proof}
Remarquons qu'un point $u$ d'un schéma $U$ est le point générique d'une
composante irréductible de $U$ si et seulement si $\dim(\mathcal{O}_{U, u}) = 0$
(Propriétés, lemme \ref{properties-lemma-generic-point}).
Le résultat découle donc de la définition de la codimension d'un
point dans $X$ (définition \ref{spaces-properties-definition-dimension-local-ring}).
\end{proof}

\begin{lemma}
\label{spaces-properties-lemma-codimension-0-points-dense}
Soit $S$ un schéma et soit $X$ un espace algébrique sur $S$.
L'ensemble des points de codimension $0$ de $X$ est dense dans $|X|$.
\end{lemma}

\begin{proof}
Si $U$ est un schéma, l'ensemble des points génériques de ses composantes
irréductibles est dense dans $U$ (cela vaut pour tout espace topologique quasi-sobre).
Ainsi, si $U \to X$ est un morphisme \'etale surjectif, l'ensemble
des points de codimension $0$ de $X$ est l'image d'une partie dense de
$|U|$ (lemme \ref{spaces-properties-lemma-codimension-0-points}). Puisque $|X|$ possède la
topologie quotient induite par $|U| \to |X|$, on conclut.
\end{proof}





\section{Espaces réduits}
\label{spaces-properties-section-reduced}

\noindent
Nous avons déjà défini les espaces algébriques réduits à la
section \ref{spaces-properties-section-types-properties}.
Nous démontrons ici seulement quelques lemmes simples sur les espaces
algébriques réduits.

\begin{lemma}
\label{spaces-properties-lemma-subspace-induced-topology}
Soit $S$ un schéma. Soit $Z \to X$ une immersion d'espaces algébriques.
Alors $|Z| \to |X|$ est un homéomorphisme de $|Z|$ sur une partie localement fermée
de $|X|$.
\end{lemma}

\begin{proof}
Soit $U$ un schéma et soit $U \to X$ un morphisme \'etale surjectif.
Alors $Z \times_X U \to U$ est une immersion de schémas, donc induit un
homéomorphisme de $|Z \times_X U|$ sur une partie localement fermée $T'$
de $|U|$. D'après le lemme \ref{spaces-properties-lemma-points-cartesian}, la partie
$T'$ est l'image réciproque de l'image $T$ de $|Z| \to |X|$.
L'application $|Z| \to |X|$ est injective, car la transformation de
foncteurs $Z \to X$ est injective, voir
Espaces, section \ref{spaces-section-Zariski}. D'après
Topologie, lemme \ref{topology-lemma-open-morphism-quotient-topology}
on voit que $T$ est localement fermée dans $|X|$. De plus, l'application continue
$|Z| \to T$ est un homéomorphisme, puisque l'application $|Z \times_X U| \to T'$
est un homéomorphisme et que $|Z \times_X U| \to |Z|$ est submersive.
\end{proof}

\noindent
Le lemme suivant nous aidera à construire des sous-espaces
(localement) fermés.

\begin{lemma}
\label{spaces-properties-lemma-subspaces-presentation}
Soit $S$ un schéma. Soit $j : R \to U \times_S U$ une relation d'équivalence \'etale.
Soit $X = U/R$ l'espace algébrique associé
(Espaces, théorème \ref{spaces-theorem-presentation}). Il existe une
bijection canonique
$$
\begin{gathered}
\text{Sous-schémas localement fermés stables par }R : Z'\text{ de }U\\[2pt]
\leftrightarrow\\[2pt]
\text{sous-espaces localement fermés }Z\text{ de }X
\end{gathered}
$$
De plus, $Z \to X$ est fermé (resp.\ ouvert) si et seulement si
$Z' \to U$ est fermé (resp.\ ouvert).
\end{lemma}

\begin{proof}
Notons $\varphi : U \to X$ l'application canonique. La bijection envoie
$Z \to X$ sur $Z' = Z \times_X U \to U$. Il résulte immédiatement de la définition
que $Z' \to U$ est une immersion, resp.\ une immersion fermée, resp.\ une
immersion ouverte, si tel est le cas de $Z \to X$. Il est aussi clair que $Z'$ est stable par $R$
(voir Groupoïdes, définition \ref{groupoids-definition-invariant-open}).

\medskip\noindent
Réciproquement, supposons que $Z' \to U$ soit une immersion stable par $R$.
Soit $R'$ la restriction de $R$ à $Z'$, voir
Groupoïdes, définition \ref{groupoids-definition-restrict-groupoid}.
Puisque $R' = R \times_{s, U} Z' = Z' \times_{U, t} R$ dans ce cas,
on voit que $R'$ est une relation d'équivalence \'etale sur $Z'$. D'après
Espaces, théorème \ref{spaces-theorem-presentation}, on voit que
$Z = Z'/R'$ est un espace algébrique. Par construction, on a
$U \times_X Z = Z'$ ; ainsi, $U \times_X Z \to U$ est une immersion.
Notons que la propriété ``immersion'' est préservée par changement de base
et locale pour la topologie fppf sur la base (voir Espaces, section \ref{spaces-section-lists}).
De plus, les immersions sont séparées et localement quasi-finies (voir
Schémas, lemme \ref{schemes-lemma-immersions-monomorphisms}
et
Morphismes, lemme \ref{morphisms-lemma-immersion-locally-quasi-finite}).
Par conséquent, d'après Compléments sur les morphismes, lemme
\ref{more-morphisms-lemma-separated-locally-quasi-finite-morphisms-fppf-descend}
les immersions satisfont à la descente pour les recouvrements fppf. Toutes les hypothèses de
Espaces,
lemme \ref{spaces-lemma-morphism-sheaves-with-P-effective-descent-etale}
sont donc satisfaites pour $Z \to X$, $\mathcal{P}=$``immersion'',
et le morphisme \'etale surjectif $U \to X$. On conclut que $Z \to X$
est représentable et est une immersion, ce qui est la
définition d'un sous-espace (voir
Espaces, définition \ref{spaces-definition-immersion}).

\medskip\noindent
Il est clair que ces constructions sont inverses l'une de l'autre, d'où le résultat.
\end{proof}

\begin{lemma}
\label{spaces-properties-lemma-reduced-closed-subspace}
Soit $S$ un schéma.
Soit $X$ un espace algébrique sur $S$.
Soit $T \subset |X|$ une partie fermée.
Il existe un unique sous-espace fermé $Z \subset X$ possédant
les propriétés suivantes : (a) on a $|Z| = T$, et (b) $Z$ est réduit.
\end{lemma}

\begin{proof}
Soit $U \to X$ un morphisme \'etale surjectif, où $U$ est un schéma.
Posons $R = U \times_X U$, de sorte que $X = U/R$, voir
Espaces, lemme \ref{spaces-lemma-space-presentation}.
Comme d'habitude, on note $s, t : R \to U$ les deux morphismes de projection.
D'après le lemme \ref{spaces-properties-lemma-points-presentation},
on voit que $T$ correspond à une partie fermée $T' \subset |U|$ telle
que $s^{-1}(T') = t^{-1}(T')$.
Soit $Z' \subset U$ le sous-schéma porté par $T'$ et muni de la structure réduite induite.
Dans ce cas, les produits fibrés
$Z' \times_{U, t} R$ et $Z' \times_{U, s} R$ sont des sous-schémas fermés
de $R$
(Schémas, lemme \ref{schemes-lemma-base-change-immersion})
qui sont \'etales sur $Z'$
(Morphismes, lemme \ref{morphisms-lemma-base-change-etale}),
et sont donc réduits
(car la propriété d'être réduit est locale pour la topologie \'etale, voir la
remarque \ref{spaces-properties-remark-list-properties-local-etale-topology}).
Comme ils ont le même espace topologique sous-jacent (voir ci-dessus),
on conclut que $Z' \times_{U, t} R = Z' \times_{U, s} R$.
On peut donc appliquer le lemme \ref{spaces-properties-lemma-subspaces-presentation}
pour obtenir un sous-espace fermé $Z \subset X$ dont l'image réciproque sur $U$ est $Z'$.
Par construction, $|Z| = T$ et $Z$ est réduit. Cela prouve l'existence.
Nous omettons la démonstration de l'unicité.
\end{proof}

\begin{lemma}
\label{spaces-properties-lemma-map-into-reduction}
Soit $S$ un schéma.
Soient $X$, $Y$ des espaces algébriques sur $S$.
Soit $Z \subset X$ un sous-espace fermé.
Supposons $Y$ réduit.
Un morphisme $f : Y \to X$ se factorise par $Z$ si et seulement si
$f(|Y|) \subset |Z|$.
\end{lemma}

\begin{proof}
Supposons $f(|Y|) \subset |Z|$. Choisissons un diagramme
$$
\xymatrix{
V \ar[d]_b \ar[r]_h & U \ar[d]^a \\
Y \ar[r]^f & X
}
$$
où $U$, $V$ sont des schémas et les flèches verticales sont surjectives et
\'etales. Le schéma $V$ est réduit, voir le
lemme \ref{spaces-properties-lemma-type-property}.
Par conséquent, $h$ se factorise par $a^{-1}(Z)$ d'après
Schémas, lemme \ref{schemes-lemma-map-into-reduction}.
Ainsi, $a \circ h$ se factorise par $Z$.
Puisque $Z \subset X$ est un sous-faisceau et que $V \to Y$ est une surjection de faisceaux
sur $(\Sch/S)_{fppf}$, on conclut que $Y \to X$ se factorise
par $Z$.
\end{proof}

\begin{definition}
\label{spaces-properties-definition-reduced-induced-space}
Soit $S$ un schéma et soit $X$ un espace algébrique sur $S$.
Soit $Z \subset |X|$ une partie fermée.
Une {\it structure d'espace algébrique sur $Z$} est donnée par un sous-espace fermé
$Z'$ de $X$ tel que $|Z'|$ soit égal à $Z$.
La {\it structure d'espace algébrique réduite induite}
sur $Z$ est celle qui a été construite dans le
lemme \ref{spaces-properties-lemma-reduced-closed-subspace}.
La {\it réduction $X_{red}$ de $X$} est la structure d'espace
algébrique réduite induite sur $|X|$.
\end{definition}











\section{Le lieu schématique}
\label{spaces-properties-section-schematic}

\noindent
Tout espace algébrique possède un plus grand sous-espace ouvert qui est un
schéma ; cela est plus ou moins clair, mais nous en donnons aussi la preuve ci-dessous.
Bien entendu, ce sous-espace peut être vide, par exemple si
$X = \mathbf{A}^1_{\mathbf{Q}}/\mathbf{Z}$ (le
contre-exemple universel). En revanche, si $X$ est par exemple quasi-séparé,
alors ce plus grand sous-espace ouvert est en fait dense dans $X$ !

\begin{lemma}
\label{spaces-properties-lemma-subscheme}
Soit $S$ un schéma.
Soit $X$ un espace algébrique sur $S$.
Il existe un plus grand sous-espace ouvert $X' \subset X$ qui est un schéma.
\end{lemma}

\begin{proof}
Soit $U \to X$ un morphisme \'etale surjectif, où $U$ est un schéma.
Soit $R = U \times_X U$. Les sous-espaces ouverts de $X$ correspondent biunivoquement
avec les sous-schémas ouverts de $U$ qui sont stables par $R$. Il existe donc un
ensemble de tels sous-espaces. Soient $X_i$, $i \in I$, les sous-espaces ouverts
de $X$ qui sont des schémas, c'est-à-dire qui sont représentables. Considérons le
sous-espace ouvert $X' \subset X$ dont l'ensemble de points sous-jacent est
l'ouvert $\bigcup |X_i|$ de $|X|$. D'après le
lemme \ref{spaces-properties-lemma-characterize-surjective},
on voit que
$$
\coprod X_i \longrightarrow X'
$$
est une application surjective de faisceaux sur $(\Sch/S)_{fppf}$.
Mais puisque chaque $X_i \to X'$ est représentable par des immersions ouvertes,
on voit qu'en fait l'application est surjective pour la topologie de Zariski.
En effet, si $T \to X'$ est un morphisme d'un schéma
vers $X'$, alors $X_i \times_{X'} T$ est un sous-schéma ouvert de $T$.
On peut donc appliquer
Schémas, lemme \ref{schemes-lemma-glue-functors}
pour voir que $X'$ est un schéma.
\end{proof}

\noindent
Dans la suite de cette section, on dit qu'un sous-espace ouvert
$X'$ d'un espace algébrique $X$ est {\it dense} si la partie ouverte
correspondante $|X'| \subset |X|$ est dense.

\begin{lemma}
\label{spaces-properties-lemma-quasi-separated-finite-etale-cover-dense-open-scheme}
Soit $S$ un schéma. Soit $X$ un espace algébrique sur $S$.
S'il existe un morphisme fini, \'etale et surjectif
$U \to X$ où $U$ est un schéma quasi-séparé, alors
il existe un sous-espace ouvert dense $X'$ de $X$ qui est un schéma.
Plus précisément, tout point $x \in |X|$ de codimension $0$ dans $X$
appartient à $X'$.
\end{lemma}

\begin{proof}
Soit $X' \subset X$ le sous-espace ouvert maximal qui est un schéma
(lemme \ref{spaces-properties-lemma-subscheme}).
Soit $x \in |X|$ un point de codimension $0$ dans $X$.
D'après le lemme \ref{spaces-properties-lemma-codimension-0-points-dense},
il suffit de montrer que $x \in X'$.
Soit $U \to X$ comme dans l'énoncé du lemme.
Écrivons $R = U \times_X U$ et notons $s, t : R \to U$ les projections usuelles.
Notons que $s, t$ sont surjectifs, finis et \'etales.
D'après le lemme \ref{spaces-properties-lemma-finite-fibres-presentation},
la fibre de $|U| \to |X|$ au-dessus de $x$ est finie, disons
$\{\eta_1, \ldots, \eta_n\}$. D'après le lemme \ref{spaces-properties-lemma-codimension-0-points},
chaque $\eta_i$ est le point générique d'une composante irréductible de $U$.
D'après Propriétés, lemme \ref{properties-lemma-maximal-points-affine},
on peut trouver un ouvert affine $W \subset U$ contenant
$\{\eta_1, \ldots, \eta_n\}$
(c'est ici que l'on utilise le fait que $U$ est quasi-séparé). D'après
Groupoïdes, lemme \ref{groupoids-lemma-find-invariant-affine}
on peut supposer que $W$ est stable par $R$.
Puisque $W \subset U$ est un ouvert affine stable par $R$, la restriction
$R_W$ de $R$ à $W$ est égale à $R_W = s^{-1}(W) = t^{-1}(W)$ (voir
Groupoïdes, définition \ref{groupoids-definition-invariant-open}
et la discussion qui la suit). En particulier, les morphismes $R_W \to W$ sont
eux aussi finis et \'etales. Il en résulte que $R_W$ est affine.
On voit donc que $W/R_W$ est un schéma, d'après
Groupoïdes, proposition \ref{groupoids-proposition-finite-flat-equivalence}.
D'autre part, $W/R_W$ est un sous-espace ouvert de $X$ d'après
Espaces, lemme \ref{spaces-lemma-finding-opens}, et il contient
$x$ par construction.
\end{proof}

\noindent
Nous étendrons la proposition suivante au cas des
espaces algébriques décents dans
Espaces algébriques décents, théorème
\ref{decent-spaces-theorem-decent-open-dense-scheme}.

\begin{proposition}
\label{spaces-properties-proposition-locally-quasi-separated-open-dense-scheme}
Soit $S$ un schéma. Soit $X$ un espace algébrique sur $S$. Si $X$ est
localement quasi-séparé pour la topologie de Zariski (par exemple si $X$ est quasi-séparé), alors
il existe un sous-espace ouvert dense $X'$ de $X$ qui est un schéma.
Plus précisément, tout point $x \in |X|$ de codimension $0$ dans $X$
appartient à $X'$.
\end{proposition}

\begin{proof}
La question est locale sur $X$ d'après le lemme \ref{spaces-properties-lemma-subscheme}.
Ainsi, d'après le lemme \ref{spaces-properties-lemma-quasi-separated-quasi-compact-pieces},
on peut supposer qu'il existe un schéma affine $U$ et un morphisme
surjectif, quasi-compact et \'etale $U \to X$.
De plus, $U \to X$ est séparé (lemme \ref{spaces-properties-lemma-separated-cover}).
Posons $R = U \times_X U$ et notons $s, t : R \to U$ les projections
usuelles. Alors $s, t$ sont surjectifs, quasi-compacts, séparés et
\'etales. Par conséquent, $s, t$ sont aussi quasi-finis et ont des fibres finies
(Morphismes,
lemmes \ref{morphisms-lemma-etale-locally-quasi-finite},
\ref{morphisms-lemma-quasi-finite-locally-quasi-compact}, et
\ref{morphisms-lemma-quasi-finite}).
D'après Morphismes, lemme \ref{morphisms-lemma-generically-finite},
pour tout $\eta \in U$ qui est le point générique d'une
composante irréductible de $U$, il existe un voisinage ouvert
$V \subset U$ de $\eta$ tel que $s^{-1}(V) \to V$ soit fini. D'après
Descente, lemme \ref{descent-lemma-descending-property-finite},
la propriété d'être fini est locale pour la topologie fpqc (et en particulier \'etale) sur le but.
On peut donc appliquer
Compléments sur les groupoïdes, lemme \ref{more-groupoids-lemma-property-invariant},
qui affirme que le plus grand ouvert $W \subset U$ au-dessus duquel $s$ est
fini est stable par $R$. D'après ce qui précède, $W$ contient tout
point générique d'une composante irréductible de $U$.
La restriction $R_W$ de $R$ à $W$ est égale à $R_W = s^{-1}(W) = t^{-1}(W)$
(voir Groupoïdes, définition \ref{groupoids-definition-invariant-open}
et la discussion qui la suit).
Par construction, $s_W, t_W : R_W \to W$ sont finis \'etales.
Considérons le sous-espace ouvert $X' = W/R_W \subset X$ (voir
Espaces, lemme \ref{spaces-lemma-finding-opens}).
Par construction, l'application d'inclusion $X' \to X$
induit une bijection sur les points de codimension $0$.
On est ainsi ramené au
lemme \ref{spaces-properties-lemma-quasi-separated-finite-etale-cover-dense-open-scheme}.
\end{proof}




\section{Obtention d'un schéma}
\label{spaces-properties-section-getting-a-scheme}

\noindent
Nous avons utilisé dans la section précédente le fait que le quotient $U/R$ d'un
schéma affine $U$ par une relation d'équivalence $R$ est un schéma si les
morphismes $s, t : R \to U$ sont finis \'etales. C'est un cas particulier
du résultat suivant.

\begin{proposition}
\label{spaces-properties-proposition-finite-flat-equivalence-global}
Soit $S$ un schéma.
Soit $(U, R, s, t, c)$ un groupoïde en schémas sur $S$.
Supposons que
\begin{enumerate}
\item $s, t : R \to U$ soient finis localement libres,
\item $j = (t, s)$ soit une relation d'équivalence, et
\item toute partie fermée non vide $Z$ de $U$ contienne un point
$u$ dont la classe d'équivalence pour $R$, $t(s^{-1}(\{u\}))$, soit contenue
dans un ouvert affine de $U$\footnote{Soit $E \subset U$ la
partie des points $u$ tels que $t(s^{-1}(\{u\}))$
soit contenu dans un ouvert affine de $U$. La condition (3) est satisfaite
si $E = U$, ou si tout point de type fini de $U$ appartient à $E$, ou
si tout $u \in U$ se spécialise en un point de $E$.}.
\end{enumerate}
Alors il existe un morphisme fini localement libre $U \to M$
de schémas sur $S$ tel que $R = U \times_M U$ et tel que $M$
représente le faisceau quotient $U/R$ pour la topologie fppf.
\end{proposition}

\begin{proof}
Par l'hypothèse (3) et
Groupoïdes, lemme \ref{groupoids-lemma-find-invariant-affine}
nous pouvons trouver un recouvrement ouvert $U = \bigcup U_i$ tel que chaque $U_i$
soit stable par $R$ et ouvert affine dans $U$. Posons $R_i = R|_{U_i}$.
Considérons les faisceaux fppf $F = U/R$ et $F_i = U_i/R_i$.
D'après Espaces, lemme \ref{spaces-lemma-finding-opens}, les morphismes
$F_i \to F$ sont représentables et sont des immersions ouvertes.
D'après Groupoïdes, proposition \ref{groupoids-proposition-finite-flat-equivalence},
les faisceaux $F_i$ sont représentables par des schémas affines.
Si $T$ est un schéma et $T \to F$ un morphisme, alors $V_i = F_i \times_F T$
est ouvert dans $T$, et nous affirmons que $T = \bigcup V_i$. En effet,
localement pour la topologie fppf sur $T$, nous pouvons relever $T \to F$ en un morphisme
$f : T \to U$ et, dans ce cas, $f^{-1}(U_i) \subset V_i$.
Nous en concluons que $F$ est représentable par un schéma, voir
Schémas, lemme \ref{schemes-lemma-glue-functors}.
\end{proof}

\noindent
Par exemple, si $U$ est isomorphe à un sous-schéma localement fermé d'un
schéma affine ou à un sous-schéma localement fermé de
$\text{Proj}(A)$ pour un anneau gradué $A$, alors la troisième hypothèse est satisfaite d'après
Propriétés, lemme \ref{properties-lemma-ample-finite-set-in-affine}.
En particulier, nous pouvons appliquer ce résultat aux actions libres de groupes finis et de
schémas en groupes finis sur des schémas quasi-affines ou quasi-projectifs.
Par exemple, le quotient $X/G$ d'une variété quasi-projective
$X$ par l'action libre d'un groupe fini $G$ est un schéma. Voici un
énoncé détaillé.

\begin{lemma}
\label{spaces-properties-lemma-quotient-scheme}
Soit $S$ un schéma. Soit $G \to S$ un schéma en groupes. Soit $X \to S$ un
morphisme de schémas. Soit $a : G \times_S X \to X$ une action. Supposons que
\begin{enumerate}
\item $G \to S$ soit fini localement libre,
\item l'action $a$ soit libre,
\item $X \to S$ soit affine, ou quasi-affine, ou projectif, ou
quasi-projectif, ou que $X$ soit isomorphe à un sous-schéma ouvert d'un
schéma affine, ou que $X$ soit isomorphe à un sous-schéma ouvert de $\text{Proj}(A)$
pour un anneau gradué $A$, ou que $G \to S$ soit radiciel.
\end{enumerate}
Alors le faisceau quotient $X/G$ pour la topologie fppf est un schéma et $X \to X/G$
est un torseur sous $G$ pour la topologie fppf.
\end{lemma}

\begin{proof}
Montrons d'abord que $X/G$ est un schéma. Puisque l'action est libre, le morphisme
$j = (a, \text{pr}) : G \times_S X \to X \times_S X$
est un monomorphisme et donc une relation d'équivalence, voir
Groupoïdes, lemme \ref{groupoids-lemma-free-action}. Les morphismes
$s, t : G \times_S X \to X$
sont finis localement libres, puisque nous avons supposé que $G \to S$ est fini localement
libre. Pour conclure, il suffit maintenant de vérifier la dernière hypothèse de la
proposition \ref{spaces-properties-proposition-finite-flat-equivalence-global}.
Puisque l'action de $G$ est définie sur $S$, il suffit de montrer que
toute partie finie d'une fibre de $X \to S$ est contenue dans un
ouvert affine de $X$. Si $X$ est isomorphe à un sous-schéma ouvert d'un
schéma affine ou à un sous-schéma ouvert de $\text{Proj}(A)$
pour un anneau gradué $A$, cela résulte de
Propriétés, lemme \ref{properties-lemma-ample-finite-set-in-affine}.
Si $X \to S$ est affine, ou quasi-affine, ou projectif, ou
quasi-projectif, nous pouvons remplacer $S$ par un ouvert affine et nous
nous ramenons au cas que nous venons de traiter. Si $G \to S$ est radiciel,
alors les orbites des points de $X$ sous l'action de $G$ sont réduites à un point
et la condition est trivialement satisfaite. Nous omettons quelques détails.

\medskip\noindent
Pour voir que $X \to X/G$ est un torseur sous $G$ pour la topologie fppf
(Groupoïdes, définition \ref{groupoids-definition-principal-homogeneous-space})
il faut montrer que $G \times_S X \to X \times_{X/G} X$
est un isomorphisme et que $X \to X/G$ admet localement des sections pour la topologie fppf.
La seconde assertion résulte du fait que $X \to X/G$ est surjectif
comme morphisme de faisceaux fppf (par construction). La première résulte de
l'isomorphisme $R = U \times_M U$ dans la conclusion de la
proposition \ref{spaces-properties-proposition-finite-flat-equivalence-global}
(notons que $R = G \times_S X$ dans notre cas).
\end{proof}

\begin{lemma}
\label{spaces-properties-lemma-quotient-separated}
Notations et hypothèses comme dans la
proposition \ref{spaces-properties-proposition-finite-flat-equivalence-global}. Alors
\begin{enumerate}
\item si $U$ est quasi-séparé sur $S$, alors $U/R$ est quasi-séparé
sur $S$,
\item si $U$ est quasi-séparé, alors $U/R$ est quasi-séparé,
\item si $U$ est séparé sur $S$, alors $U/R$ est séparé sur $S$,
\item si $U$ est séparé, alors $U/R$ est séparé, et
\item d'autres énoncés pourront être ajoutés ici.
\end{enumerate}
Des résultats analogues valent dans la situation du Lemme \ref{spaces-properties-lemma-quotient-scheme}.
\end{lemma}

\begin{proof}
Puisque $M$ représente le faisceau quotient, nous avons un diagramme cartésien
$$
\xymatrix{
R \ar[r]_-j \ar[d] & U \times_S U \ar[d] \\
M \ar[r] & M \times_S M
}
$$
de schémas. Puisque $U \times_S U \to M \times_S M$ est surjectif, fini et localement
libre, pour montrer que $M \to M \times_S M$ est quasi-compact, resp.\ une
immersion fermée, il suffit de montrer que $j : R \to U \times_S U$ est
quasi-compact, resp.\ une immersion fermée, voir
Descente, lemmes \ref{descent-lemma-descending-property-quasi-compact} et
\ref{descent-lemma-descending-property-closed-immersion}.
Puisque $j : R \to U \times_S U$ est un morphisme sur $U$ et que
$R$ est fini sur $U$, nous voyons que $j$ est quasi-compact dès que
le morphisme de projection $U \times_S U \to U$ est quasi-séparé
(Schémas, lemme \ref{schemes-lemma-quasi-compact-permanence}).
Puisque $j$ est un monomorphisme localement de type fini, nous voyons que
$j$ est une immersion fermée dès qu'il est propre
(Morphismes \'etales, lemme \ref{etale-lemma-characterize-closed-immersion})
ce qui est le cas dès que le morphisme de projection
$U \times_S U \to U$ est séparé
(Morphismes, lemme \ref{morphisms-lemma-image-proper-scheme-closed}).
Ceci prouve (1) et (3). Pour prouver (2) et (4), nous remplaçons $S$ par
$\Spec(\mathbf{Z})$, voir définition \ref{spaces-properties-definition-separated}.
Puisque le Lemme \ref{spaces-properties-lemma-quotient-scheme} se démontre en appliquant la
proposition \ref{spaces-properties-proposition-finite-flat-equivalence-global},
le dernier énoncé est lui aussi clair.
\end{proof}




\section{Points des espaces quasi-séparés}
\label{spaces-properties-section-points-quasi-separated}

\noindent
Les points peuvent se comporter très mal pour les espaces algébriques, au degré de généralité adopté
dans le Projet Champs. Cependant, pour les espaces quasi-séparés, leur comportement
est essentiellement semblable à celui des points des schémas. Nous démontrons quelques résultats
à ce sujet dans cette section ; le chapitre sur les espaces décents en contient bien davantage,
voir par exemple
Espaces algébriques décents, section \ref{decent-spaces-section-points}.

\begin{lemma}
\label{spaces-properties-lemma-quasi-separated-sober}
Soit $S$ un schéma. Soit $X$ un espace algébrique qui soit localement
quasi-séparé pour la topologie de Zariski et défini sur $S$. Alors l'espace topologique $|X|$ est sobre (voir
Topologie, définition \ref{topology-definition-generic-point}).
\end{lemma}

\begin{proof}
En combinant
Topologie, lemme \ref{topology-lemma-sober-local}
et le
lemme \ref{spaces-properties-lemma-quasi-separated-quasi-compact-pieces}
on voit que l'on peut supposer qu'il existe un schéma affine $U$
et un morphisme \'etale, surjectif et quasi-compact $U \to X$.
Posons $R = U \times_X U$, avec pour projections $s, t : R \to U$. En appliquant le
lemme \ref{spaces-properties-lemma-finite-fibres-presentation}
on voit que les fibres de $s, t$ sont finies. Il s'ensuit que toutes les hypothèses de
Topologie, lemme \ref{topology-lemma-quotient-kolmogorov}
sont satisfaites, et l'on conclut que $|X|$ est un espace de Kolmogoroff\footnote{
En fait, nous utilisons aussi ici le
Schémas, lemme \ref{schemes-lemma-scheme-sober} (sobriété des schémas),
Morphismes, les lemmes \ref{morphisms-lemma-etale-flat}
et \ref{morphisms-lemma-generalizations-lift-flat} (les générisations
se relèvent le long des morphismes \'etales), le
lemme \ref{spaces-properties-lemma-points-presentation} (les points d'un espace algébrique en
termes d'une présentation), et le
lemme \ref{spaces-properties-lemma-topology-points} (l'application quotient est ouverte).}.

\medskip\noindent
Il reste à montrer que toute partie fermée irréductible
$T \subset |X|$ possède un point générique. D'après le
lemme \ref{spaces-properties-lemma-reduced-closed-subspace}
il existe un sous-espace fermé $Z \subset X$ tel que $|Z| = T$.
Notons que $U \times_X Z \to Z$ est un morphisme \'etale, surjectif et
quasi-compact d'un schéma affine vers $Z$ ; ainsi, $Z$ est localement
quasi-séparé pour la topologie de Zariski d'après le
lemme \ref{spaces-properties-lemma-quasi-separated-quasi-compact-pieces}.
Par la
proposition \ref{spaces-properties-proposition-locally-quasi-separated-open-dense-scheme}
nous voyons qu'il existe un sous-espace ouvert dense $Z' \subset Z$
qui est un schéma. Cela signifie que $|Z'| \subset T$ est ouvert et dense.
L'espace topologique $|Z'|$ est donc irréductible, ce qui signifie que
$Z'$ est un schéma irréductible. D'après
Schémas, lemme \ref{schemes-lemma-scheme-sober}
nous concluons que $|Z'|$ est l'adhérence d'un unique point
$\eta \in |Z'| \subset T$ et donc aussi que $T = \overline{\{\eta\}}$, ce qui conclut.
\end{proof}

\begin{lemma}
\label{spaces-properties-lemma-quasi-compact-quasi-separated-spectral}
Soit $S$ un schéma. Soit $X$ un espace algébrique quasi-compact et
quasi-séparé sur $S$. L'espace topologique $|X|$ est un espace spectral.
\end{lemma}

\begin{proof}
D'après Topologie, définition \ref{topology-definition-spectral-space},
nous devons vérifier que $|X|$ est sobre et quasi-compact, qu'il possède une base
d'ouverts quasi-compacts et que l'intersection de deux
ouverts quasi-compacts quelconques est quasi-compacte. D'après le
Lemme \ref{spaces-properties-lemma-quasi-separated-sober}, nous voyons que $|X|$ est sobre.
D'après le Lemme \ref{spaces-properties-lemma-quasi-compact-space}, nous voyons que $|X|$ est quasi-compact.
D'après le Lemme \ref{spaces-properties-lemma-quasi-compact-affine-cover}, il existe un schéma affine
$U$ et un morphisme \'etale surjectif $f : U \to X$.
Puisque $|f| : |U| \to |X|$ est ouverte et continue, et que $|U|$ possède
une base d'ouverts quasi-compacts, nous concluons que $|X|$ possède une base
d'ouverts quasi-compacts. Enfin, supposons que
$A, B \subset |X|$ soient des ouverts quasi-compacts. Alors $A = |X'|$ et $B = |X''|$
pour des sous-espaces ouverts $X', X'' \subset X$ (Lemme \ref{spaces-properties-lemma-open-subspaces}),
et nous pouvons choisir des schémas affines $V$ et $W$ et des morphismes
\'etales surjectifs $V \to X'$ et $W \to X''$
(Lemme \ref{spaces-properties-lemma-quasi-compact-affine-cover}).
Alors $A \cap B$ est l'image de
$|V \times_X W| \to |X|$ (Lemme \ref{spaces-properties-lemma-points-cartesian}).
Puisque $V \times_X W$ est quasi-compact car $X$ est quasi-séparé
(Lemme \ref{spaces-properties-lemma-characterize-quasi-separated})
nous concluons que $A \cap B$ est quasi-compacte, ce qui achève la démonstration.
\end{proof}

\noindent
Le lemme suivant permet de prouver qu'un
espace algébrique est isomorphe au spectre d'un corps.

\begin{lemma}
\label{spaces-properties-lemma-point-like-spaces}
Soit $S$ un schéma. Soit $k$ un corps.
Soit $X$ un espace algébrique sur $S$ et supposons qu'il existe
un morphisme \'etale surjectif $\Spec(k) \to X$.
Si $X$ est quasi-séparé, alors $X \cong \Spec(k')$,
où $k/k'$ est une extension finie séparable.
\end{lemma}

\begin{proof}
Posons $R = \Spec(k) \times_X \Spec(k)$ ; nous avons donc un
diagramme de produit fibré
$$
\xymatrix{
R \ar[r]_-s \ar[d]_-t & \Spec(k) \ar[d] \\
\Spec(k) \ar[r] & X
}
$$
D'après
Espaces, lemme \ref{spaces-lemma-space-presentation}
nous savons que $X = \Spec(k)/R$ est le faisceau quotient.
Comme $\Spec(k) \to X$ est \'etale, les morphismes $s$ et $t$ sont
\'etales. Ainsi, $R = \coprod_{i \in I} \Spec(k_i)$ est une somme
disjointe de spectres de corps, et $s$ et $t$
induisent tous deux des extensions finies séparables $s, t : k \subset k_i$, voir
Morphismes, lemme \ref{morphisms-lemma-etale-over-field}.
Comme
$$
R = \Spec(k) \times_X \Spec(k)
= (\Spec(k) \times_S \Spec(k)) \times_{X \times_S X, \Delta} X
$$
et que $\Delta$ est quasi-compact par hypothèse, nous concluons que
$R \to \Spec(k) \times_S \Spec(k)$ est quasi-compact.
Ainsi, $R$ est quasi-compact puisque $\Spec(k) \times_S \Spec(k)$ est
affine. Nous concluons que $I$ est fini. Il s'ensuit
que $s$ et $t$ sont des morphismes finis localement libres. D'après
Groupoïdes, proposition \ref{groupoids-proposition-finite-flat-equivalence}
nous concluons que $\Spec(k)/R$ est
représenté par $\Spec(k')$, l'inclusion $k' \subset k$ étant finie localement libre,
où
$$
k' = \{x \in k \mid s_i(x) = t_i(x)\text{ pour tout }i \in I\}
$$
Il est facile de voir que $k'$ est un corps.
\end{proof}

\begin{remark}
\label{spaces-properties-remark-cannot-decide-yet}
Le Lemme \ref{spaces-properties-lemma-point-like-spaces} vaut pour les espaces algébriques décents, voir
Espaces algébriques décents, lemme \ref{decent-spaces-lemma-decent-point-like-spaces}.
En fait, un espace algébrique décent à un seul point est un schéma, voir
Espaces algébriques décents, lemme \ref{decent-spaces-lemma-when-field}.
Cela vaut aussi lorsque $X$ est localement séparé, car un
espace algébrique localement séparé est décent, voir
Espaces algébriques décents, lemme \ref{decent-spaces-lemma-locally-separated-decent}.
\end{remark}







\section{Morphismes \'etales d'espaces algébriques}
\label{spaces-properties-section-etale-morphisms}

\noindent
Cette section devrait en réalité figurer dans le chapitre consacré aux morphismes d'espaces
algébriques, mais nous avons besoin de la notion d'espace algébrique \'etale sur un autre
afin de définir le petit site \'etale d'un espace algébrique.
Nous devons donc effectuer quelques travaux préliminaires sur les morphismes \'etales de schémas vers des
espaces algébriques et sur les morphismes \'etales entre espaces algébriques.
Pour davantage de résultats sur les morphismes \'etales d'espaces algébriques, voir
Morphismes d'espaces, section \ref{spaces-morphisms-section-etale}.

\begin{lemma}
\label{spaces-properties-lemma-etale-over-space}
Soit $S$ un schéma.
Soit $X$ un espace algébrique sur $S$.
Soient $U$, $U'$ des schémas sur $S$.
\begin{enumerate}
\item Si $U \to U'$ est un morphisme \'etale de schémas, et
si $U' \to X$ est un morphisme \'etale de $U'$ vers $X$, alors la
composée $U \to X$ est un morphisme \'etale de $U$ vers $X$.
\item Si $\varphi : U \to X$ et $\varphi' : U' \to X$ sont des
morphismes \'etales vers $X$, et si $\chi : U \to U'$ est un
morphisme de schémas tel que $\varphi = \varphi' \circ \chi$,
alors $\chi$ est un morphisme \'etale de schémas.
\item Si $\chi : U \to U'$ est un morphisme \'etale surjectif
de schémas et si $\varphi' : U' \to X$ est un morphisme tel que
$\varphi = \varphi' \circ \chi$ soit \'etale, alors $\varphi'$
est \'etale.
\end{enumerate}
\end{lemma}

\begin{proof}
Rappelons que la définition d'un morphisme \'etale d'un schéma vers un
espace algébrique est celle de
Espaces, définition
\ref{spaces-definition-relative-representable-property}
puisque tout morphisme d'un schéma vers un espace algébrique
est représentable.

\medskip\noindent
L'assertion (1) du lemme résulte de ceci, du fait que
les morphismes \'etales sont stables par composition
(Morphismes, lemme
\ref{morphisms-lemma-composition-etale})
et des lemmes
d'Espaces
\ref{spaces-lemma-composition-representable-transformations-property} et
\ref{spaces-lemma-morphism-schemes-gives-representable-transformation-property}
(qui sont formels).

\medskip\noindent
Pour prouver l'assertion (2), choisissons un schéma $W$ sur $S$ et un
morphisme \'etale surjectif $W \to X$. Considérons le changement de base
$\chi_W : W \times_X U \to W \times_X U'$ de $\chi$.
Comme $W \times_X U$ et $W \times_X U'$ sont \'etales sur $W$, nous concluons que
$\chi_W$ est \'etale d'après
Morphismes, lemme \ref{morphisms-lemma-etale-permanence}.
D'autre part, dans le diagramme commutatif
$$
\xymatrix{
W \times_X U \ar[r] \ar[d] & W \times_X U' \ar[d] \\
U \ar[r] & U'
}
$$
les deux flèches verticales sont \'etales et surjectives.
Ainsi, d'après
Descente, lemme \ref{descent-lemma-syntomic-smooth-etale-permanence}
nous concluons que $U \to U'$ est \'etale.

\medskip\noindent
Pour prouver l'assertion (3), choisissons un schéma $W$ sur $S$ et un morphisme $W \to X$.
Comme ci-dessus, considérons le diagramme
$$
\xymatrix{
W \times_X U \ar[r] \ar[d] & W \times_X U' \ar[d] \ar[r] & W \ar[d] \\
U \ar[r] & U' \ar[r] & X
}
$$
Nous savons que $W \times_X U \to W \times_X U'$ est \'etale surjectif
(comme changement de base de $U \to U'$)
et que $W \times_X U \to W$ est \'etale. Ainsi, $W \times_X U' \to W$
est \'etale d'après Descente, lemme
\ref{descent-lemma-syntomic-smooth-etale-permanence}. Par définition,
cela signifie que $\varphi'$ est \'etale.
\end{proof}

\begin{definition}
\label{spaces-properties-definition-etale}
Soit $S$ un schéma.
Un morphisme $f : X \to Y$ entre espaces algébriques sur $S$ est
dit {\it \'etale} si et seulement si, pour tout morphisme \'etale
$\varphi : U \to X$, où $U$ est un schéma, la composée
$f \circ \varphi$ est elle aussi \'etale.
\end{definition}

\noindent
Si $X$ et $Y$ sont des schémas, cela coïncide avec la notion usuelle de
morphisme \'etale de schémas. En fait, dès que $X \to Y$ est un morphisme
représentable d'espaces algébriques, cela coïncide avec la notion définie dans
Espaces, définition \ref{spaces-definition-relative-representable-property}.
Cela résulte en combinant le lemme \ref{spaces-properties-lemma-etale-local} ci-dessous avec
Espaces, lemme \ref{spaces-lemma-representable-morphisms-spaces-property}.

\begin{lemma}
\label{spaces-properties-lemma-etale-local}
Soit $S$ un schéma.
Soit $f : X \to Y$ un morphisme d'espaces algébriques sur $S$.
Les assertions suivantes sont équivalentes :
\begin{enumerate}
\item $f$ est \'etale,
\item il existe un morphisme \'etale surjectif $\varphi : U \to X$,
où $U$ est un schéma, tel que la composée $f \circ \varphi$ soit
\'etale (en tant que morphisme d'espaces algébriques),
\item il existe un morphisme \'etale surjectif $\psi : V \to Y$,
où $V$ est un schéma, tel que le changement de base $V \times_Y X \to V$
soit \'etale (en tant que morphisme d'espaces algébriques),
\item il existe un diagramme commutatif
$$
\xymatrix{
U \ar[d] \ar[r] & V \ar[d] \\
X \ar[r] & Y
}
$$
où $U$, $V$ sont des schémas, les flèches verticales sont \'etales, la
flèche verticale gauche est surjective et la flèche horizontale est \'etale.
\end{enumerate}
\end{lemma}

\begin{proof}
Montrons que (4) implique (1). Supposons donné un diagramme comme en (4).
Soit $W \to X$ un morphisme \'etale, où $W$ est un schéma. Alors
$W \times_X U \to U$ est \'etale. Ainsi, $W \times_X U \to V$ est \'etale
comme composée des morphismes \'etales de schémas $W \times_X U \to U$
et $U \to V$. Par conséquent, $W \times_X U \to Y$ est \'etale d'après le
Lemme \ref{spaces-properties-lemma-etale-over-space} (1). Comme, de plus,
la projection $W \times_X U \to W$ est surjective et \'etale, nous concluons
du Lemme \ref{spaces-properties-lemma-etale-over-space} (3) que $W \to Y$ est \'etale.

\medskip\noindent
Montrons que (1) implique (4). Supposons (1). Choisissons un diagramme commutatif
$$
\xymatrix{
U \ar[d] \ar[r] & V \ar[d] \\
X \ar[r] & Y
}
$$
où $U \to X$ et $V \to Y$ sont surjectifs et \'etales, voir
Espaces, lemme \ref{spaces-lemma-lift-morphism-presentations}.
Par hypothèse, le morphisme $U \to Y$ est \'etale ;
donc $U \to V$ est \'etale d'après le Lemme \ref{spaces-properties-lemma-etale-over-space} (2).

\medskip\noindent
Nous omettons la preuve que (2) et (3) sont aussi équivalents à (1).
\end{proof}

\begin{lemma}
\label{spaces-properties-lemma-composition-etale}
La composée de deux morphismes \'etales d'espaces algébriques
est \'etale.
\end{lemma}

\begin{proof}
Cela résulte immédiatement de la définition.
\end{proof}

\begin{lemma}
\label{spaces-properties-lemma-base-change-etale}
Le changement de base d'un morphisme \'etale d'espaces algébriques
par un morphisme quelconque d'espaces algébriques est \'etale.
\end{lemma}

\begin{proof}
Soit $X \to Y$ un morphisme \'etale d'espaces algébriques sur $S$.
Soit $Z \to Y$ un morphisme d'espaces algébriques.
Choisissons un schéma $U$ et un morphisme \'etale surjectif $U \to X$.
Choisissons un schéma $W$ et un morphisme \'etale surjectif $W \to Z$.
Alors $U \to Y$ est \'etale ; ainsi, dans le diagramme
$$
\xymatrix{
W \times_Y U \ar[d] \ar[r] & W \ar[d] \\
Z \times_Y X \ar[r] & Z
}
$$
la flèche horizontale supérieure est \'etale.
De plus, la flèche verticale gauche est surjective
et \'etale (nous omettons la vérification). Nous concluons donc que la flèche
horizontale inférieure est \'etale d'après le Lemme \ref{spaces-properties-lemma-etale-local}.
\end{proof}

\begin{lemma}
\label{spaces-properties-lemma-etale-permanence}
Soit $S$ un schéma. Soient $X, Y, Z$ des espaces algébriques.
Soient $g : X \to Z$, $h : Y \to Z$ des morphismes \'etales, et soit
$f : X \to Y$ un morphisme tel que $h \circ f = g$.
Alors $f$ est \'etale.
\end{lemma}

\begin{proof}
Choisissons un diagramme commutatif
$$
\xymatrix{
U \ar[d] \ar[r]_\chi & V \ar[d] \\
X \ar[r] & Y
}
$$
où $U \to X$ et $V \to Y$ sont surjectifs et \'etales, voir
Espaces, lemme \ref{spaces-lemma-lift-morphism-presentations}.
Par hypothèse, les morphismes $\varphi : U \to X \to Z$ et
$\psi : V \to Y \to Z$ sont \'etales. De plus, $\psi \circ \chi = \varphi$
par notre hypothèse sur $f, g, h$.
Ainsi, $U \to V$ est \'etale d'après le Lemme \ref{spaces-properties-lemma-etale-over-space},
assertion (2).
\end{proof}

\begin{lemma}
\label{spaces-properties-lemma-etale-open}
Soit $S$ un schéma.
Si $X \to Y$ est un morphisme \'etale d'espaces algébriques sur $S$,
alors l'application associée $|X| \to |Y|$ entre espaces topologiques
est ouverte.
\end{lemma}

\begin{proof}
Cela résulte clairement du diagramme du
lemme \ref{spaces-properties-lemma-etale-local}
et du
lemme \ref{spaces-properties-lemma-topology-points}.
\end{proof}

\noindent
Enfin, voici un joli lemme. Il est faux qu'un espace algébrique
muni d'un morphisme \'etale vers un schéma soit nécessairement un schéma, voir
Espaces, exemple \ref{spaces-example-non-representable-descent}.
Mais cela est vrai si le but est le spectre d'un corps.

\begin{lemma}
\label{spaces-properties-lemma-etale-over-field-scheme}
Soit $S$ un schéma. Soit $X \to \Spec(k)$
un morphisme \'etale sur $S$, où $k$ est un corps.
Alors $X$ est un schéma.
\end{lemma}

\begin{proof}
Soit $U$ un schéma affine et soit $U \to X$ un morphisme \'etale. D'après la
définition \ref{spaces-properties-definition-etale}
nous voyons que $U \to \Spec(k)$ est un morphisme \'etale.
Ainsi, $U = \coprod_{i = 1, \ldots, n} \Spec(k_i)$
est une somme disjointe finie de spectres d'extensions finies séparables
$k_i$ de $k$, voir
Morphismes, lemme \ref{morphisms-lemma-etale-over-field}.
Le morphisme $R = U \times_X U \to U \times_{\Spec(k)} U$ est un monomorphisme,
et $U \times_{\Spec(k)} U$ est aussi une somme disjointe finie de
spectres d'extensions finies séparables de $k$. Ainsi, d'après
Schémas, lemme \ref{schemes-lemma-mono-towards-spec-field}
nous voyons que $R$ est de même une somme disjointe finie de
spectres d'extensions finies séparables de $k$.
Ainsi, $U$ et $R$ sont affines, et
les deux projections $R \to U$ sont finies localement libres.
Par conséquent, $U/R$ est un schéma d'après
Groupoïdes, proposition \ref{groupoids-proposition-finite-flat-equivalence}.
D'après
Espaces, lemme \ref{spaces-lemma-finding-opens}
il est aussi un sous-espace ouvert de $X$. D'après le
lemme \ref{spaces-properties-lemma-subscheme}
nous concluons que $X$ est un schéma.
\end{proof}













\section{Espaces et recouvrements fpqc}
\label{spaces-properties-section-fpqc}

\noindent
Soit $S$ un schéma.
Un espace algébrique sur $S$ est défini comme un faisceau pour la topologie fppf muni de
propriétés supplémentaires. Il n'est donc pas immédiatement clair qu'il vérifie
la condition de faisceau pour la topologie fpqc (voir
Topologies, définition \ref{topologies-definition-sheaf-property-fpqc}).
Dans cette section, nous donnons l'argument de Gabber qui le démontre.
Toutefois, lorsque nous disons que l'espace algébrique $X$ vérifie la
condition de faisceau pour la topologie fpqc, nous ne considérons en réalité que les
recouvrements fpqc $\{f_i : T_i \to T\}_{i \in I}$ tels que $T, T_i$ soient des
objets du grand site $(\Sch/S)_{fppf}$ (conformément à
nos conventions, voir la section \ref{spaces-properties-section-conventions}).

\begin{proposition}[Gabber]
\label{spaces-properties-proposition-sheaf-fpqc}
Soit $S$ un schéma. Soit $X$ un espace algébrique sur $S$. Alors
$X$ vérifie la condition de faisceau pour la topologie fpqc.
\end{proposition}

\begin{proof}
Puisque $X$ est un faisceau pour la topologie de Zariski, il suffit de montrer
ce qui suit. Étant donné un morphisme plat surjectif de schémas affines
$f : T' \to T$, on a :
$X(T)$ est le noyau du couple d'applications $X(T') \to X(T' \times_T T')$.
Voir Topologies, lemme \ref{topologies-lemma-sheaf-property-fpqc}
(nous omettons ici un petit argument, car le lemme cité
est formulé pour des foncteurs définis sur la catégorie de tous les schémas).

\medskip\noindent
Soient $a, b : T \to X$ deux morphismes tels que $a \circ f = b \circ f$.
Il faut montrer que $a = b$. Considérons le produit fibré
$$
E = X \times_{\Delta_{X/S}, X \times_S X, (a, b)} T.
$$
D'après Espaces, lemme \ref{spaces-lemma-properties-diagonal},
le morphisme $\Delta_{X/S}$ est un monomorphisme représentable. Ainsi,
$E \to T$ est un monomorphisme de schémas. Notre hypothèse
$a \circ f = b \circ f$ implique que $T' \to T$ se factorise (uniquement) par $E$.
Considérons le diagramme commutatif
$$
\xymatrix{
T' \times_T E \ar[r] \ar[d] & E \ar[d] \\
T' \ar[r] \ar@/^5ex/[u] \ar[ru] & T
}
$$
Puisque la projection $T' \times_T E \to T'$ est un monomorphisme
admettant une section, nous concluons que c'est un isomorphisme. Nous en déduisons que
$E \to T$ est un isomorphisme d'après
Descente, lemme \ref{descent-lemma-descending-property-isomorphism}.
Cela signifie que $a = b$, comme voulu.

\medskip\noindent
Soit ensuite $c : T' \to X$ un morphisme tel que les deux composées
$T' \times_T T' \to T' \to X$ soient égales. Il faut trouver un morphisme
$a : T \to X$ dont la composée avec $T' \to T$ soit $c$. Choisissons un
schéma affine $U$ et un morphisme \'etale $U \to X$ tels que l'image de
$|U| \to |X|$ contienne l'image de $|c| : |T'| \to |X|$.
C'est possible d'après les lemmes \ref{spaces-properties-lemma-topology-points} et
\ref{spaces-properties-lemma-cover-by-union-affines}, le fait qu'une somme disjointe finie de
schémas affines est affine, et le fait que $|T'|$ est quasi-compact
(nous omettons un petit argument). Puisque $U \to X$ est séparé
(lemme \ref{spaces-properties-lemma-separated-cover}), on voit que
$$
V = U \times_{X, c} T' \longrightarrow T'
$$
est un morphisme de schémas \'etale, surjectif et séparé
(pour voir qu'il est surjectif, utiliser le lemme \ref{spaces-properties-lemma-points-cartesian}
et notre choix de $U \to X$). L'égalité
$c \circ \text{pr}_0 = c \circ \text{pr}_1$ signifie que nous obtenons une
donnée de descente sur $V/T'/T$
(Descente, définition \ref{descent-definition-descent-datum}),
car
\begin{align*}
V \times_{T'} (T' \times_T T')
& =
U \times_{X, c \circ \text{pr}_0} (T' \times_T T') \\
& =
(T' \times_T T') \times_{c \circ \text{pr}_1, X} U \\
& =
(T' \times_T T') \times_{T'} V
\end{align*}
Le morphisme $V \to T'$ est ind-quasi-affine d'après
Compléments sur les morphismes, lemme
\ref{more-morphisms-lemma-etale-separated-ind-quasi-affine}
(car les morphismes \'etales sont localement quasi-finis, voir
Morphismes, lemme \ref{morphisms-lemma-etale-locally-quasi-finite}).
D'après Compléments sur les groupoïdes, lemme \ref{more-groupoids-lemma-ind-quasi-affine},
la donnée de descente est effective. Soit $W \to T$ un morphisme
tel qu'il existe un isomorphisme $\alpha : T' \times_T W \to V$
compatible avec la donnée de descente donnée sur $V$ et la donnée de descente
canonique sur $T' \times_T W$. Alors $W \to T$ est surjectif et \'etale
(Descente, lemmes \ref{descent-lemma-descending-property-surjective} et
\ref{descent-lemma-descending-property-etale}).
Considérons la composée
$$
b' : T' \times_T W \longrightarrow V = U \times_{X, c} T' \longrightarrow U
$$
Les deux composées
$b' \circ (\text{pr}_0, 1), 
b' \circ (\text{pr}_1, 1) :
(T' \times_T T') \times_T W \to T' \times_T W \to U$
coïncident par notre choix de $\alpha$ et la propriété correspondante de $c$
(calcul omis). Ainsi, $b'$ descend en un morphisme $b : W \to U$ d'après
Descente, lemme \ref{descent-lemma-fpqc-universal-effective-epimorphisms}.
Le diagramme
$$
\xymatrix{
T' \times_T W \ar[r] \ar[d] & W \ar[r]_b & U \ar[d] \\
T' \ar[rr]^c &  & X
}
$$
est commutatif. Cela signifie que nous avons démontré l'existence
d'une solution $a$ localement pour la topologie \'etale sur $T$, c'est-à-dire d'un $a' : W \to X$.
Cependant, puisque nous avons démontré l'unicité
dans le premier paragraphe, nous constatons que cette solution locale pour la topologie \'etale
satisfait à la condition de recollement, c'est-à-dire que
$\text{pr}_0^*a' = \text{pr}_1^*a'$ comme éléments de $X(W \times_T W)$.
Puisque $X$ est un faisceau pour la topologie \'etale, il existe un unique $a \in X(T)$ dont la restriction
est $a'$ sur $W$.
\end{proof}









\section{Le site \'etale d'un espace algébrique}
\label{spaces-properties-section-etale-site}

\noindent
Dans cette section, nous définissons le petit site \'etale d'un espace algébrique.
C'est l'analogue du petit site \'etale $S_\etale$ d'un schéma.
Le lemme \ref{spaces-properties-lemma-etale-over-space} implique que, dans la définition ci-dessous,
tout morphisme entre objets du site \'etale de $X$ est \'etale, et que
tout schéma \'etale sur un objet de $X_\etale$ est encore un objet de
$X_\etale$.

\begin{definition}
\label{spaces-properties-definition-etale-site}
Soit $S$ un schéma.
Soit $\Sch_{fppf}$ un grand site fppf contenant $S$,
et soit $\Sch_\etale$ le grand site \'etale correspondant
(c'est-à-dire ayant la même catégorie sous-jacente).
Soit $X$ un espace algébrique sur $S$.
Le {\it petit site \'etale $X_\etale$} de $X$ est défini comme suit :
\begin{enumerate}
\item un objet de $X_\etale$ est un morphisme $\varphi : U \to X$
où $U \in \Ob((\Sch/S)_\etale)$ est un schéma et
$\varphi$ est un morphisme \'etale,
\item un morphisme $(\varphi : U \to X) \to (\varphi' : U' \to X)$
est donné par un morphisme de schémas $\chi : U \to U'$ tel que
$\varphi = \varphi' \circ \chi$, et
\item une famille de morphismes $\{(U_i \to X) \to (U \to X)\}_{i \in I}$
de $X_\etale$ est un recouvrement si et seulement si $\{U_i \to U\}_{i \in I}$
est un recouvrement de $(\Sch/S)_\etale$.
\end{enumerate}
\end{definition}

\noindent
Une conséquence de notre choix est que le site \'etale d'un espace algébrique
n'a en général pas d'objet final ! En revanche, si $X$ est
un schéma, la définition ci-dessus coïncide avec
Topologies, définition \ref{topologies-definition-big-small-etale}.

\medskip\noindent
Le site précédent est celui que nous employons par défaut, mais nous utiliserons aussi
deux variantes. On peut considérer tous les {\it espaces algébriques}
$U$ \'etales sur $X$, ce qui donne le site
$X_{spaces, \etale}$ défini ci-dessous, ou tous les {\it schémas affines}
$U$ \'etales sur $X$, ce qui donne le site
$X_{affine, \etale}$ défini ci-dessous. La première de ces deux notions
sert à étudier la fonctorialité du petit site \'etale, voir le
lemme \ref{spaces-properties-lemma-functoriality-etale-site}.

\begin{definition}
\label{spaces-properties-definition-spaces-etale-site}
Soit $S$ un schéma.
Soit $\Sch_{fppf}$ un grand site fppf contenant $S$,
et soit $\Sch_\etale$ le grand site \'etale correspondant
(c'est-à-dire ayant la même catégorie sous-jacente).
Soit $X$ un espace algébrique sur $S$.
Le site {\it $X_{spaces, \etale}$} de $X$ est défini comme suit :
\begin{enumerate}
\item un objet de $X_{spaces, \etale}$ est un morphisme
$\varphi : U \to X$, où $U$ est un espace algébrique sur $S$ et
$\varphi$ est un morphisme \'etale d'espaces algébriques sur $S$,
\item un morphisme $(\varphi : U \to X) \to (\varphi' : U' \to X)$ de
$X_{spaces, \etale}$ est donné par un morphisme d'espaces algébriques
$\chi : U \to U'$ tel que $\varphi = \varphi' \circ \chi$, et
\item une famille de morphismes
$\{\varphi_i : (U_i \to X) \to (U \to X)\}_{i \in I}$
de $X_{spaces, \etale}$ est un recouvrement si et seulement si
$|U| = \bigcup \varphi_i(|U_i|)$.
\end{enumerate}
Comme d'habitude, nous choisissons un ensemble de recouvrements de ce type, contenant au moins
les recouvrements de $X_\etale$, comme dans
Ensembles, lemme \ref{sets-lemma-coverings-site},
afin de faire de $X_{spaces, \etale}$ un site.
\end{definition}

\noindent
Puisque le morphisme identité de $X$ est \'etale, il est clair que
$X_{spaces, \etale}$ possède un objet final.
Montrons immédiatement que le topos correspondant est équivalent au
petit topos \'etale de $X$.

\begin{lemma}
\label{spaces-properties-lemma-compare-etale-sites}
Le foncteur
$$
X_\etale \longrightarrow X_{spaces, \etale}, \quad
U/X \longmapsto U/X
$$
est un foncteur cocontinu spécial
(Sites, définition \ref{sites-definition-special-cocontinuous-functor})
et induit donc une équivalence de topos
$\Sh(X_\etale) \to \Sh(X_{spaces, \etale})$.
\end{lemma}

\begin{proof}
Il faut montrer que le foncteur satisfait aux hypothèses (1) -- (5) du
Sites, lemme \ref{sites-lemma-equivalence}.
Il est clair que le foncteur est continu et cocontinu, ce qui
prouve les hypothèses (1) et (2).
Les hypothèses (3) et (4) sont satisfaites simplement parce que le foncteur est pleinement fidèle.
L'hypothèse (5) est satisfaite parce que, par définition, un espace algébrique admet
un recouvrement par un schéma.
\end{proof}

\begin{remark}
\label{spaces-properties-remark-explain-equivalence}
Expliquons le sens du lemme \ref{spaces-properties-lemma-compare-etale-sites}.
Soit $S$ un schéma et soit $X$ un espace algébrique sur $S$.
Soit $\mathcal{F}$ un faisceau sur le petit site \'etale $X_\etale$ de
$X$. Le lemme affirme qu'il existe un unique faisceau $\mathcal{F}'$ sur
$X_{spaces, \etale}$ dont la restriction est $\mathcal{F}$ sur la
sous-catégorie $X_\etale$. Si $U \to X$ est un morphisme \'etale
d'espaces algébriques, comment calculer $\mathcal{F}'(U)$ ? Par définition
d'un espace algébrique, il existe un schéma $U'$ et un morphisme
\'etale surjectif $U' \to U$. Alors $\{U' \to U\}$ est un recouvrement de
$X_{spaces, \etale}$, et l'on obtient donc un diagramme exact
$$
\xymatrix{
\mathcal{F}'(U) \ar[r] &
\mathcal{F}(U') \ar@<1ex>[r] \ar@<-1ex>[r] &
\mathcal{F}(U' \times_U U').
}
$$
Notons que $U' \times_U U'$ est un schéma ; nous pouvons donc
écrire $\mathcal{F}$ et non $\mathcal{F}'$.
Nous voyons ainsi comment calculer $\mathcal{F}'$
à partir du faisceau $\mathcal{F}$.
\end{remark}

\begin{definition}
\label{spaces-properties-definition-affine-etale-site}
Soit $S$ un schéma.
Soit $\Sch_{fppf}$ un grand site fppf contenant $S$,
et soit $\Sch_\etale$ le grand site \'etale correspondant
(c'est-à-dire ayant la même catégorie sous-jacente).
Soit $X$ un espace algébrique sur $S$.
Le site {\it $X_{affine, \etale}$} de $X$ est défini comme suit :
\begin{enumerate}
\item un objet de $X_{affine, \etale}$ est un morphisme
$\varphi : U \to X$, où $U \in \Ob((\Sch/S)_\etale)$ est un schéma affine et
$\varphi$ est un morphisme \'etale,
\item un morphisme $(\varphi : U \to X) \to (\varphi' : U' \to X)$ de
$X_{affine, \etale}$ est donné par un morphisme de schémas
$\chi : U \to U'$ tel que $\varphi = \varphi' \circ \chi$, et
\item une famille de morphismes
$\{\varphi_i : (U_i \to X) \to (U \to X)\}_{i \in I}$
de $X_{affine, \etale}$ est un recouvrement si et seulement si
$\{U_i \to U\}$ est un recouvrement \'etale standard, voir
Topologies, définition \ref{topologies-definition-standard-etale}.
\end{enumerate}
Comme d'habitude, nous choisissons un ensemble de recouvrements de ce type, comme dans
Ensembles, lemme \ref{sets-lemma-coverings-site},
afin de faire de $X_{affine, \etale}$ un site.
\end{definition}

\begin{lemma}
\label{spaces-properties-lemma-alternative}
Soit $S$ un schéma. Soit $X$ un espace algébrique sur $S$.
Le foncteur $X_{affine, \etale} \to X_\etale$
est cocontinu spécial et induit une équivalence de topos de
$\Sh(X_{affine, \etale})$ vers $\Sh(X_\etale)$.
\end{lemma}

\begin{proof}
Omis. Indication : comparer avec la preuve de
Topologies, lemme \ref{topologies-lemma-affine-big-site-etale}.
\end{proof}

\begin{definition}
\label{spaces-properties-definition-etale-topos}
Soit $S$ un schéma. Soit $X$ un espace algébrique sur $S$.
Le {\it topos \'etale} de $X$, ou plus précisément le
{\it petit topos \'etale} de $X$, est la catégorie
$\Sh(X_\etale)$
des faisceaux d'ensembles sur $X_\etale$.
\end{definition}

\noindent
D'après le
lemme \ref{spaces-properties-lemma-compare-etale-sites},
on a
$\Sh(X_\etale) = \Sh(X_{spaces, \etale})$,
de sorte que l'on peut aussi voir ce topos comme la catégorie des faisceaux d'ensembles sur
$X_{spaces, \etale}$. De même, d'après le
lemme \ref{spaces-properties-lemma-alternative},
on voit que
$\Sh(X_\etale) = \Sh(X_{affine, \etale})$.
Ce topos est fonctoriel par rapport aux morphismes
d'espaces algébriques. Voici un énoncé précis.

\begin{lemma}
\label{spaces-properties-lemma-functoriality-etale-site}
Soit $S$ un schéma.
Soit $f : X \to Y$ un morphisme d'espaces algébriques sur $S$.
\begin{enumerate}
\item Le foncteur continu
$$
Y_{spaces, \etale} \longrightarrow X_{spaces, \etale}, \quad
V \longmapsto X \times_Y V
$$
induit un morphisme de sites
$$
f_{spaces, \etale} :
X_{spaces, \etale}
\to
Y_{spaces, \etale}.
$$
\item La règle $f \mapsto f_{spaces, \etale}$ est compatible aux
compositions, autrement dit $(f \circ g)_{spaces, \etale}
= f_{spaces, \etale} \circ g_{spaces, \etale}$ (voir
Sites, définition \ref{sites-definition-composition-morphisms-sites}).
\item Le morphisme de topos associé à $f_{spaces, \etale}$
induit, via le lemme \ref{spaces-properties-lemma-compare-etale-sites}, un morphisme de topos
$f_{small} : \Sh(X_\etale) \to \Sh(Y_\etale)$
dont la construction est compatible aux compositions.
\item Si $f$ est un morphisme représentable d'espaces algébriques,
alors $f_{small}$ provient d'un morphisme de sites
$X_\etale \to Y_\etale$,
correspondant au foncteur continu $V \mapsto X \times_Y V$.
\end{enumerate}
\end{lemma}

\begin{proof}
Montrons que le foncteur décrit en (1) satisfait aux hypothèses
de Sites, proposition \ref{sites-proposition-get-morphism}.
Il faut donc montrer que
$Y_{spaces, \etale}$ possède un objet final (à savoir $Y$) et que
le foncteur transforme celui-ci en un objet final de $X_{spaces, \etale}$
(à savoir $X$). C'est clair, car $X \times_Y Y = X$ dans toute catégorie.
Il faut ensuite montrer que $Y_{spaces, \etale}$ possède des produits fibrés.
C'est vrai puisque la catégorie des espaces algébriques possède des produits fibrés,
et puisque $V \times_Y V'$ est \'etale sur $Y$ si $V$ et $V'$ sont \'etales
sur $Y$ (voir les lemmes \ref{spaces-properties-lemma-composition-etale} et
\ref{spaces-properties-lemma-base-change-etale} ci-dessus).
La proposition s'applique donc et fournit un morphisme
de sites comme en (1).

\medskip\noindent
L'assertion (2) s'obtient en explicitant les définitions.
L'assertion (3) est claire en utilisant les équivalences pour $X$ et $Y$
du lemme \ref{spaces-properties-lemma-compare-etale-sites} ci-dessus.
L'assertion (4) en résulte car, si $f$ est représentable, les
foncteurs ci-dessus s'insèrent dans un diagramme commutatif
$$
\xymatrix{
X_\etale \ar[r] &
X_{spaces, \etale} \\
Y_\etale \ar[r] \ar[u] &
Y_{spaces, \etale} \ar[u]
}
$$
de catégories.
\end{proof}

\noindent
On peut préciser quelque peu le lemme précédent lorsqu'on décrit
la relation entre les faisceaux sur $X$ et ceux sur $Y$.
Plus précisément, on peut la formuler en termes de $f$-applications ; comparer
Faisceaux, définition \ref{sheaves-definition-f-map}, comme suit.

\begin{definition}
\label{spaces-properties-definition-f-map}
Soit $S$ un schéma.
Soit $f : X \to Y$ un morphisme d'espaces algébriques sur $S$.
Soit $\mathcal{F}$ un faisceau d'ensembles sur $X_\etale$ et
soit $\mathcal{G}$ un faisceau d'ensembles sur $Y_\etale$.
Une {\it $f$-application $\varphi : \mathcal{G} \to \mathcal{F}$}
est une collection d'applications
$\varphi_{(U, V, g)} : \mathcal{G}(V) \to \mathcal{F}(U)$
indexée par les diagrammes commutatifs
$$
\xymatrix{
U \ar[d]_g \ar[r] & X \ar[d]^f \\
V \ar[r] & Y
}
$$
où $U \in X_\etale$, $V \in Y_\etale$, telle que, pour tout
diagramme prolongé
$$
\xymatrix{
U' \ar[r] \ar[d]_{g'} & U \ar[d]_g \ar[r] & X \ar[d]^f \\
V' \ar[r] & V \ar[r] & Y
}
$$
où $V' \to V$ et $U' \to U$ sont des morphismes \'etales de schémas, le diagramme
$$
\xymatrix{
\mathcal{G}(V)
\ar[rr]_{\varphi_{(U, V, g)}}
\ar[d]_{\text{restriction de }\mathcal{G}} & &
\mathcal{F}(U)
\ar[d]^{\text{restriction de }\mathcal{F}} \\
\mathcal{G}(V')
\ar[rr]^{\varphi_{(U', V', g')}} & &
\mathcal{F}(U')
}
$$
est commutatif.
\end{definition}

\begin{lemma}
\label{spaces-properties-lemma-f-map}
Soit $S$ un schéma.
Soit $f : X \to Y$ un morphisme d'espaces algébriques sur $S$.
Soit $\mathcal{F}$ un faisceau d'ensembles sur $X_\etale$ et
soit $\mathcal{G}$ un faisceau d'ensembles sur $Y_\etale$.
Il existe des bijections canoniques entre les trois ensembles suivants :
\begin{enumerate}
\item l'ensemble des morphismes $\mathcal{G} \to f_{small, *}\mathcal{F}$ ;
\item l'ensemble des morphismes $f_{small}^{-1}\mathcal{G} \to \mathcal{F}$ ;
\item l'ensemble des $f$-applications $\varphi : \mathcal{G} \to \mathcal{F}$.
\end{enumerate}
\end{lemma}

\begin{proof}
Notons que (1) et (2) coïncident parce que les foncteurs $f_{small, *}$
et $f_{small}^{-1}$ forment un couple de foncteurs adjoints.
Supposons que $\alpha : f_{small}^{-1}\mathcal{G} \to \mathcal{F}$
soit un morphisme de faisceaux sur $X_\etale$. Considérons un diagramme
$$
\xymatrix{
U \ar[d]_g \ar[r]_{j_U} & X \ar[d]^f \\
V \ar[r]^{j_V} & Y
}
$$
comme dans la définition \ref{spaces-properties-definition-f-map}.
La commutativité du diagramme fournit également un morphisme
$g_{small}^{-1}(j_V)^{-1}\mathcal{G} \to (j_U)^{-1}\mathcal{F}$
(comparer Sites, section \ref{sites-section-localize}, pour la
description des foncteurs de localisation). Nous obtenons donc en particulier
un morphisme
$\varphi_{(U, V, g)} :
\mathcal{G}(V) = (j_V)^{-1}\mathcal{G}(V)
\to
(j_U)^{-1}\mathcal{F}(U) = \mathcal{F}(U)$.
Nous omettons de vérifier que cette règle est compatible aux
restrictions ultérieures et définit une $f$-application de $\mathcal{G}$ vers
$\mathcal{F}$.

\medskip\noindent
Réciproquement, supposons donnée une $f$-application
$\varphi = (\varphi_{(U, V, g)})$.
Désignons par $\mathcal{G}'$ (resp.\ $\mathcal{F}'$) le prolongement de
$\mathcal{G}$ (resp.\ $\mathcal{F}$) à $Y_{spaces, \etale}$
(resp.\ $X_{spaces, \etale}$), voir le
lemme \ref{spaces-properties-lemma-compare-etale-sites}.
Il faut alors construire un morphisme de faisceaux
$$
\mathcal{G}' \longrightarrow (f_{spaces, \etale})_*\mathcal{F}'
$$
Pour cela, soit $V \to Y$ un morphisme \'etale d'espaces algébriques.
Il faut construire une application d'ensembles
$$
\mathcal{G}'(V) \to \mathcal{F}'(X \times_Y V)
$$
Choisissons un morphisme \'etale surjectif $V' \to V$ avec $V'$ schéma,
puis un morphisme \'etale surjectif
$U' \to X \times_Y V'$ avec $U'$ schéma. Posons $U'' = U' \times_{X \times_Y V} U'$
et $V'' = V' \times_V V'$. Nous obtenons un morphisme de schémas $g' : U' \to V'$ ainsi qu'un morphisme de schémas
$$
g'' : U' \times_{X \times_Y V} U' \longrightarrow V' \times_V V'
$$
Considérons le diagramme suivant
$$
\xymatrix{
\mathcal{F}'(X \times_Y V) \ar[r] &
\mathcal{F}(U') \ar@<1ex>[r] \ar@<-1ex>[r] &
\mathcal{F}(U' \times_{X \times_Y V} U') \\
\mathcal{G}'(V) \ar[r] \ar@{..>}[u] &
\mathcal{G}(V') \ar@<1ex>[r] \ar@<-1ex>[r] \ar[u]_{\varphi_{(U', V', g')}} &
\mathcal{G}(V' \times_V V') \ar[u]_{\varphi_{(U'', V'', g'')}}
}
$$
La compatibilité des applications $\varphi_{...}$
aux restrictions montre que les deux carrés de droite sont commutatifs.
La définition des recouvrements de $X_{spaces, \etale}$ montre que
les lignes horizontales sont des diagrammes exacts. On obtient donc
la flèche pointillée. Nous laissons au lecteur le soin de montrer que ces
flèches sont compatibles aux applications de restriction.
\end{proof}

\noindent
Si le morphisme d'espaces algébriques $X \to Y$ est \'etale, alors le morphisme
de topos $\Sh(X_\etale) \to \Sh(Y_\etale)$
est un morphisme de localisation. Voici un énoncé précis.

\begin{lemma}
\label{spaces-properties-lemma-etale-morphism-topoi}
Soit $S$ un schéma, et soit $f : X \to Y$ un morphisme d'espaces algébriques
sur $S$. Supposons $f$ \'etale. Il existe alors un foncteur
$$
j : X_\etale \to Y_\etale, \quad
(\varphi : U \to X) \mapsto (f \circ \varphi : U \to Y)
$$
qui est cocontinu. Le morphisme de topos $f_{small}$ est le
morphisme de topos associé à $j$, voir
Sites, lemme \ref{sites-lemma-cocontinuous-morphism-topoi}.
De plus, $j$ est aussi continu ; par conséquent,
Sites, lemme \ref{sites-lemma-when-shriek},
s'applique. En particulier, $f_{small}^{-1}\mathcal{G}(U) = \mathcal{G}(jU)$
pour tout faisceau $\mathcal{G}$ sur $Y_\etale$.
\end{lemma}

\begin{proof}
Notons que notre définition même d'un morphisme \'etale d'espaces algébriques
(définition \ref{spaces-properties-definition-etale}) assure
que la règle donnée définit bien le foncteur $j$ indiqué.
Il est clair que $j$ est cocontinu et continu, simplement parce qu'un recouvrement
$\{U_i \to U\}$ de $j(\varphi : U \to X)$ dans $Y_\etale$ est
la même chose qu'un recouvrement de $(\varphi : U \to X)$ dans $X_\etale$. Il
reste à montrer que $j$ induit le même morphisme de topos que $f_{small}$.
Considérons pour cela le diagramme
$$
\xymatrix{
X_\etale \ar[r] \ar[d]^j &
X_{spaces, \etale} \ar@/_/[d]_{j_{spaces}} \\
Y_\etale \ar[r] &
Y_{spaces, \etale} \ar@/_/[u]_{v : V \mapsto X \times_Y V}
}
$$
de catégories. Ici, le foncteur $j_{spaces}$ est le prolongement évident de $j$
à la catégorie $X_{spaces, \etale}$. Le carré intérieur est donc
commutatif. En fait, $j_{spaces}$ s'identifie au
foncteur de localisation
$j_X : Y_{spaces, \etale}/X \to Y_{spaces, \etale}$
étudié dans
Sites, section \ref{sites-section-localize}.
Par conséquent, d'après
Sites, lemme \ref{sites-lemma-localize-given-products},
le foncteur cocontinu $j_{spaces}$ et le foncteur $v$ du diagramme
induisent le même morphisme de topos. D'après
Sites, lemme \ref{sites-lemma-composition-cocontinuous},
la commutativité du carré intérieur (formé de foncteurs cocontinus
entre sites) fournit un diagramme commutatif des morphismes de topos associés.
La construction de $f_{small}$ dans le
lemme \ref{spaces-properties-lemma-functoriality-etale-site} achève donc la preuve.
\end{proof}

\noindent
Le lemme précédent dit que l'image réciproque de $\mathcal{G}$ par un morphisme \'etale
$f : X \to Y$ d'espaces algébriques est simplement la restriction de $\mathcal{G}$
à la catégorie $X_\etale$. Nous emploierons souvent l'abréviation
\begin{equation}
\label{spaces-properties-equation-restrict}
\mathcal{G}|_{X_\etale} = f_{small}^{-1}\mathcal{G}
\end{equation}
pour le signaler. Notons que le foncteur
$j : X_\etale \to Y_\etale$
du lemme est alors fidèle, mais n'est pas pleinement fidèle en
général. Nous l'étudierons plus techniquement dans la
section \ref{spaces-properties-section-localize}.

\begin{lemma}
\label{spaces-properties-lemma-pushforward-etale-base-change}
Soit $S$ un schéma. Soit
$$
\xymatrix{
X' \ar[r] \ar[d]_{f'} & X \ar[d]^f \\
Y' \ar[r]^g & Y
}
$$
un carré cartésien d'espaces algébriques sur $S$. Soit
$\mathcal{F}$ un faisceau sur $X_\etale$. Si $g$ est \'etale, alors
\begin{enumerate}
\item $f'_{small, *}(\mathcal{F}|_{X'}) = (f_{small, *}\mathcal{F})|_{Y'}$
dans $\Sh(Y'_\etale)$\footnote{On a aussi
$(f')_{small}^{-1}(\mathcal{G}|_{Y'}) = (f_{small}^{-1}\mathcal{G})|_{X'}$
en raison de la commutativité du diagramme et de (\ref{spaces-properties-equation-restrict})}, et
\item si $\mathcal{F}$ est un faisceau abélien, alors
$R^if'_{small, *}(\mathcal{F}|_{X'}) = (R^if_{small, *}\mathcal{F})|_{Y'}$.
\end{enumerate}
\end{lemma}

\begin{proof}
Considérons le diagramme de foncteurs suivant
$$
\xymatrix{
X'_{spaces, \etale} \ar[r]_j &
X_{spaces, \etale} \\
Y'_{spaces, \etale} \ar[r]^j \ar[u]^{V' \mapsto V' \times_{Y'} X'} &
Y_{spaces, \etale} \ar[u]_{V \mapsto V \times_Y X}
}
$$
Les flèches horizontales sont des foncteurs de localisation et les flèches verticales induisent
des morphismes de sites. La dernière assertion de
Sites, lemme \ref{sites-lemma-localize-morphism},
donne donc (1). Pour obtenir (2), appliquons (1) à une résolution injective de $\mathcal{F}$
et utilisons que la restriction est exacte et préserve les injectifs (voir
Cohomologie sur les sites, lemme \ref{sites-cohomology-lemma-cohomology-of-open}).
\end{proof}

\noindent
Le lemme suivant dit que l'on peut considérer un faisceau sur le petit
site \'etale d'un espace algébrique comme une famille compatible de faisceaux
sur les petits sites \'etales des schémas \'etales sur cet espace. Notons bien
que tous les morphismes de comparaison $c_f$ du lemme sont des isomorphismes,
ce qui est compatible avec
Topologies, lemme \ref{topologies-lemma-characterize-sheaf-big-etale},
et avec le fait que tous les morphismes entre objets de $X_\etale$
sont \'etales.

\begin{lemma}
\label{spaces-properties-lemma-characterize-sheaf-small-etale}
Soit $S$ un schéma. Soit $X$ un espace algébrique sur $S$.
Un faisceau $\mathcal{F}$ sur $X_\etale$ équivaut aux données suivantes :
\begin{enumerate}
\item pour tout $U \in \Ob(X_\etale)$, un faisceau
$\mathcal{F}_U$ sur $U_\etale$ ;
\item pour tout $f : U' \to U$ dans $X_\etale$, un isomorphisme
$c_f : f_{small}^{-1}\mathcal{F}_U \to \mathcal{F}_{U'}$.
\end{enumerate}
Ces données sont soumises à la condition que, pour tous $f : U' \to U$
et $g : U'' \to U'$ dans $X_\etale$, la composée
$c_g \circ g_{small}^{-1} c_f$ soit égale à $c_{f \circ g}$.
\end{lemma}

\begin{proof}
On peut interpréter $g_{small}^{-1}$ comme dans le lemme \ref{spaces-properties-lemma-etale-morphism-topoi}.
Le lemme résulte alors d'un fait général sur les sites, voir
Sites, lemme \ref{sites-lemma-glue-sheaves-absolute}.
\end{proof}

\noindent
Soit $S$ un schéma. Soit $X$ un espace algébrique sur $S$.
Soit $X = U/R$ une présentation de $X$ provenant d'un morphisme
\'etale surjectif quelconque $\varphi : U \to X$, voir
Espaces, définition \ref{spaces-definition-presentation}.
On obtient en particulier un groupoïde $(U, R, s, t, c, e, i)$ tel que
$j = (t, s) : R \to U \times_S U$, voir
Groupoïdes, lemme \ref{groupoids-lemma-equivalence-groupoid}.

\begin{lemma}
\label{spaces-properties-lemma-descent-sheaf}
Avec $S$, $\varphi : U \to X$ et $(U, R, s, t, c, e, i)$ comme ci-dessus,
pour tout faisceau $\mathcal{F}$ sur $X_\etale$, le
faisceau\footnote{Dans ce lemme
et sa preuve, nous écrivons simplement $\varphi^{-1}$ au lieu de $\varphi_{small}^{-1}$,
et de même pour toutes les autres images réciproques.}
$\mathcal{G} = \varphi^{-1}\mathcal{F}$ est muni d'un
isomorphisme canonique
$$
\alpha :
t^{-1}\mathcal{G}
\longrightarrow
s^{-1}\mathcal{G}
$$
tel que le diagramme
$$
\xymatrix{
& \text{pr}_1^{-1}t^{-1}\mathcal{G} \ar[r]_-{\text{pr}_1^{-1}\alpha} &
\text{pr}_1^{-1}s^{-1}\mathcal{G} \ar@{=}[rd] & \\
\text{pr}_0^{-1}s^{-1}\mathcal{G} \ar@{=}[ru] & & &
c^{-1}s^{-1}\mathcal{G} \\
&
\text{pr}_0^{-1}t^{-1}\mathcal{G} \ar[lu]^{\text{pr}_0^{-1}\alpha} \ar@{=}[r] &
c^{-1}t^{-1}\mathcal{G} \ar[ru]_{c^{-1}\alpha}
}
$$
est commutatif. Le foncteur $\mathcal{F} \mapsto (\mathcal{G}, \alpha)$
définit une équivalence de catégories entre les faisceaux sur
$X_\etale$ et les couples $(\mathcal{G}, \alpha)$ ci-dessus.
\end{lemma}

\begin{proof}[Première preuve du lemme \ref{spaces-properties-lemma-descent-sheaf}]
Soit $\mathcal{C} = X_{spaces, \etale}$. D'après le
lemme \ref{spaces-properties-lemma-etale-morphism-topoi}
et sa preuve, on a $U_{spaces, \etale} = \mathcal{C}/U$
et le foncteur image réciproque $\varphi^{-1}$ est simplement le foncteur de restriction.
De plus, $\{U \to X\}$ est un recouvrement du site $\mathcal{C}$ et
$R = U \times_X U$. L'isomorphisme $\alpha$ est simplement l'identification
canonique
$$
\left(\mathcal{F}|_{\mathcal{C}/U}\right)|_{\mathcal{C}/U \times_X U}
=
\left(\mathcal{F}|_{\mathcal{C}/U}\right)|_{\mathcal{C}/U \times_X U}
$$
et la commutativité du diagramme est la condition de cocycle pour les données
de recollement. Ce lemme est donc un cas particulier du recollement des faisceaux, voir
Sites, section \ref{sites-section-glueing-sheaves}.
\end{proof}

\begin{proof}[Seconde preuve du lemme \ref{spaces-properties-lemma-descent-sheaf}]
L'existence de $\alpha$ vient de l'égalité
$\varphi \circ t = \varphi \circ s$ et de la fonctorialité de l'image réciproque
par rapport au morphisme, voir le
lemme \ref{spaces-properties-lemma-functoriality-etale-site}.
De la même manière, c'est-à-dire par fonctorialité de l'image réciproque, on voit
que l'isomorphisme $\alpha$ s'insère dans le diagramme commutatif.
La construction $\mathcal{F} \mapsto (\varphi^{-1}\mathcal{F}, \alpha)$
est manifestement fonctorielle en le faisceau $\mathcal{F}$.
On obtient ainsi le foncteur.

\medskip\noindent
Réciproquement, supposons que $(\mathcal{G}, \alpha)$ soit un couple.
Soit $V \to X$ un objet de $X_\etale$.
Alors le morphisme $V' = U \times_X V \to V$ est un morphisme \'etale surjectif
de schémas, et $\{V' \to V\}$ est donc un recouvrement \'etale
de $V$. Posons $\mathcal{G}' = (V' \to V)^{-1}\mathcal{G}$.
Puisque $R = U \times_X U$, avec $t = \text{pr}_0$
et $s = \text{pr}_1$, on a $V' \times_V V' = R \times_X V$,
les applications de projection $s', t' : V' \times_V V' \to V'$ étant obtenues par changement de base
à partir de $t$ et $s$. Ainsi, $\alpha$ s'induit en un isomorphisme
$\alpha' : (t')^{-1}\mathcal{G}' \to (s')^{-1}\mathcal{G}'$. Cela étant posé,
nous définissons simplement
$$
\xymatrix{
\mathcal{F}(V) \ar@{=}[r] &
\text{Noyau du couple}(\mathcal{G}(V') \ar@<1ex>[r] \ar@<-1ex>[r] &
\mathcal{G}(V' \times_V V')).
}
$$
Nous omettons de vérifier que cela définit un faisceau. Pour voir que
$\mathcal{G}(V) = \mathcal{F}(V)$ s'il existe un morphisme $V \to U$,
notons que le noyau du couple est alors
$H^0(\{V' \to V\}, \mathcal{G}) = \mathcal{G}(V)$.
\end{proof}







\section{Points du petit site \'etale}
\label{spaces-properties-section-points-small-etale-site}

\noindent
Cette section est l'analogue de
Cohomologie \'etale, section
\ref{etale-cohomology-section-stalks}.

\begin{definition}
\label{spaces-properties-definition-geometric-point}
Soit $S$ un schéma. Soit $X$ un espace algébrique sur $S$.
\begin{enumerate}
\item Un {\it point géométrique} de $X$ est un morphisme
$\overline{x} : \Spec(k) \to X$, où $k$ est un corps algébriquement
clos. Nous commettons souvent un abus de notation en
écrivant $\overline{x} = \Spec(k)$.
\item À tout point géométrique $\overline{x}$ correspond le point
« image » $x \in |X|$. On dit que $\overline{x}$ est un
{\it point géométrique au-dessus de $x$}.
\end{enumerate}
\end{definition}

\noindent
On peut prendre les fibres des faisceaux sur $X_\etale$
en des points géométriques exactement comme dans le cas
du petit site \'etale d'un schéma. Pour cela, nous
définissons comme suit la notion de voisinage \'etale.

\begin{definition}
\label{spaces-properties-definition-etale-neighbourhood}
Soit $S$ un schéma. Soit $X$ un espace algébrique sur $S$.
Soit $\overline{x}$ un point géométrique de $X$.
\begin{enumerate}
\item Un {\it voisinage \'etale} de $\overline{x}$
dans $X$ est un diagramme commutatif
$$
\xymatrix{
& U \ar[d]^\varphi \\
{\bar x} \ar[r]^{\bar x} \ar[ur]^{\bar u} & X
}
$$
où $\varphi$ est un morphisme \'etale d'espaces algébriques sur $S$.
Nous noterons $\varphi : (U, \overline{u}) \to (X, \overline{x})$
une telle situation.
\item Un {\it morphisme de voisinages \'etales}
$(U, \overline{u}) \to (U', \overline{u}')$
est un $X$-morphisme $h : U \to U'$
tel que $\overline{u}' = h \circ \overline{u}$.
\end{enumerate}
\end{definition}

\noindent
Notons que nous permettons à $U$ d'être un espace algébrique. Lorsque nous
prenons les fibres d'un faisceau sur $X_\etale$, nous devons nous restreindre aux $U$ qui
appartiennent à $X_\etale$ et ne considérons donc, dans ce cas,
que des schémas $U$. On peut aussi travailler avec le site
$X_{spaces, \etale}$ et considérer tous les voisinages \'etales. Cela
ne fait aucune différence, d'après la dernière assertion du lemme suivant.

\begin{lemma}
\label{spaces-properties-lemma-cofinal-etale}
Soit $S$ un schéma. Soit $X$ un espace algébrique sur $S$.
Soit $\overline{x}$ un point géométrique de $X$.
La catégorie des voisinages \'etales est cofiltrante. Plus précisément :
\begin{enumerate}
\item Soient $(U_i, \overline{u}_i)_{i = 1, 2}$ deux voisinages \'etales de
$\overline{x}$ dans $X$. Il existe alors un troisième voisinage \'etale
$(U, \overline{u})$ et des morphismes
$(U, \overline{u}) \to (U_i, \overline{u}_i)$, $i = 1, 2$.
\item Soient $h_1, h_2: (U, \overline{u}) \to (U', \overline{u}')$ deux
morphismes entre voisinages \'etales de $\overline{x}$. Il existe alors un
voisinage \'etale $(U'', \overline{u}'')$ et un morphisme
$h : (U'', \overline{u}'') \to (U, \overline{u})$
qui égalise $h_1$ et $h_2$, c'est-à-dire tel que
$h_1 \circ h = h_2 \circ h$.
\end{enumerate}
De plus, pour tout voisinage \'etale
$(U, \overline{u}) \to (X, \overline{x})$
il existe un morphisme de voisinages \'etales
$(U', \overline{u}') \to (U, \overline{u})$
où $U'$ est un schéma.
\end{lemma}

\begin{proof}
Pour (1), considérons le produit fibré $U = U_1 \times_X U_2$.
Il est \'etale sur $U_1$ et sur $U_2$, car les morphismes \'etales sont
stables par changement de base et par composition, voir les
lemmes \ref{spaces-properties-lemma-base-change-etale} et \ref{spaces-properties-lemma-composition-etale}.
Le morphisme $\overline{u} \to U$ défini par $(\overline{u}_1, \overline{u}_2)$
en fait un voisinage \'etale muni de morphismes vers
$U_1$ et $U_2$.

\medskip\noindent
Pour (2), définissons $U''$ comme le produit fibré
$$
\xymatrix{
U'' \ar[r] \ar[d] & U \ar[d]^{(h_1, h_2)} \\
U' \ar[r]^-\Delta & U' \times_X U'.
}
$$
Comme les composés de $\overline{u}$ et $\overline{u}'$ vers $X$ valent $\overline{x}$,
$\overline{u}'' = (\overline{u}, \overline{u}')$ est un point géométrique
de $U''$. En particulier, $U'' \not = \emptyset$.
De plus, comme $U'$ est \'etale sur $X$, le produit fibré
$U'\times_X U'$ l'est aussi (comme on l'a vu ci-dessus pour $U_1 \times_X U_2$).
La flèche verticale $(h_1, h_2)$ est donc \'etale d'après le
lemme \ref{spaces-properties-lemma-etale-permanence}.
Ainsi, $U''$ est \'etale sur $U'$ par changement de base, donc aussi
\'etale sur $X$ (puisque les composés de morphismes \'etales sont \'etales).
Par conséquent, $(U'', \overline{u}'')$ résout le problème posé en (2).

\medskip\noindent
Pour la dernière assertion, choisissons un morphisme \'etale surjectif
$U' \to U$, où $U'$ est un schéma. Alors
$U' \times_U \overline{u}$ est un schéma \'etale et surjectif sur
$\overline{u} = \Spec(k)$, où $k$ est algébriquement clos.
Il s'ensuit (voir
Morphismes, lemme \ref{morphisms-lemma-etale-over-field})
que $U' \times_U \overline{u} \to \overline{u}$ possède une section,
qui fournit le point $\overline{u}'$ voulu.
\end{proof}

\begin{lemma}
\label{spaces-properties-lemma-geometric-lift-to-usual}
Soit $S$ un schéma. Soit $X$ un espace algébrique sur $S$.
Soit $\overline{x} : \Spec(k) \to X$ un point géométrique de $X$
au-dessus de $x \in |X|$. Soit $\varphi : U \to X$ un morphisme \'etale
d'espaces algébriques et soit $u \in |U|$ tel que $\varphi(u) = x$.
Il existe alors un point géométrique
$\overline{u} : \Spec(k) \to U$ au-dessus de $u$ tel que
$\overline{x} = \varphi \circ \overline{u}$.
\end{lemma}

\begin{proof}
Choisissons un schéma affine $U'$ muni de $u' \in U'$ et un morphisme \'etale
$U' \to U$ qui envoie $u'$ sur $u$. Il suffit de prouver le lemme pour
$(U', u') \to (X, x)$. Nous pouvons donc supposer que
$U$ est un schéma, et en particulier que $U \to X$ est représentable.
Considérons alors le diagramme cartésien
$$
\xymatrix{
\Spec(k) \times_{\overline{x}, X, \varphi} U
\ar[d]_{\text{pr}_1} \ar[r]_-{\text{pr}_2} & U
\ar[d]^\varphi \\
\Spec(k) \ar[r]^-{\overline{x}} & X
}
$$
La projection $\text{pr}_1$ est le changement de base d'un morphisme \'etale ;
elle est donc \'etale, voir le
lemme \ref{spaces-properties-lemma-base-change-etale}.
Par conséquent, le schéma $\Spec(k) \times_{\overline{x}, X, \varphi} U$
est une somme disjointe de spectres d'extensions finies séparables de $k$, voir
Morphismes, lemme \ref{morphisms-lemma-etale-over-field}.
Mais $k$ est algébriquement clos, donc toutes ces extensions sont triviales.
Ainsi, $\Spec(k) \times_{\overline{x}, X, \varphi} U$
est une somme disjointe de copies de $\Spec(k)$, et chacune
correspond à un point géométrique $\overline{u}$ tel que
$\varphi \circ \overline{u} = \overline{x}$. D'après le
lemme \ref{spaces-properties-lemma-points-cartesian},
l'application
$$
|\Spec(k) \times_{\overline{x}, X, \varphi} U|
\longrightarrow
|\Spec(k)| \times_{|X|} |U|
$$
est surjective ; nous pouvons donc choisir $\overline{u}$ au-dessus de $u$.
\end{proof}

\begin{lemma}
\label{spaces-properties-lemma-geometric-lift-to-cover}
Soit $S$ un schéma. Soit $X$ un espace algébrique sur $S$.
Soit $\overline{x}$ un point géométrique de $X$.
Soit $(U, \overline{u})$ un voisinage \'etale de $\overline{x}$.
Soit $\{\varphi_i : U_i \to U\}_{i \in I}$ un recouvrement \'etale dans
$X_{spaces, \etale}$.
Il existe alors $i \in I$ et $\overline{u}_i : \overline{x} \to U_i$
tels que $\varphi_i : (U_i, \overline{u}_i) \to (U, \overline{u})$
soit un morphisme de voisinages \'etales.
\end{lemma}

\begin{proof}
Soit $u \in |U|$ l'image de $\overline{u}$.
Comme $|U| = \bigcup_{i \in I} \varphi_i(|U_i|)$, il existe
$i$ et un point $u_i \in U_i$ envoyé sur $u$. Appliquons le
lemme \ref{spaces-properties-lemma-geometric-lift-to-usual}
à $(U_i, u_i) \to (U, u)$ et à $\overline{u}$ pour
obtenir le point géométrique voulu.
\end{proof}

\begin{definition}
\label{spaces-properties-definition-stalk}
Soit $S$ un schéma. Soit $X$ un espace algébrique sur $S$.
Soit $\mathcal{F}$ un préfaisceau sur $X_\etale$.
Soit $\overline{x}$ un point géométrique de $X$.
La {\it fibre} de $\mathcal{F}$ en $\overline{x}$ est
$$
\mathcal{F}_{\bar x}
=
\colim_{(U, \overline{u})} \mathcal{F}(U)
$$
où $(U, \overline{u})$ parcourt tous les voisinages \'etales
de $\overline{x}$ dans $X$ tels que $U \in \Ob(X_\etale)$.
\end{definition}

\noindent
D'après le
lemme \ref{spaces-properties-lemma-cofinal-etale},
cette limite inductive est indexée par une catégorie filtrante,
à savoir l'opposée de la catégorie des voisinages \'etales
dans $X_\etale$. Plus précisément, le
lemme \ref{spaces-properties-lemma-cofinal-etale}
dit que l'opposée de la catégorie de tous les voisinages \'etales est filtrante,
et que la sous-catégorie pleine de ceux qui appartiennent à $X_\etale$ est
cofinale, donc elle aussi filtrante.

\medskip\noindent
Ainsi, un élément de $\mathcal{F}_{\overline{x}}$ peut être
représenté par un triplet $(U, \overline{u}, \sigma)$, où
$U \in \Ob(X_\etale)$ et $\sigma \in \mathcal{F}(U)$.
Deux triplets $(U, \overline{u}, \sigma)$, $(U', \overline{u}', \sigma')$
définissent le même élément de la fibre s'il existe un troisième
voisinage \'etale
$(U'', \overline{u}'')$, $U'' \in \Ob(X_\etale)$
et des morphismes de voisinages \'etales
$h : (U'', \overline{u}'') \to (U, \overline{u})$,
$h' : (U'', \overline{u}'') \to (U', \overline{u}')$ tels que
$h^*\sigma = (h')^*\sigma'$ dans $\mathcal{F}(U'')$. Voir
Catégories, section \ref{categories-section-directed-colimits}.

\medskip\noindent
Cela implique aussi que, si $\mathcal{F}'$ est le faisceau sur
$X_{spaces, \etale}$ correspondant à $\mathcal{F}$ sur
$X_\etale$, alors
\begin{equation}
\label{spaces-properties-equation-stalk-spaces-etale}
\mathcal{F}_{\overline{x}} = \colim_{(U, \overline{u})} \mathcal{F}'(U)
\end{equation}
où la limite inductive porte maintenant sur tous les voisinages \'etales de $\overline{x}$.
Nous passerons souvent sans autre mention du point de vue de $X_\etale$
à celui de $X_{spaces, \etale}$, et réciproquement.

\medskip\noindent
En particulier, si $\mathcal{F}$ est un préfaisceau de
groupes abéliens, d'anneaux, etc., alors $\mathcal{F}_{\overline{x}}$ est
un groupe abélien, un anneau, etc., par la construction usuelle de la
structure algébrique sur une limite inductive filtrante de tels objets.

\begin{lemma}
\label{spaces-properties-lemma-stalk-gives-point}
\begin{slogan}
Un point géométrique d'un espace algébrique donne un point de son topos \'etale.
\end{slogan}
Soit $S$ un schéma.
Soit $X$ un espace algébrique sur $S$.
Soit $\overline{x}$ un point géométrique de $X$.
Considérons le foncteur
$$
u : X_\etale \longrightarrow \textit{Ens}, \quad
U \longmapsto |U_{\overline{x}}|
$$
Alors $u$ définit un point $p$ du site $X_\etale$
(Sites, définition \ref{sites-definition-point}),
et le foncteur fibre associé $\mathcal{F} \mapsto \mathcal{F}_p$
(Sites, équation \ref{sites-equation-stalk})
est le foncteur $\mathcal{F} \mapsto \mathcal{F}_{\overline{x}}$
défini ci-dessus.
\end{lemma}

\begin{proof}
Dans la preuve du
lemme \ref{spaces-properties-lemma-geometric-lift-to-usual},
nous avons vu que le schéma $U_{\overline{x}} = \overline{x} \times_X U$
est une somme disjointe de
schémas isomorphes à $\overline{x}$. On peut donc aussi voir
$|U_{\overline{x}}|$ comme l'ensemble des points géométriques de $U$
au-dessus de $\overline{x}$, c'est-à-dire comme l'ensemble des morphismes
$\overline{u} : \overline{x} \to U$ qui s'insèrent dans le diagramme de la
définition \ref{spaces-properties-definition-etale-neighbourhood}.
Il s'ensuit que $u(X)$ est un singleton et que
$u(U \times_V W) = u(U) \times_{u(V)} u(W)$
dès que $U \to V$ et $W \to V$ sont des morphismes dans $X_\etale$.
De plus, pour un recouvrement $\{U_i \to U\}_{i \in I}$ de $X_\etale$,
l'application $\coprod u(U_i) \to u(U)$ est surjective d'après le
lemme \ref{spaces-properties-lemma-geometric-lift-to-cover}.
La
proposition \ref{sites-proposition-point-limits} de Sites
s'applique donc, si bien que $p$ est un point du site $X_\etale$.
Enfin, notre foncteur $\mathcal{F} \mapsto \mathcal{F}_{\overline{x}}$
est donné par exactement la même limite inductive que le foncteur
$\mathcal{F} \mapsto \mathcal{F}_p$ associé à $p$ dans
Sites, équation \ref{sites-equation-stalk},
ce qui prouve la dernière assertion.
\end{proof}

\begin{lemma}
\label{spaces-properties-lemma-stalk-exact}
Soit $S$ un schéma.
Soit $X$ un espace algébrique sur $S$.
Soit $\overline{x}$ un point géométrique de $X$.
\begin{enumerate}
\item Le foncteur fibre
$\textit{PAb}(X_\etale) \to \textit{Ab}$,
$\mathcal{F}  \mapsto  \mathcal{F}_{\overline{x}}$
est exact.
\item On a $(\mathcal{F}^\#)_{\overline{x}} = \mathcal{F}_{\overline{x}}$
pour tout préfaisceau d'ensembles $\mathcal{F}$ sur $X_\etale$.
\item Le foncteur
$\textit{Ab}(X_\etale) \to \textit{Ab}$,
$\mathcal{F} \mapsto \mathcal{F}_{\overline{x}}$ est exact.
\item De même, les foncteurs
$\textit{PSh}(X_\etale) \to \textit{Ens}$ et
$\Sh(X_\etale) \to \textit{Ens}$ donnés par le foncteur fibre
$\mathcal{F} \mapsto \mathcal{F}_{\overline{x}}$ sont exacts (voir
Catégories, définition \ref{categories-definition-exact})
et commutent aux colimites quelconques.
\end{enumerate}
\end{lemma}

\begin{proof}
Ce résultat découle des résultats généraux de
Modules sur les sites, section \ref{sites-modules-section-stalks}.
En effet, $\mathcal{F} \mapsto \mathcal{F}_{\overline{x}}$
provient d'un point du petit site \'etale de $X$, voir le
lemme \ref{spaces-properties-lemma-stalk-gives-point}. Voir la preuve de
Cohomologie \'etale, lemme \ref{etale-cohomology-lemma-stalk-exact},
pour une preuve directe de certaines de ces assertions dans le cadre
du petit site \'etale d'un schéma.
\end{proof}

\noindent
Nous verrons plus bas que le foncteur fibre
$\mathcal{F} \mapsto \mathcal{F}_{\overline{x}}$
est en fait l'image réciproque par le morphisme $\overline{x}$. En ce sens,
le lemme suivant généralise le lemme précédent.

\begin{lemma}
\label{spaces-properties-lemma-stalk-pullback}
Soit $S$ un schéma.
Soit $f : X \to Y$ un morphisme d'espaces algébriques sur $S$.
\begin{enumerate}
\item Le foncteur
$f_{small}^{-1} :
\textit{Ab}(Y_\etale)
\to
\textit{Ab}(X_\etale)$
est exact.
\item Le foncteur
$f_{small}^{-1} :
\Sh(Y_\etale)
\to
\Sh(X_\etale)$
est exact, c'est-à-dire qu'il commute aux limites et colimites finies, voir
Catégories, définition \ref{categories-definition-exact}.
\item Pour tout morphisme \'etale $V \to Y$ d'espaces algébriques,
on a $f_{small}^{-1}h_V = h_{X \times_Y V}$.
\item Soit $\overline{x} \to X$ un point géométrique.
Soit $\mathcal{G}$ un faisceau sur $Y_\etale$.
Il existe alors une identification canonique
$$
(f_{small}^{-1}\mathcal{G})_{\overline{x}} = \mathcal{G}_{\overline{y}}.
$$
où $\overline{y} = f \circ \overline{x}$.
\end{enumerate}
\end{lemma}

\begin{proof}
Rappelons que $f_{small}$ est défini au moyen de $f_{spaces, small}$ dans le
lemme \ref{spaces-properties-lemma-functoriality-etale-site}.
Les assertions (1), (2) et (3) découlent en général du fait que
$f_{spaces, \etale} :
X_{spaces, \etale}
\to
Y_{spaces, \etale}$
est un morphisme de sites : voir
Sites, définition \ref{sites-definition-morphism-sites},
pour (2),
Modules sur les sites, lemme \ref{sites-modules-lemma-flat-pullback-exact},
pour (1), et
Sites, lemme \ref{sites-lemma-pullback-representable-sheaf},
pour (3).

\medskip\noindent
Preuve de (4). Cette assertion est un cas particulier de
Sites, lemme \ref{sites-lemma-point-morphism-sites},
via le
lemme \ref{spaces-properties-lemma-stalk-gives-point}.
Nous en donnons aussi une preuve directe. D'après le
lemme \ref{spaces-properties-lemma-stalk-exact},
le foncteur fibre commute au passage au faisceau associé.
Soit $\mathcal{G}'$ le faisceau sur $Y_{spaces, \etale}$ dont la restriction
à $Y_\etale$ est $\mathcal{G}$.
Rappelons que $f_{spaces, \etale}^{-1}\mathcal{G}'$ est le faisceau
associé au préfaisceau
$$
U \longrightarrow \colim_{U \to X \times_Y V} \mathcal{G}'(V),
$$
voir
Sites, sections \ref{sites-section-continuous-functors} et
\ref{sites-section-functoriality-PSh}.
On a donc
\begin{align*}
(f_{spaces, \etale}^{-1}\mathcal{G}')_{\overline{x}}
& =
\colim_{(U, \overline{u})} f_{spaces, \etale}^{-1}\mathcal{G}'(U)
\\
& = \colim_{(U, \overline{u})}
\colim_{a : U \to X \times_Y V} \mathcal{G}'(V) \\
& = \colim_{(V, \overline{v})} \mathcal{G}'(V) \\
& = \mathcal{G}'_{\overline{y}}
\end{align*}
où, dans la troisième égalité, le couple $(U, \overline{u})$ et le morphisme
$a : U \to X \times_Y V$ correspondent au couple $(V, a \circ \overline{u})$.
Comme la fibre de $\mathcal{G}'$
(resp.\ de $f_{spaces, \etale}^{-1}\mathcal{G}'$)
coïncide avec celle de $\mathcal{G}$ (resp.\ de $f_{small}^{-1}\mathcal{G}$),
voir
l'équation (\ref{spaces-properties-equation-stalk-spaces-etale}),
le résultat en découle.
\end{proof}

\begin{remark}
\label{spaces-properties-remark-stalk-pullback}
Cette remarque est l'analogue de
Cohomologie \'etale, remarque \ref{etale-cohomology-remark-stalk-pullback}.
Soit $S$ un schéma.
Soit $X$ un espace algébrique sur $S$.
Soit $\overline{x} : \Spec(k) \to X$ un point géométrique de $X$.
D'après
Cohomologie \'etale,
théorème \ref{etale-cohomology-theorem-equivalence-sheaves-point},
la catégorie des faisceaux sur $\Spec(k)_\etale$ est
équivalente à la catégorie des ensembles (en associant à un faisceau ses
sections globales). Il découle alors de l'assertion (4) du
lemme \ref{spaces-properties-lemma-stalk-pullback},
appliquée au morphisme $\overline{x}$, que le foncteur
$$
\Sh(X_\etale) \longrightarrow \textit{Ens}, \quad
\mathcal{F} \longmapsto \mathcal{F}_{\overline{x}}
$$
est isomorphe au foncteur
$$
\Sh(X_\etale)
\longrightarrow
\Sh(\Spec(k)_\etale) = \textit{Ens},
\quad
\mathcal{F} \longmapsto \overline{x}^*\mathcal{F}
$$
Nous pouvons donc considérer les foncteurs fibres comme des foncteurs
d'image réciproque par des morphismes géométriques (et non simplement par
des morphismes abstraits de topos, comme dans le résultat du
lemme \ref{spaces-properties-lemma-stalk-gives-point}).
\end{remark}

\begin{remark}
\label{spaces-properties-remark-map-stalks}
Soit $S$ un schéma.
Soit $X$ un espace algébrique sur $S$.
Soit $x \in |X|$.
Nous affirmons que, pour toute paire de points géométriques $\overline{x}$ et
$\overline{x}'$ au-dessus de $x$, les foncteurs fibres sont isomorphes.
Par définition de $|X|$, on peut trouver un troisième point géométrique
$\overline{x}''$ tel qu'il existe un diagramme commutatif
$$
\xymatrix{
\overline{x}'' \ar[r] \ar[d] \ar[rd]^{\overline{x}''} &
\overline{x}' \ar[d]^{\overline{x}'} \\
\overline{x} \ar[r]^{\overline{x}} & X.
}
$$
Comme le foncteur fibre $\mathcal{F} \mapsto \mathcal{F}_{\overline{x}}$
est donné par l'image réciproque par le morphisme $\overline{x}$ (et de même
pour les autres), on conclut par fonctorialité des images réciproques.
\end{remark}




\noindent
Le théorème suivant affirme que le petit site \'etale d'un espace
algébrique a assez de points.

\begin{theorem}
\label{spaces-properties-theorem-exactness-stalks}
Soit $S$ un schéma. Soit $X$ un espace algébrique sur $S$.
Un morphisme $a : \mathcal{F} \to \mathcal{G}$ de faisceaux d'ensembles est
injectif (resp.\ surjectif) si et seulement si le morphisme sur les fibres
$a_{\overline{x}} : \mathcal{F}_{\overline{x}} \to \mathcal{G}_{\overline{x}}$
est injectif (resp.\ surjectif) pour tout point géométrique de $X$.
Une suite de faisceaux abéliens sur $X_\etale$ est exacte
si et seulement si elle est exacte sur toutes les fibres aux points géométriques de $X$.
\end{theorem}

\begin{proof}
Nous savons que le théorème est vrai lorsque $X$ est un schéma, voir
Cohomologie \'etale, théorème \ref{etale-cohomology-theorem-exactness-stalks}.
Choisissons un morphisme \'etale surjectif $f : U \to X$, où $U$ est un schéma.
Comme $\{U \to X\}$ est un recouvrement (dans $X_{spaces, \etale}$), on peut
vérifier qu'un morphisme de faisceaux est injectif ou surjectif après
restriction à $U$. Or, si $\overline{u} : \Spec(k) \to U$ est un point
géométrique de $U$, alors
$(\mathcal{F}|_U)_{\overline{u}} = \mathcal{F}_{\overline{x}}$
où $\overline{x} = f \circ \overline{u}$. (Cela résulte directement des
limites inductives qui définissent les fibres en $\overline{u}$ et $\overline{x}$, mais
cela résulte aussi du
lemme \ref{spaces-properties-lemma-stalk-pullback}.)
Le résultat pour $U$ implique donc celui pour $X$, ce qui conclut.
\end{proof}

\noindent
Le lemme suivant peut être omis en première lecture.

\begin{lemma}
\label{spaces-properties-lemma-points-small-etale-site}
Soit $S$ un schéma.
Soit $X$ un espace algébrique sur $S$.
Soit $p : \Sh(pt) \to \Sh(X_\etale)$
un point du petit topos \'etale de $X$.
Il existe alors un point géométrique $\overline{x}$ de $X$
tel que le foncteur fibre $\mathcal{F} \mapsto \mathcal{F}_p$
soit isomorphe au foncteur fibre
$\mathcal{F} \mapsto \mathcal{F}_{\overline{x}}$.
\end{lemma}

\begin{proof}
D'après
Sites, lemme \ref{sites-lemma-point-site-topos},
il y a une correspondance biunivoque entre les points du site et ceux
du topos associé. Nous pouvons donc supposer que $p$ est donné par
un foncteur $u : X_\etale \to \textit{Ens}$ qui définit un point
du site $X_\etale$.
Soit $U \in \Ob(X_\etale)$ un objet dont le morphisme structural
$j : U \to X$ est surjectif. Notons que $h_U$ est un faisceau
muni d'une surjection sur le faisceau final. Comme le foncteur fibre est exact,
on voit que $(h_U)_p = u(U)$ n'est pas vide (utiliser
Sites, lemme \ref{sites-lemma-points-recover}).
Choisissons $x \in u(U)$. D'après
Sites, lemme \ref{sites-lemma-point-localize},
on obtient un point $q : \Sh(pt) \to \Sh(U_\etale)$
tel que $p = j_{small} \circ q$, et donc, fonctoriellement,
$\mathcal{F}_p = (\mathcal{F}|_U)_q$.
D'après
Cohomologie \'etale, lemme \ref{etale-cohomology-lemma-points-small-etale-site},
il existe un point géométrique $\overline{u}$ de $U$ et un isomorphisme
fonctoriel $\mathcal{G}_q = \mathcal{G}_{\overline{u}}$
pour $\mathcal{G} \in \Sh(U_\etale)$. Posons
$\overline{x} = j \circ \overline{u}$. On voit alors que
$\mathcal{F}_{\overline{x}} \cong (\mathcal{F}|_U)_{\overline{u}}$
fonctoriellement en $\mathcal{F}$ sur $X_\etale$, d'après le
lemme \ref{spaces-properties-lemma-stalk-pullback},
ce qui conclut.
\end{proof}





\section{Supports de faisceaux abéliens}
\label{spaces-properties-section-support}

\noindent
Commençons par les supports des sections locales.

\begin{lemma}
\label{spaces-properties-lemma-support-subsheaf-final}
Soit $S$ un schéma. Soit $X$ un espace algébrique sur $S$.
Soit $\mathcal{F}$ un sous-faisceau de l'objet final du topos \'etale
de $X$ (voir
Sites, exemple \ref{sites-example-singleton-sheaf}).
Il existe alors un unique sous-espace ouvert
$W \subset X$ tel que $\mathcal{F} = h_W$.
\end{lemma}

\begin{proof}
La condition signifie que $\mathcal{F}(U)$ est un singleton ou
est vide pour tout $\varphi : U \to X$ dans $\Ob(X_{spaces, \etale})$.
En particulier, les sections locales se recollent toujours. Si
$\mathcal{F}(U) \not = \emptyset$, alors
$\mathcal{F}(\varphi(U)) \not = \emptyset$, car
$\varphi(U) \subset X$ est un sous-espace ouvert
(lemme \ref{spaces-properties-lemma-etale-open})
et
$\{\varphi : U \to \varphi(U)\}$ est un recouvrement dans $X_{spaces, \etale}$.
Il suffit de prendre
$W = \bigcup_{\varphi : U \to X, \mathcal{F}(U) \not = \emptyset} \varphi(U)$
pour conclure.
\end{proof}

\begin{lemma}
\label{spaces-properties-lemma-zero-over-image}
Soit $S$ un schéma.
Soit $X$ un espace algébrique sur $S$.
Soit $\mathcal{F}$ un faisceau abélien sur $X_{spaces, \etale}$.
Soit $\sigma \in \mathcal{F}(U)$ une section locale.
Il existe un sous-espace ouvert $W \subset U$ tel que
\begin{enumerate}
\item $W \subset U$ soit le plus grand sous-espace ouvert de $U$
tel que $\sigma|_W = 0$,
\item pour tout $\varphi : V \to U$ dans $X_{spaces, \etale}$, on ait
$$
\sigma|_V = 0 \Leftrightarrow \varphi(V) \subset W,
$$
\item pour tout point géométrique $\overline{u}$ de $U$, on ait
$$
(U, \overline{u}, \sigma) = 0\text{ dans }\mathcal{F}_{\overline{x}}
\Leftrightarrow
\overline{u} \in W
$$
où $\overline{x} = (U \to X) \circ \overline{u}$.
\end{enumerate}
\end{lemma}

\begin{proof}
Comme $\mathcal{F}$ est un faisceau pour la topologie \'etale, la restriction de
$\mathcal{F}$ à $U_{Zar}$ est un faisceau sur $U$ pour la topologie de Zariski.
Il existe donc un ouvert de Zariski $W$ ayant la propriété (1), voir
Modules, lemme \ref{modules-lemma-support-section-closed}. Soit
$\varphi : V \to U$ une flèche de $X_{spaces, \etale}$. Notons que
$\varphi(V) \subset U$ est un sous-espace ouvert
(lemme \ref{spaces-properties-lemma-etale-open})
et que $\{V \to \varphi(V)\}$ est un recouvrement \'etale. Si donc
$\sigma|_V = 0$, la condition de faisceau de $\mathcal{F}$ montre que
$\sigma|_{\varphi(V)} = 0$. Cela prouve (2).
Pour prouver (3), il faut montrer que, si $(U, \overline{u}, \sigma)$
définit l'élément nul de $\mathcal{F}_{\overline{x}}$, alors
$\overline{u} \in W$. En effet, l'hypothèse signifie
qu'il existe un morphisme de voisinages \'etales
$(V, \overline{v}) \to (U, \overline{u})$ tel que
$\sigma|_V = 0$. D'après (2), le morphisme $V \to U$ se factorise par $W$,
et donc $\overline{u} \in W$.
\end{proof}

\noindent
Soit $S$ un schéma.
Soit $X$ un espace algébrique sur $S$.
Soit $x \in |X|$.
Soit $\mathcal{F}$ un faisceau sur $X_\etale$. D'après la
remarque \ref{spaces-properties-remark-map-stalks},
la classe d'isomorphisme de la fibre du faisceau $\mathcal{F}$
en un point géométrique au-dessus de $x$ est bien définie.

\begin{definition}
\label{spaces-properties-definition-support}
Soit $S$ un schéma.
Soit $X$ un espace algébrique sur $S$.
Soit $\mathcal{F}$ un faisceau abélien sur $X_\etale$.
\begin{enumerate}
\item Le {\it support de $\mathcal{F}$} est l'ensemble des
points $x \in |X|$ tels que $\mathcal{F}_{\overline{x}} \not = 0$
pour tout (ou, ce qui revient au même, pour un) point géométrique $\overline{x}$ au-dessus de $x$.
\item Soit $\sigma \in \mathcal{F}(U)$ une section.
Le {\it support de $\sigma$} est la partie fermée $U \setminus W$, où
$W \subset U$ est le plus grand ouvert de $U$ sur lequel la restriction de $\sigma$
est nulle (voir le
lemme \ref{spaces-properties-lemma-zero-over-image}).
\end{enumerate}
\end{definition}

\begin{lemma}
\label{spaces-properties-lemma-support-section-closed}
Soit $S$ un schéma.
Soit $X$ un espace algébrique sur $S$.
Soit $\mathcal{F}$ un faisceau abélien sur $X_\etale$.
Soient $U \in \Ob(X_\etale)$ et $\sigma \in \mathcal{F}(U)$.
\begin{enumerate}
\item Le support de $\sigma$ est fermé dans $|U|$.
\item Le support de $\sigma + \sigma'$ est contenu dans la réunion des
supports de $\sigma, \sigma' \in \mathcal{F}(U)$.
\item Si $\varphi : \mathcal{F} \to \mathcal{G}$ est un morphisme de
faisceaux abéliens sur $X_\etale$, alors le support de $\varphi(\sigma)$ est
contenu dans celui de $\sigma \in \mathcal{F}(U)$.
\item Le support de $\mathcal{F}$ est la réunion des images des
supports de toutes les sections locales de $\mathcal{F}$.
\item Si $\mathcal{F} \to \mathcal{G}$ est surjectif, alors le support
de $\mathcal{G}$ est contenu dans celui de $\mathcal{F}$.
\item Si $\mathcal{F} \to \mathcal{G}$ est injectif, alors le support
de $\mathcal{F}$ est contenu dans celui de $\mathcal{G}$.
\end{enumerate}
\end{lemma}

\begin{proof}
L'assertion (1) résulte de la définition.
Les assertions (2) et (3) sont vraies parce qu'elles le sont pour les
restrictions de $\mathcal{F}$ et $\mathcal{G}$ à $U_{Zar}$, voir
Modules, lemme \ref{modules-lemma-support-section-closed}.
L'assertion (4) découle directement de l'assertion (3) du
lemme \ref{spaces-properties-lemma-zero-over-image}.
Les assertions (5) et (6) découlent des autres.
\end{proof}

\begin{lemma}
\label{spaces-properties-lemma-support-sheaf-rings-closed}
Le support d'un faisceau d'anneaux sur le petit site \'etale d'un
espace algébrique est fermé.
\end{lemma}

\begin{proof}
Cela résulte du fait que (selon nos conventions)
un anneau est $0$ si et seulement si
$1 = 0$ ; le support d'un faisceau d'anneaux
est donc celui de sa section unité.
\end{proof}






\section{Le faisceau structural d'un espace algébrique}
\label{spaces-properties-section-structure-sheaf}

\noindent
Le faisceau structural d'un espace algébrique est le faisceau d'anneaux
du lemme suivant.

\begin{lemma}
\label{spaces-properties-lemma-sheaf-condition-holds}
Soit $S$ un schéma. Soit $X$ un espace algébrique sur $S$.
La règle $U \mapsto \Gamma(U, \mathcal{O}_U)$ définit
un faisceau d'anneaux sur $X_\etale$.
\end{lemma}

\begin{proof}
Cela résulte immédiatement de la définition d'un recouvrement et de
Descente, lemme \ref{descent-lemma-sheaf-condition-holds}.
\end{proof}

\begin{definition}
\label{spaces-properties-definition-structure-sheaf}
Soit $S$ un schéma.
Soit $X$ un espace algébrique sur $S$.
Le {\it faisceau structural} de $X$
est le faisceau d'anneaux $\mathcal{O}_X$
sur le petit site \'etale $X_\etale$ décrit dans le
lemme \ref{spaces-properties-lemma-sheaf-condition-holds}.
\end{definition}

\noindent
D'après le lemme \ref{spaces-properties-lemma-characterize-sheaf-small-etale}, le faisceau
$\mathcal{O}_X$ correspond à un système de faisceaux \'etales $(\mathcal{O}_X)_U$,
où $U$ parcourt les objets de $X_\etale$. La preuve de ce lemme et
notre définition montrent clairement que l'on a simplement
$(\mathcal{O}_X)_U = \mathcal{O}_U$, où $\mathcal{O}_U$ est le faisceau
structural de $U_\etale$ introduit dans
Descente, définition \ref{descent-definition-structure-sheaf}.
En particulier, si $X$ est un schéma, on retrouve le faisceau $\mathcal{O}_X$
sur le petit site \'etale de $X$.

\medskip\noindent
Au moyen de l'équivalence
$\Sh(X_\etale) = \Sh(X_{spaces, \etale})$
du lemme \ref{spaces-properties-lemma-compare-etale-sites}, on peut également considérer $\mathcal{O}_X$
comme un faisceau d'anneaux sur $X_{spaces, \etale}$. La
remarque \ref{spaces-properties-remark-explain-equivalence}
explique comment calculer $\mathcal{O}_X(Y)$, et en particulier $\mathcal{O}_X(X)$, lorsque
$Y \to X$ est un objet de $X_{spaces, \etale}$.

\begin{lemma}
\label{spaces-properties-lemma-morphism-ringed-topoi}
Soit $S$ un schéma.
Soit $f : X \to Y$ un morphisme d'espaces algébriques sur $S$.
Il existe alors un morphisme canonique
$f^\sharp : f_{small}^{-1}\mathcal{O}_Y \to \mathcal{O}_X$ tel que
$$
(f_{small}, f^\sharp) :
(\Sh(X_\etale), \mathcal{O}_X)
\longrightarrow
(\Sh(Y_\etale), \mathcal{O}_Y)
$$
soit un morphisme de topos annelés. De plus,
\begin{enumerate}
\item La construction $f \mapsto (f_{small}, f^\sharp)$ est compatible aux
compositions.
\item Si $f$ est un morphisme de schémas, alors $f^\sharp$ est le morphisme décrit dans
Descente, remarque \ref{descent-remark-change-topologies-ringed}.
\end{enumerate}
\end{lemma}

\begin{proof}
D'après le lemme \ref{spaces-properties-lemma-f-map}, il suffit de donner une $f$-application de
$\mathcal{O}_Y$ vers $\mathcal{O}_X$. Autrement dit, pour tout
diagramme commutatif
$$
\xymatrix{
U \ar[d]_g \ar[r] & X \ar[d]^f \\
V \ar[r] & Y
}
$$
où $U \in X_\etale$, $V \in Y_\etale$, il faut donner un
morphisme d'anneaux
$
(f^\sharp)_{(U, V, g)} :
\Gamma(V, \mathcal{O}_V)
\to
\Gamma(U, \mathcal{O}_U).
$
On prend bien sûr $(f^\sharp)_{(U, V, g)} = g^\sharp$.
Il est clair que ce choix est compatible aux morphismes de restriction
et donne donc bien une $f$-application.
Nous omettons la vérification de la compatibilité aux compositions et de l'accord avec la
construction de
Descente, remarque \ref{descent-remark-change-topologies-ringed}.
\end{proof}

\begin{lemma}
\label{spaces-properties-lemma-reduced-space}
Soit $S$ un schéma. Soit $X$ un espace algébrique sur $S$.
Les conditions suivantes sont équivalentes :
\begin{enumerate}
\item $X$ est réduit,
\item pour tout $x \in |X|$, l'anneau local de $X$ en $x$ est
réduit (remarque \ref{spaces-properties-remark-list-properties-local-ring-local-etale-topology}).
\end{enumerate}
Dans ce cas, $\Gamma(X, \mathcal{O}_X)$ est un anneau réduit et,
si $f \in \Gamma(X, \mathcal{O}_X)$ vérifie $X = V(f)$, alors $f = 0$.
\end{lemma}

\begin{proof}
L'équivalence de (1) et (2) résulte de
Propriétés, lemme \ref{properties-lemma-characterize-reduced},
appliqué aux schémas affines \'etales sur $X$. Les dernières assertions
résultent du lemme cité et du fait que $\Gamma(X, \mathcal{O}_X)$ est
un sous-anneau de $\Gamma(U, \mathcal{O}_U)$ pour un certain
schéma réduit $U$ \'etale sur $X$.
\end{proof}







\section{Fibres du faisceau structural}
\label{spaces-properties-section-stalks-structure-sheaf}

\noindent
Cette section est l'analogue de
Cohomologie \'etale, section
\ref{etale-cohomology-section-stalks-structure-sheaf}.

\begin{lemma}
\label{spaces-properties-lemma-describe-etale-local-ring}
Soit $S$ un schéma.
Soit $X$ un espace algébrique sur $S$.
Soit $\overline{x}$ un point géométrique de $X$.
Soit $(U, \overline{u})$ un voisinage \'etale de $\overline{x}$,
où $U$ est un schéma. On a alors
$$
\mathcal{O}_{X, \overline{x}} =
\mathcal{O}_{U, \overline{u}} =
\mathcal{O}_{U, u}^{sh}
$$
où le membre de gauche est la fibre du faisceau structural de $X$,
et celui de droite l'hensélisé strict de l'anneau local
de $U$ en la localité $u$ de $\overline{u}$.
\end{lemma}

\begin{proof}
On sait que le faisceau structural $\mathcal{O}_U$ sur
$U_\etale$ est la restriction du faisceau structural de $X$.
La première égalité résulte donc du
lemme \ref{spaces-properties-lemma-stalk-pullback}, partie (4).
La seconde égalité est expliquée dans
Cohomologie \'etale,
lemme \ref{etale-cohomology-lemma-describe-etale-local-ring}.
\end{proof}

\begin{definition}
\label{spaces-properties-definition-etale-local-rings}
Soit $S$ un schéma.
Soit $X$ un espace algébrique sur $S$.
Soit $\overline{x}$ un point géométrique de $X$ au-dessus du point
$x \in |X|$.
\begin{enumerate}
\item L'{\it anneau local \'etale de $X$ en $\overline{x}$}
est la fibre du faisceau structural $\mathcal{O}_X$ sur $X_\etale$
en $\overline{x}$.
Notation : $\mathcal{O}_{X, \overline{x}}$.
\item L'{\it hensélisé strict de $X$ en $\overline{x}$}
est le schéma $\Spec(\mathcal{O}_{X, \overline{x}})$.
\end{enumerate}
\end{definition}

\noindent
À isomorphisme près, l'hensélisé strict de $X$ en $\overline{x}$
(comme schéma sur $X$) ne dépend que du point $x \in |X|$, et non
du choix du point géométrique au-dessus de $x$ ; voir la
remarque \ref{spaces-properties-remark-map-stalks}.

\begin{lemma}
\label{spaces-properties-lemma-etale-site-locally-ringed}
Soit $S$ un schéma.
Soit $X$ un espace algébrique sur $S$.
Le petit site \'etale $X_\etale$, muni de son
faisceau structural $\mathcal{O}_X$, est un site localement annelé ; voir
Modules sur les sites, définition \ref{sites-modules-definition-locally-ringed}.
\end{lemma}

\begin{proof}
Cela résulte du fait que les fibres
$\mathcal{O}_{X, \overline{x}}$ sont
des anneaux locaux et que $X_\etale$ a assez de points ; voir le
lemme \ref{spaces-properties-lemma-describe-etale-local-ring} et le
théorème \ref{spaces-properties-theorem-exactness-stalks}.
Voir
Modules sur les sites, lemme \ref{sites-modules-lemma-locally-ringed-stalk} et
\ref{sites-modules-lemma-ringed-stalk-not-zero}
pour le fait qu'il en résulte que le petit site \'etale est localement annelé.
\end{proof}

\begin{lemma}
\label{spaces-properties-lemma-dimension-local-ring}
Soit $S$ un schéma. Soit $X$ un espace algébrique sur $S$.
Soit $x \in |X|$ un point. Soit $d \in \{0, 1, 2, \ldots, \infty\}$.
Les conditions suivantes sont équivalentes :
\begin{enumerate}
\item la dimension de l'anneau local de $X$ en $x$
(définition \ref{spaces-properties-definition-dimension-local-ring}) est $d$ ;
\item $\dim(\mathcal{O}_{X, \overline{x}}) = d$ pour un point
géométrique $\overline{x}$ au-dessus de $x$ ;
\item $\dim(\mathcal{O}_{X, \overline{x}}) = d$ pour tout point
géométrique $\overline{x}$ au-dessus de $x$.
\end{enumerate}
\end{lemma}

\begin{proof}
L'équivalence de (2) et (3) résulte du fait que le
type d'isomorphie de $\mathcal{O}_{X, \overline{x}}$ ne dépend que
de $x \in |X|$ ; voir la remarque \ref{spaces-properties-remark-map-stalks}.
En utilisant le lemme \ref{spaces-properties-lemma-describe-etale-local-ring},
l'équivalence de (1) et (2)$+$(3) se ramène à
l'assertion suivante : pour tout anneau local $R$, on a
$\dim(R) = \dim(R^{sh})$. C'est le contenu de
Compléments d'algèbre, lemme \ref{more-algebra-lemma-henselization-dimension}.
\end{proof}

\begin{lemma}
\label{spaces-properties-lemma-dimension-decent-invariant-under-etale}
Soit $S$ un schéma. Soit $f : X \to Y$ un morphisme \'etale
d'espaces algébriques sur $S$. Soit $x \in |X|$. Alors
(1) $\dim_x(X) = \dim_{f(x)}(Y)$ et (2) la dimension de
l'anneau local de $X$ en $x$ est égale à la dimension de
l'anneau local de $Y$ en $f(x)$. Si $f$ est surjectif, alors
(3) $\dim(X) = \dim(Y)$.
\end{lemma}

\begin{proof}
Choisissons un schéma $U$, un point $u \in U$ et un morphisme \'etale
$U \to X$ qui envoie $u$ sur $x$. Le composé $U \to Y$
est encore \'etale et envoie $u$ sur $f(x)$. Les assertions (1) et (2)
en résultent, car les nombres considérés sont définis par le comportement
du schéma $U$ en $u$. Voir la
définition \ref{spaces-properties-definition-dimension-at-point} pour (1). La partie (3) est
une conséquence immédiate de (1) ; voir la définition \ref{spaces-properties-definition-dimension}.
\end{proof}

\begin{lemma}
\label{spaces-properties-lemma-reduced-local-ring}
Soit $S$ un schéma. Soit $X$ un espace algébrique sur $S$.
Soit $x \in |X|$ un point. Les conditions suivantes sont équivalentes :
\begin{enumerate}
\item l'anneau local de $X$ en $x$ est réduit
(remarque \ref{spaces-properties-remark-list-properties-local-ring-local-etale-topology}) ;
\item $\mathcal{O}_{X, \overline{x}}$ est réduit pour un point
géométrique $\overline{x}$ au-dessus de $x$ ;
\item $\mathcal{O}_{X, \overline{x}}$ est réduit pour tout point
géométrique $\overline{x}$ au-dessus de $x$.
\end{enumerate}
\end{lemma}

\begin{proof}
L'équivalence de (2) et (3) résulte du fait que le
type d'isomorphie de $\mathcal{O}_{X, \overline{x}}$ ne dépend que
de $x \in |X|$ ; voir la remarque \ref{spaces-properties-remark-map-stalks}.
En utilisant le lemme \ref{spaces-properties-lemma-describe-etale-local-ring},
l'équivalence de (1) et (2)$+$(3) se ramène à
l'assertion suivante : un anneau local est réduit si et seulement si
son hensélisé strict est réduit. C'est le contenu de
Compléments d'algèbre, lemme \ref{more-algebra-lemma-henselization-reduced}.
\end{proof}











\section{Irréductibilité locale}
\label{spaces-properties-section-irreducible-local-ring}

\noindent
Tout point d'un espace algébrique possède un anneau local \'etale bien défini, qui
correspond au hensélisé strict de l'anneau local dans le cas d'un
schéma. En général, il ne permet pas de voir combien de composantes
irréductibles d'un schéma ou d'un espace algébrique passent par le point donné :
l'anneau local \'etale ne permet de compter que le nombre de branches géométriques.

\begin{lemma}
\label{spaces-properties-lemma-irreducible-local-ring}
Soit $S$ un schéma.
Soit $X$ un espace algébrique sur $S$.
Soit $x \in |X|$ un point.
Les conditions suivantes sont équivalentes :
\begin{enumerate}
\item pour tout schéma $U$, tout morphisme \'etale $a : U \to X$ et
tout $u \in U$ tel que $a(u) = x$, l'anneau local $\mathcal{O}_{U, u}$ possède un
unique idéal premier minimal ;
\item pour tout schéma $U$, tout morphisme \'etale $a : U \to X$ et
tout $u \in U$ tel que $a(u) = x$, il existe une unique composante irréductible de $U$
passant par $u$ ;
\item pour tout schéma $U$, tout morphisme \'etale $a : U \to X$ et
tout $u \in U$ tel que $a(u) = x$, l'anneau local $\mathcal{O}_{U, u}$
est unibranche ;
\item pour tout schéma $U$, tout morphisme \'etale $a : U \to X$ et
tout $u \in U$ tel que $a(u) = x$, l'anneau local $\mathcal{O}_{U, u}$
est géométriquement unibranche ;
\item $\mathcal{O}_{X, \overline{x}}$ possède un unique idéal premier minimal
pour tout point géométrique $\overline{x}$ au-dessus de $x$.
\end{enumerate}
\end{lemma}

\begin{proof}
L'équivalence de (1) et (2) résulte du fait que les composantes
irréductibles de $U$ passant par $u$ sont en correspondance biunivoque ($1$-$1$) avec les
idéaux premiers minimaux de l'anneau local de $U$ en $u$. Soient $a : U \to X$ et
$u \in U$ comme dans (1). Alors $\mathcal{O}_{X, \overline{x}}$ est
l'hensélisé strict de $\mathcal{O}_{U, u}$ d'après le
lemme \ref{spaces-properties-lemma-describe-etale-local-ring}.
En particulier, (4) et (5) sont équivalentes d'après
Compléments d'algèbre, lemme \ref{more-algebra-lemma-geometrically-unibranch}.
L'équivalence de (2), (3) et (4) résulte de
Compléments sur les morphismes, lemme \ref{more-morphisms-lemma-nr-branches}.
\end{proof}

\begin{definition}
\label{spaces-properties-definition-unibranch}
Soit $S$ un schéma. Soit $X$ un espace algébrique sur $S$.
Soit $x \in |X|$. On dit que $X$ est {\it géométriquement unibranche
en $x$} si les conditions équivalentes du
lemme \ref{spaces-properties-lemma-irreducible-local-ring}
sont satisfaites. On dit que $X$ est {\it géométriquement unibranche} si $X$ l'est
en tout $x \in |X|$.
\end{definition}

\noindent
Cette définition est compatible avec celle donnée pour les schémas
(Propriétés, définition \ref{properties-definition-unibranch}).

\begin{lemma}
\label{spaces-properties-lemma-nr-branches-local-ring}
Soit $S$ un schéma. Soit $X$ un espace algébrique sur $S$.
Soit $x \in |X|$ un point. Soit $n \in \{1, 2, \ldots\}$ un entier.
Les conditions suivantes sont équivalentes :
\begin{enumerate}
\item pour tout schéma $U$, tout morphisme \'etale $a : U \to X$ et
tout $u \in U$ tel que $a(u) = x$, le nombre d'idéaux premiers minimaux
de l'anneau local $\mathcal{O}_{U, u}$ est $\leq n$,
et, pour au moins un choix de $U, a, u$, il vaut $n$ ;
\item pour tout schéma $U$, tout morphisme \'etale $a : U \to X$ et
tout $u \in U$ tel que $a(u) = x$, le nombre de composantes irréductibles de
$U$ passant par $u$ est $\leq n$ et, pour au moins un choix
de $U, a, u$, il vaut $n$ ;
\item pour tout schéma $U$, tout morphisme \'etale $a : U \to X$ et tout $u \in U$
tel que $a(u) = x$, le nombre de branches de $U$ en $u$ est $\leq n$,
et, pour au moins un choix de $U, a, u$, il vaut $n$ ;
\item pour tout schéma $U$, tout morphisme \'etale $a : U \to X$ et tout $u \in U$
tel que $a(u) = x$, le nombre de branches géométriques de $U$ en $u$ est $n$ ;
\item le nombre d'idéaux premiers minimaux de
$\mathcal{O}_{X, \overline{x}}$ est $n$.
\end{enumerate}
\end{lemma}

\begin{proof}
L'équivalence de (1) et (2) résulte du fait que les composantes
irréductibles de $U$ passant par $u$ sont en correspondance biunivoque ($1$-$1$) avec les
idéaux premiers minimaux de l'anneau local de $U$ en $u$. Soient $a : U \to X$ et
$u \in U$ comme dans (1). Alors $\mathcal{O}_{X, \overline{x}}$ est
l'hensélisé strict de $\mathcal{O}_{U, u}$ d'après le
lemme \ref{spaces-properties-lemma-describe-etale-local-ring}. Rappelons que
le nombre (géométrique) de branches de $U$ en $u$ est le nombre
d'idéaux premiers minimaux de l'hensélisé (strict) de $\mathcal{O}_{U, u}$.
En particulier, (4) et (5) sont équivalentes.
L'équivalence de (2), (3) et (4) résulte de
Compléments sur les morphismes, lemme \ref{more-morphisms-lemma-nr-branches}.
\end{proof}

\begin{definition}
\label{spaces-properties-definition-number-of-geometric-branches}
Soit $S$ un schéma.
Soit $X$ un espace algébrique sur $S$.
Soit $x \in |X|$. Le {\it nombre de branches géométriques de $X$ en $x$} est
égal à $n \in \mathbf{N}$ si les conditions équivalentes du
lemme \ref{spaces-properties-lemma-nr-branches-local-ring}
sont satisfaites, et à $\infty$ sinon.
\end{definition}








\section{Espaces noethériens}
\label{spaces-properties-section-noetherian}

\noindent
Nous avons déjà défini les espaces algébriques localement noethériens dans la
section \ref{spaces-properties-section-types-properties}.

\begin{definition}
\label{spaces-properties-definition-noetherian}
Soit $S$ un schéma. Soit $X$ un espace algébrique sur $S$.
On dit que $X$ est {\it noethérien} si $X$ est quasi-compact, quasi-séparé
et localement noethérien.
\end{definition}

\noindent
Notons qu'un espace algébrique noethérien $X$ n'est pas seulement quasi-compact
et localement noethérien, mais aussi quasi-séparé. Cela ne contredit pas
la définition d'un schéma noethérien, car un schéma localement noethérien
est quasi-séparé ; voir
Propriétés, lemme \ref{properties-lemma-locally-Noetherian-quasi-separated}.
Cette propriété n'est pas vraie pour les espaces algébriques. En effet,
$X = \mathbf{A}^1_k/\mathbf{Z}$, voir
Espaces, exemple \ref{spaces-example-affine-line-translation},
est localement noethérien et quasi-compact, mais non quasi-séparé
(donc non noethérien selon nos définitions).

\medskip\noindent
Une conséquence du choix fait ci-dessus est qu'un espace algébrique
de type fini sur un espace algébrique noethérien n'est pas automatiquement
noethérien ; autrement dit, l'analogue de
Morphismes, lemme \ref{morphisms-lemma-finite-type-noetherian},
n'est pas vrai. L'énoncé correct est qu'un espace algébrique de
présentation finie sur un espace algébrique noethérien est noethérien
(voir
Morphismes d'espaces,
lemme \ref{spaces-morphisms-lemma-finite-presentation-noetherian}).

\medskip\noindent
Un espace algébrique noethérien $X$ est très proche d'être un schéma.
Dans la suite de cette section, nous rassemblons quelques lemmes qui l'illustrent.

\begin{lemma}
\label{spaces-properties-lemma-Noetherian-topology}
Soit $S$ un schéma. Soit $X$ un espace algébrique sur $S$.
\begin{enumerate}
\item Si $X$ est localement noethérien, alors $|X|$ est un espace topologique
localement noethérien.
\item Si $X$ est quasi-compact et localement noethérien, alors $|X|$
est un espace topologique noethérien.
\end{enumerate}
\end{lemma}

\begin{proof}
Supposons $X$ localement noethérien.
Choisissons un schéma $U$ et un morphisme \'etale surjectif
$U \to X$. Puisque $X$ est localement noethérien, $U$ l'est aussi.
D'après
Propriétés, lemme \ref{properties-lemma-Noetherian-topology},
cela signifie que $|U|$ est un espace topologique localement noethérien.
Comme $|U| \to |X|$ est ouverte et surjective, on en conclut que
$|X|$ est localement noethérien d'après
Topologie, lemme \ref{topology-lemma-image-Noetherian}.
Cela prouve (1). Si $X$ est quasi-compact et localement noethérien,
alors $|X|$ est quasi-compact et localement noethérien. Ainsi, $|X|$
est noethérien d'après
Topologie,
lemme \ref{topology-lemma-quasi-compact-locally-Noetherian-Noetherian}.
\end{proof}

\begin{lemma}
\label{spaces-properties-lemma-Noetherian-sober}
Soit $S$ un schéma. Soit $X$ un espace algébrique sur $S$.
Si $X$ est noethérien, alors $|X|$ est un espace topologique noethérien sobre.
\end{lemma}

\begin{proof}
L'espace topologique sous-jacent à un espace algébrique quasi-séparé
est sobre ; voir le
lemme \ref{spaces-properties-lemma-quasi-separated-sober}.
Il est noethérien d'après le
lemme \ref{spaces-properties-lemma-Noetherian-topology}.
\end{proof}

\begin{lemma}
\label{spaces-properties-lemma-Noetherian-local-ring-Noetherian}
Soit $S$ un schéma. Soit $X$ un espace algébrique noethérien sur $S$.
Soit $\overline{x}$ un point géométrique de $X$. Alors
$\mathcal{O}_{X, \overline{x}}$ est un anneau local noethérien.
\end{lemma}

\begin{proof}
Choisissons un voisinage \'etale $(U, \overline{u})$ de $\overline{x}$,
où $U$ est un schéma. Alors $\mathcal{O}_{X, \overline{x}}$ est
l'hensélisé strict de l'anneau local de $U$ en $u$ ; voir le
lemme \ref{spaces-properties-lemma-describe-etale-local-ring}.
D'après notre définition des espaces noethériens, le schéma $U$ est localement noethérien.
On conclut donc par
Compléments d'algèbre, lemme \ref{more-algebra-lemma-henselization-noetherian}.
\end{proof}








\section{Espaces algébriques réguliers}
\label{spaces-properties-section-regular}

\noindent
Nous avons déjà défini les espaces algébriques réguliers à la
section \ref{spaces-properties-section-types-properties}.

\begin{lemma}
\label{spaces-properties-lemma-regular}
Soit $S$ un schéma.
Soit $X$ un espace algébrique localement noethérien sur $S$.
Les conditions suivantes sont équivalentes :
\begin{enumerate}
\item $X$ est régulier ;
\item tout anneau local \'etale $\mathcal{O}_{X, \overline{x}}$ est
régulier.
\end{enumerate}
\end{lemma}

\begin{proof}
Soit $U$ un schéma et soit $U \to X$ un morphisme \'etale surjectif.
Par hypothèse, $U$ est localement noethérien. De plus, tout anneau local \'etale
$\mathcal{O}_{X, \overline{x}}$ est le hensélisé strict d'un
anneau local de $U$, et réciproquement ; voir le
lemme \ref{spaces-properties-lemma-describe-etale-local-ring}.
Ainsi, d'après le
lemme \ref{more-algebra-lemma-henselization-regular} de Compléments d'algèbre,
nous voyons que (2) équivaut à ce que tout anneau local de $U$ soit
régulier, c'est-à-dire à ce que $U$ soit un schéma régulier (voir le
lemme \ref{properties-lemma-characterize-regular} de Propriétés).
Cela équivaut à (1) d'après la
définition \ref{spaces-properties-definition-type-property}.
\end{proof}

\noindent
Nous pouvons utiliser le lemme \ref{descent-lemma-regular-local-ring-local}
de Descente pour définir ce que signifie, pour un espace algébrique $X$, être
régulier en un point $x$.

\begin{definition}
\label{spaces-properties-definition-regular-at-point}
Soit $S$ un schéma. Soit $X$ un espace algébrique sur $S$.
Soit $x \in |X|$ un point. Nous disons que {\it $X$ est régulier en $x$}
si $\mathcal{O}_{U, u}$ est un anneau local régulier pour tout
(ou, de manière équivalente, pour un certain) couple $(a : U \to X, u)$ formé d'un
morphisme \'etale $a : U \to X$ d'un schéma vers $X$ et d'un point
$u \in U$ tel que $a(u) = x$.
\end{definition}

\noindent
Voir la définition \ref{spaces-properties-definition-property-at-point}, le
lemme \ref{spaces-properties-lemma-local-source-target-at-point} et le
lemme \ref{descent-lemma-regular-local-ring-local} de Descente.

\begin{lemma}
\label{spaces-properties-lemma-regular-at-x}
Soit $S$ un schéma. Soit $X$ un espace algébrique sur $S$.
Soit $x \in |X|$ un point. Les conditions suivantes sont équivalentes :
\begin{enumerate}
\item $X$ est régulier en $x$ ;
\item l'anneau local \'etale $\mathcal{O}_{X, \overline{x}}$ est
régulier pour tout point géométrique $\overline{x}$ au-dessus du point
$x$ (ou, de manière équivalente, pour un certain).
\end{enumerate}
\end{lemma}

\begin{proof}
Soient $U$ un schéma, $u \in U$ un point, et $a : U \to X$ un
morphisme \'etale envoyant $u$ sur $x$. Pour tout point géométrique
$\overline{x}$ de $X$ au-dessus de $x$, l'anneau local \'etale
$\mathcal{O}_{X, \overline{x}}$ est le hensélisé strict de
l'anneau local de $U$ en $u$ ; voir le
lemme \ref{spaces-properties-lemma-describe-etale-local-ring}.
On conclut donc par le
lemme \ref{more-algebra-lemma-henselization-regular} de Compléments d'algèbre.
\end{proof}

\begin{lemma}
\label{spaces-properties-lemma-regular-normal}
Tout espace algébrique régulier est normal.
\end{lemma}

\begin{proof}
Cela résulte des définitions et du cas des schémas.
Voir le lemme \ref{properties-lemma-regular-normal} de Propriétés.
\end{proof}

 








\section{Faisceaux de modules sur les espaces algébriques}
\label{spaces-properties-section-modules}

\noindent
Si $X$ est un espace algébrique, un faisceau de modules sur $X$ est
un faisceau de $\mathcal{O}_X$-modules sur le petit site \'etale de $X$,
où $\mathcal{O}_X$ est le faisceau structural de $X$. La catégorie
des faisceaux de modules est notée $\textit{Mod}(\mathcal{O}_X)$.

\medskip\noindent
Étant donné un morphisme $f : X \to Y$ d'espaces algébriques, le
lemme \ref{spaces-properties-lemma-morphism-ringed-topoi}
nous donne un morphisme de topos annelés et, par conséquent, la
définition \ref{sites-modules-definition-pushforward} de Modules sur les sites
nous donne des foncteurs image inverse et image directe bien définis
\begin{equation}
\label{spaces-properties-equation-push-pull}
f^* :
\textit{Mod}(\mathcal{O}_Y)
\longrightarrow
\textit{Mod}(\mathcal{O}_X), \quad
f_* :
\textit{Mod}(\mathcal{O}_X)
\longrightarrow
\textit{Mod}(\mathcal{O}_Y)
\end{equation}
qui sont adjoints de la manière habituelle. Si $g : Y \to Z$ est un autre morphisme
d'espaces algébriques sur $S$, alors nous avons
$(g \circ f)^* = f^* \circ g^*$ et $(g \circ f)_* = g_* \circ f_*$
simplement parce que les morphismes de topos annelés se composent de la manière
correspondante (d'après le lemme).

\begin{lemma}
\label{spaces-properties-lemma-etale-exact-pullback}
Soit $S$ un schéma.
Soit $f : X \to Y$ un morphisme \'etale d'espaces algébriques sur $S$.
Alors $f^{-1}\mathcal{O}_Y = \mathcal{O}_X$, et
$f^*\mathcal{G} = f_{small}^{-1}\mathcal{G}$ pour tout faisceau de
$\mathcal{O}_Y$-modules $\mathcal{G}$. En particulier,
$f^* : \textit{Mod}(\mathcal{O}_Y) \to \textit{Mod}(\mathcal{O}_X)$
est exact.
\end{lemma}

\begin{proof}
D'après la description de l'image réciproque dans le lemme \ref{spaces-properties-lemma-etale-morphism-topoi}
et la définition des faisceaux structuraux, il est clair que
$f_{small}^{-1}\mathcal{O}_Y = \mathcal{O}_X$. Puisque, par définition, l'image inverse est
$$
f^*\mathcal{G} =
f_{small}^{-1}\mathcal{G} \otimes_{f_{small}^{-1}\mathcal{O}_Y}
\mathcal{O}_X
$$
nous concluons que $f^*\mathcal{G} = f_{small}^{-1}\mathcal{G}$.
L'exactitude est claire, car $f_{small}^{-1}$ est exact puisque $f_{small}$
est un morphisme de topos.
\end{proof}

\noindent
Nous poursuivons l'abus de notation introduit dans
l'équation (\ref{spaces-properties-equation-restrict})
en écrivant
\begin{equation}
\label{spaces-properties-equation-restrict-modules}
\mathcal{G}|_{X_\etale}
= f^*\mathcal{G}
= f_{small}^{-1}\mathcal{G}
\end{equation}
dans la situation du lemme ci-dessus. Nous en discuterons de manière plus
technique à la
section \ref{spaces-properties-section-localize}.

\begin{lemma}
\label{spaces-properties-lemma-pushforward-etale-base-change-modules}
Soit $S$ un schéma. Soit
$$
\xymatrix{
X' \ar[r] \ar[d]_{f'} & X \ar[d]^f \\
Y' \ar[r]^g & Y
}
$$
un carré cartésien d'espaces algébriques sur $S$. Soit
$\mathcal{F} \in \textit{Mod}(\mathcal{O}_X)$. Si $g$ est \'etale, alors
$f'_*(\mathcal{F}|_{X'}) = (f_*\mathcal{F})|_{Y'}$\footnote{Nous avons aussi
$(f')^*(\mathcal{G}|_{Y'}) = (f^*\mathcal{G})|_{X'}$
par commutativité du diagramme et par (\ref{spaces-properties-equation-restrict-modules})} et
$R^if'_*(\mathcal{F}|_{X'}) = (R^if_*\mathcal{F})|_{Y'}$ dans
$\textit{Mod}(\mathcal{O}_{Y'})$.
\end{lemma}

\begin{proof}
C'est une reformulation du
lemme \ref{spaces-properties-lemma-pushforward-etale-base-change}
dans le cas des modules.
\end{proof}

\begin{lemma}
\label{spaces-properties-lemma-characterize-module-small-etale}
Soit $S$ un schéma. Soit $X$ un espace algébrique sur $S$.
Un faisceau $\mathcal{F}$ de $\mathcal{O}_X$-modules est donné par les
données suivantes :
\begin{enumerate}
\item pour tout $U \in \Ob(X_\etale)$, un faisceau
$\mathcal{F}_U$ de $\mathcal{O}_U$-modules sur $U_\etale$ ;
\item pour tout $f : U' \to U$ dans $X_\etale$, un isomorphisme
$c_f : f_{small}^*\mathcal{F}_U \to \mathcal{F}_{U'}$.
\end{enumerate}
Ces données sont soumises à la condition suivante : pour tous $f : U' \to U$
et $g : U'' \to U'$ dans $X_\etale$, le composé
$c_g \circ g_{small}^*c_f$ est égal à $c_{f \circ g}$.
\end{lemma}

\begin{proof}
Combiner les lemmes \ref{spaces-properties-lemma-etale-exact-pullback}
et \ref{spaces-properties-lemma-characterize-sheaf-small-etale}, et utiliser le fait que
tout morphisme entre objets de $X_\etale$ est un morphisme \'etale
de schémas.
\end{proof}







\section{Localisation \'etale}
\label{spaces-properties-section-localize}

\noindent
Il faut éviter à tout prix de lire cette section.

\medskip\noindent
Soit $X \to Y$ un morphisme \'etale d'espaces algébriques.
Alors $X$ est un objet de $Y_{spaces, \etale}$ et il résulte
immédiatement des définitions, voir aussi la preuve du
lemme \ref{spaces-properties-lemma-etale-morphism-topoi},
que
\begin{equation}
\label{spaces-properties-equation-localize}
X_{spaces, \etale} = Y_{spaces, \etale}/X
\end{equation}
où le membre de droite est la localisation du site
$Y_{spaces, \etale}$ en l'objet $X$ ; voir la
définition \ref{sites-definition-localize} de Sites.
De plus, cette identification est compatible avec les faisceaux structuraux d'après le
lemme \ref{spaces-properties-lemma-etale-exact-pullback}.
Par conséquent, le site annelé $(X_{spaces, \etale}, \mathcal{O}_X)$
s'identifie à la localisation du site annelé
$(Y_{spaces, \etale}, \mathcal{O}_Y)$ en l'objet $X$ :
\begin{equation}
\label{spaces-properties-equation-localize-ringed}
(X_{spaces, \etale}, \mathcal{O}_X) =
(Y_{spaces, \etale}/X, \mathcal{O}_Y|_{Y_{spaces, \etale}/X})
\end{equation}
La localisation d'un site annelé utilisée dans le membre de droite est définie dans
Modules sur les sites, à la
définition \ref{sites-modules-definition-localize-ringed-site}.

\medskip\noindent
Supposons maintenant que $X \to Y$ soit un morphisme \'etale d'espaces algébriques et que $X$ soit
un schéma. Alors $X$ est un objet de $Y_\etale$ et il s'ensuit que
\begin{equation}
\label{spaces-properties-equation-localize-at-scheme}
X_\etale = Y_\etale/X
\end{equation}
et
\begin{equation}
\label{spaces-properties-equation-localize-at-scheme-ringed}
(X_\etale, \mathcal{O}_X) =
(Y_\etale/X, \mathcal{O}_Y|_{Y_\etale/X})
\end{equation}
comme ci-dessus.

\medskip\noindent
Enfin, si $X \to Y$ est un morphisme \'etale d'espaces algébriques et si $X$ est
un schéma affine, alors $X$ est un objet de $Y_{affine, \etale}$ et
\begin{equation}
\label{spaces-properties-equation-localize-at-affine}
X_{affine, \etale} = Y_{affine, \etale}/X
\end{equation}
et
\begin{equation}
\label{spaces-properties-equation-localize-at-affine-ringed}
(X_{affine, \etale}, \mathcal{O}_X) =
(Y_{affine, \etale}/X, \mathcal{O}_Y|_{Y_{affine, \etale}/X})
\end{equation}
comme ci-dessus.

\medskip\noindent
Montrons ensuite que ces localisations sont compatibles avec les morphismes.

\begin{lemma}
\label{spaces-properties-lemma-relocalize-morphism}
Soit $S$ un schéma. Soit
$$
\xymatrix{
U \ar[d]_p \ar[r]_g & V \ar[d]^q \\
X \ar[r]^f & Y
}
$$
un diagramme commutatif d'espaces algébriques sur $S$ dans lequel $p$ et $q$ sont \'etales.
Par les identifications
(\ref{spaces-properties-equation-localize-ringed}) pour $U \to X$ et $V \to Y$,
le morphisme de topos annelés
$$
(g_{spaces, \etale}, g^\sharp) :
(\Sh(U_{spaces, \etale}), \mathcal{O}_U)
\longrightarrow
(\Sh(V_{spaces, \etale}), \mathcal{O}_V)
$$
est $2$-isomorphe au morphisme $(f_{spaces, \etale, c}, f_c^\sharp)$
construit dans
Modules sur les sites, au
lemme \ref{sites-modules-lemma-relocalize-morphism-ringed-sites},
à partir du morphisme de sites annelés
$(f_{spaces, \etale}, f^\sharp)$ et de
l'application $c : U \to V \times_Y X$ correspondant à $g$.
\end{lemma}

\begin{proof}
Le morphisme $(f_{spaces, \etale, c}, f_c^\sharp)$ est défini comme le
composé $f' \circ j$
d'un morphisme de localisation et d'une application de changement de base. De même, $g$ est le composé
$U \to V \times_Y X \to V$. Il suffit donc de démontrer
le lemme dans les deux cas suivants : (1) $f = \text{id}$, et
(2) $U = X \times_Y V$. Dans le cas (1), le morphisme $g : U \to V$ est
\'etale ; voir le
lemme \ref{spaces-properties-lemma-etale-permanence}.
Par conséquent, $(g_{spaces, \etale}, g^\sharp)$ est un morphisme de localisation
d'après la discussion autour des
équations (\ref{spaces-properties-equation-localize}) et
(\ref{spaces-properties-equation-localize-ringed}),
ce qui est exactement l'assertion du lemme dans ce cas.
Dans le cas (2), le morphisme $g_{spaces, \etale}$
provient du morphisme de sites annelés donné par le foncteur
$V_{spaces, \etale} \to U_{spaces, \etale}$,
$V'/V \mapsto V' \times_V U/U$
qui est aussi celui par lequel le morphisme $f'$ est défini ; voir le
lemme \ref{sites-lemma-localize-morphism} de Sites.
Nous omettons la vérification de $(f')^\sharp = g^\sharp$
dans ce cas (tous deux sont la restriction de $f^\sharp$
à $U_{spaces, \etale}$).
\end{proof}

\begin{lemma}
\label{spaces-properties-lemma-relocalize-morphism-at-schemes}
Mêmes notations et hypothèses que dans le
lemme \ref{spaces-properties-lemma-relocalize-morphism},
sauf que nous supposons aussi que $U$ et $V$ sont des schémas.
Par les identifications
(\ref{spaces-properties-equation-localize-at-scheme-ringed})
pour $U \to X$ et $V \to Y$, le morphisme de topos annelés
$$
(g_{small}, g^\sharp) :
(\Sh(U_\etale), \mathcal{O}_U)
\longrightarrow
(\Sh(V_\etale), \mathcal{O}_V)
$$
est $2$-isomorphe au morphisme $(f_{small, s}, f_s^\sharp)$
construit dans
Modules sur les sites, au
lemme \ref{sites-modules-lemma-relocalize-morphism-ringed-topoi},
à partir de $(f_{small}, f^\sharp)$ et de
l'application $s : h_U \to f_{small}^{-1}h_V$ correspondant à $g$.
\end{lemma}

\begin{proof}
Notons que $(g_{small}, g^\sharp)$ est $2$-isomorphe, comme
morphisme de topos annelés, au morphisme de topos annelés
associé au morphisme de sites annelés
$(g_{spaces, \etale}, g^\sharp)$. Nous concluons donc par le
lemme \ref{spaces-properties-lemma-relocalize-morphism}
et par le résultat de
Modules sur les sites donné au
lemme \ref{sites-modules-lemma-relocalize-morphism-compare}.
\end{proof}

\noindent
Enfin, nous examinons la relation entre les faisceaux d'ensembles sur le petit
site \'etale $Y_\etale$ d'un espace algébrique $Y$ et les espaces algébriques
\'etales sur $Y$. Soit $S$ un schéma et soit $Y$ un espace algébrique
sur $S$. Soit $\mathcal{F}$ un objet de $\Sh(Y_\etale)$. Considérons
le foncteur
$$
X : (\Sch/S)_{fppf}^{opp} \longrightarrow \textit{Ens}
$$
défini par la règle
$$
X(T) = \{(y, s) \mid y : T \to Y\text{ est un morphisme sur }S\text{ et }
s \in \Gamma(T, y_{small}^{-1}\mathcal{F})\}
$$
Étant donné un morphisme $g : T' \to T$, l'application de restriction envoie
$(y, s)$ sur $(y \circ g, g_{small}^{-1}s)$. Cela a un sens
puisque $y_{small} \circ g_{small} = (y \circ g)_{small}$ d'après le
lemme \ref{spaces-properties-lemma-functoriality-etale-site}. Il existe une application canonique
$X \to Y$ qui envoie le couple $(y, s)$ sur $y$.

\begin{lemma}
\label{spaces-properties-lemma-sheaf-gives-space}
Soit $S$ un schéma et soit $Y$ un espace algébrique sur $S$.
Soit $\mathcal{F}$ un faisceau d'ensembles sur $Y_\etale$.
Sous réserve d'une condition ensembliste (voir la preuve), nous avons :
\begin{enumerate}
\item le foncteur $X$ associé ci-dessus à $\mathcal{F}$ est un
espace algébrique ;
\item l'application $X \to Y$ est un morphisme \'etale d'espaces algébriques ;
\item par l'identification $\Sh(Y_\etale) = \Sh(Y_{spaces, \etale})$,
nous avons $\mathcal{F} \cong h_X$ ;
\item nous avons $\mathcal{F} \cong f_{small, !}*$. Ici, $*$ est l'objet final
de la catégorie $\Sh(X_\etale)$ et $f_{small, !}$ existe d'après le
lemme \ref{spaces-properties-lemma-etale-morphism-topoi}.
\end{enumerate}
\end{lemma}

\begin{proof}
Montrons que $X$ est un faisceau pour la topologie fppf. Supposons donc
que $\{g_i : T_i \to T\}$ soit un recouvrement de $(\Sch/S)_{fppf}$ et que
les $(y_i, s_i) \in X(T_i)$ vérifient la condition de recollement, c'est-à-dire que les
restrictions de $(y_i, s_i)$ et de $(y_j, s_j)$ à $T_i \times_T T_j$ coïncident.
Puisque $Y$ est un faisceau pour la topologie fppf, nous voyons alors que
les $y_i$ donnent naissance à un unique morphisme $y : T \to Y$ tel que
$y_i = y \circ g_i$. Nous voyons ensuite que
$y_{i, small}^{-1}\mathcal{F} = g_{i, small}^{-1}y_{small}^{-1}\mathcal{F}$.
Les sections $s_i$ se recollent donc de manière unique en une section de
$y_{small}^{-1}\mathcal{F}$ d'après le résultat donné dans
Cohomologie \'etale, lemme \ref{etale-cohomology-lemma-describe-pullback}.

\medskip\noindent
La construction qui envoie $\mathcal{F} \in \Ob(\Sh(Y_\etale))$
sur $X \in \Ob(\Sh((\Sch/S)_{fppf}))$ commute aux limites finies et
aux colimites quelconques, puisque chacun des foncteurs $y_{small}^{-1}$ possède
ces propriétés. Bien entendu, si $V \in \Ob(Y_\etale)$, alors
la construction envoie le faisceau représentable $h_V$ de $Y_\etale$
sur le foncteur représentable représenté par $V$.

\medskip\noindent
D'après le lemme \ref{sites-lemma-sheaf-coequalizer-representable} de Sites,
nous pouvons trouver un ensemble $I$ et, pour chaque $i \in I$, un objet $V_i$ de $Y_\etale$
et une surjection de faisceaux
$$
\coprod h_{V_i} \longrightarrow \mathcal{F}
$$
sur $Y_\etale$. La condition ensembliste dont nous avons besoin est que l'ensemble
d'indices $I$ ne soit pas trop grand\footnote{Il suffit que la borne supérieure des
cardinalités des fibres de $\mathcal{F}$ aux points géométriques de $Y$
soit majorée par la taille d'un certain objet de $(\Sch/S)_{fppf}$.}.
Alors $V = \coprod V_i$ est un objet de $(\Sch/S)_{fppf}$ et
donc un objet de $Y_\etale$, et nous avons une surjection
$h_V \to \mathcal{F}$.

\medskip\noindent
Observons que le produit de $h_V$ par lui-même dans $\Sh(Y_\etale)$
est $h_{V \times_Y V}$. Considérons le produit fibré
$$
h_V \times_\mathcal{F} h_V \subset h_{V \times_Y V}
$$
Il existe un sous-schéma ouvert $R$ de $V \times_Y V$ tel que
$h_V \times_\mathcal{F} h_V = h_R$ ; voir le
lemme \ref{spaces-properties-lemma-support-subsheaf-final} (un petit détail est omis).
Par le lemme de Yoneda, nous obtenons deux morphismes $s, t : R \to V$ dans $Y_\etale$
et nous trouvons un diagramme de conoyau du couple
$$
\xymatrix{
h_R \ar@<1ex>[r] \ar@<-1ex>[r] &
h_V \ar[r] &
\mathcal{F}
}
$$
dans $\Sh(Y_\etale)$. Bien entendu, les morphismes $s, t$ sont \'etales
et définissent une relation d'équivalence \'etale $(t, s) : R \to V \times_S V$.

\medskip\noindent
La discussion des deux paragraphes précédents nous donne un
diagramme de conoyau du couple
$$
\xymatrix{
R \ar@<1ex>[r] \ar@<-1ex>[r] &
V \ar[r] &
X
}
$$
dans $(\Sch/S)_{fppf}$. Ainsi, $X = V/R$ est un espace algébrique d'après le
théorème \ref{spaces-theorem-presentation} d'Espaces. Cela démontre (1).
L'assertion (2) résulte de ce que $V \to Y$ est \'etale. L'assertion (3) découle immédiatement
de la définition de $X$ et de $h_X$. Nous omettons la preuve de l'assertion (4) ;
elle résulte de l'identification du morphisme associé au foncteur cocontinu
$j$ du lemme \ref{spaces-properties-lemma-etale-morphism-topoi}
avec la description de $X_{spaces, \etale}$ comme localisation
de $Y_{spaces, \etale}$ en $X$ donnée ci-dessus, et du
lemme \ref{sites-lemma-describe-j-shriek-representable} de Sites.
\end{proof}







\section{Reconstruction des morphismes}
\label{spaces-properties-section-morphisms}

\noindent
Dans cette section, nous démontrons que le procédé qui associe à un espace algébrique
son petit topos \'etale localement annelé est pleinement fidèle en un sens
approprié ; voir le
théorème \ref{spaces-properties-theorem-fully-faithful}.

\begin{lemma}
\label{spaces-properties-lemma-morphism-locally-ringed}
Soit $S$ un schéma.
Soit $f : X \to Y$ un morphisme d'espaces algébriques sur $S$.
Le morphisme de topos annelés $(f_{small}, f^\sharp)$
associé à $f$ est un morphisme de topos localement annelés ; voir
Modules sur les sites,
la définition \ref{sites-modules-definition-morphism-locally-ringed-topoi}.
\end{lemma}

\begin{proof}
Remarquons que l'assertion a un sens puisque nous avons vu que
$(X_\etale, \mathcal{O}_{X_\etale})$ et
$(Y_\etale, \mathcal{O}_{Y_\etale})$
sont des sites localement annelés ; voir le
lemme \ref{spaces-properties-lemma-etale-site-locally-ringed}.
De plus, nous savons que $X_\etale$ a assez de points ; voir le
théorème \ref{spaces-properties-theorem-exactness-stalks}.
Il suffit donc de démontrer que $(f_{small}, f^\sharp)$
satisfait la condition (3) du
Modules sur les sites,
lemme \ref{sites-modules-lemma-locally-ringed-morphism}.
Pour le voir, prenons un point $p$ de $X_\etale$. D'après le
lemme \ref{spaces-properties-lemma-points-small-etale-site},
$p$ correspond à un point géométrique $\overline{x}$ de $X$.
D'après le
lemme \ref{spaces-properties-lemma-stalk-pullback},
le point $q = f_{small} \circ p$ correspond au
point géométrique $\overline{y} = f \circ \overline{x}$ de $Y$.
L'assertion à démontrer est donc que l'homomorphisme induit
d'anneaux locaux \'etales
$$
\mathcal{O}_{Y, \overline{y}} \longrightarrow \mathcal{O}_{X, \overline{x}}
$$
est un homomorphisme local d'anneaux locaux. On peut le démontrer directement, mais nous le déduisons plutôt
du résultat correspondant pour les schémas. Choisissons à cette fin un
diagramme commutatif
$$
\xymatrix{
U \ar[d] \ar[r]_\psi & V \ar[d] \\
X \ar[r] & Y
}
$$
où $U$ et $V$ sont des schémas et les flèches verticales sont surjectives
et \'etales (voir le
lemme \ref{spaces-lemma-lift-morphism-presentations} d'Espaces).
Choisissons un relèvement $\overline{u} : \overline{x} \to U$ (ce qui est possible d'après le
lemme \ref{spaces-properties-lemma-geometric-lift-to-cover}).
Posons $\overline{v} = \psi \circ \overline{u}$. Nous obtenons un
diagramme commutatif d'anneaux locaux \'etales
$$
\xymatrix{
\mathcal{O}_{U, \overline{u}} &
\mathcal{O}_{V, \overline{v}} \ar[l] \\
\mathcal{O}_{X, \overline{x}} \ar[u] &
\mathcal{O}_{Y, \overline{y}}. \ar[l] \ar[u]
}
$$
D'après
Cohomologie \'etale, le lemme \ref{etale-cohomology-lemma-morphism-locally-ringed},
la flèche horizontale supérieure est un homomorphisme local d'anneaux locaux. Enfin, d'après le
lemme \ref{spaces-properties-lemma-describe-etale-local-ring},
les flèches verticales sont des isomorphismes. Cela conclut.
\end{proof}

\begin{lemma}
\label{spaces-properties-lemma-2-morphism}
Soit $S$ un schéma.
Soient $X$, $Y$ des espaces algébriques sur $S$.
Soit $f : X \to Y$ un morphisme d'espaces algébriques sur $S$.
Soit $t$ un $2$-morphisme de $(f_{small}, f^\sharp)$ vers lui-même ; voir
Modules sur les sites,
la définition \ref{sites-modules-definition-2-morphism-ringed-topoi}.
Alors $t = \text{id}$.
\end{lemma}

\begin{proof}
Soient $X'$, resp.\ $Y'$, les espaces $X$, resp.\ $Y$, considérés comme espaces algébriques sur
$\Spec(\mathbf{Z})$ ; voir la
définition \ref{spaces-definition-base-change} d'Espaces.
Il ressort de la construction que $(X_{small}, \mathcal{O})$
est égal à $(X'_{small}, \mathcal{O})$, et de même pour $Y$.
Nous pouvons donc travailler avec $X'$ et $Y'$. Autrement dit, nous pouvons
supposer que $S = \Spec(\mathbf{Z})$.

\medskip\noindent
Supposons que $S = \Spec(\mathbf{Z})$ et que $f : X \to Y$ et $t$ soient comme dans
le lemme. Autrement dit, $t : f^{-1}_{small} \to f^{-1}_{small}$
est une transformation de foncteurs telle que le diagramme
$$
\xymatrix{
f_{small}^{-1}\mathcal{O}_Y
\ar[rd]_{f^\sharp}  & &
f_{small}^{-1}\mathcal{O}_Y \ar[ll]^t \ar[ld]^{f^\sharp} \\
& \mathcal{O}_X
}
$$
soit commutatif. Supposons que $V \to Y$ soit \'etale, avec $V$ affine.
Écrivons $V = \Spec(B)$. Choisissons des générateurs $b_j \in B$, $j \in J$,
de $B$ comme $\mathbf{Z}$-algèbre. Posons
$T = \Spec(\mathbf{Z}[\{x_j\}_{j \in J}])$.
Dans la suite, nous utiliserons le fait que
$\Mor_{\Sch}(U, T) = \prod_{j \in J} \Gamma(U, \mathcal{O}_U)$
pour tout schéma $U$, sans le rappeler davantage.
L'homomorphisme surjectif d'anneaux $\mathbf{Z}[x_j] \to B$, $x_j \mapsto b_j$,
correspond à une immersion fermée $V \to T$.
Nous obtenons un monomorphisme
$$
i : V \longrightarrow T_Y = T \times Y
$$
d'espaces algébriques sur $Y$. En termes de faisceaux sur $Y_\etale$,
le morphisme $i$ induit une injection de faisceaux
$h_i : h_V \to \prod_{j \in J} \mathcal{O}_Y$.
Le changement de base $i' : X \times_Y V \to T_X$ de $i$ à $X$
est lui aussi un monomorphisme
(voir le
lemme \ref{spaces-lemma-base-change-representable-transformations-property} d'Espaces).
Ainsi, $i' : X \times_Y V \to T_X$ est un monomorphisme, ce qui
signifie à son tour que
$h_{i'} : h_{X \times_Y V} \to \prod_{j \in J} \mathcal{O}_X$
est une injection de faisceaux.
Par l'identification $f_{small}^{-1}h_V = h_{X \times_Y V}$ du
lemme \ref{spaces-properties-lemma-stalk-pullback},
l'application $h_{i'}$ est égale à
$$
\xymatrix{
f_{small}^{-1}h_V \ar[r]^-{f^{-1}h_i} &
\prod_{j \in J} f_{small}^{-1}\mathcal{O}_Y
\ar[r]^{\prod f^\sharp} &
\prod_{j \in J} \mathcal{O}_X
}
$$
(vérification omise). Cela signifie que l'application
$t : f_{small}^{-1}h_V \to f_{small}^{-1}h_V$
s'insère dans le diagramme commutatif
$$
\xymatrix{
f_{small}^{-1}h_V \ar[r]^-{f^{-1}h_i} \ar[d]^t &
\prod_{j \in J} f_{small}^{-1}\mathcal{O}_Y
\ar[r]^-{\prod f^\sharp} \ar[d]^{\prod t} &
\prod_{j \in J} \mathcal{O}_X \ar[d]^{\text{id}}\\
f_{small}^{-1}h_V \ar[r]^-{f^{-1}h_i} &
\prod_{j \in J} f_{small}^{-1}\mathcal{O}_Y
\ar[r]^-{\prod f^\sharp} &
\prod_{j \in J} \mathcal{O}_X
}
$$
Le carré de droite est commutatif par l'hypothèse faite ci-dessus sur $t$.

Comme la composée des flèches horizontales est injective
d'après la discussion précédente, nous concluons que la flèche verticale de gauche
est elle aussi l'identité. Tout faisceau d'ensembles sur
$Y_\etale$ admet une surjection provenant d'une somme disjointe (énorme) de faisceaux
de la forme $h_V$, avec $V$ affine (combiner le
lemme \ref{spaces-properties-lemma-alternative}
et le
lemme \ref{sites-lemma-sheaf-coequalizer-representable} de Sites).
Nous concluons ainsi que $t : f_{small}^{-1} \to f_{small}^{-1}$
est la transformation identique, comme souhaité.
\end{proof}

\begin{lemma}
\label{spaces-properties-lemma-faithful}
Soit $S$ un schéma.
Soient $X$, $Y$ des espaces algébriques sur $S$.
Deux morphismes quelconques $a, b : X \to Y$ d'espaces algébriques sur $S$
pour lesquels il existe un $2$-isomorphisme
$(a_{small}, a^\sharp) \cong (b_{small}, b^\sharp)$
dans la $2$-catégorie des topos annelés sont égaux.
\end{lemma}

\begin{proof}
Soit $t : a_{small}^{-1} \to b_{small}^{-1}$ le $2$-isomorphisme.
Nous pouvons, de manière équivalente, considérer $t$ comme une transformation
$t : a_{spaces, \etale}^{-1} \to b_{spaces, \etale}^{-1}$
puisqu'il n'y a aucune différence entre les faisceaux sur $X_\etale$
et ceux sur $X_{spaces, \etale}$.
Choisissons un diagramme commutatif
$$
\xymatrix{
U \ar[d]_p \ar[r]_\alpha  & V \ar[d]^q \\
X \ar[r]^a & Y
}
$$
où $U$ et $V$ sont des schémas, et où $p$ et $q$ sont \'etales surjectifs.
Considérons le diagramme
$$
\xymatrix{
h_U \ar[r]_-\alpha \ar@{=}[d] & a_{spaces, \etale}^{-1}h_V \ar[d]^t \\
h_U \ar@{..>}[r] & b_{spaces, \etale}^{-1}h_V
}
$$
Puisque le faisceau $b_{spaces, \etale}^{-1}h_V$ est isomorphe à
$h_{V \times_{Y, b} X}$, nous voyons que la flèche pointillée provient d'un
morphisme de schémas
$\beta : U \to V$ qui s'insère dans un diagramme commutatif
$$
\xymatrix{
U \ar[d]_p \ar[r]_\beta  & V \ar[d]^q \\
X \ar[r]^b & Y
}
$$
Nous affirmons qu'il existe une suite de $2$-isomorphismes
\begin{align*}
(\alpha_{small}, \alpha^\sharp)
& \cong
(\alpha_{spaces, \etale}, \alpha^\sharp) \\
& \cong
(a_{spaces, \etale, c}, a_c^\sharp) \\
& \cong
(b_{spaces, \etale, d}, b_d^\sharp) \\
& \cong
(\beta_{spaces, \etale}, \beta^\sharp) \\
& \cong
(\beta_{small}, \beta^\sharp)
\end{align*}
Le premier et le dernier $2$-isomorphisme proviennent des identifications
entre les faisceaux sur $U_{spaces, \etale}$ et ceux sur
$U_\etale$, et de même pour $V$. Les deuxième et quatrième
$2$-isomorphismes sont ceux du
lemme \ref{spaces-properties-lemma-relocalize-morphism},
où $c : U \to X \times_{a, Y} V$ est induit par $\alpha$ et
$d : U \to X \times_{b, Y} V$ est induit par $\beta$.
Le $2$-isomorphisme médian provient de la transformation $t$.
Plus précisément, le foncteur $a_{spaces, \etale, c}^{-1}$ correspond
au foncteur
$$
(\mathcal{H} \to h_V) \longmapsto
(a_{spaces, \etale}^{-1}\mathcal{H}
\times_{a_{spaces, \etale}^{-1}h_V, \alpha}
h_U \to
h_U)
$$
et de même pour $b_{spaces, \etale, d}^{-1}$ ; voir le
lemme \ref{sites-lemma-relocalize-morphism} de Sites.
On utilise ici l'identification des faisceaux sur $Y_{spaces, \etale}/V$
avec les flèches $(\mathcal{H} \to h_V)$ dans $\Sh(Y_{spaces, \etale})$,
et de même pour $U/X$ ; voir le
lemme \ref{sites-lemma-essential-image-j-shriek} de Sites.
Par cette identification, le faisceau structural $\mathcal{O}_V$ correspond à la
paire $(\mathcal{O}_Y \times h_V \to h_V)$, et de même
pour $\mathcal{O}_U$ ; voir le
Modules sur les sites,
lemme \ref{sites-modules-lemma-localize-compare}.
Puisque $t$ échange $\alpha$ et $\beta$, nous voyons que $t$ induit un isomorphisme
$$
t :
a_{spaces, \etale}^{-1}\mathcal{H}
\times_{a_{spaces, \etale}^{-1}h_V, \alpha}
h_U
\longrightarrow
b_{spaces, \etale}^{-1}\mathcal{H}
\times_{b_{spaces, \etale}^{-1}h_V, \beta}
h_U
$$
au-dessus de $h_U$, fonctoriellement en $(\mathcal{H} \to h_V)$. De plus, $t$ est compatible
avec $a_c^\sharp$ et $b_d^\sharp$, puisque $t$ est
compatible avec $a^\sharp$ et $b^\sharp$ d'après notre description
ci-dessus des faisceaux structuraux $\mathcal{O}_U$ et $\mathcal{O}_V$.
Par conséquent, les morphismes de topos annelés
$(\alpha_{small}, \alpha^\sharp)$ et $(\beta_{small}, \beta^\sharp)$
sont $2$-isomorphes. D'après
Cohomologie \'etale, le lemme \ref{etale-cohomology-lemma-faithful},
nous concluons que $\alpha = \beta$ ! Puisque $p : U \to X$ est une surjection
de faisceaux, il s'ensuit que $a = b$.
\end{proof}

\noindent
Voici le résultat principal de cette section.

\begin{theorem}
\label{spaces-properties-theorem-fully-faithful}
Soient $X$, $Y$ des espaces algébriques sur $\Spec(\mathbf{Z})$.
Soit
$$
(g, g^\sharp) :
(\Sh(X_\etale), \mathcal{O}_X)
\longrightarrow
(\Sh(Y_\etale), \mathcal{O}_Y)
$$
un morphisme de topos localement annelés. Alors il existe un
unique morphisme d'espaces algébriques $f : X \to Y$ tel que
$(g, g^\sharp)$ soit isomorphe à $(f_{small}, f^\sharp)$.
Autrement dit, la construction
$$
\textit{Espaces}/\Spec(\mathbf{Z})
\longrightarrow \textit{Topos localement annelés},
\quad
X \longrightarrow (X_\etale, \mathcal{O}_X)
$$
est pleinement fidèle (les morphismes du membre de droite étant pris à $2$-isomorphisme près).
\end{theorem}

\begin{proof}
L'unicité a été démontrée au
lemme \ref{spaces-properties-lemma-faithful}.
Il suffit donc de démontrer l'existence.
Dans cette preuve, nous utiliserons librement les identifications de
l'équation (\ref{spaces-properties-equation-localize-at-scheme-ringed})
ainsi que le résultat du
lemme \ref{spaces-properties-lemma-relocalize-morphism-at-schemes}.

\medskip\noindent
Soit $U \in \Ob(X_\etale)$, soit
$V \in \Ob(Y_\etale)$,
et soit $s \in g^{-1}h_V(U)$ une section. Nous pouvons considérer
$s$ comme une application de faisceaux $s : h_U \to g^{-1}h_V$. D'après le
Modules sur les sites,
lemme \ref{sites-modules-lemma-relocalize-morphism-ringed-topoi},
nous obtenons un diagramme commutatif de morphismes de topos annelés
$$
\xymatrix{
(\Sh(X_\etale/U), \mathcal{O}_U)
\ar[rr]_-{(j, j^\sharp)} \ar[d]_{(g_s, g_s^\sharp)} & &
(\Sh(X_\etale), \mathcal{O}_X) \ar[d]^{(g, g^\sharp)} \\
(\Sh(V_\etale), \mathcal{O}_V) \ar[rr] & &
(\Sh(Y_\etale), \mathcal{O}_Y).
}
$$
D'après
Cohomologie \'etale, le théorème \ref{etale-cohomology-theorem-fully-faithful},
nous obtenons un unique morphisme de schémas $f_s : U \to V$ tel que
$(g_s, g_s^\sharp)$ soit $2$-isomorphe à $(f_{s, small}, f_s^\sharp)$.
La construction $(U, V, s) \leadsto f_s$ que nous venons de décrire satisfait
la propriété de fonctorialité suivante. Supposons donnés des morphismes
$a : U' \to U$ dans $X_\etale$ et $b : V' \to V$ dans $Y_\etale$,
ainsi qu'une application $s' : h_{U'} \to g^{-1}h_{V'}$ telle que le diagramme
$$
\xymatrix{
h_{U'} \ar[d]_a \ar[r]_{s'} & g^{-1}h_{V'} \ar[d]^{g^{-1}b} \\
h_U \ar[r]^s & g^{-1}h_V
}
$$
soit commutatif. Alors le diagramme
$$
\xymatrix{
U' \ar[r]_-{f_{s'}} \ar[d]_a & u(V') \ar[d]^{u(b)} \\
U \ar[r]^-{f_s} & u(V)
}
$$
de schémas est commutatif. En effet, la même condition
vaut pour les morphismes $(g_s, g_s^\sharp)$ construits dans
Modules sur les sites,
au lemme \ref{sites-modules-lemma-relocalize-morphism-ringed-topoi},
et l'on applique l'unicité de
Cohomologie \'etale, théorème \ref{etale-cohomology-theorem-fully-faithful}.

\medskip\noindent
Il s'agit de recoller les morphismes $f_s$ en un morphisme d'espaces
algébriques. Pour ce faire, choisissons d'abord un schéma $V$ et un morphisme \'etale
surjectif $V \to Y$. Ainsi, $h_V \to *$ est surjectif et, par conséquent,
$g^{-1}h_V \to *$ est lui aussi surjectif. Il existe donc un schéma $U$,
un morphisme \'etale surjectif $U \to X$ et un morphisme
$s : h_U \to g^{-1}h_V$. Posons ensuite $R = V \times_Y V$ et
$R' = U \times_X U$. Nous obtenons alors
$g^{-1}h_R = g^{-1}h_V \times g^{-1}h_V$, car $g^{-1}$ est exact.
Ainsi, $s$ induit un morphisme $s \times s : h_{R'} \to g^{-1}h_R$.
En appliquant les constructions précédentes, nous obtenons un
diagramme commutatif de morphismes de schémas
$$
\xymatrix{
R' \ar@<1ex>[d] \ar@<-1ex>[d] \ar[rr]_{f_{s \times s}} & &
R \ar@<1ex>[d] \ar@<-1ex>[d] \\
U \ar[rr]^{f_s} & &
V
}
$$
Comme $X = U/R'$ et $Y = V/R$ (voir le
lemme \ref{spaces-lemma-space-presentation} d'Espaces),
nous concluons que ce diagramme
définit un morphisme d'espaces algébriques $f : X \to Y$ qui s'insère
dans le diagramme commutatif évident.
Il reste à montrer que $(f_{small}, f^\sharp)$ est
$2$-isomorphe à $(g, g^\sharp)$.
Soient $t_V : f_{s, small}^{-1} \to g_s^{-1}$ et
$t_R : f_{s \times s, small}^{-1} \to g_{s \times s}^{-1}$ les
$2$-isomorphismes fournis par la construction précédente.
Soit $\mathcal{G}$ un faisceau sur $Y_\etale$. Nous voyons alors que
$t_V$ définit un isomorphisme
$$
f_{small}^{-1}\mathcal{G}|_{U_\etale}
=
f_{s, small}^{-1}\mathcal{G}|_{V_\etale}
\xrightarrow{t_V}
g_s^{-1}\mathcal{G}|_{V_\etale}
=
g^{-1}\mathcal{G}|_{U_\etale}.
$$
De plus, l'image inverse de cet isomorphisme sur $R'$ par l'une ou l'autre projection
$R' \to U$ est l'isomorphisme
$$
f_{small}^{-1}\mathcal{G}|_{R'_\etale}
=
f_{s \times s, small}^{-1}\mathcal{G}|_{R_\etale}
\xrightarrow{t_R}
g_{s \times s}^{-1}\mathcal{G}|_{R_\etale}
=
g^{-1}\mathcal{G}|_{R'_\etale}.
$$
Comme $\{U \to X\}$ est un recouvrement du site $X_{spaces, \etale}$,
cela signifie que le premier isomorphisme affiché descend en un isomorphisme
$t : f_{small}^{-1}\mathcal{G} \to g^{-1}\mathcal{G}$
de faisceaux (un petit détail est omis). L'isomorphisme est fonctoriel
en $\mathcal{G}$, puisque $t_V$ et $t_R$ sont des transformations de foncteurs.
Enfin, $t$ est compatible avec $f^\sharp$ et $g^\sharp$, puisque
$t_V$ et $t_R$ le sont (certains détails sont omis).
Cela achève la preuve du théorème.
\end{proof}

\begin{lemma}
\label{spaces-properties-lemma-isomorphism-ringed-topoi}
Soient $X$, $Y$ des espaces algébriques sur $\mathbf{Z}$. Si
$$
(g, g^\sharp) :
(\Sh(X_\etale), \mathcal{O}_X)
\longrightarrow
(\Sh(Y_\etale), \mathcal{O}_Y)
$$
est un isomorphisme de topos annelés, alors il existe un unique
morphisme d'espaces algébriques $f : X \to Y$ tel que
$(g, g^\sharp)$ soit isomorphe à $(f_{small}, f^\sharp)$ ;
de plus, $f$ est un isomorphisme d'espaces algébriques.
\end{lemma}

\begin{proof}
D'après le
théorème \ref{spaces-properties-theorem-fully-faithful},
il suffit de montrer que $(g, g^\sharp)$ est un morphisme de
topos localement annelés. D'après le
lemme \ref{sites-modules-lemma-locally-ringed-morphism} de Modules sur les sites
(et puisque le site $X_\etale$ a assez de points),
il suffit de vérifier que l'homomorphisme
$\mathcal{O}_{Y, q} \to \mathcal{O}_{X, p}$ induit par $g^\sharp$
est un homomorphisme local d'anneaux locaux, où $q = g \circ p$ et $p$ est un point quelconque de
$X_\etale$. Comme c'est un isomorphisme, cela est clair.
\end{proof}














\section{Faisceaux quasi-cohérents sur les espaces algébriques}
\label{spaces-properties-section-quasi-coherent}

\noindent
Dans les
sections \ref{descent-section-quasi-coherent-sheaves},
\ref{descent-section-quasi-coherent-cohomology} et
\ref{descent-section-quasi-coherent-sheaves-bis}
de Descente, nous avons vu que, pour un schéma $U$, il n'y a aucune différence entre un
$\mathcal{O}_U$-module quasi-cohérent sur $U$ et un
$\mathcal{O}$-module quasi-cohérent sur le petit site \'etale de $U$. La définition
suivante est donc compatible avec notre notion initiale de faisceau quasi-cohérent
sur un schéma
(Schémas, section \ref{schemes-section-quasi-coherent})
lorsqu'on l'applique à un espace algébrique représentable.

\begin{definition}
\label{spaces-properties-definition-quasi-coherent}
Soit $S$ un schéma. Soit $X$ un espace algébrique sur $S$.
Un {\it quasi-cohérent} $\mathcal{O}_X$-module
est un module quasi-cohérent sur le site annelé
$(X_\etale, \mathcal{O}_X)$ au sens de
Modules sur les sites,
la définition \ref{sites-modules-definition-site-local}.
La catégorie des faisceaux quasi-cohérents sur $X$ se note
$\QCoh(\mathcal{O}_X)$.
\end{definition}

\noindent
Remarquons que, la quasi-cohérence étant une notion intrinsèque (voir le
lemme \ref{sites-modules-lemma-special-locally-free} de Modules sur les sites),
cela revient à dire que le $\mathcal{O}_X$-module correspondant
sur $X_{spaces, \etale}$ est quasi-cohérent.

\medskip\noindent
Comme d'habitude, les faisceaux quasi-cohérents se comportent bien par image inverse.

\begin{lemma}
\label{spaces-properties-lemma-pullback-quasi-coherent}
Soit $S$ un schéma.
Soit $f : X \to Y$ un morphisme d'espaces algébriques sur $S$.
Le foncteur image inverse
$f^* : \textit{Mod}(\mathcal{O}_Y) \to \textit{Mod}(\mathcal{O}_X)$
préserve les faisceaux quasi-cohérents.
\end{lemma}

\begin{proof}
C'est un fait général ; voir le
lemme \ref{sites-modules-lemma-local-pullback} de Modules sur les sites.
\end{proof}

\noindent
Remarquons que ce foncteur image inverse coïncide avec le foncteur image inverse usuel
entre faisceaux de modules quasi-cohérents lorsque $X$ et $Y$ sont des
schémas ; voir la
proposition de Descente
\ref{descent-proposition-equivalence-quasi-coherent-functorial}.
Voici le lemme incontournable qui compare cette notion aux faisceaux quasi-cohérents
sur les objets du petit site \'etale de $X$.

\begin{lemma}
\label{spaces-properties-lemma-characterize-quasi-coherent-small-etale}
Soit $S$ un schéma. Soit $X$ un espace algébrique sur $S$.
Un $\mathcal{O}_X$-module quasi-cohérent $\mathcal{F}$
est donné par les données suivantes :
\begin{enumerate}
\item pour tout $U \in \Ob(X_\etale)$, un
$\mathcal{O}_U$-module quasi-cohérent $\mathcal{F}_U$ sur $U_\etale$ ;
\item pour tout $f : U' \to U$ dans $X_\etale$, un isomorphisme
$c_f : f_{small}^*\mathcal{F}_U \to \mathcal{F}_{U'}$.
\end{enumerate}
Ces données sont soumises à la condition que, pour tous $f : U' \to U$
et $g : U'' \to U'$ dans $X_\etale$, la composée
$c_g \circ g_{small}^*c_f$ soit égale à $c_{f \circ g}$.
\end{lemma}

\begin{proof}
Combiner les lemmes \ref{spaces-properties-lemma-pullback-quasi-coherent} et
\ref{spaces-properties-lemma-characterize-module-small-etale}.
\end{proof}

\begin{lemma}
\label{spaces-properties-lemma-stalk-quasi-coherent}
Soit $S$ un schéma.
Soit $X$ un espace algébrique sur $S$.
Soit $\mathcal{F}$ un $\mathcal{O}_X$-module quasi-cohérent.
Soit $x \in |X|$ un point et soit $\overline{x}$ un point
géométrique situé au-dessus de $x$. Enfin, soit
$\varphi : (U, \overline{u}) \to (X, \overline{x})$
un voisinage \'etale, où $U$ est un schéma.
Alors
$$
(\varphi^*\mathcal{F})_u \otimes_{\mathcal{O}_{U, u}}
\mathcal{O}_{X, \overline{x}} =
\mathcal{F}_{\overline{x}}
$$
où $u \in U$ est l'image de $\overline{u}$.
\end{lemma}

\begin{proof}
Remarquons que $\mathcal{O}_{X, \overline{x}} = \mathcal{O}_{U, u}^{sh}$ d'après le
lemme \ref{spaces-properties-lemma-describe-etale-local-ring} ;
le produit tensoriel a donc un sens. De plus, il ressort de la
définition \ref{spaces-properties-definition-stalk}
que
$$
\mathcal{F}_{\overline{x}} = \colim (\varphi^*\mathcal{F})_u
$$
où la colimite porte sur les $\varphi : (U, \overline{u}) \to (X, \overline{x})$
comme dans le lemme. Il existe donc un homomorphisme canonique du membre de gauche vers celui de droite dans
l'énoncé du lemme. Nous disposons d'une description colimite analogue pour
$\mathcal{O}_{X, \overline{x}}$
et, d'après le
lemme \ref{spaces-properties-lemma-characterize-quasi-coherent-small-etale},
nous avons
$$
((\varphi')^*\mathcal{F})_{u'} =
(\varphi^*\mathcal{F})_u \otimes_{\mathcal{O}_{U, u}} \mathcal{O}_{U', u'}
$$
dès que $(U', \overline{u}') \to (U, \overline{u})$ est un morphisme de
voisinages \'etales. Pour achever la preuve, nous utilisons le fait que
$\otimes$ commute aux colimites.
\end{proof}

\begin{lemma}
\label{spaces-properties-lemma-stalk-pullback-quasi-coherent}
Soit $S$ un schéma. Soit $f : X \to Y$ un morphisme d'espaces algébriques
sur $S$. Soit $\mathcal{G}$ un $\mathcal{O}_Y$-module quasi-cohérent.
Soit $\overline{x}$ un point géométrique de $X$ et soit
$\overline{y} = f \circ \overline{x}$ son image dans $Y$.
Alors il existe un isomorphisme canonique
$$
(f^*\mathcal{G})_{\overline{x}} =
\mathcal{G}_{\overline{y}} \otimes_{\mathcal{O}_{Y, \overline{y}}}
\mathcal{O}_{X, \overline{x}}
$$
de la fibre de l'image inverse avec le produit tensoriel de la fibre
par l'anneau local de $X$ en $\overline{x}$.
\end{lemma}

\begin{proof}
Comme $f^*\mathcal{G} =
f_{small}^{-1}\mathcal{G} \otimes_{f_{small}^{-1}\mathcal{O}_Y} \mathcal{O}_X$
cela résulte de la description des fibres des images inverses donnée au
lemme \ref{spaces-properties-lemma-stalk-pullback}
et du fait que la prise des fibres commute aux produits tensoriels.
Voici une démonstration plus directe.
Choisissons un diagramme commutatif
$$
\xymatrix{
U \ar[d]_p \ar[r]_\alpha  & V \ar[d]^q \\
X \ar[r]^f & Y
}
$$
où $U$ et $V$ sont des schémas, et où $p$ et $q$ sont \'etales surjectifs.
D'après le
lemme \ref{spaces-properties-lemma-geometric-lift-to-usual},
nous pouvons choisir un point géométrique $\overline{u}$ de $U$ tel que
$\overline{x} = p \circ \overline{u}$. Posons
$\overline{v} = \alpha \circ \overline{u}$.
Nous voyons alors que
\begin{align*}
(f^*\mathcal{G})_{\overline{x}} & =
(p^*f^*\mathcal{G})_u \otimes_{\mathcal{O}_{U, u}}
\mathcal{O}_{X, \overline{x}} \\
& = (\alpha^*q^*\mathcal{G})_u \otimes_{\mathcal{O}_{U, u}}
\mathcal{O}_{X, \overline{x}} \\
& = (q^*\mathcal{G})_v \otimes_{\mathcal{O}_{V, v}}
\mathcal{O}_{U, u} \otimes_{\mathcal{O}_{U, u}}
\mathcal{O}_{X, \overline{x}} \\
& = (q^*\mathcal{G})_v \otimes_{\mathcal{O}_{V, v}}
\mathcal{O}_{X, \overline{x}} \\
& = (q^*\mathcal{G})_v \otimes_{\mathcal{O}_{V, v}}
\mathcal{O}_{Y, \overline{y}} \otimes_{\mathcal{O}_{Y, \overline{y}}}
\mathcal{O}_{X, \overline{x}} \\
& = \mathcal{G}_{\overline{y}} \otimes_{\mathcal{O}_{Y, \overline{y}}}
\mathcal{O}_{X, \overline{x}}
\end{align*}
Nous avons utilisé ici
le lemme \ref{spaces-properties-lemma-stalk-quasi-coherent} (deux fois)
et le résultat correspondant pour les images inverses de faisceaux quasi-cohérents
sur les schémas ; voir le
lemme \ref{sheaves-lemma-stalk-pullback-modules} de Faisceaux.
\end{proof}

\begin{lemma}
\label{spaces-properties-lemma-characterize-quasi-coherent}
Soit $S$ un schéma. Soit $X$ un espace algébrique sur $S$.
Soit $\mathcal{F}$ un faisceau de $\mathcal{O}_X$-modules.
Les conditions suivantes sont équivalentes :
\begin{enumerate}
\item $\mathcal{F}$ est un $\mathcal{O}_X$-module quasi-cohérent ;
\item il existe un morphisme \'etale $f : Y \to X$
d'espaces algébriques sur $S$, avec $|f| : |Y| \to |X|$ surjectif,
tel que $f^*\mathcal{F}$ soit quasi-cohérent sur $Y$ ;
\item il existe un schéma $U$ et un morphisme \'etale surjectif
$\varphi : U \to X$ tel que $\varphi^*\mathcal{F}$ soit un
$\mathcal{O}_U$-module quasi-cohérent ;
\item pour tout schéma affine $U$ et tout morphisme \'etale $\varphi : U \to X$, la
restriction $\varphi^*\mathcal{F}$ est un $\mathcal{O}_U$-module quasi-cohérent.
\end{enumerate}
\end{lemma}

\begin{proof}
Il est clair que (1) implique (2) en considérant $\text{id}_X$.
Supposons que $f : Y \to X$ soit comme dans (2), et soit $V \to Y$ un morphisme
\'etale surjectif d'un schéma vers $Y$. Alors la composée $V \to X$ est elle aussi
\'etale surjective
et, d'après le lemme \ref{spaces-properties-lemma-pullback-quasi-coherent}, l'image inverse de $\mathcal{F}$
sur $V$ est elle aussi quasi-cohérente. Ainsi, (2) implique (3).

\medskip\noindent
Soit $U \to X$ comme dans (3). Employons l'abus de notation introduit
dans l'équation (\ref{spaces-properties-equation-restrict-modules}).
Comme $\mathcal{F}|_{U_\etale}$ est quasi-cohérent, il existe un
recouvrement \'etale $\{U_i \to U\}$ tel que
$\mathcal{F}|_{U_{i, \etale}}$ admette une présentation globale ; voir la
définition \ref{sites-modules-definition-global} de Modules sur les sites et le
lemme \ref{sites-modules-lemma-local-final-object}.
Soit $V \to X$ un objet de $X_\etale$. Puisque $U \to X$ est
surjectif et \'etale, la famille de morphismes $\{U_i \times_X V \to V\}$ est un
recouvrement \'etale
de $V$. Par les morphismes $U_i \times_X V \to U_i$, nous pouvons restreindre les
présentations globales de $\mathcal{F}|_{U_{i, \etale}}$ pour obtenir une présentation
globale de $\mathcal{F}|_{(U_i \times_X V)_\etale}$.
Le faisceau $\mathcal{F}$ sur $X_\etale$ satisfait donc la condition de la
définition \ref{sites-modules-definition-site-local} de Modules sur les sites
et est par conséquent quasi-cohérent.

\medskip\noindent
L'équivalence de (3) et (4) résulte du fait que tout schéma possède
un recouvrement par des ouverts affines.
\end{proof}

\begin{lemma}
\label{spaces-properties-lemma-properties-quasi-coherent}
Soit $S$ un schéma. Soit $X$ un espace algébrique sur $S$.
La catégorie $\QCoh(\mathcal{O}_X)$ des faisceaux quasi-cohérents sur $X$
possède les propriétés suivantes :
\begin{enumerate}
\item Toute somme directe de faisceaux quasi-cohérents est quasi-cohérente.
\item Toute colimite de faisceaux quasi-cohérents est quasi-cohérente.
\item Le noyau et le conoyau d'un morphisme de faisceaux quasi-cohérents
sont quasi-cohérents.
\item Dans une suite exacte courte de $\mathcal{O}_X$-modules
$0 \to \mathcal{F}_1 \to \mathcal{F}_2 \to \mathcal{F}_3 \to 0$
si deux termes sur trois sont quasi-cohérents, le troisième l'est aussi.
\item Le produit tensoriel de deux $\mathcal{O}_X$-modules quasi-cohérents
est quasi-cohérent.
\item Étant donnés deux $\mathcal{O}_X$-modules quasi-cohérents
$\mathcal{F}$, $\mathcal{G}$, dont $\mathcal{F}$
est de présentation finie (voir la
section \ref{spaces-properties-section-properties-modules}),
le Hom interne
$\SheafHom_{\mathcal{O}_X}(\mathcal{F}, \mathcal{G})$
est quasi-cohérent.
\end{enumerate}
\end{lemma}

\begin{proof}
Si $X$ est un schéma, c'est le
lemme \ref{descent-lemma-properties-quasi-coherent} de Descente.
Nous allons réduire le lemme à ce cas par localisation \'etale.

\medskip\noindent
Choisissons un schéma $U$ et un morphisme \'etale surjectif $\varphi : U \to X$.
Nous adopterons les notations $\textit{Mod}(\mathcal{O}_U) =
\textit{Mod}(U_\etale, \mathcal{O}_U)$ et
$\QCoh(\mathcal{O}_U) = \QCoh(U_\etale, \mathcal{O}_U)$ ; autrement
dit, bien que $U$ soit un schéma, nous considérons les modules quasi-cohérents
sur $U$ comme des modules sur le petit site \'etale de $U$.
D'après le lemme \ref{spaces-properties-lemma-pullback-quasi-coherent}, nous avons un diagramme commutatif
$$
\xymatrix{
\QCoh(\mathcal{O}_X) \ar[r]_{\varphi^*} \ar[d] &
\QCoh(\mathcal{O}_U) \ar[d] \\
\textit{Mod}(\mathcal{O}_X) \ar[r]^{\varphi^*} &
\textit{Mod}(\mathcal{O}_U)
}
$$
La flèche horizontale inférieure est le foncteur de restriction
(\ref{spaces-properties-equation-restrict-modules})
$\mathcal{G} \mapsto \mathcal{G}|_{U_\etale}$.
Ce foncteur possède à la fois un adjoint à gauche et un adjoint à droite ; voir la
section \ref{sites-modules-section-localize} de Modules sur les sites ;
il commute donc à toutes les limites et colimites.
De plus, nous savons qu'un objet de $\textit{Mod}(\mathcal{O}_X)$ appartient à
$\QCoh(\mathcal{O}_X)$ si et seulement si sa restriction à $U$ appartient à
$\QCoh(\mathcal{O}_U)$ ; voir le
lemme \ref{spaces-properties-lemma-characterize-quasi-coherent}.
Ces préliminaires étant établis, nous pouvons commencer la preuve.

\medskip\noindent
Preuve de (1). Soit $\mathcal{F}_i$, $i \in I$, une famille de
$\mathcal{O}_X$-modules quasi-cohérents. D'après la discussion précédente, nous avons
$$
\Big(\bigoplus \mathcal{F}_i\Big)|_{U_\etale} =
\bigoplus \mathcal{F}_i|_{U_\etale}
$$
Chacun des modules $\mathcal{F}_i|_{U_\etale}$ est quasi-cohérent.
La somme directe est donc quasi-cohérente d'après le cas des schémas.
Par conséquent, $\bigoplus \mathcal{F}_i$ est quasi-cohérent, puisque sa
restriction à $U$ est un module quasi-cohérent.

\medskip\noindent
Preuve de (2). Soit $\mathcal{I} \to \QCoh(\mathcal{O}_X)$,
$i \mapsto \mathcal{F}_i$, un diagramme. Alors
$$
(\colim \mathcal{F}_i)|_{U_\etale} = \colim \mathcal{F}_i|_{U_\etale}
$$
d'après la discussion précédente, et nous concluons de la même manière.

\medskip\noindent
Preuve de (3). Soit $a : \mathcal{F} \to \mathcal{F}'$
une flèche de $\QCoh(\mathcal{O}_X)$. Alors
nous avons $\Ker(a)|_{U_\etale} = \Ker(a|_{U_\etale})$
et $\Coker(a)|_{U_\etale} = \Coker(a|_{U_\etale})$,
et nous concluons de la même manière.

\medskip\noindent
Preuve de (4). La restriction
$0 \to \mathcal{F}_1|_{U_\etale} \to \mathcal{F}_2|_{U_\etale}
\to \mathcal{F}_3|_{U_\etale} \to 0$
est une suite exacte courte. Elle possède donc la propriété des deux sur trois,
et nous concluons comme précédemment.

\medskip\noindent
Preuve de (5). Soient $\mathcal{F}$ et $\mathcal{G}$ dans $\QCoh(\mathcal{O}_X)$.
Alors nous avons
$$
(\mathcal{F} \otimes_{\mathcal{O}_X} \mathcal{G})_{U_\etale} =
\mathcal{F}|_{U_\etale} \otimes_{\mathcal{O}_U} \mathcal{G}|_{U_\etale}
$$
et nous concluons comme précédemment.

\medskip\noindent
Preuve de (6). Soient $\mathcal{F}$ et $\mathcal{G}$
dans $\QCoh(\mathcal{O}_X)$, avec $\mathcal{F}$
de présentation finie. Nous avons
$$
\SheafHom_{\mathcal{O}_X}(\mathcal{F}, \mathcal{G})|_{U_\etale} =
\SheafHom_{\mathcal{O}_U}(\mathcal{F}|_{U_\etale}, \mathcal{G}|_{U_\etale})
$$
En effet, la restriction est une localisation ; voir la
section \ref{spaces-properties-section-localize}, en particulier la formule
(\ref{spaces-properties-equation-localize-at-scheme-ringed}), et la formation du Hom
interne commute à la localisation ; voir le
lemme \ref{sites-modules-lemma-internal-hom-restriction} de Modules sur les sites.
Nous concluons donc comme précédemment.
\end{proof}

\noindent
En général, l'image directe d'un faisceau quasi-cohérent
par un morphisme d'espaces algébriques n'est pas quasi-cohérente. Nous reviendrons sur cette
question dans
Morphismes d'espaces, section \ref{spaces-morphisms-section-pushforward}.




\section{Propriétés des modules}
\label{spaces-properties-section-properties-modules}

\noindent
Dans
Modules sur les sites, les sections
\ref{sites-modules-section-global},
\ref{sites-modules-section-local} et la
définition \ref{sites-modules-definition-flat},
nous avons défini plusieurs propriétés intrinsèques des
$\mathcal{O}$-modules sur tout topos annelé. Si $X$ est un espace
algébrique, nous appliquerons librement ces notions aux modules sur le site
annelé $(X_\etale, \mathcal{O}_X)$, ou, de manière équivalente, sur le site annelé
$(X_{spaces, \etale}, \mathcal{O}_X)$.

\medskip\noindent
Propriétés globales $\mathcal{P}$ :
\begin{enumerate}
\item[(a)] {\it libre} ;
\item[(b)] {\it libre de rang fini} ;
\item[(c)] {\it engendré par ses sections globales} ;
\item[(d)] {\it engendré par un nombre fini de sections globales} ;
\item[(e)] admettre une {\it présentation globale} ;
\item[(f)] admettre une {\it présentation globale finie}.
\end{enumerate}
Propriétés locales $\mathcal{P}$ :
\begin{enumerate}
\item[(g)] {\it localement libre} ;
\item[(h)] {\it localement libre de rang fini} ;
\item[(i)] {\it localement engendré par des sections} ;
\item[(j)] {\it localement engendré par $r$ sections} ;
\item[(k)] {\it de type fini} ;
\item[(l)] {\it quasi-cohérent} (voir la section \ref{spaces-properties-section-quasi-coherent}) ;
\item[(m)] {\it de présentation finie} ;
\item[(n)] {\it cohérent} ;
\item[(o)] {\it plat}.
\end{enumerate}
Voici quelques résultats qui découlent immédiatement des définitions :
\begin{enumerate}
\item Dans chaque cas, sauf pour $\mathcal{P}=$``cohérent'', la propriété
est préservée par image inverse ; voir les
lemmes \ref{sites-modules-lemma-global-pullback},
\ref{sites-modules-lemma-local-pullback} et
\ref{sites-modules-lemma-pullback-flat}.
\item Chacune des propriétés ci-dessus (y compris la cohérence) est préservée par
image inverse le long d'un morphisme \'etale d'espaces algébriques (car, dans ce
cas, l'image inverse est donnée par restriction ; voir le
lemme \ref{spaces-properties-lemma-etale-morphism-topoi}).
\item Supposons que $f : Y \to X$ soit un morphisme \'etale surjectif d'espaces
algébriques. Pour chacune des propriétés locales (g) -- (o), si
$f^*\mathcal{F}$ possède $\mathcal{P}$, alors $\mathcal{F}$ possède
$\mathcal{P}$. Cela résulte du fait que $\{Y \to X\}$ est un recouvrement de
$X_{spaces, \etale}$ et du
lemme \ref{sites-modules-lemma-local-final-object} de Modules sur les sites.
\item Si $X$ est un schéma, si $\mathcal{F}$ est un module quasi-cohérent
sur $X_\etale$ et si $\mathcal{P}$ est une propriété autre que ``cohérent'' ou
``localement libre'', alors $\mathcal{P}$ pour $\mathcal{F}$ sur $X_\etale$
équivaut à la propriété correspondante pour
$\mathcal{F}|_{X_{Zar}}$ ; autrement dit, elle correspond à $\mathcal{P}$
pour $\mathcal{F}$ considéré comme faisceau quasi-cohérent
sur le schéma $X$. Voir le
lemme \ref{descent-lemma-equivalence-quasi-coherent-properties} de Descente.
\item Si $X$ est un schéma localement noethérien et si $\mathcal{F}$ est un
module quasi-cohérent sur $X_\etale$, alors $\mathcal{F}$
est cohérent sur $X_\etale$ si et seulement si
$\mathcal{F}|_{X_{Zar}}$ est cohérent ; autrement dit, cela équivaut à ce que
le faisceau quasi-cohérent usuel sur le schéma $X$ soit
cohérent. Voir le
lemme \ref{descent-lemma-equivalence-quasi-coherent-properties} de Descente.
\end{enumerate}



\section{Modules localement projectifs}
\label{spaces-properties-section-locally-projective}

\noindent
Rappelons que, dans
Propriétés, section \ref{properties-section-locally-projective},
nous avons défini la notion de module quasi-cohérent
localement projectif.

\begin{lemma}
\label{spaces-properties-lemma-locally-projective}
Soit $S$ un schéma. Soit $X$ un espace algébrique sur $S$.
Soit $\mathcal{F}$ un $\mathcal{O}_X$-module quasi-cohérent.
Les conditions suivantes sont équivalentes :
\begin{enumerate}
\item il existe un schéma $U$ et un morphisme \'etale surjectif
$U \to X$ tels que la restriction $\mathcal{F}|_U$ soit localement projective
sur $U$ ;
\item pour tout schéma $U$ et tout morphisme \'etale
$U \to X$, la restriction $\mathcal{F}|_U$ est localement projective
sur $U$.
\end{enumerate}
\end{lemma}

\begin{proof}
Soit $U \to X$ comme dans (1), et soit $V \to X$ \'etale, où
$V$ est un schéma. Alors $\{U \times_X V \to V\}$ est un recouvrement fppf
de schémas. Ainsi, si $\mathcal{F}|_U$ est localement projectif, alors
$\mathcal{F}|_{U \times_X V}$ est localement projectif (voir le
lemme \ref{properties-lemma-locally-projective-pullback} de Propriétés)
et $\mathcal{F}|_V$ est donc localement projectif ; voir le
lemme \ref{descent-lemma-locally-projective-descends} de Descente.
\end{proof}

\begin{definition}
\label{spaces-properties-definition-locally-projective}
Soit $S$ un schéma. Soit $X$ un espace algébrique sur $S$.
Soit $\mathcal{F}$ un $\mathcal{O}_X$-module quasi-cohérent.
On dit que $\mathcal{F}$ est {\it localement projectif}
si les conditions équivalentes du
lemme \ref{spaces-properties-lemma-locally-projective}
sont satisfaites.
\end{definition}

\begin{lemma}
\label{spaces-properties-lemma-locally-projective-pullback}
Soit $S$ un schéma.
Soit $f : X \to Y$ un morphisme d'espaces algébriques sur $S$.
Soit $\mathcal{G}$ un $\mathcal{O}_Y$-module quasi-cohérent.
Si $\mathcal{G}$ est localement projectif sur $Y$, alors $f^*\mathcal{G}$
est localement projectif sur $X$.
\end{lemma}

\begin{proof}
Choisissons un morphisme \'etale surjectif $V \to Y$, où $V$ est un schéma.
Choisissons un morphisme \'etale surjectif $U \to V \times_Y X$, où
$U$ est un schéma. Notons $\psi : U \to V$ le morphisme induit.
Alors
$$
f^*\mathcal{G}|_U = \psi^*(\mathcal{G}|_V)
$$
Le lemme résulte donc de la définition et du résultat dans le
cas des schémas ; voir le
lemme \ref{properties-lemma-locally-projective-pullback} de Propriétés.
\end{proof}





\section{Faisceaux quasi-cohérents et présentations}
\label{spaces-properties-section-quasi-coherent-presentation}

\noindent
Soit $S$ un schéma. Soit $X$ un espace algébrique sur $S$.
Soit $X = U/R$ une présentation de $X$ provenant d'un morphisme
\'etale surjectif quelconque $\varphi : U \to X$ ; voir la
définition \ref{spaces-definition-presentation} d'Espaces.
Nous obtenons en particulier un groupoïde $(U, R, s, t, c)$ tel que
$j = (t, s) : R \to U \times_S U$ ; voir le
lemme \ref{groupoids-lemma-equivalence-groupoid} de Groupoïdes.
Dans la
définition \ref{groupoids-definition-groupoid-module} de Groupoïdes,
nous avons défini la notion de faisceau quasi-cohérent
sur un groupoïde arbitraire. Ces notions étant posées, nous avons
l'observation suivante.

\begin{proposition}
\label{spaces-properties-proposition-quasi-coherent}
Conservons $S$, $\varphi : U \to X$ et $(U, R, s, t, c)$ comme ci-dessus.
Pour tout $\mathcal{O}_X$-module quasi-cohérent $\mathcal{F}$, le
faisceau $\varphi^*\mathcal{F}$ est muni d'un
isomorphisme canonique
$$
\alpha : t^*\varphi^*\mathcal{F} \longrightarrow s^*\varphi^*\mathcal{F}
$$
qui satisfait aux conditions de la
définition \ref{groupoids-definition-groupoid-module} de Groupoïdes
et définit donc un faisceau quasi-cohérent sur $(U, R, s, t, c)$.
Le foncteur $\mathcal{F} \mapsto (\varphi^*\mathcal{F}, \alpha)$
définit une équivalence de catégories
$$
\begin{matrix}
\text{Catégorie des} \\
\mathcal{O}_X\text{-Modules quasi-cohérents}
\end{matrix}
\longleftrightarrow
\begin{matrix}
\text{Catégorie des Modules}\\
\text{quasi-cohérents sur }(U, R, s, t, c)
\end{matrix}
$$
\end{proposition}

\begin{proof}
Dans l'énoncé de la proposition et dans cette preuve, nous considérons un
faisceau quasi-cohérent sur un schéma comme un faisceau quasi-cohérent sur le petit
site \'etale de ce schéma. Cela est permis par les résultats de
Descente, sections \ref{descent-section-quasi-coherent-sheaves},
\ref{descent-section-quasi-coherent-cohomology} et
\ref{descent-section-quasi-coherent-sheaves-bis}.

\medskip\noindent
L'existence de $\alpha$ résulte du fait que
$\varphi \circ t = \varphi \circ s$ et de la fonctorialité de l'image inverse
par rapport au morphisme ; voir la discussion autour de
l'équation (\ref{spaces-properties-equation-push-pull}). De la même manière, c'est-à-dire par
fonctorialité de l'image inverse, nous voyons que l'isomorphisme $\alpha$ satisfait
à la condition (1) de la
définition \ref{groupoids-definition-groupoid-module} de Groupoïdes.
Pour la condition (2) de la définition, il suffit de constater que $\alpha$
est un isomorphisme, ce qui est clair. La construction
$\mathcal{F} \mapsto (\varphi^*\mathcal{F}, \alpha)$
est manifestement fonctorielle par rapport au faisceau quasi-cohérent $\mathcal{F}$.
Nous obtenons donc le foncteur de gauche à droite dans la formule
affichée de la proposition.

\medskip\noindent
Réciproquement, supposons que $(\mathcal{F}, \alpha)$ soit un faisceau
quasi-cohérent sur $(U, R, s, t, c)$. Soit $V \to X$ un objet de $X_\etale$.
Alors le morphisme $V' = U \times_X V \to V$ est un morphisme \'etale
surjectif de schémas, et $\{V' \to V\}$ est donc un recouvrement \'etale
de $V$. En outre, l'image inverse du faisceau quasi-cohérent $\mathcal{F}$
est un faisceau quasi-cohérent $\mathcal{F}'$ sur $V'$.
Comme $R = U \times_X U$, avec $t = \text{pr}_0$ et $s = \text{pr}_1$,
nous avons $V' \times_V V' = R \times_X V$, les morphismes de projection
$V' \times_V V' \to V'$ étant obtenus par changement de base à partir de $t$ et $s$. Ainsi,
l'image inverse de $\alpha$ est un isomorphisme
$\alpha' : \text{pr}_0^*\mathcal{F}' \to \text{pr}_1^*\mathcal{F}'$, et
le couple $(\mathcal{F}', \alpha')$ est une donnée de descente pour les faisceaux
quasi-cohérents relativement à $\{V' \to V\}$. D'après la
proposition
\ref{descent-proposition-fpqc-descent-quasi-coherent} de Descente,
cette donnée de descente est effective, et nous obtenons un
$\mathcal{O}_V$-module quasi-cohérent $\mathcal{F}_V$ sur $V_\etale$.
Pour voir que cela définit un faisceau quasi-cohérent sur $X_\etale$, il faut
montrer (d'après le
lemme \ref{spaces-properties-lemma-characterize-quasi-coherent-small-etale})
que, pour tout morphisme $f : V_1 \to V_2$ dans $X_\etale$,
il existe un isomorphisme canonique
$c_f : f^*\mathcal{F}_{V_2} \to \mathcal{F}_{V_1}$
compatible aux compositions de morphismes. Nous omettons la vérification.
Nous omettons aussi de vérifier que cela définit un foncteur de la
catégorie de droite vers celle de gauche, inverse
du foncteur décrit ci-dessus.
\end{proof}

\begin{proposition}
\label{spaces-properties-proposition-coherator}
Soit $S$ un schéma. Soit $X$ un espace algébrique sur $S$.
\begin{enumerate}
\item La catégorie $\QCoh(\mathcal{O}_X)$ est une catégorie
abélienne de Grothendieck. Par conséquent, $\QCoh(\mathcal{O}_X)$
possède suffisamment d'injectifs et toutes les limites.
\item Le foncteur d'inclusion
$\QCoh(\mathcal{O}_X) \to \textit{Mod}(\mathcal{O}_X)$
possède un adjoint à droite\footnote{Ce foncteur est parfois appelé
le {\it cohérateur}.}
$$
Q : \textit{Mod}(\mathcal{O}_X) \longrightarrow \QCoh(\mathcal{O}_X)
$$
tel que, pour tout faisceau quasi-cohérent $\mathcal{F}$, le morphisme d'adjonction
$Q(\mathcal{F}) \to \mathcal{F}$ soit un isomorphisme.
\end{enumerate}
\end{proposition}

\begin{proof}
Cette preuve reprend celle du cas des schémas ; voir la
proposition \ref{properties-proposition-coherator} de Propriétés.
Nous conseillons au lecteur de lire d'abord cette preuve.

\medskip\noindent
L'assertion (1) signifie que $\QCoh(\mathcal{O}_X)$ (a) admet toutes les colimites,
(b) que les colimites filtrantes y sont exactes et (c) qu'elle admet un générateur ; voir la
section \ref{injectives-section-grothendieck-conditions} d'Injectifs.
D'après le lemme \ref{spaces-properties-lemma-properties-quasi-coherent},
les colimites dans $\QCoh(\mathcal{O}_X)$ existent et coïncident
avec celles de $\textit{Mod}(\mathcal{O}_X)$. D'après le
lemme \ref{sites-modules-lemma-limits-colimits} de Modules sur les sites,
les colimites filtrantes sont exactes. Les conditions (a) et (b) sont donc satisfaites.

\medskip\noindent
Pour construire un générateur, choisissons une présentation $X = U/R$ telle que
$(U, R, s, t, c)$ soit un
groupoïde en schémas \'etale ; en particulier, $s$ et $t$ sont des morphismes plats
de schémas. Choisissons un cardinal $\kappa$ comme dans le
lemme \ref{groupoids-lemma-colimit-kappa} de Groupoïdes.
Choisissons une famille $(\mathcal{E}_t, \alpha_t)_{t \in T}$ de
modules quasi-cohérents $\kappa$-engendrés sur
$(U, R, s, t, c)$ comme dans le
lemme \ref{groupoids-lemma-set-of-iso-classes} de Groupoïdes.
Soit $\mathcal{F}_t$ le module quasi-cohérent sur $X$ qui
correspond au module quasi-cohérent $(\mathcal{E}_t, \alpha_t)$ par
l'équivalence de catégories de la
proposition \ref{spaces-properties-proposition-quasi-coherent}.
Nous voyons alors que tout module quasi-cohérent $\mathcal{H}$ est la
colimite filtrante de ses sous-modules quasi-cohérents qui sont isomorphes
à l'un des $\mathcal{F}_t$. Ainsi, $\bigoplus_t \mathcal{F}_t$ est
un générateur de $\QCoh(\mathcal{O}_X)$, ce qui établit (c).
Les assertions sur les limites et les injectifs valent dans toute catégorie
abélienne de Grothendieck ; voir le
théorème
\ref{injectives-theorem-injective-embedding-grothendieck} et le
lemme \ref{injectives-lemma-grothendieck-products} d'Injectifs.

\medskip\noindent
Preuve de (2). Pour construire $Q$, nous employons le procédé général suivant.
Étant donné un objet $\mathcal{F}$ de $\textit{Mod}(\mathcal{O}_X)$,
nous considérons le foncteur
$$
\QCoh(\mathcal{O}_X)^{opp} \longrightarrow \textit{Ens},\quad
\mathcal{G} \longmapsto \Hom_X(\mathcal{G}, \mathcal{F})
$$
Ce foncteur transforme les colimites en limites ;
il est donc représentable ; voir le
lemme \ref{injectives-lemma-grothendieck-brown} d'Injectifs.
Il existe ainsi un faisceau quasi-cohérent $Q(\mathcal{F})$
et un isomorphisme fonctoriel
$\Hom_X(\mathcal{G}, \mathcal{F}) = \Hom_X(\mathcal{G}, Q(\mathcal{F}))$
pour $\mathcal{G}$ dans $\QCoh(\mathcal{O}_X)$. D'après le lemme de Yoneda
(lemme \ref{categories-lemma-yoneda} de Catégories),
la construction $\mathcal{F} \leadsto Q(\mathcal{F})$ est
fonctorielle en $\mathcal{F}$. Par construction, $Q$ est un adjoint à
droite du foncteur d'inclusion.
Le fait que $Q(\mathcal{F}) \to \mathcal{F}$ soit un isomorphisme
lorsque $\mathcal{F}$ est quasi-cohérent résulte formellement du fait
que le foncteur d'inclusion
$\QCoh(\mathcal{O}_X) \to \textit{Mod}(\mathcal{O}_X)$
est pleinement fidèle.
\end{proof}






\section{Morphismes à valeurs dans des schémas}
\label{spaces-properties-section-morphisms-to-schemes}

\noindent
Voici l'analogue du
lemme \ref{schemes-lemma-morphism-into-affine} de Schémas.

\begin{lemma}
\label{spaces-properties-lemma-morphism-to-affine-scheme}
Soit $X$ un espace algébrique sur $\mathbf{Z}$.
Soit $T$ un schéma affine.
L'application
$$
\Mor(X, T)
\longrightarrow
\Hom(\Gamma(T, \mathcal{O}_T), \Gamma(X, \mathcal{O}_X))
$$
qui envoie $f$ sur $f^\sharp$ (sur les sections globales) est bijective.
\end{lemma}

\begin{proof}
Nous construisons l'application inverse.
Soit $\varphi : \Gamma(T, \mathcal{O}_T) \to \Gamma(X, \mathcal{O}_X)$
un homomorphisme d'anneaux. Choisissons une présentation $X = U/R$ ; voir la
définition \ref{spaces-definition-presentation} d'Espaces.
D'après le
lemme \ref{schemes-lemma-morphism-into-affine} de Schémas,
la composée
$$
\Gamma(T, \mathcal{O}_T) \to \Gamma(X, \mathcal{O}_X) \to
\Gamma(U, \mathcal{O}_U)
$$
correspond à un unique morphisme de schémas $g : U \to T$. Par le même lemme,
les deux composées $R \to U \to T$ sont égales. Nous obtenons donc un morphisme
$f : X = U/R \to T$ tel que $U \to X \to T$ soit égal à $g$. Par construction,
le diagramme
$$
\xymatrix{
\Gamma(U, \mathcal{O}_U) & \Gamma(X, \mathcal{O}_X) \ar[l] \\
& \Gamma(T, \mathcal{O}_T) \ar[lu]^{g^\sharp} \ar[u]^{\varphi}_{f^\sharp}
}
$$
est commutatif. Ainsi, $f^\sharp$ est égal à $\varphi$, car $U \to X$ est un
recouvrement \'etale et $\mathcal{O}_X$ est un faisceau sur $X_\etale$.
L'unicité de $f$ résulte de celle de $g$.
\end{proof}





\section{Quotients par des actions libres}
\label{spaces-properties-section-quotient-by-free}

\noindent
Soit $S$ un schéma.
Soit $X$ un espace algébrique sur $S$.
Soit $G$ un groupe abstrait.
Soit $a : G \to \text{Aut}(X)$ un homomorphisme, c'est-à-dire que $a$ est une
{\it action} de $G$ sur $X$. Nous dirons que l'action est {\it libre}
si, pour tout schéma $T$ sur $S$, l'action
$$
G \times X(T) \longrightarrow X(T)
$$
est libre. (Nous ne pouvons pas employer un critère comme celui du
lemme \ref{spaces-lemma-quotient} d'Espaces,
car les points peuvent ne pas avoir de corps résiduels bien définis.)
Lorsque l'action est libre, nous allons construire le quotient $X/G$
comme espace algébrique. C'est un cas particulier du
lemme général \ref{bootstrap-lemma-quotient-free-action} d'Amorçage
que nous démontrerons plus tard.

\begin{lemma}
\label{spaces-properties-lemma-quotient}
Soit $S$ un schéma.
Soit $X$ un espace algébrique sur $S$.
Soit $G$ un groupe abstrait opérant librement sur $X$.
Alors le faisceau quotient $X/G$ est un espace algébrique.
\end{lemma}

\begin{proof}
L'énoncé signifie que le faisceau $F$ associé au préfaisceau
$$
T \longmapsto X(T)/G
$$
est un espace algébrique. Pour le voir, nous allons construire une présentation.
Choisissons un schéma $U$ et un morphisme \'etale surjectif
$\varphi : U \to X$. Posons $V = \coprod_{g \in G} U$ et définissons
$\psi : V \to X$ comme $a(g) \circ \varphi$ sur la composante correspondant
à $g \in G$. Faisons agir $G$ sur $V$ en permutant les composantes : ainsi,
$g_0 \in G$ envoie la composante correspondant à $g$ sur celle correspondant
à $g_0g$ par le morphisme identité de $U$.
Alors $\psi$ est un morphisme $G$-équivariant ; nous nous ramenons donc au
cas traité dans le paragraphe suivant.

\medskip\noindent
Supposons qu'il existe une action de $G$ sur $U$ et que $U \to X$ soit surjectif,
\'etale et $G$-équivariant. Il existe alors une
action induite de $G$ sur $R = U \times_X U$, compatible avec les morphismes de
projection $t, s : R \to U$. Nous affirmons que
$$
X/G = U/\coprod\nolimits_{g \in G} R
$$
où l'application
$$
j : \coprod\nolimits_{g \in G} R
\longrightarrow
U \times_S U
$$
est donnée par $(r, g) \mapsto (t(r), g(s(r)))$. Remarquons que $j$ est un monomorphisme :
si $(t(r), g(s(r))) = (t(r'), g'(s(r')))$, alors
$t(r) = t(r')$ ; par conséquent, $r$ et $r'$ ont la même image dans $X$ par
$s$ et $t$, puis $g = g'$ (car $G$ opère librement sur $X$), puis
$s(r) = s(r')$, et enfin $r = r'$ (car $R$ est une relation d'équivalence sur $U$).
De plus, $j$ est une relation d'équivalence (détails omis).
Les deux projections $\coprod\nolimits_{g \in G} R \to U$ sont \'etales, puisque
$s$ et $t$ sont \'etales. Ainsi, $j$ est une relation d'équivalence \'etale
et $U/\coprod\nolimits_{g \in G} R$ est un espace algébrique d'après le
théorème \ref{spaces-theorem-presentation} d'Espaces.
Il existe une application
$$
U/\coprod\nolimits_{g \in G} R \longrightarrow X/G
$$
induite par l'application $U \to X$. Nous omettons la preuve qu'il s'agit d'un
isomorphisme de faisceaux.
\end{proof}










% OK, start here.
%

\chapter{Morphismes d'espaces algébriques}



\phantomsection
\label{spaces-morphisms-section-phantom}


\section{Introduction}
\label{spaces-morphisms-section-introduction}

\noindent
Dans ce chapitre, nous introduisons certains types de morphismes d'espaces
algébriques. On pourra consulter \cite{Kn}.

\medskip\noindent
Le but est d'étendre à la catégorie des espaces algébriques la définition
de chacun des types de morphismes de schémas définis dans les chapitres
Schémas et Morphismes de schémas. Chaque cas présente de légères différences
et il nous paraît préférable de les traiter séparément.





\section{Conventions}
\label{spaces-morphisms-section-conventions}

\noindent
Nous supposons une fois pour toutes que tous les schémas appartiennent
à un grand site fppf $\Sch_{fppf}$. Tous les anneaux $A$ considérés
ont en outre la propriété que $\Spec(A)$ est (isomorphe à) un
objet de ce grand site.

\medskip\noindent
Soient $S$ un schéma et $X$ un espace algébrique au-dessus de $S$.
Dans ce chapitre et le suivant, nous noterons $X \times_S X$
le produit de $X$ par lui-même (dans la catégorie des espaces
algébriques au-dessus de $S$), au lieu de $X \times X$.





\section{Propriétés des morphismes représentables}
\label{spaces-morphisms-section-representable}

\noindent
Soit $S$ un schéma.
Soit $f : X \to Y$ un morphisme représentable d'espaces algébriques. Dans
Espaces, section \ref{spaces-section-representable-properties},
nous avons défini ce que signifie pour $f$
posséder la propriété $\mathcal{P}$ lorsque $\mathcal{P}$ est une propriété
des morphismes de schémas qui
\begin{enumerate}
\item est préservée par tout changement de base,
voir Schémas, définition \ref{schemes-definition-preserved-by-base-change},
et
\item est locale sur la base pour la topologie fppf, voir
Descente, définition \ref{descent-definition-property-morphisms-local}.
\end{enumerate}
Plus précisément, dans ce cas, on dit que $f$ possède la propriété $\mathcal{P}$ si et seulement
si, pour tout schéma $U$ et tout morphisme $U \to Y$, le morphisme de schémas
$X \times_Y U \to U$ possède la propriété $\mathcal{P}$.

\medskip\noindent
D'après les listes de
Espaces, section \ref{spaces-section-lists},
ceci s'applique aux propriétés suivantes :
(1)(a) immersion fermée,
(1)(b) immersion ouverte,
(1)(c) immersion quasi-compacte,
(2) quasi-compact,
(3) universellement fermé,
(4) (quasi-)séparé,
(5) monomorphisme,
(6) surjectif,
(7) universellement injectif,
(8) affine,
(9) quasi-affine,
(10) (localement) de type fini,
(11) (localement) quasi-fini,
(12) (localement) de présentation finie,
(13) localement de type fini de dimension relative $d$,
(14) universellement ouvert,
(15) plat,
(16) syntomique,
(17) lisse,
(18) non ramifié (resp.\ G-non ramifié),
(19) \'etale,
(20) propre,
(21) fini ou entier,
(22) fini localement libre,
(23) universellement submersif,
(24) homéomorphisme universel, et
(25) immersion.

\medskip\noindent
Dans ce chapitre, nous redéfinirons ces notions pour les morphismes d'espaces
algébriques qui ne sont pas nécessairement représentables. Chaque fois, nous nous
assurerons que la nouvelle définition coïncide avec l'ancienne, afin d'éviter
toute ambiguïté.

\medskip\noindent
Remarquons que la définition ci-dessus s'applique lorsque $X$ est un schéma,
puisqu'un morphisme d'un schéma vers un espace algébrique est représentable.
Elle s'applique donc en particulier lorsque $X$ et $Y$ sont tous deux des schémas.
Dans
Espaces, lemme
\ref{spaces-lemma-morphism-schemes-gives-representable-transformation-property}
nous avons vu que, dans ce cas, les définitions
coïncident et qu'aucune ambiguïté ne se présente.

\medskip\noindent
En outre, dans
Espaces, lemme
\ref{spaces-lemma-base-change-representable-transformations-property}
nous avons vu que la propriété
ainsi définie pour les morphismes représentables d'espaces algébriques est stable par
tout changement de base par un morphisme d'espaces algébriques.
Enfin, dans
Espaces, lemmes
\ref{spaces-lemma-composition-representable-transformations-property} et
\ref{spaces-lemma-product-representable-transformations-property}
nous avons vu que, si $\mathcal{P}$ est stable par composition,
ce qui est le cas des propriétés
(1)(a), (1)(b), (1)(c), (2) -- (25), à l'exception de (13) ci-dessus, alors
les produits de morphismes représentables possèdent encore la propriété $\mathcal{P}$
et les composées de morphismes représentables possèdent encore la propriété $\mathcal{P}$.

\medskip\noindent
Nous utiliserons ces faits ci-dessous et, lorsque ce sera le cas, nous nous contenterons
de renvoyer à la présente section.















\section{Axiomes de séparation}
\label{spaces-morphisms-section-separation-axioms}

\noindent
Il est naturel d'énumérer quelques propriétés a priori de la diagonale
d'un morphisme d'espaces algébriques.

\begin{lemma}
\label{spaces-morphisms-lemma-properties-diagonal}
Soit $S$ un schéma appartenant à $\Sch_{fppf}$.
Soit $f : X \to Y$ un morphisme d'espaces algébriques au-dessus de $S$.
Soit $\Delta_{X/Y} : X \to X \times_Y X$ le morphisme diagonal.
Alors
\begin{enumerate}
\item $\Delta_{X/Y}$ est représentable,
\item $\Delta_{X/Y}$ est localement de type fini,
\item $\Delta_{X/Y}$ est un monomorphisme,
\item $\Delta_{X/Y}$ est séparé, et
\item $\Delta_{X/Y}$ est localement quasi-fini.
\end{enumerate}
\end{lemma}

\begin{proof}
Nous allons utiliser le fait que $\Delta_{X/S}$ est
représentable (par définition d'un espace algébrique) et qu'il
possède les propriétés (2) -- (5), voir
Espaces, lemme \ref{spaces-lemma-properties-diagonal}.
Remarquons que nous avons une factorisation
$$
X
\longrightarrow
X \times_Y X
\longrightarrow
X \times_S X
$$
de la diagonale $\Delta_{X/S} : X \to X \times_S X$. Puisque
$X \times_Y X \to X \times_S X$ est un monomorphisme et que
$\Delta_{X/S}$ est représentable, il s'ensuit formellement que
$\Delta_{X/Y}$ est représentable. En particulier, les autres
assertions ont désormais un sens, voir
la section \ref{spaces-morphisms-section-representable}.

\medskip\noindent
Choisissons un morphisme \'etale surjectif $U \to X$, où $U$ est un schéma.
Considérons le diagramme
$$
\xymatrix{
R = U \times_X U \ar[r] \ar[d] &
U \times_Y U \ar[d] \ar[r] &
U \times_S U \ar[d] \\
X \ar[r] & X \times_Y X \ar[r] & X \times_S X
}
$$
Les deux carrés sont cartésiens; le rectangle extérieur l'est donc aussi.
La ligne supérieure est formée de schémas et les flèches verticales
sont des morphismes \'etales surjectifs. D'après
Espaces, lemme \ref{spaces-lemma-representable-morphisms-spaces-property},
les propriétés (2) -- (5) de $\Delta_{X/Y}$ sont équivalentes à celles de
$R \to U \times_Y U$. Dans la démonstration de
Espaces, lemme \ref{spaces-lemma-properties-diagonal},
nous avons vu que $R \to U \times_S U$ possède les propriétés (2) -- (5).
Le morphisme $U \times_Y U \to U \times_S U$ est un monomorphisme
de schémas. Il en résulte que $R \to U \times_Y U$ possède les
propriétés (2) -- (5).

\medskip\noindent
En effet, pour (3), remarquons que $R \to U \times_Y U$
est un monomorphisme, puisque la composée
$R \to U \times_S U$ en est un. Pour (2), remarquons que
$R \to U \times_Y U$ est localement de type fini, car la
composée $R \to U \times_S U$ est localement de type fini
(Morphismes, lemme \ref{morphisms-lemma-permanence-finite-type}).
Un monomorphisme localement de type fini est localement quasi-fini
parce qu'il a des fibres finies
(Morphismes, lemme \ref{morphisms-lemma-finite-fibre}); d'où (5).
Un monomorphisme est séparé
(Schémas, lemme \ref{schemes-lemma-monomorphism-separated}); d'où (4).
\end{proof}

\begin{definition}
\label{spaces-morphisms-definition-separated}
Soit $S$ un schéma.
Soit $f : X \to Y$ un morphisme d'espaces algébriques au-dessus de $S$.
Soit $\Delta_{X/Y} : X \to X \times_Y X$ le morphisme diagonal.
\begin{enumerate}
\item On dit que $f$ est {\it séparé} si $\Delta_{X/Y}$ est une immersion fermée.
\item On dit que $f$ est {\it localement séparé}\footnote{Dans la littérature,
ce terme désigne souvent les morphismes quasi-séparés et localement séparés.}
si $\Delta_{X/Y}$ est une immersion.
\item On dit que $f$ est {\it quasi-séparé} si $\Delta_{X/Y}$ est quasi-compact.
\end{enumerate}
\end{definition}

\noindent
Cette définition a un sens puisque $\Delta_{X/Y}$ est représentable;
nous savons donc ce que signifie pour lui posséder l'une des propriétés
décrites dans la définition. Nous verrons plus bas
(lemme \ref{spaces-morphisms-lemma-match-separated}) que cette définition coïncide avec celles
que nous avons déjà pour les morphismes de schémas et les morphismes représentables.

\begin{lemma}
\label{spaces-morphisms-lemma-trivial-implications}
Soit $S$ un schéma. Soit $f : X \to Y$ un morphisme d'espaces algébriques
au-dessus de $S$. Si $f$ est séparé, alors $f$ est localement séparé et
$f$ est quasi-séparé.
\end{lemma}

\begin{proof}
Cela résulte du principe général de
Espaces,
lemme \ref{spaces-lemma-representable-transformations-property-implication},
puisqu'une immersion fermée de schémas est une immersion et est quasi-compacte.
\end{proof}

\begin{lemma}
\label{spaces-morphisms-lemma-base-change-separated}
Tous les axiomes de séparation énumérés dans la définition \ref{spaces-morphisms-definition-separated}
sont stables par changement de base.
\end{lemma}

\begin{proof}
Soient $f : X \to Y$ et $Y' \to Y$ des morphismes d'espaces algébriques.
Soit $f' : X' \to Y'$ le changement de base de $f$ par $Y' \to Y$. Alors
$\Delta_{X'/Y'}$ est le changement de base de $\Delta_{X/Y}$ par
le morphisme $X' \times_{Y'} X' \to X \times_Y X$. D'après les résultats de
la section \ref{spaces-morphisms-section-representable},
chacune des propriétés de la diagonale utilisées dans la
définition \ref{spaces-morphisms-definition-separated}
est stable par changement de base. D'où le lemme.
\end{proof}

\begin{lemma}
\label{spaces-morphisms-lemma-fibre-product-after-map}
\begin{slogan}
La flèche supérieure d'un « diagramme magique » d'espaces algébriques possède
de bonnes propriétés du type immersion, qui se renforcent
sous des hypothèses de séparation.
\end{slogan}
Soit $S$ un schéma. Soient $f : X \to Z$, $g : Y \to Z$ et $Z \to T$
des morphismes d'espaces algébriques au-dessus de $S$. Considérons le morphisme induit
$i : X \times_Z Y \to X \times_T Y$. Alors
\begin{enumerate}
\item $i$ est représentable, localement de type fini, localement quasi-fini,
séparé et est un monomorphisme,
\item si $Z \to T$ est localement séparé, alors $i$ est une immersion,
\item si $Z \to T$ est séparé, alors $i$ est une immersion fermée, et
\item si $Z \to T$ est quasi-séparé, alors $i$ est quasi-compact.
\end{enumerate}
\end{lemma}

\begin{proof}
Par la théorie générale des catégories, le diagramme suivant
$$
\xymatrix{
X \times_Z Y \ar[r]_i \ar[d] & X \times_T Y \ar[d] \\
Z \ar[r]^-{\Delta_{Z/T}} \ar[r] & Z \times_T Z
}
$$
est un diagramme de produit fibré. Ainsi, $i$ est le changement de base du
morphisme diagonal $\Delta_{Z/T}$. Le lemme résulte donc
du lemme \ref{spaces-morphisms-lemma-properties-diagonal} et des résultats de la
section \ref{spaces-morphisms-section-representable}.
\end{proof}

\begin{lemma}
\label{spaces-morphisms-lemma-semi-diagonal}
\begin{slogan}
Propriétés du graphe d'un morphisme d'espaces algébriques
déduites des propriétés de séparation du but.
\end{slogan}
Soit $S$ un schéma. Soit $T$ un espace algébrique au-dessus de $S$.
Soit $g : X \to Y$ un morphisme d'espaces algébriques au-dessus de $T$.
Considérons le morphisme graphe $i : X \to X \times_T Y$ de $g$. Alors
\begin{enumerate}
\item $i$ est représentable, localement de type fini, localement quasi-fini,
séparé et est un monomorphisme,
\item si $Y \to T$ est localement séparé, alors $i$ est une immersion,
\item si $Y \to T$ est séparé, alors $i$ est une immersion fermée, et
\item si $Y \to T$ est quasi-séparé, alors $i$ est quasi-compact.
\end{enumerate}
\end{lemma}

\begin{proof}
C'est le cas particulier du lemme \ref{spaces-morphisms-lemma-fibre-product-after-map}
obtenu en l'appliquant au morphisme $X = X \times_Y Y \to X \times_T Y$.
\end{proof}

\begin{lemma}
\label{spaces-morphisms-lemma-section-immersion}
Soit $S$ un schéma.
Soit $f : X \to T$ un morphisme d'espaces algébriques au-dessus de $S$.
Soit $s : T \to X$ une section de $f$ (autrement dit,
$f \circ s = \text{id}_T$). Alors
\begin{enumerate}
\item $s$ est représentable, localement de type fini, localement quasi-fini,
séparé et est un monomorphisme,
\item si $f$ est localement séparé, alors $s$ est une immersion,
\item si $f$ est séparé, alors $s$ est une immersion fermée, et
\item si $f$ est quasi-séparé, alors $s$ est quasi-compact.
\end{enumerate}
\end{lemma}

\begin{proof}
C'est le cas particulier du lemme \ref{spaces-morphisms-lemma-semi-diagonal} appliqué à
$g = s$, le morphisme considéré étant $i = s : T \to T \times_T X$.
\end{proof}

\begin{lemma}
\label{spaces-morphisms-lemma-composition-separated}
Tous les axiomes de séparation énumérés dans la définition \ref{spaces-morphisms-definition-separated}
sont stables par composition des morphismes.
\end{lemma}

\begin{proof}
Soient $f : X \to Y$ et $g : Y \to Z$ des morphismes d'espaces algébriques
auxquels s'applique l'axiome considéré.
Le morphisme diagonal $\Delta_{X/Z}$ est la composée
$$
X \longrightarrow X \times_Y X \longrightarrow X \times_Z X.
$$
Notre axiome de séparation consiste à imposer à la diagonale
une certaine propriété $\mathcal{P}$. D'après le
lemme \ref{spaces-morphisms-lemma-fibre-product-after-map} ci-dessus, nous voyons que
la seconde flèche possède aussi cette propriété. Le lemme en résulte,
puisque la composée de morphismes (représentables) possédant la propriété
$\mathcal{P}$ possède encore la propriété $\mathcal{P}$, voir
la section \ref{spaces-morphisms-section-representable}.
\end{proof}

\begin{lemma}
\label{spaces-morphisms-lemma-separated-over-separated}
Soit $S$ un schéma.
Soit $f : X \to Y$ un morphisme d'espaces algébriques au-dessus de $S$.
\begin{enumerate}
\item Si $Y$ est séparé et $f$ séparé, alors $X$ est séparé.
\item Si $Y$ est quasi-séparé et $f$ quasi-séparé, alors
$X$ est quasi-séparé.
\item Si $Y$ est localement séparé et $f$ localement séparé, alors
$X$ est localement séparé.
\item Si $Y$ est séparé au-dessus de $S$ et $f$ séparé, alors
$X$ est séparé au-dessus de $S$.
\item Si $Y$ est quasi-séparé au-dessus de $S$ et $f$ quasi-séparé, alors
$X$ est quasi-séparé au-dessus de $S$.
\item Si $Y$ est localement séparé au-dessus de $S$ et $f$ localement séparé, alors
$X$ est localement séparé au-dessus de $S$.
\end{enumerate}
\end{lemma}

\begin{proof}
Les assertions (4), (5) et (6) résultent immédiatement du
lemme \ref{spaces-morphisms-lemma-composition-separated}
et de
Espaces, définition \ref{spaces-definition-separated}.
Les assertions (1), (2) et (3) se ramènent aux assertions (4), (5) et (6) en considérant
$X$ et $Y$ comme des espaces algébriques au-dessus de $\Spec(\mathbf{Z})$, voir
Propriétés des espaces, définition \ref{spaces-properties-definition-separated}.
\end{proof}

\begin{lemma}
\label{spaces-morphisms-lemma-compose-after-separated}
Soit $S$ un schéma.
Soient $f : X \to Y$ et $g : Y \to Z$ des morphismes d'espaces algébriques au-dessus de $S$.
\begin{enumerate}
\item Si $g \circ f$ est séparé, alors $f$ l'est aussi.
\item Si $g \circ f$ est localement séparé, alors $f$ l'est aussi.
\item Si $g \circ f$ est quasi-séparé, alors $f$ l'est aussi.
\end{enumerate}
\end{lemma}

\begin{proof}
Considérons la factorisation
$$
X \to X \times_Y X \to X \times_Z X
$$
du morphisme diagonal de $g \circ f$. Dans tous les cas, le dernier morphisme
est un monomorphisme. Ainsi, pour tout schéma $T$ et tout morphisme
$T \to X \times_Y X$, on a l'égalité
$$
X \times_{(X \times_Y X)} T = X \times_{(X \times_Z X)} T.
$$
Le résultat est donc clair.
\end{proof}

\begin{lemma}
\label{spaces-morphisms-lemma-separated-implies-morphism-separated}
Soit $S$ un schéma. Soit $X$ un espace algébrique au-dessus de $S$.
\begin{enumerate}
\item Si $X$ est séparé, alors $X$ est séparé au-dessus de $S$.
\item Si $X$ est localement séparé, alors $X$ est localement séparé au-dessus de $S$.
\item Si $X$ est quasi-séparé, alors $X$ est quasi-séparé au-dessus de $S$.
\end{enumerate}
Soit $f : X \to Y$ un morphisme d'espaces algébriques au-dessus de $S$.
\begin{enumerate}
\item[(4)] Si $X$ est séparé au-dessus de $S$, alors $f$ est séparé.
\item[(5)] Si $X$ est localement séparé au-dessus de $S$, alors $f$ est localement séparé.
\item[(6)] Si $X$ est quasi-séparé au-dessus de $S$, alors $f$ est quasi-séparé.
\end{enumerate}
\end{lemma}

\begin{proof}
Les assertions (4), (5) et (6) résultent immédiatement du
lemme \ref{spaces-morphisms-lemma-compose-after-separated}
et de
Espaces, définition \ref{spaces-definition-separated}.
Les assertions (1), (2) et (3) résultent des assertions (4), (5) et (6) en
considérant $X$ et $Y$ comme des espaces algébriques au-dessus de
$\Spec(\mathbf{Z})$, voir
Propriétés des espaces, définition \ref{spaces-properties-definition-separated}.
\end{proof}

\begin{lemma}
\label{spaces-morphisms-lemma-separated-local}
Soit $S$ un schéma.
Soit $f : X \to Y$ un morphisme d'espaces algébriques sur $S$.
Soit $\mathcal{P}$ l'un quelconque des axiomes de
séparation de la définition \ref{spaces-morphisms-definition-separated}.
Les assertions suivantes sont équivalentes :
\begin{enumerate}
\item $f$ vérifie $\mathcal{P}$,
\item pour tout schéma $Z$ et tout morphisme $Z \to Y$, le
changement de base $Z \times_Y X \to Z$ de $f$ vérifie $\mathcal{P}$,
\item pour tout schéma affine $Z$ et tout morphisme $Z \to Y$, le
changement de base $Z \times_Y X \to Z$ de $f$ vérifie $\mathcal{P}$,
\item pour tout schéma affine $Z$ et tout morphisme $Z \to Y$,
l'espace algébrique $Z \times_Y X$ vérifie $\mathcal{P}$ (voir
Propriétés des espaces, définition \ref{spaces-properties-definition-separated}),
\item il existe un schéma $V$ et un morphisme \'etale surjectif
$V \to Y$ tels que le changement de base $V \times_Y X \to V$ vérifie
$\mathcal{P}$, et
\item il existe un recouvrement ouvert $Y = \bigcup Y_i$ tel que chacun
des morphismes $f^{-1}(Y_i) \to Y_i$ vérifie $\mathcal{P}$.
\end{enumerate}
\end{lemma}

\begin{proof}
Nous utiliserons à plusieurs reprises le
lemme \ref{spaces-morphisms-lemma-base-change-separated}
sans le mentionner davantage. En particulier, il est clair que
(1) implique (2) et que (2) implique (3).

\medskip\noindent
Montrons que (3) et (4) sont équivalentes. Remarquons que si $Z$ est un schéma
affine, alors le morphisme $Z \to \Spec(\mathbf{Z})$ est séparé
en tant que morphisme d'espaces algébriques sur $\Spec(\mathbf{Z})$.
Si $Z \times_Y X \to Z$ vérifie $\mathcal{P}$, alors
$Z \times_Y X \to \Spec(\mathbf{Z})$ vérifie $\mathcal{P}$
comme composé (voir le
lemme \ref{spaces-morphisms-lemma-composition-separated}). Par conséquent, l'espace
algébrique $Z \times_Y X$ vérifie $\mathcal{P}$. Réciproquement, si l'espace
algébrique $Z \times_Y X$ vérifie $\mathcal{P}$, alors
$Z \times_Y X \to \Spec(\mathbf{Z})$ vérifie $\mathcal{P}$, et
donc, d'après le
lemme \ref{spaces-morphisms-lemma-compose-after-separated},
$Z \times_Y X \to Z$ vérifie $\mathcal{P}$.

\medskip\noindent
Montrons que (3) implique (5). Supposons (3). Soit $V$ un schéma,
et soit $V \to Y$ \'etale surjectif. Nous devons montrer que
$V \times_Y X \to V$ vérifie la propriété $\mathcal{P}$. Autrement dit,
nous devons montrer que le morphisme
$$
V \times_Y X \longrightarrow
(V \times_Y X) \times_V (V \times_Y X) = V \times_Y X \times_Y X
$$
possède la propriété correspondante (c'est-à-dire qu'il est une immersion fermée,
une immersion, ou qu'il est quasi-compact). Soit $V = \bigcup V_j$ un
recouvrement de $V$ par des ouverts affines. Par hypothèse, chacun des morphismes
$$
V_j \times_Y X \longrightarrow V_j \times_Y X \times_Y X
$$
possède bien la propriété correspondante. Puisque les propriétés d'être une immersion fermée,
une immersion, une immersion quasi-compacte ou un morphisme quasi-compact sont locales
sur le but pour la topologie de Zariski, et puisque les $V_j$ recouvrent $V$, nous obtenons la conclusion voulue.

\medskip\noindent
Montrons que (5) implique (1). Soit $V \to Y$ comme en (5).
Nous avons alors le diagramme de produit fibré
$$
\xymatrix{
V \times_Y X \ar[r] \ar[d] &
X \ar[d] \\
V \times_Y X \times_Y X \ar[r] &
X \times_Y X
}
$$
Par hypothèse, la flèche verticale de gauche est une immersion fermée,
une immersion, une immersion quasi-compacte, ou est quasi-compacte. Il résulte du
lemme \ref{spaces-lemma-descent-representable-transformations-property} du chapitre Espaces
que la flèche verticale de droite est également une immersion fermée,
une immersion, une immersion quasi-compacte, ou est quasi-compacte.

\medskip\noindent
Il est clair que (1) implique (6) en prenant le recouvrement $Y = Y$.
Supposons que $Y = \bigcup Y_i$ soit comme en (6). Choisissons des schémas $V_i$ et
des morphismes \'etales surjectifs $V_i \to Y_i$. Notons que les morphismes
$V_i \times_Y X \to V_i$ vérifient $\mathcal{P}$, puisqu'ils sont des changements de base
des morphismes $f^{-1}(Y_i) \to Y_i$. Posons $V = \coprod V_i$.
Alors $V \to Y$ est un morphisme comme en (5) (détails omis). Ainsi,
(6) implique (5), ce qui termine la démonstration.
\end{proof}

\begin{lemma}
\label{spaces-morphisms-lemma-match-separated}
Soit $S$ un schéma.
Soit $f : X \to Y$ un morphisme représentable d'espaces algébriques sur $S$.
\begin{enumerate}
\item Le morphisme $f$ est localement séparé.
\item Le morphisme $f$ est (quasi-)séparé au sens de la
définition \ref{spaces-morphisms-definition-separated}
ci-dessus si et seulement si $f$ est (quasi-)séparé au sens de la
section \ref{spaces-morphisms-section-representable}.
\end{enumerate}
En particulier, si $f : X \to Y$ est un morphisme de schémas sur $S$, alors
$f$ est (quasi-)séparé au sens de la
définition \ref{spaces-morphisms-definition-separated}
si et seulement si $f$ est (quasi-)séparé en tant que morphisme de schémas.
\end{lemma}

\begin{proof}
C'est l'équivalence de (1) et (2) du
lemme \ref{spaces-morphisms-lemma-separated-local},
combinée au fait que tout morphisme de schémas est localement séparé ; voir le
lemme \ref{schemes-lemma-diagonal-immersion} du chapitre Schémas.
\end{proof}










\section{Morphismes surjectifs}
\label{spaces-morphisms-section-surjective}

\noindent
Nous avons déjà défini à la section \ref{spaces-morphisms-section-representable}
ce que signifie, pour un morphisme représentable d'espaces algébriques,
d'être surjectif.

\begin{lemma}
\label{spaces-morphisms-lemma-surjective-representable}
Soit $S$ un schéma. Soit $f : X \to Y$ un morphisme représentable
d'espaces algébriques sur $S$. Alors
$f$ est surjectif (au sens de la section \ref{spaces-morphisms-section-representable})
si et seulement si l'application $|f| : |X| \to |Y|$ est surjective.
\end{lemma}

\begin{proof}
En effet, si $f : X \to Y$ est représentable, il est surjectif si et seulement si,
pour tout schéma $T$ et tout morphisme $T \to Y$, le changement de base
$f_T : T \times_Y X \to T$ de $f$ est un morphisme surjectif de schémas,
autrement dit, si et seulement si $|f_T|$ est surjective. D'après le
lemme \ref{spaces-properties-lemma-points-cartesian} du chapitre Propriétés des espaces,
l'application $|T \times_Y X| \to |T| \times_{|Y|} |X|$ est toujours surjective.
Par conséquent, $|f_T| : |T \times_Y X| \to |T|$ est surjective si $|f| : |X| \to |Y|$
est surjective. Réciproquement, si $|f_T|$ est surjective
pour tout $T \to Y$ comme ci-dessus, alors, en prenant pour $T$ le spectre d'un
corps, nous concluons que $|X| \to |Y|$ est surjective.
\end{proof}

\noindent
Cela nous permet de poser la définition suivante.

\begin{definition}
\label{spaces-morphisms-definition-surjective}
Soit $S$ un schéma. Soit $f : X \to Y$ un morphisme d'espaces
algébriques sur $S$. Nous disons que $f$ est {\it surjectif}
si l'application $|f| : |X| \to |Y|$ induite sur les espaces topologiques sous-jacents
est surjective.
\end{definition}

\begin{lemma}
\label{spaces-morphisms-lemma-surjective-local}
Soit $S$ un schéma.
Soit $f : X \to Y$ un morphisme d'espaces algébriques sur $S$.
Les assertions suivantes sont équivalentes :
\begin{enumerate}
\item $f$ est surjectif,
\item pour tout schéma $Z$ et tout morphisme $Z \to Y$, le morphisme
$Z \times_Y X \to Z$ est surjectif,
\item pour tout schéma affine $Z$ et tout morphisme
$Z \to Y$, le morphisme $Z \times_Y X \to Z$ est surjectif,
\item il existe un schéma $V$ et un morphisme \'etale surjectif
$V \to Y$ tels que $V \times_Y X \to V$ soit un morphisme surjectif,
\item il existe un schéma $U$ et un morphisme \'etale surjectif
$\varphi : U \to X$ tels que le composé $f \circ \varphi$
soit surjectif,
\item il existe un diagramme commutatif
$$
\xymatrix{
U \ar[d] \ar[r] & V \ar[d] \\
X \ar[r] & Y
}
$$
où $U$, $V$ sont des schémas, les flèches verticales étant \'etales surjectives
et la flèche horizontale supérieure étant surjective, et
\item il existe un recouvrement ouvert $Y = \bigcup Y_i$ tel que
chacun des morphismes $f^{-1}(Y_i) \to Y_i$ est surjectif.
\end{enumerate}
\end{lemma}

\begin{proof}
Démonstration omise.
\end{proof}

\begin{lemma}
\label{spaces-morphisms-lemma-composition-surjective}
Le composé de morphismes surjectifs est surjectif.
\end{lemma}

\begin{proof}
Cela résulte immédiatement de la définition.
\end{proof}

\begin{lemma}
\label{spaces-morphisms-lemma-base-change-surjective}
Tout changement de base d'un morphisme surjectif est surjectif.
\end{lemma}

\begin{proof}
Cela résulte immédiatement du
lemme \ref{spaces-properties-lemma-points-cartesian} du chapitre Propriétés des espaces.
\end{proof}










\section{Morphismes ouverts}
\label{spaces-morphisms-section-open}

\noindent
Pour un morphisme représentable d'espaces algébriques, nous avons déjà défini (à la
section \ref{spaces-morphisms-section-representable})
ce que signifie être universellement ouvert. Avant de donner la définition
naturelle, vérifions donc qu'elle coïncide avec celle-ci dans le cas représentable.

\begin{lemma}
\label{spaces-morphisms-lemma-characterize-representable-universally-open}
Soit $S$ un schéma. Soit $f : X \to Y$ un morphisme représentable
d'espaces algébriques sur $S$. Les assertions suivantes sont équivalentes :
\begin{enumerate}
\item $f$ est universellement ouvert
(au sens de la section \ref{spaces-morphisms-section-representable}), et
\item pour tout morphisme d'espaces algébriques $Z \to Y$, l'application
$|Z \times_Y X| \to |Z|$ entre espaces topologiques est ouverte.
\end{enumerate}
\end{lemma}

\begin{proof}
Supposons (1), et soit $Z \to Y$ comme en (2). Choisissons un schéma $V$ et
un morphisme \'etale surjectif $V \to Y$. Par hypothèse, le morphisme de
schémas $V \times_Y X \to V$ est universellement ouvert. D'après la
section \ref{spaces-properties-section-points} du chapitre Propriétés des espaces,
dans le diagramme commutatif
$$
\xymatrix{
|V \times_Y X| \ar[r] \ar[d] & |Z \times_Y X| \ar[d] \\
|V| \ar[r] & |Z|
}
$$
les flèches horizontales sont ouvertes et surjectives ; de plus,
$$
|V \times_Y X| \longrightarrow |V| \times_{|Z|} |Z \times_Y X|
$$
est surjective. Puisque la flèche verticale de gauche
est ouverte, il s'ensuit que celle de droite est
ouverte. Cela démontre (2). L'implication (2) $\Rightarrow$ (1) résulte
immédiatement des définitions.
\end{proof}

\noindent
Nous pouvons donc adopter la définition naturelle suivante.

\begin{definition}
\label{spaces-morphisms-definition-open}
Soit $S$ un schéma. Soit $f : X \to Y$ un morphisme d'espaces algébriques
sur $S$.
\begin{enumerate}
\item Nous disons que $f$ est {\it ouvert} si l'application entre espaces topologiques
$|f| : |X| \to |Y|$ est ouverte.
\item Nous disons que $f$ est {\it universellement ouvert} si, pour tout morphisme
d'espaces algébriques $Z \to Y$, l'application entre espaces topologiques
$$
|Z \times_Y X| \to |Z|
$$
est ouverte, c'est-à-dire si le changement de base $Z \times_Y X \to Z$ est ouvert.
\end{enumerate}
\end{definition}

\noindent
Notons qu'un morphisme \'etale d'espaces algébriques est universellement ouvert ;
voir la
définition \ref{spaces-properties-definition-etale} du chapitre Propriétés des espaces et les
lemmes \ref{spaces-properties-lemma-etale-open} et
\ref{spaces-properties-lemma-base-change-etale}.

\begin{lemma}
\label{spaces-morphisms-lemma-base-change-universally-open}
Tout changement de base d'un morphisme universellement ouvert d'espaces algébriques
par un morphisme quelconque d'espaces algébriques est universellement ouvert.
\end{lemma}

\begin{proof}
Cela résulte immédiatement de la définition.
\end{proof}

\begin{lemma}
\label{spaces-morphisms-lemma-composition-open}
Le composé de deux morphismes (universellement) ouverts d'espaces algébriques
est (universellement) ouvert.
\end{lemma}

\begin{proof}
Démonstration omise.
\end{proof}

\begin{lemma}
\label{spaces-morphisms-lemma-universally-open-local}
Soit $S$ un schéma. Soit $f : X \to Y$ un morphisme d'espaces algébriques
sur $S$. Les assertions suivantes sont équivalentes :
\begin{enumerate}
\item $f$ est universellement ouvert,
\item pour tout schéma $Z$ et tout morphisme $Z \to Y$,
la projection $|Z \times_Y X| \to |Z|$ est ouverte,
\item pour tout schéma affine $Z$ et tout morphisme $Z \to Y$,
la projection $|Z \times_Y X| \to |Z|$ est ouverte, et
\item il existe un schéma $V$ et un morphisme \'etale surjectif
$V \to Y$ tels que $V \times_Y X \to V$ soit un morphisme universellement ouvert
d'espaces algébriques, et
\item il existe un recouvrement ouvert $Y = \bigcup Y_i$ tel que
chacun des morphismes $f^{-1}(Y_i) \to Y_i$ soit universellement ouvert.
\end{enumerate}
\end{lemma}

\begin{proof}
Nous omettons de démontrer que (1) implique (2) et que (2) implique (3).

\medskip\noindent
Supposons (3). Choisissons un morphisme \'etale surjectif $V \to Y$.
Nous allons montrer que $V \times_Y X \to V$ est un morphisme
universellement ouvert d'espaces algébriques. Soit $Z \to V$ un morphisme
d'un espace algébrique vers $V$. Soit $W \to Z$ un morphisme \'etale surjectif,
où $W = \coprod W_i$ est une somme disjointe de schémas affines ; voir le
lemme
\ref{spaces-properties-lemma-cover-by-union-affines} du chapitre Propriétés des espaces.
Nous avons alors le diagramme commutatif suivant :
$$
\xymatrix{
\coprod_i |W_i \times_Y X| \ar@{=}[r] \ar[d] &
|W \times_Y X| \ar[r] \ar[d] &
|Z \times_Y X| \ar[d] \ar@{=}[r] &
|Z \times_V (V \times_Y X)| \ar[ld] \\
\coprod |W_i| \ar@{=}[r] &
|W| \ar[r] &
|Z|
}
$$
Nous devons montrer que la flèche sud-est est ouverte. Les flèches horizontales
centrales sont surjectives et ouvertes
(lemme \ref{spaces-properties-lemma-etale-open} du chapitre Propriétés des espaces).
D'après l'hypothèse (3) et le fait que
$W_i$ est affine, les flèches verticales de gauche sont ouvertes. Par conséquent,
la flèche verticale de droite est ouverte.

\medskip\noindent
Supposons que $V \to Y$ soit comme en (4). Montrons que $f$ est universellement ouvert.
Soit $Z \to Y$ un morphisme d'espaces algébriques. Considérons le
diagramme
$$
\xymatrix{
|(V \times_Y Z) \times_V (V \times_Y X)| \ar@{=}[r] \ar[rd] &
|V \times_Y X| \ar[r] \ar[d] &
|Z \times_Y X| \ar[d] \\
 &
|V \times_Y Z| \ar[r] &
|Z|
}
$$
La flèche sud-ouest est ouverte par hypothèse. Les flèches horizontales sont
surjectives et ouvertes, car les morphismes correspondants
d'espaces algébriques sont \'etales (voir le
lemme \ref{spaces-properties-lemma-etale-open} du chapitre Propriétés des espaces).
Il s'ensuit que la flèche verticale de droite est ouverte.

\medskip\noindent
Bien entendu, (1) implique (5) en prenant le recouvrement $Y = Y$.
Supposons que $Y = \bigcup Y_i$ soit comme en (5). Alors, pour tout $Z \to Y$,
nous obtenons un recouvrement ouvert correspondant $Z = \bigcup Z_i$ tel que
le changement de base de $f$ à $Z_i$ soit ouvert. Un argument topologique élémentaire
montre alors que $Z \times_Y X \to Z$ est ouvert. Par conséquent, (1) est vérifiée.
\end{proof}

\begin{lemma}
\label{spaces-morphisms-lemma-space-over-field-universally-open}
Soit $S$ un schéma. Soit $p : X \to \Spec(k)$ un morphisme
d'espaces algébriques sur $S$, où $k$ est un corps. Alors
$p : X \to \Spec(k)$ est universellement ouvert.
\end{lemma}

\begin{proof}
Choisissons un schéma $U$ et un morphisme \'etale surjectif $U \to X$.
Le composé $U \to \Spec(k)$ est universellement ouvert (en tant que morphisme
de schémas) d'après le
lemme \ref{morphisms-lemma-scheme-over-field-universally-open} du chapitre Morphismes.
Soit $Z \to \Spec(k)$ un morphisme de schémas. Alors
$U \times_{\Spec(k)} Z \to X \times_{\Spec(k)} Z$ est surjectif ;
voir le
lemme \ref{spaces-morphisms-lemma-base-change-surjective}.
Par conséquent, la première des applications
$$
|U \times_{\Spec(k)} Z| \to |X \times_{\Spec(k)} Z| \to |Z|
$$
est surjective. Comme leur composé est ouvert d'après ce qui précède, nous en déduisons que
la seconde application est elle aussi ouverte. Ainsi, $p$ est universellement ouvert d'après le
lemme \ref{spaces-morphisms-lemma-universally-open-local}.
\end{proof}








\section{Morphismes submersifs}
\label{spaces-morphisms-section-submersive}

\noindent
Pour un morphisme représentable d'espaces algébriques, nous avons déjà défini (à la
section \ref{spaces-morphisms-section-representable})
ce que signifie être universellement submersif. Avant de donner la définition
naturelle, vérifions donc qu'elle coïncide avec celle-ci dans le cas représentable.

\begin{lemma}
\label{spaces-morphisms-lemma-characterize-representable-universally-submersive}
Soit $S$ un schéma. Soit $f : X \to Y$ un morphisme représentable
d'espaces algébriques sur $S$. Les assertions suivantes sont équivalentes :
\begin{enumerate}
\item $f$ est universellement submersif
(au sens de la section \ref{spaces-morphisms-section-representable}), et
\item pour tout morphisme d'espaces algébriques $Z \to Y$, l'application
$|Z \times_Y X| \to |Z|$ entre espaces topologiques est submersive.
\end{enumerate}
\end{lemma}

\begin{proof}
Supposons (1), et soit $Z \to Y$ comme en (2). Choisissons un schéma $V$ et
un morphisme \'etale surjectif $V \to Y$. Par hypothèse, le morphisme de
schémas $V \times_Y X \to V$ est universellement submersif. D'après la
section \ref{spaces-properties-section-points} du chapitre Propriétés des espaces,
dans le diagramme commutatif
$$
\xymatrix{
|V \times_Y X| \ar[r] \ar[d] & |Z \times_Y X| \ar[d] \\
|V| \ar[r] & |Z|
}
$$
les flèches horizontales sont ouvertes et surjectives ; de plus,
$$
|V \times_Y X| \longrightarrow |V| \times_{|Z|} |Z \times_Y X|
$$
est surjective. Puisque la flèche verticale de gauche est submersive,
il s'ensuit que celle de droite est submersive.
Cela démontre (2). L'implication (2) $\Rightarrow$ (1) résulte
immédiatement des définitions.
\end{proof}

\noindent
Nous pouvons donc adopter la définition naturelle suivante.

\begin{definition}
\label{spaces-morphisms-definition-submersive}
Soit $S$ un schéma.
Soit $f : X \to Y$ un morphisme d'espaces algébriques sur $S$.
\begin{enumerate}
\item Nous disons que $f$ est {\it submersif}\footnote{Ceci diffère profondément
de la notion de submersion de variétés différentielles.}
si l'application continue $|X| \to |Y|$ est submersive ; voir la
définition \ref{topology-definition-submersive} du chapitre Topologie.
\item Nous disons que $f$ est {\it universellement submersif} si, pour tout
morphisme d'espaces algébriques $Y' \to Y$, le changement de base
$Y' \times_Y X \to Y'$ est submersif.
\end{enumerate}
\end{definition}

\noindent
Remarquons qu'un morphisme submersif est en particulier surjectif.

\begin{lemma}
\label{spaces-morphisms-lemma-base-change-universally-submersive}
Tout changement de base d'un morphisme universellement submersif d'espaces algébriques
par un morphisme quelconque d'espaces algébriques est universellement submersif.
\end{lemma}

\begin{proof}
Cela résulte immédiatement de la définition.
\end{proof}

\begin{lemma}
\label{spaces-morphisms-lemma-composition-universally-submersive}
Le composé de deux morphismes (universellement) submersifs
d'espaces algébriques est (universellement) submersif.
\end{lemma}

\begin{proof}
Démonstration omise.
\end{proof}












\section{Morphismes quasi-compacts}
\label{spaces-morphisms-section-quasi-compact}

\noindent
D'après la section \ref{spaces-morphisms-section-representable}, nous savons ce que signifie
pour un morphisme représentable d'espaces algébriques d'être quasi-compact.
Afin de formuler la définition pour un morphisme quelconque
d'espaces algébriques, faisons l'observation suivante.

\begin{lemma}
\label{spaces-morphisms-lemma-characterize-representable-quasi-compact}
Soit $S$ un schéma.
Soit $f : X \to Y$ un morphisme représentable d'espaces algébriques sur $S$.
Les conditions suivantes sont équivalentes :
\begin{enumerate}
\item $f$ est quasi-compact
(au sens de la section \ref{spaces-morphisms-section-representable}), et
\item pour tout espace algébrique quasi-compact $Z$ et tout morphisme
$Z \to Y$, l'espace algébrique $Z \times_Y X$ est quasi-compact.
\end{enumerate}
\end{lemma}

\begin{proof}
Supposons (1), et soit $Z \to Y$ un morphisme d'espaces algébriques tel que
$Z$ soit quasi-compact. D'après
Propriétés des espaces,
définition \ref{spaces-properties-definition-quasi-compact},
il existe un schéma quasi-compact $U$ et un morphisme \'etale surjectif
$U \to Z$. Puisque $f$ est représentable et quasi-compact,
la définition montre que $U \times_Y X$ est un schéma et que
$U \times_Y X \to U$ est quasi-compact. Ainsi $U \times_Y X$ est
un schéma quasi-compact. Le morphisme $U \times_Y X \to Z \times_Y X$
est \'etale et surjectif (car c'est le changement de base du morphisme
représentable, \'etale et surjectif $U \to Z$ ; voir la
section \ref{spaces-morphisms-section-representable}).
Par définition, $Z \times_Y X$ est donc quasi-compact.

\medskip\noindent
Supposons (2). Soit $Z \to Y$ un morphisme, où $Z$ est un schéma.
Il faut montrer que $p : Z \times_Y X \to Z$ est quasi-compact.
Soit $U \subset Z$ un ouvert affine. Alors $p^{-1}(U) = U \times_Y X$
et le schéma $U \times_Y X$ est quasi-compact par l'hypothèse (2).
Ainsi $p$ est quasi-compact ; voir
Schémas, section \ref{schemes-section-quasi-compact}.
\end{proof}

\noindent
Ceci motive la définition suivante.

\begin{definition}
\label{spaces-morphisms-definition-quasi-compact}
Soit $S$ un schéma.
Soit $f : X \to Y$ un morphisme d'espaces algébriques sur $S$.
Nous disons que $f$ est {\it quasi-compact} si, pour tout espace
algébrique quasi-compact $Z$ et tout morphisme $Z \to Y$, le produit fibré
$Z \times_Y X$ est quasi-compact.
\end{definition}

\noindent
D'après le lemme \ref{spaces-morphisms-lemma-characterize-representable-quasi-compact}
ci-dessus, cette définition coïncide avec la notion déjà existante
pour les morphismes représentables d'espaces algébriques.

\begin{lemma}
\label{spaces-morphisms-lemma-quasi-compact-is-quasi-compact}
Soit $S$ un schéma. Si $f : X \to Y$ est un morphisme quasi-compact
d'espaces algébriques sur $S$, alors l'application sous-jacente
$|f| : |X| \to |Y|$ d'espaces topologiques est quasi-compacte.
\end{lemma}

\begin{proof}
Soit $V \subset |Y|$ un ouvert quasi-compact. D'après Propriétés des espaces,
lemme \ref{spaces-properties-lemma-open-subspaces},
il existe un sous-espace ouvert $Y' \subset Y$ tel que $V = |Y'|$.
Alors $Y'$ est un espace algébrique quasi-compact d'après
Propriétés des espaces, lemme \ref{spaces-properties-lemma-quasi-compact-space},
et donc $X' = Y' \times_Y X$ est un espace algébrique
quasi-compact d'après la définition \ref{spaces-morphisms-definition-quasi-compact}.
D'autre part, $X' \subset X$ est un sous-espace ouvert
(Espaces, lemme \ref{spaces-lemma-base-change-immersions})
et $|X'| = |f|^{-1}(V)$ d'après
Propriétés des espaces, lemme \ref{spaces-properties-lemma-points-cartesian}.
En appliquant de nouveau
Propriétés des espaces, lemme \ref{spaces-properties-lemma-quasi-compact-space},
on conclut que $|X'|$ est un ouvert quasi-compact de $|X|$, comme souhaité.
\end{proof}

\begin{lemma}
\label{spaces-morphisms-lemma-base-change-quasi-compact}
Tout changement de base d'un morphisme quasi-compact d'espaces algébriques
par un morphisme quelconque d'espaces algébriques est quasi-compact.
\end{lemma}

\begin{proof}
Démonstration omise. Indication : transitivité des produits fibrés.
\end{proof}

\begin{lemma}
\label{spaces-morphisms-lemma-composition-quasi-compact}
Le composé de deux morphismes quasi-compacts d'espaces algébriques
est quasi-compact.
\end{lemma}

\begin{proof}
Démonstration omise. Indication : transitivité des produits fibrés.
\end{proof}

\begin{lemma}
\label{spaces-morphisms-lemma-surjection-from-quasi-compact}
\begin{slogan}
L'image d'un espace algébrique quasi-compact par un morphisme surjectif
est quasi-compacte.
\end{slogan}
Soit $S$ un schéma.
\begin{enumerate}
\item Si $X \to Y$ est un morphisme surjectif d'espaces algébriques sur $S$,
et si $X$ est quasi-compact, alors $Y$ est quasi-compact.
\item Si
$$
\xymatrix{
X \ar[rr]_f \ar[rd]_p & &
Y \ar[dl]^q \\
& Z
}
$$
est un diagramme commutatif de morphismes d'espaces algébriques sur $S$
et si $f$ est surjectif et $p$ quasi-compact, alors $q$ est quasi-compact.
\end{enumerate}
\end{lemma}

\begin{proof}
Supposons $X$ quasi-compact et $X \to Y$ surjectif. D'après la
définition \ref{spaces-morphisms-definition-surjective},
l'application $|X| \to |Y|$ est surjective ; ainsi $Y$ est quasi-compact d'après
Propriétés des espaces, lemme \ref{spaces-properties-lemma-quasi-compact-space},
et le fait topologique que l'image d'un espace quasi-compact par une
application continue est quasi-compacte ; voir
Topologie, lemme \ref{topology-lemma-image-quasi-compact}.
Soient $f, p, q$ comme en (2).
Soit $T \to Z$ un morphisme dont la source est un espace algébrique quasi-compact.
Par hypothèse, $T \times_Z X$ est quasi-compact. D'après le
lemme \ref{spaces-morphisms-lemma-base-change-surjective},
le morphisme $T \times_Z X \to T \times_Z Y$ est surjectif.
La partie (1) montre donc que $T \times_Z Y$ est également quasi-compact.
Ainsi $q$ est quasi-compact.
\end{proof}

\begin{lemma}
\label{spaces-morphisms-lemma-descent-quasi-compact}
Soit $S$ un schéma.
Soit $f : X \to Y$ un morphisme d'espaces algébriques sur $S$.
Soit $g : Y' \to Y$ un morphisme universellement ouvert et surjectif
d'espaces algébriques tel que le changement de base $f' : X' \to Y'$ soit quasi-compact.
Alors $f$ est quasi-compact.
\end{lemma}

\begin{proof}
Soit $Z \to Y$ un morphisme d'espaces algébriques avec $Z$ quasi-compact.
Comme $g$ est universellement ouvert et surjectif, on voit que
$Y' \times_Y Z \to Z$ est ouvert et surjectif. Tout point de
$|Y' \times_Y Z|$ admettant un système fondamental de voisinages ouverts
quasi-compacts (voir
Propriétés des espaces,
lemme \ref{spaces-properties-lemma-space-locally-quasi-compact}),
on peut trouver un ouvert quasi-compact $W \subset |Y' \times_Y Z|$
dont l'image est $Z$. Notons
$f'' : W \times_Y X \to W$ le changement de base de $f'$ par $W \to Y'$.
Par hypothèse, $W \times_Y X$ est quasi-compact. Comme $W \to Z$ est surjectif,
on voit que $W \times_Y X \to Z \times_Y X$ est surjectif.
Ainsi $Z \times_Y X$ est quasi-compact d'après le
lemme \ref{spaces-morphisms-lemma-surjection-from-quasi-compact}.
Ainsi $f$ est quasi-compact.
\end{proof}

\begin{lemma}
\label{spaces-morphisms-lemma-quasi-compact-local}
\begin{slogan}
Les morphismes quasi-compacts d'espaces algébriques sont stables par
changement de base et cette propriété est locale sur la base.
\end{slogan}
Soit $S$ un schéma.
Soit $f : X \to Y$ un morphisme d'espaces algébriques sur $S$.
Les conditions suivantes sont équivalentes :
\begin{enumerate}
\item $f$ est quasi-compact,
\item pour tout schéma $Z$ et tout morphisme $Z \to Y$, le morphisme
d'espaces algébriques $Z \times_Y X \to Z$ est quasi-compact,
\item pour tout schéma affine $Z$ et tout morphisme
$Z \to Y$, l'espace algébrique $Z \times_Y X$ est quasi-compact,
\item il existe un schéma $V$ et un morphisme \'etale surjectif
$V \to Y$ tels que $V \times_Y X \to V$ soit un morphisme quasi-compact
d'espaces algébriques, et
\item il existe un morphisme \'etale surjectif
$Y' \to Y$ d'espaces algébriques tel que $Y' \times_Y X \to Y'$
soit un morphisme quasi-compact d'espaces algébriques, et
\item il existe un recouvrement de Zariski $Y = \bigcup Y_i$ tel que
chacun des morphismes $f^{-1}(Y_i) \to Y_i$ soit quasi-compact.
\end{enumerate}
\end{lemma}

\begin{proof}
Nous utiliserons le lemme \ref{spaces-morphisms-lemma-base-change-quasi-compact}
sans plus le mentionner.
Il est clair que (1) implique (2) et que (2) implique (3).
Supposons (3). Soit $Z$ un espace algébrique quasi-compact sur $S$,
et soit $Z \to Y$ un morphisme. D'après
Propriétés des espaces, lemme
\ref{spaces-properties-lemma-quasi-compact-affine-cover},
il existe un schéma affine $U$ et un morphisme \'etale surjectif
$U \to Z$. Alors $U \times_Y X \to Z \times_Y X$ est un morphisme
surjectif d'espaces algébriques ; voir le
lemme \ref{spaces-morphisms-lemma-base-change-surjective}.
Par hypothèse, $|U \times_Y X|$ est quasi-compact et son image
est $|Z \times_Y X|$ ; on en déduit que $|Z \times_Y X|$
est quasi-compact ; voir
Topologie, lemme \ref{topology-lemma-image-quasi-compact}.
Ceci prouve que (3) implique (1).

\medskip\noindent
Les implications (1) $\Rightarrow$ (4) et (4) $\Rightarrow$ (5) sont claires.
L'implication (5) $\Rightarrow$ (1) découle du
lemme \ref{spaces-morphisms-lemma-descent-quasi-compact}
et du fait qu'un morphisme \'etale d'espaces algébriques est universellement ouvert
(voir la discussion qui suit la
définition \ref{spaces-morphisms-definition-open}).

\medskip\noindent
Bien entendu, (1) implique (6) en prenant le recouvrement $Y = Y$.
Supposons que $Y = \bigcup Y_i$ soit comme en (6). Soit $Z$ affine et soit
$Z \to Y$ un morphisme. Il existe alors un recouvrement affine standard fini
$Z = Z_1 \cup \ldots \cup Z_n$ tel que chaque $Z_j \to Y$
se factorise par $Y_{i_j}$ pour un certain $i_j$. Ainsi, l'espace algébrique
$$
Z_j \times_Y X = Z_j \times_{Y_{i_j}} f^{-1}(Y_{i_j})
$$
est quasi-compact. Puisque
$Z \times_Y X = \bigcup_{j = 1, \ldots, n} Z_j \times_Y X$
est un recouvrement de Zariski, on voit que
$|Z \times_Y X| = \bigcup_{j = 1, \ldots, n} |Z_j \times_Y X|$
(voir Propriétés des espaces, lemme \ref{spaces-properties-lemma-open-subspaces})
est une réunion finie d'espaces quasi-compacts, donc est quasi-compact.
Ainsi, (6) implique (3).
\end{proof}

\noindent
Le lemme ci-dessous (ainsi que le suivant) garantit en particulier qu'un morphisme
$X \to \Spec(A)$ est quasi-compact dès que
$X$ est un espace algébrique quasi-compact.

\begin{lemma}
\label{spaces-morphisms-lemma-quasi-compact-permanence}
Soit $S$ un schéma.
Soient $f : X \to Y$ et $g : Y \to Z$ des morphismes d'espaces algébriques sur $S$.
Si $g \circ f$ est quasi-compact et si $g$ est quasi-séparé,
alors $f$ est quasi-compact.
\end{lemma}

\begin{proof}
Ceci est vrai parce que $f$ est le composé
$(1, f) : X \to X \times_Z Y \to Y$. Le premier morphisme
est quasi-compact d'après le lemme \ref{spaces-morphisms-lemma-section-immersion},
car c'est une section du morphisme quasi-séparé $X \times_Z Y \to X$
(un changement de base de $g$ ; voir le lemme \ref{spaces-morphisms-lemma-base-change-separated}).
Le second morphisme est quasi-compact puisqu'il
est le changement de base de $g \circ f$ ; voir le
lemme \ref{spaces-morphisms-lemma-base-change-quasi-compact}.
Enfin, les composés de morphismes quasi-compacts
sont quasi-compacts ; voir le lemme \ref{spaces-morphisms-lemma-composition-quasi-compact}.
\end{proof}

\begin{lemma}
\label{spaces-morphisms-lemma-quasi-compact-quasi-separated-permanence}
Soit $f : X \to Y$ un morphisme d'espaces algébriques
sur un schéma $S$.
\begin{enumerate}
\item Si $X$ est quasi-compact et $Y$ quasi-séparé, alors $f$ est
quasi-compact.
\item Si $X$ est quasi-compact et quasi-séparé et si $Y$ est quasi-séparé,
alors $f$ est quasi-compact et quasi-séparé.
\item Un produit fibré d'espaces algébriques quasi-compacts et quasi-séparés
est quasi-compact et quasi-séparé.
\end{enumerate}
\end{lemma}

\begin{proof}
La partie (1) résulte du
lemme \ref{spaces-morphisms-lemma-quasi-compact-permanence}
avec $Z = S = \Spec(\mathbf{Z})$.
La partie (2) résulte de (1) et du
lemme \ref{spaces-morphisms-lemma-compose-after-separated}.
Pour (3), soient $X \to Y$ et $Z \to Y$ des morphismes
d'espaces algébriques quasi-compacts et quasi-séparés. Alors
$X \times_Y Z \to Z$ est quasi-compact et quasi-séparé en tant que
changement de base de $X \to Y$, en vertu de (2) et des
lemmes \ref{spaces-morphisms-lemma-base-change-quasi-compact} et
\ref{spaces-morphisms-lemma-base-change-separated}.
Ainsi $X \times_Y Z$ est quasi-compact et quasi-séparé comme
espace algébrique quasi-compact et quasi-séparé sur
$Z$ ; voir les
lemmes \ref{spaces-morphisms-lemma-separated-over-separated} et
\ref{spaces-morphisms-lemma-composition-quasi-compact}.
\end{proof}


\section{Morphismes universellement fermés}
\label{spaces-morphisms-section-universally-closed}

\noindent
Pour un morphisme représentable d'espaces algébriques, nous avons déjà défini (à la
section \ref{spaces-morphisms-section-representable})
ce que signifie être universellement fermé. Avant de donner la définition
naturelle, vérifions donc qu'elle coïncide avec celle-ci dans le cas représentable.

\begin{lemma}
\label{spaces-morphisms-lemma-characterize-representable-universally-closed}
Soit $S$ un schéma. Soit $f : X \to Y$ un morphisme représentable
d'espaces algébriques sur $S$. Les conditions suivantes sont équivalentes :
\begin{enumerate}
\item $f$ est universellement fermé
(au sens de la section \ref{spaces-morphisms-section-representable}), et
\item pour tout morphisme d'espaces algébriques $Z \to Y$, le morphisme
d'espaces topologiques $|Z \times_Y X| \to |Z|$ est fermé.
\end{enumerate}
\end{lemma}

\begin{proof}
Supposons (1), et soit $Z \to Y$ comme en (2). Choisissons un schéma $V$ et
un morphisme \'etale surjectif $V \to Z$. Par hypothèse, le morphisme
de schémas $V \times_Y X \to V$ est universellement fermé. D'après
Propriétés des espaces, section \ref{spaces-properties-section-points},
dans le diagramme commutatif
$$
\xymatrix{
|V \times_Y X| \ar[r] \ar[d] & |Z \times_Y X| \ar[d] \\
|V| \ar[r] & |Z|
}
$$
les flèches horizontales sont ouvertes et surjectives et, de plus,
$$
|V \times_Y X| \longrightarrow |V| \times_{|Z|} |Z \times_Y X|
$$
est surjectif. Puisque la flèche verticale
de gauche est fermée, il s'ensuit que celle de droite est
fermée. Ceci prouve (2). L'implication (2) $\Rightarrow$ (1) résulte
immédiatement des définitions.
\end{proof}

\noindent
Nous pouvons donc adopter la définition naturelle suivante.

\begin{definition}
\label{spaces-morphisms-definition-closed}
Soit $S$ un schéma. Soit $f : X \to Y$ un morphisme d'espaces algébriques
sur $S$.
\begin{enumerate}
\item Nous disons que $f$ est {\it fermé} si l'application d'espaces
topologiques $|X| \to |Y|$ est fermée.
\item Nous disons que $f$ est {\it universellement fermé} si, pour tout morphisme
d'espaces algébriques $Z \to Y$, le morphisme d'espaces topologiques
$$
|Z \times_Y X| \to |Z|
$$
est fermé, c'est-à-dire si le changement de base $Z \times_Y X \to Z$ est fermé.
\end{enumerate}
\end{definition}

\begin{lemma}
\label{spaces-morphisms-lemma-base-change-universally-closed}
Tout changement de base d'un morphisme universellement fermé d'espaces algébriques
par un morphisme quelconque d'espaces algébriques est universellement fermé.
\end{lemma}

\begin{proof}
Cela résulte immédiatement de la définition.
\end{proof}

\begin{lemma}
\label{spaces-morphisms-lemma-composition-universally-closed}
Le composé de deux morphismes (universellement) fermés d'espaces algébriques
est (universellement) fermé.
\end{lemma}

\begin{proof}
Démonstration omise.
\end{proof}

\begin{lemma}
\label{spaces-morphisms-lemma-universally-closed-local}
Soit $S$ un schéma. Soit $f : X \to Y$ un morphisme d'espaces algébriques
sur $S$. Les conditions suivantes sont équivalentes :
\begin{enumerate}
\item $f$ est universellement fermé,
\item pour tout schéma $Z$ et tout morphisme $Z \to Y$,
la projection $|Z \times_Y X| \to |Z|$ est fermée,
\item pour tout schéma affine $Z$ et tout morphisme $Z \to Y$,
la projection $|Z \times_Y X| \to |Z|$ est fermée,
\item il existe un schéma $V$ et un morphisme \'etale surjectif
$V \to Y$ tel que $V \times_Y X \to V$ soit un morphisme universellement fermé
d'espaces algébriques, et
\item il existe un recouvrement de Zariski $Y = \bigcup Y_i$ tel que
chacun des morphismes $f^{-1}(Y_i) \to Y_i$ soit universellement fermé.
\end{enumerate}
\end{lemma}

\begin{proof}
Nous omettons la preuve que (1) implique (2) et que (2) implique (3).

\medskip\noindent
Supposons (3). Choisissons un morphisme \'etale surjectif $V \to Y$.
Nous allons montrer que $V \times_Y X \to V$ est un morphisme
universellement fermé d'espaces algébriques. Soit $Z \to V$ un morphisme
d'un espace algébrique vers $V$. Soit $W \to Z$ un morphisme \'etale surjectif
où $W = \coprod W_i$ est une somme disjointe de schémas affines ; voir
Propriétés des espaces,
lemme \ref{spaces-properties-lemma-cover-by-union-affines}.
Nous avons alors le diagramme commutatif suivant :
$$
\xymatrix{
\coprod_i |W_i \times_Y X| \ar@{=}[r] \ar[d] &
|W \times_Y X| \ar[r] \ar[d] &
|Z \times_Y X| \ar[d] \ar@{=}[r] &
|Z \times_V (V \times_Y X)| \ar[ld] \\
\coprod |W_i| \ar@{=}[r] &
|W| \ar[r] &
|Z|
}
$$
Il faut montrer que la flèche sud-est est fermée. Les flèches horizontales
centrales sont surjectives et ouvertes
(Propriétés des espaces, lemme \ref{spaces-properties-lemma-etale-open}).
L'hypothèse (3), jointe au fait que
$W_i$ est affine, montre que les flèches verticales de gauche sont fermées. Il
s'ensuit que la flèche verticale de droite est fermée.

\medskip\noindent
Supposons (4). Montrons que $f$ est universellement fermé.
Soit $Z \to Y$ un morphisme d'espaces algébriques. Considérons le
diagramme
$$
\xymatrix{
|(V \times_Y Z) \times_V (V \times_Y X)| \ar@{=}[r] \ar[rd] &
|V \times_Y X| \ar[r] \ar[d] &
|Z \times_Y X| \ar[d] \\
 &
|V \times_Y Z| \ar[r] &
|Z|
}
$$
La flèche sud-ouest est fermée par hypothèse. Les flèches horizontales sont
surjectives et ouvertes, car les morphismes correspondants d'espaces
algébriques sont \'etales (voir
Propriétés des espaces, lemme \ref{spaces-properties-lemma-etale-open}).
Il s'ensuit que la flèche verticale de droite est fermée.

\medskip\noindent
Bien entendu, (1) implique (5) en prenant le recouvrement $Y = Y$.
Supposons que $Y = \bigcup Y_i$ soit comme en (5). Pour tout $Z \to Y$,
on obtient alors un recouvrement de Zariski correspondant $Z = \bigcup Z_i$ tel que
le changement de base de $f$ à $Z_i$ soit fermé. Un argument topologique
élémentaire montre que $Z \times_Y X \to Z$ est fermé. Ainsi (1) est satisfaite.
\end{proof}

\begin{example}
\label{spaces-morphisms-example-strange-universally-closed}
Exemple étrange d'un morphisme universellement fermé.
Soit $\mathbf{Q} \subset k$ un corps de caractéristique zéro.
Soit $X = \mathbf{A}^1_k/\mathbf{Z}$ comme dans
Espaces, exemple \ref{spaces-example-affine-line-translation}.
Nous affirmons que le morphisme structural $p : X \to \Spec(k)$
est universellement fermé.
En effet, si $Z/k$ est un schéma et si $T \subset |X \times_k Z|$ est fermé,
alors $T$ correspond à une partie fermée $\mathbf{Z}$-invariante
$T' \subset |\mathbf{A}^1 \times Z|$. Il est facile de voir que
cela implique que $T'$ est l'image inverse d'une partie $T''$ de
$Z$. D'après
Morphismes, lemme \ref{morphisms-lemma-fpqc-quotient-topology},
la partie $T'' \subset Z$ est fermée.
Bien entendu, $T''$ est l'image de $T$. Le morphisme $p$ est donc universellement
fermé d'après le lemme \ref{spaces-morphisms-lemma-universally-closed-local}.
\end{example}

\begin{lemma}
\label{spaces-morphisms-lemma-universally-closed-quasi-compact}
Soit $S$ un schéma.
Tout morphisme universellement fermé d'espaces algébriques sur $S$ est quasi-compact.
\end{lemma}

\begin{proof}
Cette démonstration reprend celle du cas des schémas ; voir
Morphismes, lemme \ref{morphisms-lemma-universally-closed-quasi-compact}.
Soit $f : X \to Y$ un morphisme d'espaces algébriques sur $S$.
Supposons que $f$ ne soit pas quasi-compact.
Notre but est de montrer que $f$ n'est pas universellement fermé. D'après le
lemme \ref{spaces-morphisms-lemma-quasi-compact-local},
il existe un schéma affine $Z$ et un morphisme $Z \to Y$
tels que $Z \times_Y X \to Z$ ne soit pas quasi-compact. Pour atteindre notre but,
il suffit de montrer que $Z \times_Y X \to Z$ n'est pas universellement fermé ;
nous pouvons donc supposer que $Y = \Spec(B)$ pour un certain anneau $B$.

\medskip\noindent
Écrivons $X = \bigcup_{i \in I} X_i$, où les $X_i$ sont des sous-espaces ouverts
quasi-compacts de $X$. Par exemple, choisissons un morphisme \'etale surjectif
$U \to X$, où $U$ est un schéma, puis un recouvrement ouvert affine
$U = \bigcup U_i$, et soit $X_i \subset X$ l'image de $U_i$.
Nous utiliserons plus loin que les morphismes $X_i \to Y$ sont quasi-compacts ; voir le
lemme \ref{spaces-morphisms-lemma-quasi-compact-permanence}.
Soit $T = \Spec(B[a_i ; i \in I])$. Soit $T_i = D(a_i) \subset T$.
Soit $Z \subset T \times_Y X$ le sous-espace fermé réduit dont l'ensemble fermé
sous-jacent de points est
$|T \times_Y X| \setminus \bigcup_{i \in I} |T_i \times_Y X_i|$ ; voir
Propriétés des espaces,
lemme \ref{spaces-properties-lemma-reduced-closed-subspace}.
(Notons que $T_i \times_Y X_i$ est un sous-espace ouvert de $T \times_Y X$, car
$T_i \to T$ et $X_i \to X$ sont des immersions ouvertes ; voir
Espaces, lemmes \ref{spaces-lemma-base-change-immersions} et
\ref{spaces-lemma-composition-immersions}.) Voici un diagramme :
$$
\xymatrix{
Z \ar[r] \ar[rd] &
T \times_Y X \ar[d]^{f_T} \ar[r]_q &
X \ar[d]^f \\
& T \ar[r]^p & Y
}
$$
Il suffit de prouver que l'image $f_T(|Z|)$ n'est pas fermée dans $|T|$.

\medskip\noindent
Nous affirmons qu'il existe un point $y \in Y$ tel qu'il n'existe aucun
voisinage ouvert affine $V$ de $y$ dans $Y$ tel que $X_V$ soit quasi-compact.
Sinon, on pourrait recouvrir $Y$ par un nombre fini de tels $V$ et, pour
chaque $V$, le morphisme $X_V \to V$ serait quasi-compact d'après le
lemme \ref{spaces-morphisms-lemma-quasi-compact-permanence} ;
alors le
lemme \ref{spaces-morphisms-lemma-quasi-compact-local}
entraînerait que $f$ est quasi-compact, contradiction. Fixons un $y \in Y$ ainsi obtenu.

\medskip\noindent
Soit $t \in T$ le point au-dessus de $y$ tel que $\kappa(t) = \kappa(y)$
et tel que $a_i = 1$ dans $\kappa(t)$ pour tout $i$. Supposons $z \in |Z|$ avec
$f_T(z) = t$. Alors $q(z) \in X_i$ pour un certain $i$. Par conséquent $f_T(z) \not \in T_i$
par construction de $Z$, ce qui contredit le fait que $t \in T_i$ par
construction. Ainsi $t \in |T| \setminus f_T(|Z|)$.

\medskip\noindent
Supposons $f_T(|Z|)$ fermé dans $|T|$. Il existe alors un élément
$g \in B[a_i; i \in I]$ tel que $f_T(|Z|) \subset V(g)$ mais $t \not \in V(g)$.
L'image de $g$ dans $\kappa(t)$ est donc non nulle. En particulier, un
coefficient de $g$ a une image non nulle dans $\kappa(y)$. Ce coefficient est donc
inversible sur un voisinage ouvert affine $V$ de $y$. Soit $J$ l'ensemble fini
des $j \in I$ tels que la variable $a_j$ apparaisse dans $g$.
Puisque $X_V$ n'est pas quasi-compact et que chaque $X_{i, V}$ est quasi-compact,
on peut choisir un point $x \in |X_V| \setminus \bigcup_{j \in J} |X_{j, V}|$.
Autrement dit, $x \in |X| \setminus \bigcup_{j \in J} |X_j|$ et $x$ se trouve
au-dessus d'un certain $v \in V$. Puisque $g$ possède un coefficient inversible sur $V$,
on peut trouver un point $t' \in T$ au-dessus de $v$ tel que $t' \not \in V(g)$ et
$t' \in V(a_i)$ pour tout $i \notin J$. En effet,
$V(a_i; i \in I \setminus J) = \Spec(B[a_j; j\in J])$
et l'ensemble des points de ce schéma au-dessus de $v$ est en bijection
avec $\Spec(\kappa(v)[a_j; j \in J])$, et la restriction de $g$ est
un élément non nul de cet anneau de polynômes par construction.
Autrement dit, $t' \not \in T_i$ pour tout $i \not \in J$. D'après
Propriétés des espaces, lemme \ref{spaces-properties-lemma-points-cartesian},
on peut trouver un point $z$ de $X \times_Y T$ qui s'envoie sur $x \in X$ et sur
$t' \in T$. Puisque $x \not \in |X_j|$ pour $j \in J$ et $t' \not \in T_i$
pour $i \in I \setminus J$, on voit que $z \in |Z|$. D'autre part,
$f_T(z) = t' \not \in V(g)$, ce qui contredit $f_T(Z) \subset V(g)$.
L'hypothèse « $f_T(|Z|)$ est fermé » est donc fausse, et l'on conclut bien
que $f_T$ n'est pas fermé, comme souhaité.
\end{proof}

\noindent
L'espace d'arrivée d'un morphisme surjectif et universellement fermé issu d'un
espace algébrique séparé est séparé.

\begin{lemma}
\label{spaces-morphisms-lemma-image-universally-closed-separated}
Soit $S$ un schéma. Soit $B$ un espace algébrique sur $S$.
Soit $f : X \to Y$ un morphisme surjectif et universellement fermé
d'espaces algébriques sur $B$.
\begin{enumerate}
\item Si $X$ est quasi-séparé, alors $Y$ est quasi-séparé.
\item Si $X$ est séparé, alors $Y$ est séparé.
\item Si $X$ est quasi-séparé sur $B$, alors $Y$ est quasi-séparé sur $B$.
\item Si $X$ est séparé sur $B$, alors $Y$ est séparé sur $B$.
\end{enumerate}
\end{lemma}

\begin{proof}
Les parties (1) et (2) résultent de (3) et (4) pour
$S = B = \Spec(\mathbf{Z})$ (voir
Propriétés des espaces,
définition \ref{spaces-properties-definition-separated}).
Considérons le diagramme commutatif
$$
\xymatrix{
X \ar[d] \ar[rr]_{\Delta_{X/B}} & & X \times_B X \ar[d] \\
Y \ar[rr]^{\Delta_{Y/B}} & & Y \times_B Y
}
$$
La flèche verticale de gauche est surjective (c'est-à-dire universellement surjective).
La flèche verticale de droite est universellement fermée comme composée
des morphismes universellement fermés
$X \times_B X \to X \times_B Y \to Y \times_B Y$. Elle est donc aussi
quasi-compacte ; voir le
lemme \ref{spaces-morphisms-lemma-universally-closed-quasi-compact}.

\medskip\noindent
Supposons $X$ quasi-séparé sur $B$, c'est-à-dire  $\Delta_{X/B}$
quasi-compact. Si $Z$ est quasi-compact et si $Z \to Y \times_B Y$ est
un morphisme, alors $Z \times_{Y \times_B Y} X \to Z \times_{Y \times_B Y} Y$
est surjectif et $Z \times_{Y \times_B Y} X$ est quasi-compact d'après les remarques
ci-dessus. On en déduit que $\Delta_{Y/B}$ est quasi-compact, c'est-à-dire que $Y$
est quasi-séparé sur $B$.

\medskip\noindent
Supposons $X$ séparé sur $B$, c'est-à-dire que $\Delta_{X/B}$ est une immersion
fermée. Si $Z$ est affine et si $Z \to Y \times_B Y$ est
un morphisme, alors $Z \times_{Y \times_B Y} X \to Z \times_{Y \times_B Y} Y$
est surjectif et $Z \times_{Y \times_B Y} X \to Z$ est universellement fermé
d'après les remarques ci-dessus. On en déduit que $\Delta_{Y/B}$ est universellement fermé.
Il s'ensuit que $\Delta_{Y/B}$ est représentable, localement de type fini et est un
monomorphisme (voir le
lemme \ref{spaces-morphisms-lemma-properties-diagonal}),
et qu'il est universellement fermé, donc une immersion fermée ; voir
Morphismes \'etales,
lemme \ref{etale-lemma-characterize-closed-immersion}
(ainsi que le principe abstrait
Espaces, lemme
\ref{spaces-lemma-representable-transformations-property-implication}).
Ainsi $Y$ est séparé sur $B$.
\end{proof}













\section{Monomorphismes}
\label{spaces-morphisms-section-monomorphisms}

\noindent
Un morphisme représentable $X \to Y$ d'espaces algébriques est un monomorphisme
au sens de la section \ref{spaces-morphisms-section-representable} si, pour tout schéma
$Z$ et tout morphisme $Z \to Y$, le morphisme $Z \times_Y X \to Z$
est représentable par un monomorphisme de schémas.
Cela signifie exactement que $Z \times_Y X \to Z$
est une application injective de faisceaux sur $(\Sch/S)_{fppf}$. Puisque ceci
doit valoir pour tout $Z$ et toute application $Z \to Y$, c'est encore
équivalent à dire que l'application $X \to Y$ est une application injective de faisceaux sur
$(\Sch/S)_{fppf}$. Nous pouvons donc définir les monomorphismes parmi les morphismes
éventuellement non représentables\footnote{Nous ignorons si tout monomorphisme
d'espaces algébriques est représentable. Pour une discussion, voir
Compléments sur les morphismes d'espaces, section
\ref{spaces-more-morphisms-section-monomorphisms}.})
d'espaces algébriques comme suit.

\begin{definition}
\label{spaces-morphisms-definition-monomorphism}
Soit $S$ un schéma.
Un morphisme d'espaces algébriques sur $S$ est appelé {\it monomorphisme}
s'il est une application injective de faisceaux, c'est-à-dire un monomorphisme dans la catégorie
des faisceaux sur $(\Sch/S)_{fppf}$.
\end{definition}

\noindent
Le lemme suivant montre que cela signifie aussi qu'il est un monomorphisme
dans la catégorie des espaces algébriques sur $S$.

\begin{lemma}
\label{spaces-morphisms-lemma-monomorphism}
Soit $S$ un schéma.
Soit $j : X \to Y$ un morphisme d'espaces algébriques sur $S$.
Les conditions suivantes sont équivalentes :
\begin{enumerate}
\item $j$ est un monomorphisme (au sens de la définition \ref{spaces-morphisms-definition-monomorphism}),
\item $j$ est un monomorphisme dans la catégorie des espaces algébriques sur $S$, et
\item le morphisme diagonal $\Delta_{X/Y} : X \to X \times_Y X$ est
un isomorphisme.
\end{enumerate}
\end{lemma}

\begin{proof}
Notons que $X \times_Y X$ est à la fois le produit fibré dans la catégorie des
faisceaux sur $(\Sch/S)_{fppf}$ et le produit fibré dans la catégorie
des espaces algébriques sur $S$ ; voir
Espaces, lemme \ref{spaces-lemma-fibre-product-spaces}.
L'équivalence de (1) et (3) est une caractérisation générale
des applications injectives de faisceaux sur un site quelconque.
L'équivalence de (2) et (3) est une caractérisation des monomorphismes
dans toute catégorie possédant des produits fibrés.
\end{proof}

\begin{lemma}
\label{spaces-morphisms-lemma-monomorphism-separated}
Un monomorphisme d'espaces algébriques est séparé.
\end{lemma}

\begin{proof}
Cela vient de ce qu'un isomorphisme est une immersion fermée, compte tenu du
lemme \ref{spaces-morphisms-lemma-monomorphism} ci-dessus.
\end{proof}

\begin{lemma}
\label{spaces-morphisms-lemma-composition-monomorphism}
Un composé de monomorphismes est un monomorphisme.
\end{lemma}

\begin{proof}
C'est vrai, car un composé d'applications injectives de faisceaux est injectif.
\end{proof}

\begin{lemma}
\label{spaces-morphisms-lemma-base-change-monomorphism}
Tout changement de base d'un monomorphisme est un monomorphisme.
\end{lemma}

\begin{proof}
C'est un fait général sur les produits fibrés dans une catégorie de faisceaux.
\end{proof}

\begin{lemma}
\label{spaces-morphisms-lemma-monomorphism-local}
Soit $S$ un schéma.
Soit $f : X \to Y$ un morphisme d'espaces algébriques sur $S$.
Les conditions suivantes sont équivalentes :
\begin{enumerate}
\item $f$ est un monomorphisme,
\item pour tout schéma $Z$ et tout morphisme $Z \to Y$, le
changement de base $Z \times_Y X \to Z$ de $f$ est un monomorphisme,
\item pour tout schéma affine $Z$ et tout morphisme $Z \to Y$, le
changement de base $Z \times_Y X \to Z$ de $f$ est un monomorphisme,
\item il existe un schéma $V$ et un morphisme \'etale surjectif
$V \to Y$ tel que le changement de base $V \times_Y X \to V$ soit un
monomorphisme, et
\item il existe un recouvrement de Zariski $Y = \bigcup Y_i$ tel que chacun
des morphismes $f^{-1}(Y_i) \to Y_i$ soit un monomorphisme.
\end{enumerate}
\end{lemma}

\begin{proof}
Nous utiliserons sans plus le mentionner que tout changement de base d'un monomorphisme
est un monomorphisme ; voir le
lemme \ref{spaces-morphisms-lemma-base-change-monomorphism}.
En particulier, il est clair que
(1) $\Rightarrow$ (2) $\Rightarrow$ (3) $\Rightarrow$ (4)
(en prenant pour $V$ une somme disjointe de schémas affines \'etales sur $Y$ ; voir
Propriétés des espaces,
lemme \ref{spaces-properties-lemma-cover-by-union-affines}).
Soit $V$ un schéma et soit $V \to Y$ un morphisme \'etale surjectif.
Si $V \times_Y X \to V$ est un monomorphisme, il
s'ensuit que $X \to Y$ est un monomorphisme. En effet, étant donné un
diagramme cartésien quelconque de faisceaux
$$
\vcenter{
\xymatrix{
\mathcal{F} \ar[r]_a \ar[d]_b & \mathcal{G} \ar[d]^c \\
\mathcal{H} \ar[r]^d & \mathcal{I}
}
}
\quad
\quad
\mathcal{F} = \mathcal{H} \times_\mathcal{I} \mathcal{G}
$$
si $c$ est une surjection de faisceaux et si $a$ est injectif, alors
$d$ est également injectif. Ainsi (4) implique (1). La preuve de l'équivalence de
(5) et (1) est omise.
\end{proof}

\begin{lemma}
\label{spaces-morphisms-lemma-immersions-monomorphisms}
Une immersion d'espaces algébriques est un monomorphisme.
En particulier, toute immersion est séparée.
\end{lemma}

\begin{proof}
Soit $f : X \to Y$ une immersion d'espaces algébriques.
Pour tout morphisme $Z \to Y$ avec $Z$ représentable, le changement de
base $Z \times_Y X \to Z$ est une immersion de schémas, donc
un monomorphisme ; voir
Schémas, lemme \ref{schemes-lemma-immersions-monomorphisms}.
Ainsi $f$ est représentable et est un monomorphisme.
\end{proof}

\noindent
Nous améliorerons le lemme suivant dans
Espaces décents, lemme
\ref{decent-spaces-lemma-monomorphism-toward-disjoint-union-dim-0-rings}.

\begin{lemma}
\label{spaces-morphisms-lemma-monomorphism-toward-field}
Soit $S$ un schéma. Soit $k$ un corps et soit $Z \to \Spec(k)$
un monomorphisme d'espaces algébriques sur $S$. Alors, soit
$Z = \emptyset$, soit $Z = \Spec(k)$.
\end{lemma}

\begin{proof}
D'après les
lemmes \ref{spaces-morphisms-lemma-monomorphism-separated} et
\ref{spaces-morphisms-lemma-separated-over-separated},
$Z$ est un espace algébrique séparé. Il existe donc un
sous-espace ouvert dense $Z' \subset Z$ qui est un schéma ; voir
Propriétés des espaces, proposition
\ref{spaces-properties-proposition-locally-quasi-separated-open-dense-scheme}.
D'après
Schémas, lemme \ref{schemes-lemma-mono-towards-spec-field},
soit $Z' = \emptyset$, soit $Z' \cong \Spec(k)$.
Dans le premier cas, on conclut que $Z = \emptyset$ et, dans le
second, que $Z' = Z = \Spec(k)$,
car $Z \to \Spec(k)$ est un monomorphisme qui est un
isomorphisme au-dessus de $Z'$.
\end{proof}

\begin{lemma}
\label{spaces-morphisms-lemma-monomorphism-injective-points}
Soit $S$ un schéma. Si $X \to Y$ est un monomorphisme d'espaces algébriques
sur $S$, alors $|X| \to |Y|$ est injectif.
\end{lemma}

\begin{proof}
Cela résulte immédiatement des définitions.
\end{proof}















\section{Image directe des faisceaux quasi-cohérents}
\label{spaces-morphisms-section-pushforward}

\noindent
Nous démontrons d'abord un lemme simple qui relie l'image directe des faisceaux de modules
pour un morphisme d'espaces algébriques à l'image directe des faisceaux de modules pour
un morphisme de schémas.

\begin{lemma}
\label{spaces-morphisms-lemma-compute-pushforward}
Soit $S$ un schéma.
Soit $f : X \to Y$ un morphisme d'espaces algébriques sur $S$.
Soit $U \to X$ un morphisme \'etale surjectif d'un schéma vers $X$.
Posons $R = U \times_X U$ et notons, comme d'habitude, $t, s : R \to U$ les morphismes
de projection. Notons $a : U \to Y$ et $b : R \to Y$ les morphismes
induits. Pour tout objet $\mathcal{F}$ de $\textit{Mod}(\mathcal{O}_X)$,
il existe une suite exacte
$$
0 \to f_*\mathcal{F} \to a_*(\mathcal{F}|_U) \to b_*(\mathcal{F}|_R)
$$
où la seconde flèche est la différence $t^* - s^*$.
\end{lemma}

\begin{proof}
Nous notons encore $\mathcal{F}$ son prolongement en un faisceau de modules sur
$X_{spaces, \etale}$ ; voir
Propriétés des espaces,
remarque \ref{spaces-properties-remark-explain-equivalence}.
Soit $V \to Y$ un objet de $Y_\etale$. Alors $V \times_Y X$ est un
objet de $X_{spaces, \etale}$ et, par définition,
$f_*\mathcal{F}(V) = \mathcal{F}(V \times_Y X)$. Puisque $U \to X$ est
\'etale surjectif, on voit que $\{V \times_Y U \to V \times_Y X\}$
est un recouvrement. De plus,
$(V \times_Y U) \times_X (V \times_Y U) = V \times_Y R$.
Par conséquent, la condition de faisceau pour $\mathcal{F}$ sur
$X_{spaces, \etale}$ donne une suite exacte courte
$$
0 \to \mathcal{F}(V \times_Y X)
\to \mathcal{F}(V \times_Y U) \to \mathcal{F}(V \times_Y R)
$$
où la seconde flèche est la différence des restrictions par $t$ et par $s$.
Cette suite exacte est fonctorielle en $V$, ce qui donne le lemme.
\end{proof}

\noindent
Soit $S$ un schéma. Soit $f : X \to Y$ un morphisme quasi-compact et
quasi-séparé entre les espaces algébriques représentables $X$
et $Y$ sur $S$. D'après
Descente,
proposition \ref{descent-proposition-equivalence-quasi-coherent-functorial},
le foncteur
$f_* : \QCoh(\mathcal{O}_X) \to \QCoh(\mathcal{O}_Y)$
coïncide avec le foncteur usuel lorsque l'on considère $X$ et $Y$ comme des schémas.

\medskip\noindent
Plus généralement, supposons que $f : X \to Y$ soit un morphisme représentable,
quasi-compact et quasi-séparé d'espaces algébriques sur $S$. Soit $V$ un schéma
et soit $V \to Y$ un morphisme \'etale surjectif. Soit $U = V \times_Y X$
et soit $f' : U \to V$ le changement de base de $f$. Pour tout
$\mathcal{O}_X$-module quasi-cohérent $\mathcal{F}$, on a
\begin{equation}
\label{spaces-morphisms-equation-representable-pushforward}
f'_*(\mathcal{F}|_U) = (f_*\mathcal{F})|_V,
\end{equation}
voir
Propriétés des espaces,
lemme \ref{spaces-properties-lemma-pushforward-etale-base-change-modules}.
Comme $f' : U \to V$ est un morphisme quasi-compact et quasi-séparé
de schémas, la remarque du paragraphe précédent permet de
calculer $f'_*(\mathcal{F}|_U)$ en considérant $\mathcal{F}|_U$ comme un
faisceau quasi-cohérent sur le schéma $U$ et $f'$ comme un morphisme de schémas.
Nous l'utiliserons fréquemment sans plus le mentionner.

\medskip\noindent
Le degré de généralité suivant consiste à considérer un morphisme quelconque
quasi-compact et quasi-séparé d'espaces algébriques.

\begin{lemma}
\label{spaces-morphisms-lemma-pushforward}
Soit $S$ un schéma.
Soit $f : X \to Y$ un morphisme d'espaces algébriques sur $S$.
Si $f$ est quasi-compact et quasi-séparé, alors $f_*$ transforme les
$\mathcal{O}_X$-modules quasi-cohérents en
$\mathcal{O}_Y$-modules quasi-cohérents.
\end{lemma}

\begin{proof}
Soit $\mathcal{F}$ un faisceau quasi-cohérent sur $X$. Il faut montrer que
$f_*\mathcal{F}$ est un faisceau quasi-cohérent sur $Y$. Pour cela, il suffit
de montrer que, pour tout schéma affine $V$ et tout morphisme \'etale $V \to Y$,
la restriction de $f_*\mathcal{F}$ à $V$ est quasi-cohérente ; voir
Propriétés des espaces,
lemme \ref{spaces-properties-lemma-characterize-quasi-coherent}.
Soit $f' : V \times_Y X \to V$
le changement de base de $f$ par $V \to Y$. Notons que $f'$ est encore
quasi-compact et quasi-séparé ; voir les
lemmes \ref{spaces-morphisms-lemma-base-change-quasi-compact} et
\ref{spaces-morphisms-lemma-base-change-separated}.
D'après (\ref{spaces-morphisms-equation-representable-pushforward}),
la restriction de $f_*\mathcal{F}$ à $V$ est l'image directe par $f'_*$ de la
restriction de $\mathcal{F}$ à $V \times_Y X$. Ainsi,
on peut remplacer $f$ par $f'$ et supposer que $Y$ est un schéma affine.

\medskip\noindent
Supposons que $Y$ soit un schéma affine. Puisque $f$ est quasi-compact, $X$
est quasi-compact. On peut donc choisir un schéma affine $U$ et un morphisme
\'etale surjectif $U \to X$ ; voir
Propriétés des espaces,
lemme \ref{spaces-properties-lemma-quasi-compact-affine-cover}.
Le lemme \ref{spaces-morphisms-lemma-compute-pushforward} donne une suite exacte
$$
0 \to f_*\mathcal{F} \to a_*(\mathcal{F}|_U) \to b_*(\mathcal{F}|_R).
$$
où $R = U \times_X U$.
Comme $X \to Y$ est quasi-séparé, $R \to U \times_Y U$
est un monomorphisme quasi-compact. Il en résulte que $R$ est un schéma
quasi-compact séparé (puisque $U$ et $Y$ sont affines à ce stade).
Ainsi $a : U \to Y$ et $b : R \to Y$ sont des morphismes quasi-compacts et
quasi-séparés de schémas. Par conséquent, d'après
Descente,
proposition \ref{descent-proposition-equivalence-quasi-coherent-functorial},
les faisceaux $a_*(\mathcal{F}|_U)$ et $b_*(\mathcal{F}|_R)$
sont quasi-cohérents (voir aussi la discussion qui précède ce lemme).
Il en résulte que $f_*\mathcal{F}$ est un noyau de
modules quasi-cohérents, et est donc lui-même quasi-cohérent ; voir
Propriétés des espaces,
lemme \ref{spaces-properties-lemma-properties-quasi-coherent}.
\end{proof}

\noindent
Les images directes supérieures sont étudiées dans
Cohomologie des espaces, section
\ref{spaces-cohomology-section-higher-direct-image}.









\section{Immersions}
\label{spaces-morphisms-section-immersions}

\noindent
Les immersions ouvertes, fermées et localement fermées d'espaces algébriques ont été définies dans
Espaces, section \ref{spaces-section-Zariski}.
Plus précisément, un morphisme d'espaces algébriques est une
{\it immersion fermée} (resp. une {\it immersion ouverte}, resp.\ une {\it immersion})
s'il est représentable et est une immersion fermée (resp.\ une immersion ouverte,
resp.\ une immersion) au sens de la section \ref{spaces-morphisms-section-representable}.

\medskip\noindent
En particulier, ces types de morphismes sont stables par changement de base
et par composition de morphismes dans la catégorie des espaces
algébriques sur $S$ ; voir
Espaces, lemmes \ref{spaces-lemma-composition-immersions} et
\ref{spaces-lemma-base-change-immersions}.

\begin{lemma}
\label{spaces-morphisms-lemma-closed-immersion-local}
Soit $S$ un schéma. Soit $f : X \to Y$ un morphisme d'espaces algébriques
sur $S$. Les conditions suivantes sont équivalentes :
\begin{enumerate}
\item $f$ est une immersion fermée (resp.\ une immersion ouverte, resp.\ une immersion),
\item pour tout schéma $Z$ et tout morphisme $Z \to Y$, le morphisme
$Z \times_Y X \to Z$ est une immersion fermée (resp.\ une immersion ouverte,
resp.\ une immersion),
\item pour tout schéma affine $Z$ et tout morphisme
$Z \to Y$, le morphisme $Z \times_Y X \to Z$ est une immersion fermée
(resp.\ une immersion ouverte, resp.\ une immersion),
\item il existe un schéma $V$ et un morphisme \'etale surjectif
$V \to Y$ tel que $V \times_Y X \to V$ soit une immersion fermée
(resp.\ une immersion ouverte, resp.\ une immersion), et
\item il existe un recouvrement de Zariski $Y = \bigcup Y_i$ tel que
chacun des morphismes $f^{-1}(Y_i) \to Y_i$ soit une immersion fermée
(resp.\ une immersion ouverte, resp.\ une immersion).
\end{enumerate}
\end{lemma}

\begin{proof}
Puisque tout changement de base d'une
immersion fermée (resp.\ immersion ouverte, resp.\ immersion)
en est encore une, il est clair que (1) implique (2) et que (2) implique (3).
De plus, (3) implique (4), car on peut prendre pour $V$ une somme disjointe de
schémas affines ; voir
Propriétés des espaces,
lemme \ref{spaces-properties-lemma-cover-by-union-affines}.

\medskip\noindent
Supposons que $V \to Y$ soit comme en (4).
Soit $\mathcal{P}$ la propriété d'être une
immersion fermée (resp.\ une immersion ouverte, resp.\ une immersion)
pour les morphismes de schémas. Notons que la propriété $\mathcal{P}$
est stable par tout changement de base et locale pour la topologie fppf sur la
base (voir la section \ref{spaces-morphisms-section-representable}).
De plus, les morphismes de type $\mathcal{P}$ sont séparés et
localement quasi-finis (dans chacun des trois cas ; voir
Schémas, lemme \ref{schemes-lemma-immersions-monomorphisms}, et
Morphismes, lemme \ref{morphisms-lemma-immersion-locally-quasi-finite}).
Ainsi, d'après
Compléments sur les morphismes, lemme
\ref{more-morphisms-lemma-separated-locally-quasi-finite-morphisms-fppf-descend}
les morphismes de type $\mathcal{P}$ satisfont à la descente pour les recouvrements fppf. Ainsi
Espaces, lemme \ref{spaces-lemma-morphism-sheaves-with-P-effective-descent-etale}
s'applique, et l'on voit que $X \to Y$ est représentable et possède la propriété
$\mathcal{P}$ ; autrement dit, (1) est satisfaite.

\medskip\noindent
L'équivalence de (1) et (5) résulte du fait que
$\mathcal{P}$ est locale sur le but pour la topologie de Zariski (puisque nous avons vu
ci-dessus que $\mathcal{P}$ est en fait locale sur le but pour la topologie fppf).
\end{proof}

\begin{lemma}
\label{spaces-morphisms-lemma-immersion-permanence}
Soit $S$ un schéma.
Soient $Z \to Y \to X$ des morphismes d'espaces algébriques sur $S$.
\begin{enumerate}
\item Si $Z \to X$ est représentable, localement de type fini, localement
quasi-fini, séparé et est un monomorphisme, alors $Z \to Y$ est
représentable, localement de type fini, localement quasi-fini,
séparé et est un monomorphisme.
\item Si $Z \to X$ est une immersion et si $Y \to X$ est localement séparé,
alors $Z \to Y$ est une immersion.
\item Si $Z \to X$ est une immersion fermée et si $Y \to X$ est séparé,
alors $Z \to Y$ est une immersion fermée.
\end{enumerate}
\end{lemma}

\begin{proof}
Dans chaque cas, la preuve consiste à considérer le diagramme commutatif
$$
\xymatrix{
Z \ar[r] \ar[rd] & Y \times_X Z \ar[r] \ar[d] & Z \ar[d] \\
& Y \ar[r] & X
}
$$
où le composé des flèches horizontales supérieures est l'identité.
Démontrons (1). La première flèche horizontale est une section de
$Y \times_X Z \to Z$ ; elle est donc représentable, localement de type fini,
localement quasi-finie, séparée et est un monomorphisme d'après le
lemme \ref{spaces-morphisms-lemma-section-immersion}.
La flèche $Y \times_X Z \to Y$ est un changement de base de $Z \to X$ ;
elle est donc représentable, localement de type fini,
localement quasi-finie, séparée et est un monomorphisme
(car chacune de ces propriétés des morphismes de schémas est stable
par changement de base ; voir
Espaces, remarque \ref{spaces-remark-list-properties-stable-base-change}).
Il en va donc de même du composé (car chacune de ces propriétés des
morphismes de schémas est stable par composition ; voir Espaces, remarque
\ref{spaces-remark-list-properties-stable-composition}).
Ceci prouve (1). Les autres assertions se démontrent exactement de la même façon.
\end{proof}

\begin{lemma}
\label{spaces-morphisms-lemma-immersion-when-closed}
Soit $S$ un schéma. Soit $i : Z \to X$ une immersion d'espaces
algébriques sur $S$. Alors $|i| : |Z| \to |X|$ est un homéomorphisme sur une
partie localement fermée, et $i$ est une immersion fermée si et seulement si
l'image $|i|(|Z|) \subset |X|$ est une partie fermée.
\end{lemma}

\begin{proof}
La première assertion est le lemme
\ref{spaces-properties-lemma-subspace-induced-topology} de Propriétés des espaces.
Soit $U$ un schéma et soit $U \to X$ un morphisme \'etale surjectif.
Par hypothèse, $T = U \times_X Z$ est un schéma et le morphisme $j : T \to U$
est une immersion de schémas. D'après le lemme \ref{spaces-morphisms-lemma-closed-immersion-local},
le morphisme $i$ est une immersion fermée si et seulement si $j$ est une immersion
fermée. D'après Schémas, lemme \ref{schemes-lemma-immersion-when-closed},
cela équivaut à ce que $j(T)$ soit fermé dans $U$.
Or la partie $j(T) \subset U$ est l'image inverse de
$|i|(|Z|) \subset |X|$ ; voir
Propriétés des espaces, lemme \ref{spaces-properties-lemma-points-cartesian}.
Ceci achève la démonstration.
\end{proof}

\begin{remark}
\label{spaces-morphisms-remark-immersion}
Soit $S$ un schéma. Soit $i : Z \to X$ une immersion d'espaces
algébriques sur $S$. Puisque $i$ est un monomorphisme, on peut considérer $|Z|$ comme
une partie de $|X|$ ; c'est ce que nous ferons dans la suite de cette remarque.
Soit $\partial |Z|$ la frontière de $|Z|$ dans
l'espace topologique $|X|$. En formule,
$$
\partial |Z| = \overline{|Z|} \setminus |Z|.
$$
Soit $\partial Z$ le sous-espace fermé réduit de $X$ tel que
$|\partial Z| = \partial |Z|$,
obtenu en prenant la structure réduite induite de sous-espace fermé ; voir
Propriétés des espaces,
définition \ref{spaces-properties-definition-reduced-induced-space}.
Par construction, $|Z|$ est fermé dans
$|X| \setminus |\partial Z| = |X \setminus \partial Z|$.
Ainsi, toute immersion d'espaces algébriques peut bien se
factoriser en une immersion fermée suivie d'une immersion ouverte
(mais pas dans l'autre ordre en général ; voir
Morphismes, exemple \ref{morphisms-example-thibaut}).
\end{remark}

\begin{remark}
\label{spaces-morphisms-remark-space-structure-locally-closed-subset}
Soit $S$ un schéma. Soit $X$ un espace algébrique sur $S$.
Soit $T \subset |X|$ une partie localement fermée.
Soit $\partial T$ la frontière de $T$ dans
l'espace topologique $|X|$. En formule,
$$
\partial T = \overline{T} \setminus T.
$$
Soit $U \subset X$ le sous-espace ouvert de $X$ tel que
$|U| = |X| \setminus \partial T$ ; voir
Propriétés des espaces, lemme \ref{spaces-properties-lemma-open-subspaces}.
Soit $Z$ le sous-espace fermé réduit de $U$ tel que
$|Z| = T$, obtenu en prenant la structure réduite induite
de sous-espace fermé ; voir
Propriétés des espaces,
définition \ref{spaces-properties-definition-reduced-induced-space}.
Par construction, $Z \to U$ est une immersion fermée d'espaces algébriques
et $U \to X$ est une immersion ouverte ; par conséquent,
$Z \to X$ est une immersion d'espaces algébriques sur $S$ (voir
Espaces, lemme \ref{spaces-lemma-composition-immersions}).
Notons que $Z$ est un espace algébrique réduit et que
$|Z| = T$ comme parties de $|X|$. On dit parfois que
$Z$ est la {\it structure réduite induite de sous-espace} sur $T$.
\end{remark}

\begin{lemma}
\label{spaces-morphisms-lemma-factor-the-other-way}
Soit $S$ un schéma. Soit $Z \to X$ une immersion d'espaces algébriques sur
$S$. Supposons $Z \to X$ quasi-compact.
Il existe une factorisation $Z \to \overline{Z} \to X$ où
$Z \to \overline{Z}$ est une immersion ouverte et $\overline{Z} \to X$
est une immersion fermée.
\end{lemma}

\begin{proof}
Soit $U$ un schéma et soit $U \to X$ \'etale surjectif.
Notons comme d'habitude $R = U \times_X U$, avec pour projections
$s, t : R \to U$. Posons $T = Z \times_X U$. Soit $\overline{T} \subset U$
l'image schématique de $T \to U$. Notons que
$s^{-1}\overline{T} = t^{-1}\overline{T}$, car la formation des
images schématiques de morphismes quasi-compacts commute
au changement de base plat ; voir
Morphismes, lemme \ref{morphisms-lemma-flat-base-change-scheme-theoretic-image}.
On obtient donc un sous-espace fermé $\overline{Z} \subset X$ dont
le changement de base à $U$ est $\overline{T}$ ; voir
Propriétés des espaces, lemme
\ref{spaces-properties-lemma-subspaces-presentation}.
D'après Morphismes, lemme \ref{morphisms-lemma-quasi-compact-immersion},
le morphisme $T \to \overline{T}$
est une immersion ouverte. Il s'ensuit que $Z \to \overline{Z}$ est
une immersion ouverte, ce qui donne le résultat.
\end{proof}


















\section{Immersions fermées}
\label{spaces-morphisms-section-closed-immersions}

\noindent
Dans cette section, nous précisons certains résultats obtenus précédemment sur les
immersions d'espaces algébriques. Voir
Espaces, section \ref{spaces-section-Zariski},
et la
section \ref{spaces-morphisms-section-immersions} du présent chapitre.
Cette section est l'analogue de
Morphismes, section \ref{morphisms-section-closed-immersions},
pour les espaces algébriques.

\begin{lemma}
\label{spaces-morphisms-lemma-closed-immersion-ideals}
Soit $S$ un schéma.
Soit $X$ un espace algébrique sur $S$.
Pour toute immersion fermée $i : Z \to X$, le faisceau
$i_*\mathcal{O}_Z$ est un $\mathcal{O}_X$-module quasi-cohérent ; l'application
$i^\sharp : \mathcal{O}_X \to i_*\mathcal{O}_Z$ est surjective et son
noyau est un faisceau quasi-cohérent d'idéaux. La règle
$Z \mapsto \Ker(\mathcal{O}_X \to i_*\mathcal{O}_Z)$
définit une bijection renversant les inclusions
$$
\begin{matrix}
\text{sous-espaces fermés}\\
Z \subset X
\end{matrix}
\longrightarrow
\begin{matrix}
\text{faisceaux quasi-cohérents}\\
\text{d'idéaux }\mathcal{I} \subset \mathcal{O}_X
\end{matrix}
$$
De plus, étant donné un sous-espace fermé $Z$ correspondant au faisceau
quasi-cohérent d'idéaux $\mathcal{I} \subset \mathcal{O}_X$, un morphisme d'espaces
algébriques $h : Y \to X$ se factorise par $Z$ si et seulement si l'application
$h^*\mathcal{I} \to h^*\mathcal{O}_X = \mathcal{O}_Y$ est nulle.
\end{lemma}

\begin{proof}
Soit $U \to X$ un morphisme \'etale surjectif dont la source est un schéma.
Considérons le diagramme
$$
\xymatrix{
U \times_X Z \ar[r] \ar[d]_{i'} & Z \ar[d]^i \\
U \ar[r] & X
}
$$
D'après le
lemme \ref{spaces-morphisms-lemma-closed-immersion-local},
$i$ est une immersion fermée
si et seulement si $i'$ est une immersion fermée. D'après
Propriétés des espaces,
lemme \ref{spaces-properties-lemma-pushforward-etale-base-change-modules},
$i'_*\mathcal{O}_{U \times_X Z}$ est la restriction de
$i_*\mathcal{O}_Z$ à $U$. Ainsi, les assertions concernant
$\mathcal{O}_X \to i_*\mathcal{O}_Z$ sont équivalentes aux
assertions correspondantes concernant
$\mathcal{O}_U \to i'_*\mathcal{O}_{U \times_X Z}$.
Puisque $i'$ est une immersion fermée de schémas, ces résultats découlent de
Morphismes, lemme \ref{morphisms-lemma-closed-immersion}.

\medskip\noindent
Démontrons que, pour un faisceau quasi-cohérent
d'idéaux $\mathcal{I} \subset \mathcal{O}_X$, la formule
$$
Z(T) = \{h : T \to X \mid h^*\mathcal{I} \to \mathcal{O}_T
\text{ est nulle}\}
$$
définit un sous-espace fermé de $X$. C'est clairement un sous-foncteur de $X$.
Pour montrer que $Z \to X$ est représentable par des immersions fermées, soit
$\varphi : U \to X$ un morphisme d'un schéma vers $X$. Alors
$Z \times_X U$ est représenté par le sous-foncteur analogue de $U$ correspondant
au faisceau d'idéaux $\Im(\varphi^*\mathcal{I} \to \mathcal{O}_U)$. D'après
Propriétés des espaces,
lemme \ref{spaces-properties-lemma-pullback-quasi-coherent},
le $\mathcal{O}_U$-module $\varphi^*\mathcal{I}$ est quasi-cohérent
sur $U$, et donc $\Im(\varphi^*\mathcal{I} \to \mathcal{O}_U)$
est un faisceau quasi-cohérent d'idéaux sur $U$. D'après
Schémas, lemme \ref{schemes-lemma-characterize-closed-subspace},
on conclut que $Z \times_X U$ est représenté par le sous-schéma fermé
de $U$ associé à $\Im(\varphi^*\mathcal{I} \to \mathcal{O}_U)$.
Ainsi $Z$ est un sous-espace fermé de $X$.

\medskip\noindent
Dans la formule ci-dessus définissant $Z$, les arguments $T$ sont des schémas, puisque les espaces
algébriques sont des faisceaux sur $(\Sch/S)_{fppf}$. Nous omettons la vérification
que la même formule reste vraie lorsque $T$ est un espace algébrique.
\end{proof}

\begin{definition}
\label{spaces-morphisms-definition-inverse-image-closed-subspace}
Soit $S$ un schéma. Soit $f : Y \to X$ un morphisme d'espaces algébriques
sur $S$. Soit $Z \subset X$ un sous-espace fermé. L'
{\it image inverse $f^{-1}(Z)$ du sous-espace fermé $Z$}
est le sous-espace fermé $Z \times_X Y$ de $Y$.
\end{definition}

\noindent
Cette définition a un sens d'après le lemme \ref{spaces-morphisms-lemma-closed-immersion-local}.
Si $\mathcal{I} \subset \mathcal{O}_X$ est le faisceau quasi-cohérent
d'idéaux correspondant à $Z$ par le lemme \ref{spaces-morphisms-lemma-closed-immersion-ideals}, alors
$f^{-1}\mathcal{I}\mathcal{O}_Y = \Im(f^*\mathcal{I} \to \mathcal{O}_Y)$
est le faisceau d'idéaux correspondant à $f^{-1}(Z)$.

\begin{lemma}
\label{spaces-morphisms-lemma-closed-immersion-quasi-compact}
Une immersion fermée d'espaces algébriques est quasi-compacte.
\end{lemma}

\begin{proof}
Ceci résulte de
Schémas, lemme \ref{schemes-lemma-closed-immersion-quasi-compact},
par des principes généraux ; voir
Espaces, lemme
\ref{spaces-lemma-representable-transformations-property-implication}.
\end{proof}

\begin{lemma}
\label{spaces-morphisms-lemma-closed-immersion-separated}
Une immersion fermée d'espaces algébriques est séparée.
\end{lemma}

\begin{proof}
Ceci résulte de
Schémas, lemme \ref{schemes-lemma-immersions-monomorphisms},
par des principes généraux ; voir
Espaces, lemme
\ref{spaces-lemma-representable-transformations-property-implication}.
\end{proof}

\begin{lemma}
\label{spaces-morphisms-lemma-closed-immersion-push-pull}
Soit $S$ un schéma. Soit $i : Z \to X$ une immersion fermée d'espaces
algébriques sur $S$.
\begin{enumerate}
\item Le foncteur
$$
i_{small, *} :
\Sh(Z_\etale)
\longrightarrow
\Sh(X_\etale)
$$
est pleinement fidèle et son image essentielle est formée des faisceaux d'ensembles
$\mathcal{F}$ sur $X_\etale$ dont la restriction à $X \setminus Z$ est
isomorphe à $*$, et
\item le foncteur
$$
i_{small, *} :
\textit{Ab}(Z_\etale)
\longrightarrow
\textit{Ab}(X_\etale)
$$
est pleinement fidèle et son image essentielle est formée des faisceaux abéliens sur
$X_\etale$ dont le support est contenu dans $|Z|$.
\end{enumerate}
Dans les deux cas, $i_{small}^{-1}$ est un inverse à gauche du foncteur
$i_{small, *}$.
\end{lemma}

\begin{proof}
Soit $U$ un schéma et soit $U \to X$ \'etale surjectif.
Posons $V = Z \times_X U$. Alors $V$ est un schéma et $i' : V \to U$ est
une immersion fermée de schémas. D'après
Propriétés des espaces,
lemme \ref{spaces-properties-lemma-pushforward-etale-base-change},
pour tout faisceau $\mathcal{G}$ sur $Z$, on a
$$
(i_{small}^{-1}i_{small, *}\mathcal{G})|_V =
(i')_{small}^{-1}i'_{small, *}(\mathcal{G}|_V)
$$
D'après
Cohomologie \'etale, proposition
\ref{etale-cohomology-proposition-closed-immersion-pushforward},
l'application
$(i')_{small}^{-1}i'_{small, *}(\mathcal{G}|_V) \to \mathcal{G}|_V$
est un isomorphisme. Puisque $V \to Z$ est surjectif et \'etale, cela entraîne
que $i_{small}^{-1}i_{small, *}\mathcal{G} \to \mathcal{G}$ est un
isomorphisme. Ceci implique clairement que $i_{small, *}$ est pleinement fidèle ; voir
Sites, lemme \ref{sites-lemma-exactness-properties}.
Pour prouver l'assertion sur l'image essentielle, considérons un faisceau d'ensembles
$\mathcal{F}$ sur $X_\etale$ dont la restriction à $X \setminus Z$ est
isomorphe à $*$. Comme dans la preuve de
Cohomologie \'etale, proposition
\ref{etale-cohomology-proposition-closed-immersion-pushforward}
nous considérons le morphisme d'adjonction
$$
\mathcal{F} \longrightarrow i_{small, *}i_{small}^{-1}\mathcal{F}.
$$
Comme dans la première partie, le résultat correspondant pour le cas des
schémas montre que la restriction de ce morphisme à $U$ est un isomorphisme.
Puisque $U$ est un recouvrement \'etale de $X$, on
conclut qu'il s'agit d'un isomorphisme.
\end{proof}

\begin{lemma}
\label{spaces-morphisms-lemma-stalk-push-closed}
Soit $S$ un schéma. Soit $i : Z \to X$ une immersion fermée d'espaces
algébriques sur $S$. Soit $\overline{z}$ un point géométrique de $Z$
d'image $\overline{x}$ dans $X$. Alors
$(i_{small, *}\mathcal{F})_{\overline{x}} = \mathcal{F}_{\overline{z}}$
pour tout faisceau $\mathcal{F}$ sur $Z_\etale$.
\end{lemma}

\begin{proof}
Choisissons un voisinage \'etale $(U, \overline{u})$ de $\overline{x}$.
Alors la fibre $(i_{small, *}\mathcal{F})_{\overline{x}}$
est la fibre de $i_{small, *}\mathcal{F}|_U$ en $\overline{u}$.
D'après Propriétés des espaces,
lemme \ref{spaces-properties-lemma-pushforward-etale-base-change},
on peut remplacer $X$ par $U$ et $Z$ par $Z \times_X U$.
Alors $Z \to X$ est une immersion fermée de schémas, et le résultat est
Cohomologie \'etale, lemme
\ref{etale-cohomology-lemma-stalk-pushforward-closed-immersion}.
\end{proof}

\noindent
Le lemme suivant vaut plus généralement dans le cadre d'une immersion fermée
de topos (insérer ici une référence future).

\begin{lemma}
\label{spaces-morphisms-lemma-closed-immersion-rings}
Soit $S$ un schéma. Soit $i : Z \to X$ une immersion fermée d'espaces
algébriques sur $S$. Soit $\mathcal{A}$ un faisceau d'anneaux sur $X_\etale$.
Soit $\mathcal{B}$ un faisceau d'anneaux sur $Z_\etale$.
Soit $\varphi : \mathcal{A} \to i_{small, *}\mathcal{B}$
un homomorphisme de faisceaux d'anneaux, ce qui donne un
morphisme de topos annelés
$$
f : (\Sh(Z_\etale), \mathcal{B}) \longrightarrow (\Sh(X_\etale), \mathcal{A}).
$$
Pour un faisceau de $\mathcal{A}$-modules $\mathcal{F}$ et un
faisceau de $\mathcal{B}$-modules $\mathcal{G}$, l'application canonique
$$
\mathcal{F} \otimes_\mathcal{A} f_*\mathcal{G}
\longrightarrow
f_*(f^*\mathcal{F} \otimes_\mathcal{B} \mathcal{G})
$$
est un isomorphisme.
\end{lemma}

\begin{proof}
Cette application est l'adjointe de l'application
$$
f^*\mathcal{F} \otimes_\mathcal{B}
f^* f_*\mathcal{G} =
f^*(\mathcal{F} \otimes_\mathcal{A} f_*\mathcal{G})
\longrightarrow
f^*\mathcal{F} \otimes_\mathcal{B} \mathcal{G}
$$
provenant de $\text{id} : f^*\mathcal{F} \to f^*\mathcal{F}$
et du morphisme d'adjonction $f^* f_*\mathcal{G} \to \mathcal{G}$.
Pour voir que cette application est un isomorphisme, il suffit de le vérifier sur les fibres
(Propriétés des espaces, théorème
\ref{spaces-properties-theorem-exactness-stalks}).
Soit $\overline{z} : \Spec(k) \to Z$ un point géométrique
d'image $\overline{x} = i \circ \overline{z} : \Spec(k) \to X$.
En explicitant l'action de notre application sur les fibres, on voit qu'il
faut montrer
$$
\mathcal{F}_{\overline{x}}
\otimes_{\mathcal{A}_{\overline{x}}}
\mathcal{G}_{\overline{z}} =
(\mathcal{F}_{\overline{x}}
\otimes_{\mathcal{A}_{\overline{x}}}
\mathcal{B}_{\overline{z}}) \otimes_{\mathcal{B}_{\overline{z}}}
\mathcal{G}_{\overline{z}}
$$
ce qui est vrai. Nous avons utilisé ici que
la formation des produits tensoriels commute à celle des fibres, ainsi que le
comportement des fibres par image inverse (Propriétés des espaces,
lemme \ref{spaces-properties-lemma-stalk-pullback}) et par image directe
le long d'une immersion fermée (lemme
\ref{spaces-morphisms-lemma-stalk-push-closed}).
\end{proof}






\section{Immersions fermées et faisceaux quasi-cohérents}
\label{spaces-morphisms-section-closed-immersions-quasi-coherent}

\noindent
Cette section est l'analogue de
Morphismes, section \ref{morphisms-section-closed-immersions-quasi-coherent}.

\begin{lemma}
\label{spaces-morphisms-lemma-i-star-equivalence}
Soit $S$ un schéma. Soit $i : Z \to X$ une immersion fermée d'espaces
algébriques sur $S$. Soit $\mathcal{I} \subset \mathcal{O}_X$ le faisceau
quasi-cohérent d'idéaux définissant $Z$.
\begin{enumerate}
\item Pour tout $\mathcal{O}_X$-module $\mathcal{F}$, le morphisme d'adjonction
$\mathcal{F} \to i_*i^*\mathcal{F}$ induit un isomorphisme
$\mathcal{F}/\mathcal{I}\mathcal{F} \cong i_*i^*\mathcal{F}$.
\item Le foncteur $i^*$ est un inverse à gauche de $i_*$, c'est-à-dire que, pour tout
$\mathcal{O}_Z$-module $\mathcal{G}$, le morphisme d'adjonction
$i^*i_*\mathcal{G} \to \mathcal{G}$ est un isomorphisme.
\item Le foncteur
$$
i_* :
\QCoh(\mathcal{O}_Z)
\longrightarrow
\QCoh(\mathcal{O}_X)
$$
est exact et pleinement fidèle, et son image essentielle est formée des
$\mathcal{O}_X$-modules quasi-cohérents $\mathcal{F}$ tels que $\mathcal{I}\mathcal{F} = 0$.
\end{enumerate}
\end{lemma}

\begin{proof}
Au cours de cette démonstration, nous travaillons exclusivement avec les faisceaux sur
les petits sites étales, et nous employons $i_*, i^{-1}, \ldots$
pour désigner l'image directe et l'image inverse des faisceaux de groupes abéliens
au lieu de $i_{small, *}, i_{small}^{-1}$.

\medskip\noindent
Soit $\mathcal{F}$ un $\mathcal{O}_X$-module. Le
lemme \ref{spaces-morphisms-lemma-closed-immersion-rings}, appliqué avec
$\mathcal{A} = \mathcal{O}_X$ et
$\mathcal{G} = \mathcal{B} = \mathcal{O}_Z$, montre que
$i_*i^*\mathcal{F} = \mathcal{F} \otimes_{\mathcal{O}_X} i_*\mathcal{O}_Z$.
Le
lemme \ref{spaces-morphisms-lemma-closed-immersion-ideals}
nous donne une suite exacte courte
$$
0 \to \mathcal{I} \to \mathcal{O}_X \to i_*\mathcal{O}_Z \to 0
$$
Il résulte des propriétés du produit tensoriel que
$\mathcal{F} \otimes_{\mathcal{O}_X} i_*\mathcal{O}_Z
= \mathcal{F}/\mathcal{I}\mathcal{F}$. Ceci démontre (1) (à ceci près
que nous omettons de vérifier que le morphisme est induit par le
morphisme d'adjonction).

\medskip\noindent
Soit $\mathcal{G}$ un $\mathcal{O}_Z$-module quelconque. Le
lemme \ref{spaces-morphisms-lemma-closed-immersion-push-pull}
montre que $i^{-1}i_*\mathcal{G} = \mathcal{G}$.
Pour démontrer (2), il nous suffit donc de montrer que le morphisme canonique
$\mathcal{G} \otimes_{i^{-1}\mathcal{O}_X} \mathcal{O}_Z \to \mathcal{G}$
est un isomorphisme. Cela résulte des propriétés générales des produits tensoriels
dès que l'on montre que $i^{-1}\mathcal{O}_X \to \mathcal{O}_Z$ est surjectif. D'après le
lemme \ref{spaces-morphisms-lemma-closed-immersion-push-pull},
il suffit de prouver que
$i_*i^{-1}\mathcal{O}_X \to i_*\mathcal{O}_Z$
est surjectif. Comme le morphisme surjectif $\mathcal{O}_X \to i_*\mathcal{O}_Z$
se factorise par ce morphisme, nous obtenons (2).

\medskip\noindent
Enfin, démontrons la partie la plus intéressante du lemme, à savoir (3).
Une immersion fermée est quasi-compacte et séparée ; voir les
lemmes \ref{spaces-morphisms-lemma-closed-immersion-quasi-compact} et
\ref{spaces-morphisms-lemma-closed-immersion-separated}. Le
lemme \ref{spaces-morphisms-lemma-pushforward}
s'applique donc, et l'image directe d'un faisceau quasi-cohérent
sur $Z$ est bien un faisceau quasi-cohérent sur $X$.
Nous obtenons ainsi le foncteur
$i^{QCoh}_* : \QCoh(\mathcal{O}_Z)
\to \QCoh(\mathcal{O}_X)$.
Il résulte de (2) que $i^{QCoh}_*$ est pleinement fidèle, puisque
$i^*$ est son adjoint à gauche et que le morphisme d'adjonction de (2) est un isomorphisme.

\medskip\noindent
Décrivons maintenant l'image essentielle du
foncteur $i_*$. Il est clair que $\mathcal{I}(i_*\mathcal{G}) = 0$
pour tout $\mathcal{O}_Z$-module, puisque $\mathcal{I}$ est le noyau
du morphisme $\mathcal{O}_X \to i_*\mathcal{O}_Z$, qui est précisément celui
par lequel on munit d'une structure de $\mathcal{O}_X$-module le faisceau $i_*\mathcal{G}$.
Supposons ensuite que $\mathcal{F}$ soit un $\mathcal{O}_X$-module
quasi-cohérent tel que $\mathcal{I}\mathcal{F} = 0$.
Alors $\mathcal{F}$ est un $i_*\mathcal{O}_Z$-module,
car $i_*\mathcal{O}_Z = \mathcal{O}_X/\mathcal{I}$. En particulier,
son support est contenu dans $|Z|$. Le
lemme \ref{spaces-morphisms-lemma-closed-immersion-push-pull}
montre que $\mathcal{F} \cong i_*\mathcal{G}$ pour un certain
$\mathcal{O}_Z$-module $\mathcal{G}$. Le seul détail restant consiste
à montrer que $\mathcal{G}$ est quasi-cohérent. Or,
$\mathcal{G} \cong i^*\mathcal{F}$ par (2), et
Propriétés des espaces,
lemme \ref{spaces-properties-lemma-pullback-quasi-coherent}, donne le résultat.
\end{proof}

\noindent
Soit $i : Z \to X$ une immersion fermée d'espaces algébriques.
D'après le lemme précédent, nous noterons souvent,
par abus de notation, $\mathcal{F}$ le faisceau $i_*\mathcal{F}$ sur $X$.

\begin{lemma}
\label{spaces-morphisms-lemma-largest-quasi-coherent-subsheaf}
Soit $S$ un schéma. Soit $X$ un espace algébrique sur $S$.
Soit $\mathcal{F}$ un $\mathcal{O}_X$-module quasi-cohérent.
Soit $\mathcal{G} \subset \mathcal{F}$ un $\mathcal{O}_X$-sous-module.
Il existe un unique $\mathcal{O}_X$-sous-module quasi-cohérent
$\mathcal{G}' \subset \mathcal{G}$ ayant la propriété suivante :
pour tout $\mathcal{O}_X$-module quasi-cohérent $\mathcal{H}$, le morphisme
$$
\Hom_{\mathcal{O}_X}(\mathcal{H}, \mathcal{G}')
\longrightarrow
\Hom_{\mathcal{O}_X}(\mathcal{H}, \mathcal{G})
$$
est bijectif. En particulier, $\mathcal{G}'$ est le plus grand
$\mathcal{O}_X$-sous-module quasi-cohérent de $\mathcal{F}$ contenu dans $\mathcal{G}$.
\end{lemma}

\begin{proof}
Soit $\mathcal{G}_a$, $a \in A$, l'ensemble des
$\mathcal{O}_X$-sous-modules quasi-cohérents contenus dans $\mathcal{G}$.
Alors l'image $\mathcal{G}'$ de
$$
\bigoplus\nolimits_{a \in A} \mathcal{G}_a \longrightarrow \mathcal{F}
$$
est quasi-cohérente, car l'image d'un morphisme de faisceaux quasi-cohérents
sur $X$ est quasi-cohérente et toute somme directe de faisceaux quasi-cohérents
est quasi-cohérente ; voir le lemme de
Propriétés des espaces,
lemme \ref{spaces-properties-lemma-properties-quasi-coherent}.
Le module $\mathcal{G}'$ est contenu dans $\mathcal{G}$. C'est donc le
plus grand $\mathcal{O}_X$-sous-module quasi-cohérent contenu dans $\mathcal{G}$.

\medskip\noindent
Pour démontrer la formule, soit $\mathcal{H}$ un
$\mathcal{O}_X$-module quasi-cohérent et soit $\alpha : \mathcal{H} \to \mathcal{G}$
un morphisme de $\mathcal{O}_X$-modules. L'image du composé
$\mathcal{H} \to \mathcal{G} \to \mathcal{F}$ est quasi-cohérente,
comme image d'un morphisme de faisceaux quasi-cohérents. Elle est donc contenue
dans $\mathcal{G}'$. Par conséquent, $\alpha$ se factorise par $\mathcal{G}'$,
comme souhaité.
\end{proof}

\begin{lemma}
\label{spaces-morphisms-lemma-i-upper-shriek}
Soit $S$ un schéma.
Soit $i : Z \to X$ une immersion fermée d'espaces algébriques sur $S$.
Il existe un foncteur\footnote{Cette notation n'est vraisemblablement pas standard.}
$i^! : \QCoh(\mathcal{O}_X) \to \QCoh(\mathcal{O}_Z)$
qui est un adjoint à droite de $i_*$. (Comparer avec
Modules, lemme \ref{modules-lemma-i-star-right-adjoint}.)
\end{lemma}

\begin{proof}
Étant donné un $\mathcal{O}_X$-module quasi-cohérent $\mathcal{G}$, considérons
le sous-faisceau $\mathcal{H}_Z(\mathcal{G})$ de $\mathcal{G}$ formé des sections locales
annulées par le faisceau d'idéaux $\mathcal{I}$ définissant $Z$. D'après le
lemme \ref{spaces-morphisms-lemma-largest-quasi-coherent-subsheaf},
il existe canoniquement un plus grand $\mathcal{O}_X$-sous-module quasi-cohérent,
noté $\mathcal{H}_Z(\mathcal{G})'$. Par construction, nous avons
$$
\Hom_{\mathcal{O}_X}(i_*\mathcal{F}, \mathcal{H}_Z(\mathcal{G})')
=
\Hom_{\mathcal{O}_X}(i_*\mathcal{F}, \mathcal{G})
$$
pour tout $\mathcal{O}_Z$-module quasi-cohérent $\mathcal{F}$.
Nous pouvons donc poser $i^!\mathcal{G} = i^*(\mathcal{H}_Z(\mathcal{G})')$.
Nous omettons les détails.
\end{proof}

\noindent
Au moyen de la correspondance biunivoque ($1$-à-$1$) entre les faisceaux quasi-cohérents
d'idéaux et les sous-espaces fermés (voir le
lemme \ref{spaces-morphisms-lemma-closed-immersion-ideals}),
nous pouvons définir les intersections et réunions schématiques
de sous-espaces fermés.

\begin{definition}
\label{spaces-morphisms-definition-scheme-theoretic-intersection-union}
Soit $S$ un schéma. Soit $X$ un espace algébrique sur $S$.
Soient $Z, Y \subset X$ des sous-espaces fermés
correspondant aux faisceaux quasi-cohérents d'idéaux
$\mathcal{I}, \mathcal{J} \subset \mathcal{O}_X$.
L'{\it intersection schématique} de $Z$ et $Y$
est le sous-espace fermé de $X$ défini par $\mathcal{I} + \mathcal{J}$.
La {\it réunion schématique} de $Z$ et $Y$
est le sous-espace fermé de $X$ défini par
$\mathcal{I} \cap \mathcal{J}$.
\end{definition}

\noindent
Il est clair que la formation de l'intersection schématique
commute à la localisation étale, et qu'il en va de même pour la
réunion schématique.

\begin{lemma}
\label{spaces-morphisms-lemma-scheme-theoretic-intersection}
Soit $S$ un schéma. Soit $X$ un espace algébrique sur $S$.
Soient $Z, Y \subset X$ des sous-espaces fermés.
Soit $Z \cap Y$ l'intersection schématique de $Z$ et $Y$.
Alors $Z \cap Y \to Z$ et $Z \cap Y \to Y$ sont des immersions fermées,
et
$$
\xymatrix{
Z \cap Y \ar[r] \ar[d] & Z \ar[d] \\
Y \ar[r] & X
}
$$
est un diagramme cartésien d'espaces algébriques sur $S$, c'est-à-dire que
$Z \cap Y = Z \times_X Y$.
\end{lemma}

\begin{proof}
Les morphismes $Z \cap Y \to Z$ et $Z \cap Y \to Y$ sont des immersions fermées
d'après le lemme \ref{spaces-morphisms-lemma-closed-immersion-ideals}.
Comme la formation de l'intersection schématique commute
à la localisation étale, le cas des schémas montre que le diagramme est cartésien.
Voir le
lemme de Morphismes \ref{morphisms-lemma-scheme-theoretic-intersection}.
\end{proof}

\begin{lemma}
\label{spaces-morphisms-lemma-scheme-theoretic-union}
Soit $S$ un schéma. Soit $X$ un espace algébrique sur $S$.
Soient $Y, Z \subset X$ des sous-espaces fermés.
Soit $Y \cup Z$ la réunion schématique de $Y$ et $Z$.
Soit $Y \cap Z$ l'intersection schématique de $Y$ et $Z$.
Alors $Y \to Y \cup Z$ et $Z \to Y \cup Z$ sont des immersions fermées,
il existe une suite exacte courte
$$
0 \to \mathcal{O}_{Y \cup Z} \to \mathcal{O}_Y \times \mathcal{O}_Z
\to \mathcal{O}_{Y \cap Z} \to 0
$$
de $\mathcal{O}_X$-modules, et le diagramme
$$
\xymatrix{
Y \cap Z \ar[r] \ar[d] & Y \ar[d] \\
Z \ar[r] & Y \cup Z
}
$$
est cocartésien dans la catégorie des espaces algébriques sur $S$, c'est-à-dire que
$Y \cup Z = Y \amalg_{Y \cap Z} Z$.
\end{lemma}

\begin{proof}
Les morphismes $Y \to Y \cup Z$ et $Z \to Y \cup Z$ sont des immersions fermées
d'après le lemme \ref{spaces-morphisms-lemma-closed-immersion-ideals}. Dans la suite exacte courte,
nous employons l'équivalence du lemme \ref{spaces-morphisms-lemma-i-star-equivalence} pour considérer les
modules quasi-cohérents sur les sous-espaces fermés de $X$ comme des modules quasi-cohérents
sur $X$. Le premier morphisme de la suite provient des morphismes canoniques
$\mathcal{O}_{Y \cup Z} \to \mathcal{O}_Y$ et
$\mathcal{O}_{Y \cup Z} \to \mathcal{O}_Z$,
et le second du morphisme canonique
$\mathcal{O}_Y \to \mathcal{O}_{Y \cap Z}$ et de
l'opposé du morphisme canonique
$\mathcal{O}_Z \to \mathcal{O}_{Y \cap Z}$. Pour vérifier
l'exactitude, nous pouvons travailler localement pour la topologie étale et la déduire
du cas des schémas
(Morphismes, lemme \ref{morphisms-lemma-scheme-theoretic-union}).

\medskip\noindent
Pour montrer que le diagramme est cocartésien, soient un espace algébrique
$T$ sur $S$ et des morphismes $f : Y \to T$, $g : Z \to T$ qui coïncident sur
$Y \cap Z \to T$. Il s'agit de montrer qu'il existe un unique morphisme
$h : Y \cup Z \to T$ dont les restrictions sont $f$ et $g$.
Pour construire $h$, nous pouvons travailler localement sur $Y \cup Z$ pour la topologie étale
(car $Y \cup Z$, étant un espace algébrique, est un faisceau étale).
Nous pouvons donc supposer que $X$ est un schéma.
Dans ce cas, nous savons que $Y \cup Z$ est la somme amalgamée
de $Y$ et $Z$ le long de $Y \cap Z$ dans la catégorie des schémas
d'après Morphismes, lemme \ref{morphisms-lemma-scheme-theoretic-union}.
Choisissons un schéma $T'$ et un morphisme étale surjectif $T' \to T$.
Posons $Y' = T' \times_{T, f} Y$ et $Z' = T' \times_{T, g} Z$.
Alors $Y'$ et $Z'$ sont des schémas, et nous disposons d'un isomorphisme canonique
$\varphi : Y' \times_Y (Y \cap Z) \to Z' \times_Z (Y \cap Z)$
de schémas. D'après Compléments sur les morphismes, lemme
\ref{more-morphisms-lemma-pushout-along-closed-immersions}
la somme amalgamée $W' = Y' \amalg_{Y' \times_Y (Y \cap Z), \varphi} Z'$
existe dans la catégorie des schémas.
Le morphisme $W' \to Y \cup Z$ est étale d'après
Compléments sur les morphismes, lemme
\ref{more-morphisms-lemma-pushout-along-closed-immersions-properties-above}.
Il est surjectif puisque $Y' \to Y$ et $Z' \to Z$ le sont.
Les morphismes $f' : Y' \to T'$ et $g' : Z' \to T'$
se recollent en un unique morphisme de schémas $h' : W' \to T'$.
Par unicité, le composé $W' \to T' \to T$
descend au morphisme souhaité $h : Y \cup Z \to T$.
Nous omettons certains détails.
\end{proof}






\section{Supports de modules}
\label{spaces-morphisms-section-support}

\noindent
Dans cette section, nous rassemblons quelques résultats élémentaires sur les
supports des modules quasi-cohérents sur les espaces algébriques. Soit $X$ un
espace algébrique. Le support d'un faisceau abélien sur $X_\etale$
a été défini dans Propriétés des espaces, section
\ref{spaces-properties-section-support}.
Nous adoptons la même définition pour les supports de modules.
Le lemme suivant affirme qu'elle coïncide avec la notion
définie pour les modules quasi-cohérents sur les schémas.

\begin{lemma}
\label{spaces-morphisms-lemma-support-covering}
Soit $S$ un schéma. Soit $X$ un espace algébrique sur $S$.
Soit $\mathcal{F}$ un $\mathcal{O}_X$-module quasi-cohérent.
Soit $U$ un schéma et soit $\varphi : U \to X$ un morphisme étale.
Alors
$$
\text{Supp}(\varphi^*\mathcal{F}) = |\varphi|^{-1}(\text{Supp}(\mathcal{F}))
$$
où le membre de gauche est le support de $\varphi^*\mathcal{F}$ en tant que
module quasi-cohérent sur le schéma $U$.
\end{lemma}

\begin{proof}
Soit $u\in U$ un point (usuel) et soit $\overline{x}$ un
point géométrique au-dessus de $u$. D'après
Propriétés des espaces, lemme \ref{spaces-properties-lemma-stalk-quasi-coherent},
nous avons
$(\varphi^*\mathcal{F})_u \otimes_{\mathcal{O}_{U, u}}
\mathcal{O}_{X, \overline{x}} = \mathcal{F}_{\overline{x}}$.
Comme $\mathcal{O}_{U, u} \to \mathcal{O}_{X, \overline{x}}$
est l'hensélisé strict d'après
Propriétés des espaces, lemme
\ref{spaces-properties-lemma-describe-etale-local-ring}
il est fidèlement plat (voir
Compléments d'algèbre, lemme
\ref{more-algebra-lemma-dumb-properties-henselization}).
Par conséquent, $(\varphi^*\mathcal{F})_u = 0$ si et seulement si
$\mathcal{F}_{\overline{x}} = 0$. Cela démontre le lemme.
\end{proof}

\noindent
Pour les modules quasi-cohérents de type fini, le support est fermé,
se vérifie sur les fibres et commute au changement de base.

\begin{lemma}
\label{spaces-morphisms-lemma-support-finite-type}
Soit $S$ un schéma. Soit $X$ un espace algébrique sur $S$.
Soit $\mathcal{F}$ un $\mathcal{O}_X$-module quasi-cohérent de type fini.
Alors
\begin{enumerate}
\item Le support de $\mathcal{F}$ est fermé.
\item Pour tout point géométrique $\overline{x}$ au-dessus de $x \in |X|$, nous avons
$$
x \in \text{Supp}(\mathcal{F})
\Leftrightarrow
\mathcal{F}_{\overline{x}} \not = 0
\Leftrightarrow
\mathcal{F}_{\overline{x}} \otimes_{\mathcal{O}_{X, \overline{x}}}
\kappa(\overline{x}) \not = 0.
$$
\item Pour tout morphisme d'espaces algébriques $f : Y \to X$, l'image inverse
$f^*\mathcal{F}$ est également de type fini et nous avons
$\text{Supp}(f^*\mathcal{F}) = f^{-1}(\text{Supp}(\mathcal{F}))$.
\end{enumerate}
\end{lemma}

\begin{proof}
Choisissons un schéma $U$ et un morphisme étale surjectif $\varphi : U \to X$.
D'après le lemme \ref{spaces-morphisms-lemma-support-covering}, l'image inverse du support de
$\mathcal{F}$ est le support de $\varphi^*\mathcal{F}$, qui est fermé d'après
Morphismes, lemme \ref{morphisms-lemma-support-finite-type}.
Ainsi, (1) découle de la définition de la topologie sur $|X|$.

\medskip\noindent
La première équivalence de (2) est la définition du support.
La seconde équivalence découle du lemme de Nakayama ; voir
Algèbre, lemme \ref{algebra-lemma-NAK}.

\medskip\noindent
Soit $f : Y \to X$ comme en (3). Notons que $f^*\mathcal{F}$ est de type fini
d'après Propriétés des espaces, section
\ref{spaces-properties-section-properties-modules}.
Pour la dernière assertion, soit $\overline{y}$ un point géométrique de $Y$
d'image le point géométrique $\overline{x}$ dans $X$. Rappelons que
$$
(f^*\mathcal{F})_{\overline{y}} =
\mathcal{F}_{\overline{x}} \otimes_{\mathcal{O}_{X, \overline{x}}}
\mathcal{O}_{Y, \overline{y}},
$$
voir Propriétés des espaces, lemme
\ref{spaces-properties-lemma-stalk-pullback-quasi-coherent}.
Par conséquent, $(f^*\mathcal{F})_{\overline{y}} \otimes \kappa(\overline{y})$
est non nul si et seulement si
$\mathcal{F}_{\overline{x}} \otimes \kappa(\overline{x})$ est non nul.
D'après (2), il s'ensuit que $x \in \text{Supp}(\mathcal{F})$ si et seulement
si $y \in \text{Supp}(f^*\mathcal{F})$, ce qui constitue
l'assertion (3).
\end{proof}

\noindent
Notre objectif suivant est de montrer que le support schématique
d'un module quasi-cohérent de type fini (voir
Morphismes, définition \ref{morphisms-definition-scheme-theoretic-support})
a également un sens pour les modules quasi-cohérents de type fini sur les
espaces algébriques.

\begin{lemma}
\label{spaces-morphisms-lemma-scheme-theoretic-support}
Soit $S$ un schéma. Soit $X$ un espace algébrique sur $S$.
Soit $\mathcal{F}$ un $\mathcal{O}_X$-module quasi-cohérent de type fini.
Il existe un plus petit sous-espace fermé $i : Z \to X$ tel qu'il
existe un $\mathcal{O}_Z$-module quasi-cohérent $\mathcal{G}$ vérifiant
$i_*\mathcal{G} \cong \mathcal{F}$. De plus :
\begin{enumerate}
\item Si $U$ est un schéma et si $\varphi : U \to X$ est un morphisme étale,
alors $Z \times_X U$ est le support schématique de $\varphi^*\mathcal{F}$.
\item Le faisceau quasi-cohérent $\mathcal{G}$ est unique à isomorphisme
unique près.
\item Le faisceau quasi-cohérent $\mathcal{G}$ est de type fini.
\item Les supports de $\mathcal{G}$ et de $\mathcal{F}$ sont tous deux $|Z|$.
\end{enumerate}
\end{lemma}

\begin{proof}
Choisissons un schéma $U$ et un morphisme étale surjectif $\varphi : U \to X$.
Posons $R = U \times_X U$, de projections $s, t : R \to U$.
Soit $i' : Z' \to U$ le support schématique de $\varphi^*\mathcal{F}$,
et soit $\mathcal{G}'$ le $\mathcal{O}_{Z'}$-module quasi-cohérent de type fini
(unique à isomorphisme unique près)
tel que $i'_*\mathcal{G}' = \varphi^*\mathcal{F}$ ; voir
Morphismes, lemme \ref{morphisms-lemma-scheme-theoretic-support}.
Comme $s^*\varphi^*\mathcal{F} = t^*\varphi^*\mathcal{F}$, nous avons
$R' = s^{-1}Z' = t^{-1}Z'$ comme sous-schémas fermés de $R$, d'après
Morphismes, lemme \ref{morphisms-lemma-flat-pullback-support}.
Nous pouvons donc appliquer Propriétés des espaces, lemme
\ref{spaces-properties-lemma-subspaces-presentation}
pour obtenir un sous-espace fermé $i : Z \to X$ dont l'image inverse sur $U$ est $Z'$.
Notons $s', t' : R' \to Z'$ les projections et
$j' : R' \to R$ l'immersion fermée considérée ; nous avons
$$
j'_* (s')^*\mathcal{G}' = s^* i'_*\mathcal{G}' =
s^*\varphi^*\mathcal{F} = t^*\varphi^*\mathcal{F} =
t^*i'_*\mathcal{G}' = j'_*(t')^*\mathcal{G}'
$$
(la première et la dernière égalité résultant de Cohomologie des schémas,
lemme \ref{coherent-lemma-flat-base-change-cohomology}).
Le lemme
Morphismes, lemme \ref{morphisms-lemma-flat-pullback-support},
appliqué à $R' \to R$, fournit donc, par unicité, un isomorphisme
$\alpha : (t')^*\mathcal{G}' \to (s')^*\mathcal{G}'$
compatible avec l'isomorphisme canonique
$t^*\varphi^*\mathcal{F} = s^*\varphi^*\mathcal{F}$
par l'intermédiaire de $j'_*$. Il est clair que $\alpha$ satisfait la condition de cocycle ;
nous pouvons donc appliquer
Propriétés des espaces, Proposition
\ref{spaces-properties-proposition-quasi-coherent}
pour obtenir un module quasi-cohérent $\mathcal{G}$ sur $Z$ dont la restriction
à $Z'$ s'identifie à $\mathcal{G}'$, compatiblement à $\alpha$.
En utilisant à nouveau l'équivalence de la proposition précitée
(cette fois pour $X$), nous concluons que $i_*\mathcal{G} \cong \mathcal{F}$.

\medskip\noindent
Cela démontre l'existence. Les autres propriétés du lemme découlent
de la comparaison avec le résultat pour les schémas, au moyen du
lemme \ref{spaces-morphisms-lemma-support-covering}.
Nous omettons les démonstrations détaillées.
\end{proof}

\begin{definition}
\label{spaces-morphisms-definition-scheme-theoretic-support}
Soit $S$ un schéma. Soit $X$ un espace algébrique sur $S$.
Soit $\mathcal{F}$ un $\mathcal{O}_X$-module quasi-cohérent de type fini.
Le {\it support schématique de $\mathcal{F}$} est le sous-espace fermé
$Z \subset X$ construit dans le lemme \ref{spaces-morphisms-lemma-scheme-theoretic-support}.
\end{definition}

\noindent
Dans cette situation, nous considérons souvent $\mathcal{F}$ comme un faisceau
quasi-cohérent de type fini sur $Z$ (au moyen de l'équivalence de catégories du
lemme \ref{spaces-morphisms-lemma-i-star-equivalence}).






\section{Image schématique}
\label{spaces-morphisms-section-scheme-theoretic-image}

\noindent
Attention : une partie des résultats de cette section est d'une généralité extrême
et se comporte autrement qu'on pourrait s'y attendre.

\begin{lemma}
\label{spaces-morphisms-lemma-scheme-theoretic-image}
\begin{slogan}
L'image schématique d'un morphisme d'espaces algébriques existe.
\end{slogan}
Soit $S$ un schéma. Soit $f : X \to Y$ un morphisme d'espaces algébriques
sur $S$. Il existe un sous-espace fermé $Z \subset Y$ tel que $f$ se factorise
par $Z$ et tel que, pour tout autre sous-espace fermé $Z' \subset Y$ tel que
$f$ se factorise par $Z'$, nous ayons $Z \subset Z'$.
\end{lemma}

\begin{proof}
Posons $\mathcal{I} = \Ker(\mathcal{O}_Y \to f_*\mathcal{O}_X)$.
Si $\mathcal{I}$ est quasi-cohérent, il suffit de prendre pour $Z$ le
sous-espace fermé défini par $\mathcal{I}$ ; voir
le lemme \ref{spaces-morphisms-lemma-closed-immersion-ideals}.
En général, il faut, pour démontrer le lemme, établir qu'il existe
un plus grand faisceau quasi-cohérent d'idéaux $\mathcal{I}'$ contenu dans
$\mathcal{I}$.
Cela découle du lemme \ref{spaces-morphisms-lemma-largest-quasi-coherent-subsheaf}.
\end{proof}

\noindent
Supposons que, dans la situation du lemme \ref{spaces-morphisms-lemma-scheme-theoretic-image}
ci-dessus, $X$ et $Y$ soient représentables. Alors le sous-espace fermé $Z \subset Y$
obtenu dans le lemme coïncide avec le sous-schéma fermé $Z \subset Y$ obtenu dans
Morphismes, lemme \ref{morphisms-lemma-scheme-theoretic-image}.
En effet, les sous-espaces fermés (ou sous-schémas fermés) sont en correspondance,
en renversant les inclusions, avec les faisceaux quasi-cohérents d'idéaux sur
$Y_\etale$ et $Y$. Comme les catégories des modules quasi-cohérents
sur $Y_\etale$ et $Y$ sont les mêmes
(Propriétés des espaces, section \ref{spaces-properties-section-quasi-coherent}),
nous pouvons conclure. Ainsi, la définition suivante coïncide avec la définition
antérieure pour les morphismes de schémas.

\begin{definition}
\label{spaces-morphisms-definition-scheme-theoretic-image}
Soit $S$ un schéma. Soit $f : X \to Y$ un morphisme d'espaces algébriques
sur $S$. L'{\it image schématique} de $f$ est le plus petit sous-espace
fermé $Z \subset Y$ par lequel $f$
se factorise ; voir le lemme \ref{spaces-morphisms-lemma-scheme-theoretic-image} ci-dessus.
\end{definition}

\noindent
Nous désignons souvent simplement par $f : X \to Z$ la factorisation de $f$.
Si le morphisme $f$ n'est pas quasi-compact, alors, en général, la
construction de l'image schématique ne commute pas à la
restriction aux sous-espaces ouverts de $Y$.

\begin{lemma}
\label{spaces-morphisms-lemma-quasi-compact-scheme-theoretic-image}
Soit $S$ un schéma.
Soit $f : X \to Y$ un morphisme d'espaces algébriques sur $S$.
Soit $Z \subset Y$ l'image schématique de $f$.
Si $f$ est quasi-compact, alors
\begin{enumerate}
\item le faisceau d'idéaux
$\mathcal{I} = \Ker(\mathcal{O}_Y \to f_*\mathcal{O}_X)$
est quasi-cohérent,
\item l'image schématique $Z$ est le sous-espace fermé
correspondant à $\mathcal{I}$,
\item pour tout morphisme étale $V \to Y$, l'image schématique de
$X \times_Y V \to V$ est égale à $Z \times_Y V$, et
\item l'image $|f|(|X|) \subset |Z|$ est une partie dense de $|Z|$.
\end{enumerate}
\end{lemma}

\begin{proof}
Pour démontrer (3), il suffit de démontrer (1) et (2), puisque la
formation de $\mathcal{I}$ commute à la localisation étale.
Si (1) est vraie, nous avons établi (2) dans la preuve du lemme
\ref{spaces-morphisms-lemma-scheme-theoretic-image}. Montrons que $\mathcal{I}$ est quasi-cohérent.
Comme la propriété d'être quasi-cohérent est locale pour la topologie étale, nous pouvons
supposer que $Y$ est un schéma affine. Comme $f$ est quasi-compact,
nous pouvons trouver un schéma affine $U$ et un morphisme étale surjectif
$U \to X$. Notons $f'$ le composé $U \to X \to Y$.
Alors $f_*\mathcal{O}_X$ est un sous-faisceau de $f'_*\mathcal{O}_U$,
et donc $\mathcal{I} = \Ker(\mathcal{O}_Y \to f'_*\mathcal{O}_U)$.
D'après le lemme \ref{spaces-morphisms-lemma-pushforward},
le faisceau $f'_*\mathcal{O}_U$ est quasi-cohérent sur $Y$. Ainsi, $\mathcal{I}$
est quasi-cohérent comme noyau d'un morphisme de modules quasi-cohérents.
Enfin, la partie (4) découle des parties (1), (2) et (3), car l'idéal
$\mathcal{I}$ sera l'idéal unité en tout point de $|Y|$ qui
n'appartient pas à l'adhérence de $|f|(|X|)$.
\end{proof}

\begin{lemma}
\label{spaces-morphisms-lemma-scheme-theoretic-image-reduced}
Soit $S$ un schéma. Soit $f : X \to Y$ un morphisme d'espaces algébriques
sur $S$. Supposons que $X$ soit réduit. Alors
\begin{enumerate}
\item l'image schématique $Z$ de $f$ est l'adhérence $\overline{|f|(|X|)}$, munie de la
structure réduite induite, et
\item pour tout morphisme étale $V \to Y$, l'image schématique de
$X \times_Y V \to V$ est égale à $Z \times_Y V$.
\end{enumerate}
\end{lemma}

\begin{proof}
La partie (1) est vraie parce que l'adhérence $\overline{|f|(|X|)}$, munie de la
structure réduite induite, est le plus petit sous-espace fermé
de $Y$ par lequel $f$ se factorise ; voir
Propriétés des espaces, lemme \ref{spaces-properties-lemma-map-into-reduction}.
La partie (2) découle de (1), du fait que l'application $|V| \to |Y|$ est ouverte et du
fait que la réduction est préservée par localisation étale.
\end{proof}

\begin{lemma}
\label{spaces-morphisms-lemma-reach-points-scheme-theoretic-image}
Soit $S$ un schéma.
Soit $f : X \to Y$ un morphisme quasi-compact d'espaces algébriques sur $S$.
Soit $Z$ l'image schématique de $f$.
Soit $z \in |Z|$. Il existe un anneau de valuation $A$ de
corps des fractions $K$ et un diagramme commutatif
$$
\xymatrix{
\Spec(K) \ar[rr] \ar[d] & & X \ar[d] \ar[ld] \\
\Spec(A) \ar[r] & Z \ar[r] & Y
}
$$
tel que le point fermé de $\Spec(A)$ s'envoie sur $z$.
\end{lemma}

\begin{proof}
Choisissons un schéma affine $V$, un point $z' \in V$
et un morphisme étale $V \to Y$ envoyant $z'$ sur $z$.
Soit $Z' \subset V$ l'image schématique de $X \times_Y V \to V$.
D'après le lemme \ref{spaces-morphisms-lemma-quasi-compact-scheme-theoretic-image}, nous avons
$Z' = Z \times_Y V$. Ainsi, $z' \in Z'$.
Comme $f$ est quasi-compact et $V$ affine, nous voyons que
$X \times_Y V$ est quasi-compact. Il existe donc
un schéma affine $W$ et un morphisme étale surjectif
$W \to X \times_Y V$. Alors $Z' \subset V$ est également
l'image schématique de $W \to V$.
D'après Morphismes, lemme \ref{morphisms-lemma-reach-points-scheme-theoretic-image},
nous pouvons choisir un diagramme
$$
\xymatrix{
\Spec(K) \ar[r] \ar[d] &
W \ar[r] \ar[d] &
X \times_Y V \ar[d] \ar[r] &
X \ar[d] \\
\Spec(A) \ar[r] &
Z' \ar[r] &
V \ar[r] &
Y
}
$$
tel que le point fermé de $\Spec(A)$ s'envoie sur $z'$.
En composant avec $Z' \to Z$ et $W \to X \times_Y V \to X$,
nous obtenons une solution.
\end{proof}

\begin{lemma}
\label{spaces-morphisms-lemma-factor-factor}
Soit $S$ un schéma. Soit
$$
\xymatrix{
X_1 \ar[d] \ar[r]_{f_1} & Y_1 \ar[d] \\
X_2 \ar[r]^{f_2} & Y_2
}
$$
un diagramme commutatif d'espaces algébriques sur $S$.
Soient $Z_i \subset Y_i$, $i = 1, 2$, les
images schématiques de $f_i$. Alors le morphisme
$Y_1 \to Y_2$ induit un morphisme $Z_1 \to Z_2$ et un
diagramme commutatif
$$
\xymatrix{
X_1 \ar[r] \ar[d] & Z_1 \ar[d] \ar[r] & Y_1 \ar[d] \\
X_2 \ar[r] & Z_2 \ar[r] & Y_2
}
$$
\end{lemma}

\begin{proof}
L'image inverse schématique de $Z_2$ dans $Y_1$
est un sous-espace fermé de $Y_1$ par lequel
$f_1$ se factorise. Par conséquent, $Z_1$ y est contenu.
Cela démontre le lemme.
\end{proof}

\begin{lemma}
\label{spaces-morphisms-lemma-scheme-theoretic-image-of-partial-section}
Soit $S$ un schéma.
Soit $f : X \to Y$ un morphisme séparé d'espaces algébriques sur $S$.
Soit $V \subset Y$ un sous-espace ouvert tel que $V \to Y$ soit quasi-compact.
Soit $s : V \to X$ un morphisme tel que $f \circ s = \text{id}_V$.
Soit $Y'$ l'image schématique de $s$.
Alors $Y' \to Y$ est un isomorphisme au-dessus de $V$.
\end{lemma}

\begin{proof}
D'après le lemme \ref{spaces-morphisms-lemma-quasi-compact-permanence},
le morphisme $s : V \to X$ est quasi-compact.
La construction de l'image schématique $Y'$
de $s$ commute donc à la restriction aux ouverts, d'après le
lemme \ref{spaces-morphisms-lemma-quasi-compact-scheme-theoretic-image}.
En particulier, nous voyons que $Y' \cap f^{-1}(V)$ est
l'image schématique d'une section du morphisme séparé
$f^{-1}(V) \to V$. Comme une section d'un morphisme séparé
est une immersion fermée
(lemme \ref{spaces-morphisms-lemma-section-immersion}),
nous concluons que
$Y' \cap f^{-1}(V) \to V$ est un isomorphisme, comme souhaité.
\end{proof}








\section{Adhérence schématique et densité}
\label{spaces-morphisms-section-scheme-theoretic-closure}

\noindent
Cette section est l'analogue de
Morphismes, section \ref{morphisms-section-scheme-theoretic-closure}.

\begin{lemma}
\label{spaces-morphisms-lemma-scheme-theoretically-dense-representable}
Soit $S$ un schéma. Soit $W \subset S$ un sous-schéma ouvert
schématiquement dense
(Morphismes, définition \ref{morphisms-definition-scheme-theoretically-dense}).
Soit $f : X \to S$ un morphisme de schémas plat, localement de
présentation finie et localement quasi-fini.
Alors $f^{-1}(W)$ est schématiquement dense dans $X$.
\end{lemma}

\begin{proof}
Nous utiliserons la caractérisation de Morphismes, lemme
\ref{morphisms-lemma-characterize-scheme-theoretically-dense}.
Supposons que $V \subset X$ soit un ouvert et que $g \in \Gamma(V, \mathcal{O}_V)$
soit une fonction dont la restriction à $f^{-1}(W) \cap V$ est nulle.
Il faut montrer que $g = 0$. Supposons $g \not = 0$ afin d'obtenir une
contradiction. D'après
Compléments sur les morphismes, lemme \ref{more-morphisms-lemma-go-down-with-annihilators}
nous pouvons, quitte à rétrécir $V$, trouver un ouvert $U \subset S$ s'insérant dans un
diagramme commutatif
$$
\xymatrix{
V \ar[r] \ar[d]_\pi & X \ar[d]^f \\
U \ar[r] & S,
}
$$
un sous-faisceau quasi-cohérent $\mathcal{F} \subset \mathcal{O}_U$, un entier
$r > 0$ et un morphisme injectif de $\mathcal{O}_U$-modules
$\mathcal{F}^{\oplus r} \to \pi_*\mathcal{O}_V$
dont l'image contient $g|_V$. Soit
$(g_1, \ldots, g_r) \in \Gamma(U, \mathcal{F}^{\oplus r})$ un antécédent de $g$.
Alors $g_i|_{W \cap U} = 0$, car $g|_{f^{-1}W \cap V} = 0$.
Par conséquent, $g_i = 0$, puisque $\mathcal{F} \subset \mathcal{O}_U$ et que
$W$ est schématiquement dense dans $S$.
Il s'ensuit que $g = 0$, ce qui est la contradiction cherchée.
\end{proof}

\begin{lemma}
\label{spaces-morphisms-lemma-scheme-theoretically-dense}
Soit $S$ un schéma.
Soit $X$ un espace algébrique sur $S$.
Soit $U \subset X$ un sous-espace ouvert.
Les conditions suivantes sont équivalentes :
\begin{enumerate}
\item pour tout morphisme étale $\varphi : V \to X$ (d'espaces algébriques),
l'adhérence schématique de $\varphi^{-1}(U)$ dans $V$ est égale à $V$,
\item il existe un schéma $V$ et un morphisme étale surjectif
$\varphi : V \to X$ tels que l'adhérence schématique de
$\varphi^{-1}(U)$ dans $V$ soit égale à $V$,
\end{enumerate}
\end{lemma}

\begin{proof}
Observons que, si $V \to V'$ est un morphisme entre espaces algébriques étales
sur $X$, et si $Z \subset V$, resp.\ $Z' \subset V'$, sont les adhérences schématiques
de $U \times_X V$, resp.\ $U \times_X V'$, dans $V$, resp.\ $V'$,
alors $Z$ s'envoie dans $Z'$. Ainsi, si $V \to V'$ est étale et surjectif,
$Z = V$ entraîne $Z' = V'$. Ensuite, notons qu'un morphisme étale est
plat, localement de présentation finie et localement quasi-fini
(voir Morphismes, section \ref{morphisms-section-etale}).
Le lemme \ref{spaces-morphisms-lemma-scheme-theoretically-dense-representable}
implique donc que, si $V$ et $V'$ sont des schémas, $Z' = V'$ entraîne
$Z = V$. Un argument formel utilisant le fait que tout espace algébrique admet un
recouvrement étale par un schéma montre que (1) et (2) sont équivalentes.
\end{proof}

\noindent
Il résulte du
lemme \ref{spaces-morphisms-lemma-scheme-theoretically-dense}
que la définition suivante est compatible avec celle
du cas des schémas.

\begin{definition}
\label{spaces-morphisms-definition-scheme-theoretically-dense}
Soit $S$ un schéma.
Soit $X$ un espace algébrique sur $S$.
Soit $U \subset X$ un sous-espace ouvert.
\begin{enumerate}
\item L'image schématique du morphisme $U \to X$
est appelée l'{\it adhérence schématique de $U$ dans $X$}.
\item Nous disons que $U$ est {\it schématiquement dense dans $X$}
si les conditions équivalentes du
lemme \ref{spaces-morphisms-lemma-scheme-theoretically-dense} sont satisfaites.
\end{enumerate}
\end{definition}

\noindent
Avec cette définition, il n'est {\bf pas} vrai que $U$ soit schématiquement
dense dans $X$ si et seulement si l'adhérence schématique
de $U$ est $X$. C'est quelque peu inélégant. Nous verrons cependant que,
sous des conditions de finitude convenables, cette équivalence est vraie.

\begin{lemma}
\label{spaces-morphisms-lemma-scheme-theoretically-dense-quasi-compact}
Soit $S$ un schéma. Soit $X$ un espace algébrique sur $S$.
Soit $U \subset X$ un sous-espace ouvert.
Si $U \to X$ est quasi-compact, alors $U$
est schématiquement dense dans $X$ si et seulement si l'adhérence schématique
de $U$ dans $X$ est $X$.
\end{lemma}

\begin{proof}
Cela découle de la partie (3) du lemme \ref{spaces-morphisms-lemma-quasi-compact-scheme-theoretic-image}.
\end{proof}

\begin{lemma}
\label{spaces-morphisms-lemma-characterize-scheme-theoretically-dense}
Soit $S$ un schéma.
Soit $j : U \to X$ une immersion ouverte d'espaces algébriques sur $S$.
Alors $U$ est schématiquement dense dans $X$ si et seulement si
$\mathcal{O}_X \to j_*\mathcal{O}_U$ est injectif.
\end{lemma}

\begin{proof}
Si $\mathcal{O}_X \to j_*\mathcal{O}_U$ est injectif,
il en va de même après restriction à tout
espace algébrique $V$ étale sur $X$.
L'adhérence schématique de $U \times_X V$ dans $V$
est donc égale à $V$ ; voir la preuve du
lemme \ref{spaces-morphisms-lemma-scheme-theoretic-image}.
Réciproquement, supposons que l'adhérence schématique
de $U \times_X V$ soit égale à $V$ pour tout $V$ étale sur $X$.
Supposons que $\mathcal{O}_X \to j_*\mathcal{O}_U$ ne soit pas injectif.
Nous pouvons alors trouver un schéma affine, disons $V = \Spec(A)$, étale sur $X$,
et un élément non nul $f \in A$ tel que $f$ s'envoie sur zéro dans
$\Gamma(V \times_X U, \mathcal{O})$. Dans ce cas, l'adhérence schématique
de $V \times_X U$ dans $V$ est clairement contenue dans $\Spec(A/(f))$,
ce qui est absurde.
\end{proof}

\begin{lemma}
\label{spaces-morphisms-lemma-intersection-scheme-theoretically-dense}
Soit $S$ un schéma. Soit $X$ un espace algébrique sur $S$.
Si $U$ et $V$ sont des sous-espaces ouverts
schématiquement denses dans $X$, alors $U \cap V$ l'est aussi.
\end{lemma}

\begin{proof}
Soit $W \to X$ un morphisme étale quelconque. Considérons l'application
$\mathcal{O}(W) \to \mathcal{O}(W \times_X V)
\to \mathcal{O}(W \times_X (V \cap U))$.
D'après le lemme \ref{spaces-morphisms-lemma-characterize-scheme-theoretically-dense},
les deux applications sont injectives. Leur composée est donc injective.
Par conséquent, d'après le lemme \ref{spaces-morphisms-lemma-characterize-scheme-theoretically-dense},
$U \cap V$ est schématiquement dense dans $X$.
\end{proof}

\begin{lemma}
\label{spaces-morphisms-lemma-quasi-compact-immersion}
Soit $S$ un schéma. Soit $h : Z \to X$ une immersion d'espaces algébriques
sur $S$. Supposons que $Z \to X$ soit quasi-compact ou que $Z$ soit réduit.
Soit $\overline{Z} \subset X$ l'image schématique de $h$.
Alors le morphisme $Z \to \overline{Z}$ est une immersion ouverte
qui identifie $Z$ à un sous-espace ouvert schématiquement dense
de $\overline{Z}$. De plus, $Z$ est topologiquement
dense dans $\overline{Z}$.
\end{lemma}

\begin{proof}
Dans les deux cas, la formation de $\overline{Z}$ commute à la
localisation étale ; voir les lemmes
\ref{spaces-morphisms-lemma-quasi-compact-scheme-theoretic-image} et
\ref{spaces-morphisms-lemma-scheme-theoretic-image-reduced}.
Le présent lemme découle donc du cas des schémas ; voir
Morphismes, lemme \ref{morphisms-lemma-quasi-compact-immersion}.
\end{proof}

\begin{lemma}
\label{spaces-morphisms-lemma-equality-of-morphisms}
Soit $S$ un schéma. Soit $B$ un espace algébrique sur $S$.
Soient $f, g : X \to Y$ des morphismes d'espaces algébriques sur $B$.
Soit $U \subset X$ un sous-espace ouvert tel que
$f|_U = g|_U$. Si l'adhérence schématique de $U$
dans $X$ est $X$ et si $Y \to B$ est séparé, alors $f = g$.
\end{lemma}

\begin{proof}
Comme $Y \to B$ est séparé, le produit fibré
$Y \times_{\Delta, Y \times_B Y, (f, g)} X$ est un sous-espace fermé $Z \subset X$. 
Comme $f|_U = g|_U$, nous avons $U \subset Z$. Donc $Z = X$,
car l'hypothèse sur $U$ affirme que son adhérence schématique est $X$.
\end{proof}






\section{Morphismes dominants}
\label{spaces-morphisms-section-dominant}

\noindent
Nous reprenons la définition d'un morphisme dominant de schémas afin d'obtenir
la notion de morphisme dominant d'espaces algébriques. Nous avertissons le
lecteur que cette définition ne se comporte pas bien à moins que le morphisme
ne soit quasi-compact et que les espaces algébriques ne satisfassent certains
axiomes de séparation.

\begin{definition}
\label{spaces-morphisms-definition-dominant}
Soit $S$ un schéma. Un morphisme $f : X \to Y$ d'espaces algébriques sur $S$ est
dit {\it dominant} si l'image de $|f| : |X| \to |Y|$ est dense dans $|Y|$.
\end{definition}






\section{Morphismes universellement injectifs}
\label{spaces-morphisms-section-universally-injective}

\noindent
Nous avons déjà défini à la section \ref{spaces-morphisms-section-representable}
ce que signifie, pour un morphisme représentable d'espaces algébriques,
être universellement injectif. Pour un corps $K$ sur $S$ (rappelons que cela
signifie que l'on s'est donné un morphisme structural $\Spec(K) \to S$) et un
espace algébrique $X$ sur $S$, nous posons
$X(K) = \Mor_S(\Spec(K), X)$. Nous traduisons d'abord la
condition portant sur les morphismes représentables en une condition sur le foncteur
des points.

\begin{lemma}
\label{spaces-morphisms-lemma-universally-injective-representable}
Soit $S$ un schéma. Soit $f : X \to Y$ un morphisme représentable
d'espaces algébriques sur $S$. Alors $f$ est universellement injectif
(au sens de la section \ref{spaces-morphisms-section-representable})
si et seulement si, pour tout corps $K$, l'application $X(K) \to Y(K)$ est injective.
\end{lemma}

\begin{proof}
Nous utiliserons
Morphismes, lemme \ref{morphisms-lemma-universally-injective}
sans autre mention.
Supposons que $f$ soit universellement injectif. Alors, pour tout corps $K$ et tout
morphisme $\Spec(K) \to Y$, le morphisme de schémas
$\Spec(K) \times_Y X \to \Spec(K)$ est universellement injectif.
Il existe donc au plus une section du morphisme
$\Spec(K) \times_Y X \to \Spec(K)$. Ainsi, l'application
$X(K) \to Y(K)$ est injective. Réciproquement, supposons que, pour tout corps $K$,
l'application $X(K) \to Y(K)$ soit injective. Soit $T \to Y$ un morphisme d'un
schéma vers $Y$, et considérons le changement de base $f_T : T \times_Y X \to T$.
Pour tout corps $K$, on a
$$
(T \times_Y X)(K) = T(K) \times_{Y(K)} X(K)
$$
par définition du produit fibré ; ainsi, l'injectivité de
$X(K) \to Y(K)$ entraîne celle de
$(T \times_Y X)(K) \to T(K)$, ce qui signifie que $f_T$ est universellement injectif,
comme souhaité.
\end{proof}

\noindent
Traduisons maintenant l'injectivité de l'application sur les points à valeurs
dans des corps en une condition plus géométrique.

\begin{lemma}
\label{spaces-morphisms-lemma-universally-injective}
Soit $S$ un schéma.
Soit $f : X \to Y$ un morphisme d'espaces algébriques sur $S$.
Les conditions suivantes sont équivalentes :
\begin{enumerate}
\item l'application $X(K) \to Y(K)$ est injective pour tout corps $K$ sur $S$ ;
\item pour tout morphisme $Y' \to Y$ d'espaces algébriques sur $S$,
l'application induite $|Y' \times_Y X| \to |Y'|$ est injective ; et
\item le morphisme diagonal $X \to X \times_Y X$ est surjectif.
\end{enumerate}
\end{lemma}

\begin{proof}
Supposons (1). Soit $g : Y' \to Y$ un morphisme d'espaces
algébriques, et notons $f' : Y' \times_Y X \to Y'$ le changement de base de $f$.
Soient $K_i$, $i = 1, 2$, des corps et soient
$\varphi_i : \Spec(K_i) \to Y' \times_Y X$ des morphismes
tels que $f' \circ \varphi_1$ et $f' \circ \varphi_2$ définissent le
même élément de $|Y'|$. Par définition, cela signifie qu'il existe un
corps $\Omega$ et des plongements $\alpha_i : K_i \subset \Omega$ tels que
les deux morphismes
$f' \circ \varphi_i \circ \alpha_i : \Spec(\Omega) \to Y'$ soient égaux.
Voici le diagramme commutatif correspondant :
$$
\xymatrix{
\Spec(\Omega) \ar@/_5ex/[ddrr] \ar[rd]^{\alpha_1} \ar[r]_{\alpha_2} &
\Spec(K_2) \ar[rd]^{\varphi_2} \\
& \Spec(K_1) \ar[r]^{\varphi_1} &
Y' \times_Y X  \ar[d]^{f'} \ar[r]^{g'} &
 X \ar[d]^f \\
& & Y' \ar[r]^g & Y.
}
$$
En particulier, les composés $g \circ f' \circ \varphi_i \circ \alpha_i$
sont égaux. D'après (1), cela implique que les morphismes
$g' \circ \varphi_i \circ \alpha_i$ sont égaux, où $g' : Y' \times_Y X \to X$
est la projection. Par la propriété universelle du produit fibré, nous en déduisons
que les morphismes
$\varphi_i \circ \alpha_i : \Spec(\Omega) \to Y' \times_Y X$ sont
égaux. Autrement dit, $\varphi_1$ et $\varphi_2$ définissent le même point
de $Y' \times_Y X$. Nous en concluons que (2) est vérifiée.

\medskip\noindent
Supposons (2). Soit $K$ un corps sur $S$, et soient $a, b : \Spec(K) \to X$
deux morphismes tels que $f \circ a = f \circ b$. Notons
$c : \Spec(K) \to Y$ leur valeur commune. Par hypothèse,
$|\Spec(K) \times_{c, Y} X| \to |\Spec(K)|$ est injective.
Cela signifie qu'il existe un corps $\Omega$ et des plongements
$\alpha_i : K \to \Omega$ tels que
$$
\xymatrix{
\Spec(\Omega) \ar[r]_{\alpha_1} \ar[d]_{\alpha_2} &
\Spec(K) \ar[d]^a \\
\Spec(K) \ar[r]^-b &
\Spec(K) \times_{c, Y} X
}
$$
soit commutatif. En composant avec la projection sur $\Spec(K)$,
on voit que $\alpha_1 = \alpha_2$. Notons leur valeur commune $\alpha$.
Alors la famille $\{\alpha : \Spec(\Omega) \to \Spec(K)\}$
est couvrante pour la topologie fpqc de $\Spec(K)$ et les deux morphismes
$a, b$ deviennent égaux après ce changement de base. D'après
Propriétés des espaces, Proposition
\ref{spaces-properties-proposition-sheaf-fpqc},
on a $a = b$. Nous en concluons que (1) est vérifiée.

\medskip\noindent
Supposons (3). Soient $x, x' \in |X|$ deux points tels que
$f(x) = f(x')$ dans $|Y|$. D'après
Propriétés des espaces, lemme \ref{spaces-properties-lemma-points-cartesian},
il existe $x'' \in |X \times_Y X|$ dont les projections
sont $x$ et $x'$. Par hypothèse et d'après
Propriétés des espaces,
lemme \ref{spaces-properties-lemma-characterize-surjective},
il existe $x''' \in |X|$ tel que $\Delta_{X/Y}(x''') = x''$.
Ainsi $x = x'$. Autrement dit, $f$ est injectif sur les points.
Puisque la condition (3) est stable par changement de base, on voit que $f$
vérifie (2).

\medskip\noindent
Supposons (2). Alors, en particulier, $|X \times_Y X| \to |X|$ est injective,
ce qui entraîne immédiatement que
$|\Delta_{X/Y}| : |X| \to |X \times_Y X|$ est surjective ; cela
entraîne que $\Delta_{X/Y}$ est surjectif d'après
Propriétés des espaces,
lemme \ref{spaces-properties-lemma-characterize-surjective}.
\end{proof}

\noindent
D'après les deux lemmes précédents, la définition suivante est compatible avec
la notion déjà définie de morphisme représentable universellement injectif
d'espaces algébriques.

\begin{definition}
\label{spaces-morphisms-definition-universally-injective}
Soit $S$ un schéma. Soit $f : X \to Y$ un morphisme d'espaces
algébriques sur $S$. Nous disons que $f$ est {\it universellement injectif} si,
pour tout morphisme $Y' \to Y$, l'application induite
$|Y' \times_Y X| \to |Y'|$ est injective.
\end{definition}

\noindent
Précisément, cela signifie qu'une, et donc chacune, des conditions équivalentes du
lemme \ref{spaces-morphisms-lemma-universally-injective}
est satisfaite.

\begin{remark}
\label{spaces-morphisms-remark-universally-injective-not-separated}
Un morphisme universellement injectif de schémas est séparé ; voir
Morphismes, lemme
\ref{morphisms-lemma-universally-injective-separated}.
Ce n'est pas le cas pour les morphismes d'espaces algébriques.
En effet, l'espace algébrique
$X = \mathbf{A}^1_k/\{x \sim -x \mid x \not = 0\}$
construit dans
Espaces, exemple \ref{spaces-example-affine-line-involution}
est muni d'un morphisme $X \to \mathbf{A}^1_k$ qui envoie
le point de coordonnée $x$ sur le point de coordonnée $x^2$.
C'est un isomorphisme hors de $0$, et il existe un unique point
de $X$ au-dessus de $0$. Comme $X$ n'est pas séparé, il s'agit d'un morphisme
universellement injectif d'espaces algébriques qui n'est pas séparé.
\end{remark}

\begin{lemma}
\label{spaces-morphisms-lemma-base-change-universally-injective}
Tout changement de base d'un morphisme universellement injectif est universellement injectif.
\end{lemma}

\begin{proof}
Omis. Indication : c'est formel.
\end{proof}

\begin{lemma}
\label{spaces-morphisms-lemma-universally-injective-local}
Soit $S$ un schéma.
Soit $f : X \to Y$ un morphisme d'espaces algébriques sur $S$.
Les conditions suivantes sont équivalentes :
\begin{enumerate}
\item $f$ est universellement injectif ;
\item pour tout schéma $Z$ et tout morphisme $Z \to Y$, le morphisme
$Z \times_Y X \to Z$ est universellement injectif ;
\item pour tout schéma affine $Z$ et tout morphisme
$Z \to Y$, le morphisme $Z \times_Y X \to Z$ est universellement injectif ;
\item il existe un schéma $Z$ et un morphisme surjectif
$Z \to Y$ tels que $Z \times_Y X \to Z$ soit universellement injectif ; et
\item il existe un recouvrement de Zariski $Y = \bigcup Y_i$ tel que
chacun des morphismes $f^{-1}(Y_i) \to Y_i$ soit universellement injectif.
\end{enumerate}
\end{lemma}

\begin{proof}
Nous utiliserons sans autre mention dans cette démonstration le fait que la propriété
d'être universellement injectif est préservée par changement de base
(lemme \ref{spaces-morphisms-lemma-base-change-universally-injective}).
Il est clair que (1) $\Rightarrow$ (2) $\Rightarrow$ (3)
$\Rightarrow$ (4).

\medskip\noindent
Supposons donné $g : Z \to Y$ comme dans (4). Soit $y : \Spec(K) \to Y$ un
morphisme du spectre d'un corps vers $Y$. Par hypothèse, on
peut trouver une extension de corps $\alpha : K \subset K'$ et un morphisme
$z : \Spec(K') \to Z$ tels que $y \circ \alpha = g \circ z$
(avec l'abus de notation évident). Par hypothèse, le
morphisme $Z \times_Y X \to Z$ est universellement injectif ; il existe donc
au plus un
relèvement de $g \circ z : \Spec(K') \to Y$ en un morphisme vers $X$.
Comme la famille $\{\alpha : \Spec(K') \to \Spec(K)\}$ est
couvrante pour la topologie fpqc, il existe au plus un relèvement de
$y : \Spec(K) \to Y$ en un morphisme vers $X$ ; voir
Propriétés des espaces, Proposition
\ref{spaces-properties-proposition-sheaf-fpqc}.
Ainsi, (1) est vérifiée.

\medskip\noindent
Nous omettons de vérifier que (5) est équivalente à (1).
\end{proof}

\begin{lemma}
\label{spaces-morphisms-lemma-composition-universally-injective}
Tout composé de morphismes universellement injectifs est universellement injectif.
\end{lemma}

\begin{proof}
Omis.
\end{proof}











\section{Morphismes affines}
\label{spaces-morphisms-section-affine}

\noindent
Nous avons déjà défini à la section \ref{spaces-morphisms-section-representable}
ce que signifie, pour un morphisme représentable d'espaces algébriques,
être affine.

\begin{lemma}
\label{spaces-morphisms-lemma-affine-representable}
Soit $S$ un schéma. Soit $f : X \to Y$ un morphisme représentable
d'espaces algébriques sur $S$. Alors
$f$ est affine (au sens de la section \ref{spaces-morphisms-section-representable})
si et seulement si, pour tout schéma affine $Z$
et tout morphisme $Z \to Y$, le schéma $X \times_Y Z$ est affine.
\end{lemma}

\begin{proof}
Cela résulte directement de la définition d'un morphisme affine de schémas
(Morphismes, définition \ref{morphisms-definition-affine}).
\end{proof}

\noindent
Nous pouvons ainsi poser la définition suivante.

\begin{definition}
\label{spaces-morphisms-definition-affine}
Soit $S$ un schéma.
Soit $f : X \to Y$ un morphisme d'espaces algébriques sur $S$.
Nous disons que $f$ est {\it affine} si, pour tout schéma affine $Z$ et tout
morphisme $Z \to Y$, l'espace algébrique $X \times_Y Z$ est représentable
par un schéma affine.
\end{definition}

\begin{lemma}
\label{spaces-morphisms-lemma-affine-local}
Soit $S$ un schéma.
Soit $f : X \to Y$ un morphisme d'espaces algébriques sur $S$.
Les conditions suivantes sont équivalentes :
\begin{enumerate}
\item $f$ est représentable et affine ;
\item $f$ est affine ;
\item pour tout schéma affine $V$ et tout morphisme étale $V \to Y$,
le schéma $X \times_Y V$ est affine ;
\item il existe un schéma $V$ et un morphisme étale surjectif
$V \to Y$ tels que $V \times_Y X \to V$ soit affine ; et
\item il existe un recouvrement de Zariski $Y = \bigcup Y_i$ tel
que chacun des morphismes $f^{-1}(Y_i) \to Y_i$ soit affine.
\end{enumerate}
\end{lemma}

\begin{proof}
Il est clair que (1) implique (2), que (2) implique (3), et que
(3) implique (4) en prenant pour $V$ une somme disjointe de schémas affines
étales sur $Y$ ; voir Propriétés des espaces,
lemme \ref{spaces-properties-lemma-cover-by-union-affines}.
Supposons $V \to Y$ comme dans (4). Pour tout ouvert affine $W$ de $V$,
$W \times_Y X$ est alors un ouvert affine de $V \times_Y X$. Par
Propriétés des espaces, lemme \ref{spaces-properties-lemma-subscheme},
on en déduit que $V \times_Y X$ est un schéma. De plus, le morphisme
$V \times_Y X \to V$ est affine. Nous pouvons donc appliquer
Espaces,
lemme \ref{spaces-lemma-morphism-sheaves-with-P-effective-descent-etale},
car la classe des morphismes affines satisfait à toutes les propriétés
requises (voir
Morphismes, lemmes \ref{morphisms-lemma-base-change-affine} et
Descente, lemmes \ref{descent-lemma-descending-property-affine}
et \ref{descent-lemma-affine}). Ce lemme donne
que $f$ est représentable et affine, c'est-à-dire que (1) est vérifiée.

\medskip\noindent
L'équivalence de (1) et (5) résulte de ce que la propriété
d'être affine est locale sur le but pour la topologie de Zariski (la référence précédente montre
qu'elle est en fait locale sur le but pour la topologie fpqc).
\end{proof}

\begin{lemma}
\label{spaces-morphisms-lemma-composition-affine}
Tout composé de morphismes affines est affine.
\end{lemma}

\begin{proof}
Omis. Indication : transitivité des produits fibrés.
\end{proof}

\begin{lemma}
\label{spaces-morphisms-lemma-base-change-affine}
Tout changement de base d'un morphisme affine est affine.
\end{lemma}

\begin{proof}
Omis. Indication : transitivité des produits fibrés.
\end{proof}

\begin{lemma}
\label{spaces-morphisms-lemma-closed-immersion-affine}
Toute immersion fermée est affine.
\end{lemma}

\begin{proof}
Cela résulte immédiatement de l'énoncé correspondant pour les morphismes de
schémas ; voir Morphismes, lemme \ref{morphisms-lemma-closed-immersion-affine}.
\end{proof}

\begin{lemma}
\label{spaces-morphisms-lemma-affine-equivalence-algebras}
Soit $S$ un schéma. Soit $X$ un espace algébrique sur $S$.
Il existe une anti-équivalence de catégories
$$
\begin{matrix}
\text{espaces algébriques} \\
\text{affines sur }X
\end{matrix}
\longleftrightarrow
\begin{matrix}
\text{faisceaux quasi-cohérents} \\
\text{d'}\mathcal{O}_X\text{-algèbres}
\end{matrix}
$$
qui associe à $f : Y \to X$ le faisceau $f_*\mathcal{O}_Y$.
De plus, cette équivalence est compatible avec tout changement de base.
\end{lemma}

\begin{proof}
Ce lemme est l'analogue de Morphismes, lemme
\ref{morphisms-lemma-affine-equivalence-algebras}.
Soit $\mathcal{A}$ un faisceau quasi-cohérent de $\mathcal{O}_X$-algèbres.
Nous allons construire un morphisme affine d'espaces algébriques
$\pi : Y = \underline{\Spec}_X(\mathcal{A}) \to X$ tel que
$\pi_*\mathcal{O}_Y \cong \mathcal{A}$. Pour cela, choisissons un schéma
$U$ et un morphisme étale surjectif $\varphi : U \to X$. Posons comme d'habitude
$R = U \times_X U$, avec projections $s, t : R \to U$. Notons
$\psi : R \to X$ le composé $\psi = \varphi \circ s = \varphi \circ t$.
D'après le lemme précité, il existe des morphismes affines de schémas
$\pi_0 : V \to U$ et $\pi_1 : W \to R$ tels que
$\pi_{0, *}\mathcal{O}_V \cong \varphi^*\mathcal{A}$ et
$\pi_{1, *}\mathcal{O}_W \cong \psi^*\mathcal{A}$.
Puisque la construction est compatible avec le changement de base, il existe
des morphismes $s', t' : W \to V$ tels que les diagrammes
$$
\vcenter{
\xymatrix{
W \ar[r]_{s'} \ar[d] & V \ar[d] \\
R \ar[r]^s & U
}
}
\quad\text{et}\quad
\vcenter{
\xymatrix{
W \ar[r]_{t'} \ar[d] & V \ar[d] \\
R \ar[r]^t & U
}
}
$$
soient cartésiens. Il s'ensuit que $s', t'$ sont étales. Ce qui précède entraîne
formellement que $(t', s') : W \to V \times_S V$ est un
monomorphisme. Nous omettons de vérifier que $W \to V \times_S V$ est une
relation d'équivalence (indication : considérer l'image inverse de $\mathcal{A}$
sur $U \times_X U \times_X U = R \times_{s, U, t} R$).
Le faisceau quotient $Y = V/W$
est un espace algébrique ; voir
Espaces, théorème \ref{spaces-theorem-presentation}.
D'après Groupoïdes, lemme \ref{groupoids-lemma-criterion-fibre-product},
on a $Y \times_X U \cong V$. Ainsi, $Y \to X$ est affine par le
lemme \ref{spaces-morphisms-lemma-affine-local}. Enfin, l'isomorphisme
$$
(Y \times_X U \to U)_*\mathcal{O}_{Y \times_X U} =
\pi_{0, *}\mathcal{O}_V \cong \varphi^*\mathcal{A}
$$
est compatible avec les isomorphismes de recollement, d'où
$(Y \to X)_*\mathcal{O}_Y \cong \mathcal{A}$ d'après
Propriétés des espaces, Proposition
\ref{spaces-properties-proposition-quasi-coherent}.
Nous omettons de vérifier que cette construction est compatible avec tout
changement de base.
\end{proof}

\begin{definition}
\label{spaces-morphisms-definition-relative-spec}
Soit $S$ un schéma. Soit $X$ un espace algébrique sur $S$.
Soit $\mathcal{A}$ un faisceau quasi-cohérent de
$\mathcal{O}_X$-algèbres. Le {\it spectre relatif de $\mathcal{A}$ sur
$X$}, ou simplement le {\it spectre de $\mathcal{A}$ sur $X$}, est le
morphisme affine $\underline{\Spec}(\mathcal{A}) \to X$
correspondant à $\mathcal{A}$ par l'équivalence de catégories du
lemme \ref{spaces-morphisms-lemma-affine-equivalence-algebras}.
\end{definition}

\noindent
La formation du spectre relatif commute à tout changement de base.

\begin{remark}
\label{spaces-morphisms-remark-factorization-quasi-compact-quasi-separated}
Soit $S$ un schéma. Soit $f : Y \to X$ un morphisme quasi-compact et
quasi-séparé d'espaces algébriques sur $S$. Alors
$f$ admet une factorisation canonique
$$
Y \longrightarrow \underline{\Spec}_X(f_*\mathcal{O}_Y) \longrightarrow X
$$
Cette écriture a un sens car $f_*\mathcal{O}_Y$ est quasi-cohérent
par le lemme \ref{spaces-morphisms-lemma-pushforward}. Le morphisme
$Y \to \underline{\Spec}_X(f_*\mathcal{O}_Y)$ provient du
morphisme canonique de $\mathcal{O}_Y$-algèbres
$f^*f_*\mathcal{O}_Y \to \mathcal{O}_Y$, auquel correspond
un morphisme canonique
$Y \to Y \times_X \underline{\Spec}_X(f_*\mathcal{O}_Y)$ sur $Y$ (voir
le lemme \ref{spaces-morphisms-lemma-affine-equivalence-algebras}), d'où une factorisation
de $f$ comme ci-dessus.
\end{remark}

\begin{lemma}
\label{spaces-morphisms-lemma-affine-equivalence-modules}
Soit $S$ un schéma.
Soit $f : Y \to X$ un morphisme affine d'espaces algébriques sur $S$.
Posons $\mathcal{A} = f_*\mathcal{O}_Y$.
Le foncteur $\mathcal{F} \mapsto f_*\mathcal{F}$ induit
une équivalence de catégories
$$
\left\{
\begin{matrix}
\text{catégorie des}\\
\mathcal{O}_Y\text{-modules quasi-cohérents}
\end{matrix}
\right\}
\longrightarrow
\left\{
\begin{matrix}
\text{catégorie des}\\
\mathcal{A}\text{-modules quasi-cohérents}
\end{matrix}
\right\}
$$
De plus, un $\mathcal{A}$-module est
quasi-cohérent comme $\mathcal{O}_X$-module si et seulement s'il
est quasi-cohérent comme $\mathcal{A}$-module.
\end{lemma}

\begin{proof}
Omis.
\end{proof}

\begin{lemma}
\label{spaces-morphisms-lemma-affine-permanence}
Soit $S$ un schéma. Soit $B$ un espace algébrique sur $S$.
Supposons que $g : X \to Y$ soit un morphisme d'espaces algébriques sur $B$.
\begin{enumerate}
\item Si $X$ est affine sur $B$ et si $\Delta : Y \to Y \times_B Y$ est affine,
alors $g$ est affine.
\item Si $X$ est affine sur $B$ et si $Y$ est séparé sur $B$,
alors $g$ est affine.
\item Tout morphisme d'un schéma affine vers un espace algébrique dont la
diagonale sur $\mathbf{Z}$ est affine (au sens de Propriétés des espaces, définition
\ref{spaces-properties-definition-separated}) est affine.
\item Tout morphisme d'un schéma affine vers un espace algébrique séparé
est affine.
\end{enumerate}
\end{lemma}

\begin{proof}
Démonstration de (1). Le changement de base $X \times_B Y \to Y$ est affine par le
lemme \ref{spaces-morphisms-lemma-base-change-affine}.
Le morphisme $(1, g) : X \to X \times_B Y$ est le changement de base de
$Y \to Y \times_B Y$ par le morphisme $X \times_B Y \to Y \times_B Y$.
Il est donc affine par le
lemme \ref{spaces-morphisms-lemma-base-change-affine}.
Tout composé de morphismes affines est affine
(voir le lemme \ref{spaces-morphisms-lemma-composition-affine}), d'où (1).
L'assertion (2) résulte de (1), car une immersion fermée est affine
(voir le lemme \ref{spaces-morphisms-lemma-closed-immersion-affine}) et dire que $Y/B$ est séparé
signifie que $\Delta$ est une immersion fermée. Les assertions (3) et (4) sont des cas
particuliers de (1) et (2).
\end{proof}

\begin{lemma}
\label{spaces-morphisms-lemma-Artinian-affine}
Soit $S$ un schéma. Soit $X$ un espace algébrique quasi-séparé sur $S$.
Soit $A$ un anneau artinien. Tout morphisme $\Spec(A) \to X$ est affine.
\end{lemma}

\begin{proof}
Soit $U \to X$ un morphisme étale, avec $U$ affine. Pour démontrer le
lemme, il suffit de montrer que $\Spec(A) \times_X U$ est affine ; voir le
lemme \ref{spaces-morphisms-lemma-affine-local}. Comme $X$ est quasi-séparé, le
schéma $\Spec(A) \times_X U$ est quasi-compact. De plus, le
morphisme de projection $\Spec(A) \times_X U \to \Spec(A)$ est étale.
Ce morphisme a donc des fibres discrètes finies ; en outre,
la topologie de $\Spec(A)$ est discrète. Ainsi, $\Spec(A) \times_X U$
est un schéma dont l'espace topologique
sous-jacent est un ensemble discret fini. On conclut par
Schémas, lemme \ref{schemes-lemma-scheme-finite-discrete-affine}.
\end{proof}












\section{Morphismes quasi-affines}
\label{spaces-morphisms-section-quasi-affine}

\noindent
Nous avons déjà défini à la section \ref{spaces-morphisms-section-representable}
ce que signifie, pour un morphisme représentable d'espaces algébriques,
être quasi-affine.

\begin{lemma}
\label{spaces-morphisms-lemma-quasi-affine-representable}
Soit $S$ un schéma. Soit $f : X \to Y$ un morphisme représentable
d'espaces algébriques sur $S$. Alors
$f$ est quasi-affine (au sens de la section \ref{spaces-morphisms-section-representable})
si et seulement si, pour tout schéma affine $Z$
et tout morphisme $Z \to Y$, le schéma $X \times_Y Z$ est quasi-affine.
\end{lemma}

\begin{proof}
Cela résulte directement de la définition d'un morphisme quasi-affine
de schémas
(Morphismes, définition \ref{morphisms-definition-quasi-affine}).
\end{proof}

\noindent
Nous pouvons ainsi poser la définition suivante.

\begin{definition}
\label{spaces-morphisms-definition-quasi-affine}
Soit $S$ un schéma.
Soit $f : X \to Y$ un morphisme d'espaces algébriques sur $S$.
Nous disons que $f$ est {\it quasi-affine} si, pour tout schéma affine $Z$ et tout
morphisme $Z \to Y$, l'espace algébrique $X \times_Y Z$ est représentable
par un schéma quasi-affine.
\end{definition}

\begin{lemma}
\label{spaces-morphisms-lemma-quasi-affine-local}
Soit $S$ un schéma.
Soit $f : X \to Y$ un morphisme d'espaces algébriques sur $S$.
Les conditions suivantes sont équivalentes :
\begin{enumerate}
\item $f$ est représentable et quasi-affine ;
\item $f$ est quasi-affine ;
\item il existe un schéma $V$ et un morphisme étale surjectif
$V \to Y$ tels que $V \times_Y X \to V$ soit quasi-affine ; et
\item il existe un recouvrement de Zariski $Y = \bigcup Y_i$ tel
que chacun des morphismes $f^{-1}(Y_i) \to Y_i$ soit quasi-affine.
\end{enumerate}
\end{lemma}

\begin{proof}
Il est clair que (1) implique (2) et que (2) implique (3) en prenant
pour $V$ une somme disjointe de schémas affines étales sur $Y$ ; voir
Propriétés des espaces,
lemme \ref{spaces-properties-lemma-cover-by-union-affines}.
Supposons $V \to Y$ comme dans (3). Pour tout ouvert affine $W$ de $V$,
$W \times_Y X$ est alors un ouvert quasi-affine de $V \times_Y X$. Par
Propriétés des espaces, lemme \ref{spaces-properties-lemma-subscheme},
on en déduit que $V \times_Y X$ est un schéma. De plus, le morphisme
$V \times_Y X \to V$ est quasi-affine. Nous pouvons donc appliquer
Espaces,
lemme \ref{spaces-lemma-morphism-sheaves-with-P-effective-descent-etale},
car la classe des morphismes quasi-affines satisfait à toutes les
propriétés requises (voir
Morphismes, lemmes \ref{morphisms-lemma-base-change-quasi-affine} et
Descente, lemmes \ref{descent-lemma-descending-property-quasi-affine}
et \ref{descent-lemma-quasi-affine}). Ce lemme donne
que $f$ est représentable et quasi-affine, c'est-à-dire que (1) est vérifiée.

\medskip\noindent
L'équivalence de (1) et (4) résulte de ce que la propriété
d'être quasi-affine est locale sur le but pour la topologie de Zariski (la référence précédente
montre qu'elle est en fait locale sur le but pour la topologie fpqc).
\end{proof}

\begin{lemma}
\label{spaces-morphisms-lemma-composition-quasi-affine}
Tout composé de morphismes quasi-affines est quasi-affine.
\end{lemma}

\begin{proof}
Omis.
\end{proof}

\begin{lemma}
\label{spaces-morphisms-lemma-base-change-quasi-affine}
Tout changement de base d'un morphisme quasi-affine est quasi-affine.
\end{lemma}

\begin{proof}
Omis.
\end{proof}

\begin{lemma}
\label{spaces-morphisms-lemma-characterize-quasi-affine}
Soit $S$ un schéma.
Un morphisme quasi-compact et quasi-séparé d'espaces algébriques
$f : Y \to X$ est quasi-affine si et seulement si la factorisation canonique
$Y \to \underline{\Spec}_X(f_*\mathcal{O}_Y)$
(remarque \ref{spaces-morphisms-remark-factorization-quasi-compact-quasi-separated})
est une immersion ouverte.
\end{lemma}

\begin{proof}
Soit $U \to X$ un morphisme étale surjectif, où $U$ est un schéma.
Comme on peut vérifier que $f$ est quasi-affine après le changement de base à
$U$ (lemme \ref{spaces-morphisms-lemma-quasi-affine-local}), que $f_*\mathcal{O}_Y|_U$
est égal à $(Y \times_X U \to U)_*\mathcal{O}_{Y \times_X U}$
(Propriétés des espaces, lemme
\ref{spaces-properties-lemma-pushforward-etale-base-change-modules}), et
que la formation du spectre relatif commute au changement de base
(lemme \ref{spaces-morphisms-lemma-affine-equivalence-algebras}),
on voit que l'assertion se ramène au cas où $X$ est un schéma.
Si $X$ est un schéma et si $f$ est quasi-affine ou si
$Y \to \underline{\Spec}_X(f_*\mathcal{O}_Y)$ est une immersion ouverte,
alors $Y$ est également un schéma. On se ramène donc à
Morphismes, lemme \ref{morphisms-lemma-characterize-quasi-affine}.
\end{proof}







\section{Propriétés de morphismes locales sur la source et le but pour la topologie étale}
\label{spaces-morphisms-section-local-source-target}

\noindent
Étant donnée une propriété de morphismes de schémas qui est {\it locale sur la
source et le but pour la topologie étale}, voir
Descente, définition \ref{descent-definition-local-source-target}
nous pouvons l'utiliser pour définir une propriété correspondante
de morphismes d'espaces algébriques, en imposant l'une ou l'autre des
conditions équivalentes du lemme ci-dessous.

\begin{lemma}
\label{spaces-morphisms-lemma-local-source-target}
Soit $\mathcal{P}$ une propriété de morphismes de schémas
qui est locale sur la source et le but pour la topologie étale.
Soit $S$ un schéma.
Soit $f : X \to Y$ un morphisme d'espaces algébriques sur $S$.
Considérons des diagrammes commutatifs
$$
\xymatrix{
U \ar[d]_a \ar[r]_h & V \ar[d]^b \\
X \ar[r]^f & Y
}
$$
où $U$ et $V$ sont des schémas et où les flèches verticales sont étales.
Les conditions suivantes sont équivalentes :
\begin{enumerate}
\item pour tout diagramme comme ci-dessus, le morphisme $h$ possède la propriété
$\mathcal{P}$ ;
\item pour un certain diagramme comme ci-dessus dans lequel $a : U \to X$ est
surjectif, le morphisme $h$ possède la propriété $\mathcal{P}$.
\end{enumerate}
Si $X$ et $Y$ sont représentables, cela équivaut encore à ce que
$f$ (considéré comme morphisme de schémas) possède la propriété $\mathcal{P}$.
Si, de plus, $\mathcal{P}$ est préservée par tout changement de base et
locale sur la base pour la topologie fppf, cela équivaut encore, pour les
morphismes représentables $f$, à ce que $f$ possède la propriété $\mathcal{P}$
au sens de la section \ref{spaces-morphisms-section-representable}.
\end{lemma}

\begin{proof}
Montrons l'équivalence de (1) et (2).
L'implication (1) $\Rightarrow$ (2) est immédiate (compte tenu de
Espaces, lemme \ref{spaces-lemma-lift-morphism-presentations}).
Supposons que
$$
\xymatrix{
U \ar[d] \ar[r]_h & V \ar[d] \\
X \ar[r]^f & Y
}
\quad\quad
\xymatrix{
U' \ar[d] \ar[r]_{h'} & V' \ar[d] \\
X \ar[r]^f & Y
}
$$
soient deux diagrammes comme dans le lemme. Supposons que $U \to X$ soit
surjectif et que $h$ possède la propriété $\mathcal{P}$. Pour montrer que (2)
implique (1), il faut prouver que $h'$ possède $\mathcal{P}$. Considérons
pour cela le diagramme
$$
\xymatrix{
U \ar[d]_h &
U \times_X U' \ar[l] \ar[d]^{(h, h')} \ar[r] &
U' \ar[d]^{h'} \\
V &
V \times_Y V' \ar[l] \ar[r] &
V'
}
$$
D'après
Descente, lemme \ref{descent-lemma-local-source-target-characterize}
le fait que $h$ possède $\mathcal{P}$ entraîne que $(h, h')$ possède $\mathcal{P}$ ;
comme $U \times_X U' \to U'$ est surjectif, cela entraîne (par le même
lemme) que $h'$ possède $\mathcal{P}$.

\medskip\noindent
Si $X$ et $Y$ sont représentables, alors
Descente, lemme \ref{descent-lemma-local-source-target-characterize}
s'applique et montre que (1) et (2) équivalent au fait que $f$ possède
$\mathcal{P}$.

\medskip\noindent
Enfin, supposons que $f$ soit représentable, que $U, V, a, b, h$ soient
comme dans la partie (2) du lemme et que $\mathcal{P}$ soit préservée par
tout changement de base. Il faut montrer que, pour tout schéma
$Z$ et tout morphisme $Z \to Y$, le changement de base $Z \times_Y X \to Z$
possède la propriété $\mathcal{P}$. Considérons le diagramme
$$
\xymatrix{
Z \times_Y U \ar[d] \ar[r] &
Z \times_Y V \ar[d] \\
Z \times_Y X \ar[r] &
Z
}
$$
Notons que la flèche horizontale supérieure est un changement de base de $h$
et possède donc la propriété $\mathcal{P}$. La flèche verticale gauche est étale
et surjective, tandis que la flèche verticale droite est étale. Ainsi,
Descente, lemme \ref{descent-lemma-local-source-target-characterize}
intervient une fois encore et montre que $Z \times_Y X \to Z$
possède la propriété $\mathcal{P}$.
\end{proof}

\begin{definition}
\label{spaces-morphisms-definition-P}
Soit $S$ un schéma.
Soit $\mathcal{P}$ une propriété de morphismes de schémas
qui est locale sur la source et le but pour la topologie étale.
On dit qu'un morphisme $f : X \to Y$ d'espaces algébriques sur $S$
{\it possède la propriété $\mathcal{P}$} si les conditions équivalentes du
lemme \ref{spaces-morphisms-lemma-local-source-target}
sont satisfaites.
\end{definition}

\noindent
Voici deux remarques immédiates.

\begin{remark}
\label{spaces-morphisms-remark-composition-P}
Soit $S$ un schéma. Soit $\mathcal{P}$ une propriété de morphismes de schémas
qui est locale sur la source et le but pour la topologie étale. Supposons de plus
que $\mathcal{P}$ soit stable par composition. Alors la classe des morphismes
d'espaces algébriques possédant la propriété $\mathcal{P}$ est stable par composition.
\end{remark}

\begin{remark}
\label{spaces-morphisms-remark-base-change-P}
Soit $S$ un schéma. Soit $\mathcal{P}$ une propriété de morphismes de schémas
qui est locale sur la source et le but pour la topologie étale. Supposons de plus
que $\mathcal{P}$ soit stable par changement de base. Alors la classe des morphismes
d'espaces algébriques possédant la propriété $\mathcal{P}$ est stable par changement de base.
\end{remark}

\noindent
Étant donnée une propriété de morphismes de germes de schémas qui est {\it locale
sur la source et le but pour la topologie étale}, voir
Descente, définition \ref{descent-definition-local-source-target-at-point}
nous pouvons l'utiliser pour définir une propriété correspondante de morphismes
d'espaces algébriques en un point, en imposant l'une ou l'autre des
conditions équivalentes du lemme ci-dessous.

\begin{lemma}
\label{spaces-morphisms-lemma-local-source-target-at-point}
Soit $\mathcal{Q}$ une propriété de morphismes de germes qui est
locale sur la source et le but pour la topologie étale.
Soit $S$ un schéma.
Soit $f : X \to Y$ un morphisme d'espaces algébriques sur $S$.
Soit $x \in |X|$ un point de $X$.
Considérons les diagrammes
$$
\xymatrix{
U \ar[d]_a \ar[r]_h & V \ar[d]^b \\
X \ar[r]^f & Y
}
\quad\quad
\xymatrix{
u \ar[d] \ar[r] & v \ar[d] \\
x \ar[r] & y
}
$$
où $U$ et $V$ sont des schémas, où $a, b$ sont étales et où $u, v, x, y$
sont des points des espaces correspondants.
Les conditions suivantes sont équivalentes :
\begin{enumerate}
\item pour tout diagramme comme ci-dessus, on a $\mathcal{Q}((U, u) \to (V, v))$ ;
\item pour un certain diagramme comme ci-dessus, on a $\mathcal{Q}((U, u) \to (V, v))$.
\end{enumerate}
Si $X$ et $Y$ sont représentables, cela équivaut encore
à $\mathcal{Q}((X, x) \to (Y, y))$.
\end{lemma}

\begin{proof}
Omis. Indication : la démonstration est très semblable à celle du
lemme \ref{spaces-morphisms-lemma-local-source-target}.
\end{proof}

\begin{definition}
\label{spaces-morphisms-definition-P-at-point}
Soit $\mathcal{Q}$ une propriété de morphismes de germes
de schémas qui est locale sur la source et le but pour la topologie étale.
Soit $S$ un schéma.
Étant donnés un morphisme $f : X \to Y$ d'espaces algébriques sur $S$ et
un point $x \in |X|$, on dit que $f$
{\it possède la propriété $\mathcal{Q}$ en $x$} si les conditions équivalentes du
lemme \ref{spaces-morphisms-lemma-local-source-target-at-point}
sont satisfaites.
\end{definition}

\noindent
Il ne faut pas utiliser aveuglément le lemme suivant pour passer d'une propriété
de morphismes à une propriété de morphismes en un point. Par exemple, si
$\mathcal{P}$ est la propriété d'être plat, alors la propriété
$\mathcal{Q}$ du lemme suivant signifie « $f$ est plat dans un voisinage ouvert
de $x$ », ce qui n'est pas équivalent à « $f$ est plat en $x$ ».

\begin{lemma}
\label{spaces-morphisms-lemma-local-source-target-global-implies-local}
Soit $\mathcal{P}$ une propriété de morphismes de schémas
qui est locale sur la source et le but pour la topologie étale.
Considérons la propriété $\mathcal{Q}$ de morphismes
de germes associée à $\mathcal{P}$ dans
Descente, lemme \ref{descent-lemma-local-source-target-global-implies-local}.
Alors :
\begin{enumerate}
\item $\mathcal{Q}$ est locale sur la source et le but pour la topologie étale ;
\item étant donnés un morphisme d'espaces algébriques $f : X \to Y$ et $x \in |X|$,
les conditions suivantes sont équivalentes :
\begin{enumerate}
\item $f$ possède $\mathcal{Q}$ en $x$ ;
\item il existe un voisinage ouvert $X' \subset X$ de $x$
tel que $X' \to Y$ possède $\mathcal{P}$ ;
\end{enumerate}
\item étant donné un morphisme d'espaces algébriques $f : X \to Y$,
les conditions suivantes sont équivalentes :
\begin{enumerate}
\item $f$ possède $\mathcal{P}$ ;
\item pour tout $x \in |X|$, le morphisme $f$ possède $\mathcal{Q}$ en $x$.
\end{enumerate}
\end{enumerate}
\end{lemma}

\begin{proof}
Voir
Descente, lemme \ref{descent-lemma-local-source-target-global-implies-local}
pour (1). L'implication (2)(a) $\Rightarrow$ (2)(b) s'obtient en posant
$X' = a(U) \subset X$ dans un diagramme comme dans le
lemme \ref{spaces-morphisms-lemma-local-source-target-at-point}.
L'implication (2)(b) $\Rightarrow$ (2)(a) est claire.
L'équivalence de (3)(a) et (3)(b) résulte de l'énoncé correspondant
pour les morphismes de schémas ; voir
Descente, lemme \ref{descent-lemma-local-source-target-local-implies-global}.
\end{proof}

\begin{remark}
\label{spaces-morphisms-remark-when-apply}
Nous appliquerons le
lemme \ref{spaces-morphisms-lemma-local-source-target-global-implies-local}
ci-dessus à tous les cas énumérés dans
Descente, remarque \ref{descent-remark-list-local-source-target}
sauf « plat ». Dans chaque cas, nous le ferons en disant que
$f$ possède la propriété $\mathcal{P}$ en $x$ si $f$ possède
$\mathcal{P}$ dans un voisinage de $x$.
\end{remark}




\section{Morphismes de type fini}
\label{spaces-morphisms-section-finite-type}

\noindent
La propriété « localement de type fini » des morphismes de schémas
est locale sur la source et le but pour la topologie étale ; voir
Descente, remarque \ref{descent-remark-list-local-source-target}.
Elle est aussi stable par changement de base et locale sur le but pour
la topologie fpqc ; voir Morphismes, lemme \ref{morphisms-lemma-base-change-finite-type}, et
Descente, lemmes \ref{descent-lemma-descending-property-locally-finite-type}.
Ainsi, d'après le
lemme \ref{spaces-morphisms-lemma-local-source-target}
ci-dessus, nous pouvons définir comme suit ce que signifie, pour un morphisme
d'espaces algébriques, être localement de type fini ;
cette définition coïncide avec la notion déjà définie à la
section \ref{spaces-morphisms-section-representable}
lorsque le morphisme est représentable.

\begin{definition}
\label{spaces-morphisms-definition-locally-finite-type}
Soit $S$ un schéma.
Soit $f : X \to Y$ un morphisme d'espaces algébriques sur $S$.
\begin{enumerate}
\item On dit que $f$ est
{\it localement de type fini} si les conditions équivalentes du
lemme \ref{spaces-morphisms-lemma-local-source-target}
sont satisfaites pour
$\mathcal{P} = \text{localement de type fini}$.
\item Soit $x \in |X|$. On dit que $f$ est {\it de type fini en $x$}
s'il existe un voisinage ouvert $X' \subset X$ de $x$ tel
que $f|_{X'} : X' \to Y$ soit localement de type fini.
\item On dit que $f$ est
{\it de type fini} s'il est localement de type fini et quasi-compact.
\end{enumerate}
\end{definition}

\noindent
Considérons l'espace algébrique $\mathbf{A}^1_k/\mathbf{Z}$ de
Espaces, exemple \ref{spaces-example-affine-line-translation}.
Le morphisme $\mathbf{A}^1_k/\mathbf{Z} \to \Spec(k)$
est de type fini.

\begin{lemma}
\label{spaces-morphisms-lemma-composition-finite-type}
Le composé de morphismes de type fini est de type fini.
Il en est de même des morphismes localement de type fini.
\end{lemma}

\begin{proof}
Voir la remarque \ref{spaces-morphisms-remark-composition-P} et
Morphismes, lemme \ref{morphisms-lemma-composition-finite-type}.
\end{proof}

\begin{lemma}
\label{spaces-morphisms-lemma-base-change-finite-type}
Tout changement de base d'un morphisme de type fini est de type fini.
Il en est de même des morphismes localement de type fini.
\end{lemma}

\begin{proof}
Voir la remarque \ref{spaces-morphisms-remark-base-change-P} et
Morphismes, lemme \ref{morphisms-lemma-base-change-finite-type}.
\end{proof}

\begin{lemma}
\label{spaces-morphisms-lemma-finite-type-local}
Soit $S$ un schéma.
Soit $f : X \to Y$ un morphisme d'espaces algébriques sur $S$.
Les conditions suivantes sont équivalentes :
\begin{enumerate}
\item $f$ est localement de type fini ;
\item pour tout $x \in |X|$, le morphisme $f$ est de type fini en $x$ ;
\item pour tout schéma $Z$ et tout morphisme $Z \to Y$, le morphisme
$Z \times_Y X \to Z$ est localement de type fini ;
\item pour tout schéma affine $Z$ et tout morphisme
$Z \to Y$, le morphisme $Z \times_Y X \to Z$ est localement de type fini ;
\item il existe un schéma $V$ et un morphisme étale surjectif
$V \to Y$ tels que $V \times_Y X \to V$ soit localement de type fini ;
\item il existe un schéma $U$ et un morphisme étale surjectif
$\varphi : U \to X$ tels que le composé $f \circ \varphi$
soit localement de type fini ;
\item pour tout diagramme commutatif
$$
\xymatrix{
U \ar[d] \ar[r] & V \ar[d] \\
X \ar[r] & Y
}
$$
où $U$, $V$ sont des schémas et où les flèches verticales sont étales,
la flèche horizontale supérieure est localement de type fini ;
\item il existe un diagramme commutatif
$$
\xymatrix{
U \ar[d] \ar[r] & V \ar[d] \\
X \ar[r] & Y
}
$$
où $U$, $V$ sont des schémas, où les flèches verticales sont étales, où $U \to X$
est surjectif et où la flèche horizontale supérieure est localement de type fini ;
\item il existe des recouvrements de Zariski $Y = \bigcup_{i \in I} Y_i$
et $f^{-1}(Y_i) = \bigcup X_{ij}$ tels que
chaque morphisme $X_{ij} \to Y_i$ soit localement de type fini.
\end{enumerate}
\end{lemma}

\begin{proof}
Chacune des conditions (2), (3), (4), (5), (6), (7) et (9)
implique directement la condition (8). Par exemple, si
(5) est satisfaite, nous pouvons choisir un schéma $V$ qui soit une somme
disjointe de schémas affines et un morphisme étale surjectif $V \to Y$
(voir Propriétés des espaces, lemme
\ref{spaces-properties-lemma-cover-by-union-affines}).
Alors $V \times_Y X \to V$ est localement de type fini d'après (5).
Choisissons un schéma $U$ et un morphisme étale surjectif
$U \to V \times_Y X$. Alors $U \to V$ est localement de type
fini d'après le lemme \ref{spaces-morphisms-lemma-composition-finite-type}.
La condition (8) est donc satisfaite.

\medskip\noindent
Les conditions (1), (7) et (8) sont équivalentes par définition.

\medskip\noindent
Pour terminer la démonstration, montrons que (1) implique toutes les conditions
(2), (3), (4), (5), (6) et (9). Pour (2), c'est immédiat.
Pour (3), (4), (5) et (9), cela résulte du fait que la propriété d'être
localement de type fini est préservée par changement de base ; voir le
lemme \ref{spaces-morphisms-lemma-base-change-finite-type}.
Pour (6), prenons une présentation étale surjective $U \to X$ par un schéma ; le morphisme composé est alors localement de type fini.
\end{proof}

\begin{lemma}
\label{spaces-morphisms-lemma-locally-finite-type-locally-noetherian}
Soit $S$ un schéma.
Soit $f : X \to Y$ un morphisme d'espaces algébriques sur $S$.
Si $f$ est localement de type fini et si $Y$ est localement noethérien,
alors $X$ est localement noethérien.
\end{lemma}

\begin{proof}
Soit
$$
\xymatrix{
U \ar[d] \ar[r] & V \ar[d] \\
X \ar[r] & Y
}
$$
un diagramme commutatif où $U$, $V$ sont des schémas et où les flèches verticales
sont étales surjectives. Si $f$ est localement de type fini, alors
$U \to V$ est localement de type fini. Si $Y$ est localement noethérien, alors
$V$ est localement noethérien. D'après
Morphismes, lemme \ref{morphisms-lemma-finite-type-noetherian}
on voit que $U$ est localement noethérien, ce qui signifie que $X$ est localement
noethérien.
\end{proof}

\begin{lemma}
\label{spaces-morphisms-lemma-permanence-finite-type}
Soit $S$ un schéma.
Soient $f : X \to Y$, $g : Y \to Z$ des morphismes d'espaces algébriques sur $S$.
Si $g \circ f : X \to Z$ est localement de type fini, alors $f : X \to Y$
est localement de type fini.
\end{lemma}

\begin{proof}
On peut trouver un diagramme
$$
\xymatrix{
U \ar[r] \ar[d] & V \ar[r] \ar[d] & W \ar[d] \\
X \ar[r] & Y \ar[r] & Z
}
$$
où $U$, $V$, $W$ sont des schémas et où les flèches verticales sont étales et surjectives ;
voir
Espaces, lemme \ref{spaces-lemma-lift-morphism-presentations}.
On peut alors utiliser le
lemme \ref{spaces-morphisms-lemma-finite-type-local}
et
Morphismes, lemme \ref{morphisms-lemma-permanence-finite-type}
pour conclure.
\end{proof}

\begin{lemma}
\label{spaces-morphisms-lemma-immersion-locally-finite-type}
Une immersion est localement de type fini.
\end{lemma}

\begin{proof}
Cela résulte du principe général
Espaces, lemme
\ref{spaces-lemma-representable-transformations-property-implication}
et
Morphismes, lemme \ref{morphisms-lemma-immersion-locally-finite-type}.
\end{proof}







\section{Points et points géométriques}
\label{spaces-morphisms-section-points-fields}

\noindent
Dans cette section, nous faisons quelques remarques sur les points et les points géométriques (voir
Propriétés des espaces,
définition \ref{spaces-properties-definition-geometric-point}).
Une manière de concevoir un point géométrique de $X$ consiste à
considérer un point géométrique $\overline{s} : \Spec(k) \to S$ de $S$
et un relèvement de $\overline{s}$ en un morphisme $\overline{x}$ vers $X$.
On a le diagramme
$$
\xymatrix{
\Spec(k) \ar[r]_-{\overline{x}} \ar[rd]_{\overline{s}} & X \ar[d] \\
& S.
}
$$
Nous disons souvent « soit $k$ un corps algébriquement clos sur $S$ »
pour indiquer que $\Spec(k)$ est muni d'un morphisme
$\Spec(k) \to S$. Dans cette situation, nous écrivons
$$
X(k) = \Mor_S(\Spec(k), X) =
\{\overline{x} \in X\text{ situé au-dessus de }\overline{s}\}
$$
pour l'ensemble des points à valeurs dans $k$ de $X$. Dans ce cas, l'application
$X(k) \to |X|$ est à valeurs dans la partie $|X_s| \subset |X|$.
Ici, $X_s = \Spec(\kappa(s)) \times_S X$, où $s \in S$
est le point correspondant à $\overline{s}$. Comme
$\Spec(\kappa(s)) \to S$ est un monomorphisme, son changement de base
$X_s \to X$ est aussi un monomorphisme, et $|X_s|$ est bien une partie de $|X|$.

\begin{lemma}
\label{spaces-morphisms-lemma-locally-finite-type-surjective-geometric-points}
Soit $S$ un schéma. Soit $f : X \to Y$ un morphisme d'espaces algébriques
sur $S$. Supposons que $f$ soit localement de type fini. Les conditions suivantes sont équivalentes :
\begin{enumerate}
\item $f$ est surjectif ;
\item pour tout corps algébriquement clos $k$ sur $S$, l'application induite
$X(k) \to Y(k)$ est surjective.
\end{enumerate}
\end{lemma}

\begin{proof}
Choisissons un diagramme
$$
\xymatrix{
U \ar[d] \ar[r] & V \ar[d] \\
X \ar[r] & Y
}
$$
où $U$, $V$ sont des schémas sur $S$ et où les flèches verticales sont surjectives
et étales ; voir
Espaces, lemme \ref{spaces-lemma-lift-morphism-presentations}.
Comme $f$ est localement de type fini, on voit que $U \to V$ est localement de
type fini.

\medskip\noindent
Supposons (1) et soit $\overline{y} \in Y(k)$. Alors $U \to Y$ est
surjectif et localement de type fini d'après les
lemmes \ref{spaces-morphisms-lemma-composition-surjective} et
\ref{spaces-morphisms-lemma-composition-finite-type}.
Posons $Z = U \times_{Y, \overline{y}} \Spec(k)$. C'est un schéma.
La projection $Z \to \Spec(k)$
est surjective et localement de type fini d'après les
lemmes \ref{spaces-morphisms-lemma-base-change-surjective} et
\ref{spaces-morphisms-lemma-base-change-finite-type}.
Il résulte de
Variétés, lemme \ref{varieties-lemma-locally-finite-type-Jacobson}
que $Z$ possède un point à valeurs dans $k$, noté $\overline{z}$. L'image
$\overline{x} \in X(k)$ de $\overline{z}$ s'envoie sur $\overline{y}$,
comme voulu.

\medskip\noindent
Supposons (2). D'après
Propriétés des espaces,
lemme \ref{spaces-properties-lemma-characterize-surjective}
il suffit de montrer que $|X| \to |Y|$ est surjective.
Soit $y \in |Y|$. Choisissons $u \in U$ s'envoyant sur $y$.
Soit $k \supset \kappa(u)$ une clôture algébrique. Notons
$\overline{u} \in U(k)$ le point correspondant et
$\overline{y} \in Y(k)$ son image. Par hypothèse, il existe
$\overline{x} \in X(k)$ s'envoyant sur $\overline{y}$.
Il est alors clair que l'image $x \in |X|$ de $\overline{x}$
s'envoie sur $y$.
\end{proof}

\noindent
Pour énoncer le lemme suivant, introduisons la
notation suivante. Étant donné un schéma $T$, posons
$$
\lambda(T) = \sup\{\aleph_0, |\kappa(t)| ; t \in T\}.
$$
Autrement dit, $\lambda(T)$ est le plus petit cardinal infini qui majore
toutes les cardinalités des corps résiduels de $T$. Notons que, si $R$
est un anneau, la cardinalité de tout corps résiduel $\kappa(\mathfrak p)$
de $R$ est majorée par celle de $R$ (détails omis).
Il s'ensuit que $\lambda(T) \leq \text{size}(T)$, où
$\text{size}(T)$ est la taille du schéma $T$ introduite dans
Ensembles, section \ref{sets-section-categories-schemes}.
Si $L/K$ est une extension de corps de type fini, alors
$|K| \leq |L| \leq \max\{\aleph_0, |K|\}$. Il s'ensuit que, si $T' \to T$
est un morphisme de schémas localement de type
fini, alors $\lambda(T') \leq \lambda(T)$ ; si $T' \to T$ est en outre
surjectif, il y a égalité. Supposons ensuite que $S$ soit un schéma
et que $X$ soit un espace algébrique sur $S$. Dans ce cas, posons
$$
\lambda(X) := \lambda(U)
$$
où $U$ est un schéma quelconque sur $S$ muni d'un morphisme étale surjectif
vers $X$. Cette définition est indépendante du choix de $U$ :
en effet, étant donnés deux tels schémas $U$ et $U'$, le produit fibré
$U \times_X U'$ est un schéma qui admet un morphisme étale surjectif
vers $U$ comme vers $U'$, d'où $\lambda(U) = \lambda(U \times_X U') =
\lambda(U')$ d'après ce qui précède.

\begin{lemma}
\label{spaces-morphisms-lemma-large-enough}
Soit $S$ un schéma. Soient $X$, $Y$ des espaces algébriques sur $S$.
\begin{enumerate}
\item Lorsque $k$ parcourt les corps algébriquement clos sur $S$,
les points géométriques $\overline{y} \in Y(k)$ recouvrent tout $|Y|$.
\item Lorsque $k$ parcourt les corps algébriquement clos sur $S$ tels que
$|k| \geq \lambda(Y)$ et $|k| > \lambda(X)$, les points géométriques
$\overline{y} \in Y(k)$ recouvrent tout $|Y|$.
\item Pour tout point géométrique
$\overline{s} : \Spec(k) \to S$ tel que
$k$ soit de cardinalité $> \lambda(X)$, l'application
$$
X(k) \longrightarrow |X_s|
$$
est surjective.
\item Soit $X \to Y$ un morphisme d'espaces algébriques sur $S$.
Pour tout point géométrique $\overline{s} : \Spec(k) \to S$ tel que
$k$ soit de cardinalité $> \lambda(X)$, l'application
$$
X(k) \longrightarrow |X| \times_{|Y|} Y(k)
$$
est surjective.
\item Soit $X \to Y$ un morphisme d'espaces algébriques sur $S$.
Les conditions suivantes sont équivalentes :
\begin{enumerate}
\item l'application $X \to Y$ est surjective ;
\item pour tout corps algébriquement clos $k$ sur $S$ tel que
$|k| > \lambda(X)$ et $|k| \geq \lambda(Y)$, l'application $X(k) \to Y(k)$
est surjective.
\end{enumerate}
\end{enumerate}
\end{lemma}

\begin{proof}
Pour démontrer (1), choisissons un morphisme étale surjectif $V \to Y$, où
$V$ est un schéma. Pour chaque $v \in V$, choisissons une clôture algébrique
$\kappa(v) \subset k_v$. Considérons les morphismes
$\overline{y}_v : \Spec(k_v) \to V \to Y$.
Par construction de $|Y|$, ceux-ci recouvrent $|Y|$.

\medskip\noindent
Pour démontrer (2), nous utiliserons les deux faits suivants, dont nous omettons les démonstrations :
(\romannumeral1) si $K$ est un corps et $\overline{K}$ une clôture algébrique, alors
$|\overline{K}| \leq \max\{\aleph_0, |K|\}$ ;
(\romannumeral2) pour tout corps algébriquement clos $k$ et tout cardinal
$\aleph$ tel que $\aleph \geq |k|$, il existe une extension de corps
algébriquement clos $k'/k$ telle que $|k'| = \aleph$.
Posons maintenant $\aleph = \max\{\lambda(X), \lambda(Y)\}^+$.
Ici, $\lambda^+ > \lambda$ désigne le cardinal immédiatement supérieur ; voir
Ensembles, section \ref{sets-section-cardinals}.
D'après (\romannumeral1),
les corps $k_v$ construits au premier paragraphe
de la démonstration ont tous une cardinalité majorée par $\lambda(Y)$. Par suite,
d'après (\romannumeral2), il existe des extensions $k_v \subset k'_v$ telles que
$|k'_v| = \aleph$. Les morphismes
$\overline{y}'_v : \Spec(k'_v) \to Y$ recouvrent $|Y|$, comme voulu.
Pour achever réellement la démonstration de (2), il faut vérifier que les schémas
$\Spec(k'_v)$ sont (isomorphes à) des objets de $\Sch_{fppf}$,
car nos conventions imposent que tous les schémas soient des objets de
$\Sch_{fppf}$ ; le reste de ce paragraphe peut être omis par
qui ne s'intéresse pas aux considérations ensemblistes. Par
construction, il existe un objet $T$ de $\Sch_{fppf}$ tel que
$\lambda(X)$ et $\lambda(Y)$ soient majorés par $\text{size}(T)$.
D'après notre construction de la catégorie $\Sch_{fppf}$ dans
Topologies, définitions \ref{topologies-definition-big-fppf-site}
comme la catégorie $\Sch_\alpha$ construite dans
Ensembles, lemme \ref{sets-lemma-construct-category}
on voit que tout schéma dont la taille est $\leq \text{size}(T)^+$
est isomorphe à un objet de $\Sch_{fppf}$. Voir
l'expression de la fonction $Bound$ dans
Ensembles, équation (\ref{sets-equation-bound}).
Comme $\aleph \leq \text{size}(T)^+$, on conclut.

\medskip\noindent
Dans la partie (3), la notation $X_s$ désigne
le produit fibré $\Spec(\kappa(s)) \times_S X$, où $s \in S$
est le point correspondant à $\overline{s}$. La partie (3) résulte donc
de (4) en prenant $Y = \Spec(\kappa(s))$.

\medskip\noindent
Démontrons (4). Soit $X \to Y$ un morphisme d'espaces algébriques sur $S$.
Soit $k$ un corps algébriquement clos sur $S$, de cardinalité
$> \lambda(X)$. Soient $\overline{y} \in Y(k)$ et $x \in |X|$ ayant pour image le
même élément $y$ de $|Y|$.
Il faut trouver $\overline{x} \in X(k)$ s'envoyant sur $x$ et $\overline{y}$.
Choisissons un diagramme commutatif
$$
\xymatrix{
U \ar[d] \ar[r] & V \ar[d] \\
X \ar[r] & Y
}
$$
où $U$, $V$ sont des schémas sur $S$ et où les flèches verticales sont surjectives
et étales ; voir
Espaces, lemme \ref{spaces-lemma-lift-morphism-presentations}.
Choisissons $u \in |U|$ s'envoyant sur $x$, et notons $v \in |V|$ son image.
Nous considérerons $u = \Spec(\kappa(u))$ et
$v = \Spec(\kappa(v))$ comme des schémas.
Notons que $V \times_Y \Spec(k)$ est un schéma étale sur $k$.
C'est donc une somme disjointe de spectres d'extensions finies
séparables de $k$ ; voir
Morphismes, lemme \ref{morphisms-lemma-etale-over-field}.
Comme $v$ s'envoie sur $y$, on voit que $v \times_Y \Spec(k)$ est un
schéma non vide. Comme $v \to V$ est un monomorphisme, on voit que
$v \times_Y \Spec(k) \to V \times_Y \Spec(k)$
est un monomorphisme. Ainsi, $v \times_Y \Spec(k)$ est une somme
disjointe de spectres d'extensions finies séparables de $k$, d'après
Schémas, lemme \ref{schemes-lemma-mono-towards-spec-field}.
On en déduit que le morphisme
$v \times_Y \Spec(k) \to \Spec(k)$
possède une section, c'est-à-dire qu'il existe un morphisme
$\overline{v} : \Spec(k) \to V$ situé au-dessus de $v$ et de
$\overline{y}$. Considérons enfin le schéma
$$
u \times_{V, \overline{v}} \Spec(k)
=
\Spec(\kappa(u) \otimes_{\kappa(v)} k)
$$
où $\kappa(v) \to k$ est le morphisme de corps qui définit le morphisme
$\overline{v}$. Comme, par hypothèse, la cardinalité de $k$ est strictement supérieure à celle
de $\kappa(u)$, on peut appliquer
Algèbre, lemme \ref{algebra-lemma-base-change-Jacobson}
pour voir que tout idéal maximal
$\mathfrak m \subset \kappa(u) \otimes_{\kappa(v)} k$
a un corps résiduel algébrique sur $k$, donc égal à $k$.
Un tel idéal maximal fournit donc un morphisme
$\overline{u} : \Spec(k) \to U$ situé au-dessus de $u$ et s'envoyant sur
$\overline{v}$. Le composé $\Spec(k) \to U \to X$
est le point géométrique recherché $\overline{x} \in X(k)$.
Ceci achève la démonstration de (4).

\medskip\noindent
La partie (5) est une conséquence formelle des parties (2) et (4) et de
Propriétés des espaces,
lemme \ref{spaces-properties-lemma-characterize-surjective}.
\end{proof}








\section{Points de type fini}
\label{spaces-morphisms-section-points-finite-type}

\noindent
Soit $S$ un schéma. Soit $X$ un espace algébrique sur $S$.
Un point de type fini $x \in |X|$ est un point qui peut être représenté par
un morphisme $\Spec(k) \to X$ localement de type fini.
Les points de type fini remplacent avantageusement les points fermés pour les espaces
algébriques et les champs algébriques. Par exemple, il y en a toujours « assez ».

\begin{lemma}
\label{spaces-morphisms-lemma-point-finite-type}
Soit $S$ un schéma. Soit $X$ un espace algébrique sur $S$.
Soit $x \in |X|$. Les conditions suivantes sont équivalentes :
\begin{enumerate}
\item Il existe un morphisme $\Spec(k) \to X$
qui est localement de type fini et représente $x$.
\item Il existe un schéma $U$, un point fermé $u \in U$ et un morphisme étale
$\varphi : U \to X$ tels que $\varphi(u) = x$.
\end{enumerate}
\end{lemma}

\begin{proof}
Soient $u \in U$ et $U \to X$ comme en (2). Alors
$\Spec(\kappa(u)) \to U$ est de type fini, et $U \to X$ est
représentable et localement de type fini (d'après le principe général
Espaces, lemme
\ref{spaces-lemma-representable-transformations-property-implication}
et
Morphismes, lemmes \ref{morphisms-lemma-etale-locally-finite-presentation} et
\ref{morphisms-lemma-finite-presentation-finite-type}). On voit donc
que (1) est satisfaite d'après le
lemme \ref{spaces-morphisms-lemma-composition-finite-type}.

\medskip\noindent
Réciproquement, supposons que $\Spec(k) \to X$ soit localement de type fini et
représente $x$. Soit $U \to X$ un morphisme étale surjectif, où $U$
est un schéma. Par hypothèse, $U \times_X \Spec(k) \to U$ est localement
de type fini. Choisissons un point de type fini $v$ de
$U \times_X \Spec(k)$ (il en existe au moins un ; voir
Morphismes, lemme \ref{morphisms-lemma-identify-finite-type-points}).
D'après
Morphismes, lemme \ref{morphisms-lemma-finite-type-points-morphism}
l'image $u \in U$ de $v$ est un point de type fini de $U$.
Par conséquent, d'après
Morphismes, lemme \ref{morphisms-lemma-identify-finite-type-points}
quitte à rétrécir $U$, on peut supposer que $u$ est un point fermé de $U$, c'est-à-dire que
(2) est satisfaite.
\end{proof}

\begin{definition}
\label{spaces-morphisms-definition-finite-type-point}
Soit $S$ un schéma. Soit $X$ un espace algébrique sur $S$.
On dit qu'un point $x \in |X|$ est un {\it point de type fini}\footnote{C'est un
léger abus de langage : il serait peut-être plus correct de dire
« point localement de type fini ».} si les conditions équivalentes du
lemme \ref{spaces-morphisms-lemma-point-finite-type}
sont satisfaites. On note $X_{\text{ft-pts}}$ l'ensemble des points de type fini
de $X$.
\end{definition}

\noindent
On peut décrire comme suit l'ensemble des points de type fini.

\begin{lemma}
\label{spaces-morphisms-lemma-identify-finite-type-points}
Soit $S$ un schéma. Soit $X$ un espace algébrique sur $S$. On a
$$
X_{\text{ft-pts}} =
\bigcup\nolimits_{\varphi : U \to X\text{ étale }} |\varphi|(U_0)
$$
où $U_0$ est l'ensemble des points fermés de $U$.
Ici, on peut faire parcourir à $U$ tous les schémas étales sur $X$, ou tous les
schémas affines étales sur $X$.
\end{lemma}

\begin{proof}
Cela résulte immédiatement du
lemme \ref{spaces-morphisms-lemma-point-finite-type}.
\end{proof}

\begin{lemma}
\label{spaces-morphisms-lemma-finite-type-points-morphism}
Soit $S$ un schéma.
Soit $f : X \to Y$ un morphisme d'espaces algébriques sur $S$.
Si $f$ est localement de type fini, alors
$f(X_{\text{ft-pts}}) \subset Y_{\text{ft-pts}}$.
\end{lemma}

\begin{proof}
Prenons $x \in X_{\text{ft-pts}}$. Représentons $x$ par un morphisme localement de type fini
$x : \Spec(k) \to X$. Alors $f \circ x$ est localement de type fini d'après le
lemme \ref{spaces-morphisms-lemma-composition-finite-type}.
Ainsi, $f(x) \in Y_{\text{ft-pts}}$.
\end{proof}

\begin{lemma}
\label{spaces-morphisms-lemma-finite-type-points-surjective-morphism}
Soit $S$ un schéma.
Soit $f : X \to Y$ un morphisme d'espaces algébriques sur $S$.
Si $f$ est localement de type fini et surjectif, alors
$f(X_{\text{ft-pts}}) = Y_{\text{ft-pts}}$.
\end{lemma}

\begin{proof}
On a $f(X_{\text{ft-pts}}) \subset Y_{\text{ft-pts}}$ d'après le
lemme \ref{spaces-morphisms-lemma-finite-type-points-morphism}.
Soit $y \in |Y|$ un point de type fini. Représentons $y$ par un morphisme
$\Spec(k) \to Y$ localement de type fini.
Comme $f$ est surjectif, l'espace algébrique
$X_k = \Spec(k) \times_Y X$ est non vide et possède donc
un point de type fini $x \in |X_k|$ d'après le
lemme \ref{spaces-morphisms-lemma-identify-finite-type-points}.
Or $X_k \to X$ est un morphisme localement de type fini comme changement de base
de $\Spec(k) \to Y$
(lemme \ref{spaces-morphisms-lemma-base-change-finite-type}).
Ainsi, l'image de $x$ dans $X$ est un point de type fini d'après le
lemme \ref{spaces-morphisms-lemma-finite-type-points-morphism}
et s'envoie sur $y$ par construction.
\end{proof}

\begin{lemma}
\label{spaces-morphisms-lemma-enough-finite-type-points}
Soit $S$ un schéma. Soit $X$ un espace algébrique sur $S$.
Pour toute partie localement fermée $T \subset |X|$, on a
$$
T \not = \emptyset
\Rightarrow
T \cap X_{\text{ft-pts}} \not = \emptyset.
$$
En particulier, pour toute partie fermée $T \subset |X|$,
$T \cap X_{\text{ft-pts}}$ est dense dans $T$.
\end{lemma}

\begin{proof}
Soit $i : Z \to X$ la structure réduite induite de sous-espace sur $T$ ; voir la
remarque \ref{spaces-morphisms-remark-space-structure-locally-closed-subset}.
Toute immersion est localement de type fini ; voir le
lemme \ref{spaces-morphisms-lemma-immersion-locally-finite-type}.
Ainsi, d'après le
lemme \ref{spaces-morphisms-lemma-finite-type-points-morphism},
on a $Z_{\text{ft-pts}} \subset X_{\text{ft-pts}} \cap T$.
Enfin, tout schéma affine non vide $U$ muni d'un morphisme étale vers
$Z$ possède au moins un point fermé. Ainsi, $Z$ possède au moins un
point de type fini d'après le
lemme \ref{spaces-morphisms-lemma-identify-finite-type-points}.
Le lemme en résulte.
\end{proof}

\noindent
Voici une autre caractérisation, plus technique, d'un point de type fini
d'un espace algébrique.

\begin{lemma}
\label{spaces-morphisms-lemma-point-finite-type-monomorphism}
Soit $S$ un schéma. Soit $X$ un espace algébrique sur $S$.
Soit $x \in |X|$. Les conditions suivantes sont équivalentes :
\begin{enumerate}
\item $x$ est un point de type fini ;
\item il existe un espace algébrique $Z$ dont l'espace topologique sous-jacent
$|Z|$ est un singleton, et un morphisme $f : Z \to X$ localement de type
fini tel que $\{x\} = |f|(|Z|)$ ;
\item il existe un espace algébrique $Z$ et un morphisme $f : Z \to X$
possédant les propriétés suivantes :
\begin{enumerate}
\item il existe un morphisme étale surjectif
$z : \Spec(k) \to Z$, où $k$ est un corps ;
\item $f$ est localement de type fini ;
\item $f$ est un monomorphisme ;
\item $x = f(z)$.
\end{enumerate}
\end{enumerate}
\end{lemma}

\begin{proof}
Supposons que $x$ soit un point de type fini. Choisissons un schéma affine $U$,
un point fermé $u \in U$ et un morphisme étale $\varphi : U \to X$
avec $\varphi(u) = x$ ; voir le
lemme \ref{spaces-morphisms-lemma-identify-finite-type-points}.
Posons comme d'habitude $u = \Spec(\kappa(u))$. Les morphismes de projection
$u \times_X u \to u$ sont les composés
$$
u \times_X u \to u \times_X U \to u \times_X X = u
$$
où la première flèche est une immersion fermée (un changement de base de
$u \to U$) et la seconde est étale (un changement de base du morphisme étale
$U \to X$). Ainsi, $u \times_X U$ est une somme disjointe de spectres
d'extensions finies séparables de $k$ (voir
Morphismes, lemme \ref{morphisms-lemma-etale-over-field})
et, par conséquent, le sous-schéma fermé $u \times_X u$ est une somme disjointe de
spectres d'extensions finies séparables de $k$ ; ainsi, $u \times_X u \to u$ est étale. D'après
Espaces, théorème \ref{spaces-theorem-presentation}
on voit que $Z = u/u \times_X u$ est un espace algébrique. Par construction,
le diagramme
$$
\xymatrix{
u \ar[d] \ar[r] & U \ar[d] \\
Z \ar[r] & X
}
$$
est commutatif et ses flèches verticales sont étales. Ainsi, $Z \to X$ est localement
de type fini (voir le
lemme \ref{spaces-morphisms-lemma-finite-type-local}).
Par construction, le morphisme $Z \to X$ est un monomorphisme et
l'image de $u$ est $x$. La condition (3) est donc satisfaite.

\medskip\noindent
Il est clair que (3) implique (2).
Si (2) est satisfaite, alors $x$ est un point de type fini de $X$ d'après le
lemme \ref{spaces-morphisms-lemma-finite-type-points-morphism}
(et le
lemme \ref{spaces-morphisms-lemma-enough-finite-type-points}
montre que $Z_{\text{ft-pts}}$ est non vide, c'est-à-dire que l'unique point de
$Z$ est un point de type fini de $Z$).
\end{proof}





\section{Espaces de Nagata}
\label{spaces-morphisms-section-nagata}

\noindent
Voir Propriétés des espaces, section
\ref{spaces-properties-section-types-properties} pour la définition
d'un espace algébrique de Nagata.

\begin{lemma}
\label{spaces-morphisms-lemma-finite-type-nagata}
Soit $S$ un schéma.
Soit $f : X \to Y$ un morphisme d'espaces algébriques sur $S$.
Si $Y$ est de Nagata et si $f$ est localement de type fini, alors $X$ est de Nagata.
\end{lemma}

\begin{proof}
Soit $V$ un schéma et soit $V \to Y$ un morphisme étale surjectif.
Soit $U$ un schéma et soit $U \to X \times_Y V$ un morphisme étale
surjectif. Si $Y$ est de Nagata, alors $V$ est un schéma de Nagata.
Si $X \to Y$ est localement de type fini, alors $U \to V$ est localement
de type fini. Ainsi, $U$ est un schéma de Nagata d'après
Morphismes, lemme \ref{morphisms-lemma-finite-type-nagata}.
Alors $X$ est de Nagata par définition.
\end{proof}

\begin{lemma}
\label{spaces-morphisms-lemma-ubiquity-nagata}
Les types suivants d'espaces algébriques sont de Nagata.
\begin{enumerate}
\item Tout espace algébrique localement de type fini sur un schéma de Nagata.
\item Tout espace algébrique localement de type fini sur un corps.
\item Tout espace algébrique localement de type fini sur un
anneau local noethérien complet.
\item Tout espace algébrique localement de type fini sur $\mathbf{Z}$.
\item Tout espace algébrique localement de type fini sur un anneau de Dedekind de
caractéristique nulle.
\item Et ainsi de suite.
\end{enumerate}
\end{lemma}

\begin{proof}
La première assertion résulte du lemme \ref{spaces-morphisms-lemma-finite-type-nagata}.
Il en va donc de même des autres ; voir
Morphismes, lemme \ref{morphisms-lemma-ubiquity-nagata}.
\end{proof}








\section{Morphismes quasi-finis}
\label{spaces-morphisms-section-quasi-finite}

\noindent
La propriété « localement quasi-fini » des morphismes de schémas est
locale sur la source et le but pour la topologie étale ; voir
Descente, remarque \ref{descent-remark-list-local-source-target}.
Elle est aussi stable par changement de base et locale sur le but pour la
topologie fpqc ; voir Morphismes, lemme \ref{morphisms-lemma-base-change-quasi-finite}, et
Descente, lemme \ref{descent-lemma-descending-property-quasi-finite}.
Ainsi, d'après le
lemme \ref{spaces-morphisms-lemma-local-source-target}
ci-dessus, on peut définir comme suit ce que signifie, pour un morphisme
d'espaces algébriques, d'être localement quasi-fini ; cette définition
coïncide avec la notion déjà définie dans la
section \ref{spaces-morphisms-section-representable}
lorsque le morphisme est représentable.

\begin{definition}
\label{spaces-morphisms-definition-locally-quasi-finite}
Soit $S$ un schéma.
Soit $f : X \to Y$ un morphisme d'espaces algébriques sur $S$.
\begin{enumerate}
\item On dit que $f$ est
{\it localement quasi-fini} si les conditions équivalentes du
lemme \ref{spaces-morphisms-lemma-local-source-target} sont satisfaites pour
$\mathcal{P} = \text{localement quasi-fini}$.
\item Soit $x \in |X|$. On dit que $f$ est {\it quasi-fini en $x$}
s'il existe un voisinage ouvert $X' \subset X$ de $x$ tel
que $f|_{X'} : X' \to Y$ soit localement quasi-fini.
\item Un morphisme d'espaces algébriques $f : X \to Y$ est
{\it quasi-fini} s'il est localement quasi-fini et quasi-compact.
\end{enumerate}
\end{definition}

\noindent
La dernière partie est compatible avec la notion de morphisme quasi-fini pour les morphismes
de schémas d'après
Morphismes, lemme \ref{morphisms-lemma-quasi-finite-locally-quasi-compact}.

\begin{lemma}
\label{spaces-morphisms-lemma-base-change-quasi-finite-locus}
Soit $S$ un schéma. Soient $f : X \to Y$ et $g : Y' \to Y$ des morphismes
d'espaces algébriques sur $S$. Notons $f' : X' \to Y'$ le changement de base
de $f$ par $g$. Notons $g' : X' \to X$ la projection.
Supposons que $f$ soit localement de type fini.
Soient $W \subset |X|$, resp.\ $W' \subset |X'|$,
l'ensemble des points où $f$, resp.\ $f'$, est quasi-fini.
\begin{enumerate}
\item $W \subset |X|$ et $W' \subset |X'|$ sont ouverts ;
\item $W' = (g')^{-1}(W)$, c'est-à-dire que la formation du lieu où
$f$ est quasi-fini commute aux changements de base ;
\item le changement de base d'un morphisme localement quasi-fini est
localement quasi-fini ;
\item le changement de base d'un morphisme quasi-fini est quasi-fini.
\end{enumerate}
\end{lemma}

\begin{proof}
Choisissons un schéma $V$ et un morphisme étale surjectif $V \to Y$.
Choisissons un schéma $U$ et un morphisme étale surjectif $U \to V \times_Y X$.
Choisissons un schéma $V'$ et un morphisme étale surjectif $V' \to Y' \times_Y V$.
Posons $U' = V' \times_V U$, de sorte que $U' \to X'$ soit lui aussi un morphisme
étale surjectif. On a le diagramme
$$
\vcenter{
\xymatrix{
U' \ar[d] \ar[r] & U \ar[d] \\
V' \ar[r] & V
}
}
\quad\text{au-dessus de}\quad
\vcenter{
\xymatrix{
X' \ar[d] \ar[r] & X \ar[d] \\
Y' \ar[r] & Y
}
}
$$
Choisissons $u \in |U|$ d'image $x \in |X|$.
La propriété d'être « localement quasi-fini » est locale pour la topologie étale sur
la source et le but ; voir
Descente, remarque \ref{descent-remark-list-local-source-target}.
Les lemmes \ref{spaces-morphisms-lemma-local-source-target-at-point} et
\ref{spaces-morphisms-lemma-local-source-target-global-implies-local} s'appliquent donc, et
on voit que $f : X \to Y$ est quasi-fini en $x$
si et seulement si $U \to V$ est quasi-fini en $u$.
Il en va de même de $f' : X' \to Y'$ et du morphisme $U' \to V'$.
Ainsi, les assertions (1), (2) et (3) se ramènent à
Morphismes, lemmes \ref{morphisms-lemma-base-change-quasi-finite} et
\ref{morphisms-lemma-quasi-finite-points-open}.
L'assertion (4) résulte de (3) et du lemme \ref{spaces-morphisms-lemma-base-change-quasi-compact}.
\end{proof}

\begin{lemma}
\label{spaces-morphisms-lemma-composition-quasi-finite}
Le composé de morphismes quasi-finis est quasi-fini.
Il en va de même pour « localement quasi-fini ».
\end{lemma}

\begin{proof}
Voir la remarque \ref{spaces-morphisms-remark-composition-P} et
Morphismes, lemme \ref{morphisms-lemma-composition-quasi-finite}.
\end{proof}

\begin{lemma}
\label{spaces-morphisms-lemma-base-change-quasi-finite}
Un changement de base d'un morphisme quasi-fini est quasi-fini.
Il en va de même pour « localement quasi-fini ».
\end{lemma}

\begin{proof}
Conséquence immédiate du lemme \ref{spaces-morphisms-lemma-base-change-quasi-finite-locus}.
\end{proof}

\noindent
Le lemme suivant caractérise les morphismes localement quasi-finis comme les
morphismes localement de type fini à « fibres discrètes ».
Cependant, cela ne revient pas à demander que $|X| \to |Y|$ ait
des fibres discrètes, comme le montre la discussion de
Exemples, section
\ref{examples-section-specializations-fibre-etale}.

\begin{lemma}
\label{spaces-morphisms-lemma-locally-quasi-finite}
Soit $S$ un schéma. Soit $f : X \to Y$ un morphisme d'espaces algébriques.
Supposons que $f$ soit localement de type fini. Les conditions suivantes sont équivalentes :
\begin{enumerate}
\item $f$ est localement quasi-fini ;
\item pour tout morphisme $\Spec(k) \to Y$, où $k$ est un corps,
l'espace $|X_k|$ est discret. Ici, $X_k = \Spec(k) \times_Y X$.
\end{enumerate}
\end{lemma}

\begin{proof}
Supposons que $f$ soit localement quasi-fini. Soit $\Spec(k) \to Y$ comme
en (2). Choisissons un morphisme étale surjectif $U \to X$, où $U$ est un schéma.
Alors $U_k = \Spec(k) \times_Y U \to X_k$ est un
morphisme étale d'espaces algébriques d'après
Propriétés des espaces, lemme \ref{spaces-properties-lemma-base-change-etale}.
D'après le
lemme \ref{spaces-morphisms-lemma-base-change-quasi-finite},
on voit que $X_k \to \Spec(k)$ est localement quasi-fini. Par
définition, cela signifie que $U_k \to \Spec(k)$ est localement
quasi-fini. Ainsi, $|U_k|$ est discret d'après
Morphismes, lemme \ref{morphisms-lemma-locally-quasi-finite-fibres}.
Comme $|U_k| \to |X_k|$ est surjective et ouverte, on en conclut que $|X_k|$
est discret.

\medskip\noindent
Réciproquement, supposons (2). Choisissons un morphisme étale surjectif $V \to Y$,
où $V$ est un schéma. Choisissons un morphisme étale surjectif
$U \to V \times_Y X$, où $U$ est un schéma.
Remarquons que $U \to V$ est localement de type fini puisque $f$ l'est.
On a le diagramme
$$
\xymatrix{
U \ar[r] \ar[rd] & X \times_Y V \ar[d] \ar[r] & V \ar[d] \\
& X \ar[r] & Y
}
$$
Si $f$ n'est pas localement quasi-fini, alors $U \to V$ ne l'est pas.
Il existe donc une spécialisation $u \leadsto u'$ pour certains $u, u' \in U$
au-dessus d'un même point $v \in V$ ; voir
Morphismes, lemme \ref{morphisms-lemma-quasi-finite-at-point-characterize}.
Nous affirmons que $u, u'$ n'ont pas la même image dans
$X_v = \Spec(\kappa(v)) \times_Y X$,
ce qui contredira, comme voulu, l'hypothèse que $|X_v|$ est discret.
Posons $d = \text{trdeg}_{\kappa(v)}(\kappa(u))$ et
$d' = \text{trdeg}_{\kappa(v)}(\kappa(u'))$.
On voit alors que $d > d'$ d'après
Morphismes, lemme \ref{morphisms-lemma-dimension-fibre-specialization}.
Remarquons que $U_v$ (la fibre de $U \to V$ au-dessus de $v$) est le produit fibré de
$U$ et de $X_v$ au-dessus de $X \times_Y V$ ; ainsi, $U_v \to X_v$ est étale (comme
changement de base du morphisme étale $U \to X \times_Y V$). Si
$u, u' \in U_v$ s'envoient sur le même élément de $|X_v|$, alors il existe un point
$r \in R_v = U_v \times_{X_v} U_v$ tel que $t(r) = u$ et $s(r) = u'$ ; voir
Propriétés des espaces, lemme \ref{spaces-properties-lemma-points-cartesian}.
Remarquons que $s, t : R_v \to U_v$ sont des morphismes étales de schémas
sur $\kappa(v)$ ; ainsi, les extensions $\kappa(u) \subset \kappa(r) \supset \kappa(u')$
sont finies séparables entre corps sur $\kappa(v)$ (voir
Morphismes, lemme \ref{morphisms-lemma-etale-over-field}).
On en conclut que les degrés de transcendance sont égaux.
Cette contradiction achève la démonstration.
\end{proof}

\begin{lemma}
\label{spaces-morphisms-lemma-quasi-finite-local}
Soit $S$ un schéma.
Soit $f : X \to Y$ un morphisme d'espaces algébriques sur $S$.
Les conditions suivantes sont équivalentes :
\begin{enumerate}
\item $f$ est localement quasi-fini ;
\item pour tout $x \in |X|$, le morphisme $f$ est quasi-fini en $x$ ;
\item pour tout schéma $Z$ et tout morphisme $Z \to Y$, le morphisme
$Z \times_Y X \to Z$ est localement quasi-fini ;
\item pour tout schéma affine $Z$ et tout morphisme
$Z \to Y$, le morphisme $Z \times_Y X \to Z$ est localement quasi-fini ;
\item il existe un schéma $V$ et un morphisme étale surjectif
$V \to Y$ tels que $V \times_Y X \to V$ soit localement quasi-fini ;
\item il existe un schéma $U$ et un morphisme étale surjectif
$\varphi : U \to X$ tels que le composé $f \circ \varphi$
soit localement quasi-fini ;
\item pour tout diagramme commutatif
$$
\xymatrix{
U \ar[d] \ar[r] & V \ar[d] \\
X \ar[r] & Y
}
$$
où $U$ et $V$ sont des schémas et où les flèches verticales sont étales,
la flèche horizontale supérieure est localement quasi-finie ;
\item il existe un diagramme commutatif
$$
\xymatrix{
U \ar[d] \ar[r] & V \ar[d] \\
X \ar[r] & Y
}
$$
où $U$ et $V$ sont des schémas, où les flèches verticales sont étales et où
$U \to X$ est surjectif, tel que la flèche horizontale supérieure soit
localement quasi-finie ;
\item il existe des recouvrements de Zariski $Y = \bigcup_{i \in I} Y_i$
et $f^{-1}(Y_i) = \bigcup X_{ij}$ tels que
chaque morphisme $X_{ij} \to Y_i$ soit localement quasi-fini.
\end{enumerate}
\end{lemma}

\begin{proof}
La démonstration est omise.
\end{proof}

\begin{lemma}
\label{spaces-morphisms-lemma-immersion-quasi-finite}
Une immersion est localement quasi-finie.
\end{lemma}

\begin{proof}
La démonstration est omise.
\end{proof}

\begin{lemma}
\label{spaces-morphisms-lemma-permanence-quasi-finite}
Soit $S$ un schéma.
Soient $X \to Y \to Z$ des morphismes d'espaces algébriques sur $S$.
Si $X \to Z$ est localement quasi-fini, alors $X \to Y$
est localement quasi-fini.
\end{lemma}

\begin{proof}
Choisissons un diagramme commutatif
$$
\xymatrix{
U \ar[d] \ar[r] & V \ar[d] \ar[r] & W \ar[d] \\
X \ar[r] & Y \ar[r] & Z
}
$$
dont les flèches verticales sont étales et surjectives. (Voir
Espaces, lemme \ref{spaces-lemma-lift-morphism-presentations}.)
Appliquons
Morphismes, lemme \ref{morphisms-lemma-permanence-quasi-finite}
à la ligne supérieure.
\end{proof}

\begin{lemma}
\label{spaces-morphisms-lemma-quasi-finite-at-a-finite-number-of-points}
Soit $S$ un schéma. Soit $f : X \to Y$ un morphisme de type fini
d'espaces algébriques sur $S$. Soit $y \in |Y|$.
Il n'y a qu'un nombre fini
de points de $|X|$ au-dessus de $y$ où $f$ est quasi-fini.
\end{lemma}

\begin{proof}
Choisissons un corps $k$ et un morphisme $\Spec(k) \to Y$ dans la classe
d'équivalence déterminée par $y$. La fibre $X_k = \Spec(k) \times_Y X$ est un
espace algébrique de type fini sur un corps, donc en particulier quasi-compact.
L'application $|X_k| \to |X|$ a pour image la fibre de $|X| \to |Y|$
au-dessus de $y$ (Propriétés des espaces, lemme
\ref{spaces-properties-lemma-points-cartesian}).
De plus, l'ensemble des points où $X_k \to \Spec(k)$ est
quasi-fini s'envoie sur l'ensemble des points au-dessus de $y$ où
$f$ est quasi-fini d'après le lemme \ref{spaces-morphisms-lemma-base-change-quasi-finite-locus}.
Choisissons un schéma affine $U$ et un morphisme étale surjectif $U \to X_k$
(Propriétés des espaces, lemme
\ref{spaces-properties-lemma-quasi-compact-affine-cover}).
Alors $U \to \Spec(k)$ est un morphisme de type fini et il n'y a qu'un
nombre fini de points où ce morphisme est quasi-fini ;
voir Morphismes, lemme
\ref{morphisms-lemma-quasi-finite-at-a-finite-number-of-points}.
Comme $X_k \to \Spec(k)$ est quasi-fini en un point $x'$ si et seulement
si ce point est l'image d'un point de $U$ où $U \to \Spec(k)$ est
quasi-fini, on conclut.
\end{proof}

\begin{lemma}
\label{spaces-morphisms-lemma-monomorphism-loc-finite-type-loc-quasi-finite}
Soit $S$ un schéma.
Soit $f : X \to Y$ un morphisme d'espaces algébriques sur $S$.
Si $f$ est localement de type fini et est un monomorphisme, alors $f$
est séparé et localement quasi-fini.
\end{lemma}

\begin{proof}
Un monomorphisme est séparé ; voir le
lemme \ref{spaces-morphisms-lemma-monomorphism-separated}.
D'après le
lemme \ref{spaces-morphisms-lemma-quasi-finite-local},
il suffit de démontrer le lemme après avoir effectué un changement de base
par $Z \to Y$, avec $Z$ affine. On peut donc supposer que $Y$ est un
schéma affine. Choisissons un schéma affine $U$ et un morphisme étale
$U \to X$. Comme $X \to Y$ est localement de type fini, le morphisme
de schémas affines $U \to Y$ est de type fini.
Comme $X \to Y$ est un monomorphisme, on a $U \times_X U = U \times_Y U$.
En particulier, les applications $U \times_Y U \to U$ sont étales.
Soit $y \in Y$. Alors soit $U_y$ est vide, soit
$\Spec(\kappa(u)) \times_{\Spec(\kappa(y))} U_y$
est isomorphe à la fibre de $U \times_Y U \to U$ au-dessus de $u$ pour
un certain $u \in U$ au-dessus de $y$. Il s'ensuit que les fibres de
$U \to Y$ sont des ensembles discrets finis (car $U \times_Y U \to U$
est un morphisme étale de schémas affines ; voir
Morphismes, lemme \ref{morphisms-lemma-etale-over-field}).
Ainsi, $U \to Y$ est quasi-fini ; voir
Morphismes, lemme \ref{morphisms-lemma-quasi-finite-at-point-characterize}.
Comme $U \to X$ était un morphisme étale quelconque, avec $U$ affine,
il en résulte que $X \to Y$ est localement quasi-fini.
\end{proof}













\section{Morphismes de présentation finie}
\label{spaces-morphisms-section-finite-presentation}

\noindent
La propriété « localement de présentation finie » des morphismes de schémas est
locale sur la source et le but pour la topologie étale ; voir
Descente, remarque \ref{descent-remark-list-local-source-target}.
Elle est aussi stable par changement de base et locale sur le but pour la topologie fpqc ; voir
Morphismes, lemme \ref{morphisms-lemma-base-change-finite-presentation}, et
Descente,
lemme \ref{descent-lemma-descending-property-locally-finite-presentation}.
Ainsi, d'après le
lemme \ref{spaces-morphisms-lemma-local-source-target}
ci-dessus, on peut définir comme suit ce que signifie, pour un morphisme
d'espaces algébriques, d'être localement de présentation
finie ; cette définition coïncide avec la notion déjà définie dans la
section \ref{spaces-morphisms-section-representable}
lorsque le morphisme est représentable.

\begin{definition}
\label{spaces-morphisms-definition-locally-finite-presentation}
Soit $S$ un schéma.
Soit $f : X \to Y$ un morphisme d'espaces algébriques sur $S$.
\begin{enumerate}
\item On dit que $f$ est {\it localement de présentation finie} si
les conditions équivalentes du
lemme \ref{spaces-morphisms-lemma-local-source-target}
sont satisfaites pour $\mathcal{P} =$ « localement de présentation finie ».
\item Soit $x \in |X|$. On dit que $f$ est de {\it présentation finie en $x$}
s'il existe un voisinage ouvert $X' \subset X$ de $x$ tel
que $f|_{X'} : X' \to Y$ soit localement de
présentation finie\footnote{Il semblerait maladroit de dire « localement de présentation finie
en $x$ », mais la terminologie retenue peut prêter à confusion : en effet,
« de présentation finie en $x$ » {\bf ne signifie pas} qu'il existe
un voisinage ouvert $X' \subset X$ tel que $f|_{X'}$ soit de présentation
finie.}.
\item Un morphisme d'espaces algébriques $f : X \to Y$ est
{\it de présentation finie}
s'il est localement de présentation finie, quasi-compact et
quasi-séparé.
\end{enumerate}
\end{definition}

\noindent
Remarquons qu'un morphisme de présentation finie n'est {\bf pas} simplement un morphisme
quasi-compact et localement de présentation finie.

\begin{lemma}
\label{spaces-morphisms-lemma-composition-finite-presentation}
Le composé de morphismes de présentation finie est de présentation finie.
Il en va de même pour « localement de présentation finie ».
\end{lemma}

\begin{proof}
Voir la remarque \ref{spaces-morphisms-remark-composition-P} et
Morphismes, lemme \ref{morphisms-lemma-composition-finite-presentation}.
On utilise aussi le résultat pour les morphismes quasi-compacts et quasi-séparés
(lemmes \ref{spaces-morphisms-lemma-composition-quasi-compact} et
\ref{spaces-morphisms-lemma-composition-separated}).
\end{proof}

\begin{lemma}
\label{spaces-morphisms-lemma-base-change-finite-presentation}
Un changement de base d'un morphisme de présentation finie est de présentation finie.
Il en va de même pour « localement de présentation finie ».
\end{lemma}

\begin{proof}
Voir la remarque \ref{spaces-morphisms-remark-base-change-P} et
Morphismes, lemme \ref{morphisms-lemma-base-change-finite-presentation}.
On utilise aussi le résultat pour les morphismes quasi-compacts et quasi-séparés
(lemmes \ref{spaces-morphisms-lemma-base-change-quasi-compact} et
\ref{spaces-morphisms-lemma-base-change-separated}).
\end{proof}

\begin{lemma}
\label{spaces-morphisms-lemma-finite-presentation-local}
Soit $S$ un schéma.
Soit $f : X \to Y$ un morphisme d'espaces algébriques sur $S$.
Les conditions suivantes sont équivalentes :
\begin{enumerate}
\item $f$ est localement de présentation finie ;
\item pour tout $x \in |X|$, le morphisme $f$ est de présentation finie en $x$ ;
\item pour tout schéma $Z$ et tout morphisme $Z \to Y$, le morphisme
$Z \times_Y X \to Z$ est localement de présentation finie ;
\item pour tout schéma affine $Z$ et tout morphisme
$Z \to Y$, le morphisme $Z \times_Y X \to Z$ est localement de présentation finie ;
\item il existe un schéma $V$ et un morphisme étale surjectif
$V \to Y$ tels que $V \times_Y X \to V$ soit
localement de présentation finie ;
\item il existe un schéma $U$ et un morphisme étale surjectif
$\varphi : U \to X$ tels que le composé $f \circ \varphi$
soit localement de présentation finie ;
\item pour tout diagramme commutatif
$$
\xymatrix{
U \ar[d] \ar[r] & V \ar[d] \\
X \ar[r] & Y
}
$$
où $U$ et $V$ sont des schémas et où les flèches verticales sont étales,
la flèche horizontale supérieure est localement de présentation finie ;
\item il existe un diagramme commutatif
$$
\xymatrix{
U \ar[d] \ar[r] & V \ar[d] \\
X \ar[r] & Y
}
$$
où $U$ et $V$ sont des schémas, où les flèches verticales sont étales et où
$U \to X$ est surjectif, tel que la flèche horizontale supérieure soit
localement de présentation finie ;
\item il existe des recouvrements de Zariski $Y = \bigcup_{i \in I} Y_i$
et $f^{-1}(Y_i) = \bigcup X_{ij}$ tels que
chaque morphisme $X_{ij} \to Y_i$ soit localement de présentation finie.
\end{enumerate}
\end{lemma}

\begin{proof}
La démonstration est omise.
\end{proof}

\begin{lemma}
\label{spaces-morphisms-lemma-finite-presentation-finite-type}
Un morphisme localement de présentation finie est localement de type fini.
Un morphisme de présentation finie est de type fini.
\end{lemma}

\begin{proof}
Soit $f : X \to Y$ un morphisme d'espaces algébriques localement de
présentation finie. Cela signifie qu'il existe un diagramme comme dans le
lemme \ref{spaces-morphisms-lemma-local-source-target},
où $h$ est localement de présentation finie et où $a$ est une flèche verticale surjective. D'après
Morphismes, lemme \ref{morphisms-lemma-finite-presentation-finite-type}
$h$ est localement de type fini.
Ainsi, $X \to Y$ est localement de type fini par définition.
Si $f$ est de présentation finie, il est quasi-compact et
il en résulte que $f$ est de type fini.
\end{proof}

\begin{lemma}
\label{spaces-morphisms-lemma-finite-presentation-noetherian}
Soit $S$ un schéma.
Soit $f : X \to Y$ un morphisme d'espaces algébriques sur $S$.
Si $f$ est de présentation finie et si $Y$ est noethérien,
alors $X$ est noethérien.
\end{lemma}

\begin{proof}
Supposons que $f$ soit de présentation finie et que $Y$ soit noethérien. D'après les
lemmes \ref{spaces-morphisms-lemma-finite-presentation-finite-type} et
\ref{spaces-morphisms-lemma-locally-finite-type-locally-noetherian}
on voit que $X$ est localement noethérien. Comme $f$ est quasi-compact
et que $Y$ est quasi-compact, on voit que $X$ est quasi-compact.
Comme $f$ est de présentation finie, il est quasi-séparé (voir la
définition \ref{spaces-morphisms-definition-locally-finite-presentation}),
et comme $Y$ est noethérien, il est quasi-séparé (voir
Propriétés des espaces,
définition \ref{spaces-properties-definition-noetherian}).
Ainsi, $X$ est quasi-séparé d'après le
lemme \ref{spaces-morphisms-lemma-separated-over-separated}.
On a donc vérifié les trois conditions de
Propriétés des espaces,
définition \ref{spaces-properties-definition-noetherian}
et cela conclut.
\end{proof}

\begin{lemma}
\label{spaces-morphisms-lemma-noetherian-finite-type-finite-presentation}
Soit $S$ un schéma.
Soit $f : X \to Y$ un morphisme d'espaces algébriques sur $S$.
\begin{enumerate}
\item Si $Y$ est localement noethérien et si $f$ est localement de type fini,
alors $f$ est localement de présentation finie.
\item Si $Y$ est localement noethérien et si $f$ est de type fini et quasi-séparé,
alors $f$ est de présentation finie.
\end{enumerate}
\end{lemma}

\begin{proof}
Supposons que $f : X \to Y$ soit localement de type fini et que $Y$ soit localement noethérien.
Cela signifie qu'il existe un diagramme comme dans le
lemme \ref{spaces-morphisms-lemma-local-source-target},
où $h$ est localement de type fini et où $a$ est une flèche verticale surjective. D'après
Morphismes, lemme
\ref{morphisms-lemma-noetherian-finite-type-finite-presentation}
$h$ est localement de présentation finie.
Ainsi, $X \to Y$ est localement de présentation finie par définition.
Ceci démontre (1).
Si $f$ est de type fini et quasi-séparé, alors il est aussi
quasi-compact et quasi-séparé, et (2) en résulte immédiatement.
\end{proof}

\begin{lemma}
\label{spaces-morphisms-lemma-finite-presentation-quasi-compact-quasi-separated}
Soit $S$ un schéma. Soit $Y$ un espace algébrique sur $S$ qui est
quasi-compact et quasi-séparé. Si $X$ est de présentation finie sur
$Y$, alors $X$ est quasi-compact et quasi-séparé.
\end{lemma}

\begin{proof}
La démonstration est omise.
\end{proof}

\begin{lemma}
\label{spaces-morphisms-lemma-finite-presentation-permanence}
Soit $S$ un schéma.
Soient $f : X \to Y$ et $Y \to Z$ des morphismes d'espaces algébriques sur $S$.
Si $X$ est localement de présentation finie sur $Z$ et si
$Y$ est localement de type fini sur $Z$, alors $f$ est localement
de présentation finie.
\end{lemma}

\begin{proof}
Choisissons un schéma $W$ et un morphisme étale surjectif $W \to Z$.
Choisissons ensuite un schéma $V$ et un morphisme étale surjectif $V \to W \times_Z Y$.
Choisissons enfin un schéma $U$ et un morphisme étale surjectif
$U \to V \times_Y X$. Par définition, $U$ est localement de présentation finie
sur $W$ et $V$ est localement de type fini sur $W$. D'après
Morphismes, lemme \ref{morphisms-lemma-finite-presentation-permanence}
le morphisme $U \to V$ est localement de présentation finie.
Ainsi, $f$ est localement de présentation finie.
\end{proof}

\begin{lemma}
\label{spaces-morphisms-lemma-diagonal-morphism-finite-type}
Soit $S$ un schéma. Soit $f : X \to Y$ un morphisme d'espaces algébriques
sur $S$, de diagonale $\Delta : X \to X \times_Y X$. Si $f$ est localement de
type fini, alors $\Delta$ est localement de présentation finie. Si $f$ est
quasi-séparé et localement de type fini, alors $\Delta$ est de
présentation finie.
\end{lemma}

\begin{proof}
Remarquons que $\Delta$ est un morphisme sur $X$ (par la seconde
projection $X \times_Y X \to X$). Supposons $f$ localement de type fini.
Remarquons que $X$ est de présentation finie sur $X$ et que $X \times_Y X$ est
de type fini sur $X$ (d'après le lemme \ref{spaces-morphisms-lemma-base-change-finite-type}).
La première assertion résulte donc du
lemme \ref{spaces-morphisms-lemma-finite-presentation-permanence}.
La seconde assertion résulte de la première, des définitions et
du fait qu'un morphisme diagonal est séparé
(lemme \ref{spaces-morphisms-lemma-properties-diagonal}).
\end{proof}

\begin{lemma}
\label{spaces-morphisms-lemma-open-immersion-locally-finite-presentation}
Une immersion ouverte d'espaces algébriques est localement de présentation finie.
\end{lemma}

\begin{proof}
Une immersion ouverte est représentable par définition ; on peut donc
utiliser le principe général
Espaces,
lemme \ref{spaces-lemma-representable-transformations-property-implication}
et
Morphismes,
lemme \ref{morphisms-lemma-open-immersion-locally-finite-presentation}.
\end{proof}

\begin{lemma}
\label{spaces-morphisms-lemma-closed-immersion-finite-presentation}
Une immersion fermée $i : Z \to X$ est de présentation finie si et seulement si
le faisceau quasi-cohérent d'idéaux associé
$\mathcal{I} = \Ker(\mathcal{O}_X \to i_*\mathcal{O}_Z)$
est de type fini (comme $\mathcal{O}_X$-module).
\end{lemma}

\begin{proof}
Soit $U$ un schéma et soit $U \to X$ un morphisme étale surjectif.
D'après le lemme \ref{spaces-morphisms-lemma-finite-presentation-local},
on voit que $i' : Z \times_X U \to U$ est de présentation finie si et
seulement si $i$ l'est. D'après Propriétés des espaces, section
\ref{spaces-properties-section-properties-modules}
on voit que $\mathcal{I}$ est de type fini si et seulement si
$\mathcal{I}|_U = \Ker(\mathcal{O}_U \to i'_*\mathcal{O}_{Z \times_X U})$
l'est. Le résultat découle donc du cas des schémas ; voir Morphismes,
lemme \ref{morphisms-lemma-closed-immersion-finite-presentation}.
\end{proof}







\section{Parties constructibles}
\label{spaces-morphisms-section-constructible}

\noindent
Cette section fait suite à
Propriétés des espaces, section \ref{spaces-properties-section-constructible}.

\begin{lemma}
\label{spaces-morphisms-lemma-inverse-image-constructible}
Soit $S$ un schéma.
Soit $f : X \to Y$ un morphisme d'espaces algébriques sur $S$.
Soit $E \subset |Y|$ une partie.
Si $E$ est localement constructible pour la topologie étale dans $Y$, alors
$f^{-1}(E)$ est localement constructible pour la topologie étale dans $X$.
\end{lemma}

\begin{proof}
Choisissons un schéma $V$ et un morphisme étale surjectif $\varphi : V \to Y$.
Choisissons un schéma $U$ et un morphisme étale surjectif
$U \to V \times_Y X$. Alors $U \to X$ est étale surjectif,
et l'image inverse de $f^{-1}(E)$ dans $U$ est l'image
inverse de $\varphi^{-1}(E)$ par $U \to V$. Le lemme découle donc
du cas des schémas pour $U \to V$
(Morphismes, lemme \ref{morphisms-lemma-inverse-image-constructible})
et de la définition (Propriétés des espaces, définition
\ref{spaces-properties-definition-locally-constructible}).
\end{proof}

\begin{theorem}[Théorème de Chevalley]
\label{spaces-morphisms-theorem-chevalley}
Soit $S$ un schéma.
Soit $f : X \to Y$ un morphisme d'espaces algébriques sur $S$.
Supposons que $f$ soit quasi-compact et localement de présentation finie.
Alors l'image de toute partie de $|X|$ localement constructible pour la topologie étale est
une partie de $|Y|$ localement constructible pour la topologie étale.
\end{theorem}

\begin{proof}
Soit $E \subset |X|$ localement constructible pour la topologie étale.
Soit $V \to Y$ un morphisme étale, avec $V$ affine.
Il suffit de montrer que l'image inverse de $f(E)$ dans
$V$ est constructible ; voir Propriétés des espaces, définition
\ref{spaces-properties-definition-locally-constructible}.
Comme $f$ est quasi-compact, $V \times_Y X$ est un espace algébrique quasi-compact.
Choisissons un schéma affine $U$ et un morphisme étale surjectif
$U \to V \times_Y X$ (Propriétés des espaces, lemme
\ref{spaces-properties-lemma-quasi-compact-affine-cover}).
D'après Propriétés des espaces, lemme \ref{spaces-properties-lemma-points-cartesian},
l'image inverse de $f(E)$ dans $V$ est l'image par $U \to V$
de l'image inverse de $E$ dans $U$.
Le résultat découle donc du cas des schémas ; voir
Morphismes, lemme \ref{morphisms-lemma-chevalley}.
\end{proof}







\section{Morphismes plats}
\label{spaces-morphisms-section-flat}

\noindent
La propriété « plat » des morphismes de schémas est
locale sur la source et le but pour la topologie étale ; voir
Descente, remarque \ref{descent-remark-list-local-source-target}.
Elle est aussi stable par changement de base et locale sur le but pour la topologie fpqc ; voir
Morphismes, lemme \ref{morphisms-lemma-base-change-flat} et
Descente, lemme \ref{descent-lemma-descending-property-flat}.
Ainsi, d'après le
lemme \ref{spaces-morphisms-lemma-local-source-target}
ci-dessus, on peut définir comme suit la notion de morphisme plat d'espaces
algébriques ; cette définition coïncide avec la notion déjà définie dans la
section \ref{spaces-morphisms-section-representable}
lorsque le morphisme est représentable.

\begin{definition}
\label{spaces-morphisms-definition-flat}
Soit $S$ un schéma.
Soit $f : X \to Y$ un morphisme d'espaces algébriques sur $S$.
\begin{enumerate}
\item On dit que $f$ est {\it plat} si les conditions équivalentes du
lemme \ref{spaces-morphisms-lemma-local-source-target} sont satisfaites pour
$\mathcal{P} =$ « plat ».
\item Soit $x \in |X|$. On dit que $f$ est {\it plat en $x$} si les
conditions équivalentes du
lemme \ref{spaces-morphisms-lemma-local-source-target-at-point}
sont satisfaites pour $\mathcal{Q} =$ « l'application induite sur les anneaux locaux est plate ».
\end{enumerate}
Remarquons que la seconde partie a un sens d'après
Descente, lemme \ref{descent-lemma-flat-at-point}.
\end{definition}

\noindent
Vérifions rapidement la cohérence de cette définition.

\begin{lemma}
\label{spaces-morphisms-lemma-flat-is-flat-at-all-points}
Soit $S$ un schéma. Soit $f : X \to Y$ un morphisme d'espaces algébriques
sur $S$. Alors $f$ est plat si et seulement si $f$ est plat en tout point de $|X|$.
\end{lemma}

\begin{proof}
Choisissons un diagramme commutatif
$$
\xymatrix{
U \ar[d]_a \ar[r]_h & V \ar[d]^b \\
X \ar[r]^f & Y
}
$$
où $U$ et $V$ sont des schémas, où les flèches verticales sont étales et où
$a$ est surjectif. Par définition, $f$ est plat si et seulement si $h$ est plat
(définition \ref{spaces-morphisms-definition-P}).
Par définition, $f$ est plat en $x \in |X|$ si et seulement si $h$ est plat
en un certain (ce qui équivaut à tout) $u \in U$ s'envoyant sur $x$
(définition \ref{spaces-morphisms-definition-P-at-point}).
Le lemme résulte donc du fait qu'un morphisme de schémas
est plat si et seulement s'il est plat en tout point de la source
(Morphismes, définition \ref{morphisms-definition-flat}).
\end{proof}

\begin{lemma}
\label{spaces-morphisms-lemma-composition-flat}
Le composé de morphismes plats est plat.
\end{lemma}

\begin{proof}
Voir la remarque \ref{spaces-morphisms-remark-composition-P} et
Morphismes, lemme \ref{morphisms-lemma-composition-flat}.
\end{proof}

\begin{lemma}
\label{spaces-morphisms-lemma-base-change-flat}
Le changement de base d'un morphisme plat est plat.
\end{lemma}

\begin{proof}
Voir la remarque \ref{spaces-morphisms-remark-base-change-P} et
Morphismes, lemme \ref{morphisms-lemma-base-change-flat}.
\end{proof}

\begin{lemma}
\label{spaces-morphisms-lemma-flat-local}
Soit $S$ un schéma.
Soit $f : X \to Y$ un morphisme d'espaces algébriques sur $S$.
Les conditions suivantes sont équivalentes :
\begin{enumerate}
\item $f$ est plat ;
\item pour tout $x \in |X|$, le morphisme $f$ est plat en $x$ ;
\item pour tout schéma $Z$ et tout morphisme $Z \to Y$, le morphisme
$Z \times_Y X \to Z$ est plat ;
\item pour tout schéma affine $Z$ et tout morphisme
$Z \to Y$, le morphisme $Z \times_Y X \to Z$ est plat ;
\item il existe un schéma $V$ et un morphisme étale surjectif
$V \to Y$ tels que $V \times_Y X \to V$ soit plat ;
\item il existe un schéma $U$ et un morphisme étale surjectif
$\varphi : U \to X$ tels que le composé $f \circ \varphi$
soit plat ;
\item pour tout diagramme commutatif
$$
\xymatrix{
U \ar[d] \ar[r] & V \ar[d] \\
X \ar[r] & Y
}
$$
où $U$ et $V$ sont des schémas et où les flèches verticales sont étales,
la flèche horizontale supérieure est plate ;
\item il existe un diagramme commutatif
$$
\xymatrix{
U \ar[d] \ar[r] & V \ar[d] \\
X \ar[r] & Y
}
$$
où $U$ et $V$ sont des schémas, où les flèches verticales sont étales et où
$U \to X$ est surjectif, tel que la flèche horizontale supérieure soit plate ;
\item il existe des recouvrements de Zariski $Y = \bigcup Y_i$ et
$f^{-1}(Y_i) = \bigcup X_{ij}$ tels que
chaque morphisme $X_{ij} \to Y_i$ soit plat.
\end{enumerate}
\end{lemma}

\begin{proof}
La démonstration est omise.
\end{proof}

\begin{lemma}
\label{spaces-morphisms-lemma-fppf-open}
Un morphisme plat localement de présentation finie est universellement ouvert.
\end{lemma}

\begin{proof}
Soit $f : X \to Y$ un morphisme plat et localement de présentation finie
d'espaces algébriques sur $S$. Choisissons un diagramme
$$
\xymatrix{
U \ar[r]_\alpha \ar[d] & V \ar[d] \\
X \ar[r] & Y
}
$$
où $U$ et $V$ sont des schémas et où les flèches verticales sont surjectives et
étales ; voir
Espaces, lemme \ref{spaces-lemma-lift-morphism-presentations}.
D'après les
lemmes \ref{spaces-morphisms-lemma-flat-local} et \ref{spaces-morphisms-lemma-finite-presentation-local},
le morphisme $\alpha$ est plat et localement de présentation finie.
Ainsi, d'après
Morphismes, lemme \ref{morphisms-lemma-fppf-open}
on voit que $\alpha$ est universellement ouvert.
Ainsi, $X \to Y$ est universellement ouvert d'après le
lemme \ref{spaces-morphisms-lemma-universally-open-local}.
\end{proof}

\begin{lemma}
\label{spaces-morphisms-lemma-fpqc-quotient-topology}
Soit $S$ un schéma.
Soit $f : X \to Y$ un morphisme plat, quasi-compact et surjectif
d'espaces algébriques sur $S$.
Une partie $T \subset |Y|$ est ouverte (resp.\ fermée) si et seulement si
$f^{-1}(T)$ est ouverte (resp.\ fermée) dans $|X|$.
Autrement dit, $f$ est submersif, et même universellement submersif.
\end{lemma}

\begin{proof}
Choisissons des schémas affines $V_i$ et des morphismes étales $V_i \to Y$ tels que
$g : V = \coprod V_i \to Y$ soit surjectif ; voir
Propriétés des espaces,
lemme \ref{spaces-properties-lemma-cover-by-union-affines}.
Pour tout $i$, l'espace algébrique $V_i \times_Y X$ est quasi-compact.
On peut donc trouver un schéma affine $U_i$ et un morphisme étale surjectif
$U_i \to V_i \times_Y X$ ; voir
Propriétés des espaces,
lemme \ref{spaces-properties-lemma-quasi-compact-affine-cover}.
Alors le composé $U_i \to V_i \times_Y X \to V_i$ est un morphisme surjectif
et plat entre schémas affines.
Bien sûr, $U = \coprod U_i \to X$ est alors étale surjectif,
et $U \to V \times_Y X$ est surjectif. De plus, le morphisme $U \to V$ est la
somme disjointe des morphismes $U_i \to V_i$. Ainsi, $U \to V$ est surjectif,
quasi-compact et plat. Considérons le diagramme
$$
\xymatrix{
U \ar[r] \ar[d] & X \ar[d] \\
V \ar[r] & Y
}
$$
Par définition de la topologie sur $|Y|$, la partie $T$ est fermée
(resp.\ ouverte) si et seulement si $g^{-1}(T) \subset |V|$ est fermée
(resp.\ ouverte). Il en va de même de
$f^{-1}(T)$ et de son image inverse dans $|U|$.
Comme $U \to V$ est quasi-compact, surjectif et plat, on conclut par
Morphismes, lemme \ref{morphisms-lemma-fpqc-quotient-topology}.
\end{proof}

\begin{lemma}
\label{spaces-morphisms-lemma-flat-at-point-etale-local-rings}
Soit $S$ un schéma.
Soit $f : X \to Y$ un morphisme d'espaces algébriques sur $S$.
Soit $\overline{x}$ un point géométrique de $X$ au-dessus du point
$x \in |X|$. Posons $\overline{y} = f \circ \overline{x}$. Les conditions suivantes
sont équivalentes :
\begin{enumerate}
\item $f$ est plat en $x$ ;
\item l'application entre les anneaux locaux étales
$\mathcal{O}_{Y, \overline{y}} \to \mathcal{O}_{X, \overline{x}}$
est plate.
\end{enumerate}
\end{lemma}

\begin{proof}
Choisissons un diagramme commutatif
$$
\xymatrix{
U \ar[d]_a \ar[r]_h & V \ar[d]^b \\
X \ar[r]^f & Y
}
$$
où $U$ et $V$ sont des schémas, où $a, b$ sont étales, et où
$u \in U$ s'envoie sur $x$. On peut trouver un point géométrique
$\overline{u} : \Spec(k) \to U$ au-dessus de $u$ tel que
$\overline{x} = a \circ \overline{u}$ ; voir
Propriétés des espaces, lemme
\ref{spaces-properties-lemma-geometric-lift-to-usual}.
Posons $\overline{v} = h \circ \overline{u}$, d'image $v \in V$.
On sait que
$$
\mathcal{O}_{X, \overline{x}} = \mathcal{O}_{U, u}^{sh}
\quad\text{et}\quad
\mathcal{O}_{Y, \overline{y}} = \mathcal{O}_{V, v}^{sh}
$$
voir
Propriétés des espaces, lemme
\ref{spaces-properties-lemma-describe-etale-local-ring}.
On obtient un diagramme commutatif
$$
\xymatrix{
\mathcal{O}_{U, u} \ar[r] &
\mathcal{O}_{X, \overline{x}} \\
\mathcal{O}_{V, v} \ar[u] \ar[r] &
\mathcal{O}_{Y, \overline{y}} \ar[u]
}
$$
formé d'anneaux locaux et dont les flèches horizontales sont plates. Il faut montrer que la
flèche verticale gauche est plate si et seulement si la flèche verticale droite l'est.
Algèbre, lemme \ref{algebra-lemma-flatness-descends-more-general}
nous dit que $\mathcal{O}_{U, u}$ est plat sur $\mathcal{O}_{V, v}$
si et seulement si $\mathcal{O}_{X, \overline{x}}$ est plat sur
$\mathcal{O}_{V, v}$. Le résultat découle donc de
Compléments sur la platitude, lemme \ref{flat-lemma-flat-up-down-henselization}.
\end{proof}

\begin{lemma}
\label{spaces-morphisms-lemma-flat-morphism-sites}
Soit $S$ un schéma.
Soit $f : X \to Y$ un morphisme d'espaces algébriques sur $S$.
Alors $f$ est plat si et seulement si le morphisme de sites
$
(f_{small}, f^\sharp) :
(X_\etale, \mathcal{O}_X)
\to
(Y_\etale, \mathcal{O}_Y)
$
associé à $f$ est plat.
\end{lemma}

\begin{proof}
La platitude de $(f_{small}, f^\sharp)$ est définie au moyen de la
platitude de $\mathcal{O}_X$ comme $f_{small}^{-1}\mathcal{O}_Y$-module.
Elle se vérifie sur les fibres ; voir
Modules sur les sites, lemme \ref{sites-modules-lemma-check-flat-stalks}
et
Propriétés des espaces, théorème \ref{spaces-properties-theorem-exactness-stalks}.
Or on a déjà vu que la platitude de $f$ se vérifie sur les fibres ; voir le
lemme \ref{spaces-morphisms-lemma-flat-at-point-etale-local-rings}.
\end{proof}

\begin{lemma}
\label{spaces-morphisms-lemma-flat-pullback-support}
Soit $S$ un schéma. Soit $f : Y \to X$ un morphisme d'espaces algébriques
sur $S$. Soit $\mathcal{F}$ un $\mathcal{O}_X$-module quasi-cohérent
de type fini de support schématique $Z \subset X$.
Si $f$ est plat, alors $f^{-1}(Z)$ est le support schématique de
$f^*\mathcal{F}$.
\end{lemma}

\begin{proof}
En utilisant la caractérisation du support schématique
donnée dans le lemme \ref{spaces-morphisms-lemma-scheme-theoretic-support}
et la caractérisation des morphismes plats au moyen de
recouvrements étales donnée dans le lemme \ref{spaces-morphisms-lemma-flat-local},
on se ramène au cas des schémas, qui est
Morphismes, lemme \ref{morphisms-lemma-flat-pullback-support}.
\end{proof}

\begin{lemma}
\label{spaces-morphisms-lemma-flat-morphism-scheme-theoretically-dense-open}
Soit $S$ un schéma.
Soit $f : X \to Y$ un morphisme plat d'espaces algébriques sur $S$.
Soit $V \to Y$ une immersion ouverte quasi-compacte. Si $V$
est schématiquement dense dans $Y$, alors $f^{-1}V$
est schématiquement dense dans $X$.
\end{lemma}

\begin{proof}
En utilisant la caractérisation des ouverts schématiquement denses
du lemme \ref{spaces-morphisms-lemma-scheme-theoretically-dense}
et la caractérisation des morphismes plats au moyen de
recouvrements étales donnée dans le lemme \ref{spaces-morphisms-lemma-flat-local},
on se ramène au cas des schémas, qui est
Morphismes, lemme
\ref{morphisms-lemma-flat-morphism-scheme-theoretically-dense-open}.
\end{proof}

\begin{lemma}
\label{spaces-morphisms-lemma-flat-base-change-scheme-theoretic-image}
Soit $S$ un schéma. Soit $f : X \to Y$ un morphisme plat d'espaces algébriques
sur $S$. Soit $g : V \to Y$ un morphisme quasi-compact d'espaces algébriques.
Soit $Z \subset Y$ l'image schématique de $g$, et soit $Z' \subset X$
l'image schématique du changement de base $V \times_Y X \to X$.
Alors $Z' = f^{-1}Z$.
\end{lemma}

\begin{proof}
Soit $Y' \to Y$ un morphisme étale surjectif tel que $Y'$ soit une
somme disjointe de schémas affines (Propriétés des espaces,
lemme \ref{spaces-properties-lemma-cover-by-union-affines}).
Soit $X' \to X \times_Y Y'$ un morphisme étale surjectif tel
que $X'$ soit une somme disjointe de schémas affines.
D'après le lemme \ref{spaces-morphisms-lemma-flat-local}, le morphisme $X' \to Y'$ est plat.
Posons $V' = V \times_Y Y'$.
D'après le lemme \ref{spaces-morphisms-lemma-quasi-compact-scheme-theoretic-image},
l'image inverse de $Z$ dans $Y'$ est l'image schématique
de $V' \to Y'$, et l'image inverse de $Z'$ dans $X'$
est l'image schématique de $V' \times_{Y'} X' \to X'$.
Comme $X' \to X$ est étale surjectif, il suffit de démontrer
le résultat pour les morphismes $X' \to Y'$ et $V' \to Y'$.
On peut donc supposer que $X$ et $Y$ sont des schémas affines.
Dans ce cas, $V$ est un espace algébrique quasi-compact.
Choisissons un schéma affine $W$ et un morphisme étale surjectif
$W \to V$ (Propriétés des espaces, lemme
\ref{spaces-properties-lemma-quasi-compact-affine-cover}).
Il est clair que l'image schématique de $V \to Y$
coïncide avec l'image schématique de $W \to Y$, et
de même pour $V \times_Y X \to X$ et $W \times_Y X \to X$.
On se ramène ainsi au cas des schémas, qui est
Morphismes, lemme
\ref{morphisms-lemma-flat-base-change-scheme-theoretic-image}.
\end{proof}





\section{Modules plats}
\label{spaces-morphisms-section-flat-modules}

\noindent
Dans cette section, nous définissons ce que signifie, pour un module, être plat
en un point. Nous utiliserons pour cela la notion de fibre d'un faisceau sur
le petit site étale $X_\etale$ d'un espace algébrique ; voir
Propriétés des espaces, définition \ref{spaces-properties-definition-stalk}.

\begin{lemma}
\label{spaces-morphisms-lemma-flat-at-point}
Soit $S$ un schéma. Soit $f : X \to Y$ un morphisme d'espaces
algébriques sur $S$. Soit $\mathcal{F}$ un faisceau quasi-cohérent sur $X$.
Soit $x \in |X|$. Les conditions suivantes sont équivalentes :
\begin{enumerate}
\item il existe un diagramme commutatif
$$
\xymatrix{
U \ar[d]_a \ar[r]_h & V \ar[d]^b \\
X \ar[r]^f & Y
}
$$
où $U$ et $V$ sont des schémas, où $a, b$ sont étales et où
$u \in U$ s'envoie sur $x$, tel que le module $a^*\mathcal{F}$ soit plat en $u$ sur $V$ ;
\item la fibre $\mathcal{F}_{\overline{x}}$ est plate sur
l'anneau local étale $\mathcal{O}_{Y, \overline{y}}$,
où $\overline{x}$ est un point géométrique quelconque au-dessus de
$x$ et $\overline{y} = f \circ \overline{x}$.
\end{enumerate}
\end{lemma}

\begin{proof}
Au cours de cette démonstration, fixons un point géométrique
$\overline{x} : \Spec(k) \to X$ au-dessus de $x$ et
notons $\overline{y} = f \circ \overline{x}$ son image dans $Y$.
Étant donné un diagramme comme en (1), on peut trouver un point géométrique
$\overline{u} : \Spec(k) \to U$ au-dessus de $u$ tel que
$\overline{x} = a \circ \overline{u}$ ; voir
Propriétés des espaces, lemme
\ref{spaces-properties-lemma-geometric-lift-to-usual}.
Posons $\overline{v} = h \circ \overline{u}$, d'image $v \in V$.
On sait que
$$
\mathcal{O}_{X, \overline{x}} = \mathcal{O}_{U, u}^{sh}
\quad\text{et}\quad
\mathcal{O}_{Y, \overline{y}} = \mathcal{O}_{V, v}^{sh}
$$
voir
Propriétés des espaces, lemme
\ref{spaces-properties-lemma-describe-etale-local-ring}.
On obtient un diagramme commutatif
$$
\xymatrix{
\mathcal{O}_{U, u} \ar[r] &
\mathcal{O}_{X, \overline{x}} \\
\mathcal{O}_{V, v} \ar[u] \ar[r] &
\mathcal{O}_{Y, \overline{y}} \ar[u]
}
$$
d'anneaux locaux. Enfin, on a
$$
\mathcal{F}_{\overline{x}} =
(a^*\mathcal{F})_u \otimes_{\mathcal{O}_{U, u}}
\mathcal{O}_{X, \overline{x}}
$$
d'après
Propriétés des espaces, lemme \ref{spaces-properties-lemma-stalk-quasi-coherent}.
Ainsi,
Algèbre, lemme \ref{algebra-lemma-flatness-descends-more-general}
nous dit que $(a^*\mathcal{F})_u$ est plat sur $\mathcal{O}_{V, v}$
si et seulement si $\mathcal{F}_{\overline{x}}$ est plate sur $\mathcal{O}_{V, v}$.
Le résultat découle donc de
Compléments sur la platitude, lemme \ref{flat-lemma-flat-up-down-henselization}.
\end{proof}

\begin{definition}
\label{spaces-morphisms-definition-flat-module}
Soit $S$ un schéma.
Soit $f : X \to Y$ un morphisme d'espaces algébriques sur $S$.
Soit $\mathcal{F}$ un faisceau quasi-cohérent sur $X$.
\begin{enumerate}
\item Soit $x \in |X|$. On dit que $\mathcal{F}$ est {\it plate en $x$ sur $Y$}
si les conditions équivalentes du
lemme \ref{spaces-morphisms-lemma-flat-at-point}
sont satisfaites.
\item On dit que $\mathcal{F}$ est {\it plate sur $Y$} si $\mathcal{F}$ est
plate sur $Y$ en tout $x \in |X|$.
\end{enumerate}
\end{definition}

\noindent
Cette définition posée, on dispose de l'inévitable lemme de changement de base.
Ce lemme entraîne que la formation du lieu de platitude d'un faisceau
quasi-cohérent commute au changement de base plat.

\begin{lemma}
\label{spaces-morphisms-lemma-base-change-module-flat}
Soit $S$ un schéma. Soit
$$
\xymatrix{
X' \ar[d]_{f'} \ar[r]_{g'} & X \ar[d]^f \\
Y' \ar[r]^g & Y
}
$$
un diagramme cartésien d'espaces algébriques sur $S$. Soit $x' \in |X'|$
d'image $x \in |X|$. Soit $\mathcal{F}$ un faisceau quasi-cohérent
sur $X$, et posons $\mathcal{F}' = (g')^*\mathcal{F}$.
\begin{enumerate}
\item Si $\mathcal{F}$ est plate en $x$ sur $Y$,
alors $\mathcal{F}'$ est plate en $x'$ sur $Y'$.
\item Si $g$ est plat en $f'(x')$ et si
$\mathcal{F}'$ est plate en $x'$ sur $Y'$, alors
$\mathcal{F}$ est plate en $x$ sur $Y$.
\end{enumerate}
En particulier, si $\mathcal{F}$ est plate sur $Y$, alors
$\mathcal{F}'$ est plate sur $Y'$.
\end{lemma}

\begin{proof}
Choisissons un schéma $V$ et un morphisme étale surjectif $V \to Y$.
Choisissons un schéma $U$ et un morphisme étale surjectif $U \to V \times_Y X$.
Choisissons un schéma $V'$ et un morphisme étale surjectif $V' \to V \times_Y Y'$.
Alors $U' = V' \times_V U$ est un schéma muni d'un morphisme étale
surjectif $U' = V' \times_V U \to Y' \times_Y X = X'$. Choisissons $u' \in U'$
s'envoyant sur $x' \in |X'|$. On peut alors vérifier la platitude de
$\mathcal{F}'$ en $x'$ sur $Y'$ au moyen de la platitude de
$\mathcal{F}'|_{U'}$ en $u'$ sur $V'$. Le lemme découle donc de
Compléments sur les morphismes, lemme \ref{more-morphisms-lemma-flat-locus-base-change}.
\end{proof}

\noindent
Le lemme suivant traite de la « composition » des morphismes plats en
termes de modules. Il montre aussi que la platitude satisfait une sorte de
descente du haut vers le bas.

\begin{lemma}
\label{spaces-morphisms-lemma-composition-module-flat}
Soit $S$ un schéma. Soient $X \to Y \to Z$ des morphismes d'espaces algébriques
sur $S$. Soit $\mathcal{F}$ un faisceau quasi-cohérent sur $X$.
Soit $x \in |X|$, d'image $y \in |Y|$.
\begin{enumerate}
\item Si $\mathcal{F}$ est plate en $x$ sur $Y$ et si
$Y$ est plat en $y$ sur $Z$, alors $\mathcal{F}$ est plate en
$x$ sur $Z$.
\item Soit $x : \Spec(K) \to X$ un représentant de $x$. Si
\begin{enumerate}
\item $\mathcal{F}$ est plate en $x$ sur $Y$ ;
\item $x^*\mathcal{F} \not = 0$ ;
\item $\mathcal{F}$ est plate en $x$ sur $Z$,
\end{enumerate}
alors $Y$ est plat en $y$ sur $Z$.
\item Soit $\overline{x}$ un point géométrique de $X$ au-dessus de $x$,
d'image $\overline{y}$ dans $Y$. Si $\mathcal{F}_{\overline{x}}$ est un
$\mathcal{O}_{Y, \overline{y}}$-module fidèlement plat et si
$\mathcal{F}$ est plate en $x$ sur $Z$, alors
$Y$ est plat en $y$ sur $Z$.
\end{enumerate}
\end{lemma}

\begin{proof}
Choisissons $\overline{x}$ et $\overline{y}$ comme dans (3), et notons
$\overline{z}$ le point géométrique de $Z$ ainsi induit. Au moyen de la
caractérisation de la platitude donnée dans les
lemmes \ref{spaces-morphisms-lemma-flat-at-point} et
\ref{spaces-morphisms-lemma-flat-at-point-etale-local-rings}
le lemme se ramène à une question purement algébrique sur l'homomorphisme
d'anneaux locaux $\mathcal{O}_{Z, \overline{z}} \to \mathcal{O}_{Y, \overline{y}}$
et sur le module $\mathcal{F}_{\overline{x}}$.
La partie (1) découle de
Algèbre, lemme \ref{algebra-lemma-composition-flat}.
Remarquons que la condition (2)(b) garantit que
$\mathcal{F}_{\overline{x}}/
\mathfrak m_{\overline{y}} \mathcal{F}_{\overline{x}}$
est non nul. Ainsi, (2)(a) $+$ (2)(b) entraînent que $\mathcal{F}_{\overline{x}}$
est un $\mathcal{O}_{Y, \overline{y}}$-module fidèlement plat ; voir
Algèbre, lemme \ref{algebra-lemma-ff}.
Ainsi, (2) est un cas particulier de (3).
Enfin, (3) découle de
Algèbre, lemme \ref{algebra-lemma-flat-permanence}.
\end{proof}

\noindent
Il arrive que le changement de base se fasse « en haut ». Voici un énoncé précis.

\begin{lemma}
\label{spaces-morphisms-lemma-flat-permanence}
Soit $S$ un schéma. Soient $f : X \to Y$, $g : Y \to Z$ des morphismes
d'espaces algébriques sur $S$. Soit $\mathcal{G}$ un faisceau quasi-cohérent sur $Y$.
Soit $x \in |X|$, d'image $y \in |Y|$.
Si $f$ est plat en $x$, alors
$$
\mathcal{G}\text{ est plate sur }Z\text{ en }y
\Leftrightarrow
f^*\mathcal{G}\text{ est plate sur }Z\text{ en }x.
$$
En particulier, si $f$ est surjectif et plat, alors
$\mathcal{G}$ est plate sur $Z$ si et seulement si
$f^*\mathcal{G}$ est plate sur $Z$.
\end{lemma}

\begin{proof}
Choisissons un point géométrique $\overline{x}$ de $X$, et notons
$\overline{y}$ son image dans $Y$ et $\overline{z}$ son image dans $Z$.
Au moyen de la caractérisation de la platitude donnée dans les
lemmes \ref{spaces-morphisms-lemma-flat-at-point} et
\ref{spaces-morphisms-lemma-flat-at-point-etale-local-rings}
et de la description de la fibre de $f^*\mathcal{G}$ en $\overline{x}$ dans
Propriétés des espaces,
lemme \ref{spaces-properties-lemma-stalk-pullback-quasi-coherent}
le lemme se ramène à une question purement algébrique sur les homomorphismes
d'anneaux locaux
$\mathcal{O}_{Z, \overline{z}} \to \mathcal{O}_{Y, \overline{y}}
\to \mathcal{O}_{X, \overline{x}}$
et sur le module $\mathcal{G}_{\overline{y}}$.
Cet énoncé algébrique est
Algèbre, lemme \ref{algebra-lemma-flatness-descends-more-general}.
\end{proof}

\begin{lemma}
\label{spaces-morphisms-lemma-pf-flat-module-open}
Soit $S$ un schéma.
Soit $f : X \to Y$ un morphisme d'espaces algébriques sur $S$.
Soit $\mathcal{F}$ un $\mathcal{O}_X$-module quasi-cohérent.
Supposons que $f$ soit localement de présentation finie, que $\mathcal{F}$ soit de
type fini, que $X = \text{Supp}(\mathcal{F})$ et que
$\mathcal{F}$ soit plate sur $Y$. Alors $f$ est universellement ouvert.
\end{lemma}

\begin{proof}
Choisissons un morphisme étale surjectif $\varphi : V \to Y$, où $V$ est un schéma.
Choisissons un morphisme étale surjectif $U \to V \times_Y X$, où $U$ est un schéma.
Il suffit alors de démontrer le lemme pour $U \to V$ et le
$\mathcal{O}_U$-module quasi-cohérent $\mathcal{F}|_U$.
Le lemme découle donc du cas des schémas ; voir
Morphismes, lemme \ref{morphisms-lemma-pf-flat-module-open}.
\end{proof}






\section{Platitude générique}
\label{spaces-morphisms-section-generic-flatness}

\noindent
Cette section est l'analogue de
Morphismes, section \ref{morphisms-section-generic-flatness}.

\begin{proposition}
\label{spaces-morphisms-proposition-generic-flatness-reduced}
Soit $S$ un schéma.
Soit $f : X \to Y$ un morphisme d'espaces algébriques sur $S$.
Soit $\mathcal{F}$ un faisceau quasi-cohérent de $\mathcal{O}_X$-modules.
Supposons que
\begin{enumerate}
\item $Y$ soit réduit ;
\item $f$ soit de type fini ;
\item $\mathcal{F}$ soit un $\mathcal{O}_X$-module de type fini.
\end{enumerate}
Alors il existe un sous-espace ouvert dense $W \subset Y$ tel que
le changement de base $X_W \to W$ de $f$ soit plat, localement de présentation finie et
quasi-compact, et tel que $\mathcal{F}|_{X_W}$ soit plate sur $W$ et de
présentation finie sur $\mathcal{O}_{X_W}$.
\end{proposition}

\begin{proof}
Soit $V$ un schéma et soit $V \to Y$ un morphisme étale surjectif.
Posons $X_V = V \times_Y X$ et notons $\mathcal{F}_V$ la restriction de
$\mathcal{F}$ à $X_V$. Supposons le résultat vrai pour le morphisme
$X_V \to V$ et le faisceau $\mathcal{F}_V$. Il existe alors un sous-schéma ouvert
$V' \subset V$ tel que $X_{V'} \to V'$ soit plat et de présentation finie,
et que $\mathcal{F}_{V'}$ soit un $\mathcal{O}_{X_{V'}}$-module de
présentation finie plat sur $V'$. Soit $W \subset Y$ l'image du
morphisme étale $V' \to Y$ ; voir
Propriétés des espaces, lemme \ref{spaces-properties-lemma-etale-image-open}.
Alors $V' \to W$ est un morphisme étale surjectif ; on voit donc que
$X_W \to W$ est plat, localement de présentation finie et quasi-compact d'après les
lemmes \ref{spaces-morphisms-lemma-finite-presentation-local},
\ref{spaces-morphisms-lemma-flat-local} et
\ref{spaces-morphisms-lemma-quasi-compact-local}.
D'après la discussion de
Propriétés des espaces, section
\ref{spaces-properties-section-properties-modules}
on voit que $\mathcal{F}_W$ est de présentation finie comme
$\mathcal{O}_{X_W}$-module et, d'après le
lemme \ref{spaces-morphisms-lemma-base-change-module-flat},
que $\mathcal{F}_W$ est plate sur $W$. Cet argument
ramène la proposition au cas où $Y$ est un schéma.

\medskip\noindent
Supposons que l'on sache démontrer la proposition lorsque $Y$ est un schéma affine.
Soit $f : X \to Y$ un morphisme de type fini d'espaces algébriques
sur $S$, où $Y$ est un schéma, et soit $\mathcal{F}$ un $\mathcal{O}_X$-module
quasi-cohérent de type fini. Choisissons un recouvrement ouvert affine
$Y = \bigcup V_j$. Par hypothèse, on peut trouver des ouverts denses $W_j \subset V_j$
tels que $X_{W_j} \to W_j$ soit plat, localement de présentation finie et
quasi-compact, et tels que $\mathcal{F}|_{X_{W_j}}$ soit plate sur $W_j$
et de présentation finie comme $\mathcal{O}_{X_{W_j}}$-module. Dans cette
situation, il suffit de prendre $W = \bigcup W_j$. On ramène ainsi
la proposition au cas où $Y$ est un schéma affine.

\medskip\noindent
Soit $Y$ un schéma affine sur $S$, soit
$f : X \to Y$ un morphisme de type fini d'espaces algébriques
sur $S$, et soit $\mathcal{F}$ un $\mathcal{O}_X$-module quasi-cohérent
de type fini. Comme $f$ est de type fini,
il est quasi-compact ; par suite, $X$ est quasi-compact. On peut donc trouver
un schéma affine $U$ et un morphisme étale surjectif $U \to X$ ; voir
Propriétés des espaces, lemme
\ref{spaces-properties-lemma-quasi-compact-affine-cover}.
Notons que $U \to Y$ est de type fini (c'est précisément ce que signifie ici
que $f$ est de type fini). On peut donc appliquer
Morphismes, Proposition \ref{morphisms-proposition-generic-flatness-reduced}
pour voir qu'il existe un ouvert dense $W \subset Y$ tel que
$U_W \to W$ soit plat et de présentation finie, et tel que
$\mathcal{F}|_{U_W}$ soit plate sur $W$ et de présentation finie
comme $\mathcal{O}_{U_W}$-module. D'après nos définitions, cela signifie que
le changement de base $X_W \to W$ de $f$ est plat, localement de présentation finie
et quasi-compact, et que $\mathcal{F}|_{X_W}$ est plate sur $W$ et de
présentation finie sur $\mathcal{O}_{X_W}$.
\end{proof}

\noindent
On ne peut renforcer le résultat de la proposition précédente en exigeant que
$X_W \to W$ soit de présentation finie, car
$\mathbf{A}^1_{\mathbf{Q}}/\mathbf{Z} \to \Spec(\mathbf{Q})$
fournit un contre-exemple. Le problème est que le morphisme diagonal
$\Delta_{X/Y}$ peut ne pas être quasi-compact, autrement dit que $f$ peut ne pas être
quasi-séparé. C'est manifestement le seul problème.

\begin{proposition}
\label{spaces-morphisms-proposition-generic-flatness-reduced-quasi-separated}
Soit $S$ un schéma.
Soit $f : X \to Y$ un morphisme d'espaces algébriques sur $S$.
Soit $\mathcal{F}$ un faisceau quasi-cohérent de $\mathcal{O}_X$-modules.
Supposons que
\begin{enumerate}
\item $Y$ soit réduit ;
\item $f$ soit quasi-séparé ;
\item $f$ soit de type fini ;
\item $\mathcal{F}$ soit un $\mathcal{O}_X$-module de type fini.
\end{enumerate}
Alors il existe un sous-espace ouvert dense $W \subset Y$ tel que
le changement de base $X_W \to W$ de $f$ soit plat et de présentation finie,
et tel que $\mathcal{F}|_{X_W}$ soit plate sur $W$ et de
présentation finie sur $\mathcal{O}_{X_W}$.
\end{proposition}

\begin{proof}
Cela résulte immédiatement de la
Proposition \ref{spaces-morphisms-proposition-generic-flatness-reduced}
et du fait que « de présentation finie » $=$
« localement de présentation finie » $+$ « quasi-compact » $+$
« quasi-séparé ».
\end{proof}















\section{Dimension relative}
\label{spaces-morphisms-section-relative-dimension}

\noindent
Dans cette section, nous définissons la dimension relative
d'un morphisme d'espaces algébriques en un point, ainsi que quelques
propriétés étroitement liées.

\begin{definition}
\label{spaces-morphisms-definition-dimension-fibre}
Soit $S$ un schéma.
Soit $f : X \to Y$ un morphisme d'espaces algébriques sur $S$.
Soit $x \in |X|$.
Soient $d, r \in \{0, 1, 2, \ldots, \infty\}$.
\begin{enumerate}
\item On dit que la
{\it dimension de l'anneau local de la fibre de $f$ en $x$} vaut $d$
si les conditions équivalentes du
lemme \ref{spaces-morphisms-lemma-local-source-target-at-point}
sont satisfaites pour la propriété
$\mathcal{P}_d$ décrite dans
Descente, lemme \ref{descent-lemma-dimension-local-ring-fibre}.
\item On dit que le
{\it degré de transcendance de $x/f(x)$} vaut $r$
si les conditions équivalentes du
lemme \ref{spaces-morphisms-lemma-local-source-target-at-point}
sont satisfaites pour la propriété
$\mathcal{P}_r$ décrite dans
Descente, lemme \ref{descent-lemma-transcendence-degree-at-point}.
\item On dit que
{\it $f$ est de dimension relative $d$ en $x$}
si les conditions équivalentes du
lemme \ref{spaces-morphisms-lemma-local-source-target-at-point}
sont satisfaites pour la propriété
$\mathcal{P}_d$ décrite dans
Descente, lemme \ref{descent-lemma-dimension-at-point}.
\end{enumerate}
\end{definition}

\noindent
Précisons ce que cela signifie. Choisissons des
diagrammes
$$
\xymatrix{
U \ar[d]_a \ar[r]_h & V \ar[d]^b \\
X \ar[r]^f & Y
}
\quad\quad
\xymatrix{
u \ar[d] \ar[r] & v \ar[d] \\
x \ar[r] & y
}
$$
comme dans le
lemme \ref{spaces-morphisms-lemma-local-source-target-at-point}.
Alors
$$
\begin{matrix}
\text{dimension relative de }f\text{ en }x & = &
\dim_u (U_v) \\
\text{dimension de l'anneau local de la fibre de }f\text{ en }x & = &
\dim(\mathcal{O}_{U_v, u})\\
\text{degré de transcendance de }x/f(x) & = &
\text{trdeg}_{\kappa(v)}(\kappa(u))
\end{matrix}
$$
Notons que si $Y = \Spec(k)$ est le spectre d'un corps, alors
la dimension relative de $X/Y$ en $x$ est égale à $\dim_x(X)$,
le degré de transcendance de $x/f(x)$ est le degré de transcendance
sur $k$, et la dimension de l'anneau local de la fibre de $f$
en $x$ est simplement la dimension de l'anneau local en $x$ ; autrement dit, les
notions relatives deviennent alors des notions absolues.

\begin{definition}
\label{spaces-morphisms-definition-relative-dimension}
Soit $S$ un schéma.
Soit $f : X \to Y$ un morphisme d'espaces algébriques sur $S$.
Soit $d \in \{0, 1, 2, \ldots\}$.
\begin{enumerate}
\item On dit que $f$ est de {\it dimension relative $\leq d$} si
$f$ est de dimension relative $\leq d$ en tout $x \in |X|$.
\item On dit que $f$ est de {\it dimension relative $d$} si
$f$ est de dimension relative $d$ en tout $x \in |X|$.
\end{enumerate}
\end{definition}

\noindent
Être de dimension relative {\it égale} à $d$ signifie, en gros, que toutes les
fibres non vides sont équidimensionnelles de dimension $d$.

\begin{lemma}
\label{spaces-morphisms-lemma-compare-tr-deg}
Soit $S$ un schéma. Soient $X \to Y \to Z$ des morphismes d'espaces
algébriques sur $S$. Soit $x \in |X|$, et soient $y \in |Y|$, $z \in |Z|$
ses images. Supposons que $X \to Y$ soit localement quasi-fini et que $Y \to Z$ soit
localement de type fini. Alors le degré de transcendance de $x/z$
est égal au degré de transcendance de $y/z$.
\end{lemma}

\begin{proof}
On peut choisir des diagrammes commutatifs
$$
\xymatrix{
U \ar[d] \ar[r] & V \ar[d] \ar[r] & W \ar[d] \\
X \ar[r] & Y \ar[r] & Z
}
\quad\quad
\xymatrix{
u \ar[d] \ar[r] & v \ar[d] \ar[r] & w \ar[d] \\
x \ar[r] & y \ar[r] & z
}
$$
où $U, V, W$ sont des schémas et où les flèches verticales sont étales.
Par définition, le morphisme $U \to V$ est localement quasi-fini,
ce qui entraîne que $\kappa(v) \subset \kappa(u)$ est finie ; voir
Morphismes, lemme \ref{morphisms-lemma-residue-field-quasi-finite}.
Le résultat est donc clair.
\end{proof}

\begin{lemma}
\label{spaces-morphisms-lemma-jacobson-finite-type-points}
Soit $S$ un schéma. Soit $f : X \to Y$ un morphisme d'espaces algébriques
sur $S$. Si $f$ est localement de type fini, si $Y$ est de Jacobson
(Propriétés des espaces, remarque
\ref{spaces-properties-remark-list-properties-local-etale-topology}),
et si $x \in |X|$ est un point de type fini de $X$,
alors le degré de transcendance de $x/f(x)$ vaut $0$.
\end{lemma}

\begin{proof}
Choisissons un schéma $V$ et un morphisme étale surjectif $V \to Y$.
Choisissons un schéma $U$ et un morphisme étale surjectif $U \to X \times_Y V$.
D'après le lemme \ref{spaces-morphisms-lemma-finite-type-points-surjective-morphism},
on peut trouver un point de type fini $u \in U$ s'envoyant sur $x$.
Après avoir rétréci $U$, on peut supposer que $u \in U$ est fermé
(Morphismes, lemme \ref{morphisms-lemma-identify-finite-type-points}).
Soit $v \in V$ l'image de $u$. D'après
Morphismes, lemme \ref{morphisms-lemma-jacobson-finite-type-points}
l'extension $\kappa(u)/\kappa(v)$ est finie.
Cela achève la démonstration.
\end{proof}

\begin{lemma}
\label{spaces-morphisms-lemma-rel-dimension-dimension}
Soit $S$ un schéma. Soit $f : X \to Y$ un morphisme d'espaces algébriques
localement noethériens sur $S$, plat, localement de type fini et de
dimension relative $d$. Pour tout point $x$ de $|X|$, d'image
$y$ dans $|Y|$, on a $\dim_x(X) = \dim_y(Y) + d$.
\end{lemma}

\begin{proof}
Par définition de la dimension d'un espace algébrique
en un point (Propriétés des espaces, définition
\ref{spaces-properties-definition-dimension-at-point})
et par définition de la dimension relative $d$,
on se ramène à l'énoncé correspondant pour les schémas
(Morphismes, lemme \ref{morphisms-lemma-rel-dimension-dimension}).
\end{proof}







\section{Morphismes et dimensions des fibres}
\label{spaces-morphisms-section-dimension-fibres}

\noindent
Cette section est l'analogue de
Morphismes, section \ref{morphisms-section-dimension-fibres}.
Les formulations de cette section sont un peu malcommodes, puisque
nous ne disposons pas d'anneaux locaux des espaces algébriques en leurs points.

\begin{lemma}
\label{spaces-morphisms-lemma-dimension-fibre-at-a-point}
Soit $S$ un schéma.
Soit $f : X \to Y$ un morphisme d'espaces algébriques sur $S$.
Soit $x \in |X|$.
Supposons que $f$ soit localement de type fini.
Alors
$$
\begin{matrix}
\text{dimension relative de }f\text{ en }x \\
= \\
\text{dimension de l'anneau local de la fibre de }f\text{ en }x \\
+ \\
\text{degré de transcendance de }x/f(x)
\end{matrix}
$$
où les notations sont celles de la
définition \ref{spaces-morphisms-definition-dimension-fibre}.
\end{lemma}

\begin{proof}
Cela résulte immédiatement de
Morphismes, lemme \ref{morphisms-lemma-dimension-fibre-at-a-point}
appliqué à $h : U \to V$ et $u \in U$
comme dans le
lemme \ref{spaces-morphisms-lemma-local-source-target-at-point}.
\end{proof}

\begin{lemma}
\label{spaces-morphisms-lemma-dimension-fibre-at-a-point-additive}
Soit $S$ un schéma.
Soient $f : X \to Y$ et $g : Y \to Z$ des morphismes d'espaces algébriques sur $S$.
Soit $x \in |X|$ et posons $y = f(x)$.
Supposons que $f$ et $g$ soient localement de type fini.
Alors :
\begin{enumerate}
\item
$$
\begin{matrix}
\text{dimension relative de }g \circ f\text{ en }x \\
\leq \\
\text{dimension relative de }f\text{ en }x \\
+ \\
\text{dimension relative de }g\text{ en }y
\end{matrix}
$$
\item il y a égalité dans (1) si, pour un certain morphisme $\Spec(k) \to Z$
depuis le spectre d'un corps représentant la classe de $g(f(x)) = g(y)$,
le morphisme $X_k \to Y_k$ est plat en $x$, par exemple si $f$ est plat en $x$ ;
\item
$$
\begin{matrix}
\text{degré de transcendance de }x/g(f(x)) \\
= \\
\text{degré de transcendance de }x/f(x) \\
+ \\
\text{degré de transcendance de }f(x)/g(f(x))
\end{matrix}
$$
\end{enumerate}
\end{lemma}

\begin{proof}
Choisissons un diagramme
$$
\xymatrix{
U \ar[d] \ar[r] & V \ar[d] \ar[r] & W \ar[d] \\
X \ar[r] & Y \ar[r] & Z
}
$$
où $U, V, W$ sont des schémas et où les flèches verticales sont étales et surjectives. (Voir
Espaces, lemme \ref{spaces-lemma-lift-morphism-presentations}.)
Choisissons $u \in U$ s'envoyant sur $x$. Notons $v, w$ les images
de $u$ dans $V, W$.
Appliquons
Morphismes, lemme \ref{morphisms-lemma-dimension-fibre-at-a-point-additive}
à la ligne supérieure et aux points $u, v, w$. Nous omettons les détails.
\end{proof}

\begin{lemma}
\label{spaces-morphisms-lemma-dimension-fibre-after-base-change}
Soit $S$ un schéma. Soit
$$
\xymatrix{
X' \ar[r]_{g'} \ar[d]_{f'} & X \ar[d]^f \\
Y' \ar[r]^g & Y
}
$$
un diagramme de produit fibré d'espaces algébriques sur $S$.
Soit $x' \in |X'|$. Posons $x = g'(x')$. Supposons que $f$ soit localement de type fini.
Alors :
\begin{enumerate}
\item
$$
\begin{matrix}
\text{dimension relative de }f\text{ en }x \\
= \\
\text{dimension relative de }f'\text{ en }x'
\end{matrix}
$$
\item on a
$$
\begin{matrix}
\text{dimension de l'anneau local de la fibre de }f'\text{ en }x' \\
- \\
\text{dimension de l'anneau local de la fibre de }f\text{ en }x \\
= \\
\text{degré de transcendance de }x/f(x) \\
- \\
\text{degré de transcendance de }x'/f'(x')
\end{matrix}
$$
et cette valeur commune est $\geq 0$ ;
\item étant donnés $x$ et $y' \in |Y'|$ s'envoyant sur un même $y \in |Y|$,
on peut choisir $x'$ de sorte que l'entier de (2) soit $0$.
\end{enumerate}
\end{lemma}

\begin{proof}
Choisissons un morphisme étale surjectif $V \to Y$, où $V$ est un schéma.
Choisissons un morphisme étale surjectif $U \to V \times_Y X$, où $U$ est un schéma.
Choisissons un morphisme étale surjectif $V' \to V \times_Y Y'$, où $V'$ est un schéma.
Posons $U' = V' \times_V U$.
Alors le morphisme induit $U' \to X'$ est lui aussi étale et
surjectif (nous omettons l'argument). Choisissons $u' \in U'$
s'envoyant sur $x'$. Les parties (1) et (2) résultent alors de l'application de
Morphismes, lemme \ref{morphisms-lemma-dimension-fibre-after-base-change}
au diagramme de schémas faisant intervenir $U', U, V', V$ et le point $u'$.
Pour démontrer (3), choisissons d'abord $v \in V$ s'envoyant sur $y$.
Puis, à l'aide de Propriétés des espaces, lemme
\ref{spaces-properties-lemma-points-cartesian}
nous pouvons choisir $v' \in V'$ s'envoyant sur $y'$ et $v$, ainsi que
$u \in U$ s'envoyant sur $x$ et $v$. Enfin, d'après
Morphismes, lemme \ref{morphisms-lemma-dimension-fibre-after-base-change}
nous pouvons choisir $u' \in U'$ s'envoyant sur $v'$ et $u$ de sorte que
l'entier soit nul. Il suffit alors de prendre pour $x' \in |X'|$ l'image de $u'$.
\end{proof}

\begin{lemma}
\label{spaces-morphisms-lemma-openness-bounded-dimension-fibres}
Soit $S$ un schéma.
Soit $f : X \to Y$ un morphisme d'espaces algébriques sur $S$.
Soit $n \geq 0$. Supposons que $f$ soit localement de type fini.
La partie
$$
W_n = \{x \in |X|
\text{ tels que la dimension relative de }f\text{ en } x \leq n\}
$$
est ouverte dans $|X|$.
\end{lemma}

\begin{proof}
Choisissons un diagramme
$$
\xymatrix{
U \ar[r]_h \ar[d]_a & V \ar[d] \\
X \ar[r] & Y
}
$$
où $U$ et $V$ sont des schémas et où les flèches verticales sont surjectives et
étales ; voir
Espaces, lemme \ref{spaces-lemma-lift-morphism-presentations}.
D'après
Morphismes, lemme \ref{morphisms-lemma-openness-bounded-dimension-fibres}
la partie $U_n$ des points où $h$ est de dimension relative
$\leq n$ est ouverte dans $U$. D'après notre définition de la dimension relative
des morphismes d'espaces algébriques en leurs points, on voit que
$U_n = a^{-1}(W_n)$.
Le lemme résulte de la définition de la topologie sur $|X|$.
\end{proof}

\begin{lemma}
\label{spaces-morphisms-lemma-openness-bounded-dimension-fibres-finite-presentation}
Soit $S$ un schéma.
Soit $f : X \to Y$ un morphisme d'espaces algébriques sur $S$.
Soit $n \geq 0$. Supposons que $f$ soit localement de présentation finie.
L'ouvert
$$
W_n = \{x \in |X|
\text{ tels que la dimension relative de }f\text{ en } x \leq n\}
$$
du lemme \ref{spaces-morphisms-lemma-openness-bounded-dimension-fibres}
est rétrocompact dans $|X|$. (Voir
Topologie, définition \ref{topology-definition-quasi-compact}.)
\end{lemma}

\begin{proof}
Choisissons un diagramme
$$
\xymatrix{
U \ar[r]_h \ar[d]_a & V \ar[d] \\
X \ar[r] & Y
}
$$
où $U$ et $V$ sont des schémas et où les flèches verticales sont surjectives et
étales ; voir
Espaces, lemme \ref{spaces-lemma-lift-morphism-presentations}.
Dans la démonstration du
lemme \ref{spaces-morphisms-lemma-openness-bounded-dimension-fibres},
nous avons vu que $a^{-1}(W_n) = U_n$ est la partie correspondante
pour le morphisme $h$. D'après
Morphismes, lemme
\ref{morphisms-lemma-openness-bounded-dimension-fibres-finite-presentation}
on voit que $U_n$ est rétrocompact dans $U$.
Le lemme résulte de la définition de la topologie sur $|X|$ ; comparer à
Propriétés des espaces,
lemme \ref{spaces-properties-lemma-space-locally-quasi-compact}
et à sa démonstration.
\end{proof}

\begin{lemma}
\label{spaces-morphisms-lemma-locally-quasi-finite-rel-dimension-0}
Soit $S$ un schéma.
Soit $f : X \to Y$ un morphisme d'espaces algébriques sur $S$.
Supposons que $f$ soit localement de type fini.
Alors $f$ est localement quasi-fini si et seulement si $f$ est de
dimension relative $0$ en tout $x \in |X|$.
\end{lemma}

\begin{proof}
Choisissons un diagramme
$$
\xymatrix{
U \ar[r]_h \ar[d]_a & V \ar[d] \\
X \ar[r] & Y
}
$$
où $U$ et $V$ sont des schémas et où les flèches verticales sont surjectives et
étales ; voir
Espaces, lemme \ref{spaces-lemma-lift-morphism-presentations}.
Les définitions entraînent que
$h$ est localement quasi-fini si et seulement si $f$ l'est,
et que $f$ est de dimension relative $0$ en tout $x \in |X|$ si et
seulement si $h$ est de dimension relative $0$ en tout $u \in U$.
Le résultat découle donc du résultat pour $h$, à savoir
Morphismes, lemme \ref{morphisms-lemma-locally-quasi-finite-rel-dimension-0}.
\end{proof}

\begin{lemma}
\label{spaces-morphisms-lemma-locally-finite-type-quasi-finite-part}
Soit $S$ un schéma.
Soit $f : X \to Y$ un morphisme d'espaces algébriques sur $S$.
Supposons que $f$ soit localement de type fini.
Alors il existe un sous-espace ouvert canonique $X' \subset X$ tel que
$f|_{X'} : X' \to Y$ soit localement quasi-fini, et tel que la
dimension relative de $f$ en tout $x \in |X|$, $x \not \in |X'|$, soit
$\geq 1$. La formation de $X'$ commute à tout changement de base.
\end{lemma}

\begin{proof}
Combiner les
lemmes \ref{spaces-morphisms-lemma-openness-bounded-dimension-fibres},
\ref{spaces-morphisms-lemma-locally-quasi-finite-rel-dimension-0} et
\ref{spaces-morphisms-lemma-dimension-fibre-after-base-change}.
\end{proof}

\begin{lemma}
\label{spaces-morphisms-lemma-quasi-finite-at-point}
Soit $S$ un schéma. Considérons un diagramme cartésien
$$
\xymatrix{
X \ar[d] & F \ar[l]^p \ar[d] \\
Y & \Spec(k) \ar[l]
}
$$
où $X \to Y$ est un morphisme d'espaces algébriques sur $S$
localement de type fini et où $k$ est un corps sur $S$.
Soit $z \in |F|$ tel que $\dim_z(F) = 0$. Alors, après avoir remplacé $X$
par un sous-espace ouvert contenant $p(z)$, le morphisme
$$
X \longrightarrow Y
$$
est localement quasi-fini.
\end{lemma}

\begin{proof}
Soit $X' \subset X$ le sous-espace ouvert sur lequel $f$ est localement quasi-fini,
construit dans le
lemme \ref{spaces-morphisms-lemma-locally-finite-type-quasi-finite-part}.
Comme la formation de $X'$ commute à tout changement de base, on voit
que $z \in X' \times_Y \Spec(k)$. Le lemme est donc clair.
\end{proof}





\section{La formule de dimension}
\label{spaces-morphisms-section-dimension-formula}

\noindent
L'analogue de la formule de dimension
(Morphismes, lemme \ref{morphisms-lemma-dimension-formula})
est un peu délicat à formuler, car il faudrait définir
les espaces algébriques intègres (ce que nous ferons plus loin), ainsi que
les espaces algébriques universellement caténaires. La
version suivante est toutefois immédiate.

\begin{lemma}
\label{spaces-morphisms-lemma-dimension-formula-general}
Soit $S$ un schéma. Soit $f : X \to Y$ un morphisme d'espaces
algébriques sur $S$. Supposons que $Y$ soit localement noethérien et que $f$ soit localement
de type fini. Soit $x \in |X|$, d'image $y \in |Y|$.
Alors
\begin{align*}
& \text{la dimension de l'anneau local de }X\text{ en }x \leq \\
& \text{la dimension de l'anneau local de }Y\text{ en }y + E - \\
& \text{ le degré de transcendance de }x/y
\end{align*}
où $E$ est le maximum des degrés de transcendance de $\xi/f(\xi)$,
lorsque $\xi \in |X|$ parcourt les points se spécialisant en $x$ en lesquels
l'anneau local de $X$ est de dimension $0$.
\end{lemma}

\begin{proof}
Choisissons un schéma affine $V$, un morphisme étale $V \to Y$ et un point
$v \in V$ s'envoyant sur $y$. Choisissons un schéma affine $U$, un morphisme étale
$U \to X \times_Y V$ et un point $u \in U$ s'envoyant sur $v$ dans $V$ et sur $x$
dans $X$. En explicitant la définition \ref{spaces-morphisms-definition-dimension-fibre} et
Propriétés des espaces, définition
\ref{spaces-properties-definition-dimension-local-ring}
il faut montrer que
$$
\dim(\mathcal{O}_{U, u}) \leq
\dim(\mathcal{O}_{V, v}) + E - \text{trdeg}_{\kappa(v)}(\kappa(u))
$$
Soit $\xi_U \in U$ un point générique d'une composante irréductible de
$U$ contenant $u$. Alors $\xi_U$ s'envoie sur un point $\xi \in |X|$
figurant dans la liste utilisée pour définir la quantité $E$ et, en fait,
chaque $\xi$ intervenant dans la définition de $E$ s'obtient ainsi
(nous omettons ce petit détail). En particulier, il n'y a qu'un nombre fini
de tels $\xi$, et l'on peut prendre le maximum (c'est donc bien
un maximum et non une borne supérieure).
Le degré de transcendance de $\xi$ sur $f(\xi)$ est
$\text{trdeg}_{\kappa(\xi_V)}(\kappa(\xi_U))$, où $\xi_V \in V$
est l'image de $\xi_U$. Le lemme découle donc de
Morphismes, lemme \ref{morphisms-lemma-dimension-formula-general}.
\end{proof}

\begin{lemma}
\label{spaces-morphisms-lemma-alteration-dimension-general}
Soit $S$ un schéma.
Soit $f : X \to Y$ un morphisme d'espaces algébriques sur $S$. Supposons
que $Y$ soit localement noethérien et que $f$ soit localement de type fini.
Alors
$$
\dim(X) \leq \dim(Y) + E
$$
où $E$ est la borne supérieure des degrés de transcendance de
$\xi/f(\xi)$, lorsque $\xi$ parcourt les points en
lesquels l'anneau local de $X$ est de dimension $0$.
\end{lemma}

\begin{proof}
Conséquence immédiate du lemme \ref{spaces-morphisms-lemma-dimension-formula-general}
et de Propriétés des espaces, lemme \ref{spaces-properties-lemma-dimension}.
\end{proof}







\section{Morphismes syntomiques}
\label{spaces-morphisms-section-syntomic}

\noindent
La propriété « syntomique » des morphismes de schémas est
locale sur la source et le but pour la topologie étale ; voir
Descente, remarque \ref{descent-remark-list-local-source-target}.
Elle est aussi stable par changement de base et locale sur le but pour la topologie fpqc ; voir
Morphismes, lemme \ref{morphisms-lemma-base-change-syntomic} et
Descente, lemme \ref{descent-lemma-descending-property-syntomic}.
Ainsi, d'après le
lemme \ref{spaces-morphisms-lemma-local-source-target}
ci-dessus, on peut définir comme suit la notion de morphisme syntomique d'espaces
algébriques ; cette définition coïncide avec la notion déjà définie dans la
section \ref{spaces-morphisms-section-representable}
lorsque le morphisme est représentable.

\begin{definition}
\label{spaces-morphisms-definition-syntomic}
Soit $S$ un schéma.
Soit $f : X \to Y$ un morphisme d'espaces algébriques sur $S$.
\begin{enumerate}
\item On dit que $f$ est {\it syntomique} si les conditions équivalentes du
lemme \ref{spaces-morphisms-lemma-local-source-target}
sont satisfaites pour $\mathcal{P} =$ « syntomique ».
\item Soit $x \in |X|$. On dit que $f$ est {\it syntomique en $x$} s'il
existe un voisinage ouvert $X' \subset X$ de $x$ tel
que $f|_{X'} : X' \to Y$ soit syntomique.
\end{enumerate}
\end{definition}

\begin{lemma}
\label{spaces-morphisms-lemma-composition-syntomic}
Le composé de morphismes syntomiques est syntomique.
\end{lemma}

\begin{proof}
Voir la remarque \ref{spaces-morphisms-remark-composition-P} et
Morphismes, lemme \ref{morphisms-lemma-composition-syntomic}.
\end{proof}

\begin{lemma}
\label{spaces-morphisms-lemma-base-change-syntomic}
Le changement de base d'un morphisme syntomique est syntomique.
\end{lemma}

\begin{proof}
Voir la remarque \ref{spaces-morphisms-remark-base-change-P} et
Morphismes, lemme \ref{morphisms-lemma-base-change-syntomic}.
\end{proof}

\begin{lemma}
\label{spaces-morphisms-lemma-syntomic-local}
Soit $S$ un schéma.
Soit $f : X \to Y$ un morphisme d'espaces algébriques sur $S$.
Les conditions suivantes sont équivalentes :
\begin{enumerate}
\item $f$ est syntomique ;
\item pour tout $x \in |X|$, le morphisme $f$ est syntomique en $x$ ;
\item pour tout schéma $Z$ et tout morphisme $Z \to Y$, le morphisme
$Z \times_Y X \to Z$ est syntomique ;
\item pour tout schéma affine $Z$ et tout morphisme
$Z \to Y$, le morphisme $Z \times_Y X \to Z$ est syntomique ;
\item il existe un schéma $V$ et un morphisme étale surjectif
$V \to Y$ tels que $V \times_Y X \to V$ soit syntomique ;
\item il existe un schéma $U$ et un morphisme étale surjectif
$\varphi : U \to X$ tels que le composé $f \circ \varphi$
soit syntomique ;
\item pour tout diagramme commutatif
$$
\xymatrix{
U \ar[d] \ar[r] & V \ar[d] \\
X \ar[r] & Y
}
$$
où $U$ et $V$ sont des schémas et où les flèches verticales sont étales,
la flèche horizontale supérieure est syntomique ;
\item il existe un diagramme commutatif
$$
\xymatrix{
U \ar[d] \ar[r] & V \ar[d] \\
X \ar[r] & Y
}
$$
où $U$ et $V$ sont des schémas, où les flèches verticales sont étales et où
$U \to X$ est surjectif, tel que la flèche horizontale supérieure soit syntomique ;
\item il existe des recouvrements de Zariski $Y = \bigcup_{i \in I} Y_i$
et $f^{-1}(Y_i) = \bigcup X_{ij}$ tels que
chaque morphisme $X_{ij} \to Y_i$ soit syntomique.
\end{enumerate}
\end{lemma}

\begin{proof}
La démonstration est omise.
\end{proof}

\begin{lemma}
\label{spaces-morphisms-lemma-syntomic-locally-finite-presentation}
Un morphisme syntomique est localement de présentation finie.
\end{lemma}

\begin{proof}
Cela résulte immédiatement du cas des schémas
(Morphismes, lemme \ref{morphisms-lemma-syntomic-locally-finite-presentation}).
\end{proof}

\begin{lemma}
\label{spaces-morphisms-lemma-syntomic-flat}
Un morphisme syntomique est plat.
\end{lemma}

\begin{proof}
Cela résulte immédiatement du cas des schémas
(Morphismes, lemme \ref{morphisms-lemma-syntomic-flat}).
\end{proof}

\begin{lemma}
\label{spaces-morphisms-lemma-syntomic-open}
Un morphisme syntomique est universellement ouvert.
\end{lemma}

\begin{proof}
Combiner les
lemmes \ref{spaces-morphisms-lemma-syntomic-locally-finite-presentation},
\ref{spaces-morphisms-lemma-syntomic-flat} et
\ref{spaces-morphisms-lemma-fppf-open}.
\end{proof}





\section{Morphismes lisses}
\label{spaces-morphisms-section-smooth}

\noindent
La propriété « lisse » des morphismes de schémas est
locale sur la source et le but pour la topologie étale ; voir
Descente, remarque \ref{descent-remark-list-local-source-target}.
Elle est aussi stable par changement de base et locale sur le but pour la topologie fpqc ; voir
Morphismes, lemme \ref{morphisms-lemma-base-change-smooth}
et
Descente, lemme \ref{descent-lemma-descending-property-smooth}.
Ainsi, d'après le
lemme \ref{spaces-morphisms-lemma-local-source-target}
ci-dessus, on peut définir comme suit la notion de morphisme lisse d'espaces
algébriques ; cette définition coïncide avec la notion déjà définie dans la
section \ref{spaces-morphisms-section-representable}
lorsque le morphisme est représentable.

\begin{definition}
\label{spaces-morphisms-definition-smooth}
Soit $S$ un schéma.
Soit $f : X \to Y$ un morphisme d'espaces algébriques sur $S$.
\begin{enumerate}
\item On dit que $f$ est {\it lisse} si les conditions équivalentes du
lemme \ref{spaces-morphisms-lemma-local-source-target} sont satisfaites pour
$\mathcal{P} =$ « lisse ».
\item Soit $x \in |X|$. On dit que $f$ est {\it lisse en $x$} s'il existe
un voisinage ouvert $X' \subset X$ de $x$ tel que $f|_{X'} : X' \to Y$
soit lisse.
\end{enumerate}
\end{definition}

\begin{lemma}
\label{spaces-morphisms-lemma-composition-smooth}
Le composé de morphismes lisses est lisse.
\end{lemma}

\begin{proof}
Voir la remarque \ref{spaces-morphisms-remark-composition-P} et
Morphismes, lemme \ref{morphisms-lemma-composition-smooth}.
\end{proof}

\begin{lemma}
\label{spaces-morphisms-lemma-base-change-smooth}
Le changement de base d'un morphisme lisse est lisse.
\end{lemma}

\begin{proof}
Voir la remarque \ref{spaces-morphisms-remark-base-change-P} et
Morphismes, lemme \ref{morphisms-lemma-base-change-smooth}.
\end{proof}

\begin{lemma}
\label{spaces-morphisms-lemma-smooth-local}
Soit $S$ un schéma.
Soit $f : X \to Y$ un morphisme d'espaces algébriques sur $S$.
Les conditions suivantes sont équivalentes :
\begin{enumerate}
\item $f$ est lisse ;
\item pour tout $x \in |X|$, le morphisme $f$ est lisse en $x$ ;
\item pour tout schéma $Z$ et tout morphisme $Z \to Y$, le morphisme
$Z \times_Y X \to Z$ est lisse ;
\item pour tout schéma affine $Z$ et tout morphisme
$Z \to Y$, le morphisme $Z \times_Y X \to Z$ est lisse ;
\item il existe un schéma $V$ et un morphisme étale surjectif
$V \to Y$ tels que $V \times_Y X \to V$ soit un morphisme lisse ;
\item il existe un schéma $U$ et un morphisme étale surjectif
$\varphi : U \to X$ tels que le composé $f \circ \varphi$
soit lisse ;
\item pour tout diagramme commutatif
$$
\xymatrix{
U \ar[d] \ar[r] & V \ar[d] \\
X \ar[r] & Y
}
$$
où $U$ et $V$ sont des schémas et où les flèches verticales sont étales,
la flèche horizontale supérieure est lisse ;
\item il existe un diagramme commutatif
$$
\xymatrix{
U \ar[d] \ar[r] & V \ar[d] \\
X \ar[r] & Y
}
$$
où $U$ et $V$ sont des schémas, où les flèches verticales sont étales et où
$U \to X$ est surjectif, tel que la flèche horizontale supérieure soit lisse ;
\item il existe des recouvrements de Zariski $Y = \bigcup_{i \in I} Y_i$
et $f^{-1}(Y_i) = \bigcup X_{ij}$ tels que
chaque morphisme $X_{ij} \to Y_i$ soit lisse.
\end{enumerate}
\end{lemma}

\begin{proof}
La démonstration est omise.
\end{proof}

\begin{lemma}
\label{spaces-morphisms-lemma-smooth-locally-finite-presentation}
Un morphisme lisse d'espaces algébriques est localement de présentation finie.
\end{lemma}

\begin{proof}
Soit $X \to Y$ un morphisme lisse d'espaces algébriques. Par
définition, cela signifie qu'il existe un diagramme comme dans le
lemme \ref{spaces-morphisms-lemma-local-source-target},
où $h$ est lisse et où $a$ est une flèche verticale surjective. D'après
Morphismes, lemme \ref{morphisms-lemma-smooth-locally-finite-presentation}
$h$ est localement de présentation finie. Ainsi, $X \to Y$ est
localement de présentation finie par définition.
\end{proof}

\begin{lemma}
\label{spaces-morphisms-lemma-smooth-locally-finite-type}
Un morphisme lisse d'espaces algébriques est localement de type fini.
\end{lemma}

\begin{proof}
Combiner les
lemmes \ref{spaces-morphisms-lemma-smooth-locally-finite-presentation} et
\ref{spaces-morphisms-lemma-finite-presentation-finite-type}.
\end{proof}

\begin{lemma}
\label{spaces-morphisms-lemma-smooth-flat}
Un morphisme lisse d'espaces algébriques est plat.
\end{lemma}

\begin{proof}
Soit $X \to Y$ un morphisme lisse d'espaces algébriques. Par
définition, cela signifie qu'il existe un diagramme comme dans le
lemme \ref{spaces-morphisms-lemma-local-source-target},
où $h$ est lisse et où $a$ est une flèche verticale surjective. D'après
Morphismes, lemme \ref{morphisms-lemma-smooth-flat},
$h$ est plat. Ainsi, $X \to Y$ est plat par définition.
\end{proof}

\begin{lemma}
\label{spaces-morphisms-lemma-smooth-syntomic}
Un morphisme lisse d'espaces algébriques est syntomique.
\end{lemma}

\begin{proof}
Soit $X \to Y$ un morphisme lisse d'espaces algébriques. Par
définition, cela signifie qu'il existe un diagramme comme dans le
lemme \ref{spaces-morphisms-lemma-local-source-target},
où $h$ est lisse et où $a$ est une flèche verticale surjective. D'après
Morphismes, lemme \ref{morphisms-lemma-smooth-syntomic}
$h$ est syntomique. Ainsi, $X \to Y$ est syntomique par définition.
\end{proof}

\begin{lemma}
\label{spaces-morphisms-lemma-where-smooth}
Soit $S$ un schéma. Soit $f : X \to Y$ un morphisme d'espaces
algébriques sur $S$. Il existe un plus grand sous-espace ouvert $U \subset X$
tel que $f|_U : U \to Y$ soit lisse. De plus, la formation de
cet ouvert commute au changement de base par
\begin{enumerate}
\item les morphismes plats et
localement de présentation finie ;
\item les morphismes plats, pourvu que $f$ soit localement de présentation finie.
\end{enumerate}
\end{lemma}

\begin{proof}
L'existence de $U$ résulte du fait que la propriété
d'être lisse est locale sur la source pour la topologie de Zariski (et même étale) ; voir le
lemme \ref{spaces-morphisms-lemma-smooth-local}. De plus, ce lemme permet
de ramener les propriétés (1) et (2) au cas
des morphismes de schémas. Le cas des schémas est traité dans
Morphismes, lemme \ref{morphisms-lemma-set-points-where-fibres-smooth}.
Certains détails sont omis.
\end{proof}

\begin{lemma}
\label{spaces-morphisms-lemma-smoothness-dimension-spaces}
Soient $X$ et $Y$ des espaces algébriques localement noethériens sur un schéma
$S$, et soit $f : X \to Y$ un morphisme lisse.
Pour tout point $x \in |X|$ d'image $y \in |Y|$, on a
$$
\dim_x(X) = \dim_y(Y) + \dim_x(X_y)
$$
où $\dim_x(X_y)$ est la dimension relative de $f$ en $x$ au sens de la
définition \ref{spaces-morphisms-definition-dimension-fibre}.
\end{lemma}

\begin{proof}
Par définition de la dimension d'un espace algébrique
en un point (Propriétés des espaces, définition
\ref{spaces-properties-definition-dimension-at-point}),
on se ramène à l'énoncé correspondant pour les schémas
(Morphismes, lemme \ref{morphisms-lemma-smoothness-dimension}).
\end{proof}



\section{Morphismes non ramifiés}
\label{spaces-morphisms-section-unramified}

\noindent
La propriété « non ramifié » (resp.\ « G-non ramifié »)
des morphismes de schémas est locale sur la source et le but pour la topologie étale ; voir
Descente, remarque \ref{descent-remark-list-local-source-target}.
Elle est aussi stable par changement de base et locale sur le but pour la topologie fpqc ; voir
Morphismes, lemme \ref{morphisms-lemma-base-change-unramified} et
Descente, lemme \ref{descent-lemma-descending-property-unramified}.
Ainsi, d'après le
lemme \ref{spaces-morphisms-lemma-local-source-target}
ci-dessus, on peut définir comme suit la notion de morphisme non ramifié
(resp.\ de morphisme G-non ramifié) d'espaces algébriques ;
cette définition coïncide avec la notion déjà définie dans la
section \ref{spaces-morphisms-section-representable}
lorsque le morphisme est représentable.

\begin{definition}
\label{spaces-morphisms-definition-unramified}
Soit $S$ un schéma.
Soit $f : X \to Y$ un morphisme d'espaces algébriques sur $S$.
\begin{enumerate}
\item On dit que $f$ est {\it non ramifié} si les conditions équivalentes du
lemme \ref{spaces-morphisms-lemma-local-source-target}
sont satisfaites pour $\mathcal{P} = \text{non ramifié}$.
\item Soit $x \in |X|$. On dit que $f$ est {\it non ramifié en $x$} s'il
existe un voisinage ouvert $X' \subset X$ de $x$ tel que
$f|_{X'} : X' \to Y$ soit non ramifié.
\item On dit que $f$ est {\it G-non ramifié} si les conditions équivalentes du
lemme \ref{spaces-morphisms-lemma-local-source-target}
sont satisfaites pour $\mathcal{P} = \text{G-non ramifié}$.
\item Soit $x \in |X|$. On dit que $f$ est {\it G-non ramifié en $x$} s'il
existe un voisinage ouvert $X' \subset X$ de $x$ tel que
$f|_{X'} : X' \to Y$ soit G-non ramifié.
\end{enumerate}
\end{definition}

\noindent
En raison du lemme suivant, nous ne développerons désormais que la théorie
des morphismes non ramifiés et, chaque fois que nous voudrons utiliser un morphisme
G-non ramifié, nous dirons simplement « un morphisme non ramifié localement de
présentation finie ».

\begin{lemma}
\label{spaces-morphisms-lemma-unramified-G-unramified}
Soit $S$ un schéma.
Soit $f : X \to Y$ un morphisme d'espaces algébriques sur $S$.
Alors $f$ est G-non ramifié si et seulement si $f$ est non ramifié et
localement de présentation finie.
\end{lemma}

\begin{proof}
Considérons un diagramme quelconque comme dans le
lemme \ref{spaces-morphisms-lemma-local-source-target}.
L'assertion revient alors à dire que le morphisme $h$ est
G-non ramifié si et seulement s'il est non ramifié et localement de
présentation finie. Cela résulte immédiatement de
Morphismes, définition \ref{morphisms-definition-unramified}.
\end{proof}

\begin{lemma}
\label{spaces-morphisms-lemma-composition-unramified}
Le composé de morphismes non ramifiés est non ramifié.
\end{lemma}

\begin{proof}
Voir la remarque \ref{spaces-morphisms-remark-composition-P} et
Morphismes, lemme \ref{morphisms-lemma-composition-unramified}.
\end{proof}

\begin{lemma}
\label{spaces-morphisms-lemma-base-change-unramified}
Le changement de base d'un morphisme non ramifié est non ramifié.
\end{lemma}

\begin{proof}
Voir la remarque \ref{spaces-morphisms-remark-base-change-P} et
Morphismes, lemme \ref{morphisms-lemma-base-change-unramified}.
\end{proof}

\begin{lemma}
\label{spaces-morphisms-lemma-unramified-local}
Soit $S$ un schéma.
Soit $f : X \to Y$ un morphisme d'espaces algébriques sur $S$.
Les conditions suivantes sont équivalentes :
\begin{enumerate}
\item $f$ est non ramifié ;
\item pour tout $x \in |X|$, le morphisme $f$ est non ramifié en $x$ ;
\item pour tout schéma $Z$ et tout morphisme $Z \to Y$, le morphisme
$Z \times_Y X \to Z$ est non ramifié ;
\item pour tout schéma affine $Z$ et tout morphisme
$Z \to Y$, le morphisme $Z \times_Y X \to Z$ est non ramifié ;
\item il existe un schéma $V$ et un morphisme étale surjectif
$V \to Y$ tels que $V \times_Y X \to V$ soit un morphisme non ramifié ;
\item il existe un schéma $U$ et un morphisme étale surjectif
$\varphi : U \to X$ tels que le composé $f \circ \varphi$
soit non ramifié ;
\item pour tout diagramme commutatif
$$
\xymatrix{
U \ar[d] \ar[r] & V \ar[d] \\
X \ar[r] & Y
}
$$
où $U$ et $V$ sont des schémas et où les flèches verticales sont étales,
la flèche horizontale supérieure est non ramifiée ;
\item il existe un diagramme commutatif
$$
\xymatrix{
U \ar[d] \ar[r] & V \ar[d] \\
X \ar[r] & Y
}
$$
où $U$ et $V$ sont des schémas, où les flèches verticales sont étales et où
$U \to X$ est surjectif, tel que la flèche horizontale supérieure soit non ramifiée ;
\item il existe des recouvrements de Zariski $Y = \bigcup_{i \in I} Y_i$
et $f^{-1}(Y_i) = \bigcup X_{ij}$ tels que
chaque morphisme $X_{ij} \to Y_i$ soit non ramifié.
\end{enumerate}
\end{lemma}

\begin{proof}
La démonstration est omise.
\end{proof}


\begin{lemma}
\label{spaces-morphisms-lemma-unramified-locally-finite-type}
Un morphisme non ramifié d'espaces algébriques est localement de type fini.
\end{lemma}

\begin{proof}
Au moyen d'un diagramme comme dans le
lemme \ref{spaces-morphisms-lemma-local-source-target},
l'assertion se ramène à
Morphismes, lemme \ref{morphisms-lemma-unramified-locally-finite-type}.
\end{proof}

\begin{lemma}
\label{spaces-morphisms-lemma-unramified-quasi-finite}
Si $f$ est non ramifié en $x$, alors $f$ est quasi-fini en $x$.
En particulier, un morphisme non ramifié est localement quasi-fini.
\end{lemma}

\begin{proof}
Au moyen d'un diagramme comme dans le
lemme \ref{spaces-morphisms-lemma-local-source-target},
l'assertion se ramène à
Morphismes, lemme \ref{morphisms-lemma-unramified-quasi-finite}.
\end{proof}

\begin{lemma}
\label{spaces-morphisms-lemma-immersion-unramified}
Une immersion d'espaces algébriques est non ramifiée.
\end{lemma}

\begin{proof}
Soit $i : X \to Y$ une immersion d'espaces algébriques. Choisissons un schéma
$V$ et un morphisme étale surjectif $V \to Y$. Alors $V \times_Y X \to V$
est une immersion de schémas, donc est non ramifié (voir
Morphismes, lemmes \ref{morphisms-lemma-open-immersion-unramified} et
\ref{morphisms-lemma-closed-immersion-unramified}).
Ainsi, $i$ est non ramifié par définition.
\end{proof}

\begin{lemma}
\label{spaces-morphisms-lemma-diagonal-unramified-morphism}
Soit $S$ un schéma.
Soit $f : X \to Y$ un morphisme d'espaces algébriques sur $S$.
\begin{enumerate}
\item Si $f$ est non ramifié, alors le morphisme diagonal
$\Delta_{X/Y} : X \to X \times_Y X$ est une immersion ouverte.
\item Si $f$ est localement de type fini
et si $\Delta_{X/Y}$ est une immersion ouverte, alors $f$ est non ramifié.
\end{enumerate}
\end{lemma}

\begin{proof}
On sait dans tous les cas que $\Delta_{X/Y}$ est un monomorphisme représentable ; voir le
lemme \ref{spaces-morphisms-lemma-properties-diagonal}.
Choisissons un schéma $V$ et un morphisme étale surjectif $V \to Y$.
Choisissons un schéma $U$ et un morphisme étale surjectif $U \to X \times_Y V$.
Considérons le diagramme commutatif
$$
\xymatrix{
U \ar[d] \ar[rr]_-{\Delta_{U/V}} & &
U \times_V U \ar[d] \ar[r] &
V \ar[d]^{\Delta_{V/Y}} \\
X \ar[rr]^-{\Delta_{X/Y}} & &
X \times_Y X \ar[r] &
V \times_Y V
}
$$
dont le carré de droite est cartésien. La flèche verticale de gauche est étale surjective.
La flèche verticale de droite est étale en tant que morphisme entre des schémas
étales sur $Y$ ; voir
Propriétés des espaces,
lemme \ref{spaces-properties-lemma-etale-permanence}.
Par conséquent, la flèche verticale du milieu est également étale (mais elle n'est pas
nécessairement surjective).

\medskip\noindent
Supposons que $f$ soit non ramifié. Alors $U \to V$ est non ramifié, donc
$\Delta_{U/V}$ est une immersion ouverte d'après
Morphismes, lemme \ref{morphisms-lemma-diagonal-unramified-morphism}.
En considérant le carré de gauche du diagramme ci-dessus, on conclut que
$\Delta_{X/Y}$ est un morphisme étale ; voir
Propriétés des espaces,
lemme \ref{spaces-properties-lemma-etale-local}.
Ainsi, $\Delta_{X/Y}$ est un monomorphisme étale représentable, ce qui
implique que c'est une immersion ouverte d'après
\'Morphismes étales, théorème \ref{etale-theorem-etale-radicial-open}.
(Voir aussi
Espaces, lemme
\ref{spaces-lemma-representable-transformations-property-implication}
pour le passage du langage des schémas à celui des foncteurs.)

\medskip\noindent
Supposons que $f$ soit localement de type fini et que $\Delta_{X/Y}$
soit une immersion ouverte. Cela implique que $U \to V$ est aussi localement de type
fini (par définition d'un morphisme d'espaces algébriques qui est
localement de type fini). En considérant le diagramme ci-dessus,
on conclut que $\Delta_{U/V}$ est étale en tant que morphisme entre des
schémas étales sur $X \times_Y X$ ; voir
Propriétés des espaces,
lemme \ref{spaces-properties-lemma-etale-permanence}.
Mais, puisque $\Delta_{U/V}$ est la diagonale d'un morphisme entre schémas,
on voit que c'est en tout cas une immersion ; voir
Schémas, lemme \ref{schemes-lemma-diagonal-immersion}.
Il s'agit donc d'une immersion ouverte, et l'on conclut
que $U \to V$ est non ramifié d'après
Morphismes, lemme \ref{morphisms-lemma-diagonal-unramified-morphism}.
Cela signifie à son tour que $f$ est non ramifié par définition.
\end{proof}

\begin{lemma}
\label{spaces-morphisms-lemma-where-unramified}
Soit $S$ un schéma. Considérons un diagramme commutatif
$$
\xymatrix{
X \ar[rr]_f \ar[rd]_p & & Y \ar[ld]^q \\
& Z
}
$$
d'espaces algébriques sur $S$. Supposons que $X \to Z$ soit localement de
type fini. Alors il existe un sous-espace ouvert $U(f) \subset X$
tel que $|U(f)| \subset |X|$ soit l'ensemble des points où $f$ est non ramifié.
De plus, pour tout morphisme d'espaces algébriques $Z' \to Z$, si
$f' : X' \to Y'$ est le changement de base de $f$ par $Z' \to Z$, alors
$U(f')$ est l'image inverse de $U(f)$ par la projection $X' \to X$.
\end{lemma}

\begin{proof}
Ce lemme est l'analogue de
Morphismes, lemme \ref{morphisms-lemma-set-points-where-fibres-unramified}
et nous allons en fait l'en déduire. D'après la
définition \ref{spaces-morphisms-definition-unramified},
l'ensemble $\{x \in |X| : f \text{ est non ramifié en }x\}$ est
ouvert dans $X$. Il suffit donc de démontrer la dernière assertion. D'après le
lemme \ref{spaces-morphisms-lemma-permanence-finite-type},
le morphisme $X \to Y$ est localement de type fini. D'après le
lemme \ref{spaces-morphisms-lemma-base-change-finite-type},
le morphisme $X' \to Y'$ est localement de type fini.

\medskip\noindent
Choisissons un schéma $W$ et un morphisme étale surjectif $W \to Z$.
Choisissons un schéma $V$ et un morphisme étale surjectif $V \to W \times_Z Y$.
Choisissons un schéma $U$ et un morphisme étale surjectif $U \to V \times_Y X$.
Enfin, choisissons un schéma $W'$ et un morphisme étale surjectif
$W' \to W \times_Z Z'$.
Posons $V' = W' \times_W V$ et $U' = W' \times_W U$ ; on obtient ainsi des
morphismes étales surjectifs $V' \to Y'$ et $U' \to X'$.
Nous utiliserons sans autre commentaire le fait qu'un morphisme étale d'espaces algébriques
induit une application ouverte entre les espaces topologiques associés (voir
Propriétés des espaces, lemme
\ref{spaces-properties-lemma-etale-open}).
Combiné avec le
lemme \ref{spaces-morphisms-lemma-unramified-local}, cela
implique que $U(f)$ est l'image dans $|X|$ de l'ensemble $T$ des points de $U$
où le morphisme $U \to V$ est non ramifié. De même, $U(f')$ est l'image
dans $|X'|$ de l'ensemble $T'$ des points de $U'$ où le morphisme $U' \to V'$
est non ramifié. Or, par construction, le diagramme
$$
\xymatrix{
U' \ar[r] \ar[d] & U \ar[d] \\
V' \ar[r] & V
}
$$
est cartésien (dans la catégorie des schémas). Le lemme précité
Morphismes, lemme \ref{morphisms-lemma-set-points-where-fibres-unramified}
montre donc que $T'$ est l'image inverse de $T$. Puisque
$|U'| \to |X'|$ est surjective, cela démontre le lemme.
\end{proof}

\begin{lemma}
\label{spaces-morphisms-lemma-permanence-unramified}
Soit $S$ un schéma.
Soient $X \to Y \to Z$ des morphismes d'espaces algébriques sur $S$.
Si $X \to Z$ est non ramifié, alors $X \to Y$ est non ramifié.
\end{lemma}

\begin{proof}
Choisissons un diagramme commutatif
$$
\xymatrix{
U \ar[d] \ar[r] & V \ar[d] \ar[r] & W \ar[d] \\
X \ar[r] & Y \ar[r] & Z
}
$$
dont les flèches verticales sont étales et surjectives. (Voir
Espaces, lemme \ref{spaces-lemma-lift-morphism-presentations}.)
Appliquer
Morphismes, lemme \ref{morphisms-lemma-unramified-permanence}
à la ligne supérieure.
\end{proof}














\section{Morphismes étales}
\label{spaces-morphisms-section-etale}

\noindent
La notion de morphisme étale d'espaces algébriques a été définie dans
Propriétés des espaces, définition \ref{spaces-properties-definition-etale}.
Voici ce que signifie, pour un morphisme, être étale en un point.

\begin{definition}
\label{spaces-morphisms-definition-etale}
Soit $S$ un schéma.
Soit $f : X \to Y$ un morphisme d'espaces algébriques sur $S$.
Soit $x \in |X|$. On dit que $f$ est {\it étale en $x$} s'il
existe un voisinage ouvert $X' \subset X$ de $x$ tel que
$f|_{X'} : X' \to Y$ soit étale.
\end{definition}

\begin{lemma}
\label{spaces-morphisms-lemma-etale-local}
Soit $S$ un schéma.
Soit $f : X \to Y$ un morphisme d'espaces algébriques sur $S$.
Les conditions suivantes sont équivalentes :
\begin{enumerate}
\item $f$ est étale ;
\item pour tout $x \in |X|$, le morphisme $f$ est étale en $x$ ;
\item pour tout schéma $Z$ et tout morphisme $Z \to Y$, le morphisme
$Z \times_Y X \to Z$ est étale ;
\item pour tout schéma affine $Z$ et tout morphisme
$Z \to Y$, le morphisme $Z \times_Y X \to Z$ est étale ;
\item il existe un schéma $V$ et un morphisme étale surjectif
$V \to Y$ tels que $V \times_Y X \to V$ soit un morphisme étale ;
\item il existe un schéma $U$ et un morphisme étale surjectif
$\varphi : U \to X$ tels que le composé $f \circ \varphi$
soit étale ;
\item pour tout diagramme commutatif
$$
\xymatrix{
U \ar[d] \ar[r] & V \ar[d] \\
X \ar[r] & Y
}
$$
où $U$ et $V$ sont des schémas et où les flèches verticales sont étales,
la flèche horizontale supérieure est étale ;
\item il existe un diagramme commutatif
$$
\xymatrix{
U \ar[d] \ar[r] & V \ar[d] \\
X \ar[r] & Y
}
$$
où $U$ et $V$ sont des schémas, où les flèches verticales sont étales et où
$U \to X$ est surjectif, tel que la flèche horizontale supérieure soit étale ;
\item il existe des recouvrements de Zariski $Y = \bigcup Y_i$ et
$f^{-1}(Y_i) = \bigcup X_{ij}$ tels que chaque morphisme
$X_{ij} \to Y_i$ soit étale.
\end{enumerate}
\end{lemma}

\begin{proof}
Combiner
Propriétés des espaces, lemmes
\ref{spaces-properties-lemma-etale-local},
\ref{spaces-properties-lemma-base-change-etale} et
\ref{spaces-properties-lemma-composition-etale}.
Certains détails sont omis.
\end{proof}

\begin{lemma}
\label{spaces-morphisms-lemma-composition-etale}
Le composé de deux morphismes étales d'espaces algébriques
est étale.
\end{lemma}

\begin{proof}
Il s'agit d'une reprise de
Propriétés des espaces, lemme \ref{spaces-properties-lemma-composition-etale}.
\end{proof}

\begin{lemma}
\label{spaces-morphisms-lemma-base-change-etale}
Le changement de base d'un morphisme étale d'espaces algébriques
par un morphisme quelconque d'espaces algébriques est étale.
\end{lemma}

\begin{proof}
Il s'agit d'une reprise de
Propriétés des espaces, lemme \ref{spaces-properties-lemma-base-change-etale}.
\end{proof}

\begin{lemma}
\label{spaces-morphisms-lemma-etale-locally-quasi-finite}
Un morphisme étale d'espaces algébriques est localement quasi-fini.
\end{lemma}

\begin{proof}
Soit $X \to Y$ un morphisme étale d'espaces algébriques ; voir
Propriétés des espaces, définition \ref{spaces-properties-definition-etale}.
D'après
Propriétés des espaces, lemme \ref{spaces-properties-lemma-etale-local}
on voit que cela signifie qu'il existe un diagramme comme dans le
lemme \ref{spaces-morphisms-lemma-local-source-target},
où $h$ est étale et où $a$ est une flèche verticale surjective. D'après
Morphismes, lemme \ref{morphisms-lemma-etale-locally-quasi-finite}
$h$ est localement quasi-fini. Ainsi, $X \to Y$ est localement quasi-fini
par définition.
\end{proof}

\begin{lemma}
\label{spaces-morphisms-lemma-etale-smooth}
Un morphisme étale d'espaces algébriques est lisse.
\end{lemma}

\begin{proof}
La démonstration est identique à celle du
lemme \ref{spaces-morphisms-lemma-etale-locally-quasi-finite}.
Elle utilise le fait qu'un morphisme étale de schémas est lisse
(par définition d'un morphisme étale de schémas).
\end{proof}

\begin{lemma}
\label{spaces-morphisms-lemma-etale-flat}
Un morphisme étale d'espaces algébriques est plat.
\end{lemma}

\begin{proof}
La démonstration est identique à celle du
lemme \ref{spaces-morphisms-lemma-etale-locally-quasi-finite}.
Elle utilise
Morphismes, lemme \ref{morphisms-lemma-etale-flat}.
\end{proof}

\begin{lemma}
\label{spaces-morphisms-lemma-etale-locally-finite-presentation}
\begin{slogan}
Étale implique localement de présentation finie.
\end{slogan}
Un morphisme étale d'espaces algébriques est localement de présentation finie.
\end{lemma}

\begin{proof}
La démonstration est identique à celle du
lemme \ref{spaces-morphisms-lemma-etale-locally-quasi-finite}.
Elle utilise
Morphismes, lemme \ref{morphisms-lemma-etale-locally-finite-presentation}.
\end{proof}

\begin{lemma}
\label{spaces-morphisms-lemma-etale-locally-finite-type}
Un morphisme étale d'espaces algébriques est localement de type fini.
\end{lemma}

\begin{proof}
Un morphisme étale est localement de présentation finie,
et un morphisme localement de présentation finie est localement de type fini ;
voir les
lemmes \ref{spaces-morphisms-lemma-etale-locally-finite-presentation} et
\ref{spaces-morphisms-lemma-finite-presentation-finite-type}.
\end{proof}

\begin{lemma}
\label{spaces-morphisms-lemma-etale-unramified}
Un morphisme étale d'espaces algébriques est non ramifié.
\end{lemma}

\begin{proof}
La démonstration est identique à celle du
lemme \ref{spaces-morphisms-lemma-etale-locally-quasi-finite}.
Elle utilise
Morphismes, lemme \ref{morphisms-lemma-etale-smooth-unramified}.
\end{proof}

\begin{lemma}
\label{spaces-morphisms-lemma-etale-permanence}
Soit $S$ un schéma. Soient $X, Y$ des espaces algébriques étales sur
un espace algébrique $Z$. Tout morphisme $X \to Y$ sur $Z$ est étale.
\end{lemma}

\begin{proof}
Il s'agit d'une reprise de
Propriétés des espaces, lemme \ref{spaces-properties-lemma-etale-permanence}.
\end{proof}

\begin{lemma}
\label{spaces-morphisms-lemma-unramified-flat-lfp-etale}
Un morphisme d'espaces algébriques localement de présentation finie,
plat et non ramifié est étale.
\end{lemma}

\begin{proof}
Soit $X \to Y$ un morphisme d'espaces algébriques localement de présentation
finie, plat et non ramifié. D'après
Propriétés des espaces, lemme \ref{spaces-properties-lemma-etale-local}
on voit que cela signifie qu'il existe un diagramme comme dans le
lemme \ref{spaces-morphisms-lemma-local-source-target},
où $h$ est localement de présentation finie, plat et non ramifié,
et où $a$ est une flèche verticale surjective. D'après
Morphismes, lemme \ref{morphisms-lemma-flat-unramified-etale}
$h$ est étale. Ainsi, $X \to Y$ est étale par définition.
\end{proof}






\section{Morphismes propres}
\label{spaces-morphisms-section-proper}

\noindent
La notion de morphisme propre joue un rôle important en géométrie algébrique.
Voici la définition d'un morphisme propre d'espaces algébriques.

\begin{definition}
\label{spaces-morphisms-definition-proper}
Soit $S$ un schéma.
Soit $f : X \to Y$ un morphisme d'espaces algébriques sur $S$.
On dit que $f$ est {\it propre} si $f$ est séparé, de type fini et
universellement fermé.
\end{definition}

\begin{lemma}
\label{spaces-morphisms-lemma-proper-local}
Soit $S$ un schéma. Soit $f : X \to Y$ un morphisme d'espaces algébriques
sur $S$. Les conditions suivantes sont équivalentes :
\begin{enumerate}
\item $f$ est propre ;
\item pour tout schéma $Z$ et tout morphisme $Z \to Y$,
la projection $Z \times_Y X \to Z$ est propre ;
\item pour tout schéma affine $Z$ et tout morphisme $Z \to Y$,
la projection $Z \times_Y X \to Z$ est propre ;
\item il existe un schéma $V$ et un morphisme étale surjectif
$V \to Y$ tel que $V \times_Y X \to V$ soit propre ;
\item il existe un recouvrement de Zariski $Y = \bigcup Y_i$ tel que
chacun des morphismes $f^{-1}(Y_i) \to Y_i$ soit propre.
\end{enumerate}
\end{lemma}

\begin{proof}
Combiner les lemmes \ref{spaces-morphisms-lemma-separated-local},
\ref{spaces-morphisms-lemma-finite-type-local},
\ref{spaces-morphisms-lemma-quasi-compact-local} et
\ref{spaces-morphisms-lemma-universally-closed-local}.
\end{proof}

\begin{lemma}
\label{spaces-morphisms-lemma-base-change-proper}
Un changement de base d'un morphisme propre est propre.
\end{lemma}

\begin{proof}
Voir les
lemmes \ref{spaces-morphisms-lemma-base-change-separated},
\ref{spaces-morphisms-lemma-base-change-finite-type} et
\ref{spaces-morphisms-lemma-base-change-universally-closed}.
\end{proof}

\begin{lemma}
\label{spaces-morphisms-lemma-composition-proper}
Un composé de morphismes propres est propre.
\end{lemma}

\begin{proof}
Voir les
lemmes \ref{spaces-morphisms-lemma-composition-separated},
\ref{spaces-morphisms-lemma-composition-finite-type} et
\ref{spaces-morphisms-lemma-composition-universally-closed}.
\end{proof}

\begin{lemma}
\label{spaces-morphisms-lemma-closed-immersion-proper}
Une immersion fermée d'espaces algébriques est un morphisme propre
d'espaces algébriques.
\end{lemma}

\begin{proof}
Comme une immersion fermée est représentable par définition, cela résulte de
Espaces,
lemme \ref{spaces-lemma-representable-transformations-property-implication}
et du résultat correspondant pour les morphismes de schémas ; voir
Morphismes, lemme \ref{morphisms-lemma-closed-immersion-proper}.
\end{proof}

\begin{lemma}
\label{spaces-morphisms-lemma-universally-closed-permanence}
Soit $S$ un schéma.
Considérons un diagramme commutatif d'espaces algébriques
$$
\xymatrix{
X \ar[rr] \ar[rd] & &
Y \ar[ld] \\
& B &
}
$$
sur $S$.
\begin{enumerate}
\item Si $X \to B$ est universellement fermé et si $Y \to B$ est
séparé, alors le morphisme $X \to Y$ est universellement fermé.
En particulier, l'image de $|X|$ dans $|Y|$ est fermée.
\item Si $X \to B$ est propre et si $Y \to B$ est séparé, alors
le morphisme $X \to Y$ est propre.
\end{enumerate}
\end{lemma}

\begin{proof}
Supposons que $X \to B$ soit universellement fermé et que $Y \to B$ soit séparé.
Factorisons le morphisme en $X \to X \times_B Y \to Y$.
Le premier morphisme est une immersion fermée ; voir le
lemme \ref{spaces-morphisms-lemma-semi-diagonal} ;
il est donc universellement fermé.
La projection $X \times_B Y \to Y$ est le changement de base
d'un morphisme universellement fermé et est donc
universellement fermée ; voir le
lemme \ref{spaces-morphisms-lemma-base-change-universally-closed}.
Ainsi, $X \to Y$ est universellement fermé en tant que composé
de morphismes universellement fermés ; voir le
lemme \ref{spaces-morphisms-lemma-composition-universally-closed}.
Cela démontre (1). Pour en déduire (2), combiner (1) avec les
lemmes \ref{spaces-morphisms-lemma-compose-after-separated},
\ref{spaces-morphisms-lemma-quasi-compact-permanence} et
\ref{spaces-morphisms-lemma-permanence-finite-type}.
\end{proof}

\begin{lemma}
\label{spaces-morphisms-lemma-image-proper-is-proper}
Soit $S$ un schéma. Soit $B$ un espace algébrique sur $S$.
Soit $f : X \to Y$ un morphisme d'espaces algébriques sur $B$.
Si $X$ est universellement fermé sur $B$ et si $f$ est surjectif, alors
$Y$ est universellement fermé sur $B$. En particulier, si $Y$ est aussi
séparé et de type fini sur $B$, alors $Y$ est propre sur $B$.
\end{lemma}

\begin{proof}
Supposons que $X$ soit universellement fermé et que $f$ soit surjectif.
Notons $p : X \to B$ et $q : Y \to B$ les morphismes structuraux.
Soit $B' \to B$ un morphisme d'espaces algébriques sur $S$.
Le changement de base $f' : X_{B'} \to Y_{B'}$ est surjectif
(lemme \ref{spaces-morphisms-lemma-base-change-surjective}), et le changement de
base $p' : X_{B'} \to B'$ est fermé.
Si $T \subset Y_{B'}$ est fermé, alors $(f')^{-1}(T) \subset X_{B'}$
est fermé ; par conséquent, $p'((f')^{-1}(T)) = q'(T)$ est fermé.
Ainsi, $q'$ est fermé.
\end{proof}

\begin{lemma}
\label{spaces-morphisms-lemma-scheme-theoretic-image-is-proper}
Soit $S$ un schéma. Soit
$$
\xymatrix{
X \ar[rr]_h \ar[rd]_f & & Y \ar[ld]^g \\
& B
}
$$
un diagramme commutatif de morphismes d'espaces algébriques sur $S$.
Supposons que
\begin{enumerate}
\item $X \to B$ soit un morphisme propre ;
\item $Y \to B$ soit séparé et localement de type fini.
\end{enumerate}
Alors l'image schématique $Z \subset Y$ de $h$
est propre sur $B$ et $X \to Z$ est surjectif.
\end{lemma}

\begin{proof}
L'image schématique de $h$ est construite à la section
\ref{spaces-morphisms-section-scheme-theoretic-image}.
Observons que $h$ est quasi-compact
(lemme \ref{spaces-morphisms-lemma-quasi-compact-quasi-separated-permanence})
et que, par conséquent, $|h|(|X|) \subset |Z|$
est dense (lemme \ref{spaces-morphisms-lemma-quasi-compact-scheme-theoretic-image}).
D'autre part, $|h|(|X|)$ est fermé dans $|Y|$
(lemme \ref{spaces-morphisms-lemma-universally-closed-permanence})
et donc $X \to Z$ est surjectif.
Ainsi, $Z \to B$ est propre (lemme \ref{spaces-morphisms-lemma-image-proper-is-proper}).
\end{proof}

\begin{lemma}
\label{spaces-morphisms-lemma-separated-diagonal-proper}
Soit $S$ un schéma.
Soit $f : X \to Y$ un morphisme d'espaces algébriques sur $S$.
Les conditions suivantes sont équivalentes :
\begin{enumerate}
\item $f$ est séparé ;
\item $\Delta_{X/Y} : X \to X \times_Y X$ est universellement fermé ;
\item $\Delta_{X/Y} : X \to X \times_Y X$ est propre.
\end{enumerate}
\end{lemma}

\begin{proof}
L'implication (1) $\Rightarrow$ (3) résulte du
lemme \ref{spaces-morphisms-lemma-closed-immersion-proper}.
Nous utiliserons
Espaces, lemme
\ref{spaces-lemma-representable-transformations-property-implication}
sans autre mention dans la suite de la démonstration.
Rappelons que $\Delta_{X/Y}$ est un monomorphisme
représentable qui est localement de type fini ; voir le
lemme \ref{spaces-morphisms-lemma-properties-diagonal}.
Puisque propre $\Rightarrow$ universellement fermé pour les morphismes de schémas,
on conclut que (3) implique (2).
Si $\Delta_{X/Y}$ est universellement fermé, alors
\'Morphismes étales,
lemme \ref{etale-lemma-characterize-closed-immersion}
implique que c'est une immersion fermée. Ainsi, (2) $\Rightarrow$ (1),
ce qui achève la démonstration.
\end{proof}





\section{Critères valuatifs}
\label{spaces-morphisms-section-valuative}

\noindent
Cette section introduit les notions de base sur les critères valuatifs pour les morphismes
d'espaces algébriques. Voici une liste de références vers des résultats complémentaires :
\begin{enumerate}
\item le critère valuatif de fermeture universelle se trouve à la
section \ref{spaces-morphisms-section-valuative-criterion-universally-closed} ;
\item le critère valuatif de séparation se trouve à la
section \ref{spaces-morphisms-section-valuative-separatedness} ;
\item le critère valuatif de propreté se trouve à la
section \ref{spaces-morphisms-section-valuative-criterion-properness} ;
\item des réciproques supplémentaires se trouvent dans
Espaces décents, section
\ref{decent-spaces-section-valuative-criterion-universally-closed}
et Espaces décents, lemme
\ref{decent-spaces-lemma-re-characterize-universally-closed} ;
\item dans le cas noethérien, il suffit de vérifier le critère pour les
anneaux de valuation discrète, comme le montrent Cohomologie des espaces, section
\ref{spaces-cohomology-section-Noetherian-valuative-criterion} et
Limites d'espaces, section
\ref{spaces-limits-section-Noetherian-valuative-criterion} ;
\item des versions affinées des critères valuatifs dans le cas noethérien
se trouvent dans Limites d'espaces, section
\ref{spaces-limits-section-refined-valuative-criteria}.
\end{enumerate}
Nous énonçons d'abord formellement la définition, puis nous expliquons en quoi elle
diffère du cas des morphismes de schémas.

\begin{definition}
\label{spaces-morphisms-definition-valuative-criterion}
Soit $S$ un schéma.
Soit $f : X \to Y$ un morphisme d'espaces algébriques sur $S$.
On dit que $f$ {\it satisfait la partie d'unicité du critère valuatif}
si, pour tout diagramme commutatif en traits pleins
$$
\xymatrix{
\Spec(K) \ar[r] \ar[d] & X \ar[d] \\
\Spec(A) \ar[r] \ar@{-->}[ru] & Y
}
$$
où $A$ est un anneau de valuation de corps des fractions $K$, il existe
au plus une flèche pointillée (sans en exiger l'existence).
On dit que $f$ {\it satisfait la partie d'existence du critère valuatif}
si, pour tout diagramme en traits pleins comme ci-dessus, il existe une extension
de corps $K'/K$, un anneau de valuation $A' \subset K'$ qui domine
$A$ et un morphisme $\Spec(A') \to X$ tels que le diagramme
suivant soit commutatif :
$$
\xymatrix{
\Spec(K') \ar[r] \ar[d] & \Spec(K) \ar[r] & X \ar[d] \\
\Spec(A') \ar[r] \ar[rru] & \Spec(A) \ar[r] & Y
}
$$
On dit que $f$ {\it satisfait le critère valuatif}
si $f$ satisfait à la fois les parties d'existence et d'unicité.
\end{definition}

\noindent
La formulation de la partie d'existence du critère valuatif est
légèrement différente pour les morphismes d'espaces algébriques, puisqu'il peut être
nécessaire d'étendre le corps des fractions de l'anneau de valuation.
En pratique, cette différence ne joue presque jamais de rôle.
\begin{enumerate}
\item La vérification de la partie d'unicité du critère valuatif ne
fait jamais intervenir d'extension du corps des fractions ; elle est donc exactement la même
que dans le cas des schémas.
\item Il est nécessaire d'autoriser en général des extensions de corps ; voir
l'exemple \ref{spaces-morphisms-example-finite-separable-needed}.
\item Pour les morphismes d'espaces algébriques, il suffit toujours de
prendre une extension finie séparable $K'/K$ dans la partie d'existence
du critère valuatif ; voir le lemme \ref{spaces-morphisms-lemma-finite-separable-enough}.
\item Si $f : X \to Y$ est un morphisme séparé d'espaces algébriques, on
peut toujours prendre $K = K'$ lors de la vérification de la partie d'existence du
critère valuatif ; voir le lemme \ref{spaces-morphisms-lemma-usual-enough}.
\item Pour un morphisme quasi-compact et quasi-séparé
$f : X \to Y$, on obtient l'équivalence entre
« $f$ est séparé et universellement fermé » et « $f$ satisfait
le critère valuatif usuel » ; voir le
lemme \ref{spaces-morphisms-lemma-characterize-separated-and-universally-closed}.
Le critère valuatif de propreté est le critère usuel ; voir le
lemme \ref{spaces-morphisms-lemma-characterize-proper}.
\end{enumerate}
Comme première étape de la théorie, montrons que le critère est identique
à celui formulé pour les morphismes de schémas
lorsque le morphisme d'espaces algébriques est représentable.

\begin{lemma}
\label{spaces-morphisms-lemma-valuative-criterion-representable}
Soit $S$ un schéma.
Soit $f : X \to Y$ un morphisme d'espaces algébriques sur $S$.
Supposons que $f$ soit représentable. Les conditions suivantes sont équivalentes :
\begin{enumerate}
\item $f$ satisfait la partie d'existence du critère valuatif
au sens de la définition \ref{spaces-morphisms-definition-valuative-criterion} ;
\item pour tout diagramme commutatif en traits pleins
$$
\xymatrix{
\Spec(K) \ar[r] \ar[d] & X \ar[d] \\
\Spec(A) \ar[r] \ar@{-->}[ru] & Y
}
$$
où $A$ est un anneau de valuation de corps des fractions $K$, il existe
une flèche pointillée, c'est-à-dire que $f$ satisfait la partie d'existence du critère
valuatif au sens de
Schémas, définition \ref{schemes-definition-valuative-criterion}.
\end{enumerate}
\end{lemma}

\begin{proof}
Il suffit de montrer que, pour tout diagramme commutatif de la forme
$$
\xymatrix{
\Spec(K') \ar[r] \ar[d] & \Spec(K) \ar[r] & X \ar[d] \\
\Spec(A') \ar[r] \ar[rru]^\varphi & \Spec(A) \ar[r] & Y
}
$$
comme dans la définition \ref{spaces-morphisms-definition-valuative-criterion}, on peut aussi
trouver un morphisme $\Spec(A) \to X$ qui s'insère dans le diagramme.
Posons $X_A = \Spec(A) \times_Y X$. Comme $f$ est représentable, on voit
que $X_A$ est un schéma. Le morphisme $\varphi$ donne un morphisme
$\varphi' : \Spec(A') \to X_A$. Soit $x \in X_A$ l'image du
point fermé par $\varphi' : \Spec(A') \to X_A$. On dispose alors du
diagramme commutatif d'anneaux suivant :
$$
\xymatrix{
K' & K \ar[l] & \mathcal{O}_{X_A, x} \ar[l] \ar[lld] \\
A' \ar[u] & A \ar[l] & A \ar[l] \ar[u]
}
$$
Puisque $A$ est un anneau de valuation et que $A'$ domine $A$, on voit
que $K \cap A' = A$. L'image de l'homomorphisme d'anneaux $\mathcal{O}_{X_A, x} \to K$
est donc contenue dans $A$. On obtient ainsi un morphisme $\Spec(A) \to X_A$ (voir
Schémas, section \ref{schemes-section-points})
comme souhaité.
\end{proof}

\begin{lemma}
\label{spaces-morphisms-lemma-finite-separable-enough}
Soit $S$ un schéma.
Soit $f : X \to Y$ un morphisme d'espaces algébriques sur $S$.
Les conditions suivantes sont équivalentes :
\begin{enumerate}
\item $f$ satisfait la partie d'existence du critère valuatif
au sens de la définition \ref{spaces-morphisms-definition-valuative-criterion} ;
\item $f$ satisfait la partie d'existence du critère valuatif
au sens de la définition \ref{spaces-morphisms-definition-valuative-criterion}, modifiée en
exigeant que l'extension $K'/K$ soit finie séparable.
\end{enumerate}
\end{lemma}

\begin{proof}
Il faut montrer que (1) implique (2). Supposons donné un diagramme
$$
\xymatrix{
\Spec(K') \ar[r] \ar[d] & \Spec(K) \ar[r] & X \ar[d] \\
\Spec(A') \ar[r] \ar[rru] & \Spec(A) \ar[r] & Y
}
$$
comme dans la définition \ref{spaces-morphisms-definition-valuative-criterion}, avec $K \subset K'$
arbitraire. Choisissons un schéma $U$ et un morphisme étale surjectif $U \to X$.
Alors
$$
\Spec(A') \times_X U \longrightarrow \Spec(A')
$$
est étale surjectif. Soit $p$ un point de $\Spec(A') \times_X U$
qui s'envoie sur le point fermé de $\Spec(A')$. Soit $p' \leadsto p$
une générisation de $p$ qui s'envoie sur le point générique de $\Spec(A')$.
Une telle générisation existe parce que les générisations se relèvent le long des morphismes
plats de schémas ; voir
Morphismes, lemme \ref{morphisms-lemma-generalizations-lift-flat}.
Alors $p'$ correspond à un point du schéma $\Spec(K') \times_X U$.
Notons que
$$
\Spec(K') \times_X U
=
\Spec(K') \times_{\Spec(K)} (\Spec(K) \times_X U)
$$
Par conséquent, $p'$ s'envoie sur un point $q' \in \Spec(K) \times_X U$ dont le
corps résiduel est une extension finie séparable de $K$. Enfin,
$p' \leadsto p$ s'envoie sur une spécialisation $u' \leadsto u$ dans le
schéma $U$. Avec toutes ces notations, on obtient le diagramme d'anneaux
suivant :
$$
\xymatrix{
\kappa(p') & & \kappa(q') \ar[ll] & \kappa(u') \ar[l] \\
& \mathcal{O}_{\Spec(A') \times_X U, p} \ar[lu] & &
\mathcal{O}_{U, u} \ar[ll] \ar[u] \\
K' \ar[uu] & A' \ar[l] \ar[u] & A \ar[l] \ar'[u][uu]
}
$$
Cela signifie que l'anneau $B \subset \kappa(q')$ engendré par
les images de $A$ et de $\mathcal{O}_{U, u}$ s'envoie dans un sous-anneau
de $\kappa(p')$ contenu dans l'image $B'$ de
$\mathcal{O}_{\Spec(A') \times_X U, p} \to \kappa(p')$.
Notons que $B'$ est un anneau local. Soit $\mathfrak m \subset B'$
son idéal maximal. Par construction, $A \cap \mathfrak m$
(resp.\ $\mathcal{O}_{U, u} \cap \mathfrak m$, resp.\ $A' \cap \mathfrak m$)
est l'idéal maximal de $A$ (resp.\ de $\mathcal{O}_{U, u}$, resp.\ de $A'$).
Posons $\mathfrak q = B \cap \mathfrak m$. C'est un
idéal premier tel que $A \cap \mathfrak q$ soit l'idéal maximal de $A$.
Ainsi, $B_{\mathfrak q} \subset \kappa(q')$ est un anneau local qui domine
$A$. D'après
Algèbre, lemme \ref{algebra-lemma-dominate}
on peut trouver un anneau de valuation $A_1 \subset \kappa(q')$
de corps des fractions $\kappa(q')$
qui domine $B_{\mathfrak q}$. L'homomorphisme (local) d'anneaux
$\mathcal{O}_{U, u} \to A_1$ donne un morphisme
$\Spec(A_1) \to U \to X$
tel que le diagramme
$$
\xymatrix{
\Spec(\kappa(q')) \ar[r] \ar[d] & \Spec(K) \ar[r] & X \ar[d] \\
\Spec(A_1) \ar[r] \ar[rru] & \Spec(A) \ar[r] & Y
}
$$
soit commutatif. Puisque le corps des fractions de $A_1$ est $\kappa(q')$
et que $\kappa(q')/K$ est une extension finie
séparable par construction, le lemme est démontré.
\end{proof}

\begin{lemma}
\label{spaces-morphisms-lemma-push-down-solution}
Soit $S$ un schéma. Soit $f : X \to Y$ un morphisme séparé d'espaces
algébriques sur $S$. Supposons donné un diagramme
$$
\xymatrix{
\Spec(K') \ar[r] \ar[d] & \Spec(K) \ar[r] & X \ar[d] \\
\Spec(A') \ar[r] \ar[rru] & \Spec(A) \ar[r] \ar@{-->}[ru] & Y
}
$$
comme dans la définition \ref{spaces-morphisms-definition-valuative-criterion}, avec $K \subset K'$
arbitraire. Alors il existe une flèche pointillée qui rend le diagramme commutatif.
\end{lemma}

\begin{proof}
Il faut montrer que l'on peut trouver un morphisme $\Spec(A) \to X$ qui s'insère
dans le diagramme.

\medskip\noindent
Considérons le changement de base $X_A = \Spec(A) \times_Y X$ de $X$.
Alors $X_A \to \Spec(A)$ est un morphisme séparé d'espaces
algébriques (lemme \ref{spaces-morphisms-lemma-base-change-separated}). En changeant la base de
tous les morphismes du diagramme ci-dessus, on obtient
$$
\xymatrix{
\Spec(K') \ar[r] \ar[d] & \Spec(K) \ar[r] & X_A \ar[d] \\
\Spec(A') \ar[r] \ar[rru] & \Spec(A) \ar@{=}[r] & \Spec(A)
}
$$
On peut donc remplacer $X$ par $X_A$, supposer que $Y = \Spec(A)$ et
que l'on dispose d'un diagramme comme ci-dessus. On remplace alors $X$ par
un sous-espace ouvert quasi-compact contenant l'image de $|\Spec(A')| \to |X|$.

\medskip\noindent
Le morphisme $\Spec(A') \to X$ est quasi-compact d'après le
lemme \ref{spaces-morphisms-lemma-quasi-compact-permanence}.
Soit $Z \subset X$ l'image schématique de
$\Spec(A') \to X$. Alors $Z$ est un espace algébrique réduit
(lemme \ref{spaces-morphisms-lemma-scheme-theoretic-image-reduced}),
quasi-compact (en tant que sous-espace fermé de $X$) et séparé
(en tant que sous-espace fermé de $X$) sur $A$.
Considérons le changement de base
$$
\Spec(K') = \Spec(A') \times_{\Spec(A)} \Spec(K) \to
X \times_{\Spec(A)} \Spec(K) = X_K
$$
du morphisme $\Spec(A') \to X$ par le morphisme plat de schémas
$\Spec(K) \to \Spec(A)$. D'après le
lemme \ref{spaces-morphisms-lemma-flat-base-change-scheme-theoretic-image},
on voit que l'image schématique de ce morphisme
est le changement de base $Z_K$ de $Z$. D'autre part, par
hypothèse (c'est-à-dire d'après le diagramme commutatif ci-dessus),
ce morphisme se factorise par un morphisme
$\Spec(K) \to Z_K$ qui est une section du morphisme structural
$Z_K \to \Spec(K)$. Comme $Z_K$ est séparé, cette section est une
immersion fermée (lemme \ref{spaces-morphisms-lemma-section-immersion}).
On conclut que $Z_K = \Spec(K)$.

\medskip\noindent
Soit $V \to Z$ un morphisme étale surjectif avec $V$ un
schéma affine (Propriétés des espaces, lemme
\ref{spaces-properties-lemma-quasi-compact-affine-cover}).
Écrivons $V = \Spec(B)$. Alors $V \times_Z \Spec(A') = \Spec(C)$
est affine puisque $Z$ est séparé. Notons que $B \to C$ est injectif,
car $V$ est l'image schématique de $V \times_Z \Spec(A') \to V$
d'après le lemme \ref{spaces-morphisms-lemma-quasi-compact-scheme-theoretic-image}.
D'autre part, $A' \to C$ est étale, car il correspond
au changement de base de $V \to Z$.
Puisque $A'$ est un $A$-module sans torsion,
la platitude de $A' \to C$ implique que $C$ est un
$A$-module sans torsion ; ainsi, $B$ est un $A$-module sans torsion.
Notons qu'être sans torsion comme $A$-module équivaut
à être plat (Compléments d'algèbre, lemme
\ref{more-algebra-lemma-valuation-ring-torsion-free-flat}).
Ensuite, écrivons
$$
V \times_Z V = \Spec(B')
$$
Notons que les deux homomorphismes d'anneaux $B \to B'$ sont étales puisque $V \to Z$
est étale. L'homomorphisme canonique surjectif $B \otimes_A B \to B'$
devient un isomorphisme après tensorisation par $K$ sur $A$,
car $Z_K = \Spec(K)$. Cependant, $B \otimes_A B$ est sans torsion
comme $A$-module d'après les remarques précédentes. Ainsi, $B' = B \otimes_A B$.
Il en résulte que le changement de base de l'homomorphisme d'anneaux
$A \to B$ par l'homomorphisme fidèlement plat $A \to B$
est étale (notons que $\Spec(B) \to \Spec(A)$ est surjectif,
puisque $Z \to \Spec(A)$ est surjectif). Ainsi, $A \to B$ est
étale (Descente, lemme \ref{descent-lemma-descending-property-etale}),
autrement dit, $V \to \Spec(A)$ est étale.
Comme $V \times_Z V = V \times_{\Spec(A)} V$,
on conclut que $Z = \Spec(A)$ en tant qu'espaces algébriques
(par exemple d'après Espaces, lemme \ref{spaces-lemma-space-presentation}),
ce qui achève la démonstration.
\end{proof}

\begin{lemma}
\label{spaces-morphisms-lemma-usual-enough}
Soit $S$ un schéma. Soit $f : X \to Y$ un morphisme séparé d'espaces
algébriques sur $S$. Les conditions suivantes sont équivalentes :
\begin{enumerate}
\item $f$ satisfait la partie d'existence du critère valuatif
au sens de la définition \ref{spaces-morphisms-definition-valuative-criterion} ;
\item pour tout diagramme commutatif en traits pleins
$$
\xymatrix{
\Spec(K) \ar[r] \ar[d] & X \ar[d] \\
\Spec(A) \ar[r] \ar@{-->}[ru] & Y
}
$$
où $A$ est un anneau de valuation de corps des fractions $K$, il existe
une flèche pointillée, c'est-à-dire que $f$ satisfait la partie d'existence du critère
valuatif au sens de
Schémas, définition \ref{schemes-definition-valuative-criterion}.
\end{enumerate}
\end{lemma}

\begin{proof}
Il faut montrer que (1) implique (2). Supposons donné un diagramme commutatif
$$
\xymatrix{
\Spec(K) \ar[r] \ar[d] & X \ar[d] \\
\Spec(A) \ar[r] & Y
}
$$
comme dans la partie (2). D'après (1), il existe un diagramme commutatif
$$
\xymatrix{
\Spec(K') \ar[r] \ar[d] & \Spec(K) \ar[r] & X \ar[d] \\
\Spec(A') \ar[r] \ar[rru] & \Spec(A) \ar[r] & Y
}
$$
comme dans la définition \ref{spaces-morphisms-definition-valuative-criterion}, avec $K \subset K'$
arbitraire. D'après le lemme \ref{spaces-morphisms-lemma-push-down-solution}, on peut trouver un morphisme
$\Spec(A) \to X$ qui s'insère dans le diagramme, c'est-à-dire que (2) est satisfaite.
\end{proof}

\begin{example}
\label{spaces-morphisms-example-finite-separable-needed}
Considérons l'espace algébrique $X$ construit dans
Espaces, exemple \ref{spaces-example-non-representable-descent}.
Rappelons qu'il s'agit d'une forme tordue galoisienne de la droite affine à origine dédoublée.
La torsion galoisienne est relative à une extension galoisienne de degré deux
$k'/k$ de corps. Elle est ainsi munie d'un morphisme
$$
\pi : X \longrightarrow S = \mathbf{A}^1_k
$$
qui est quasi-compact. Affirmons que $\pi$ est universellement fermé.
En effet, après le changement de base par $\Spec(k') \to \Spec(k)$,
le morphisme $\pi$ s'identifie au morphisme
$$
\text{droite affine à origine dédoublée}
\longrightarrow
\text{droite affine}
$$
qui est universellement fermé (certains détails sont omis). Puisque le morphisme
$\Spec(k') \to \Spec(k)$ est universellement fermé et
surjectif, une chasse au diagramme montre que $\pi$ est universellement fermé.
D'autre part, considérons le diagramme
$$
\xymatrix{
\Spec(k((x))) \ar[r] \ar[d] & X \ar[d]^\pi \\
\Spec(k[[x]]) \ar[r] \ar@{..>}[ru] & \mathbf{A}^1_k
}
$$
Puisque l'unique point de $X$ au-dessus de $0 \in \mathbf{A}^1_k$
correspond à un monomorphisme $\Spec(k') \to X$,
il est clair qu'il ne peut exister de flèche pointillée ! Cela montre qu'une
extension finie séparable de corps est nécessaire en général.
\end{example}

\begin{lemma}
\label{spaces-morphisms-lemma-base-change-valuative-criteria}
Le changement de base d'un morphisme d'espaces algébriques qui satisfait la
partie d'existence (resp.\ la partie d'unicité) du critère valuatif
par un morphisme quelconque d'espaces algébriques satisfait la
partie d'existence (resp.\ la partie d'unicité) du critère valuatif.
\end{lemma}

\begin{proof}
Soit $f : X \to Y$ un morphisme d'espaces algébriques sur le schéma $S$.
Soit $Z \to Y$ un morphisme quelconque d'espaces algébriques sur $S$.
Considérons un diagramme commutatif en traits pleins de la forme suivante :
$$
\xymatrix{
\Spec(K) \ar[r] \ar[d] & Z \times_Y X \ar[r] \ar[d] & X \ar[d] \\
\Spec(A) \ar[r] \ar@{-->}[ru] \ar@{-->}[rru] & Z \ar[r] & Y
}
$$
L'ensemble des flèches pointillées nord-ouest qui rendent le diagramme commutatif
est en correspondance biunivoque avec l'ensemble des flèches pointillées ouest-nord-ouest
qui le rendent commutatif. Cela démontre le lemme dans le cas de
« l'unicité ». Pour la partie d'existence, supposons que $f$ satisfasse la partie d'existence
du critère valuatif. Si l'on se donne un diagramme commutatif en traits pleins
comme ci-dessus, il existe par hypothèse une extension de corps $K'/K$,
un anneau de valuation $A' \subset K'$ qui domine $A$ et
un morphisme $\Spec(A') \to X$ qui s'insère dans le diagramme commutatif
suivant :
$$
\xymatrix{
\Spec(K') \ar[r] \ar[d] &
\Spec(K) \ar[r] & Z \times_Y X \ar[r] & X \ar[d] \\
\Spec(A') \ar[r] \ar[rrru] & \Spec(A) \ar[r] & Z \ar[r] & Y
}
$$
Et, d'après les remarques précédentes, la flèche oblique correspond à une flèche
$\Spec(A') \to Z \times_Y X$, comme souhaité.
\end{proof}

\begin{lemma}
\label{spaces-morphisms-lemma-composition-valuative-criteria}
Le composé de deux morphismes d'espaces algébriques qui satisfont la
partie d'existence (resp.\ la partie d'unicité) du critère valuatif
satisfait la partie d'existence (resp.\ la partie d'unicité) du critère
valuatif.
\end{lemma}

\begin{proof}
Soient $f : X \to Y$ et $g : Y \to Z$ des morphismes d'espaces algébriques sur le
schéma $S$. Considérons un diagramme commutatif en traits pleins de la forme suivante :
$$
\xymatrix{
\Spec(K) \ar[dd] \ar[r] & X \ar[d]^f \\
& Y \ar[d]^g \\
\Spec(A) \ar[r] \ar@{-->}[ru] \ar@{-->}[ruu] & Z
}
$$
Si l'on dispose de la partie d'unicité pour $g$, il existe au
plus une flèche pointillée nord-ouest qui rende le diagramme commutatif.
Si l'on dispose aussi de la partie d'unicité pour $f$, il existe
au plus une flèche pointillée nord-nord-ouest qui rende le diagramme
commutatif. Dans le cas de l'existence, la démonstration résulte de l'examen
du diagramme suivant :
$$
\xymatrix{
\Spec(K'') \ar[r] \ar[dd] &
\Spec(K') \ar[r] &
\Spec(K) \ar[r] &
X \ar[d]^f \\
& & & Y \ar[d]^g \\
\Spec(A'') \ar[r] \ar[rrruu] &
\Spec(A') \ar[r] \ar[rru] &
\Spec(A) \ar[r] &
Z
}
$$
En effet, la partie d'existence pour $g$ fournit l'extension $K'$, l'anneau
de valuation $A'$ et la flèche $\Spec(A') \to Y$, après quoi
la partie d'existence pour $f$ fournit l'extension $K''$, l'anneau
de valuation $A''$ et la flèche $\Spec(A'') \to X$.
\end{proof}






\section{Critère valuatif de fermeture universelle}
\label{spaces-morphisms-section-valuative-criterion-universally-closed}

\noindent
La partie d'existence du critère valuatif entraîne qu'un morphisme
quasi-compact est universellement fermé ; voir le
lemme \ref{spaces-morphisms-lemma-quasi-compact-existence-universally-closed}.
Dans le cas des schémas, il s'agit d'une équivalence,
mais cet énoncé est faux pour les morphismes d'espaces algébriques.
L'exemple \ref{spaces-morphisms-example-strange-universally-closed}
montre que $\mathbf{A}^1_k/\mathbf{Z} \to \Spec(k)$ est universellement
fermé, alors qu'il est facile de voir que la partie d'existence
du critère valuatif n'est pas satisfaite. Nous revenons sur ce sujet dans
Espaces décents, section
\ref{decent-spaces-section-valuative-criterion-universally-closed}
et y montrons que la réciproque vaut lorsque la source du morphisme est un espace décent
(voir aussi
Espaces décents,
lemme \ref{decent-spaces-lemma-re-characterize-universally-closed},
pour une version relative).

\begin{lemma}
\label{spaces-morphisms-lemma-quasi-compact-existence-universally-closed}
Soit $S$ un schéma.
Soit $f : X \to Y$ un morphisme d'espaces algébriques sur $S$.
Supposons que
\begin{enumerate}
\item $f$ soit quasi-compact, et que
\item $f$ satisfasse la partie d'existence du critère valuatif.
\end{enumerate}
Alors $f$ est universellement fermé.
\end{lemma}

\begin{proof}
D'après les lemmes \ref{spaces-morphisms-lemma-base-change-quasi-compact}
et \ref{spaces-morphisms-lemma-base-change-valuative-criteria},
les propriétés (1) et (2) sont préservées par
tout changement de base. D'après le lemme \ref{spaces-morphisms-lemma-universally-closed-local},
il suffit de montrer que $|T \times_Y X| \to |T|$ est fermé
chaque fois que $T$ est un schéma affine sur $S$ muni d'un morphisme vers $Y$. Il
suffit donc de prouver ceci : si $Y$ est un schéma affine, si $f : X \to Y$ est quasi-compact
et satisfait la partie d'existence du critère valuatif, alors
$f : |X| \to |Y|$ est fermé. Dans cette situation, $X$ est un espace
algébrique quasi-compact. D'après
Propriétés des espaces,
lemme \ref{spaces-properties-lemma-quasi-compact-affine-cover},
il existe un schéma affine $U$ et un morphisme étale surjectif
$\varphi : U \to X$. Soit $T \subset |X|$ fermé. L'image inverse
$\varphi^{-1}(T) \subset U$ est fermée ; c'est donc l'ensemble des points
d'un sous-schéma fermé affine $Z \subset U$. Ainsi, d'après
Algèbre, lemme \ref{algebra-lemma-image-stable-specialization-closed}
la partie $f(T) = f(\varphi(|Z|)) \subset |Y|$ est fermée dès qu'elle est
stable par spécialisation.

\medskip\noindent
Soit $y' \leadsto y$ une spécialisation dans $Y$ telle que $y' \in f(T)$.
Choisissons un point $x' \in T \subset |X|$ qui s'envoie sur $y'$ par $f$.
On peut représenter $x'$ par un morphisme $\Spec(K) \to X$
pour un certain corps $K$. On obtient donc le diagramme suivant
$$
\xymatrix{
\Spec(K) \ar[r]_-{x'} \ar[d] & X \ar[d]^f \\
\Spec(\mathcal{O}_{Y, y}) \ar[r] & Y,
}
$$
voir
Schémas, section \ref{schemes-section-points}
pour l'existence de la flèche verticale de gauche.
Choisissons un anneau de valuation $A \subset K$ qui domine l'image de
l'homomorphisme d'anneaux $\mathcal{O}_{Y, y} \to K$ (c'est possible puisque
cette image est un anneau local qui n'est pas un corps, car $y' \not = y$ ; voir
Algèbre, lemme \ref{algebra-lemma-dominate}).
Par hypothèse, il existe une extension de corps $K'/K$, un
anneau de valuation $A' \subset K'$ qui domine $A$, et un morphisme
$\Spec(A') \to X$ qui complète le diagramme commutatif.
Comme $A'$ domine $A$ et que $A$ domine $\mathcal{O}_{Y, y}$,
le point fermé de $\Spec(A')$ s'envoie sur
un point $x \in X$ tel que $f(x) = y$ et qui est une spécialisation de $x'$.
Ainsi $x \in T$, puisque $T$ est fermé, et donc $y \in f(T)$, comme voulu.
\end{proof}

\noindent
Le lemme suivant sera généralisé dans Espaces décents, lemme
\ref{decent-spaces-lemma-re-characterize-universally-closed}.

\begin{lemma}
\label{spaces-morphisms-lemma-characterize-universally-closed-quasi-separated}
Soit $S$ un schéma. Soit $f : X \to Y$ un
morphisme d'espaces algébriques sur $S$.
\begin{enumerate}
\item Si $f$ est quasi-séparé et universellement fermé, alors
$f$ satisfait la partie d'existence du critère valuatif.
\item Si $f$ est quasi-compact et quasi-séparé, alors
$f$ est universellement fermé si et seulement si la partie d'existence du
critère valuatif est satisfaite.
\end{enumerate}
\end{lemma}

\begin{proof}
Si (1) est vrai, sa combinaison avec le
lemme \ref{spaces-morphisms-lemma-quasi-compact-existence-universally-closed}
donne (2). Supposons que $f$ soit quasi-séparé et universellement fermé.
Considérons un diagramme
$$
\xymatrix{
\Spec(K) \ar[r] \ar[d] & X \ar[d] \\
\Spec(A) \ar[r] & Y
}
$$
comme dans la définition \ref{spaces-morphisms-definition-valuative-criterion}.
Un argument formel montre que l'existence du diagramme souhaité
$$
\xymatrix{
\Spec(K') \ar[r] \ar[d] & \Spec(K) \ar[r] & X \ar[d] \\
\Spec(A') \ar[r] \ar[rru] & \Spec(A) \ar[r] & Y
}
$$
résulte de l'existence d'un tel diagramme pour le morphisme $X_A \to \Spec(A)$.
Puisque les propriétés d'être quasi-séparé et universellement fermé sont préservées
par changement de base, le lemme résulte de ce qui est démontré au paragraphe suivant.

\medskip\noindent
Considérons un diagramme en traits pleins
$$
\xymatrix{
\Spec(K) \ar[r]_-x \ar[d] & X \ar[d]^f \\
\Spec(A) \ar@{=}[r] \ar@{..>}[ru] & \Spec(A)
}
$$
où $A$ est un anneau de valuation de corps des fractions $K$.
D'après le lemme \ref{spaces-morphisms-lemma-quasi-compact-permanence} et puisque
$f$ est quasi-séparé, le
morphisme $x$ est quasi-compact.
Comme $f$ est universellement fermé,
$|f|(\overline{\{x\}})$ est en particulier fermé dans $\Spec(A)$.
Puisque cette image contient le point générique de $\Spec(A)$,
il existe un point $x' \in |X|$, dans l'adhérence
de $x$, qui s'envoie sur le point fermé de $\Spec(A)$.
D'après le lemme \ref{spaces-morphisms-lemma-reach-points-scheme-theoretic-image},
on peut trouver un diagramme commutatif
$$
\xymatrix{
\Spec(K') \ar[r] \ar[d] & \Spec(K) \ar[d] \\
\Spec(A') \ar[r] & X
}
$$
tel que le point fermé de $\Spec(A')$ s'envoie sur $x' \in |X|$.
Il s'ensuit que $\Spec(A') \to \Spec(A)$ envoie le point fermé
sur le point fermé ; autrement dit, $A'$ domine $A$, ce qui achève la preuve.
\end{proof}

\begin{lemma}
\label{spaces-morphisms-lemma-characterize-universally-closed-separated}
Soit $S$ un schéma. Soit $f : X \to Y$ un morphisme d'espaces algébriques
sur $S$. Supposons que $f$ soit quasi-compact et séparé. Alors les assertions suivantes
sont équivalentes :
\begin{enumerate}
\item $f$ est universellement fermé ;
\item la partie d'existence du critère valuatif est satisfaite
au sens de la définition \ref{spaces-morphisms-definition-valuative-criterion} ;
\item pour tout diagramme commutatif en traits pleins
$$
\xymatrix{
\Spec(K) \ar[r] \ar[d] & X \ar[d] \\
\Spec(A) \ar[r] \ar@{-->}[ru] & Y
}
$$
où $A$ est un anneau de valuation de corps des fractions $K$, il existe
une flèche pointillée ; autrement dit, $f$ satisfait la partie d'existence du critère
valuatif au sens de
Schémas, définition \ref{schemes-definition-valuative-criterion}.
\end{enumerate}
\end{lemma}

\begin{proof}
Puisque $f$ est séparé, les assertions (2) et (3) sont équivalentes d'après le
lemme \ref{spaces-morphisms-lemma-usual-enough}.
L'équivalence de (3) et (1) résulte du
lemme \ref{spaces-morphisms-lemma-characterize-universally-closed-quasi-separated}.
\end{proof}

\begin{lemma}
\label{spaces-morphisms-lemma-lift-valuation-ring-through-flat-morphism}
Soit $S$ un schéma. Soit $f : X \to Y$ un morphisme plat
d'espaces algébriques sur $S$.
Soit $\Spec(A) \to Y$ un morphisme, où $A$ est un
anneau de valuation. Si le point fermé de $\Spec(A)$ s'envoie sur un
point de $|Y|$ appartenant à l'image de $|X| \to |Y|$, alors il existe
un diagramme commutatif
$$
\xymatrix{
\Spec(A') \ar[r] \ar[d] & X \ar[d] \\
\Spec(A) \ar[r] & Y
}
$$
où $A \to A'$ est une extension d'anneaux de valuation
(Compléments d'algèbre, définition
\ref{more-algebra-definition-extension-valuation-rings}).
\end{lemma}

\begin{proof}
Le changement de base $X_A \to \Spec(A)$ est plat
(lemme \ref{spaces-morphisms-lemma-base-change-flat}) et le point fermé de
$\Spec(A)$ appartient à l'image de $|X_A| \to |\Spec(A)|$
(Propriétés des espaces, lemme \ref{spaces-properties-lemma-points-cartesian}).
On peut donc supposer que $Y = \Spec(A)$. Soit $U \to X$ un morphisme
étale surjectif, où $U$ est un schéma. Soit $u \in U$ au-dessus
du point fermé de $\Spec(A)$. Considérons l'homomorphisme local plat
$A \to B = \mathcal{O}_{U, u}$. D'après
Algèbre, lemme \ref{algebra-lemma-ff-rings},
il existe un idéal premier $\mathfrak q \subset B$ tel que
$\mathfrak q$ soit au-dessus de $(0) \subset A$. D'après
Algèbre, lemme \ref{algebra-lemma-dominate},
on peut trouver un anneau de valuation $A' \subset \kappa(\mathfrak q)$
qui domine $B/\mathfrak q$. Le morphisme induit
$\Spec(A') \to U \to X$ résout le problème
posé par le lemme.
\end{proof}

\begin{lemma}
\label{spaces-morphisms-lemma-refined-valuative-criterion-universally-closed}
Soit $S$ un schéma. Soient $f : X \to Y$ et $h : U \to X$ des morphismes
d'espaces algébriques sur $S$. Si
\begin{enumerate}
\item $f$ et $h$ sont quasi-compacts,
\item $|h|(|U|)$ est dense dans $|X|$,
\end{enumerate}
et si, pour tout diagramme commutatif en traits pleins
$$
\xymatrix{
\Spec(K) \ar[r] \ar[d] & U \ar[r] & X \ar[d] \\
\Spec(A) \ar[rr] \ar@{-->}[rru] & & Y
}
$$
où $A$ est un anneau de valuation de corps des fractions $K$,
\begin{enumerate}
\item[(3)] il existe au plus une flèche pointillée rendant le diagramme
commutatif, et
\item[(4)] il existe une extension de corps $K'/K$, un
anneau de valuation $A' \subset K'$ qui domine $A$ et un morphisme
$\Spec(A') \to X$ tels que le diagramme suivant soit commutatif
$$
\xymatrix{
\Spec(K') \ar[r] \ar[d] & \Spec(K) \ar[r] & U \ar[r] & X \ar[d] \\
\Spec(A') \ar[r] \ar[rrru] & \Spec(A) \ar[rr] & & Y
}
$$
\end{enumerate}
alors $f$ est universellement fermé. Si, de plus,
\begin{enumerate}
\item[(5)] $f$ est quasi-séparé,
\end{enumerate}
alors $f$ est séparé et universellement fermé.
\end{lemma}

\begin{proof}
Supposons (1), (2), (3) et (4).
Nous allons vérifier la partie d'existence du critère valuatif pour $f$,
ce qui entraînera que $f$ est universellement fermé d'après le
lemme \ref{spaces-morphisms-lemma-quasi-compact-existence-universally-closed}.
Pour cela, considérons un diagramme commutatif
\begin{equation}
\label{spaces-morphisms-equation-start-with}
\vcenter{
\xymatrix{
\Spec(K) \ar[r] \ar[d] & X \ar[d] \\
\Spec(A) \ar[r] & Y
}
}
\end{equation}
où $A$ est un anneau de valuation et $K$ le corps des fractions de $A$.
Notons que, les anneaux de valuation et les corps étant réduits, on peut
remplacer $U$, $X$ et $Y$ par leurs réductions respectives d'après
Propriétés des espaces, lemme \ref{spaces-properties-lemma-map-into-reduction}.
Dans ce cas, l'hypothèse que $h(U)$ est dense signifie que
l'image schématique de $h : U \to X$ est $X$ ; voir le
lemme \ref{spaces-morphisms-lemma-scheme-theoretic-image-reduced}.

\medskip\noindent
Réduction au cas où $Y$ est affine. Choisissons un morphisme étale
$\Spec(R) \to Y$ tel que l'image du point fermé de $\Spec(A)$
appartienne à $\Im(|\Spec(R)| \to |Y|)$. D'après le
lemme \ref{spaces-morphisms-lemma-lift-valuation-ring-through-flat-morphism},
on peut trouver un homomorphisme local $A \to A'$ d'anneaux de valuation
et un morphisme $\Spec(A') \to \Spec(R)$ s'insérant dans un
diagramme commutatif
$$
\xymatrix{
\Spec(A') \ar[r] \ar[d] & \Spec(R) \ar[d] \\
\Spec(A) \ar[r] & Y
}
$$
Puisque la définition \ref{spaces-morphisms-definition-valuative-criterion}
autorise les extensions d'anneaux de valuation,
on peut clairement remplacer $A$ par $A'$, $Y$ par $\Spec(R)$,
$X$ par $X \times_Y \Spec(R)$ et $U$ par $U \times_Y \Spec(R)$.

\medskip\noindent
Nous supposons désormais que $Y = \Spec(R)$ est un schéma affine.
Soit $\Spec(B) \to X$ un morphisme étale dont la source est un schéma affine
tel que le morphisme $\Spec(K) \to X$ appartienne à l'image de
$|\Spec(B)| \to |X|$. Puisque l'on peut remplacer $K$ par une extension
$K' \supset K$ et $A$ par un anneau de valuation $A' \subset K'$
qui domine $A$ (un tel anneau existe d'après
Algèbre, lemme \ref{algebra-lemma-dominate}),
on peut supposer que le morphisme $\Spec(K) \to X$ se factorise par $\Spec(B)$
(par définition de $|X|$). Autrement dit, on peut considérer $K$ comme une $B$-algèbre.
Choisissons une algèbre polynomiale $P$ sur $B$ et une surjection de $B$-algèbres
$P \to K$. Alors $\Spec(P) \to X$ est plat, comme composée
$\Spec(P) \to \Spec(B) \to X$. Par conséquent, l'image schématique
du morphisme $U \times_X \Spec(P) \to \Spec(P)$ est $\Spec(P)$ d'après le
lemme \ref{spaces-morphisms-lemma-flat-base-change-scheme-theoretic-image}.
D'après le lemme \ref{spaces-morphisms-lemma-reach-points-scheme-theoretic-image},
on peut trouver un diagramme commutatif
$$
\xymatrix{
\Spec(K') \ar[r] \ar[d] & U \times_X \Spec(P) \ar[d] \\
\Spec(A') \ar[r] & \Spec(P)
}
$$
où $A'$ est un anneau de valuation et $K'$ le corps des fractions de $A'$ ;
le point fermé de $\Spec(A')$ s'envoie sur
$\Spec(K) \subset \Spec(P)$. Autrement dit, il existe un homomorphisme de $B$-algèbres
$\varphi : K \to A'/\mathfrak m_{A'}$. Choisissons un anneau de valuation
$A'' \subset A'/\mathfrak m_{A'}$ qui domine $\varphi(A)$ et dont le
corps des fractions est $K'' = A'/\mathfrak m_{A'}$
(Algèbre, lemme \ref{algebra-lemma-dominate}). Posons
$$
C = \{\lambda \in A' \mid \lambda \bmod \mathfrak m_{A'} \in A''\}
$$
qui est un anneau de valuation d'après
Algèbre, lemme \ref{algebra-lemma-stack-valuation-rings}.
Comme $C$ est une $R$-algèbre de corps des fractions $K'$, on obtient un diagramme
commutatif en traits pleins
$$
\xymatrix{
\Spec(K'_1) \ar@{-->}[r] \ar@{-->}[d] &
\Spec(K') \ar[r] \ar[d] & U \ar[r] & X \ar[d] \\
\Spec(C_1) \ar@{-->}[r] \ar@{-->}[rrru] & \Spec(C) \ar[rr] & & Y
}
$$
comme dans l'énoncé du lemme. L'hypothèse (4) fournit donc
$C \to C_1$ et les flèches pointillées qui rendent le diagramme commutatif.
Soit $A_1' = (C_1)_\mathfrak p$ le localisé de $C_1$
en un idéal premier $\mathfrak p \subset C_1$ situé au-dessus de
$\mathfrak m_{A'} \subset C$. Comme $C \to C_1$ est plat d'après
Compléments d'algèbre, lemme
\ref{more-algebra-lemma-valuation-ring-torsion-free-flat}
un tel idéal premier $\mathfrak p$ existe d'après
Algèbre, lemmes \ref{algebra-lemma-local-flat-ff} et
\ref{algebra-lemma-ff-rings}.
Notons que $A'$ est le localisé de $C$ en $\mathfrak m_{A'}$
et que $A'_1$ est un anneau de valuation
(Algèbre, lemme \ref{algebra-lemma-make-valuation-rings}).
Autrement dit, $A' \to A'_1$ est un homomorphisme local d'anneaux de
valuation. L'hypothèse (3) implique que le diagramme
$$
\xymatrix{
\Spec(A'_1) \ar[r] \ar[d] & \Spec(C_1) \ar[r] & X \\
\Spec(A') \ar[r] & \Spec(P) \ar[r] & \Spec(B) \ar[u]
}
$$
est commutatif. Ainsi, le morphisme induit par $\Spec(C_1) \to X$
sur $\Spec(C_1/\mathfrak p)$ a pour restriction la composée
$$
\Spec(\kappa(\mathfrak p)) \to
\Spec(A'/\mathfrak m_{A'}) = \Spec(K'') \to
\Spec(K) \to X
$$
au point générique de $\Spec(C_1/\mathfrak p)$. De plus,
$C_1/\mathfrak p$ est un anneau de valuation
(Algèbre, lemme \ref{algebra-lemma-make-valuation-rings})
qui domine $A''$, lequel domine $A$.
Ainsi, le morphisme $\Spec(C_1/\mathfrak p) \to X$ fournit la solution exigée par la
partie d'existence du critère valuatif pour le diagramme
(\ref{spaces-morphisms-equation-start-with}), comme voulu.

\medskip\noindent
Supposons enfin que (5) soit également satisfaite ; autrement dit, que le morphisme
$\Delta : X \to X \times_Y X$ soit quasi-compact. Dans ce cas,
les hypothèses (1) -- (4) sont satisfaites pour $h$ et $\Delta$. La première
partie de la preuve montre donc que $\Delta$ est universellement fermé.
D'après le lemme \ref{spaces-morphisms-lemma-separated-diagonal-proper}, on en conclut que
$f$ est séparé.
\end{proof}

















\section{Critère valuatif de séparation}
\label{spaces-morphisms-section-valuative-separatedness}

\noindent
Nous démontrons d'abord une réciproque, puis nous énonçons le critère.

\begin{lemma}
\label{spaces-morphisms-lemma-separated-implies-valuative}
Soit $S$ un schéma.
Soit $f : X \to Y$ un morphisme d'espaces algébriques sur $S$.
Si $f$ est séparé, alors $f$ satisfait la partie
d'unicité du critère valuatif.
\end{lemma}

\begin{proof}
Considérons un diagramme comme dans la définition \ref{spaces-morphisms-definition-valuative-criterion}.
Supposons qu'il existe deux morphismes distincts
$a, b : \Spec(A) \to X$ qui complètent ce diagramme.
Soit $Z \subset \Spec(A)$ l'égalisateur de $a$ et $b$.
Alors $Z = \Spec(A) \times_{(a, b), X \times_Y X, \Delta} X$.
Si $f$ est séparé, alors $\Delta$ est une immersion fermée, et
C'est un sous-schéma fermé de $\Spec(A)$. Par hypothèse, il contient
le point générique de $\Spec(A)$. Comme $A$ est intègre,
cela implique $Z = \Spec(A)$. Ainsi $a = b$, comme voulu.
\end{proof}

\begin{lemma}[Critère valuatif de séparation]
\label{spaces-morphisms-lemma-valuative-criterion-separatedness}
Soit $S$ un schéma.
Soit $f : X \to Y$ un morphisme d'espaces algébriques sur $S$.
Supposons que
\begin{enumerate}
\item le morphisme $f$ soit quasi-séparé, et que
\item le morphisme $f$ satisfasse la partie d'unicité
du critère valuatif.
\end{enumerate}
Alors $f$ est séparé.
\end{lemma}

\begin{proof}
L'hypothèse (1) signifie que $\Delta_{X/Y}$ est quasi-compact.
Nous affirmons que le morphisme
$\Delta_{X/Y} : X \to X \times_Y X$ satisfait la partie d'existence
du critère valuatif.
Considérons un diagramme commutatif en traits pleins
$$
\xymatrix{
\Spec(K) \ar[r] \ar[d] & X \ar[d] \\
\Spec(A) \ar[r] \ar@{-->}[ru] & X \times_Y X
}
$$
La flèche horizontale inférieure correspond à une
paire de morphismes $a, b : \Spec(A) \to X$ sur $Y$.
D'après l'hypothèse (2), on a $a = b$. Prendre $a$ pour flèche
pointillée convient. Par conséquent, le
lemme \ref{spaces-morphisms-lemma-quasi-compact-existence-universally-closed}
s'applique, et $\Delta_{X/Y}$ est universellement fermé.
Puisque $\Delta_{X/Y}$ est toujours localement de type fini et
séparé, il résulte de
Compléments sur les morphismes, lemme \ref{more-morphisms-lemma-characterize-finite}
que $\Delta_{X/Y}$ est un morphisme fini (on peut aussi utiliser le
principe général de
Espaces, lemme
\ref{spaces-lemma-representable-transformations-property-implication}).
À ce stade, $\Delta_{X/Y}$ est un monomorphisme représentable et fini,
donc une immersion fermée d'après
Morphismes, lemme \ref{morphisms-lemma-finite-monomorphism-closed}.
\end{proof}

\begin{lemma}
\label{spaces-morphisms-lemma-characterize-separated-and-universally-closed}
Soit $S$ un schéma. Soit $f : X \to Y$ un morphisme d'espaces algébriques
sur $S$. Supposons que $f$ soit quasi-compact et quasi-séparé. Alors les
assertions suivantes sont équivalentes :
\begin{enumerate}
\item $f$ est séparé et universellement fermé ;
\item le critère valuatif est satisfait au sens de la définition
\ref{spaces-morphisms-definition-valuative-criterion},
\item pour tout diagramme commutatif en traits pleins
$$
\xymatrix{
\Spec(K) \ar[r] \ar[d] & X \ar[d] \\
\Spec(A) \ar[r] \ar@{-->}[ru] & Y
}
$$
où $A$ est un anneau de valuation de corps des fractions $K$, il existe
une unique flèche pointillée ; autrement dit, $f$ satisfait le critère
valuatif au sens de
Schémas, définition \ref{schemes-definition-valuative-criterion}.
\end{enumerate}
\end{lemma}

\begin{proof}
Puisque $f$ est quasi-séparé, la partie d'unicité
du critère valuatif implique que $f$ est séparé
(lemme \ref{spaces-morphisms-lemma-valuative-criterion-separatedness}).
Réciproquement, si $f$ est séparé, il satisfait la
partie d'unicité du critère valuatif
(lemme \ref{spaces-morphisms-lemma-separated-implies-valuative}).
Cela étant, dans chacun des trois cas, le
morphisme $f$ est séparé et satisfait la partie d'unicité
du critère valuatif. Dans ce cas, le lemme est une conséquence
formelle du
lemme \ref{spaces-morphisms-lemma-characterize-universally-closed-separated}.
\end{proof}



\section{Critère valuatif de propreté}
\label{spaces-morphisms-section-valuative-criterion-properness}

\noindent
Voici l'énoncé.

\begin{lemma}[Critère valuatif de propreté]
\label{spaces-morphisms-lemma-characterize-proper}
Soit $S$ un schéma. Soit $f : X \to Y$ un morphisme d'espaces algébriques
sur $S$. Supposons que $f$ soit de type fini et quasi-séparé. Alors les
assertions suivantes sont équivalentes :
\begin{enumerate}
\item $f$ est propre ;
\item le critère valuatif est satisfait au sens de la définition
\ref{spaces-morphisms-definition-valuative-criterion},
\item pour tout diagramme commutatif en traits pleins
$$
\xymatrix{
\Spec(K) \ar[r] \ar[d] & X \ar[d] \\
\Spec(A) \ar[r] \ar@{-->}[ru] & Y
}
$$
où $A$ est un anneau de valuation de corps des fractions $K$, il existe
une unique flèche pointillée ; autrement dit, $f$ satisfait le critère
valuatif au sens de
Schémas, définition \ref{schemes-definition-valuative-criterion}.
\end{enumerate}
\end{lemma}

\begin{proof}
Cela résulte formellement du
lemme \ref{spaces-morphisms-lemma-characterize-separated-and-universally-closed}
et des définitions.
\end{proof}





\section{Morphismes entiers et finis}
\label{spaces-morphisms-section-integral}

\noindent
Nous avons déjà défini dans la section \ref{spaces-morphisms-section-representable}
ce que signifie, pour un morphisme représentable d'espaces algébriques,
être entier (resp.\ fini).

\begin{lemma}
\label{spaces-morphisms-lemma-integral-representable}
Soit $S$ un schéma. Soit $f : X \to Y$ un morphisme représentable
d'espaces algébriques sur $S$. Alors
$f$ est entier, resp.\ fini
(au sens de la section \ref{spaces-morphisms-section-representable}),
si et seulement si, pour tout schéma affine $Z$
et tout morphisme $Z \to Y$, le schéma $X \times_Y Z$ est affine et
entier, resp.\ fini, sur $Z$.
\end{lemma}

\begin{proof}
Cela résulte directement de la définition d'un morphisme entier (resp.\ fini)
de schémas
(Morphismes, définition \ref{morphisms-definition-integral}).
\end{proof}

\noindent
Cela permet de poser la définition suivante.

\begin{definition}
\label{spaces-morphisms-definition-integral}
Soit $S$ un schéma.
Soit $f : X \to Y$ un morphisme d'espaces algébriques sur $S$.
\begin{enumerate}
\item Nous disons que $f$ est {\it entier} si, pour tout schéma affine $Z$
et tout morphisme $Z \to Y$, l'espace algébrique $X \times_Y Z$ est
représentable par un schéma affine entier sur $Z$.
\item Nous disons que $f$ est {\it fini} si, pour tout schéma affine $Z$
et tout morphisme $Z \to Y$, l'espace algébrique $X \times_Y Z$ est
représentable par un schéma affine fini sur $Z$.
\end{enumerate}
\end{definition}

\begin{lemma}
\label{spaces-morphisms-lemma-integral-local}
Soit $S$ un schéma.
Soit $f : X \to Y$ un morphisme d'espaces algébriques sur $S$.
Les assertions suivantes sont équivalentes :
\begin{enumerate}
\item $f$ est représentable et entier (resp.\ fini) ;
\item $f$ est entier (resp.\ fini) ;
\item il existe un schéma $V$ et un morphisme étale surjectif
$V \to Y$ tel que $V \times_Y X \to V$ soit entier (resp. fini) ;
\item il existe un recouvrement de Zariski $Y = \bigcup Y_i$ tel que
chacun des morphismes $f^{-1}(Y_i) \to Y_i$ soit entier (resp.\ fini).
\end{enumerate}
\end{lemma}

\begin{proof}
Il est clair que (1) implique (2) et que (2) implique (3) en prenant pour
$V$ une somme disjointe de schémas affines étales sur $Y$ ; voir
Propriétés des espaces,
lemme \ref{spaces-properties-lemma-cover-by-union-affines}.
Supposons que $V \to Y$ soit comme en (3). Pour tout ouvert affine $W$ de $V$,
$W \times_Y X$ est alors un ouvert affine de $V \times_Y X$. D'après
Propriétés des espaces, lemme \ref{spaces-properties-lemma-subscheme}
on en conclut que $V \times_Y X$ est un schéma. De plus, le morphisme
$V \times_Y X \to V$ est affine. On peut donc appliquer
Espaces,
lemme \ref{spaces-lemma-morphism-sheaves-with-P-effective-descent-etale}
car la classe des morphismes entiers (resp.\ finis)
satisfait toutes les propriétés requises (voir
Morphismes, lemmes \ref{morphisms-lemma-base-change-finite} et
Descente, lemmes
\ref{descent-lemma-descending-property-integral},
\ref{descent-lemma-descending-property-finite},
et \ref{descent-lemma-affine}).
La conclusion de ce lemme est que $f$ est représentable
et entier (resp.\ fini), c'est-à-dire que (1) est satisfaite.

\medskip\noindent
L'équivalence de (1) et (4) résulte du fait que la propriété d'être
entier (resp.\ fini) est locale pour la topologie de Zariski sur le but (la
référence ci-dessus montre qu'en fait, la propriété d'être entier ou fini est
locale pour la topologie fpqc sur le but).
\end{proof}

\begin{lemma}
\label{spaces-morphisms-lemma-composition-integral}
La composée de morphismes entiers (resp.\ finis) est entière
(resp.\ finie).
\end{lemma}

\begin{proof}
Omis.
\end{proof}

\begin{lemma}
\label{spaces-morphisms-lemma-base-change-integral}
Tout changement de base d'un morphisme entier (resp.\ fini) est entier
(resp.\ fini).
\end{lemma}

\begin{proof}
Omis.
\end{proof}

\begin{lemma}
\label{spaces-morphisms-lemma-finite-integral}
Un morphisme fini d'espaces algébriques est entier.
Un morphisme entier d'espaces algébriques
qui est localement de type fini est fini.
\end{lemma}

\begin{proof}
Dans les deux cas, le morphisme est représentable, et la condition se vérifie
après un changement de base par un schéma affine muni d'un morphisme vers $Y$ ; voir le
lemme \ref{spaces-morphisms-lemma-integral-local}. Le résultat découle donc du
lemme analogue pour les schémas ; voir
Morphismes, lemme \ref{morphisms-lemma-finite-integral}.
\end{proof}

\begin{lemma}
\label{spaces-morphisms-lemma-integral-universally-closed}
Soit $S$ un schéma.
Soit $f : X \to Y$ un morphisme d'espaces algébriques sur $S$.
Les assertions suivantes sont équivalentes :
\begin{enumerate}
\item $f$ est entier ;
\item $f$ est affine et universellement fermé.
\end{enumerate}
\end{lemma}

\begin{proof}
Dans les deux cas, le morphisme est représentable, et la condition se vérifie
après un changement de base par un schéma affine muni d'un morphisme vers $Y$ ; voir les
lemmes \ref{spaces-morphisms-lemma-integral-local},
\ref{spaces-morphisms-lemma-affine-local} et
\ref{spaces-morphisms-lemma-universally-closed-local}.
Le résultat découle donc de
Morphismes, lemme \ref{morphisms-lemma-integral-universally-closed}.
\end{proof}

\begin{lemma}
\label{spaces-morphisms-lemma-finite-quasi-finite}
Un morphisme fini d'espaces algébriques est quasi-fini.
\end{lemma}

\begin{proof}
Soit $f : X \to Y$ un morphisme d'espaces algébriques.
D'après la
définition \ref{spaces-morphisms-definition-integral}
et les
lemmes \ref{spaces-morphisms-lemma-quasi-compact-local} et
\ref{spaces-morphisms-lemma-quasi-finite-local}
les deux propriétés se vérifient après changement de base à un schéma affine sur $Y$ ;
autrement dit, on peut supposer $Y$ affine.
Si $f$ est fini, alors $X$ est un schéma.
Le résultat découle donc du résultat correspondant pour les schémas ; voir
Morphismes, lemme \ref{morphisms-lemma-finite-quasi-finite}.
\end{proof}

\begin{lemma}
\label{spaces-morphisms-lemma-finite-proper}
Soit $S$ un schéma. Soit $f : X \to Y$ un morphisme d'espaces algébriques sur
$S$. Les assertions suivantes sont équivalentes :
\begin{enumerate}
\item $f$ est fini ;
\item $f$ est affine et propre.
\end{enumerate}
\end{lemma}

\begin{proof}
Dans les deux cas, le morphisme est représentable, et la condition se vérifie
après changement de base à un schéma affine muni d'un morphisme vers $Y$ ; voir les
lemmes \ref{spaces-morphisms-lemma-integral-local}, \ref{spaces-morphisms-lemma-affine-local} et
\ref{spaces-morphisms-lemma-proper-local}. Le résultat découle donc de
Morphismes, lemme \ref{morphisms-lemma-finite-proper}.
\end{proof}

\begin{lemma}
\label{spaces-morphisms-lemma-closed-immersion-finite}
Une immersion fermée est finie (et a fortiori entière).
\end{lemma}

\begin{proof}
Omis.
\end{proof}

\begin{lemma}
\label{spaces-morphisms-lemma-finite-union-finite}
Soit $S$ un schéma.
Soient $X_i \to Y$, $i = 1, \ldots, n$, des morphismes finis
d'espaces algébriques sur $S$.
Alors $X_1 \amalg \ldots \amalg X_n \to Y$ est également fini.
\end{lemma}

\begin{proof}
Cela résulte du cas des schémas
(Morphismes, lemme \ref{morphisms-lemma-finite-union-finite})
par localisation étale.
\end{proof}

\begin{lemma}
\label{spaces-morphisms-lemma-finite-permanence}
Soit $S$ un schéma.
Soient $f : X \to Y$ et $g : Y \to Z$ des morphismes d'espaces algébriques sur $S$.
\begin{enumerate}
\item Si $g \circ f$ est fini et si $g$ est séparé, alors $f$ est fini.
\item Si $g \circ f$ est entier et si $g$ est séparé, alors $f$ est entier.
\end{enumerate}
\end{lemma}

\begin{proof}
Supposons que $g \circ f$ soit fini (resp.\ entier) et que $g$ soit séparé.
Le changement de base $X \times_Z Y \to Y$ est fini (resp.\ entier) d'après le
lemme \ref{spaces-morphisms-lemma-base-change-integral}.
Le morphisme $X \to X \times_Z Y$ est
une immersion fermée puisque $Y \to Z$ est séparé ; voir le
lemme \ref{spaces-morphisms-lemma-section-immersion}.
Une immersion fermée est finie (resp.\ entière) ;
voir le lemme \ref{spaces-morphisms-lemma-closed-immersion-finite}.
La composée de morphismes finis (resp.\ entiers) est finie
(resp.\ entière) ;
voir le lemme \ref{spaces-morphisms-lemma-composition-integral}. D'où le résultat.
\end{proof}





\section{Morphismes finis localement libres}
\label{spaces-morphisms-section-finite-locally-free}

\noindent
Nous avons déjà défini dans la section \ref{spaces-morphisms-section-representable}
ce que signifie, pour un morphisme représentable d'espaces algébriques,
être fini localement libre.

\begin{lemma}
\label{spaces-morphisms-lemma-finite-locally-free-representable}
Soit $S$ un schéma. Soit $f : X \to Y$ un morphisme représentable
d'espaces algébriques sur $S$. Alors $f$ est fini localement libre
(au sens de la section \ref{spaces-morphisms-section-representable})
si et seulement si $f$ est affine et si le faisceau $f_*\mathcal{O}_X$ est
un $\mathcal{O}_Y$-module fini localement libre.
\end{lemma}

\begin{proof}
Supposons que $f$ soit fini localement libre (au sens de la
section \ref{spaces-morphisms-section-representable}). Cela signifie que,
pour tout morphisme $V \to Y$ dont la source est un schéma, le
changement de base $f' : V \times_Y X \to V$ est un morphisme fini localement libre
de schémas. Cela signifie à son tour (par définition d'un morphisme fini
localement libre de schémas) que
$f'_*\mathcal{O}_{V \times_Y X}$
est un $\mathcal{O}_V$-module fini localement libre. On peut choisir $V \to Y$
étale et surjectif. D'après
Propriétés des espaces,
lemme \ref{spaces-properties-lemma-pushforward-etale-base-change-modules}
la restriction de $f_*\mathcal{O}_X$ à $V$ est
finie localement libre. Par conséquent, d'après
Modules sur les sites, lemme \ref{sites-modules-lemma-local-final-object}
appliqué au faisceau $f_*\mathcal{O}_X$ sur $Y_{spaces, \etale}$,
on en conclut que $f_*\mathcal{O}_X$ est fini localement libre.

\medskip\noindent
Réciproquement, supposons que $f$ soit affine et que $f_*\mathcal{O}_X$ soit un
$\mathcal{O}_Y$-module fini localement libre. Soit $V$ un schéma et soit
$V \to Y$ un morphisme étale surjectif. De nouveau, d'après
Propriétés des espaces,
lemme \ref{spaces-properties-lemma-pushforward-etale-base-change-modules}
$f'_*\mathcal{O}_{V \times_Y X}$ est fini localement libre.
Ainsi, $f' : V \times_Y X \to V$ est fini localement libre (puisqu'il est aussi affine).
D'après
Espaces,
lemme \ref{spaces-lemma-morphism-sheaves-with-P-effective-descent-etale}
on en conclut que $f$ est fini localement libre (utiliser
Morphismes, lemme \ref{morphisms-lemma-base-change-finite-locally-free}
et Descente, lemmes \ref{descent-lemma-descending-property-finite-locally-free}
et \ref{descent-lemma-affine}). D'où le résultat.
\end{proof}

\noindent
Cela permet de poser la définition suivante.

\begin{definition}
\label{spaces-morphisms-definition-finite-locally-free}
Soit $S$ un schéma.
Soit $f : X \to Y$ un morphisme d'espaces algébriques sur $S$.
Nous disons que $f$ est {\it fini localement libre} si $f$ est affine
et si $f_*\mathcal{O}_X$ est un $\mathcal{O}_Y$-module fini localement libre.
Dans ce cas, nous disons que $f$ est
de {\it rang} ou de {\it degré} $d$
si le faisceau $f_*\mathcal{O}_X$ est fini localement libre de rang $d$.
\end{definition}

\begin{lemma}
\label{spaces-morphisms-lemma-finite-locally-free-local}
Soit $S$ un schéma.
Soit $f : X \to Y$ un morphisme d'espaces algébriques sur $S$.
Les assertions suivantes sont équivalentes :
\begin{enumerate}
\item $f$ est représentable et fini localement libre ;
\item $f$ est fini localement libre ;
\item il existe un schéma $V$ et un morphisme étale surjectif
$V \to Y$ tel que $V \times_Y X \to V$ soit fini localement libre ;
\item il existe un recouvrement de Zariski $Y = \bigcup Y_i$ tel que
chaque morphisme $f^{-1}(Y_i) \to Y_i$ soit fini localement libre.
\end{enumerate}
\end{lemma}

\begin{proof}
Il est clair que (1) implique (2) et que (2) implique (3) en prenant pour
$V$ une somme disjointe de schémas affines étales sur $Y$ ; voir
Propriétés des espaces,
lemme \ref{spaces-properties-lemma-cover-by-union-affines}.
Supposons que $V \to Y$ soit comme en (3). Pour tout ouvert affine $W$ de $V$,
$W \times_Y X$ est alors un ouvert affine de $V \times_Y X$. D'après
Propriétés des espaces, lemme \ref{spaces-properties-lemma-subscheme}
on en conclut que $V \times_Y X$ est un schéma. De plus, le morphisme
$V \times_Y X \to V$ est affine. On peut donc appliquer
Espaces,
lemme \ref{spaces-lemma-morphism-sheaves-with-P-effective-descent-etale}
car la classe des morphismes finis localement libres
satisfait toutes les propriétés requises (voir
Morphismes, lemme \ref{morphisms-lemma-base-change-finite-locally-free}
et Descente, lemmes \ref{descent-lemma-descending-property-finite-locally-free}
et \ref{descent-lemma-affine}).
La conclusion de ce lemme est que $f$ est représentable
et fini localement libre, c'est-à-dire que (1) est satisfaite.

\medskip\noindent
L'équivalence de (1) et (4) résulte du fait que la propriété d'être
fini localement libre est locale pour la topologie de Zariski sur le but (la référence ci-dessus montre
qu'en fait, la propriété d'être fini localement libre est locale pour la topologie fpqc sur le but).
\end{proof}

\begin{lemma}
\label{spaces-morphisms-lemma-composition-finite-locally-free}
La composée de morphismes finis localement libres est finie localement libre.
\end{lemma}

\begin{proof}
Omis.
\end{proof}

\begin{lemma}
\label{spaces-morphisms-lemma-base-change-finite-locally-free}
Tout changement de base d'un morphisme fini localement libre est fini localement libre.
\end{lemma}

\begin{proof}
Omis.
\end{proof}

\begin{lemma}
\label{spaces-morphisms-lemma-finite-flat}
Soit $S$ un schéma.
Soit $f : X \to Y$ un morphisme d'espaces algébriques sur $S$.
Les assertions suivantes sont équivalentes :
\begin{enumerate}
\item $f$ est fini localement libre ;
\item $f$ est fini, plat et localement de présentation finie.
\end{enumerate}
Si $Y$ est localement noethérien, ces assertions sont aussi équivalentes à
\begin{enumerate}
\item[(3)] $f$ est fini et plat.
\end{enumerate}
\end{lemma}

\begin{proof}
Dans chacun des trois cas, le morphisme est représentable et la propriété
se vérifie après changement de base par un morphisme étale surjectif
$V \to Y$ ; voir les
lemmes \ref{spaces-morphisms-lemma-integral-local},
\ref{spaces-morphisms-lemma-finite-locally-free-local},
\ref{spaces-morphisms-lemma-flat-local} et
\ref{spaces-morphisms-lemma-finite-presentation-local}.
Si $Y$ est localement noethérien, alors $V$ est localement noethérien.
Le résultat découle donc du résultat correspondant
pour les schémas ; voir
Morphismes, lemme \ref{morphisms-lemma-finite-flat}.
\end{proof}





\section{Applications rationnelles}
\label{spaces-morphisms-section-rational-maps}

\noindent
Cette section est l'analogue de
Morphismes, section \ref{morphisms-section-rational-maps}.
Nous utiliserons sans autre mention le fait que l'intersection d'ouverts denses
d'un espace topologique est un ouvert dense.

\begin{definition}
\label{spaces-morphisms-definition-rational-map}
Soit $S$ un schéma. Soient $X$, $Y$ des espaces algébriques sur $S$.
\begin{enumerate}
\item Soient $f : U \to Y$, $g : V \to Y$ des morphismes d'espaces algébriques
sur $S$, définis sur des sous-espaces ouverts denses $U$, $V$ de $X$. On dit que $f$ est
{\it équivalent} à $g$ si $f|_W = g|_W$ pour un sous-espace ouvert dense
$W \subset U \cap V$.
\item Une {\it application rationnelle de $X$ dans $Y$}
est une classe d'équivalence pour la relation d'équivalence définie en (1).
\item Étant donnés des morphismes $X \to B$ et $Y \to B$ d'espaces algébriques sur $S$,
on dit qu'une application rationnelle de $X$ dans $Y$ est une
{\it $B$-application rationnelle de $X$ dans $Y$}
s'il existe un représentant $f : U \to Y$ de la classe d'équivalence
qui soit un morphisme au-dessus de $B$.
\end{enumerate}
\end{definition}

\noindent
On dira que deux morphismes $f$, $g$ comme en (1) de la définition
définissent la même application rationnelle, plutôt que de les dire équivalents.
Dans de nombreux cas considérés par la suite, les espaces algébriques
$X$ et $Y$ contiendront des sous-espaces ouverts denses $X'$ et $Y'$ qui
sont des schémas. Alors une application rationnelle de $X$ dans $Y$ est la même chose
qu'une $S$-application rationnelle de $X'$ dans $Y'$ au sens de
Morphismes, définition \ref{spaces-morphisms-definition-rational-map}.
Toute la théorie développée pour les schémas s'applique alors.

\begin{definition}
\label{spaces-morphisms-definition-rational-function}
Soit $S$ un schéma. Soit $X$ un espace algébrique sur $S$. Une
{\it fonction rationnelle sur $X$} est une application rationnelle de $X$ dans $\mathbf{A}^1_S$.
\end{definition}

\noindent
En considérant la discussion qui suit
Morphismes, définition \ref{morphisms-definition-rational-function},
on voit qu'il s'agit de la même notion lorsque
$X$ est un schéma.

\medskip\noindent
Rappelons l'identification canonique
$$
\Mor_S(T, \mathbf{A}^1_S) = \Mor(T, \mathbf{A}^1_\mathbf{Z}) =
\Gamma(T, \mathcal{O}_T)
$$
pour tout schéma $T$ sur $S$ ; voir
Schémas, exemple \ref{schemes-example-global-sections}.
Ainsi, les lois d'addition et de multiplication font de $\mathbf{A}^1_S$ un anneau dans la
catégorie des schémas sur $S$. Autrement dit, l'addition
et la multiplication définissent des morphismes
$$
+ : \mathbf{A}^1_S \times_S \mathbf{A}^1_S \to \mathbf{A}^1_S
\quad\text{et}\quad
* : \mathbf{A}^1_S \times_S \mathbf{A}^1_S \to \mathbf{A}^1_S
$$
qui satisfont aux axiomes de l'addition et de la multiplication dans un anneau
(commutatif, d'élément unité $1$, comme toujours). L'ensemble des applications rationnelles
à valeurs dans $\mathbf{A}^1_S$ est donc lui aussi muni d'une structure naturelle d'anneau.

\begin{definition}
\label{spaces-morphisms-definition-ring-of-rational-functions}
Soit $S$ un schéma.
Soit $X$ un espace algébrique sur $S$.
L'{\it anneau des fonctions rationnelles sur $X$}
est l'anneau $R(X)$ dont les éléments sont les fonctions rationnelles, muni de
l'addition et de la multiplication qui viennent d'être décrites.
\end{definition}

\noindent
Nous définirons plus loin le corps des fonctions d'un espace algébrique intègre ; voir
Espaces sur les corps, section \ref{spaces-over-fields-section-integral-spaces}.

\begin{definition}
\label{spaces-morphisms-definition-domain-of-definition}
Soit $S$ un schéma. Soit $\varphi$ une application rationnelle entre deux
espaces algébriques $X$ et $Y$ sur $S$. On dit que
$\varphi$ est {\it définie au point $x \in |X|$} s'il existe un
représentant $(U, f)$ de $\varphi$ tel que $x \in |U|$. Le
{\it domaine de définition} de $\varphi$ est l'ensemble des points où
$\varphi$ est définie.
\end{definition}

\noindent
Le domaine de définition est considéré comme un sous-espace ouvert de $X$ grâce à
Propriétés des espaces, lemme \ref{spaces-properties-lemma-open-subspaces}.
Avec cette définition, il n'est pas vrai en général que $\varphi$ possède un
représentant défini sur tout son domaine de définition.

\begin{lemma}
\label{spaces-morphisms-lemma-rational-map-from-reduced-to-separated}
Soit $S$ un schéma. Soient $X$ et $Y$ des espaces algébriques sur $S$.
Supposons $X$ réduit et $Y$ séparé sur $S$. Soit
$\varphi$ une application rationnelle de $X$ dans $Y$, de domaine de définition
$U \subset X$. Il existe alors un unique morphisme $f : U \to Y$
d'espaces algébriques qui représente $\varphi$.
\end{lemma}

\begin{proof}
Soient $(V, g)$ et $(V', g')$ des représentants de $\varphi$. Alors
Les morphismes $g, g'$ coïncident sur un sous-espace ouvert dense $W \subset V \cap V'$.
D'autre part, l'égalisateur $E$ de $g|_{V \cap V'}$ et $g'|_{V \cap V'}$
est un sous-espace fermé de $V \cap V'$, car il s'obtient par changement de base
de $\Delta : Y \to Y \times_S Y$ par le morphisme
$V \cap V' \to Y \times_S Y$ défini par $g|_{V \cap V'}$ et $g'|_{V \cap V'}$.
Or $W \subset E$ implique que $|E| = |V \cap V'|$. Comme $V \cap V'$
est réduit, on en déduit $E = V \cap V'$ schématiquement, c'est-à-dire
$g|_{V \cap V'} = g'|_{V \cap V'}$ ; voir
Propriétés des espaces, lemme \ref{spaces-properties-lemma-map-into-reduction}.
On peut donc recoller les représentants $g : V \to Y$ de $\varphi$
en un morphisme $f : U \to Y$, puisque $\coprod V \to U$ est un morphisme surjectif de
faisceaux fppf et que
$\coprod_{V, V'} V \cap V' = (\coprod V) \times_U (\coprod V)$.
\end{proof}

\noindent
En général, la composition d'applications rationnelles n'a pas de sens. En effet,
l'image d'un représentant de la première application rationnelle peut
ne pas rencontrer le domaine de définition de la seconde.
Toutefois, si les espaces sont irréductibles et si l'on considère
des applications rationnelles dominantes, on peut les composer.

\begin{definition}
\label{spaces-morphisms-definition-dominant-rational}
Soit $S$ un schéma. Soient $X$ et $Y$ des espaces algébriques sur $S$.
Supposons $|X|$ et $|Y|$ irréductibles. Une application rationnelle de $X$ dans $Y$
est dite {\it dominante} si tout représentant $f : U \to Y$ est un morphisme dominant
au sens de la définition \ref{spaces-morphisms-definition-dominant}.
\end{definition}

\noindent
On peut composer une application rationnelle dominante
$\varphi$ entre des espaces algébriques irréductibles
$X$ et $Y$ avec une application rationnelle arbitraire
$\psi$ de $Y$ dans $Z$. Choisissons en effet des
représentants $f : U \to Y$, avec $|U| \subset |X|$ ouvert dense,
et $g : V \to Z$, avec $|V| \subset |Y|$ ouvert dense. Alors
$W = |f|^{-1}(V) \subset |X|$ est ouvert non vide (car
l'image de $|f|$ est dense et rencontre donc l'ouvert non vide $V$),
donc dense puisque $|X|$ est irréductible. On définit $\psi \circ \varphi$
comme la classe d'équivalence de $g \circ f|_W : W \to Z$. Nous omettons la vérification
que cette définition ne dépend pas des choix.

\medskip\noindent
On obtient ainsi une catégorie dont les objets sont les espaces algébriques
irréductibles sur $S$ et dont les morphismes sont les applications rationnelles dominantes.

\begin{definition}
\label{spaces-morphisms-definition-birational-spaces}
Soit $S$ un schéma. Soient $X$ et $Y$ des espaces algébriques
sur $S$, avec $|X|$ et $|Y|$ irréductibles.
On dit que $X$ et $Y$ sont {\it birationnels} si $X$ et $Y$ sont isomorphes
dans la catégorie des espaces algébriques irréductibles sur $S$
et des applications rationnelles dominantes.
\end{definition}

\noindent
Si $X$ et $Y$ sont des espaces algébriques irréductibles birationnels,
l'ensemble des applications rationnelles de $X$ dans $Z$ est en bijection
avec l'ensemble des applications rationnelles de $Y$ dans $Z$, pour tout espace
algébrique $Z$ (fonctoriellement en $Z$). Pour des espaces algébriques irréductibles « généraux »,
ce n'est là qu'une définition possible. On pourrait aussi
exiger que $X$ et $Y$ aient des anneaux de fonctions
rationnelles isomorphes ; ces deux notions sont parfois équivalentes
(insérer ici un renvoi futur).

\begin{lemma}
\label{spaces-morphisms-lemma-birational}
Soit $S$ un schéma. Soient $X$ et $Y$ des espaces
algébriques sur $S$, avec $|X|$ et $|Y|$ irréductibles.
Alors $X$ et $Y$ sont birationnels si et seulement s'il existe
des sous-espaces ouverts non vides $U \subset X$ et $V \subset Y$
qui sont isomorphes comme espaces algébriques sur $S$.
\end{lemma}

\begin{proof}
Supposons $X$ et $Y$ birationnels. Soient $f : U \to Y$ et $g : V \to X$
des représentants d'applications rationnelles dominantes inverses de $X$ dans $Y$ et de $Y$ dans $X$.
Après avoir rétréci $U$, on peut supposer que $f : U \to Y$ se factorise par $V$.
Comme $g \circ f$ est l'identité en tant qu'application rationnelle dominante,
la composée $U \to V \to X$ est l'identité sur un ouvert dense de $U$.
En remplaçant $U$ par un ouvert plus petit, on peut donc supposer que
$U \to V \to X$ est l'inclusion de $U$ dans $X$.
Par symétrie, il existe un sous-espace ouvert $V' \subset V$
tel que $g|_{V'} : V' \to X$ se factorise par $U \subset X$
et que $V' \to U \to Y$ soit l'identité.
L'image réciproque de $|V'|$ par $|U| \to |V|$ est un ouvert de $|U|$,
donc elle est égale à $|U'|$ pour un sous-espace ouvert $U' \subset U$ ; voir
Propriétés des espaces, lemme \ref{spaces-properties-lemma-open-subspaces}.
Alors $U' \subset U \to V$ se factorise en $U' \to V'$.
De même, $V' \to U$ se factorise en $V' \to U'$.
On vérifie que $U' \to V'$
et $V' \to U'$ sont des morphismes
d'espaces algébriques sur $S$ inverses l'un de l'autre,
ce qui achève la démonstration.
\end{proof}









\section{Normalisation relative des espaces algébriques}
\label{spaces-morphisms-section-normalization-X-in-Y}

\noindent
Cette section est l'analogue de
Morphismes, section \ref{morphisms-section-normalization-X-in-Y}.

\begin{lemma}
\label{spaces-morphisms-lemma-integral-closure}
Soit $S$ un schéma. Soit $X$ un espace algébrique sur $S$.
Soit $\mathcal{A}$ un faisceau quasi-cohérent de $\mathcal{O}_X$-algèbres.
Il existe un sous-faisceau quasi-cohérent de $\mathcal{O}_X$-algèbres
$\mathcal{A}' \subset \mathcal{A}$ tel que,
pour tout objet affine $U$ de $X_\etale$, l'anneau
$\mathcal{A}'(U) \subset \mathcal{A}(U)$ soit
la clôture intégrale de $\mathcal{O}_X(U)$ dans $\mathcal{A}(U)$.
\end{lemma}

\begin{proof}
Soit $U$ un objet de $X_\etale$. Alors $U$ est un schéma. Notons
$\mathcal{A}|_U$ la restriction au site de Zariski. Le faisceau
$\mathcal{A}|_U$ est un faisceau quasi-cohérent de $\mathcal{O}_U$-algèbres ;
on peut donc appliquer Morphismes, lemme \ref{morphisms-lemma-integral-closure},
pour obtenir une sous-algèbre quasi-cohérente $\mathcal{A}'_U \subset \mathcal{A}|_U$
telle que la valeur de $\mathcal{A}'_U$ sur tout ouvert affine $W \subset U$
est celle indiquée dans l'énoncé du lemme. Si $f : U' \to U$ est un morphisme
dans $X_\etale$, alors $\mathcal{A}|_{U'} = f^*(\mathcal{A}|_U)$,
où $f^*$ désigne l'image inverse par le morphisme $f$ pour la topologie de Zariski ;
cela résulte de la quasi-cohérence de $\mathcal{A}$ (voir l'introduction de
Propriétés des espaces, section \ref{spaces-properties-section-quasi-coherent},
ainsi que les renvois à la discussion du chapitre sur la
descente des schémas). Comme $f$ est étale, le lemme
Compléments sur les morphismes,
\ref{more-morphisms-lemma-integral-closure-smooth-pullback},
donne un isomorphisme canonique
$f^*(\mathcal{A}'_U) = \mathcal{A}'_{U'}$.
Cela montre aussitôt que l'on obtient un sous-préfaisceau
$\mathcal{A}' \subset \mathcal{A}$ de $\mathcal{O}_X$-algèbres sur $X_\etale$
qui est un faisceau pour la topologie de Zariski
et prend les valeurs voulues sur les objets affines.
Le fait que chaque $\mathcal{A}'_U$ soit quasi-cohérent
sur le schéma $U$ et que, pour $f : U' \to U$ étale, on ait
$\mathcal{A}'_{U'} = f^*(\mathcal{A}'_U)$ implique en outre que $\mathcal{A}'$
est quasi-cohérent sur $X_\etale$ (il s'agit en effet d'une propriété locale,
et les références données ci-dessus décrivent précisément de cette manière les modules quasi-cohérents
sur $U_\etale$).
\end{proof}

\begin{definition}
\label{spaces-morphisms-definition-integral-closure}
Soit $S$ un schéma. Soit $X$ un espace algébrique sur $S$.
Soit $\mathcal{A}$ un faisceau quasi-cohérent de $\mathcal{O}_X$-algèbres.
La {\it clôture intégrale de $\mathcal{O}_X$ dans $\mathcal{A}$} est la
$\mathcal{O}_X$-sous-algèbre quasi-cohérente $\mathcal{A}' \subset \mathcal{A}$
construite dans le lemme \ref{spaces-morphisms-lemma-integral-closure} ci-dessus.
\end{definition}

\noindent
Nous appliquerons notamment ceci lorsque $\mathcal{A} = f_*\mathcal{O}_Y$
pour un morphisme quasi-compact et quasi-séparé d'espaces algébriques
$f : Y \to X$ (voir le lemme \ref{spaces-morphisms-lemma-pushforward}). On peut alors prendre
le spectre relatif de la $\mathcal{O}_X$-algèbre quasi-cohérente
(lemme \ref{spaces-morphisms-lemma-affine-equivalence-algebras}) pour obtenir la
normalisation de $X$ dans $Y$.

\begin{definition}
\label{spaces-morphisms-definition-normalization-X-in-Y}
Soit $S$ un schéma. Soit $f : Y \to X$ un morphisme quasi-compact et quasi-séparé
d'espaces algébriques sur $S$. Soit $\mathcal{O}'$ la fermeture
intégrale de $\mathcal{O}_X$ dans $f_*\mathcal{O}_Y$. La {\it normalisation de
$X$ dans $Y$} est le morphisme d'espaces algébriques
$$
\nu : X' = \underline{\Spec}_X(\mathcal{O}') \to X
$$
sur $S$. Elle est munie d'une factorisation naturelle
$$
Y \xrightarrow{f'} X' \xrightarrow{\nu} X
$$
du morphisme initial $f$.
\end{definition}

\noindent
Pour obtenir cette factorisation, on utilise la remarque
\ref{spaces-morphisms-remark-factorization-quasi-compact-quasi-separated}
et la fonctorialité de la construction $\underline{\Spec}$.

\begin{lemma}
\label{spaces-morphisms-lemma-properties-normalization}
Soit $S$ un schéma. Soit $f : Y \to X$ un morphisme quasi-compact et quasi-séparé
d'espaces algébriques sur $S$. Soit $Y \to X' \to X$ la
normalisation de $X$ dans $Y$.
\begin{enumerate}
\item Si $W \to X$ est un morphisme étale d'espaces algébriques sur $S$,
alors $W \times_X X'$ est la normalisation de $W$ dans $W \times_X Y$.
\item Si $Y$ et $X$ sont représentables, alors $X'$ est représentable
et canoniquement isomorphe à la normalisation du schéma $X$
dans le schéma $Y$ construite dans
Morphismes, section \ref{morphisms-section-normalization}.
\end{enumerate}
\end{lemma}

\begin{proof}
Il résulte immédiatement de la construction que la formation
de la normalisation de $X$ dans $Y$ commute au changement de base
étale, ce qui prouve (1). D'autre part, si
$X$ et $Y$ sont des schémas, alors, pour tout ouvert affine $U \subset X$,
$f_*\mathcal{O}_Y(U) = \mathcal{O}_Y(f^{-1}(U))$ ; par conséquent,
$\nu^{-1}(U)$ est le spectre du même anneau que celui
obtenu dans la construction correspondante pour les schémas.
\end{proof}

\noindent
Voici une caractérisation de cette construction.

\begin{lemma}
\label{spaces-morphisms-lemma-characterize-normalization}
Soit $S$ un schéma. Soit $f : Y \to X$ un morphisme quasi-compact et quasi-séparé
d'espaces algébriques sur $S$. La factorisation $f = \nu \circ f'$,
où $\nu : X' \to X$ est la normalisation de $X$ dans $Y$, est caractérisée
par les deux propriétés suivantes :
\begin{enumerate}
\item le morphisme $\nu$ est entier ;
\item pour toute factorisation $f = \pi \circ g$, où $\pi : Z \to X$
est entier, il existe un diagramme commutatif
$$
\xymatrix{
Y \ar[d]_{f'} \ar[r]_g & Z \ar[d]^\pi \\
X' \ar[ru]^h \ar[r]^\nu & X
}
$$
pour un unique morphisme $h : X' \to Z$.
\end{enumerate}
De plus, dans (2), le morphisme $h : X' \to Z$ est la normalisation de
$Z$ dans $Y$.
\end{lemma}

\begin{proof}
Soit $\mathcal{O}' \subset f_*\mathcal{O}_Y$ la clôture intégrale de
$\mathcal{O}_X$ définie comme dans la définition \ref{spaces-morphisms-definition-normalization-X-in-Y}.
Le morphisme $\nu$ est entier par construction, ce qui prouve (1).
Supposons donnée une factorisation $f = \pi \circ g$ avec $\pi : Z \to X$
entier comme en (2). D'après la définition \ref{spaces-morphisms-definition-integral},
$\pi$ est affine ; ainsi, $Z$ est le spectre relatif
d'un faisceau quasi-cohérent de $\mathcal{O}_X$-algèbres $\mathcal{B}$.
Le morphisme $g : Y \to Z$ correspond à un morphisme de $\mathcal{O}_X$-algèbres
$\chi : \mathcal{B} \to f_*\mathcal{O}_Y$. Comme $\mathcal{B}(U)$ est
entière sur $\mathcal{O}_X(U)$ pour tout objet affine $U$ étale sur $X$
(d'après la définition \ref{spaces-morphisms-definition-integral}),
le lemme \ref{spaces-morphisms-lemma-integral-closure} montre
que $\chi(\mathcal{B}) \subset \mathcal{O}'$.
Par fonctorialité du spectre relatif,
le lemme \ref{spaces-morphisms-lemma-affine-equivalence-algebras}
fournit alors un unique morphisme
$h : X' \to Z$. Nous omettons la vérification de la commutativité du diagramme.

\medskip\noindent
Il est clair que (1) et (2) caractérisent la
factorisation $f = \nu \circ f'$, puisqu'ils la caractérisent
comme objet initial d'une catégorie. Le morphisme $h$ de (2)
est entier d'après le lemme \ref{spaces-morphisms-lemma-finite-permanence}.
Étant donnée une factorisation $g = \pi' \circ g'$ avec $\pi' : Z' \to Z$
entier, on obtient une factorisation $f = (\pi \circ \pi') \circ g'$ et
un morphisme $h' : X' \to Z'$. L'unicité implique que
$\pi' \circ h' = h$. La caractérisation (1), (2) s'applique donc
au morphisme $h : X' \to Z$, ce qui donne la dernière assertion du lemme.
\end{proof}

\begin{lemma}
\label{spaces-morphisms-lemma-normalization-in-reduced}
Soit $S$ un schéma. Soit $f : Y \to X$ un morphisme quasi-compact et
quasi-séparé d'espaces algébriques sur $S$.
Soit $X' \to X$ la normalisation de $X$ dans $Y$.
Si $Y$ est réduit, alors $X'$ l'est aussi.
\end{lemma}

\begin{proof}
Cela résulte du fait qu'un sous-anneau d'un anneau réduit est réduit.
Nous omettons quelques détails.
\end{proof}

\begin{lemma}
\label{spaces-morphisms-lemma-normalization-generic}
Soit $S$ un schéma. Soit $f : Y \to X$ un morphisme quasi-compact et quasi-séparé
de schémas. Soit $X' \to X$ la normalisation de $X$ dans $Y$.
Si $x' \in |X'|$ est un point de codimension $0$
(Propriétés des espaces, définition
\ref{spaces-properties-definition-dimension-local-ring}), alors
$x'$ est l'image d'un point $y \in |Y|$ de codimension $0$.
\end{lemma}

\begin{proof}
D'après le lemme \ref{spaces-morphisms-lemma-properties-normalization} et les définitions, on peut
supposer que $X = \Spec(A)$ est affine. Alors $X' = \Spec(A')$, où $A'$ est
la clôture intégrale de $A$ dans $\Gamma(Y, \mathcal{O}_Y)$, et $x'$ correspond
à un idéal premier minimal de $A'$. Choisissons un morphisme étale surjectif
$V \to Y$, où $V = \Spec(B)$ est affine. Alors
$A' \to B$ est injectif ; tout idéal premier minimal de $A'$ est donc
l'image d'un idéal premier minimal de $B$ ; voir
Algèbre, lemme \ref{algebra-lemma-injective-minimal-primes-in-image}.
Le résultat en découle.
\end{proof}

\begin{lemma}
\label{spaces-morphisms-lemma-normalization-in-disjoint-union}
Soit $S$ un schéma.
Soit $f : Y \to X$ un morphisme quasi-compact et quasi-séparé
d'espaces algébriques sur $S$.
Supposons que $Y = Y_1 \amalg Y_2$ soit la somme disjointe de deux
espaces algébriques.
Écrivons $f_i = f|_{Y_i}$. Soit $X_i'$ la normalisation de $X$ dans $Y_i$.
Alors $X_1' \amalg X_2'$ est la normalisation de $X$ dans $Y$.
\end{lemma}

\begin{proof}
Démonstration omise.
\end{proof}

\begin{lemma}
\label{spaces-morphisms-lemma-normalization-in-universally-closed}
Soit $S$ un schéma.
Soit $f : X \to Y$ un morphisme quasi-compact, quasi-séparé et
universellement fermé d'espaces algébriques sur $S$.
Alors $f_*\mathcal{O}_X$ est entière sur $\mathcal{O}_Y$. Autrement
dit, la normalisation de $Y$ dans $X$ est égale à la factorisation
$$
X \longrightarrow \underline{\Spec}_Y(f_*\mathcal{O}_X)
\longrightarrow Y
$$
de la remarque \ref{spaces-morphisms-remark-factorization-quasi-compact-quasi-separated}.
\end{lemma}

\begin{proof}
La question est locale sur $Y$ pour la topologie étale ; on peut donc se ramener au
cas où $Y = \Spec(R)$ est affine. Soit $h \in \Gamma(X, \mathcal{O}_X)$.
Il faut montrer que $h$ satisfait une équation unitaire sur $R$. Considérons $h$
comme un morphisme dans le diagramme commutatif suivant :
$$
\xymatrix{
X \ar[rr]_h \ar[rd]_f & & \mathbf{A}^1_Y \ar[ld] \\
& Y &
}
$$
Soit $Z \subset \mathbf{A}^1_Y$ l'image schématique de $h$ ;
voir la définition \ref{spaces-morphisms-definition-scheme-theoretic-image}.
Le morphisme $h$ est quasi-compact, car $f$ est quasi-compact et
$\mathbf{A}^1_Y \to Y$ est séparé ; voir
le lemme \ref{spaces-morphisms-lemma-quasi-compact-permanence}.
D'après le lemme \ref{spaces-morphisms-lemma-quasi-compact-scheme-theoretic-image}, le
morphisme $X \to Z$ a une image dense sur les espaces topologiques sous-jacents. D'après
le lemme \ref{spaces-morphisms-lemma-universally-closed-permanence}, le morphisme
$X \to Z$ est fermé. Ainsi $h(X) = Z$ ensemblistement.
On peut donc appliquer
le lemme \ref{spaces-morphisms-lemma-image-proper-is-proper}
pour conclure que $Z \to Y$ est universellement fermé (et même propre).
Comme $Z \subset \mathbf{A}^1_Y$, le morphisme $Z \to Y$ est affine
et propre, donc entier d'après le lemme \ref{spaces-morphisms-lemma-integral-universally-closed}.
En écrivant $\mathbf{A}^1_Y = \Spec(R[T])$, on en déduit que
l'idéal $I \subset R[T]$ de $Z$ contient un polynôme unitaire
$P(T) \in R[T]$. Ainsi $P(h) = 0$, ce qui prouve le résultat.
\end{proof}

\begin{lemma}
\label{spaces-morphisms-lemma-normalization-in-integral}
Soit $S$ un schéma. Soit $f : Y \to X$ un morphisme entier
d'espaces algébriques sur $S$.
Alors la normalisation de $X$ dans $Y$ est égale à $Y$.
\end{lemma}

\begin{proof}
D'après le lemme \ref{spaces-morphisms-lemma-integral-universally-closed}, c'est un cas particulier du
lemme \ref{spaces-morphisms-lemma-normalization-in-universally-closed}.
\end{proof}

\begin{lemma}
\label{spaces-morphisms-lemma-nagata-normalization-finite}
Soit $S$ un schéma. Soit $f : X \to Y$ un morphisme d'espaces algébriques
sur $S$. Supposons que
\begin{enumerate}
\item $Y$ soit de Nagata ;
\item $f$ soit quasi-séparé et de type fini ;
\item $X$ soit réduit.
\end{enumerate}
Alors la normalisation $\nu : Y' \to Y$ de $Y$ dans $X$ est finie.
\end{lemma}

\begin{proof}
La question est locale sur $Y$ pour la topologie étale ; voir
le lemme \ref{spaces-morphisms-lemma-properties-normalization}.
On peut donc supposer que $Y = \Spec(R)$ est affine.
Alors $R$ est un anneau de Nagata noethérien,
et il faut montrer que la clôture intégrale de
$R$ dans $\Gamma(X, \mathcal{O}_X)$ est finie sur $R$.
Comme $f$ est quasi-compact, l'espace $X$ est quasi-compact.
Choisissons un schéma affine $U$ et un morphisme étale surjectif
$U \to X$ (Propriétés des espaces, lemme
\ref{spaces-properties-lemma-quasi-compact-affine-cover}).
Alors $\Gamma(X, \mathcal{O}_X) \subset \Gamma(U, \mathcal{O}_U)$.
Comme $R$ est noethérien, il suffit de montrer que
la clôture intégrale de $R$ dans $\Gamma(U, \mathcal{O}_U)$
est finie sur $R$. Puisque $U \to Y$ est de type fini,
cela résulte de
Morphismes, lemme \ref{morphisms-lemma-nagata-normalization-finite}.
\end{proof}







\section{Normalisation}
\label{spaces-morphisms-section-normalization}

\noindent
Cette section est l'analogue de
Morphismes, section \ref{morphisms-section-normalization}.

\begin{lemma}
\label{spaces-morphisms-lemma-prepare-normalization}
Soit $S$ un schéma. Soit $X$ un espace algébrique sur $S$.
Les conditions suivantes sont équivalentes :
\begin{enumerate}
\item il existe un morphisme étale surjectif $U \to X$, où $U$
est un schéma tel que tout ouvert quasi-compact de $U$ ait
un nombre fini de composantes irréductibles ;
\item pour tout schéma $U$ et tout morphisme étale
$U \to X$, tout ouvert quasi-compact de $U$ a un nombre fini de
composantes irréductibles ;
\item pour tout espace algébrique quasi-compact $Y$ étale sur $X$,
l'ensemble des points de codimension $0$ de $Y$ (Propriétés des espaces,
définition \ref{spaces-properties-definition-dimension-local-ring})
est fini ;
\item pour tout espace algébrique quasi-compact $Y$ étale sur $X$,
l'espace $|Y|$ a un nombre fini de composantes irréductibles.
\end{enumerate}
Si $X$ est représentable, cela signifie que tout ouvert quasi-compact de $X$
a un nombre fini de composantes irréductibles.
\end{lemma}

\begin{proof}
L'équivalence de (1) et (2), ainsi que la dernière assertion, résultent de
Descente, lemme \ref{descent-lemma-locally-finite-nr-irred-local-fppf}, et de
Propriétés des espaces, lemme \ref{spaces-properties-lemma-type-property}.
Il est clair que (4) implique (1) et (2) en ne considérant que les $Y$
qui sont des schémas. De même, (3) implique (1) et (2), puisque, pour un schéma,
les points de codimension $0$ sont les points génériques de ses
composantes irréductibles ; voir par exemple
Propriétés des espaces, lemme \ref{spaces-properties-lemma-codimension-0-points}.

\medskip\noindent
Réciproquement, supposons (2) et soit $Y \to X$ un morphisme étale d'espaces
algébriques, avec $Y$ quasi-compact. On peut alors choisir un schéma
affine $V$ et un morphisme étale surjectif $V \to Y$
(Propriétés des espaces, lemme
\ref{spaces-properties-lemma-quasi-compact-affine-cover}).
Comme $V$ a un nombre fini de composantes irréductibles d'après (2) et que
$|V| \to |Y|$ est surjectif et continu, on conclut que
$|Y|$ a un nombre fini de composantes irréductibles d'après Topologie, lemme
\ref{topology-lemma-surjective-continuous-irreducible-components}.
La condition (4) est donc satisfaite. De même, d'après
Propriétés des espaces, lemme \ref{spaces-properties-lemma-codimension-0-points},
les images des points génériques des composantes irréductibles de $V$
sont les points de codimension $0$ de $Y$ ; ils sont donc
en nombre fini, ce qui prouve (3).
\end{proof}

\begin{lemma}
\label{spaces-morphisms-lemma-locally-noetherian-can-be-normalized}
Soit $S$ un schéma. Soit $X$ un espace algébrique localement noethérien
sur $S$. Alors $X$ satisfait aux conditions équivalentes du
lemme \ref{spaces-morphisms-lemma-prepare-normalization}.
\end{lemma}

\begin{proof}
Si $U \to X$ est étale et si $U$ est un schéma, alors $U$ est un schéma
localement noethérien ; voir Propriétés des espaces, section
\ref{spaces-properties-section-types-properties}.
L'ensemble des composantes irréductibles d'un schéma localement noethérien est
localement fini (Diviseurs, lemme \ref{divisors-lemma-components-locally-finite}).
On en conclut que $X$ satisfait à la condition (2) du lemme.
\end{proof}

\begin{lemma}
\label{spaces-morphisms-lemma-flat-image-point-codimension-0}
Soit $S$ un schéma. Soit $f : X \to Y$ un morphisme plat d'espaces
algébriques sur $S$. Alors, pour $x \in |X|$, on a : $x$ est de codimension $0$ dans
$X \Rightarrow f(x)$ est de codimension $0$ dans $Y$.
\end{lemma}

\begin{proof}
À l'aide de Propriétés des espaces, lemme
\ref{spaces-properties-lemma-codimension-0-points}, et par localisation étale,
on se ramène au cas d'un morphisme de schémas et de points génériques
de composantes irréductibles. Le résultat découle alors de ce que les générisations
se relèvent le long des morphismes plats de schémas ; voir
Morphismes, lemme \ref{morphisms-lemma-generalizations-lift-flat}.
\end{proof}

\begin{lemma}
\label{spaces-morphisms-lemma-flat-locally-finite-type-points-codimension-0}
Soit $S$ un schéma. Soit $f : X \to Y$ un morphisme d'espaces
algébriques sur $S$. Supposons $f$ plat et localement de type fini, et
supposons que $Y$ satisfasse aux conditions équivalentes du
lemme \ref{spaces-morphisms-lemma-prepare-normalization}.
Alors $X$ satisfait aux conditions équivalentes du
lemme \ref{spaces-morphisms-lemma-prepare-normalization} et, pour $x \in |X|$, on a :
$x$ est de codimension $0$ dans $X \Rightarrow f(x)$ est de codimension $0$ dans $Y$.
\end{lemma}

\begin{proof}
La dernière assertion résulte du
lemme \ref{spaces-morphisms-lemma-flat-image-point-codimension-0}.
Choisissons un morphisme étale surjectif $V \to Y$, où $V$ est un schéma.
Choisissons un morphisme étale surjectif $U \to X \times_Y V$, où $U$
est un schéma. Il suffit de montrer que tout ouvert quasi-compact
de $U$ a un nombre fini de composantes irréductibles.
Nous utiliserons sans autre mention les résultats de
Propriétés des espaces, lemme \ref{spaces-properties-lemma-codimension-0-points}.
D'après ce qui précède,
les points de codimension $0$ de $U$ se trouvent au-dessus de
points de codimension $0$ de $V$, lesquels sont localement en nombre fini par
hypothèse. Il suffit donc de montrer que, pour $v \in V$ de codimension $0$,
les points de codimension $0$ de la fibre schématique $U_v = U \times_V v$
sont localement en nombre fini. Cela résulte de ce que $U_v$ est un schéma
localement de type fini sur $\kappa(v)$, donc localement noethérien ;
et l'on peut appliquer le lemme \ref{spaces-morphisms-lemma-locally-noetherian-can-be-normalized}
par exemple.
\end{proof}

\begin{lemma}
\label{spaces-morphisms-lemma-normalization}
Soit $S$ un schéma. Pour tout espace algébrique $X$ sur $S$ satisfaisant
aux conditions équivalentes du lemme \ref{spaces-morphisms-lemma-prepare-normalization},
il existe un morphisme d'espaces algébriques
$$
\nu_X : X^\nu \longrightarrow X
$$
ayant les propriétés suivantes :
\begin{enumerate}
\item si $X$ satisfait aux conditions équivalentes du
lemme \ref{spaces-morphisms-lemma-prepare-normalization}, alors
$X^\nu$ est normal et $\nu_X$ est entier ;
\item si $X$ est un schéma dont tout ouvert quasi-compact a un nombre fini
de composantes irréductibles, alors $\nu_X : X^\nu \to X$ est la
normalisation de $X$ construite dans
Morphismes, section \ref{morphisms-section-normalization} ;
\item si $f : X \to Y$ est un morphisme d'espaces algébriques sur $S$
dont la source et le but satisfont aux conditions équivalentes du
lemme \ref{spaces-morphisms-lemma-prepare-normalization}, et si tout point de codimension $0$
de $X$ est envoyé par $f$ sur un point de codimension $0$ de $Y$, alors
il existe un unique morphisme $f^\nu : X^\nu \to Y^\nu$ d'espaces
algébriques sur $S$ tel que $\nu_Y \circ f^\nu = f \circ \nu_X$ ;
\item si $f : X \to Y$ est un morphisme étale ou lisse d'espaces
algébriques et si $Y$ satisfait aux conditions équivalentes du
lemme \ref{spaces-morphisms-lemma-prepare-normalization}, alors les hypothèses de (3)
sont satisfaites et le morphisme $f^\nu$ induit un isomorphisme
$X^\nu  \to X \times_Y Y^\nu$.
\end{enumerate}
\end{lemma}

\begin{proof}
Considérons la catégorie $\mathcal{C}$ dont les objets sont les
schémas $U$ sur $S$ tels que tout ouvert quasi-compact de
$U$ ait un nombre fini de composantes irréductibles, et dont les morphismes
sont les morphismes $g : U \to V$ de schémas sur $S$ tels
que tout point générique d'une composante irréductible de $U$
soit envoyé sur le point générique d'une composante irréductible de $V$.
Nous avons déjà montré que :
\begin{enumerate}
\item[(a)] pour $U \in \Ob(\mathcal{C})$, on dispose d'un
morphisme de normalisation $\nu_U : U^\nu \to U$ comme dans
Morphismes, définition \ref{morphisms-definition-normalization} ;
\item[(b)] pour $U \in \Ob(\mathcal{C})$, le morphisme $\nu_U$
est entier et $U^\nu$ est un schéma normal ; voir
Morphismes, lemme \ref{morphisms-lemma-normalization-normal} ;
\item[(c)] pour tout $g : U \to V \in \text{Arrows}(\mathcal{C})$,
il existe un unique morphisme $g^\nu : U^\nu \to V^\nu$
tel que $\nu_V \circ g^\nu = g \circ \nu_U$ ; voir
Morphismes, lemme \ref{morphisms-lemma-normalization-normal},
partie (4), appliqué à la composée $U^\nu \to U \to V$ ;
\item[(d)] si $V \in \Ob(\mathcal{C})$ et si $g : U \to V$ est
étale ou lisse, alors $U \in \Ob(\mathcal{C})$ et
$g \in \text{Arrows}(\mathcal{C})$, et le morphisme $g^\nu$
induit un isomorphisme $U^\nu \to U \times_V V^\nu$ ; voir le
lemme \ref{spaces-morphisms-lemma-flat-locally-finite-type-points-codimension-0}
et Compléments sur les morphismes, lemme
\ref{more-morphisms-lemma-normalization-and-smooth}.
\end{enumerate}
Il s'agit d'étendre cette construction à la catégorie correspondante
des espaces algébriques $X$ sur $S$.

\medskip\noindent
Soit $X$ un espace algébrique sur $S$ satisfaisant
aux conditions équivalentes du lemme \ref{spaces-morphisms-lemma-prepare-normalization}.
Soit $U \to X$ un morphisme étale surjectif, où $U$ est un schéma.
Posons $R = U \times_X U$, avec les projections $s, t : R \to U$, et
$j = (t, s) : R \to U \times_S U$, de sorte que $X = U/R$ ; voir
Espaces, lemme \ref{spaces-lemma-space-presentation}.
Observons que $U$ et $R$ sont des objets de $\mathcal{C}$
en vertu de nos hypothèses sur $X$, et que les morphismes
$s$ et $t$ sont des morphismes étales de schémas sur $S$.
D'après (a), on dispose des morphismes de normalisation
$\nu_U : U^\nu \to U$ et $\nu_R : R^\nu \to R$ ;
d'après (d), on dispose de morphismes $s^\nu : R^\nu \to U^\nu$ et
$t^\nu : R^\nu \to U^\nu$ qui définissent des
isomorphismes $R^\nu \to R \times_{s, U} U^\nu$ et
$R^\nu \to U^\nu \times_{U, t} R$. Il en résulte que
$s^\nu$ et $t^\nu$ sont étales (car ils sont isomorphes
à des changements de base de morphismes étales). Le morphisme induit
$j^\nu = (t^\nu, s^\nu) : R^\nu \to U^\nu \times_S U^\nu$
est un monomorphisme, car il est égal à la composée
\begin{align*}
R^\nu
& \to
(U^\nu \times_{U, t} R) \times_R (R \times_{s, U} U^\nu) \\
& =
U^\nu \times_{U, t} R \times_{s, U} U^\nu \\
& \xrightarrow{j}
U^\nu \times_U (U \times_S U) \times_U U^\nu \\
& =
U^\nu \times_S U^\nu
\end{align*}
Le premier morphisme est le morphisme diagonal de $\nu_R$.
(Cela montre que $R^\nu$ est un sous-schéma de la restriction de
$R$ à $U^\nu$.) Un calcul formel de produits fibrés
utilisant la propriété (d) montre que
$R^\nu \times_{s^\nu, U^\nu, t^\nu} R^\nu$ est la normalisation
de $R \times_{s, U, t} R$. Le morphisme étale
$c : R \times_{s, U, t} R \to R$ se prolonge donc de manière unique en $c^\nu$ d'après (d).
Le morphisme $c^\nu$ est compatible avec la projection
$\text{pr}_{13} : U^\nu \times_S U^\nu \times_S U^\nu \to U^\nu \times_S U^\nu$.
De même, il existe un morphisme $i^\nu : R^\nu \to R^\nu$
compatible avec le morphisme $U^\nu \times_S U^\nu \to U^\nu \times_S U^\nu$
qui échange les facteurs, et un morphisme $e^\nu : U^\nu \to R^\nu$
compatible avec le morphisme diagonal $U^\nu \to U^\nu \times_S U^\nu$.
Il en résulte que $j^\nu : R^\nu \to U^\nu \times_S U^\nu$
est une relation d'équivalence étale.
On peut alors poser, et l'on pose, $X^\nu = U^\nu/R^\nu$
(Espaces, théorème \ref{spaces-theorem-presentation}).
On voit ensuite que $U^\nu = X^\nu \times_X U$ d'après
Groupoïdes, lemme \ref{groupoids-lemma-criterion-fibre-product}.

\medskip\noindent
Voici ce que nous avons montré au paragraphe précédent : pour tout espace
algébrique $X$ sur $S$ satisfaisant aux conditions équivalentes du
lemme \ref{spaces-morphisms-lemma-prepare-normalization}, si l'on choisit un morphisme étale
surjectif $g : U \to X$, où $U$ est un schéma, on obtient
un diagramme cartésien
$$
\xymatrix{
X^\nu \ar[d]_{\nu_X} & U^\nu \ar[l]^{g^\nu} \ar[d]^{\nu_U} \\
X & U \ar[l]_g
}
$$
d'espaces algébriques. Il en résulte immédiatement que $X^\nu$ est un espace
algébrique normal et que $\nu_X$ est un morphisme entier. Cela prouve
la partie (1) du lemme.

\medskip\noindent
Nous montrerons plus bas que le morphisme $\nu_X : X^\nu \to X$
est, à isomorphisme unique près, indépendant du choix de $g$ ;
mais, pour l'instant, si $X$ est un schéma, choisissons $\text{id} : X \to X$,
ce qui établit aussitôt la partie (2) du lemme.

\medskip\noindent
Il reste à prouver les parties (3) et (4). Soient $g : U \to X$,
$\nu_X : X^\nu \to X$ et $g^\nu : U^\nu \to X^\nu$ comme ci-dessus.
Soit $Z$ un schéma normal, et soient $h : Z \to U$ et $a : Z \to X^\nu$
des morphismes sur $S$ tels que $g \circ h = \nu_X \circ a$ et que
toute composante irréductible de $Z$ domine une composante irréductible de $U$
(par $h$). D'après
Morphismes, lemme \ref{morphisms-lemma-normalization-normal}, partie (4),
on obtient un unique morphisme $h^\nu : Z \to U^\nu$ tel
que $h = \nu_U \circ h^\nu$. Voici le diagramme :
$$
\xymatrix{
X^\nu \ar[d]_{\nu_X} &
U^\nu \ar[l]^{g^\nu} \ar[d]^{\nu_U} &
Z \ar[l]^{h^\nu} \ar@/_1em/[ll]_a \ar[dl]^h
\\
X & U \ar[l]_g
}
$$
Observons que $a = g^\nu \circ h^\nu$. En effet, puisque le carré de sommets
$X^\nu$, $X$, $U^\nu$, $U$ est cartésien, cela résulte
immédiatement de l'unicité de $h^\nu$ une fois $h$ donné.
Autrement dit, étant donné $h : Z \to U$ comme ci-dessus (mais non $a$),
il existe un unique morphisme $a : Z \to X^\nu$ tel que
$\nu_X \circ a = g \circ h$.

\medskip\noindent
Soit $f : X \to Y$ comme dans la partie (3) de l'énoncé du lemme.
Supposons choisis des morphismes étales surjectifs $U \to X$ et
$V \to Y$, où $U$ et $V$ sont des schémas, tels que $f$ se relève en un
morphisme $g : U \to V$. Alors $g \in \text{Arrows}(\mathcal{C})$
et l'on obtient un unique morphisme $g^\nu : U^\nu \to V^\nu$ compatible
avec $\nu_U$ et $\nu_V$. Les deux morphismes
$$
R^\nu = U^\nu \times_{X^\nu} U^\nu
\to U^\nu \to V^\nu \to Y^\nu
$$
sont donc égaux d'après les observations du paragraphe précédent (appliquées
avec $Y$ à la place de $X$).
Comme $X^\nu$ est construit comme le quotient de
$U^\nu$ par $R^\nu$, on obtient ainsi un unique
morphisme $f^\nu : X^\nu \to Y^\nu$ comme dans (3).

\medskip\noindent
Pour voir que la construction de $X^\nu$ est indépendante du choix
du morphisme étale surjectif $g : U \to X$, appliquons la construction du
paragraphe précédent à $\text{id} : X \to X$ et à un morphisme
$U' \to U$ entre revêtements étales de $X$. Cela suffit, car,
étant donnés deux revêtements étales de $X$, il en existe un troisième qui
les domine tous deux. On vérifie que le morphisme entre les deux
normalisations construites à l'aide de $U' \to X$ et de $U \to X$
devient un isomorphisme après changement de base à $U'$ et qu'il était donc
déjà un isomorphisme. Nous omettons les détails.

\medskip\noindent
Nous omettons la démonstration de (4), qui est analogue ; indication : utiliser la partie (d) ci-dessus.
\end{proof}

\noindent
Cela conduit à la définition suivante.

\begin{definition}
\label{spaces-morphisms-definition-normalization}
Soit $S$ un schéma. Soit $X$ un espace algébrique sur $S$ satisfaisant aux
conditions équivalentes du lemme \ref{spaces-morphisms-lemma-prepare-normalization}.
On appelle {\it normalisation} de $X$ le morphisme
$$
\nu_X : X^\nu \longrightarrow X
$$
construit dans le lemme \ref{spaces-morphisms-lemma-normalization}.
\end{definition}

\noindent
La définition s'applique aux espaces algébriques localement
noethériens ; voir le lemme \ref{spaces-morphisms-lemma-locally-noetherian-can-be-normalized}.
Habituellement, la normalisation n'est définie que pour les espaces algébriques
réduits. Avec la définition ci-dessus, la normalisation de $X$ est la même
que celle de la réduction $X_{red}$ de $X$.

\begin{lemma}
\label{spaces-morphisms-lemma-normalization-reduced}
Soit $S$ un schéma. Soit $X$ un espace algébrique sur $S$ satisfaisant aux
conditions équivalentes du lemme \ref{spaces-morphisms-lemma-prepare-normalization}.
Le morphisme de normalisation $\nu$ se factorise par la réduction $X_{red}$,
et $X^\nu \to X_{red}$ est la normalisation de $X_{red}$.
\end{lemma}

\begin{proof}
On peut vérifier ceci localement sur $X$ pour la topologie étale et se ramener ainsi au cas
des schémas, qui est
Morphismes, lemme \ref{morphisms-lemma-normalization-reduced}.
Nous omettons quelques détails.
\end{proof}

\begin{lemma}
\label{spaces-morphisms-lemma-normalization-normal}
Soit $S$ un schéma. Soit $X$ un espace algébrique sur $S$ satisfaisant aux
conditions équivalentes du lemme \ref{spaces-morphisms-lemma-prepare-normalization}.
\begin{enumerate}
\item La normalisation $X^\nu$ est normale.
\item Le morphisme $\nu : X^\nu \to X$ est entier et surjectif.
\item L'application $|\nu| : |X^\nu| \to |X|$ induit une bijection entre
les ensembles de points de codimension $0$ (Propriétés des espaces,
définition \ref{spaces-properties-definition-dimension-local-ring}).
\item Soit $f : Z \to X$ un morphisme. Supposons que $Z$ soit un espace algébrique normal
et que, pour $z \in |Z|$, on ait : $z$ est de codimension $0$ dans
$Z \Rightarrow f(z)$ est de codimension $0$ dans $X$. Alors
il existe une unique factorisation $Z \to X^\nu \to X$.
\end{enumerate}
\end{lemma}

\begin{proof}
Les propriétés (1), (2) et (3) résultent des énoncés correspondants
pour les schémas (Morphismes, lemme \ref{morphisms-lemma-normalization-normal}),
combinés au fait qu'un point d'un schéma est un point
générique d'une composante irréductible si et seulement si la dimension
de son anneau local est nulle
(Propriétés, lemme \ref{properties-lemma-generic-point}).

\medskip\noindent
Soit $Z \to X$ un morphisme comme en (4). Soit $U$ un schéma, et soit
$U \to X$ un morphisme étale surjectif. Choisissons un schéma $V$ et
un morphisme étale surjectif $V \to U \times_X Z$. La condition sur les
points de codimension $0$ assure que $V \to U$ envoie les points génériques des
composantes irréductibles de $V$ sur des points génériques de composantes
irréductibles de $U$. On obtient donc une unique factorisation
$V \to U^\nu \to U$ d'après
Morphismes, lemme \ref{morphisms-lemma-normalization-normal}.
L'unicité garantit que les deux applications
$$
V \times_{U \times_X Z} V \to V \to U^\nu
$$
coïncident, car elles donnent l'unique factorisation de l'application
$V \times_{U \times_X Z} V \to V \to U$.
Comme l'espace algébrique $U \times_X Z$ est égal au quotient
$V/V \times_{U \times_X Z} V$
(voir Espaces, section \ref{spaces-section-presentations}),
on obtient un morphisme canonique
$U \times_X Z \to U^\nu$. Voici le diagramme :
$$
\xymatrix{
U \times_X Z \ar[r] \ar[d] & U^\nu \ar[r] \ar[d] & U \ar[d] \\
Z \ar@/_/[rr] \ar@{..>}[r] & X^\nu \ar[r] & X
}
$$
Pour obtenir la flèche pointillée, remarquons que la construction de la
flèche $U \times_X Z \to U^\nu$ est fonctorielle par rapport au morphisme étale $U \to X$
(nous omettons la formulation précise et la démonstration).
Posons donc $R = U \times_X U$, avec les projections $s, t : R \to U$ ;
on obtient alors un morphisme $R \times_X Z \to R^\nu$ compatible avec
$s, t : R \to U$ et $s^\nu, t^\nu : R^\nu \to U^\nu$.
Rappelons que $X^\nu = U^\nu/R^\nu$ ; voir la démonstration du
lemme \ref{spaces-morphisms-lemma-normalization}. Comme $X = U/R$, un argument simple
de théorie des faisceaux
montre que $Z = (U \times_X Z)/(R \times_X Z)$. Ainsi, les morphismes
$U \times_X Z \to U^\nu$ et $R \times_X Z \to R^\nu$ définissent un
morphisme $Z \to X^\nu$, comme voulu.
\end{proof}

\begin{lemma}
\label{spaces-morphisms-lemma-nagata-normalization}
Soit $S$ un schéma. Soit $X$ un espace algébrique de Nagata sur $S$.
La normalisation $\nu : X^\nu \to X$ est un morphisme fini.
\end{lemma}

\begin{proof}
Comme $X$ est de Nagata, il est localement noethérien,
et la définition \ref{spaces-morphisms-definition-normalization} s'applique.
Par la construction de $X^\nu$ dans le lemme \ref{spaces-morphisms-lemma-normalization},
on se ramène immédiatement au cas des schémas, qui est
Morphismes, lemme \ref{morphisms-lemma-nagata-normalization}.
\end{proof}












\section{Morphismes séparés et localement quasi-finis}
\label{spaces-morphisms-section-schemehood}

\noindent
Dans cette section, nous montrons qu'un espace algébrique qui est localement
quasi-fini et séparé sur un schéma est représentable. Il en résulte
qu'un morphisme séparé et localement quasi-fini est
représentable (voir le
lemme \ref{spaces-morphisms-lemma-locally-quasi-finite-separated-representable}).
Mais d'abord...\ un lemme (qui sera rendu caduc par la
proposition \ref{spaces-morphisms-proposition-locally-quasi-finite-separated-over-scheme}).

\begin{lemma}
\label{spaces-morphisms-lemma-neighbourhood-scheme}
Soit $S$ un schéma. Considérons un diagramme commutatif
$$
\xymatrix{
V' \ar[r] \ar[rd] & T' \times_T X \ar[r] \ar[d] & X \ar[d] \\
& T' \ar[r] & T
}
$$
d'espaces algébriques sur $S$. Supposons que
\begin{enumerate}
\item $T' \to T$ soit un morphisme étale de schémas affines,
\item $X \to T$ soit un morphisme séparé et localement quasi-fini,
\item $V'$ soit un sous-espace ouvert de $T' \times_T X$, et
\item $V' \to T'$ soit quasi-affine.
\end{enumerate}
Alors l'image $U$ de $V'$ dans $X$ est un sous-espace ouvert
quasi-compact et représentable de $X$.
\end{lemma}

\begin{proof}
Commençons par quelques observations immédiates.
Le lemme \ref{spaces-morphisms-lemma-quasi-affine-local} montre que $V'$ est représentable.
Il est aussi quasi-compact (comme schéma quasi-affine sur un schéma affine ; voir
Morphismes, lemme \ref{morphisms-lemma-quasi-affine-separated}).
Puisque $T' \times_T X \to X$ est étale
(Propriétés des espaces, lemme \ref{spaces-properties-lemma-base-change-etale}),
l'application $|T' \times_T X| \to |X|$ est ouverte ; voir
Propriétés des espaces, lemme \ref{spaces-properties-lemma-etale-open}.
Soit $U \subset X$ le sous-espace ouvert correspondant à l'image de
$|V'|$ ; voir
Propriétés des espaces, lemme \ref{spaces-properties-lemma-open-subspaces}.
Comme $|V'|$ est quasi-compact, $|U|$ l'est aussi ; par conséquent,
$U$ est un espace algébrique quasi-compact d'après
Propriétés des espaces, lemme \ref{spaces-properties-lemma-quasi-compact-space}.

\medskip\noindent
D'après
Morphismes,
lemme \ref{morphisms-lemma-locally-quasi-finite-qc-source-universally-bounded},
le morphisme $T' \to T$ est universellement borné. On peut donc raisonner par
récurrence sur l'entier $n$ qui borne le degré des fibres de $T' \to T$ ; voir
Morphismes, lemme \ref{morphisms-lemma-etale-universally-bounded}
pour une description de cet entier dans le cas d'un morphisme étale.
Si $n = 1$, alors $T' \to T$ est une immersion ouverte (voir
Morphismes étales, théorème \ref{etale-theorem-etale-radicial-open}),
et le résultat est clair. Supposons $n > 1$.

\medskip\noindent
Considérons le schéma affine $T'' = T' \times_T T'$.
Comme $T' \to T$ est étale, on a une décomposition (en sous-schémas affines
ouverts et fermés) $T'' = \Delta(T') \amalg T^*$. En effet, $\Delta = \Delta_{T'/T}$
est ouverte d'après
Morphismes, lemme \ref{morphisms-lemma-diagonal-unramified-morphism},
et fermée parce que $T' \to T$ est séparé en tant que morphisme de schémas affines.
Par changement de base, les degrés des fibres de la seconde projection
$\text{pr}_1 : T' \times_T T' \to T'$ sont bornés par $n$ ; voir
Morphismes, lemme \ref{morphisms-lemma-base-change-universally-bounded}.
D'autre part, $\text{pr}_1|_{\Delta(T')} : \Delta(T') \to T'$ est
un isomorphisme, et chacune de ses fibres a exactement un point.
En appliquant donc
Morphismes, lemme \ref{morphisms-lemma-etale-universally-bounded},
on conclut que les degrés des fibres de la restriction
$\text{pr}_1|_{T^*} : T^* \to T'$ sont bornés par $n - 1$.
L'hypothèse de récurrence appliquée au diagramme
$$
\xymatrix{
p_0^{-1}(V') \cap X^* \ar[r] \ar[rd] &
X^* \ar[r]_{p_1|_{X^*}} \ar[d] &
X' \ar[d] \\
& T^* \ar[r]^{\text{pr}_1|_{T^*}} & T'
}
$$
montre que $p_1(p_0^{-1}(V') \cap X^*)$
est un schéma quasi-compact. Nous avons posé ici
$X'' = T'' \times_T X$, $X^* = T^* \times_T X$ et $X' = T' \times_T X$,
et $p_0, p_1 : X'' \to X'$ sont les changements de base de $\text{pr}_0, \text{pr}_1$.
La plupart des hypothèses du lemme entraînent,
par changement de base, l'hypothèse correspondante pour le diagramme ci-dessus.
Par exemple, $p_0^{-1}(V') = T'' \times_{T'} V'$
est un schéma quasi-affine sur $T''$, car il s'obtient par changement de base. Nous
omettons quelques vérifications.

\medskip\noindent
D'après
Propriétés des espaces, lemme \ref{spaces-properties-lemma-subscheme},
on en conclut que
$$
p_1(p_0^{-1}(V')) =
V' \cup p_1(p_0^{-1}(V') \cap X^*)
$$
est un schéma quasi-compact. De plus, il est clair que
$p_1(p_0^{-1}(V'))$ est l'image réciproque du
sous-espace ouvert quasi-compact $U \subset X$ considéré dans le
premier paragraphe de la démonstration. Autrement dit, $T' \times_T U$ est un schéma !
Notons que $T' \times_T U$ est quasi-compact,
séparé et localement quasi-fini sur $T'$, puisque
$T' \times_T X \to T'$ est localement quasi-fini et séparé,
étant obtenu par changement de base du morphisme initial $X \to T$ (voir les
lemmes \ref{spaces-morphisms-lemma-base-change-separated} et
\ref{spaces-morphisms-lemma-base-change-quasi-finite}).
Il résulte alors de
Compléments sur les morphismes,
lemme \ref{more-morphisms-lemma-quasi-finite-separated-quasi-affine},
que $T' \times_T U \to T'$ est quasi-affine.

\medskip\noindent
D'après
Descente, lemme \ref{descent-lemma-descent-data-sheaves},
ceci fournit sur $T' \times_T U / T'$ une donnée de descente
relative au recouvrement étale $\{T' \to W\}$, où $W \subset T$
est l'image du morphisme $T' \to T$.
Comme $T' \times_T U$ est quasi-affine sur $T'$, il résulte de
Descente, lemme \ref{descent-lemma-quasi-affine},
que cette donnée est effective ; la dernière partie de
Descente, lemme \ref{descent-lemma-descent-data-sheaves},
implique alors que $U$ est un schéma, comme voulu.
Nous omettons quelques détails mineurs.
\end{proof}

\begin{proposition}
\label{spaces-morphisms-proposition-locally-quasi-finite-separated-over-scheme}
Soit $S$ un schéma.
Soit $f : X \to T$ un morphisme d'espaces algébriques sur $S$.
Supposons que
\begin{enumerate}
\item $T$ soit représentable,
\item $f$ soit localement quasi-fini, et
\item $f$ soit séparé.
\end{enumerate}
Alors $X$ est représentable.
\end{proposition}

\begin{proof}
Soit $T = \bigcup T_i$ un recouvrement ouvert affine du schéma $T$.
Si l'on montre que les sous-espaces ouverts $X_i = f^{-1}(T_i)$ sont
représentables, alors $X$ est représentable ; voir
Propriétés des espaces, lemme \ref{spaces-properties-lemma-subscheme}.
Notons que $X_i = T_i \times_T X$ et que les propriétés d'être localement quasi-fini
et d'être séparé sont toutes deux stables par changement de base ; voir les
lemmes \ref{spaces-morphisms-lemma-base-change-separated} et
\ref{spaces-morphisms-lemma-base-change-quasi-finite}.
On peut donc supposer que $T$ est un schéma affine.

\medskip\noindent
D'après
Propriétés des espaces,
lemme \ref{spaces-properties-lemma-union-of-quasi-compact},
il existe un recouvrement de Zariski $X = \bigcup X_i$
tel que chacun des $X_i$ possède un recouvrement étale surjectif par
un schéma affine. D'après
Propriétés des espaces, lemme \ref{spaces-properties-lemma-subscheme},
il suffit encore de démontrer la proposition pour chaque $X_i$.
On peut donc supposer qu'il existe un schéma affine $U$ et un
morphisme étale surjectif $U \to X$. On se ramène ainsi à la
situation du paragraphe suivant.

\medskip\noindent
Supposons donnés
$$
U \longrightarrow X \longrightarrow T
$$
où $U$ et $T$ sont des schémas affines, $U \to X$ est étale surjectif, et
$X \to T$ est séparé et localement quasi-fini. D'après les
lemmes \ref{spaces-morphisms-lemma-etale-locally-quasi-finite} et
\ref{spaces-morphisms-lemma-composition-quasi-finite}
le morphisme $U \to T$ est localement quasi-fini.
Comme $U$ et $T$ sont affines, il est quasi-fini.
Posons $R = U \times_X U$. Alors $X = U/R$ ; voir
Espaces, lemme \ref{spaces-lemma-space-presentation}.
Puisque $X \to T$ est séparé, le
morphisme $R \to U \times_T U$ est une immersion fermée ; voir le
lemme \ref{spaces-morphisms-lemma-fibre-product-after-map}.
En particulier, $R$ est lui aussi un schéma affine.
Comme $U \to X$ est étale, les morphismes de projection
$t, s : R \to U$ sont eux aussi étales. En particulier, $s$ et $t$ sont
quasi-finis, plats et de présentation finie (voir
Morphismes, lemmes \ref{morphisms-lemma-etale-locally-quasi-finite},
\ref{morphisms-lemma-etale-flat} et
\ref{morphisms-lemma-etale-locally-finite-presentation}).

\medskip\noindent
Soit $(U, R, s, t, c)$ le groupoïde associé à la relation
d'équivalence étale $R$ sur $U$. Soit $u \in U$ un point, et
notons $p \in T$ son image. Nous allons appliquer
Compléments sur les groupoïdes,
lemme \ref{more-groupoids-lemma-quasi-finite-over-base-j-proper},
au groupoïde $(U, R, s, t, c)$ sur le schéma $T$, avec les
points $p$ et $u$ ci-dessus.
D'après le paragraphe précédent, toutes les
hypothèses (1) -- (7) de ce lemme sont satisfaites.
On obtient donc un voisinage étale
$(T', p') \to (T, p)$ et des décompositions en unions disjointes
$$
U_{T'} = U' \amalg W, \quad
R_{T'} = R' \amalg W'
$$
et un point $u' \in U'$ satisfaisant les conclusions
(a), (b), (c), (d), (e), (f), (g) et (h) du lemme précité des
Compléments sur les groupoïdes,
lemme \ref{more-groupoids-lemma-quasi-finite-over-base-j-proper}.
On peut supposer, et l'on suppose, que $T'$ est affine (quitte à rétrécir $T'$).
La conclusion (h) entraîne que $R' = U' \times_{X_{T'}} U'$, les morphismes de
projection s'identifiant aux restrictions de $s'$ et $t'$.
Ainsi, le quintuplet $(U', R', s'|_{R'}, t'|_{R'}, c'|_{R' \times_{t', U', s'} R'})$ de la
conclusion (g) est une relation d'équivalence étale. D'après
Espaces, lemme \ref{spaces-lemma-finding-opens},
on conclut que $U'/R'$ est un sous-espace ouvert de $X_{T'}$. Par la conclusion (d),
les schémas $U'$ et $R'$ sont affines, et les morphismes
$s'|_{R'}, t'|_{R'}$ sont finis étales. La
proposition \ref{groupoids-proposition-finite-flat-equivalence} du chapitre Groupoïdes
s'applique donc et montre que $U'/R'$ est un schéma affine.

\medskip\noindent
On conclut que, pour toute paire de points $(u, p)$ comme ci-dessus, on peut
trouver un voisinage étale $(T', p') \to (T, p)$ tel que
$\kappa(p) = \kappa(p')$, ainsi qu'un point $u' \in U_{T'}$ au-dessus de $u$,
de telle sorte que l'image $x'$ de $u'$ dans $|X_{T'}|$ possède un voisinage ouvert
$V'$ dans $X_{T'}$ qui soit un schéma affine. En appliquant le
lemme \ref{spaces-morphisms-lemma-neighbourhood-scheme},
on obtient un sous-espace ouvert $W \subset X$ qui est un schéma et
qui contient $x$ (l'image de $u$ dans $|X|$).
Comme cela vaut pour tout $x$, on voit que $X$
est un schéma d'après
Propriétés des espaces, lemme \ref{spaces-properties-lemma-subscheme}.
La démonstration est achevée.
\end{proof}



\section{Applications}
\label{spaces-morphisms-section-applications}

\noindent
Une autre démonstration du lemme suivant consiste à le voir comme une conséquence
du théorème principal de Zariski pour les morphismes (non représentables) d'espaces algébriques,
tel qu'il est étudié dans Compléments sur les morphismes d'espaces, section
\ref{spaces-more-morphisms-section-structure-quasi-finite}.
En effet, Compléments sur les morphismes d'espaces, lemme
\ref{spaces-more-morphisms-lemma-quasi-finite-separated-quasi-affine}
entraîne qu'un morphisme quasi-fini et séparé d'espaces algébriques
est quasi-affine, donc représentable.

\begin{lemma}
\label{spaces-morphisms-lemma-locally-quasi-finite-separated-representable}
Soit $S$ un schéma. Soit $f : X \to Y$ un morphisme d'espaces
algébriques sur $S$. Si $f$ est localement quasi-fini et séparé, alors
$f$ est représentable.
\end{lemma}

\begin{proof}
Cela résulte immédiatement de la
proposition \ref{spaces-morphisms-proposition-locally-quasi-finite-separated-over-scheme}
et du fait que les propriétés d'être localement quasi-fini et séparé sont
préservées par tout changement de base ; voir les
lemmes \ref{spaces-morphisms-lemma-base-change-quasi-finite} et
\ref{spaces-morphisms-lemma-base-change-separated}.
\end{proof}

\begin{lemma}
\label{spaces-morphisms-lemma-etale-universally-injective-open}
\begin{slogan}
Les applications étales et universellement injectives sont des immersions ouvertes.
\end{slogan}
Soit $S$ un schéma. Soit $f : X \to Y$ un morphisme étale et universellement
injectif d'espaces algébriques sur $S$. Alors $f$ est une
immersion ouverte.
\end{lemma}

\begin{proof}
Soit $T \to Y$ un morphisme d'un schéma vers $Y$.
Il suffit de montrer que $X \times_Y T \to T$ est une immersion ouverte.
Comme les propriétés d'être étale et universellement injectif sont
stables par changement de base (voir les
lemmes \ref{spaces-morphisms-lemma-base-change-etale} et
\ref{spaces-morphisms-lemma-base-change-universally-injective})
on peut supposer que $Y$ est un schéma. Notons que la
diagonale $\Delta_{X/Y} : X \to X \times_Y X$ est étale, est un monomorphisme et
est surjective d'après le
lemme \ref{spaces-morphisms-lemma-universally-injective}.
On voit donc que $\Delta_{X/Y}$ est un isomorphisme (voir
Espaces, lemme
\ref{spaces-lemma-surjective-flat-locally-finite-presentation}),
et, en particulier, que $X$ est séparé sur $Y$.
Il s'ensuit que $X$ est lui aussi un schéma, d'après la
proposition \ref{spaces-morphisms-proposition-locally-quasi-finite-separated-over-scheme}.
Enfin, $X \to Y$ est une immersion ouverte d'après le théorème fondamental
sur les morphismes étales de schémas ; voir
Morphismes étales, théorème \ref{etale-theorem-etale-radicial-open}.
\end{proof}



\section{Théorème principal de Zariski (cas représentable)}
\label{spaces-morphisms-section-Zariski}

\noindent
Voici la version que l'on peut démontrer en utilisant le fait que la normalisation
commute à la localisation étale. Avant de pouvoir démontrer des versions
plus puissantes (pour les morphismes non représentables), il nous faut
développer davantage d'outils. Voir
Compléments sur les morphismes d'espaces, section
\ref{spaces-more-morphisms-section-structure-quasi-finite}.

\begin{lemma}
\label{spaces-morphisms-lemma-finite-type-separated}
Soit $S$ un schéma. Soit $f : X \to Y$ un morphisme d'espaces
algébriques sur $S$, représentable, de type fini et séparé.
Soit $Y'$ la normalisation de $Y$ dans $X$. On a le diagramme :
$$
\xymatrix{
X \ar[rd]_f \ar[rr]_{f'} & & Y' \ar[ld]^\nu \\
& Y &
}
$$
Alors il existe un sous-espace ouvert $U' \subset Y'$ tel que
\begin{enumerate}
\item $(f')^{-1}(U') \to U'$ soit un isomorphisme, et
\item $(f')^{-1}(U') \subset X$ soit l'ensemble des points où
$f$ est quasi-fini.
\end{enumerate}
\end{lemma}

\begin{proof}
Soit $W \to Y$ un morphisme étale surjectif, où $W$ est un schéma.
Alors $W \times_Y X$ est lui aussi un schéma. D'après le
lemme \ref{spaces-morphisms-lemma-properties-normalization},
l'espace algébrique $W \times_Y Y'$ est représentable et constitue
la normalisation du schéma $W$ dans le schéma $W \times_Y X$. On a le diagramme
$$
\xymatrix{
W \times_Y X \ar[rd]_{(1, f)} \ar[rr]_{(1, f')} & &
W \times_Y Y' \ar[ld]^{(1, \nu)} \\
& W &
}
$$
D'après Compléments sur les morphismes, lemme \ref{more-morphisms-lemma-finite-type-separated},
l'énoncé du lemme vaut sur $W$. Soit $V' \subset W \times_Y Y'$
le sous-schéma ouvert tel que
\begin{enumerate}
\item $(1, f')^{-1}(V') \to V'$ soit un isomorphisme, et
\item $(1, f')^{-1}(V') \subset W \times_Y X$ soit l'ensemble des points où
$(1, f)$ est quasi-fini.
\end{enumerate}
D'après le lemme \ref{spaces-morphisms-lemma-locally-finite-type-quasi-finite-part},
il existe un plus grand ouvert $U \subset X$ constitué de points où $f$
est quasi-fini, et $W \times_Y U = (1, f')^{-1}(V')$.
Le morphisme $f'|_U : U \to Y'$ est une immersion ouverte
d'après le lemme \ref{spaces-morphisms-lemma-closed-immersion-local},
car son changement de base à $W$ est l'isomorphisme $(1, f')^{-1}(V') \to V'$
suivi de l'immersion ouverte $V' \to W \times_Y Y'$.
Poser $U' = \Im(U \to Y')$ achève la démonstration
(nous omettons la vérification de $(f')^{-1}(U') = U$).
\end{proof}

\noindent
Dans le lemme suivant, on peut supprimer l'hypothèse de
représentabilité, puisque nous avons montré qu'un morphisme
localement quasi-fini et séparé est représentable.

\begin{lemma}
\label{spaces-morphisms-lemma-quasi-finite-separated-quasi-affine}
Soit $S$ un schéma.
Soit $f : X \to Y$ un morphisme d'espaces algébriques sur $S$.
Supposons $f$ quasi-fini et séparé.
Soit $Y'$ la normalisation de $Y$ dans $X$.
On a le diagramme :
$$
\xymatrix{
X \ar[rd]_f \ar[rr]_{f'} & & Y' \ar[ld]^\nu \\
& Y &
}
$$
Alors $f'$ est une immersion ouverte quasi-compacte et $\nu$ est entier.
En particulier, $f$ est quasi-affine.
\end{lemma}

\begin{proof}
D'après le lemme \ref{spaces-morphisms-lemma-locally-quasi-finite-separated-representable},
le morphisme $f$ est représentable. On peut donc appliquer le
lemme \ref{spaces-morphisms-lemma-finite-type-separated}. Il existe ainsi un sous-espace
ouvert $U' \subset Y'$ tel que
$(f')^{-1}(U') = X$ (!) et que $X \to U'$ soit un isomorphisme ! Autrement
dit, $f'$ est une immersion ouverte. Notons que $f'$ est quasi-compact, car
$f$ est quasi-compact et $\nu : Y' \to Y$ est séparé
(lemme \ref{spaces-morphisms-lemma-quasi-compact-permanence}).
Ainsi, pour tout schéma affine $Z$ et tout morphisme $Z \to Y$, le
produit fibré $Z \times_Y X$ est un sous-schéma ouvert quasi-compact
du schéma affine $Z \times_Y Y'$. Par définition, $f$ est donc
quasi-affine.
\end{proof}








\section{Homéomorphismes universels}
\label{spaces-morphisms-section-universal-homeomorphisms}

\noindent
La classe des homéomorphismes universels de schémas est stable par
composition et par tout changement de base, et elle est locale pour la topologie fppf sur le but. Voir
Morphismes, lemmes
\ref{morphisms-lemma-composition-universal-homeomorphism} et
\ref{morphisms-lemma-base-change-universal-homeomorphism}, ainsi que
Descente, lemme \ref{descent-lemma-descending-property-universal-homeomorphism}.
Ainsi, en appliquant à cette notion la discussion de la
section \ref{spaces-morphisms-section-representable},
on sait ce que signifie, pour un morphisme représentable
d'espaces algébriques, être un homéomorphisme universel.

\begin{lemma}
\label{spaces-morphisms-lemma-universal-homeomorphism-representable}
Soit $S$ un schéma.
Soit $f : X \to Y$ un morphisme représentable d'espaces algébriques sur $S$.
Alors $f$ est un homéomorphisme universel
(au sens de la section \ref{spaces-morphisms-section-representable}) si et seulement si,
pour tout morphisme d'espaces algébriques $Z \to Y$, l'application obtenue par
changement de base $Z \times_Y X \to Z$ induit un homéomorphisme
$|Z \times_Y X| \to |Z|$.
\end{lemma}

\begin{proof}
Si, pour tout morphisme d'espaces algébriques $Z \to Y$, l'application obtenue par
changement de base $Z \times_Y X \to Z$ induit un homéomorphisme
$|Z \times_Y X| \to |Z|$, il en va de même lorsque $Z$ est un schéma,
ce qui entraîne formellement que $f$ est un homéomorphisme universel au
sens de la
section \ref{spaces-morphisms-section-representable}.
Réciproquement, si $f$ est un homéomorphisme universel au sens de la
section \ref{spaces-morphisms-section-representable},
alors $X \to Y$ est entier, universellement injectif et surjectif
(d'après Espaces, lemme
\ref{spaces-lemma-representable-transformations-property-implication}
et
Morphismes, lemme \ref{morphisms-lemma-universal-homeomorphism}).
Ainsi, $f$ est universellement fermé ; voir le
lemme \ref{spaces-morphisms-lemma-integral-universally-closed} ;
il est aussi universellement injectif et (universellement) surjectif, c'est-à-dire que
$f$ est un homéomorphisme universel.
\end{proof}

\begin{definition}
\label{spaces-morphisms-definition-universal-homeomorphism}
Soit $S$ un schéma.
Un morphisme $f : X \to Y$ d'espaces algébriques sur $S$
est appelé un {\it homéomorphisme universel}
si et seulement si, pour tout morphisme d'espaces algébriques $Z \to Y$,
le changement de base $Z \times_Y X \to Z$ induit un homéomorphisme
$|Z \times_Y X| \to |Z|$.
\end{definition}

\noindent
Cette définition est compatible avec la définition antérieure pour les
morphismes représentables d'espaces algébriques, d'après le
lemme \ref{spaces-morphisms-lemma-universal-homeomorphism-representable}.
Pour les morphismes d'espaces algébriques, les homéomorphismes
universels ne sont pas toujours entiers.

\begin{example}
\label{spaces-morphisms-example-universal-homeomorphism}
Poursuivons la
remarque \ref{spaces-morphisms-remark-universally-injective-not-separated}.
Considérons l'espace algébrique
$X = \mathbf{A}^1_k/\{x \sim -x \mid x \not = 0\}$.
Il existe des morphismes
$$
\mathbf{A}^1_k \longrightarrow X \longrightarrow \mathbf{A}^1_k
$$
tels que la première flèche soit étale surjective, la seconde
universellement injective, et que leur composée soit l'application $x \mapsto x^2$.
La composée est donc universellement fermée. Il s'ensuit que
l'application $X \to \mathbf{A}^1_k$ est un homéomorphisme universel, mais que
$X \to \mathbf{A}^1_k$ n'est pas séparé.
\end{example}

\noindent
Soit $S$ un schéma.
Soit $f : X \to Y$ un homéomorphisme universel d'espaces algébriques
sur $S$. Alors $f$ est universellement fermé, donc quasi-compact ; voir le
lemme \ref{spaces-morphisms-lemma-universally-closed-quasi-compact}.
Cependant, $f$ n'est pas nécessairement séparé (voir l'exemple ci-dessus), ni même
quasi-séparé : on peut prendre pour exemple l'espace affine de dimension infinie
$\mathbf{A}^\infty = \Spec(k[x_1, x_2, \ldots])$ modulo
la relation d'équivalence donnée par le changement de signe d'un nombre fini de
coordonnées non nulles (nous omettons les détails).

\medskip\noindent
Énonçons d'abord les lemmes usuels.

\begin{lemma}
\label{spaces-morphisms-lemma-base-change-universal-homeomorphism}
Le changement de base d'un homéomorphisme universel d'espaces algébriques
par un morphisme quelconque d'espaces algébriques est un homéomorphisme universel.
\end{lemma}

\begin{proof}
Cela résulte immédiatement de la définition.
\end{proof}

\begin{lemma}
\label{spaces-morphisms-lemma-composition-universal-homeomorphism}
La composée de deux homéomorphismes universels
d'espaces algébriques est un homéomorphisme universel.
\end{lemma}

\begin{proof}
Démonstration omise.
\end{proof}

\begin{lemma}
\label{spaces-morphisms-lemma-reduction-universal-homeomorphism}
Soit $S$ un schéma. Soit $X$ un espace algébrique sur $S$.
L'immersion fermée canonique $X_{red} \to X$ (voir
Propriétés des espaces, définition
\ref{spaces-properties-definition-reduced-induced-space})
est un homéomorphisme universel.
\end{lemma}

\begin{proof}
Démonstration omise.
\end{proof}

\noindent
Nous plaçons ici le résultat suivant, faute pour le moment d'un meilleur
emplacement.

\begin{lemma}
\label{spaces-morphisms-lemma-integral-universally-injective-push-pull}
Soit $S$ un schéma. Soit $f : Y \to X$ un morphisme universellement injectif
et entier d'espaces algébriques sur $S$.
\begin{enumerate}
\item Le foncteur
$$
f_{small, *} : \Sh(Y_\etale) \longrightarrow \Sh(X_\etale)
$$
est pleinement fidèle, et son image essentielle est formée des faisceaux d'ensembles
$\mathcal{F}$ sur $X_\etale$ dont la restriction à $|X| \setminus f(|Y|)$
est isomorphe à $*$ ;
\item le foncteur
$$
f_{small, *} : \textit{Ab}(Y_\etale) \longrightarrow \textit{Ab}(X_\etale)
$$
est pleinement fidèle, et son image essentielle est formée des faisceaux abéliens sur
$X_\etale$ dont le support est contenu dans $f(|Y|)$.
\end{enumerate}
Dans les deux cas, $f_{small}^{-1}$ est un inverse à gauche du foncteur $f_{small, *}$.
\end{lemma}

\begin{proof}
Comme $f$ est entier, il est universellement fermé
(lemme \ref{spaces-morphisms-lemma-integral-universally-closed}).
En particulier, $f(|Y|)$ est une partie fermée de $|X|$,
et les énoncés ont un sens.
La suite de la démonstration est identique à celle du
lemme \ref{spaces-morphisms-lemma-closed-immersion-push-pull},
si ce n'est que l'on utilise
Cohomologie étale, proposition
\ref{etale-cohomology-proposition-integral-universally-injective-pushforward}
au lieu de
Cohomologie étale, proposition
\ref{etale-cohomology-proposition-closed-immersion-pushforward}.
\end{proof}



























% OK, start here.
%

\chapter{Espaces algébriques décents}



\phantomsection
\label{decent-spaces-section-phantom}


\section{Introduction}
\label{decent-spaces-section-introduction}

\noindent
Dans ce chapitre, nous étudions les propriétés « locales » des espaces
algébriques généraux, c'est-à-dire de ceux qui ne sont pas quasi-séparés.
Les espaces algébriques quasi-séparés sont étudiés dans \cite{Kn}.
Il apparaît que des phénomènes essentiellement nouveaux se produisent,
notamment en ce qui concerne les points et leurs spécialisations, pour les
espaces algébriques plus généraux. D'autre part, dans la plupart des résultats
fondamentaux sur les espaces algébriques, il n'est pas nécessaire de se
préoccuper de ces phénomènes. C'est pourquoi nous avons choisi de réunir cette
matière dans un chapitre séparé, après l'exposé habituel de la théorie.



\section{Conventions}
\label{decent-spaces-section-conventions}

\noindent
Nous supposons une fois pour toutes que tous les schémas sont contenus dans
un grand site fppf $\Sch_{fppf}$. De plus, tout anneau $A$ considéré est tel
que $\Spec(A)$ est (isomorphe à) un objet de ce grand site.

\medskip\noindent
Soit $S$ un schéma et soit $X$ un espace algébrique sur $S$.
Dans ce chapitre et le suivant, nous noterons $X \times_S X$ le produit de $X$
par lui-même (dans la catégorie des espaces algébriques sur $S$), au lieu de
$X \times X$.



\section{Fibres universellement bornées}
\label{decent-spaces-section-universally-bounded}

\noindent
Nous examinons brièvement ce que signifie, pour un morphisme d'un schéma vers
un espace algébrique, le fait d'avoir des fibres universellement bornées. Voir
Morphismes, section \ref{morphisms-section-universally-bounded}
pour les définitions et résultats analogues concernant les morphismes de
schémas.

\begin{definition}
\label{decent-spaces-definition-universally-bounded}
Soit $S$ un schéma. Soit $X$ un espace algébrique sur $S$, et soit
$U$ un schéma sur $S$. Soit $f : U \to X$ un morphisme sur $S$.
Nous dirons que {\it les fibres de $f$ sont universellement
bornées}\footnote{Cette terminologie n'est probablement pas standard.}
s'il existe un entier $n$ tel que, pour tout corps $k$ et tout morphisme
$\Spec(k) \to X$, le produit fibré $\Spec(k) \times_X U$ soit un schéma fini
sur $k$, dont le degré sur $k$ est $\leq n$.
\end{definition}

\noindent
Cette définition a un sens, car le produit fibré
$\Spec(k) \times_X U$ est un schéma. De plus, si $X$ est un schéma,
on retrouve la notion de
Morphismes, définition \ref{morphisms-definition-universally-bounded}
en vertu de
Morphismes, lemme \ref{morphisms-lemma-characterize-universally-bounded}.

\begin{lemma}
\label{decent-spaces-lemma-composition-universally-bounded}
Soit $S$ un schéma. Soit $X$ un espace algébrique sur $S$.
Soit $V \to U$ un morphisme de schémas sur $S$, et soit
$U \to X$ un morphisme de $U$ vers $X$. Si les fibres de
$V \to U$ et de $U \to X$ sont universellement bornées, alors celles de
$V \to X$ le sont aussi.
\end{lemma}

\begin{proof}
Soit $n$ un entier qui convient à $V \to U$, et soit $m$ un entier qui
convient à $U \to X$ dans la
Définition \ref{decent-spaces-definition-universally-bounded}.
Soit $\Spec(k) \to X$ un morphisme, où $k$ est un corps.
Considérons les morphismes
$$
\Spec(k) \times_X V
\longrightarrow
\Spec(k) \times_X U
\longrightarrow
\Spec(k).
$$
Par hypothèse, le schéma $\Spec(k) \times_X U$ est fini de degré au plus $m$
sur $k$, et $n$ borne le degré des fibres du premier morphisme. Par
Morphismes, lemme \ref{morphisms-lemma-composition-universally-bounded},
nous en déduisons que $\Spec(k) \times_X V$ est fini sur $k$ de degré au plus
$nm$.
\end{proof}

\begin{lemma}
\label{decent-spaces-lemma-base-change-universally-bounded}
Soit $S$ un schéma.
Soit $Y \to X$ un morphisme représentable d'espaces algébriques sur $S$.
Soit $U \to X$ un morphisme d'un schéma vers $X$.
Si les fibres de $U \to X$ sont universellement bornées, alors celles de
$U \times_X Y \to Y$ le sont aussi.
\end{lemma}

\begin{proof}
Cela résulte immédiatement de la définition et des propriétés des produits
fibrés. (Notons que $U \times_X Y$ est un schéma puisque nous avons supposé
$Y \to X$ représentable, si bien que la définition s'applique.)
\end{proof}

\begin{lemma}
\label{decent-spaces-lemma-descent-universally-bounded}
Soit $S$ un schéma. Soit $g : Y \to X$ un morphisme représentable d'espaces
algébriques sur $S$. Soit $f : U \to X$ un morphisme d'un schéma vers $X$.
Soit $f' : U \times_X Y \to Y$ le changement de base de $f$.
Si
$$
\Im(|f| : |U| \to |X|) \subset \Im(|g| : |Y| \to |X|)
$$
et si $f'$ a des fibres universellement bornées, alors $f$ a des fibres
universellement bornées.
\end{lemma}

\begin{proof}
Soit $n \geq 0$ un entier qui borne les degrés des produits fibrés
$\Spec(k) \times_Y (U \times_X Y)$ au sens de la
Définition \ref{decent-spaces-definition-universally-bounded} appliquée au morphisme $f'$.
Nous affirmons que $n$ convient aussi à $f$. En effet, supposons que
$x : \Spec(k) \to X$ soit un morphisme depuis le spectre d'un corps.
Alors, soit $\Spec(k) \times_X U$ est vide (et il n'y a rien à démontrer),
soit $x$ appartient à l'image de $|f|$. D'après
Propriétés des espaces,
lemme \ref{spaces-properties-lemma-points-cartesian}
et l'hypothèse du lemme, cela signifie qu'il existe une extension de corps
$k'/k$ et un diagramme commutatif
$$
\xymatrix{
\Spec(k') \ar[r] \ar[d] & Y \ar[d] \\
\Spec(k) \ar[r] & X
}
$$
On obtient alors
$$
\Spec(k') \times_Y (U \times_X Y) =
\Spec(k') \times_{\Spec(k)} (\Spec(k) \times_X U)
$$
Comme le schéma $\Spec(k') \times_Y (U \times_X Y)$ est supposé fini de degré
$\leq n$ sur $k'$, il en résulte que $\Spec(k) \times_X U$ est lui aussi fini
de degré $\leq n$ sur $k$, comme voulu. (Nous omettons quelques détails.)
\end{proof}

\begin{lemma}
\label{decent-spaces-lemma-universally-bounded-permanence}
Soit $S$ un schéma. Soit $X$ un espace algébrique sur $S$.
Considérons un diagramme commutatif
$$
\xymatrix{
U \ar[rd]_g \ar[rr]_f & & V \ar[ld]^h \\
& X &
}
$$
où $U$ et $V$ sont des schémas. Si $g$ a des fibres universellement bornées
et si $f$ est surjectif et plat, alors $h$ a lui aussi des fibres
universellement bornées.
\end{lemma}

\begin{proof}
Supposons que $g$ ait des fibres universellement bornées et que $f$ soit
surjectif et plat. Soit $n \geq 0$ un entier qui borne les degrés des schémas
$\Spec(k) \times_X U$ au sens de la
Définition \ref{decent-spaces-definition-universally-bounded}.
Nous affirmons que $n$ convient aussi à $h$.
Soit $\Spec(k) \to X$ un morphisme du spectre d'un corps vers $X$.
Considérons le morphisme de schémas
$$
\Spec(k) \times_X U \longrightarrow \Spec(k) \times_X V
$$
Il est plat et surjectif. Par hypothèse, le schéma de gauche est fini de degré
$\leq n$ sur $\Spec(k)$. Il résulte de
Morphismes, lemme \ref{morphisms-lemma-universally-bounded-permanence}
que le degré du schéma de droite est lui aussi majoré par $n$, comme voulu.
\end{proof}

\begin{lemma}
\label{decent-spaces-lemma-universally-bounded-finite-fibres}
Soit $S$ un schéma.
Soit $X$ un espace algébrique sur $S$, et soit $U$ un schéma sur $S$.
Soit $\varphi : U \to X$ un morphisme sur $S$.
Si les fibres de $\varphi$ sont universellement bornées, il existe un entier
$n$ tel que chaque fibre de $|U| \to |X|$ compte au plus $n$ éléments.
\end{lemma}

\begin{proof}
L'entier $n$ de la Définition \ref{decent-spaces-definition-universally-bounded} convient.
En effet, choisissons $x \in |X|$. Représentons $x$ par un morphisme
$x : \Spec(k) \to X$. Nous obtenons alors un diagramme commutatif
$$
\xymatrix{
\Spec(k) \times_X U \ar[r] \ar[d] & U \ar[d] \\
\Spec(k) \ar[r]^x & X
}
$$
qui montre (au moyen de
Propriétés des espaces,
lemme \ref{spaces-properties-lemma-points-cartesian})
que l'image réciproque de $x$ dans $|U|$ est l'image de la flèche horizontale
supérieure. Puisque $\Spec(k) \times_X U$ est fini de degré $\leq n$ sur $k$,
il possède au plus $n$ points.
\end{proof}











\section{Conditions de finitude et points}
\label{decent-spaces-section-points-monomorphisms}

\noindent
Dans cette section, nous approfondissons la question de savoir quand des
points peuvent être représentés par des monomorphismes de spectres de corps
dans l'espace.

\begin{remark}
\label{decent-spaces-remark-recall}
Avant de démontrer le lemme suivant, rappelons quelques faits sur les
morphismes étales de schémas :
\begin{enumerate}
\item Un morphisme étale est plat; les générisations se relèvent donc le long
d'un morphisme étale
(Morphismes, lemmes \ref{morphisms-lemma-etale-flat}
et \ref{morphisms-lemma-generalizations-lift-flat}).
\item Un morphisme étale est non ramifié; un morphisme non ramifié est
localement quasi-fini; ses fibres sont donc discrètes
(Morphismes, lemmes \ref{morphisms-lemma-flat-unramified-etale},
\ref{morphisms-lemma-unramified-quasi-finite}, et
\ref{morphisms-lemma-quasi-finite-at-point-characterize}).
\item Un morphisme étale quasi-compact est quasi-fini et possède en particulier
des fibres finies
(Morphismes, lemmes \ref{morphisms-lemma-quasi-finite-locally-quasi-compact} et
\ref{morphisms-lemma-quasi-finite}).
\item Un schéma étale sur un corps $k$ est une réunion disjointe de spectres
d'extensions finies séparables de $k$
(Morphismes, lemme \ref{morphisms-lemma-etale-over-field}).
\end{enumerate}
Pour une étude générale des morphismes étales, voir
\'Morphismes étales, section \ref{etale-section-etale-morphisms}.
\end{remark}

\begin{lemma}
\label{decent-spaces-lemma-U-finite-above-x}
Soit $S$ un schéma. Soit $X$ un espace algébrique sur $S$.
Soit $x \in |X|$. Les conditions suivantes sont équivalentes :
\begin{enumerate}
\item il existe une famille de schémas $U_i$ et de morphismes étales
$\varphi_i : U_i \to X$ telle que
$\coprod \varphi_i : \coprod U_i \to X$ soit surjectif,
et telle que, pour tout $i$, la fibre de
$|U_i| \to |X|$ au-dessus de $x$ soit finie, et
\item pour tout schéma affine $U$ et tout morphisme étale $\varphi : U \to X$,
la fibre de $|U| \to |X|$ au-dessus de $x$ est finie.
\end{enumerate}
\end{lemma}

\begin{proof}
L'implication (2) $\Rightarrow$ (1) est immédiate.
Soit $\varphi_i : U_i \to X$ une famille de morphismes étales comme en (1).
Soit $\varphi : U \to X$ un morphisme étale d'un schéma affine vers $X$.
Considérons les diagrammes de produits fibrés
$$
\xymatrix{
U \times_X U_i \ar[r]_-{p_i} \ar[d]_{q_i} & U_i \ar[d]^{\varphi_i} \\
U \ar[r]^\varphi & X
}
\quad \quad
\xymatrix{
\coprod U \times_X U_i \ar[r]_-{\coprod p_i} \ar[d]_{\coprod q_i} &
\coprod U_i \ar[d]^{\coprod \varphi_i} \\
U \ar[r]^\varphi & X
}
$$
Comme $q_i$ est étale, il est ouvert (voir la Remarque \ref{decent-spaces-remark-recall}).
De plus, le morphisme $\coprod q_i$ est surjectif.
Il existe donc un nombre fini d'indices $i_1, \ldots, i_n$ et des ouverts
quasi-compacts $W_{i_j} \subset U \times_X U_{i_j}$ dont les images recouvrent
$U$.
Le morphisme $p_i$ est étale, donc localement quasi-fini (voir la remarque sur
les morphismes étales ci-dessus). Nous pouvons donc appliquer
Morphismes, lemme
\ref{morphisms-lemma-locally-quasi-finite-qc-source-universally-bounded}
pour voir que les fibres de
$p_{i_j}|_{W_{i_j}} : W_{i_j} \to U_{i_j}$ sont finies.
D'après
Propriétés des espaces,
lemme \ref{spaces-properties-lemma-points-cartesian}
et l'hypothèse sur $\varphi_i$, nous en concluons que la fibre de $\varphi$
au-dessus de $x$ est finie. Autrement dit, (2) est vérifiée.
\end{proof}

\begin{lemma}
\label{decent-spaces-lemma-R-finite-above-x}
Soit $S$ un schéma. Soit $X$ un espace algébrique sur $S$.
Soit $x \in |X|$. Les conditions suivantes sont équivalentes :
\begin{enumerate}
\item il existe un schéma $U$, un morphisme étale
$\varphi : U \to X$, et des points $u, u' \in U$ dont l'image est
$x$ tels que, si l'on pose $R = U \times_X U$, la fibre de
$$
|R| \to |U| \times_{|X|} |U|
$$
au-dessus de $(u, u')$ soit finie,
\item pour tout schéma $U$, tout morphisme étale $\varphi : U \to X$ et
tous points $u, u' \in U$ dont l'image est
$x$, si l'on pose $R = U \times_X U$, la fibre de
$$
|R| \to |U| \times_{|X|} |U|
$$
au-dessus de $(u, u')$ est finie,
\item il existe un morphisme $\Spec(k) \to X$, où $k$ est un corps,
dans la classe d'équivalence de $x$, tel que les projections
$\Spec(k) \times_X \Spec(k) \to \Spec(k)$ soient
étales et quasi-compactes, et
\item il existe un monomorphisme $\Spec(k) \to X$, où $k$ est un corps,
dans la classe d'équivalence de $x$.
\end{enumerate}
\end{lemma}

\begin{proof}
Supposons (1), c'est-à-dire soit $\varphi : U \to X$ un morphisme étale
d'un schéma vers $X$, et soient $u, u'$ des points de $U$ situés au-dessus de
$x$ tels que la fibre de $|R| \to |U| \times_{|X|} |U|$ au-dessus de $(u, u')$
soit un ensemble fini. Dans cette démonstration, nous considérons un point
$u = \Spec(\kappa(u))$ comme un schéma. Notons que $u \to U$, $u' \to U$ sont
des monomorphismes (voir
Schémas, lemme \ref{schemes-lemma-injective-points-surjective-stalks}); ainsi,
$u \times_X u' \to R = U \times_X U$ est un monomorphisme.
Dans ce langage, l'hypothèse signifie exactement que
$u \times_X u'$ est un schéma dont l'espace topologique sous-jacent possède
un nombre fini de points.
Soit $\psi : W \to X$ un morphisme étale d'un schéma vers $X$.
Soient $w, w' \in W$ des points de $W$ dont l'image est $x$.
Nous devons montrer que $w \times_X w'$ est un schéma dont l'espace topologique
sous-jacent possède un nombre fini de points.
Considérons le diagramme de produit fibré
$$
\xymatrix{
W \times_X U \ar[r]_p \ar[d]_q & U \ar[d]^\varphi \\
W \ar[r]^\psi & X
}
$$
Comme $x$ est l'image de $u$ et de $u'$, nous pouvons choisir des points
$\tilde w, \tilde w'$ de $W \times_X U$ tels que $q(\tilde w) = w$,
$q(\tilde w') = w'$, $u = p(\tilde w)$ et $u' = p(\tilde w')$; voir
Propriétés des espaces,
lemme \ref{spaces-properties-lemma-points-cartesian}.
Comme $p$ et $q$ sont étales, les extensions de corps
$\kappa(w) \subset \kappa(\tilde w) \supset \kappa(u)$ et
$\kappa(w') \subset \kappa(\tilde w') \supset \kappa(u')$ sont
finies séparables; voir la Remarque \ref{decent-spaces-remark-recall}.
Nous obtenons alors un diagramme commutatif
$$
\xymatrix{
w \times_X w' \ar[d] &
\tilde w \times_X \tilde w' \ar[l] \ar[d] \ar[r] &
u \times_X u' \ar[d] \\
w \times_X w' &
\tilde w \times_S \tilde w' \ar[l] \ar[r] &
u \times_S u'
}
$$
où les carrés sont des carrés de produits fibrés. Les morphismes
horizontaux inférieurs sont étales et quasi-compacts, car tout schéma de la
forme $\Spec(k) \times_S \Spec(k')$ est affine, et d'après nos observations
ci-dessus sur les extensions de corps.
Nous voyons donc que les flèches horizontales supérieures sont étales et
quasi-compactes, et qu'elles ont par conséquent des fibres finies.
Nous avons vu plus haut que $|u \times_X u'|$ est fini; nous en concluons que
$|w \times_X w'|$ est fini. Autrement dit, (2) est vérifiée.

\medskip\noindent
Supposons (2). Soit $U \to X$ un morphisme étale d'un schéma $U$
tel que $x$ appartienne à l'image de $|U| \to |X|$. Soit $u \in U$ un
point dont l'image est $x$. Nous avons alors vu au paragraphe précédent que
$u = \Spec(\kappa(u)) \to X$ a la propriété que
$u \times_X u$ possède un espace topologique sous-jacent fini. D'autre part,
les projections $u \times_X u \to u$ sont les composés
$$
u \times_X u \longrightarrow
u \times_X U \longrightarrow
u \times_X X = u,
$$
c'est-à-dire les composés d'un monomorphisme (le changement de base du
monomorphisme $u \to U$) et d'un morphisme étale (le changement de base du
morphisme étale $U \to X$). Ainsi, $u \times_X U$ est une réunion disjointe
de spectres de corps qui sont des extensions finies séparables de $\kappa(u)$
(voir la Remarque \ref{decent-spaces-remark-recall}). Puisque $u \times_X u$ est fini, son
image dans $u \times_X U$ est une réunion disjointe finie de spectres de corps
qui sont des extensions finies séparables de $\kappa(u)$. Par
Schémas, lemme \ref{schemes-lemma-mono-towards-spec-field},
nous en concluons que $u \times_X u$ est une réunion disjointe finie de
spectres de corps qui sont des extensions finies séparables de $\kappa(u)$.
Autrement dit, nous voyons que $u \times_X u \to u$ est quasi-compact et
étale. Cela signifie que (3) est vérifiée.

\medskip\noindent
Démontrons que (3) implique (4). Soit $\Spec(k) \to X$ un morphisme
du spectre d'un corps vers $X$, dans la classe d'équivalence de $x$,
tel que les deux projections
$t, s : R = \Spec(k) \times_X \Spec(k)  \to \Spec(k)$
soient quasi-compactes et étales.
Cela signifie en particulier que $R$ est une relation d'équivalence étale
sur $\Spec(k)$.
Par Espaces, théorème \ref{spaces-theorem-presentation},
nous savons que le faisceau quotient
$X' = \Spec(k)/R$ est un espace algébrique. Par
Groupoïdes, lemme \ref{groupoids-lemma-quotient-groupoid-restrict},
le morphisme $X' \to X$ est un monomorphisme.
Comme $s, t$ sont quasi-compacts, nous voyons que $R$ est quasi-compact;
Propriétés des espaces,
lemme \ref{spaces-properties-lemma-point-like-spaces}
s'applique donc à $X'$, et nous voyons que
$X' = \Spec(k')$ pour un certain corps $k'$. Nous obtenons ainsi une
factorisation
$$
\Spec(k) \longrightarrow
\Spec(k') \longrightarrow X
$$
qui montre que $\Spec(k') \to X$ est un monomorphisme dont l'image est
$x \in |X|$. Autrement dit, (4) est vérifiée.

\medskip\noindent
Enfin, démontrons que (4) implique (1). Soit $\Spec(k) \to X$
un monomorphisme, où $k$ est un corps, dans la classe d'équivalence de $x$.
Soit $U \to X$ un morphisme étale surjectif d'un schéma $U$ vers $X$.
Soit $u \in U$ un point au-dessus de $x$. Puisque $\Spec(k) \times_X u$
est non vide et que $\Spec(k) \times_X u \to u$ est un monomorphisme,
nous en concluons que $\Spec(k) \times_X u = u$ (voir
Schémas, lemme \ref{schemes-lemma-mono-towards-spec-field}).
Le morphisme $u \to U \to X$ se factorise donc par $\Spec(k) \to X$; voici
le diagramme
$$
\xymatrix{
u \ar[r] \ar[d] & U \ar[d] \\
\Spec(k) \ar[r] & X
}
$$
Comme la flèche verticale de droite est étale, il en résulte que
$\kappa(u)/k$ est une extension finie séparable. Par conséquent,
$$
u \times_X u = u \times_{\Spec(k)} u
$$
est un schéma fini, et le résultat découle de l'analyse de la propriété (1)
faite au premier paragraphe de cette démonstration.
\end{proof}

\begin{lemma}
\label{decent-spaces-lemma-weak-UR-finite-above-x}
Soit $S$ un schéma. Soit $X$ un espace algébrique sur $S$.
Soit $x \in |X|$.
Soit $U$ un schéma et soit $\varphi : U \to X$ un morphisme étale.
Les conditions suivantes sont équivalentes :
\begin{enumerate}
\item $x$ appartient à l'image de $|U| \to |X|$ et, si l'on pose
$R = U \times_X U$, les fibres des deux applications
$$
|U| \longrightarrow |X|
\quad\text{et}\quad
|R| \longrightarrow |X|
$$
au-dessus de $x$ sont finies,
\item il existe un monomorphisme $\Spec(k) \to X$, où $k$ est un corps,
dans la classe d'équivalence de $x$, et
le produit fibré $\Spec(k) \times_X U$ est
un schéma fini non vide sur $k$.
\end{enumerate}
\end{lemma}

\begin{proof}
Supposons (1). Cette condition implique clairement la première condition du
Lemme \ref{decent-spaces-lemma-R-finite-above-x}; nous obtenons donc un monomorphisme
$\Spec(k) \to X$ dans la classe de $x$. En prenant le produit fibré,
nous voyons que $\Spec(k) \times_X U \to \Spec(k)$ est un schéma
étale sur $\Spec(k)$ ayant un nombre fini de points, donc un schéma fini
non vide sur $k$, c'est-à-dire que (2) est vérifiée.

\medskip\noindent
Supposons (2). Par hypothèse, $x$ appartient à l'image de
$|U| \to |X|$. La finitude de la fibre de
$|U| \to |X|$ au-dessus de $x$ est claire, puisque cette fibre est égale à
$|\Spec(k) \times_X U|$ d'après
Propriétés des espaces,
lemme \ref{spaces-properties-lemma-points-cartesian}.
La finitude de la fibre de $|R| \to |X|$ au-dessus de $x$ est également
claire, puisqu'elle est égale à l'ensemble sous-jacent au schéma
$$
(\Spec(k) \times_X U) \times_{\Spec(k)} (\Spec(k) \times_X U)
$$
qui est fini sur $k$. Ainsi, (1) est vérifiée.
\end{proof}

\begin{lemma}
\label{decent-spaces-lemma-UR-finite-above-x}
Soit $S$ un schéma. Soit $X$ un espace algébrique sur $S$.
Soit $x \in |X|$. Les conditions suivantes sont équivalentes :
\begin{enumerate}
\item pour tout schéma affine $U$ et tout morphisme étale
$\varphi : U \to X$, si l'on pose $R = U \times_X U$, les fibres des deux
applications
$$
|U| \longrightarrow |X|
\quad\text{et}\quad
|R| \longrightarrow |X|
$$
au-dessus de $x$ sont finies,
\item il existe des schémas $U_i$ et des morphismes étales
$U_i \to X$ tels que $\coprod U_i \to X$ soit surjectif et que, pour tout
$i$, si l'on pose $R_i = U_i \times_X U_i$, les fibres des deux applications
$$
|U_i| \longrightarrow |X|
\quad\text{et}\quad
|R_i| \longrightarrow |X|
$$
au-dessus de $x$ soient finies,
\item il existe un monomorphisme $\Spec(k) \to X$, où $k$ est un corps,
dans la classe d'équivalence de $x$, et, pour tout schéma affine $U$ et tout
morphisme étale $U \to X$, le produit fibré $\Spec(k) \times_X U$ est
un schéma fini sur $k$,
\item il existe un monomorphisme quasi-compact $\Spec(k) \to X$,
où $k$ est un corps dans la classe d'équivalence de $x$,
\item il existe un morphisme quasi-compact $\Spec(k) \to X$,
où $k$ est un corps dans la classe d'équivalence de $x$, et
\item tout morphisme $\Spec(k) \to X$, où $k$ est un corps dans la
classe d'équivalence de $x$, est quasi-compact.
\end{enumerate}
\end{lemma}

\begin{proof}
L'équivalence de (1) et (3) résulte de l'application du
Lemme \ref{decent-spaces-lemma-weak-UR-finite-above-x}
à tout morphisme étale $U \to X$ avec $U$ affine.
Il est clair que (3) implique (2).
Supposons que $U_i \to X$ et $R_i$ soient comme en (2). Le
Lemme \ref{decent-spaces-lemma-U-finite-above-x}
nous permet de conclure que, pour tout schéma affine $U$ et tout morphisme
étale $U \to X$, la fibre de $|U| \to |X|$ au-dessus de $x$ est finie.
Écrivons cette fibre $\{u_1, \ldots, u_n\}$. Comme la condition (1) du
Lemme \ref{decent-spaces-lemma-R-finite-above-x}
s'applique à $U_i \to X$ pour un certain $i$ tel que $x$ appartienne à
l'image de $|U_i| \to |X|$, nous voyons que la fibre de
$|R = U \times_X U| \to |U| \times_{|X|} |U|$
est finie au-dessus de $(u_a, u_b)$, où $a, b \in \{1, \ldots, n\}$.
La fibre de $|R| \to |X|$ au-dessus de $x$ est donc finie.
Nous voyons ainsi que (1) est vérifiée. À ce stade, nous savons que
(1), (2) et (3) sont équivalentes.

\medskip\noindent
Si (4) est vérifiée, alors, pour tout schéma affine $U$ et tout morphisme
étale $U \to X$, le schéma $\Spec(k) \times_X U$ est d'une part
étale sur $k$ (donc une réunion disjointe de spectres d'extensions finies
séparables de $k$, d'après la
Remarque \ref{decent-spaces-remark-recall})
et d'autre part quasi-compact sur $U$ (donc quasi-compact).
Nous voyons ainsi que (3) est vérifiée.
Réciproquement, si $U_i \to X$ est comme en (2) et si $\Spec(k) \to X$
est un monomorphisme comme en (3), alors
$$
\coprod \Spec(k) \times_X U_i
\longrightarrow
\coprod U_i
$$
est quasi-compact (car, au-dessus de chaque $U_i$, nous voyons que
$\Spec(k) \times_X U_i$ est une réunion disjointe finie de spectres de corps).
Ainsi, $\Spec(k) \to X$ est quasi-compact d'après
Morphismes d'espaces, lemme \ref{spaces-morphisms-lemma-quasi-compact-local}.

\medskip\noindent
Il est immédiat que (4) implique (5). Réciproquement, soit
$\Spec(k) \to X$ un morphisme quasi-compact dans la classe d'équivalence de
$x$. Soit $U \to X$ un morphisme étale avec $U$ affine. Considérons le
produit fibré
$$
\xymatrix{
F \ar[r] \ar[d] & U \ar[d] \\
\Spec(k) \ar[r] & X
}
$$
Alors $F \to U$ est quasi-compact, donc $F$ est quasi-compact.
D'autre part, $F \to \Spec(k)$ est étale; ainsi, $F$ est une
réunion disjointe finie de spectres d'extensions finies séparables de $k$
(Remarque \ref{decent-spaces-remark-recall}). Comme l'image de $|F| \to |U|$
est la fibre de $|U| \to |X|$ au-dessus de $x$ (Propriétés des espaces, lemme
\ref{spaces-properties-lemma-points-cartesian}), nous en concluons que
la fibre de $|U| \to |X|$ au-dessus de $x$ est finie. Le schéma
$F \times_{\Spec(k)} F$ est aussi une réunion finie de spectres de corps,
car il est lui aussi quasi-compact et étale sur $\Spec(k)$.
Il existe un monomorphisme
$F \times_X F \to F \times_{\Spec(k)} F$; par conséquent, $F \times_X F$ est
une réunion disjointe finie de spectres de corps
(Schémas, lemme \ref{schemes-lemma-mono-towards-spec-field}).
L'image de $F \times_X F \to U \times_X U = R$ est donc finie.
Comme cette image est la fibre de $|R| \to |X|$ au-dessus de $x$ d'après
Propriétés des espaces, lemme \ref{spaces-properties-lemma-points-cartesian},
nous en concluons que (1) est vérifiée. À ce stade, nous savons que
(1) -- (5) sont équivalentes.

\medskip\noindent
Il est clair que (6) implique (5). Réciproquement, supposons que
$\Spec(k) \to X$ soit comme en (4) et soit $\Spec(k') \to X$ un autre
morphisme, où $k'$ est un corps dans la classe d'équivalence de $x$. Par
Propriétés des espaces, lemme
\ref{spaces-properties-lemma-equivalence-class-point-monomorphism},
nous avons une factorisation $\Spec(k') \to \Spec(k) \to X$ du morphisme
donné. C'est un composé de morphismes quasi-compacts, donc un morphisme
quasi-compact (Morphismes d'espaces,
lemme \ref{spaces-morphisms-lemma-composition-quasi-compact}), comme voulu.
\end{proof}

\begin{lemma}
\label{decent-spaces-lemma-U-universally-bounded}
Soit $S$ un schéma. Soit $X$ un espace algébrique sur $S$.
Les conditions suivantes sont équivalentes :
\begin{enumerate}
\item il existe des schémas $U_i$ et des morphismes étales
$U_i \to X$ tels que $\coprod U_i \to X$ soit surjectif et que
chaque $U_i \to X$ ait des fibres universellement bornées, et
\item pour tout schéma affine $U$ et tout morphisme étale $\varphi : U \to X$,
les fibres de $U \to X$ sont universellement bornées.
\end{enumerate}
\end{lemma}

\begin{proof}
L'implication (2) $\Rightarrow$ (1) est immédiate.
Supposons (1). Soit $(\varphi_i : U_i \to X)_{i \in I}$ une famille de
morphismes étales de schémas vers $X$, qui recouvre $X$, telle que
chaque $\varphi_i$ ait des fibres universellement bornées.
Soit $\psi : U \to X$ un morphisme étale d'un schéma affine vers $X$.
Pour tout $i$, considérons le diagramme de produit fibré
$$
\xymatrix{
U \times_X U_i \ar[r]_{p_i} \ar[d]_{q_i} & U_i \ar[d]^{\varphi_i} \\
U \ar[r]^\psi & X
}
$$
Comme $q_i$ est étale, il est ouvert (voir la Remarque \ref{decent-spaces-remark-recall}).
De plus, nous avons $U = \bigcup \Im(q_i)$, puisque la famille
$(\varphi_i)_{i \in I}$ est surjective. Comme $U$ est affine, donc
quasi-compact, nous pouvons trouver un nombre fini d'indices
$i_1, \ldots, i_n \in I$ et des ouverts quasi-compacts
$W_j \subset U \times_X U_{i_j}$ tels que
$U = \bigcup q_{i_j}(W_j)$.
Le morphisme $p_{i_j}$ est étale, donc localement quasi-fini
(voir la remarque sur les morphismes étales ci-dessus). Nous pouvons donc
appliquer Morphismes, lemme
\ref{morphisms-lemma-locally-quasi-finite-qc-source-universally-bounded}
pour voir que les fibres de $p_{i_j}|_{W_j} : W_j \to U_{i_j}$ sont
universellement bornées. Par le
Lemme \ref{decent-spaces-lemma-composition-universally-bounded},
nous voyons donc que les fibres de $W_j \to X$ sont universellement bornées.
Ainsi, $\coprod_{j = 1, \ldots, n} W_j \to X$ a lui aussi des fibres
universellement bornées. Comme $\coprod_{j = 1, \ldots, n} W_j \to X$ se
factorise par le morphisme étale surjectif
$\coprod q_{i_j}|_{W_j} : \coprod_{j = 1, \ldots, n} W_j \to U$, nous
voyons que les fibres de $U \to X$ sont universellement bornées d'après le
Lemme \ref{decent-spaces-lemma-universally-bounded-permanence}.
Autrement dit, (2) est vérifiée.
\end{proof}

\begin{lemma}
\label{decent-spaces-lemma-characterize-very-reasonable}
Soit $S$ un schéma.
Soit $X$ un espace algébrique sur $S$.
Les conditions suivantes sont équivalentes :
\begin{enumerate}
\item il existe un recouvrement de Zariski $X = \bigcup X_i$ et, pour
chaque $i$, un schéma $U_i$ ainsi qu'un morphisme étale surjectif
quasi-compact $U_i \to X_i$, et
\item il existe des schémas $U_i$ et des morphismes étales $U_i \to X$
tels que les projections $U_i \times_X U_i \to U_i$ soient quasi-compactes
et que $\coprod U_i \to X$ soit surjectif.
\end{enumerate}
\end{lemma}

\begin{proof}
Si (1) est vérifiée, alors les morphismes $U_i \to X_i \to X$ sont étales
(combiner Morphismes, lemme \ref{morphisms-lemma-composition-etale}
et
Espaces, lemmes
\ref{spaces-lemma-composition-representable-transformations-property} et
\ref{spaces-lemma-morphism-schemes-gives-representable-transformation-property}
).
De plus, puisque $U_i \times_X U_i = U_i \times_{X_i} U_i$,
les deux projections $U_i \times_X U_i \to U_i$ sont quasi-compactes.

\medskip\noindent
Si (2) est vérifiée, soit $X_i \subset X$ le sous-espace ouvert correspondant
à l'image de l'application ouverte $|U_i| \to |X|$; voir
Propriétés des espaces,
lemme \ref{spaces-properties-lemma-etale-image-open}.
Les morphismes $U_i \to X_i$ sont surjectifs.
Ainsi, $U_i \to X_i$ est étale surjectif, et les projections
$U_i \times_{X_i} U_i \to U_i$ sont quasi-compactes, car
$U_i \times_{X_i} U_i = U_i \times_X U_i$. Par conséquent, d'après
Espaces, lemme \ref{spaces-lemma-representable-morphisms-spaces-property},
les morphismes $U_i \to X_i$ sont quasi-compacts.
\end{proof}








\section{Conditions sur les espaces algébriques}
\label{decent-spaces-section-conditions}

\noindent
Dans cette section, nous étudions les relations entre diverses conditions
naturelles sur les espaces algébriques rencontrées plus haut. Veuillez lire
la section \ref{decent-spaces-section-reasonable-decent}
pour vous faire une idée du sens de ces conditions.

\begin{lemma}
\label{decent-spaces-lemma-bounded-fibres}
Soit $S$ un schéma. Soit $X$ un espace algébrique sur $S$.
Considérons les conditions suivantes sur $X$ :
\begin{itemize}
\item[] $(\alpha)$ Pour tout $x \in |X|$, les conditions équivalentes du
Lemme \ref{decent-spaces-lemma-U-finite-above-x}
sont vérifiées.
\item[] $(\beta)$ Pour tout $x \in |X|$, les conditions équivalentes du
Lemme \ref{decent-spaces-lemma-R-finite-above-x}
sont vérifiées.
\item[] $(\gamma)$ Pour tout $x \in |X|$, les conditions équivalentes du
Lemme \ref{decent-spaces-lemma-UR-finite-above-x}
sont vérifiées.
\item[] $(\delta)$ Les conditions équivalentes du
Lemme \ref{decent-spaces-lemma-U-universally-bounded}
sont vérifiées.
\item[] $(\epsilon)$ Les conditions équivalentes du
Lemme \ref{decent-spaces-lemma-characterize-very-reasonable}
sont vérifiées.
\item[] $(\zeta)$ L'espace $X$ est localement quasi-séparé pour la topologie
de Zariski.
\item[] $(\eta)$ L'espace $X$ est quasi-séparé.
\item[] $(\theta)$ L'espace $X$ est représentable, c'est-à-dire que $X$ est un
schéma.
\item[] $(\iota)$ L'espace $X$ est un schéma quasi-séparé.
\end{itemize}
Nous avons
$$
\xymatrix{
& (\theta) \ar@{=>}[rd] & & & &  \\
(\iota) \ar@{=>}[ru] \ar@{=>}[rd] & &
(\zeta) \ar@{=>}[r] &
(\epsilon) \ar@{=>}[r] &
(\delta) \ar@{=>}[r] &
(\gamma) \ar@{<=>}[r] & (\alpha) + (\beta) \\
& (\eta) \ar@{=>}[ru] & & & &
}
$$
\end{lemma}

\begin{proof}
L'implication $(\gamma) \Leftrightarrow (\alpha) + (\beta)$ est immédiate.
Les implications du losange de gauche résultent clairement des définitions.

\medskip\noindent
Supposons $(\zeta)$, c'est-à-dire que $X$ soit localement quasi-séparé pour
la topologie de Zariski. Alors $(\epsilon)$ est vérifiée d'après
Propriétés des espaces, lemme
\ref{spaces-properties-lemma-quasi-separated-quasi-compact-pieces}.

\medskip\noindent
Supposons $(\epsilon)$. D'après le
Lemme \ref{decent-spaces-lemma-characterize-very-reasonable},
il existe un recouvrement ouvert de Zariski $X = \bigcup X_i$ tel que, pour
chaque $i$, il existe un schéma $U_i$ et un morphisme étale surjectif
quasi-compact $U_i \to X_i$. Fixons un $i$ et un sous-schéma ouvert affine
$W \subset U_i$. Il suffit de montrer que $W \to X$ a des fibres
universellement bornées, puisque la famille de tous ces morphismes
$W \to X$ recouvre alors $X$.
Pour cela, considérons le diagramme
$$
\xymatrix{
W \times_X U_i \ar[r]_-p \ar[d]_q & U_i \ar[d] \\
W \ar[r] & X
}
$$
Puisque $W \to X$ se factorise par $X_i$, nous voyons que
$W \times_X U_i = W \times_{X_i} U_i$; par conséquent, $q$ est quasi-compact.
Comme $W$ est affine, il en résulte que le schéma $W \times_X U_i$
est quasi-compact. Nous pouvons donc appliquer
Morphismes, lemme
\ref{morphisms-lemma-locally-quasi-finite-qc-source-universally-bounded},
et nous concluons que $p$ a des fibres universellement bornées. Du
Lemme \ref{decent-spaces-lemma-descent-universally-bounded},
nous concluons que $W \to X$ a lui aussi des fibres universellement bornées.

\medskip\noindent
Supposons $(\delta)$. Soit $U$ un schéma affine, et soit $U \to X$ un morphisme
étale. Par hypothèse, les fibres du morphisme $U \to X$ sont universellement
bornées. Il en va donc de même des fibres des deux projections
$R = U \times_X U \to U$; voir le
Lemme \ref{decent-spaces-lemma-base-change-universally-bounded}.
Et, d'après le
Lemme \ref{decent-spaces-lemma-composition-universally-bounded},
les fibres de $R \to X$ sont elles aussi universellement bornées.
Par conséquent, pour tout $x \in X$, les fibres de $|U| \to |X|$ et de
$|R| \to |X|$ au-dessus de $x$ sont finies; voir le
Lemme \ref{decent-spaces-lemma-universally-bounded-finite-fibres}.
Autrement dit, les conditions équivalentes du
Lemme \ref{decent-spaces-lemma-UR-finite-above-x}
sont vérifiées. Cela prouve $(\delta) \Rightarrow (\gamma)$.
\end{proof}

\begin{lemma}
\label{decent-spaces-lemma-properties-local}
Soit $S$ un schéma.
Soit $\mathcal{P}$ l'une des propriétés
$(\alpha)$, $(\beta)$, $(\gamma)$, $(\delta)$, $(\epsilon)$, $(\zeta)$ ou
$(\theta)$ des espaces algébriques énumérées dans le
Lemme \ref{decent-spaces-lemma-bounded-fibres}.
Alors, si $X$ est un espace algébrique sur $S$ et si
$X = \bigcup X_i$ est un recouvrement ouvert de Zariski tel que chaque
$X_i$ possède $\mathcal{P}$, alors $X$ possède $\mathcal{P}$.
\end{lemma}

\begin{proof}
Soit $X$ un espace algébrique sur $S$, et soit $X = \bigcup X_i$ un
recouvrement ouvert de Zariski tel que chaque $X_i$ possède $\mathcal{P}$.

\medskip\noindent
Le cas $\mathcal{P} = (\alpha)$. La condition $(\alpha)$ pour $X_i$
signifie que, pour tout $x \in |X_i|$, tout schéma affine $U$ et tout
morphisme étale $\varphi : U \to X_i$, la fibre de
$\varphi : |U| \to |X_i|$ au-dessus de $x$ est finie. Considérons
$x \in X$, un schéma affine $U$ et un morphisme étale $U \to X$.
Puisque $X = \bigcup X_i$ est un recouvrement ouvert de Zariski, il existe
un recouvrement ouvert affine fini $U = U_1 \cup \ldots \cup U_n$ tel que
chaque $U_j \to X$ se factorise par un certain $X_{i_j}$. Par hypothèse,
les fibres de $|U_j | \to |X_{i_j}|$ au-dessus de $x$ sont finies pour
$j = 1, \ldots, n$. Cela signifie clairement que la fibre de
$|U| \to |X|$ au-dessus de $x$ est finie.
Ceci prouve le résultat pour $(\alpha)$.

\medskip\noindent
Le cas $\mathcal{P} = (\beta)$. La condition $(\beta)$ pour $X_i$ signifie
que tout $x \in |X_i|$ est représenté par un monomorphisme du spectre
d'un corps vers $X_i$. Il en va donc de même pour $X$, puisque
$X_i \to X$ est un monomorphisme et que $X = \bigcup X_i$.

\medskip\noindent
Le cas $\mathcal{P} = (\gamma)$.
Notons que $(\gamma) = (\alpha) + (\beta)$ d'après le
Lemme \ref{decent-spaces-lemma-bounded-fibres}; le lemme pour $(\gamma)$ résulte donc
des cas traités ci-dessus.

\medskip\noindent
Le cas $\mathcal{P} = (\delta)$. La condition $(\delta)$ pour $X_i$ signifie
qu'il existe des schémas $U_{ij}$ et des morphismes étales
$U_{ij} \to X_i$ à fibres universellement bornées qui recouvrent $X_i$.
Ces schémas fournissent aussi un morphisme étale surjectif
$\coprod U_{ij} \to X$, et $U_{ij} \to X$ a encore des fibres
universellement bornées.

\medskip\noindent
Le cas $\mathcal{P} = (\epsilon)$. La condition $(\epsilon)$ pour $X_i$
signifie que l'on peut trouver un ensemble $J_i$ et des morphismes
$\varphi_{ij} : U_{ij} \to X_i$ tels que chaque $\varphi_{ij}$ soit étale,
que les deux projections $U_{ij} \times_{X_i} U_{ij} \to U_{ij}$ soient
quasi-compactes et que $\coprod_{j \in J_i} U_{ij} \to X_i$ soit surjectif.
Dans ce cas, les composés $U_{ij} \to X_i \to X$ sont étales
(combiner
Morphismes, lemmes
\ref{morphisms-lemma-composition-etale} et
\ref{morphisms-lemma-open-immersion-etale}
et
Espaces, lemmes
\ref{spaces-lemma-composition-representable-transformations-property} et
\ref{spaces-lemma-morphism-schemes-gives-representable-transformation-property}
).
Comme $X_i \subset X$ est un sous-espace, nous voyons que
$U_{ij} \times_{X_i} U_{ij} = U_{ij} \times_X U_{ij}$; la condition sur
les produits fibrés est donc préservée. Et il est clair que
$\coprod_{i, j} U_{ij} \to X$ est surjectif. Ainsi, $X$
satisfait $(\epsilon)$.

\medskip\noindent
Le cas $\mathcal{P} = (\zeta)$. La condition $(\zeta)$ pour $X_i$ signifie
que $X_i$ est localement quasi-séparé pour la topologie de Zariski.
Il est immédiat que $X$ est alors localement quasi-séparé pour la
topologie de Zariski.

\medskip\noindent
Pour $(\theta)$, voir
Propriétés des espaces,
lemme \ref{spaces-properties-lemma-subscheme}.
\end{proof}

\begin{lemma}
\label{decent-spaces-lemma-representable-properties}
Soit $S$ un schéma. Soit $\mathcal{P}$ l'une des propriétés
$(\beta)$, $(\gamma)$, $(\delta)$, $(\epsilon)$ ou
$(\theta)$ des espaces algébriques énumérées dans le
Lemme \ref{decent-spaces-lemma-bounded-fibres}.
Soient $X$, $Y$ des espaces algébriques sur $S$.
Soit $X \to Y$ un morphisme représentable.
Si $Y$ possède la propriété $\mathcal{P}$, il en est de même de $X$.
\end{lemma}

\begin{proof}
Supposons que $f : X \to Y$ soit un morphisme représentable d'espaces
algébriques et que $Y$ possède $\mathcal{P}$. Soit $x \in |X|$, et posons
$y = f(x) \in |Y|$.

\medskip\noindent
Le cas $\mathcal{P} = (\beta)$. La condition $(\beta)$ pour $Y$ signifie
qu'il existe un monomorphisme $\Spec(k) \to Y$ représentant $y$.
Le produit fibré $X_y = \Spec(k) \times_Y X$ est un schéma, et
$x$ correspond à un point de $X_y$, c'est-à-dire à un monomorphisme
$\Spec(k') \to X_y$. Comme $X_y \to X$ est lui aussi un monomorphisme,
nous voyons que $x$ est représenté par le monomorphisme
$\Spec(k') \to X_y \to X$.
Autrement dit, $(\beta)$ est vérifiée pour $X$.

\medskip\noindent
Le cas $\mathcal{P} = (\gamma)$. Puisque $(\gamma) \Rightarrow (\beta)$,
nous avons vu au paragraphe précédent que $y$ et $x$ peuvent être
représentés par des monomorphismes comme dans le diagramme suivant
$$
\xymatrix{
\Spec(k') \ar[r]_-x \ar[d] & X \ar[d] \\
\Spec(k) \ar[r]^-y & Y
}
$$
De plus, par définition de la propriété $(\gamma)$ au moyen du
Lemme \ref{decent-spaces-lemma-UR-finite-above-x} (2),
il existe des schémas $V_i$ et des morphismes étales $V_i \to Y$ tels que
$\coprod V_i \to Y$ soit surjectif et que, pour chaque $i$, en posant
$R_i = V_i \times_Y V_i$, les fibres des deux applications
$$
|V_i| \longrightarrow |Y|
\quad\text{et}\quad
|R_i| \longrightarrow |Y|
$$
au-dessus de $y$ soient finies. Cela signifie que les schémas
$(V_i)_y$ et $(R_i)_y$ sont des schémas finis sur $y = \Spec(k)$.
Comme $X \to Y$ est représentable, les produits fibrés
$U_i = V_i \times_Y X$ sont des schémas. Les morphismes $U_i \to X$
sont étales, et $\coprod U_i \to X$ est surjectif. Enfin, pour chaque $i$,
nous avons
$$
(U_i)_x =
(V_i \times_Y X)_x =
(V_i)_y \times_{\Spec(k)} \Spec(k')
$$
et
$$
(U_i \times_X U_i)_x =
\left((V_i \times_Y X) \times_X (V_i \times_Y X)\right)_x =
(R_i)_y \times_{\Spec(k)} \Spec(k')
$$
de sorte que ces schémas sont finis sur $k'$, comme changements de base des
schémas finis $(V_i)_y$ et $(R_i)_y$. Il s'ensuit que $(\gamma)$ est
vérifiée pour $X$, de nouveau par la seconde condition du
Lemme \ref{decent-spaces-lemma-UR-finite-above-x}.

\medskip\noindent
Le cas $\mathcal{P} = (\delta)$. Soit $V \to Y$ un morphisme étale, où
$V$ est un schéma affine. Puisque $Y$ possède la propriété $(\delta)$,
ce morphisme a des fibres universellement bornées. D'après le
Lemme \ref{decent-spaces-lemma-base-change-universally-bounded},
le changement de base $V \times_Y X \to X$ a lui aussi des fibres
universellement bornées. La première partie du
Lemme \ref{decent-spaces-lemma-U-universally-bounded}
s'applique donc, et nous voyons que $X$ possède lui aussi la propriété
$(\delta)$.

\medskip\noindent
Le cas $\mathcal{P} = (\epsilon)$. Nous utiliserons à plusieurs reprises
Espaces, lemme
\ref{spaces-lemma-base-change-representable-transformations-property}.
Soient $V_i \to Y$ comme dans le
Lemme \ref{decent-spaces-lemma-characterize-very-reasonable} (2).
Posons $U_i = X \times_Y V_i$. Les morphismes $U_i \to X$ sont étales,
et $\coprod U_i \to X$ est surjectif. Puisque
$U_i \times_X U_i = X \times_Y (V_i \times_Y V_i)$, nous voyons que les
projections $U_i \times_X U_i \to U_i$ sont les changements de base des
projections $V_i \times_Y V_i \to V_i$, et sont donc elles aussi
quasi-compactes. Ainsi, $X$ satisfait la condition (2) du
Lemme \ref{decent-spaces-lemma-characterize-very-reasonable}.

\medskip\noindent
Le cas $\mathcal{P} = (\theta)$. Dans ce cas, le résultat est
Catégories, lemme \ref{categories-lemma-representable-over-representable}.
\end{proof}












\section{Espaces algébriques raisonnables et décents}
\label{decent-spaces-section-reasonable-decent}

\noindent
Dans le
Lemme \ref{decent-spaces-lemma-bounded-fibres},
nous avons rencontré plusieurs conditions sur les espaces algébriques liées
au comportement des morphismes étales de schémas affines vers $X$
et à l'existence de recouvrements étales particuliers de $X$ par des
schémas. Nous récapitulons ici les différents types de conditions :
$$
\boxed{
\begin{matrix}
(\alpha) & \text{les fibres des morphismes étales issus d'affines
sont finies} \\
(\beta) & \text{les points proviennent de monomorphismes
de spectres de corps} \\
(\gamma) & \text{les points proviennent de monomorphismes quasi-compacts de
spectres de corps} \\
(\delta) & \text{les fibres des morphismes étales issus d'affines sont
universellement bornées} \\
(\epsilon) & \text{recouvrement par des morphismes étales issus de schémas,
quasi-compacts sur leur image}
\end{matrix}
}
$$

\medskip\noindent
Les conditions de la définition suivante ne sont pas exactement des conditions
sur la diagonale de $X$, mais ce sont, en un certain sens, des conditions de
séparation sur $X$.

\begin{definition}
\label{decent-spaces-definition-very-reasonable}
Soit $S$ un schéma.
Soit $X$ un espace algébrique sur $S$.
\begin{enumerate}
\item Nous disons que $X$ est {\it décent} si, pour tout point $x \in X$, les
conditions équivalentes du
Lemme \ref{decent-spaces-lemma-UR-finite-above-x}
sont vérifiées, autrement dit si la propriété $(\gamma)$ du
Lemme \ref{decent-spaces-lemma-bounded-fibres}
est vérifiée.
\item Nous disons que $X$ est {\it raisonnable} si les conditions équivalentes
du
Lemme \ref{decent-spaces-lemma-U-universally-bounded}
sont vérifiées, autrement dit si la propriété $(\delta)$ du
Lemme \ref{decent-spaces-lemma-bounded-fibres}
est vérifiée.
\item Nous disons que $X$ est {\it très raisonnable} si les conditions
équivalentes du
Lemme \ref{decent-spaces-lemma-characterize-very-reasonable}
sont vérifiées, c'est-à-dire si la propriété $(\epsilon)$ du
Lemme \ref{decent-spaces-lemma-bounded-fibres}
est vérifiée.
\end{enumerate}
\end{definition}

\noindent
Nous avons les implications suivantes entre ces conditions sur les espaces
algébriques :
$$
\xymatrix{
\text{représentable} \ar@{=>}[rd] & & & \\
 & \text{très raisonnable} \ar@{=>}[r] &
\text{raisonnable} \ar@{=>}[r] &
\text{décent} \\
\text{quasi-séparé} \ar@{=>}[ru] & & &
}
$$
La notion d'espace algébrique très raisonnable est obsolète.
Elle a été introduite parce que cette hypothèse était nécessaire pour
démontrer certains résultats qui le sont maintenant pour la classe des
espaces décents.
La classe des espaces décents est la plus grande classe d'espaces $X$ pour
lesquels on dispose d'une bonne relation entre la topologie de $|X|$ et les
propriétés de $X$ lui-même.

\begin{example}
\label{decent-spaces-example-not-decent}
L'espace algébrique $\mathbf{A}^1_{\mathbf{Q}}/\mathbf{Z}$ construit dans
Espaces, exemple \ref{spaces-example-affine-line-translation}
n'est pas décent, car son ``point générique'' ne peut pas être représenté par
un monomorphisme du spectre d'un corps.
\end{example}

\begin{remark}
\label{decent-spaces-remark-reasonable}
Les espaces algébriques raisonnables sont techniquement plus faciles à manier
que les espaces algébriques très raisonnables. Par exemple, si $X \to Y$ est
un morphisme quasi-compact, étale et surjectif d'espaces algébriques et si
$X$ est raisonnable, alors $Y$ l'est aussi; voir le
Lemme \ref{decent-spaces-lemma-descent-conditions}.
Mais nous ne savons pas si cela reste vrai pour la propriété
``très raisonnable''. Nous donnons plus bas une autre propriété technique
que possèdent les espaces algébriques raisonnables.
\end{remark}

\begin{lemma}
\label{decent-spaces-lemma-fun-property-reasonable}
Soit $S$ un schéma.
Soit $X$ un espace algébrique raisonnable quasi-compact sur $S$.
Il existe alors un système inductif filtrant d'espaces algébriques
quasi-compacts et quasi-séparés $X_i$ tel que $X = \colim_i X_i$
(la colimite étant prise dans la catégorie des faisceaux). De plus, on peut
faire en sorte que
\begin{enumerate}
\item pour tout schéma quasi-compact $T$ sur $S$, on ait
$\colim X_i(T) = X(T)$,
\item les morphismes de transition $X_i \to X_{i'}$ du système et les
coprojections $X_i \to X$ soient surjectifs et étales, et
\item si $X$ est un schéma, les espaces algébriques $X_i$ soient des schémas,
et les morphismes de transition $X_i \to X_{i'}$ et les coprojections
$X_i \to X$ soient des isomorphismes locaux.
\end{enumerate}
\end{lemma}

\begin{proof}
Esquissons la démonstration. D'après
Propriétés des espaces, lemme
\ref{spaces-properties-lemma-quasi-compact-affine-cover},
nous avons $X = U/R$, avec $U$ affine.
Dans ce cas, être raisonnable signifie que $U \to X$ est
universellement borné. Il existe donc un entier $N$ tel que les ``fibres'' de
$U \to X$ soient de degré au plus $N$; voir la
Définition \ref{decent-spaces-definition-universally-bounded}.
Désignons par $s, t : R \to U$ et $c : R \times_{s, U, t} R \to R$ les
applications structurales du groupoïde.

\medskip\noindent
Affirmation : pour tout ouvert quasi-compact $A \subset R$, il existe un ouvert
$R' \subset R$ tel que
\begin{enumerate}
\item $A \subset R'$,
\item $R'$ soit quasi-compact, et
\item $(U, R', s|_{R'}, t|_{R'}, c|_{R' \times_{s, U, t} R'})$ soit
un groupoïde en schémas.
\end{enumerate}
Notons que le morphisme $e : U \to R$ est ouvert, puisqu'il est une section du
morphisme étale $s : R \to U$; voir
\'Morphismes étales, Proposition \ref{etale-proposition-properties-sections}.
De plus, $U$ est affine, donc quasi-compact. Nous pouvons ainsi remplacer $A$
par $A \cup e(U) \subset R$ et supposer que $A$ contient $e(U)$. Définissons
ensuite par récurrence $A^1 = A$ et
$$
A^n = c(A^{n - 1} \times_{s, U, t} A) \subset R
$$
pour $n \geq 2$. Un raisonnement par récurrence montre que $A^n$ est
quasi-compact pour tout $n \geq 2$, comme image du produit fibré quasi-compact
$A^{n - 1} \times_{s, U, t} A$. Si $k$ est un corps algébriquement clos sur
$S$ et si nous considérons les points à valeurs dans $k$, alors
$$
A^n(k) = \left\{(u, u') \in U(k) \times U(k)
:
\begin{matrix}
\text{il existe } u = u_1, u_2, \ldots, u_{n + 1} = u' \in U(k)
\text{ tels que} \\
(u_i , u_{i + 1}) \in A \text{ pour tout }i = 1, \ldots, n.
\end{matrix}
\right\}
$$
Mais, comme les fibres de $U(k) \to X(k)$ sont de cardinal au plus $N$, dès
que $n > N$, un terme se répète dans la suite ci-dessus et nous pouvons la
raccourcir. Cela prouve que $A^N = A^n$ pour tout $n \geq N$.
Il en résulte que $R' = A^N$ fournit le groupoïde en schémas
$(U, R', s|_{R'}, t|_{R'}, c|_{R' \times_{s, U, t} R'})$, ce qui démontre
l'affirmation.

\medskip\noindent
Considérons le morphisme de faisceaux sur $(\Sch/S)_{fppf}$
$$
\colim_{R' \subset R} U/R' \longrightarrow U/R
$$
où $R' \subset R$ parcourt les sous-schémas ouverts quasi-compacts de $R$
qui définissent des relations d'équivalence étales comme ci-dessus.
Chacun des quotients $U/R'$ est un espace algébrique
(voir Espaces, théorème \ref{spaces-theorem-presentation}).
Comme $R'$ est quasi-compact et $U$ affine, le morphisme
$R' \to U \times_{\Spec(\mathbf{Z})} U$ est quasi-compact;
par conséquent, $U/R'$ est quasi-séparé. Enfin, si $T$ est un schéma
quasi-compact, alors
$$
\colim_{R' \subset R} U(T)/R'(T) \longrightarrow U(T)/R(T)
$$
est une bijection, car tout morphisme de $T$ dans $R$ se factorise par l'une
des sous-relations ouvertes $R'$ grâce à l'affirmation ci-dessus. Cela
implique clairement que la colimite des faisceaux $U/R'$ est $U/R$. Autrement
dit, l'espace algébrique $X = U/R$ est la colimite des espaces algébriques
quasi-séparés $U/R'$.

\medskip\noindent
Les propriétés (1) et (2) résultent de la discussion précédente.
Si $X$ est un schéma, choisir pour $U$ une réunion disjointe finie d'ouverts
affines de $X$ donne (3).
Nous omettons les détails.
\end{proof}

\begin{lemma}
\label{decent-spaces-lemma-representable-named-properties}
Soit $S$ un schéma. Soient $X$, $Y$ des espaces algébriques sur $S$.
Soit $X \to Y$ un morphisme représentable.
Si $Y$ est décent (resp.\ raisonnable), alors $X$ l'est aussi.
\end{lemma}

\begin{proof}
C'est une reformulation du Lemme \ref{decent-spaces-lemma-representable-properties}.
\end{proof}

\begin{lemma}
\label{decent-spaces-lemma-etale-named-properties}
Soit $S$ un schéma. Soit $X \to Y$ un morphisme étale d'espaces
algébriques sur $S$. Si $Y$ est décent, resp.\ raisonnable,
alors $X$ l'est aussi.
\end{lemma}

\begin{proof}
Soit $U$ un schéma affine et $U \to X$ un morphisme étale.
Posons $R = U \times_X U$ et $R' = U \times_Y U$. Notons que
$R \to R'$ est un monomorphisme.

\medskip\noindent
Soit $x \in |X|$. Pour montrer que $X$ est décent, nous devons montrer que
les fibres de $|U| \to |X|$ et de $|R| \to |X|$ au-dessus de $x$ sont finies.
Mais si $Y$ est décent, les fibres de $|U| \to |Y|$ et de
$|R'| \to |Y|$ sont finies. D'où le résultat pour ``décent''.

\medskip\noindent
Pour montrer que $X$ est raisonnable, nous devons montrer que les fibres de
$U \to X$ sont universellement bornées. Or, si $Y$ est raisonnable, les fibres
de $U \to Y$ sont universellement bornées, ce qui implique immédiatement la
même propriété pour les fibres de $U \to X$.
D'où le résultat pour ``raisonnable''.
\end{proof}







\section{Points et spécialisations}
\label{decent-spaces-section-specializations}

\noindent
Il existe un morphisme étale d'espaces algébriques $f : X \to Y$ et une
spécialisation non triviale entre des points d'une fibre de
$|f| : |X| \to |Y|$; voir
Exemples, lemme \ref{examples-lemma-specializations-fibre-etale}.
Si la source du morphisme est un schéma, on peut exclure ce phénomène en
imposant la condition ($\alpha$) à $Y$.

\begin{lemma}
\label{decent-spaces-lemma-no-specializations-map-to-same-point}
Soit $S$ un schéma.
Soit $X$ un espace algébrique sur $S$.
Soit $U \to X$ un morphisme étale d'un schéma vers $X$.
Supposons que $u, u' \in |U|$ aient pour image le même point $x$ de $|X|$ et
que $u' \leadsto u$. Si le couple $(X, x)$ satisfait les conditions
équivalentes du
Lemme \ref{decent-spaces-lemma-U-finite-above-x},
alors $u = u'$.
\end{lemma}

\begin{proof}
Supposons que le couple $(X, x)$ satisfasse les conditions équivalentes du
Lemme \ref{decent-spaces-lemma-U-finite-above-x}.
Soit $U$ un schéma, soit $U \to X$ un morphisme étale, et supposons que
$u, u' \in |U|$ aient pour image $x$ dans $|X|$ et que
$u' \leadsto u$. Nous pouvons remplacer, et remplaçons, $U$ par un voisinage
affine de $u$. Soient $t, s : R = U \times_X U \to U$
les projections étales.

\medskip\noindent
Choisissons un point $r \in R$ tel que $t(r) = u$ et $s(r) = u'$.
C'est possible d'après
Propriétés des espaces,
lemme \ref{spaces-properties-lemma-points-presentation}.
Comme les générisations se relèvent le long du morphisme étale $t$
(Remarque \ref{decent-spaces-remark-recall}), nous pouvons trouver une spécialisation
$r' \leadsto r$ telle que $t(r') = u'$. Posons $u'' = s(r')$. Alors
$u'' \leadsto u'$.
Nous pouvons donc recommencer et trouver $r'' \leadsto r'$ tel que
$t(r'') = u''$. Posons $u''' = s(r'')$, et ainsi de suite.
Voici le diagramme :
$$
\xymatrix{
& r'' \ar[rd]^s \ar[ld]_t \ar@{~>}[d] & \\
u'' \ar@{~>}[d] & r' \ar[rd]^s \ar[ld]_t \ar@{~>}[d] & u''' \ar@{~>}[d] \\
u' \ar@{~>}[d] & r \ar[rd]^s \ar[ld]_t & u'' \ar@{~>}[d] \\
u & & u'
}
$$
Dans la Remarque \ref{decent-spaces-remark-recall}, nous avons vu qu'il n'y a pas de
spécialisations entre les points des fibres du morphisme étale $s$. Ainsi,
si $u^{(n + 1)} = u^{(n)}$ pour un certain $n$, alors aussi
$r^{(n)} = r^{(n - 1)}$, puis, en appliquant $t$,
$u^{(n)} = u^{(n - 1)}$. Toute la tour est alors constante; en particulier,
$u = u'$. Nous voyons donc que, si $u \not = u'$, toutes les spécialisations
sont strictes et que $\{u, u', u'', \ldots\}$ est un ensemble infini de points
de $U$ dont l'image est le point $x$ de $|X|$. Comme nous avons choisi $U$
affine, cela contredit la seconde partie du
Lemme \ref{decent-spaces-lemma-U-finite-above-x}, comme voulu.
\end{proof}

\begin{lemma}
\label{decent-spaces-lemma-no-specializations-map-to-same-point-Noetherian}
Soit $S$ un schéma. Soit $X$ un espace algébrique sur $S$.
Soit $U \to X$ un morphisme étale d'un schéma vers $X$.
Supposons que $u, u' \in |U|$ aient pour image le même point $x$ de $|X|$ et
que $u' \leadsto u$. Si $X$ est localement noethérien, alors $u = u'$.
\end{lemma}

\begin{proof}
La discussion de Schémas, section \ref{schemes-section-points}
montre que $\mathcal{O}_{U, u'}$ est un localisé de
l'anneau local noethérien $\mathcal{O}_{U, u}$.
D'après Propriétés des espaces, lemme
\ref{spaces-properties-lemma-pre-dimension-local-ring},
nous avons $\dim(\mathcal{O}_{U, u}) = \dim(\mathcal{O}_{U, u'})$.
La théorie de la dimension des anneaux locaux noethériens permet de conclure
que $u = u'$.
\end{proof}

\begin{lemma}
\label{decent-spaces-lemma-specialization}
Soit $S$ un schéma.
Soit $X$ un espace algébrique sur $S$.
Soient $x, x' \in |X|$ et supposons que $x' \leadsto x$, c'est-à-dire que $x$
est une spécialisation de $x'$.
Supposons que le couple $(X, x')$ satisfasse les conditions équivalentes
du Lemme \ref{decent-spaces-lemma-UR-finite-above-x}.
Alors, pour tout morphisme étale $\varphi : U \to X$ d'un schéma $U$ et tout
$u \in U$ tel que $\varphi(u) = x$, il existe un point $u'\in U$ tel que
$u' \leadsto u$ et $\varphi(u') = x'$.
\end{lemma}

\begin{proof}
Nous pouvons remplacer $U$ par un voisinage ouvert affine de $u$.
Nous pouvons donc supposer que $U$ est affine. Comme $x$ appartient à l'image
du morphisme ouvert $|U| \to |X|$, il en est de même de $x'$. Nous pouvons
ainsi remplacer $X$ par le sous-espace ouvert de Zariski correspondant à
l'image de $|U| \to |X|$; voir
Propriétés des espaces,
lemme \ref{spaces-properties-lemma-etale-image-open}.
Autrement dit, nous pouvons supposer que
$U \to X$ est surjectif et étale.
Soient $s, t : R = U \times_X U \to U$ les projections.
Comme $(X, x')$ satisfait par hypothèse les conditions équivalentes du
Lemme \ref{decent-spaces-lemma-UR-finite-above-x},
les fibres de $|U| \to |X|$ et de $|R| \to |X|$
au-dessus de $x'$ sont finies. Écrivons $\{u'_1, \ldots, u'_n\} \subset U$ et
$\{r'_1, \ldots, r'_m\} \subset R$ pour les images inverses respectives
de $\{x'\}$.
Considérons les fermés
$$
T = \overline{\{u'_1\}} \cup \ldots \cup \overline{\{u'_n\}} \subset |U|,
\quad
T' = \overline{\{r'_1\}} \cup \ldots \cup \overline{\{r'_m\}} \subset |R|.
$$
Il est clair que $s(T') \subset T$. Comme $R$ est une relation d'équivalence,
nous avons aussi $t(T') = s(T')$, puisque l'ensemble $\{r_j'\}$ est invariant
par la symétrie de $R$, par construction. Soit $w \in T$ un point quelconque.
Alors $u'_i \leadsto w$ pour un certain $i$. Choisissons $r \in R$
tel que $s(r) = w$. Comme les générisations se relèvent le long de
$s : R \to U$ (voir la Remarque \ref{decent-spaces-remark-recall}), nous pouvons trouver
$r' \leadsto r$ tel que $s(r') = u_i'$. Alors $r' = r'_j$ pour un certain $j$,
et nous en déduisons que $w \in s(T')$. Ainsi,
$T = s(T') = t(T')$ est une partie saturée pour la relation d'équivalence
$|R|$, fermée dans $|U|$. Cela signifie que $T$ est l'image inverse d'une
partie fermée (!) $T'' = \varphi(T)$ de $|X|$; voir
Propriétés des espaces,
lemmes \ref{spaces-properties-lemma-points-presentation} et
\ref{spaces-properties-lemma-topology-points}.
Par conséquent, $T'' = \overline{\{x'\}}$.
Ainsi, $T$ contient un point $u_1$ d'image $x$, puisque $x \in T''$.
Autrement dit, pour un certain $i$, il existe une spécialisation
$u'_i \leadsto u_1$ dont l'image est la spécialisation donnée
$x' \leadsto x$.

\medskip\noindent
Pour achever la preuve, choisissons un point $r \in R$ tel que
$s(r) = u$ et $t(r) = u_1$ (à l'aide de
Propriétés des espaces,
lemme \ref{spaces-properties-lemma-points-cartesian}).
Comme les générisations se relèvent le long de $t$ et que
$u'_i \leadsto u_1$, nous pouvons trouver une spécialisation
$r' \leadsto r$ telle que $t(r') = u'_i$.
Posons $u' = s(r')$. Alors $u' \leadsto u$ et $\varphi(u') = x'$, comme voulu.
\end{proof}

\begin{lemma}
\label{decent-spaces-lemma-generalizations-lift-flat}
Soit $S$ un schéma. Soit $f : Y \to X$ un morphisme plat d'espaces algébriques
sur $S$. Soient $x, x' \in |X|$ et supposons que $x' \leadsto x$, c'est-à-dire
que $x$ est une spécialisation de $x'$. Supposons que le couple $(X, x')$
satisfasse les conditions équivalentes du
Lemme \ref{decent-spaces-lemma-UR-finite-above-x} (par exemple si
$X$ est décent, $X$ est quasi-séparé ou $X$ est représentable).
Alors, pour tout $y \in |Y|$ tel que $f(y) = x$, il existe un point
$y' \in |Y|$ tel que $y' \leadsto y$ et $f(y') = x'$.
\end{lemma}

\begin{proof}
(L'assertion entre parenthèses résulte de la définition des espaces décents
et des implications entre les différentes conditions de séparation
mentionnées dans la section \ref{decent-spaces-section-reasonable-decent}.)
Choisissons un schéma $V$ et un morphisme étale surjectif $V \to Y$.
Choisissons $v \in V$ d'image $y$. Il suffit alors de démontrer le lemme pour
$V \to X$. Nous pouvons donc supposer que $Y$ est un schéma.
Choisissons un schéma $U$ et un morphisme étale surjectif $U \to X$.
Choisissons $u \in U$ d'image $x$. D'après le
Lemme \ref{decent-spaces-lemma-specialization}, nous pouvons choisir $u' \leadsto u$
d'image $x'$. D'après Propriétés des espaces,
lemme \ref{spaces-properties-lemma-points-cartesian},
nous pouvons choisir $z \in U \times_X Y$ d'image $y$ et $u$.
Nous nous ramenons ainsi au cas du morphisme plat de schémas
$U \times_X Y \to U$. La conclusion résulte alors de
Morphismes, lemme \ref{morphisms-lemma-generalizations-lift-flat}.
\end{proof}






\section{Stratification des espaces algébriques par des schémas}
\label{decent-spaces-section-stratifications}

\noindent
Dans cette section, nous montrons qu'un espace algébrique quasi-compact et
quasi-séparé admet une stratification finie par des sous-espaces localement
fermés dont chacun est un schéma, et telle que le recollement des strates
s'effectue par des carrés distingués élémentaires. Nous établissons d'abord
un résultat un peu plus faible pour les espaces algébriques raisonnables.

\begin{lemma}
\label{decent-spaces-lemma-quasi-compact-reasonable-stratification}
Soit $S$ un schéma. Soit $W \to X$ un morphisme d'un schéma $W$
vers un espace algébrique $X$, plat, localement de présentation finie,
séparé, localement quasi-fini et à fibres universellement bornées. Il existe
des sous-espaces fermés réduits
$$
\emptyset = Z_{-1} \subset Z_0 \subset Z_1 \subset Z_2 \subset
\ldots \subset Z_n = X
$$
tels que, si $X_r = Z_r \setminus Z_{r - 1}$, la stratification
$X = \coprod_{r = 0, \ldots, n} X_r$ est caractérisée par la propriété
universelle suivante : étant donné $g : T \to X$, la projection
$W \times_X T \to T$ est finie localement libre de degré $r$ si et seulement si
$g(|T|) \subset |X_r|$.
\end{lemma}

\begin{proof}
Soit $n$ un entier qui majore les degrés des fibres de $W \to X$.
Choisissons un schéma $U$ et un morphisme étale surjectif $U \to X$.
Appliquons Compléments sur les morphismes, lemme
\ref{more-morphisms-lemma-stratify-flat-fp-lqf-universally-bounded}
à $W \times_X U \to U$. Nous obtenons des parties fermées
$$
\emptyset = Y_{-1} \subset Y_0 \subset Y_1 \subset Y_2 \subset
\ldots \subset Y_n = U
$$
caractérisées par la propriété de l'énoncé pour le morphisme
$W \times_X U \to U$. Manifestement, la formation de ces parties fermées
commute au changement de base. Posant $R = U \times_X U$, et notant
$s, t : R \to U$ les projections, nous concluons que
$$
s^{-1}(Y_r) = t^{-1}(Y_r)
$$
en tant que parties fermées de $R$. Autrement dit, les parties fermées
$Y_r \subset U$ sont saturées pour la relation $R$. Cela signifie que
$|Y_r|$ est l'image inverse d'une partie fermée $Z_r \subset |X|$.
Notons encore $Z_r \subset X$ le sous-espace algébrique muni de la structure
réduite induite; voir Propriétés des espaces, définition
\ref{spaces-properties-definition-reduced-induced-space}.

\medskip\noindent
Soit $g : T \to X$ un morphisme d'espaces algébriques. Choisissons un schéma
$V$ et un morphisme étale surjectif $V \to T$. Pour démontrer la dernière
assertion du lemme, il suffit de la démontrer pour la composée
$V \to X$ (d'après notre définition des morphismes finis localement libres;
voir Morphismes d'espaces, section
\ref{spaces-morphisms-section-finite-locally-free}).
De même, le morphisme de schémas $W \times_X V \to V$ est fini
localement libre de degré $r$ si et seulement si le morphisme de schémas
$$
W \times_X (U \times_X V)
\longrightarrow
U \times_X V
$$
est fini localement libre de degré $r$ (voir
Descente, lemme \ref{descent-lemma-descending-property-finite-locally-free}).
Par construction, cela se produit si et seulement si
$|U \times_X V| \to |U|$ a son image contenue dans $|Y_r|$, ce qui équivaut
à ce que $|V| \to |X|$ ait son image contenue dans $|Z_r|$.
\end{proof}

\begin{lemma}
\label{decent-spaces-lemma-stratify-flat-fp-lqf}
Soit $S$ un schéma. Soit $W \to X$ un morphisme d'un schéma $W$
vers un espace algébrique $X$, plat, localement de présentation finie,
séparé et localement quasi-fini. Il existe alors des sous-espaces ouverts
$$
X = X_0 \supset X_1 \supset X_2 \supset \ldots
$$
tels qu'un morphisme $\Spec(k) \to X$, où $k$ est un corps,
se factorise par $X_d$ si et seulement si
$W \times_X \Spec(k)$ est de degré $\geq d$ sur $k$.
\end{lemma}

\begin{proof}
Choisissons un schéma $U$ et un morphisme étale surjectif $U \to X$.
Appliquons Compléments sur les morphismes, lemme
\ref{more-morphisms-lemma-stratify-flat-fp-lqf}
à $W \times_X U \to U$. Nous obtenons des sous-schémas ouverts
$$
U = U_0 \supset U_1 \supset U_2 \supset \ldots
$$
caractérisés par la propriété de l'énoncé pour le morphisme
$W \times_X U \to U$. Manifestement, la formation de ces sous-schémas ouverts
commute au changement de base. Posons $R = U \times_X U$ et notons
$s, t : R \to U$ les projections. Nous concluons que
$$
s^{-1}(U_d) = t^{-1}(U_d)
$$
en tant que sous-schémas ouverts de $R$. Autrement dit, les sous-schémas
ouverts $U_d \subset U$ sont saturés pour la relation $R$. Cela signifie que
$U_d$ est l'image inverse d'un sous-espace ouvert $X_d \subset X$
(Propriétés des espaces, lemme
\ref{spaces-properties-lemma-subspaces-presentation}).
\end{proof}

\begin{lemma}
\label{decent-spaces-lemma-filter-quasi-compact}
Soit $S$ un schéma. Soit $X$ un espace algébrique quasi-compact sur $S$.
Il existe des sous-espaces ouverts
$$
\ldots \subset U_4 \subset U_3 \subset U_2 \subset U_1 = X
$$
ayant les propriétés suivantes :
\begin{enumerate}
\item si l'on pose $T_p = U_p \setminus U_{p + 1}$ (muni de la structure
réduite induite), il existe un schéma séparé $V_p$ et un morphisme étale
surjectif $f_p : V_p \to U_p$ tels que $f_p^{-1}(T_p) \to T_p$ soit un
isomorphisme,
\item si $x \in |X|$ peut être représenté par un morphisme quasi-compact
$\Spec(k) \to X$ dont la source est le spectre d'un corps, alors
$x \in T_p$ pour un certain $p$.
\end{enumerate}
\end{lemma}

\begin{proof}
D'après Propriétés des espaces, lemme
\ref{spaces-properties-lemma-quasi-compact-affine-cover},
nous pouvons choisir un schéma affine $U$ et un morphisme étale surjectif
$U \to X$. Pour $p \geq 0$, posons
$$
W_p = U \times_X \ldots \times_X U \setminus
\text{toutes les diagonales}
$$
où le produit fibré comporte $p$ facteurs. Puisque $U$ est séparé,
le morphisme $U \to X$ est séparé et tous les produits fibrés
$U \times_X \ldots \times_X U$ sont des schémas séparés. Comme $U \to X$
est séparé, la diagonale $U \to U \times_X U$ est une immersion fermée.
Comme $U \to X$ est étale, la diagonale $U \to U \times_X U$ est une
immersion ouverte; voir Morphismes d'espaces, lemmes
\ref{spaces-morphisms-lemma-etale-unramified} et
\ref{spaces-morphisms-lemma-diagonal-unramified-morphism}.
De même, tous les morphismes diagonaux sont des immersions ouvertes et
fermées, et $W_p$ est un sous-schéma ouvert et fermé de
$U \times_X \ldots \times_X U$. De plus, le morphisme
$$
U \times_X \ldots \times_X U \longrightarrow
U \times_{\Spec(\mathbf{Z})} \ldots \times_{\Spec(\mathbf{Z})} U
$$
est localement quasi-fini et séparé (Morphismes d'espaces,
lemme \ref{spaces-morphisms-lemma-fibre-product-after-map}),
et son but est un schéma affine. Toute partie finie de
$U \times_X \ldots \times_X U$ est donc contenue dans un ouvert affine;
voir Compléments sur les morphismes, lemme
\ref{more-morphisms-lemma-separated-locally-quasi-finite-over-affine}.
Il en va donc de même pour $W_p$.
Le groupe symétrique $S_p$ opère librement sur $W_p$ au-dessus de $X$
(puisque nous avons retranché le lieu des points fixes de
$U \times_X \ldots \times_X U$). D'après ce qui précède et
Propriétés des espaces, Proposition
\ref{spaces-properties-proposition-finite-flat-equivalence-global},
le quotient $V_p = W_p/S_p$ est un schéma. Comme l'action de $S_p$ sur
$W_p$ est au-dessus de $X$, il existe un morphisme $V_p \to X$.
Puisque $W_p \to X$ est étale et que $W_p \to V_p$ est étale surjectif,
il s'ensuit que $V_p \to X$ est lui aussi étale; voir
Propriétés des espaces, lemme \ref{spaces-properties-lemma-etale-local}.
Notons que $V_p$ est un schéma séparé d'après
Propriétés des espaces, lemme
\ref{spaces-properties-lemma-quotient-separated}.

\medskip\noindent
Soit $U_p \subset X$ le sous-espace ouvert image de $V_p \to X$.
Par construction, un morphisme $\Spec(k) \to X$, où $k$ est
algébriquement clos, se factorise par $U_p$ si et seulement si
$U \times_X \Spec(k)$ possède $\geq p$ points; comme d'habitude, notons que
$U \times_X \Spec(k)$ est, en tant que schéma, une réunion disjointe
(éventuellement infinie) de copies de $\Spec(k)$; voir
remarque \ref{decent-spaces-remark-recall}. Il s'ensuit que les $U_p$ forment une filtration
de $X$ comme dans l'énoncé du lemme. De plus, notre morphisme
$\Spec(k) \to X$ se factorise par $T_p$ si et seulement si
$U \times_X \Spec(k)$ possède exactement $p$ points. Dans ce cas, nous
voyons que $V_p \times_X \Spec(k)$ possède exactement un point.
Posons $Z_p = f_p^{-1}(T_p) \subset V_p$. C'est un sous-schéma fermé de
$V_p$. Alors $Z_p \to T_p$ est un morphisme étale d'espaces algébriques qui
induit une bijection sur les points à valeurs dans $k$ pour tout corps
algébriquement clos $k$. Précisément, cela implique que $Z_p \to T_p$ est
universellement injectif, donc une immersion ouverte d'après
Morphismes d'espaces, lemme
\ref{spaces-morphisms-lemma-etale-universally-injective-open},
et par conséquent un isomorphisme. Cela démontre (1).

\medskip\noindent
Soit $x : \Spec(k) \to X$ un morphisme quasi-compact, où $k$ est un corps.
Alors la composée $\Spec(\overline{k}) \to \Spec(k) \to X$ est elle aussi
quasi-compacte (Morphismes d'espaces, lemme
\ref{spaces-morphisms-lemma-composition-quasi-compact}).
Dans ce cas, le schéma $U \times_X \Spec(\overline{k})$ est quasi-compact.
Comme nous avons vu ci-dessus qu'il est une réunion disjointe de copies de
$\Spec(\overline{k})$, nous constatons qu'il a un nombre fini de points.
Si ce nombre est $p$, alors $x \in T_p$, comme voulu, ce qui achève la preuve.
\end{proof}

\begin{lemma}
\label{decent-spaces-lemma-filter-reasonable}
Soit $S$ un schéma. Soit $X$ un espace algébrique quasi-compact et raisonnable
sur $S$. Il existe un entier $n$ et des sous-espaces ouverts
$$
\emptyset = U_{n + 1} \subset
U_n \subset U_{n - 1} \subset \ldots \subset U_1 = X
$$
ayant la propriété suivante : si l'on pose $T_p = U_p \setminus U_{p + 1}$
(muni de la structure réduite induite), il existe un schéma séparé $V_p$ et
un morphisme étale surjectif $f_p : V_p \to U_p$ tels que
$f_p^{-1}(T_p) \to T_p$ soit un isomorphisme.
\end{lemma}

\begin{proof}
La démonstration de ce lemme est identique à celle du
Lemme \ref{decent-spaces-lemma-filter-quasi-compact}.
Soit $n$ un entier qui majore les degrés des fibres de $U \to X$;
un tel entier existe puisque $X$ est raisonnable, voir
définition \ref{decent-spaces-definition-very-reasonable}.
Nous constatons alors que $U_{n + 1} = \emptyset$, ce qui achève la preuve.
\end{proof}

\begin{lemma}
\label{decent-spaces-lemma-stratify-reasonable}
Soit $S$ un schéma. Soit $X$ un espace algébrique quasi-compact et raisonnable
sur $S$. Il existe un entier $n$ et des sous-espaces ouverts
$$
\emptyset = U_{n + 1} \subset
U_n \subset U_{n - 1} \subset \ldots \subset U_1 = X
$$
tels que chaque $T_p = U_p \setminus U_{p + 1}$ (muni de la structure réduite
induite) soit un schéma.
\end{lemma}

\begin{proof}
Conséquence immédiate du Lemme \ref{decent-spaces-lemma-filter-reasonable}.
\end{proof}

\noindent
Le résultat suivant est presque identique à
\cite[Proposition 5.7.8]{GruRay}.

\begin{lemma}
\label{decent-spaces-lemma-filter-quasi-compact-quasi-separated}
\begin{reference}
Ce résultat est presque identique à \cite[Proposition 5.7.8]{GruRay}.
\end{reference}
Soit $X$ un espace algébrique quasi-compact et quasi-séparé sur
$\Spec(\mathbf{Z})$. Il existe un entier $n$ et des sous-espaces ouverts
$$
\emptyset = U_{n + 1} \subset
U_n \subset U_{n - 1} \subset \ldots \subset U_1 = X
$$
ayant la propriété suivante : si l'on pose $T_p = U_p \setminus U_{p + 1}$
(muni de la structure réduite induite), il existe un schéma $V_p$
quasi-compact et séparé et un morphisme étale surjectif
$f_p : V_p \to U_p$ tels que $f_p^{-1}(T_p) \to T_p$ soit un isomorphisme.
\end{lemma}

\begin{proof}
La démonstration de ce lemme est identique à celle du
Lemme \ref{decent-spaces-lemma-filter-quasi-compact}.
Remarquons qu'un espace quasi-séparé est raisonnable; voir
lemme \ref{decent-spaces-lemma-bounded-fibres} et
définition \ref{decent-spaces-definition-very-reasonable}.
Nous obtenons donc $U_{n + 1} = \emptyset$, comme dans le
Lemme \ref{decent-spaces-lemma-filter-reasonable}.
À la fin de l'argument, ajoutons que, puisque $X$ est quasi-séparé,
les schémas $U \times_X \ldots \times_X U$ sont tous quasi-compacts.
Les schémas $W_p$ sont donc quasi-compacts. Par conséquent, les quotients
$V_p = W_p/S_p$ par le groupe symétrique $S_p$ sont des schémas
quasi-compacts.
\end{proof}

\noindent
Le lemme suivant devrait probablement se trouver ailleurs.

\begin{lemma}
\label{decent-spaces-lemma-locally-constructible}
Soit $S$ un schéma. Soit $X$ un espace algébrique quasi-séparé sur $S$.
Soit $E \subset |X|$ une partie. Alors $E$ est étale-localement constructible
(Propriétés des espaces, définition
\ref{spaces-properties-definition-locally-constructible})
si et seulement si $E$ est une partie localement constructible de
l'espace topologique $|X|$
(Topologie, définition \ref{topology-definition-constructible}).
\end{lemma}

\begin{proof}
Supposons que $E \subset |X|$ soit une partie localement constructible de
l'espace topologique $|X|$. Soit $f : U \to X$ un morphisme étale, où $U$
est un schéma. Nous devons montrer que $f^{-1}(E)$ est localement constructible
dans $U$. La question est locale sur $U$ et $X$; nous pouvons donc supposer
que $X$ est quasi-compact, que $E \subset |X|$ est constructible et que $U$
est affine. Dans ce cas, $U \to X$ est quasi-compact, donc
$f : |U| \to |X|$ est quasi-compact. Remarquons que les ouverts rétrocompacts
de $|X|$, resp.\ $U$, ne sont autres que les ouverts quasi-compacts de
$|X|$, resp.\ $U$; voir
Topologie, lemme \ref{topology-lemma-topology-quasi-separated-scheme}.
Ainsi, $f^{-1}(E)$ est constructible d'après Topologie, lemme
\ref{topology-lemma-inverse-images-constructibles}.

\medskip\noindent
Réciproquement, supposons que $E$ soit étale-localement constructible.
Nous voulons montrer que $E$ est localement constructible dans l'espace
topologique $|X|$. La question est locale sur $X$; nous pouvons donc supposer
que $X$ est à la fois quasi-compact et quasi-séparé. Montrons que, dans ce cas,
$E$ est constructible dans $|X|$. Choisissons des sous-espaces ouverts
$$
\emptyset = U_{n + 1} \subset
U_n \subset U_{n - 1} \subset \ldots \subset U_1 = X
$$
et des morphismes étales surjectifs $f_p : V_p \to U_p$ induisant des
isomorphismes $f_p^{-1}(T_p) \to T_p = U_p \setminus U_{p + 1}$, où $V_p$
est un schéma quasi-compact et séparé, comme dans le
Lemme \ref{decent-spaces-lemma-filter-quasi-compact-quasi-separated}.
Par définition, l'image inverse $E_p \subset V_p$ de $E$ est localement
constructible dans $V_p$. Alors $E_p$ est constructible dans $V_p$ d'après
Propriétés, lemme
\ref{properties-lemma-constructible-quasi-compact-quasi-separated}.
Ainsi, $E_p \cap |f_p^{-1}(T_p)| = E \cap |T_p|$ est constructible dans
$|T_p|$ d'après
Topologie, lemme \ref{topology-lemma-intersect-constructible-with-closed}
(remarquons que $V_p \setminus f_p^{-1}(T_p)$ est quasi-compact, car c'est
l'image inverse de l'espace quasi-compact $U_{p + 1}$ par le morphisme
quasi-compact $f_p$). Par conséquent,
$$
E = (|T_n| \cap E) \cup (|T_{n - 1}| \cap E) \cup \ldots \cup
(|T_1| \cap E)
$$
est constructible d'après
Topologie, lemme \ref{topology-lemma-collate-constructible-from-constructible}.
Nous utilisons ici le fait que $|T_p|$ est constructible dans $|X|$, ce qui
résulte clairement de ce qui précède.
\end{proof}





\section{Recouvrement entier par un schéma}
\label{decent-spaces-section-integral-cover}

\noindent
Nous démontrons ici que, pour tout espace algébrique $X$ quasi-compact et
quasi-séparé, il existe un schéma $Y$ et un morphisme entier surjectif
$Y \to X$. Après avoir développé une théorie des limites d'espaces
algébriques, nous démontrerons qu'on peut obtenir un morphisme fini; voir
Limites d'espaces, section \ref{spaces-limits-section-finite-cover}.

\begin{lemma}
\label{decent-spaces-lemma-extend-integral-morphism}
Soit $S$ un schéma. Soit $j : V \to Y$ une immersion ouverte quasi-compacte
d'espaces algébriques sur $S$. Soit $\pi : Z \to V$ un morphisme entier.
Il existe alors un morphisme entier $\nu : Y' \to Y$ tel que $Z$ soit
$V$-isomorphe à l'image inverse de $V$ dans $Y'$.
\end{lemma}

\begin{proof}
Puisque $j$ et $\pi$ sont tous deux quasi-compacts et séparés, il en va de
même de $j \circ \pi$. Soit $\nu : Y' \to Y$ la normalisation de $Y$ dans
$Z$; voir Morphismes d'espaces, section
\ref{spaces-morphisms-section-normalization-X-in-Y}.
Bien entendu, $\nu$ est entier; voir
Morphismes d'espaces, lemme
\ref{spaces-morphisms-lemma-characterize-normalization}.
La dernière assertion résulte formellement de
Morphismes d'espaces, lemmes
\ref{spaces-morphisms-lemma-properties-normalization} et
\ref{spaces-morphisms-lemma-normalization-in-integral}.
\end{proof}

\begin{lemma}
\label{decent-spaces-lemma-there-is-a-scheme-integral-over}
Soit $S$ un schéma. Soit $X$ un espace algébrique quasi-compact et
quasi-séparé sur $S$.
\begin{enumerate}
\item Il existe un morphisme entier surjectif $Y \to X$, où $Y$ est un
schéma,
\item étant donné un morphisme étale surjectif $U \to X$, on peut choisir
$Y \to X$ de telle sorte que, pour tout $y \in Y$, il existe un voisinage
ouvert $V \subset Y$ tel que $V \to X$ se factorise par $U$.
\end{enumerate}
\end{lemma}

\begin{proof}
La partie (1) est le cas particulier de la partie (2) où $U = X$.
Choisissons un morphisme étale surjectif $U' \to U$, où $U'$ est un schéma.
Il est clair que nous pouvons remplacer $U$ par $U'$ et donc supposer que
$U$ est un schéma. Puisque $X$ est quasi-compact, il existe un nombre fini
d'ouverts affines $U_i \subset U$ tels que $U' = \coprod U_i \to X$ soit
surjectif. En remplaçant encore $U$ par $U'$, nous voyons que nous pouvons
supposer $U$ affine. Puisque $X$ est quasi-séparé, donc raisonnable, il existe
un entier $d$ qui majore le degré des fibres géométriques de $U \to X$
(voir Lemme \ref{decent-spaces-lemma-bounded-fibres}).
Nous démontrerons le lemme par récurrence sur $d$ pour tout schéma $U$
quasi-compact et séparé muni d'un morphisme étale surjectif vers $X$.
Si $d = 1$, alors $U = X$ et le résultat vaut avec $Y = U$.
Supposons $d > 1$.

\medskip\noindent
Nous appliquons Morphismes d'espaces, lemme
\ref{spaces-morphisms-lemma-quasi-finite-separated-quasi-affine}
et obtenons une factorisation
$$
\xymatrix{
U \ar[rr]_j \ar[rd] & & Y \ar[ld]^\pi \\
& X
}
$$
où $\pi$ est entier et $j$ une immersion ouverte quasi-compacte. Nous
pouvons supposer, et nous le faisons, que $j(U)$ est schématiquement dense
dans $Y$. Alors $U \times_X Y$ est un schéma quasi-compact et séparé
(puisqu'il est entier sur $U$), et nous avons
$$
U \times_X Y = U \amalg W
$$
Ici, le premier terme est l'image de $U \to U \times_X Y$
(qui est fermée d'après
Morphismes d'espaces, lemme \ref{spaces-morphisms-lemma-semi-diagonal},
et ouverte parce que le morphisme en question est étale entre des espaces
algébriques étales sur $Y$), et le second terme est le complémentaire
(ouvert et fermé). L'image $V \subset Y$ de $W$ est un sous-espace ouvert
contenant $Y \setminus U$.

\medskip\noindent
Le morphisme étale $W \to Y$ a des fibres géométriques de cardinalité $< d$.
En effet, c'est clair par inspection pour les points géométriques de
$U \subset Y$. Puisque $|U| \subset |Y|$ est dense, cela vaut pour tous les
points géométriques de $Y$ d'après le
Lemme \ref{decent-spaces-lemma-quasi-compact-reasonable-stratification}
(le degré des fibres d'un morphisme étale quasi-compact n'augmente pas par
spécialisation). Nous pouvons donc appliquer l'hypothèse de récurrence à
$W \to V$ et trouver un morphisme entier surjectif $Z \to V$, où $Z$ est un
schéma, qui, localement pour la topologie de Zariski, se factorise par $W$.
Choisissons une factorisation $Z \to Z' \to Y$ dans laquelle $Z' \to Y$ est
entier et $Z \to Z'$ une immersion ouverte
(Lemme \ref{decent-spaces-lemma-extend-integral-morphism}).
Après avoir remplacé $Z'$ par l'adhérence schématique de $Z$ dans $Z'$,
nous pouvons supposer que $Z$ est schématiquement dense dans $Z'$.
Après ce remplacement, nous avons $Z' \times_Y V = Z$. Enfin, soit
$T \subset Y$ le complémentaire $Y \setminus V$, muni de la structure de
sous-espace fermé induite. Considérons le morphisme
$$
Z' \amalg T \longrightarrow X
$$
Par construction, c'est un morphisme entier surjectif.
Puisque $T \subset U$, il est clair que le morphisme $T \to X$ se factorise
par $U$. Par ailleurs, soit $z \in Z'$ un point. Si $z \not \in Z$, alors $z$
s'envoie sur un point de $Y \setminus V \subset U$, et nous trouvons un
voisinage ouvert de $z$ sur lequel le morphisme se factorise par $U$.
Si $z \in Z$, nous disposons d'un voisinage ouvert de $z$ dans $Z$
(qui est aussi un voisinage ouvert de $z$ dans $Z'$) sur lequel le morphisme
se factorise par $W \subset U \times_X Y$, donc par $U$.
\end{proof}

\begin{lemma}
\label{decent-spaces-lemma-there-is-a-scheme-integral-over-refined}
Soit $S$ un schéma. Soit $X$ un espace algébrique quasi-compact et
quasi-séparé sur $S$ tel que $|X|$ ait un nombre fini de composantes
irréductibles.
\begin{enumerate}
\item Il existe un morphisme entier surjectif $f : Y \to X$, où $Y$ est un
schéma, tel que $f$ soit fini étale au-dessus d'un ouvert dense quasi-compact
$U \subset X$,
\item étant donné un morphisme étale surjectif $V \to X$, on peut choisir
$Y \to X$ de telle sorte que, pour tout $y \in Y$, il existe un voisinage
ouvert $W \subset Y$ tel que $W \to X$ se factorise par $V$.
\end{enumerate}
\end{lemma}

\begin{proof}
La démonstration est, à peu de chose près, la même que celle du
Lemme \ref{decent-spaces-lemma-there-is-a-scheme-integral-over}, avec des précisions
techniques supplémentaires pour obtenir l'ouvert dense quasi-compact $U$
(et, malheureusement, des changements de notation pour suivre $U$).

\medskip\noindent
La partie (1) est le cas particulier de la partie (2) où $V = X$.

\medskip\noindent
Démonstration de (2). Choisissons un morphisme étale surjectif $V' \to V$,
où $V'$ est un schéma. Il est clair que nous pouvons remplacer $V$ par $V'$
et donc supposer que $V$ est un schéma. Puisque $X$ est quasi-compact, il
existe un nombre fini d'ouverts affines $V_i \subset V$ tels que
$V' = \coprod V_i \to X$ soit surjectif. En remplaçant encore $V$ par $V'$,
nous voyons que nous pouvons supposer $V$ affine. Puisque $X$ est
quasi-séparé, donc raisonnable, il existe un entier $d$ qui majore le degré
des fibres géométriques de $V \to X$
(voir Lemme \ref{decent-spaces-lemma-bounded-fibres}).

\medskip\noindent
Par récurrence sur $d \geq 1$, nous allons établir l'hypothèse de récurrence
$(H_d)$ suivante :
\begin{itemize}
\item pour tout espace algébrique quasi-compact et quasi-séparé $X$ ayant
un nombre fini de composantes irréductibles, pour tout $m \geq 0$, pour tous
schémas quasi-compacts et séparés $V_j$, $j = 1, \ldots, m$, et pour tous
morphismes étales $\varphi_j : V_j \to X$, $j = 1, \ldots, m$, tels que $d$
majore le degré des fibres géométriques de $\varphi_j : V_j\to X$ et que
$\varphi = \coprod \varphi_j : V = \coprod V_j \to X$ soit surjectif,
l'énoncé du lemme vaut pour $\varphi : V \to X$.
\end{itemize}
Si $d = 1$, alors chaque $\varphi_j$ est une immersion ouverte. Ainsi, $X$
est un schéma et le résultat vaut avec $Y = V$.
Supposons $d > 1$, supposons $(H_{d - 1})$, et soient
$m$, $\varphi_j : V_j \to X$, $j = 1, \ldots, m$ comme dans $(H_d)$.

\medskip\noindent
Soient $\eta_1, \ldots, \eta_n \in |X|$ les points génériques des
composantes irréductibles de $|X|$. D'après
Propriétés des espaces, Proposition
\ref{spaces-properties-proposition-locally-quasi-separated-open-dense-scheme},
il existe un sous-schéma ouvert $U \subset X$ tel que
$\eta_1, \ldots, \eta_n \in U$. En remplaçant $U$ par un ouvert plus petit,
nous pouvons supposer $U$ affine et, d'après
Morphismes, lemme \ref{morphisms-lemma-generically-finite},
nous pouvons supposer que chaque $\varphi_j : V_j \to X$ est fini étale
au-dessus de $U$. Bien entendu, $U$ est quasi-compact et dense dans $X$,
et $\varphi_j^{-1}(U)$ est dense dans $V_j$. En particulier, chaque $V_j$
a un nombre fini de composantes irréductibles.

\medskip\noindent
Fixons $j \in \{1, \ldots, m\}$.
Comme dans Morphismes d'espaces, lemme
\ref{spaces-morphisms-lemma-quasi-finite-separated-quasi-affine},
notons $Y_j$ la normalisation de $X$ dans $V_j$. Nous obtenons une
factorisation
$$
\xymatrix{
V_j \ar[rr] \ar[rd]_{\varphi_j} & & Y_j \ar[ld]^{\pi_j} \\
& X
}
$$
où $\pi_j$ est entier et $V_j \to Y_j$ une immersion ouverte
quasi-compacte. Puisque $Y_j$ est la normalisation de $X$ dans $V_j$, les
Morphismes d'espaces, lemmes
\ref{spaces-morphisms-lemma-properties-normalization} et
\ref{spaces-morphisms-lemma-normalization-in-integral}
montrent que $\varphi_j^{-1}(U) \to \pi_j^{-1}(U)$ est un isomorphisme.
Ainsi, $\pi_j$ est fini étale au-dessus de $U$. Remarquons que $V_j$ est
schématiquement dense dans $Y_j$, puisque $Y_j$ est la normalisation de $X$
dans $V_j$ (cela résulte de la caractérisation de la normalisation relative
dans Morphismes d'espaces, lemme
\ref{spaces-morphisms-lemma-characterize-normalization}). Puisque $V_j$ est
quasi-compact, nous voyons que $|V_j| \subset |Y_j|$ est dense; voir
Morphismes d'espaces, section
\ref{spaces-morphisms-section-scheme-theoretic-closure}
(et en particulier Morphismes d'espaces, lemme
\ref{spaces-morphisms-lemma-quasi-compact-immersion}).
Il s'ensuit que $|Y_j|$ a un nombre fini de composantes irréductibles.
Alors $V_j \times_X Y_j$ est un schéma quasi-compact et séparé
(puisqu'il est fini sur $V_j$), et
$$
V_j \times_X Y_j = V_j \amalg W_j
$$
Ici, le premier terme est l'image de $V_j \to V_j \times_X Y_j$
(qui est fermée d'après
Morphismes d'espaces, lemme \ref{spaces-morphisms-lemma-semi-diagonal},
et ouverte parce que le morphisme en question est étale entre des espaces
algébriques étales sur $Y_j$), et le second terme est le complémentaire
(ouvert et fermé).

\medskip\noindent
Le morphisme étale $W_j \to Y_j$ a des fibres géométriques de cardinalité
$< d$. En effet, c'est clair par inspection pour les points géométriques de
$V_j \subset Y_j$. Puisque $|V_j| \subset |Y_j|$ est dense, cela vaut pour
tous les points géométriques de $Y_j$ d'après le
Lemme \ref{decent-spaces-lemma-quasi-compact-reasonable-stratification}
(le degré des fibres d'un morphisme étale quasi-compact n'augmente pas par
spécialisation). En appliquant $(H_{d - 1})$ à
$V_j \amalg W_j \to Y_j$, nous trouvons un morphisme entier surjectif
$Y_j' \to Y_j$, avec $Y_j'$ un schéma, qui se factorise localement pour la
topologie de Zariski par $V_j \amalg W_j$ et qui est fini étale au-dessus
d'un ouvert dense quasi-compact $U_j \subset Y_j$. En remplaçant $U$ par un
ouvert plus petit, nous pouvons supposer, et nous le faisons, que
$\pi_j^{-1}(U) \subset U_j$ (nous pouvons choisir, et nous choisissons, le
même $U$ pour tout $j$; nous omettons certains détails).

\medskip\noindent
Nous affirmons que
$$
Y = \coprod\nolimits_{j = 1, \ldots, m} Y'_j \longrightarrow X
$$
est la solution de notre problème. Premièrement, ce morphisme est entier,
car, sur chaque composante, nous avons la composition $Y'_j \to Y_j \to X$ de
morphismes entiers (Morphismes d'espaces, lemme
\ref{spaces-morphisms-lemma-composition-integral}). Deuxièmement, ce morphisme
se factorise localement pour la topologie de Zariski par
$V = \coprod V_j$, car nous avons vu plus haut que chaque
$Y'_j \to Y_j$ se factorise localement pour la topologie de Zariski par
$V_j \amalg W_j = V_j \times_X Y_j$. Enfin, puisque $Y'_j \to Y_j$ et
$Y_j \to X$ sont tous deux finis étales au-dessus de $U$, leur composition
l'est aussi. Ceci achève la démonstration.
\end{proof}



















\section{Lieu schématique}
\label{decent-spaces-section-schematic}

\noindent
Dans cette section, nous démontrons qu'un espace algébrique décent possède un
sous-espace ouvert dense qui est un schéma. Nous le démontrons d'abord pour
les espaces algébriques raisonnables.


\begin{proposition}
\label{decent-spaces-proposition-reasonable-open-dense-scheme}
Soit $S$ un schéma. Soit $X$ un espace algébrique sur $S$.
Si $X$ est raisonnable, alors il existe un sous-espace ouvert dense de $X$
qui est un schéma.
\end{proposition}

\begin{proof}
D'après Propriétés des espaces,
lemme \ref{spaces-properties-lemma-subscheme},
la question est locale sur $X$. Nous pouvons donc supposer qu'il existe un
schéma affine $U$ et un morphisme étale surjectif $U \to X$
(Propriétés des espaces, lemme
\ref{spaces-properties-lemma-cover-by-union-affines}).
Soit $n$ un entier qui majore les degrés des fibres de $U \to X$;
un tel entier existe puisque $X$ est raisonnable, voir la
Définition \ref{decent-spaces-definition-very-reasonable}.
Nous allons montrer par récurrence sur $n$ que, dès que
\begin{enumerate}
\item $U \to X$ est un morphisme étale surjectif dont les fibres sont de
degré $\leq n$, et
\item $U$ est isomorphe à un sous-schéma localement fermé d'un schéma affine,
\end{enumerate}
le lieu schématique est dense dans $X$.

\medskip\noindent
Soit $X_n \subset X$ le sous-espace ouvert complémentaire du sous-espace
fermé $Z_{n - 1} \subset X$ construit dans le
Lemme \ref{decent-spaces-lemma-quasi-compact-reasonable-stratification}
à l'aide du morphisme $U \to X$.
Soit $U_n \subset U$ l'image inverse de $X_n$. Alors
$U_n \to X_n$ est fini localement libre de degré $n$.
Par conséquent, $X_n$ est un schéma d'après
Propriétés des espaces, Proposition
\ref{spaces-properties-proposition-finite-flat-equivalence-global}
(et le fait que tout ensemble fini de points de $U_n$ est contenu dans un
ouvert affine de $U_n$; voir
Propriétés, lemme \ref{properties-lemma-ample-finite-set-in-affine}).

\medskip\noindent
Soit $X' \subset X$ le sous-espace ouvert tel que $|X'|$ soit l'intérieur de
$|Z_{n - 1}|$ dans $|X|$ (voir
Topologie, définition \ref{topology-definition-nowhere-dense}).
Soit $U' \subset U$ son image inverse. Alors $U' \to X'$ est étale surjectif
et le degré de ses fibres est majoré par $n - 1$. Par récurrence, nous voyons
que le lieu schématique de $X'$ est un ouvert dense $X'' \subset X'$.
La topologie élémentaire montre alors que $X'' \cup X_n \subset X$ est
ouvert et dense, ce qui conclut la preuve.
\end{proof}

\begin{theorem}[David Rydh]
\label{decent-spaces-theorem-decent-open-dense-scheme}
Soit $S$ un schéma. Soit $X$ un espace algébrique sur $S$.
Si $X$ est décent, alors il existe un sous-espace ouvert dense de $X$ qui est
un schéma.
\end{theorem}

\begin{proof}
Supposons que $X$ soit un espace algébrique décent pour lequel le théorème
est faux. D'après Propriétés des espaces, lemme
\ref{spaces-properties-lemma-subscheme}, il existe un plus grand sous-espace
ouvert $X' \subset X$ qui est un schéma. Puisque $X'$ n'est pas dense dans
$X$, il existe un sous-espace ouvert $X'' \subset X$ tel que
$|X''| \cap |X'| = \emptyset$. En remplaçant $X$ par $X''$, nous obtenons un
espace algébrique décent non vide $X$ qui ne contient {\it aucun}
sous-espace ouvert qui soit un schéma.

\medskip\noindent
Choisissons un schéma affine non vide $U$ et un morphisme étale $U \to X$.
Nous pouvons remplacer, et remplaçons, $X$ par le sous-espace ouvert
correspondant à l'image de $|U| \to |X|$. Considérons la suite décroissante
de sous-espaces ouverts
$$
X = X_0 \supset X_1 \supset X_2 \ldots
$$
construite dans le Lemme \ref{decent-spaces-lemma-stratify-flat-fp-lqf}
pour le morphisme $U \to X$. Remarquons que $X_0 = X_1$, puisque $U \to X$
est surjectif. Soit $U = U_0 = U_1 \supset U_2 \ldots$ la suite induite de
sous-schémas ouverts de $U$.

\medskip\noindent
Choisissons un ouvert affine non vide $V_1 \subset U_1$ (par exemple
$V_1 = U_1$). Par récurrence, nous allons construire une suite d'ouverts
affines non vides $V_1 \supset V_2 \supset \ldots$ tels que $V_n \subset U_n$.
En effet, une fois $V_1, \ldots, V_{n - 1}$ construits, nous pouvons toujours
choisir $V_n$, sauf si $V_{n - 1} \cap U_n = \emptyset$. Mais si
$V_{n - 1} \cap U_n = \emptyset$, alors le sous-espace ouvert
$X' \subset X$ tel que
$|X'| = \Im(|V_{n - 1}| \to |X|)$ est contenu dans $|X| \setminus |X_n|$.
Par conséquent, $V_{n - 1} \to X'$ est un morphisme étale dont les fibres ont
un degré majoré par $n - 1$. Autrement dit, $X'$ est raisonnable (par
définition); ainsi, $X'$ contient un sous-schéma ouvert non vide d'après la
Proposition \ref{decent-spaces-proposition-reasonable-open-dense-scheme}.
C'est une contradiction, qui montre que nous pouvons choisir $V_n$.

\medskip\noindent
D'après Limites, lemme \ref{limits-lemma-limit-nonempty},
la limite $V_\infty = \lim V_n$ est un schéma non vide. Choisissons un
morphisme $\Spec(k) \to V_\infty$. La composée
$\Spec(k) \to V_\infty \to U \to X$ a, par construction, son image contenue
dans tous les $X_d$. Autrement dit, la fibre $U \times_X \Spec(k)$ est de
degré infini, ce qui contredit la définition d'un espace décent. Cette
contradiction achève la démonstration du théorème.
\end{proof}

\begin{lemma}
\label{decent-spaces-lemma-when-quotient-scheme-at-point}
Soit $S$ un schéma. Soit $X \to Y$ un morphisme surjectif, fini et
localement libre d'espaces algébriques sur $S$. Pour $y \in |Y|$, les
conditions suivantes sont équivalentes :
\begin{enumerate}
\item $y$ appartient au lieu schématique de $Y$ ;
\item il existe un ouvert affine $U \subset X$ contenant l'image inverse
de $y$.
\end{enumerate}
\end{lemma}

\begin{proof}
Si $y \in Y$ appartient au lieu schématique, il possède un voisinage ouvert
affine $V \subset Y$, et l'image inverse $U$ de $V$ dans $X$ est ouverte et
finie sur $V$, donc affine. Ainsi, (1) implique (2).

\medskip\noindent
Réciproquement, supposons donné $U \subset X$ comme en (2).
Posons $R = X \times_Y X$ et notons $s, t : R \to X$ les projections.
Considérons $Z = R \setminus (s^{-1}(U) \cap t^{-1}(U))$.
C'est une partie fermée de $R$. L'image $t(Z)$ est une partie fermée
de $X$ que l'on peut décrire informellement comme l'ensemble des
points de $X$ qui sont $R$-équivalents à un point de
$X \setminus U$. Ainsi, $U' = X \setminus t(Z)$ est stable par $R$ et
constitue un sous-espace ouvert de $X$ contenu dans $U$ et contenant
la fibre de $X \to Y$ au-dessus de $y$. Comme $X \to Y$ est ouvert
(Morphismes d'espaces, lemme \ref{spaces-morphisms-lemma-fppf-open}),
l'image de $U'$ est un sous-espace ouvert $V' \subset Y$.
Puisque $U'$ est stable par $R$ et que $R = X \times_Y X$, on voit que $U'$
est l'image inverse de $V'$ (utiliser
Propriétés des espaces, lemme \ref{spaces-properties-lemma-points-cartesian}).
En remplaçant $Y$ par $V'$ et $X$ par $U'$, on voit que l'on peut supposer
que $X$ est un schéma isomorphe à un sous-schéma ouvert d'un schéma affine.

\medskip\noindent
Supposons que $X$ soit un schéma isomorphe à un sous-schéma ouvert d'un
schéma affine. Dans ce cas, le faisceau quotient fppf $X/R$ est un schéma,
voir Propriétés des espaces, Proposition
\ref{spaces-properties-proposition-finite-flat-equivalence-global}.
Puisque $Y$ est un faisceau pour la topologie fppf, on obtient un morphisme
canonique $X/R \to Y$ par lequel se factorise $X \to Y$. Comme $X \to Y$ est
surjectif, fini et localement libre, il est surjectif comme morphisme de
faisceaux
(Espaces, lemme \ref{spaces-lemma-surjective-flat-locally-finite-presentation}).
On en conclut que $X/R \to Y$ est surjectif comme morphisme de faisceaux.
D'autre part, puisque $R = X \times_Y X$ comme faisceaux, on en conclut que
$X/R \to Y$ est injectif comme morphisme de faisceaux. Ainsi, $X/R \to Y$
est un isomorphisme, et l'on voit que $Y$ est représentable.
\end{proof}

\noindent
À ce stade, nous disposons de plusieurs méthodes différentes pour démontrer
le lemme suivant.

\begin{lemma}
\label{decent-spaces-lemma-finite-etale-cover-dense-open-scheme}
Soit $S$ un schéma. Soit $X$ un espace algébrique sur $S$.
S'il existe un morphisme fini, étale et surjectif
$U \to X$, où $U$ est un schéma, alors il existe un sous-espace ouvert dense
de $X$ qui est un schéma.
\end{lemma}

\begin{proof}[Première démonstration]
Le morphisme $U \to X$ est fini localement libre. Il existe donc une
décomposition de $X$ en sous-espaces ouverts et fermés $X_d \subset X$ telle
que $U \times_X X_d \to X_d$ soit fini localement libre de degré $d$.
On peut donc supposer que $U \to X$ est fini localement libre de degré $d$.
Dans ce cas, soient $U_i \subset U$, $i \in I$, les ouverts affines.
Pour tout $i$, le morphisme $U_i \to X$ est étale et ses fibres sont
universellement bornées (à savoir, majorées par $d$).
Autrement dit, $X$ est raisonnable, et le résultat découle de la
Proposition \ref{decent-spaces-proposition-reasonable-open-dense-scheme}.
\end{proof}

\begin{proof}[Deuxième démonstration]
La question est locale sur $X$
(Propriétés des espaces, lemme \ref{spaces-properties-lemma-subscheme}) ;
on peut donc supposer $X$ quasi-compact. Alors $U$ est quasi-compact.
Il existe alors un sous-schéma ouvert dense $W \subset U$ qui est séparé
(Propriétés, lemme
\ref{properties-lemma-quasi-compact-dense-open-separated}).
Posons $Z = U \setminus W$.
Soit $R = U \times_X U$, et soient $s, t : R \to U$ les projections.
Alors $t^{-1}(Z)$ est une partie rare de $R$
(Topologie, lemme \ref{topology-lemma-open-inverse-image-closed-nowhere-dense})
et donc $\Delta = s(t^{-1}(Z))$ est une partie stable par $R$, fermée et rare
dans $U$
(Morphismes, lemme \ref{morphisms-lemma-image-nowhere-dense-finite}).
Soit $u \in U \setminus \Delta$ un point générique d'une composante
irréductible. Puisque ces points sont denses dans $U \setminus \Delta$
et que $\Delta$ est rare, il suffit de montrer que l'image
$x \in X$ de $u$ appartient au lieu schématique de $X$.
Remarquons que $t(s^{-1}(\{u\})) \subset W$ est un
ensemble fini de points génériques de composantes irréductibles de $W$
(comparer avec
Propriétés des espaces, lemme
\ref{spaces-properties-lemma-codimension-0-points}).
D'après Propriétés, lemme \ref{properties-lemma-maximal-points-affine},
on peut trouver un ouvert affine $V \subset W$ tel que
$t(s^{-1}(\{u\})) \subset V$. Puisque $t(s^{-1}(\{u\}))$ est la fibre
de $|U| \to |X|$ au-dessus de $x$, on conclut par le
Lemme \ref{decent-spaces-lemma-when-quotient-scheme-at-point}.
\end{proof}

\begin{proof}[Troisième démonstration]
(Cette démonstration est essentiellement identique à la deuxième, mais
utilise moins de renvois.)
Supposons que $X$ soit un espace algébrique, que $U$ soit un schéma et que
$U \to X$ soit un morphisme fini, étale et surjectif. Écrivons
$R = U \times_X U$ et notons, comme d'habitude,
$s, t : R \to U$ les projections. Remarquons que $s, t$ sont surjectifs,
finis et étales. Assertion : la réunion des ouverts affines stables par $R$
et contenus dans $U$ est topologiquement dense dans $U$.

\medskip\noindent
Démonstration de l'assertion. Soit $W \subset U$ un ouvert affine.
Posons $W' = t(s^{-1}(W)) \subset U$. Comme $s^{-1}(W)$ est affine
(donc quasi-compact), on voit que $W' \subset U$ est un ouvert
quasi-compact. D'après
Propriétés, lemme \ref{properties-lemma-quasi-compact-dense-open-separated},
il existe un ouvert dense $W'' \subset W'$ qui est un schéma séparé.
Posons $\Delta' = W' \setminus W''$. C'est une partie fermée rare de $W'$.
Comme $t|_{s^{-1}(W)} : s^{-1}(W) \to W'$ est ouvert (car il est étale),
on voit que l'image inverse
$(t|_{s^{-1}(W)})^{-1}(\Delta') \subset s^{-1}(W)$
est une partie fermée rare (voir
Topologie, lemme \ref{topology-lemma-open-inverse-image-closed-nowhere-dense}).
Il résulte alors de
Morphismes, lemme \ref{morphisms-lemma-image-nowhere-dense-finite}, que
$$
\Delta = s\left((t|_{s^{-1}(W)})^{-1}(\Delta')\right)
$$
est une partie fermée rare de $W$. Choisissons un point $\eta \in W$,
$\eta \not \in \Delta$, qui soit un point générique d'une composante
irréductible de $W$ (et donc de $U$). Par les choix effectués ci-dessus,
l'ensemble fini
$t(s^{-1}(\{\eta\})) = \{\eta_1, \ldots, \eta_n\}$
est contenu dans le schéma séparé $W''$.
Notons que les fibres de $s$ sont des espaces discrets finis et que les
générisations se relèvent le long du morphisme étale $t$, voir
Morphismes, lemmes \ref{morphisms-lemma-etale-flat}
et \ref{morphisms-lemma-generalizations-lift-flat}.
On voit ainsi que chaque $\eta_i$ est un point générique d'une composante
irréductible de $W''$. Par conséquent, d'après
Propriétés, lemme \ref{properties-lemma-maximal-points-affine},
on peut trouver un ouvert affine $V \subset W''$ tel que
$\{\eta_1, \ldots, \eta_n\} \subset V$.
D'après
Groupoïdes, lemme \ref{groupoids-lemma-find-invariant-affine},
cela implique que $\eta$ appartient à un sous-schéma ouvert affine stable
par $R$ de $U$. L'assertion en résulte, puisque $W$ a été choisi comme un
ouvert affine arbitraire de $U$ et que l'ensemble des points génériques
des composantes irréductibles de $W \setminus \Delta$ est dense dans $W$.

\medskip\noindent
L'assertion permet d'achever la démonstration. En effet, si $W \subset U$ est
un ouvert affine stable par $R$, alors la restriction $R_W$ de $R$ à $W$
vérifie $R_W = s^{-1}(W) = t^{-1}(W)$ (voir
Groupoïdes, définition \ref{groupoids-definition-invariant-open}
et la discussion qui suit). En particulier, les morphismes $R_W \to W$
sont eux aussi finis et étales. Il s'ensuit notamment que $R_W$ est affine.
Ainsi, $W/R_W$ est un schéma, d'après
Groupoïdes, Proposition \ref{groupoids-proposition-finite-flat-equivalence}.
D'autre part, $W/R_W$ est un sous-espace ouvert de $X$, d'après
Espaces, lemme \ref{spaces-lemma-finding-opens}.
Par conséquent, l'existence d'un ensemble dense de points contenus dans des
ouverts affines stables par $R$ de $U$ implique bien que le lieu schématique
de $X$ (voir Propriétés des espaces, lemme
\ref{spaces-properties-lemma-subscheme}) est un ouvert dense dans $X$.
\end{proof}













\section{Corps résiduels et anneaux locaux henséliens}
\label{decent-spaces-section-residue-fields-henselian-local-rings}

\noindent
Pour un espace algébrique décent, nous pouvons définir le corps résiduel et
l'anneau local hensélien en un point. Par exemple, le lemme suivant montre que
le corps résiduel d'un point d'un espace décent est bien défini.

\begin{lemma}
\label{decent-spaces-lemma-decent-points-monomorphism}
Soit $S$ un schéma. Soit $X$ un espace algébrique sur $S$.
Considérons l'application
$$
\{\Spec(k) \to X \text{ monomorphisme où }k\text{ est un corps}\}
\longrightarrow
|X|
$$
Cette application est toujours injective. Si $X$ est décent, elle est
bijective.
\end{lemma}

\begin{proof}
Nous avons vu dans
Propriétés des espaces,
lemme \ref{spaces-properties-lemma-points-monomorphism}
que cette application est injective en général.
D'après le Lemme \ref{decent-spaces-lemma-bounded-fibres}, elle est surjective lorsque $X$
est décent (on peut même dire que cela fait partie de la définition d'un
espace décent).
\end{proof}

\noindent
Soit $S$ un schéma. Soit $X$ un espace algébrique sur $S$.
Si un point $x \in |X|$ peut être représenté par un monomorphisme
$\Spec(k) \to X$, alors le corps $k$ est unique à isomorphisme unique près.
Pour un espace algébrique décent, un tel monomorphisme existe pour tout point
d'après le Lemme \ref{decent-spaces-lemma-decent-points-monomorphism} ; la définition
suivante a donc un sens.

\begin{definition}
\label{decent-spaces-definition-residue-field}
Soit $S$ un schéma. Soit $X$ un espace algébrique décent sur $S$.
Soit $x \in |X|$. Le {\it corps résiduel de $X$ au point $x$}
est l'unique corps $\kappa(x)$ muni d'un
monomorphisme $\Spec(\kappa(x)) \to X$ représentant $x$.
\end{definition}

\noindent
Soit $S$ un schéma. Soit $f : X \to Y$ un morphisme d'espaces algébriques
décents sur $S$. Soit $x \in |X|$ un point.
Posons $y = f(x) \in |Y|$. Alors la composée $\Spec(\kappa(x)) \to Y$
appartient à la classe d'équivalence qui définit $y$ et se factorise donc par
$\Spec(\kappa(y)) \to Y$. Autrement dit, on obtient un diagramme commutatif
$$
\xymatrix{
\Spec(\kappa(x)) \ar[r]_-x \ar[d] & X \ar[d]^f \\
\Spec(\kappa(y)) \ar[r]^-y & Y
}
$$
Le morphisme vertical de gauche correspond à un homomorphisme de corps
$\kappa(y) \to \kappa(x)$. Nous l'appellerons souvent simplement
l'homomorphisme induit par $f$.

\begin{lemma}
\label{decent-spaces-lemma-identifies-residue-fields}
Soit $S$ un schéma. Soit $f : X \to Y$ un morphisme d'espaces algébriques
décents sur $S$. Soit $x \in |X|$ un point
d'image $y = f(x) \in |Y|$.
Les conditions suivantes sont équivalentes :
\begin{enumerate}
\item $f$ induit un isomorphisme $\kappa(y) \to \kappa(x)$ ;
\item le morphisme induit $\Spec(\kappa(x)) \to Y$ est un monomorphisme.
\end{enumerate}
\end{lemma}

\begin{proof}
Cela résulte immédiatement de la discussion précédente.
\end{proof}

\noindent
Le lemme suivant montre que l'anneau local hensélien d'un point
d'un espace algébrique décent est bien défini.

\begin{lemma}
\label{decent-spaces-lemma-decent-space-elementary-etale-neighbourhood}
Soit $S$ un schéma. Soit $X$ un espace algébrique décent sur $S$.
Pour tout point $x \in |X|$, il existe un morphisme étale
$$
(U, u) \longrightarrow (X, x)
$$
où $U$ est un schéma affine, $u$ est l'unique point de $U$ au-dessus de
$x$, et l'homomorphisme induit $\kappa(x) \to \kappa(u)$ est un isomorphisme.
\end{lemma}

\begin{proof}
On peut supposer que $X$ est quasi-compact en remplaçant $X$ par un ouvert
quasi-compact contenant $x$. Rappelons que $x$ peut être représenté par un
(mono)morphisme quasi-compact provenant du spectre d'un corps (par définition
des espaces décents). Le lemme résulte donc du
Lemme \ref{decent-spaces-lemma-filter-quasi-compact}.
\end{proof}

\begin{definition}
\label{decent-spaces-definition-elemenary-etale-neighbourhood}
Soit $S$ un schéma. Soit $X$ un espace algébrique sur $S$.
Soit $x \in X$ un point. Un {\it voisinage étale élémentaire}
est un morphisme étale $(U, u) \to (X, x)$ où $U$ est un schéma,
$u \in U$ est un point au-dessus de $x$, et le morphisme
$u = \Spec(\kappa(u)) \to X$ est un monomorphisme.
Un {\it morphisme de voisinages étales élémentaires}
$(U, u) \to (U', u')$ est, par définition, un morphisme $U \to U'$
au-dessus de $X$ qui envoie $u$ sur $u'$.
\end{definition}

\noindent
Si $X$ n'est pas décent, la catégorie des voisinages étales élémentaires
peut être vide.

\begin{lemma}
\label{decent-spaces-lemma-elementary-etale-neighbourhoods}
Soit $S$ un schéma. Soit $X$ un espace algébrique décent sur $S$.
Soit $x$ un point de $X$.
La catégorie des voisinages étales élémentaires de $(X, x)$
est cofiltrante (voir
Catégories, définition \ref{categories-definition-codirected}).
\end{lemma}

\begin{proof}
La catégorie est non vide d'après le
Lemme \ref{decent-spaces-lemma-decent-space-elementary-etale-neighbourhood}.
Supposons donnés deux voisinages étales élémentaires
$(U_i, u_i) \to (X, x)$.
Considérons alors $U = U_1 \times_X U_2$. Puisque
$\Spec(\kappa(u_i)) \to X$, $i = 1, 2$, sont tous deux des monomorphismes
dans la classe de $x$ (Lemme \ref{decent-spaces-lemma-identifies-residue-fields}),
on voit que
$$
u = \Spec(\kappa(u_1)) \times_X \Spec(\kappa(u_2))
$$
est le spectre d'un corps $\kappa(u)$ tel que les homomorphismes induits
$\kappa(u_i) \to \kappa(u)$ sont des isomorphismes. Alors
$u \to U$ est un point de $U$, et l'on voit que $(U, u) \to (X, x)$ est un
voisinage étale élémentaire qui domine $(U_i, u_i)$.
Si $a, b : (U_1, u_1) \to (U_2, u_2)$ sont deux morphismes entre
nos voisinages étales élémentaires, considérons le schéma
$$
U = U_1 \times_{(a, b), (U_2 \times_X U_2), \Delta} U_2
$$
D'après Propriétés des espaces, lemme
\ref{spaces-properties-lemma-etale-permanence},
$U \to X$ est étale. De plus, exactement de la même manière que précédemment,
on voit que $U$ possède un point $u$
tel que $(U, u) \to (X, x)$ est un voisinage
étale élémentaire. Enfin, $U \to U_1$ égalise $a$ et $b$,
ce qui achève la démonstration.
\end{proof}

\begin{definition}
\label{decent-spaces-definition-henselian-local-ring}
Soit $S$ un schéma. Soit $X$ un espace algébrique décent sur $S$.
Soit $x \in |X|$. L'{\it anneau local hensélien de $X$ en $x$} est
$$
\mathcal{O}_{X, x}^h = \colim \Gamma(U, \mathcal{O}_U)
$$
où la limite inductive est prise sur les voisinages étales élémentaires
$(U, u) \to (X, x)$.
\end{definition}

\noindent
Voici l'analogue du
Propriétés des espaces, lemme
\ref{spaces-properties-lemma-describe-etale-local-ring}.

\begin{lemma}
\label{decent-spaces-lemma-describe-henselian-local-ring}
Soit $S$ un schéma. Soit $X$ un espace algébrique décent sur $S$.
Soit $x \in |X|$. Soit $(U, u) \to (X, x)$ un voisinage
étale élémentaire. Alors
$$
\mathcal{O}_{X, x}^h = \mathcal{O}_{U, u}^h
$$
Autrement dit : l'anneau local hensélien de $X$ en $x$
est égal au hensélisé $\mathcal{O}_{U, u}^h$
de l'anneau local $\mathcal{O}_{U, u}$ de $U$ en $u$.
\end{lemma}

\begin{proof}
Puisque la catégorie des voisinages étales élémentaires de $(X, x)$
est cofiltrante (Lemme \ref{decent-spaces-lemma-elementary-etale-neighbourhoods}),
on voit que
la catégorie des voisinages étales élémentaires de $(U, u)$
est initiale dans
la catégorie des voisinages étales élémentaires de $(X, x)$.
L'égalité résulte alors du
Compléments sur les morphismes, lemme \ref{more-morphisms-lemma-describe-henselization}
et du
Catégories, lemme \ref{categories-lemma-cofinal}
(initiale devient cofinale, car la limite inductive
qui définit les anneaux locaux henséliens est prise sur la
catégorie opposée de la catégorie des voisinages
étales élémentaires).
\end{proof}

\begin{lemma}
\label{decent-spaces-lemma-henselian-local-ring-strict}
Soit $S$ un schéma. Soit $X$ un espace algébrique décent sur $S$.
Soit $\overline{x}$ un point géométrique de $X$ au-dessus de $x \in |X|$.
L'anneau local étale $\mathcal{O}_{X, \overline{x}}$ de $X$ en $\overline{x}$
(Propriétés des espaces, définition
\ref{spaces-properties-definition-etale-local-rings})
est le hensélisé strict
de l'anneau local hensélien $\mathcal{O}_{X, x}^h$ de $X$ en $x$.
\end{lemma}

\begin{proof}
Cela résulte du Lemme \ref{decent-spaces-lemma-describe-henselian-local-ring},
de Propriétés des espaces, lemme
\ref{spaces-properties-lemma-describe-etale-local-ring}
et du fait que $(R^h)^{sh} = R^{sh}$
pour un anneau local $(R, \mathfrak m, \kappa)$ et une
clôture séparable donnée $\kappa^{sep}$ de $\kappa$.
Cette égalité résulte du
Algèbre, lemme \ref{algebra-lemma-uniqueness-henselian}.
\end{proof}

\begin{lemma}
\label{decent-spaces-lemma-residue-field-henselian-local-ring}
Soit $S$ un schéma. Soit $X$ un espace algébrique décent sur $S$.
Soit $x \in |X|$. Le corps résiduel de
l'anneau local hensélien de $X$ en $x$
(Définition \ref{decent-spaces-definition-henselian-local-ring})
est le corps résiduel de $X$ en $x$
(Définition \ref{decent-spaces-definition-residue-field}).
\end{lemma}

\begin{proof}
Choisissons un voisinage étale élémentaire $(U, u) \to (X, x)$.
Alors $\kappa(u) = \kappa(x)$ et
$\mathcal{O}_{X, x}^h = \mathcal{O}_{U, u}^h$
(Lemme \ref{decent-spaces-lemma-describe-henselian-local-ring}).
Le corps résiduel de $\mathcal{O}_{U, u}^h$
est $\kappa(u)$ d'après Algèbre, lemme \ref{algebra-lemma-henselization}
(le résultat de ce lemme est la construction/définition
du hensélisé d'un anneau local, voir
Algèbre, définition \ref{algebra-definition-henselization}).
\end{proof}

\begin{remark}
\label{decent-spaces-remark-functoriality-henselian-local-ring}
Soit $S$ un schéma. Soit $f : X \to Y$ un morphisme
d'espaces algébriques décents sur $S$. Soit $x \in |X|$
d'image $y \in |Y|$. Choisissons des points géométriques compatibles
$(\overline{x} \to X, \overline{y} = f(\overline{x}) \to Y)$
au-dessus de $(x, y)$.
Choisissons un voisinage étale élémentaire $(V, v) \to (Y, y)$ (ce qui est
possible d'après le
Lemme \ref{decent-spaces-lemma-decent-space-elementary-etale-neighbourhood}).
Alors $V \times_Y X$ est un espace algébrique étale sur $X$
qui possède un unique point $x'$ dont les images sont $x$ dans $X$
et $v$ dans $V$.
(Les détails sont omis; on utilise le fait que tous les points peuvent être
représentés par des monomorphismes provenant de spectres de corps.)
Choisissons un voisinage étale élémentaire $(U, u) \to (V \times_Y X, x')$.
Nous obtenons alors le diagramme commutatif suivant
$$
\xymatrix{
\Spec(\mathcal{O}_{X, \overline{x}}) \ar[r] \ar[d] &
\Spec(\mathcal{O}_{X, x}^h) \ar[r] \ar[d] &
\Spec(\mathcal{O}_{U, u}) \ar[r] \ar[d] &
U \ar[r] \ar[d] &
X \ar[d] \\
\Spec(\mathcal{O}_{Y, \overline{y}}) \ar[r] &
\Spec(\mathcal{O}_{Y, y}^h) \ar[r] &
\Spec(\mathcal{O}_{V, v}) \ar[r] &
V \ar[r] &
Y
}
$$
Cela provient des identifications
$\mathcal{O}_{X, \overline{x}} = \mathcal{O}_{U, u}^{sh}$,
$\mathcal{O}_{X, x}^h = \mathcal{O}_{U, u}^h$,
$\mathcal{O}_{Y, \overline{y}} = \mathcal{O}_{V, v}^{sh}$,
$\mathcal{O}_{Y, y}^h = \mathcal{O}_{V, v}^h$
données par le
Lemme \ref{decent-spaces-lemma-describe-henselian-local-ring}
et par
Propriétés des espaces, lemme
\ref{spaces-properties-lemma-describe-etale-local-ring},
ainsi que de la fonctorialité de l'hensélisation (stricte)
étudiée dans Algèbre, sections \ref{algebra-section-ind-etale} et
\ref{algebra-section-henselization}.
\end{remark}









\section{Points des espaces décents}
\label{decent-spaces-section-points}

\noindent
Dans cette section, nous démontrons quelques propriétés des points des
espaces algébriques décents. Le lemme suivant montre que la spécialisation
des points se comporte bien dans les espaces algébriques décents.
Espaces, exemple \ref{spaces-example-infinite-product}
montre que ceci n'est {\bf pas} vrai en général.

\begin{lemma}
\label{decent-spaces-lemma-decent-no-specializations-map-to-same-point}
Soit $S$ un schéma. Soit $X$ un espace algébrique décent sur $S$.
Soit $U \to X$ un morphisme \'etale d'un schéma vers $X$.
Si $u, u' \in |U|$ ont même image dans $|X|$ et si
$u' \leadsto u$, alors $u = u'$.
\end{lemma}

\begin{proof}
Combiner les Lemmes \ref{decent-spaces-lemma-bounded-fibres} et
\ref{decent-spaces-lemma-no-specializations-map-to-same-point}.
\end{proof}

\begin{lemma}
\label{decent-spaces-lemma-decent-specialization}
Soit $S$ un schéma. Soit $X$ un espace algébrique décent sur $S$.
Soient $x, x' \in |X|$ et supposons que $x' \leadsto x$, c'est-à-dire que
$x$ est une spécialisation de $x'$. Alors, pour tout morphisme \'etale
$\varphi : U \to X$ d'un schéma $U$ vers X et tout $u \in U$ tel que
$\varphi(u) = x$, il existe un point $u'\in U$ tel que
$u' \leadsto u$ et $\varphi(u') = x'$.
\end{lemma}

\begin{proof}
Combiner les Lemmes \ref{decent-spaces-lemma-bounded-fibres} et
\ref{decent-spaces-lemma-specialization}.
\end{proof}

\begin{lemma}
\label{decent-spaces-lemma-kolmogorov}
Soit $S$ un schéma. Soit $X$ un espace algébrique décent sur $S$.
Alors $|X|$ est un espace de Kolmogorov (voir
Topologie, définition \ref{topology-definition-generic-point}).
\end{lemma}

\begin{proof}
Soient $x_1, x_2 \in |X|$ tels que $x_1 \leadsto x_2$ et
$x_2 \leadsto x_1$. Nous devons montrer que $x_1 = x_2$. Choisissons un
schéma $U$ et un morphisme \'etale $U \to X$ tels que $x_1, x_2$
appartiennent tous deux à l'image de $|U| \to |X|$. D'après le
Lemme \ref{decent-spaces-lemma-decent-specialization}, nous pouvons trouver une spécialisation
$u_1 \leadsto u_2$ dans $U$ qui s'envoie sur $x_1 \leadsto x_2$.
D'après le Lemme \ref{decent-spaces-lemma-decent-specialization}, il existe
$u_2' \leadsto u_1$ qui s'envoie sur $x_2 \leadsto x_1$. Ainsi,
$u_2' \leadsto u_2$ est une spécialisation entre des points de $U$ qui
s'envoient sur le même point de $X$, à savoir $x_2$. Cela n'est possible que
si $u_2' = u_2$, voir le
Lemme \ref{decent-spaces-lemma-decent-no-specializations-map-to-same-point}. Par conséquent,
$u_1 = u_2$ également, comme voulu.
\end{proof}

\begin{proposition}
\label{decent-spaces-proposition-reasonable-sober}
Soit $S$ un schéma. Soit $X$ un espace algébrique décent sur $S$.
Alors l'espace topologique $|X|$ est sobre (voir
Topologie, définition \ref{topology-definition-generic-point}).
\end{proposition}

\begin{proof}
Nous avons vu au Lemme \ref{decent-spaces-lemma-kolmogorov} que $|X|$ est un espace de
Kolmogorov. Il reste donc à montrer que toute partie fermée irréductible
$T \subset |X|$ possède un point générique. D'après
Propriétés des espaces,
lemme \ref{spaces-properties-lemma-reduced-closed-subspace},
il existe un sous-espace fermé $Z \subset X$ tel que $|Z| = T$.
Par définition, cela signifie que $Z \to X$ est un morphisme représentable
d'espaces algébriques. Ainsi, $Z$ est un espace algébrique décent
d'après le Lemme \ref{decent-spaces-lemma-representable-properties}. Le
Théorème \ref{decent-spaces-theorem-decent-open-dense-scheme}
montre qu'il existe un sous-espace ouvert dense $Z' \subset Z$ qui est un
schéma. Cela signifie que $|Z'| \subset T$ est ouvert dense.
L'espace topologique $|Z'|$ est donc irréductible, ce qui signifie que
$Z'$ est un schéma irréductible. D'après
Schémas, lemme \ref{schemes-lemma-scheme-sober},
on conclut que $|Z'|$ est l'adhérence d'un seul point $\eta \in T$,
et donc aussi que $T = \overline{\{\eta\}}$, ce qui achève la preuve.
\end{proof}

\noindent
Pour les espaces algébriques décents, la dimension se comporte comme prévu.

\begin{lemma}
\label{decent-spaces-lemma-dimension-decent-space}
Soit $S$ un schéma. La dimension telle qu'elle est définie dans
Propriétés des espaces, section \ref{spaces-properties-section-dimension}
se comporte bien sur les espaces algébriques décents $X$ au-dessus de $S$.
\begin{enumerate}
\item Si $x \in |X|$, alors $\dim_x(|X|) = \dim_x(X)$, et
\item $\dim(|X|) = \dim(X)$.
\end{enumerate}
\end{lemma}

\begin{proof}
Démonstration de (1). Choisissons un schéma $U$ muni d'un point $u \in U$
et un morphisme \'etale $h : U \to X$ qui envoie $u$ sur $x$.
Par définition, la dimension de $X$ en $x$ est $\dim_u(|U|)$.
Nous pouvons donc choisir $U$ tel que $\dim_x(X) = \dim(|U|)$.
Soit $d$ un entier. Si $\dim(U) \geq d$, il existe alors
une suite de spécialisations non triviales
$u_d \leadsto \ldots \leadsto u_0$ dans $U$. En prenant l'image,
nous trouvons une suite correspondante
$h(u_d) \leadsto \ldots \leadsto h(u_0)$
dont chacune des spécialisations est non triviale d'après le
Lemme \ref{decent-spaces-lemma-decent-no-specializations-map-to-same-point}.
Ainsi, l'image de $|U|$ dans $|X|$ est de dimension au moins $d$.
Réciproquement, supposons que $x_d \leadsto \ldots \leadsto x_0$ soit une
suite de spécialisations dans $|X|$ telle que $x_0$ appartienne à l'image de
$|U| \to |X|$. Nous pouvons alors la relever en une suite de spécialisations
dans $U$ d'après le Lemme \ref{decent-spaces-lemma-decent-specialization}.

\medskip\noindent
La partie (2) est une conséquence immédiate de la partie (1),
de Topologie, lemme \ref{topology-lemma-dimension-supremum-local-dimensions},
et de Propriétés des espaces, section
\ref{spaces-properties-section-dimension}.
\end{proof}

\begin{lemma}
\label{decent-spaces-lemma-dimension-local-ring-quasi-finite}
Soit $S$ un schéma. Soit $X \to Y$ un morphisme localement quasi-fini
d'espaces algébriques sur $S$. Soit $x \in |X|$ d'image $y \in |Y|$.
Alors la dimension de l'anneau local de $Y$ en $y$ est $\geq$ celle de
l'anneau local de $X$ en $x$.
\end{lemma}

\begin{proof}
La définition de la dimension de l'anneau local d'un point d'un espace
algébrique est donnée dans Propriétés des espaces, définition
\ref{spaces-properties-definition-dimension-local-ring}.
Choisissons un morphisme \'etale $(V, v) \to (Y, y)$ où $V$ est un schéma.
Choisissons un morphisme \'etale $U \to V \times_Y X$ et un point $u \in U$
qui s'envoie sur $x \in |X|$ et sur $v \in V$. Alors $U \to V$ est localement
quasi-fini et nous devons démontrer que
$$
\dim(\mathcal{O}_{V, v}) \geq \dim(\mathcal{O}_{U, u})
$$
C'est Algèbre, lemme \ref{algebra-lemma-dimension-inequality-quasi-finite}.
\end{proof}

\begin{lemma}
\label{decent-spaces-lemma-dimension-quasi-finite}
Soit $S$ un schéma. Soit $X \to Y$ un morphisme localement quasi-fini
d'espaces algébriques sur $S$. Alors $\dim(X) \leq \dim(Y)$.
\end{lemma}

\begin{proof}
Cela résulte du Lemme \ref{decent-spaces-lemma-dimension-local-ring-quasi-finite}
et de Propriétés des espaces, lemme \ref{spaces-properties-lemma-dimension}.
\end{proof}

\noindent
Le lemme suivant est très légèrement plus fort que
Propriétés des espaces,
lemme \ref{spaces-properties-lemma-point-like-spaces}.
Nous améliorerons ce lemme au Lemme \ref{decent-spaces-lemma-when-field}.

\begin{lemma}
\label{decent-spaces-lemma-decent-point-like-spaces}
Soit $S$ un schéma. Soit $k$ un corps. Soit $X$ un espace algébrique
sur $S$ et supposons qu'il existe un morphisme \'etale surjectif
$\Spec(k) \to X$. Si $X$ est décent, alors $X \cong \Spec(k')$,
où $k/k'$ est une extension finie séparable.
\end{lemma}

\begin{proof}
L'hypothèse implique que $|X| = \{x\}$ ne possède qu'un point. Comme
$X$ est décent, il existe un monomorphisme quasi-compact
$\Spec(k') \to X$ dont l'image est $x$. Alors la projection
$U = \Spec(k') \times_X \Spec(k) \to \Spec(k)$
est un monomorphisme, d'où $U = \Spec(k)$, voir
Schémas, lemme \ref{schemes-lemma-mono-towards-spec-field}.
Par conséquent, la projection $\Spec(k) = U \to \Spec(k')$ est \'etale,
ce qui achève la preuve.
\end{proof}






\section{Espaces réduits à point unique}
\label{decent-spaces-section-singleton}

\noindent
Un espace {\it à point unique} est un espace algébrique $X$ tel que $|X|$
est réduit à un point. De tels espaces peuvent en fait être plus
intéressants que le seul spectre d'un corps, voir
Espaces, exemple \ref{spaces-example-Qbar}.
Nous développons un peu de formalisme afin de pouvoir en parler.

\begin{lemma}
\label{decent-spaces-lemma-flat-cover-by-field}
Soit $S$ un schéma. Soit $Z$ un espace algébrique sur $S$.
Soit $k$ un corps et soit $\Spec(k) \to Z$ un morphisme surjectif et plat.
Alors tout morphisme $\Spec(k') \to Z$, où $k'$ est un corps, est
surjectif et plat.
\end{lemma}

\begin{proof}
Considérons le carré cartésien
$$
\xymatrix{
T \ar[d] \ar[r] & \Spec(k) \ar[d] \\
\Spec(k') \ar[r] & Z
}
$$
Notons que $T \to \Spec(k')$ est plat et surjectif, de sorte que $T$
n'est pas vide. D'autre part, $T \to \Spec(k)$ est plat puisque
$k$ est un corps. Ainsi, $T \to Z$ est plat et surjectif.
Il résulte de
Morphismes d'espaces, lemme \ref{spaces-morphisms-lemma-flat-permanence}
que $\Spec(k') \to Z$ est plat. Ce morphisme est surjectif puisque
l'hypothèse entraîne que $|Z|$ est réduit à un point.
\end{proof}

\begin{lemma}
\label{decent-spaces-lemma-unique-point}
Soit $S$ un schéma.
Soit $Z$ un espace algébrique sur $S$. Les conditions suivantes sont
équivalentes :
\begin{enumerate}
\item $Z$ est réduit et $|Z|$ est réduit à un point,
\item il existe un morphisme plat surjectif $\Spec(k) \to Z$
où $k$ est un corps, et
\item il existe un morphisme localement de type fini, surjectif et plat
$\Spec(k) \to Z$, où $k$ est un corps.
\end{enumerate}
\end{lemma}

\begin{proof}
Supposons (1). Soit $W$ un schéma et
soit $W \to Z$ un morphisme \'etale surjectif. Alors $W$ est
un schéma réduit. Soit $\eta \in W$ un point générique d'une composante
irréductible de $W$. Comme $W$ est réduit, on a
$\mathcal{O}_{W, \eta} = \kappa(\eta)$. Il s'ensuit que le morphisme canonique
$\eta = \Spec(\kappa(\eta)) \to W$ est plat. On voit que le composé
$\eta \to Z$ est plat (voir
Morphismes d'espaces, lemme \ref{spaces-morphisms-lemma-composition-flat}).
Il est aussi surjectif puisque $|Z|$ est réduit à un point. Autrement dit,
(2) est satisfaite.

\medskip\noindent
Supposons (2). Soit $W$ un schéma et
soit $W \to Z$ un morphisme \'etale surjectif. Choisissons un corps
$k$ et un morphisme plat surjectif $\Spec(k) \to Z$.
Alors $W \times_Z \Spec(k)$ est un schéma \'etale sur $k$.
Ainsi, $W \times_Z \Spec(k)$ est une réunion disjointe de spectres de corps
(voir la Remarque \ref{decent-spaces-remark-recall}), donc est en particulier réduit. Comme
$W \times_Z \Spec(k) \to W$
est plat et surjectif, on conclut que $W$ est réduit
(Descente, lemme \ref{descent-lemma-descend-reduced}).
Autrement dit, (1) est satisfaite.

\medskip\noindent
Il est clair que (3) implique (2). Enfin, supposons (2). Choisissons un schéma
affine non vide $W$ et un morphisme \'etale $W \to Z$. Choisissons un point
fermé $w \in W$ et posons $k = \kappa(w)$. Le composé
$$
\Spec(k) \xrightarrow{w} W \longrightarrow Z
$$
est localement de type fini d'après
Morphismes d'espaces, lemmes
\ref{spaces-morphisms-lemma-composition-finite-type} et
\ref{spaces-morphisms-lemma-etale-locally-finite-type}.
Il est également plat et surjectif d'après le
Lemme \ref{decent-spaces-lemma-flat-cover-by-field}.
Ainsi, (3) est satisfaite.
\end{proof}

\noindent
Le lemme suivant isole une classe d'espaces algébriques à point unique
légèrement meilleure que celle du lemme précédent.

\begin{lemma}
\label{decent-spaces-lemma-unique-point-better}
Soit $S$ un schéma. Soit $Z$ un espace algébrique sur $S$.
Les conditions suivantes sont équivalentes :
\begin{enumerate}
\item $Z$ est réduit, localement noethérien, et $|Z|$
est réduit à un point, et
\item il existe un morphisme localement de présentation finie, surjectif et
plat $\Spec(k) \to Z$, où $k$ est un corps.
\end{enumerate}
\end{lemma}

\begin{proof}
Supposons (2) satisfaite. D'après le
Lemme \ref{decent-spaces-lemma-unique-point},
on voit que $Z$ est réduit et que $|Z|$ est réduit à un point.
Soit $W$ un schéma et soit $W \to Z$ un morphisme \'etale surjectif.
Choisissons un corps $k$ et un morphisme localement de présentation finie,
surjectif et plat $\Spec(k) \to Z$.
Alors $W \times_Z \Spec(k)$ est un schéma
\'etale sur $k$, donc une réunion disjointe de spectres de corps
(voir la Remarque \ref{decent-spaces-remark-recall}),
donc localement noethérien. Comme $W \times_Z \Spec(k) \to W$
est plat, surjectif et localement de présentation finie, on voit
que $\{W \times_Z \Spec(k) \to W\}$ est un recouvrement fppf,
et on conclut que $W$ est localement noethérien
(Descente, lemme
\ref{descent-lemma-Noetherian-local-fppf}).
Autrement dit, (1) est satisfaite.

\medskip\noindent
Supposons (1). Choisissons un schéma affine non vide $W$ et un morphisme
\'etale $W \to Z$. Choisissons un point fermé $w \in W$ et posons
$k = \kappa(w)$. Comme $W$ est localement noethérien, le morphisme
$w : \Spec(k) \to W$ est de présentation finie, voir
Morphismes, lemme \ref{morphisms-lemma-closed-immersion-finite-presentation}.
Par conséquent, le composé
$$
\Spec(k) \xrightarrow{w} W \longrightarrow Z
$$
est localement de présentation finie d'après
Morphismes d'espaces, lemmes
\ref{spaces-morphisms-lemma-composition-finite-presentation} et
\ref{spaces-morphisms-lemma-etale-locally-finite-presentation}.
Il est aussi plat et surjectif d'après le
Lemme \ref{decent-spaces-lemma-flat-cover-by-field}.
Ainsi, (2) est satisfaite.
\end{proof}

\begin{lemma}
\label{decent-spaces-lemma-monomorphism-into-point}
Soit $S$ un schéma.
Soit $Z' \to Z$ un monomorphisme d'espaces algébriques sur $S$.
Supposons qu'il existe un corps $k$ et un morphisme localement de présentation
finie, surjectif et plat $\Spec(k) \to Z$. Alors ou bien $Z'$
est vide, ou bien $Z' = Z$.
\end{lemma}

\begin{proof}
Nous pouvons supposer que $Z'$ est non vide. Dans ce cas, le
produit fibré $T = Z' \times_Z \Spec(k)$
est non vide, voir
Propriétés des espaces, lemme \ref{spaces-properties-lemma-points-cartesian}.
Maintenant, $T$ est un espace algébrique et la projection
$T \to \Spec(k)$ est un monomorphisme. Ainsi, $T = \Spec(k)$, voir
Morphismes d'espaces, lemme
\ref{spaces-morphisms-lemma-monomorphism-toward-field}.
On conclut que $\Spec(k) \to Z$ se factorise par $Z'$.
Mais comme $\Spec(k) \to Z$ est surjectif, plat et localement de présentation
finie, on voit que $\Spec(k) \to Z$ est surjectif comme application de
faisceaux sur $(\Sch/S)_{fppf}$ (voir
Espaces, remarque \ref{spaces-remark-warning}),
et on conclut que $Z' = Z$.
\end{proof}

\noindent
Le lemme suivant affirme qu'à chaque point d'un espace algébrique on
peut associer un espace algébrique canonique, réduit, localement noethérien
et à point unique.

\begin{lemma}
\label{decent-spaces-lemma-find-singleton-from-point}
Soit $S$ un schéma. Soit $X$ un espace algébrique sur $S$.
Soit $x \in |X|$. Alors il existe un unique monomorphisme
$Z \to X$ d'espaces algébriques
sur $S$ tel que $Z$ soit un espace algébrique satisfaisant aux conditions
équivalentes du
Lemme \ref{decent-spaces-lemma-unique-point-better}
et tel que l'image de $|Z| \to |X|$ soit $\{x\}$.
\end{lemma}

\begin{proof}
Choisissons un schéma $U$ et un morphisme \'etale surjectif $U \to X$.
Posons $R = U \times_X U$ ; alors $X = U/R$ est une présentation (voir
Espaces, section \ref{spaces-section-presentations}).
Posons
$$
U' = \coprod\nolimits_{u \in U\text{ au-dessus de }x} \Spec(\kappa(u)).
$$
Le morphisme canonique $U' \to U$ est un monomorphisme. Posons
$$
R' = U' \times_X U' = R \times_{(U \times_S U)} (U' \times_S U').
$$
Comme $U' \to U$ est un monomorphisme, on voit que les projections
$s', t' : R' \to U'$ se factorisent en un monomorphisme suivi d'un
morphisme \'etale. Ainsi, puisque $U'$ est une réunion disjointe de spectres
de corps, en utilisant
la Remarque \ref{decent-spaces-remark-recall},
et
Schémas, lemme \ref{schemes-lemma-mono-towards-spec-field},
on conclut que $R'$ est une réunion disjointe de spectres de corps et que
les morphismes $s', t' : R' \to U'$ sont \'etales. Ainsi,
$Z = U'/R'$ est un espace algébrique d'après
Espaces, théorème \ref{spaces-theorem-presentation}.
Comme $R'$ est la restriction de $R$ induite par $U' \to U$, on voit que
$Z \to X$ est un monomorphisme d'après
Groupoïdes, lemme \ref{groupoids-lemma-quotient-groupoid-restrict}.
Puisque $Z \to X$ est un monomorphisme, on voit que $|Z| \to |X|$ est
injective, voir Morphismes d'espaces, lemme
\ref{spaces-morphisms-lemma-monomorphism-injective-points}.
D'après
Propriétés des espaces, lemme \ref{spaces-properties-lemma-points-cartesian},
on voit que
$$
|U'| = |Z \times_X U'| \to |Z| \times_{|X|} |U'|
$$
est surjective, ce qui implique (par notre choix de $U'$) que
l'image de $|Z| \to |X|$ est $\{x\}$. On conclut que $|Z|$ est réduit
à un point. Enfin, par construction, $U'$ est localement noethérien et réduit,
c'est-à-dire que $Z$ satisfait aux conditions équivalentes du
Lemme \ref{decent-spaces-lemma-unique-point-better}.

\medskip\noindent
Prouvons l'unicité de $Z \to X$. Supposons que
$Z' \to X$ soit un second monomorphisme d'espaces algébriques ayant ces
propriétés. Alors les projections
$$
Z' \longleftarrow Z' \times_X Z \longrightarrow Z
$$
sont des monomorphismes. L'espace algébrique du milieu est non vide d'après
Propriétés des espaces,
lemme \ref{spaces-properties-lemma-points-cartesian}.
Par conséquent, les deux projections sont des isomorphismes d'après le
Lemme \ref{decent-spaces-lemma-monomorphism-into-point},
ce qui achève la preuve.
\end{proof}

\noindent
Nous introduisons la terminologie suivante, qui préfigure
les gerbes résiduelles que nous introduirons plus tard, voir
Propriétés des champs, définition
\ref{stacks-properties-definition-residual-gerbe}.

\begin{definition}
\label{decent-spaces-definition-residual-space}
Soit $S$ un schéma.
Soit $X$ un espace algébrique sur $S$. Soit $x \in |X|$.
L'
{\it espace résiduel de $X$ en $x$}\footnote{Notation non standard.}
est le monomorphisme $Z_x \to X$ construit dans le
Lemme \ref{decent-spaces-lemma-find-singleton-from-point}.
\end{definition}

\noindent
En particulier, nous savons que $Z_x$ est un espace algébrique
localement noethérien, réduit et à point unique,
et qu'il existe un corps et un morphisme surjectif, plat et localement
de présentation finie
$$
\Spec(k) \longrightarrow Z_x.
$$
L'espace résiduel est souvent donné par un monomorphisme
provenant du spectre d'un corps.

\begin{lemma}
\label{decent-spaces-lemma-residual-space-monomorphism}
Soit $S$ un schéma. Soit $X$ un espace algébrique sur $S$. Soit $x \in |X|$.
L'espace résiduel $Z_x$ de $X$ en $x$ est isomorphe au spectre d'un corps
si et seulement si $x$ peut être représenté par un monomorphisme
$\Spec(k) \to X$, où $k$ est un corps. Si $X$ est décent, cela vaut pour
tout $x \in |X|$.
\end{lemma}

\begin{proof}
Puisque $Z_x \to X$ est un monomorphisme, si $Z_x = \Spec(k)$ pour un corps
$k$, alors $x$ est représenté par le monomorphisme $\Spec(k) = Z_x \to X$.
Réciproquement, si $\Spec(k) \to X$ est un monomorphisme représentant $x$,
alors $Z_x \times_X \Spec(k) \to \Spec(k)$ est un monomorphisme dont la source
est non vide d'après Propriétés des espaces, lemme
\ref{spaces-properties-lemma-points-cartesian}.
Ainsi, $Z_x \times_X \Spec(k) = \Spec(k)$ d'après Morphismes d'espaces, lemme
\ref{spaces-morphisms-lemma-monomorphism-toward-field}.
On obtient donc un monomorphisme $\Spec(k) \to Z_x$. C'est
un isomorphisme d'après le Lemme \ref{decent-spaces-lemma-monomorphism-into-point}.
La dernière assertion découle du Lemme \ref{decent-spaces-lemma-decent-points-monomorphism}.
\end{proof}

\noindent
L'espace résiduel est un espace algébrique régulier d'après le lemme suivant.

\begin{lemma}
\label{decent-spaces-lemma-residual-space-regular}
Un espace algébrique $Z$ réduit, localement noethérien et à point unique
est régulier.
\end{lemma}

\begin{proof}
Soit $Z$ un espace algébrique réduit, localement noethérien et à point unique
sur un schéma $S$. Soit $W \to Z$ un morphisme \'etale surjectif, où $W$
est un schéma. Soit $k$ un corps et soit $\Spec(k) \to Z$
un morphisme surjectif, plat et localement de présentation finie (voir le
Lemme \ref{decent-spaces-lemma-unique-point-better}).
Le schéma $T = W \times_Z \Spec(k)$ est
\'etale sur $k$, donc en particulier régulier, voir la
Remarque \ref{decent-spaces-remark-recall}.
Comme $T \to W$ est localement de présentation finie, plat et surjectif,
il s'ensuit que $W$ est régulier, voir
Descente, lemme \ref{descent-lemma-descend-regular}.
Par définition, cela signifie que $Z$ est régulier.
\end{proof}

\begin{lemma}
\label{decent-spaces-lemma-factor-through-residual-space-Noetherian}
Soit $S$ un schéma. Soit $f : Y \to X$ un morphisme d'espaces
algébriques sur $S$. Soit $x \in |X|$ un point. Supposons que
\begin{enumerate}
\item $|f|(|Y|)$ soit contenu dans $\{x\} \subset |X|$,
\item $Y$ soit réduit, et
\item $X$ soit localement noethérien.
\end{enumerate}
Alors $f$ se factorise par l'espace résiduel $Z_x$ de $X$ en $x$.
\end{lemma}

\begin{proof}
Remarque préliminaire : puisque $Z_x \to X$ est un monomorphisme, il suffit
de trouver un morphisme \'etale surjectif $Y' \to Y$ tel que $Y' \to X$
se factorise par $Z_x$. Remarquons ici que $Y'$ est également réduit.

\medskip\noindent
Soit $U$ un schéma affine et soit $U \to X$ un morphisme \'etale
tel que $x$ appartienne à l'image de $|U| \to |X|$. Comme $X$ est localement
noethérien, $U$ est un schéma affine noethérien. En vertu de (1),
on voit que $Y' = U \times_X Y \to Y$ est un morphisme \'etale surjectif.
Notons $E \subset |U|$ l'ensemble des points d'image $x$.
Il n'existe aucune spécialisation non triviale entre les éléments de $E$, voir
le Lemme \ref{decent-spaces-lemma-no-specializations-map-to-same-point-Noetherian}.
Le morphisme $Y' \to U$ envoie $|Y'|$ dans $E$. Par notre construction
de $Z_x$ dans la preuve du Lemme \ref{decent-spaces-lemma-find-singleton-from-point},
nous savons que $\coprod_{u \in E} u \to X$ se factorise par $Z_x$.
Il suffit donc de démontrer que $Y' \to U$ se factorise par
$\coprod_{u \in E} u \to U$. Après avoir remplacé $Y'$ par un recouvrement
\'etale par un schéma (ce que permet notre remarque préliminaire), cela résulte
de Morphismes, lemme \ref{morphisms-lemma-factor-through-set-of-points}.
\end{proof}

\begin{lemma}
\label{decent-spaces-lemma-factor-through-residual-space-decent}
Soit $S$ un schéma. Soit $f : Y \to X$ un morphisme d'espaces
algébriques sur $S$. Soit $x \in |X|$ un point. Supposons que
\begin{enumerate}
\item $|f|(|Y|)$ soit contenu dans $\{x\} \subset |X|$,
\item $Y$ soit réduit, et
\item $x$ puisse être représenté par un monomorphisme quasi-compact
$x : \Spec(k) \to X$, où $k$ est un corps (par exemple si $X$ est décent).
\end{enumerate}
Alors $f$ se factorise par
l'espace résiduel $Z_x = \Spec(k)$ de $X$ en $x$.
\end{lemma}

\begin{proof}
D'après le Lemme \ref{decent-spaces-lemma-residual-space-monomorphism}, on a
$Z_x = \Spec(k)$.

\medskip\noindent
Remarque préliminaire : puisque $\Spec(k) \to X$ est un monomorphisme,
il suffit de trouver
un morphisme \'etale surjectif $Y' \to Y$ tel que $Y' \to X$
se factorise par $Z_x$. Remarquons ici que $Y'$ est également réduit.

\medskip\noindent
Après avoir remplacé $X$ par un voisinage ouvert quasi-compact de $x$,
nous pouvons supposer $X$ quasi-compact. D'après le
Lemme \ref{decent-spaces-lemma-filter-quasi-compact}, $x$
est un point de $T \subset U \subset X$, où $T \to U$ (resp.\ $U \to X$)
est une immersion fermée (resp.\ ouverte), et $T$ est un schéma.
D'après Propriétés des espaces, lemme
\ref{spaces-properties-lemma-factor-through-open-subspace},
$f$ se factorise par $U$ ; nous pouvons donc supposer $U = X$. Alors $f$
se factorise par $T$ puisque $Y$ est réduit, voir
Propriétés des espaces, lemme \ref{spaces-properties-lemma-map-into-reduction}.
Nous pouvons donc supposer que $X = T$ est un schéma.
D'après notre remarque préliminaire, nous pouvons aussi supposer que $Y$
est un schéma. Nous sommes ainsi ramenés à Morphismes, lemme
\ref{morphisms-lemma-factor-through-point}.
\end{proof}

\begin{example}
\label{decent-spaces-example-does-not-factor-through-residual-gerbe}
Voici un contre-exemple aux
Lemmes \ref{decent-spaces-lemma-factor-through-residual-space-Noetherian} et
\ref{decent-spaces-lemma-factor-through-residual-space-decent}
lorsque $X$ n'est ni localement noethérien ni décent.
Soit $k$ un corps. Soit $G$ un groupe profini infini.
Soit $Y$ le groupe $G$ vu comme un $k$-schéma en groupes affine de
dimension zéro, c'est-à-dire
$Y = \Spec(\text{applications localement constantes } G \to k)$.
Soit $\Gamma$ le groupe $G$ vu comme un $k$-schéma en groupes discret,
opérant sur $Y$ par translations. Posons $X = Y/\Gamma$.
C'est un espace algébrique réduit à un point, muni de la projection
$q : Y \to X$.
Soit $e \in G$ l'élément neutre (n'importe quel élément conviendrait), et
regardons-le comme un $k$-point de $Y$. On obtient un $k$-point
$x : \Spec(k) \to X$ qui est un monomorphisme, puisqu'il est une section de
$X \to \Spec(k)$.
Nous affirmons que (bien que $Y$ soit affine et réduit et que
$|X| = \{x\}$) le morphisme $q$ ne se factorise par aucun morphisme
$\Spec(K) \to X$, où $K$ est un corps. Sinon, il se factoriserait par
$x$ d'après Propriétés des espaces, lemme
\ref{spaces-properties-lemma-equivalence-class-point-monomorphism}.
Or, le changement de base de $q$ par $x$ est $\Gamma \to \Spec(k)$, la
projection $\Gamma \to Y$ étant l'application d'orbite
$g \mapsto g \cdot e$. Cette dernière n'a pas de section, d'où l'assertion.
\end{example}

\begin{lemma}
\label{decent-spaces-lemma-residual-space-closed}
Soit $S$ un schéma. Soit $X$ un espace algébrique sur $S$. Soit
$x \in |X|$, et soit $Z_x \subset X$ son espace résiduel. Supposons que $X$
soit localement noethérien. Alors $x$ est un point fermé de $|X|$ si et
seulement si le morphisme $Z_x \to X$ est une immersion fermée.
\end{lemma}

\begin{proof}
Si $Z_x \to X$ est une immersion fermée, alors $x$ est un point fermé de
$|X|$, voir Morphismes d'espaces, lemme
\ref{spaces-morphisms-lemma-immersion-when-closed}.
Réciproquement, supposons que $x$ soit un point fermé de $|X|$.
Soit $Z \subset X$ le sous-espace fermé réduit tel que $|Z| = \{x\}$
(Propriétés des espaces,
lemme \ref{spaces-properties-lemma-reduced-closed-subspace}).
Alors $Z$ est localement noethérien d'après Morphismes d'espaces, lemmes
\ref{spaces-morphisms-lemma-immersion-locally-finite-type} et
\ref{spaces-morphisms-lemma-locally-finite-type-locally-noetherian}.
Puisque $Z$ est également réduit et que $|Z| = \{x\}$, $Z = Z_x$ est
l'espace résiduel par définition.
\end{proof}











\section{Espaces décents}
\label{decent-spaces-section-decent}

\noindent
Dans cette section, nous rassemblons quelques faits utiles sur les espaces
décents.

\begin{lemma}
\label{decent-spaces-lemma-locally-Noetherian-decent-quasi-separated}
Tout espace algébrique décent localement noethérien est quasi-séparé.
\end{lemma}

\begin{proof}
En effet, soit $X$ un espace algébrique (sur un schéma de base quelconque,
par exemple sur $\mathbf{Z}$) qui est décent et localement noethérien.
Soient $U \to X$ et $V \to X$ des morphismes étales, où $U$ et $V$ sont
des schémas affines. Nous devons montrer que $W = U \times_X V$ est
quasi-compact (Propriétés des espaces, lemme
\ref{spaces-properties-lemma-characterize-quasi-separated}).
Puisque $X$ est localement noethérien, les schémas $U$, $V$ sont noethériens
et $W$ est localement noethérien. Puisque $X$ est décent, les fibres du
morphisme $W \to U$ sont finies. En effet, on peut représenter tout
$x \in |X|$ par un monomorphisme quasi-compact $\Spec(k) \to X$.
Alors $U_k$ et $V_k$ sont des réunions disjointes finies de spectres
d'extensions finies séparables de $k$ (Remarque \ref{decent-spaces-remark-recall}), et
nous voyons que $W_k = U_k \times_{\Spec(k)} V_k$ est fini.
Soit $n$ le degré maximal d'une fibre de $W \to U$ en un point générique
d'une composante irréductible de $U$. Considérons la stratification
$$
U = U_0 \supset U_1 \supset U_2 \supset \ldots
$$
associée à $W \to U$ dans
Compléments sur les morphismes, lemme \ref{more-morphisms-lemma-stratify-flat-fp-lqf}.
Par le choix de $n$ ci-dessus, nous concluons que $U_{n + 1}$ est vide.
Nous voyons donc que les fibres de $W \to U$ sont universellement bornées.
Nous pouvons alors appliquer Compléments sur les morphismes, lemme
\ref{more-morphisms-lemma-stratify-flat-fp-lqf-universally-bounded}
pour trouver une stratification
$$
\emptyset = Z_{-1} \subset Z_0 \subset Z_1 \subset Z_2 \subset
\ldots \subset Z_n = U
$$
par des parties fermées telle que, si $S_r = Z_r \setminus Z_{r - 1}$,
le morphisme $W \times_U S_r \to S_r$ soit fini localement libre.
Puisque $U$ est noethérien, les schémas $S_r$ sont noethériens, donc les
schémas $W \times_U S_r$ sont noethériens, et par conséquent
$W = \coprod W \times_U S_r$ est quasi-compact, comme voulu.
\end{proof}

\begin{lemma}
\label{decent-spaces-lemma-when-field}
Soit $S$ un schéma. Soit $X$ un espace algébrique décent sur $S$.
\begin{enumerate}
\item Si $|X|$ est réduit à un point, alors $X$ est un schéma.
\item Si $|X|$ est réduit à un point et que $X$ est réduit, alors
$X \cong \Spec(k)$ pour un corps $k$.
\end{enumerate}
\end{lemma}

\begin{proof}
Supposons que $|X|$ soit réduit à un point. Il résulte immédiatement du
Théorème \ref{decent-spaces-theorem-decent-open-dense-scheme} que $X$ est un schéma,
mais nous pouvons aussi raisonner directement comme suit.
Choisissons un schéma affine $U$ et un morphisme étale surjectif $U \to X$.
Posons $R = U \times_X U$. Alors $U$ et $R$ ont un nombre fini de points
d'après le Lemme \ref{decent-spaces-lemma-UR-finite-above-x} (et la définition d'un espace
décent). Tous ces points sont fermés dans $U$ et $R$ d'après le
Lemme \ref{decent-spaces-lemma-decent-no-specializations-map-to-same-point}.
Il s'ensuit que $U$ et $R$ sont des schémas affines.
Nous pouvons restreindre $U$ à un espace réduit à un point. Alors $U$ est
le spectre d'un anneau local hensélien, voir
Algèbre, lemme \ref{algebra-lemma-local-dimension-zero-henselian}.
Les projections $R \to U$ sont étales, donc finies et étales puisque
$U$ est le spectre d'un anneau local hensélien de dimension $0$, voir
Algèbre, lemme \ref{algebra-lemma-characterize-henselian}.
Il s'ensuit que $X$ est un schéma d'après
Groupoïdes, Proposition \ref{groupoids-proposition-finite-flat-equivalence}.

\medskip\noindent
L'assertion (2) découle de (1) et du fait qu'un schéma réduit dont l'espace
sous-jacent est réduit à un point est le spectre d'un corps.
\end{proof}

\begin{remark}
\label{decent-spaces-remark-one-point-decent-scheme}
Nous verrons dans
Limites d'espaces, lemme \ref{spaces-limits-lemma-reduction-scheme}
qu'un espace algébrique dont la réduction est un schéma est un schéma.
\end{remark}

\begin{lemma}
\label{decent-spaces-lemma-algebraic-residue-field-extension-closed-point}
Soit $S$ un schéma. Soit $X$ un espace algébrique décent sur $S$.
Considérons un diagramme commutatif
$$
\xymatrix{
\Spec(k) \ar[rr] \ar[rd] & & X \ar[ld] \\
& S
}
$$
Supposons que le point image $s \in S$ de $\Spec(k) \to S$ soit un point
fermé et que $\kappa(s) \subset k$ soit une extension algébrique.
Alors le point image $x$ de $\Spec(k) \to X$ est un point fermé de $|X|$.
\end{lemma}

\begin{proof}
Supposons que $x \leadsto x'$ pour un $x' \in |X|$. Choisissons un
morphisme étale $U \to X$, où $U$ est un schéma, et un point $u' \in U$
qui s'envoie sur $x'$. Choisissons une spécialisation $u \leadsto u'$ dans
$U$, où $u$ s'envoie sur $x$ dans $X$, voir le
Lemme \ref{decent-spaces-lemma-decent-specialization}.
Alors $u$ est l'image d'un point $w$ du schéma
$W = \Spec(k) \times_X U$. Puisque la projection $W \to \Spec(k)$ est étale,
nous voyons que $\kappa(w) \supset k$ est une extension finie. Ainsi
$\kappa(w) \supset \kappa(s)$ est algébrique. Ainsi
$\kappa(u) \supset \kappa(s)$ est algébrique. Donc $u$ est un point fermé de
$U$ d'après Morphismes, lemme
\ref{morphisms-lemma-algebraic-residue-field-extension-closed-point-fibre}.
Ainsi $u = u'$, d'où $x = x'$.
\end{proof}

\begin{lemma}
\label{decent-spaces-lemma-finite-residue-field-extension-finite}
Soit $S$ un schéma. Soit $X$ un espace algébrique décent sur $S$.
Considérons un diagramme commutatif
$$
\xymatrix{
\Spec(k) \ar[rr] \ar[rd] & & X \ar[ld] \\
& S
}
$$
Supposons que le point image $s \in S$ de $\Spec(k) \to S$ soit un point
fermé et que l'extension de corps $k/\kappa(s)$ soit finie.
Alors $\Spec(k) \to X$ est un morphisme fini. Si $\kappa(s) = k$,
alors $\Spec(k) \to X$ est une immersion fermée.
\end{lemma}

\begin{proof}
D'après le Lemme \ref{decent-spaces-lemma-algebraic-residue-field-extension-closed-point},
le point image $x \in |X|$ est fermé. Soit $Z \subset X$ le sous-espace
fermé réduit tel que $|Z| = \{x\}$ (Propriétés des espaces,
lemme \ref{spaces-properties-lemma-reduced-closed-subspace}).
Remarquons que $Z$ est un espace algébrique décent d'après le
Lemme \ref{decent-spaces-lemma-representable-named-properties}.
D'après le Lemme \ref{decent-spaces-lemma-when-field}, nous voyons que
$Z = \Spec(k')$ pour un corps $k'$. Bien sûr,
$k \supset k' \supset \kappa(s)$. Alors $\Spec(k) \to Z$ est un morphisme
fini de schémas, et $Z \to X$ est un morphisme fini puisqu'il s'agit d'une
immersion fermée. Par conséquent, $\Spec(k) \to X$ est fini (Morphismes
d'espaces, lemme \ref{spaces-morphisms-lemma-composition-integral}).
Si $k = \kappa(s)$, alors $\Spec(k) = Z$ et $\Spec(k) \to X$ est une
immersion fermée.
\end{proof}

\begin{lemma}
\label{decent-spaces-lemma-decent-space-closed-point}
Soit $S$ un schéma. Supposons que $X$ soit un espace algébrique décent sur
$S$. Soit $x \in |X|$ un point fermé. Alors $x$ peut être représenté par une
immersion fermée $i : \Spec(k) \to X$ provenant du spectre d'un corps.
\end{lemma}

\begin{proof}
Nous savons que $x$ peut être représenté par un monomorphisme quasi-compact
$i : \Spec(k) \to X$, où $k$ est un corps
(Définition \ref{decent-spaces-definition-very-reasonable}).
Soit $U \to X$ un morphisme étale, où $U$ est un schéma affine.
Comme $x$ est fermé et que $X$ est décent, la fibre $F$ de
$|U| \to |X|$ au-dessus de $x$ est formée de points fermés
(Lemme \ref{decent-spaces-lemma-decent-no-specializations-map-to-same-point}).
Puisque $i$ est un monomorphisme, il en va de même de
$U_k = U \times_X \Spec(k) \to U$. En particulier, l'application
$|U_k| \to F$ est injective. Comme $U_k$ est quasi-compact et étale sur un
corps, nous voyons que $U_k$ est une réunion disjointe finie de spectres de
corps (Remarque \ref{decent-spaces-remark-recall}). Écrivons
$U_k = \Spec(k_1) \amalg \ldots \amalg \Spec(k_r)$.
Puisque $\Spec(k_i) \to U$ est un monomorphisme, nous voyons que son image
$u_i$ a pour corps résiduel $\kappa(u_i) = k_i$.
Comme $u_i \in F$ est un point fermé, nous concluons que le morphisme
$\Spec(k_i) \to U$ est une immersion fermée. Puisque les $u_i$ sont deux à
deux distincts, $U_k \to U$ est une immersion fermée. Par conséquent, $i$
est une immersion fermée (Morphismes d'espaces, lemme
\ref{spaces-morphisms-lemma-closed-immersion-local}). Cela achève la preuve.
\end{proof}


\section{Espaces localement séparés}
\label{decent-spaces-section-locally-separated}

\noindent
Il se trouve qu'un espace algébrique localement séparé est décent.

\begin{lemma}
\label{decent-spaces-lemma-infinite-number}
Soit $A$ un anneau. Soit $k$ un corps. Soit $\mathfrak p_n$, $n \geq 1$,
une suite d'idéaux premiers deux à deux distincts de $A$. De plus, pour chaque
$n$, soit $k \to \kappa(\mathfrak p_n)$ un plongement. Alors l'adhérence
de l'image de
$$
\coprod\nolimits_{n \not = m}
\Spec(\kappa(\mathfrak p_n) \otimes_k \kappa(\mathfrak p_m))
\longrightarrow
\Spec(A \otimes A)
$$
rencontre la diagonale.
\end{lemma}

\begin{proof}
Posons $k_n = \kappa(\mathfrak p_n)$. Nous pouvons supposer que
$A = \prod k_n$. Notons $x_n = \Spec(k_n)$ le point ouvert et fermé
correspondant à $A \to k_n$. Alors $\Spec(A) = Z \amalg \{x_n\}$, où $Z$
est une partie fermée non vide. En effet, $Z = V(e_n; n \geq 1)$, où $e_n$
est l'idempotent de $A$ correspondant au facteur $k_n$, et $Z$ est non vide
car l'idéal engendré par les $e_n$ n'est pas égal à $A$. Nous allons montrer
que l'adhérence de l'image contient $\Delta(Z)$. Le noyau de l'application
$$
(\prod k_n) \otimes_k (\prod k_m)
\longrightarrow
\prod\nolimits_{n \not = m} k_n \otimes_k k_m
$$
est l'idéal engendré par les $e_n \otimes e_n$, $n \geq 1$.
Par conséquent, l'adhérence de l'image du morphisme induit sur les spectres
est $V(e_n \otimes e_n; n \geq 1)$, dont l'intersection avec
$\Delta(\Spec(A))$ est $\Delta(Z)$. Il suffit donc de montrer que
$$
\coprod\nolimits_{n \not = m} \Spec(k_n \otimes_k k_m)
\longrightarrow
\Spec(\prod\nolimits_{n \not = m} k_n \otimes_k k_m)
$$
a une image dense. Cela résulte du fait que la famille d'homomorphismes
d'anneaux
$\prod_{n \not = m} k_n \otimes_k k_m \to k_n \otimes_k k_m$
est conjointement injective.
\end{proof}

\begin{lemma}[David Rydh]
\label{decent-spaces-lemma-locally-separated-decent}
Un espace algébrique localement séparé est décent.
\end{lemma}

\begin{proof}
Soit $S$ un schéma et soit $X$ un espace algébrique localement séparé sur
$S$. Nous pouvons supposer que $S = \Spec(\mathbf{Z})$, voir
Propriétés des espaces, Définition
\ref{spaces-properties-definition-separated}.
Les produits fibrés sans indication seront pris au-dessus de $\mathbf{Z}$.
Soit $x \in |X|$. Choisissons un schéma $U$, un morphisme étale
$U \to X$ et un point $u \in U$ qui s'envoie sur $x$ dans $|X|$.
Comme d'habitude, nous identifions $u = \Spec(\kappa(u))$.
Puisque $X$ est localement séparé, le morphisme
$$
u \times_X u \to u \times u
$$
est une immersion (Morphismes d'espaces, lemme
\ref{spaces-morphisms-lemma-fibre-product-after-map}).
Par conséquent, Compléments sur les groupoïdes, lemme
\ref{more-groupoids-lemma-locally-closed-image-is-closed}
nous dit que c'est une immersion fermée (en utilisant
Schémas, lemme \ref{schemes-lemma-immersion-when-closed}).
Puisque $u \times_X u \to u \times_X U$ est un monomorphisme (changement de
base de $u \to U$) et que $u \times_X U \to u$ est étale, nous concluons que
$u \times_X u$ est une réunion disjointe de spectres de corps
(voir la Remarque \ref{decent-spaces-remark-recall} et
Schémas, lemme \ref{schemes-lemma-mono-towards-spec-field}).
Puisqu'elle est aussi fermée dans le schéma affine $u \times u$, nous
concluons que $u \times_X u$ est une réunion disjointe finie de spectres de
corps. Ainsi, $x$ peut être représenté par un monomorphisme
$\Spec(k) \to X$, où $k$ est un corps, voir le
Lemme \ref{decent-spaces-lemma-R-finite-above-x}.

\medskip\noindent
Ensuite, soit $U = \Spec(A)$ un schéma affine et soit $U \to X$ un morphisme
étale. Pour terminer la preuve, il suffit de montrer que
$F = U \times_X \Spec(k)$ est fini. Écrivons
$F = \coprod_{i \in I} \Spec(k_i)$ comme la réunion disjointe de spectres
d'extensions finies séparables de $k$. Nous devons montrer que $I$ est fini.
Posons $R = U \times_X U$. Puisque $X$ est localement séparé, le morphisme
$j : R \to U \times U$ est une immersion. Soit $U' \subset U \times U$ un
ouvert tel que $j$ se factorise par une immersion fermée $j' : R \to U'$.
Soit $e : U \to R$ l'application diagonale. Comme $e$ est un morphisme entre
schémas étales sur $U$ tel que $\Delta = j \circ e$ soit une immersion fermée,
nous concluons que $R = e(U) \amalg W$ pour un sous-schéma ouvert et fermé
$W \subset R$. Puisque $j'$ est une immersion fermée, nous concluons que
$j'(W) \subset U'$ est fermé et disjoint de $j'(e(U))$. Par conséquent,
$\overline{j(W)} \cap \Delta(U) = \emptyset$ dans $U \times U$.
Remarquons que $W$ contient $\Spec(k_i \otimes_k k_{i'})$ pour tous
$i \not = i'$, $i, i' \in I$. D'après le Lemme
\ref{decent-spaces-lemma-infinite-number}, nous concluons que $I$ est fini, comme voulu.
\end{proof}








\section{Critère valuatif}
\label{decent-spaces-section-valuative-criterion-universally-closed}

\noindent
Pour un morphisme quasi-compact de source un espace décent, le critère
valuatif est une condition nécessaire pour que le morphisme soit
universellement fermé.


\begin{proposition}
\label{decent-spaces-proposition-characterize-universally-closed}
Soit $S$ un schéma. Soit $f : X \to Y$ un morphisme d'espaces algébriques
sur $S$. Supposons que $f$ soit quasi-compact et que $X$ soit décent.
Alors $f$ est
universellement fermé si et seulement si la partie d'existence du critère
valuatif est satisfaite.
\end{proposition}

\begin{proof}
Dans
Morphismes d'espaces,
lemme \ref{spaces-morphisms-lemma-quasi-compact-existence-universally-closed},
nous avons vu l'une des implications.
Pour démontrer l'autre, supposons $f$ universellement fermé. Soit
$$
\xymatrix{
\Spec(K) \ar[r] \ar[d] & X \ar[d] \\
\Spec(A) \ar[r] & Y
}
$$
un diagramme comme dans
Morphismes d'espaces,
définition \ref{spaces-morphisms-definition-valuative-criterion}.
Posons $X_A = \Spec(A) \times_Y X$ ; on a alors
$$
\xymatrix{
\Spec(K) \ar[r] \ar[rd] & X_A \ar[d] \\
 & \Spec(A)
}
$$
D'après
Morphismes d'espaces,
lemme \ref{spaces-morphisms-lemma-base-change-quasi-compact},
$X_A \to \Spec(A)$ est quasi-compact. Comme $X_A \to X$
est représentable, $X_A$ est lui aussi décent ; voir le
Lemme \ref{decent-spaces-lemma-representable-properties}.
De plus, puisque $f$ est universellement fermé, $X_A \to \Spec(A)$
est universellement fermé.
On peut donc remplacer, et l'on remplace, $X$ par $X_A$ et $Y$ par $\Spec(A)$.

\medskip\noindent
Soit $x' \in |X|$ la classe d'équivalence du morphisme
$\Spec(K) \to X$. Soit $y \in |Y| = |\Spec(A)|$
le point fermé. Posons $y' = f(x')$ ; c'est le point générique de
$\Spec(A)$. Comme $f$ est universellement fermé,
$f(\overline{\{x'\}})$ contient $\overline{\{y'\}}$, donc
contient $y$. Soit $x \in \overline{\{x'\}}$ un point tel que
$f(x) = y$. Soit $U$ un schéma, et soit $\varphi : U \to X$
un morphisme étale tel qu'il existe un point $u \in U$ avec
$\varphi(u) = x$. D'après le
Lemme \ref{decent-spaces-lemma-specialization}
et l'hypothèse selon laquelle $X$ est décent,
il existe une spécialisation $u' \leadsto u$ dans $U$ telle que
$\varphi(u') = x'$.
Cela signifie qu'il existe une extension de corps commune
$K \subset K' \supset \kappa(u')$ telle que le diagramme
$$
\xymatrix{
\Spec(K') \ar[r] \ar[d] & U \ar[d] \\
\Spec(K) \ar[r] \ar[rd] & X \ar[d] \\
 & \Spec(A)
}
$$
soit commutatif. On obtient le diagramme commutatif d'anneaux suivant :
$$
\xymatrix{
K' & \mathcal{O}_{U, u} \ar[l] \\
K \ar[u] & \\
 & A \ar[lu] \ar[uu]
}
$$
D'après
Algèbre, lemme \ref{algebra-lemma-dominate},
on peut trouver un anneau de valuation $A' \subset K'$ qui domine l'image de
$\mathcal{O}_{U, u}$ dans $K'$. Comme, par construction, $\mathcal{O}_{U, u}$
domine $A$, on voit que $A'$ domine également $A$. On obtient donc un
diagramme
semblable au deuxième diagramme de
Morphismes d'espaces,
définition \ref{spaces-morphisms-definition-valuative-criterion},
ce qui démontre la proposition.
\end{proof}







\section{Conditions relatives}
\label{decent-spaces-section-relative-conditions}

\noindent
Il s'agit d'une section technique (encore une) consacrée à des conditions
sur les espaces algébriques qui concernent les points. Il est probablement
préférable de sauter cette section en première lecture.

\begin{definition}
\label{decent-spaces-definition-relative-conditions}
Soit $S$ un schéma. On dit qu'un espace algébrique $X$ sur $S$
{\it possède la propriété $(\beta)$} si $X$ possède la propriété
correspondante
du Lemme \ref{decent-spaces-lemma-bounded-fibres}.
Soit $f : X \to Y$ un morphisme d'espaces algébriques sur $S$.
\begin{enumerate}
\item On dit que $f$ {\it possède la propriété $(\beta)$} si, pour tout schéma
$T$ et tout morphisme $T \to Y$, le produit fibré $T \times_Y X$ possède la
propriété $(\beta)$.
\item On dit que $f$ est {\it décent} si, pour tout schéma $T$ et tout
morphisme $T \to Y$, le produit fibré $T \times_Y X$ est un espace
algébrique décent.
\item On dit que $f$ est {\it raisonnable} si, pour tout schéma $T$ et tout
morphisme $T \to Y$, le produit fibré $T \times_Y X$ est un espace
algébrique raisonnable.
\item On dit que $f$ est {\it très raisonnable} si, pour tout schéma $T$ et
tout morphisme $T \to Y$, le produit fibré $T \times_Y X$ est un espace
algébrique très raisonnable.
\end{enumerate}
\end{definition}

\noindent
Nous renvoyons à la Remarque \ref{decent-spaces-remark-very-reasonable} pour une discussion
informelle. On verra que la classe des morphismes très raisonnables n'est pas
très utile, mais que celles des morphismes décents et raisonnables le sont.

\begin{lemma}
\label{decent-spaces-lemma-properties-trivial-implications}
Soit $S$ un schéma.
Soit $f : X \to Y$ un morphisme d'espaces algébriques sur $S$.
On a les implications suivantes entre les conditions imposées à $f$ :
$$
\xymatrix{
\text{représentable} \ar@{=>}[rd] & & & & \\
& \text{très raisonnable} \ar@{=>}[r] & \text{raisonnable} \ar@{=>}[r] &
\text{décent} \ar@{=>}[r] & (\beta) \\
\text{quasi-séparé} \ar@{=>}[ru] & & & &
}
$$
\end{lemma}

\begin{proof}
Cela résulte immédiatement des définitions, du
Lemme \ref{decent-spaces-lemma-bounded-fibres}
et de
Morphismes d'espaces,
Lemme \ref{spaces-morphisms-lemma-separated-local}.
\end{proof}

\noindent
Voici une autre vérification élémentaire.

\begin{lemma}
\label{decent-spaces-lemma-property-for-morphism-out-of-property}
Soit $S$ un schéma. Soit $f : X \to Y$ un morphisme d'espaces
algébriques sur $S$. Si $X$ est décent (resp.\ raisonnable, resp.\ possède la
propriété $(\beta)$ du Lemme \ref{decent-spaces-lemma-bounded-fibres}), alors $f$ est décent
(resp.\ raisonnable, resp.\ possède la propriété $(\beta)$).
\end{lemma}

\begin{proof}
Soit $T$ un schéma et soit $T \to Y$ un morphisme. Alors $T \to Y$ est
représentable, donc le changement de base $T \times_Y X \to X$ est
représentable. Par conséquent, si $X$ est décent (ou raisonnable), il en est
de
même de $T \times_Y X$ ; voir le
Lemme \ref{decent-spaces-lemma-representable-named-properties}.
De même, pour la propriété $(\beta)$, voir le
Lemme \ref{decent-spaces-lemma-representable-properties}.
\end{proof}

\begin{lemma}
\label{decent-spaces-lemma-base-change-relative-conditions}
Le fait de posséder la propriété $(\beta)$, d'être décent ou d'être
raisonnable est stable par changement de base arbitraire.
\end{lemma}

\begin{proof}
Cela découle immédiatement de la définition.
\end{proof}

\begin{lemma}
\label{decent-spaces-lemma-property-over-property}
Soit $S$ un schéma.
Soit $f : X \to Y$ un morphisme d'espaces algébriques sur $S$.
Soit $\omega \in \{\beta, \text{décent}, \text{raisonnable}\}$.
Supposons que $Y$ possède la propriété $(\omega)$ et que $f : X \to Y$
possède $(\omega)$. Alors $X$ possède $(\omega)$.
\end{lemma}

\begin{proof}
Démontrons le lemme dans le cas où $\omega = \beta$. Il faut alors montrer
que tout $x \in |X|$ est représenté par un monomorphisme du spectre d'un
corps vers $X$. Posons $y = f(x) \in |Y|$. Par hypothèse, il existe un corps
$k$ et un monomorphisme $\Spec(k) \to Y$ représentant $y$.
Alors $x$ correspond à un point $x'$ de $\Spec(k) \times_Y X$.
Par hypothèse, $x'$ est représenté par un monomorphisme
$\Spec(k') \to \Spec(k) \times_Y X$. Il est clair que le composé
$\Spec(k') \to X$ est un monomorphisme représentant $x$.

\medskip\noindent
Démontrons le lemme dans le cas où $\omega = \text{décent}$.
Soient $x \in |X|$ et $y = f(x) \in |Y|$. D'après le paragraphe précédent,
on peut choisir un diagramme
$$
\xymatrix{
\Spec(k') \ar[r]_x \ar[d] & X \ar[d]^f \\
\Spec(k) \ar[r]^y & Y
}
$$
dont les flèches horizontales sont des monomorphismes. Puisque $Y$ est décent,
le morphisme $y$ est quasi-compact. Puisque $f$ est décent, l'espace
algébrique $\Spec(k) \times_Y X$ est décent. Le monomorphisme
$\Spec(k') \to \Spec(k) \times_Y X$ est donc quasi-compact.
Le monomorphisme $x : \Spec(k') \to X$ est alors quasi-compact en tant que
composé de morphismes quasi-compacts (utiliser
Morphismes d'espaces, lemmes
\ref{spaces-morphisms-lemma-base-change-quasi-compact} et
\ref{spaces-morphisms-lemma-composition-quasi-compact}).
Comme le point $x$ était arbitraire, il s'ensuit que $X$ est décent.

\medskip\noindent
Démontrons le lemme dans le cas où $\omega = \text{raisonnable}$.
Choisissons $V \to Y$ étale, avec $V$ un schéma affine.
Choisissons $U \to V \times_Y X$ étale, avec $U$ un schéma affine.
Par hypothèse, les fibres de $V \to Y$ sont universellement bornées.
D'après le
Lemme \ref{decent-spaces-lemma-base-change-universally-bounded},
les fibres du morphisme $V \times_Y X \to X$ sont universellement bornées.
L'hypothèse sur $f$ montre que les fibres de $U \to V \times_Y X$ sont
universellement bornées. D'après le
Lemme \ref{decent-spaces-lemma-composition-universally-bounded},
les fibres du composé $U \to X$ sont universellement bornées.
Il existe donc suffisamment de morphismes étales $U \to X$ provenant de
schémas dont les fibres sont universellement bornées, et l'on conclut que
$X$ est raisonnable.
\end{proof}

\begin{lemma}
\label{decent-spaces-lemma-composition-relative-conditions}
Le fait de posséder la propriété $(\beta)$, d'être décent ou d'être
raisonnable est stable par composition.
\end{lemma}

\begin{proof}
Soit $\omega \in \{\beta, \text{décent}, \text{raisonnable}\}$.
Soient $f : X \to Y$ et $g : Y \to Z$ des morphismes d'espaces algébriques
sur le schéma $S$. Supposons que $f$ et $g$ possèdent tous deux la propriété
$(\omega)$. Il faut alors montrer que, pour tout schéma $T$ et tout morphisme
$T \to Z$, l'espace $T \times_Z X$ possède $(\omega)$. D'après le
Lemme \ref{decent-spaces-lemma-base-change-relative-conditions},
on se ramène à l'assertion suivante : si $Y$ est un espace algébrique qui
possède la propriété $(\omega)$ et si $f : X \to Y$ est un morphisme qui
possède $(\omega)$, alors $X$ possède $(\omega)$.
C'est le contenu du Lemme \ref{decent-spaces-lemma-property-over-property}.
\end{proof}

\begin{lemma}
\label{decent-spaces-lemma-fibre-product-conditions}
Soit $S$ un schéma. Soient $f : X \to Y$ et $g : Z \to Y$ des morphismes
d'espaces algébriques sur $S$. Si $X$ et $Z$ sont décents
(resp.\ raisonnables, resp.\ possèdent la propriété
$(\beta)$ du Lemme \ref{decent-spaces-lemma-bounded-fibres}), il en va de même de
$X \times_Y Z$.
\end{lemma}

\begin{proof}
En effet, d'après le Lemme \ref{decent-spaces-lemma-property-for-morphism-out-of-property},
le morphisme $X \to Y$ possède la propriété considérée. Son changement de
base $X \times_Y Z \to Z$ possède alors cette propriété d'après le
Lemme \ref{decent-spaces-lemma-base-change-relative-conditions}.
Enfin, le Lemme \ref{decent-spaces-lemma-property-over-property} implique que
$X \times_Y Z$ possède cette propriété.
\end{proof}

\begin{lemma}
\label{decent-spaces-lemma-descent-conditions}
Soit $S$ un schéma.
Soit $f : X \to Y$ un morphisme d'espaces algébriques sur $S$.
Soit $\mathcal{P} \in \{(\beta), \text{décent}, \text{raisonnable}\}$.
Supposons que
\begin{enumerate}
\item $f$ soit quasi-compact,
\item $f$ soit étale,
\item $|f| : |X| \to |Y|$ soit surjective, et
\item l'espace algébrique $X$ possède la propriété $\mathcal{P}$.
\end{enumerate}
Alors $Y$ possède la propriété $\mathcal{P}$.
\end{lemma}

\begin{proof}
Démontrons ceci dans le cas où $\mathcal{P} = (\beta)$. Soit $y \in |Y|$
un point. Il faut montrer que $y$ peut être représenté par un monomorphisme
provenant d'un corps. Choisissons un point $x \in |X|$ tel que $f(x) = y$.
Par hypothèse, on peut représenter $x$ par un monomorphisme
$\Spec(k) \to X$, où $k$ est un corps. D'après le
Lemme \ref{decent-spaces-lemma-R-finite-above-x},
il suffit de montrer que les projections
$\Spec(k) \times_Y \Spec(k) \to \Spec(k)$
sont étales et quasi-compactes. On peut factoriser la première projection en
$$
\Spec(k) \times_Y \Spec(k)
\longrightarrow
\Spec(k) \times_Y X
\longrightarrow
\Spec(k)
$$
Le premier morphisme est un monomorphisme, et le second est étale et
quasi-compact. D'après
Propriétés des espaces,
Lemme \ref{spaces-properties-lemma-etale-over-field-scheme},
$\Spec(k) \times_Y X$ est un schéma. C'est donc une réunion disjointe finie
de spectres d'extensions finies séparables de $k$. D'après
Schémas, lemme \ref{schemes-lemma-mono-towards-spec-field},
la première flèche identifie $\Spec(k) \times_Y \Spec(k)$ à une réunion
disjointe finie de spectres d'extensions finies séparables de $k$.
Le morphisme de projection est donc étale et quasi-compact.

\medskip\noindent
Démontrons ceci dans le cas où $\mathcal{P} = \text{décent}$.
Le premier paragraphe de la preuve montre déjà que tout $y \in |Y|$ peut être
représenté par un monomorphisme $y : \Spec(k) \to Y$. Fixons un tel $y$.
Choisissons un schéma affine $U$ et un morphisme étale $U \to X$ tels que
l'image de $|U| \to |Y|$ contienne $y$. D'après le
Lemme \ref{decent-spaces-lemma-UR-finite-above-x},
il suffit de montrer que $U_y$ est un schéma fini sur $k$. Le produit fibré
$X_y = \Spec(k) \times_Y X$ est un espace algébrique quasi-compact et étale
sur $k$. D'après
Propriétés des espaces,
Lemme \ref{spaces-properties-lemma-etale-over-field-scheme},
c'est donc un schéma. C'est une réunion disjointe finie de spectres
d'extensions finies séparables de $k$. Écrivons
$X_y = \{x_1, \ldots, x_n\}$, où $x_i$ est donné par
$x_i : \Spec(k_i) \to X$ avec $[k_i : k] < \infty$. Par hypothèse, $X$ est
décent, donc les schémas $U_{x_i} = \Spec(k_i) \times_X U$ sont finis sur
$k_i$. Enfin, $U_y = \coprod U_{x_i}$ comme schéma, d'où l'on conclut que
$U_y$ est fini sur $k$, comme voulu.

\medskip\noindent
Démontrons ceci dans le cas où $\mathcal{P} = \text{raisonnable}$.
Choisissons un schéma affine $V$ et un morphisme étale $V \to Y$.
Il faut montrer que les fibres de $V \to Y$ sont universellement bornées.
L'espace algébrique $V \times_Y X$ est quasi-compact.
On peut donc trouver un schéma affine $W$ et un morphisme étale surjectif
$W \to V \times_Y X$ ; voir
Propriétés des espaces,
Lemme \ref{spaces-properties-lemma-quasi-compact-affine-cover}.
Voici un diagramme, dont les flèches pleines sont données :
$$
\xymatrix{
W \ar[r]  \ar[rd] &
V \times_Y X \ar[r] \ar[d] &
X \ar[d]_f & \Spec(k) \ar@{..>}[l]^x \ar@{..>}[ld]^y \\
 & V \ar[r] & Y
}
$$
Les fibres du morphisme $W \to X$ sont universellement bornées, puisque
l'espace $X$ est raisonnable. Soit $n$ un entier qui majore les degrés des
fibres de $W \to X$. Nous affirmons que le même entier majore les fibres de
$V \to Y$. En effet, soit $y \in |Y|$ un point. Il existe alors un point
$x \in |X|$ tel que $f(x) = y$ (voir ci-dessus).
On peut donc trouver un corps $k$ et des morphismes $x, y$ donnés par les
flèches pointillées du diagramme ci-dessus. En particulier, on obtient un
morphisme étale surjectif
$$
\Spec(k) \times_{x, X} W
\to
\Spec(k) \times_{x, X} (V \times_Y X) = \Spec(k) \times_{y, Y} V
$$
ce qui montre que le degré de $\Spec(k) \times_{y, Y} V$ sur $k$ est
inférieur ou égal au degré de $\Spec(k) \times_{x, X} W$ sur $k$, donc qu'il
est $\leq n$, et cela donne le résultat. (Cette dernière partie de
l'argument est la même que dans la preuve du
Lemme \ref{decent-spaces-lemma-descent-universally-bounded}.
Malheureusement, ce lemme n'est pas assez général, car il ne s'applique
qu'aux morphismes représentables.)
\end{proof}

\begin{lemma}
\label{decent-spaces-lemma-relative-conditions-local}
Soit $S$ un schéma.
Soit $f : X \to Y$ un morphisme d'espaces algébriques sur $S$.
Soit $\mathcal{P} \in \{(\beta), \text{décent}, \text{raisonnable}, \text{très raisonnable}\}$.
Les conditions suivantes sont équivalentes :
\begin{enumerate}
\item $f$ possède $\mathcal{P}$,
\item pour tout schéma affine $Z$ et tout morphisme $Z \to Y$, le changement
de base $Z \times_Y X \to Z$ de $f$ possède $\mathcal{P}$,
\item pour tout schéma affine $Z$ et tout morphisme $Z \to Y$, l'espace
algébrique $Z \times_Y X$ possède $\mathcal{P}$, et
\item il existe un recouvrement de Zariski $Y = \bigcup Y_i$ tel que chaque
morphisme $f^{-1}(Y_i) \to Y_i$ possède $\mathcal{P}$.
\end{enumerate}
Si $\mathcal{P} \in \{(\beta), \text{décent}, \text{raisonnable}\}$, ces conditions sont
aussi
équivalentes à la suivante :
\begin{enumerate}
\item[(5)] il existe un schéma $V$ et un morphisme étale surjectif
$V \to Y$ tels que le changement de base $V \times_Y X \to V$ possède
$\mathcal{P}$.
\end{enumerate}
\end{lemma}

\begin{proof}
Les implications (1) $\Rightarrow$ (2) $\Rightarrow$ (3) $\Rightarrow$ (4)
sont immédiates.
Pour voir que (3) $\Rightarrow$ (1), procédons comme suit.
Soit $Z \to Y$ un morphisme dont la source est un schéma sur $S$.
Considérons l'espace algébrique $Z \times_Y X$. Si l'on suppose (3), alors,
pour tout ouvert affine $W \subset Z$, le sous-espace ouvert
$W \times_Y X$ de $Z \times_Y X$ possède la propriété $\mathcal{P}$.
D'après le
Lemme \ref{decent-spaces-lemma-properties-local},
l'espace $Z \times_Y X$ possède donc la propriété $\mathcal{P}$, ce qui donne
(1). Un argument analogue, que nous omettons, montre que (4) implique (1).

\medskip\noindent
L'implication (1) $\Rightarrow$ (5) est immédiate. Soit $V \to Y$ un morphisme
étale provenant d'un schéma comme dans (5). Soit $Z$ un schéma affine et soit
$Z \to Y$ un morphisme. Considérons le diagramme
$$
\xymatrix{
Z \times_Y V \ar[r]_q \ar[d]_p & V \ar[d] \\
Z \ar[r] & Y
}
$$
Puisque $p$ est étale, donc ouvert, on peut choisir un nombre fini de
sous-schémas ouverts affines $W_i \subset Z \times_Y V$ tels que
$Z = \bigcup p(W_i)$. Considérons le diagramme commutatif
$$
\xymatrix{
V \times_Y X \ar[d] &
(\coprod W_i) \times_Y X \ar[l] \ar[d] \ar[r] &
Z \times_Y X \ar[d] \\
V &
\coprod W_i \ar[l] \ar[r] &
Z
}
$$
On sait que $V \times_Y X$ possède la propriété $\mathcal{P}$. D'après le
Lemme \ref{decent-spaces-lemma-representable-properties},
$(\coprod W_i) \times_Y X$ possède la propriété $\mathcal{P}$.
Remarquons que le morphisme
$(\coprod W_i) \times_Y X \to Z \times_Y X$ est étale et quasi-compact,
puisqu'il est le changement de base de $\coprod W_i \to Z$.
Le Lemme \ref{decent-spaces-lemma-descent-conditions}
montre donc que $Z \times_Y X$ possède la propriété $\mathcal{P}$.
\end{proof}

\begin{remark}
\label{decent-spaces-remark-very-reasonable}
Une description informelle des propriétés $(\beta)$, décent, raisonnable et
très raisonnable a été donnée dans la section \ref{decent-spaces-section-reasonable-decent}.
Un morphisme possède l'une de ces propriétés si, de manière très
approximative,
ses fibres possèdent les propriétés correspondantes.
La propriété d'être décent est utile pour démontrer des résultats sur les
spécialisations de points de $|X|$. La propriété d'être raisonnable est un peu
plus forte et se prête assez facilement aux raisonnements techniques.
\end{remark}

\noindent
Voici un lemme promis plus haut qui utilise les morphismes décents.

\begin{lemma}
\label{decent-spaces-lemma-re-characterize-universally-closed}
Soit $S$ un schéma.
Soit $f : X \to Y$ un morphisme d'espaces algébriques sur $S$.
Supposons $f$ quasi-compact et décent.
(C'est par exemple le cas si $f$ est représentable ou quasi-séparé ; voir le
Lemme \ref{decent-spaces-lemma-properties-trivial-implications}.)
Alors $f$ est universellement fermé si et seulement si la partie d'existence
du critère valuatif est satisfaite.
\end{lemma}

\begin{proof}
Dans
Morphismes d'espaces,
Lemme \ref{spaces-morphisms-lemma-quasi-compact-existence-universally-closed},
nous avons démontré que tout morphisme quasi-compact qui satisfait la partie
d'existence du critère valuatif est universellement fermé.
Pour démontrer la réciproque, supposons que $f$ soit universellement fermé.
Dans la preuve de la
Proposition \ref{decent-spaces-proposition-characterize-universally-closed},
nous avons vu qu'il suffit de montrer, pour tout anneau de valuation $A$ et
tout morphisme $\Spec(A) \to Y$, que le changement de base
$f_A : X_A \to \Spec(A)$ satisfait la partie d'existence du critère valuatif.
Par définition, l'espace algébrique $X_A$ possède la propriété $(\gamma)$ ; la
Proposition \ref{decent-spaces-proposition-characterize-universally-closed}
s'applique donc au morphisme $f_A$, ce qui achève la preuve.
\end{proof}





\section{Points des fibres}
\label{decent-spaces-section-points-fibres}

\noindent
Soit $S$ un schéma. Considérons un diagramme cartésien
\begin{equation}
\label{decent-spaces-equation-points-fibres}
\xymatrix{
W \ar[r]_q \ar[d]_p & Z \ar[d]^g \\
X \ar[r]^f & Y
}
\end{equation}
d'espaces algébriques sur $S$. Soient $x \in |X|$ et $z \in |Z|$
des points qui s'envoient sur le même point $y \in |Y|$. On peut se demander :
quand l'ensemble
\begin{equation}
\label{decent-spaces-equation-fibre}
F_{x, z} = \{ w \in |W| \text{ tel que }p(w) = x\text{ et }q(w) = z\}
\end{equation}
est-il fini ?

\begin{example}
\label{decent-spaces-example-schemes}
Si $X, Y, Z$ sont des schémas, alors l'ensemble $F_{x, z}$ s'identifie à
l'ensemble sous-jacent au spectre de
$\kappa(x) \otimes_{\kappa(y)} \kappa(z)$
(Schémas, lemme \ref{schemes-lemma-points-fibre-product}). On obtient donc un
ensemble fini si l'extension $\kappa(y) \subset \kappa(x)$ est finie ou si
l'extension $\kappa(y) \subset \kappa(z)$ est finie. C'est en particulier
toujours le cas si $g$ est quasi-fini en $z$ (Morphismes, lemme
\ref{morphisms-lemma-residue-field-quasi-finite}).
\end{example}

\begin{example}
\label{decent-spaces-example-not-finite}
Soit $K$ un corps de caractéristique $0$ muni d'un automorphisme $\sigma$
d'ordre infini. Posons $Y = \Spec(K)/\mathbf{Z}$ et
$X = \mathbf{A}^1_K/\mathbf{Z}$, où $\mathbf{Z}$ agit sur $K$ par $\sigma$
et sur $\mathbf{A}^1_K = \Spec(K[t])$ par $t \mapsto t + 1$.
Soit $Z = \Spec(K)$. Alors $W = \mathbf{A}^1_K$. Voici le diagramme :
$$
\xymatrix{
\mathbf{A}^1_K \ar[r]_q \ar[d]_p & \Spec(K) \ar[d]^g \\
\mathbf{A}^1_K/\mathbf{Z} \ar[r]^f & \Spec(K)/\mathbf{Z}
}
$$
Prenons pour $x$ le point correspondant à $t = 0$ et pour $z$ l'unique point
de $\Spec(K)$. On voit alors que $F_{x, z} = \mathbf{Z}$ en tant qu'ensemble.
\end{example}

\begin{lemma}
\label{decent-spaces-lemma-surjective-on-fibres}
Dans la situation de (\ref{decent-spaces-equation-points-fibres}), si $Z' \to Z$ est un
morphisme et si $z' \in |Z'|$ s'envoie sur $z$, alors l'application induite
$F_{x, z'} \to F_{x, z}$ est surjective.
\end{lemma}

\begin{proof}
Posons $W' = X \times_Y Z' = W \times_Z Z'$. Alors l'application
$|W'| \to |W| \times_{|Z|} |Z'|$ est surjective d'après
Propriétés des espaces, lemme
\ref{spaces-properties-lemma-points-cartesian}.
D'où la surjectivité de $F_{x, z'} \to F_{x, z}$.
\end{proof}

\begin{lemma}
\label{decent-spaces-lemma-qf-and-qc-finite-fibre}
Dans le diagramme (\ref{decent-spaces-equation-points-fibres}), l'ensemble
(\ref{decent-spaces-equation-fibre}) est fini si $f$ est de type fini et si $f$ est
quasi-fini en $x$.
\end{lemma}

\begin{proof}
Le morphisme $q$ est quasi-fini en tout $w \in F_{x, z}$ ; voir
Morphismes d'espaces, lemme
\ref{spaces-morphisms-lemma-base-change-quasi-finite-locus}.
Le lemme résulte donc de
Morphismes d'espaces, lemme
\ref{spaces-morphisms-lemma-quasi-finite-at-a-finite-number-of-points}.
\end{proof}

\begin{lemma}
\label{decent-spaces-lemma-decent-finite-fibre}
Dans le diagramme (\ref{decent-spaces-equation-points-fibres}), l'ensemble
(\ref{decent-spaces-equation-fibre}) est fini si $y$ peut être représenté par un
monomorphisme $\Spec(k) \to Y$, où $k$ est un corps, et si $g$ est quasi-fini
en $z$. (Cas particulier : $Y$ est décent et $g$ est étale.)
\end{lemma}

\begin{proof}
En appliquant deux fois le Lemme \ref{decent-spaces-lemma-surjective-on-fibres},
on peut remplacer $Z$ par $Z_k = \Spec(k) \times_Y Z$ et
$X$ par $X_k = \Spec(k) \times_Y X$. On peut aussi remplacer, et l'on
remplace, $Y$ par $\Spec(k)$. Remarquons que $Z_k \to \Spec(k)$ est quasi-fini
en $z$ d'après Morphismes d'espaces, lemme
\ref{spaces-morphisms-lemma-base-change-quasi-finite-locus}.
Choisissons un schéma $V$, un point $v \in V$ et un morphisme étale
$V \to Z_k$ qui envoie $v$ sur $z$. Choisissons un schéma $U$, un point
$u \in U$ et un morphisme étale $U \to X_k$ qui envoie $u$ sur $x$.
D'après le Lemme \ref{decent-spaces-lemma-surjective-on-fibres},
il suffit encore de montrer que $F_{u, v}$ est fini dans le diagramme
$$
\xymatrix{
U \times_{\Spec(k)} V \ar[r] \ar[d] & V \ar[d] \\
U \ar[r] & \Spec(k)
}
$$
Le morphisme $V \to \Spec(k)$ est quasi-fini en $v$
(cela résulte de la discussion générale de
Morphismes d'espaces, section
\ref{spaces-morphisms-section-local-source-target},
et de la définition de la quasi-finitude en un point).
La finitude résulte alors de l'Exemple \ref{decent-spaces-example-schemes}.
La remarque entre parenthèses de l'énoncé découle du fait que les points des
espaces décents sont représentés par des monomorphismes provenant de corps,
et du fait qu'un morphisme étale d'espaces algébriques est localement
quasi-fini.
\end{proof}

\begin{lemma}
\label{decent-spaces-lemma-topology-fibre}
\begin{slogan}
Les fibres des points à valeurs dans un corps des espaces algébriques ont la
topologie de Zariski attendue.
\end{slogan}
Soit $S$ un schéma. Soit $f : X \to Y$ un morphisme d'espaces algébriques
sur $S$.
Soit $y \in |Y|$ et supposons que $y$ soit représenté par un monomorphisme
quasi-compact $\Spec(k) \to Y$. Alors $|X_k| \to |X|$ est un homéomorphisme
sur $f^{-1}(\{y\}) \subset |X|$, muni de la topologie induite.
\end{lemma}

\begin{proof}
Nous utiliserons
Propriétés des espaces, lemme \ref{spaces-properties-lemma-etale-open},
ainsi que
Morphismes d'espaces, lemme
\ref{spaces-morphisms-lemma-monomorphism-injective-points},
sans autre mention.
Soit $V \to Y$ un morphisme étale, avec $V$ affine, tel qu'il existe un point
$v \in V$ qui s'envoie sur $y$. Comme $\Spec(k) \to Y$ est quasi-compact,
seul un nombre fini de points de $V$ s'envoient sur $y$
(Lemme \ref{decent-spaces-lemma-UR-finite-above-x}). Après avoir rétréci
$V$, on peut supposer que $v$ est le seul. Choisissons un schéma $U$ et un
morphisme étale surjectif $U \to X$.
Considérons le diagramme commutatif
$$
\xymatrix{
U \ar[d] & U_V \ar[l] \ar[d] & U_v \ar[l] \ar[d] \\
X \ar[d] & X_V \ar[l] \ar[d] & X_v \ar[l] \ar[d] \\
Y & V \ar[l] & v \ar[l]
}
$$
Comme $U_v \to U_V$ identifie $U_v$ à une partie de $U_V$ munie de la
topologie induite (Schémas, lemme
\ref{schemes-lemma-fibre-topological}),
et comme $|U_V| \to |X_V|$ et $|U_v| \to |X_v|$ sont surjectives et ouvertes,
on voit que $|X_v| \to |X_V|$ est un homéomorphisme sur son image, munie de la
topologie induite.
D'autre part, l'image inverse de $f^{-1}(\{y\})$ par l'application ouverte
$|X_V| \to |X|$ est égale à $|X_v|$.
On en conclut que $|X_v| \to f^{-1}(\{y\})$ est ouverte.
Le morphisme $X_v \to X$ se factorise par $X_k$, et l'application
$|X_k| \to |X|$ est injective, d'image $f^{-1}(\{y\})$, d'après
Propriétés des espaces, lemme
\ref{spaces-properties-lemma-points-cartesian}. En utilisant
$|X_v| \to |X_k| \to f^{-1}(\{y\})$, le lemme résulte de la surjectivité de
$X_v \to X_k$.
\end{proof}

\begin{lemma}
\label{decent-spaces-lemma-conditions-on-point-on-space-over-field}
Soit $X$ un espace algébrique localement de type fini sur un corps $k$.
Soit $x \in |X|$. Considérons les conditions suivantes :
\begin{enumerate}
\item $\dim_x(|X|) = 0$,
\item $x$ est fermé dans $|X|$ et, si $x' \leadsto x$  dans $|X|$, alors
$x' = x$,
\item $x$ est un point isolé de $|X|$,
\item $\dim_x(X) = 0$,
\item $X \to \Spec(k)$ est quasi-fini en $x$.
\end{enumerate}
Alors (2), (3), (4) et (5) sont équivalentes.
Si $X$ est décent, (1) est équivalente aux autres.
\end{lemma}

\begin{proof}
Les assertions (4) et (5) sont par exemple équivalentes d'après
Morphismes d'espaces, lemmes
\ref{spaces-morphisms-lemma-locally-finite-type-quasi-finite-part} et
\ref{spaces-morphisms-lemma-quasi-finite-at-point}.

\medskip\noindent
Soit $U \to X$ un morphisme étale, où $U$ est un schéma affine, et soit
$u \in U$ un point qui s'envoie sur $x$. De plus, si $x$ est un point fermé,
par exemple dans les cas (2) ou (3), on peut supposer, et l'on suppose, que
$u$ est fermé. Remarquons que $\dim_u(U) = \dim_x(X)$ par définition, et que
ce nombre est égal à $\dim(\mathcal{O}_{U, u})$ si $u$ est fermé ; voir
Algèbre, lemme
\ref{algebra-lemma-dimension-closed-point-finite-type-field}.

\medskip\noindent
Si $\dim_x(X) > 0$ et si $u$ est fermé, les arguments précédents permettent de
choisir une spécialisation non triviale $u' \leadsto u$ dans $U$.
Alors le degré de transcendance de $\kappa(u')$ sur $k$ est strictement
supérieur à celui de $\kappa(u)$ sur $k$. Il s'ensuit que les images $x$ et
$x'$ dans $X$ sont distinctes, car les degrés de transcendance de $x/k$ et de
$x'/k$ sont bien définis ; voir Morphismes d'espaces, définition
\ref{spaces-morphisms-definition-dimension-fibre}.
Cela s'applique notamment dans les cas (2) et (3), et l'on conclut que (2) et
(3) impliquent (4).

\medskip\noindent
Réciproquement, si $X \to \Spec(k)$ est localement quasi-fini en $x$, alors
$U \to \Spec(k)$ est localement quasi-fini en $u$, de sorte que $u$ est un
point isolé de $U$
(Morphismes, lemme \ref{morphisms-lemma-quasi-finite-at-point-characterize}).
Comme $|U| \to |X|$ est continue et ouverte, il s'ensuit que (5) implique (2)
et (3).

\medskip\noindent
Supposons que $X$ soit décent et que (1) soit satisfaite. Alors
$\dim_x(X) = \dim_x(|X|)$ d'après le
Lemme \ref{decent-spaces-lemma-dimension-decent-space}, ce qui achève la preuve.
\end{proof}

\begin{lemma}
\label{decent-spaces-lemma-conditions-on-space-over-field}
Soit $X$ un espace algébrique localement de type fini sur un corps $k$.
Considérons les conditions suivantes :
\begin{enumerate}
\item $|X|$ est un ensemble fini,
\item $|X|$ est un espace discret,
\item $\dim(|X|) = 0$,
\item $\dim(X) = 0$,
\item $X \to \Spec(k)$ est localement quasi-fini,
\end{enumerate}
Alors (2), (3), (4) et (5) sont équivalentes.
Si $X$ est décent, (1) implique les autres.
\end{lemma}

\begin{proof}
Les assertions (4) et (5) sont par exemple équivalentes d'après
Morphismes d'espaces, lemme
\ref{spaces-morphisms-lemma-locally-finite-type-quasi-finite-part}.

\medskip\noindent
Soit $U \to X$ un morphisme étale surjectif, où $U$ est un schéma.

\medskip\noindent
Si $\dim(U) > 0$, choisissons une spécialisation non triviale
$u \leadsto u'$ dans $U$. Le degré de transcendance de $\kappa(u)$ sur $k$
est strictement supérieur à celui de $\kappa(u')$ sur $k$.
Il s'ensuit que les images $x$ et $x'$ dans $X$ sont distinctes, car les
degrés
de transcendance de $x/k$ et de $x'/k$ sont bien définis ; voir
Morphismes d'espaces, définition
\ref{spaces-morphisms-definition-dimension-fibre}.
On conclut que (2) et (3) impliquent (4).

\medskip\noindent
Réciproquement, si $X \to \Spec(k)$ est localement quasi-fini, alors $U$ est
localement noethérien
(Morphismes, lemme \ref{morphisms-lemma-finite-type-noetherian}),
de dimension $0$
(Morphismes, lemme
\ref{morphisms-lemma-locally-quasi-finite-rel-dimension-0}),
et est donc une réunion disjointe de spectres d'anneaux locaux artiniens
(Propriétés, lemme \ref{properties-lemma-locally-Noetherian-dimension-0}).
Ainsi, $U$ est un espace topologique discret et, comme $|U| \to |X|$ est
continue et ouverte, il en va de même de $|X|$.
Autrement dit, (4) implique (2) et (3).

\medskip\noindent
Supposons que $X$ soit décent et que (1) soit satisfaite. On peut alors
choisir $U$ affine ci-dessus. Les fibres de $|U| \to |X|$ sont finies
(cela fait partie
de la propriété qui définit les espaces décents). Ainsi, $U$ est un schéma de
type fini sur $k$ qui possède un nombre fini de points. Ainsi, $U$ est
quasi-fini
sur $k$ (Morphismes, lemme \ref{morphisms-lemma-finite-fibre}),
ce qui signifie par définition que $X \to \Spec(k)$ est localement quasi-fini.
\end{proof}

\begin{lemma}
\label{decent-spaces-lemma-conditions-on-point-in-fibre-and-qf}
Soit $S$ un schéma. Soit $f : X \to Y$ un morphisme d'espaces algébriques
sur $S$, localement de type fini. Soit $x \in |X|$, d'image
$y \in |Y|$. Munissons $F = f^{-1}(\{y\})$ de la topologie induite par $|X|$.
Soit $k$ un corps et soit $\Spec(k) \to Y$ un morphisme appartenant à la
classe d'équivalence qui définit $y$. Posons $X_k = \Spec(k) \times_Y X$.
Soit $\tilde x \in |X_k|$ un point qui s'envoie sur $x \in |X|$.
Considérons les conditions suivantes :
\begin{enumerate}
\item
\label{decent-spaces-item-fibre-at-x-dim-0}
$\dim_x(F) = 0$,
\item
\label{decent-spaces-item-isolated-in-fibre}
$x$ est isolé dans $F$,
\item
\label{decent-spaces-item-no-specializations-in-fibre}
$x$ est fermé dans $F$ et, si $x' \leadsto x$ dans $F$, alors $x = x'$,
\item
\label{decent-spaces-item-dimension-top-k-fibre}
$\dim_{\tilde x}(|X_k|) = 0$,
\item
\label{decent-spaces-item-isolated-in-k-fibre}
$\tilde x$ est isolé dans $|X_k|$,
\item
\label{decent-spaces-item-no-specializations-in-k-fibre}
$\tilde x$ est fermé dans $|X_k|$ et, si $\tilde x' \leadsto \tilde x$
dans $|X_k|$, alors $\tilde x = \tilde x'$,
\item
\label{decent-spaces-item-k-fibre-at-x-dim-0}
$\dim_{\tilde x}(X_k) = 0$,
\item
\label{decent-spaces-item-quasi-finite-at-x}
$f$ est quasi-fini en $x$.
\end{enumerate}
On a alors
$$
\xymatrix{
(\ref{decent-spaces-item-dimension-top-k-fibre}) \ar@{=>}[r]_{f\text{ décent}} &
(\ref{decent-spaces-item-isolated-in-k-fibre}) \ar@{<=>}[r] &
(\ref{decent-spaces-item-no-specializations-in-k-fibre}) \ar@{<=>}[r] &
(\ref{decent-spaces-item-k-fibre-at-x-dim-0}) \ar@{<=>}[r] &
(\ref{decent-spaces-item-quasi-finite-at-x})
}
$$
Si $Y$ est décent, les conditions (\ref{decent-spaces-item-isolated-in-fibre}) et
(\ref{decent-spaces-item-no-specializations-in-fibre}) sont équivalentes entre elles ainsi
qu'aux conditions
(\ref{decent-spaces-item-isolated-in-k-fibre}),
(\ref{decent-spaces-item-no-specializations-in-k-fibre}),
(\ref{decent-spaces-item-k-fibre-at-x-dim-0}) et
(\ref{decent-spaces-item-quasi-finite-at-x}).
Si $Y$ et $X$ sont décents, toutes les conditions sont équivalentes.
\end{lemma}

\begin{proof}
D'après le Lemme \ref{decent-spaces-lemma-conditions-on-point-on-space-over-field}, les
conditions (\ref{decent-spaces-item-isolated-in-k-fibre}),
(\ref{decent-spaces-item-no-specializations-in-k-fibre}) et
(\ref{decent-spaces-item-k-fibre-at-x-dim-0}) sont équivalentes entre elles, ainsi qu'à la
condition que $X_k \to \Spec(k)$ soit quasi-fini en $\tilde x$.
Par conséquent, d'après Morphismes d'espaces, lemme
\ref{spaces-morphisms-lemma-base-change-quasi-finite-locus},
elles sont aussi équivalentes à (\ref{decent-spaces-item-quasi-finite-at-x}).
Si $f$ est décent, alors $X_k$ est un espace algébrique décent, et le
Lemme \ref{decent-spaces-lemma-conditions-on-point-on-space-over-field}
montre que (\ref{decent-spaces-item-dimension-top-k-fibre}) implique
(\ref{decent-spaces-item-isolated-in-k-fibre}).

\medskip\noindent
Si $Y$ est décent, on peut choisir un monomorphisme quasi-compact
$\Spec(k') \to Y$ appartenant à la classe d'équivalence de $y$. Dans ce cas,
le Lemme \ref{decent-spaces-lemma-topology-fibre}
nous dit que $|X_{k'}| \to F$ est un homéomorphisme.
Combiné aux arguments précédents, cela entraîne les autres assertions du
lemme ; les détails sont omis.
\end{proof}

\begin{lemma}
\label{decent-spaces-lemma-conditions-on-fibre-and-qf}
Soit $S$ un schéma. Soit $f : X \to Y$ un morphisme d'espaces algébriques
sur $S$, localement de type fini. Soit $y \in |Y|$. Soit $k$ un corps et soit
$\Spec(k) \to Y$ un morphisme appartenant à la classe d'équivalence qui
définit $y$. Posons $X_k = \Spec(k) \times_Y X$ et munissons
$F = f^{-1}(\{y\})$ de la topologie induite par $|X|$. Considérons les
conditions suivantes :
\begin{enumerate}
\item
\label{decent-spaces-item-fibre-finite}
$F$ est fini,
\item
\label{decent-spaces-item-fibre-discrete}
$F$ est un espace topologique discret,
\item
\label{decent-spaces-item-fibre-no-specializations}
$\dim(F) = 0$,
\item
\label{decent-spaces-item-k-fibre-finite}
$|X_k|$ est un ensemble fini,
\item
\label{decent-spaces-item-k-fibre-discrete}
$|X_k|$ est un espace discret,
\item
\label{decent-spaces-item-k-fibre-no-specializations}
$\dim(|X_k|) = 0$,
\item
\label{decent-spaces-item-k-fibre-dim-0}
$\dim(X_k) = 0$,
\item
\label{decent-spaces-item-quasi-finite-at-points-fibre}
$f$ est quasi-fini en tout point de $|X|$ situé au-dessus de $y$.
\end{enumerate}
On a alors
$$
\xymatrix{
(\ref{decent-spaces-item-fibre-finite}) &
(\ref{decent-spaces-item-k-fibre-finite}) \ar@{=>}[l] \ar@{=>}[r]_{f\text{ décent}} &
(\ref{decent-spaces-item-k-fibre-discrete}) \ar@{<=>}[r] &
(\ref{decent-spaces-item-k-fibre-no-specializations}) \ar@{<=>}[r] &
(\ref{decent-spaces-item-k-fibre-dim-0}) \ar@{<=>}[r] &
(\ref{decent-spaces-item-quasi-finite-at-points-fibre})
}
$$
Si $Y$ est décent, les conditions (\ref{decent-spaces-item-fibre-discrete}) et
(\ref{decent-spaces-item-fibre-no-specializations}) sont équivalentes entre elles ainsi
qu'aux conditions (\ref{decent-spaces-item-k-fibre-discrete}),
(\ref{decent-spaces-item-k-fibre-no-specializations}), (\ref{decent-spaces-item-k-fibre-dim-0}) et
(\ref{decent-spaces-item-quasi-finite-at-points-fibre}).
Si $Y$ et $X$ sont décents, (\ref{decent-spaces-item-fibre-finite}) implique toutes les
autres conditions.
\end{lemma}

\begin{proof}
D'après le Lemme \ref{decent-spaces-lemma-conditions-on-space-over-field}, les conditions
(\ref{decent-spaces-item-k-fibre-discrete}),
(\ref{decent-spaces-item-k-fibre-no-specializations}) et (\ref{decent-spaces-item-k-fibre-dim-0})
sont équivalentes entre elles, ainsi qu'à la condition que
$X_k \to \Spec(k)$ soit localement quasi-fini.
Par conséquent, d'après Morphismes d'espaces, lemme
\ref{spaces-morphisms-lemma-base-change-quasi-finite-locus},
elles sont aussi équivalentes à (\ref{decent-spaces-item-quasi-finite-at-points-fibre}).
Si $f$ est décent, alors $X_k$ est un espace algébrique décent, et le
Lemme \ref{decent-spaces-lemma-conditions-on-space-over-field}
montre que (\ref{decent-spaces-item-k-fibre-finite}) implique
(\ref{decent-spaces-item-k-fibre-discrete}).

\medskip\noindent
L'application $|X_k| \to F$ est surjective d'après
Propriétés des espaces, lemme
\ref{spaces-properties-lemma-points-cartesian},
et l'on obtient
(\ref{decent-spaces-item-k-fibre-finite}) $\Rightarrow$ (\ref{decent-spaces-item-fibre-finite}).

\medskip\noindent
Si $Y$ est décent, on peut choisir un monomorphisme quasi-compact
$\Spec(k') \to Y$ appartenant à la classe d'équivalence de $y$. Dans ce cas,
le Lemme \ref{decent-spaces-lemma-topology-fibre}
nous dit que $|X_{k'}| \to F$ est un homéomorphisme.
Combiné aux arguments précédents, cela entraîne les autres assertions du
lemme ; les détails sont omis.
\end{proof}







\section{Monomorphismes}
\label{decent-spaces-section-monomorphisms}

\noindent
Voici un autre cas où les monomorphismes sont représentables.
On trouvera davantage d'informations dans Morphismes d'espaces, compléments,
section \ref{spaces-more-morphisms-section-monomorphisms}.

\begin{lemma}
\label{decent-spaces-lemma-monomorphism-toward-disjoint-union-dim-0-rings}
Soit $S$ un schéma. Soit $Y$ une réunion disjointe de spectres d'anneaux
locaux de dimension zéro sur $S$.
Soit $f : X \to Y$ un monomorphisme d'espaces algébriques sur $S$.
Alors $f$ est représentable, c'est-à-dire que $X$ est un schéma.
\end{lemma}

\begin{proof}
On se ramène immédiatement au cas où $Y = \Spec(A)$ et où
$A$ est un anneau local de dimension zéro, c'est-à-dire où
$\Spec(A) = \{\mathfrak m_A\}$
est réduit à un seul point. Si $X = \emptyset$, il n'y a rien à démontrer.
Sinon, choisissons un schéma affine non vide $U = \Spec(B)$ et un morphisme
étale $U \to X$. Comme $|X|$ est réduit à un seul point (en tant que partie de
$|Y|$ ; voir Morphismes d'espaces, lemme
\ref{spaces-morphisms-lemma-monomorphism-injective-points}),
on voit que $U \to X$ est surjectif. Remarquons que
$U \times_X U = U \times_Y U = \Spec(B \otimes_A B)$.
Les homomorphismes d'anneaux $B \to B \otimes_A B$ sont donc étales.
Puisque
$$
(B \otimes_A B)/\mathfrak m_A(B \otimes_A B)
=
(B/\mathfrak m_AB) \otimes_{A/\mathfrak m_A} (B/\mathfrak m_AB)
$$
on voit que
$B/\mathfrak m_AB \to (B \otimes_A B)/\mathfrak m_A(B \otimes_A B)$
est plat, et même libre de rang égal à la dimension de
$B/\mathfrak m_AB$ comme espace vectoriel sur $A/\mathfrak m_A$. Comme
$B \to B \otimes_A B$ est étale, cela n'est possible que si cette dimension
est finie (voir par exemple
Morphismes, lemmes \ref{morphisms-lemma-etale-universally-bounded} et
\ref{morphisms-lemma-locally-quasi-finite-qc-source-universally-bounded}).
Tout idéal premier de $B$ est situé au-dessus de $\mathfrak m_A$ (l'unique
idéal premier de $A$). Par conséquent,
$\Spec(B) = \Spec(B/\mathfrak m_A B)$ comme espace topologique, et cet espace
est un ensemble discret fini, puisque $B/\mathfrak m_A B$ est un anneau
artinien ; voir
Algèbre, lemmes \ref{algebra-lemma-finite-dimensional-algebra} et
\ref{algebra-lemma-artinian-finite-length}.
Tous les idéaux premiers de $B$ sont donc maximaux, et
$B = B_1 \times \ldots \times B_n$ est un produit d'un nombre fini d'anneaux
locaux de dimension zéro ; voir
Algèbre, lemme \ref{algebra-lemma-product-local}.
Ainsi, $B \to B \otimes_A B$ est fini étale, car tous les anneaux locaux
$B_i$ sont henséliens d'après
Algèbre, lemme \ref{algebra-lemma-local-dimension-zero-henselian}.
Enfin, $X$ est un schéma affine d'après
Groupoïdes, proposition \ref{groupoids-proposition-finite-flat-equivalence}.
\end{proof}








\section{Points génériques}
\label{decent-spaces-section-generic-points}

\noindent
Cette section prolonge
Propriétés des espaces, section
\ref{spaces-properties-section-generic-points}.

\begin{lemma}
\label{decent-spaces-lemma-decent-generic-points}
Soit $S$ un schéma. Soit $X$ un espace algébrique décent sur $S$.
Soit $x \in |X|$. Les conditions suivantes sont équivalentes :
\begin{enumerate}
\item $x$ est un point générique d'une composante irréductible de $|X|$,
\item pour tout morphisme étale $(Y, y) \to (X, x)$ d'espaces algébriques
pointés, $y$ est un point générique d'une composante irréductible de $|Y|$,
\item pour un certain morphisme étale $(Y, y) \to (X, x)$ d'espaces
algébriques pointés, $y$ est un point générique d'une composante irréductible
de $|Y|$,
\item l'anneau local de $X$ en $x$ est de dimension zéro, et
\item $x$ est un point de codimension $0$ sur $X$.
\end{enumerate}
\end{lemma}

\begin{proof}
Les conditions (4) et (5) sont équivalentes par définition pour tout espace
algébrique ; voir Propriétés des espaces, définition
\ref{spaces-properties-definition-dimension-local-ring}.
Remarquons que tout $Y$ comme en (2) et (3) est décent d'après le
Lemme \ref{decent-spaces-lemma-etale-named-properties}.
Il suffit donc de démontrer l'équivalence de (1) et (4) : l'équivalence avec
(2) et (3) en découle, puisque la dimension de l'anneau local de $Y$ en $y$
est égale à celle de l'anneau local de $X$ en $x$.
Soit $f : U \to X$ un morphisme étale provenant d'un schéma affine et soit
$u \in U$ un point qui s'envoie sur $x$.

\medskip\noindent
Supposons (1). Soit $u' \leadsto u$ une spécialisation dans $U$.
Alors $f(u') = f(u) = x$. D'après le
Lemme \ref{decent-spaces-lemma-decent-no-specializations-map-to-same-point},
on a $u' = u$. Ainsi, $u$ est un point générique d'une composante irréductible
de $U$. Par conséquent, $\dim(\mathcal{O}_{U, u}) = 0$, et (4) est satisfaite.

\medskip\noindent
Supposons (4). Le point $x$ appartient à une composante irréductible
$T \subset |X|$. Puisque $|X|$ est sobre
(Proposition \ref{decent-spaces-proposition-reasonable-sober}),
$T$ possède un point générique $x'$. Bien sûr, $x' \leadsto x$.
On peut relever cette spécialisation en une spécialisation
$u' \leadsto u$ dans $U$
(Lemme \ref{decent-spaces-lemma-decent-specialization}). Cela contredit l'hypothèse
$\dim(\mathcal{O}_{U, u}) = 0$, sauf si $u' = u$, c'est-à-dire si $x' = x$.
\end{proof}

\begin{lemma}
\label{decent-spaces-lemma-codimension-local-ring}
Soit $S$ un schéma. Soit $X$ un espace algébrique décent sur $S$.
Soit $T \subset |X|$ une partie fermée irréductible. Soit $\xi \in T$
son point générique (Proposition \ref{decent-spaces-proposition-reasonable-sober}).
Alors $\text{codim}(T, |X|)$
(Topologie, définition \ref{topology-definition-codimension})
est la dimension de l'anneau local de $X$ en $\xi$
(Propriétés des espaces, définition
\ref{spaces-properties-definition-dimension-local-ring}).
\end{lemma}

\begin{proof}
Choisissons un schéma $U$, un point $u \in U$ et un morphisme étale
$U \to X$ qui envoie $u$ sur $\xi$. Alors toute suite de spécialisations non
triviales $\xi_e \leadsto \ldots \leadsto \xi_0 = \xi$ peut se relever en une
suite $u_e \leadsto \ldots \leadsto u_0 = u$ dans $U$ d'après le
Lemme \ref{decent-spaces-lemma-decent-specialization}.
Réciproquement, toute suite de spécialisations non triviales
$u_e \leadsto \ldots \leadsto u_0 = u$ dans $U$ s'envoie sur une suite de
spécialisations non triviales
$\xi_e \leadsto \ldots \leadsto \xi_0 = \xi$ d'après le
Lemme \ref{decent-spaces-lemma-decent-no-specializations-map-to-same-point}.
Comme $|X|$ et $U$ sont des espaces topologiques sobres, on conclut que la
codimension de $T$ dans $|X|$ est égale à celle de $\overline{\{u\}}$ dans
$U$. Le lemme se ramène ainsi au cas des schémas, qui est
Propriétés, lemme \ref{properties-lemma-codimension-local-ring}.
\end{proof}

\begin{lemma}
\label{decent-spaces-lemma-get-reasonable}
Soit $S$ un schéma. Soit $X$ un espace algébrique sur $S$. Supposons que
\begin{enumerate}
\item tout schéma quasi-compact et étale sur $X$ possède un nombre fini de
composantes irréductibles, et
\item tout $x \in |X|$ de codimension $0$ sur $X$ puisse être représenté
par un monomorphisme $\Spec(k) \to X$.
\end{enumerate}
Alors $X$ est un espace algébrique raisonnable.
\end{lemma}

\begin{proof}
Soit $U$ un schéma affine et soit $a : U \to X$ un morphisme étale.
Il faut montrer que les fibres de $a$ sont universellement bornées. D'après
l'hypothèse (1), le schéma $U$ possède un nombre fini de composantes
irréductibles. Soient $u_1, \ldots, u_n \in U$ les points génériques de ces
composantes irréductibles. Soit $\{x_1, \ldots, x_m\} \subset |X|$ l'image de
$\{u_1, \ldots, u_n\}$. Chaque $x_j$ est un point de codimension $0$.
D'après l'hypothèse (2), on peut choisir un monomorphisme
$\Spec(k_j) \to X$ qui représente $x_j$. D'après Propriétés des espaces, lemme
\ref{spaces-properties-lemma-codimension-0-points}, on a
$$
U \times_X \Spec(k_j) = \coprod\nolimits_{a(u_i) = x_j} \Spec(\kappa(u_i))
$$
C'est un schéma fini sur $\Spec(k_j)$, de degré
$d_j = \sum_{a(u_i) = x_j} [\kappa(u_i) : k_j]$. Posons $n = \max d_j$.

\medskip\noindent
Remarquons que $a$ est séparé
(Propriétés des espaces, lemme
\ref{spaces-properties-lemma-separated-cover}).
Considérons la stratification
$$
X = X_0 \supset X_1 \supset X_2 \supset \ldots
$$
associée à $U \to X$ dans le Lemme \ref{decent-spaces-lemma-stratify-flat-fp-lqf}.
Le choix de $n$ ci-dessus montre que $X_{n + 1}$ est vide.
En effet, sinon, $a^{-1}(X_{n + 1})$ serait un ouvert non vide de $U$, et
contiendrait donc l'un des $u_i$. Cela signifierait que $X_{n + 1}$ contient
$x_j = a(u_i)$, ce qui est impossible.
On voit donc que les fibres de $U \to X$ sont universellement bornées
(en fait, par l'entier $n$).
\end{proof}

\begin{lemma}
\label{decent-spaces-lemma-finitely-many-irreducible-components}
Soit $S$ un schéma. Soit $X$ un espace algébrique sur $S$.
Les conditions suivantes sont équivalentes :
\begin{enumerate}
\item $X$ est décent et $|X|$ possède un nombre fini de composantes
irréductibles,
\item tout schéma quasi-compact et étale sur $X$ possède un nombre fini de
composantes irréductibles, il existe un nombre fini de $x \in |X|$ de
codimension $0$ sur $X$, et chacun d'eux peut être représenté par un
monomorphisme $\Spec(k) \to X$,
\item il existe un ouvert dense $X' \subset X$ qui est un schéma, tel que
$X'$ possède un nombre fini de composantes irréductibles de points génériques
$\{x'_1, \ldots, x'_m\}$ et que le morphisme $x'_j \to X$ soit quasi-compact
pour $j = 1, \ldots, m$.
\end{enumerate}
De plus, si ces conditions sont satisfaites, alors $X$ est raisonnable et les
points $x'_j \in |X|$ sont les points génériques des composantes irréductibles
de $|X|$.
\end{lemma}

\begin{proof}
Dans la preuve, nous utiliserons sans autre mention Propriétés des espaces,
lemme \ref{spaces-properties-lemma-codimension-0-points}.
Supposons (1). D'après le
Théorème \ref{decent-spaces-theorem-decent-open-dense-scheme},
$X$ possède un sous-schéma ouvert dense $X'$.
Comme l'adhérence d'une composante irréductible de $|X'|$ est une composante
irréductible de $|X|$, on voit que $|X'|$ possède un nombre fini de
composantes irréductibles. Ainsi, (3) est satisfaite.

\medskip\noindent
Supposons que $X' \subset X$ soit comme en (3). Soit
$\{x'_1, \ldots, x'_m\}$ l'ensemble des points génériques des composantes
irréductibles de $X'$. Soit $a : U \to X$ un morphisme étale, où $U$ est un
schéma quasi-compact. Pour démontrer (2), il suffit de montrer que $U$ possède
un nombre fini de composantes irréductibles dont les points génériques sont
situés au-dessus de $\{x'_1, \ldots, x'_m\}$. Il suffit de le démontrer pour
les membres d'un recouvrement ouvert affine fini de $U$ ; on peut donc
supposer, et l'on suppose, que $U$ est affine.
Remarquons que $U' = a^{-1}(X') \subset U$ est un ouvert dense.
Puisque $U' \to X'$ est un morphisme étale de schémas, les points génériques
des composantes irréductibles de $U'$ sont les points situés au-dessus de
$\{x'_1, \ldots, x'_m\}$. Comme $x'_j \to X$ est quasi-compact, seul un nombre
fini de points de $U$ sont situés au-dessus de $x'_j$
(Lemme \ref{decent-spaces-lemma-UR-finite-above-x}). Ainsi, $U'$ possède un nombre fini de
composantes irréductibles, ce qui implique que les adhérences de ces
composantes irréductibles sont les composantes irréductibles de $U$.
Ainsi, (2) est satisfaite.

\medskip\noindent
Supposons (2). Le Lemme \ref{decent-spaces-lemma-get-reasonable} implique alors (1) ainsi
que l'assertion finale. (On utilise aussi qu'un espace algébrique raisonnable
est décent ; voir la discussion qui suit la
Définition \ref{decent-spaces-definition-very-reasonable}.)
\end{proof}



\section{Morphismes génériquement finis}
\label{decent-spaces-section-generically-finite}

\noindent
Cette section traite, pour les morphismes d'espaces algébriques, les résultats
exposés dans Morphismes, section
\ref{morphisms-section-generically-finite}, et dans
Variétés, section \ref{varieties-section-generically-finite},
pour les morphismes de schémas.

\begin{lemma}
\label{decent-spaces-lemma-generically-finite}
Soit $S$ un schéma. Soit $f : X \to Y$ un morphisme d'espaces algébriques
sur $S$. Supposons que $f$ soit quasi-séparé et de type fini.
Soit $y \in |Y|$ un point de codimension $0$ sur $Y$.
Les conditions suivantes sont équivalentes :
\begin{enumerate}
\item l'espace $|X_k|$ est fini, où $\Spec(k) \to Y$ représente $y$,
\item $X \to Y$ est quasi-fini en tout point de $|X|$ situé au-dessus de $y$,
\item il existe un sous-espace ouvert $Y' \subset Y$, avec $y \in |Y'|$,
tel que $Y' \times_Y X \to Y'$ soit fini.
\end{enumerate}
Si $Y$ est décent, ces conditions sont aussi équivalentes à
\begin{enumerate}
\item[(4)] l'ensemble $f^{-1}(\{y\})$ est fini.
\end{enumerate}
\end{lemma}

\begin{proof}
L'équivalence de (1) et (2) résulte du
Lemme \ref{decent-spaces-lemma-conditions-on-fibre-and-qf}
(et du fait qu'un morphisme quasi-séparé est décent, d'après le
Lemme \ref{decent-spaces-lemma-properties-trivial-implications}).

\medskip\noindent
Supposons satisfaites les conditions équivalentes (1) et (2). Choisissons un
schéma affine $V$ et un morphisme étale $V \to Y$ qui envoie un point
$v \in V$ sur $y$. Alors $v$ est un point générique d'une composante
irréductible de $V$, d'après Propriétés des espaces, lemme
\ref{spaces-properties-lemma-codimension-0-points}.
Choisissons un schéma affine $U$ et un morphisme étale surjectif
$U \to V \times_Y X$. Alors $U \to V$ est de type fini.
Le morphisme $U \to V$ est quasi-fini en tout point situé au-dessus de $v$,
d'après (2). Il s'ensuit que la fibre de $U \to V$ au-dessus de $v$ est finie
(Morphismes, lemme
\ref{morphisms-lemma-quasi-finite-at-a-finite-number-of-points}). D'après
Morphismes, lemme \ref{morphisms-lemma-generically-finite},
après avoir rétréci $V$, on peut supposer que $U \to V$ est fini.
Posons
$$
R = U \times_{V \times_Y X} U.
$$
Puisque $f$ est quasi-séparé, $V \times_Y X$ est quasi-séparé, et $R$ est
donc un schéma quasi-compact. De plus, le morphisme $R \to V$ est quasi-fini,
comme composé du morphisme étale $R \to U$ et du morphisme fini $U \to V$.
On peut donc appliquer de nouveau Morphismes, lemme
\ref{morphisms-lemma-generically-finite} et, après avoir encore rétréci $V$,
supposer que $R \to V$ est fini. Cela implique bien sûr que les deux
projections $R \to U$ sont finies étales. Il s'ensuit que
$U/R = V \times_Y X$ est un schéma affine ; voir Groupoïdes, proposition
\ref{groupoids-proposition-finite-flat-equivalence}.
D'après Morphismes, lemme \ref{morphisms-lemma-image-proper-is-proper},
on conclut que $V \times_Y X \to V$ est propre, puis, d'après
Morphismes, lemme \ref{morphisms-lemma-finite-proper},
que $V \times_Y X \to V$ est fini. Enfin, soit $Y' \subset Y$ le
sous-espace ouvert de $Y$
correspondant à l'image de $|V| \to |Y|$.
D'après Morphismes d'espaces, lemme
\ref{spaces-morphisms-lemma-integral-local}, on conclut que
$Y' \times_Y X \to Y'$ est fini, puisque son changement de base à $V$ est
fini et que $V \to Y'$ est un morphisme étale surjectif.

\medskip\noindent
Si $Y$ est décent et $f$ quasi-séparé, alors $X$ est également décent ;
utiliser les Lemmes \ref{decent-spaces-lemma-properties-trivial-implications} et
\ref{decent-spaces-lemma-property-over-property}.
Le Lemme \ref{decent-spaces-lemma-conditions-on-fibre-and-qf} montre donc que (4) implique
(1) et (2). Réciproquement, (2) implique (4) d'après Morphismes d'espaces,
lemme \ref{spaces-morphisms-lemma-quasi-finite-at-a-finite-number-of-points}.
\end{proof}

\begin{lemma}
\label{decent-spaces-lemma-generically-finite-reprise}
Soit $S$ un schéma. Soit $f : X \to Y$ un morphisme d'espaces algébriques
sur $S$. Supposons que $f$ soit quasi-séparé et localement de type fini,
et que $Y$ soit quasi-séparé. Soit $y \in |Y|$ un point de codimension $0$
sur $Y$. Les conditions suivantes sont équivalentes :
\begin{enumerate}
\item l'ensemble $f^{-1}(\{y\})$ est fini,
\item l'espace $|X_k|$ est fini, où $\Spec(k) \to Y$ représente $y$,
\item il existe des sous-espaces ouverts $X' \subset X$ et $Y' \subset Y$
tels que $f(X') \subset Y'$, $y \in |Y'|$ et
$f^{-1}(\{y\}) \subset |X'|$, et tels que
$f|_{X'} : X' \to Y'$ soit fini.
\end{enumerate}
\end{lemma}

\begin{proof}
Les espaces algébriques quasi-séparés étant décents, l'équivalence de (1)
et (2) résulte du Lemme \ref{decent-spaces-lemma-conditions-on-fibre-and-qf}.
Pour montrer que (1) et (2) impliquent (3), on peut, et l'on va, remplacer
$Y$ par un ouvert quasi-compact contenant $y$.
Comme $f^{-1}(\{y\})$ est fini, on peut trouver un sous-espace ouvert
quasi-compact $X' \subset X$ contenant la fibre.
La restriction $f|_{X'} : X' \to Y$ est quasi-compacte et quasi-séparée
d'après Morphismes d'espaces, lemme
\ref{spaces-morphisms-lemma-quasi-compact-quasi-separated-permanence}
(c'est ici qu'on utilise que $Y$ est quasi-séparé).
En appliquant le Lemme \ref{decent-spaces-lemma-generically-finite} à
$f|_{X'} : X' \to Y$, on obtient (3).
Nous omettons la démonstration de l'implication de (3) vers (2).
\end{proof}

\begin{lemma}
\label{decent-spaces-lemma-quasi-finiteness-over-generic-point}
Soit $S$ un schéma. Soit $f : X \to Y$ un morphisme d'espaces algébriques
sur $S$. Supposons que $f$ soit localement de type fini.
Notons $X^0 \subset |X|$, respectivement $Y^0 \subset |Y|$, l'ensemble des
points de codimension $0$ de $X$, respectivement de $Y$. Soit $y \in Y^0$.
Les conditions suivantes sont équivalentes :
\begin{enumerate}
\item $f^{-1}(\{y\}) \subset X^0$,
\item $f$ est quasi-fini en tout point situé au-dessus de $y$,
\item $f$ est quasi-fini en tout $x \in X^0$ situé au-dessus de $y$.
\end{enumerate}
\end{lemma}

\begin{proof}
Soit $V$ un schéma et soit $V \to Y$ un morphisme étale surjectif.
Soit $U$ un schéma et soit $U \to V \times_Y X$ un morphisme étale
surjectif. Alors $f$ est quasi-fini en l'image $x$ d'un point $u \in U$
si et seulement si $U \to V$ est quasi-fini en $u$. De plus, $x \in X^0$
si et seulement si $u$ est le point générique d'une composante irréductible
de $U$ (Propriétés des espaces, lemme
\ref{spaces-properties-lemma-codimension-0-points}).
Le lemme se ramène donc au cas du morphisme $U \to V$, c'est-à-dire à
Morphismes, lemme
\ref{morphisms-lemma-quasi-finiteness-over-generic-point}.
\end{proof}

\begin{lemma}
\label{decent-spaces-lemma-finite-over-dense-open}
Soit $S$ un schéma. Soit $f : X \to Y$ un morphisme d'espaces algébriques
sur $S$. Supposons que $f$ soit localement de type fini.
Notons $X^0 \subset |X|$, respectivement $Y^0 \subset |Y|$, l'ensemble des
points de codimension $0$ de $X$, respectivement de $Y$. Supposons que
\begin{enumerate}
\item $Y$ soit décent,
\item $X^0$ et $Y^0$ soient finis et $f^{-1}(Y^0) = X^0$,
\item soit $f$ soit quasi-compact, soit $f$ soit séparé.
\end{enumerate}
Alors il existe un ouvert dense $V \subset Y$ tel que
$f^{-1}(V) \to V$ soit fini.
\end{lemma}

\begin{proof}
D'après les Lemmes \ref{decent-spaces-lemma-finitely-many-irreducible-components} et
\ref{decent-spaces-lemma-decent-generic-points}, on peut supposer que $Y$ est un schéma
ayant un nombre fini de composantes irréductibles. En rétrécissant encore,
on peut supposer que $Y$ est un schéma affine irréductible de point générique
$y$. La fibre de $f$ au-dessus de $y$ est alors finie.

\medskip\noindent
Supposons que $f$ soit quasi-compact et $Y$ affine irréductible. Alors $X$
est quasi-compact, et l'on peut choisir un schéma affine $U$ et un morphisme
étale surjectif $U \to X$. Le morphisme $U \to Y$ est alors de type fini,
et la fibre de $U \to Y$ au-dessus de $y$ est l'ensemble $U^0$ des points
génériques des composantes irréductibles de $U$ (Propriétés des espaces,
lemme \ref{spaces-properties-lemma-codimension-0-points}).
Ainsi, $U^0$ est fini
(Morphismes, lemme
\ref{morphisms-lemma-quasi-finite-at-a-finite-number-of-points}),
et, après avoir rétréci $Y$, on peut supposer que $U \to Y$ est fini
(Morphismes, lemme \ref{morphisms-lemma-generically-finite}).
Considérons ensuite $R = U \times_X U$. Comme la projection
$s : R \to U$ est étale, on a $R^0 = s^{-1}(U^0)$, qui est situé au-dessus
de $y$. Puisque $R \to U \times_Y U$ est un monomorphisme, on conclut que
$R^0$ est fini, car $U \times_Y U \to Y$ est fini.
De plus, $R$ est séparé
(Propriétés des espaces, lemme \ref{spaces-properties-lemma-separated-cover}).
On peut donc rétrécir encore une fois $Y$ pour se ramener au cas où $R$ est
fini sur $Y$
(Morphismes, lemme \ref{morphisms-lemma-finite-over-dense-open}).
Il s'ensuit que $X = U/R$ est fini sur $Y$, par les mêmes arguments que dans
la démonstration du Lemme \ref{decent-spaces-lemma-generically-finite}
(ou bien on peut simplement appliquer ce lemme, puisqu'il en résulte
immédiatement que $X$ est également quasi-séparé).

\medskip\noindent
Supposons que $f$ soit séparé et $Y$ affine irréductible. Choisissons
$V \subset Y$ et $U \subset X$ comme dans le
Lemme \ref{decent-spaces-lemma-generically-finite-reprise}.
Puisque $f|_U : U \to V$ est fini, $U \subset f^{-1}(V)$ est à la fois fermé
et ouvert
(Morphismes d'espaces, lemmes
\ref{spaces-morphisms-lemma-universally-closed-permanence} et
\ref{spaces-morphisms-lemma-finite-proper}).
Ainsi, $f^{-1}(V) = U \amalg W$ pour un certain sous-espace ouvert $W$ de
$X$. Mais, comme $U$ contient tous les points de codimension $0$ de $X$,
on conclut que $W = \emptyset$
(Propriétés des espaces, lemme
\ref{spaces-properties-lemma-codimension-0-points-dense}),
comme souhaité.
\end{proof}






\section{Morphismes birationnels}
\label{decent-spaces-section-birational}

\noindent
La définition suivante d'un morphisme birationnel d'espaces algébriques
semble être celle qui se rapproche le plus de notre définition
(Morphismes, définition \ref{morphisms-definition-birational})
d'un morphisme birationnel de schémas.

\begin{definition}
\label{decent-spaces-definition-birational}
Soit $S$ un schéma. Soient $X$ et $Y$ des espaces algébriques sur $S$.
Supposons que $X$ et $Y$ soient décents et que $|X|$ et $|Y|$ possèdent
un nombre fini de composantes irréductibles. On dit qu'un morphisme
$f : X \to Y$ est {\it birationnel} si
\begin{enumerate}
\item $|f|$ induit une bijection entre l'ensemble des points génériques
des composantes irréductibles de $|X|$ et l'ensemble des points génériques
des composantes irréductibles de $|Y|$, et
\item pour tout point générique $x \in |X|$ d'une composante irréductible,
l'homomorphisme d'anneaux locaux
$\mathcal{O}_{Y, f(x)} \to \mathcal{O}_{X, x}$ est un isomorphisme
(voir la précision ci-dessous).
\end{enumerate}
\end{definition}

\noindent
Précision : puisque $X$ et $Y$ sont décents, les espaces topologiques
$|X|$ et $|Y|$ sont sobres
(Proposition \ref{decent-spaces-proposition-reasonable-sober}).
La condition (1) a donc bien un sens. De plus, puisque nous avons supposé que
$|X|$ et $|Y|$ ont un nombre fini de composantes irréductibles, les points
génériques $x_1, \ldots, x_n \in |X|$, respectivement
$y_1, \ldots, y_n \in |Y|$, appartiennent à tout ouvert dense de $|X|$,
respectivement de $|Y|$. En particulier, ils appartiennent au lieu schématique
de $X$, respectivement de $Y$, d'après le
Théorème \ref{decent-spaces-theorem-decent-open-dense-scheme}.
On peut donc définir $\mathcal{O}_{X, x_i}$, respectivement
$\mathcal{O}_{Y, y_i}$, comme l'anneau local de ce schéma en $x_i$,
respectivement en $y_i$.

\medskip\noindent
On en conclut que, si le morphisme $f : X \to Y$ est birationnel, il existe
des sous-espaces ouverts denses $X' \subset X$ et $Y' \subset Y$ tels que
\begin{enumerate}
\item $f(X') \subset Y'$,
\item $X'$ et $Y'$ soient représentables, et
\item $f|_{X'} : X' \to Y'$ soit birationnel au sens de Morphismes,
définition \ref{morphisms-definition-birational}.
\end{enumerate}
Nous exigeons toutefois que $X$ et $Y$ soient décents et possèdent un nombre
fini de composantes irréductibles. Le
Lemme \ref{decent-spaces-lemma-finitely-many-irreducible-components} donne d'autres
caractérisations des espaces algébriques décents ayant un nombre fini de
composantes irréductibles. Dans la plupart des cas, les morphismes
birationnels sont des isomorphismes au-dessus d'ouverts denses.

\begin{lemma}
\label{decent-spaces-lemma-birational-dominant}
Soit $S$ un schéma. Soit $f : X \to Y$ un morphisme d'espaces algébriques
sur $S$, supposés décents et ayant un nombre fini de composantes
irréductibles. Si $f$ est birationnel, alors $f$ est dominant.
\end{lemma}

\begin{proof}
Cela résulte immédiatement des définitions. Voir Morphismes d'espaces,
définition \ref{spaces-morphisms-definition-dominant}.
\end{proof}

\begin{lemma}
\label{decent-spaces-lemma-birational-generic-fibres}
Soit $S$ un schéma. Soit $f : X \to Y$ un morphisme birationnel d'espaces
algébriques sur $S$, supposés décents et ayant un nombre fini de composantes
irréductibles. Si $y \in |Y|$ est le point générique d'une composante
irréductible, alors le changement de base
$X \times_Y \Spec(\mathcal{O}_{Y, y}) \to
\Spec(\mathcal{O}_{Y, y})$ est un isomorphisme.
\end{lemma}

\begin{proof}
Soient $X' \subset X$ et $Y' \subset Y$ les sous-espaces ouverts maximaux
qui sont représentables ; voir le
Lemme \ref{decent-spaces-lemma-finitely-many-irreducible-components}.
D'après le Lemme \ref{decent-spaces-lemma-quasi-finiteness-over-generic-point},
la fibre de $f$ au-dessus de $y$ est constituée de points de codimension $0$
de $X$ et est donc contenue dans $X'$. Ainsi,
$X \times_Y \Spec(\mathcal{O}_{Y, y}) =
X' \times_{Y'} \Spec(\mathcal{O}_{Y', y})$, et le résultat découle de
Morphismes, lemme \ref{morphisms-lemma-birational-generic-fibres}.
\end{proof}

\begin{lemma}
\label{decent-spaces-lemma-birational-birational}
Soit $S$ un schéma. Soit $f : X \to Y$ un morphisme birationnel d'espaces
algébriques sur $S$, supposés décents et ayant un nombre fini de composantes
irréductibles. Supposons satisfaite l'une des conditions suivantes :
\begin{enumerate}
\item $f$ est localement de type fini et $Y$ est réduit
(par exemple intègre),
\item $f$ est localement de présentation finie.
\end{enumerate}
Alors il existe des ouverts denses $U \subset X$ et $V \subset Y$ tels que
$f(U) \subset V$ et que $f|_U : U \to V$ soit un isomorphisme.
\end{lemma}

\begin{proof}
D'après le Lemme \ref{decent-spaces-lemma-finitely-many-irreducible-components},
on peut supposer que $X$ et $Y$ sont des schémas. Dans ce cas, le résultat est
Morphismes, lemme \ref{morphisms-lemma-birational-birational}.
\end{proof}

\begin{lemma}
\label{decent-spaces-lemma-birational-isomorphism-over-dense-open}
Soit $S$ un schéma. Soit $f : X \to Y$ un morphisme birationnel d'espaces
algébriques sur $S$, supposés décents et ayant un nombre fini de composantes
irréductibles. Supposons que
\begin{enumerate}
\item soit $f$ soit quasi-compact, soit $f$ soit séparé, et
\item soit $f$ soit localement de type fini et $Y$ soit réduit, soit $f$
soit localement de présentation finie.
\end{enumerate}
Alors il existe un ouvert dense $V \subset Y$ tel que
$f^{-1}(V) \to V$ soit un isomorphisme.
\end{lemma}

\begin{proof}
D'après le Lemme \ref{decent-spaces-lemma-finitely-many-irreducible-components},
on peut supposer que $Y$ est un schéma. D'après le
Lemme \ref{decent-spaces-lemma-finite-over-dense-open}, on peut supposer que $f$ est fini.
Alors $X$ est aussi un schéma, et le résultat découle de Morphismes,
lemme \ref{morphisms-lemma-birational-isomorphism-over-dense-open}.
\end{proof}

\begin{lemma}
\label{decent-spaces-lemma-birational-etale-localization}
Soit $S$ un schéma. Soit $f : X \to Y$ un morphisme d'espaces algébriques
sur $S$, supposés décents et ayant un nombre fini de composantes
irréductibles. Si $f$ est birationnel et si $V \to Y$ est un morphisme étale
avec $V$ affine, alors $X \times_Y V$ est décent, possède un nombre fini de
composantes irréductibles et $X \times_Y V \to V$ est birationnel.
\end{lemma}

\begin{proof}
L'espace algébrique $U = X \times_Y V$ est décent
(Lemme \ref{decent-spaces-lemma-etale-named-properties}).
Les points génériques de $V$ et de $U$ sont les éléments de $|V|$ et de $|U|$
qui sont situés au-dessus des points génériques de $|Y|$ et de $|X|$
(Lemme \ref{decent-spaces-lemma-decent-generic-points}).
Comme $Y$ est décent, on en conclut que $V$ ne possède qu'un nombre fini de
points génériques. Soit $\xi \in |X|$ un point générique d'une composante
irréductible. D'après la discussion qui suit la
Définition \ref{decent-spaces-definition-birational}, on a un carré cartésien
$$
\xymatrix{
\Spec(\mathcal{O}_{X, \xi}) \ar[d] \ar[r] & X \ar[d] \\
\Spec(\mathcal{O}_{Y, f(\xi)}) \ar[r] & Y
}
$$
dont les morphismes horizontaux sont des monomorphismes qui identifient les
anneaux locaux, et dont la flèche verticale de gauche est un isomorphisme.
Il en résulte que, dans le diagramme
$$
\xymatrix{
\Spec(\mathcal{O}_{X, \xi}) \times_X U \ar[d] \ar[r] & U \ar[d] \\
\Spec(\mathcal{O}_{Y, f(\xi)}) \times_Y V \ar[r] & V
}
$$
la flèche verticale de gauche est un isomorphisme. Les flèches horizontales
ont leur image contenue dans le lieu schématique de $U$ et de $V$ et
identifient les anneaux locaux (nous omettons certains détails). Comme les
images des flèches horizontales sont les points de $|U|$, respectivement de
$|V|$, situés au-dessus de $\xi$, respectivement de $f(\xi)$, on conclut.
\end{proof}

\begin{lemma}
\label{decent-spaces-lemma-birational-induced-morphism-normalizations}
Soit $S$ un schéma. Soit $f : X \to Y$ un morphisme birationnel entre des
espaces algébriques sur $S$, supposés décents et ayant un nombre fini de
composantes irréductibles. Alors les normalisations $X^\nu \to X$ et
$Y^\nu \to Y$ existent, et il existe un diagramme commutatif
$$
\xymatrix{
X^\nu \ar[r] \ar[d] &  Y^\nu \ar[d] \\
X \ar[r] & Y
}
$$
d'espaces algébriques sur $S$. Le morphisme $X^\nu \to Y^\nu$ est
birationnel.
\end{lemma}

\begin{proof}
D'après le Lemme \ref{decent-spaces-lemma-finitely-many-irreducible-components},
$X$ et $Y$ satisfont aux conditions équivalentes de Morphismes d'espaces,
lemme \ref{spaces-morphisms-lemma-prepare-normalization}, et les
normalisations sont définies. D'après Morphismes d'espaces, lemme
\ref{spaces-morphisms-lemma-normalization-normal}, l'espace algébrique
$X^\nu$ est normal et envoie les points de codimension $0$ sur des points de
codimension $0$. Comme $f$ envoie les points de codimension $0$ sur des
points de codimension $0$ (ils coïncident avec les points génériques sur les
espaces décents d'après le Lemme \ref{decent-spaces-lemma-decent-generic-points}),
Morphismes d'espaces, lemme
\ref{spaces-morphisms-lemma-normalization-normal}, fournit une factorisation
du composé $X^\nu \to X \to Y$ par $Y^\nu$.

\medskip\noindent
Remarquons que $X^\nu$ et $Y^\nu$ sont décents, par exemple d'après le
Lemme \ref{decent-spaces-lemma-representable-named-properties}.
De plus, les applications $X^\nu \to X$ et $Y^\nu \to Y$ induisent des
bijections sur les composantes irréductibles (voir les références ci-dessus) ;
ainsi, $X^\nu$ et $Y^\nu$ possèdent tous deux un nombre fini de composantes
irréductibles, et l'application $X^\nu \to Y^\nu$ induit une bijection entre
leurs points génériques. Pour montrer que $X^\nu \to Y^\nu$ est birationnel,
il suffit donc de montrer qu'il induit un isomorphisme sur les anneaux locaux
en ces points. Pour cela, on peut remplacer $X$ et $Y$ par des voisinages
ouverts de leurs points génériques ; on peut donc supposer que $X$ et $Y$
sont des schémas affines irréductibles de points génériques $x$ et $y$.
Puisque $f$ est birationnel, l'homomorphisme
$\mathcal{O}_{Y, y} \to \mathcal{O}_{X, x}$ est un isomorphisme.
Soient $x^\nu \in X^\nu$ et $y^\nu \in Y^\nu$ les points situés au-dessus
de $x$ et de $y$. Par construction de la normalisation, on a
$\mathcal{O}_{X^\nu, x^\nu} = \mathcal{O}_{X, x}/\mathfrak m_x$,
et de même sur $Y$. Ainsi, l'homomorphisme
$\mathcal{O}_{Y^\nu, y^\nu} \to \mathcal{O}_{X^\nu, x^\nu}$
est lui aussi un isomorphisme.
\end{proof}

\begin{lemma}
\label{decent-spaces-lemma-finite-birational-over-normal}
Soit $S$ un schéma. Soit $f : X \to Y$ un morphisme d'espaces algébriques
sur $S$. Supposons que
\begin{enumerate}
\item $X$ et $Y$ soient décents et aient un nombre fini de composantes
irréductibles,
\item $f$ soit intégral et birationnel,
\item $Y$ soit normal, et
\item $X$ soit réduit.
\end{enumerate}
Alors $f$ est un isomorphisme.
\end{lemma}

\begin{proof}
Soit $V \to Y$ un morphisme étale, avec $V$ affine. Il suffit de montrer que
$U = X \times_Y V \to V$ est un isomorphisme. D'après le
Lemme \ref{decent-spaces-lemma-birational-etale-localization} et sa démonstration,
$U$ et $V$ sont décents et possèdent un nombre fini de composantes
irréductibles, et $U \to V$ est birationnel. D'après Propriétés, lemme
\ref{properties-lemma-normal-locally-finite-nr-irreducibles},
$V$ est une réunion disjointe finie de schémas intègres.
On peut donc supposer que $V$ est intègre. Comme $f$ est birationnel, $U$ est
irréductible et réduit, c'est-à-dire intègre
(remarquer que $U$ est un schéma, puisque $f$ est intégral, donc
représentable). On peut ainsi supposer que $X$ et $Y$ sont des schémas
intègres, et le résultat découle du cas des schémas ; voir Morphismes, lemme
\ref{morphisms-lemma-finite-birational-over-normal}.
\end{proof}

\begin{lemma}
\label{decent-spaces-lemma-normalization-normal}
Soit $S$ un schéma. Soit $f : X \to Y$ un morphisme intégral birationnel
d'espaces algébriques décents sur $S$, ayant un nombre fini de composantes
irréductibles. Il existe alors une factorisation $Y^\nu \to X \to Y$, et
$Y^\nu \to X$ est la normalisation de $X$.
\end{lemma}

\begin{proof}
Considérons l'application $X^\nu \to Y^\nu$ du
Lemme \ref{decent-spaces-lemma-birational-induced-morphism-normalizations}.
Cette application est intégrale d'après Morphismes d'espaces, lemme
\ref{spaces-morphisms-lemma-finite-permanence}.
Elle est donc un isomorphisme d'après le
Lemme \ref{decent-spaces-lemma-finite-birational-over-normal}.
\end{proof}




\section{Espaces de Jacobson}
\label{decent-spaces-section-jacobson}

\noindent
Nous avons défini la propriété de Jacobson pour les espaces algébriques dans
Propriétés des espaces, remarque
\ref{spaces-properties-remark-list-properties-local-etale-topology}.
Pour les espaces algébriques représentables, elle coïncide avec la propriété
étudiée dans Propriétés, section \ref{properties-section-jacobson}.
Le rapport entre la propriété de Jacobson et le comportement de l'espace
topologique $|X|$ n'est pas évident pour un espace algébrique général $X$.
Cependant, un espace algébrique décent $X$ (par exemple quasi-séparé ou
localement séparé) est de Jacobson si et seulement si $|X|$ est de Jacobson
(voir le Lemme \ref{decent-spaces-lemma-decent-Jacobson}).

\begin{lemma}
\label{decent-spaces-lemma-Jacobson-universally-Jacobson}
Soit $S$ un schéma. Soit $X$ un espace algébrique de Jacobson sur $S$.
Tout espace algébrique localement de type fini sur $X$ est de Jacobson.
\end{lemma}

\begin{proof}
Soit $U \to X$ un morphisme étale surjectif, où $U$ est un schéma.
Alors $U$ est de Jacobson (par définition), et, pour tout morphisme de schémas
$V \to U$ localement de type fini, $V$ est de Jacobson d'après le résultat
correspondant pour les schémas (Morphismes, lemme
\ref{morphisms-lemma-Jacobson-universally-Jacobson}).
Ainsi, si $Y \to X$ est un morphisme d'espaces algébriques localement de type
fini, en posant $V = U \times_X Y$, on voit par définition que $Y$ est de
Jacobson.
\end{proof}

\begin{lemma}
\label{decent-spaces-lemma-Jacobson-ft-points-lift-to-closed}
Soit $S$ un schéma. Soit $X$ un espace algébrique de Jacobson sur $S$.
Soient $x \in X_{\text{ft-pts}}$ et $g : W \to X$ un morphisme localement de
type fini, où $W$ est un schéma. Si $x \in \Im(|g|)$, alors il existe un point
fermé de $W$ qui s'envoie sur $x$.
\end{lemma}

\begin{proof}
Soit $U \to X$ un morphisme étale, où $U$ est un schéma, et soit $u \in U$
un point fermé qui s'envoie sur $x$ ; voir Morphismes d'espaces, lemme
\ref{spaces-morphisms-lemma-identify-finite-type-points}.
Remarquons que $W$, $W \times_X U$ et $U$ sont des schémas de Jacobson d'après
le Lemme \ref{decent-spaces-lemma-Jacobson-universally-Jacobson}.
Sur ces schémas, les points de type fini sont donc exactement les points
fermés, d'après Morphismes, lemme
\ref{morphisms-lemma-jacobson-finite-type-points}.
L'image inverse $T \subset W \times_X U$ de $u$ est une partie fermée non vide
(puisque $x$ appartient à l'image de $W \to X$).
D'après Morphismes, lemme \ref{morphisms-lemma-enough-finite-type-points},
il existe un point fermé $t$ de $W \times_X U$ qui s'envoie sur $u$.
Comme $W \times_X U \to W$ est localement de type fini, l'image de $t$ dans
$W$ est fermée d'après Morphismes, lemme
\ref{morphisms-lemma-jacobson-finite-type-points}.
\end{proof}

\begin{lemma}
\label{decent-spaces-lemma-decent-Jacobson-ft-pts}
Soit $S$ un schéma. Soit $X$ un espace algébrique décent de Jacobson sur $S$.
Alors $X_{\text{ft-pts}} \subset |X|$ est l'ensemble des points fermés.
\end{lemma}

\begin{proof}
Si $x \in |X|$ est fermé, on peut représenter $x$ par une immersion fermée
$\Spec(k) \to X$ ; voir le Lemme \ref{decent-spaces-lemma-decent-space-closed-point}.
Ainsi, $x$ est certainement un point de type fini.

\medskip\noindent
Réciproquement, soit $x \in |X|$ un point de type fini. On sait que $x$ peut
être représenté par un monomorphisme quasi-compact $\Spec(k) \to X$, où $k$
est un corps (Définition \ref{decent-spaces-definition-very-reasonable}). D'autre part, par
définition, il existe un morphisme $\Spec(k') \to X$ localement de type fini
qui représente $x$
(Morphismes, définition \ref{morphisms-definition-finite-type-point}).
On obtient une factorisation $\Spec(k') \to \Spec(k) \to X$.
Soit $U \to X$ un morphisme étale quelconque, avec $U$ affine, et considérons
les morphismes
$$
\Spec(k') \times_X U \to \Spec(k) \times_X U \to U.
$$
Le schéma quasi-compact $\Spec(k) \times_X U$ est étale sur $\Spec(k)$ ;
c'est donc une réunion disjointe finie de spectres de corps
(Remarque \ref{decent-spaces-remark-recall}). De plus, le premier morphisme est surjectif
et localement de type fini
(Morphismes, lemme \ref{morphisms-lemma-permanence-finite-type}) ;
il est donc surjectif sur les points de type fini
(Morphismes, lemme
\ref{morphisms-lemma-finite-type-points-surjective-morphism}), et le composé
(qui est localement de type fini) envoie les points de type fini sur des
points fermés, puisque $U$ est de Jacobson
(Morphismes, lemme \ref{morphisms-lemma-jacobson-finite-type-points}).
Ainsi, l'image de $\Spec(k) \times_X U \to U$ est un ensemble fini de points
fermés, donc fermé. Comme cela vaut pour tout $U$ affine et tout morphisme
étale $U \to X$, on conclut que $x \in |X|$ est fermé.
\end{proof}

\begin{lemma}
\label{decent-spaces-lemma-decent-Jacobson}
Soit $S$ un schéma. Soit $X$ un espace algébrique décent sur $S$.
Alors $X$ est de Jacobson si et seulement si $|X|$ est de Jacobson.
\end{lemma}

\begin{proof}
Supposons que $X$ soit de Jacobson et que $T \subset |X|$ soit une partie
fermée. D'après Morphismes d'espaces, lemme
\ref{spaces-morphisms-lemma-enough-finite-type-points}, on voit que
$T \cap X_{\text{ft-pts}}$ est dense dans $T$.
D'après le Lemme \ref{decent-spaces-lemma-decent-Jacobson-ft-pts},
$X_{\text{ft-pts}}$ est l'ensemble des points fermés de $|X|$.
Ainsi, $|X|$ est bien de Jacobson.

\medskip\noindent
Supposons que $|X|$ soit de Jacobson. Soit $f : U \to X$ un morphisme étale,
où $U$ est un schéma affine. Il faut montrer que $U$ est de Jacobson.
Si $x \in |X|$ est fermé, la fibre $F = f^{-1}(\{x\})$ est une partie finie
(par définition de décent) et fermée (par construction de la topologie sur
$|X|$) de $U$. Comme il n'y a pas de spécialisations entre les points de $F$
(Lemme \ref{decent-spaces-lemma-decent-no-specializations-map-to-same-point}),
on conclut que tout point de $F$ est fermé dans $U$.
Si $U$ n'est pas de Jacobson, il existe un point non fermé $u \in U$ tel que
$\{u\}$ soit localement fermé (Topologie, lemme
\ref{topology-lemma-non-jacobson-Noetherian-characterize}).
Nous allons montrer que $f(u) \in |X|$ est fermé ; d'après ce qui précède,
$u$ serait alors fermé dans $U$, contradiction qui achèvera la démonstration.
Pour cela, on peut remplacer $U$ par un voisinage ouvert affine de $u$.
On peut donc supposer que $\{u\}$ est fermé dans $U$.
Posons $R = U \times_X U$, avec projections $s, t : R \to U$.
Alors $s^{-1}(\{u\}) = \{r_1, \ldots, r_m\}$ est fini
(par définition des espaces décents). Après avoir remplacé $U$ par un
voisinage ouvert affine plus petit de $u$, on peut supposer que
$t(r_j) = u$ pour $j = 1, \ldots, m$. Il s'ensuit que $\{u\}$ est une partie
fermée $R$-invariante de $U$. Ainsi, $\{f(u)\}$ est une partie localement
fermée de $X$, puisqu'elle est fermée dans l'ouvert $|f|(|U|)$ de $|X|$.
Comme $|X|$ est de Jacobson, on conclut que $f(u)$ est fermé dans $|X|$,
comme souhaité.
\end{proof}

\begin{lemma}
\label{decent-spaces-lemma-punctured-spec}
Soit $S$ un schéma. Soit $X$ un espace algébrique décent localement
noethérien sur $S$. Soit $x \in |X|$. Alors
$$
W = \{x' \in |X| : x' \leadsto x,\ x' \not = x\}
$$
est un espace topologique noethérien, spectral, sobre et de Jacobson.
\end{lemma}

\begin{proof}
On peut effectuer le remplacement par tout sous-espace ouvert contenant $x$.
On peut donc supposer que $X$ est quasi-compact.
Alors $|X|$ est un espace topologique noethérien
(Propriétés des espaces, lemme
\ref{spaces-properties-lemma-Noetherian-topology}).
Ainsi, $W$ est un espace topologique noethérien
(Topologie, lemme \ref{topology-lemma-Noetherian}).

\medskip\noindent
En combinant le Lemme \ref{decent-spaces-lemma-locally-Noetherian-decent-quasi-separated}
et Propriétés des espaces, lemme
\ref{spaces-properties-lemma-quasi-compact-quasi-separated-spectral},
on voit que $|X|$ est un espace topologique spectral.
D'après Topologie, lemme \ref{topology-lemma-make-spectral-space},
$W \cup \{x\}$ est un espace topologique spectral.
Or $W$ est un ouvert quasi-compact de $W \cup \{x\}$ ; $W$ est donc spectral
d'après Topologie, lemme \ref{topology-lemma-spectral-sub}.

\medskip\noindent
Soit $E \subset W$ une partie fermée irréductible. Si $Z \subset |X|$ est
l'adhérence de $E$, alors $x \in Z$. Il existe un unique point générique
$\eta \in Z$ d'après la Proposition \ref{decent-spaces-proposition-reasonable-sober}.
Bien sûr, $\eta \in W$, donc $\eta \in E$. On conclut que $E$ possède un
unique point générique, c'est-à-dire que $W$ est sobre.

\medskip\noindent
Soit $x' \in W$ un point tel que $\{x'\}$ soit localement fermé dans $W$.
Pour achever la démonstration, il faut montrer que $x'$ est un point fermé de
$W$. Sinon, il existe une spécialisation non triviale
$x' \leadsto x'_1$ dans $W$. Soient $U$ un schéma affine, $u \in U$ un point
et $U \to X$ un morphisme étale envoyant $u$ sur $x$. D'après le
Lemme \ref{decent-spaces-lemma-decent-specialization}, on peut choisir des spécialisations
$u' \leadsto u'_1 \leadsto u$ qui s'envoient sur
$x' \leadsto x'_1 \leadsto x$.
Soit $\mathfrak p' \subset \mathcal{O}_{U, u}$ l'idéal premier correspondant
à $u'$. L'existence de ces spécialisations implique que
$\dim(\mathcal{O}_{U, u}/\mathfrak p') \geq 2$.
Ainsi, tout ouvert non vide de
$\Spec(\mathcal{O}_{U, u}/\mathfrak p')$ est infini d'après Algèbre, lemme
\ref{algebra-lemma-Noetherian-local-domain-dim-2-infinite-opens}.
D'après le Lemme \ref{decent-spaces-lemma-decent-no-specializations-map-to-same-point},
on obtient une application continue
$$
\Spec(\mathcal{O}_{U, u}/\mathfrak p')
\setminus \{\mathfrak m_u/\mathfrak p'\}
\longrightarrow
W.
$$
Comme le point générique du membre de gauche s'envoie sur $x'$, l'image est
contenue dans $\overline{\{x'\}}$. On en conclut que l'image inverse de
$\{x'\}$ par la flèche affichée est un ouvert non vide, donc infini.
Cependant, les fibres de $U \to X$ sont finies puisque $X$ est décent ; la
fibre au-dessus de $x'$ serait donc infinie. Cette contradiction achève la
démonstration.
\end{proof}






\section{Irréductibilité locale}
\label{decent-spaces-section-irreducible-local-ring}

\noindent
Nous avons déjà défini le nombre géométrique de branches d'un espace
algébrique en un point dans Propriétés des espaces, section
\ref{spaces-properties-section-irreducible-local-ring}.
Le nombre de branches d'un espace algébrique en un point ne peut être défini
que pour les espaces algébriques décents.

\begin{lemma}
\label{decent-spaces-lemma-irreducible-local-ring}
Soit $S$ un schéma. Soit $X$ un espace algébrique décent sur $S$.
Soit $x \in |X|$ un point. Les conditions suivantes sont équivalentes :
\begin{enumerate}
\item pour tout voisinage étale élémentaire $(U, u) \to (X, x)$,
l'anneau local $\mathcal{O}_{U, u}$ possède un unique idéal premier minimal,
\item pour tout voisinage étale élémentaire $(U, u) \to (X, x)$,
il existe une unique composante irréductible de $U$ passant par $u$,
\item pour tout voisinage étale élémentaire $(U, u) \to (X, x)$,
l'anneau local $\mathcal{O}_{U, u}$ est unibranche,
\item l'anneau local hensélien $\mathcal{O}_{X, x}^h$ possède un unique idéal
premier minimal.
\end{enumerate}
\end{lemma}

\begin{proof}
L'équivalence de (1) et (2) résulte du fait que les composantes irréductibles
de $U$ passant par $u$ sont en correspondance bijective ($1$-$1$) avec les
premiers minimaux de l'anneau local de $U$ en $u$. L'anneau
$\mathcal{O}_{X, x}^h$ est l'hensélisé de $\mathcal{O}_{U, u}$ ; voir la
discussion qui suit la Définition \ref{decent-spaces-definition-henselian-local-ring}.
En particulier, (3) et (4) sont équivalentes d'après Compléments d'algèbre,
lemme \ref{more-algebra-lemma-unibranch}.
L'équivalence de (2) et (3) découle de Compléments sur les morphismes, lemme
\ref{more-morphisms-lemma-nr-branches}.
\end{proof}

\begin{definition}
\label{decent-spaces-definition-unibranch}
Soit $S$ un schéma. Soit $X$ un espace algébrique décent sur $S$.
Soit $x \in |X|$. On dit que $X$ est {\it unibranche en $x$} si les
conditions équivalentes du Lemme \ref{decent-spaces-lemma-irreducible-local-ring}
sont satisfaites. On dit que $X$ est {\it unibranche} si $X$ est unibranche
en tout $x \in |X|$.
\end{definition}

\noindent
Cette définition est compatible avec celle donnée pour les schémas
(Propriétés, définition \ref{properties-definition-unibranch}).

\begin{lemma}
\label{decent-spaces-lemma-nr-branches-local-ring}
Soit $S$ un schéma. Soit $X$ un espace algébrique décent sur $S$.
Soit $x \in |X|$ un point. Soit $n \in \{1, 2, \ldots\}$ un entier.
Les conditions suivantes sont équivalentes :
\begin{enumerate}
\item pour tout voisinage étale élémentaire $(U, u) \to (X, x)$, le nombre
d'idéaux premiers minimaux de l'anneau local $\mathcal{O}_{U, u}$ est
$\leq n$, et, pour au moins un choix de $(U, u)$, il est égal à $n$,
\item pour tout voisinage étale élémentaire $(U, u) \to (X, x)$, le nombre
de composantes irréductibles de $U$ passant par $u$ est $\leq n$, et, pour
au moins un choix de $(U, u)$, il est égal à $n$,
\item pour tout voisinage étale élémentaire $(U, u) \to (X, x)$, le nombre
de branches de $U$ en $u$ est $\leq n$, et, pour au moins un choix de
$(U, u)$, il est égal à $n$,
\item le nombre d'idéaux premiers minimaux de
$\mathcal{O}_{X, x}^h$ est $n$.
\end{enumerate}
\end{lemma}

\begin{proof}
L'équivalence de (1) et (2) résulte du fait que les composantes irréductibles
de $U$ passant par $u$ sont en correspondance bijective ($1$-$1$) avec les
premiers minimaux de l'anneau local de $U$ en $u$.
L'anneau $\mathcal{O}_{X, x}^h$ est l'hensélisé de $\mathcal{O}_{U, u}$ ;
voir la discussion qui suit la
Définition \ref{decent-spaces-definition-henselian-local-ring}.
En particulier, (3) et (4) sont équivalentes d'après Compléments d'algèbre,
lemme \ref{more-algebra-lemma-unibranch}.
L'équivalence de (2) et (3) découle de Compléments sur les morphismes, lemme
\ref{more-morphisms-lemma-nr-branches}.
\end{proof}

\begin{definition}
\label{decent-spaces-definition-number-of-branches}
Soit $S$ un schéma. Soit $X$ un espace algébrique décent sur $S$.
Soit $x \in |X|$. Le {\it nombre de branches de $X$ en $x$} est égal soit à
$n \in \mathbf{N}$ si les conditions équivalentes du
Lemme \ref{decent-spaces-lemma-nr-branches-local-ring} sont satisfaites, soit à $\infty$.
\end{definition}






\section{Espaces algébriques caténaires}
\label{decent-spaces-section-catenary}

\noindent
Cette section étend aux espaces algébriques les résultats de Propriétés,
section \ref{properties-section-catenary}, et de Morphismes, section
\ref{morphisms-section-universally-catenary}.

\begin{definition}
\label{decent-spaces-definition-catenary}
Soit $S$ un schéma. Soit $X$ un espace algébrique décent sur $S$.
On dit que $X$ est {\it caténaire} si $|X|$ est caténaire
(Topologie, définition \ref{topology-definition-catenary}).
\end{definition}

\noindent
Si $X$ est représentable, cette définition équivaut à la notion
correspondante pour le schéma qui représente $X$.

\begin{lemma}
\label{decent-spaces-lemma-scheme-with-dimension-function}
Soit $S$ un schéma localement noethérien et universellement caténaire.
Soit $\delta : S \to \mathbf{Z}$ une fonction de dimension.
Soit $X$ un espace algébrique décent sur $S$ tel que le morphisme structural
$X \to S$ soit localement de type fini. Soit
$\delta_X : |X| \to \mathbf{Z}$ l'application qui envoie $x$ sur
$\delta(f(x))$ augmenté du degré de transcendance de $x/f(x)$.
Alors $\delta_X$ est une fonction de dimension sur $|X|$.
\end{lemma}

\begin{proof}
Soit $\varphi : U \to X$ un morphisme étale surjectif, où $U$ est un schéma.
Alors la fonction $\delta_U$, définie de manière analogue, est une fonction
de dimension sur $U$ d'après Morphismes, lemme
\ref{morphisms-lemma-dimension-function-propagates}.
D'autre part, par définition du degré de transcendance relatif dans
Morphismes d'espaces, définition
\ref{spaces-morphisms-definition-dimension-fibre}, on a
$\delta_U(u) = \delta_X(\varphi(u))$.

\medskip\noindent
Soit $x \leadsto x'$ une spécialisation de points de $|X|$.
D'après le Lemme \ref{decent-spaces-lemma-decent-specialization}, on peut trouver une
spécialisation $u \leadsto u'$ de points de $U$ telle que
$\varphi(u) = x$ et $\varphi(u') = x'$. De plus, $x = x'$ si et seulement si
$u = u'$ ; voir le
Lemme \ref{decent-spaces-lemma-decent-no-specializations-map-to-same-point}.
Le fait que $\delta_U$ soit une fonction de dimension implique donc que
$\delta_X$ en est une ; voir Topologie, définition
\ref{topology-definition-dimension-function}.
\end{proof}

\begin{lemma}
\label{decent-spaces-lemma-universally-catenary-scheme}
Soit $S$ un schéma localement noethérien et universellement caténaire.
Soit $X$ un espace algébrique sur $S$ tel que $X$ soit décent et que le
morphisme structural $X \to S$ soit localement de type fini. Alors $X$ est
caténaire.
\end{lemma}

\begin{proof}
La question est locale sur $S$ (utiliser Topologie, lemme
\ref{topology-lemma-catenary}).
On peut donc supposer que $S$ possède une fonction de dimension ; voir
Topologie, lemme \ref{topology-lemma-locally-dimension-function}.
D'après le Lemme \ref{decent-spaces-lemma-scheme-with-dimension-function},
$|X|$ possède alors une fonction de dimension.
Comme $|X|$ est sobre (Proposition \ref{decent-spaces-proposition-reasonable-sober}),
on conclut que $|X|$ est caténaire d'après Topologie, lemme
\ref{topology-lemma-dimension-function-catenary}.
\end{proof}

\noindent
D'après le Lemme \ref{decent-spaces-lemma-universally-catenary-scheme}, la définition
suivante est compatible avec la notion déjà existante pour les espaces
algébriques représentables.

\begin{definition}
\label{decent-spaces-definition-universally-catenary}
Soit $S$ un schéma. Soit $X$ un espace algébrique décent et localement
noethérien sur $S$. On dit que $X$ est {\it universellement caténaire} si,
pour tout morphisme $Y \to X$ d'espaces algébriques qui est localement de
type fini, avec $Y$ décent, l'espace algébrique $Y$ est caténaire.
\end{definition}

\noindent
Si $X$ est un espace algébrique, la condition « $X$ est décent et localement
noethérien » équivaut à la condition « $X$ est quasi-séparé et localement
noethérien ». C'est le
Lemme \ref{decent-spaces-lemma-locally-Noetherian-decent-quasi-separated}.
Une autre manière de comprendre la définition précédente est donc que $X$
est universellement caténaire si et seulement si $Y$ est caténaire pour tout
morphisme $Y \to X$ quasi-séparé et localement de type fini.

\begin{lemma}
\label{decent-spaces-lemma-universally-catenary}
Soit $S$ un schéma. Soit $X$ un espace algébrique décent, localement
noethérien et universellement caténaire sur $S$. Alors tout espace algébrique
décent localement de type fini sur $X$ est universellement caténaire.
\end{lemma}

\begin{proof}
Cela résulte formellement des définitions et du fait que les composés de
morphismes localement de type fini sont localement de type fini
(Morphismes d'espaces, lemme
\ref{spaces-morphisms-lemma-composition-finite-type}).
\end{proof}

\begin{lemma}
\label{decent-spaces-lemma-check-dimension-function-finite-cover}
Soit $S$ un schéma. Soit $f : Y \to X$ un morphisme fini surjectif d'espaces
algébriques décents et localement noethériens. Soit
$\delta : |X| \to \mathbf{Z}$ une fonction. Si $\delta \circ |f|$ est une
fonction de dimension, alors $\delta$ est une fonction de dimension.
\end{lemma}

\begin{proof}
Soit $x \leadsto x'$, avec $x \not = x'$, une spécialisation dans $|X|$.
Choisissons $y \in |Y|$ tel que $|f|(y) = x$. Comme $|f|$ est fermée
(Morphismes d'espaces, lemme \ref{spaces-morphisms-lemma-finite-proper}),
il existe une spécialisation $y \leadsto y'$ telle que $|f|(y') = x'$.
On conclut donc que
$\delta(x) = \delta(|f|(y)) > \delta(|f|(y')) = \delta(x')$
(voir Topologie, définition \ref{topology-definition-dimension-function}).
Si $x \leadsto x'$ est une spécialisation immédiate, alors
$y \leadsto y'$ en est aussi une : en effet, si
$y \leadsto y'' \leadsto y'$, alors $|f|(y'')$ est nécessairement égal à
$x$ ou à $x'$, et il n'existe aucune spécialisation non triviale entre des
points d'une même fibre de $|f|$, d'après le
Lemme \ref{decent-spaces-lemma-conditions-on-fibre-and-qf}.
\end{proof}

\noindent
La discussion se poursuivra dans Compléments sur les morphismes d'espaces,
section \ref{spaces-more-morphisms-section-catenary}.













\bibliography{my}
\bibliographystyle{amsalpha}
\end{document}
